\documentclass[11pt,a4paper,oneside]{report}
\usepackage[T1]{fontenc}
\usepackage[utf8]{inputenc}
\usepackage{lmodern,amsmath,amssymb,amsthm,mathtools}
\usepackage[margin=26mm,headheight=14pt]{geometry}
\usepackage{microtype,booktabs,longtable,array,enumitem,xcolor}
\usepackage[colorlinks=true,linkcolor=blue!45!black,citecolor=blue!45!black,urlcolor=blue!45!black]{hyperref}
\setlength{\emergencystretch}{4em}
\setcounter{tocdepth}{0}
\setlist{itemsep=3pt,topsep=5pt}
\allowdisplaybreaks[1]
\newtheorem{theorem}{Theorem}[section]
\newtheorem{lemma}[theorem]{Lemma}
\newtheorem{proposition}[theorem]{Proposition}
\newtheorem{corollary}[theorem]{Corollary}
\theoremstyle{definition}\newtheorem{definition}[theorem]{Definition}
\theoremstyle{remark}\newtheorem{remark}[theorem]{Remark}
\newtheorem{example}[theorem]{Example}
\newcommand{\R}{\mathbb R}\newcommand{\Q}{\mathbb Q}
\newcommand{\Z}{\mathbb Z}\newcommand{\C}{\mathbb C}\newcommand{\N}{\mathbb N}
\DeclareMathOperator{\per}{per}\DeclareMathOperator{\Tr}{Tr}
\DeclareMathOperator{\supp}{supp}\DeclareMathOperator{\rank}{rank}
\DeclareMathOperator{\vol}{vol}\DeclareMathOperator{\conv}{conv}
\makeatletter
\def\@makechapterhead#1{\vspace*{8pt}{\parindent\z@\raggedright\normalfont
\ifnum\c@secnumdepth>\m@ne\Large\bfseries\thechapter\quad\fi
\Large\bfseries #1\par\nobreak\vskip18pt}}
\def\@makeschapterhead#1{\vspace*{8pt}{\parindent\z@\raggedright
\normalfont\Large\bfseries #1\par\nobreak\vskip18pt}}
\makeatother
\hypersetup{pdftitle={Local Certificates, Stability, and Concavity: Integrated Research from Seven Rounds},pdfauthor={Keio}}
% R1E1 private definitions
\theoremstyle{plain}








\newtheorem{srcAAtheorem}{Theorem}[chapter]
\newtheorem{srcAAlemma}[srcAAtheorem]{Lemma}
\newtheorem{srcAAcorollary}[srcAAtheorem]{Corollary}
\newtheorem{srcAAproposition}[srcAAtheorem]{Proposition}
\newcommand{\srcAASn}{S_n}




% R1E2 private definitions
\theoremstyle{plain}





\newtheorem{srcABtheorem}{Theorem}[chapter]
\newtheorem{srcABlemma}[srcABtheorem]{Lemma}
\newtheorem{srcABcorollary}[srcABtheorem]{Corollary}
\newtheorem{srcABproposition}[srcABtheorem]{Proposition}
\newcommand{\srcABout}{N^+}
\newcommand{\srcABinn}{N^-}
\newcommand{\srcABsecond}{N^{++}}
\newcommand{\srcABdefi}{\operatorname{def}}




% R1E3 private definitions
\theoremstyle{plain}








\newtheorem{srcACtheorem}{Theorem}[chapter]
\newtheorem{srcAClemma}[srcACtheorem]{Lemma}
\newtheorem{srcACproposition}[srcACtheorem]{Proposition}
\newcommand{\srcACqb}[2]{\genfrac{[}{]}{0pt}{}{#1}{#2}_{q}}
\newcommand{\srcACLL}{\mathcal L}
\newcommand{\srcACNN}{\mathbb N_0}




% R1E4 private definitions
\theoremstyle{plain}







\newtheorem{srcADtheorem}{Theorem}[chapter]
\newtheorem{srcADlemma}[srcADtheorem]{Lemma}
\newtheorem{srcADproposition}[srcADtheorem]{Proposition}
\newtheorem{srcADcorollary}[srcADtheorem]{Corollary}
\theoremstyle{srcADremark}\newtheorem{srcADremark}[srcADtheorem]{Remark}
\newcommand{\srcADdd}{\mathrm{d}}




% R1E5 private definitions
\theoremstyle{plain}











\newtheorem{srcAEtheorem}{Theorem}[chapter]
\newtheorem{srcAElemma}[srcAEtheorem]{Lemma}
\newtheorem{srcAEcorollary}[srcAEtheorem]{Corollary}
\newtheorem{srcAEproposition}[srcAEtheorem]{Proposition}
\theoremstyle{srcAEremark}\newtheorem{srcAEremark}[srcAEtheorem]{Remark}
\DeclareMathOperator{\srcAEper}{per}
\DeclareMathOperator{\srcAEtr}{tr}
\DeclareMathOperator{\srcAEdiag}{diag}
\newcommand{\srcAEC}{\mathbb C}
\newcommand{\srcAER}{\mathbb R}
\newcommand{\srcAEcalM}{\mathcal M}




% R1E6 private definitions
\theoremstyle{plain}





\newtheorem{srcAFtheorem}{Theorem}[section]
\newtheorem{srcAFlemma}[srcAFtheorem]{Lemma}
\newtheorem{srcAFproposition}[srcAFtheorem]{Proposition}
\newtheorem{srcAFcorollary}[srcAFtheorem]{Corollary}
\theoremstyle{srcAFdefinition}\newtheorem{srcAFdefinition}[srcAFtheorem]{Definition}
\newcommand{\srcAFE}{\mathcal E}
\newcommand{\srcAFZ}{\mathbb Z}
\newcommand{\srcAFR}{\mathbb R}
\newcommand{\srcAFone}{\mathbf 1}




% R2E1 private definitions
\theoremstyle{plain}







\newtheorem{srcAGtheorem}{Theorem}[chapter]
\newtheorem{srcAGlemma}[srcAGtheorem]{Lemma}
\newtheorem{srcAGcorollary}[srcAGtheorem]{Corollary}
\newtheorem{srcAGproposition}[srcAGtheorem]{Proposition}
\newcommand{\srcAGZp}{\mathbb Z_p}
\newcommand{\srcAGQp}{\mathbb Q_p}
\newcommand{\srcAGvp}{v_p}




% R2E2 private definitions
\theoremstyle{plain}



\newtheorem{srcAHtheorem}{Theorem}[chapter]\newtheorem{srcAHlemma}[srcAHtheorem]{Lemma}\newtheorem{srcAHcorollary}[srcAHtheorem]{Corollary}
\newcommand{\srcAHD}{\Delta}\newcommand{\srcAHR}{\mathbb R}




% R2E3 private definitions
\theoremstyle{plain}








\newtheorem{srcAItheorem}{Theorem}[section]
\newtheorem{srcAIlemma}[srcAItheorem]{Lemma}
\newtheorem{srcAIcorollary}[srcAItheorem]{Corollary}
\newtheorem{srcAIproposition}[srcAItheorem]{Proposition}
\theoremstyle{srcAIremark}\newtheorem{srcAIremark}[srcAItheorem]{Remark}
\DeclareMathOperator{\srcAIper}{per}
\newcommand{\srcAIQ}{\mathcal Q}
\newcommand{\srcAIh}{\mathrm d}




% R2E4 private definitions
\theoremstyle{plain}







\newtheorem{srcAJtheorem}{Theorem}[chapter]
\newtheorem{srcAJlemma}[srcAJtheorem]{Lemma}
\newtheorem{srcAJproposition}[srcAJtheorem]{Proposition}
\newtheorem{srcAJcorollary}[srcAJtheorem]{Corollary}
\theoremstyle{srcAJremark}\newtheorem{srcAJremark}[srcAJtheorem]{Remark}
\newcommand{\srcAJdd}{\mathrm{d}}




% R2E5 private definitions
\theoremstyle{plain}









\newtheorem{srcAKtheorem}{Theorem}[chapter]
\newtheorem{srcAKlemma}[srcAKtheorem]{Lemma}
\newtheorem{srcAKcorollary}[srcAKtheorem]{Corollary}
\newtheorem{srcAKproposition}[srcAKtheorem]{Proposition}
\theoremstyle{srcAKdefinition}\newtheorem{srcAKdefinition}[srcAKtheorem]{Definition}
\DeclareMathOperator{\srcAKper}{per}
\DeclareMathOperator{\srcAKdiag}{diag}
\newcommand{\srcAKC}{\mathbb C}
\newcommand{\srcAKR}{\mathbb R}
\newcommand{\srcAKPP}{P}
\newcommand{\srcAKQQ}{Q}
\newcommand{\srcAKhh}{h}




% R2E6 private definitions
\theoremstyle{plain}






\newtheorem{srcALtheorem}{Theorem}[section]
\newtheorem{srcALlemma}[srcALtheorem]{Lemma}
\newtheorem{srcALproposition}[srcALtheorem]{Proposition}
\newtheorem{srcALcorollary}[srcALtheorem]{Corollary}
\newcommand{\srcALE}{\mathcal E}
\newcommand{\srcALZ}{\mathbb Z}
\newcommand{\srcALR}{\mathbb R}
\newcommand{\srcALemin}{\mathrm E_{\min}}
\newcommand{\srcALQ}{Q}




% R3E1 private definitions
\theoremstyle{plain}






\newtheorem{srcAMtheorem}{Theorem}[chapter]
\newtheorem{srcAMlemma}[srcAMtheorem]{Lemma}
\newtheorem{srcAMcorollary}[srcAMtheorem]{Corollary}
\newcommand{\srcAMvp}{v_p}




% R3E2 private definitions
\theoremstyle{plain}







\newtheorem{srcANtheorem}{Theorem}[chapter]\newtheorem{srcANlemma}[srcANtheorem]{Lemma}\newtheorem{srcANcorollary}[srcANtheorem]{Corollary}\newtheorem{srcANproposition}[srcANtheorem]{Proposition}
\newcommand{\srcANStr}{\operatorname{Str}}\newcommand{\srcANdefi}{\operatorname{def}}\newcommand{\srcANnorm}[1]{\left\|#1\right\|}




% R3E3 private definitions
\theoremstyle{plain}







\newtheorem{srcAOtheorem}{Theorem}[section]
\newtheorem{srcAOlemma}[srcAOtheorem]{Lemma}
\newtheorem{srcAOcorollary}[srcAOtheorem]{Corollary}
\DeclareMathOperator{\srcAOper}{per}
\newcommand{\srcAOQ}{\mathcal Q}\newcommand{\srcAOdd}{\mathrm d}




% R3E4 private definitions
\theoremstyle{plain}







\newtheorem{srcAPtheorem}{Theorem}[section]
\newtheorem{srcAPlemma}[srcAPtheorem]{Lemma}
\newtheorem{srcAPproposition}[srcAPtheorem]{Proposition}
\newtheorem{srcAPcorollary}[srcAPtheorem]{Corollary}
\theoremstyle{srcAPremark}\newtheorem{srcAPremark}[srcAPtheorem]{Remark}
\DeclareMathOperator{\srcAPper}{per}
\DeclareMathOperator{\srcAPRea}{Re}
\newcommand{\srcAPh}{\operatorname{h}}
\newcommand{\srcAPBS}{\mathcal B}
\newcommand{\srcAPC}{\mathbb C}
\newcommand{\srcAPR}{\mathbb R}




% R3E5 private definitions
\theoremstyle{plain}










\newtheorem{srcAQtheorem}{Theorem}[section]
\newtheorem{srcAQlemma}[srcAQtheorem]{Lemma}
\newtheorem{srcAQproposition}[srcAQtheorem]{Proposition}
\newtheorem{srcAQcorollary}[srcAQtheorem]{Corollary}
\theoremstyle{srcAQremark}
\newtheorem{srcAQremark}[srcAQtheorem]{Remark}
\DeclareMathOperator{\srcAQper}{per}
\DeclareMathOperator{\srcAQsgn}{sgn}
\DeclareMathOperator{\srcAQRea}{Re}
\DeclareMathOperator{\srcAQIma}{Im}
\newcommand{\srcAQSch}{\mathcal S}
\newcommand{\srcAQId}{I}
\newcommand{\srcAQC}{\mathbb C}
\newcommand{\srcAQG}{\mathfrak S}
\newcommand{\srcAQClass}{\mathcal C}




% R3E6 private definitions
\theoremstyle{plain}






\newtheorem{srcARtheorem}{Theorem}[section]
\newtheorem{srcARlemma}[srcARtheorem]{Lemma}
\newtheorem{srcARproposition}[srcARtheorem]{Proposition}
\newtheorem{srcARcorollary}[srcARtheorem]{Corollary}
\newcommand{\srcARE}{\mathcal E}
\newcommand{\srcARZ}{\mathbb Z}
\newcommand{\srcARR}{\mathbb R}
\newcommand{\srcARF}{\mathcal F}
\newcommand{\srcARrate}{\lambda}




% R4E1 private definitions
\theoremstyle{plain}




\newtheorem{srcAStheorem}{Theorem}[chapter]\newtheorem{srcASlemma}[srcAStheorem]{Lemma}\newtheorem{srcAScorollary}[srcAStheorem]{Corollary}
\newcommand{\srcASvp}{v_p}\newcommand{\srcASZp}{\mathbb Z_p}



% R4E2 private definitions
\theoremstyle{plain}






\newtheorem{srcATtheorem}{Theorem}[chapter]\newtheorem{srcATlemma}[srcATtheorem]{Lemma}\newtheorem{srcATcorollary}[srcATtheorem]{Corollary}\newtheorem{srcATproposition}[srcATtheorem]{Proposition}
\newcommand{\srcATdefi}{\operatorname{def}}\newcommand{\srcATnorm}[1]{\left\|#1\right\|}



% R4E3 private definitions
\theoremstyle{plain}







\newtheorem{srcAUtheorem}{Theorem}[section]\newtheorem{srcAUlemma}[srcAUtheorem]{Lemma}
\newtheorem{srcAUproposition}[srcAUtheorem]{Proposition}\newtheorem{srcAUcorollary}[srcAUtheorem]{Corollary}
\DeclareMathOperator{\srcAUper}{per}\DeclareMathOperator{\srcAUsgn}{sgn}
\newcommand{\srcAUSch}{\mathcal S}\newcommand{\srcAURefl}{\mathcal R}




% R4E4 private definitions
\theoremstyle{plain}






\newtheorem{srcAVtheorem}{Theorem}[section]
\newtheorem{srcAVlemma}[srcAVtheorem]{Lemma}
\newtheorem{srcAVproposition}[srcAVtheorem]{Proposition}
\newtheorem{srcAVcorollary}[srcAVtheorem]{Corollary}
\theoremstyle{srcAVremark}\newtheorem{srcAVremark}[srcAVtheorem]{Remark}
\DeclareMathOperator{\srcAVper}{per}
\DeclareMathOperator{\srcAVsgn}{sgn}
\newcommand{\srcAVSch}{\mathcal S}
\newcommand{\srcAVDef}{\mathcal D}
\newcommand{\srcAVC}{\mathbb C}




% R4E5 private definitions
\theoremstyle{plain}








\newtheorem{srcAWtheorem}{Theorem}[section]
\newtheorem{srcAWlemma}[srcAWtheorem]{Lemma}
\newtheorem{srcAWproposition}[srcAWtheorem]{Proposition}
\newtheorem{srcAWcorollary}[srcAWtheorem]{Corollary}
\theoremstyle{srcAWremark}\newtheorem{srcAWremark}[srcAWtheorem]{Remark}
\DeclareMathOperator{\srcAWper}{per}\DeclareMathOperator{\srcAWsgn}{sgn}
\DeclareMathOperator{\srcAWRea}{Re}\DeclareMathOperator{\srcAWIma}{Im}
\newcommand{\srcAWSch}{\mathcal S}\newcommand{\srcAWG}{\mathfrak S}
\newcommand{\srcAWClass}{\mathcal C}\newcommand{\srcAWC}{\mathbb C}
\newcommand{\srcAWDef}{\Delta}\newcommand{\srcAWId}{I}




% R4E6 private definitions
\theoremstyle{plain}






\newtheorem{srcAXtheorem}{Theorem}[section]
\newtheorem{srcAXlemma}[srcAXtheorem]{Lemma}
\newtheorem{srcAXproposition}[srcAXtheorem]{Proposition}
\newtheorem{srcAXcorollary}[srcAXtheorem]{Corollary}
\newtheorem{srcAXremark}[srcAXtheorem]{Remark}
\newcommand{\srcAXZ}{\mathbb Z}
\newcommand{\srcAXR}{\mathbb R}
\newcommand{\srcAXQ}{\mathbb Q}
\newcommand{\srcAXone}{\mathbf 1}
\newcommand{\srcAXdivg}{\operatorname{div}}
\newcommand{\srcAXsupp}{\operatorname{supp}}




% R5E1 private definitions
\theoremstyle{plain}




\newtheorem{srcAYtheorem}{Theorem}[chapter]\newtheorem{srcAYlemma}[srcAYtheorem]{Lemma}\newtheorem{srcAYcorollary}[srcAYtheorem]{Corollary}
\newcommand{\srcAYZp}{\mathbb Z_p}




% R5E2 private definitions
\theoremstyle{plain}





\newtheorem{srcAZtheorem}{Theorem}[chapter]\newtheorem{srcAZlemma}[srcAZtheorem]{Lemma}\newtheorem{srcAZcorollary}[srcAZtheorem]{Corollary}
\newcommand{\srcAZdefi}{\operatorname{def}}\newcommand{\srcAZnorm}[1]{\left\|#1\right\|}



% R5E3 private definitions
\theoremstyle{plain}








\newtheorem{srcBAtheorem}{Theorem}[section]\newtheorem{srcBAlemma}[srcBAtheorem]{Lemma}
\newtheorem{srcBAproposition}[srcBAtheorem]{Proposition}\newtheorem{srcBAcorollary}[srcBAtheorem]{Corollary}
\DeclareMathOperator{\srcBAarea}{area}\DeclareMathOperator{\srcBArank}{rank}
\newcommand{\srcBAZ}{\mathbb Z}\newcommand{\srcBAR}{\mathbb R}




% R5E4 private definitions
\theoremstyle{plain}






\newtheorem{srcBBtheorem}{Theorem}[section]
\newtheorem{srcBBlemma}[srcBBtheorem]{Lemma}
\newtheorem{srcBBproposition}[srcBBtheorem]{Proposition}
\newtheorem{srcBBcorollary}[srcBBtheorem]{Corollary}
\theoremstyle{srcBBremark}\newtheorem{srcBBremark}[srcBBtheorem]{Remark}
\DeclareMathOperator{\srcBBper}{per}
\newcommand{\srcBBSch}{\mathcal S}\newcommand{\srcBBDef}{\mathcal D}\newcommand{\srcBBC}{\mathbb C}




% R5E5 private definitions
\theoremstyle{plain}








\newtheorem{srcBCtheorem}{Theorem}[section]
\newtheorem{srcBClemma}[srcBCtheorem]{Lemma}
\newtheorem{srcBCproposition}[srcBCtheorem]{Proposition}
\newtheorem{srcBCcorollary}[srcBCtheorem]{Corollary}
\theoremstyle{srcBCremark}\newtheorem{srcBCremark}[srcBCtheorem]{Remark}
\DeclareMathOperator{\srcBCper}{per}\DeclareMathOperator{\srcBCRea}{Re}
\newcommand{\srcBCSch}{\mathcal S}\newcommand{\srcBCDef}{\Delta}
\newcommand{\srcBCG}{\mathfrak S}\newcommand{\srcBCId}{I}\newcommand{\srcBCC}{\mathbb C}




% R5E6 private definitions
\theoremstyle{plain}






\newtheorem{srcBDtheorem}{Theorem}[section]
\newtheorem{srcBDlemma}[srcBDtheorem]{Lemma}
\newtheorem{srcBDproposition}[srcBDtheorem]{Proposition}
\newtheorem{srcBDcorollary}[srcBDtheorem]{Corollary}
\newtheorem{srcBDremark}[srcBDtheorem]{Remark}
\newcommand{\srcBDZ}{\mathbb Z}
\newcommand{\srcBDR}{\mathbb R}
\newcommand{\srcBDQ}{\mathbb Q}
\newcommand{\srcBDone}{\mathbf1}
\newcommand{\srcBDHess}{\operatorname{Hess}}
\newcommand{\srcBDrank}{\operatorname{rank}}
\newcommand{\srcBDconv}{\operatorname{conv}}




% R6E1 private definitions
\theoremstyle{plain}





\newtheorem{srcBEtheorem}{Theorem}[chapter]\newtheorem{srcBElemma}[srcBEtheorem]{Lemma}\newtheorem{srcBEcorollary}[srcBEtheorem]{Corollary}
\newcommand{\srcBEZp}{\mathbb Z_p}




% R6E2 private definitions
\theoremstyle{plain}







\newtheorem{srcBFtheorem}{Theorem}[chapter]\newtheorem{srcBFlemma}[srcBFtheorem]{Lemma}\newtheorem{srcBFcorollary}[srcBFtheorem]{Corollary}
\newcommand{\srcBFdefi}{\operatorname{def}}\newcommand{\srcBFnorm}[1]{\left\|#1\right\|}



% R6E3 private definitions
\theoremstyle{plain}







\newtheorem{srcBGtheorem}{Theorem}[section]\newtheorem{srcBGlemma}[srcBGtheorem]{Lemma}\newtheorem{srcBGcorollary}[srcBGtheorem]{Corollary}\newtheorem{srcBGproposition}[srcBGtheorem]{Proposition}
\newcommand{\srcBGZ}{\mathbb Z}\newcommand{\srcBGR}{\mathbb R}\DeclareMathOperator{\srcBGvol}{vol}



% R6E4 private definitions
\theoremstyle{plain}






\newtheorem{srcBHtheorem}{Theorem}[section]
\newtheorem{srcBHlemma}[srcBHtheorem]{Lemma}
\newtheorem{srcBHproposition}[srcBHtheorem]{Proposition}
\newtheorem{srcBHcorollary}[srcBHtheorem]{Corollary}
\DeclareMathOperator{\srcBHper}{per}
\DeclareMathOperator{\srcBHtr}{tr}
\newcommand{\srcBHC}{\mathbb C}
\newcommand{\srcBHE}{\mathcal E}
\newcommand{\srcBHLr}{\mathcal R}




% R6E5 private definitions
\theoremstyle{plain}








\newtheorem{srcBItheorem}{Theorem}[section]
\newtheorem{srcBIlemma}[srcBItheorem]{Lemma}
\newtheorem{srcBIproposition}[srcBItheorem]{Proposition}
\newtheorem{srcBIcorollary}[srcBItheorem]{Corollary}
\theoremstyle{srcBIremark}\newtheorem{srcBIremark}[srcBItheorem]{Remark}
\DeclareMathOperator{\srcBIper}{per}\DeclareMathOperator{\srcBIRea}{Re}
\newcommand{\srcBIC}{\mathbb C}\newcommand{\srcBIId}{I}\newcommand{\srcBISch}{\mathcal S}




% R6E6 private definitions
\theoremstyle{plain}





\newtheorem{srcBJtheorem}{Theorem}[section]
\newtheorem{srcBJlemma}[srcBJtheorem]{Lemma}
\newtheorem{srcBJproposition}[srcBJtheorem]{Proposition}
\newtheorem{srcBJcorollary}[srcBJtheorem]{Corollary}
\newcommand{\srcBJZ}{\mathbb Z}\newcommand{\srcBJR}{\mathbb R}\newcommand{\srcBJone}{\mathbf1}
\newcommand{\srcBJHess}{\operatorname{Hess}}\newcommand{\srcBJrank}{\operatorname{rank}}
\newcommand{\srcBJconv}{\operatorname{conv}}





% R7E1 private definitions
\theoremstyle{plain}





\newtheorem{srcBKtheorem}{Theorem}[chapter]\newtheorem{srcBKlemma}[srcBKtheorem]{Lemma}\newtheorem{srcBKcorollary}[srcBKtheorem]{Corollary}
\newcommand{\srcBKZp}{\mathbb Z_p}\newcommand{\srcBKZloc}{\mathbb Z_{(p)}}




% R7E2 private definitions
\theoremstyle{plain}






\newtheorem{srcBLtheorem}{Theorem}[chapter]\newtheorem{srcBLlemma}[srcBLtheorem]{Lemma}\newtheorem{srcBLcorollary}[srcBLtheorem]{Corollary}\newtheorem{srcBLproposition}[srcBLtheorem]{Proposition}
\newcommand{\srcBLStr}{\operatorname{Str}}\newcommand{\srcBLdefi}{\operatorname{def}}\newcommand{\srcBLnorm}[1]{\left\|#1\right\|}



% R7E3 private definitions
\theoremstyle{plain}






\newtheorem{srcBMtheorem}{Theorem}[chapter]\newtheorem{srcBMproposition}[srcBMtheorem]{Proposition}\newtheorem{srcBMlemma}[srcBMtheorem]{Lemma}\newtheorem{srcBMcorollary}[srcBMtheorem]{Corollary}
\newcommand{\srcBMLL}{\mathcal L}\newcommand{\srcBMT}{\mathcal T_q}\newcommand{\srcBMbr}[1]{[#1]_q}




% R7E4 private definitions
\theoremstyle{plain}






\newtheorem{srcBNtheorem}{Theorem}[section]
\newtheorem{srcBNlemma}[srcBNtheorem]{Lemma}
\newtheorem{srcBNproposition}[srcBNtheorem]{Proposition}
\newtheorem{srcBNcorollary}[srcBNtheorem]{Corollary}
\DeclareMathOperator{\srcBNper}{per}
\DeclareMathOperator{\srcBNRea}{Re}
\newcommand{\srcBNC}{\mathbb C}
\newcommand{\srcBNSch}{\mathcal S}
\newcommand{\srcBNId}{I}
\newcommand{\srcBNK}{\mathcal K}




% R7E5 private definitions
\theoremstyle{plain}








\newtheorem{srcBOtheorem}{Theorem}[section]
\newtheorem{srcBOlemma}[srcBOtheorem]{Lemma}
\newtheorem{srcBOproposition}[srcBOtheorem]{Proposition}
\newtheorem{srcBOcorollary}[srcBOtheorem]{Corollary}
\theoremstyle{srcBOremark}\newtheorem{srcBOremark}[srcBOtheorem]{Remark}
\DeclareMathOperator{\srcBOper}{per}\DeclareMathOperator{\srcBOtr}{tr}
\newcommand{\srcBOC}{\mathbb C}\newcommand{\srcBOR}{\mathbb R}
\newcommand{\srcBOId}{I}\newcommand{\srcBOF}{\mathrm F}




% R7E6 private definitions
\theoremstyle{plain}





\newtheorem{srcBPtheorem}{Theorem}[section]
\newtheorem{srcBPlemma}[srcBPtheorem]{Lemma}
\newtheorem{srcBPproposition}[srcBPtheorem]{Proposition}
\newtheorem{srcBPcorollary}[srcBPtheorem]{Corollary}
\newcommand{\srcBPZ}{\mathbb Z}\newcommand{\srcBPR}{\mathbb R}\newcommand{\srcBPone}{\mathbf1}
\newcommand{\srcBPvol}{\operatorname{vol}}\newcommand{\srcBPlin}{\operatorname{lin}}
\newcommand{\srcBPdiam}{\operatorname{diam}}\newcommand{\srcBPrank}{\operatorname{rank}}
\newcommand{\srcBPcR}{\mathcal R}\newcommand{\srcBPcN}{\mathcal N}





\begin{document}
\begin{titlepage}
\centering
\vspace*{20mm}
{\Huge\bfseries Local Certificates,\par Stability, and Concavity\par}
\vspace{10mm}
{\Large An Integrated Mathematical Manuscript\par}
\vspace{5mm}
{\large Seven Rounds, Forty-Two Source Developments,\par and Further Consequences\par}
\vspace{15mm}
{\Large Keio\par}
\vspace{6mm}
21 October 2026
\vfill
\begin{minipage}{.88\textwidth}\small
\textbf{Abstract.}
We integrate forty-two mathematical developments on local valuations,
strong second neighborhoods, gamma quotients, Gaussian-binomial columns,
permanent inequalities, and lattice-point concavity. Full source arguments
are retained in a single notation-isolated manuscript. We prove additional
consequences along six dependency chains: effective rational mixed moments
for translated Gaussian sums; same-order prime regular quasirandom
tournaments with opposite strong-vertex densities and an induced-sampling
obstruction; an enlarged radial permanent-on-top domain with row bound
$1/23$; explicit diagonal budgets for mixed permanent certificates;
concavity and transition curvature of an optimal gamma exponent together
with reciprocal Gaussian moment theorems; and seed-independent Ewald
tower rates with compatible half-facet corrections. The statements are
restricted theorems with complete arguments, not resolutions of the
unrestricted conjectures discussed in the source literature. Source
attribution, overlap, computational checks, and remaining limitations are
recorded explicitly. The single \texttt{.tex} file also carries a recoverable
verification supplement.
\end{minipage}
\vspace{10mm}
\end{titlepage}
\tableofcontents

\chapter{Scope, provenance, and the integrated results}
\label{chap:scope}

This manuscript combines the mathematical contents of the supplied
\emph{Math Tournament, Rounds 01--07} archive. Each round contains six
research developments and one assessment: there are forty-two research
texts and seven editorial assessments. The research texts are preserved
as forty-two source chapters, organized by mathematical subject. Their
proofs, examples, and mathematical appendices remain available, with
chapter-local notation and collision-free references. The assessments
inform attribution and corrections; rankings are not mathematical results
and are not reproduced as additional research chapters.

Chapter~2 contains the additional consequences obtained during this
integration. Every such statement identifies the source results it uses
and supplies its further argument. The expression \emph{additional result}
means additional to the supplied forty-two texts. It is not a certification
of historical novelty. The primary-source checks are selective, and this
edition is a research manuscript rather than a peer-reviewed publication.

\section{What the combination adds}

\begin{longtable}{@{}p{.20\textwidth}p{.72\textwidth}@{}}
\toprule
Chain & Additional conclusion and location\\
\midrule\endhead
Local valuations & All finite translated mixed moments of the specified
Gaussian square sums have effectively computable rational limits.
The summatory error is logarithmic to an explicit power; for $p$-spaced
shifts it depends on the largest moment order, not their sum.
Theorem~\ref{syn:arith:mixed}.\\[5pt]
Tournament structure & Same-order prime regular quasirandom pairs have
strong proportions $1$ and tending to $0$, while almost every deficient
root in the latter has diverging deficiency. Uniform induced samples
cannot reliably estimate that proportion at the stated logarithmic
scale. Section~\ref{syn:tour:section}.\\[5pt]
Operator inequalities & The archived radial permanent-on-top row domain
increases from $1/64$ to $1/23$, with an explicit spectral gap and a
componentwise correlation-mask bound.
Theorem~\ref{syn:op:radius} and Corollary~\ref{syn:op:mask}.\\[5pt]
Mixed certificates & Structural, reflected-companion, and weak-row
permanent certificates may be glued using computable diagonal costs.
Their local deficits transfer quantitatively to the full matrix.
Log-permanent strict convexity also transfers in the weak-piece variables,
even with unbounded global row sums.
Theorems~\ref{syn:struct:allocation}, \ref{syn:struct:deficits},
and~\ref{syn:struct:hessian}.\\[5pt]
Special functions & The optimal gamma exponent is increasing and concave,
$C^1$ but not $C^2$ at its explicit transition. Positive reciprocal powers
of four normalized Gaussian iterates are infinitely divisible moment
sequences with strictly positive Hankel minors of every order.
Theorems~\ref{syn:spec:gamma} and~\ref{syn:spec:moments}.\\[5pt]
Lattice geometry & Every fixed monotone seed has the same asymptotic
minimum rate in the permitted simplex-tower class. Canonical half-facet
corrections pull back exactly under coherent series extension, admit
countable barycenters, and require at most thirty edge terms in nullity
three. Section~\ref{syn:geom:section}.\\
\bottomrule
\end{longtable}

These are separate mathematical closures inside specified classes.
There is no claim that the six subjects are instances of one theorem.
Their shared organizational principle is to retain a local certificate
while controlling the effect of combination: separated valuation roots,
bounded tournament surgery, connected cycle expansions, diagonal
allocation, positive moment representations, or lattice-preserving fiber
maps. Each combination requires its own hypotheses and proof.

\section{Attribution and corrections that affect interpretation}

\begin{enumerate}
\item \textbf{Previously refuted conjectures remain refuted.}
The universal strong Seymour conjectures are false: the second version
of Bai--Li--Park records a 36-vertex counterexample
\cite{R7E2:BLP}. The positive and extremal tournament results here concern
restricted classes. They do not settle the regular-existence question or
the public Nine-Mass conjecture. Likewise, unrestricted permanent-on-top
is false; its validity in certified subclasses is the relevant question.

\item \textbf{Shared results are counted once.}
Round 2 entries 03 and 05 independently develop the complex-cactus
conclusion. Round 6 entries 04 and 05 independently develop the scalar
row-local contraction conclusion. Both developments of each result are
retained, with joint-credit notes. The full operator result in Round 7
is a distinct strengthening; the $1/23$ enlargement belongs to this
integration. Statements in an earlier chapter about a subsequent target
are historical descriptions, and must be read with the later chapters.

\item \textbf{A proof in the archive does not establish priority.}
The fixed gamma parameter $a=6$ has an earlier-dated public proof claim,
identified in the Round 2 source discussion. The general halfway-lattice
Euler--Maclaurin mechanism has a 2025 thesis antecedent reported by
Chu--Karshon; see Section~\ref{syn:geom:framework}. Those facts prevent
counting these mechanisms as newly closed public problems in this edition.
The source arguments remain mathematically useful and are retained.

\item \textbf{A printed conjectural assertion is identified precisely.}
The arithmetic chapters prove the conditional congruence printed as
Sondow--MacMillan Conjecture 12 and the endpoint assertion printed as
Amdeberhan--Moll--Sharma--Vignat Conjecture 6.1, under their stated
hypotheses. They do not solve the Erd\H{o}s--Moser equation. The
additional effective mixed-moment theorem concerns polynomial Gaussian
side-length sequences, and does not remove the height obstruction for
the different exponential exponent-variable model.
\end{enumerate}

\section{Problems that this manuscript leaves open}

The conclusions should not be extended beyond the hypotheses proved in
their chapters. In particular, this manuscript leaves unresolved the
unrestricted higher-power correlations in the exponential valuation
model; regular strong Seymour existence and Nine-Mass; coefficientwise
infinite $q$-log-concavity of all Gaussian columns; general Chollet-type
permanent inequalities; the exact Ewald tower rate and the unrestricted
Ewald-point minimum; and general uncorrected barycentric log-concavity
for unimodular fibers. These are limits of the present arguments; the
source-status review is not an exhaustive survey of all subsequent work.

Positive-real $q$ specializations do not imply coefficientwise polynomial
positivity. A scalar permanent inequality does not imply the full Schur
operator inequality. A corrected lattice count is not the original count.
The fixed-seed Ewald comparison does not permit growing-dimensional
seeds. These distinctions are essential to the additional theorems.

\section{Reading and verification}

The next chapter can be read as the main synthesis paper. Its source
dependencies are included afterwards in six subject parts, so readers
can inspect every inherited argument without opening another archive.
The final chapter gives the complete source-to-chapter map and describes
the verification supplement embedded in the TeX file. Classical external
theorems remain cited in the usual mathematical manner.

The integration was prepared with AI assistance. Source arguments and
new deductions were reviewed by subject, and selected combinations were
checked independently in a second pass. Exact finite checks supplement
the written proofs. They are not formal proof-assistant certificates,
and bounded tests alone do not establish an infinite statement.


\chapter{Additional consequences from combining the source results}

% Integrated synthesis: arithmetic
\section{Arithmetic synthesis: effective joint moments of Gaussian sums}
\label{syn:arith:section}

The seven Entry 01 manuscripts develop two related, but different, valuation
problems.  In Rounds 1--3 the exponent is fixed while a power-sum endpoint or
block length varies.  In Rounds 4--6 the exponent varies in
$T_p(m)=\sum_{a=1}^{p+1}a^m$.  Round 7 returns to a polynomial variable, the
side length of a Gaussian square sum.  This distinction determines what can
be combined: a polynomial height estimate supplies effective logarithmic
control which the exponential model does not presently supply.

\subsection{What the seven papers establish}

Write $S_n(k)=\sum_{a=1}^k a^n$.  The following statements summarize the
arithmetic dependencies of the retained papers; they do not replace their
proofs.
\begin{enumerate}
\item Round 1 proves the implication printed as Conjecture 12 of
Sondow--MacMillan, \emph{Integers} 11 (2011), A34, and strengthens its
conclusion from modulus $p^3$ to $p^4$ for $p\ge5$.  Its hypothesis remains
$S_n(k)\equiv(k+1)^n\pmod{k^2}$.  The example $(n,k,p)=(12,42,7)$ rules out
a uniform replacement of $p^4$ by $p^5$.  This is a power-sum congruence,
not a solution of the Erd\H{o}s--Moser equation.
\item Round 2 describes three simple $p$-adic branches of
$S_n(pa)-aS_n(p)$, for even $n\ge4$, $p\ge7$, and its stated Bernoulli
valuation conditions.  The branch near $-1$ is shifted by the endpoint
term.  The $p=5,n=24$ repeated-root example prevents extending its
simple-root argument unchanged to every prime.
\item Round 3 proves endpoint stability at regular primes and, more
generally, in the stated locally nonirregular Bernoulli residue classes.
It establishes the assertion printed as Conjecture 6.1 of
Amdeberhan--Moll--Sharma--Vignat (2025).  The regular-prime valuation
formulas themselves are credited to the earlier literature; the strict
endpoint gap and local formulation are the relevant deductions here.
\item Rounds 4--5 classify the exponent tree using the Wilson quotient
$W_p=((p-1)!+1)/p$.  Outside $W_p\equiv1\pmod p$, one isometric
$p$-adic root controls the unit-base sum.  The full sum has a possible
endpoint boundary set.  The classical $p$-adic Subspace Theorem proves
that this set is finite, ineffectively, and proves the ordinary mean.
It does not prove that the set is empty.
\item Round 6 reduces distinct-shift correlations to one-root moments
and explicit maximal spikes.  All first-power distinct-shift correlations
converge, but higher powers require stronger control of those spikes.
Its explicit countermodel shows that a convergent first moment and the
correct geometric level densities do not suffice.
\item Round 7 gives a polynomial normalization for selected Gaussian
exponent families at every side length divisible by $p$, and uses
polynomial height to prove effective univariate moments of every order.
The prime-side formulas of Kalinin--Zottor retain their original
attribution (arXiv:\allowbreak2602.00206v3, 28 July 2026).
\end{enumerate}

The additional result below combines the separated-root idea of Round 6
with the effective height estimate of Round 7.  It proves all mixed
moments for finitely many translated Gaussian side-length sequences,
including repeated powers.  This statement is additional to the seven
frozen manuscripts.  No claim of historical priority is made.

\subsection{An effective finite-root summation principle}

Throughout this section, $v=v_p$, and integration over $\mathbb Z_p$ uses
Haar probability measure.  Values at a finite set of roots do not affect
these integrals.

\begin{lemma}[Polynomial singularities with logarithmic height]
\label{syn:arith:finite-roots}
Let $\Gamma\subset\mathbb Z_p$ be finite.  Suppose that a function $F$,
defined away from $\Gamma$, has the following description for some integer
$H\ge1$:
\begin{enumerate}
\item the balls $v(x-\gamma)\ge H$, $\gamma\in\Gamma$, are disjoint;
\item off those balls, $F$ depends only on $x\bmod p^H$ and takes rational
values;
\item on the ball about $\gamma$, $F(x)=Q_\gamma(v(x-\gamma))$, where
$Q_\gamma\in\mathbb Q[Y]$ has degree at most $K$, with $K\ge1$.
\end{enumerate}
Assume, effectively, that for $1\le m\le X$ and $m\ne\gamma$,
\[
 v(m-\gamma)\le C_0+C_1\log(2+X)\qquad(\gamma\in\Gamma).
\]
One may omit any fixed finite set $E$ of integer indices containing
$\Gamma\cap\mathbb Z_{>0}$.  Then
\begin{equation}
\label{syn:arith:finite-root-sum}
 \sum_{\substack{1\le m\le X\\m\notin E}}F(m)
 =X\int_{\mathbb Z_p}F(x)\,dx+O\bigl((\log(2+X))^K\bigr).
\end{equation}
The integral is rational.  All constants and the integral are effective
when the finite residue data, polynomials, and height constants are
effective.
\end{lemma}

\begin{proof}
Define the bounded, $p^H$-periodic function $F_H$ by replacing the value
on each root ball by $Q_\gamma(H)$.  Choose an effectively computable
integer $J=O(\log(2+X))$ exceeding $H$ and the assumed height bound, with
$p^J\ge X$.  For every retained index,
\[
 F(m)=F_H(m)+\sum_{\gamma\in\Gamma}\sum_{h=H+1}^{J}
 \bigl(Q_\gamma(h)-Q_\gamma(h-1)\bigr)
 \mathbf1_{v(m-\gamma)\ge h}.
\]
The count of integers $m\le X$ in the indicated root ball differs from
$X/p^h$ by at most one.  Omitting $E$ adds at most $|E|$ to this error.
The bounded periodic part has an $O(1)$ summation error.  Because
$Q_\gamma(h)-Q_\gamma(h-1)=O(h^{K-1})$, the total counting error is
$O(J^K)$.  The uncounted integral tail is bounded by
\[
 X\sum_{h>J}O(h^{K-1})p^{-h}=O(J^{K-1}).
\]
This proves \eqref{syn:arith:finite-root-sum}.  Its integral is a rational
finite residue average plus finitely many series of the form
$\sum_{h>H}P(h)p^{-h}$ with $P\in\mathbb Q[Y]$.  Such series are rational,
as follows by applying powers of $z\,d/dz$ to $z^{H+1}/(1-z)$ and setting
$z=1/p$.  These operations also give an effective computation.
\end{proof}

This lemma records explicitly the extra hypothesis missing from the
general exponential-root problem: a uniform logarithmic height bound
along the ordinary integers.  The residue geometry alone does not give it.

\subsection{The Gaussian valuation model}

For an odd prime $p$, put
\[
 G_n(N)=\sum_{a,b=1}^{N-1}(a+bi)^n,
 \qquad v_p(x+iy)=\min(v_p(x),v_p(y)).
\]
Fix
\[
 n=L+r,\quad L\ge4,\quad 4\mid L,\quad p-1\mid L,\quad 0\le r\le3,
\]
and set
\[
 b=r+v_p\binom nr,\qquad d=r+1,\qquad
 \delta=\mathbf1_{\{p\equiv3\pmod4,\ p^2-1\mid L\}}.
\]
The normalization proved in Round 7 is
\begin{equation}
\label{syn:arith:H}
 H(u)=\frac{G_n(pu)}{\binom nr p^r((1+i)/2)^r u^{r+1}}
 \in\mathbb Z_{(p)}[u],\qquad H(u)\equiv2-\delta u\pmod p.
\end{equation}
When $\delta=1$, it is an isometry of $\mathbb Z_p$ and has a unique root
$\beta\equiv2\pmod p$.  Consequently, for nonzero Gaussian sums,
\begin{equation}
\label{syn:arith:model}
 V(u):=v_p(G_n(pu))=b+d\,v(u)+\delta v(u-\beta).
\end{equation}
If $\delta=0$, every expression involving $\beta$ is to be omitted.
There is at most one positive integer zero of $G_n(pu)$, namely $u=\beta$
when that root is a positive integer.

\begin{theorem}[All translated mixed moments]
\label{syn:arith:mixed}
Let $t_1,\ldots,t_q$ be distinct nonnegative integers, let
$k_1,\ldots,k_q\ge1$, and put $K_0=k_1+\cdots+k_q$.  Omit the finite set
of $u\ge1$ for which any $G_n(p(u+t_i))$ vanishes.  Then
\begin{equation}
\label{syn:arith:joint-sum}
 \sum_{\substack{1\le u\le X\\G_n(p(u+t_i))\ne0\ (1\le i\le q)}}
 \prod_{i=1}^q V(u+t_i)^{k_i}
 =X\,\mathcal M+O_{n,p,\mathbf t,\mathbf k}
       \bigl((\log(2+X))^{K_0}\bigr),
\end{equation}
where the following integral is rational and effectively computable:
\begin{equation}
\label{syn:arith:joint-integral}
 \mathcal M=\int_{\mathbb Z_p}
 \prod_{i=1}^q
 \bigl[b+d\,v(x+t_i)+\delta v(x+t_i-\beta)\bigr]^{k_i}\,dx.
\end{equation}
If all $t_i$ are divisible by $p$, the error exponent $K_0$ can be
replaced by $K=\max_i k_i$.  In this case set $s_i=t_i/p$ and define
\[
 \mathcal J_p(\mathbf s,\mathbf k;a,c)=
 \int_{\mathbb Z_p}\prod_{i=1}^q
       [a+c(1+v(y+s_i))]^{k_i}\,dy.
\]
Then the constant has the root-independent formula
\begin{equation}
\label{syn:arith:constant}
 \mathcal M=\frac{p-1-\delta}{p}b^{K_0}
 +\frac1p\mathcal J_p(\mathbf s,\mathbf k;b,d)
 +\frac\delta p\mathcal J_p(\mathbf s,\mathbf k;b,1).
\end{equation}
\end{theorem}

\begin{proof}
The possible centers are
\[
 \Gamma=\{-t_i:1\le i\le q\}
 \ \cup\ \{\beta-t_i:1\le i\le q,\ \delta=1\},
\]
with coincident centers counted once.  Choose $H$ strictly larger than
the valuation of every nonzero difference of centers.  On a ball about
one center, every distance to another center is constant.  Thus the
integrand in \eqref{syn:arith:joint-integral} is a rational polynomial
in the single unbounded depth, of degree at most $K_0$.  Off these
balls it is a bounded function of the residue modulo $p^H$.

These finite data can be determined effectively.  Differences within
either family are ordinary integers.  In the exceptional case, a
cross-family difference has the form $\beta-a$ with $a\in\mathbb Z$;
the isometry gives $v(\beta-a)=v(H(a))$.  The exact rational evaluation
$H(a)$ both detects equality and, when nonzero, determines the required
depth.  Hensel lifting then provides every finite residue of $\beta$
needed in the residue table.  Possible integer zeros are decidable by
the rational-root test applied after clearing denominators in $H$.

For the effective height bound, an integer center gives
$v(u+t_i)\le\log_p(X+t_i)$ for positive $u+t_i$.  Choose an integer
$D_0$ prime to $p$ such that $Q=D_0H\in\mathbb Z[u]$, and write
$D=\deg Q$ and $C=\sum_j|[u^j]Q|$.  At a retained index,
\[
 p^{v(u+t_i-\beta)}=p^{v(Q(u+t_i))}
 \le |Q(u+t_i)|\le C(1+X+t_i)^D.
\]
This is the required logarithmic bound.  Lemma
\ref{syn:arith:finite-roots} proves \eqref{syn:arith:joint-sum}, including
the effective rationality assertion and the treatment of omitted zeros.

If $p\mid t_i$ for every $i$, the first family of centers is $0$ modulo
$p$ and the second is $2$ modulo $p$.  Hence there are no cross-family
coincidences, and each center is singular for just one factor of the
integrand.  Its local polynomial degree is at most $k_i$, which proves
the sharper error.  Finally, outside these one or two residue classes
every factor equals $b$.  On $x=py$, formula
\eqref{syn:arith:model} gives the first $\mathcal J_p$ term of
\eqref{syn:arith:constant}.  On $x=\beta+py$ it gives the second.
Each change of variables contributes Haar measure $1/p$.  This also
shows why the unspecified digits of $\beta$ disappear from the constant.
\end{proof}

For arbitrary shifts the proof gives a sometimes better exponent than
$K_0$: take the largest sum of the $k_i$ belonging to factors singular at
one common center.  In particular, without cross-family coincidences
the exponent is $\max_i k_i$.  The weaker uniform exponent in
\eqref{syn:arith:joint-sum} avoids imposing any unproved irrationality
condition on $\beta$.

\subsection{An explicit covariance law}

\begin{corollary}[Pair correlations at $p$-spaced shifts]
\label{syn:arith:covariance}
Let $t>0$ be divisible by $p$, and write $h=v_p(t)\ge1$.  With the same
zero-omission convention, the mean of $V(u)$ is
$\mu=b+(d+\delta)/(p-1)$ and
\begin{equation}
\label{syn:arith:pair}
 \sum_{\substack{1\le u\le X\\G_n(pu)G_n(p(u+t))\ne0}}
 V(u)V(u+t)=X M_h+O_{n,p,t}(\log(2+X)),
\end{equation}
where
\[
 M_h=b^2+\frac{2b(d+\delta)}{p-1}
       +(d^2+\delta)\frac{(p+1)(1-p^{-h})}{(p-1)^2}.
\]
The limiting covariance is therefore
\begin{equation}
\label{syn:arith:cov}
 \frac{(d^2+\delta)(p+1)(1-p^{-h})-(d+\delta)^2}{(p-1)^2}.
\end{equation}
\end{corollary}

\begin{proof}
For a nonzero integer $t$ with $v_p(t)=h$,
\begin{align*}
 C_h&:=\int_{\mathbb Z_p}v(x)v(x+t)\,dx\\
 &=\sum_{j=0}^{h-1}\frac{p-1}{p^{j+1}}j^2
   +p^{-h}\left(h^2+\frac{2h}{p-1}\right)
 =\frac{(p+1)(1-p^{-h})}{(p-1)^2}.
\end{align*}
Indeed, depths below $h$ agree, while on $p^h\mathbb Z_p$ the two
rescaled roots are distinct modulo $p$, so their positive-depth events
are disjoint.  Cross terms between an origin root and a $\beta$ root
also vanish because $t\equiv0$ and $\beta\equiv2\pmod p$.  Expanding
\eqref{syn:arith:joint-integral} gives
$b^2+2b(d+\delta)/(p-1)+(d^2+\delta)C_h$.
The theorem supplies the error $O(\log(2+X))$.  Subtracting $\mu^2$
gives \eqref{syn:arith:cov}.
\end{proof}

For example, at $(p,n)=(3,8)$ one has $(b,d,\delta)=(0,1,1)$.
For the side-length shift $u\mapsto u+3$ the product mean is $4/3$ and
the covariance is $1/3$.  At $(p,n)=(5,8)$ the corresponding shift
$u\mapsto u+5$ gives product mean $3/10$ and covariance $19/80$.
These are exact asymptotic constants, not values inferred from sampled
averages.

\subsection{Scope of the added closure}

The new theorem removes the higher-moment obstruction for every finite
collection of translates of the Gaussian family in
\eqref{syn:arith:H}.  In particular it covers every $n\ge4$ at $p=3,5$;
the elementary exponents $0\le n\le3$ follow separately from
$v_p(G_n(pu))=n(1+v_p(u))$.  At larger primes it retains the exponent
and side-length restrictions of Round 7.  It does not classify exponents
outside those families.

It also does not supply a logarithmic estimate for the exponential
root in $T_p(m)$, make the Subspace Theorem effective, or prove the
boundary set of Rounds 4--5 empty.  Accordingly, neither the unrestricted
valuation-tree question nor the higher-power correlations of that
exponential model are declared closed.  The two previously printed
conjectural implications proved in Rounds 1 and 3 should be described
as proved assertions with identified sources; an exhaustive historical
priority audit has not been performed.  The additional mixed-moment
theorem is a rigorous consequence of the retained Gaussian
normalization and the elementary finite-root argument proved here.


% Integrated synthesis: tournaments
\section{Tournament synthesis: prime quasirandom pairs and a sampling obstruction}
\label{syn:tour:section}

The seven Entry 02 manuscripts concern the matching strengthening of the
second-neighborhood property.  Their cumulative conclusion is a collection of
restricted resilience and spectral theorems, followed by a prime
regularization construction.  They do not prove universal strong existence in
regular tournaments.  In particular, the general tournament conjecture must
not be presented as open or as a new target: Bai--Li--Park's second arXiv
version records a 36-vertex counterexample \cite{R7E2:BLP}.
The ordinary second-neighborhood
property is different and holds at every vertex of a regular tournament.

Here is the additional consequence obtained by combining the constant-eleven
ladder of Round 3, the arbitrary-fiber coded surgery of Round 7, and the
finite-field Paley construction.  There are \emph{prime regular tournaments
of the same order}, both spectrally quasirandom, whose strong-vertex
proportions tend respectively to one and zero.  In the latter tournament
almost every root has deficiency tending to infinity.  The pair gives a
quantitative obstruction to estimating the strong proportion from an induced
vertex sample, even with primeness and regularity promised.  This last
restriction and the same-order comparison are not asserted in any of the
seven source manuscripts.  Historical priority beyond this source collection
is not claimed.

\subsection{Definitions and the cumulative results}

Let $T$ be a regular tournament of odd order $n=2d+1$.  Its cyclic link at
$v$ is the bipartite graph $B_v$ with shores $N^+(v),N^-(v)$ and edges
$x y$ whenever $x\to y$.  Write
\[
 \delta_T(v)=d-\nu(B_v),\qquad
 \operatorname{Str}(T)=\{v:\delta_T(v)=0\},\qquad
 \rho(T)=|\operatorname{Str}(T)|/n.
\]
Every inneighbor of $v$ is an exact second outneighbor: otherwise it would
dominate $v$ and all $d$ outneighbors of $v$, contradicting regularity.
Consequently $\delta_T(v)=0$ is precisely the strong Seymour property.
Let $K_T$ be the real skew-adjacency matrix, with entry $1$ on an arc,
$-1$ on its reverse, and zero diagonal; put $\sigma(T)=\|K_T\|_2$.
We call a sequence spectrally quasirandom when $\sigma(T)=o(n)$.
A module is a vertex set whose members have identical orientations toward
each outside vertex.  A tournament is prime when it has no module of size
between two and $n-1$.

\begin{center}
\begin{tabular}{p{0.08\linewidth}p{0.84\linewidth}}
\hline
Round & Established contribution, with its restriction \\
\hline
1 & A global deficiency budget for deleting $d_0$ arcs and reversing
     $r_0$ arcs of a doubly regular tournament: total deficiency at most
     $d_0+3r_0$, with at least $n-d_0-2r_0$ strong roots. \\
2 & Local resilience of doubly regular tournaments; the asymptotically
     sharp deletion threshold is $n/8+o(n)$.  Mixed edits have a separate
     sufficient bound and degree restriction. \\
3 & A regular constant-eleven ladder, exact substitution identities, and
     nonprime quasirandom examples with vanishing strong density;
     induced-sample nonestimability in the regular class. \\
4 & The Hall compression bound
     $\sigma^2\ge 2(n-1)(r+1)-1$ when a root has deficiency $r$,
     with matching leading spectral constructions. \\
5 & If $g=2r+2$ and $\sigma^2<(g+1)n-(4g^2-g+1)$, at most
     $\lfloor n/g\rfloor$ roots have deficiency at least $r$;
     exact attainment at infinitely many specified orders. \\
6 & Excluding homogeneous directed triangles, $\sigma^2<6n-116$
     forces at least $\lceil4n/5\rceil$ strong roots.  Prime regular
     examples attain the limiting proportion $4/5$. \\
7 & Coded triangle surgery makes arbitrary amplified regular fibers prime,
     retaining quantitative Hall deficiency.  This gives prime quasirandom
     examples with vanishing strong density. \\
\hline
\end{tabular}
\end{center}
The fixed-$\sigma/\sqrt n$ statements in Rounds 4--6 and the
$\sigma=o(n)$ constructions in Rounds 3 and 7 concern different scales.
There is no contradiction between their positive-density and vanishing-density
conclusions.

\subsection{The two attributed construction inputs}

The following statements are inputs proved in the source chapters, rather
than new claims of this synthesis.

\begin{lemma}[Round 3 ladder and substitution identities]
\label{syn:tour:inputs}
For every odd $h\ge13$ there is a regular tournament $H_h$ of order $h$
with exactly eleven strong roots and deficiency one at every other root.
Let $C_k$ be the cyclic regular tournament of odd order $k$; every root of
$C_k$ is strong.  For regular tournaments $Q,F$, with $|F|=m$, substitution
satisfies
\[
 \delta_{Q[F]}(u,v)=m\delta_Q(u)+\delta_F(v).
\]
If $P_p$ is the Paley tournament for a prime $p\equiv3\pmod4$, then
$\sigma(P_p[F])=m\sqrt p$.
\end{lemma}

The ladder is explicitly constructed and certified in Round 3, Sections
2--3 (Chapter~\ref{source:R3E2}).  For the identity, the link of a substituted root is the disjoint
union of its internal link and the $m$-fold blowup of its quotient link:
there are no cross edges between these two parts.  Maximum matchings in
the blowup have size $m$ times the original maximum matching size, by
Hall's deficiency formula.  The norm identity follows by splitting the
vertex space into vectors constant on each fiber and vectors summing to
zero on each fiber.

\begin{lemma}[Round 7 coded surgery, in its arbitrary-fiber form]
\label{syn:tour:surgery}
Let $p\ge127$ be prime with $p\equiv3\pmod4$, and let $F$ be any regular
tournament of odd order $m\ge25$.  Put $r=\lceil\sqrt m\rceil$.
There is a prime regular tournament $\mathcal C_p(F)$ of order $pm$ such that
\begin{align}
 \sigma(\mathcal C_p(F))&\le m\sqrt p+8r+4,
       \label{syn:tour:surgerynorm}\\
 \delta_{\mathcal C_p(F)}(i,v)&\ge\delta_F(v)-4r-1
       \quad\hbox{for every root }(i,v).
       \label{syn:tour:surgerydef}
\end{align}
\end{lemma}

This is the combined content of Round 7's lemmas entitled ``Regularity and
perturbation size,'' ``All modules are destroyed,'' and ``Deficiency loss,
including auxiliary roots''
(Lemmas~\ref{R7E2:lem:edit}--\ref{R7E2:lem:loss}).
Its statement is stronger than simply
substituting into the final Round 7 theorem: it allows the amplifier size
to be chosen afresh.  For clarity, its mechanism is recalled here.
Choose a directed Paley triangle $(0,1,b)$ whose three signed difference
classes are distinct.  Their $p$ translates have disjoint quotient edges.
Index a fiber by $0\le x<m$ and set
$a(x)=\lfloor x/r\rfloor$, $c(x)=x\bmod r$.  For every $i,x$ reverse
\[
 (i,x)\to(i+1,a(x))\to(i+b,c(x))\to(i,x).
\]
Injectivity of $(a(x),c(x))$ makes all edited edges disjoint.  Each root
participates in at most $1+2r$ triangles, so its incident edit degree is
at most $4r+2$ and the skew perturbation norm is at most $8r+4$.
Degrees are preserved.  A proper module would have size less than $3m$,
by the module inequality proved below.  Paley row separation gives more
than $3m$ remaining distinguishers for vertices in different fibers.
Within a fiber the distinct pair codes provide an outside distinguisher
for every pair.  Hence no proper nontrivial module survives.

Finally choose an internal Hall set and remove its at most $r$ sources
with indices below $r$.  The remaining sources have only their primary
role.  Their new external targets lie in a set of size $r$, apart from
the at most $2r+1$ new inneighbors of the root.  The total Hall loss is
at most $r+r+(2r+1)=4r+1$.  In particular this bound applies also to
roots used as auxiliaries; omitting those roots would invalidate the
count conclusion below.  The source chapter supplies all finite-order
inequalities needed in the primeness argument.

\subsection{A new same-order prime comparison}

\begin{theorem}[Prime quasirandom pairs with diverging Hall deficiency]
\label{syn:tour:pair}
Let $a\ge1$ be an integer, and let $p\ge127$ be prime with
$p\equiv3\pmod4$.  Put $n=p^{2a+3}$.  There are prime regular tournaments
$T^+_{p,a}$ and $T^-_{p,a}$, both of order $n$, such that
\begin{align}
 \rho(T^+_{p,a})&=1,&
 \sigma(T^+_{p,a})&=\sqrt n,\label{syn:tour:positive}\\
 \rho(T^-_{p,a})&\le 11p^{-a},&
 \sigma(T^-_{p,a})&\le p^{2a+5/2}+8p^{a+1}+4.
       \label{syn:tour:negative}
\end{align}
At least $(1-11p^{-a})n$ roots of $T^-_{p,a}$ have deficiency at least
\begin{equation}
 p^{a+2}-4p^{a+1}-1.
 \label{syn:tour:largegap}
\end{equation}
For each fixed $a$, as $p\to\infty$ through these primes, both sequences
are spectrally quasirandom, their strong proportions tend to one and zero,
and the bound in \eqref{syn:tour:largegap} tends to infinity.
\end{theorem}

\begin{proof}
For the negative member use the Round 3 seed $H_{p^a}$ and choose
\[
 k=p^{a+2},\qquad F=H_{p^a}[C_k],\qquad
 m=p^{2a+2},\qquad r=p^{a+1}.
\]
The substitution identity shows that every vertex above a failed seed
root has deficiency $k$, and exactly $11k$ fiber vertices lie above
strong seed roots.  Apply Lemma~\ref{syn:tour:surgery} to $F$.
It gives a prime regular tournament of order $pm=n$ and the claimed
norm bound.  For every failed-seed lift it gives deficiency at least
$k-4r-1>0$.  There are $(p^a-11)kp=(1-11p^{-a})n$ such vertices.
Only the remaining $11kp=11p^{a+3}$ vertices can be strong.  This proves
all assertions concerning $T^-_{p,a}$, without assuming that any strong
seed root remains strong after surgery.

For the positive member take the Paley tournament over the finite field
$\mathbb F_q$, where $q=n=p^{2a+3}\equiv3\pmod4$.  Orient $x\to y$
when $y-x$ is a nonzero square.  The quadratic character $\chi$ has
$\chi(-1)=-1$ and, for $t\ne0$,
\[
 \sum_{x\in\mathbb F_q}\chi(x)\chi(x+t)=-1.
\]
For completeness, after scaling this follows by counting solutions of
$z^2=x(x+1)$.  The equivalent equation
$(2x+1-2z)(2x+1+2z)=1$ has $q-1$ solutions, whereas the character
count is $q+\sum_x\chi(x(x+1))$.  Thus
$K\mathbf1=0$ and $KK^{\mathsf T}=qI-J$, and the norm is $\sqrt q$.
The tournament is doubly regular; every cyclic link is a balanced
$(q+1)/4$-regular bipartite graph.  Such a graph has a perfect matching
by Hall's theorem: edges leaving any source set number $((q+1)/4)|S|$
and can enter at most $((q+1)/4)|N(S)|$ positions.  Every root is strong.

It remains to check primeness, rather than infer it from quasirandomness.
In any regular tournament, a proper nontrivial module $M$ of size $t$
induces a regular tournament, so $t$ is odd and $t\ge3$.  Exactly half
the outside vertices are dominated by $M$, and the other half dominate
it.  An outside vertex therefore has at least $t$ neighbors in one
direction, giving $t\le(n-1)/2$.  Moreover $K\mathbf1_M$ is zero on
$M$ and is $\pm t$ outside, whence
\begin{equation}
 t(n-t)\le\sigma(T)^2.
 \label{syn:tour:module}
\end{equation}
For the Paley tournament the right side is $n$, whereas the allowed
range gives $t(n-t)\ge3(n-3)>n$ at the present orders.  This is
impossible, proving primeness.

Finally
\[
 \frac{\sigma(T^-_{p,a})}{n}
 \le p^{-1/2}+8p^{-a-2}+4p^{-2a-3}\longrightarrow0,
 \qquad
 \frac{\sigma(T^+_{p,a})}{n}=n^{-1/2}\longrightarrow0.
\]
The strong-density and deficiency limits follow from the explicit bounds.
\end{proof}

\begin{corollary}[Absolute counts on prime quasirandom sequences]
\label{syn:tour:absolutecount}
For every $\eta>0$ there is a sequence of prime regular quasirandom
tournaments with $|\operatorname{Str}(T)|=O(n^{1/2+\eta})$, with all but
$o(n)$ roots having deficiency tending to infinity.  These tournaments
can be paired, at each order, with prime regular all-strong tournaments.
\end{corollary}
\begin{proof}
Choose a fixed integer $a\ge1$ so large that
$3/(4a+6)\le\eta$.  The negative member of
Theorem~\ref{syn:tour:pair} has at most
$11p^{a+3}=11n^{1/2+3/(4a+6)}$ strong roots.  All other assertions
are part of that theorem.
\end{proof}
This corollary is a quantitative consequence of the chosen parameters;
it does not give a bounded strong count or growth slower than every
prescribed divergent function.

\subsection{A quantitative obstruction to induced-sample estimation}

Let $\mu_{T,s}$ be the distribution of the labelled tournament induced
by an ordered uniformly chosen $s$-tuple of distinct vertices of $T$.
Only the adjacency pattern of the sample is observed, not the names of
the sampled vertices; the total order $n$ may additionally be known.
Write $(n)_s=n(n-1)\cdots(n-s+1)$.

\begin{lemma}[Counting bound, with its sampling normalization]
\label{syn:tour:counting}
For two tournaments $T_1,T_2$ of the same order $n$, put
$\sigma_i=\|K_{T_i}\|_2$.  If $2\le s\le n$ and $c=\binom s2$, then
\begin{equation}
 d_{\mathrm{TV}}(\mu_{T_1,s},\mu_{T_2,s})
 \le 2^{c-2}c(\sigma_1+\sigma_2+2)
        \frac{n^{s-1}}{(n)_s}.
 \label{syn:tour:tv}
\end{equation}
\end{lemma}
\begin{proof}
For a labelled tournament $L$ on $s$ vertices let $N_L(T)$ count its
injective labelled copies.  With $A=(J-I+K_T)/2$,
$\|A-J/2\|_2\le(\sigma(T)+1)/2$.  Telescope the product defining
$N_L(T)$ against the constant product $2^{-c}$.  In the term associated
with one edge, fix all coordinates except its two endpoints.  The
remaining factors have the form $c_0 f(i)g(j)$, where
$0\le c_0,f,g\le1$.  The corresponding bilinear error is at most
$(\sigma(T)+1)n/2$.  Summing over the other $n^{s-2}$ coordinates and
the $c$ edges yields
\[
 |N_L(T)-2^{-c}n^s|
 \le \tfrac12 c(\sigma(T)+1)n^{s-1}.
\]
The adjacency product vanishes on tuples with a repeated coordinate,
so the count is indeed injective.  Apply this inequality to $T_1,T_2$,
cancel their common central term, divide by $(n)_s$, sum over the
$2^c$ labelled tournaments, and multiply by one half as required in
total variation distance.  This gives \eqref{syn:tour:tv}.
\end{proof}

\begin{theorem}[Nonestimability inside the prime regular class]
\label{syn:tour:sampling}
Fix an integer $a\ge1$ and $0<\varepsilon<1$.  Along the pairs of
Theorem~\ref{syn:tour:pair}, let $s=s_p\ge2$ satisfy
\begin{equation}
 s_p^2\le(1-\varepsilon)\log_2p
       =\frac{1-\varepsilon}{2a+3}\log_2n.
 \label{syn:tour:samplesize}
\end{equation}
Then
\[
 d_{\mathrm{TV}}(\mu_{T^+_{p,a},s_p},
                 \mu_{T^-_{p,a},s_p})\longrightarrow0.
\]
Consequently no estimator based only on such a uniform induced sample,
even with $n$ known and prime regularity promised, can estimate $\rho(T)$
to additive error $1/4$ with probability at least $2/3$ on every member
of both sequences for all sufficiently large $p$.  The conclusion also
holds for every fixed sample size.
\end{theorem}
\begin{proof}
The norm bounds give $(\sigma_++\sigma_-+2)/n=O(p^{-1/2})$.
Also $s=O(\sqrt{\log p})$ and $n=p^{2a+3}$, so
$n^s/(n)_s=1+o(1)$.  Since $c\le s^2/2$, condition
\eqref{syn:tour:samplesize} implies
\[
 2^c p^{-1/2}\le p^{-\varepsilon/2}.
\]
Lemma~\ref{syn:tour:counting} therefore bounds the total variation
distance by $O((\log p)p^{-\varepsilon/2})$, which tends to zero.

For sufficiently large $p$, the two true proportions are $1$ and at
most $11p^{-a}<1/4$.  If an estimator achieved the stated accuracy on
both inputs, declaring the positive member precisely when the estimate
exceeds $5/8$ would be a test correct with probability at least $2/3$
on each input.  Its acceptance probabilities would differ by at least
$1/3$, contradicting convergence of the sample laws in total variation.
An independent random seed cannot increase total variation, so the
argument includes randomized estimators.
\end{proof}

\subsection{What this closes, and what it leaves open}

The combination closes a precise restricted question: imposing primeness,
regularity, and asymptotic quasirandomness does \emph{not} make the strong
proportion recoverable from bounded induced samples.  The obstruction
persists for the slowly growing samples in
\eqref{syn:tour:samplesize}, and is not caused merely by deficiency-one
roots.  The new all-strong comparison has exactly the same order as the
deficient construction, so knowing the order does not remove the
obstruction.

These statements do not settle whether every regular tournament has a
strong root, or the proposed nine-strong-root lower bound in the public
Nine-Mass notes \cite{R7E2:G}.  They do not
construct a regular tournament with no strong roots, assert bounded
strong counts in prime examples, or give an optimal sample complexity.
The statistical result concerns uniform induced-vertex sampling; it is
not an adaptive adjacency-query lower bound.  The positive and negative
members are not asserted to have small edit distance from each other,
nor to lie in a fixed $\sigma/\sqrt n$ window.  The remaining universal
lower-bound problem cannot be closed by replacing these explicit
restrictions with an unrestricted claim.


% Integrated synthesis: operator
\section{An enlarged row domain for the operator theorem}
\label{syn:op:section}

This section improves a numerical domain in the archived operator argument.
Round 5 entries 04 and 05 independently supplied the same word-energy and
least-edge estimates. Round 6 entries 04 and 05 independently removed
remote scalar decorations by connected expansions. Round 7 entry 04
retained the order of overlapping cycle operators through heaps and
proved radial permanent-on-top monotonicity for row bound $1/64$.
We keep that mechanism, expose its adjustable parameters, and obtain
$1/23$. The improvement is relative to this archive; historical priority
for the optimized constant is not asserted. The scalar-to-operator step
and the heap method remain credited to Round 7 entry 04.

For a Hermitian diagonal-one matrix $A=(a_{ij})$ of order $n$, put
\[
 w_{ij}=|a_{ij}|\ (i\ne j),\qquad
 q=\max_i\sum_{j\ne i}w_{ij},\qquad
 H_A=\sum_{i<j}w_{ij}^2(I-L_{(ij)}).
\]
Here $L_\pi$ is the left regular representation of $S_n$. Write
\[
 A_t=I+t(A-I),\quad S(t)=\mathcal S(A_t),\quad
 P(t)=\operatorname{per}(A_t),\quad T(t)=S(t)/P(t),
\]
where $\mathcal S(A)_{\sigma,\tau}=\prod_i a_{\sigma(i),\tau(i)}$.
Let $K_A=\ker H_A$ and let $\Pi_A$ be its orthogonal projection.
If $K_A^\perp\ne\{0\}$, denote the least positive eigenvalue of $H_A$
by $\delta_A$. All these operators act on $\mathbb C^{n!}$.

\begin{theorem}[An enlarged radial POT domain]\label{syn:op:radius}
If $q\le1/23$, then $A_t\succ0$ for $0\le t\le1$ and
\begin{equation}\label{syn:op:monotone}
 0\prec T(t)\preceq T(s)\preceq I\qquad(0\le s\le t\le1).
\end{equation}
In particular, $\mathcal S(A)\preceq\operatorname{per}(A)I$.
For $t>0$ define
\[
 J(t)=\frac{tP'(t)}{P(t)}I-tS'(t)S(t)^{-1}
      =-tT'(t)T(t)^{-1}.
\]
Then
\begin{equation}\label{syn:op:accretive}
 \frac1{100}t^2H_A\preceq\operatorname{Re}J(t)
 \preceq\frac{399}{100}t^2H_A,
 \qquad \operatorname{Re}J=(J+J^*)/2.
\end{equation}
Moreover,
\begin{equation}\label{syn:op:gap}
 T(t)\preceq\Pi_A+\exp(-\delta_A t^2/200)(I-\Pi_A)
\end{equation}
when $K_A^\perp\ne\{0\}$. For $q\le1/24$, the constants
$1/100,399/100,200$ in these two displays may respectively be replaced
by $1/10,39/10,20$.
If the support components have orders $n_1,\ldots,n_k$, the top
eigenspace of $\mathcal S(A)$ is exactly $K_A$, of dimension
$n!/\prod_jn_j!$. At $A=I$ this is the whole space, and no positive
gap is asserted.
\end{theorem}

\subsection{The parameterized estimate and its proof}
The following version makes explicit which constants can be changed.
It also records the ordering issue that is essential to the argument.

\begin{lemma}[Adjustable heap bound]\label{syn:op:parameter}
Let $a,b,C>0$, set $c=a+b$, and suppose
\begin{equation}\label{syn:op:parameters}
 r=qe^c<1,\qquad \frac{r^2}{1-r}\le a,
 \qquad L^2\le Ce^{bL}\quad(L\ge0).
\end{equation}
Define
\begin{equation}\label{syn:op:eta}
 \eta(x)=C\left[e^{2b}(e^{2a}-1)
       +2e^{2c}\bigl((1-xe^c)^{-2}-1\bigr)\right].
\end{equation}
Then, for $0<t\le1$ and every vector $v$,
\begin{equation}\label{syn:op:form}
 \left|\langle v,(J(t)-2t^2H_A)v\rangle\right|
 \le\eta(tq)t^2\langle v,H_Av\rangle.
\end{equation}
\end{lemma}

\begin{proof}
Use directed simple cycles $\gamma$ of length $\ell_\gamma\ge2$
as polymers, modulo cyclic rotation. A transposition is counted once;
the two orientations of a longer cycle are distinct. Two polymers are
incompatible exactly when their vertex sets meet. Write
$z_\gamma=\prod_{i\in V(\gamma)}a_{\gamma(i),i}$,
$\rho_\gamma=|z_\gamma|$, and $U_\gamma=L_\gamma$.
Compatible cycle operators commute. Permutation decomposition gives
\[
 S(t)=\sum_{\mathcal F\text{ compatible}}
             \prod_{\gamma\in\mathcal F}t^{\ell_\gamma}z_\gamma U_\gamma.
\]

We recall and justify the two jointly credited Round 5 estimates.
For self-adjoint unitaries $V_1,\ldots,V_m$, set
$E_j=\langle v,(I-V_j)v\rangle$. Telescoping the word, and then
telescoping its adjoint prefixes in reverse order, gives
\[
 |\langle v,(I-V_1\cdots V_m)v\rangle|
 \le\sum_jE_j+2\sum_{i<j}\sqrt{E_iE_j}
 \le m\sum_jE_j.
\]
Indeed $\|(I-V_j)v\|^2=2E_j$ and the norm of an adjoint prefix
minus the identity applied to $v$ is at most the sum of the
corresponding $\sqrt{2E_i}$. This bound controls complex quadratic
forms, not only their real parts.

Delete a selected least-weight edge from every long cycle and let
$F_\gamma$ be the remaining path edges. They give a word for the
cycle with $\ell_\gamma-1$ transpositions. For a two-cycle retain its
single edge. If $e$ is retained and $w_e>0$, then
\begin{equation}\label{syn:op:charge}
 \sum_{\substack{\ell_\gamma=\ell\\e\in F_\gamma}}
 \frac{\rho_\gamma}{w_e^2}\le2(\ell-1)q^{\ell-2}\qquad(\ell\ge3).
\end{equation}
To see this, remove $e$ and the chosen least edge $f$. Their factor
after division is $w_f/w_e\le1$. The remaining two paths have total
length $\ell-2$ and determine the cycle and $f$. For each of the
$\ell-1$ possible splits of this length, extend the paths to arbitrary
walks starting at the endpoints of $e$; their total weight is at most
$q^{\ell-2}$. The two orientations supply the factor two.
Ties are resolved once. No zero edge is divided by. Also, for every
vertex $i$,
\begin{equation}\label{syn:op:mass}
 \sum_{\gamma\ni i}\rho_\gamma e^{c\ell_\gamma}
 \le\sum_{\ell\ge2}(qe^c)^\ell=\frac{r^2}{1-r}.
\end{equation}
Each cycle through $i$ is a distinct closed walk from $i$, and
nonnegative walk weights have row sums bounded by powers of $q$.

A heap is a word class in which only adjacent compatible letters
may be interchanged. Its occurrences form a partially ordered set,
with every earlier incompatible occurrence ordered below the later
one. A pyramid has one minimal occurrence, its root. Put
\[
 M=\sum_{\mathcal F\text{ compatible}}(-1)^{|\mathcal F|}
                     \prod_{\gamma\in\mathcal F}x_\gamma,
 \qquad Dx_\gamma=\ell_\gamma x_\gamma.
\]
The formal identities used in Round 7 entry 04 are
\begin{equation}\label{syn:op:heap}
 M^{-1}=\sum_{h\text{ heap}}x(h),\qquad
 -(DM)M^{-1}=\sum_{h\text{ pyramid}}\ell_{\mathrm{root}(h)}x(h).
\end{equation}
For completeness, a compatible left factor of a fixed heap is a subset
of its minimal occurrences. Its coefficient in the first product is
$\sum_{F\subseteq\min h}(-1)^{|F|}$, which cancels unless the heap
is empty. In the differentiated product it is
\[
 \sum_{F\subseteq\min h}(-1)^{|F|}\sum_{\gamma\in F}\ell_\gamma
 =-\sum_{\gamma\in\min h}\ell_\gamma
                    (1-1)^{|\min h|-1}.
\]
Only a pyramid survives. In particular the derivative is on the left
and the inverse on the right; this is not a commutative operator
logarithm.

We spell out the scalar majorant. For a tuple of $k$ polymers let
$\phi$ be the sum of $(-1)^{|E|}$ over connected spanning subgraphs
of its incompatibility graph. Repetitions of polymer types are allowed.
The compatible-family polynomial $Z(u)$ satisfies the formal identity
\[
 \log Z(u)=\sum_{k\ge1}\frac1{k!}\sum_{\gamma_1,\ldots,\gamma_k}
                 \phi(\boldsymbol\gamma)\prod_j u_{\gamma_j}.
\]
This follows by expanding the pair factors that exclude incompatible
pairs and partitioning graphs into connected components. The sign is
$(-1)^{k-1}\phi\ge0$, and $|\phi|$ is at most the number of spanning
trees. One elementary justification is deletion--contraction of the
coefficient of $x$ in the chromatic polynomial: after multiplication
by $(-1)^{k-1}$ its recurrence is additive, starts nonnegative, and is
bounded by the spanning-tree recurrence (merging parallel edges can
only decrease the tree count).

Set $u_\gamma=t^{\ell_\gamma}\rho_\gamma e^{b\ell_\gamma}$.
The mass bound and a union bound over the root's vertices imply
\[
 \sum_{\delta\not\sim\gamma}u_\delta e^{a\ell_\delta}
 \le a\ell_\gamma.
\]
The rooted absolute connected sum therefore obeys
\begin{equation}\label{syn:op:root}
 R_\gamma:=\sum_{k\ge1}\frac1{(k-1)!}
   \sum_{\gamma_2,\ldots,\gamma_k}
    |\phi(\gamma,\gamma_2,\ldots,\gamma_k)|
                    \prod_{j=2}^k u_{\gamma_j}
 \le e^{a\ell_\gamma}.
\end{equation}
Here is the needed convergence argument. Majorize by labeled rooted
trees. Their height-truncated generating functions start at one and
satisfy $R_\gamma^{(h+1)}=\exp(\sum_{\delta\not\sim\gamma}
u_\delta R_\delta^{(h)})$. Induction bounds them by
$e^{a\ell_\gamma}$; monotone convergence gives
\eqref{syn:op:root}. Unordered child branches account for their
factorials, including repeated labels. The absolute connected series
is finite. The sign assertion makes $-\log Z(-u)$ coefficientwise
nonnegative; its exponential has finite coefficient mass and equals
the scalar heap series by \eqref{syn:op:heap}. Thus that heap series
also converges. It dominates the operator inverse series in norm,
and the exponential length margin dominates its differentiated series.
All rearrangements below are consequently absolutely convergent,
uniformly for $0\le t\le1$.

Substitute $x_\gamma=-t^{\ell_\gamma}z_\gamma U_\gamma$ into
\eqref{syn:op:heap}, and subtract the corresponding scalar identity.
Writing $L(h)=\sum_{\gamma\in h}\ell_\gamma$, we obtain
\[
 J(t)=\sum_{h\text{ pyramid}}(-1)^{|h|-1}\ell_{\mathrm{root}(h)}
                  t^{L(h)}z(h)(I-U(h)).
\]
Singleton transpositions contribute exactly $2t^2H_A$.
Abelianizing the scalar version of the derivative identity shows,
coefficientwise for each multiset of occurrences, that
\[
 \sum_{h\text{ pyramid}}\ell_{\mathrm{root}(h)}u(h)
 =\sum_{k\ge1}\frac1{k!}\sum_{\boldsymbol\gamma}
       |\phi(\boldsymbol\gamma)|
       \left(\sum_j\ell_{\gamma_j}\right)\prod_j u_{\gamma_j}.
\]
This scalar equality is used only after each ordered operator word has
been bounded. In that bound the word length is at most $L$, and
the edge energy is $\sum_j\sum_{e\in F_{\gamma_j}}E_e$, where
$E_e=\langle v,(I-L_e)v\rangle$. The derivative supplies a further
factor $L$. Absorb $L^2$ by $Ce^{bL}$ and mark the occurrence
supplying the edge. Symmetry changes $1/k!$ into $1/(k-1)!$.
After removing the singleton two-cycle, the absolute error is at most
\begin{align*}
 &Ce^{2b}\sum_e t^2w_e^2E_e(R_e-1)\\
 &\hspace{10mm}+C\sum_{\ell_\gamma\ge3}
 t^{\ell_\gamma}\rho_\gamma e^{b\ell_\gamma}R_\gamma
                         \sum_{e\in F_\gamma}E_e.
\end{align*}
The first line is bounded with $R_e-1\le e^{2a}-1$.
For the second, \eqref{syn:op:charge} bounds the coefficient at
$t^2w_e^2E_e$ by
\[
 2C\sum_{\ell\ge3}(\ell-1)e^{c\ell}(tq)^{\ell-2}
 =2Ce^{2c}\bigl((1-tqe^c)^{-2}-1\bigr).
\]
This proves \eqref{syn:op:form} and \eqref{syn:op:eta}.
\end{proof}

\subsection{Exact constants and the positive Sylvester inverse}
Take
\[
 a=\frac1{128},\qquad b=\frac58,\qquad c=\frac{81}{128},
 \qquad C=\frac75.
\]
The length budget is valid because
\[
 \sup_{L\ge0}L^2e^{-5L/8}=\frac{256}{25e^2}
 <\frac{147456}{105625}<\frac75;
\]
the first strict inequality uses $e>1+1+1/2+1/6+1/24=65/24$.
All other exponential bounds needed here have rational certificates:
\begin{equation}\label{syn:op:exponential}
 e^{1/64}-1\le\frac1{63},\quad
 e^{5/4}<\frac72,\quad e^{81/128}<\frac{19}{10}.
\end{equation}
The first follows by bounding the exponential series by a geometric
series. For the latter two use
\[
 e^x\le\sum_{j=0}^m\frac{x^j}{j!}
       +\frac{x^{m+1}}{(m+1)!}\frac1{1-x/(m+2)}\quad(0\le x<m+2).
\]
For $(x,m)=(5/4,5)$ this upper bound is
$11837573/3391488<7/2$; for $(x,m)=(81/128,2)$ it is
$26596603/14123008<19/10$.

For $q\le1/23$, the convergence mass is
\[
 r\le\frac{19}{230},\qquad
 \frac{r^2}{1-r}\le\frac{361}{48530}<\frac1{128}.
\]
Moreover \eqref{syn:op:eta} and \eqref{syn:op:exponential} give
\begin{equation}\label{syn:op:rational23}
 \eta(tq)\le\frac75\left[
 \frac7{126}+\frac{361}{50}
          \left(\left(\frac{230}{211}\right)^2-1\right)\right]
 =\frac{99177386}{50086125}<\frac{199}{100}<2.
\end{equation}
For $q\le1/24$ the corresponding two bounds are
\begin{equation}\label{syn:op:rational24}
 \frac{r^2}{1-r}\le\frac{361}{53040}<\frac1{128},\qquad
 \eta(tq)\le\frac{103876556}{54946125}<\frac{19}{10}.
\end{equation}
Taking real parts in Lemma~\ref{syn:op:parameter} proves the stated
accretivity bounds.

\begin{proof}[Completion of Theorem~\ref{syn:op:radius}]
The quadratic-form row estimate gives $A_t\succeq(1-tq)I\succ0$.
Its tensor power is positive definite, as is its principal Schur
submatrix $S(t)$. The constant-vector eigenvalue $P(t)$ is positive,
so $T(t)\succ0$ and $T(0)=I$.
Put $R=T(t)^{-1}$ and $Y=-T'(t)$. The defining identity for $J$
gives the Sylvester equation
\[
 RY+YR=\frac2t\operatorname{Re}J(t),\qquad
 Y=\frac2t\int_0^\infty
       e^{-sR}\operatorname{Re}J(t)e^{-sR}\,ds\succeq0.
\]
The integral converges since $R\succ0$. Differentiating its integrand
verifies the equation; uniqueness follows by diagonalizing $R$, since
each entry is multiplied by $r_i+r_j>0$. This positive inverse is
the reason accretivity implies monotonicity. Multiplying a Loewner
inequality by a noncommuting $T$ would not establish the conclusion.
Integration and continuity at zero give \eqref{syn:op:monotone}.

Every summand of $H_A$ is positive semidefinite. Its kernel consists
of vectors invariant under all support-edge transpositions. Those
transpositions generate $G=\prod_jS_{n_j}$: a path expresses any
transposition within a support component. Therefore $K_A$ consists
of functions constant on the right cosets $G\sigma$, giving the
stated dimension. Every supported permutation belongs to $G$, so
$T(t)$ acts as the identity on $K_A$; its orthogonal complement is
invariant by Hermiticity.

Let $\mu(t)$ be the largest eigenvalue of $T(t)$ on that complement.
For a maximizing unit eigenvector $v$, $T^{-1}v=\mu^{-1}v$ and
\[
 -\frac{t}{\mu(t)}\langle v,T'(t)v\rangle
 =\langle v,\operatorname{Re}J(t)v\rangle
 \ge\frac1{100}t^2\delta_A.
\]
The maximum of the smooth Rayleigh quotients is locally Lipschitz;
its right upper derivative is the maximum of the corresponding
derivatives over maximizing eigenvectors. Hence almost everywhere
$(\log\mu)'\le-\delta_A t/100$. Since $\mu(0)=1$, integration
gives \eqref{syn:op:gap}. The improved coefficient $1/10$ under
$q\le1/24$ gives the alternative denominator $20$. Strictness on
the complement for $t>0$ proves the exact top-eigenspace assertion.
\end{proof}

\subsection{A componentwise spectral mask consequence}
The next formula combines the optimized radial gap with the Schur
product transfer used in Round 4 entries 04 and 05. Its component
factor is useful when support is disconnected.

\begin{corollary}[Mask contraction from a full operator gap]
\label{syn:op:mask}
Assume $q\le1/23$ and $A\ne I$. Let $C_1,\ldots,C_k$ be its
support components and let $B$ be any correlation matrix of the same
order. Put
\[
 M_A(B)=\prod_{j=1}^k
          \frac{\operatorname{per}(B[C_j])}{|C_j|!},
 \qquad \rho=\exp(-\delta_A/200).
\]
Then
\begin{equation}\label{syn:op:maskbound}
 \frac{\operatorname{per}(A\circ B)}{\operatorname{per}(A)}
 \le \rho+(1-\rho)M_A(B)\le1.
\end{equation}
For $q\le1/24$ use $\rho=\exp(-\delta_A/20)$.
Equality $\operatorname{per}(A\circ B)=\operatorname{per}(A)$
holds exactly when $B[C_j]$ has rank one for every component.
This equality characterization was already available on the larger
scalar row domain of Round 6; it is not counted as a new closure.
\end{corollary}
\begin{proof}
The projector is the group average
$\Pi_A=|G|^{-1}\sum_{\pi\in G}L_\pi$.
The positive semidefinite matrix $\mathcal S(B)$ has diagonal one,
and $L_\pi\circ\mathcal S(B)=d_B(\pi)L_\pi$, where
$d_B(\pi)=\prod_i b_{\pi(i),i}$. Thus the row sum of
$\Pi_A\circ\mathcal S(B)$ is
\[
 |G|^{-1}\sum_{\pi\in G}d_B(\pi)
 =\prod_j\frac{\operatorname{per}(B[C_j])}{|C_j|!}=M_A(B).
\]
Take the Schur product of \eqref{syn:op:gap} at $t=1$ with
$\mathcal S(B)$, and evaluate on the normalized constant vector.
This gives the first inequality in \eqref{syn:op:maskbound}.
Each component permanent is nonnegative by its tensor Gram formula
and at most $|C_j|!$ since every permutation monomial has modulus
at most one. Therefore $0\le M_A(B)\le1$.

If the permanent ratio is one, $\rho<1$ forces $M_A(B)=1$.
Every component factor then equals one. Equality in the real-part
bound for the sum of its permutation monomials forces every monomial
to be one; in particular $|b_{ij}|^2=1$ for each pair in that
component. Its unit Gram vectors are collinear, so its rank is one.
Conversely these rank-one blocks give a diagonal unitary gauge on
each component, which preserves the permanent of $A$.
\end{proof}

\subsection{What the intersection with the Hessian theorem does establish}
Round 7 entry 05 proves, independently, that $f(A)=\log\operatorname{per}(A)$
on the diagonal-one Hermitian row ball $q(A)\le1/6$ satisfies
\[
 \frac1{20}\|H\|_F^2\le D^2f(A)[H,H]
                  \le\frac{39}{20}\|H\|_F^2
\]
for every Hermitian zero-diagonal direction $H$. At row bound
$1/12$ its factors improve to $7/10$ and $13/10$.
Its two-root cluster estimate is scalar; it does not itself prove
an order-$n!$ operator inequality. Its exact rational endpoint bounds
are $319789/341936<19/20$ and
$178721191/671304725<3/10$ for the Hessian error from the Frobenius
metric. Those constants have been independently recalculated.

Consequently on the common row ball $q\le1/23$ both the improved
operator theorem and the stronger $7/10,13/10$ Hessian comparison
hold. The former applies to arbitrary correlation masks; the latter
controls affine interpolation and, by Jensen's inequality, mixtures
of rank-one phase masks. The intersection does not justify extending
the operator radius to $1/6$, extending the Hessian conclusion to
every strictly diagonally dominant matrix, or asserting monotonicity
under arbitrary independent changes of individual entries.
In particular the archived exact example
$A=I_3-(2/5)(J_3-I_3)$ has row bound $4/5<1$ but has negative
log-permanent second derivative $-7050/28561$ in direction
$-(J_3-I_3)$.

The concrete additional closure here is the enlargement from $1/64$
to $1/23$ for the full radial operator statement, with explicit gaps
and a componentwise mask formula. It uses no total squared-mass
bound and no minimum nonzero edge size. It does not resolve a
universal permanent conjecture; unrestricted permanent-on-top is
already false. The structural synthesis can also use this enlarged
row domain as one permitted piece in an explicitly allocated
positive semidefinite one-point gluing construction.


% Integrated synthesis: structural
\section{Mixed permanent certificates with an explicit diagonal budget}
\label{syn:struct:section}

The structural and row-local results can be combined at a cut vertex,
provided the shared diagonal is allocated before testing the local
conditions.  This qualification matters: a Gram splitting preserves
positive semidefiniteness, but need not preserve a normalized row bound
or simultaneous positivity of a reflected companion.  We give a
checkable allocation criterion, a quantitative loss formula, and an
equality statement for the resulting mixed class.  These are consequences
of the proved constituent certificates, not a solution of unrestricted
permanent-on-top or of the general Chollet problem.

For a Hermitian matrix $D$ write $P(D)=\operatorname{per}D$ and
$h(D)=\prod_i d_{ii}$, with empty products and empty permanents equal to
one.  Its Schur power is
\[
 \mathcal S(D)_{\sigma,\tau}=\prod_i d_{\sigma(i),\tau(i)}.
\]
Call $D$ \emph{hereditarily POT} if every principal submatrix $F$ satisfies
$0\preceq\mathcal S(F)\preceq P(F)I$.
Set $\mathcal R(D)=2\operatorname{diag}(D)-D$.

The inputs used below are explicit.  The structural theorem
\ref{R3E5:thm:main} covers PSD matrices whose support blocks are cycles,
phase-balanced blocks, or bipartite blocks with no edge vertex-disjoint
from a simple cycle.  Its operator gluing identity is
Lemma~\ref{R3E5:lem:glue}.  The reflected-companion theorem
\ref{R4E3:reflection} supplies the upper bound whenever
$\mathcal R(D)\succeq0$ and the support has no edge vertex-disjoint from
any simple cycle; imposing $D\succeq0$ as well supplies the lower bound.
The improved row-local theorem~\ref{syn:op:radius} supplies full POT at
the normalized absolute row bound $1/23$.  All three properties are
hereditary.  The common cactus result of Round~2 entries~03 and~05 is
one predecessor result, with both proofs retained in this volume.

\subsection{An allocation criterion}

Partition an index set as $\{v\}\sqcup U_1\sqcup\cdots\sqcup U_k$ and
consider the Hermitian matrix
\begin{equation}\label{syn:struct:matrix}
 A=\begin{pmatrix}
 a&z_1^*&z_2^*&\cdots&z_k^*\\
 z_1&B_1&0&\cdots&0\\
 z_2&0&B_2&\cdots&0\\
 \vdots&\vdots&\vdots&\ddots&\vdots\\
 z_k&0&0&\cdots&B_k
 \end{pmatrix},\qquad
 A_j(t)=\begin{pmatrix}t&z_j^*\\z_j&B_j\end{pmatrix}.
\end{equation}
Assume each $B_j$ has positive diagonal and $z_j\ne0$; components with
$z_j=0$ can instead be treated as separate direct summands.  For each
$j$ choose one of the following three certificates.

\begin{enumerate}
\item \emph{Structural certificate.}
Assume $B_j\succeq0$, $z_j\in\operatorname{ran}B_j$, and that the
off-diagonal support and phases of $A_j(t)$ satisfy the structural
block hypothesis stated above.  Put
\begin{equation}\label{syn:struct:costS}
 \tau_j=z_j^*B_j^{\dagger}z_j.
\end{equation}
Here $B_j^{\dagger}$ is the Moore--Penrose inverse.

\item \emph{Reflected-companion certificate.}
Assume $B_j\succeq0$, $\mathcal R(B_j)\succeq0$, and
$z_j\in\operatorname{ran}B_j\cap\operatorname{ran}\mathcal R(B_j)$.
Assume the support of $A_j(t)$ has no edge vertex-disjoint from any
simple cycle.  Put
\begin{equation}\label{syn:struct:costR}
 \tau_j=\max\{z_j^*B_j^{\dagger}z_j,
              z_j^*\mathcal R(B_j)^{\dagger}z_j\}.
\end{equation}

\item \emph{Row-local certificate.}
Write $b_i=(B_j)_{ii}>0$ for $i\in U_j$ and set $q_0=1/23$.  Define
\[
 r_i=\sum_{\substack{\ell\in U_j\\\ell\ne i}}
       \frac{|(B_j)_{i\ell}|}{\sqrt{b_i b_\ell}},
 \qquad s_j=\sum_{i\in U_j}\frac{|(z_j)_i|}{\sqrt{b_i}}.
\]
Assume $r_i<q_0$ for every $i$, and put
\begin{equation}\label{syn:struct:costW}
 \tau_j=\max\left\{\left(\frac{s_j}{q_0}\right)^2,
       \max_{i\in U_j}
       \frac{|(z_j)_i|^2}{b_i(q_0-r_i)^2}\right\}.
\end{equation}
No separate PSD assumption on $B_j$ is needed in this case.
\end{enumerate}

\begin{theorem}[Mixed allocation and full POT]
\label{syn:struct:allocation}
If the chosen finite costs satisfy $\sum_j\tau_j\le a$, then the matrix
$A$ in \eqref{syn:struct:matrix} is PSD and hereditarily POT.  In
particular, for every same-order PSD matrix $D$,
\begin{equation}\label{syn:struct:BS}
 0\le P(A\circ D)\le P(A)h(D)\le P(A)P(D),
 \qquad P(A\circ\overline A)\le h(A)P(A).
\end{equation}
More generally, any finite iterated one-point sum of hereditarily POT
pieces is hereditarily POT, provided that each piece is certified after
its own share of every common diagonal has been specified.
\end{theorem}

\begin{proof}
Choose $t_j\ge\tau_j$ with $\sum_jt_j\le a$ and write
$t_0=a-\sum_jt_j\ge0$.  Each $\tau_j$ is positive because $z_j\ne0$.
For a PSD matrix $B$ and $z\in\operatorname{ran}B$, the generalized
Schur complement says
\[
 \begin{pmatrix}t&z^*\\z&B\end{pmatrix}\succeq0
 \quad\Longleftrightarrow\quad t\ge z^*B^{\dagger}z.
\]
For completeness, write $z=B^{1/2}w$ with
$w=B^{\dagger/2}z$; completing the square gives
$\|B^{1/2}y+xw\|^2+(t-\|w\|^2)|x|^2$.
Thus the structural cost makes $A_j(t_j)$ PSD and the structural
theorem makes it hereditarily POT.  The same calculation applied also
to $\mathcal R(B_j)$ shows that the reflected cost makes both
$A_j(t_j)$ and $\mathcal R(A_j(t_j))$ PSD.  The sign reversal of $z_j$
does not change its quadratic cost.  The reflected-companion theorem
and inheritance under principal deletion give hereditary POT.

For the row-local cost, the normalized root row sum is $s_j/\sqrt{t_j}$,
and the normalized row sum at $i\in U_j$ is
$r_i+|(z_j)_i|/\sqrt{b_i t_j}$.  Both are at most $q_0$ by
\eqref{syn:struct:costW}.  A Hermitian matrix with diagonal one and
absolute off-diagonal row sums at most $q_0<1$ is positive definite:
its quadratic form is at least $(1-q_0)\|x\|^2$.
Theorem~\ref{syn:op:radius} therefore applies.  Deleting indices can
only reduce these row sums, so it also proves hereditary POT.

Here is the operator gluing calculation, including the residual
diagonal $t_0$.  Regard permutations of a subset $V$ as permutations
of the whole index set fixing its complement, and write
\[
 \widehat{\mathcal S}_V(F)=
 \sum_{\pi\in S_V}\left(\prod_{i\in V}f_{\pi(i),i}\right)L_\pi
\]
on the full regular representation.  Restriction to subgroup cosets
makes this a direct sum of copies of $\mathcal S(F)$, so it has the
same lower and upper operator bounds.  Operators supported on disjoint
subsets commute.  A permutation cycle cannot visit two different
$U_j$'s, since it would have to pass through $v$ twice.  Expanding
according to the cycle containing $v$, and splitting its fixed-point
weight as $a=t_0+\sum_jt_j$, gives
\begin{align}
 \mathcal S(A)
 &=t_0\prod_j\widehat{\mathcal S}(B_j)
   +\sum_j\widehat{\mathcal S}(A_j(t_j))
                      \prod_{\ell\ne j}\widehat{\mathcal S}(B_\ell),
       \label{syn:struct:operator}\\
 P(A)&=t_0\prod_j P(B_j)
   +\sum_j P(A_j(t_j))\prod_{\ell\ne j}P(B_\ell).
       \label{syn:struct:permanent}
\end{align}
Within each product the supports are disjoint.  Products of commuting
PSD operators bounded by their scalar permanents are bounded by the
corresponding product of permanents.  Hence
\eqref{syn:struct:operator} lies between zero and
\eqref{syn:struct:permanent} times the identity.
The matrix $A$ itself is a sum of the embedded PSD pieces and the
nonnegative scalar at $v$, hence is PSD.

For a principal submatrix retaining $v$, repeat this calculation with
the restricted pieces.  For one deleting $v$, use the direct sum of
principal submatrices of the $B_j$.  This proves the hereditary claim
and also proves the iterated assertion by induction.
Finally, Schur multiplying $P(A)I-\mathcal S(A)\succeq0$ by
$\mathcal S(D)\succeq0$ and testing on the constant vector gives
$P(A\circ D)\le P(A)h(D)$.  The permanent Hadamard inequality for PSD
matrices gives $h(D)\le P(D)$, and $D=\overline A$ gives the last
assertion of \eqref{syn:struct:BS}.
\end{proof}

The distinction between \eqref{syn:struct:costR} and two separate
global PSD tests is essential.  The simultaneous-allocation example
in Round~4 entry~03 has a root diagonal $2/5$, four root edges $1/4$,
leaf diagonals one, and two disjoint leaf edges $1/2$ and $-1/2$.
Both global matrices $A$ and $\mathcal R(A)$ are positive definite,
but the two rooted triangles each require a simultaneous cost $1/4$.
Their total $1/2$ exceeds $2/5$.  Our hypothesis explicitly checks this
budget and makes no assertion that ordinary Gram splitting preserves
all three local certificates.

\subsection{Quantitative loss and equality after gluing}

\begin{theorem}[Transfer of local deficits]
\label{syn:struct:deficits}
Let $A$ be an allocated one-point sum as in
Theorem~\ref{syn:struct:allocation}, and let $H$ be any correlation
matrix on its full index set.  Put $V_j=\{v\}\cup U_j$ and
\[
 P_j=P(A_j(t_j)),\quad R_j=P(B_j),\quad
 \widetilde P_j=P(A_j(t_j)\circ H[V_j]),\quad
 \Delta_j=P_j-\widetilde P_j.
\]
Then
\begin{equation}\label{syn:struct:loss}
 P(A)-P(A\circ H)
 \ge\sum_j\Delta_j\prod_{\ell\ne j}R_\ell.
\end{equation}
Consequently every proved local lower bound for $\Delta_j$ can be
inserted simultaneously in this formula.  If all pieces have positive
diagonal and connected support and each satisfies
\begin{equation}\label{syn:struct:strictlocal}
 P(A_j(t_j)\circ K)=P_j
 \quad\Longleftrightarrow\quad \operatorname{rank}K=1
 \quad\text{for correlation matrices }K,
\end{equation}
then
\begin{equation}\label{syn:struct:equality}
 P(A\circ H)=P(A)
 \quad\Longleftrightarrow\quad \operatorname{rank}H=1.
\end{equation}
\end{theorem}

\begin{proof}
Set $\widetilde R_j=P(B_j\circ H[U_j])$.
The local POT inequalities imply
$0\le\widetilde P_j\le P_j$ and
$0\le\widetilde R_j\le R_j$.
Apply \eqref{syn:struct:permanent} to $A$ and to $A\circ H$; the
residual root weight stays $t_0$ because $h_{vv}=1$.
For each summand,
\[
 P_j\prod_{\ell\ne j}R_\ell
 -\widetilde P_j\prod_{\ell\ne j}\widetilde R_\ell
 =\Delta_j\prod_{\ell\ne j}R_\ell
 +\widetilde P_j\left(\prod_{\ell\ne j}R_\ell
              -\prod_{\ell\ne j}\widetilde R_\ell\right).
\]
All omitted terms, including the residual-root term, are nonnegative.
This proves \eqref{syn:struct:loss}.
Positive diagonal and PSD give $R_j\ge h(B_j)>0$.
Thus equality forces $\Delta_j=0$ for every $j$ and hence every
$H[V_j]$ has rank one.  In a Gram representation of the correlation
matrix $H$, all unit vectors in $V_j$ are collinear with the common
vector at $v$.  Therefore all vectors are collinear and $H$ has rank
one.  Conversely a rank-one correlation mask changes $A$ only by a
diagonal unitary congruence, preserving its permanent.
\end{proof}

Several exact local certificates can now be mixed in a single bound.
For a piece with positive diagonal write $C_j$ for its correlation
normalization, so that
$\Delta_j=h(A_j(t_j))[P(C_j)-P(C_j\circ H[V_j])]$.
For a cycle of length at least four, the structural scalar certificate
\eqref{R3E5:eq:scalar} gives
\begin{equation}\label{syn:struct:cycleloss}
 P(C_j)-P(C_j\circ H[V_j])
 \ge2\sum_{\substack{M\text{ a matching}\\|M|=2}}
       \prod_{e\in M}|(C_j)_e|^2
       \left(1-\prod_{e\in M}|H_e|^2\right).
\end{equation}
For a reflected-companion book with two hubs and independent leaves,
write its normalized hub-to-leaf entries as $u_p,v_p$.
Corollary~\ref{R4E3:equality} gives
\begin{equation}\label{syn:struct:bookloss}
 P(C_j)-P(C_j\circ H[V_j])
 \ge2\sum_{p\ne q}|u_p|^2|v_q|^2
       (1-|H_{1p}|^2|H_{2q}|^2).
\end{equation}
The indices $p,q$ in the mask denote the corresponding leaf vertices.
For a connected row-local piece of order $m_j$, let $\delta_j>0$ be
the least positive eigenvalue of
\[
 \sum_{r<s}|(C_j)_{rs}|^2(I-L_{(rs)}).
\]
The projector-mask consequence, Theorem~\ref{syn:op:mask}, gives
\begin{equation}\label{syn:struct:weakloss}
 \frac{\Delta_j}{P_j}
 \ge (1-e^{-\delta_j/200})
       \left(1-\frac{P(H[V_j])}{m_j!}\right).
\end{equation}
All these bounds can be substituted in \eqref{syn:struct:loss} at
once.  Equivalently, with
$w_j=P_j\prod_{\ell\ne j}R_\ell/P(A)$, the normalized loss is at
least $\sum_j w_j\Delta_j/P_j$; the nonnegative weights sum to at
most one, the remaining weight corresponding to $t_0$.

For a cycle of length at least four with every edge nonzero, every
edge has a disjoint edge, so vanishing of
\eqref{syn:struct:cycleloss} forces $|H_e|=1$ on every edge.
The same conclusion follows from \eqref{syn:struct:bookloss} for a
book with at least two leaves and all hub-to-leaf edges nonzero.
In each case connectivity then forces rank one.  A connected
phase-balanced piece also has property \eqref{syn:struct:strictlocal}:
after gauging it to an entrywise nonnegative matrix, each permutation
weight is nonnegative, the real part of each mask monomial is at most
one, and the positive transposition weights force $|H_e|=1$ on every
support edge when the loss is zero.
Finally, \eqref{syn:struct:weakloss} gives the property for a connected
row-local piece.  Indeed $P(K)\le m!$ for every order-$m$ correlation
matrix $K$, and equality forces all transposition monomials
$|k_{rs}|^2$ to equal one, hence $\operatorname{rank}K=1$.
Thus these strict piece types may be combined freely in
\eqref{syn:struct:equality}.  No strictness assertion is needed for
other pieces in Theorem~\ref{syn:struct:allocation}.

\subsection{A Hessian consequence on the weak pieces}

The scalar Hessian theorem of Round~7 entry~05 can also be transported
through the exact gluing formula.  This yields convexity in selected
matrix variables on a mixed class that need not have a small global
row sum.

\begin{theorem}[Strict convexity in the weak-piece variables]
\label{syn:struct:hessian}
Fix an allocated one-point decomposition with $t_0\ge0$ and all local
diagonals positive.  Hold the structural and reflected-companion
pieces fixed.  Let a nonempty set $W$ of row-local pieces vary, with
their individual diagonals fixed and their correlation normalizations
having absolute off-diagonal row sums at most $1/23$.
Zero off-diagonal entries are permitted throughout this variable
domain.  Then $\log P(A)$ is strictly convex as a function of all
these weak-piece off-diagonal entries, viewed as real variables.

More precisely, write $D_j=\operatorname{diag}(A_j)$ for the fixed
positive diagonal matrix of a weak piece.  For Hermitian zero-diagonal
directions $K_j$ in its original entries, set
$H_j=D_j^{-1/2}K_jD_j^{-1/2}$.
Along the induced affine direction $K$ of the glued matrix,
\begin{equation}\label{syn:struct:hessianbound}
 D^2\log P(A)[K,K]\ge\frac7{10}\sum_{j\in W}
 \left[w_j\|H_j\|_F^2
       +(1-w_j)\|H_j[U_j]\|_F^2\right],
 \qquad
 w_j=\frac{P_j\prod_{\ell\ne j}R_\ell}{P(A)}>0.
\end{equation}
The norm in this formula is therefore the norm of the normalized
local direction, not that of the unnormalized $K_j$.
\end{theorem}

\begin{proof}
All pieces and deletions are PSD with positive diagonal, so the
permanent Hadamard inequality gives $P_j\ge h(A_j)>0$ and
$R_j\ge h(B_j)>0$.  In particular $P(A)>0$ by
\eqref{syn:struct:permanent}.  Put $f_j=\log P_j$ and
$g_j=\log R_j$.  That identity writes $\log P(A)$ as the logarithm
of a sum of exponentials of
\[
 F_j=f_j+\sum_{\ell\ne j}g_\ell\quad(1\le j\le k),
 \qquad F_0=\log t_0+\sum_\ell g_\ell.
\]
The term $F_0$ is omitted if $t_0=0$; no logarithm of zero is taken.
Let $\omega_\alpha=e^{F_\alpha}/\sum_\beta e^{F_\beta}$ be the
weights of the retained terms, so $\omega_j=w_j>0$ and their sum
is one.  Along any real affine line, direct differentiation gives
\[
 \bigl(\log\sum_\alpha e^{F_\alpha}\bigr)''
 =\sum_\alpha\omega_\alpha F_\alpha''
   +\sum_\alpha\omega_\alpha
       \left(F_\alpha'-\sum_\beta\omega_\beta F_\beta'\right)^2
 \ge\sum_\alpha\omega_\alpha F_\alpha''.
\]
Theorem~\ref{R7E5:thm:main} gives Hessian lower bound
$(7/10)\|H\|_F^2$ for $\log\operatorname{per}$ on the
correlation row ball of radius $1/12$.  It applies both to each weak
piece here and to its root deletion, because $1/23<1/12$ and
principal deletion reduces row sums.  Fixed diagonal congruence
changes the logarithm of a permanent only by a constant.  Hence
\[
 f_j''\ge\tfrac7{10}\|H_j\|_F^2,\qquad
 g_j''\ge\tfrac7{10}\|H_j[U_j]\|_F^2\qquad(j\in W).
\]
The function $f_j$ occurs only in $F_j$, while $g_j$ occurs in every
other retained term.  Summing with the weights therefore gives
\eqref{syn:struct:hessianbound}.  A nonzero combined direction has
$H_j\ne0$ for some weak piece, and $w_j>0$, so its lower bound is
strictly positive.  The product of the local row balls is convex;
integrating this positive second derivative along a nonconstant
segment proves strict convexity.  Neither a small global row sum
nor a Hessian claim in directions changing the fixed strong pieces
is used.
\end{proof}

There is a uniform consequence when the weak pieces' entries incident
to $v$ are also fixed.  Then $H_j$ has zero root row and column, so
the two norms in \eqref{syn:struct:hessianbound} agree and the weights
cancel.  The Hessian lower bound becomes
$(7/10)\sum_{j\in W}\|H_j\|_F^2$, independently of the number or
strength of the fixed pieces.  In particular, for two such matrices
$A^{(0)},A^{(1)}$, let $C_j^{(0)},C_j^{(1)}$ be their normalized
weak pieces and $A^{(t)}=(1-t)A^{(0)}+tA^{(1)}$.  Integration gives
\begin{equation}\label{syn:struct:interiorjensen}
 (1-t)\log P(A^{(0)})+t\log P(A^{(1)})-\log P(A^{(t)})
 \ge\frac7{20}t(1-t)\sum_{j\in W}
       \|C_j^{(1)}-C_j^{(0)}\|_F^2.
\end{equation}
This is a dimension-independent strong-convexity estimate for the weak
interior entries.  For unrestricted weak-piece directions, the
weights in \eqref{syn:struct:hessianbound} are retained; no universal
positive lower bound for all root-edge weights is claimed.

\subsection{A class beyond the separate global hypotheses}

Let $W$ be the order-five correlation matrix with all upper-triangular
off-diagonal entries $\varepsilon=1/256$, except
$w_{12}=i\varepsilon$, and lower-triangular entries given by
Hermitian symmetry.  Its normalized row sums are $4\varepsilon=1/64$,
so it is positive definite and row-locally POT.  It is not
phase-balanced: the triangle on $1,2,3$ has a nonreal cycle product.
Its support $K_5$ is neither a cycle nor bipartite, and has an edge
vertex-disjoint from a triangle.  Thus this block is outside both
structural support hypotheses used above.

Let $C$ be the order-four correlation matrix on a cycle, with each
edge of modulus $1/3$ and one directed edge equal to $i/3$.
It is positive definite by strict diagonal dominance and is a
structurally certified frustrated cycle.  Glue $k\ge1$ copies of
$W$ and one copy of $C$ at a common vertex, adding their original
diagonals there.  The resulting matrix $A_k$ has order $4k+4$ and
root diagonal $k+1$.  Each allocation is exactly one, so
Theorem~\ref{syn:struct:allocation} proves hereditary full POT and
\eqref{syn:struct:BS}, and
Theorem~\ref{syn:struct:deficits} gives the mixed quantitative bound
and rank-one equality characterization.

In the correlation normalization of $A_k$, the cycle vertex opposite
the root has absolute row sum $2/3$, already outside the global
$1/23$ row-local theorem.  The root row sum is
\[
 \frac{4k\varepsilon+2/3}{\sqrt{k+1}},
\]
which tends to infinity.  Moreover, the six edges among the four
nonroot vertices of each copy of $W$ retain modulus $\varepsilon$,
so the total squared edge mass is at least $6k\varepsilon^2$ and
also tends to infinity.  The nonbalanced $K_5$ blocks prevent the
global structural block theorem from applying; they also violate
the reflected theorem's edge--cycle support condition.  This gives
a concrete unbounded-order mixed family beyond those separate global
criteria.  It establishes an additional certified class and loss
formula, without claiming an optimal class or historical priority.
With the cycle held fixed, Theorem~\ref{syn:struct:hessian} also
proves strict convexity of $\log P(A_k)$ in the entries of the weak
blocks throughout their local row balls, despite these unbounded
global quantities.


% Integrated synthesis: special
\section{Special functions and Gaussian columns: two further consequences}
\label{syn:spec:section}

The four contributions in this branch address two distinct problems.
Chapters~\ref{source:R1E4} and~\ref{source:R2E4} study a logarithmic gamma
quotient; Chapters~\ref{source:R1E3} and~\ref{source:R7E3} study
Gaussian-binomial columns.  Their combination does not identify
these problems with one another.  It does give two useful extensions: the
concavity and exact transition regularity of an optimal gamma exponent, and
strict total positivity of all positive reciprocal powers of the first four
normalized Gaussian iterates.  The proofs below are deductions from the
structural theorems proved in those contributions.  No historical priority is
claimed for these deductions.

\subsection{A concave optimal exponent with an exact transition}

For $a>0$ and $x>1$, write
\[
 F(a,x)=\frac{\log\Gamma(x+1)}{\log((x^2+a)/(x+a))},\qquad
 M(a)=\inf_{x>1}F(a,x),\qquad p=1-\gamma.
\]
Thus $M(a)$ is the largest exponent for which
\[
 \Gamma(x+1)\geq\left(\frac{x^2+a}{x+a}\right)^{M(a)}\quad(x>1).
\]
The classification in Chapter~\ref{source:R2E4} uses
\[
 \Delta=3-\zeta(2)-2\gamma>0,\qquad
 a_* =\frac{\zeta(2)-\gamma}{\Delta},
\]
and proves that $M(a)=p(a+1)$ for $a\leq a_*$.  For $a>a_*$,
the infimum is attained at a unique $b_a>1$, with
$F_x(a,b_a)=0$ and $F_{xx}(a,b_a)>0$; $b_a$ is analytic and increasing.
Put
\[
 R_* =\frac{4p}{a_*+1}+6(\zeta(2)-1)
       -2(a_*+1)(\zeta(3)-1)>0.
\]
The endpoint identities proved there are
\begin{equation}\label{syn:spec:endpoint}
 F(a,1)=p(a+1),\quad F_x(a,1)=\frac{\Delta}{2}(a_*-a),\quad
 F_{xx}(a_*,1)=\frac{R_*}{3}.
\end{equation}
Here $F$ extends real analytically to a neighborhood of $(a_*,1)$:
the numerator and denominator have simple zeros at $x=1$.

\begin{theorem}[Concavity and transition of the best exponent]
\label{syn:spec:gamma}
The function $M$ is positive, strictly increasing, concave, and continuously
differentiable on $(0,\infty)$.  It is affine on $(0,a_*]$, analytic and
strictly concave on $(a_*,\infty)$, and is not twice differentiable at $a_*$.
More precisely, as $\varepsilon\downarrow0$,
\begin{equation}\label{syn:spec:Mexpansion}
 M(a_*+\varepsilon)
 =p(a_*+1+\varepsilon)
   -\frac{3\Delta^2}{8R_*}\varepsilon^2+O(\varepsilon^3).
\end{equation}
Consequently
\[
 M'(a_*\pm)=p,\qquad M''(a_*-)=0,\qquad
 M''(a_*+)=-\frac{3\Delta^2}{4R_*}.
\]
On $(a_*,\infty)$ the positive derivative $M'$ decreases strictly to zero.
\end{theorem}

\begin{proof}
Fix $x>1$, set $d=x^2-x>0$, and write $t=(a+x)/d$.  Then
\[
 F(a,x)=\log\Gamma(x+1)\,h(t),\qquad
 h(t)=\frac{1}{\log(1+1/t)}.
\]
Since $\Gamma(x+1)>1$ for $x>1$, differentiation shows $F_a>0$ and
\[
 h''(t)=
 \frac{2-(2t+1)\log(1+1/t)}
 {t^2(t+1)^2\log^3(1+1/t)}<0.
\]
Indeed, for $u>0$,
\[
 \log(1+u)>\frac{2u}{2+u},
\]
because the difference vanishes at zero and its derivative is
$u^2/((1+u)(2+u)^2)>0$.  Thus $F_{aa}(a,x)<0$ for every $x>1$.
An infimum of concave functions is concave: for $0<\lambda<1$,
\begin{align*}
 M(\lambda a+(1-\lambda)b)
 &\geq\inf_{x>1}\{\lambda F(a,x)+(1-\lambda)F(b,x)\}\\
 &\geq\lambda M(a)+(1-\lambda)M(b).
\end{align*}

For $a>a_*$ the analytic minimizer and the envelope identity give
\begin{equation}\label{syn:spec:envelope}
 M'(a)=F_a(a,b_a)>0,\qquad
 M''(a)=F_{aa}(a,b_a)
       -\frac{F_{ax}(a,b_a)^2}{F_{xx}(a,b_a)}<0.
\end{equation}
The second formula follows by differentiating $F_x(a,b_a)=0$ and then
$M'(a)=F_a(a,b_a)$.  In particular, the strict concavity is not inferred
merely from taking an infimum of strictly concave functions.

For the transition, put $y=b_{a_*+\varepsilon}-1$.  The implicit function
theorem and \eqref{syn:spec:endpoint} give
\[
 y=\frac{3\Delta}{2R_*}\varepsilon+O(\varepsilon^2).
\]
Taylor expansion at $x=1$ yields
\begin{align*}
 F(a_*+\varepsilon,1+y)
 &=p(a_*+1+\varepsilon)-\frac{\Delta}{2}\varepsilon y
   +\frac{R_*}{6}y^2+O(\varepsilon^3)\\
 &=p(a_*+1+\varepsilon)
   -\frac{3\Delta^2}{8R_*}\varepsilon^2+O(\varepsilon^3).
\end{align*}
All remainders come from analytic functions near $(a_*,1)$; hence they may
also be differentiated with the corresponding loss of order.  This proves
the stated one-sided derivatives and $C^1$ regularity.  Positivity and
strict increase across the transition follow from the displayed formulas
and \eqref{syn:spec:envelope}.

Finally $M(a)=O(\sqrt a\log a)$: evaluate its defining infimum at
$x=\sqrt a$ for $a>1$ and use Stirling's estimate.  Concavity implies
\[
 0<M'(a)\leq\frac{M(a)-M(a/2)}{a/2}
            \leq\frac{2M(a)}a\longrightarrow0.
\]
Together with strict concavity this proves the last assertion.
\end{proof}

\begin{corollary}[Interpolation of sharp gamma powers]
\label{syn:spec:jensen}
For $a,b>0$ and $0<\lambda<1$,
\[
 M(\lambda a+(1-\lambda)b)
 \geq\lambda M(a)+(1-\lambda)M(b).
\]
For distinct $a,b$, equality holds exactly when $a,b\leq a_*$.
Thus a convex combination of two established optimal exponents is a valid
lower bound for the optimal exponent at the interpolated parameter, with
a strict gain whenever the parameter interval enters $(a_*,\infty)$.
\end{corollary}
\begin{proof}
The inequality was proved above.  A concave differentiable function has
equality at an interior point of a chord only if it is affine on that
chord's entire interval.  Theorem~\ref{syn:spec:gamma} gives precisely the
claimed affine intervals.
\end{proof}

This completes an additional regularity question for the optimal exponent;
it is not a new solution of the fixed-$a=6$ source question.  The latter
already has a public proof claim identified in the Round~2 bibliography.
The threshold $a_*$ itself is also a previously published constant.

\subsection{A uniform reciprocal-moment theorem for Gaussian iterates}

Extend sequences by zero at negative indices and define
\[
 (\mathcal L a)_n=a_n^2-a_{n-1}a_{n+1},\qquad
 a_n(q)=\genfrac{[}{]}{0pt}{}{n+k}{k}_q\quad(n\geq0).
\]
Fix $k\geq3$ and $r\in\{0,1,2,3\}$.  Put $e_r=2^r-1$ and, for $q\ne1$,
\[
 \eta_r(q)=
 \begin{cases}e_r,&0<q<1,\\ k2^r-e_r,&q>1.\end{cases}
\]
Define the positive normalized iterate
\begin{equation}\label{syn:spec:A}
 A_{r,n}(q)=
 \begin{cases}
 q^{-\eta_r(q)n}(\mathcal L^r a(q))_n,&q\ne1,\\
 (\mathcal L^r a(1))_n,&q=1.
 \end{cases}
\end{equation}
We shall use the following explicit consequence of the rectangular
identities and fractional-shift factorization proved in
Chapter~\ref{source:R7E3}, especially its Theorem~\ref{R7E3:main}.

For $Q=\min(q,q^{-1})\leq1$, let
\[
 M_j=\{i+\ell-1:1\leq i\leq j,\ 1\leq\ell\leq k+1-j\},
 \qquad C=2M_1+M_2+M_4-\{1,k\}.
\]
Multisets are counted with multiplicity; $M_4$ is empty if $k=3$.
Write $[u]_Q=(1-Q^u)/(1-Q)$, with $[u]_1=u$, and set
$R=[3]_Q[k-2]_Q/[k]_Q$.  Let $\rho\in(0,1)$ be the unique solution of
\[
 Q^\rho+Q^{k+1-\rho}
 =1+Q^{k+1}-\frac{R}{R+1}(1-Q)(1-Q^k)
\]
when $Q<1$; at $Q=1$ use
$\rho(k+1-\rho)=Rk/(R+1)$.  Define
\begin{equation}\label{syn:spec:E}
 E_0=M_1,\quad E_1=M_2,\quad E_2=M_1+M_3,\quad
 E_3=C+\{\rho,k+1-\rho\}.
\end{equation}
Each $E_r$ is a nonempty multiset of positive real numbers.  Then
\begin{equation}\label{syn:spec:product}
 A_{r,n}(q)=\prod_{c\in E_r}\frac{[n+c]_Q}{[c]_Q}.
\end{equation}
For $q\leq1$ this follows directly from the first-three-iterate
factorizations, evaluating their positive scalar factors at $n=0$.
Indeed $(\mathcal L^r a)_0=1$ for every $r$.  For $q>1$ use
$a_n(q)=q^{kn}a_n(q^{-1})$ and the elementary identity
$\mathcal L(c^na_n)=c^{2n}(\mathcal L a)_n$; iteration gives exactly
the exponent in \eqref{syn:spec:A}.

\begin{theorem}[All positive powers, all Hankel orders]
\label{syn:spec:moments}
For every $q>0$, $k\geq3$, $r\in\{0,1,2,3\}$ and $\theta>0$, the sequence
\[
 b_n=A_{r,n}(q)^{-\theta}\quad(n\geq0)
\]
is the moment sequence of a probability measure $\mu_\theta$ on $[0,1]$
with infinite support in $(0,1]$.  These measures form a multiplicative
convolution semigroup in $\theta$.  Their Hankel matrices are strictly
totally positive: for every $m\geq1$ and strictly increasing nonnegative
integer sequences $i_1<\cdots<i_m$, $j_1<\cdots<j_m$,
\begin{equation}\label{syn:spec:minor}
 \det(b_{i_u+j_v})_{u,v=1}^m>0.
\end{equation}
Moreover the continuous interpolation
$b(x)=\prod_{c\in E_r}([c]_Q/[x+c]_Q)^\theta$ is completely monotone
and logarithmically completely monotone for $x\geq0$.
\end{theorem}

\begin{proof}
Suppose first that $0<Q<1$ and put $h=-\log Q$.  Expanding $-\log(1-z)$
in its absolutely convergent power series gives $b(x)=e^{-\theta\Phi(x)}$,
where
\begin{equation}\label{syn:spec:levy}
 \Phi(x)=\sum_{\ell=1}^{\infty}
 \frac{\sum_{c\in E_r}Q^{\ell c}}{\ell}(1-e^{-h\ell x}).
\end{equation}
The sum of its positive coefficients is finite.  Thus $b(x)$ is the
Laplace transform of a compound Poisson random variable $X_\theta$,
with jump sizes $h\ell$ and rates
$\theta\sum_c Q^{\ell c}/\ell$.  This assertion can alternatively be
read directly from the positive probability generating function
\begin{equation}\label{syn:spec:weights}
 \sum_{j\geq0}w_jz^j
 =\prod_{c\in E_r}\left(\frac{1-Q^c}{1-Q^cz}\right)^\theta,
 \qquad w_j>0,\qquad\sum_{j\geq0}w_j=1.
\end{equation}
The binomial series proves positivity and normalization.  Consequently
\[
 \mu_\theta=\sum_{j\geq0}w_j\delta_{Q^j}.
\]

If $Q=1$, each factor $(c/(x+c))^\theta$ is the Laplace transform of a
gamma random variable with shape $\theta$ and rate $c$.  Take independent
variables for the members of $E_r$, with multiplicity, and let $X_\theta$
be their sum.  Its density is positive on $(0,\infty)$, and the law of
$e^{-X_\theta}$ gives the required measure.  Equivalently,
\begin{equation}\label{syn:spec:gammalevy}
 \Phi(x)=\int_0^\infty(1-e^{-xt})
                  \frac{\sum_{c\in E_r}e^{-ct}}{t}\,dt.
\end{equation}
The identity follows by differentiation in $x$ and normalization at zero.
In both cases the exponent is additive in $\theta$, which proves the
convolution-semigroup assertion.  The representations show that
$(-1)^j b^{(j)}(x)\geq0$ for $j\geq0$ and
$(-1)^j(\log b)^{(j)}(x)\geq0$ for $j\geq1$; all differentiations are
justified by the finite exponential moments near $x=0$.

For completeness, the determinant assertion needs more than positive
semidefiniteness.  Expanding the determinant and symmetrizing the integrals
gives the continuous Cauchy--Binet identity
\[
 \det(b_{i_u+j_v})=
 \frac1{m!}\int_{[0,1]^m}
 \det(t_v^{i_u})_{u,v=1}^m\,
 \det(t_v^{j_u})_{u,v=1}^m\prod_{v=1}^m d\mu_\theta(t_v).
\]
For $0<t_1<\cdots<t_m$, both generalized Vandermonde determinants are
positive.  One elementary justification divides each alternant by the
ordinary Vandermonde determinant: the quotient is a Schur polynomial,
which has nonnegative integer coefficients and is nonzero at positive
arguments.  The integrand is symmetric, nonnegative everywhere, and
positive when all coordinates are distinct and positive.  Infinite
positive support supplies such tuples with positive product measure
(use disjoint neighborhoods of $m$ support points), proving
\eqref{syn:spec:minor}.
\end{proof}

\begin{corollary}[Infinite log-convexity of the reciprocals]
\label{syn:spec:infinite}
Define $(\mathcal D b)_n=b_nb_{n+2}-b_{n+1}^2$.  Every iterate
$\mathcal D^s b$, $s\geq0$, is a Hausdorff moment sequence of a finite
positive measure with infinite positive support.  In particular,
$(\mathcal D^s b)_n>0$ for every $s\geq1$ and $n\geq0$.
\end{corollary}
\begin{proof}
For any finite positive representing measure $\mu$ on $[0,1]$,
\[
 b_nb_{n+2}-b_{n+1}^2
 =\frac12\int\!\int (xy)^n(x-y)^2\,d\mu(x)d\mu(y).
\]
Thus push forward the finite positive measure
$\tfrac12(x-y)^2\,d\mu(x)d\mu(y)$ by $(x,y)\mapsto xy$.
Its support remains infinite in $(0,1]$: fix a positive support point
$x_0$ and vary $y$ through infinitely many other positive support points.
Every $x_0y$ is in the new support, as follows from neighborhoods separated
from the diagonal, and these products are distinct.  Induction proves
the assertion and strict positivity of all moments.
\end{proof}

\begin{corollary}[A sharp leading term for every Hankel determinant]
\label{syn:spec:asymptotic}
For $q\ne1$, let $Q=\min(q,q^{-1})$ and use the weights
\eqref{syn:spec:weights}.  For fixed $m\geq1$ put
\[
 H_m(n)=\det(b_{n+i+j})_{i,j=0}^{m-1},\qquad
 s_m=\frac{m(m-1)}2,
\]
\[
 C_m=\left(\prod_{j=0}^{m-1}w_j\right)
      \prod_{0\leq i<j<m}(Q^j-Q^i)^2>0.
\]
Then for every integer $n\geq0$,
\begin{equation}\label{syn:spec:bound}
 0\leq\frac{H_m(n)}{C_mQ^{ns_m}}-1
 \leq\left(\frac{H_m(0)}{C_m}-1\right)Q^n.
\end{equation}
In particular $H_m(n)\sim C_mQ^{ns_m}$ as $n\to\infty$.
\end{corollary}
\begin{proof}
Apply Cauchy--Binet first to finite atomic truncations, then pass to the
limit using entrywise convergence of moments and nonnegative summands:
\[
 H_m(n)=\sum_{0\leq j_1<\cdots<j_m}
 \left(\prod_{u=1}^m w_{j_u}\right)Q^{n\sum_u j_u}
 \prod_{u<v}(Q^{j_v}-Q^{j_u})^2.
\]
The unique smallest sum of indices is $s_m$, attained at
$(j_1,\ldots,j_m)=(0,\ldots,m-1)$.  Every other sum is at least $s_m+1$.
Their coefficients add to $H_m(0)-C_m$, so subtraction of the leading
term proves \eqref{syn:spec:bound}.
\end{proof}

\subsection{What is and is not closed by this synthesis}

The gamma theorem adds global concavity, sharp transition curvature, and
an interpolation principle for the optimal exponent.  The Gaussian theorem
adds arbitrary Hankel orders, an explicit determinant asymptotic, and
unlimited log-convexity depth after taking normalized reciprocal powers.
These are exact consequences with complete proofs, not experimental
extrapolations.

They do not settle coefficientwise infinite $q$-log-concavity of the
original Gaussian columns.  Round~1 gives coefficientwise certificates
for $k=3$ through depth six and $k=4,5,6$ through depth four.  Round~7 gives
fourth-fold positivity for every positive real $q$ and every column;
positive values for $q>0$ do not force nonnegative polynomial
coefficients.  The Hankel results concern a different operator
$\mathcal D$ applied to reciprocal powers, not further $\mathcal L$
iterates of the original sequence.  Also, the fourth normalized
interpolation polynomial can have nonreal roots, so the three-step
factorization is not an unrestricted root-preservation principle.
At $q=1$, infinite log-concavity of ordinary Pascal columns was already
proved by Br\"and\'en--Chasse; it must not be counted as a new closure.


% Integrated synthesis: geometry
% Integration fragment. All labels and bibliography keys use syn:geom:.
\section{Geometric synthesis: seed independence and functorial half-facet corrections}
\label{syn:geom:section}

The nine geometry manuscripts form two different chains. The first concerns symmetric lattice points in smooth reflexive polytopes; the second concerns barycentric inequalities for unimodular fibers. They do not solve one common conjecture. Within each chain, however, their conclusions combine into further precise results. We prove that the asymptotic tower minimum is independent of every fixed seed, that half-facet corrections commute exactly with coherent series extensions, and that the corrected inequalities extend to countable distributions in every nullity. We also obtain an ambient-rank-independent bound on the size of the correction formula in nullity three. These are deductions from the preserved manuscripts and classical inputs, not claims of independently established historical priority.

\subsection{One asymptotic Ewald rate for every fixed seed}
\label{syn:geom:seeds}
For a lattice polytope $B\subset\mathbb R^d$ put
\[
 \mathcal E(B)=B\cap(-B)\cap\mathbb Z^d,
 \qquad e(B)=|\mathcal E(B)|.
\]
A monotone polytope is a smooth reflexive lattice polytope. Fix such a seed $B$ of positive dimension $d$. At any surviving facet $\ell(x)\le1$, permit the attachment
\begin{equation}\label{syn:geom:attachment}
 S_{n,q}(B,\ell)=\{(x,y):x\in B,\ y_j\ge-1\ (1\le j\le n),\
             \textstyle\sum_jy_j+q\ell(x)\le1\},
 \qquad n\ge1,\quad 0\le q\le n,
\end{equation}
with integral $n,q$. Facets may be used repeatedly and the attachments may branch. Let $M_B(d+m)$ be the minimum Ewald count among these towers with $m$ added coordinates. Let $M_I(m+1)$ denote the corresponding minimum for the segment $I=[-1,1]$. These minima exist: there is at least one tower in every added dimension and the possible counts are positive integers.

\begin{theorem}[Independence of the fixed seed]\label{syn:geom:seedtheorem}
For every fixed positive-dimensional monotone seed $B$ and every integer $m\ge0$,
\begin{equation}\label{syn:geom:comparison}
 \frac{M_I(m+1)}3\ \le\ M_B(d+m)\ \le\ e(B)M_I(m+1).
\end{equation}
Consequently the limit
\[
 \lim_{D\to\infty}M_B(D)^{1/D}=\Lambda_*
\]
exists, is independent of $B$, and equals the segment-tower limit. In particular the inherited bounds are
\begin{equation}\label{syn:geom:window}
 \frac{203}{100}\le\Lambda_*
 \le\bigl(19195+4\sqrt{2297138}\bigr)^{1/14}
 =2.0627969707\ldots.
\end{equation}
The same limit holds if the seed may be chosen from any fixed nonempty finite collection of positive-dimensional monotone polytopes.
\end{theorem}

\begin{proof}
We first recall and justify the precise combinatorial data needed for grafting. An attachment creates facet heights
\[
 h_0=\sum_jy_j+q\ell(x),\qquad h_j=-y_j\quad(1\le j\le n),
 \qquad \sum_{j=0}^nh_j=q\ell(x).
\]
For an Ewald point all these heights belong to $\{-1,0,1\}$. Every assignment satisfying the displayed sum determines exactly one choice of the new coordinates. Each later attachment depends on its predecessors only through its parent facet height. Organize the construction by the birth of each facet. Descendants at different facet nodes use disjoint new coordinates; repeated direct attachments at a node form parallel collections.

For a descendant collection write $(u,v)$ for the number of completions at parent heights $1$ and $0$. Negating the descendant coordinates proves that the count at height $-1$ is also $u$. An empty collection has $(u,v)=(1,1)$. Parallel collections multiply both coordinates. A simplex node with child collections $(u_i,v_i)$ has
\begin{equation}\label{syn:geom:branch}
 u=[z^q]\prod_{i=0}^n(v_i+u_i(z+z^{-1})),\qquad
 v=[z^0]\prod_{i=0}^n(v_i+u_i(z+z^{-1})).
\end{equation}
Every collection has $0<u\le v$. To see this without an assumption about global facet slices, normalize each factor by $v_i$. A symmetric sequence $(c_k)$ decreasing for $k\ge0$, multiplied by $1+r(z+z^{-1})$ with $0\le r\le1$, has coefficients $(d_k)$ satisfying
\[
 d_k-d_{k+1}=r(c_{k-1}-c_{k+2})+(1-r)(c_k-c_{k+1})\ge0.
\]
For $k=0$ use $c_{-1}=c_1\ge c_2$. Induction proves the assertion, with positivity because $q\le n<n+1$. Parallel multiplication preserves it. This is the conditional-count lemma of Round 3, entry 06, restated to make the new grafting argument complete.

Take any tower $P$ over $B$. For each original facet of $B$, let $(u_j,v_j)$ be the collection of all its historical descendants. Put
\[
 U=\prod_j u_j,\qquad V=\prod_j v_j.
\]
The seed origin has every facet height zero, so exactly $V$ Ewald points of $P$ lie over it. In particular $e(P)\ge V$. Graft all these descendant collections, in parallel, onto the facet $x=1$ of a new segment. Every step is again an allowed operation \eqref{syn:geom:attachment}, with the same twist and number of new coordinates. The resulting segment tower $Q$ has dimension $m+1$ and
\[
 e(Q)=2U+V\le3V\le3e(P).
\]
Taking the minimum over $P$ proves the first inequality of \eqref{syn:geom:comparison}.

Conversely, take a segment tower $Q$ with $m$ added coordinates. Combine its two root-facet collections. Evenness in the parent height means that their conditional counts are again a pair $(U,V)$ with $U\le V$, and $e(Q)=2U+V$. Choose any facet of $B$ and graft both collections in parallel onto it. It does not matter that the original segment facets had opposite signs: the completion counts are even. Over each seed Ewald point the number of completions is either $U$ or $V$, hence is at most $V$. The resulting tower $P$ has
\[
 e(P)\le e(B)V\le e(B)e(Q).
\]
Minimizing over $Q$ proves the upper comparison.

For completeness, the segment-tower limit exists because that class is closed under Cartesian products: after constructing one factor, an $S_{1,0}$ attachment creates an independent segment, on which the second factor's history can be performed. Thus $M_I(a+b)\le M_I(a)M_I(b)$. The subadditivity lemma applied to $\log M_I$ gives a limit of $M_I(D)^{1/D}$. Constants $3,e(B)$ disappear after taking logarithms and dividing \eqref{syn:geom:comparison} by $d+m$; also $(m+1)/(d+m)\to1$. This proves seed independence. Round 3 supplies the lower bound in \eqref{syn:geom:window}; Round 2 supplies its explicit upper construction. Finally, the minimum of finitely many sequences with the same logarithmic rate has that rate, proving the last statement.
\end{proof}

The theorem closes an exact fixed-seed question within this construction class. It does not identify $\Lambda_*$, close the interval \eqref{syn:geom:window}, or give a lower bound for arbitrary monotone polytopes. Allowing growing-dimensional seeds changes the assertion and is not justified by the comparison.

\subsection{The half-facet framework and its source status}
\label{syn:geom:framework}
Let $M\in\mathbb Z^{r\times n}$ have full row rank, all nonsingular maximal minors equal to $\pm1$, and nullity $d=n-r\ge1$. Redundant equations can first be removed using the lattice normalization proved in the source manuscripts. For parameters $\theta=(b,\ell,u)$ with nonempty real fiber set
\[
 P_M(\theta)=\{x:Mx=b,\ \ell\le x\le u\}.
\]
All bounds are finite here. The nonnegative setting is also allowed when $M$ is pointed, meaning $\ker M\cap\mathbb R_{\ge0}^n=\{0\}$. Choose integral kernel coordinates, so that $\ker M\cap\mathbb Z^n$ becomes $\mathbb Z^d$. The half-relaxed fiber is
\[
 Q_M(\theta)=\{x:Mx=b,\ \ell-\tfrac12\mathbf1\le x\le u+\tfrac12\mathbf1\}.
\]
An inequality constant on the affine fiber can be omitted from this formula while retaining feasibility of the original parameters; on the domain being used this changes no points. In the nonnegative case the half-relaxation is $x_j\ge-1/2$. Use lattice-normalized volume and define, at feasible integral parameters,
\begin{equation}\label{syn:geom:Rdef}
 F_M(\theta)=\#(P_M(\theta)\cap\mathbb Z^n),\quad
 V_M(\theta)=\operatorname{vol}_d Q_M(\theta),\quad
 R_M(\theta)=F_M(\theta)-V_M(\theta).
\end{equation}

Round 7, entry 06, makes this correction explicit using the local Euler--Maclaurin formula of Berline--Vergne \cite{R7E6:BV}. It proves that $R_M$ is piecewise polynomial of degree at most $d-2$ for $d\ge2$, uniformly bounded by a constant times $\max(1,\operatorname{diam}Q_M)^{d-2}$ in fixed kernel coordinates. The correction is zero for $d=1$. Odd-codimension local constant terms vanish. In dimension three its formula is
\begin{equation}\label{syn:geom:edges}
 R_M=\frac1{24}\sum_{E\text{ edge of }Q_M}
 \ell_{\mathbb Z}(E)(a_E\cdot b_E)
 \left(\|a_E\|^{-2}+\|b_E\|^{-2}\right),
\end{equation}
where $a_E,b_E$ are the two primitive outward facet normals incident to the edge, expressed in the chosen kernel coordinates. Nonsimple vertices and parameters on chamber walls are included in that argument. Local coefficients depend on the chosen metric; their total in \eqref{syn:geom:Rdef} does not.

There is a relevant antecedent beyond the references printed in Round 7. Chu--Karshon \cite[Remark 1.1 and reference 5]{R5E3:CK} report an Euler--Maclaurin analogue for unimodular halfway-lattice polytopes with the Todd function replaced by the $\widehat A$ function. Their cited source is Natalia Juliana Lozano Duarte, \emph{Khovanskii--Pukhlikov formula for half-way polytopes}, M.Sc. thesis, supervised by Alexander Caviedes Castro, Pontificia Universidad Javeriana, 2025. This audit verified that attribution in Chu--Karshon, but did not obtain the thesis itself. Accordingly no historical novelty is claimed for the general halfway-lattice mechanism or parity phenomenon. The following results concern its consequences for the tournament's particular parameter families.

\subsection{Exact compatibility with coherent series extensions}
\label{syn:geom:series}
Fix positive integers $k_e$. Replace each variable $x_e$ by a chain $x_{e,1},\ldots,x_{e,k_e}$, with equations
\begin{equation}\label{syn:geom:serieseq}
 M(x_{e,1})_{e=1}^n=b,\qquad
 x_{e,j+1}-x_{e,j}=\kappa_{e,j}\quad(1\le j<k_e).
\end{equation}
Write $M^{\mathrm{ser}}$ for this fixed equation matrix. Let
\[
 s_{e,1}=0,\quad s_{e,j}=\sum_{h<j}\kappa_{e,h},\quad
 L_e=\max_j(\ell_{e,j}-s_{e,j}),\quad
 U_e=\min_j(u_{e,j}-s_{e,j}).
\]
Bounds and internal balances may all vary.

\begin{theorem}[Series pullback of the canonical correction]\label{syn:geom:seriespullback}
The matrix $M^{\mathrm{ser}}$ is unimodular and has the same nullity $d$ as $M$. At every feasible integral boxed parameter vector,
\begin{align}
 F_{M^{\mathrm{ser}}}(b,\kappa,\ell,u)&=F_M(b,L,U),\label{syn:geom:seriesF}\\
 V_{M^{\mathrm{ser}}}(b,\kappa,\ell,u)&=V_M(b,L,U),\label{syn:geom:seriesV}\\
 R_{M^{\mathrm{ser}}}(b,\kappa,\ell,u)&=R_M(b,L,U).\label{syn:geom:seriesR}
\end{align}
For nonnegative fibers put $L_e=\max_j(-s_{e,j})$. If $M$ is pointed, the analogous identities are
\begin{equation}\label{syn:geom:seriesnonnegative}
 F_{M^{\mathrm{ser}}}(b,\kappa)=F_M(b-ML),\qquad
 V_{M^{\mathrm{ser}}}(b,\kappa)=V_M(b-ML),\qquad
 R_{M^{\mathrm{ser}}}(b,\kappa)=R_M(b-ML).
\end{equation}
In particular an arbitrary subdivision requires no new local correction computation. A piecewise-affine nullity-three formula remains piecewise affine after refining by the finitely many max/min choices.
\end{theorem}

\begin{proof}
The difference equations give $x_{e,j}=y_e+s_{e,j}$, where $y_e=x_{e,1}$. They are independent, and their common homogeneous solution space consists of chains constant along each $e$. Hence $M^{\mathrm{ser}}$ has rank $r+\sum_e(k_e-1)$ and nullity $d$.

Here is also a maximal-minor verification of unimodularity. A square maximal column submatrix uses all equation rows. If fewer than $k_e-1$ columns are selected from one chain, its $k_e-1$ difference rows force singularity. Otherwise each chain contributes either $k_e-1$ or $k_e$ columns; precisely $r$ chains contribute all their columns. For a chain with $k_e-1$ selected columns, its difference-row block is a tree-incidence minor of determinant $\pm1$. Eliminate that block. For a fully selected chain, use its difference equations to leave one surviving variable; its column in the old rows is the original column $M_e$. The remaining determinant is a maximal minor of $M$, up to sign. Thus every nonsingular maximal minor is $\pm1$.

The map $y\mapsto(x_{e,j})=(y_e+s_{e,j})$ is an affine bijection from the base fiber $My=b$, $L\le y\le U$ to the series fiber. For integral parameters it maps the two affine lattices bijectively, proving \eqref{syn:geom:seriesF}. On the homogeneous spaces its linear part repeats coordinates along chains and identifies the kernel lattices bijectively. It therefore preserves lattice-normalized volume, even though it need not preserve Euclidean volume in their original ambient spaces.

Half-relaxing every chain bound gives precisely
\[
 \max_j(\ell_{e,j}-s_{e,j}-\tfrac12)=L_e-\tfrac12,
 \qquad
 \min_j(u_{e,j}-s_{e,j}+\tfrac12)=U_e+\tfrac12.
\]
Thus the same map identifies the two half-relaxed bodies. This proves \eqref{syn:geom:seriesV}; subtracting gives \eqref{syn:geom:seriesR}. No invariance assertion for individual metric-dependent local coefficients is needed.

For nonnegative fibers, $y_e\ge L_e$. Translate by the integral vector $L$ to obtain $z=y-L\ge0$, $Mz=b-ML$. Half-relaxation gives $z_e\ge-1/2$ exactly. The same lattice and volume arguments prove \eqref{syn:geom:seriesnonnegative}. Pointedness is preserved because a homogeneous chain is constant and nonnegative precisely when its base value is nonnegative. Finally max/min functions have finitely many affine pieces; composition with the base correction proves the claimed degree statement.
\end{proof}

In the $A_3$ case, Round 6, entry 03, provides a concrete function $R(a,b,c)$ satisfying
\[
 K(a,b,c)-R(a,b,c)=\operatorname{vol}P(a+\tfrac32,b+2,c+\tfrac32).
\]
Theorem~\ref{syn:geom:seriespullback} transfers that \emph{same} correction, evaluated at the effective margin $b-ML$, to every coherent series extension, with all internal balances variable. This is stronger bookkeeping than merely reapplying a general existence theorem to the larger matrix, while preserving the distinction between corrected cube-root concavity and uncorrected log-concavity.

\subsection{Countable barycenters in every nullity}
\label{syn:geom:countable}
The following passage to countable support combines the conditional-tail argument of Round 5, entry 03, with the all-nullity volume extension of Round 7. It does not require a separate count-integrability hypothesis.

\begin{theorem}[Countable corrected Brunn--Minkowski]\label{syn:geom:countablethm}
Fix a unimodular $M$ of nullity $d\ge1$, with finite variable boxes, or with nonnegative fibers and $M$ pointed. Let $\theta_i$ be a finite or countable family of feasible integral parameter vectors and let $\mu_i>0$, $\sum_i\mu_i=1$. Assume
\[
 \sum_i\mu_i\|\theta_i\|<\infty,
 \qquad \bar\theta=\sum_i\mu_i\theta_i\text{ is integral}.
\]
Then the central fiber is nonempty and
\begin{equation}\label{syn:geom:countableBM}
 \bigl(F_M(\bar\theta)-R_M(\bar\theta)\bigr)^{1/d}
 \ge\sum_i\mu_i\bigl(F_M(\theta_i)-R_M(\theta_i)\bigr)^{1/d}.
\end{equation}
The series on the right converges to a finite value. The assertion applies in particular to all the series systems of Theorem~\ref{syn:geom:seriespullback}, including variable capacities and internal balances.
\end{theorem}

\begin{proof}
Let $\mathcal D$ be the real feasible parameter set. It is a closed convex polyhedron: it is the projection of the finite system consisting of the equalities and bounds defining the fiber. In the nonnegative case it is the finitely generated column cone of $M$. The expectation of an integrable random vector in $\mathcal D$ remains in $\mathcal D$, as follows by taking expectations in each of its defining affine inequalities and equalities. Thus $\bar\theta\in\mathcal D$.

Integral feasibility implies an integer point here. Indeed normalize a maximal minor to obtain a kernel parametrization $x=x^0(b)+Kt$ with $x^0$ integral on integral margins and the nonzero maximal minors of $K$ equal to $\pm1$. A vertex of the nonempty bounded feasible fiber is obtained from independent active integer bounds and is integral. This also applies to a pointed nonnegative fiber. Hence $F_M(\bar\theta)>0$.

On all real feasible parameters define
\[
 G(\theta)=\operatorname{vol}_dQ_M(\theta)^{1/d}.
\]
It is finite and nonnegative. The affine dependence of the equations and relaxed inequalities implies Minkowski containment under averaging, in fixed kernel coordinates. Classical Brunn--Minkowski gives concavity of $G$ on $\mathcal D$. In fact each half-relaxed body has positive $d$-volume, since any point of the original fiber strictly satisfies all inequalities that are nonconstant on the fiber.

For a genuinely infinite support, set $q_N=\sum_{i>N}\mu_i$. Its real conditional-tail mean
\[
 \tau_N=q_N^{-1}\sum_{i>N}\mu_i\theta_i
\]
exists by the first-moment hypothesis and belongs to $\mathcal D$. Finite Jensen, applied to $\theta_1,\ldots,\theta_N,\tau_N$, gives
\[
 G(\bar\theta)\ge\sum_{i\le N}\mu_iG(\theta_i)+q_NG(\tau_N)
 \ge\sum_{i\le N}\mu_iG(\theta_i).
\]
The partial sums increase and are bounded by the finite number $G(\bar\theta)$. Taking their limit proves convergence and \eqref{syn:geom:countableBM}, using \eqref{syn:geom:Rdef}. The tail mean need not be integral: the real volume extension is exactly what makes this argument valid.
\end{proof}

This theorem does not imply the uncorrected inequality asked in Leake--Mohammadi Yekta Question 1.7 \cite{R4E6:LMY}. The correction can have either sign, and its removal is not justified by taking limits in \eqref{syn:geom:countableBM}.

\subsection{A bounded formula size for all nullity-three systems}
\label{syn:geom:bounded}
\begin{corollary}[At most thirty edge terms]\label{syn:geom:30edges}
Every nonempty half-relaxed fiber of a unimodular matrix of nullity three, with finite boxes or pointed nonnegative constraints, has at most twelve facets, twenty vertices, and thirty edges in kernel coordinates. Its canonical correction \eqref{syn:geom:edges} therefore uses at most thirty terms, independently of the rank and the number of original variables. Opposite kernel directions are allowed.
\end{corollary}

\begin{proof}
Choose a maximal minor of $M$ and normalize it to the identity. Its kernel is represented by $K=(-N,I_3)^{\mathsf T}$, with all nonsingular three-row minors $\pm1$, as follows by the complementary-minor identity. Every nonzero row of $K$ is primitive: extend it to a nonsingular three-row minor, whose determinant is $\pm1$.

Remove zero rows, normalize signs, and retain one representative of each proportionality class. The rank-three direction certificate proved in Round 5, entry 06, and reproduced in Round 6, entry 06, identifies this simple unimodular row system, after an integral change of coordinates, with a subset of
\[
 e_1,e_2,e_3,e_1-e_2,e_1-e_3,e_2-e_3.
\]
That certificate enumerates all ten additional candidate directions beside a chosen basis, excludes every four-element candidate set, and prints integral basis transformations for every maximal admissible set. We use that preserved finite proof, rather than an assumption that every fiber is an $A_3$ nonnegative fiber.

There are thus at most six unoriented normal directions, or twelve oriented directions. Inequalities with the same primitive oriented normal reduce to the tightest one; adding $1/2$ to every bound commutes with taking this minimum. The half-relaxed body is bounded and full-dimensional, even if the original fiber was lower-dimensional. It has at most $f=12$ facets. If $v,e,f$ are its face numbers, Euler's relation and minimum vertex degree three give
\[
 v-e+f=2,\qquad 3v\le2e,
 \qquad v\le2f-4\le20,\quad e\le3f-6\le30.
\]
No simplicity assumption is required. The edge formula of Round 7 now proves the conclusion. Opposite normals correspond to distinct oriented facets and were explicitly included in the bound twelve.
\end{proof}

The bound controls the number of geometric terms after compression, not the bit complexity of constructing the original kernel or reading a large input. It also does not turn the corrected result into an uncorrected theorem for mixed nullity-three systems.

\subsection{What is additionally closed, and what is retained as open}
The seed comparison \eqref{syn:geom:comparison}, the series pullback identities \eqref{syn:geom:seriesF}--\eqref{syn:geom:seriesnonnegative}, the countable all-nullity inequality \eqref{syn:geom:countableBM}, and the thirty-term bound are the additional deductions in this synthesis. The source proofs retain credit for their input theorems. The unrestricted Ewald-point minimum, the exact tower rate, general uncorrected unimodular barycentric log-concavity, the optimal discrete uncorrected $A_3$ power, and arbitrary weighted or boxed $K_4$ barycentric inequalities remain unresolved by these arguments.




\part{Power sums and local valuation laws}

\chapter[An elementary proof of a cubic power-sum conjecture and a sharp fourth-power strengthening]{An elementary proof of a cubic power-sum conjecture and a sharp fourth-power strengthening}
\label{source:R1E1}
\noindent\textit{Source development: Round 1, entry 01. Original mathematical arguments retained; notation is local to this chapter.}\par\medskip


\begin{quote}\small\noindent\textbf{Chapter abstract.} 
We give a self-contained proof of the assertion published as Conjecture 12 by Sondow and MacMillan in 2011: if $S_n(k)\equiv(k+1)^n\pmod{k^2}$, then $S_n(k)\equiv(k/p)S_n(p)\pmod{p^3}$ for every prime $p\mid k$, where $S_n(x)=\sum_{j=1}^xj^n$. For $p\ge5$ we obtain the stronger modulus $p^4$ under weaker hypotheses. The example $(n,k,p)=(12,42,7)$ shows that a universal replacement by $p^5$ is false. The proofs use elementary binomial expansion and pairing; their ingredients are standard. \textbf{Status:} the stated subsidiary assertion is proved here, but historical priority and its present open status have not been established by an exhaustive literature review. No solution of the Erd\H{o}s--Moser equation or improvement of its known global bounds is claimed.
\end{quote}

\section{The precise target and scope}
Write $S_r(x)=\sum_{j=1}^x j^r$ for positive integers $r,x$. The Erd\H{o}s--Moser equation is $S_n(k)=(k+1)^n$; the conjecture is that its only positive-integer solution is $(n,k)=(1,2)$. Our target is not that equation. It is the following explicitly stated conjecture in the published paper of Sondow and MacMillan \cite[Conjecture 12, p.~7]{R1E1:SM}; it appears as Conjecture 2 in their 2010 preprint \cite{R1E1:SMpre}.

\begin{srcAAtheorem}\label{R1E1:main}
If $n,k$ are positive integers and
\begin{equation}\label{R1E1:square}
 S_n(k)\equiv(k+1)^n\pmod{k^2},
\end{equation}
then, for each prime divisor $p$ of $k$,
\begin{equation}\label{R1E1:cube}
 S_n(k)\equiv\frac{k}{p}S_n(p)\pmod{p^3}.
\end{equation}
For every such prime $p\ge5$, the conclusion holds modulo $p^4$.
\end{srcAAtheorem}

The source was published on 27 May 2011 and calls the cubic assertion a conjecture. Searches during this experiment did not locate a subsequent explicit resolution, but that negative search result is not evidence of priority. Related power-sum pairing congruences occur in Me\v{s}trovi\'{c} \cite{R1E1:M}; our use of pairing is not presented as a new method. In particular, the contribution evaluated here is the complete deduction of the stated implication, the treatment of the exceptional prime $3$, and the stronger modulus with its sharpness example.

\section{Elementary local facts}
For a prime $p$ put $A_r(p)=\sum_{j=1}^{p-1}j^r$. The familiar finite-field sum is
\begin{equation}\label{R1E1:field}
 A_r(p)\equiv\begin{cases}-1&(p-1)\mid r,\\0&(p-1)\nmid r\end{cases}\pmod p \qquad(r\ge1).
\end{equation}
Indeed, summing powers of a generator of the cyclic group $\mathbb F_p^\times$ gives either $p-1$ or a geometric series whose ratio differs from $1$.

\begin{srcAAlemma}\label{R1E1:necessary}
If $S_n(k)\equiv(k+1)^n\pmod k$, then $k$ is squarefree and, for every prime $p\mid k$, writing $a=k/p$, one has
\begin{equation}\label{R1E1:basic}
 (p-1)\mid n,\qquad a\equiv-1\pmod p.
\end{equation}
\end{srcAAlemma}
\begin{proof}
Splitting $1,\ldots,k$ into $a$ blocks of length $p$ gives $S_n(k)\equiv a A_n(p)\pmod p$. The assumed congruence makes this $1$ modulo $p$. Formula \eqref{R1E1:field} therefore forces $(p-1)\mid n$ and $a\equiv-1\pmod p$. In particular $p\nmid a$, so $p^2\nmid k$.
\end{proof}
This is the elementary condition already recorded in \cite[Theorem 3]{R1E1:SM}; the proof was included to make the argument self-contained.

For $h=0,\ldots,a-1$ set
\[
 H_i(a)=\sum_{h=0}^{a-1}h^i,\qquad
 H_1(a)+H_2(a)=\frac{a(a^2-1)}3.
\]
The last identity follows from $H_1=a(a-1)/2$ and $H_2=a(a-1)(2a-1)/6$.

\begin{srcAAlemma}[Paired sums]\label{R1E1:pair}
Let $p$ be an odd prime and $r\ge3$ odd. Then
\begin{equation}\label{R1E1:pair2}
 2A_r(p)\equiv rpA_{r-1}(p)\pmod{p^2},
\end{equation}
and
\begin{equation}\label{R1E1:pair3}
 2A_r(p)\equiv rpA_{r-1}(p)-\binom r2p^2 A_{r-2}(p)\pmod{p^3}.
\end{equation}
\end{srcAAlemma}
\begin{proof}
The substitution $j\mapsto p-j$ permutes $1,\ldots,p-1$. Expand $(p-j)^r$ by the binomial theorem, use the oddness of $r$, and add $A_r(p)$ to both sides. Terms containing $p^2$ or $p^3$, respectively, may be discarded in the two displayed congruences.
\end{proof}
Fractions with denominator prime to the modulus below mean modular inverses. All formulas involving a denominator $3$ at the prime $3$ are handled as integer identities before reduction.

\section{The fourth-power result away from $2$ and $3$}
\begin{srcAAtheorem}\label{R1E1:fourth}
Let $p\ge5$ be prime, $n\ge4$ be even with $(p-1)\mid n$, and $a\ge1$. Then
\begin{equation}\label{R1E1:defect4}
 S_n(pa)-aS_n(p)
 \equiv\frac{n(n-1)}6p^2a(a^2-1)A_{n-2}(p)\pmod{p^4}.
\end{equation}
In particular, if $a\equiv-1\pmod p$, this difference is divisible by $p^4$.
\end{srcAAtheorem}
\begin{proof}
Multiples of $p$ contribute zero modulo $p^4$, since $n\ge4$. Expanding the other terms block by block yields, with $D=S_n(pa)-aS_n(p)$,
\begin{align}\label{R1E1:expand4}
 D\equiv{}&npH_1 A_{n-1}(p)
 +\binom n2p^2H_2 A_{n-2}(p)\notag\\
 &+\binom n3p^3H_3 A_{n-3}(p)\pmod{p^4}.
\end{align}
The exponent $n-3$ is positive and odd, so $A_{n-3}(p)\equiv0\pmod p$ by \eqref{R1E1:field}. Thus the last term vanishes. Applying \eqref{R1E1:pair3} to $r=n-1$ gives
\[
 A_{n-1}(p)\equiv\frac{n-1}2pA_{n-2}(p)\pmod{p^3},
\]
because the remaining term contains $p^2 A_{n-3}(p)$. Substitution into \eqref{R1E1:expand4} gives
\[
 D\equiv\frac{n(n-1)}2p^2(H_1+H_2)A_{n-2}(p)\pmod{p^4},
\]
which is \eqref{R1E1:defect4}. Finally, $(p-1)\mid n$ and $p\ge5$ imply $(p-1)\nmid n-2$, so $p\mid A_{n-2}(p)$. If $a\equiv-1\pmod p$, also $p\mid a^2-1$. Both factors together prove the last assertion; division by $6$ is legitimate for $p\ge5$.
\end{proof}
The square-modulus hypothesis of Theorem~\ref{R1E1:main} is unnecessary for this part: its reduction modulo $k$ already supplies the needed conditions.

\section{The exceptional prime and the complete proof}
\begin{srcAAlemma}\label{R1E1:three}
For even $n\ge4$ and every positive integer $a$,
\begin{equation}\label{R1E1:defect3}
 S_n(3a)-aS_n(3)\equiv
 -\frac{3n(n-1)}2a(a^2-1)\pmod{27}.
\end{equation}
If, in addition, $a\equiv-1\pmod3$ and
\begin{equation}\label{R1E1:local9}
 S_n(3a)\equiv(3a+1)^n\pmod9,
\end{equation}
then the difference in \eqref{R1E1:defect3} vanishes modulo $27$.
\end{srcAAlemma}
\begin{proof}
Here $A_{n-2}(3)\equiv-1\pmod3$ and, by \eqref{R1E1:pair2},
$A_{n-1}(3)\equiv-3(n-1)/2\pmod9$. Multiples of $3$ contribute zero modulo $27$. The two-term block expansion therefore gives
\[
 D\equiv-\frac{9n(n-1)}2(H_1+H_2)
 =-\frac{3n(n-1)}2a(a^2-1)\pmod{27}.
\]
The equality uses the exact integer identity $3(H_1+H_2)=a(a^2-1)$; no division by $3$ in a congruence has been made.

For the second assertion, the same expansion modulo $9$ gives $S_n(3a)\equiv aS_n(3)\pmod9$, because $3\mid A_{n-1}(3)$. Write $n=2t$. Then $2^n=4^t\equiv1+3t\pmod9$, so
\[
 S_n(3)\equiv2-3n\pmod9.
\]
Using \eqref{R1E1:local9} yields $a(2-6n)\equiv1\pmod9$. Write $a=-1+3c$, as permitted by hypothesis. Substitution and division of the resulting integer divisibility by $3$ give
\[
 2c+2n-1\equiv0\pmod3,
 \qquad c\equiv2n-1\pmod3.
\]
Hence $3\mid n(n-1)c$, since $3\mid n(n-1)(2n-1)$ for every integer $n$. As $a+1=3c$, it follows that $9\mid n(n-1)(a+1)$. Formula \eqref{R1E1:defect3} is consequently zero modulo $27$.
\end{proof}

\begin{proof}[Proof of Theorem~\ref{R1E1:main}]
If $k=1$ the assertion is vacuous. Lemma~\ref{R1E1:necessary} applies to \eqref{R1E1:square}. If $n$ is odd it forces $k\in\{1,2\}$. For $k=2$ and odd $n\ge3$, the two sides of \eqref{R1E1:square} are $1$ and $3$ modulo $4$, respectively, a contradiction. For $n=1,k=2$, the required block identity is an equality.

If $n=2$, Lemma~\ref{R1E1:necessary} restricts prime divisors of $k$ to $2,3$ and makes $k$ squarefree; hence $k\in\{1,2,3,6\}$. Direct reduction gives $S_2(3)=14\not\equiv16\pmod9$ and $S_2(6)=91\not\equiv49\pmod{36}$. Only $k=1,2$ remain, and the assertion is again vacuous or an equality.

It remains to take $n\ge4$ even. For $p=2$, every odd $j$ has $j^n\equiv1\pmod8$, and every even $j$ has $j^n\equiv0\pmod8$. Thus $S_n(2a)\equiv a\equiv aS_n(2)\pmod8$. For $p=3$, the assumptions of Lemma~\ref{R1E1:three} follow from Lemma~\ref{R1E1:necessary} and \eqref{R1E1:square}. For $p\ge5$, Theorem~\ref{R1E1:fourth} gives even the asserted fourth-power strengthening.
\end{proof}

\section{Sharpness and a cubic criterion}
\begin{srcAAproposition}\label{R1E1:sharp}
The modulus $p^4$ in Theorem~\ref{R1E1:main} cannot universally be replaced by $p^5$, even for $p\ge5$. The hypothesis \eqref{R1E1:square} holds for $(n,k)=(12,42)$, but
\[
 S_{12}(42)-6S_{12}(7)\equiv2401\pmod{16807}.
\]
\end{srcAAproposition}
\begin{proof}
Exact modular summation gives the following small certificates:
\[
 S_{12}(42)\equiv43^{12}\equiv505\pmod{1764},
\]
\[
 S_{12}(42)\equiv15646\pmod{16807},\qquad
 S_{12}(7)\equiv10611\pmod{16807}.
\]
Thus $15646-6\cdot10611=-48020\equiv2401\pmod{16807}$.
All these finite equalities may be verified by repeated multiplication; the accompanying standard-library Python script independently reproduces them. Since $2401=7^4$ and $16807=7^5$, the conclusion follows.
\end{proof}

For completeness, Theorem~\ref{R1E1:main} also gives the explicit type of cubic criterion sought after Conjecture 12 in \cite{R1E1:SM}. Define, for each prime $p$,
\[
 q_p(j)=\frac{j^{p-1}-1}{p},\quad
 Q_p=\sum_{j=1}^{p-1}q_p(j),\quad R_p=\sum_{j=1}^{p-1}q_p(j)^2.
\]
\begin{srcAAcorollary}
Suppose $n\ge4$ is even and \eqref{R1E1:square} holds. For $p\mid k$ write $a=k/p$ and $t=n/(p-1)$. Then $S_n(k)\equiv(k+1)^n\pmod{k^3}$ if and only if, for every such $p$,
\begin{equation}\label{R1E1:criterion}
 a\left(p-1+ptQ_p+p^2\binom t2 R_p\right)
 \equiv1+npa+\binom n2p^2a^2\pmod{p^3}.
\end{equation}
\end{srcAAcorollary}
\begin{proof}
Fermat's theorem makes $q_p(j)$ integral. Expanding $(1+pq_p(j))^t$ and summing gives
$S_n(p)\equiv p-1+ptQ_p+p^2\binom t2R_p\pmod{p^3}$; the endpoint $p^n$ vanishes. Apply Theorem~\ref{R1E1:main}, expand $(1+pa)^n$, and use squarefreeness of $k$ and the Chinese remainder theorem. This proof also covers $p=2$, where $Q_2=R_2=0$.
\end{proof}

\section{Verification, limitations, and research assessment}
The file \texttt{verify.py} uses only exact Python integers. It checks the $3$-adic defect on $4\le n\le100$ even and $1\le a\le100$, the fourth-power identity for primes $5\le p\le31$, four positive multiples of $p-1$ and $1\le a\le50$, and Theorem~\ref{R1E1:main} for all $1\le k\le400$, $1\le n\le80$ satisfying its premise. It also reproduces Proposition~\ref{R1E1:sharp}. These checks are diagnostic, not substitutes for the general proof.

An initial route was to seek a global obstruction to the Erd\H{o}s--Moser equation by lifting its modular conditions. This did not produce a contradiction: a block-reduction identity is not the original equation, and congruences alone do not bound all possible $k$. The alternative developed here isolates and proves the explicitly sourced subsidiary assertion. The failed attempt to lift universally to $p^5$ instead supplied the sharpness certificate.

This is a short elementary number-theory note, not an Annals-level breakthrough. The weakest claim is bibliographic novelty, not an identified mathematical gap: the literature search was limited, and both power-sum pairing and block expansion are classical. Me\v{s}trovi\'{c}'s 2014 paper \cite{R1E1:M}, including its Theorem 1 and Proposition 7, establishes related congruences. It does not explicitly state the implication proved here in the portions inspected. A specialist literature check is necessary before any claim of a first proof. The full Erd\H{o}s--Moser conjecture remains untouched by the argument.



\chapter[Three p-adic branches of power-sum block defects and an endpoint obstruction to naive lifting]{Three $p$-adic branches of power-sum block defects and an endpoint obstruction to naive lifting}
\label{source:R2E1}
\noindent\textit{Source development: Round 2, entry 01. Original mathematical arguments retained; notation is local to this chapter.}\par\medskip


\begin{quote}\small\noindent\textbf{Chapter abstract.} 
\textbf{Status: an all-precision local classification under explicit Bernoulli-valuation hypotheses, not a solution of the Erd\H{o}s--Moser equation; priority unverified.}
Let $S_n(x)$ be the power-sum polynomial and $D_{n,p}(a)=S_n(pa)-aS_n(p)$. For primes $p\ge7$ and even $n\ge4$, we prove that, when $v_p(B_{n-2})\in\{-1,0\}$, the normalized defect has exactly three simple $p$-adic root branches, at $0$, $1$, and a shifted root $\alpha\equiv-1\pmod p$. An explicit factorization determines $v_p(D_{n,p}(a))$ for every $a\in\mathbb Z_p$ and every congruence precision. The normalization is $p^e$, where $e=3+v_p(n(n-1))+v_p(B_{n-2})$, and $v_p(\alpha+1)=n-e$. The same conclusion covers $v_p(B_{n-2})=1$ when $(p-1)\mid n$. Centering the sum restores an exact root at $-1$; the omitted endpoint explains why merely forcing $a$ closer to $-1$ eventually stops increasing the uncentered valuation. A repeated-root example at $p=5$ marks a genuine limitation of the argument.
\end{quote}

\section{Context, provenance, and the new question}
For positive integers $n,x$, let $S_n(x)=\sum_{j=1}^x j^n$, extended to $x\in\mathbb Q_p$ by its rational polynomial. Sondow--MacMillan \cite[Conjecture 12]{R1E1:SM} asked whether
\[
 S_n(k)\equiv(k+1)^n\pmod{k^2}
 \quad\Longrightarrow\quad
 S_n(k)\equiv(k/p)S_n(p)\pmod{p^3}\quad(p\mid k).
\]
Round 1 Entry 01 \cite{R2E1:R1} supplied a proof of that assertion and the stronger modulus $p^4$ for $p\ge5$, with a $p^5$ counterexample. We do not count those results again. The original source states a conjecture in 2011; neither round establishes first-proof priority or an exhaustive present-day status assessment.

This note asks a different, structural question: for fixed $n,p$, what determines \emph{every} precision to which the block-reduction congruence holds as $a=k/p$ varies? Its answer below replaces a single uniform modulus by a complete valuation formula. The result is not a claim that this exact classification was itself a published open conjecture. It is a local extension motivated by the sourced problem. The Erd\H{o}s--Moser equality remains outside the theorem.

The general theory of valuations of power sums and their Bernoulli connections is classical; see Sondow--Tsukerman \cite{R2E1:ST} and the recent study \cite{R2E1:AMSV}. Our polynomial is a \emph{difference} of two such sums. We make no novelty claim for Faulhaber's formula, Bernoulli denominator estimates, simple-root lifting, or polynomial factorization. The proposed increment is their uniform coefficient control for this defect, its exact three-branch classification, and its quantified endpoint shift.

\section{Main statements}
Let $p\ge7$ be prime and $n\ge4$ even. Write
\[
 s=v_p(n(n-1)),\qquad b=v_p(B_{n-2}),\qquad
 e=3+s+b,\qquad
 u=p^{-s-b}\frac{n(n-1)}6B_{n-2}.
\]
We use $B_0=1$, $B_1=-1/2$, $B_2=1/6$, and $v_p(0)=+\infty$. The pair $(n,p)$ will satisfy one of the following explicit hypotheses:
\begin{align}
 &b\in\{-1,0\};\label{R2E1:caseA}\\
 &(p-1)\mid n\quad\hbox{and}\quad b=1.\label{R2E1:caseB}
\end{align}
In either case $e$ is an integer, $u\in\srcAGZp^\times$, and $n>e$.
Define
\[
 D(a)=S_n(pa)-aS_n(p),\qquad
 \widetilde D(a)=D(a)-\frac{p^n}{2}(a^n-a).
\]

\begin{srcAGtheorem}[Exact factorization and valuation]\label{R2E1:main}
Under \eqref{R2E1:caseA} or \eqref{R2E1:caseB}, there exist polynomials $V,U\in\srcAGZp[a]$ and a unique $\alpha\in-1+p\srcAGZp$ such that
\begin{align}
 \widetilde D(a)&=p^e a(a^2-1)V(a),\label{R2E1:centerfact}\\
 D(a)&=p^e a(a-1)(a-\alpha)U(a),\label{R2E1:rawfact}\\
 V(a)&\equiv U(a)\equiv u\pmod p\quad\hbox{as polynomials}.\label{R2E1:units}
\end{align}
In particular, $U$ and $V$ take unit values everywhere on $\srcAGZp$, and
\begin{align}
 v_p(\widetilde D(a))&=e+v_p(a)+v_p(a-1)+v_p(a+1),\label{R2E1:centerval}\\
 v_p(D(a))&=e+v_p(a)+v_p(a-1)+v_p(a-\alpha)\label{R2E1:rawval}
\end{align}
for every $a\in\srcAGZp$. Moreover,
\begin{equation}\label{R2E1:shift}
 v_p(\alpha+1)=n-e,\qquad
 \alpha\equiv-1-\frac{p^{n-e}}{2u}\pmod{p^{n-e+1}}.
\end{equation}
These formulas include the infinite valuations at the roots.
\end{srcAGtheorem}

\begin{srcAGcorollary}[All congruence levels]\label{R2E1:levels}
For every integer $h\ge1$, exactly three residue classes modulo $p^h$ satisfy
\[
 D(a)\equiv0\pmod{p^{e+h}}:
 \qquad a\equiv0,\ 1,\ \alpha\pmod{p^h}.
\]
Consequently this set of nonnegative integers has natural density $3/p^h$. In the branch $a\equiv-1\pmod p$, precisely one residue class modulo $p^h$ survives.
\end{srcAGcorollary}

\begin{srcAGcorollary}[Endpoint transition]\label{R2E1:transition}
Put $h_0=n-e$. Suppose $a\equiv-1\pmod p$ and $t=v_p(a+1)<\infty$. Then
\[
 v_p(D(a))=
 \begin{cases}
 e+t,&t<h_0,\\
 n,&t>h_0.
 \end{cases}
\]
At the boundary $a=-1+p^{h_0}z$, $z\in\srcAGZp^\times$, one has
\begin{equation}\label{R2E1:boundary}
 p^{-n}D(a)\equiv1+2uz\pmod p.
\end{equation}
Further divisibility is therefore controlled by the shifted root, not by increasing $v_p(a+1)$ indefinitely.
\end{srcAGcorollary}
In particular there are arbitrarily large positive integers with arbitrarily large $v_p(a+1)$ but with $v_p(D(a))=n$. On the other hand the $\alpha$ branch provides positive integers with unbounded $v_p(D(a))$.

\section{Uniform coefficient control}
We state precisely the classical Bernoulli fact used in the proof. The von Staudt--Clausen theorem gives, for positive even $r$,
\[
 B_r+\sum_{\ell-1\mid r}\frac1\ell\in\mathbb Z,
\]
where $\ell$ runs over primes. Thus $v_p(B_r)\ge-1$, and it equals $-1$ exactly when $(p-1)\mid r$; otherwise $B_r$ is $p$-integral. Also $B_0=1$. This formulation is recorded in \cite[Section 5]{R2E1:ST}; it is a classical dependency, not a result of this experiment.

The generating function $t/(e^t-1)=\sum B_jt^j/j!$, multiplied by the finite geometric sum of exponentials, gives the power-sum polynomial. For even $n$ it reads
\begin{equation}\label{R2E1:faulhaber}
 S_n(x)=\frac{x^n}{2}
  +\sum_{\substack{1\le d\le n+1\\d\text{ odd}}}
     \frac1d\binom n{d-1}B_{n+1-d}x^d.
\end{equation}
This also shows directly that $S_n(x)-x^n/2$ is an odd polynomial. Subtracting the block-linear term cancels the term $d=1$ and gives the exact identity
\begin{equation}\label{R2E1:exact}
 D(a)=\frac{p^n}{2}(a^n-a)
 +\sum_{\substack{3\le d\le n+1\\d\text{ odd}}}
     c_dp^d(a^d-a),\qquad
 c_d=\frac1d\binom n{d-1}B_{n+1-d}.
\end{equation}
In particular $c_3=n(n-1)B_{n-2}/6$.

\begin{srcAGlemma}[Coefficient bounds]\label{R2E1:coeff}
For $d\ge5$ odd occurring in \eqref{R2E1:exact},
\[
 v_p(p^dc_d)\ge s+4.
\]
If $(p-1)\mid n$, the stronger lower bound $s+5$ holds for every such $d$. The endpoint coefficient $p^n/2$ has valuation at least $s+4$ in general and at least $s+5$ when $(p-1)\mid n$.
\end{srcAGlemma}
\begin{proof}
The exact identity
\[
 c_d=\frac{n(n-1)}{d(d-1)(d-2)}
       \binom{n-2}{d-3}B_{n+1-d}
\]
gives
\begin{equation}\label{R2E1:estimate}
 v_p(p^dc_d)\ge s+d-1-v_p(d(d-1)(d-2)).
\end{equation}
At most one of three consecutive integers is divisible by $p\ge7$, so the last valuation is at most $\lfloor\log_p d\rfloor$. When $d=5$, it is zero, giving $s+4$. For $d\ge7$, the elementary inequality $\lfloor\log_p d\rfloor\le d-6$ gives at least $s+5$. For example the latter inequality follows from $7^{d-5}>d$ for $d\ge7$, proved by induction.

If $(p-1)\mid n$, then $n\ge6$ and $n-4$ is positive and not divisible by $p-1$. Consequently $B_{n-4}$ is $p$-integral. For $d=5$ we may omit the loss of one in \eqref{R2E1:estimate}, giving $s+5$ also in this case.

Finally at most one of $n,n-1$ is divisible by $p$, and $s\le\lfloor\log_p n\rfloor$. This implies $s\le n-4$ for $n\ge4$ and $s\le n-5$ for $n\ge6$: check $4\le n\le6$ directly and use the preceding logarithm bound for $n\ge7$. The endpoint assertions follow.
\end{proof}

\begin{srcAGcorollary}[A cubic reduction]\label{R2E1:cubic}
Under the hypotheses of Theorem~\ref{R2E1:main}, both $F=p^{-e}D$ and $\widetilde F=p^{-e}\widetilde D$ belong to $\srcAGZp[a]$, and
\begin{equation}\label{R2E1:reduction}
 F(a)\equiv\widetilde F(a)\equiv u(a^3-a)\pmod p.
\end{equation}
Moreover $D(0)=D(1)=0$ and $D(-1)=p^n$.
\end{srcAGcorollary}
\begin{proof}
The $d=3$ term has coefficient $p^3c_3=p^e u$. In case \eqref{R2E1:caseA}, $e\le s+3$, so Lemma~\ref{R2E1:coeff} puts every other coefficient in $p^{e+1}\srcAGZp$. In case \eqref{R2E1:caseB}, $e=s+4$, and the stronger bound does the same. This proves \eqref{R2E1:reduction} and also $n>e$. The three evaluations follow immediately from \eqref{R2E1:exact}, because $n$ is even and every $d$ in the sum is odd.
\end{proof}
This reduction is coefficientwise, not merely an equality of polynomial functions on the residue field. That distinction will matter at the excluded small prime.

\section{From the cubic reduction to exact valuations}
\begin{srcAGlemma}[Simple-root lifting, with proof]\label{R2E1:hensel}
Let $f\in\srcAGZp[a]$. A residue $r$ with $f(r)\equiv0\pmod p$ and $f'(r)\not\equiv0\pmod p$ lifts uniquely to a root in $r+p\srcAGZp$. At each precision $p^h$ there is precisely one root residue above $r$.
\end{srcAGlemma}
\begin{proof}
If $r_h$ is a root modulo $p^h$, write its next lift as $r_h+tp^h$, with $t\in\{0,\ldots,p-1\}$. Taylor expansion gives
\[
 f(r_h+tp^h)\equiv f(r_h)+tp^hf'(r_h)\pmod{p^{h+1}}.
\]
The unit derivative makes the equation for $t$ unique. Starting at $h=1$ gives compatible residues, hence a unique $p$-adic integer; polynomial continuity makes it a root. The same digit argument proves uniqueness.
\end{proof}

\begin{proof}[Proof of Theorem~\ref{R2E1:main}]
For the centered polynomial, \eqref{R2E1:exact} shows that $\widetilde F$ vanishes at $0,1,-1$, so monic polynomial division gives
$\widetilde F=a(a^2-1)V$ with $V\in\srcAGZp[a]$. Reducing this identity and \eqref{R2E1:reduction} modulo $p$ shows $V\equiv u$ as a polynomial.

For $F$, the reduction $u(a^3-a)$ has exactly three roots in $\mathbb F_p$, namely $0,1,-1$. They are simple: the derivatives there are $-u,2u,2u$. The first two are already exact roots. Lemma~\ref{R2E1:hensel} gives a unique third root $\alpha\equiv-1\pmod p$. Dividing successively by the monic factors $a$, $a-1$, $a-\alpha$ preserves coefficient integrality and yields \eqref{R2E1:rawfact}. Reduction again gives $U\equiv u$; in particular $U(a)$ and $V(a)$ are units for every $a\in\srcAGZp$. Taking valuations proves \eqref{R2E1:centerval}--\eqref{R2E1:rawval}.

Evaluate \eqref{R2E1:rawfact} at $a=-1$. Since $D(-1)=p^n$,
\[
 p^{n-e}=-2(\alpha+1)U(-1).
\]
The unit property proves the exact valuation and the congruence in \eqref{R2E1:shift}.
\end{proof}

\begin{proof}[Proofs of Corollaries~\ref{R2E1:levels} and \ref{R2E1:transition}]
The three roots $0,1,\alpha$ occupy distinct residue classes modulo $p$. Hence, for any $a\in\srcAGZp$, at most one of the three factors in \eqref{R2E1:rawval} has positive valuation. Divisibility by $p^{e+h}$ is therefore equivalent to congruence to one root modulo $p^h$. Periodicity gives the asserted density.

In the branch $a\equiv-1\pmod p$, the factors $a$ and $a-1$ are units. Compare $a+1$ with $\alpha+1$ using \eqref{R2E1:shift}. If their valuations differ, the valuation of their difference is the smaller one, giving the two cases. At equality, Taylor expansion of $F$ at $-1$ gives
\[
 F(-1+p^{h_0}z)\equiv p^{h_0}(1+2uz)\pmod{p^{h_0+1}},
\]
since $h_0\ge1$ and $F'(-1)\equiv2u\pmod p$. This is \eqref{R2E1:boundary}.
\end{proof}

\section{Examples, a genuine exceptional case, and interpretation}
For $p=7,n=6$, one has $s=b=0$, $e=3$, $u=-1/6$, and the exact normalized polynomial is
\[
 F(a)=7^3(a^7-a)+\frac{7^3}{2}(a^6-a)
       +\frac{7^2}{2}(a^5-a)-\frac16(a^3-a).
\]
Thus $F\equiv a^3-a\pmod7$. The shifted root satisfies
$\alpha\equiv-1+3\cdot7^3\pmod{7^4}$, or $\alpha\equiv1028\pmod{2401}$.
The following exact checks illustrate both valuation regimes and the shifted branch:
\begin{center}
\begin{tabular}{rrrrr}
\toprule
$a$&$6$&$48$&$342$&$2400$\\
$v_7(a+1)$&$1$&$2$&$3$&$4$\\
$v_7(D(a))$&$4$&$5$&$6$&$6$\\
\bottomrule
\end{tabular}
\end{center}
In contrast, $a=1028$ gives $v_7(D(a))=8$. This finite equality is a diagnostic example, not a premise of the theorem. The compatible root residues modulo $7^h$, $h=1,\ldots,8$, are
\[
 6,\ 48,\ 342,\ 1028,\ 1028,\ 68256,\ 421203,\ 1244746.
\]
Case \eqref{R2E1:caseB} is nonempty: $(p,n)=(7,30)$ has $b=1$, $s=0$, $e=4$. Thus its entire $a\equiv-1\pmod7$ branch automatically has defect divisible by $7^5$, and the theorem classifies every higher lift.

The restriction $p\ge7$ cannot simply be erased from the proof. For $p=5,n=24$, exact Bernoulli arithmetic gives
\begin{equation}\label{R2E1:five}
 5^{-3}D_{24,5}(a)\equiv4a+2a^3+4a^5
 =4a(a^2-1)^2\pmod5.
\end{equation}
Indeed $B_{22}=854513/138$ makes the cubic coefficient $2$ modulo $5$, while the $d=5$ term has coefficient
$5\binom{24}{4}B_{20}\equiv4\pmod5$, since $5B_{20}\equiv-1\pmod5$. Higher terms vanish modulo $5$ after normalization; indeed $d-1-\lfloor\log_5d\rfloor\ge5$ for odd $d\ge7$, and the three consecutive denominator factors contain at most one multiple of $5$. Thus the roots $\pm1$ are multiple, and the simple-root argument is unavailable despite the unit cubic coefficient. We do not assert that no alternative classification exists there.

The theorem concerns a necessary auxiliary local structure, not the simultaneous Erd\H{o}s--Moser equation. Arbitrarily deep solutions to the block congruence at one fixed prime do not yield global solutions, but they do rule out the hope of obtaining a universal local contradiction merely by repeatedly demanding this block reduction at that prime.

\section{Audit, limitations, and exact delta from Round 1}
\paragraph{What is new within the experiment.}
Round 1 Entry 01, Theorem 4, controlled one fixed modulus by a paired-sum expansion. Its Proposition 6 ruled out a uniform fifth-power extension. Here the valuation scale incorporates $v_p(n(n-1))$ and the Bernoulli valuation, and the conclusion specifies every precision, the unique displaced branch, its first nonzero displacement digit, and the saturation phenomenon. The centered exact factorization is also new to this submission. These are not re-statements of the Round 1 cubic theorem. No theorem from Entries 02--06 is used; their frozen papers and the assessment were read as required by the cumulative protocol.

\paragraph{Verification.}
The standalone \texttt{verify.py} uses rational arithmetic for Bernoulli numbers and defect coefficients, and exact modular arithmetic for root lifting. For even $4\le n\le120$ and $p\in\{7,11,13,17,19,23,29,31\}$ it checks 447 instances of \eqref{R2E1:caseA} and four additional instances of \eqref{R2E1:caseB}, the coefficientwise cubic reduction, 4510 truncated valuation comparisons, and 639 independent direct-sum evaluations. It checks \eqref{R2E1:five} and the displayed examples. The general proof does not rely on these finite tests.

\paragraph{Failed alternatives and remaining gap.}
A blanket formula using $a+1$ in place of $a-\alpha$ fails because $D(-1)=p^n$, not zero. Centering repairs that false symmetry, but restoring the endpoint shifts the root. Extending the same simple-root theorem to $p=5$ fails at \eqref{R2E1:five}. For larger positive $b$ not covered above, higher-degree terms can compete with the cubic after normalization. Their Newton polygons and multiple roots require further analysis. Those singular cases, and actual global consequences for the original Diophantine equation, are not solved here.

\paragraph{Novelty status.}
The inspected primary literature includes the original conjecture \cite{R1E1:SM}, the valuation/denominator results in \cite{R2E1:ST}, and the 2025 study \cite{R2E1:AMSV}. Searches for power-sum block defects and $p$-adic root classifications did not locate the precise theorem above. This is limited evidence, not a priority certification. The method is an elementary application of classical $p$-adic polynomial theory once the coefficient bounds are established. It is a potentially useful specialized structural result, not an Annals-level claim.



\chapter[Endpoint stability of power-sum valuations at regular and locally nonirregular primes]{Endpoint stability of power-sum valuations at regular and locally nonirregular primes}
\label{source:R3E1}
\noindent\textit{Source development: Round 3, entry 01. Original mathematical arguments retained; notation is local to this chapter.}\par\medskip


\begin{quote}\small\noindent\textbf{Chapter abstract.} 
We prove the endpoint-removal assertion printed as Conjecture 6.1 in Amdeberhan--Moll--Sharma--Vignat (2025): at a regular prime dividing $N$, the sums through $N-1$ and $N$ have the same valuation. For exponents at least two, a strict valuation gap makes cancellation impossible; exponent one is handled separately. We give complete coefficient-domination arguments for the underlying valuation formulas, including the small primes, and extend the conclusion to irregular primes whenever the relevant exponent lies outside their irregular Bernoulli residue classes. The regular-prime valuation formulas themselves already appear in the cited paper. Thus the contribution is a transparent closure of its printed assertion and a local refinement, not a claim to invent those formulas or a deep new general method. Historical first-proof priority remains unverified.
\end{quote}

\section{The printed assertion and precise scope}
Put
\[
 S_n(N)=\sum_{j=1}^N j^n,\qquad n,N\ge1.
\]
An odd prime $p$ is regular if it divides no numerator of $B_r$ for even $2\le r\le p-3$. As customary in this context, include $2$ among the regular primes. The Bernoulli convention is $B_0=1$, $B_1=-1/2$.

Amdeberhan, Moll, Sharma and Vignat \cite[Conjecture 6.1, p.~128]{R3E1:AMSV} state the assertion
\begin{equation}\label{R3E1:target}
 v_p(S_n(pm-1))=v_p(S_n(pm))
 \qquad(p\text{ regular};\ m,n\ge1).
\end{equation}
Their notation is $S_k(n)=\sum_{j=1}^k j^n$; we transpose the roles of the subscript and argument to retain the notation of our earlier rounds. Their preceding Theorems 6.1--6.3 already give the regular-prime valuation formulas used below. We do not claim those formulas as new. The current journal PDF, dated 2025 and inspected on 20 October 2026, still labels \eqref{R3E1:target} a conjecture. Targeted searches did not locate a subsequent resolution, but do not establish its present open status or first-proof priority.

A strict inequality is useful here. From $X-Y=Z$ and $v_p(X)\le v_p(Y)$ one cannot in general conclude $v_p(Z)=v_p(X)$: equality permits cancellation. For example, $X=p$, $Y=p(1+p)$ have equal valuations but their difference has valuation two. The weak-bound equivalence asserted after the printed conjecture therefore needs an additional argument. We prove a \emph{strict} bound when $n\ge2$; the remaining exponent has an explicit formula.

\begin{srcAMtheorem}[Endpoint stability]\label{R3E1:main}
For every regular prime $p$, every positive integer $N$ divisible by $p$, and every $n\ge1$,
\[
 v_p(S_n(N-1))=v_p(S_n(N)).
\]
Moreover, for $n\ge2$, writing $s=v_p(N)$, one has
\begin{equation}\label{R3E1:gap}
 v_p(S_n(N))<ns.
\end{equation}
Consequently, for $n\ge2$, the valuation is unchanged on replacing $S_n(N)$ by $S_n(N)+cN^n$, for any $c\in\mathbb Z_p$.
\end{srcAMtheorem}
The last sentence is just a consequence of the strict gap, not an independent discovery.

We formulate the larger local scope before proving the theorem. For an odd prime $p$ and positive even $k$, call $k$ \emph{admissible at $p$} if either $(p-1)\mid k$, or its unique even residue $r\in\{2,4,\ldots,p-3\}$ modulo $p-1$ satisfies $p\nmid\operatorname{num}(B_r)$. Thus every even $k$ is admissible at a regular prime; admissibility does not require the prime itself to be regular.

\begin{srcAMtheorem}[Local valuation formulas]\label{R3E1:local}
Let $p$ be odd, $p\mid N$, and $s=v_p(N)$. For even $n\ge2$ admissible at $p$,
\begin{equation}\label{R3E1:even}
 v_p(S_n(N))=
 \begin{cases}s-1,&(p-1)\mid n,\\s+v_p(n),&(p-1)\nmid n.\end{cases}
\end{equation}
For odd $n\ge3$ with $n-1$ admissible at $p$,
\begin{equation}\label{R3E1:odd}
 v_p(S_n(N))=
 \begin{cases}2s+v_p(n)-1,&(p-1)\mid n-1,\\
 2s+v_p(n(n-1)),&(p-1)\nmid n-1.
 \end{cases}
\end{equation}
In all these cases \eqref{R3E1:gap} and endpoint stability hold. For $p=2$ and even $N$, without any exponent restriction other than $n\ge2$,
\begin{equation}\label{R3E1:two}
 v_2(S_n(N))=
 \begin{cases}s-1,&n\text{ even},\\2s-2,&n\text{ odd}.
 \end{cases}
\end{equation}
\end{srcAMtheorem}
In the odd-prime cases we prove the more precise leading-term statements
\[
 \frac{S_n(N)}{NB_n}\in1+p\mathbb Z_p\quad(n\text{ even}),\qquad
 \frac{S_n(N)}{(n/2)B_{n-1}N^2}\in1+p\mathbb Z_p\quad(n\text{ odd}).
\]
These ratios are well-defined because the indicated even Bernoulli numbers are nonzero.

\section{Classical inputs and coefficient bookkeeping}
We use exactly two classical Bernoulli facts. The von Staudt--Clausen theorem gives, for positive even $k$,
\begin{equation}\label{R3E1:vsc}
 v_p(B_k)\ge-1,\qquad
 v_p(B_k)=-1\ \Longleftrightarrow\ (p-1)\mid k.
\end{equation}
The first Kummer congruence gives, when $k\equiv r\pmod{p-1}$ and neither is divisible by $p-1$,
\begin{equation}\label{R3E1:kummer}
 \frac{B_k}{k}\equiv\frac{B_r}{r}\pmod p.
\end{equation}
Both fractions are $p$-integral. Precise forms of these facts are recorded in \cite[Lemmas 2--3]{R3E1:BM}. In particular, admissibility and \eqref{R3E1:kummer} yield
\begin{equation}\label{R3E1:unit}
 v_p(B_k)=v_p(k)\quad\text{if }(p-1)\nmid k\text{ and }k\text{ is admissible}.
\end{equation}
For $p=3$ this second case never occurs. No assertion about the distribution of regular primes is required.

Faulhaber's formula, obtained by multiplying the Bernoulli generating function by the finite geometric sum of exponentials, is
\begin{equation}\label{R3E1:faul}
 S_n(N)=\frac{N^n}{2}+
 \sum_{\substack{1\le d\le n+1\\n+1-d\text{ even}}}c_dN^d,
 \qquad c_d=\frac1d\binom n{d-1}B_{n+1-d},\quad n\ge2.
\end{equation}
The summation includes $B_0=1$, and contains only odd $d$ if $n$ is even, only even $d$ if $n$ is odd. Useful identities are
\begin{align}
 c_d&=\frac{n}{d(d-1)}\binom{n-1}{d-2}B_{n+1-d}\quad(d\ge2),\label{R3E1:factorone}\\
 c_d&=\frac{n(n-1)}{d(d-1)(d-2)}\binom{n-2}{d-3}B_{n+1-d}\quad(d\ge3).\label{R3E1:factortwo}
\end{align}
The binomial coefficients are integers, including at $d=n+1$.

Here are elementary bounds used below. Two consecutive positive integers have at most one multiple of an odd prime; three consecutive integers have at most one multiple of a prime $p\ge5$. Therefore the valuation of the relevant product is at most $\lfloor\log_p d\rfloor$. The inequalities
\begin{align*}
 \lfloor\log_3 d\rfloor&\le d-3 &&(d\ge4),\\
 \lfloor\log_5 d\rfloor&\le d-3 &&(d\ge3),\\
 \lfloor\log_5 d\rfloor&\le d-4 &&(d\ge4)
\end{align*}
follow respectively from $3^{d-2}>d$, $5^{d-2}>d$, and $5^{d-3}>d$, each proved from its initial value by induction. These deliberately coarse bounds suffice. For an odd prime $p$,
\[
 v_p(n)\le n-2\quad(n\ge2).
\]
For $p\ge5$ and $n\ge3$, at most one of $n,n-1$ is divisible by $p$, giving
\begin{equation}\label{R3E1:endpointbound}
 v_p(n(n-1))\le\lfloor\log_5n\rfloor\le n-3.
\end{equation}
The final inequality holds directly at $n=3$ and by the bounds above thereafter.

\section{Proof of the local formulas}
\begin{proof}[Even exponents, odd primes]
The first term in the sum \eqref{R3E1:faul} is $NB_n$. Suppose first $(p-1)\mid n$. Its valuation is $s-1$ by \eqref{R3E1:vsc}. Every other summand has odd $d\ge3$ and
\[
 v_p(c_dN^d)\ge ds-1-v_p(d)>s-1,
\]
because $(d-1)s>v_p(d)$. The endpoint term $N^n/2$ has valuation $ns>s-1$. Thus the first term is the unique least-valuation term.

Otherwise $p\ge5$ and \eqref{R3E1:unit} gives $v_p(NB_n)=s+t$, where $t=v_p(n)$. Using \eqref{R3E1:factorone}, every odd $d\ge3$ satisfies
\[
 v_p(c_dN^d)\ge ds+t-1-\lfloor\log_p d\rfloor\ge s+t+1.
\]
The endpoint valuation $ns$ exceeds $s+t$ since $t\le n-2$. Again the first term is uniquely least. This proves \eqref{R3E1:even} and its leading-term congruence.
\end{proof}

\begin{proof}[Odd exponents, odd primes]
The first summand in \eqref{R3E1:faul} is $(n/2)B_{n-1}N^2$. Write $t=v_p(n)$. If $(p-1)\mid n-1$, its valuation is $2s+t-1$. For even $d\ge4$, \eqref{R3E1:factorone} gives
\[
 v_p(c_dN^d)\ge ds+t-1-\lfloor\log_p d\rfloor
 \ge2s+t.
\]
The endpoint has valuation $ns>2s+t-1$, because $t\le n-2$ and $s\ge1$.

If $(p-1)\nmid n-1$, then $p\ge5$. Put $w=v_p(n(n-1))$; the leading valuation is $2s+w$. Using \eqref{R3E1:factortwo} and the three-consecutive-factor bound gives, for even $d\ge4$,
\[
 v_p(c_dN^d)\ge ds+w-1-\lfloor\log_p d\rfloor
 \ge2s+w+1.
\]
By \eqref{R3E1:endpointbound}, $ns>2s+w$. This proves \eqref{R3E1:odd} and the claimed congruence. In particular the argument includes $p=3$ in the first case; no denominator factor has been discarded there.
\end{proof}

\begin{proof}[The prime two]
For even $n$, the leading term $NB_n$ has valuation $s-1$. Other odd-degree terms have valuation at least $ds-1>s-1$, since $d$ is odd and hence a unit at two. The endpoint has valuation $ns-1>s-1$.

For odd $n\ge3$, the leading term $(n/2)B_{n-1}N^2$ has valuation $2s-2$, since $n$ is odd. Other even-degree terms have $d\ge4$ and valuation at least
\[
 ds-1-v_2(d)>2s-2.
\]
Indeed $(d-2)s+1>v_2(d)$, as $v_2(d)\le d-2$ for $d\ge4$. The endpoint has valuation $ns-1>2s-2$. This proves \eqref{R3E1:two}.
\end{proof}

\section{The strict gap and the printed conjecture}
\begin{proof}[Proof of the endpoint assertions]
The preceding proofs already compared each odd-prime leading term strictly against the endpoint $N^n/2$. Because two is a unit, this proves \eqref{R3E1:gap}. Alternatively substitute the explicit formulas and use \eqref{R3E1:endpointbound} and $v_p(n)\le n-2$. For $p=2$, the values $s-1$ and $2s-2$ are strictly less than $ns$ for the stated exponents.

If $v_p(X)<v_p(Y)$, factoring the smaller power of $p$ shows that $X-Y$ has valuation $v_p(X)$. Apply this with $X=S_n(N)$ and $Y=N^n$ to obtain endpoint stability for $n\ge2$, and with $Y=-cN^n$ for the perturbation assertion. A regular odd prime makes every even exponent admissible, so this proves the full stated regular-prime case except $n=1$.

Finally,
\[
 S_1(N)=\frac{N(N+1)}2,\qquad S_1(N-1)=\frac{N(N-1)}2.
\]
If $p\mid N$, both $N-1$ and $N+1$ are units at $p$. The two valuations are therefore $v_p(N)-v_p(2)$, proving the remaining case, including $p=2$. Taking $N=pm$ establishes exactly \eqref{R3E1:target}.
\end{proof}

\section{Provenance, checks, and limitations}
\paragraph{Prior work and the actual increment.}
The regular-prime versions of \eqref{R3E1:even}--\eqref{R3E1:odd} are already encoded in \cite[Theorems 6.1--6.3]{R3E1:AMSV}. We do not regard rederiving them as a new theorem. Our proof explicitly supplies their coefficient domination, treats exponent one separately, and extracts the strict bound needed for the assertion still printed as a conjecture. The exponent-local admissibility version removes the unnecessary global regularity assumption. Baoulina--Moree \cite[Section 2 and proof of Theorem 6]{R3E1:BM} use the same classical Kummer--Faulhaber divisibility strategy in studying integer ratios of consecutive power sums. This is an important methodological antecedent, not a technique invented here. We do not solve their integer-ratio problem. Sondow--Tsukerman \cite{R2E1:ST} is another relevant primary source for power-sum valuation formulas.

\paragraph{Delta from the experiment's prior rounds.}
Round 1 Entry 01 treated a cubic block congruence; Round 2 Entry 01 classified the roots of $S_n(pa)-aS_n(p)$ under Bernoulli conditions. Neither theorem is used here. The shared Bernoulli bookkeeping is rederived. All six frozen Round 2 papers and their assessment were read; no competitor theorem is needed. Previously corrected source-status claims concerning strong Seymour vertices and the fixed-parameter gamma question are not revived.

\paragraph{What the theorem does not imply.}
At an irregular residue the displayed regular-residue formulas can fail: for $p=N=37$, $n=32$, direct exact summation gives $v_{37}(S_{32}(37))=2$, not $1$. This is not a counterexample to endpoint stability; it is an obstruction to extending the present formula without additional information. The unrestricted irregular-residue endpoint question is not settled here. We make no claim about Erd\H{o}s--Moser or a broad new principle in analytic number theory.

\paragraph{Verification.}
The companion standard-library \texttt{verify.py} completes 12,345 coefficient audits, 8,932 direct-sum audits and 9,296 endpoint checks. It constructs rational Bernoulli numbers, checks Kummer valuation consequences in the admissible cases, and checks every coefficient against the proposed unique least term. It also computes the defining power sums independently by modular exponentiation, verifying both endpoint valuations. Its bounded ranges are primes at most $101$, exponents at most $160$, and $v_p(N)=1,2,3$ for coefficient audits; direct sums use exponents at most $60$ and $v_p(N)\le2$. It separately tests exponent one and the irregular example. These are diagnostics only: the all-exponent proof is above.

\paragraph{Status and failed alternatives.}
Merely quoting the source's weak valuation bound would leave a cancellation gap. Reusing the Round 2 three-root theorem would address a different polynomial and does not resolve the endpoint question. Dropping the Bernoulli admissibility condition breaks the exact formula, as the example shows. The defensible result is a proof of the printed regular-prime assertion, together with a local sufficient condition, using established arithmetic. A broader literature review may reveal an earlier implicit or explicit proof. No first-proof, major-breakthrough, or exceptional-journal claim is made.



\chapter[A Wilson-quotient dichotomy for exponent valuation trees]{A Wilson-quotient dichotomy for exponent valuation trees}
\label{source:R4E1}
\noindent\textit{Source development: Round 4, entry 01. Original mathematical arguments retained; notation is local to this chapter.}\par\medskip

\begin{quote}\small\noindent\textbf{Chapter abstract.} 
For an odd prime $p$, consider $T_p(n)=\sum_{a=1}^{p+1}a^n$ as the positive exponent $n$ varies. We prove a dichotomy governed by the Wilson quotient modulo $p$. If that quotient is $1$, the valuation is always zero or one, with an exact elementary classification. Otherwise a single $p$-adic root controls every divisibility level outside an explicitly bounded initial segment. This gives exact finite counting formulas and geometric valuation densities. The specialization $p=5$ proves an eventual version of the continuing-branch assertion in Amdeberhan--Moll--Sharma--Vignat's Conjecture 5.1, but does not establish its unrestricted all-exponent version. We isolate the nonunit summand obstruction and prove exponential separation of its possible exceptional exponents.
\end{quote}

\section{Statement and scope}
Write $v_p(0)=+\infty$, and let
\[
 W_p=\frac{(p-1)!+1}{p},\qquad U_p=\{1,\ldots,p-1,p+1\},\qquad
 g_p(m)=\frac1p\sum_{a\in U_p}a^{(p-1)m}\quad(m\geq0).
\]
The last expression is an integer, since its $p$ summands are each $1$ modulo $p$.

\begin{srcAStheorem}\label{R4E1:main}
Let $p$ be an odd prime and $n\geq1$.
\begin{enumerate}
\item If $p-1\nmid n$, then $v_p(T_p(n))=0$.
\item If $W_p\equiv1\pmod p$ and $p-1\mid n$, then $v_p(T_p(n))=1$.
\item If $W_p\not\equiv1\pmod p$, then $g_p$ extends to an isometric bijection of $\srcASZp$ to itself. Its unique zero $\alpha_p$ satisfies
\[
 \alpha_p\equiv-(W_p-1)^{-1}\pmod p,\qquad
 v_p(g_p(x))=v_p(x-\alpha_p)\quad(x\in\srcASZp).
\]
For every $h\geq1$ and every $n\geq h$,
\begin{equation}\label{R4E1:tail}
 p^h\mid T_p(n)\quad\Longleftrightarrow\quad
 n=(p-1)m\text{ and }m\equiv\alpha_p\pmod {p^{h-1}}.
\end{equation}
\end{enumerate}
\end{srcAStheorem}

Thus in the third case, at every height there is exactly one continuing residue class on the tail of the exponent sequence. The cutoff $n\geq h$ is part of the theorem, not a suppressed technicality.

For $p=5$, $W_5=5$, so this case applies to $T_5(n)=S_6(n)$ in the notation of \cite{R3E1:AMSV}. Section 5.5 of that paper develops a $5$-adic exponent tree and states Conjecture 5.1 about its indefinitely continuing branch. Our result proves a precise tail version at every height, together with an all-prime criterion and distribution results. It does \emph{not} prove that every positive exponent is described by a single uncorrected root formula. The printed conjecture contains notation/indexing inconsistencies; we use its surrounding continuing-branch description rather than silently treating those inconsistencies as a new conjecture.

\section{Elementary lifting proof}
For $a$ prime to $p$, put $q_p(a)=(a^{p-1}-1)/p$, and set $A_p=\sum_{a\in U_p}q_p(a)$.
\begin{srcASlemma}\label{R4E1:lerch}
$A_p\equiv W_p-1\pmod p$.
\end{srcASlemma}
\begin{proof}
Multiplication modulo $p^2$ gives
\[
 1+p\sum_{a=1}^{p-1}q_p(a)
 \equiv\prod_{a=1}^{p-1}(1+pq_p(a))
 =((p-1)!)^{p-1}
 =(-1+pW_p)^{p-1}\equiv1+pW_p\pmod {p^2}.
\]
Also $(p+1)^{p-1}\equiv1-p\pmod {p^2}$, so $q_p(p+1)\equiv-1\pmod p$. This proves the assertion. The underlying sum-of-Fermat-quotients identity is classical (Lerch); no novelty is claimed for it.
\end{proof}

\begin{srcASlemma}\label{R4E1:lift}
If $b\geq0$, $r\geq0$, $t\geq1$ and $p\nmid t$, then
\begin{equation}\label{R4E1:isom}
 g_p(b+p^rt)-g_p(b)\equiv p^rtA_p\pmod {p^{r+1}}.
\end{equation}
Moreover $g_p(m)\equiv1+mA_p\pmod p$ for every $m\geq0$.
\end{srcASlemma}
\begin{proof}
For any integer $q$, the binomial theorem gives
\[
 (1+pq)^{p^rt}\equiv1+p^{r+1}tq\pmod {p^{r+2}}.
\]
For completeness, the case $r=0$ follows by discarding terms of degree at least two. Raising a congruence $1+p^{r+1}tq$ modulo $p^{r+2}$ to the $p$th power gives the next case: all terms of degree at least two have valuation at least $r+3$, including when $r=0$, because $p$ is odd. This proves the formula by induction. Apply it to $q=q_p(a)$, multiply by $a^{(p-1)b}\equiv1\pmod p$, sum over $U_p$, and divide by $p$. The second assertion follows directly by expanding $(1+pq_p(a))^m$ modulo $p^2$, using $g_p(0)=1$.
\end{proof}

\begin{proof}[Proof of Theorem \ref{R4E1:main}]
The finite-field power sum $\sum_{a=1}^{p-1}a^n$ is $0$ modulo $p$ unless $p-1\mid n$, in which case it is $-1$. One proof uses a generator of the cyclic group $\mathbb F_p^\times$ and the geometric-series formula. Adding $p^n$ and $(p+1)^n$ proves the first assertion.

Suppose $n=(p-1)m$, $m\geq1$. Then
\begin{equation}\label{R4E1:actual}
 T_p(n)=p\bigl(g_p(m)+p^{n-1}\bigr).
\end{equation}
Since $n\geq2$, Lemmas \ref{R4E1:lerch} and \ref{R4E1:lift} give the second assertion when $A_p\equiv0\pmod p$.

If $A_p$ is a unit, (\ref{R4E1:isom}) shows for distinct nonnegative integers $b,m$ that
\[
 v_p(g_p(m)-g_p(b))=v_p(m-b).
\]
Nonnegative integers are dense in $\srcASZp$, and completeness therefore gives a unique isometric extension. To check surjectivity without any analytic theorem, start with $g_p(m)\equiv1+A_pm\pmod p$. There is exactly one residue mapping to any prescribed residue modulo $p$. Given a solution modulo $p^r$, its $p$ lifts have images differing by $p^rjA_p$ modulo $p^{r+1}$, by (\ref{R4E1:isom}); the lift with any prescribed next image digit is unique. The compatible residues converge to the unique preimage of any element of $\srcASZp$. Taking the image zero gives $\alpha_p$, its initial digit, and the valuation identity. Finally, if $n\geq h$, the second term in (\ref{R4E1:actual}) vanishes modulo $p^{h-1}$. The identity for $g_p$ then proves (\ref{R4E1:tail}).
\end{proof}

\section{Exact counting and densities}
Assume in this section that $W_p\not\equiv1\pmod p$. Let $a_r$ be the least nonnegative residue of $\alpha_p$ modulo $p^r$, with $a_0=0$. Write
\[
 N_h(X)=\#\{1\leq n\leq X:p^h\mid T_p(n)\},\quad
 H=\left\lceil\frac h{p-1}\right\rceil,\quad M=\left\lfloor\frac X{p-1}\right\rfloor.
\]
\begin{srcAScorollary}\label{R4E1:count}
For $X\geq1$ and $h\geq1$,
\begin{align*}
 N_h(X)={}&\#\{H\leq m\leq M:m\equiv a_{h-1}\pmod {p^{h-1}}\}\\
 &+\sum_{1\leq m< H,\ m\leq M}\boldsymbol1_{p^h\mid T_p((p-1)m)}.
\end{align*}
In particular,
\[
 \left|N_h(X)-\frac{X}{(p-1)p^{h-1}}\right|
 \leq1+\left\lfloor\frac{h-1}{p-1}\right\rfloor.
\]
The natural density of $\{n:v_p(T_p(n))=k\}$ is $p^{-k}$ for each $k\geq1$, and that of valuation zero is $(p-2)/(p-1)$.
\end{srcAScorollary}
\begin{proof}
Split the multiples of $p-1$ at $h$ and apply (\ref{R4E1:tail}) to the upper part. Without the cutoff, one residue class of $m$ modulo $p^{h-1}$ gives a progression of $n$ of step $(p-1)p^{h-1}$; its count differs from $X$ divided by that step by at most one. Altering its membership at the at most $H-1=\lfloor(h-1)/(p-1)\rfloor$ omitted multiples proves the bound. For each fixed $h$, divide by $X$ and let $X$ tend to infinity. Subtracting the densities for heights $k$ and $k+1$ gives $p^{-k}$. The zero case follows from Theorem \ref{R4E1:main}(1).
\end{proof}

When $W_p\equiv1$, the exact classification instead gives densities $1/(p-1)$ at valuation one and $(p-2)/(p-1)$ at zero, with no higher valuations. Examples are $p=3,11,107$, by direct factorial calculation; no assertion about the infinitude of such primes is involved.

For $p=5$, the nonzero densities are $5^{-k}$ and the height-$h$ density is $1/(4\cdot5^{h-1})$. The first root residues, at moduli $5,25,125,625,3125,15625$, are respectively
\[
 1,\quad6,\quad6,\quad131,\quad1381,\quad1381.
\]
These are computable certificates, not evidence replacing the lifting proof.

\section{What prevents the unrestricted root formula?}
\begin{srcAStheorem}\label{R4E1:boundary}
Assume $W_p\not\equiv1\pmod p$. For $m\geq1$ put $n=(p-1)m$ and $t=1+v_p(m-\alpha_p)$. Then $t$ is finite, and
\[
 v_p(T_p(n))=\begin{cases}t,&t<n,\\ n,&t>n,\end{cases}
 \qquad v_p(T_p(n))\geq n\quad\text{if }t=n.
\]
Define the candidate obstruction set
\[
 E_p=\{m\geq1:1+v_p(m-\alpha_p)\geq(p-1)m\}.
\]
Any two members $m_1<m_2$ satisfy
\[
 m_2-m_1\geq p^{(p-1)m_1-1}.
\]
\end{srcAStheorem}
\begin{proof}
For an ordinary positive integer $m$, $g_p(m)>0$, hence is not zero, so its valuation is finite. Formula (\ref{R4E1:actual}) and the valuation of a sum of two terms of unequal valuations give the three alternatives. If both $m_i$ belong to $E_p$, both $m_i-\alpha_p$ are divisible in $\srcASZp$ by $p^{(p-1)m_1-1}$. Their positive integer difference is therefore at least that power.
\end{proof}

The distinction is substantive: the term $p^n$ is zero modulo $p^h$ only for $h\leq n$. Dropping it at every precision is unjustified. The theorem neither proves $E_p$ empty nor proves it finite, and at $t=n$ additional cancellation is possible. Exponential separation alone does not resolve those questions. The finite correction in Corollary \ref{R4E1:count} remains exact even if such exceptions occur. Our argument proves the entire distribution law without resolving this boundary problem. We do not infer moment convergence from pointwise density formulas.

\section{Provenance, novelty limits, and reproducibility}
The motivating source is the published 2025 version of \cite{R3E1:AMSV}, inspected on 20 October 2026, especially Section 5.5, Lemmas 5.2--5.4 and Conjecture 5.1. That source already proves the initial two continuing digits for $p=5$. It also attributes many results to Sharma's 2022 thesis, whose full text was not obtained in this audit. Consequently neither priority nor the continued open status of the unrestricted assertion is certified here. Searches for the power-sum exponent tree and its Wilson-quotient connection did not locate this exact all-prime dichotomy, but an unsuccessful search is not a novelty proof. The Fermat-quotient congruence in Lemma \ref{R4E1:lerch} is classical; \cite{R4E1:Sondow} discusses this tradition. The lifting mechanism is elementary Hensel-type reasoning, not a new general lifting method.

Relative to this entrant's Round 3, the variable has changed: that paper varied the endpoint with a fixed exponent and used regular-prime Bernoulli valuations. This paper fixes endpoint $p+1$, varies the exponent, imposes no regularity condition, and gives a Wilson-quotient criterion, counting law and explicit obstruction. It is not another version of the Round 3 endpoint domination theorem. The Round 2 roots concerned a different block-defect polynomial; none of its results is used here. Other entrants' frozen Round 3 results were read as required, but are not used in this proof.

The supplied standard-library Python verifier checks the Fermat/Wilson identity, lifting congruences, direct modular sums, the tail equivalence, finite counting corrections, and root digits. All computations, including the discrepancy check after clearing denominators, use exact integers. Finite absence of boundary candidates is reported only as a finite check. Run \texttt{python3 verify.py}; \texttt{build.sh} rebuilds the PDF in the tournament's TeX environment.



\chapter[Finite boundary exceptions and mean valuations of power sums]{Finite boundary exceptions and mean valuations of power sums}
\label{source:R5E1}
\noindent\textit{Source development: Round 5, entry 01. Original mathematical arguments retained; notation is local to this chapter.}\par\medskip

\begin{quote}\small\noindent\textbf{Chapter abstract.} 
Let $T_p(n)=\sum_{a=1}^{p+1}a^n$ for an odd prime $p$. Combining the preceding round's Wilson-quotient exponent tree with the classical $p$-adic Subspace Theorem, we prove that its boundary exceptions are finite and that its ordinary mean valuation exists. The mean is $1/(p-1)$ when the Wilson quotient is $1$ modulo $p$, and $p/(p-1)^2$ otherwise. In the latter case the summatory valuation equals its linear main term plus the largest individual valuation, with an $O_p(\log X)$ remainder. The finiteness and convergence are ineffective. An explicit $p$-adic root model with the same geometric level densities and a linear pointwise bound nevertheless has two different Ces\`aro limit points, showing why a tail argument is required. The Subspace Theorem and its power-sum consequence are known inputs, not discoveries; no all-exponent resolution of the motivating valuation-tree conjecture is claimed.
\end{quote}

\section{Main results and the increment}
For an odd prime $p$, let $T_p(n)=\sum_{b=1}^{p+1}b^n$ and set
\begin{gather*}
 a=p-1,\qquad W_p=\frac{(p-1)!+1}{p},\qquad V_p(n)=v_p(T_p(n)),\\
 S_p(X)=\sum_{1\le n\le X}V_p(n),\qquad
 M_p(X)=\max_{1\le n\le X}V_p(n).
\end{gather*}
We use $v_p(p)=1$ and $v_p(0)=+\infty$. All power sums appearing in these definitions are positive integers.

\begin{srcAYtheorem}[Unconditional ordinary mean]\label{R5E1:main}
For every odd prime $p$,
\begin{equation}\label{R5E1:mean}
 \lim_{X\to\infty}\frac{S_p(X)}X=\mu_p,
 \qquad
 \mu_p=\begin{cases}
 \dfrac1{p-1},&W_p\equiv1\pmod p,\\[4pt]
 \dfrac p{(p-1)^2},&W_p\not\equiv1\pmod p.
 \end{cases}
\end{equation}
In the first case $S_p(X)=\lfloor X/(p-1)\rfloor$ exactly. In the second case,
\begin{equation}\label{R5E1:maxidentity}
 S_p(X)=\mu_pX+M_p(X)+O_p(\log(2+X)),\qquad M_p(X)=o(X).
\end{equation}
No effective modulus of convergence is asserted.
\end{srcAYtheorem}

Thus, for example, the mean at $p=5$ is $5/16$. Equivalently,
\[
 v_p\left(\prod_{n=1}^N T_p(n)\right)=\mu_pN+o(N).
\]
The product reformulation is only a consequence, not an additional discovery.

To state the structural improvement, put
\[
 U_p=\{1,2,\ldots,p-1,p+1\},\qquad
 F_p(m)=\sum_{b\in U_p}b^{(p-1)m},\qquad g_p(m)=F_p(m)/p.
\]
When $W_p\not\equiv1\pmod p$, the earlier elementary lifting argument gives a unique $\alpha_p\in\srcAYZp$ with
\begin{equation}\label{R5E1:root}
 v_p(g_p(m))=v_p(m-\alpha_p)\quad(m\ge0).
\end{equation}
We reprove this input below.

\begin{srcAYtheorem}[Finite boundary set]\label{R5E1:finite}
Suppose $W_p\not\equiv1\pmod p$. The set
\begin{equation}\label{R5E1:E}
 E_p=\{m\ge1:1+v_p(m-\alpha_p)\ge(p-1)m\}
\end{equation}
is finite. Outside this finite set,
\begin{equation}\label{R5E1:eventual}
 V_p((p-1)m)=1+v_p(m-\alpha_p).
\end{equation}
In particular, there is one cutoff in the exponent beyond which the root formula is valid at every precision simultaneously. Neither that cutoff nor emptiness of $E_p$ is proved effectively.
\end{srcAYtheorem}

Round 4 \cite{R5E1:R4} proved geometric densities at each fixed height and an exact tail criterion restricted to $n\ge h$. It proved exponential separation of $E_p$, but not its finiteness, and deliberately did not infer mean convergence. Theorems \ref{R5E1:main} and \ref{R5E1:finite} close those two particular gaps using a classical deep input. They do not prove the unrestricted, all-positive-exponent continuing-branch assertion motivating Conjecture 5.1 of \cite{R3E1:AMSV}. A finite exceptional set is not an empty exceptional set.

\section{The inherited elementary root classification}
The following is the exact Round 4 input, supplied with proof for standalone use. The lifting method is classical; compare the Fermat-quotient criterion in \cite[Theorem 1]{R5E1:BD}.

\begin{srcAYlemma}\label{R5E1:lifting}
If $p-1\nmid n$, then $V_p(n)=0$. If $W_p\equiv1\pmod p$, then $V_p(n)=1$ for every positive multiple of $p-1$. Otherwise $g_p$ extends to an isometric bijection $\srcAYZp\to\srcAYZp$ with unique zero $\alpha_p$, satisfying \eqref{R5E1:root} and
\[
 \alpha_p\equiv-(W_p-1)^{-1}\pmod p.
\]
For $m\ge1$, let $t=1+v_p(m-\alpha_p)$ and $n=(p-1)m$. Then $t$ is finite and
\[
 V_p(n)=t\quad(t<n),\qquad V_p(n)=n\quad(t>n),\qquad
 V_p(n)\ge n\quad(t=n).
\]
\end{srcAYlemma}
\begin{proof}
The finite-field power sum over $\mathbb F_p^\times$ is zero unless $p-1\mid n$, when it is $-1$. Adding the terms $p^n$ and $(p+1)^n$ proves the zero-valuation assertion.

Put $q_p(b)=(b^{p-1}-1)/p$ and $A=\sum_{b\in U_p}q_p(b)$. Multiplication modulo $p^2$ gives
\[
 1+p\sum_{b=1}^{p-1}q_p(b)
 \equiv((p-1)!)^{p-1}=(-1+pW_p)^{p-1}
 \equiv1+pW_p\pmod {p^2}.
\]
Also $q_p(p+1)\equiv-1\pmod p$, so $A\equiv W_p-1\pmod p$. This is the classical Fermat/Wilson quotient relation.

For $r\ge0$, $p\nmid u$ and $u>0$, binomial expansion gives
\[
 (1+pq)^{p^ru}\equiv1+p^{r+1}uq\pmod {p^{r+2}}.
\]
The case $r=0$ is immediate; raising to the $p$th power proves the induction step, since $p$ is odd. Therefore, for $b\ge0$,
\begin{equation}\label{R5E1:increment}
 g_p(b+p^ru)-g_p(b)\equiv p^ruA\pmod {p^{r+1}},
 \qquad g_p(m)\equiv1+mA\pmod p.
\end{equation}
If $A$ is a unit, the first congruence gives
$v_p(g_p(m)-g_p(b))=v_p(m-b)$. Density and completeness extend this isometry to $\srcAYZp$. The second congruence gives a unique preimage of each residue modulo $p$; the first gives exactly one lift with any prescribed next image digit. This proves surjectivity and uniqueness at every precision, hence the root assertion. For an ordinary $m\ge0$, $g_p(m)>0$, so the root is not any nonnegative integer and the valuation in \eqref{R5E1:root} is finite.

Finally,
\begin{equation}\label{R5E1:nonunit}
 T_p((p-1)m)=F_p(m)+p^{(p-1)m}
             =p\bigl(g_p(m)+p^{(p-1)m-1}\bigr).
\end{equation}
If $A\equiv0\pmod p$, $g_p(m)\equiv1\pmod p$ and $(p-1)m\ge2$, giving valuation one. In the unit case, compare the valuations of the two terms in \eqref{R5E1:nonunit}. Unequal valuations give their minimum; equal valuations give the stated lower bound.
\end{proof}

\section{The known Diophantine input, with its normalization checked}
We use the following special case of Schlickewei's $p$-adic Subspace Theorem. It is stated explicitly in this normalization in Evertse's author-hosted notes \cite[Theorem 1.1, with equal numbers of forms and variables]{R5E1:Evertse}.

\paragraph{Classical Subspace Theorem (input).}
Let $S$ be a finite set of places of $\mathbb Q$ containing infinity. At each $v\in S$ choose $d$ linearly independent rational linear forms $L_{1,v},\ldots,L_{d,v}$ in $d\ge2$ variables. For each $\delta>0$, the primitive integer vectors $x\ne0$ satisfying
\begin{equation}\label{R5E1:ST}
 \prod_{v\in S}\prod_{i=1}^d |L_{i,v}(x)|_v
 \le \|x\|_\infty^{-\delta}
\end{equation}
lie in finitely many proper rational linear subspaces. Here $|q|_q=q^{-1}$ for a prime $q$, and $\|x\|_\infty=\max_i|x_i|$. This theorem is invoked, not reproved; its subspaces are not effectively determined by the version used.

\begin{srcAYlemma}[Standard power-sum consequence]\label{R5E1:sublinear}
Let $1=b_1<b_2<\cdots<b_d$ be positive integers, $d\ge2$, and let $p$ be any prime. Then
\[
 v_p\left(\sum_{i=1}^d b_i^k\right)=o(k)\qquad(k\to\infty).
\]
This is a standard application of the classical input, not a new valuation estimate claimed by this work.
\end{srcAYlemma}
\begin{proof}
Take $x(k)=(1,b_2^k,\ldots,b_d^k)$, which is primitive. Let $S$ consist of infinity, $p$, and all prime divisors of $b_2\cdots b_d$. At places other than $p$, use all coordinate forms. At $p$, replace the first coordinate form by $X_1+\cdots+X_d$ and keep $X_2,\ldots,X_d$. These forms are linearly independent, with coefficient determinant one. The product formula for each coordinate, and $x_1=1$, give the exact identity
\begin{equation}\label{R5E1:product}
 \prod_{v\in S}\prod_{i=1}^d|L_{i,v}(x(k))|_v
 =\left|\sum_{i=1}^d b_i^k\right|_p.
\end{equation}
For a prescribed $\varepsilon>0$, a valuation at least $\varepsilon k$ makes this at most $p^{-\varepsilon k}=\|x(k)\|_\infty^{-\delta}$ with
$\delta=\varepsilon\log p/\log b_d>0$.

Each of the finitely many exceptional subspaces has a nonzero rational linear equation $\sum_i c_iX_i=0$. If $j$ is the largest index with $c_j\ne0$, division by $b_j^k$ shows
\[
 \sum_{i\le j}c_i(b_i/b_j)^k\longrightarrow c_j\ne0.
\]
Thus only finitely many $x(k)$ lie in that subspace. There are consequently only finitely many $k$ with valuation at least $\varepsilon k$. Since $\varepsilon$ was arbitrary and all valuations are nonnegative, the conclusion follows. This also checks directly the nondegeneracy issue; no unproved assertion about zeros of general recurrence sequences is used.
\end{proof}

\begin{proof}[Proof of Theorem \ref{R5E1:finite}]
Apply Lemma \ref{R5E1:sublinear} to the distinct positive integer bases
$\{b^{p-1}:b\in U_p\}$, which include $1$. It gives
\[
 1+v_p(m-\alpha_p)=v_p(F_p(m))=o(m).
\]
For all sufficiently large $m$, this is strictly less than $(p-1)m$. Thus \eqref{R5E1:E} is finite, and Lemma \ref{R5E1:lifting} proves \eqref{R5E1:eventual}. Independently, applying Lemma \ref{R5E1:sublinear} to $1,2,\ldots,p+1$ gives $V_p(n)=o(n)$. No finite computation is being used to locate the exceptional subspaces or the eventual cutoff.
\end{proof}

\section{An exact one-root summation estimate}
The next lemma is elementary and is useful independently of power sums. Its explicit largest-term contribution is more informative than merely assuming uniform integrability.

\begin{srcAYlemma}[Largest-term formula]\label{R5E1:summation}
Let $p$ be prime and $\alpha\in\srcAYZp$ not be a positive integer. Put
\[
 f(m)=1+v_p(m-\alpha),\quad
 B(N)=\max_{1\le m\le N}f(m),\quad c=\frac p{p-1}.
\]
For every integer $N\ge1$ and every integer $H\ge1$ with $p^H>N$,
\begin{equation}\label{R5E1:explicit}
 \left|\sum_{m=1}^N f(m)-cN-B(N)\right|
 \le 2H+cNp^{-H}.
\end{equation}
In particular the remainder is $O_p(\log(2+N))$, uniformly in $\alpha$.
\end{srcAYlemma}
\begin{proof}
For $h\ge1$, the condition $f(m)\ge h$ specifies one residue class modulo $p^{h-1}$. Its count in $[1,N]$ differs from $N/p^{h-1}$ by at most one. Thus
\[
 \sum_{m=1}^N\min(f(m),H)
 =cN(1-p^{-H})+R,\qquad |R|\le H.
\]
There is at most one $m\in[1,N]$ with $f(m)>H$, because two such integers would differ by a nonzero multiple of $p^H>N$. Hence the contribution above the truncation is exactly $(B(N)-H)_+$. Subtracting $cN+B(N)$ leaves
$R-cNp^{-H}-\min(B(N),H)$, which proves \eqref{R5E1:explicit}.
\end{proof}

\begin{srcAYcorollary}[An exact mean criterion]\label{R5E1:criterion}
Under the hypotheses of Lemma \ref{R5E1:summation},
\[
 \frac1N\sum_{m=1}^Nf(m)\longrightarrow\frac p{p-1}
 \quad\Longleftrightarrow\quad f(m)=o(m).
\]
\end{srcAYcorollary}
\begin{proof}
The lemma equates convergence to the stated mean with $B(N)=o(N)$. That is equivalent to $f(m)=o(m)$: one direction uses $f(N)\le B(N)$; for the other, fix a cutoff beyond which $f(m)\le\varepsilon m$, and separate the bounded finite initial maximum. The content particular to this setting is the largest-term identity, not this elementary equivalence of pointwise and maximum sublinearity.
\end{proof}

\begin{proof}[Proof of Theorem \ref{R5E1:main}]
The case $W_p\equiv1\pmod p$ follows from Lemma \ref{R5E1:lifting}. In the other case let $f(m)=1+v_p(m-\alpha_p)$ and define the root model on all exponents by
\[
 R_p(n)=\begin{cases}f(m),&n=(p-1)m,\\0,&p-1\nmid n.\end{cases}
\]
By Theorem \ref{R5E1:finite}, $R_p(n)$ differs from $V_p(n)$ at only finitely many exponents. Therefore their summatory functions differ by $O_p(1)$, and their running maxima also differ by $O_p(1)$: the latter follows from the maximum of the finitely many absolute differences. Lemma \ref{R5E1:summation}, with $N=\lfloor X/(p-1)\rfloor$, gives
\[
 \sum_{n\le X}R_p(n)
 =\frac p{p-1}\left\lfloor\frac X{p-1}\right\rfloor
   +\max_{n\le X}R_p(n)+O_p(\log(2+X)).
\]
This proves \eqref{R5E1:maxidentity}'s first assertion. Lemma \ref{R5E1:sublinear} gives $V_p(n)=o(n)$, hence $M_p(X)=o(X)$ by the finite-initial-segment argument in Corollary \ref{R5E1:criterion}. Dividing proves \eqref{R5E1:mean}. The finitely many corrections make the implied constant for the actual sums ineffective; the model bound \eqref{R5E1:explicit} itself is fully explicit.
\end{proof}

\section{Why the arithmetic input cannot be replaced by densities}
The following explicit model provides an obstruction to an invalid shortcut. It is not claimed to arise from $T_p$ or to contradict any result above.

\begin{srcAYtheorem}[Geometric densities with nonconvergent means]\label{R5E1:example}
For every prime $p$, define integers
\[
 m_1=p,\qquad m_{j+1}=m_j+p^{m_j},\qquad
 \beta=\lim_{j\to\infty}m_j\quad\text{in }\srcAYZp,
 \qquad f_\beta(m)=1+v_p(m-\beta).
\]
The limit is not a positive integer. For every $m\ge1$,
$1\le f_\beta(m)\le m+1$, and for every fixed $k\ge1$ the natural density of $f_\beta(m)=k$ is $(p-1)/p^k$. Nevertheless,
\begin{equation}\label{R5E1:oscillation}
 \liminf_{N\to\infty}\frac1N\sum_{m\le N}f_\beta(m)=\frac p{p-1},
 \qquad
 \limsup_{N\to\infty}\frac1N\sum_{m\le N}f_\beta(m)=\frac p{p-1}+1.
\end{equation}
\end{srcAYtheorem}
\begin{proof}
The differences $m_{j+1}-m_j$ have valuations $m_j$ strictly tending to infinity, so the limit exists and
$v_p(\beta-m_j)=m_j$. For $m_j<m<m_{j+1}$, one has
$0<m-m_j<p^{m_j}$, hence $v_p(m-m_j)<m_j$. The unequal-valuation rule now gives
\[
 v_p(m-\beta)=v_p(m-m_j)<m_j.
\]
For $1\le m<p$, the valuation is zero. These formulas show both that $\beta$ cannot equal a positive integer and that $v_p(m-\beta)\le m$. They also show exactly
\[
 \max_{1\le m\le N} f_\beta(m)=m_j+1
 \quad\text{whenever }m_j\le N<m_{j+1}.
\]
For each height $h$, membership is one residue class modulo $p^{h-1}$, so subtraction of consecutive height densities gives $(p-1)/p^k$. Lemma \ref{R5E1:summation} applies. The displayed running maximum divided by $N$ has limsup one, attained at $N=m_j$, and liminf zero, attained at $N=m_{j+1}-1$, since $m_{j+1}=m_j+p^{m_j}$. The $O(\log N)/N$ remainder vanishes, proving \eqref{R5E1:oscillation}.
\end{proof}

This example has the same one-root residue geometry as our arithmetic model, and even a global linear bound, but no ordinary mean. Thus the use of Subspace-Theorem sublinearity is not a disguised appeal to geometric densities alone. We make no claim that this elementary lacunary construction or the general phenomenon is historically new.

\section{Source audit, verification, and remaining questions}
The current published AMSV source \cite{R3E1:AMSV} was rechecked on 20 October 2026. Its Section 5.5 and Conjecture 5.1 motivate the $p=5$ problem; this paper does not announce full closure or certify continued open status. Sharma's 2022 thesis, cited there, was not obtained in full. The earlier notation/indexing caveat remains in force.

The central new ingredient relative to Round 4 is an \emph{application} of known Diophantine approximation, not a new Subspace Theorem. Evertse's stated normalization was checked directly: equal numbers of forms and variables, primitive vectors, normalized prime absolute values, and the exact product \eqref{R5E1:product}. The dominant-base argument checks all exceptional subspaces. The standard $o(k)$ consequence is explicitly credited rather than relabeled as a new power-sum theorem. The published Bayarmagnai--Delger paper \cite{R5E1:BD}, including its final 2022 version, was additionally inspected: its Fermat-quotient lifting criterion concerns unboundedness, reinforcing that this mechanism is established prior art. No historical priority is claimed for the present combination, mean formula, or model construction on the basis of unsuccessful searches. All six frozen Round 4 manuscripts and their assessment, and the Round 3 pool, were read; no competitor result is a mathematical dependency here.

The verifier uses exact integer and rational arithmetic to check the lifting/root model against full power sums in bounded ranges, the product-form normalization, the explicit bound \eqref{R5E1:explicit}, and the lacunary example. It labels finite averages as diagnostics, not convergence proofs. The Subspace Theorem cannot be certified by finite sampling, and no output of the program is presented as an effective eventual cutoff.

The remaining gap is now sharply smaller than in Round 4: $E_p$ is finite, but its elements have not been effectively bounded or shown absent. There is no proved $O(\log n)$ bound for $V_p(n)$ here, no effective error rate in \eqref{R5E1:mean}, and no higher-moment convergence theorem. The first-moment proof does not silently extend to those stronger statements. An effective bound on the exceptional exponents, or a stronger Diophantine approximation bound for $\alpha_p$, is the next substantive target.



\chapter[Distinct-shift correlations and maximal-spike formulas for power-sum valuations]{Distinct-shift correlations and maximal-spike formulas for power-sum valuations}
\label{source:R6E1}
\noindent\textit{Source development: Round 6, entry 01. Original mathematical arguments retained; notation is local to this chapter.}\par\medskip

\begin{quote}\small\noindent\textbf{Chapter abstract.} 
For each odd prime, we prove convergence of every fixed distinct-shift correlation of the valuations of $\sum_{a=1}^{p+1}a^n$, with an explicitly computable rational limit. On the active exponent subsequence the pair correlation at distance $d\ne0$ is $(p+1)(p-p^{-v_p(d)})/(p-1)^2$ in the nonconstant Wilson-quotient case. The mechanism is a general maximal-spike expansion for products of powers of separated $p$-adic root valuations. It isolates every possible higher-moment obstruction in a positive sum of powers of individual running maxima. The preceding round's classical Subspace-Theorem input controls precisely the distinct-shift, degree-one case. A new explicit root model has convergent first mean but nonconvergent second mean, so that input does not justify higher moments. The arithmetic cutoff and convergence remain ineffective; no unrestricted valuation-tree conjecture or higher-moment convergence is claimed.
\end{quote}

\section{Arithmetic correlations and the exact increment}
For an odd prime $p$ write
\[
 a=p-1,\quad T_p(n)=\sum_{b=1}^{p+1}b^n,\quad V_p(n)=v_p(T_p(n)),\quad
 A_p(m)=V_p(am),\quad W_p=\frac{(p-1)!+1}{p}.
\]
Valuations are normalized by $v_p(p)=1$. Let $t_1,\ldots,t_r$ be distinct nonnegative integers, and $d_1,\ldots,d_r$ positive integers. Define
\[
 D=\max_i d_i,\qquad c_i=\prod_{j\ne i}(1+v_p(t_i-t_j))^{d_j},\qquad
 B_i(N)=\max_{1\le m\le N} A_p(m+t_i).
\]
An empty product is one. Let $dx$ denote probability Haar measure on $\srcBEZp$, and put
\begin{equation}\label{R6E1:integral}
 I_p(\mathbf t,\mathbf d)=\int_{\srcBEZp}\prod_{i=1}^r(1+v_p(x+t_i))^{d_i}\,dx.
\end{equation}
The values at its finitely many singular points do not affect this integral.

\begin{srcBEtheorem}[Mixed-moment expansion]\label{R6E1:main}
Suppose $W_p\not\equiv1\pmod p$. For fixed $\mathbf t,\mathbf d$ as above,
\begin{equation}\label{R6E1:expansion}
 \sum_{m=1}^N\prod_i A_p(m+t_i)^{d_i}
 =N I_p(\mathbf t,\mathbf d)+\sum_i c_i B_i(N)^{d_i}
       +O_{p,\mathbf t,\mathbf d}(\log^D(2+N)).
\end{equation}
The integral is finite, rational, and computable by the finite formula in Section~\ref{R6E1:finiteformula}. The implied constant for these arithmetic sums can be ineffective. In particular,
\begin{equation}\label{R6E1:correlation}
 \lim_{N\to\infty}\frac1N\sum_{m=1}^N\prod_i A_p(m+t_i)
 =I_p(\mathbf t,\mathbf 1).
\end{equation}
For arbitrary positive $\mathbf d$, convergence to the integral in \eqref{R6E1:expansion} holds if and only if $B_i(N)^{d_i}=o(N)$ for every $i$.
If $W_p\equiv1\pmod p$, then $A_p(m)=1$ for all $m\ge1$, so every displayed mixed average is exactly one instead.
\end{srcBEtheorem}

\begin{srcBEcorollary}[All exponent shifts and explicit pairs]\label{R6E1:pairs}
Let $h_1,\ldots,h_r$ be distinct nonnegative integers. The average
$X^{-1}\sum_{n\le X}\prod_i V_p(n+h_i)$ converges. Its limit is zero if the $h_i$ are not all congruent modulo $a$. Otherwise write $h_i=h_*+at_i$, where $h_*=\min_i h_i$. The limit is $1/a$ in the constant Wilson case, and $I_p(\mathbf t,\mathbf1)/a$ otherwise.
In the nonconstant case, for $d\ge1$ and $s=v_p(d)$,
\begin{equation}\label{R6E1:pairconstant}
 \lim_{N\to\infty}\frac1N\sum_{m=1}^N A_p(m)A_p(m+d)
 =\frac{(p+1)(p-p^{-s})}{(p-1)^2}.
\end{equation}
For example, at $p=5$ the mean of $V_5(n)V_5(n+4)$ is $3/8$, whereas that of $V_5(n)V_5(n+1)$ is zero.
\end{srcBEcorollary}

Round 4 entry 01 established the root classification and fixed-height densities; Round 5 entry 01 proved finite boundary exceptions, the ordinary first mean, and its largest-term correction \cite{R6E1:R4,R6E1:R5}. These are inherited inputs, reproved below. The increment here is the separated-root reduction, the general polynomial maximal-spike expansion, and unconditional correlations of \emph{arbitrarily many distinct shifts}. These are not ordinary higher moments at a coincident shift. Our explicit square-moment obstruction in Section~\ref{R6E1:obstruction} distinguishes them.

The published source \cite[Section 5.5]{R3E1:AMSV} motivates the $p=5$ exponent tree, but these statistical results do not settle its all-exponent assertion. No claim of current open status or historical priority is based on unsuccessful searches. All six frozen Round 5 papers and their assessment were read; none of the other five papers is a mathematical dependency here.

\section{Separated singularities reduce to one-root moments}
We first prove a statement independent of power sums. Fix a prime $p$ and $\alpha\in\srcBEZp$ such that $\alpha$ is not a positive integer. All offsets below are nonnegative, so
\[
 f_i(m)=1+v_p(m+t_i-\alpha)
\]
is finite for positive integers $m$. Let $b_i(N)=\max_{m\le N}f_i(m)$.

\begin{srcBElemma}[One-root polynomial summation]\label{R6E1:single}
For an integer $k\ge1$, put
\[
 \mu_{p,k}=\int_{\srcBEZp}(1+v_p(x))^k\,dx
 =\sum_{j\ge1}(j^k-(j-1)^k)p^{-(j-1)}.
\]
For every such root $\gamma$ and $f(m)=1+v_p(m-\gamma)$, with finite positive-integer values,
\begin{equation}\label{R6E1:singleformula}
 \sum_{m=1}^N f(m)^k=N\mu_{p,k}+\bigl(\max_{m\le N}f(m)\bigr)^k
                 +O_{p,k}(\log^k(2+N)),
\end{equation}
with a constant independent of $\gamma$.
\end{srcBElemma}
\begin{proof}
Choose $H\ge1$ minimally with $p^H>N$ and write $\Delta_j=j^k-(j-1)^k$. Each condition $f(m)\ge j$ is one residue class modulo $p^{j-1}$, whose count differs from $N p^{-(j-1)}$ by at most one. Hence
\[
 \sum_{m\le N}\min(f(m),H)^k
 =N\sum_{j=1}^H\Delta_j p^{-(j-1)}+R,\qquad |R|\le H^k.
\]
There is at most one integer with $f(m)>H$, since its residue class has modulus $p^H>N$. If $b=\max f(m)$, the omitted contribution is exactly $(b^k-H^k)_+$. The difference from $b^k$ has absolute value at most $H^k$. Finally,
\[
 N\sum_{j>H}\Delta_j p^{-(j-1)}=O_{p,k}(H^{k-1})
\]
because $Np^{-H}<1$ and a polynomial times a geometric tail is bounded by a constant times its initial polynomial value. These estimates prove the assertion. The series for $\mu_{p,k}$ is the layer formula for the geometric valuation distribution, and is rational by finite differentiation of a geometric series.
\end{proof}

\begin{srcBElemma}[Periodic remainder]\label{R6E1:periodic}
Let $r\ge2$, $H=1+\max_{i\ne j}v_p(t_i-t_j)$ and $P=p^H$. On $\srcBEZp$ away from the roots $\alpha-t_i$, the function
\begin{equation}\label{R6E1:remainder}
 R(x)=\prod_i(1+v_p(x+t_i-\alpha))^{d_i}
          -\sum_i c_i(1+v_p(x+t_i-\alpha))^{d_i}
\end{equation}
extends to a bounded function constant on residue classes modulo $P$.
\end{srcBElemma}
\begin{proof}
The root classes $x\equiv\alpha-t_i\pmod {P}$ are distinct. Off all of them every valuation is less than $H$ and is determined by the residue class. In the class of root $i$, for $j\ne i$ the unequal-valuation rule gives
\[
 v_p(x+t_j-\alpha)=v_p(t_j-t_i)<H.
\]
Thus the product in \eqref{R6E1:remainder} is exactly $c_i(1+v_p(x+t_i-\alpha))^{d_i}$. This cancels its corresponding term, and the remaining terms are constant on that class. Defining the value at the root by that constant proves the claim.
\end{proof}

\begin{srcBEtheorem}[General root-model expansion]\label{R6E1:model}
With the preceding definitions, for every $r\ge1$,
\begin{equation}\label{R6E1:modelidentity}
 \sum_{m\le N}\prod_i f_i(m)^{d_i}
 =N I_p(\mathbf t,\mathbf d)+\sum_i c_i b_i(N)^{d_i}
      +O_{p,\mathbf t,\mathbf d}(\log^D(2+N)).
\end{equation}
The constant is uniform in $\alpha$. Consequently the average tends to the integral if and only if every $b_i(N)^{d_i}=o(N)$, equivalently $f_i(m)=o(m^{1/d_i})$ for every $i$.
\end{srcBEtheorem}
\begin{proof}
For $r=1$ this is Lemma~\ref{R6E1:single}. For $r\ge2$, Lemma~\ref{R6E1:periodic} gives
\[
 \sum_{m\le N}R(m)=N\int R(x)\,dx+O_{p,\mathbf t,\mathbf d}(1).
\]
Indeed the sum over each full block of $P$ integers is the sum over all residue classes. The bound on $R$ depends only on the offset distances and powers, not on the translation $\alpha$. Add Lemma~\ref{R6E1:single} for each summand $c_i f_i^{d_i}$. Translation invariance of Haar measure identifies the total main term with \eqref{R6E1:integral}. This also proves finiteness of that integral. Every $c_i$ is strictly positive, so the sum of maximum corrections is $o(N)$ exactly when each correction is. Finally a finite-initial-segment argument shows that $\max_{m\le N}f_i(m)^{d_i}=o(N)$ is equivalent to $f_i(m)^{d_i}=o(m)$.
\end{proof}

\section{Finite evaluation of the rational constants}\label{R6E1:finiteformula}
For $r\ge2$ use $H,P$ from Lemma~\ref{R6E1:periodic}; for $r=1$ take $H=1$ and $P=p$. Put $v_p(0)=+\infty$ when taking a minimum, and define
\[
 J=\frac1P\sum_{x=0}^{P-1}\prod_i\bigl(1+\min\{v_p(x+t_i),H\}\bigr)^{d_i},\qquad
 Q_{p,k}(H)=\sum_{j=H+2}^{\infty}\Delta_j p^{-(j-1)}.
\]
Then the exact formula is
\begin{equation}\label{R6E1:finiteconstant}
 I_p(\mathbf t,\mathbf d)=J+\sum_i c_i Q_{p,d_i}(H).
\end{equation}
To see this, cap each valuation at $H$. Only one factor can exceed the cap at any point, and all other factors there are the constants giving $c_i$. Integrate the difference by layers. The infinite-looking $Q$ is explicitly rational: $\Delta_j$ is a polynomial of degree $k-1$ in $j$, and each series $\sum_{\ell\ge0}\ell^u z^\ell$ equals $(z\,d/dz)^u(1-z)^{-1}$ at $z=1/p$. Thus \eqref{R6E1:finiteconstant} is a finite algorithm using rational arithmetic. In particular $Q_{p,1}(H)=p^{-H}/(p-1)$.

For the pair formula write $f_0(x)=1+v_p(x)$ and $f_d(x)=1+v_p(x+d)$, and expand each in layers indexed by $i,j\ge0$. The two residue conditions are compatible exactly when $\min(i,j)\le s=v_p(d)$; their intersection then has measure $p^{-\max(i,j)}$. Therefore
\[
 \int f_0f_d=\sum_{i,j\ge0}p^{-\max(i,j)}
             -\sum_{i,j\ge s+1}p^{-\max(i,j)}
 =\frac{p(p+1)}{(p-1)^2}-\frac{(p+1)p^{-s}}{(p-1)^2}.
\]
Both geometric double sums can be evaluated by grouping equal maxima; the first has $2k+1$ pairs of maximum $k$. This proves \eqref{R6E1:pairconstant} once the arithmetic reduction is established.

\section{Inherited arithmetic reduction, with complete arguments}
The statements in this section are the Round 4 lifting theorem and Round 5 finite-exception theorem \cite{R6E1:R4,R6E1:R5}, not new contributions. Put
\[
 U_p=\{1,\ldots,p-1,p+1\},\qquad
 F_p(m)=\sum_{b\in U_p}b^{am},\qquad g_p(m)=F_p(m)/p.
\]
\begin{srcBElemma}\label{R6E1:arithmetic}
If $a\nmid n$, then $V_p(n)=0$. If $W_p\equiv1\pmod p$, then $A_p(m)=1$ for all positive $m$. Otherwise there is a unique $\alpha_p\in\srcBEZp$, not a nonnegative integer, with
\[
 v_p(g_p(m))=v_p(m-\alpha_p)\quad(m\ge0).
\]
Moreover $f(m)=1+v_p(m-\alpha_p)=o(m)$, and $A_p(m)=f(m)$ outside a finite set.
\end{srcBElemma}
\begin{proof}
The finite-field unit power sum is zero unless $a\mid n$, when it is $-1$; the additional $(p+1)^n$ contributes one modulo $p$. This proves the first assertion.
Let $q_p(b)=(b^a-1)/p$ and $C=\sum_{b\in U_p}q_p(b)$. Multiplication modulo $p^2$ gives
\[
 1+p\sum_{b=1}^{p-1}q_p(b)\equiv((p-1)!)^{p-1}
 \equiv(-1+pW_p)^{p-1}\equiv1+pW_p\pmod {p^2}.
\]
Since $q_p(p+1)\equiv-1\pmod p$, $C\equiv W_p-1\pmod p$. Binomial expansion, followed by induction on $r$, gives
$(1+pq)^{p^ru}\equiv1+p^{r+1}uq\pmod {p^{r+2}}$ for positive $u$ prime to $p$. Consequently
\[
 g_p(b+p^ru)-g_p(b)\equiv p^ru C\pmod {p^{r+1}},\qquad
 g_p(m)\equiv1+mC\pmod p.
\]
If $C$ is a unit, this is an isometry on nonnegative integers and extends by completeness to an isometry of $\srcBEZp$. The second congruence gives a unique preimage of each image residue modulo $p$, and the first gives exactly one lift for each successive image digit. The extension is therefore bijective. Its unique zero $\alpha_p$ satisfies the stated valuation identity. Positivity of $g_p(m)$ excludes a nonnegative integer zero. Finally
\begin{equation}\label{R6E1:boundary}
 A_p(m)=v_p(F_p(m)+p^{am}).
\end{equation}
If $C\equiv0\pmod p$, then $F_p(m)$ has valuation one, whereas $am\ge2$, proving the constant case.

For completeness the classical Diophantine input used in Round 5 is stated in its exact normalization. Schlickewei's $p$-adic Subspace Theorem, as in Evertse \cite[Theorem 1.1]{R5E1:Evertse}, says: for finitely many places $S$ including infinity, and $d$ independent rational linear forms at each place, primitive integer vectors $x$ satisfying
\[
 \prod_{v\in S}\prod_{i=1}^d |L_{i,v}(x)|_v\le\|x\|_\infty^{-\delta},\quad\delta>0,
\]
lie in finitely many proper rational subspaces. We invoke this classical theorem, not a new or effective version.
Its standard consequence is
\begin{equation}\label{R6E1:sublinear}
 v_p\left(\sum_{i=1}^d b_i^k\right)=o(k)
 \quad\text{for }1=b_1<b_2<\cdots<b_d\text{ positive integers}.
\end{equation}
Indeed use the primitive vector $(1,b_2^k,\ldots,b_d^k)$; let $S$ include $p$, infinity, and all primes dividing the bases. At all places use coordinate forms, except replace the first form at $p$ by their sum. The product of their normalized absolute values is exactly $|\sum b_i^k|_p$ by the product formula. Valuation at least $\varepsilon k$ gives the theorem's inequality with $\delta=\varepsilon\log p/\log b_d$. Any proper rational subspace has a nonzero equation $\sum c_i b_i^k=0$, which has only finitely many solutions $k$: divide by the largest base with nonzero coefficient and let $k$ tend to infinity. Thus there are finitely many exceptions for each $\varepsilon>0$, proving \eqref{R6E1:sublinear}.
Apply this to the bases $\{b^a:b\in U_p\}$. It gives
$f(m)=v_p(F_p(m))=o(m)$. Eventually $f(m)<am$, so the unequal-valuation rule in \eqref{R6E1:boundary} gives $A_p(m)=f(m)$. The exceptional set is finite but not effectively located by this argument.
\end{proof}

\begin{proof}[Proof of Theorem~\ref{R6E1:main} and Corollary~\ref{R6E1:pairs}]
In the nonconstant case apply Theorem~\ref{R6E1:model} with $\alpha=\alpha_p$. Each shifted arithmetic sequence differs from its root model at only finitely many $m$, so their product sums differ by a fixed constant. For sufficiently large $N$ their maxima are in fact identical: the root-model maxima tend to infinity, since every prescribed residue class is hit by positive integers, while the finitely many altered values are bounded. Thus replacing $b_i(N)^{d_i}$ by $B_i(N)^{d_i}$ preserves \eqref{R6E1:expansion} after enlarging an ineffective constant to cover the finite initial segment. This eventual equality, rather than merely an $O(1)$ difference of maxima, is needed when $d_i>1$.
Since $f(m+t_i)=o(m)$, every $B_i(N)=o(N)$. Setting all powers to one proves \eqref{R6E1:correlation}; positivity of all $c_i$ proves the general criterion. The constant Wilson case is Lemma~\ref{R6E1:arithmetic}.
For exponent shifts, simultaneous nonzero valuations require $a\mid n+h_i$ for every $i$, so incongruent shifts give an identically zero product. For congruent shifts write $n+h_*=am$ and apply \eqref{R6E1:correlation} on the resulting interval of $m$. Its lower endpoint is fixed and its length is $X/a+O(1)$, giving the claimed factor $1/a$. The pair integral was evaluated above. Since $W_5=5\not\equiv1\pmod5$, the stated examples follow.
\end{proof}

\section{First-mean convergence does not imply second-mean convergence}\label{R6E1:obstruction}
The following model is not asserted to be a power-sum root. It strengthens the preceding round's density-only obstruction by satisfying precisely the sublinearity needed for its first-mean theorem.
\begin{srcBEtheorem}[Square-root spikes]\label{R6E1:counterexample}
For any prime $p$, define
\[
 m_1=p^4,\qquad e_j=\lfloor\sqrt{m_j}\rfloor,\qquad
 m_{j+1}=m_j+p^{e_j},\qquad \beta=\lim_j m_j\text{ in }\srcBEZp.
\]
Then $\beta$ is not a positive integer, and $f(m)=1+v_p(m-\beta)$ satisfies $f(m)=o(m)$. Its first mean is $p/(p-1)$, but
\[
 \liminf_{N\to\infty}\frac1N\sum_{m\le N}f(m)^2=\frac{p(p+1)}{(p-1)^2},\qquad
 \limsup_{N\to\infty}\frac1N\sum_{m\le N}f(m)^2=\frac{p(p+1)}{(p-1)^2}+1.
\]
All fixed distinct-shift products of $f$ nevertheless have the limits in \eqref{R6E1:correlation}.
\end{srcBEtheorem}
\begin{proof}
For every integer $e\ge4$, $2^e\ge2e+1$ (start at $e=4$ and induct). Since $m_j\ge e_j^2$ and $e_j\ge4$, this implies $m_{j+1}\ge(e_j+1)^2$. Thus $e_j$ strictly increases, the $p$-adic limit exists, and $v_p(\beta-m_j)=e_j$.
For $m_j<m<m_{j+1}$, $0<m-m_j<p^{e_j}$, so
\[
 v_p(m-\beta)=v_p(m-m_j)<e_j.
\]
For $1\le m<m_1$, the same argument gives $v_p(m-\beta)=v_p(m-m_1)<4\le e_1$. These formulas show that $\beta$ is not a positive integer, $f(m)\le1+\sqrt m$ for $m\ge m_1$, and exactly
\[
 B(N)=\max_{m\le N} f(m)=e_j+1\qquad(m_j\le N<m_{j+1}).
\]
Hence $B(N)/N\to0$, while $B(N)^2/N$ has limsup one (at $N=m_j$, since $(e_j+1)^2/m_j\to1$) and liminf zero (at $N=m_{j+1}-1$, since $m_j/p^{\lfloor\sqrt{m_j}\rfloor}\to0$). Lemma~\ref{R6E1:single} for $k=1,2$ proves the asserted limits, with $\mu_{p,2}=p(p+1)/(p-1)^2$. Theorem~\ref{R6E1:model} and sublinearity prove every distinct-shift assertion.
\end{proof}

\section{Audit, scope, and remaining arithmetic obstruction}
The arithmetic source \cite{R3E1:AMSV}, its published 2025 version, was checked again on 20 October 2026; its notation/indexing caveat from earlier rounds is preserved. Sharma's thesis cited there was not obtained in full. The published Fermat-quotient lifting antecedent \cite[Theorem 1]{R5E1:BD} and Evertse's primary normalization were rechecked. Lifting, the Subspace Theorem, and its $o(k)$ consequence are established inputs, not discoveries. The elementary separated-singularity argument is supplied in full; no historical priority claim is made for its general principle or for these applications merely because searches did not locate the same formulas.

The new general formula is useful beyond an abstract uniform-integrability restatement: it gives explicit positive maximum corrections, exact rational limiting constants, and a pointwise Diophantine threshold for every prescribed mixed moment. Nevertheless, it does not improve the known arithmetic estimate $f(m)=o(m)$ to $o(m^{1/2})$. Thus ordinary square moments of the actual power sums remain unproved here. The example above shows that sublinearity and the full geometric level geometry alone cannot fill that gap. There is no claimed effective convergence rate, effective exceptional cutoff, empty exceptional set, or resolution of an unrelated major open problem.

The verifier checks exact periodic cancellation, rational integral evaluation and pair constants, one-root polynomial error bounds, finite power-sum/root comparisons, and the square-root-spike formulas. It uses no floating-point evidence as proof. Numerical mixed averages are labeled diagnostics. No finite verification establishes the ineffective Subspace-Theorem cutoff. Further progress requires stronger approximation bounds for the specific roots $\alpha_p$, or a separate arithmetic way to control their maximal square spikes.



\chapter[Gaussian square sums at composite side lengths: exact valuations and effective moments]{Gaussian square sums at composite side lengths: exact valuations and effective moments}
\label{source:R7E1}
\noindent\textit{Source development: Round 7, entry 01. Original mathematical arguments retained; notation is local to this chapter.}\par\medskip

\begin{quote}\small\noindent\textbf{Chapter abstract.} 
Let $G_n(N)=\sum_{a,b=1}^{N-1}(a+bi)^n$. For every odd prime $p$, we determine its $p$-adic valuation at all side lengths divisible by $p$ in the infinite exponent families $n=L+r$, where $4\mid L$, $p-1\mid L$, $L>0$, and $0\le r\le3$. The generic valuation is explicit. An extra depth occurs only when $p\equiv3\pmod4$ and $p^2-1\mid L$, and is exactly the distance to one simple algebraic $p$-adic root in the side-length variable. The four roots belonging to $L,L+1,L+2,L+3$ agree modulo $p^3$. At $p=5$ this gives a closed formula for every exponent and every side length divisible by five; at $p=3$ it gives a complete one-root classification. Polynomial height bounds make every ordinary moment effective, with summatory error $O(\log^k X)$, after omitting a possible zero value. The prime-side formulas and finite-ring ingredients have prior antecedents. The contribution is the composite-side classification and its effective statistical consequence, not a newly solved unrestricted conjecture.
\end{quote}

\section{The exact classification}
For integers $N\ge2$ and $n\ge0$ put
\[
 G_n(N)=\sum_{a=1}^{N-1}\sum_{b=1}^{N-1}(a+bi)^n\in\mathbb Z[i].
\]
For a Gaussian rational $z=x+iy$ define $v_p(z)=\min(v_p(x),v_p(y))$, with $v_p(0)=+\infty$. At split primes this is a divisibility order, not a valuation on the entire Gaussian field; we will only multiply by Gaussian units and rational scalars, where the required additivity is valid. Set $c=(1+i)/2$, a unit in $\srcBKZloc[i]$ for every odd $p$.

Fix an odd prime $p$ and write
\begin{equation}\label{R7E1:family}
 n=L+r,\qquad L\ge4,\quad 4\mid L,\quad p-1\mid L,\quad 0\le r\le3.
\end{equation}
Define
\[
 \delta=\begin{cases}1,&p\equiv3\pmod4\text{ and }p^2-1\mid L,\\0,&\text{otherwise.}\end{cases}
\]

\begin{srcBKtheorem}[A normalized polynomial and its unique exceptional branch]\label{R7E1:main}
Under \eqref{R7E1:family}, the expression
\begin{equation}\label{R7E1:Hdef}
 H_{n,p}(u)=\frac{G_n(pu)}{\binom nr p^r c^r u^{r+1}}
\end{equation}
extends from nonzero integers to a polynomial in $\srcBKZloc[u]$. Its coefficientwise reduction is
\begin{equation}\label{R7E1:reduction}
 H_{n,p}(u)\equiv2-\delta u\pmod p.
\end{equation}
If $\delta=0$, it is a unit at every $u\in\srcBKZp$. If $\delta=1$, it has a unique zero $\beta_{n,p}\in\srcBKZp$, with $\beta_{n,p}\equiv2\pmod p$, and
\begin{equation}\label{R7E1:isometry}
 v_p(H_{n,p}(u))=v_p(u-\beta_{n,p})\qquad(u\in\srcBKZp).
\end{equation}
Consequently, for every positive integer $N$ with $s=v_p(N)\ge1$,
\begin{equation}\label{R7E1:valuation}
 v_p(G_n(N))=s(r+1)-1+v_p\binom nr+
 \begin{cases}0,&\delta=0,\\v_p(N/p-\beta_{n,p}),&\delta=1.\end{cases}
\end{equation}
The equality includes $+\infty$ if a sum vanishes. In the exceptional case all extra depth is confined to $N/p\equiv2\pmod p$, in particular to $s=1$. The root can be computed uniquely to any prescribed precision by successive lifting.
\end{srcBKtheorem}

\begin{srcBKcorollary}[Four consecutive exponents share three root digits]\label{R7E1:coherence}
Fix $L$ satisfying \eqref{R7E1:family}. The four normalized polynomials satisfy
\[
 H_{L+r,p}(u)\equiv H_{L,p}(u)\pmod {p^3}\qquad(0\le r\le3).
\]
If $\delta=1$, their respective roots are congruent modulo $p^3$.
\end{srcBKcorollary}

\begin{srcBKcorollary}[Complete small-prime families]\label{R7E1:small}
For every $N$ divisible by $5$, $s=v_5(N)$, and every $n\ge4$, writing $n=4t+r$ with $0\le r\le3$,
\[
 v_5(G_n(N))=s(r+1)-1+v_5\binom nr.
\]
For $p=3$, the same generic formula holds for every $n\ge4$, except that when $8\mid n-r$ it has the extra term $v_3(N/3-\beta_{n,3})$ of Theorem~\ref{R7E1:main}. Thus it always holds without extra depth if $9\mid N$ or $N\equiv3\pmod9$.
For $0\le n\le3$ and every odd $p\mid N$, one has $v_p(G_n(N))=n\,v_p(N)$.
\end{srcBKcorollary}
For example, $G_8(6)=8\,534\,460$ has $3$-adic order one, although the generic value for $n=8$ and $v_3(N)=1$ is zero. This is a genuine composite-side exception, not an extension of the prime-side formula without qualification.

\paragraph{Target, prior work, and exact increment.}
The latest Kalinin--Zottor version (v3, 28 July 2026) \cite{R7E1:KZ}, studies the prime-side specialization $G_n(p)$ and includes complete formulas at $p=3,5$, finite-field congruences, and a depth-five congruence at the special exponent $n=p$. The previous versions' advertised conjectures must not be treated as still open: v3 already proves the relevant assertions. The present theorem changes the variable to \emph{all side lengths $N$ divisible by $p$}, with an exact algebraic branch and four-exponent coherence. It does not claim a new proof of their complete prime-side $p=3,5$ formulas or cover all exponents at every larger prime. Its families have period $\operatorname{lcm}(4,p-1)$ and contain four residue classes per period.

Their Section 5 already uses fourfold centering, and their Lemma 11 isolates the unique minimum term for the prime-side small-prime formulas. Those methodological antecedents are credited, not newly claimed. The Faulhaber expansion, von Staudt--Clausen denominator theorem, finite-field power sums, and simple-root lifting are classical tools. Finite-ring Gaussian power sums have earlier general treatments \cite{R7E1:AGO,R7E1:GO}. Historical priority of the precise composite-side combination is not certified. All six frozen Round 6 papers and the assessment were read. Earlier Entry 01 root-summation arguments \cite{R7E1:R6} motivate the effective moment application below, but their exponential roots and ineffective Subspace-Theorem step are not used here.

\section{Polynomial integrality and finite-ring reduction}
Let $S_j(Y-1)$ denote the polynomial agreeing with $\sum_{a=1}^{Y-1}a^j$ at positive integers $Y$. In particular $S_0(Y-1)=Y-1$.

\begin{srcBKlemma}\label{R7E1:faulhaber}
For $j\ge1$, $S_j(pu-1)\in\srcBKZloc[u]$, its constant term is zero, and
\begin{equation}\label{R7E1:Smod}
 S_j(pu-1)\equiv-u\,\mathbf1_{p-1\mid j}\pmod p.
\end{equation}
Also $S_0(pu-1)\equiv-1\pmod p$.
\end{srcBKlemma}
\begin{proof}
Use Bernoulli numbers with $B_1=-1/2$. Faulhaber's polynomial identity is
\[
 S_j(Y-1)=\sum_{k=1}^{j+1}\frac1k\binom j{k-1}B_{j-k+1}Y^k.
\]
The classical von Staudt--Clausen theorem implies $v_p(B_m)\ge-1$ for all $m\ge0$, and $pB_j\equiv-1\pmod p$ precisely when $j>0$ and $p-1\mid j$; otherwise $pB_j\equiv0$. These are the only Bernoulli inputs. After $Y=pu$, the coefficient with $k\ge2$ has valuation at least $k-v_p(k)-1\ge1$ for odd $p$. Indeed $v_p(k)\le k-2$ for $k\ge2$. The $k=1$ coefficient is $pB_j$. This proves integrality, the zero constant, and \eqref{R7E1:Smod}.
\end{proof}

Define the centered polynomial moments
\[
 M_j(Y)=\sum_{a,b=1}^{Y-1}\bigl(a-Y/2+i(b-Y/2)\bigr)^j
\]
first at positive integers and then by polynomial continuation. The square of centered points is invariant under multiplication by $i$ and under conjugation. Thus, as polynomial identities,
\begin{equation}\label{R7E1:symmetry}
 M_j=0\quad(4\nmid j),\qquad M_{4t}\in\mathbb Q[Y],\qquad M_0(Y)=(Y-1)^2.
\end{equation}
For example, multiplication by $i$ sends $(a,b)$ to $(Y-b,a)$, still in the square. Polynomial equality follows from equality at infinitely many integers. The exact translation identity is
\begin{equation}\label{R7E1:translation}
 G_n(Y)=\sum_{j=0}^n\binom nj(cY)^{n-j}M_j(Y).
\end{equation}

\begin{srcBKlemma}[The leading centered polynomial]\label{R7E1:center}
For every $j\ge0$, $G_j(pu),M_j(pu)\in\srcBKZloc[i][u]$. For $j>0$ each has zero constant term. For $L$ satisfying \eqref{R7E1:family},
\begin{equation}\label{R7E1:Mmod}
 M_L(pu)\equiv G_L(pu)\equiv2u-\delta u^2\pmod p.
\end{equation}
In particular $M_L(pu)/u\in\srcBKZloc[u]$.
\end{srcBKlemma}
\begin{proof}
The exact binomial factorization is
\[
 G_j(Y)=\sum_{h=0}^j\binom jh i^h S_{j-h}(Y-1)S_h(Y-1).
\]
Lemma~\ref{R7E1:faulhaber} gives coefficient integrality. For $j>0$ every summand has at least one positive-index power sum, proving the zero constant. Translating in the opposite direction expresses $M_j(pu)$ as a sum of $G_h(pu)$ times powers of $-cpu$, so it too is integral and has zero constant for $j>0$. All translated terms except $G_j(pu)$ are divisible by $p$.

For $L>0$ divisible by both $4$ and $p-1$, the two endpoint terms reduce to $2u$. Every interior term reduces to $u^2\binom Lh i^h$ when both $h$ and $L-h$ are divisible by $p-1$, and to zero otherwise. Their coefficient sum is exactly
\[
 \sum_{a,b\in\mathbb F_p}(a+bi)^L\quad\text{in }\mathbb F_p[i].
\]
Indeed the endpoint factors involve $\sum_{a\in\mathbb F_p}1=p=0$, while a positive power sum is $-1$ or zero according to divisibility by $p-1$. If $p\equiv3\pmod4$, this ring is the field $\mathbb F_{p^2}$, so its sum is $-1$ exactly when $p^2-1\mid L$, and zero otherwise, by cyclicity of its multiplicative group. If $p\equiv1\pmod4$, the ring is $\mathbb F_p\times\mathbb F_p$. In either coordinate each field element occurs $p$ times as the pair varies, so the sum is zero. This proves \eqref{R7E1:Mmod}, including the split case. The reality assertion follows from \eqref{R7E1:symmetry}.
\end{proof}

\section{Normalization, coherence, and exact lifting}
\begin{srcBKlemma}[A small binomial-denominator bound]\label{R7E1:binomial}
For $k\ge4$ and $0\le r\le3$,
\[
 v_p\binom{k+r}{r}\le\lfloor\log_p(k+r)\rfloor\le k-3.
\]
Consequently, whenever $k\le n-r$,
\begin{equation}\label{R7E1:ratio}
 p^k\frac{\binom n{r+k}}{\binom nr}\in p^3\srcBKZloc.
\end{equation}
\end{srcBKlemma}
\begin{proof}
For $r=0$ the first valuation is zero. For $r=1,2$, at most one of the $r$ consecutive numerator factors is divisible by odd $p$, and $r!$ is a unit. For $r=3,p\ge5$ the same argument applies. For $r=3,p=3$, exactly one of the three numerator factors is divisible by three, and dividing by $3!$ only reduces its valuation. This proves the first bound. The elementary inequality $3^{k-2}>k+3$ for $k\ge4$ follows at $k=4$ and by induction, and gives the second. Finally,
\[
 \frac{\binom n{r+k}}{\binom nr}=\frac{\binom{n-r}k}{\binom{r+k}r},
\]
whose numerator is integral. This proves \eqref{R7E1:ratio}.
\end{proof}

\begin{proof}[Polynomial assertions and Corollary~\ref{R7E1:coherence}]
Only $j=L-4\ell$ survive in \eqref{R7E1:translation}. Dividing it by the expression in \eqref{R7E1:Hdef} gives the leading term $M_L(pu)/u$. Each remaining term, with $k=4\ell\ge4$, is
\begin{equation}\label{R7E1:higher}
 p^k c^k\frac{\binom n{r+k}}{\binom nr}\,
                 u^{k-1}M_{L-k}(pu).
\end{equation}
Here $c^k=(-1/4)^\ell$ is a rational $p$-adic unit; $u^{k-1}$ has a nonnegative exponent even for the final term $L-k=0$. Lemmas~\ref{R7E1:center} and \ref{R7E1:binomial} show that every coefficient in \eqref{R7E1:higher} lies in $p^3\srcBKZloc$. Thus $H_{n,p}$ is a real rational polynomial with $p$-integral coefficients, and
\begin{equation}\label{R7E1:common}
 H_{L+r,p}(u)\equiv M_L(pu)/u\pmod {p^3}.
\end{equation}
Lemma~\ref{R7E1:center} proves \eqref{R7E1:reduction}, and \eqref{R7E1:common} proves polynomial coherence.
\end{proof}

\begin{proof}[Lifting and the valuation formula]
If $\delta=0$, \eqref{R7E1:reduction} is the constant unit two. If $\delta=1$, coefficientwise reduction is $2-u$. For any $u,v\in\srcBKZp$, polynomial division gives
\[
 H(u)-H(v)=(u-v)(-1+pQ(u,v)),\qquad Q\in\srcBKZp[u,v].
\]
The second factor is always a unit. Hence $H$ is an isometry. Its unique zero modulo $p$ is two; if $a_j$ is a zero modulo $p^j$, the Taylor congruence
\[
 H(a_j+p^jt)\equiv H(a_j)+p^jtH'(a_j)\pmod {p^{j+1}}
\]
gives exactly one next digit $t\in\{0,\ldots,p-1\}$ because $H'(a_j)\equiv-1$. Completeness gives the root $\beta$ and the isometry gives \eqref{R7E1:isometry}. Polynomial coherence and isometry imply equality of the four root residues modulo $p^3$.

For $u=N/p$, $v_p(u)=s-1$. Multiplication by $c^r$ preserves Gaussian divisibility because its inverse belongs to $\srcBKZloc[i]$. Taking orders in \eqref{R7E1:Hdef} gives
\[
 v_p(G_n(N))=v_p\binom nr+r+(r+1)(s-1)+v_p(H(u)),
\]
which is \eqref{R7E1:valuation}. Since $\beta\equiv2$, its extra term can be positive only for unit $u\equiv2$. This also proves the small-prime assertions for $n\ge4$: every multiple of four is divisible by $p-1$ at $p=3,5$, and $\delta=0$ at five. For $n\le3$, \eqref{R7E1:translation} has only $M_0$, giving $G_n(N)=(cN)^n(N-1)^2$. Since $p\mid N$, the last factor is a unit. This completes all the stated classifications.
\end{proof}

\section{Effective laws for every ordinary moment}
The new side-length roots are algebraic, in contrast with the exponential-sequence roots of the earlier rounds. This permits an elementary effective height estimate, without the Subspace Theorem.

Fix $n,p$ as in \eqref{R7E1:family}, and set
\[
 b=r+v_p\binom nr,\qquad d=r+1,\qquad V(u)=v_p(G_n(pu)).
\]
If $G_n(pu)=0$, omit that index in sums of valuations. There is at most one such positive integer $u$: in the constant case $H$ is never zero, and in the exceptional case its unique $p$-adic zero is $\beta$. Whether such an integer zero exists is decidable by the rational-root test after clearing denominators. The constant coefficient of $H$ is a unit, so it is nonzero and the test is finite.

\begin{srcBKtheorem}[Effective all-moment convergence]\label{R7E1:moments}
Define the rational function
\begin{equation}\label{R7E1:pgf}
 \Phi(z)=z^b\left\{\frac{p-1-\delta}{p}
 +\frac{p-1}{p}\left(\frac{z^d}{p-z^d}
                      +\delta\frac{z}{p-z}\right)\right\}.
\end{equation}
For every fixed integer $k\ge1$, put
$\mu_k=[(z\,d/dz)^k\Phi(z)]_{z=1}\in\mathbb Q$. Then
\begin{equation}\label{R7E1:momentsum}
 \sum_{\substack{1\le u\le X\\G_n(pu)\ne0}} V(u)^k
 =\mu_k X+O_{n,p,k}(\log^k(2+X)),
\end{equation}
with an effectively computable implied constant. In particular
\begin{equation}\label{R7E1:mean}
 \mu_1=r+v_p\binom nr+\frac{r+1+\delta}{p-1}.
\end{equation}
The coefficients of $\Phi$ are the natural-density distribution of these valuations, with any omitted zero having density zero.
\end{srcBKtheorem}
\begin{proof}
The exact formula is
\begin{equation}\label{R7E1:tworoot}
 V(u)=b+d\,v_p(u)+\delta\,v_p(u-\beta)
\end{equation}
whenever the sum is nonzero. In the exceptional case the two roots are in distinct residue classes modulo $p$, since $\beta\equiv2$. Their positive-depth events are disjoint. Each root has exact-depth density $(p-1)/p^{h+1}$ at depth $h\ge1$. Off these one or two root classes the density is $(p-1-\delta)/p$. Summing their geometric generating series gives \eqref{R7E1:pgf}, whose coefficients are nonnegative and sum to one. It converges near $z=1$, and differentiation gives the candidate moments and \eqref{R7E1:mean}.

To justify ordinary moments effectively, choose an integer $D_0>0$, prime to $p$, clearing all denominators of $H$, and write $Q=D_0H\in\mathbb Z[u]$, of degree $D$. Let $C\ge1$ be the sum of the absolute values of its coefficients. At any nonzero integer value and $1\le u\le X$,
\[
 p^{v_p(H(u))}=p^{v_p(Q(u))}\le |Q(u)|\le C(1+X)^D.
\]
Thus every finite extra depth is at most $\log_p C+D\log_p(1+X)$. Also $v_p(u)\le\log_p X$. Both bounds are explicit in the computed polynomial. Choose an integer $J=O_{n,p}(\log(2+X))$ exceeding these bounds and $2\log_p(2+X)$.

For a root $\gamma\in\srcBKZp$, the count of $u\in[1,X]$ with $v_p(u-\gamma)\ge h$ differs from $X/p^h$ by at most one. Removing a possible integer zero changes it by at most one more. By disjointness in \eqref{R7E1:tworoot}, its $k$th power has the layer decomposition
\[
 b^k+\sum_{h\ge1}\bigl((b+dh)^k-(b+d(h-1))^k\bigr)
                     \mathbf1_{v_p(u)\ge h}
 +\delta\sum_{h\ge1}\bigl((b+h)^k-(b+h-1)^k\bigr)
                     \mathbf1_{v_p(u-\beta)\ge h}.
\]
On the retained indices the sums stop at $J$. Replacing each residue count by $X/p^h$ costs at most a constant times $(b+dJ)^k=O_{n,p,k}(\log^k(2+X))$, by telescoping the positive layer increments. The omitted infinite geometric tails cost $O_{n,p,k}(XJ^{k-1}p^{-J})$, which is smaller. Removing the possible zero also costs $O(b^k)$ in the constant layer. This proves \eqref{R7E1:momentsum}. All constants follow from finite polynomial arithmetic and explicit geometric tails. The layer-counting mechanism is the earlier Entry 01 method \cite{R7E1:R6}; the new arithmetic input is the effective polynomial bound, not a new general summation principle.
\end{proof}
For instance, the mean at $(p,n)=(3,8)$ is one, and at $(p,n)=(3,9)$ it is $9/2$. At $(p,n)=(5,8)$ it is $1/4$. These refer to $N=pu$ with $u$ uniformly averaged; they are not the prime-side sequence studied in the source or the exponent averages studied in earlier rounds.

\section{Scope, source audit, and verification}
The formula applies exactly to \eqref{R7E1:family}; exponents outside these four residue families at larger primes may have other Bernoulli or root behavior and are not classified. It also requires $p\mid N$. The simple exceptional root is algebraic and effectively liftable, but no claim that it is rational, irrational, or never a positive integer is needed. A possible zero is explicitly omitted from moments; assigning an infinite valuation to it and then averaging would be invalid.

The latest primary Gaussian source v3 was inspected rather than relying on a search summary of v1. Its prime-side $p=3,5$ formulas and its resolved older conjectures retain their source credit. The general finite-ring context is credited separately. The composite-side theorem uses the same classical finite-field mechanism, but its polynomial normalization, exact side-length branch and coherence are not asserted in that inspected source. This comparison is not an exhaustive novelty certification. The earlier source-status corrections concerning strong Seymour and gamma problems are untouched; neither topic is used here.

A bounded attempt at weighted three-normal polygon barycentric concavity found neither a proof nor a counterexample and was abandoned. It is not a result. The present work instead obtains a complete effective classification in a different Gaussian family. The verifier constructs the polynomials by independent Faulhaber arithmetic, checks all coefficient integrality and reductions in stated ranges, compares against direct Gaussian sums, lifts roots, checks four-exponent coherence, and tests exact moment layer counts. Finite averages are only diagnostics; the universal conclusions follow from the proofs above. No external posting, submission, or author contact occurred.



\part{Strong second neighborhoods in tournaments}

\chapter[A Deficiency Budget for Strong Seymour Vertices under Tournament Perturbations]{A Deficiency Budget for Strong Seymour Vertices under Tournament Perturbations}
\label{source:R1E2}
\noindent\textit{Source development: Round 1, entry 02. Original mathematical arguments retained; notation is local to this chapter.}\par\medskip


\begin{quote}\small\noindent\textbf{Chapter abstract.} 
Bai, Li and Park conjectured in July 2026 that every oriented graph has a vertex whose out-neighbors can be matched by directed arcs into its exact second out-neighborhood. We prove a perturbative special case, not a universal result. The originally cited universal strong conjectures were disproved in version 2 of the source; the original submission missed that update. If $D$ is obtained from an $n$-vertex doubly regular tournament by deleting $d$ arcs and reversing $r$ other arcs, then $D$ has at least $n-d-2r$ such vertices. Moreover, the sum, over all vertices, of the unmatched out-neighbor deficiencies is at most $d+3r$. For deletions alone this becomes a total deficiency budget of $d$. The proof combines a regular bipartite matching-deficiency inequality with exact directed-triangle incidence counting. The constants in the uniform deficiency budget are sharp already on three vertices; no large-order optimality is asserted. These are elementary special-case results whose publication novelty has not been established. Exact finite checks support, but are not used in, the proofs.
\end{quote}

\section{The target and the scope}
All graphs here are finite. An oriented graph has no loops, multiple arcs, or pairs of opposite arcs. For a vertex $x$ in an oriented graph $D$, put
\[
 \srcABsecond_D(x)=\left(\bigcup_{u\in\srcABout_D(x)}\srcABout_D(u)\right)
 \setminus(\srcABout_D(x)\cup\{x\}).
\]
The exclusion of first out-neighbors is essential. Let $\nu_D(x)$ be the maximum cardinality of a matching consisting of arcs from $\srcABout_D(x)$ to $\srcABsecond_D(x)$, and define
\[
 \eta_D(x)=|\srcABout_D(x)|-\nu_D(x).
\]
Thus $\eta_D(x)$ is a nonnegative integer. A \emph{strong Seymour vertex} is a vertex with $\eta_D(x)=0$; a sink qualifies, with the empty matching.

Bai, Li and Park~\cite{R1E2:BLP}, Conjecture 1.2, ask whether every oriented graph has a strong Seymour vertex. They separately state the tournament case as Conjecture 1.3 in version 1. Their July 2026 preprint proves existence when the minimum out-degree is at most five, and for $5$-anti-transitive oriented graphs. \textbf{Status correction (20 October 2026):} version 2 of the same paper, posted on 24 July 2026, reports a $36$-vertex counterexample to both strong conjectures in Remark 3.1, attributed to David Dzitsoev. The original submission of this note inspected version 1 only and missed that update. Its current-open-problem framing is withdrawn. The perturbation theorem proved below is unaffected: it identifies a positive subclass for a property whose universal version is false. The ordinary (unstrengthened) Seymour conjecture is a different question.

A tournament $T$ of order $n=4q+3$ is \emph{doubly regular} if every vertex has out-degree $2q+1$ and every two distinct vertices have exactly $q$ common out-neighbors. Set
\[
 a=q+1=\frac{n+1}{4},\qquad b=2q+1=\frac{n-1}{2}=2a-1.
\]
Doubly regular tournaments are a classical family; Paley tournaments of prime order $p\equiv3\pmod4$ provide examples~\cite{R1E2:GN}. We do not claim their unperturbed matching property as a new result. The contribution examined here is a quantitative transfer estimate under arbitrary specified arc deletions and reversals.

\begin{srcABtheorem}[Perturbation and deficiency budget]\label{R1E2:thm:main}
Let $T$ be a doubly regular tournament on $n$ vertices. Let $F_0,F_1$ be disjoint subsets of its arc set, with $|F_0|=d$ and $|F_1|=r$. Obtain $D$ by deleting the arcs in $F_0$ and reversing those in $F_1$. Let $H$ be the set of heads, in $T$, of arcs in $F_1$. Then
\begin{align}
 \sum_{x\in V(T)}\eta_D(x)&\le d+3r,\label{R1E2:eq:mass}\\
 |\{x:\eta_D(x)=0\}|&\ge n-|H|-d-r\ge n-d-2r.\label{R1E2:eq:count}
\end{align}
In particular, $D$ has a strong Seymour vertex whenever $d+2r<n$. If $d+r<a$, every vertex outside $H$ is strong.
\end{srcABtheorem}
An inequality with a negative right-hand side in~\eqref{R1E2:eq:count} is simply vacuous. All edits are specified relative to $T$; repeated edits to one arc are not counted as separate operations.

\section{A matching-deficiency inequality}
We give the matching argument explicitly to isolate the only combinatorial ingredient.

\begin{srcABlemma}\label{R1E2:lem:matching}
Let $B$ be an $a$-regular bipartite graph with parts $L,R$ of equal size $b$, where $a\ge1$. Remove $c$ of its edges, obtaining $B'$. A maximum matching in $B'$ leaves at most $\lfloor c/a\rfloor$ vertices of $L$ unmatched.
\end{srcABlemma}
\begin{proof}
Take a maximum matching $M$, and suppose it leaves $z>0$ vertices of $L$ unmatched. Start alternating paths from all unmatched vertices of $L$, traversing unmatched edges from left to right and matched edges from right to left. Let $S\subseteq L$ and $W\subseteq R$ be the reached sets. No vertex of $W$ is unmatched, since an alternating path to such a vertex would augment $M$. Also $N_{B'}(S)=W$, and the matching pairs $W$ bijectively with the matched vertices of $S$. Hence $|S|-|W|=z$. In $B$ there are $a|S|$ edges incident with $S$, of which at most $a|W|$ end in $W$. All the remaining edges were removed. Consequently $c\ge a(|S|-|W|)=az$. The case $z=0$ is immediate.
\end{proof}

\begin{srcABlemma}[Cyclic links]\label{R1E2:lem:link}
For $x\in V(T)$, form the bipartite graph $B_x$ with left part $L_x=\srcABout_T(x)$ and right part $R_x=\srcABinn_T(x)$, with an edge $uv$ when $u\to v$ in $T$. Then $B_x$ is $a$-regular on two parts of size $b$. Every arc of $T$ appears as an edge of exactly $a$ of these bipartite graphs.
\end{srcABlemma}
\begin{proof}
For $u\in L_x$, exactly $q$ of its out-neighbors are common out-neighbors with $x$, and therefore lie in $L_x$. Since $u$ does not dominate $x$, it has $2q+1-q=q+1=a$ out-neighbors in $R_x$. For $v\in R_x$, exactly $q$ of the vertices in $L_x$ are out-neighbors of $v$. Thus precisely $b-q=a$ vertices in $L_x$ dominate $v$.

An arc $u\to v$ belongs to $B_x$ exactly when $v\to x\to u\to v$ is a directed triangle. Among the $b$ out-neighbors of $v$, exactly $q$ are also out-neighbors of $u$, and all the other $b-q=a$ dominate $u$. These are exactly the possible $x$.
\end{proof}

\section{Proof of the perturbation theorem}
Write $F=F_0\cup F_1$, and define the integer loads
\[
 c_x=|F\cap E(B_x)|,\qquad
 h_x=|\{u:(u,x)\in F_1\}|.
\]
Here $E(B_x)$ is identified with its corresponding directed arcs in $T$. By Lemma~\ref{R1E2:lem:link},
\begin{equation}\label{R1E2:eq:incidence}
 \sum_x c_x=a(d+r),\qquad \sum_x h_x=r,\qquad H=\{x:h_x>0\}.
\end{equation}

\begin{srcABproposition}[Local certificate]\label{R1E2:prop:local}
For every vertex $x$,
\begin{equation}\label{R1E2:eq:local}
 \eta_D(x)\le \left\lfloor\frac{c_x}{a}\right\rfloor+2h_x.
\end{equation}
In particular, $h_x=0$ and $c_x<a$ certify that $x$ is strong.
\end{srcABproposition}
\begin{proof}
Delete from $B_x$ every edge corresponding to an arc in $F$. By Lemma~\ref{R1E2:lem:matching}, the remaining graph has a matching $M$ leaving at most $\lfloor c_x/a\rfloor$ vertices of $L_x$ unmatched.

The current out-neighborhood of $x$ consists of a subset $L'_x\subseteq L_x$, together with a set $A_x\subseteq R_x$ of size $h_x$: these are precisely the original in-neighbors whose arcs into $x$ were reversed. First restrict $M$ to tails in $L'_x$. This does not increase the number of unmatched vertices among those retained tails. Next remove matches whose head is in $A_x$; at most $h_x$ matches are removed. Every remaining head lies in $R_x\setminus A_x$, which is disjoint from $\srcABout_D(x)\cup\{x\}$, and the corresponding arc remains in $D$. Since its tail belongs to $\srcABout_D(x)$, that head is in $\srcABsecond_D(x)$. Thus the retained arcs are a valid second-neighborhood matching in $D$. Finally, leave the $h_x$ new out-neighbors in $A_x$ unmatched. This constructs a matching with deficiency at most the right-hand side of~\eqref{R1E2:eq:local}.
\end{proof}

Summing~\eqref{R1E2:eq:local} and using~\eqref{R1E2:eq:incidence} proves~\eqref{R1E2:eq:mass}:
\[
 \sum_x\eta_D(x)\le\sum_x\left\lfloor\frac{c_x}{a}\right\rfloor+2r
 \le d+r+2r.
\]
For the counting bound, a nonstrong vertex outside $H$ has $c_x\ge a$. Consequently there are at most $d+r$ such vertices. At most $|H|\le r$ other vertices exist in $H$, proving~\eqref{R1E2:eq:count}. If $d+r<a$, then $c_x\le d+r<a$ for every $x$, proving the final assertion of Theorem~\ref{R1E2:thm:main}. A sometimes stronger, directly computable version is
\begin{equation}\label{R1E2:eq:refined}
 |\{x:\eta_D(x)=0\}|\ge
 n-|H|-\left\lfloor\frac{\sum_{x\notin H}c_x}{a}\right\rfloor.
\end{equation}

\begin{srcABcorollary}[Deletion-only neighborhood]
Deleting $d$ arbitrary arcs from $T$ produces an oriented graph with total deficiency at most $d$ and at least $n-d$ strong Seymour vertices. In particular, fewer than $n$ deletions preserve existence. Fewer than $(n+1)/4$ deletions preserve the property at every vertex.
\end{srcABcorollary}

\section{Sharp local thresholds and limitations}
The local loss threshold $a$ cannot be increased for a prescribed vertex. Fix $x$ and $u\in L_x$, and delete all $a$ arcs from $u$ into $R_x$. No arc at $x$ was changed, so its first out-neighborhood remains $L_x$. The vertex $u$ has no remaining arc to any vertex outside $L_x\cup\{x\}$. It therefore cannot be matched into $\srcABsecond_D(x)$, and $x$ is no longer strong. This proves sharpness of the local certificate, not optimality of the global existence threshold.

The coefficients of $d$ and $r$ in the uniform deficiency budget~\eqref{R1E2:eq:mass} are individually sharp. On the directed triangle, deleting one arc produces a directed path whose vertex deficiencies are $0,1,0$ in path order. Reversing one arc produces a transitive tournament, whose deficiencies in decreasing out-degree order are $2,1,0$. These have total deficiencies $1$ and $3$, respectively. They also attain the respective lower bounds $n-d$ and $n-2r$ in~\eqref{R1E2:eq:count}. No asymptotic sharpness claim follows from these examples.

An initial approach excluded \emph{both} endpoints of each reversed arc. That gives only the weaker count $n-d-3r$. The useful refinement is asymmetric: an outgoing reversal removes a demanded tail, whereas an incoming reversal creates a new demanded tail and can remove a permitted target. Proposition~\ref{R1E2:prop:local} separates these two effects.

The main obstruction to extending the proof to arbitrary tournaments is precise. The argument requires many matchings, or the equivalent regular-link deficiency bound, for the \emph{original} tournament. For a general tournament a cyclic link need not be balanced or regular, and may already lack a complete matching. A transfer theorem cannot produce this missing base certificate. Replacing double regularity merely by regularity is not justified. Nor is monotonicity under arbitrary arc deletion asserted: the fixed-vertex example above directly disproves such monotonicity. The results here concern the stated perturbation class only. The universal assertions for tournaments and oriented graphs are false by the cited version-2 counterexample; no general rescue of them is claimed.

\section{Exact checks, reproducibility, and novelty status}
The standalone standard-library Python script accompanying this note constructs prime-order Paley tournaments using modular squares, checks double regularity with exact set counts, and computes maximum matchings by augmenting paths. It verifies the local and global bounds, not just existence. The deterministic run (seed $20261020$) gives:
\begin{center}
\begin{tabular}{lrr}
\toprule
Test & Orders & Number of cases\\
\midrule
All deletion/reversal patterns & $3$ & $27$\\
All patterns on at most three arcs & $7$ & $11{,}523$\\
Seeded arbitrary mixed edits & $11,19,31,43$ & $2{,}000$\\
\bottomrule
\end{tabular}
\end{center}
The script also records an explicit $31$-vertex example with $20$ deletions and $5$ reversals. It has minimum out-degree $13$ and $27$ strong vertices. The theorem guarantees at least one. This example is not covered by the source's minimum-out-degree-at-most-five sufficient condition; this observation does not claim exclusion from every previously known special class. The exact edit list and all test counts are emitted as JSON. There is no floating-point spectral computation. These finite tests are debugging evidence, not a replacement for the proof.

The unperturbed regular-link fact and the matching argument are standard. The particular mixed-edit bound and total-deficiency formulation were derived in this experiment. Searches for strong Seymour vertices together with double regularity, deletion, and perturbation did not locate the same statement. That is limited novelty evidence, not certification of a new publishable theorem. The mathematical scope is an elementary quantitative theorem for a positive subclass of a property whose universal conjecture has been disproved, not a general solution or a claim of exceptional journal-level significance. Independent expert review and a broader literature comparison remain necessary.



\chapter[Locally Resilient Strong Second Neighborhoods in Doubly Regular Tournaments]{Locally Resilient Strong Second Neighborhoods in Doubly Regular Tournaments}
\label{source:R2E2}
\noindent\textit{Source development: Round 2, entry 02. Original mathematical arguments retained; notation is local to this chapter.}\par\medskip

\begin{quote}\small\noindent\textbf{Chapter abstract.} 
We prove local edit resilience for strong second-neighborhood matchings in doubly regular tournaments. If the maximum incident edit degree is at most $n/8-O(\sqrt n)$, every vertex whose edited outdegree is at most $(n-1)/2$ admits such a matching. This permits a quadratic total number of deletions and reversals, unlike the global edit-budget result of our Round 1 entry. For deletions alone, all vertices retain the property below $(n+1)/8-\sqrt n/4$. The leading constant $1/8$ is optimal for preserving the property at every vertex: a balanced Hall obstruction can be created with incident deletion degree $n/8+O(\sqrt{n\log n})$. These are restricted structural theorems, not a solution of Seymour's conjecture. In particular, the stronger universal conjecture discussed in Round 1 was disproved in version 2 of its cited source; we explicitly correct that source-status error.
\end{quote}
\section{Scope, source status, and main results}
An oriented graph has no loops, opposite arc pairs, or multiple arcs. Write $N^+(x)$ for its outneighbors, and $N^{++}(x)$ for vertices reachable from $x$ by two arcs but not in $N^+(x)\cup\{x\}$. A vertex is \emph{strong} if there is a matching of directed arcs from $N^+(x)$ into $N^{++}(x)$ saturating $N^+(x)$. This implies the ordinary Seymour inequality $|N^{++}(x)|\ge |N^+(x)|$.

Bai--Li--Park's July 20 version proposed universal strong-vertex conjectures. Their July 24 version, Remark 3.1, reports Dzitsoev's 36-vertex tournament with no strong vertex \cite{R2E2:BLP}. Thus those strong conjectures are false. The ordinary second-neighborhood conjecture remains open in that source. Our Round 1 paper incorrectly treated the first-version strong conjecture as currently open. Its restricted matching theorem is unaffected, but that framing is withdrawn. The present paper proves a resilient positive subclass, not a universal conjecture.

A doubly regular tournament (DRT) has $n=4a-1$ vertices, every outdegree $b=2a-1$, and every pair has $a-1$ common outneighbors. Paley tournaments give an infinite family \cite{R1E2:GN}. A set of edits changes distinct original arcs, each by deletion or reversal. Its \emph{incident degree} $\srcAHD$ is the maximum number of changed original arcs touching a vertex. Unchanged arcs retain their orientation.

\begin{srcAHtheorem}[Mixed-edit resilience]\label{R2E2:thm:mixed}
Let $T$ be a DRT of order $n\ge127$, and obtain $D$ by edits of incident degree $\srcAHD$. Put
\[
 \tau_n=\frac{n-1}{8}-\frac{n}{\sqrt{2(n-1)}}.
\]
If $\srcAHD\le\tau_n$, every vertex of $D$ with outdegree at most $b=(n-1)/2$ is strong. In particular $D$ has at least $\lceil n/(\srcAHD+1)\rceil$ strong vertices. If all edits are reversals, the strong vertices are exactly the vertices of outdegree at most $b$.
\end{srcAHtheorem}
\begin{srcAHtheorem}[Deletion resilience and sharp leading constant]\label{R2E2:thm:delete}
For a DRT $T$ of order $n=4a-1$, every vertex remains strong after deletions of incident degree
\[
 \srcAHD\le L_n:=\frac{n+1}{8}-\frac{\sqrt n}{4}.
\]
Let $R_{\rm del}(T)$ be the minimum incident degree of a deletion set after which at least one vertex is not strong. Then
\[
 L_n<R_{\rm del}(T)\le
 \left\lceil\frac{a^2}{2a-1}+\sqrt{\frac a2\log(4(2a-1))}\right\rceil.
\]
Consequently $R_{\rm del}(T)/n\to1/8$ along any sequence of DRTs of unbounded order.
\end{srcAHtheorem}
The upper bound destroys the property at one vertex, not necessarily at every vertex. It is not a counterexample to the ordinary Seymour conjecture. Local resilience and its characteristic half-degree phenomenon are classical themes \cite{R2E2:SV}; we claim no novelty for that general methodology. The specific strong-neighborhood conclusions and explicit edit bounds above are the research contribution proposed here; priority is not certified by our limited literature search.

\section{Two spectral facts}
Let $A$ be the adjacency matrix of $T$ and $K=A-A^{\mathsf T}$. The DRT parameters imply
\begin{equation}\label{R2E2:eq:skew}
 K\mathbf1=0,\qquad KK^{\mathsf T}=nI-J.
\end{equation}
Indeed $A+A^{\mathsf T}=J-I$ and $AA^{\mathsf T}=aI+(a-1)J$; expanding $K=2A-J+I$ gives the claim.
\begin{srcAHlemma}[Skew mixing]\label{R2E2:lem:mix}
For disjoint vertex sets $S,U$ of sizes $s,t$, respectively,
\[
 e_T(S,U)\ge\frac{st-\sqrt{st(n-s-t)}}2.
\]
\end{srcAHlemma}
\begin{proof}
Set $u=\mathbf1_S-(s/n)\mathbf1$ and $v=\mathbf1_U-(t/n)\mathbf1$. Since $K$ is skew, $u^{\mathsf T}Ku=0$. Projecting $v$ off $u$ and using \eqref{R2E2:eq:skew} bounds
\[
 |u^{\mathsf T}Kv|\le\sqrt n
 \sqrt{\|u\|^2\|v\|^2-(u^{\mathsf T}v)^2}
 =\sqrt{st(n-s-t)}.
\]
(The zero-vector cases are immediate.) The left bilinear form equals $2e_T(S,U)-st$.
\end{proof}
For a vertex $x$, let $B_x$ be the bipartite graph with sides $L_0=N_T^+(x)$ and $R_0=N_T^-(x)$, using arcs directed from $L_0$ to $R_0$ as edges. Both sides have size $b$, and every degree is $a$. For example, an outneighbor $u$ of $x$ has $a-1$ outneighbors in $L_0$, hence $a$ in $R_0$; the other degree follows by the analogous common-inneighbor count. These are the cyclic links used in Round 1, rederived here.
\begin{srcAHlemma}\label{R2E2:lem:link}
The nonconstant singular values of $B_x$ are at most $\sqrt n/2$.
\end{srcAHlemma}
\begin{proof}
For zero-sum vectors on the two sides, the all-ones block vanishes in $B_x=(J_{L_0,R_0}+K_{L_0,R_0})/2$. Extend the vectors by zero to all vertices and apply \eqref{R2E2:eq:skew}. Regularity separates the constant singular direction, whose singular value is $a$.
\end{proof}
\section{Mixed edits}
\begin{proof}[Proof of Theorem \ref{R2E2:thm:mixed}]
Fix $x$ with $L=N_D^+(x)$ of size at most $b$, and put $R=V\setminus(L\cup\{x\})$, so $|R|\ge b\ge|L|$. Every vertex $u\ne x$ sends at least $a-1$ original arcs into $R_0$, and every vertex $v\ne x$ receives at least $a-1$ original arcs from $L_0$. This follows from the DRT common-neighbor counts, allowing the loss of $x$ itself. At most $\srcAHD$ vertices of $R_0$ leave $R$, and at most $\srcAHD$ vertices of $L_0$ leave $L$. At most $\srcAHD$ incident arcs change at each other vertex. Thus the directed bipartite graph from $L$ to $R$ has minimum degree on both sides at least
\[
 \delta=a-1-2\srcAHD.
\]
If Hall's condition fails, there are $S\subseteq L$, $W=N_D^+(S)\cap R$, with $w=|W|<|S|=s_0$. Put $U=R\setminus W$, $t_0=|U|$. There are no arcs $S\to U$ in $D$ and $s_0+t_0>|R|\ge b$. The degree bound gives $w\ge\delta$, hence $s_0\ge\delta+1$. Each vertex of $U$ receives its arcs from $L\setminus S$, so $|L|-s_0\ge\delta$ and
$t_0=|R|-w\ge |R|-s_0+1\ge\delta+1$.

Let $s=\min(s_0,t_0)$ and $t=\max(s_0,t_0)$, without changing the direction of the cut. Every original arc from $S$ to $U$ must have changed, and therefore $e_T(S,U)\le\srcAHD s$. But Lemma \ref{R2E2:lem:mix} implies
\begin{equation}\label{R2E2:eq:cut}
 \frac{e_T(S,U)}s\ge\frac t2-\frac12\sqrt{\frac{t(n-s-t)}s}.
\end{equation}
Write $\tau_n=b/4-n/(2\sqrt b)$. Under the hypothesis,
$s\ge\delta+1\ge n/\sqrt b+1/2>16$ for $n\ge127$.
If $s\ge b/4$, then $t\ge(b+1)/2$ and $t(n-s-t)\le n^2/4$. Hence \eqref{R2E2:eq:cut} is at least
$(b+1)/4-n/(2\sqrt b)=\tau_n+1/4>\srcAHD$.
If $16\le s<b/4$, then $t>3b/4$, and \eqref{R2E2:eq:cut} is greater than
$3b/8-n/16=b/4-1/16>\tau_n\ge\srcAHD$.
Both cases contradict the changed-arc bound. Hall therefore gives a matching saturating $L$. Each matched target is outside $L\cup\{x\}$ and reached via a matched arc, so lies in the exact second neighborhood.

Let $\ell$ be the number of vertices with outdegree at most $b$. Every degree is at least $b-\srcAHD$, while the sum of degrees is at most $nb$. Consequently
$(n-\ell)(b+1)+\ell(b-\srcAHD)\le nb$, giving $\ell\ge n/(\srcAHD+1)$. Finally, in a tournament a vertex of degree greater than $b$ has fewer than its outdegree many vertices outside its first neighborhood, so cannot be strong.
\end{proof}
\section{Deletion-only lower bound}
\begin{srcAHlemma}[Robust regular bipartite matching]\label{R2E2:lem:robust}
Let $B$ be an $a$-regular bipartite graph with equal side sizes $b$, whose nonconstant singular values are at most $\lambda<a$. Removing a subgraph of maximum degree at most $(a-\lambda)/2$ leaves a perfect matching.
\end{srcAHlemma}
\begin{proof}
A Hall failure supplies opposite-side sets of sizes $s,t$ with $s+t>b$ and no remaining edges between them. Exchange side names to assume $s\le t$, and write $x=s/b$, $y=t/b$. Centering their indicator vectors gives
\[
 \frac{e_B(S,U)}s\ge ay-\lambda\sqrt{\frac{y(1-x)(1-y)}x}.
\]
Since $y\ge\max\{x,1-x\}\ge1/2$, one has $y(1-y)\le x(1-x)$, so the square root is at most $1-x$. If $x\le1/2$, then $y>1-x$ and the bound is greater than $(a-\lambda)(1-x)\ge(a-\lambda)/2$. If $x>1/2$, the bound is at least $ax-\lambda(1-x)>(a-\lambda)/2$. Removing all these edges requires maximum degree greater than $(a-\lambda)/2$, a contradiction.
\end{proof}
\begin{proof}[Lower bound in Theorem \ref{R2E2:thm:delete}]
For $n>3$, use Lemmas \ref{R2E2:lem:link} and \ref{R2E2:lem:robust} with $\lambda=\sqrt n/2<a$. Every cyclic link retains a perfect matching under $\srcAHD\le(a-\sqrt n/2)/2=L_n$. Restrict that matching to the surviving outneighbors of its root. Its targets were old inneighbors and cannot become outneighbors under deletion, so the restricted matching certifies strength. For $n=3$, the hypothesis forces $\srcAHD=0$, and the conclusion follows directly from the directed triangle. Since every deletion set below or at $L_n$ is safe, $R_{\rm del}(T)>L_n$.
\end{proof}
\section{A matching obstruction at the sharp scale}
\begin{srcAHlemma}[Sampling bound]\label{R2E2:lem:tail}
If $X$ counts marked elements in a uniform $k$-subset of a $b$-element set with marked proportion $p$, then
$\Pr(X\ge kp+u)\le\exp(-2u^2/k)$ for $u\ge0$.
\end{srcAHlemma}
\begin{proof}
For $t\ge0$, the exponential moment is the normalized elementary symmetric polynomial of degree $k$ in $b$ nonnegative numbers, each equal to $1$ or $e^t$. This is at most their arithmetic mean to the power $k$. To see the inequality, averaging two variables with fixed sum increases every elementary symmetric polynomial: the terms depending on both have nonnegative coefficient and their product increases. Iterating such averaging converges to all variables equal. Thus
$\mathbb E e^{tX}\le(1-p+pe^t)^k$.
The second derivative of $\log(1-p+pe^t)$ is at most $1/4$, so
$\log(1-p+pe^t)-pt\le t^2/8$. Markov's inequality followed by $t=4u/k$ gives the assertion.
\end{proof}
\begin{proof}[Upper bound in Theorem \ref{R2E2:thm:delete}]
Fix $x$, and independently choose uniform $a$-subsets $S\subseteq L_0$ and $U\subseteq R_0$. Delete every arc of the cyclic link from $S$ to $U$, and no other arcs. The first neighborhood of $x$ stays $L_0$, while the $a$ vertices of $S$ have at most $b-a=a-1$ possible targets in $R_0\setminus U$. Thus $x$ is not strong.

For each vertex in $L_0$, its number of neighbors in $U$ is hypergeometric with sample size $a$, marked proportion $a/b$, and mean $\mu=a^2/b$. The analogous assertion holds for each vertex of $R_0$ with respect to $S$. Put $u=\sqrt{(a/2)\log(4b)}$. Lemma \ref{R2E2:lem:tail} and a union bound over $2b$ vertices show that the probability any of these potential deleted degrees is at least $\mu+u$ is at most $1/2$. Hence there exists a choice with maximum deleted degree at most $\lceil\mu+u\rceil$. All other vertices have deleted degree zero. This proves the claimed upper bound. Since $a=(n+1)/4$, both bounds divided by $n$ tend to $1/8$.
\end{proof}
\section{Verification, provenance, and limitations}
The companion standard-library Python script checks DRT parameters and the integer identity $KK^{\mathsf T}=nI-J$ for Paley tournaments, computes maximum matchings by augmenting paths, and validates edited neighborhoods directly. It also implements the published 36-vertex counterexample, independently checking that no vertex is strong. Numerical experiments are consistency checks only; no theorem rests on random testing. Exact small cyclic-link Hall-cut enumeration is explicitly distinguished from general deletion optimality.

Concretely, all 27 sampled edit instances at orders $127,191,263$ passed, with maximum deletion degrees $13,20,28$ and mixed/reversal degrees $7,13,21$, respectively. Deleting the undirected cyclic distances $1,\ldots,6$ from the order-$127$ Paley tournament removes $762$ arcs, has incident degree $12$, and preserves all $127$ strong vertices. A sampled balanced Hall cut at the same order gives an explicit degree-$19$ deletion set destroying strength at its root. Exhaustive cyclic-link cut enumeration at orders $7,11,19$ gives cut optima $1,2,3$ over $15,210,43758$ cuts, respectively. Full edit lists and matching deficiencies are recorded in \texttt{checks.json}.

The new contribution relative to Round 1 is local, rather than global, resilience, its mixed-edit degree characterization, and an asymptotically sharp deletion threshold. Round 1's regular cyclic links are reused and reproved. Its global-budget theorem is not counted again. No result from another entrant is a proof dependency. Standard dependencies are Hall's theorem, elementary singular-value estimates, and the elementary sampling argument proved above. The most delicate proof step is the mixed-edit Hall-barrier reduction: both edited sides must retain minimum degree despite movements of the root's neighborhood. The two independent losses of at most $\srcAHD$ are retained explicitly.

The broad strong conjecture cannot be a target because it is false. Trying to preserve every vertex under reversals also fails: a vertex whose degree rises above $b$ cannot be strong. Our low-degree characterization isolates the exact unavoidable obstruction in the proved regime. We do not determine the lower-order term of $R_{\rm del}$, nor establish a sharp mixed-edit threshold. Possible next directions are an $O(\sqrt n)$ upper error via discrepancy theory, a sharper mixed-edit analysis using actual neighborhood movement, and extension to almost doubly regular tournaments with quantified spectral and degree errors. None is claimed here.


\chapter[Arbitrarily Sparse Strong Seymour Sets in Regular Quasirandom Tournaments]{Arbitrarily Sparse Strong Seymour Sets in Regular Quasirandom Tournaments}
\label{source:R3E2}
\noindent\textit{Source development: Round 3, entry 02. Original mathematical arguments retained; notation is local to this chapter.}\par\medskip

\begin{quote}\small\noindent\textbf{Chapter abstract.} 
We construct pairs of regular tournaments with identical order and identical skew-adjacency operator norm, one with every vertex strong and the other with an arbitrarily slowly growing strong-vertex set. Both sequences are quasirandom, including convergence of every fixed labelled tournament density to its random value. The two members differ by reversals of vanishing maximum incident degree relative to their order, and every failed vertex has matching deficiency exactly one. Explicitly, for a prime $p\equiv3\pmod4$ and odd $k\ge13$, the pair has order $pk$, skew norm $k\sqrt p$, and respectively $pk$ and $11p$ strong vertices. The strong-vertex density is consequently not estimable from a bounded uniform induced-vertex sample. A fully certified constant-eleven ladder and Paley substitution prove the result. This is an obstruction to replacing the local structure in the preceding round's resilience theorem by global spectral quasirandomness; it neither refutes regular-tournament existence nor resolves ordinary Seymour's conjecture. The ladder operation and substitution viewpoint have a public antecedent, explicitly credited; the constant eleven is not an extremal improvement over that source's claimed nine.
\end{quote}
\section{Precise statements and current scope}
For an oriented graph, $N^{++}(v)$ consists of vertices reached in two steps but outside $N^+(v)\cup\{v\}$. A vertex is \emph{strong} if the arcs from $N^+(v)$ to $N^{++}(v)$ contain a matching saturating $N^+(v)$. Write $\srcANStr(T)$ for the strong set. For a regular tournament of odd order $n$, every degree is $d=(n-1)/2$, and
\begin{equation}\label{R3E2:eq:exactsecond}
 N^{++}(v)=N^-(v).
\end{equation}
Indeed an inneighbor not reached from any outneighbor would dominate $v$ and all its $d$ outneighbors, contradicting regularity. Let $B_v$ be the bipartite arc graph from $N^+(v)$ to $N^-(v)$, and define $\srcANdefi_T(v)=d-\nu(B_v)$, where $\nu$ is maximum matching size.

The skew adjacency matrix $K_T$ has entry $1$ on an arc, $-1$ on its reverse, and zero on the diagonal. Its Euclidean operator norm is denoted $\srcANnorm{K_T}$. We will prove quasirandomness directly in both cut and finite-pattern formulations, without using an equivalence theorem as a black box.

\begin{srcANtheorem}[An exact pair]\label{R3E2:thm:pair}
For every prime $p\equiv3\pmod4$ and every odd $k\ge13$, there are regular tournaments $T^-_{p,k}$ and $T^+_{p,k}$ on the same $n=pk$ vertices with
\[
 \srcANnorm{K_{T^-}}=\srcANnorm{K_{T^+}}=k\sqrt p,
 \qquad |\srcANStr(T^-)|=11p,\qquad |\srcANStr(T^+)|=pk.
\]
Every nonstrong vertex of $T^-$ has deficiency exactly one. The pair differs only inside $p$ disjoint $k$-vertex sets, hence by at most $p\binom{k}{2}$ arc reversals with maximum incident edit degree at most $k-1$.
\end{srcANtheorem}
\begin{srcANcorollary}[Arbitrarily slow strong counts]\label{R3E2:cor:slow}
Let $f:\mathbb N\to(0,\infty)$ be any nondecreasing function tending to infinity. There exist $n_j\to\infty$ and pairs $T_j^-,T_j^+$ of regular tournaments of order $n_j$ such that
\[
 |\srcANStr(T_j^-)|\le f(n_j),\qquad |\srcANStr(T_j^+)|=n_j,
 \qquad \frac{\srcANnorm{K_{T_j^\pm}}}{n_j}\longrightarrow0.
\]
The relative maximum incident edit degree between the pair tends to zero. Both sequences have every fixed labelled tournament density converging to its random value, while $|\srcANStr(T_j^-)|/n_j\to0$. The construction does not assert a bounded strong count along a quasirandom sequence.
\end{srcANcorollary}
For a simple explicit choice, take $k=2^p+1$ and primes $p\equiv3\pmod4$ tending to infinity. Then $n=p(2^p+1)$ and $|\srcANStr(T^-)|=11p\le11\log_2 n$, while all but these vertices have deficiency one. Choosing much larger $k$ yields the full quantifier in Corollary~\ref{R3E2:cor:slow}.

\paragraph{Status and attribution.}
The universal strong conjectures are false in Bai--Li--Park v2, which reports Dzitsoev's counterexample \cite{R2E2:BLP}. The correction made in Round 2 remains in force. Ordinary Seymour's conjecture is different, and all vertices of our regular examples satisfy its inequality by \eqref{R3E2:eq:exactsecond}. The public repository \cite{R3E2:G} proposes a regular-tournament lower bound of nine and gives a constant-nine construction. Its definitions agree with ours; it labels the lower bound open and its proofs unreviewed. We neither settle that problem nor improve nine. Its two-vertex ladder and substitution factorization motivated our construction. Rather than assume its constant-nine theorem, we provide a different seed, full matching and Hall certificates, and all necessary arguments. The proposed increment is the spectral/quasirandom obstruction and its arbitrarily-slow and edit-instability forms. Historical priority remains unverified.

\section{A fully certified constant-eleven family}
We first construct regular tournaments $F_k$ with exactly eleven strong vertices for every odd $k\ge13$. This auxiliary count is deliberately not advertised as an extremal record.

\subsection{The seed}
Let the eleven vertices be $0,1,\ldots,10$, and put
$S=\{1,2,3\}$, $U=\{4,5,6\}$, $W=\{7,8,9,10\}$.
Orient $S$ and $U$ as cycles in their listed orders; orient
\[
 U\to0\to S\to W\to U,\qquad U\to S.
\]
Inside $W$ use $7\to8\to9\to10\to7$, together with $7\to9$ and $8\to10$.
Finally $0\to7,8$ and $9,10\to0$. This specifies every pair exactly once. Each vertex has outdegree five. The vertex $0$ is not strong: its source set $S$ has exactly the two targets $9,10$ in $N^-(0)$.

Every other vertex is strong, certified by Table~\ref{R3E2:tab:seed}. Each displayed pair is a matching arc; the five tails and five heads are exactly the two shores of the corresponding root's link. Thus no search or classification of all small tournaments is a premise.
\begin{table}[ht]
\centering\small
\begin{tabular}{cl}
\toprule Root & Matching pairs $(\text{tail},\text{head})$\\\midrule
1 & $(2,3),(7,6),(8,5),(9,4),(10,0)$\\
2 & $(3,1),(7,6),(8,5),(9,4),(10,0)$\\
3 & $(1,2),(7,6),(8,5),(9,4),(10,0)$\\
4 & $(0,8),(1,10),(2,9),(3,7),(5,6)$\\
5 & $(0,8),(1,10),(2,9),(3,7),(6,4)$\\
6 & $(0,8),(1,10),(2,9),(3,7),(4,5)$\\
7 & $(4,3),(5,2),(6,1),(8,10),(9,0)$\\
8 & $(4,3),(5,2),(6,1),(9,0),(10,7)$\\
9 & $(0,8),(4,3),(5,2),(6,1),(10,7)$\\
10 & $(0,8),(4,3),(5,2),(6,1),(7,9)$\\\bottomrule
\end{tabular}
\caption{Ten perfect link matchings in the regular eleven-vertex seed.}\label{R3E2:tab:seed}
\end{table}
\subsection{The extension and old-root matchings}
Start with $P_0=S\cup U$ and $Q_0=\{0\}\cup W$, of sizes six and five. At step $i\ge1$ add $x_i,y_i$ and orient
\begin{equation}\label{R3E2:eq:ladder}
 P_{i-1}\to x_i\to Q_{i-1}\to y_i\to P_{i-1},
 \qquad x_i\to y_i,
\end{equation}
then set $P_i=P_{i-1}\cup\{x_i\}$ and $Q_i=Q_{i-1}\cup\{y_i\}$.
The form of this extension is inherited from \cite{R3E2:G}; its use with this seed is verified below. Every old vertex gains one outgoing and one incoming arc. At step $i$, each new vertex has outdegree $5+i$. Thus after $r$ steps the tournament is regular of order $k=11+2r$ and degree $5+r$.

All eleven old vertices are strong once $r\ge1$:
\begin{itemize}
\item For roots $1,\ldots,6$, append $(x_j,y_j)$, $1\le j\le r$, to the base matching.
\item For roots $7,\ldots,10$, remove $(4,3)$ from the base matching and insert
\[
 (4,x_r),\quad(y_1,3),\quad(y_j,x_{j-1})\quad(2\le j\le r).
\]
\item At root $0$, use
\[
 (1,x_r),(2,10),(3,9),(7,6),(8,5),(y_1,4),
 \quad(y_j,x_{j-1})\ (2\le j\le r).
\]
\end{itemize}
Each is a set of $5+r$ disjoint valid link arcs by \eqref{R3E2:eq:ladder}. In particular the original failed root becomes strong at the first step; the assertion is not that its failure persists.

\subsection{Every new root has deficiency exactly one}
For $1\le i\le r$, use $X_{>i}=\{x_j:i<j\le r\}$ and $Y_{>i}=\{y_j:i<j\le r\}$. The following are exact Hall neighborhoods inside the root's inshore:
\begin{equation}\label{R3E2:eq:witness}
\begin{array}{c|c|c}
\text{root}&\text{source set}&\text{its neighborhood in }N^-(\text{root})\\\hline
x_i&W\cup X_{>i}&U\cup Y_{>i}\\
y_i&U\cup Y_{>i}&\{0\}\cup\{x_j:i\le j\le r\}.
\end{array}
\end{equation}
These identities follow by separating seed vertices, earlier ladder vertices, and later ladder vertices in \eqref{R3E2:eq:ladder}. The first source and target sizes are $4+r-i$ and $3+r-i$; the second are $3+r-i$ and $2+r-i$. Hence no new root is strong.

For completeness, explicit size-$4+r$ matchings give the reverse bound on deficiency. At $x_i$ use
\begin{align*}
 &(7,4),(8,5),(9,6),(0,1),(y_1,2),\\
 &(y_j,x_{j-1})\ (2\le j\le i),\qquad
 (x_j,y_j)\ (i<j\le r).
\end{align*}
At $y_i$ use
\begin{align*}
 &(4,0),(5,x_r),(1,7),(2,8),(3,9),\\
 &(x_j,y_j)\ (1\le j<i),\qquad
 (y_j,x_{j-1})\ (i<j\le r).
\end{align*}
The index ranges are empty when appropriate. These are valid disjoint link arcs, each list having $5+(r-1)=4+r$ pairs. We have proved:
\begin{srcANproposition}\label{R3E2:prop:fiber}
For every odd $k\ge13$, the explicitly defined $F_k$ is regular, has exactly eleven strong vertices, and every other vertex has deficiency one.
\end{srcANproposition}

\section{Substitution separates the matching problem}
For tournaments $Q$ and $F$, define $Q[F]$ on $V(Q)\times V(F)$: arcs between different fibers follow $Q$, and arcs inside each fiber follow $F$.
\begin{srcANlemma}\label{R3E2:lem:sub}
If $Q,F$ are regular of orders $m,k$, then $Q[F]$ is regular of order $mk$, and
\[
 (q,v)\in\srcANStr(Q[F])\quad\Longleftrightarrow\quad
 q\in\srcANStr(Q)\ \hbox{and}\ v\in\srcANStr(F).
\]
If every root of $Q$ is strong, then
$\srcANdefi_{Q[F]}((q,v))=\srcANdefi_F(v)$.
\end{srcANlemma}
\begin{proof}
The outdegree is $k(m-1)/2+(k-1)/2=(mk-1)/2$, so exact second neighborhoods equal in-neighborhoods in all three tournaments.
At root $(q,v)$, separate the link sources and targets into vertices in its own fiber and vertices outside it. There are \emph{no link arcs between these parts}. An internal source cannot beat an external target, since that target's quotient vertex dominates $q$. An external source cannot beat an internal target, since $q$ dominates its quotient vertex. Thus the link is a disjoint union of the internal link $B_v(F)$ and the $k$-fold bipartite blowup of $B_q(Q)$.

If $B_q(Q)$ has a perfect matching, match the corresponding fibers by arbitrary bijections to obtain a perfect matching of the outer part. Conversely, a deficient Hall set in $B_q(Q)$ yields a deficient Hall set in its blowup by taking all $k$ copies of each vertex. Hence the outer part is perfectly matchable exactly when the quotient link is. Matchings of disjoint components have additive maximum size, proving the final assertion as well.
\end{proof}
The factorization viewpoint appears in \cite{R3E2:G}; the proof above states explicitly why the exact second-neighborhood requirement causes no extra targets or cross-matchings.

\section{The Paley quotient and its spectrum}
For a prime $p\equiv3\pmod4$, let $P_p$ have vertex set $\mathbb F_p$ and arc $i\to j$ when $j-i$ is a nonzero square. Write $\chi$ for the quadratic character, with $\chi(0)=0$. Its skew matrix is $K_{ij}=\chi(j-i)$.
\begin{srcANlemma}\label{R3E2:lem:paley}
The tournament $P_p$ is regular, all its vertices are strong, and
\[
 K\mathbf1=0,\qquad KK^{\mathsf T}=pI-J,
 \qquad\srcANnorm K=\sqrt p.
\]
\end{srcANlemma}
\begin{proof}
Since $-1$ is a nonsquare, the orientation is well defined; exactly $(p-1)/2$ differences are squares. For $h\ne0$,
$\sum_t\chi(t)\chi(t+h)=-1$: after scaling, count solutions of $z^2=t(t+1)$. The equation
$(2t+1)^2-(2z)^2=1$ has $p-1$ solutions, parametrized by a nonzero first factor in its difference-of-squares factorization. The same count is $p+\sum_t\chi(t(t+1))$. This proves the identity and hence $KK^{\mathsf T}=pI-J$; row sums vanish.

Expanding the two square indicators, with $t=i,j$ excluded, shows any distinct vertices have $(p-3)/4$ common outneighbors and the same number of common inneighbors. Put $b=(p-1)/2$ and $a=(p+1)/4$. The link at any root has $b$ vertices on each side and degree $a$: an outneighbor has $a-1$ outgoing arcs inside the outshore and therefore $a$ to the inshore; the opposite degree follows from the common-inneighbor count. An $a$-regular balanced bipartite graph has a perfect matching, because edge counting gives $a|S|\le a|N(S)|$ for each subset $S$. Hall's theorem proves strength. The norm follows from the displayed matrix identity.
\end{proof}
The DRT link observation and skew identity were used in our Round 2 entry \cite{R3E2:R2}; they are repeated here, not counted again as new results.

\begin{srcANlemma}[Exact substitution norm]\label{R3E2:lem:norm}
For any regular $k$-vertex tournament $F$ and the quotient $P_p$,
\[
 \srcANnorm{K_{P_p[F]}}=k\sqrt p.
\]
\end{srcANlemma}
\begin{proof}
The matrix is
\[
 K_{P_p[F]}=K_{P_p}\otimes J_k+I_p\otimes K_F.
\]
Decompose Euclidean space orthogonally into vectors constant on every fiber and vectors summing to zero on every fiber. Regularity gives $K_F\mathbf1=0$, so both spaces are invariant. On the first, the operator is $kK_{P_p}$; on the second, it is a direct sum of copies of $K_F$. Hence its norm is $\max\{k\sqrt p,\srcANnorm{K_F}\}$. The maximum absolute row and column sums of $K_F$ are $k-1$, giving $\srcANnorm{K_F}\le k-1<k\sqrt p$.
\end{proof}

\begin{proof}[Proof of Theorem~\ref{R3E2:thm:pair}]
Take $T^-_{p,k}=P_p[F_k]$. Lemmas~\ref{R3E2:lem:sub}--\ref{R3E2:lem:norm} and Proposition~\ref{R3E2:prop:fiber} prove its assertions. For the positive member use $T^+_{p,k}=P_p[C_k]$, where $C_k$ is the cyclic regular tournament: writing $k=2d+1$, its arcs are the differences $1,\ldots,d$ modulo $k$. At root zero the matching $a\mapsto a+d$, $1\le a\le d$, proves strength; translation handles all roots. The same norm lemma applies. Between corresponding fibers the two tournaments agree, so all changes are internal and the edit bounds follow.
\end{proof}

\section{Cut discrepancy, every fixed pattern, and the slow quantifier}
\begin{srcANlemma}[Direct spectral counting]\label{R3E2:lem:count}
Let a tournament $T$ have order $n$ and skew norm $\sigma$. For disjoint $S,U$,
\[
 \left|e_T(S,U)-\frac{|S||U|}{2}\right|
 \le\frac\sigma2\sqrt{|S||U|}.
\]
If $H$ is a fixed labelled tournament on $h$ vertices and $N_H(T)$ counts injective arc-preserving maps, then
\begin{equation}\label{R3E2:eq:patterns}
 \left|N_H(T)-2^{-\binom h2}n^h\right|
 \le\binom h2\frac{\sigma+1}{2}\,n^{h-1}.
\end{equation}
\end{srcANlemma}
\begin{proof}
The cut assertion is the operator bound on $\mathbf1_S^{\mathsf T}K_T\mathbf1_U/2$. Let $A$ be the zero-one adjacency matrix. Since $A=(J-I+K_T)/2$, its difference $E=A-J/2$ has norm at most $(\sigma+1)/2$.
Write $N_H(T)$ as a sum over all $h$-tuples of the product of the adjacency entries specified by $H$. A noninjective tuple has a zero factor, so this really counts injective maps. Telescope this product against the product in which every factor is $1/2$, replacing one edge factor at a time. Fix all coordinates except the two endpoints of that edge. Every other factor is either constant or a function of just one of those two coordinates; their product has form $c f(i)g(j)$, with $0\le c,f,g\le1$. The sum of the changed term is bounded in absolute value by
$\srcANnorm E\srcANnorm f\srcANnorm g\le(\sigma+1)n/2$.
Summing over $n^{h-2}$ choices of the fixed coordinates and the $\binom h2$ edge replacements gives \eqref{R3E2:eq:patterns}.
\end{proof}
Thus $\sigma/n\to0$ implies uniform cut discrepancy $o(n^2)$ and the usual random densities of every fixed tournament. These are standard quasirandom features \cite{R3E2:CG,R3E2:KS}; the quantitative estimates needed here have been proved directly.

\begin{proof}[Proof of Corollary~\ref{R3E2:cor:slow}]
Choose increasing primes $p_j\equiv3\pmod4$. There are infinitely many: if their list were finite, $4$ times their product minus one would have a new prime divisor congruent to three modulo four. For each $j$, choose an odd $k_j\ge13$ so large that $k_j\ge j$, $n_j=p_jk_j$ exceeds $n_{j-1}$, and $f(n_j)\ge11p_j$. This is possible because $f(n)\to\infty$.
Theorem~\ref{R3E2:thm:pair} supplies the pair and the desired strong counts. Its norm ratio is exactly $1/\sqrt{p_j}\to0$, its relative incident edit degree is at most $1/p_j$, and Lemma~\ref{R3E2:lem:count} supplies every fixed density. Also $11p_j/n_j=11/k_j\to0$. The strong counts themselves equal $11p_j$ and diverge; no bounded-count conclusion has been smuggled into the arbitrary-slow statement.
\end{proof}

\begin{srcANcorollary}[Failure of bounded induced-sample estimation]\label{R3E2:cor:sampling}
There is no fixed integer $h$ and estimator based only on the induced tournament on $h$ uniformly sampled distinct vertices that, for every sufficiently large regular tournament, estimates $|\srcANStr(T)|/|T|$ to additive error $1/4$ with success probability at least $2/3$. The conclusion allows randomized estimators and knowledge of the total order. The obstructing pairs can both belong to quasirandom sequences.
\end{srcANcorollary}
\begin{proof}
For the pair of order $pk$, couple the sample by choosing the same ordered distinct vertices in each graph. If no two lie in one fiber, the induced tournaments coincide. A given pair of sampled vertices lies in one fiber with probability $(k-1)/(pk-1)\le1/p$. The union bound therefore bounds the total variation distance between the two induced-sample laws by $\binom h2/p$. The same bound holds after any randomized processing of that sample, by using the same internal randomness in the coupling.

Choose primes $p\to\infty$ and odd $k\to\infty$. The true densities are $1$ and $11/k$. For $k>44$, their additive-$1/4$ success intervals are disjoint. If both success probabilities were at least $2/3$, the event that the estimate belongs to the first interval would have probabilities at least $2/3$ and at most $1/3$ in the two cases. This requires total variation distance at least $1/3$, contradicting $\binom h2/p\to0$. The order supplied to the estimator is identical within each pair. Quasirandomness follows from the already proved norm identity.
\end{proof}
This is specifically an obstruction to uniform induced-vertex sampling, not a lower bound for arbitrary adaptive adjacency-query algorithms. Equivalently, the normalized strong count is discontinuous at the random tournament limit even when restricted to regular tournaments.

\section{What this rules out, and what it leaves open}
Round 2 proved deletion resilience of order $n/8$ for doubly regular base tournaments. That class has much more local structure than a regular tournament with $\srcANnorm K=o(n)$. The present pair shows that degree regularity, the global skew norm, every fixed tournament density, and a vanishing-density edit comparison cannot control the number of vertices with a \emph{perfect} second-neighborhood matching. The norm is identical within each pair, not merely asymptotically comparable. Indeed
\[
 |\srcANStr(T^-_{p,k})|\left(\frac{\srcANnorm K}{pk}\right)^2=11.
\]
No converse lower estimate of this form is asserted.

The failures are microscopic: every bad vertex misses exactly one matching edge, while the ordinary second-neighborhood inequality holds with equality at every vertex. Large homogeneous fibers are invisible to fixed-pattern limits when their relative size tends to zero, but their internal Hall obstructions cannot be repaired by any external matching. This explains the distinction rather than suggesting a counterexample to ordinary Seymour.

The construction does not show that a particular codegree hypothesis is logically necessary; other sufficient local conditions may exist. It does not disprove existence of a strong vertex in every regular tournament, the proposed nine-vertex lower bound, or an estimate that quantitatively depends on the spectral convergence rate and further local data. The public constant-nine construction would improve our auxiliary constant if imported as a dependency; eleven is used to keep the full argument auditable from a short seed.

\paragraph{Verification and provenance.}
All frozen Round 2 papers were read; no competitor theorem is a dependency. The exact verifier checks 2,080 ladder-root certificates through $r=40$, independently recomputes all matchings through $r=12$, and checks five paired products through order $143$. It audits integer tensor identities, both invariant subspaces, edit bounds and cut inequalities. A separate bounded reconstruction confirms the public constant-nine examples at orders $15,17,19,21$ only. Infinite claims rest on the written proofs, not these finite diagnostics.

\paragraph{Limits and next directions.}
The attempted global-quasirandom extension led to this obstruction, not a resolution of regular strong existence. A bounded seed-partition exploration supplied the final explicit certificates. Primary-source searches on 20 October 2026 did not locate the precise pair theorem, but do not certify priority. Future questions include lower bounds incorporating local link expansion and sparse strong sets at the stronger spectral scale $O(\sqrt n)$. No named open conjecture is claimed solved.



\chapter[The Sharp Spectral Scale for\ Strong Seymour Matching Deficiency]{The Sharp Spectral Scale for\ Strong Seymour Matching Deficiency}
\label{source:R4E2}
\noindent\textit{Source development: Round 4, entry 02. Original mathematical arguments retained; notation is local to this chapter.}\par\medskip

\begin{quote}\small\noindent\textbf{Chapter abstract.} 
For a regular tournament of order $n$, a vertex whose cyclic-neighborhood matching has deficiency $r\ge1$ forces the skew-adjacency norm to satisfy
$\srcATnorm K^2\ge2(n-1)(r+1)-1$.
In particular, $\srcATnorm K^2<4n-5$ makes every vertex strong Seymour.
The leading constant is optimal for every fixed deficiency: for each $r\ge1$ and every sufficiently large prime $n\equiv3\pmod4$, we construct a regular tournament with maximum deficiency exactly $r$ and
$\srcATnorm K\le\sqrt{(2r+2)n}+140(r+1)$.
A four-cell skew compression proves the obstruction; a bounded cluster planted into a Paley tournament, followed by bounded-load two-arc degree repair, proves sharpness. This resolves the spectral threshold to its leading constant, not the regular strong-existence or nine-vertex lower-bound questions. Historical novelty remains unverified.
\end{quote}

\section{Statements and the exact increment}
A tournament is \emph{regular} when every vertex has outdegree and indegree $d=(n-1)/2$. Its skew-adjacency matrix $K$ has $K_{xy}=1$ for $x\to y$, $-1$ for $y\to x$, and zero diagonal. The norm is its Euclidean operator norm. At a root $v$, let $B_v$ be the bipartite graph of arcs from $N^+(v)$ to $N^-(v)$, and set
\[
 \srcATdefi_T(v)=d-\nu(B_v),\qquad D(T)=\max_v\srcATdefi_T(v).
\]
Here $\nu$ is maximum matching size. Every inneighbor is an exact second outneighbor: otherwise it dominates $v$ and all $d$ outneighbors, contradicting regularity. Thus $v$ is \emph{strong Seymour} exactly when $\srcATdefi_T(v)=0$.

\begin{srcATtheorem}[Universal obstruction]\label{R4E2:thm:lower}
If a regular tournament $T$ of order $n=2d+1$ has a root of deficiency $r\ge1$, then
\begin{equation}\label{R4E2:eq:lower}
 \srcATnorm{K_T}^2\ge 2(n-1)(r+1)-1.
\end{equation}
Consequently, $\srcATnorm{K_T}^2<4n-5$ implies that every vertex is strong. More generally, $\srcATnorm{K_T}^2<2(n-1)(r+1)-1$ implies $D(T)<r$ for each integer $r\ge1$.
\end{srcATtheorem}

\begin{srcATtheorem}[Asymptotically matching construction]\label{R4E2:thm:upper}
Fix an integer $r\ge1$ and put $g=2r+2$. For every prime $n\equiv3\pmod4$ with $n\ge512g$, a deterministic construction produces a regular tournament $T$ of order $n$ with a specified root $v$ such that
\begin{equation}\label{R4E2:eq:upper}
 \srcATdefi_T(v)\ge r,\qquad \srcATnorm{K_T}\le\sqrt{gn}+70g.
\end{equation}
If $n\ge10000g^3$, then $\srcATdefi_T(v)=D(T)=r$.
\end{srcATtheorem}
Thus, for every fixed $r$, the infimum asymptotic norm ratio permitting deficiency at least $r$ is exactly $\sqrt{2r+2}$. More precisely, if $M_r(n)$ is the minimum norm over regular order-$n$ tournaments with $D(T)\ge r$, and $M_r(n)=+\infty$ when there are none, then
\begin{equation}\label{R4E2:eq:liminf}
 \liminf_{\substack{n\to\infty\\ n\text{ odd}}}\frac{M_r(n)}{\sqrt n}=\sqrt{2r+2}.
\end{equation}
The upper bound is along the indicated primes, not an assertion for every large odd order. There are infinitely many such primes: four times the product of a purported complete finite list, minus one, has a new prime divisor congruent to three modulo four.

\paragraph{Target status and provenance.}
The unrestricted strong Seymour assertions are false in Bai--Li--Park v2, which records Dzitsoev's counterexample \cite{R2E2:BLP}. The current public notes \cite{R4E2:G} formulate a still-open regular-tournament lower bound of nine strong vertices, while labeling their own proofs unreviewed. We do not resolve either regular existence or that stronger count. Ordinary Seymour is different and holds at every vertex of every example here.
Our Round 3 paper \cite{R4E2:R3} showed that regular quasirandom tournaments can have arbitrarily sparse strong sets, but used a growing substitution scale. Its final section explicitly left the stronger $O(\sqrt n)$ scale for further study. Here the main increment is the universal deficiency-dependent lower bound and its matching leading-constant construction. The Paley identity is reused from Rounds 2--3 and reproved; neither the public ladder nor our previous constant-eleven family is needed. All six frozen Round 3 papers and their assessment were read; no competitor theorem is a dependency. These are proposed spectral results with conditional novelty, not a claim to settle a pre-existing named conjecture.

\section{A four-cell compression of a Hall obstruction}
We recall the matching-deficiency form of Hall's theorem:
\begin{equation}\label{R4E2:eq:hall}
 d-\nu(B)=\max_{S\subseteq L}(|S|-|N_B(S)|)
\end{equation}
for a bipartite graph with both shores of size $d$. One inequality follows since any matching leaves at least $|S|-|N(S)|$ vertices unmatched. For the reverse inequality, start with a maximum matching, and follow alternating paths from its unmatched left vertices. If the reached left and right sets are $S,N(S)$, every reached right vertex is matched (otherwise augment), its matched partner is reached, and all reached left vertices except the original unmatched ones are matched into that reached right set. Their size difference is exactly the number of unmatched left vertices. This proves \eqref{R4E2:eq:hall}.

Fix a root of deficiency $r\ge1$, and choose $S\subseteq N^+(v)$ attaining \eqref{R4E2:eq:hall}. Set
\[
 U=N^-(v)\setminus N_{B_v}(S),\quad s=|S|,\quad t=|U|,
 \quad W=V(T)\setminus(\{v\}\cup S\cup U),\quad w=|W|.
\]
Then $U\to v\to S$, $U\to S$, and
\begin{equation}\label{R4E2:eq:sizes}
 s+t=d+r,\qquad w=d-r.
\end{equation}
Both $s,t$ are positive. Summing indegrees over $S$, and outdegrees over $U$, respectively, gives
\[
 s(t+1)+\binom s2\le sd,\qquad t(s+1)+\binom t2\le td.
\]
Consequently
\begin{equation}\label{R4E2:eq:range}
 s+2t\le2d-1,\quad 2s+t\le2d-1,
 \qquad 2r+1\le s,t\le d-r-1.
\end{equation}
In particular $d\ge3r+2$ and $w\ge\max(s,t)+1>0$. All divisions below therefore have positive denominators.

Compress $K$ to the orthonormal indicators of the four cells $\{v\},S,U,W$. The resulting real skew matrix $Q$ has upper-triangular entries, in that order,
\begin{equation}\label{R4E2:eq:Q}\begin{gathered}
 Q_{12}=\sqrt s,\quad Q_{13}=-\sqrt t,\quad
 Q_{14}=\frac{t-s}{\sqrt w},\\ Q_{23}=-\sqrt{st},\quad
 Q_{24}=\frac{\sqrt s(t+1)}{\sqrt w},\quad
 Q_{34}=-\frac{\sqrt t(s+1)}{\sqrt w}.
\end{gathered}\end{equation}
Indeed the first three known inter-cell arc sums determine the others because $K\mathbf1=0$; the sum within each cell is zero by skew symmetry. Also $Q(1,\sqrt s,\sqrt t,\sqrt w)^{\mathsf T}=0$. A real skew matrix has even rank, so this nonzero $4\times4$ matrix has rank two. Its two nonzero singular values are equal. Hence
\begin{align}
 \srcATnorm K^2\ge\srcATnorm Q^2
 &=\frac12\srcATnorm Q_{\mathrm F}^2\notag\\
 &=s+t+st+\frac{(t-s)^2+s(t+1)^2+t(s+1)^2}{w}\notag\\
 &=\frac{(n-1)st+n(s+t)}{n-1-s-t}.\label{R4E2:eq:compression}
\end{align}
The last equality is a direct polynomial simplification.

\begin{proof}[Proof of Theorem~\ref{R4E2:thm:lower}]
For the fixed sum $s+t=d+r$, the product $s(d+r-s)$ is concave in $s$. On the interval \eqref{R4E2:eq:range}, its minimum is attained at either endpoint and equals
\[
 st\ge(2r+1)(d-r-1).
\]
Substitute this and \eqref{R4E2:eq:sizes} into \eqref{R4E2:eq:compression}. The result simplifies exactly to
\[
 \srcATnorm K^2\ge
 \frac{2d(2r+1)(d-r-1)+(2d+1)(d+r)}{d-r}
 =4d(r+1)-1,
\]
which is \eqref{R4E2:eq:lower}. Applying this to a deficient root proves the stated consequences. The finite inequality is not claimed attained; its leading constant will be sharp.
\end{proof}

\section{Paley input and a bounded-load degree-repair lemma}
For a prime $n\equiv3\pmod4$, orient $x\to y$ on $\mathbb F_n$ when $y-x$ is a nonzero quadratic residue. Let $P$ denote its skew matrix. The elementary identities we need are
\begin{equation}\label{R4E2:eq:paley}
 P\mathbf1=0,\qquad P^2=J-nI,\qquad \srcATnorm P=\sqrt n.
\end{equation}
For completeness, writing $P_{xy}=\chi(y-x)$, the off-diagonal correlation identity is $\sum_z\chi(z)\chi(z+h)=-1$ for $h\ne0$. After scaling it follows by counting $b^2=a(a+1)$: the equation $(2a+1)^2-(2b)^2=1$ has $n-1$ solutions, whereas the same count is $n+\sum_a\chi(a(a+1))$. Diagonal correlations are $n-1$, yielding \eqref{R4E2:eq:paley}. The row sum is zero because half the nonzero elements are squares. Expanding two character indicators also gives $(n-3)/4$ common outneighbors of each pair. It follows that for each ordered distinct pair $x,y$, the number of intermediate vertices $z$ with $x\to z\to y$ is at least $(n-3)/4$: if $x\to y$ it is $d-1-(n-3)/4$, and otherwise it is $d-(n-3)/4$.

\begin{srcATlemma}[Degree repair inside a large Paley principal tournament]\label{R4E2:lem:repair}
Let $R$ be the complement of a $g$-vertex set in a Paley tournament of order $n\ge512g$, where $g\ge4$. Suppose integer target outdegrees on $R$ have the same total as the current induced tournament and differ from its outdegrees by at most $g$ at each vertex. There is a sequence of directed two-arc path reversals attaining those degrees, with each vertex used as an endpoint at most $g$ times and as an intermediate vertex at most $16g$ times. The final change in the skew matrix has norm at most $66g$.
\end{srcATlemma}
\begin{proof}
Put $m=n-g$ and $L=16g$. While degrees differ, choose an excess vertex $x$ and a deficit vertex $y$. Seek a current path $x\to z\to y$ in $R$, with $z$ used fewer than $L$ times as an intermediate. Reverse its two arcs. This reduces the excess at $x$ and deficit at $y$ by one and does not change the degree of $z$.

At every stage an endpoint has been used at most its initial absolute discrepancy, at most $g$ times, because endpoint moves only reduce that discrepancy. Each vertex consequently has at most $g+2L=33g$ incident reversal occurrences. The initial number of length-two paths from $x$ to $y$ inside $R$ is at least $(n-3)/4-g$. Changes on arcs incident to $x$ or $y$ can invalidate at most $66g$ of these candidates. There are at most $gm$ transfers altogether, since their number is the initial total positive discrepancy, at most $gm$. Therefore at most $gm/L=m/16$ vertices have reached their intermediate capacity. The number of admissible candidates is at least
\begin{equation}\label{R4E2:eq:candidates}
 \frac{n-3}{4}-g-66g-\frac{n-g}{16}>0
 \qquad(n\ge512g).
\end{equation}
The induction continues until the target degrees are attained. Arcs may be reversed more than once; counting occurrences is a valid upper bound for the final changed incidence. A reversed arc changes each of its two skew entries by magnitude two, so every absolute row and column sum of the final difference is at most $66g$. The operator norm is at most the geometric mean of these two maximum sums, giving the claim.
\end{proof}

\section{Planting a prescribed Hall deficit without changing the order}
Fix $r\ge1$, put $g=2r+2$, and take a Paley tournament on $n=2d+1\ge512g$ vertices. Choose one vertex $a$, and choose $g-1$ of its inneighbors as a set $Z$. We \emph{reuse}, rather than delete or add, these vertices as the planted cluster
\[
 G=\{a\}\cup Z=\{v\}\cup S,\qquad v=a,\quad S=Z.
\]
Let $R$ be the remaining $n-g$ vertices. All $d$ original outneighbors of $a$ lie in $R$. Choose $g-1$ of them; let $C$ be $r+1$ of those chosen vertices and let $J_0$ be the other $r$. Put
\[
 O=N_P^+(a)\setminus(C\cup J_0),\qquad
 I=(N_P^-(a)\setminus Z)\cup J_0.
\]
These partition $R$ into sizes
\begin{equation}\label{R4E2:eq:parts}
 |O|=d-g+1=d-2r-1,\quad |C|=r+1,\quad |I|=d-r-1.
\end{equation}
Keep the Paley tournament inside $R$. Inside $G$, orient $v\to S$, and make $S$ a cyclic regular tournament of odd order $2r+1$. Prescribe cross arcs by
\begin{equation}\label{R4E2:eq:plant}
 v\to O,\quad C\cup I\to v,\qquad
 S\to O\cup C,\quad I\to S.
\end{equation}
The degree of $v$ is $(g-1)+(d-g+1)=d$, and that of every $s\in S$ is $r+|O|+|C|=d$. Thus every planted vertex is already regular.

For $x\in R$, let $c_x$ be its number of wins against $G$ in \eqref{R4E2:eq:plant}, namely $0$ on $O$, $1$ on $C$, and $g$ on $I$. Its desired degree inside $R$ is $d-c_x$. Its current Paley degree there is $d-b_x$, where $b_x\in\{0,\ldots,g\}$ was its original number of wins against $G$. The discrepancy is $c_x-b_x$, of magnitude at most $g$. The target degrees have the correct sum: all $g$ planted vertices already have total outdegree $d$, so their wins against $R$ total $gd-\binom g2=gm/2$. Hence $\sum_{x\in R}c_x=gm/2$, and
\[
 \sum_{x\in R}(d-c_x)=m(d-g/2)=\binom m2.
\]
Lemma~\ref{R4E2:lem:repair} makes the whole tournament regular using changes entirely inside $R$.

None of those changes affects the Hall witness. At $v$, the set $S$ is contained in $N^+(v)$, while $N^-(v)=C\cup I$. Its exact neighborhood in the cyclic link is $C$. Therefore
\begin{equation}\label{R4E2:eq:planteddef}
 \srcATdefi_T(v)\ge |S|-|C|=(2r+1)-(r+1)=r.
\end{equation}
The repair is not asserted to preserve every other matching; the spectral bound below controls their maximum possible deficiency instead.

\section{The sharp norm bound for the planted operator}
Let $H=P[R,R]$, $m=|R|$, and let $b\in\{1,-1\}^{R}$ be the original row of $a$ restricted to $R$, considered as a column. Thus $\srcATnorm b=\sqrt m$. The identity $P^2=J-nI$, restricted to rows $R$ and column $a$, gives
\begin{equation}\label{R4E2:eq:Hb}
 Hb=P[R,Z]P[Z,a]-\mathbf1_R,
 \qquad \srcATnorm{Hb}\le g\sqrt m.
\end{equation}
Indeed $P[R,a]=-b$, and there are $g-1$ terms in the product on the right, each a signed $\{1,-1\}$ column. Also $\srcATnorm H\le\sqrt n$, since it is a principal compression of $P$.

Consider first the auxiliary real skew matrix
\[
 B_0=\begin{pmatrix}0&\mathbf1_g b^{\mathsf T}\\-b\mathbf1_g^{\mathsf T}&H\end{pmatrix}.
\]
It need not itself be a tournament matrix; it is only an operator comparison. Set $e=\mathbf1_g/\sqrt g$ and $u=b/\sqrt m$. On the plane spanned by $e,u$, the cross block is a skew $2\times2$ matrix of norm $\sqrt{gm}$. On $u^\perp\subset\mathbb R^R$, retain the compression of $H$, whose norm is at most $\sqrt n$; on $e^\perp\subset\mathbb R^G$ use zero. These three blocks define an orthogonal direct-sum skew operator. Its difference from $B_0$ consists only of the couplings between $u$ and $u^\perp$ from $H$, and has norm $\srcATnorm{Hu}\le g$ by \eqref{R4E2:eq:Hb}. Since $gm\ge n$ at our orders,
\begin{equation}\label{R4E2:eq:basebound}
 \srcATnorm{B_0}\le\sqrt{gm}+g\le\sqrt{gn}+g.
\end{equation}
This decomposition is important: a triangle inequality between the entire cross block and $H$ would lose the sharp leading constant.

Insert the actual tournament inside $G$, costing norm at most $g-1$ by its absolute row sums. Relative to the identical-row cross block in $B_0$, \eqref{R4E2:eq:plant} changes only arcs to $C\cup J_0$: the root changes $g-1$ arcs, and each member of $S$ changes $r$ arcs. Each outside vertex is incident with at most $g$ changed arcs. Thus the cross correction has norm at most $2g$. Finally degree repair costs at most $66g$. Combining the bounds gives
\begin{equation}\label{R4E2:eq:totalbound}
 \srcATnorm{K_T}\le\sqrt{gn}+g+(g-1)+2g+66g
 <\sqrt{gn}+70g.
\end{equation}

\begin{proof}[Completion of Theorem~\ref{R4E2:thm:upper} and \eqref{R4E2:eq:liminf}]
Regularity and \eqref{R4E2:eq:planteddef} have already been proved, and \eqref{R4E2:eq:totalbound} proves the norm assertion. To justify the exact-deficiency conclusion, for $n\ge10000g^3$ we have
\[
 \frac{140g\sqrt{gn}+4900g^2+g+3}{n}
 \le\frac75+\frac{49}{100g}+\frac{g+3}{10000g^3}<2
 \quad(g\ge4).
\]
Consequently
\[
 (\sqrt{gn}+70g)^2<(g+2)(n-1)-1.
\]
Theorem~\ref{R4E2:thm:lower} applied with deficiency $r+1$ rules out any root of deficiency at least $r+1$, because $2(r+2)=g+2$. Together with \eqref{R4E2:eq:planteddef}, this proves $\srcATdefi_T(v)=D(T)=r$.
The universal lower bound divided by $n$ has limit $2r+2$; the constructed upper bound divided by $\sqrt n$ tends to $\sqrt{2r+2}$ along the primes. This proves \eqref{R4E2:eq:liminf}.
\end{proof}

\section{Interpretation, verification, and remaining gaps}
For $r=1$, the sharp leading threshold is $2\sqrt n$: below $\sqrt{4n-5}$ every root is strong, whereas examples with a failed root have norm at most $2\sqrt n+280$. Paley itself has norm $\sqrt n$. Thus the positive spectral criterion is nonvacuous, and failure can occur at a constant multiple of the minimum spectral scale, not merely at the weaker $o(n)$ scale of Round 3. For each fixed $r$, the examples are quasirandom since $\srcATnorm K/n\to0$; the cut and fixed-pattern estimates in Round 3, Lemma 7, establish this directly. We do not count that consequence again as a new quasirandomness theorem.

This paper controls the \emph{worst local deficiency}, not the number of strong vertices. It does not extend Round 3's sparse-strong-count conclusion to the $O(\sqrt n)$ scale, nor prove regular strong existence or the public nine-vertex lower bound. It does not establish an exact finite-order extremum or an every-order sharp construction. The additive constant $70g$ and the explicit order cutoffs are deliberately loose. The construction permits large changes incident with the planted cluster; only the subsequent outside degree repair has bounded incident load.

\paragraph{Checks.}
The standard-library verifier performs 180,075 exact rational checks of the compression identity and extremal parameter bound. It checks the repair-availability and exact-deficiency cutoff inequalities, constructs four examples at orders $127,251,2063,3083$, and independently computes maximum root matchings. Their root deficiencies are respectively $1,2,1,2$. The first two are below the general repair cutoff and are diagnostics of the successful algorithm, not evidence that the proved cutoff was unnecessary in every case. It verifies every degree, every pair orientation, the exact Hall neighborhood, reversal loads, and the integer identity \eqref{R4E2:eq:Hb}. No floating eigenvalue is a premise; the operator estimate is the written decomposition above. The maximum tested incident reversal counts are $26,69,131,198$, all within $33g$. The finite tests do not replace the all-order path-availability proof.

\paragraph{Failed routes and novelty limitations.}
Compressing only the Hall source and forbidden target against the constant vector gives a weaker leading obstruction and misses the root's contribution. Treating the Paley remainder and planted cross block by a crude norm triangle inequality also misses sharpness. The successful four-cell compression and the nearly invariant two-dimensional spike avoid these losses. Earlier substitution examples cannot attain the new deficiency-one constant because their nontrivial failed fiber is much larger than four. Searches of primary sources on 20 October 2026 did not identify this precise deficiency-dependent sharp threshold, but absence of a hit is not historical priority. Openly disjoint directed triangle fans at one prescribed root are the equivalent matching problem; they are not the vertex-disjoint triangle-factor problem or general arc-disjoint triangle packing. The latter known results do not automatically imply this theorem. No external submission or author contact was made.

Useful next directions are a sharp additive error, all-order constructions, and a lower bound on the \emph{number} of strong roots under a fixed $\srcATnorm K/\sqrt n$ constraint. These remain questions rather than claimed results.



\chapter[A Sharp Three-Quarter Law and Spectral Tail Bounds for Strong Seymour Matchings]{A Sharp Three-Quarter Law and Spectral Tail Bounds for Strong Seymour Matchings}
\label{source:R5E2}
\noindent\textit{Source development: Round 5, entry 02. Original mathematical arguments retained; notation is local to this chapter.}\par\medskip

\begin{quote}\small\noindent\textbf{Chapter abstract.} 
We pass from a spectral bound on the worst matching deficiency to a sharp bound on how many roots can be deficient. Let $T$ be a regular tournament of order $n$, with skew-adjacency matrix $K$. For $r\ge1$, put $g=2r+2$. If $\srcAZnorm K^2<(g+1)n-(4g^2-g+1)$, at most $\lfloor n/g\rfloor$ vertices have cyclic-link matching deficiency at least $r$. In particular, $\srcAZnorm K^2<5n-61$ forces at least $\lceil3n/4\rceil$ strong Seymour vertices. Every bound is attained at infinitely many orders: for all sufficiently large Paley primes $p$, an explicit regular tournament of order $gp-1$ has exactly $p-1$ vertices of deficiency $r$, no larger deficiency, and norm at most $g\sqrt p+2g-1$. For $r=1$ its strong count is exactly $3p$ and its norm ratio tends to the preceding round's sharp constant $2$. The proof identifies homogeneous local obstructions, packs their near-identical skew rows, and couples many obstructions through disjoint link matchings. These are restricted spectral count theorems, not a resolution of regular strong existence or the proposed nine-vertex lower bound.
\end{quote}

\section{The count theorem and its scope}
A regular tournament of odd order $n=2d+1$ has every outdegree and indegree equal to $d$. Let $B_v$ be the bipartite arc graph from $N^+(v)$ to $N^-(v)$ and define
\[
 \srcAZdefi_T(v)=d-\nu(B_v),\qquad
 b_r(T)=|\{v:\srcAZdefi_T(v)\ge r\}|.
\]
Here $\nu$ denotes maximum matching size. Every inneighbor is an exact second outneighbor: an unreached inneighbor would dominate $v$ and all its $d$ outneighbors, impossible in a regular tournament. Thus a vertex is \emph{strong Seymour} precisely when its deficiency is zero. The real skew-adjacency matrix $K$ has entry $1$ on an arc, $-1$ on its reverse and zero diagonal; its norm is the Euclidean operator norm.

\begin{srcAZtheorem}[A finite-order spectral tail bound]\label{R5E2:thm:count}
Let $r\ge1$, $g=2r+2$, and $C_g=4g^2-g+1$. Every regular tournament satisfying
\begin{equation}\label{R5E2:eq:window}
 \srcAZnorm K^2<(g+1)n-C_g
\end{equation}
has
\begin{equation}\label{R5E2:eq:count}
 b_r(T)\le\left\lfloor\frac ng\right\rfloor.
\end{equation}
Every vertex counted in $b_r(T)$ has deficiency exactly $r$. It admits a $g$-vertex witness consisting of itself and a homogeneous regular $(g-1)$-vertex subtournament, which it either dominates or is dominated by. Witnesses for distinct counted vertices are disjoint.
\end{srcAZtheorem}
A set is \emph{homogeneous} when every vertex outside it either dominates the whole set or is dominated by the whole set. No assumption about the existence of such sets is imposed in the theorem.

\begin{srcAZtheorem}[Exact attainment at infinitely many orders]\label{R5E2:thm:construction}
Fix $r\ge1$, put $g=2r+2$, and let $p\equiv3\pmod4$ be prime with $p\ge16g^2$. There is an explicitly constructible regular tournament $T_{p,r}$ of order $n=gp-1$ such that
\begin{equation}\label{R5E2:eq:examples}
 \srcAZnorm{K_{T_{p,r}}}\le g\sqrt p+2g-1,\qquad
 \max_v\srcAZdefi_{T_{p,r}}(v)=r,\qquad b_r(T_{p,r})=p-1=\lfloor n/g\rfloor.
\end{equation}
It satisfies \eqref{R5E2:eq:window}. Thus \eqref{R5E2:eq:count} is attained at infinitely many orders for every fixed $r$.
\end{srcAZtheorem}
The count bound is universal over regular tournaments of every order; construction attainment is asserted only on the stated prime-derived orders, not on every large odd order.

\begin{srcAZcorollary}[The sharp three-quarter law]\label{R5E2:cor:threequarters}
If $\srcAZnorm K^2<5n-61$, at least $\lceil3n/4\rceil$ vertices are strong. For every prime $p\equiv3\pmod4$ with $p\ge256$, there is a regular tournament of order $4p-1$ satisfying that spectral inequality with exactly $3p=\lceil3n/4\rceil$ strong vertices. Along these examples, $\srcAZnorm K/\sqrt n\to2$.
\end{srcAZcorollary}
More generally, among sequences with $\srcAZnorm K/\sqrt n\to\sqrt{2r+2}$, the largest possible limiting proportion of vertices of deficiency at least $r$ is $1/(2r+2)$. Theorem~\ref{R5E2:thm:count} gives the upper bound eventually; Theorem~\ref{R5E2:thm:construction} attains it. The lower spectral estimate recalled below supplies the matching limit of the constructed norm ratio.

\paragraph{Status, attribution, and the actual increment.}
The general strong conjectures are false in Bai--Li--Park v2 \cite{R2E2:BLP}; ordinary Seymour is a different statement and holds at every vertex here. The public Nine-Mass notes \cite{R5E2:G} ask whether every regular tournament of order at least nine has at least nine strong vertices. That source labels its results unreviewed and the lower-bound problem open. We prove a much larger count only under a restrictive spectral hypothesis, not the unrestricted assertion.
Round 4 Entry 02 \cite{R5E2:R4} proved the sharp leading spectral scale permitting a \emph{single} root of any fixed positive deficiency, and explicitly left the number of strong roots at this scale unresolved. We reuse its Hall compression with attribution and a self-contained proof. The new ingredients are the rigidity of low-spectrum Hall witnesses, their disjointness, and a simultaneous extremal construction attaining the count bound. Round 3's sparse-count construction at the weaker $o(n)$ scale \cite{R5E2:R3} does not imply these results. All six frozen Round 4 papers and their assessment, as well as the Round 3 pool, were read; no competitor theorem is used. Historical priority remains unverified.

\section{Inherited Hall compression and new rigidity}
We first record precisely the earlier input. The deficiency form of Hall's theorem is
\[
 d-\nu(B)=\max_{S\subseteq L}(|S|-|N_B(S)|)
\]
for balanced shores of size $d$. To recall the reverse inequality, take a maximum matching and the alternating paths from its unmatched left vertices. Their reached left set $S$ and reached right set $N(S)$ differ in size by exactly the number of unmatched left vertices, because an unmatched reached right vertex would give an augmenting path. The other inequality follows directly by counting matched vertices in $S$.

\begin{srcAZlemma}[Round 4 Hall compression, reproved]\label{R5E2:lem:compression}
If a root has deficiency $h\ge1$, choose a maximizing Hall set $S\subseteq N^+(v)$ and put $U=N^-(v)\setminus N_{B_v}(S)$. Writing $s=|S|$, $t=|U|$, one has
\begin{equation}\label{R5E2:eq:ranges}
 s+t=d+h,\qquad 2h+1\le s,t\le d-h-1,
\end{equation}
and
\begin{equation}\label{R5E2:eq:compress}
 \srcAZnorm K^2\ge\frac{(n-1)st+n(s+t)}{n-1-s-t}
 \ge2(n-1)(h+1)-1.
\end{equation}
\end{srcAZlemma}
\begin{proof}
The Hall equality gives $s+t=d+h$. The known arcs are $U\to v\to S$ and $U\to S$. Sum indegrees on $S$ and outdegrees on $U$ to get
\[
 s(t+1)+\binom s2\le sd,\qquad t(s+1)+\binom t2\le td.
\]
These give \eqref{R5E2:eq:ranges}. The remaining cell $W$ has size $w=d-h\ge\max(s,t)+1>0$. Compress $K$ to the normalized indicators of $\{v\},S,U,W$. Its six upper entries are
\[
 \sqrt s,\quad-\sqrt t,\quad\frac{t-s}{\sqrt w},\quad
 -\sqrt{st},\quad\frac{\sqrt s(t+1)}{\sqrt w},\quad
 -\frac{\sqrt t(s+1)}{\sqrt w}.
\]
The entries involving $W$ follow from zero row sums. The vector $(1,\sqrt s,\sqrt t,\sqrt w)^{\mathsf T}$ lies in the compressed matrix's kernel. Its even rank is therefore two. Its squared norm is half its squared Frobenius norm, namely the first expression in \eqref{R5E2:eq:compress}. At fixed sum $d+h$, the product $st$ is minimized at the endpoints of \eqref{R5E2:eq:ranges}. Its minimum is $(2h+1)(d-h-1)$. Substitution yields $4d(h+1)-1$, the second bound. This is the originating Round 4 calculation, not new credit here.
\end{proof}

\begin{srcAZlemma}[Rigidity below the next spectral level]\label{R5E2:lem:rigidity}
Under \eqref{R5E2:eq:window}, any root of deficiency at least $r$ has deficiency $r$, and some maximizing Hall witness satisfies $\min(s,t)=2r+1=g-1$. If $s=g-1$, its set $S$ is a homogeneous regular tournament of order $g-1$, dominated by $v$. The skew rows of the $g$ vertices $\{v\}\cup S$ have pairwise squared distances at most $4g-2$. The case $t=g-1$ gives the reversed configuration and the same row-distance bound.
\end{srcAZlemma}
\begin{proof}
Deficiency $h\ge r+1$ would imply
$\srcAZnorm K^2\ge(g+2)(n-1)-1$, which exceeds $(g+1)n-C_g$ because their difference is $n+C_g-g-3>0$. Thus $h=r$.

If both $s,t\ge2r+2$, then $d\ge3r+4$. Concavity of the fixed-sum product gives $st\ge(2r+2)(d-r-2)$. The first bound in \eqref{R5E2:eq:compress} consequently gives
\begin{align*}
 \srcAZnorm K^2
 &\ge(4r+6)d-4r-7-\frac{4r^2+6r}{d-r}\\
 &>(2r+3)n-(8r+10)
   =(g+1)n-(4g+2).
\end{align*}
Here $d-r\ge2r+4$ makes the fraction strictly less than $2r$. Since $C_g>4g+2$ for $g\ge4$, this contradicts \eqref{R5E2:eq:window}. Thus one shore size is $g-1$.

Suppose $s=2r+1$, so $t=d-r-1$. Every member of $S$ already has $t+1$ incoming arcs from $U\cup\{v\}$. Its internal indegree is at most $r$, but their average is exactly $r$, so each internal indegree equals $r$. There can be no incoming arc from $W$ to $S$. Hence $S\to W$, $U\cup\{v\}\to S$, and $S$ is a regular homogeneous set. Exactly $r+1$ vertices of $W$ dominate $v$, since $v$ has $d$ inneighbors and $|U|=d-r-1$. On the other external vertices, the row of $v$ agrees with every row of $S$.

For two distinct rows from $S$, all outside entries agree, and their squared distance is at most $4((g-1)-2)+2=4g-10$: there are at most $g-3$ other internal positions with difference two, and the two diagonal positions contribute one each. For $v$ and $x\in S$, the two diagonal positions contribute two, the other $S$ positions contribute $4r$, and the $r+1$ exceptional external positions contribute $4(r+1)$. The total is $8r+6=4g-2$.

If $t=g-1$, reverse every arc. The cyclic link is transposed, deficiency and skew norm are unchanged, and $U$ is a source-shore Hall set with the same deficiency in the reversed tournament. Indeed its neighborhood is contained in the complement of $S$, and maximal deficiency forces equality. The preceding argument applies, then reversing back gives the claimed sink configuration.
\end{proof}
The homogeneity conclusion is exact, not an assertion that the whole $g$-set is a module. Its root differs from the homogeneous $(g-1)$-set on exactly $r+1$ external vertices.

\section{Packing witnesses by skew-row variance}
\begin{srcAZlemma}[Near-row packing]\label{R5E2:lem:packing}
Let $K$ be the skew matrix of any tournament of order $n$. If there are $q$ distinct rows $a_0,\ldots,a_{q-1}$ such that $\srcAZnorm{a_i-a_0}^2\le c$ for every $i>0$, then
\begin{equation}\label{R5E2:eq:variance}
 \srcAZnorm K^2\ge q(n-1)-(q-1)c.
\end{equation}
\end{srcAZlemma}
\begin{proof}
Every row has squared norm $n-1$. If $\bar a=q^{-1}\sum_i a_i$, then
\[
 \left\|\frac1{\sqrt q}\sum_i a_i\right\|^2
 =q(n-1)-\sum_i\srcAZnorm{a_i-\bar a}^2
 \ge q(n-1)-\sum_i\srcAZnorm{a_i-a_0}^2.
\]
The last inequality uses the minimizing property of the mean, obtained by expanding squared distances. The normalized sum of rows is the image under $K^{\mathsf T}$ of a unit normalized indicator. This proves \eqref{R5E2:eq:variance}.
\end{proof}

\begin{proof}[Proof of Theorem~\ref{R5E2:thm:count}]
For every root counted in $b_r(T)$, choose a witness $G_v$ of size $g$ from Lemma~\ref{R5E2:lem:rigidity}. If two distinct witness sets overlap, choose a common vertex and $g$ other vertices in their union. Every chosen row is within squared distance $4g-2$ of the common row. Lemma~\ref{R5E2:lem:packing} with $q=g+1$ gives
\[
 \srcAZnorm K^2\ge(g+1)(n-1)-g(4g-2)
 =(g+1)n-C_g,
\]
contradicting \eqref{R5E2:eq:window}. Thus two witnesses are either identical or disjoint.

An identical witness cannot belong to two different roots. In its induced tournament each certifying root must be either a source or a sink, and its deletion must leave a regular tournament of order $2r+1$. Two sources or two sinks are impossible. If there were one of each, deleting the source would leave a sink in a regular tournament of positive degree $r$, again impossible. Therefore witnesses for different counted roots are disjoint, and $g b_r(T)\le n$. The exact deficiency and structural assertions were established in Lemma~\ref{R5E2:lem:rigidity}.
\end{proof}

\begin{srcAZcorollary}[A module-free strengthening for strong vertices]\label{R5E2:cor:prime}
A regular tournament with no homogeneous directed triangle has every vertex strong whenever $\srcAZnorm K^2<5n-18$.
\end{srcAZcorollary}
\begin{proof}
A deficient root cannot have deficiency at least two, since Lemma~\ref{R5E2:lem:compression} would give $\srcAZnorm K^2\ge6n-7$. For deficiency one, if $s,t\ge4$, the first computation in Lemma~\ref{R5E2:lem:rigidity} gives $\srcAZnorm K^2>5n-18$. Otherwise the equality case gives a homogeneous regular three-set, which is a directed triangle. Both possibilities contradict the hypotheses.
\end{proof}
This includes prime tournaments, but only under the spectral condition. It does not prove regular strong existence for every prime tournament.

\section{Many simultaneous obstructions at the sharp spectral scale}
\subsection{The Paley quotient and disjoint cyclic-link matchings}
Let $P$ be the Paley skew matrix of prime order $p\equiv3\pmod4$: $P_{ij}=\chi(j-i)$ with quadratic character $\chi(0)=0$. Put $h=(p-1)/2$. We use
\begin{equation}\label{R5E2:eq:paley}
 P\mathbf1=0,\qquad PP^{\mathsf T}=pI-J,\qquad\srcAZnorm P=\sqrt p,
\end{equation}
and the fact that every pair has $(p-3)/4$ common outneighbors and the same number of common inneighbors. These standard Paley facts were used in Rounds 2--4. A short proof of the correlation identity is to count solutions of $z^2=x(x+1)$: the factorization $(2x+1)^2-(2z)^2=1$ gives $p-1$ solutions, while summing quadratic characters gives $p+\sum_x\chi(x(x+1))$. Thus the character correlation is $-1$. Expanding the two square indicators gives the common-neighbor counts and \eqref{R5E2:eq:paley}.

Fix a quotient vertex $a$, and set $O=N_P^+(a)$, $I=N_P^-(a)$. The bipartite graph of arcs $O\to I$ is balanced, of degree
\[
 k=\frac{p+1}{4}.
\]
For example an $i\in O$ has $h$ outneighbors, $(p-3)/4$ of them in $O$, and does not beat $a$, leaving $k$ in $I$; common inneighbors give the opposite degree. A positive-degree regular balanced bipartite graph has a perfect matching by edge counting and Hall's theorem. Deleting such a matching preserves regularity. Therefore it has $k$ disjoint perfect matchings. Choose $g/2=r+1$ of them, denoted $M_1,\ldots,M_{g/2}:O\to I$. Our order cutoff makes $k\ge g/2$. A lexicographically specified matching algorithm makes this choice deterministic if desired.

\subsection{Fibers, degrees, and exact deficient neighborhoods}
Replace $a$ by a cyclic regular tournament $F_a$ of order $g-1=2r+1$. For every $i\ne a$, use a fiber
\[
 F_i=\{v_i\}\cup S_i,\qquad |S_i|=g-1,
\]
where $S_i$ is cyclic regular of degree $r$. For $i\in O$, orient $v_i\to S_i$; for $i\in I$, orient $S_i\to v_i$. Between fibers use the Paley orientation. The total order is
\[
 n=(g-1)+g(p-1)=gp-1,
 \qquad d=(n-1)/2=gh-1+g/2.
\]
This initial tournament is not regular. Its $S_i$ vertices and its $F_a$ vertices already have degree $d$:
\[
 \begin{array}{c|c|c}
 \text{vertex type}&\text{external outdegree}&\text{internal outdegree}\\\hline
 F_a&gh&r\\
 S_i,\ i\in O&gh&r\\
 S_i,\ i\in I&gh-1&r+1.
 \end{array}
\]
Here $d=gh+r$ since $r=g/2-1$. Each $v_i$ for $i\in O$ has degree $d+g/2$, and each $v_j$ for $j\in I$ has degree $d-g/2$.

For each matching arc $i\to M_\ell(i)$, reverse the arc
\begin{equation}\label{R5E2:eq:reversal}
 v_i\longrightarrow v_{M_\ell(i)}.
\end{equation}
The matchings are edge-disjoint, and each root occurs exactly $g/2$ times. This corrects every root degree and changes no other degree. The resulting tournament is regular.

At a source root $v_i$, $i\in O$, the exact neighborhood of $S_i$ in $N^-(v_i)$ is
\begin{equation}\label{R5E2:eq:hallsource}
 N_{B_{v_i}}(S_i)=\{v_{M_\ell(i)}:1\le\ell\le g/2\}.
\end{equation}
Before the reversals, no external inneighbor of $v_i$ is beaten by $S_i$; afterward the only new inneighbors are these roots, and all of $S_i$ beats them. Thus $|S_i|-|N(S_i)|=g-1-g/2=r$.
At a sink root $v_j$, $j\in I$, the exact set of vertices of $N^+(v_j)$ sending arcs into the target set $S_j$ is
\begin{equation}\label{R5E2:eq:hallsink}
 \{v_{M_\ell^{-1}(j)}:1\le\ell\le g/2\}.
\end{equation}
This is the right-shore version of the same Hall deficiency. Hence all $p-1$ distinguished roots have deficiency at least $r$. The triples, or larger odd sets, $S_i$ remain homogeneous, because the only cross-fiber edits joined distinguished roots.

\subsection{The spectral estimate and the complete cutoff}
Let $K_0$ contain only the cross-fiber skew entries, with zero blocks on the fibers. On the subspace constant on each fiber it is unitarily represented by $DPD$, where $D$ is diagonal with entries $\sqrt{|F_i|}$. It is zero on the orthogonal complement. Since every fiber size is at most $g$,
\[
 \srcAZnorm{K_0}=\srcAZnorm{DPD}\le\srcAZnorm D^2\srcAZnorm P\le g\sqrt p.
\]
The internal-fiber skew matrices, whether regular or not, have norm at most $g-1$ by their absolute row and column sums. The reversal matrix has maximum absolute row and column sum $2(g/2)=g$. Therefore
\begin{equation}\label{R5E2:eq:normbound}
 \srcAZnorm{K_{T_{p,r}}}\le g\sqrt p+(g-1)+g=g\sqrt p+2g-1.
\end{equation}
No regular-substitution matching factorization is used: the even fibers are deliberately nonregular before the degree correction.

\begin{proof}[Proof of Theorem~\ref{R5E2:thm:construction} and Corollary~\ref{R5E2:cor:threequarters}]
It remains to verify the spectral window including all additive terms. With $x=\sqrt p$, its right side minus the square of \eqref{R5E2:eq:normbound} is
\begin{align*}
 &(g+1)(gp-1)-C_g-(g\sqrt p+2g-1)^2\\
 &\hspace{1cm}=g\{x^2-(4g-2)x\}-(8g^2-4g+3).
\end{align*}
At $x=4g$ this equals $4g-3>0$, and its derivative is $g(2x-4g+2)>0$ for $x\ge4g$. Thus $p\ge16g^2$ implies the strict window \eqref{R5E2:eq:window}.

The $p-1$ distinguished roots are deficient by at least $r$, while Theorem~\ref{R5E2:thm:count} permits at most $\lfloor(gp-1)/g\rfloor=p-1$ such roots and excludes deficiency above $r$. Therefore these and only these roots have deficiency at least $r$, all with exact deficiency $r$. Other vertices have deficiency less than $r$; for $r>1$ we do not assert here that they are all strong. This proves Theorem~\ref{R5E2:thm:construction}. There are infinitely many primes congruent to three modulo four, as follows from the usual product-minus-one argument.

For $r=1$, the number of nonstrong vertices is precisely $p-1$, so the strong count is $gp-1-(p-1)=3p$. The universal count theorem gives $n-\lfloor n/4\rfloor=\lceil3n/4\rceil$. Finally Lemma~\ref{R5E2:lem:compression} gives the lower norm ratio $\srcAZnorm K/\sqrt n\ge\sqrt{4-5/n}$ in these deficient examples; \eqref{R5E2:eq:normbound} gives the upper limit $2$. The same argument with general $r$ gives norm ratio tending to $\sqrt g$.
\end{proof}

\section{What has changed, and what remains unresolved}
There is now a sharp count result at the first spectral threshold: the previous round forced every root strong below approximately $2\sqrt n$, and the present construction realizes a quarter of all roots failing with norm ratio tending to $2$ from the allowed side. Within the larger finite window $\srcAZnorm K^2<5n-61$, more than a quarter cannot fail. The constant $61$ is a convenient rigorous overlap allowance, not claimed optimal. Nor is the coefficient $5$ claimed to be the greatest possible spectral window for the same three-quarter conclusion.

For larger $r$, the theorem concerns the \emph{tail} of the deficiency distribution, not the total number of strong roots. The extremal examples have a positive density of prescribed failures, unlike the single planted obstruction of Round 4, and use matchings to correct all their degrees simultaneously. The norm estimate is at the same sharp leading scale, with a different construction and a new count conclusion. Reusing the earlier Hall calculation is essential and explicitly credited; it is not itself a fifth-round discovery.

The hypotheses are much stronger than regularity or ordinary quasirandomness. We do not settle the public nine-vertex lower bound, regular strong existence, the exact finite minimum at every order, or the full tradeoff between strong density and arbitrary fixed $\srcAZnorm K/\sqrt n$. The module-free corollary does not resolve the prime-tournament case without its spectral bound. The known universal strong counterexample and the distinction from ordinary Seymour remain in force.

\paragraph{Reproducibility.}
The supplied standard-library verifier checks the rigidity parameter inequalities exactly, constructs disjoint Paley link matchings, and builds the finite tournaments from scratch. It verifies every degree and pair orientation, every prescribed Hall neighborhood, all within-witness row distances, and the exact radical-free cutoff comparison. Independent maximum-matching computations test all roots at small orders and a stated sample at the larger orders. The all-order count and spectral conclusions follow from the proofs above, not finite matching tests or numerical eigenvalues. The exact output and tested orders are retained in \texttt{verification.txt}; no floating eigenvalue is used as a certificate.

\paragraph{Failed avenues and source discipline.}
An initial construction with a five-vertex exceptional fiber raised the available leading norm bound to the next spectral scale. Replacing that fiber by a smaller odd fiber, and choosing opposite root orientations on the two quotient shores, gives the successful degree balance. Crude independent planting and repair would not have controlled a positive density of failures at the desired constant. Searches of current primary sources on 20 October 2026 did not locate the precise spectral count theorem, but do not establish priority. The public module and substitution viewpoint is an antecedent \cite{R5E2:G}; its ladder theorem is not a premise of this paper. No external posting, submission, or author contact occurred.

Next questions include improving the finite overlap constant, describing equality structures beyond the constructed orders, and extending the count tradeoff beyond the present spectral windows. These are not claimed results.



\chapter[A Sharp Four-Fifths Law for Prime Regular Tournaments]{A Sharp Four-Fifths Law for Prime Regular Tournaments}
\label{source:R6E2}
\noindent\textit{Source development: Round 6, entry 02. Original mathematical arguments retained; notation is local to this chapter.}\par\medskip

\begin{quote}\small\noindent\textbf{Chapter abstract.} 
We determine a new sharp spectral density regime for strong Seymour matchings after excluding homogeneous directed triangles. Every regular tournament of order $n$ with this exclusion and skew norm squared less than $6n-116$ has at least $\lceil4n/5\rceil$ strong vertices. This includes prime tournaments. The proportion is asymptotically optimal even in the prime class: for all sufficiently large Paley primes $p$, we construct prime regular tournaments of order $5p$ with strong proportion tending to $4/5$ and skew norm divided by $\sqrt{5p}$ tending to $\sqrt5$. Combined with the preceding round's all-strong bound below $5n-18$, this also identifies the sharp leading spectral onset of matching failure in prime regular tournaments. The construction uses many four-cycle reversals through a sublinear reservoir, followed by bounded-load triangle reversals that destroy every remaining module. It is not a substitution example: the final tournaments are prime. Universal strong existence and the public nine-vertex lower-bound question remain unresolved.
\end{quote}

\section{The two results and the precise increment}
A tournament is \emph{regular} when every outdegree and indegree is $d=(n-1)/2$. A set $H$ is \emph{homogeneous} if each outside vertex either dominates all of $H$ or is dominated by all of $H$. A tournament is \emph{prime} if its only homogeneous sets have sizes $0,1,n$. Our exclusion below means no homogeneous three-set inducing a directed cycle, not absence of directed triangles.

For a root $v$, let $B_v$ be the bipartite arc graph from $N^+(v)$ to $N^-(v)$, and put
\[
 \srcBFdefi_T(v)=d-\nu(B_v),\qquad b(T)=|\{v:\srcBFdefi_T(v)>0\}|.
\]
Every inneighbor is an exact second outneighbor: an unreached one would dominate $v$ and all its $d$ outneighbors. Thus deficiency zero is exactly the strong Seymour matching condition. Write $K_{xy}=1$ for $x\to y$, $K_{xy}=-1$ for the reverse, and $K_{xx}=0$; $\srcBFnorm K$ is its Euclidean operator norm.

\begin{srcBFtheorem}[A universal finite-order four-fifths bound]\label{R6E2:thm:count}
Let $T$ be a regular tournament of order $n$ having no homogeneous directed triangle. If
\begin{equation}\label{R6E2:eq:window}
 \srcBFnorm K^2<6n-116,
\end{equation}
then
\[
 b(T)\le\lfloor n/5\rfloor,
 \qquad |\operatorname{Str}(T)|\ge\lceil4n/5\rceil.
\]
Every failed root has deficiency exactly one. It has a five-vertex Hall witness, and witnesses belonging to distinct failed roots are disjoint.
\end{srcBFtheorem}

\begin{srcBFtheorem}[Asymptotic attainment by prime regular tournaments]\label{R6E2:thm:construction}
For every prime $p\equiv3\pmod4$ with $p\ge10^{24}$, there exists a prime regular tournament $T_p$ of order $n=5p$ such that
\begin{equation}\label{R6E2:eq:construction}
 \srcBFnorm{K_{T_p}}\le5\sqrt p+8001p^{1/4},\qquad
 4p\le|\operatorname{Str}(T_p)|\le4p+2p^{3/4}.
\end{equation}
Every failed root has deficiency one. In particular, along these orders,
\[
 \frac{|\operatorname{Str}(T_p)|}{5p}\longrightarrow\frac45,
 \qquad \frac{\srcBFnorm{K_{T_p}}}{\sqrt{5p}}\longrightarrow\sqrt5.
\]
The construction satisfies \eqref{R6E2:eq:window}. The lower count in Theorem~\ref{R6E2:thm:count} is universal at every regular order; the attainment statement is asymptotic on the specified prime-derived orders, not exact at every order.
\end{srcBFtheorem}
The large cutoff is a convenient sufficient value, not an optimized bound. The finite diagnostics below also construct much smaller prime examples, without invoking this cutoff theorem.

\paragraph{The preceding threshold, credited rather than reclaimed.}
Round 5 Entry 02 \cite[Corollary 7]{R6E2:R5} proved that no homogeneous directed triangle together with $\srcBFnorm K^2<5n-18$ forces \emph{every} root strong. We reproduce that implication below. Theorem~\ref{R6E2:thm:construction} shows its leading coefficient $5$ is sharp, even for prime tournaments. More generally, among prime regular sequences with norm ratio tending to $\sqrt5$, the minimum possible limiting strong proportion is $4/5$: Theorem~\ref{R6E2:thm:count} gives the eventual lower bound, and Theorem~\ref{R6E2:thm:construction} attains its limit.

\paragraph{Target, source corrections, and novelty limits.}
The universal strong conjectures are false in Bai--Li--Park v2 \cite{R2E2:BLP}; ordinary Seymour is different and holds at every root of a regular tournament. The current public Nine-Mass notes \cite{R6E2:G} label the regular nine-strong-vertex lower-bound problem open and unreviewed, and emphasize prime instances after module reduction. Neither unrestricted existence nor that lower bound is settled here. We study a restricted spectral extremal problem within that program.

Round 4 Entry 02 \cite{R6E2:R4} supplied the Hall compression and the single-deficiency spectral scale. Round 5 Entry 02 \cite{R6E2:R5} supplied the first density theorem, near-row packing, and simultaneous even-fiber examples with homogeneous odd obstructions. They are genuine premises, reproduced with attribution. The new increment is a \emph{prime} extremal family and a sharp $4/5$ density regime after removing homogeneous triangles. Its four-element Hall shores are not modules. Its reservoir surgery differs from the previous disjoint-matching correction, and destroys all nontrivial modules rather than relying on them. All six frozen Round 5 manuscripts and their assessment were read; no live Round 6 opponent manuscript was inspected. Targeted primary-source searches on 20 October 2026 did not establish historical priority.

\section{Hall compression and five-vertex witnesses}
\begin{srcBFlemma}[Inherited Hall compression]\label{R6E2:lem:hall}
If a root has deficiency $h\ge1$, there are sets $S\subseteq N^+(v)$ and $U\subseteq N^-(v)$, with $U\to S$, whose sizes satisfy
\begin{equation}\label{R6E2:eq:ranges}
 s+t=d+h,\quad 2h+1\le s,t\le d-h-1,
 \quad w:=n-1-s-t=d-h\ge\max(s,t)+1>0.
\end{equation}
Moreover
\begin{equation}\label{R6E2:eq:hall}
 \srcBFnorm K^2\ge\frac{(n-1)st+n(s+t)}{n-1-s-t}
 \ge2(n-1)(h+1)-1.
\end{equation}
\end{srcBFlemma}
\begin{proof}
The deficiency version of Hall's theorem says $h=\max_S(|S|-|N_{B_v}(S)|)$. It follows, for example, by taking alternating paths from the unmatched left vertices of a maximum matching; their reached shores differ by exactly the number of unmatched vertices. Choose a maximizing set and put $U=N^-(v)\setminus N_{B_v}(S)$. Thus $s+t=d+h$. Summing indegrees on $S$ and outdegrees on $U$ gives
\[
 s(t+1)+\binom s2\le sd,\qquad t(s+1)+\binom t2\le td.
\]
These imply \eqref{R6E2:eq:ranges}. Compress $K$ to normalized indicators of $\{v\},S,U,W$, where $W$ is the remaining set. Its six upper entries are
\[
 \sqrt s,\ -\sqrt t,\ (t-s)/\sqrt w,\ -\sqrt{st},\
 \sqrt s(t+1)/\sqrt w,\ -\sqrt t(s+1)/\sqrt w.
\]
Zero row sums determine the last three cell interactions involving $W$. The vector $(1,\sqrt s,\sqrt t,\sqrt w)^{\mathsf T}$ is in the kernel, so this nonzero real skew matrix has rank two. Its squared norm is half its squared Frobenius norm, the first expression in \eqref{R6E2:eq:hall}. At fixed sum $d+h$, the product $st$ is minimized at the allowed endpoints, giving $(2h+1)(d-h-1)$. Substitution gives $4d(h+1)-1$. This is the Round 4 calculation, not a new result.
\end{proof}

\begin{srcBFlemma}[New four-shore rigidity]\label{R6E2:lem:rigidity}
Under the hypotheses of Theorem~\ref{R6E2:thm:count}, a failed root has deficiency one and a maximizing Hall pair with $\min(s,t)=4$. If $s=4$, its witness $G_v=\{v\}\cup S$ has these properties:
\begin{enumerate}
\item $v$ dominates $S$, every internal indegree in $S$ is at most two, and exactly two arcs enter $S$ from $W$ in total.
\item Squared distances between distinct full skew rows in $G_v$ are at most $22$ (and at most $14$ for two rows in $S$).
\end{enumerate}
For $t=4$ the reversed statements hold.
\end{srcBFlemma}
\begin{proof}
Deficiency at least two would force $\srcBFnorm K^2\ge6n-7$, contradicting \eqref{R6E2:eq:window}. Thus $h=1$, $s+t=d+1$, and $s,t\ge3$. If $s=3$, then $t=d-2$. Each member of $S$ already has $d-1$ incoming arcs from $U\cup\{v\}$, so its internal indegree is at most one. The average is one, hence all equal one and no arc enters from $W$. Thus $S$ is a homogeneous directed triangle. The same conclusion holds if $t=3$ after reversal. These cases are excluded.

If $s,t\ge5$, then $d\ge9$ and the fixed-sum product is at least $5(d-4)$. Compression gives
\[
 \srcBFnorm K^2\ge12d-25-\frac{24}{d-1}
 =6n-31-\frac{24}{d-1}\ge6n-34,
\]
again contradicting \eqref{R6E2:eq:window}. Consequently the smaller shore has size four.

Suppose $s=4$, so $t=d-3$. For $x\in S$ let $a_x$ be its internal indegree. Its incoming degree from $W$ is exactly $2-a_x$, hence $a_x\le2$ and these incoming degrees sum to $8-6=2$. The root has exactly three inneighbors in $W$. The rows of $v,x$ agree on $U$. Their two diagonal positions contribute two to squared distance; other positions in $S$ contribute $4a_x$; the differing positions in $W$ number at most $3+(2-a_x)$. Thus
\[
 \srcBFnorm{K_{v,*}-K_{x,*}}^2\le2+4a_x+4(5-a_x)=22.
\]
For $x,y\in S$, among the other two positions of $S$ at most $a_x+a_y-1$ can differ: the edge $xy$ contributes one to their total indegrees. Their differing external positions number at most $(2-a_x)+(2-a_y)$. The squared distance is therefore at most
\[
 2+4(a_x+a_y-1)+4(4-a_x-a_y)=14.
\]
Reversal transposes the cyclic link and preserves deficiency and norm. If $t=4$, $U$ is a maximizing left Hall set in the reversed tournament: its neighborhood is contained in the complement of $S$, and maximal deficiency forces equality. This proves the reversed case.
\end{proof}

\begin{srcBFlemma}[Near-row packing, from Round 5]\label{R6E2:lem:packing}
If $q$ distinct skew rows $a_0,\ldots,a_{q-1}$ of an order-$n$ tournament satisfy $\srcBFnorm{a_i-a_0}^2\le c$ for $i>0$, then
\[
 \srcBFnorm K^2\ge q(n-1)-(q-1)c.
\]
\end{srcBFlemma}
\begin{proof}
Every row has squared norm $n-1$. With $\bar a=q^{-1}\sum a_i$, the mean minimizes the sum of squared distances, so
\[
 \left\|q^{-1/2}\sum_i a_i\right\|^2
 =q(n-1)-\sum_i\srcBFnorm{a_i-\bar a}^2
 \ge q(n-1)-\sum_i\srcBFnorm{a_i-a_0}^2.
\]
The left side is at most $\srcBFnorm K^2$, being the squared norm of the image of a unit vector under $K^{\mathsf T}$.
\end{proof}

\begin{proof}[Proof of Theorem~\ref{R6E2:thm:count}]
Choose a five-set witness for every failed root. If two different witness sets intersect, their union contains six vertices, including a common vertex. Every chosen row lies within squared distance $22$ of that common row. Lemma~\ref{R6E2:lem:packing} gives $\srcBFnorm K^2\ge6(n-1)-5\cdot22=6n-116$, a contradiction.

An identical five-set cannot certify two roots. Each certifying root is a source or a sink in its induced five-set; two sources or two sinks are impossible. If it had both, the complement of its source would contain a vertex of indegree three, contradicting Lemma~\ref{R6E2:lem:rigidity}. Thus witnesses for distinct roots are disjoint. There are at most $\lfloor n/5\rfloor$ of them.
\end{proof}

\begin{srcBFcorollary}[The inherited onset bound]\label{R6E2:cor:onset}
A regular tournament without a homogeneous directed triangle has every root strong if $\srcBFnorm K^2<5n-18$.
\end{srcBFcorollary}
\begin{proof}
This is Round 5, Corollary 7. To reproduce it, deficiency at least two gives $6n-7$. For deficiency one, the preceding argument excludes a shore of size three. If $s,t\ge4$, then $d\ge7$ and compression gives
\[
 \srcBFnorm K^2\ge10d-11-\frac{10}{d-1}
 =5n-16-\frac{10}{d-1}>5n-18.
\]
Every possible failure is excluded.
\end{proof}

\section{A small reservoir in a Paley tournament}
Let $P$ be the Paley skew matrix of prime order $p\equiv3\pmod4$, with $P_{ij}=\chi(j-i)$ for the quadratic character and $\chi(0)=0$. We use
\begin{equation}\label{R6E2:eq:paley}
 P\mathbf1=0,\qquad PP^{\mathsf T}=pI-J,\qquad\srcBFnorm P=\sqrt p.
\end{equation}
For completeness, the character correlation is $-1$: the equation $z^2=x(x+1)$ has $p-1$ solutions by $(2x+1-2z)(2x+1+2z)=1$, but its number of solutions is $p+\sum_x\chi(x(x+1))$. Translation and scaling give all off-diagonal correlations. The diagonal and row-sum assertions are immediate. In particular, for disjoint quotient sets $A,B$,
\begin{equation}\label{R6E2:eq:mixing}
 e(A,B)\ge\frac{|A||B|-\sqrt{p|A||B|}}2.
\end{equation}
Indeed $2e(A,B)-|A||B|=\mathbf1_A^{\mathsf T}P\mathbf1_B$, whose absolute value is at most $\sqrt{p|A||B|}$.

\begin{srcBFlemma}[Reservoir and two bounded-load schedules]\label{R6E2:lem:reservoir}
Let $p\ge10^{24}$ be as above, and put $m=p^{3/4}$ and $L=\lceil1000p^{1/4}\rceil$. There is a quotient set $R$ with
\[
 |R|\le2m,\qquad |N^+(i)\cap R|,\ |N^-(i)\cap R|\ge m/3
 \quad\text{for every }i.
\]
It supports the following schedules, where all directions refer to the original Paley tournament.
\begin{enumerate}
\item For each $i\notin R$, choose $j_i,k_i\in R$ with $i\to j_i\to k_i\to i$. No unordered pair $\{j_i,k_i\}$ is reused, and every reservoir vertex occurs as $j_i$ or $k_i$ at most $L$ times.
\item For each $i\in R$, choose a directed triangle $i\to j\to k\to i$ wholly in $R$. No edge is reused between these triangles, and every vertex belongs to at most $1001$ chosen triangles.
\end{enumerate}
\end{srcBFlemma}
\begin{proof}
Include each quotient vertex in $R$ independently with probability $\rho=p^{-1/4}$. Its expected size is $m$, and each in- or out-neighborhood intersection has mean $\mu=(p-1)\rho/2\ge0.49m$. The elementary Chernoff bounds
\[
 \Pr(X\ge2\mathbb EX)\le e^{-\mathbb EX/3},\qquad
 \Pr(X\le(1-\delta)\mathbb EX)\le e^{-\delta^2\mathbb EX/2}
\]
for sums of independent indicators follow from exponential Markov and
$\mathbb Ee^{tX}\le\exp((e^t-1)\mathbb EX)$, optimizing $t$ and using the elementary logarithmic bounds. With $\delta=1/4$, the threshold $m/3$ is below $(3/4)\mu$, so a union bound gives failure probability at most
\[
 e^{-m/3}+2p e^{-m/100}<1.
\]
For our cutoff the last strict inequality follows, for instance, from $m/100>\log(8p)$ and $m/3>\log4$; use $\log(8p)\le\sqrt{8p}$ and $p^{1/4}\ge10^6$. Fix a successful reservoir.

For any $i$, set $A=N^+(i)\cap R$, $B=N^-(i)\cap R$. They are disjoint, each of size at least $m/3$. Since $m\ge6\sqrt p$, \eqref{R6E2:eq:mixing} gives at least $|A||B|/4\ge m^2/36$ pairs $(j,k)$ with $j\to k$. These are candidate directed triangles through $i$.

For the first schedule, process all $i\notin R$ greedily. Before any step at most $2p$ reservoir incidences have been used. Vertices already of load $L$ number at most $2p/L\le m/500$. Discard candidates meeting them, at most $|R|m/500\le m^2/250$ pairs. Discard the at most $p$ previously used reservoir pairs. A candidate remains because
\begin{equation}\label{R6E2:eq:firstbudget}
 m^2/36-m^2/250-p>0.
\end{equation}
Choose one and increase the two loads. The load cap and unused-pair condition follow.

For the second schedule, process every $i\in R$ once as its prescribed vertex; write $M=1000$. Auxiliary vertices are chosen only while their current triangle load is below $M$. A vertex's single prescribed turn can increase its load by one beyond that cap, so all loads are at most $M+1$. Total load before any step is at most $3|R|\le6m$, hence at most $6m/M$ vertices currently have load at least $M$. Discarding pairs meeting them costs at most $12m^2/M=3m^2/250$. At most $2(M+1)=2002$ previously used edges touch the prescribed vertex $i$, so excluding pairs using one of those edges costs at most $2002|R|\le4004m$. Finally exclude candidates whose edge $jk$ was previously used, at most $3|R|\le6m$ pairs. There is still a choice, since
\begin{equation}\label{R6E2:eq:secondbudget}
 m^2(1/36-3/250)-4010m>0.
\end{equation}
Both budget inequalities follow directly from $p\ge10^{24}$. The auxiliary cap and the distinct-edge rule give all assertions. The two schedules need not have disjoint quotient edge sets; later they use different vertex roles.
\end{proof}

\section{Regularity-preserving surgery with many failed roots}
Start with the regular substitution of a cyclic five-tournament in each Paley vertex. Label a fiber
\[
 F_i=\{(i,0),(i,1),(i,2),(i,3),(i,4)\},
\]
and orient $(i,a)\to(i,b)$ when $b-a\pmod5$ is $1$ or $2$. Between fibers use Paley directions. The initial order is $5p$ and every degree is $(5p-1)/2$.

For each $i\notin R$, use the first schedule and reverse all arcs of the following two directed four-cycles:
\begin{equation}\label{R6E2:eq:gadget}
 ((i,0),(j_i,1),(k_i,0),(i,3)),\qquad
 ((i,0),(j_i,2),(k_i,0),(i,4)).
\end{equation}
A listed cyclic order means that the last vertex also points to the first before reversal. This holds because $i\to j_i\to k_i\to i$, and in the initial five-cycle both $(i,3),(i,4)$ point to $(i,0)$. The two four-cycles share vertices but not edges. Across different $i$, edges between a nonreservoir and a reservoir fiber cannot coincide, and edges between reservoir fibers cannot coincide because the pairs $\{j_i,k_i\}$ are distinct. Thus every four-cycle is still directed when reversed, in any order.

For each triangle $i\to j\to k\to i$ in the second schedule, reverse the directed triangle on
\begin{equation}\label{R6E2:eq:break}
 (i,4),(j,4),(k,4).
\end{equation}
These edges are distinct and were not touched in \eqref{R6E2:eq:gadget}: the earlier reservoir-to-reservoir edges had roles $1$ or $2$ at one end and role $0$ at the other. Reversing a directed cycle changes no degree. The final tournament $T_p$ is therefore regular.

\begin{srcBFlemma}[Exact Hall neighborhoods]\label{R6E2:lem:failures}
Every $v_i=(i,0)$ with $i\notin R$ fails, with Hall set $S_i=F_i\setminus\{v_i\}$ satisfying
\begin{equation}\label{R6E2:eq:exacttargets}
 N_{B_{v_i}}(S_i)=\{(j_i,1),(j_i,2),(k_i,0)\}.
\end{equation}
Consequently $b(T_p)\ge p-|R|\ge p-2p^{3/4}$.
\end{srcBFlemma}
\begin{proof}
The internal reversals make $v_i$ dominate all four vertices of $S_i$. Its only new external inneighbors are $(j_i,1),(j_i,2)$, both still dominated by all of $S_i$. Its old external inneighbors were all in quotient inneighbor fibers. Among them, only $(k_i,0)$ becomes beaten by a member of $S_i$: it becomes beaten by $(i,3),(i,4)$. No other gadget or later triangle touches $F_i$, since their auxiliary fibers lie in $R$. This proves the exact equality \eqref{R6E2:eq:exacttargets}, not merely an upper bound, and the Hall deficiency is $4-3=1$.
\end{proof}

\begin{srcBFlemma}[Spectral and edit bounds]\label{R6E2:lem:norm}
Every skew row differs from its initial row in at most
\[
 D=4L+2002
\]
positions, and
\[
 \srcBFnorm{K_{T_p}}\le5\sqrt p+8L+4004\le5\sqrt p+8001p^{1/4}.
\]
In particular $\srcBFnorm{K_{T_p}}^2<26p$.
\end{srcBFlemma}
\begin{proof}
The initial matrix is $P\otimes J_5+I_p\otimes C$, where $C$ is the cyclic five-tournament skew matrix. On the fiber-constant space the norm is $5\sqrt p$; on its orthogonal complement it is $\srcBFnorm C\le4$. Thus its total norm is $5\sqrt p$.

A nonreservoir row is changed at most four times. Each appearance of a reservoir vertex in the first schedule changes at most four incident edges of any one row of its fiber, hence at most $4L$ total. The second schedule changes at most $2(1001)=2002$ incident edges. Since no edge is reversed twice, the total skew perturbation has maximum absolute row and column sum at most $2D$. Its operator norm is at most $2D$, proving the first norm bound.

Put $x=p^{1/4}\ge10^6$. Since $L\le1000x+1$ and $4012\le x$, the extra term is at most $8001x$. Also $5+8001/x<5.01$, so its squared norm is less than $5.01^2p<26p$.
\end{proof}

\section{Why the final tournaments really are prime}
\begin{srcBFlemma}[Modules in a regular tournament]\label{R6E2:lem:module}
A proper homogeneous set $H$ with $2\le h:=|H|<n$ in a regular tournament has odd size, induces a regular tournament, and satisfies
\[
 h\le(n-1)/2,\qquad h(n-h)\le\srcBFnorm K^2.
\]
\end{srcBFlemma}
\begin{proof}
All its vertices have identical external outdegree; regularity forces a constant internal outdegree, necessarily $(h-1)/2$. This proves oddness and internal regularity. An outside vertex either has at least $h$ outneighbors or at least $h$ inneighbors, giving $h\le(n-1)/2$. The vector $K\mathbf1_H$ is zero on $H$ and equals $\pm h$ off $H$, so its squared norm divided by $\srcBFnorm{\mathbf1_H}^2=h$ is $h(n-h)$. This is bounded by $\srcBFnorm K^2$.
\end{proof}

\begin{srcBFlemma}[Primeness of the construction]\label{R6E2:lem:prime}
The tournament $T_p$ has no nontrivial homogeneous set.
\end{srcBFlemma}
\begin{proof}
Suppose $H$ were one. By Lemmas~\ref{R6E2:lem:norm}--\ref{R6E2:lem:module}, its odd size is at most five. Indeed if $7\le h\le(5p-1)/2$, monotonicity of $h(5p-h)$ on that range gives
\[
 \srcBFnorm K^2\ge7(5p-7)=35p-49>26p,
\]
a contradiction.

No such $H$ can meet two different fibers. For distinct quotient vertices $i,j$, \eqref{R6E2:eq:paley} gives squared distance $2p$ between their Paley rows. Their two diagonal positions contribute two; hence they disagree in exactly $(p-1)/2$ third-fiber positions. Two initial full rows from $F_i,F_j$ therefore disagree on $5(p-1)/2$ positions in those third fibers. After surgery at most $2D$ of these disagreements can disappear. With $x=p^{1/4}\ge10^6$ one has $D\le4001x$, so
\[
 5(p-1)/2-2D\ge(5/2)(x^4-1)-8002x>5.
\]
But rows belonging to a homogeneous set of size at most five must agree at every position outside that set. This contradiction puts $H$ inside one fiber.

Its possible sizes are now three or five. In a reservoir fiber the induced tournament remains the regular cyclic five-tournament, which has no homogeneous regular three-set: such a set would violate the first bound of Lemma~\ref{R6E2:lem:module} applied to that five-tournament. In a nonreservoir fiber, vertex $0$ is a source, so a regular three-set cannot contain it. Among $\{1,2,3,4\}$ the only directed triangles are $\{1,2,4\}$ and $\{1,3,4\}$. The remaining vertex $3$, respectively $2$, sees those triples nonuniformly. Thus neither is internally homogeneous. This excludes size three in every fiber.

Finally no entire five-fiber is homogeneous. For $i\notin R$, the outside vertex $(k_i,0)$ is beaten by $(i,3),(i,4)$ but still beats $(i,0),(i,1),(i,2)$. For $i\in R$, its role-$4$ vertex participates in at least one triangle reversal \eqref{R6E2:eq:break}. At the opposite endpoint of such a reversed edge, the orientation toward $(i,4)$ differs from its orientation toward $(i,3)$, which was untouched on every reservoir-to-reservoir edge. No edge was reversed twice. Hence this fiber too is seen nonuniformly by an outside vertex. This excludes the last possibility and proves primeness.
\end{proof}

\begin{proof}[Proof of Theorem~\ref{R6E2:thm:construction}]
Regularity, the norm estimate, and primeness were proved above. Since $26p<30p-116=6n-116$, Theorem~\ref{R6E2:thm:count} applies and gives $b(T_p)\le p$, with every failed root of deficiency one. Lemma~\ref{R6E2:lem:failures} gives $b(T_p)\ge p-2p^{3/4}$. Subtracting from $5p$ proves both count bounds.

There is at least one failed root, since $p-2p^{3/4}>0$. Corollary~\ref{R6E2:cor:onset} therefore gives $\srcBFnorm K^2\ge5n-18$, which together with \eqref{R6E2:eq:construction} proves the norm-ratio limit. The strong-count limit follows from the two count bounds. There are infinitely many primes congruent to three modulo four by the elementary product-minus-one argument, so this yields an infinite sequence. The reservoir existence argument and the greedy schedules specify a finite construction; they make no claim to a small practical cutoff or to exact equality $|\operatorname{Str}|=4p$.
\end{proof}

\section{Verification, failed routes, and remaining questions}
The verifier independently builds Paley fibers and all directed-cycle reversals at four smaller prime orders. It checks every degree and pair orientation, every designated Hall neighborhood, all five-witness row distances, and the absence of a homogeneous set through an exact distinguishing-digraph test. It computes every cyclic-link maximum matching for the first two examples and a stated sample for the larger two. It separately checks $77{,}815$ Hall parameter cases and exact polynomial certificates for the large-cutoff budgets. The finite examples do not satisfy the deliberately huge sufficient cutoff; their directly checked properties are not presented as applications of that cutoff theorem.

For clarity, the exact prime test is as follows. For each anchor $x$, form a directed graph on $V\setminus\{x\}$, with $y\to z$ when $z$ distinguishes $x$ and $y$. A proper homogeneous set containing $x$ and another vertex is exactly a proper nonempty set closed under these implications, with $x$ adjoined. Thus checking strong connectivity of this graph for every anchor excludes every nontrivial module. Forward and reverse reachability are checked using integer bitsets. This finite certificate is independent of the asymptotic spectral argument proving primeness.

\paragraph{What failed and what changed.}
Simply removing the old homogeneous triples cannot preserve the first spectral scale: the inherited onset bound forbids it. A first five-fiber construction used disjoint quotient triangles as reservoirs, giving only a $1/15$ failure density. Reusing a \emph{sublinear} reservoir with bounded load raises the failure density to $1/5$ while changing the skew norm by only $O(p^{1/4})$. Letting every fiber simultaneously act as a root and a repair reservoir introduced extra Hall targets and destroyed the four-against-three certificate; the separated roles in \eqref{R6E2:eq:gadget} avoid that error. Finally, leaving even one fiber untouched would invalidate primeness, which is why the second schedule is part of the proof, not cosmetic cleanup.

The constant $116$, the cutoff $10^{24}$, and the spectral window coefficient $6$ are not claimed optimal. Exact finite attainment of $\lceil4n/5\rceil$, optimal density at larger norm ratios, and unrestricted prime regular strong existence remain open in this paper. The result proves that primeness alone does not make \emph{every} root strong, even at a controlled pseudorandom spectral scale; it does not produce a prime tournament with no strong root. No unrestricted Nine-Mass consequence or prize-level claim is made. Public source versions were checked on 20 October 2026; historical priority remains unverified. No external posting, submission, or author contact occurred.



\chapter[Prime Regularization Preserving Strong Seymour Failures and Quasirandom Tournaments of Vanishing Strong Density]{Prime Regularization Preserving Strong Seymour Failures and Quasirandom Tournaments of Vanishing Strong Density}
\label{source:R7E2}
\noindent\textit{Source development: Round 7, entry 02. Original mathematical arguments retained; notation is local to this chapter.}\par\medskip

\begin{quote}\small\noindent\textbf{Chapter abstract.} 
We give a degree-preserving, explicitly coded surgery that makes a Paley substitution prime while losing only $O(\sqrt m)$ from each internal matching deficiency, where $m$ is the fiber order. Applied after amplification, it preserves every failed root of any regular seed tournament. Specifically, a regular $h$-vertex seed $H$ yields, for every prime $p\equiv3\pmod4$, $p\ge127$, a prime regular tournament of order $25h^2p$, skew norm at most $25h^2\sqrt p+40h+4$, and strong proportion at most $|\srcBLStr(H)|/h$. Thus there are prime regular quasirandom tournaments with strong proportion tending to zero; almost all their roots have matching deficiency tending to infinity. The earlier constant-eleven ladder is an attributed input, not a new extremal count. The new ingredient is simultaneous destruction of every module without repairing any amplified failure. A conditional reduction to prime regular quasirandom counterexamples does not settle strong-vertex existence or the public nine-vertex lower-bound problem.
\end{quote}

\section{Statements, definitions, and the increment}
All tournaments here are finite. A \emph{module} is a vertex set $M$ such that each vertex outside $M$ either dominates all of $M$ or is dominated by all of $M$. A tournament is \emph{prime} if it has no module with $1<|M|<n$. For a regular tournament of order $n=2d+1$, the exact second out-neighborhood is
\begin{equation}\label{R7E2:eq:second}N^{++}(v)=N^-(v).
\end{equation}
Indeed an inneighbor not reached from $N^+(v)$ would dominate $v$ and its $d$ outneighbors. Write $B_v$ for the bipartite arc graph from $N^+(v)$ to $N^-(v)$, and put
$\srcBLdefi_T(v)=d-\nu(B_v)$, where $\nu$ is maximum matching size. A root is \emph{strong} precisely when this deficiency is zero. Hall's deficiency formula is
\begin{equation}\label{R7E2:eq:hall}
\srcBLdefi_T(v)=\max_{S\subseteq N^+(v)}\bigl(|S|-|N_{B_v}(S)|\bigr).
\end{equation}
For completeness, add $t$ universal target vertices to a balanced bipartite graph, where $t$ is this maximum. Hall's condition then holds, giving a matching that leaves at most $t$ original sources unmatched; the reverse inequality follows from every source subset. This proves the formula.

The skew adjacency matrix $K_T$ has entries $1$ on arcs, $-1$ on their reverses, and zero diagonal; its Euclidean operator norm is denoted $\srcBLnorm{K_T}$.

\begin{srcBLtheorem}[Prime regularization of an arbitrary seed]\label{R7E2:thm:main}
Let $H$ be any regular tournament of odd order $h\ge1$. For every prime $p\equiv3\pmod4$, $p\ge127$, there is an explicitly specified prime regular tournament $\mathcal R_p(H)$ on vertices $(i,v,a)$, where $i\in\mathbb F_p$, $v\in V(H)$, and $0\le a<25h$, such that
\begin{align}
 n&=25h^2p, \qquad \srcBLnorm{K_{\mathcal R_p(H)}}\le25h^2\sqrt p+40h+4,\\
 \srcBLdefi_{\mathcal R_p(H)}(i,v,a)&\ge25h\srcBLdefi_H(v)-20h-1.\label{R7E2:eq:main_def}
\end{align}
In particular every vertex above a failed seed root is failed, and
\begin{equation}\label{R7E2:eq:density}
 \frac{|\srcBLStr(\mathcal R_p(H))|}{25h^2p}\le\frac{|\srcBLStr(H)|}{h}.
\end{equation}
There is no assertion of equality in this count or preservation of strong roots.
\end{srcBLtheorem}

\begin{srcBLcorollary}[Prime, quasirandom, and almost everywhere highly deficient]\label{R7E2:cor:sparse}
For every odd $h\ge13$ and every Paley prime $p\ge127$, there is a prime regular tournament $T_{h,p}$ of order $25h^2p$ satisfying
\[
 \srcBLnorm{K_{T_{h,p}}}\le25h^2\sqrt p+40h+4,
 \qquad |\srcBLStr(T_{h,p})|\le275hp.
\]
At least a proportion $1-11/h$ of its roots have deficiency at least $5h-1$. Consequently, along any sequence $h\to\infty$, $p\to\infty$ with these conditions, the strong proportion tends to zero, the displayed deficiency lower bound tends to infinity for almost all roots, and $\srcBLnorm K/n\to0$. Every fixed labelled tournament density converges to its random value.
\end{srcBLcorollary}

\begin{srcBLcorollary}[A conditional existence reduction]\label{R7E2:cor:reduction}
If any regular tournament has no strong root, then there is a sequence of prime regular tournaments with no strong root and $\srcBLnorm K/n\to0$. Therefore, to prove existence of a strong root in every regular tournament, it suffices to rule out such prime regular quasirandom counterexample sequences.
\end{srcBLcorollary}
This is an implication, not a proof that either side has counterexamples or that regular strong existence holds. It does not transfer a fixed absolute lower bound such as nine, since the construction increases order and may replicate strong roots.

\paragraph{Prior results and current status.}
Round 3 entry 02 \cite{R7E2:R3} proved vanishing strong density in regular quasirandom substitutions, whose large homogeneous fibers remained essential. Round 6 entry 02 \cite{R7E2:R6} constructed prime regular tournaments with asymptotic strong proportion $4/5$ at the sharp leading spectral scale $\sqrt{5n}$ for its restricted window. Neither result supplied vanishing density in the prime class. Here the new coded surgery applies to \emph{any} regular fiber, preserves all sufficiently large deficiencies, and removes every module. We reuse and prove the elementary Paley spectrum, substitution separation, the module norm inequality, and the required constant-eleven ladder obstruction. These inputs are not additional discoveries.

Bai--Li--Park v2 \cite{R7E2:BLP} reports counterexamples to the universal strong conjectures; this correction to the original Round 1 framing remains in force. Ordinary Seymour is a different statement and holds at every regular root by \eqref{R7E2:eq:second}. The public notes \cite{R7E2:G}, dated 23 July 2026, label the regular nine-vertex lower bound open and their work unreviewed. Their ladder/substitution viewpoint and constant-nine construction predate this experiment. We use the independently certified eleven-root family from Round 3 instead, without claiming eleven is best. The source versions were checked again on 20 October 2026. Targeted primary-source searches do not certify historical priority for the new regularization theorem.

\section{Two elementary Paley facts}
Let $P_p$ be the tournament on $\mathbb F_p$ with $i\to j$ when $j-i$ is a nonzero square. Write $\chi$ for the quadratic character, with $\chi(0)=0$. Since $p\equiv3\pmod4$, this is a regular tournament and its skew matrix is $K_{ij}=\chi(j-i)$.
\begin{srcBLlemma}\label{R7E2:lem:paley}
The matrix of $P_p$ satisfies $K\mathbf1=0$ and $KK^{\mathsf T}=pI-J$. Distinct rows differ in exactly $(p-1)/2$ positions outside their two diagonal positions. For $p\ge127$ there exists $b$ such that
\[
 0\to1\to b\to0,
 \qquad \{\pm1,\pm b,\pm(b-1)\}\text{ consists of six distinct elements}.
\]
The $p$ translated directed triangles $(i,i+1,i+b)$ have pairwise disjoint unordered edge sets. Each quotient vertex occurs once in each of their three roles.
\end{srcBLlemma}
\begin{proof}
For $t\ne0$, $\sum_x\chi(x)\chi(x+t)=-1$. To see this after scaling, count solutions of $y^2=x(x+1)$: the equation $(2x+1-2y)(2x+1+2y)=1$ has $p-1$ solutions, while the character count is $p+\sum_x\chi(x(x+1))$. This gives the matrix identity and row separation. The norm is consequently $\sqrt p$.

The number of $b$ satisfying $\chi(b)=-1$, $\chi(b-1)=1$ is $(p+1)/4$. Indeed summing $(1-\chi(b))(1+\chi(b-1))/4$ over all $b$ gives $(p+1)/4$, using the character correlation; the summands at $b=0,1$ are zero. Apart from the already excluded $0,1$, a repeated signed difference can occur only for $b=-1,2,1/2$. At most three candidates are excluded, so a valid $b$ exists. Fix, for definiteness, the least representative in $\{0,\ldots,p-1\}$ satisfying these conditions.

The three unordered edge types have differences $\pm1,\pm b,\pm(b-1)$, respectively, and so cannot coincide. Within one type the translations are distinct, since $p$ is odd. Translations also prove the role counts.
\end{proof}

For regular tournaments $Q,F$, the substitution $Q[F]$ replaces every quotient vertex by a copy of $F$, with all between-fiber arcs following $Q$.
\begin{srcBLlemma}[Attributed substitution identities]\label{R7E2:lem:sub}
If $Q,F$ have odd orders $q,m$, their substitution is regular. At root $(u,v)$ its link is the disjoint union of the internal link at $v$ and the $m$-fold bipartite blowup of the link at $u$. Hence
\[
 \srcBLdefi_{Q[F]}(u,v)=m\srcBLdefi_Q(u)+\srcBLdefi_F(v).
\]
For a regular $m$-vertex $F$, $\srcBLnorm{K_{P_p[F]}}=m\sqrt p$.
\end{srcBLlemma}
\begin{proof}
Degrees are $m(q-1)/2+(m-1)/2$. An internal source cannot beat an external target, and an external source cannot beat an internal target, so the two link parts have no cross arcs. A bipartite $m$-fold blowup has maximum matching size $m$ times that of the original graph: lifting a matching proves one inequality, and a maximum-deficiency Hall set lifted to all its copies proves the other by \eqref{R7E2:eq:hall}. This proves the deficiency identity.

The skew matrix is $K_{P_p}\otimes J_m+I_p\otimes K_F$. The subspaces of vectors constant on each fiber and vectors summing to zero on each fiber are invariant and orthogonal. The two norms are $m\sqrt p$ and $\srcBLnorm{K_F}\le m-1$, respectively. Regularity gives $K_F\mathbf1=0$ and justifies this decomposition.
\end{proof}
The exact link separation was proved in Round 3 entry 02, with the public substitution viewpoint credited there; it is repeated rather than newly claimed.

\section{Coded triangle surgery on arbitrary regular fibers}
Fix an odd $m\ge25$, a regular tournament $F$ on $\{0,\ldots,m-1\}$, a Paley prime $p\ge127$, and $b$ as in Lemma~\ref{R7E2:lem:paley}. Put
\[
 r=\lceil\sqrt m\rceil,\quad A=\{0,\ldots,r-1\},\quad
 a(x)=\lfloor x/r\rfloor,\quad c(x)=x\bmod r.
\]
Both $a(x),c(x)$ belong to $A$, and $x\mapsto(a(x),c(x))$ is injective. Start with $P_p[F]$, writing its vertices as $(i,x)$. For every $i$ and every $x$, reverse the directed triangle
\begin{equation}\label{R7E2:eq:switch}
 (i,x)\ \longrightarrow\ (i+1,a(x))\ \longrightarrow\ (i+b,c(x))\ \longrightarrow\ (i,x).
\end{equation}
Call the result $\mathcal C_p(F)$.

\begin{srcBLlemma}[Regularity and perturbation size]\label{R7E2:lem:edit}
All triangles in \eqref{R7E2:eq:switch} have disjoint edge sets. The resulting tournament is regular, has no within-fiber edits, and its maximum incident edit degree is at most $D=4r+2$. Thus
\begin{equation}\label{R7E2:eq:norm}
 \srcBLnorm{K_{\mathcal C_p(F)}}\le m\sqrt p+8r+4.
\end{equation}
\end{srcBLlemma}
\begin{proof}
Different quotient triangles have disjoint edge sets by Lemma~\ref{R7E2:lem:paley}. Within one quotient triangle, primary-to-auxiliary edges have different primary endpoint $x$; auxiliary-to-auxiliary edges are distinct by injectivity of $(a(x),c(x))$. Thus every reversed triangle is directed in the initial tournament and no edge is reversed twice. A directed triangle reversal preserves every degree.

Each vertex is primary once. It can be the middle auxiliary vertex in at most $r$ triangles, because each value of $a(x)$ occurs at most $r$ times, and the last auxiliary vertex in at most $r$ triangles for the same reason with $c(x)$. Each occurrence changes two incident edges, proving $D$. A skew perturbation has absolute row and column sums at most $2D$, so its Euclidean norm is at most $2D$ (by the square root of the product of those two bounds). Lemma~\ref{R7E2:lem:sub} gives \eqref{R7E2:eq:norm}.
\end{proof}

\begin{srcBLlemma}[All modules are destroyed]\label{R7E2:lem:prime}
The tournament $\mathcal C_p(F)$ is prime.
\end{srcBLlemma}
\begin{proof}
We first bound the size of any proper nontrivial module $M$ in a regular $n$-vertex tournament. All its vertices have the same outside outdegree, so its induced tournament is regular; in particular its size $t$ is odd. Summing the regular degree equations over $M$ shows that exactly half the outside vertices are dominated by $M$. Both outside orientations therefore occur, and an outside vertex incident to all $t$ module vertices in one direction gives $t\le(n-1)/2$. Moreover $K\mathbf1_M$ is zero on $M$ and equals $\pm t$ outside, so
\begin{equation}\label{R7E2:eq:module}
 t(n-t)\le\srcBLnorm K^2,\qquad t<\frac{2\srcBLnorm K^2}{n}.
\end{equation}
The inequality and degree facts were used in Round 6 entry 02 and are reproved here.

For odd $m\ge25$, $4r+2\le m$. For $r=5$ this is immediate; for $r\ge6$, $m\ge(r-1)^2+1\ge4r+2$. Thus \eqref{R7E2:eq:norm} gives $\srcBLnorm K\le m(\sqrt p+2)$. Since $p\ge127$, $8/\sqrt p+8/p<1$, and \eqref{R7E2:eq:module} implies
\begin{equation}\label{R7E2:eq:smallmodule}t<3m.
\end{equation}

Two vertices in different fibers initially have different orientations to exactly $m(p-1)/2$ vertices in third fibers, by Paley row separation. Edits at either endpoint remove at most $2D\le2m$ of these distinguishing vertices. They still have at least $m(p-5)/2>3m$ distinguishing vertices. If they belonged to one module of size less than $3m$, every distinguishing vertex would have to lie in that module, a contradiction. Hence any proper nontrivial module lies in a single fiber.

Finally take distinct $(i,x),(i,y)$ in one fiber. If $a(x)\ne a(y)$, the external vertex $(i+1,a(x))$ has its orientation to $(i,x)$ reversed and its orientation to $(i,y)$ unchanged. If $a(x)=a(y)$, then $c(x)\ne c(y)$, and $(i+b,c(x))$ distinguishes them instead. No other gadget can change these edges because quotient edge sets are disjoint. Thus every pair in a fiber has an external distinguisher, excluding every nontrivial subset of that fiber as a module. This also excludes the entire fiber.
\end{proof}

\begin{srcBLlemma}[Deficiency loss, including auxiliary roots]\label{R7E2:lem:loss}
For every root $(i,v)$, including indices $v\in A$ used in auxiliary roles,
\begin{equation}\label{R7E2:eq:loss}
 \srcBLdefi_{\mathcal C_p(F)}(i,v)\ge\srcBLdefi_F(v)-4r-1.
\end{equation}
\end{srcBLlemma}
\begin{proof}
Choose a maximum-deficiency Hall source set $S$ in the internal link of $v$ in $F$, and let $N$ be its internal target neighborhood. Delete the at most $r$ vertices whose indices belong to $A$, obtaining $S'=S\setminus A$. Internal arcs are unaltered, so $S'$ remains a subset of the root's outshore and its internal target neighborhood is contained in $N$.

Every vertex of $S'$ participates only in its \emph{own primary} triangle: all auxiliary roles use indices in $A$. Its only gained outgoing edge is to a vertex in
$C_i=\{(i+b,z):0\le z<r\}$. Its other changed edge is a \emph{lost} outgoing edge and cannot create a target. Consequently, among vertices that were already in the root's external inshore before surgery, new targets from $S'$ lie in $C_i$, at most $r$ vertices. Before surgery there were no external targets from an internal source to an external inneighbor.

The root itself participates in at most $1+2r$ directed triangles, even when it has an auxiliary role. Each reversal changes exactly one of its outgoing edges to incoming. Therefore at most $2r+1$ vertices enter its inshore. Count all these as potential additional targets, regardless of their orientation to $S'$. No other source of new targets exists. Thus the new Hall neighborhood of $S'$ has size at most $|N|+r+(2r+1)$, and
\[
 |S'|-|N_{B_{(i,v)}}(S')|\ge |S|-|N|-r-r-(2r+1)
 =\srcBLdefi_F(v)-4r-1.
\]
Hall's formula proves the claim, including when its right side is nonpositive.
\end{proof}

\section{Amplification and the prime reduction}
Let $C_k$ denote the cyclic regular tournament of odd order $k=2e+1$, with arcs given by differences $1,\ldots,e$. Every root is strong: at zero use the matching $a\mapsto a+e$, $1\le a\le e$, and translate it.

\begin{proof}[Proof of Theorem~\ref{R7E2:thm:main}]
Take $k=25h$ and $F=H[C_k]$, labelling its vertices by $x=kv+a$. It has order $m=25h^2\ge25$ and $r=5h$. Lemma~\ref{R7E2:lem:sub} gives the exact amplification
$\srcBLdefi_F(v,a)=25h\srcBLdefi_H(v)$.
Apply the coded construction and Lemmas~\ref{R7E2:lem:edit}--\ref{R7E2:lem:loss}. The degree, prime and norm claims follow, as does \eqref{R7E2:eq:main_def}. If a seed root fails, its deficiency is at least one, so the new deficiency is at least $5h-1>0$. Only the $25hp$ vertices above each originally strong root can possibly be strong, proving \eqref{R7E2:eq:density}.
\end{proof}

\begin{proof}[Proof of Corollary~\ref{R7E2:cor:reduction}]
Fix a hypothetical regular seed with no strong root, and let $p\equiv3\pmod4$ tend to infinity through primes. Every root of $\mathcal R_p(H)$ fails by Theorem~\ref{R7E2:thm:main}, while
\begin{equation}\label{R7E2:eq:ratio}
 \frac{\srcBLnorm K}{n}\le\frac1{\sqrt p}+\frac{40h+4}{25h^2p}\longrightarrow0.
\end{equation}
There are infinitely many such primes: if there were finitely many, four times their product minus one would have a new prime divisor congruent to three modulo four. No counterexample seed is supplied or asserted to exist.
\end{proof}
In particular, the infimum of the strong proportion over prime regular tournaments equals that over all regular tournaments: one inequality is inclusion, and the opposite follows by applying \eqref{R7E2:eq:density} to each regular seed. The following attributed explicit seeds show that both infima are zero.

\section{A self-contained sparse seed and quasirandomness}
We reproduce only the part of the Round 3 constant-eleven family needed here: regularity and failure of all new ladder roots. The exact eleven count and explicit matchings for the old roots were established there, but the upper bound below does not require those matchings. The two-vertex ladder operation itself has the public antecedent \cite{R7E2:G}.

Start with vertices $0,1,\ldots,10$, put $U=\{4,5,6\}$, $V=\{1,2,3\}$, $W=\{7,8,9,10\}$, and orient $U,V$ cyclically in their displayed orders. Orient
\[
 U\to0\to V\to W\to U,\qquad U\to V.
\]
On $W$ use $7\to8\to9\to10\to7$ and $7\to9$, $8\to10$. Finally set $0\to7,8$ and $9,10\to0$. Every vertex has outdegree five, as direct counting by these four parts shows.

Initially set $P_0=U\cup V$, $Q_0=\{0\}\cup W$. For $j=1,\ldots,\ell$ add $x_j,y_j$ with
\begin{equation}\label{R7E2:eq:ladder}
 P_{j-1}\to x_j\to Q_{j-1}\to y_j\to P_{j-1},\qquad x_j\to y_j,
\end{equation}
then put $P_j=P_{j-1}\cup\{x_j\}$, $Q_j=Q_{j-1}\cup\{y_j\}$. Every old vertex gains one incoming and one outgoing edge; each new root has outdegree $5+j$. The result $H_h$ is regular of order $h=11+2\ell$.

For $1\le j\le\ell$, direct application of \eqref{R7E2:eq:ladder} gives the following exact Hall neighborhoods in the relevant root's inshore:
\[
\begin{array}{c|c|c}
\text{root}&\text{source set}&\text{target neighborhood}\\\hline
x_j&W\cup\{x_t:t>j\}&U\cup\{y_t:t>j\}\\
y_j&U\cup\{y_t:t>j\}&\{0\}\cup\{x_t:t\ge j\}.
\end{array}
\]
For example, for $x_j$ the old inshore is $P_{j-1}$; $W$ beats precisely $U$ there, later $x_t$ beat no additional earlier $P$ vertices, and the later inshore vertices $y_t$ are all reached from $W$. For $y_j$, the old inshore is $Q_{j-1}\cup\{x_j\}$; $U$ reaches $0,x_j$, later $y_t$ add no earlier $Q$ targets, and all later inshore $x_t$ are reached from $U$. These statements account separately for seed, earlier and later ladder vertices. The source size exceeds the target size by one in both rows. Hence all $2\ell$ new roots fail and $|\srcBLStr(H_h)|\le11$ for every odd $h\ge13$.

\begin{proof}[Proof of Corollary~\ref{R7E2:cor:sparse}]
Use $T_{h,p}=\mathcal R_p(H_h)$. The count, deficiency and norm claims follow from Theorem~\ref{R7E2:thm:main}. Formula~\eqref{R7E2:eq:ratio} tends to zero whenever $p\to\infty$, uniformly over $h\ge1$; $h\to\infty$ is needed separately for the strong density to vanish and the deficiency lower bound to diverge.

For clarity, spectral quasirandomness implies the fixed-pattern statement directly. If $A$ is the zero-one adjacency matrix of an $n$-vertex tournament, $A=(J-I+K)/2$, so $\srcBLnorm{A-J/2}\le(\srcBLnorm K+1)/2$. For a fixed labelled tournament $L$ on $s$ vertices, telescope its adjacency-product count against the constant product $2^{-\binom s2}$. Fix all coordinates except the endpoints of the edge being replaced; the remaining factors separate as $c f(i)g(j)$ with $0\le c,f,g\le1$. The resulting bilinear form has absolute value at most $(\srcBLnorm K+1)n/2$. Sum over the $n^{s-2}$ other coordinates and all $\binom s2$ replacements to obtain
\[
 \left|N_L(T)-2^{-\binom s2}n^s\right|
 \le\binom s2\frac{\srcBLnorm K+1}{2}n^{s-1}.
\]
Repeated-coordinate tuples have a zero adjacency factor, so $N_L$ counts injective maps. Dividing by $n^s$ proves the claim. This standard counting argument, already supplied in Round 3, is not a new contribution here.
\end{proof}

\section{Scope, verification, and remaining questions}
The construction proves vanishing \emph{density}, not an arbitrarily slowly growing absolute strong count. Its bound is $275hp$ at order $25h^2p$; unlike the earlier nonprime family it does not assert domination by every prescribed divergent function. Nor does it give a fixed $O(\sqrt n)$ spectral constant along the vanishing-density sequence: the principal ratio $\srcBLnorm K/\sqrt n$ is of order $5h$. There is therefore no conflict with the prior sharp density windows at bounded spectral constants.

The essential distinction from an ordinary substitution is that the final tournament has no nontrivial module of any size. The quantitative all-root loss bound handles auxiliary roots as well as primary-only roots, and the amplification preserves every failed seed root. It need not preserve a strong root, an exact deficiency, or an exact number of strong roots. The conditional existence reduction leaves the existence question entirely open. No external posting, submission or author contact occurred.

The exact verifier checks translated triangle schedules, edge disjointness, pair codes and degree budgets; independently tests all modules on a practical finite construction; checks the seed regularity and exact ladder Hall neighborhoods; and verifies the amplified Hall certificates using an exact compressed representation of the much larger constructions. Finite diagnostics are not the proof of primeness or the infinite theorem. The full argument above covers all stated orders, including the explicit cutoff $p\ge127$. Attempts to assign distinct private auxiliary vertices to each fiber vertex incurred a linear target cost and failed to preserve the certificates; sharing a square-root-sized pair code and deleting all auxiliary-role sources removes that obstruction. Questions left open include sharper regularization cost, absolute sparse strong counts in prime tournaments, and universal lower bounds for regular strong sets.



\part{Special functions and Gaussian columns}

\chapter[A unique-minimum theorem for a logarithmic gamma quotient]{A unique-minimum theorem for a logarithmic gamma quotient}
\label{source:R1E4}
\noindent\textit{Source development: Round 1, entry 04. Original mathematical arguments retained; notation is local to this chapter.}\par\medskip
\begin{quote}\small\textbf{Integrated source-status correction.}
The fixed-parameter $a=6$ assertion has an earlier public proof claim,
as documented in the following parameter-classification chapter.
This chapter is retained for its complete argument; it is not counted
as a first solution of a currently open problem.\end{quote}


\begin{quote}\small\noindent\textbf{Chapter abstract.} 
This note gives an analytic proof of the unique-minimum assertion posed as an open problem by Bulboac\u{a} and Zayed in January 2026 for the continuous extension of
$\log\Gamma(x+1)/\log((x^2+6)/(x+6))$ to $(-1,\infty)$.
Its derivative is strictly negative before one point $b\in(1,2)$ and strictly positive afterwards, and the minimum is nondegenerate. A parameter theorem also classifies the quotient on $(1,\infty)$ for $0<a\leq8$. The proof uses elementary polygamma bounds and explicit polynomial-positivity certificates; no numerical root finding is needed. The left-interval argument is included in full. \textbf{Status:} a proposed, fully written proof of a narrowly scoped published open problem, not a claim of an Annals-level breakthrough. Priority beyond the sources inspected has not been established.
\end{quote}

\section{Problem, provenance, and statements}
Let $\gamma$ be Euler's constant, $\psi=\Gamma'/\Gamma$, and
$z_2=\zeta(2)=\pi^2/6$. For $a>0$ and $x>1$, put
\begin{equation}\label{R1E4:eq:family}
 F_a(x)=\frac{f(x)}{g_a(x)},\qquad
 f(x)=\log\Gamma(x+1),\qquad
 g_a(x)=\log\frac{x^2+a}{x+a}.
\end{equation}
The natural extension at $1$ is $F_a(1)=(a+1)(1-\gamma)$.

Section 3 of Bulboac\u{a}--Zayed \cite{R1E4:BZ}, published on 22 January 2026, asks for an analytic proof of the unique-minimum behavior of $F_6$ on $(-1,\infty)$. Their Theorem 1 proves decrease on $(-1,1)$; their evidence for the remaining behavior is numerical. We prove the full derivative-sign assertion below, including both removable points. The derivative-ratio strategy and the rational logarithm bound used below already occur in that work; the new step is the right-half-line positivity certificate.

\begin{srcADtheorem}\label{R1E4:thm:main}
Define
\[
 \mathcal F(x)=
 \begin{cases}
 \displaystyle\frac{\log\Gamma(x+1)}{\log((x^2+6)/(x+6))},&x\ne0,1,\\[4pt]
 6\gamma,&x=0,\\
 7(1-\gamma),&x=1.
 \end{cases}
\]
Then $\mathcal F$ is real analytic on $(-1,\infty)$. There is exactly one
$b\in(-1,\infty)$ with $\mathcal F'(b)=0$, and
\[
 1<b<2,\qquad
 \mathcal F'(x)<0\quad(-1<x<b),\qquad
 \mathcal F'(x)>0\quad(x>b),\qquad
 \mathcal F''(b)>0.
\]
In particular, $b$ is its unique global minimizer.
\end{srcADtheorem}

The argument on the right of $1$ yields a family result. The constant in the following statement is \emph{not new}: it is the sharp decreasing threshold on $(0,1)$ in Kup\'an--Sz\'asz \cite[Theorem 2]{R1E4:KS}; see also Yang--Tian \cite{R1E4:YT}. Our assertion concerns the other side of $1$.

\begin{srcADtheorem}\label{R1E4:thm:family}
Set
\[
 a_* =\frac{z_2-\gamma}{3-z_2-2\gamma}.
\]
Then $0<a_*<6$. For $0<a\leq8$ the following hold.
\begin{enumerate}
\item If $a\leq a_*$, then $F_a'(x)>0$ for every $x>1$.
\item If $a>a_*$, there is exactly one $b_a>1$ such that
$F_a'<0$ on $(1,b_a)$ and $F_a'>0$ on $(b_a,\infty)$.
Moreover, $F_a''(b_a)>0$.
\end{enumerate}
The upper restriction $8$ is a sufficient range for the proof, not an asserted optimal parameter bound.
\end{srcADtheorem}

\section{Elementary analytic bounds}
\begin{srcADlemma}\label{R1E4:lem:bounds}
For $t>0$,
\begin{align}
 \frac1t+\frac1{2t^2}&<\psi'(t)
 <\frac1t+\frac1{2t^2}+\frac1{6t^3},\label{R1E4:eq:trigamma}\\
 -\psi''(t)&<\frac1{t^2}+\frac1{t^3}+\frac1{2t^4},\label{R1E4:eq:tetragamma}\\
 \log t-\frac1{2t}-\frac1{12t^2}&<\psi(t)
 <\log t-\frac1{2t}.\label{R1E4:eq:digamma}
\end{align}
Also, $\gamma<289/500$ and $z_2<33/20$.
\end{srcADlemma}
\begin{proof}
The usual differentiated Euler product gives
\[
 \psi'(t)=\sum_{k=0}^{\infty}\frac1{(t+k)^2}
 =\int_0^{\infty}e^{-tu}\frac{u}{1-e^{-u}}\,\srcADdd u.
\]
Differentiation under this integral is justified by exponential decay and the bounded behavior of its kernel at zero. For $u>0$,
\begin{equation}\label{R1E4:eq:kernel}
 1+\frac u2<\frac{u}{1-e^{-u}}<1+\frac u2+\frac{u^2}{12}.
\end{equation}
For the first inequality, multiplication by $e^u-1$ gives the positive series
\[
 (u/2-1)e^u+1+u/2
 =\sum_{k=3}^{\infty}\frac{k-2}{2k!}u^k>0.
\]
For the second, the corresponding difference is
\[
 (1-u/2+u^2/12)e^u-(1+u/2+u^2/12)
 =\sum_{k=5}^{\infty}\frac{(k-3)(k-4)}{12k!}u^k>0.
\]
Integrating \eqref{R1E4:eq:kernel}, and then its product with $u$, proves
\eqref{R1E4:eq:trigamma}--\eqref{R1E4:eq:tetragamma}. The elementary sum--integral comparison for the digamma series gives $\psi(t)-\log t\to0$ as $t\to\infty$. Thus integrating \eqref{R1E4:eq:trigamma} from $t$ to infinity proves \eqref{R1E4:eq:digamma}.

For explicit rational constants, the convergent series
\[
 \log2=2\sum_{j=0}^{\infty}\frac{1}{(2j+1)3^{2j+1}}
 >2\sum_{j=0}^{3}\frac{1}{(2j+1)3^{2j+1}}>\frac{693}{1000}
\]
and \eqref{R1E4:eq:digamma} at $t=2$ give
\[
 1-\gamma=\psi(2)>\log2-\frac{13}{48}
 >\frac{211}{500}.
\]
Also, the integral test gives
\[
 z_2<\sum_{n=1}^{10}\frac1{n^2}+\frac1{10}
 =\frac{2095337}{1270080}<\frac{33}{20}.
\]
\end{proof}

\section{The derivative-ratio certificate on the right half-line}
In this section $a>0$ is fixed and $x\geq1$. Write
\[
 Q=(x^2+a)(x+a),\qquad D=x^2+2ax-a,\qquad
 X=Q'D-QD',\qquad Y=QD.
\]
Then $Q,D>0$, $g_a'=D/Q>0$, and $\psi(x+1)\geq\psi(2)>0$.
Define
\[
 U=\frac{f'}{g_a'}=\frac{Q\psi(x+1)}D,
 \qquad A=X\psi(x+1)+Y\psi'(x+1).
\]
Direct differentiation gives
\begin{equation}\label{R1E4:eq:Uprime}
 U'=\frac{A}{D^2},\qquad
 A'=X'\psi(x+1)+2Q'D\psi'(x+1)+QD\psi''(x+1).
\end{equation}
The cancellation $X+Y'=2Q'D$ is useful here.

\begin{srcADproposition}\label{R1E4:prop:A}
If $0<a\leq8$, then $A'(x)>0$ for every $x\geq1$.
\end{srcADproposition}
\begin{proof}
Put $t=x+1$. Applying Lemma~\ref{R1E4:lem:bounds} to \eqref{R1E4:eq:Uprime} gives
\begin{equation}\label{R1E4:eq:Aprimebound}
 A'>X'\psi(x+1)+\frac{D\,R_a(x)}{2t^4},
 \quad
 R_a(x)=2Q'(2t^3+t^2)-Q(2t^2+2t+1).
\end{equation}
All signs of the multipliers used in this bound are positive.
With $u=x-1\geq0$, exact expansion gives
\begin{align*}
 X'(u+1)={}&4u^3+(12a+12)u^2+(4a^2+16a+12)u+4a+4,\\
 R_a(u+1)={}&10u^5+(6a+86)u^4+(48a+287)u^3\\
 &+(-2a^2+143a+463)u^2+(-10a^2+189a+359)u\\
 &+(a+1)(107-13a).
\end{align*}
For $0<a\leq8$ every displayed coefficient is positive: in particular,
$143-2a>0$, $189-10a>0$, and $107-13a\geq3$.
Consequently both terms on the right of \eqref{R1E4:eq:Aprimebound} are positive.
\end{proof}

\begin{proof}[Proof of Theorem~\ref{R1E4:thm:family}]
Let $\delta=3-z_2-2\gamma$. Lemma~\ref{R1E4:lem:bounds} gives
$\delta>97/500>0$ and $z_2-\gamma>0$. Moreover,
\[
 7z_2+11\gamma<7\frac{33}{20}+11\frac{289}{500}
 =\frac{4477}{250}<18,
\]
so $a_*<6$. At $x=1$,
\begin{equation}\label{R1E4:eq:A1}
 A(1)=(a+1)^2\bigl((a+1)(z_2-1)-(2a-1)(1-\gamma)\bigr)
 =(a+1)^2\delta(a_*-a).
\end{equation}
If $a\leq a_*$, Proposition~\ref{R1E4:prop:A} shows that $U'(x)>0$ for $x>1$.
If $a>a_*$, $A(1)<0$ while $A(x)>0$ for sufficiently large $x$, because
$X$ has leading term $x^4$ and $Y,\psi,\psi'$ are positive there.
Thus there is a unique $r_a>1$ at which $U'$ changes from negative to positive.

To transfer this sign pattern to the quotient without appealing to a monotone l'H\^opital theorem, put
\[
 W=g_aU-f.
\]
This function is regular at $1$ with $W(1)=0$, and
\begin{equation}\label{R1E4:eq:W}
 W'=g_aU',\qquad F_a'=\frac{g_a'}{g_a^2}W \quad (x>1).
\end{equation}
If $a\leq a_*$, then $W'>0$, hence $F_a'>0$ on $(1,\infty)$.
Otherwise $W$ decreases strictly from zero to $W(r_a)<0$, then increases strictly.
It eventually becomes positive: Lemma~\ref{R1E4:lem:bounds}, integrated once, gives
\[
 f=x\log x-x+O(\log x),\quad f'=\log x+O(x^{-1}),\quad
 g_a=\log x+O(x^{-1}),\quad g_a'=x^{-1}+O(x^{-2}),
\]
whence
\begin{equation}\label{R1E4:eq:Wasymp}
 \frac{W(x)}{x(\log x)^2}\longrightarrow1.
\end{equation}
Therefore $W$ has exactly one zero $b_a>r_a$. Equation~\eqref{R1E4:eq:W} gives the required derivative signs. At that zero,
\[
 F_a''(b_a)=\frac{g_a'(b_a)}{g_a(b_a)^2}W'(b_a)>0.
\]
\end{proof}

\section{A self-contained left-interval certificate}
Here $a=6$ and $-1<x<1$. The above polynomial definitions specialize to
\begin{align*}
 X&=x^4+24x^3+48x^2-144x-468,\\
 Y&=(x+6)(x^2+6)(x^2+12x-6).
\end{align*}
Let $\alpha=\sqrt{42}-6\in(0,1)$. The sign of $Y$ is negative to the left of $\alpha$ and positive to its right. Also $X<0$ on $(-1,1)$, since the sum of the absolute values of its nonconstant terms is at most $217<468$.

\begin{srcADlemma}\label{R1E4:lem:left}
For $-1<x<1$,
\[
 A(x)=X(x)\psi(x+1)+Y(x)\psi'(x+1)<0.
\]
\end{srcADlemma}
\begin{proof}
The rational logarithm estimate
\begin{equation}\label{R1E4:eq:logbound}
 \log(1+x)\geq L(x):=\frac{x(6+5x)}{2(x+3)(x+1)}\qquad(x>-1)
\end{equation}
follows from equality at zero and
\[
 \frac{\srcADdd}{\srcADdd x}\bigl(\log(1+x)-L(x)\bigr)
 =\frac{x^3}{(x+1)^2(x+3)^2}.
\]
Put $t=x+1$. For $x\leq\alpha$, Lemma~\ref{R1E4:lem:bounds}, $X<0$ and $Y\leq0$ give
\begin{align}
 A(x)&<X\left(L(x)-\frac1{2t}-\frac1{12t^2}\right)
       +Y\left(\frac1t+\frac1{2t^2}\right)
 =\frac{N(x)}{12(x+1)^2(x+3)},\label{R1E4:eq:Nbound}\\
 N(x)&=42x^7+1050x^6+4781x^5+4527x^4-10128x^3\notag\\
 &\hspace{35mm}-7992x^2+7596x-1836.\notag
\end{align}
For $x>\alpha$, use the upper, instead of lower, trigamma bound to obtain
\begin{align}
 A(x)&<X\left(L(x)-\frac1{2t}-\frac1{12t^2}\right)
       +Y\left(\frac1t+\frac1{2t^2}+\frac1{6t^3}\right)
 =\frac{E(x)}{12(x+1)^3(x+3)},\label{R1E4:eq:Ebound}\\
 E(x)&=42x^8+1092x^7+5833x^6+9350x^5-5349x^4\notag\\
 &\hspace{24mm}-17544x^3+828x^2+7704x-3132.\notag
\end{align}
Here are exact, finite sign certificates. For $u>0$,
\begin{align*}
 -(1+u)^7N\left(\frac{-u}{1+u}\right)
 ={}&6542u^7+52807u^6+162848u^5+251061u^4\\
 &+208032u^3+92124u^2+20448u+1836,\\
 -(1+u)^7N\left(\frac{u}{1+u}\right)
 ={}&1960u^7+23555u^6+46942u^5+28245u^4\\
 &+408u^3+972u^2+5256u+1836,\\
 -(1+u)^8E\left(\frac{u}{1+u}\right)
 ={}&1176u^8+34468u^7+87143u^6+76678u^5\\
 &+30249u^4+26184u^3+32940u^2+17352u+3132.
\end{align*}
Every coefficient is positive. Hence $N<0$ on $(-1,1)$ (including zero directly) and $E<0$ on $(0,1)$.
The denominators in \eqref{R1E4:eq:Nbound}--\eqref{R1E4:eq:Ebound} are positive, proving the claim.
\end{proof}

\section{Completion of the unique-minimum theorem}
\begin{proof}[Proof of Theorem~\ref{R1E4:thm:main}]
The only zeros of $g_6$ on $(-1,\infty)$ are $0$ and $1$, and they are simple. The numerator vanishes at both points, so their singularities are analytically removable. Taylor expansion gives
\begin{align}
 \mathcal F(x)&=6\gamma+\left(\frac{13\gamma}{2}-3z_2\right)x+O(x^2)
 \qquad(x\to0),\label{R1E4:eq:zero}\\
 \mathcal F(1+y)&=7(1-\gamma)
 +\frac{7z_2+11\gamma-18}{2}y+O(y^2)
 \qquad(y\to0).\label{R1E4:eq:one}
\end{align}
The coefficient in \eqref{R1E4:eq:one} is negative by the rational estimates used above.
The coefficient in \eqref{R1E4:eq:zero} is negative because
$\gamma<289/500$ and $z_2>1+1/4+1/9=49/36$.

It remains to check strict derivative negativity away from the removable points on $(-1,1)$; mere strict monotonicity would not by itself exclude stationary points.
On $(-1,\alpha)$ and $(\alpha,1)$ define $U=f'/g_6'$ and $W=g_6U-f$ as before.
Lemma~\ref{R1E4:lem:left} yields $U'<0$. On $(-1,0)$, we have $g_6>0$, so $W'<0$ and $W(0)=0$ imply $W>0$.
On $(0,\alpha)$, we have $g_6<0$, so $W'>0$ and again $W>0$.
As $g_6'<0$ on these two intervals, equation~\eqref{R1E4:eq:W} gives $\mathcal F'<0$ there.
On $(\alpha,1)$, $g_6<0$ gives $W'>0$, and $W(1)=0$ gives $W<0$.
Now $g_6'>0$, so again $\mathcal F'<0$.
At $x=\alpha$ itself, $D=Y=0$, $X=-QD'<0$, and $A=X\psi(\alpha+1)<0$; thus $\psi(\alpha+1)>0$.
Since $g_6(\alpha)<0$ and $g_6'(\alpha)=0$,
\[
 \mathcal F'(\alpha)=\frac{\psi(\alpha+1)}{g_6(\alpha)}<0.
\]
Together with \eqref{R1E4:eq:zero}--\eqref{R1E4:eq:one}, this settles the entire left interval and both junctions.

Theorem~\ref{R1E4:thm:family}, applied to $a=6>a_*$, gives exactly one critical point $b>1$ and the asserted signs and nondegeneracy on the right. Finally this point lies below $2$. Indeed,
\[
 W(2)=\frac{40}{11}\left(\frac32-\gamma\right)\log\frac54-\log2.
\]
The logarithm series gives $\log(5/4)>2/9$ and
\[
 \log2<2\left(\frac13+\frac{1}{3}\frac{(1/3)^3}{1-(1/3)^2}\right)
 =\frac{25}{36}<\frac7{10}.
\]
Consequently
\[
 W(2)>\frac{40}{11}\frac{461}{500}\frac29-\frac7{10}
 =\frac{223}{4950}>0.
\]
Since $W$ has just one zero after its initial decrease, $b<2$ follows.
\end{proof}

\section{Scope, reproducibility, and next questions}
The note proves the qualitative assertion requested in \cite[Section 3]{R1E4:BZ}. The certified enclosure given here is only $1<b<2$; the decimal approximation reported in that source is not used in the proof. The argument also gives the parameter classification in Theorem~\ref{R1E4:thm:family}, with the already known constant $a_*$ credited to the earlier literature. It does not settle the behavior for all $a>8$, nor claim optimality of that cutoff.

The initial attempt to use only the crude bound
$-\psi''(t)\leq2\psi'(t)/t$ did not produce a positive derivative certificate near $x=1$. Keeping the next elementary kernel terms in Lemma~\ref{R1E4:lem:bounds} removes that obstruction. The extra polynomial checks in the left interval were developed to avoid importing a computer-assisted sign claim from the published proof.

The companion script \texttt{verify.py} checks all polynomial identities and rational sign bounds exactly using SymPy; it also provides optional floating-point diagnostics, explicitly separated from proof. No computer algebra output is a missing premise: all identities and positive coefficients needed for the theorem are printed above. A further round could seek a parameter classification beyond $8$, a sharper analytically certified enclosure of $b$, or monotonicity of $b_a$ with respect to $a$. These are proposed directions, not results claimed here.

A targeted source search on 20 October 2026 inspected \cite{R1E4:BZ,R1E4:KS,R1E4:YT} and searched the exact title and the source's reported critical-point decimal. No earlier proof of the present right-half-line assertion was identified. This is limited novelty evidence, not a comprehensive priority certification. The problem is a specialized special-function monotonicity problem; its resolution alone should not be described as a major-field breakthrough.

\clearpage


\chapter[A global parameter classification and moving minima for a logarithmic gamma quotient]{A global parameter classification and moving minima  for a logarithmic gamma quotient}
\label{source:R2E4}
\noindent\textit{Source development: Round 2, entry 04. Original mathematical arguments retained; notation is local to this chapter.}\par\medskip


\begin{quote}\small\noindent\textbf{Chapter abstract.} 
For every $a>0$, we classify the logarithmic quotient
$F_a(x)=\log\Gamma(x+1)/\log((x^2+a)/(x+a))$ on $x>1$.
It is strictly increasing up to a known sharp parameter threshold; above that threshold it has exactly one nondegenerate minimum. This removes the bounded-parameter restriction in the author's Round 1 result. The minimizing point is an analytic, strictly increasing function of the parameter, emerges regularly from the endpoint, and has an explicit two-term large-parameter asymptotic. We also obtain the asymptotic of the optimal exponent in a gamma inequality. The new global step is a normalized cubic whose three coefficient functions are proved increasing by exact positive-polynomial certificates. \textbf{Status:} a proposed complete structural extension of a narrowly scoped published problem, with fully written proofs and reproducible exact algebra; the fixed-parameter $a=6$ case has a public proof claim bearing an earlier manuscript date and is not claimed as new here.
\end{quote}

\section{Context and main results}
Write $\gamma$ for Euler's constant, $z_2=\zeta(2)$, and $\psi=\Gamma'/\Gamma$.
Throughout the paper,
\begin{equation}\label{R2E4:eq:family}
 f(x)=\log\Gamma(x+1),\qquad
 g_a(x)=\log\frac{x^2+a}{x+a},\qquad F_a(x)=\frac{f(x)}{g_a(x)}
 \quad (x>1,\ a>0).
\end{equation}
The continuous extension is $F_a(1)=(a+1)(1-\gamma)$.
Kup\'an--Sz\'asz \cite[Theorem 2]{R2E4:KS} identified
\begin{equation}\label{R2E4:eq:threshold}
 \delta=3-z_2-2\gamma>0,\qquad a_* =\frac{z_2-\gamma}{\delta}
\end{equation}
as the sharp decreasing threshold for this family on $(0,1)$; see also Yang--Tian \cite{R2E4:YT}. The constant $a_*$ is not new.
Bulboac\u{a}--Zayed \cite[Section 3]{R2E4:BZ}, published on 22 January 2026, posed the unique-minimum question for $a=6$ on $(-1,\infty)$. Round 1 of this experiment \cite{R2E4:R1} gave a full proof for $a=6$ and a classification on $(1,\infty)$ for $0<a\leq8$.
A late literature check also located a public manuscript by Salazar and Pudim AI \cite{R2E4:SP}, bearing the version date 23 June 2026, whose APP-0029 and APP-0034 state the same fixed-$a=6$ theorem with a ratio-derivative argument. Its public commit date has not been independently verified; nevertheless, no priority is claimed here for that case. The present paper treats \emph{every} positive parameter on the right half-line and studies how its minimum moves. We do not claim that the authors of \cite{R2E4:BZ} posed the full parameter question, or that a literature search establishes absolute novelty.

\newpage
\begin{srcAJtheorem}[Global parameter classification]\label{R2E4:thm:classification}
For all $a>0$, exactly the following alternatives hold.
\begin{enumerate}
\item If $a\leq a_*$, then $F_a'(x)>0$ for all $x>1$.
\item If $a>a_*$, there is exactly one $b_a>1$ with
\[
 F_a'(x)<0\ (1<x<b_a),\qquad F_a'(b_a)=0,\qquad
 F_a'(x)>0\ (x>b_a),\qquad F_a''(b_a)>0.
\]
\end{enumerate}
In particular, the transition parameter is sharp and $0<a_*<6$.
\end{srcAJtheorem}

\begin{srcAJtheorem}[Motion and endpoint emergence]\label{R2E4:thm:motion}
The function $a\mapsto b_a$ is real analytic and strictly increasing from
$(a_*,\infty)$ onto $(1,\infty)$. Put $p=1-\gamma$, $q=z_2-1$, and
\begin{equation}\label{R2E4:eq:Rstar}
 R_* =\frac{4p}{a_*+1}+6q-2(a_*+1)(\zeta(3)-1)>0.
\end{equation}
Then, as $a\downarrow a_*$,
\begin{equation}\label{R2E4:eq:birth}
 b_a=1+\frac{3\delta}{2R_*}(a-a_*)+O((a-a_*)^2).
\end{equation}
The minimum value $m(a)=F_a(b_a)$ is also real analytic and strictly increasing on $(a_*,\infty)$.
\end{srcAJtheorem}

\begin{srcAJtheorem}[Large-parameter asymptotics]\label{R2E4:thm:asymptotics}
There is a unique $c>1$ satisfying
\begin{equation}\label{R2E4:eq:c}
 (1+c^2)\log(1+c^2)=2c^2.
\end{equation}
With
\[
 \kappa=\frac{2c(1+c^2)}{c^2-1},\qquad J(c)=\frac{c}{\log(1+c^2)},
\]
one has, as $a\to\infty$,
\begin{align}
 \frac{b_a}{\sqrt a}&=c-\frac{\kappa}{\log a}+O((\log a)^{-2}),\label{R2E4:eq:basy}\\
 m(a)&=J(c)\sqrt a\left(\frac12\log a+\log c-1\right)
       +O\left(\frac{\sqrt a}{\log a}\right).\label{R2E4:eq:masy}
\end{align}
\end{srcAJtheorem}

\begin{srcAJcorollary}[Optimal power]\label{R2E4:cor:power}
The largest exponent $M(a)$ for which
\[
 \Gamma(x+1)\geq \left(\frac{x^2+a}{x+a}\right)^{M(a)}\qquad(x>1)
\]
holds is $(a+1)(1-\gamma)$ if $a\leq a_*$, and $m(a)$ otherwise.
In the first case the inequality is strict for $x>1$ and sharp as $x\downarrow1$; in the second, equality holds exactly at $x=b_a$.
\end{srcAJcorollary}

\section{Elementary bounds and a digamma ratio}
The elementary estimates used here were also used in Round 1; their proofs are repeated to make the paper independent.
\begin{srcAJlemma}\label{R2E4:lem:bounds}
For $t>0$,
\begin{align}
 \frac1t+\frac1{2t^2}&<\psi'(t)<\frac1t+\frac1{2t^2}+\frac1{6t^3},\label{R2E4:eq:trigamma}\\
 -\psi''(t)&<\frac1{t^2}+\frac1{t^3}+\frac1{2t^4},\label{R2E4:eq:tetragamma}\\
 \log t-\frac1{2t}-\frac1{12t^2}&<\psi(t)<\log t-\frac1{2t}.\label{R2E4:eq:digamma}
\end{align}
Also $\gamma<289/500$ and $z_2<33/20$.
\end{srcAJlemma}
\begin{proof}
The differentiated Euler product gives
\[
 \psi'(t)=\sum_{k=0}^{\infty}(t+k)^{-2}
 =\int_0^\infty e^{-tu}\frac{u}{1-e^{-u}}\,\srcAJdd u.
\]
For $u>0$ the kernel lies strictly between $1+u/2$ and
$1+u/2+u^2/12$. Indeed, multiplying the two positive differences by
$e^u-1$ gives, respectively,
\[
 \sum_{k=3}^{\infty}\frac{k-2}{2k!}u^k,
 \qquad
 \sum_{k=5}^{\infty}\frac{(k-3)(k-4)}{12k!}u^k.
\]
Integrating these bounds, and their product with $u$, proves
\eqref{R2E4:eq:trigamma}--\eqref{R2E4:eq:tetragamma}. Differentiation under the integral is justified by local boundedness at zero and exponential decay at infinity.
The digamma series and integral comparison give $\psi(t)-\log t\to0$;
integrating \eqref{R2E4:eq:trigamma} from $t$ to infinity proves \eqref{R2E4:eq:digamma}.
Finally,
\[
 \log2>2\sum_{j=0}^{3}\frac1{(2j+1)3^{2j+1}}>\frac{693}{1000},
 \quad \psi(2)>\log2-\frac{13}{48}>\frac{211}{500},
\]
so $\gamma<289/500$. The integral test gives
\[
 z_2<\sum_{n=1}^{10}\frac1{n^2}+\frac1{10}
 =\frac{2095337}{1270080}<\frac{33}{20}.
\]
\end{proof}

\begin{srcAJlemma}\label{R2E4:lem:h}
For $x\geq1$, put
\[
 h(x)=\frac{\psi(x+1)}{\psi'(x+1)},\quad h_0(x)=\frac32x-\frac9{10},\quad
 K(x)=\frac{2(x+1)^2+2(x+1)+1}{(x+1)^2(2x+3)}.
\]
Then
\begin{equation}\label{R2E4:eq:hbound}
 h>h_0>0,\qquad h'<1+Kh,\qquad V:=2h-2x+1>x-\frac45>0.
\end{equation}
\end{srcAJlemma}
\begin{proof}
Set $L(x)=x(6+5x)/(2(x+3)(x+1))$.
The identity
\[
 (\log(1+x)-L(x))'=\frac{x^3}{(x+1)^2(x+3)^2}
\]
and equality at zero imply $\log(1+x)\geq L(x)$ for $x\geq0$.
The lower bound for $\psi$ and the upper bound for $\psi'$ in Lemma~\ref{R2E4:lem:bounds}, whose rational numerator is positive for $x\geq1$, give
\[
 h(x)>H(x):=\frac{(x+1)(30x^3+60x^2+11x-21)}{2(x+3)(6x^2+15x+10)}.
\]
Writing $u=x-1\geq0$, direct simplification gives
\[
 H-h_0=\frac{60u^4+249u^3+214u^2-84u+56}
 {10(u+4)(6(u+1)^2+15(u+1)+10)}>0.
\]
The quadratic $214u^2-84u+56$ is positive because its discriminant is negative.
Moreover, $h'=1+kh$, where $k=-\psi''(x+1)/\psi'(x+1)$.
The upper bound \eqref{R2E4:eq:tetragamma} and lower bound \eqref{R2E4:eq:trigamma} give $k<K$. The assertions follow.
\end{proof}

\section{A normalized cubic valid for every parameter}
For $x\geq1$, define
\[
 Q=(x^2+a)(x+a),\quad D=x^2+2ax-a,\quad X=Q'D-QD',\quad Y=QD.
\]
Here $Q,D>0$, and $g_a'=D/Q>0$. Put
\begin{equation}\label{R2E4:eq:U}
 U=\frac{f'}{g_a'}=\frac{Q\psi(x+1)}D,\qquad
 A=X\psi(x+1)+Y\psi'(x+1),\qquad U'=\frac{A}{D^2}.
\end{equation}
The decisive step is not to prove $A'>0$ for all parameters. Instead we normalize its cubic dependence on $a$.

\begin{srcAJproposition}\label{R2E4:prop:normalized}
For each fixed $a>0$, the function
\begin{equation}\label{R2E4:eq:S}
 S_a(x):=\frac{A(x)}{\psi'(x+1)V(x)}
       =-a^3+p_2(x)a^2+p_1(x)a+p_0(x)
\end{equation}
is strictly increasing on $[1,\infty)$. In fact all three coefficient functions $p_j$ are strictly increasing.
\end{srcAJproposition}
\begin{proof}
Direct expansion gives
\begin{align*}
 X&=-2a^3+(2x^2-4x-1)a^2+(4x^3-4x^2)a+x^4,\\
 Y&=(2x-1)a^3+(2x^3+2x^2-x)a^2+3x^4a+x^5.
\end{align*}
Thus $p_j=N_j/V$, with
\begin{align}
 N_0&=x^4(h+x),\notag\\
 N_1&=x^2\bigl((4x-4)h+3x^2\bigr),\label{R2E4:eq:N}\\
 N_2&=(2x^2-4x-1)h+x(2x^2+2x-1).\notag
\end{align}
Regard $N_j=N_j(x,h)$ as a polynomial in two independent variables and set
\[
 T_j(x,h,h')=(\partial_xN_j+h'\partial_hN_j)V-N_j(2h'-2).
\]
Then $V^2p_j'=T_j$. The coefficient of $h'$ in $T_j$ is, respectively,
\begin{equation}\label{R2E4:eq:slopes}
 -x^4(4x-1),\qquad -2x^2(7x^2-6x+2),\qquad -8x^3+6x^2-1.
\end{equation}
All are strictly negative for $x\geq1$. Consequently Lemma~\ref{R2E4:lem:h} yields
\[
 T_j(x,h,h')>B_j(x,h):=T_j(x,h,1+K(x)h).
\]
Here is a fully rational certificate that $B_j>0$. Write $u=x-1\geq0$ and $v=h-h_0>0$, and let
\[
 P_j(u,v)=\sum_{m=0}^2\sum_{k\geq0} C_{j,m,k}u^kv^m.
\]
The finite coefficient lists in Table~\ref{R2E4:tab:certificate} give the identities
\begin{equation}\label{R2E4:eq:certidentity}
 B_0=\frac{(u+1)^3P_0}{50(u+2)^2(2u+5)},\qquad
 B_1=\frac{(u+1)P_1}{25(u+2)^2(2u+5)},\qquad
 B_2=\frac{P_2}{50(u+2)^2(2u+5)}.
\end{equation}
All coefficients are nonnegative, and $C_{j,0,0}>0$. Thus $B_j>0$, whence $p_j'>0$ and $S_a'>0$ for every $a>0$. The identities in this proof are polynomial arithmetic; the accompanying script checks them over the rationals.
\end{proof}

\begin{table}[ht]
\centering
\footnotesize
\setlength{\tabcolsep}{5pt}
\begin{tabular}{cc l}
\toprule
$j$ & $m$ & $(C_{j,m,0},C_{j,m,1},\ldots)$ in increasing powers of $u$\\
\midrule
0&0&$(510,6597,15297,12768,4590,600)$\\
0&1&$(15650,46590,44240,17260,2400)$\\
0&2&$(8000,11200,5200,800)$\\[2pt]
1&0&$(270,5221,18479,26887,17807,5530,650)$\\
1&1&$(12850,56470,96230,73190,25440,3300)$\\
1&2&$(4000,17600,19400,8200,1200)$\\[2pt]
2&0&$(30,2855,10782,13852,7768,2040,200)$\\
2&1&$(10050,45100,74440,55040,18760,2400)$\\
2&2&$(0,8000,11200,5200,800)$\\
\bottomrule
\end{tabular}
\caption{Exact positive-coefficient certificate for \eqref{R2E4:eq:certidentity}.}\label{R2E4:tab:certificate}
\end{table}

\begin{proof}[Proof of Theorem~\ref{R2E4:thm:classification}]
Lemma~\ref{R2E4:lem:bounds} gives $\delta>97/500>0$; also $z_2-\gamma>0$.
Since $7z_2+11\gamma<4477/250<18$, we have $a_*<6$.
At $x=1$, direct evaluation gives
\begin{equation}\label{R2E4:eq:Aone}
 A(1)=(a+1)^2\bigl((a+1)(z_2-1)-(2a-1)(1-\gamma)\bigr)
     =(a+1)^2\delta(a_*-a).
\end{equation}
For each fixed $a$, $A(x)$ is positive for sufficiently large $x$: the polynomial $X$ has leading term $x^4$, and $Y,\psi,\psi'$ are positive there.
Because the normalizing factor in \eqref{R2E4:eq:S} is positive and $S_a'>0$,
$U'$ is positive for all $x>1$ if $a\leq a_*$. Otherwise it changes sign exactly once, from negative to positive, at a point $r_a>1$.

For completeness, the sign-transfer step used in Round 1 is repeated.
Set $W=g_aU-f$. It is regular at $1$, with $W(1)=0$, and
\begin{equation}\label{R2E4:eq:W}
 W'=g_aU',\qquad F_a'=\frac{g_a'}{g_a^2}W\quad(x>1).
\end{equation}
If $a\leq a_*$ then $W'>0$, so $F_a'>0$.
Otherwise $W$ strictly decreases to a negative value at $r_a$ and then strictly increases. It eventually becomes positive. Indeed, Lemma~\ref{R2E4:lem:bounds} and its integral imply, for fixed $a$,
\[
 f=x\log x-x+O(\log x),\quad f'=\log x+O(x^{-1}),\quad
 g_a=\log x+O(x^{-1}),\quad g_a'=x^{-1}+O(x^{-2}),
\]
so $W/[x(\log x)^2]\to1$. Hence $W$ has exactly one zero $b_a>r_a$, and \eqref{R2E4:eq:W} proves the derivative signs. At that point
$F_a''(b_a)=g_a'(b_a)W'(b_a)/g_a(b_a)^2>0$.
\end{proof}

\section{Parameter monotonicity and the endpoint}
\begin{proof}[Proof of Theorem~\ref{R2E4:thm:motion}, except surjectivity]
Since $f,g_a>0$ for $x>1$, put
\[
 H(a,x)=\partial_x\log F_a(x)=\frac{f'}f-\frac{g_a'}{g_a}.
\]
At $x=b_a$, one has $H_x=F_a''(b_a)/F_a(b_a)>0$.
All functions are real analytic, so the implicit function theorem gives a locally analytic branch $b_a$; uniqueness joins these branches globally.

To determine its direction, write $P=a+x$, $Q_1=a+x^2$. A logarithmic-mean identity gives
\begin{equation}\label{R2E4:eq:mean}
 \partial_a\log g_a=-M(a,x),\qquad
 M=\frac{1/P-1/Q_1}{\log(Q_1/P)}
   =\int_0^1P^{-s}Q_1^{-(1-s)}\,\srcAJdd s.
\end{equation}
Since $P_x=1$ and $(Q_1)_x=2x$ are positive, differentiation under this integral gives $M_x<0$. Commuting the mixed derivatives shows
\[
 \partial_a(g_a'/g_a)=-M_x>0,\qquad H_a=M_x<0,
 \qquad b_a'=-H_a/H_x>0.
\]
In particular $b_a$ has a finite limit $\ell\geq1$ as $a\downarrow a_*$. If $\ell>1$, continuity would give $F_{a_*}'(\ell)=0$, contradicting Theorem~\ref{R2E4:thm:classification}. Thus $\ell=1$.

The quotient extends real analytically in $(a,x)$ across $x=1$, because $f$ and $g_a$ have simple zeros there and $g_a'(1)>0$.
From \eqref{R2E4:eq:U} and \eqref{R2E4:eq:Aone}, or by Taylor division,
\begin{equation}\label{R2E4:eq:endpointderivatives}
 F_a'(1)=\frac12U'(1)=\frac\delta2(a_*-a).
\end{equation}
At $a=a_*$ one has $U'(1)=0$, and Taylor division then gives
\[
 F_{a_*}''(1)=\frac13 U''(1)=\frac{A'(1)}{3(a_*+1)^2}=\frac{R_*}{3}.
\]
To evaluate $A'$, differentiate its definition and use $X+Y'=2Q'D$:
\[
 A'=X'\psi+2Q'D\psi'+QD\psi''.
\]
At $1$, $X'=4(a+1)$, $Q'=3(a+1)$, $Q=(a+1)^2$, and $D=a+1$; also $\psi''(2)=-2(\zeta(3)-1)$. These identities give \eqref{R2E4:eq:Rstar}. Its positivity does not require a decimal estimate: at $(a_*,1)$, \eqref{R2E4:eq:S} has $S=0$ and $S'>0$, so $A'(1)>0$.
The analytic implicit function theorem applied to $F_a'(x)=0$ at $(a_*,1)$ now proves \eqref{R2E4:eq:birth}.

Finally $m(a)$ is analytic, and differentiation at its critical point gives
\[
 m'(a)=\partial_aF_a(b_a)>0,
\]
because $f>0$ and $\partial_ag_a=1/(x^2+a)-1/(x+a)<0$ for $x>1$.
Surjectivity follows from the large-parameter estimate proved next.
\end{proof}

\section{The moving minimum on the square-root scale}
\begin{proof}[Proof of Theorem~\ref{R2E4:thm:asymptotics}]
Let $L=\log a$. For $z$ in any fixed compact subinterval of $(0,\infty)$ and $x=z\sqrt a$, the elementary estimates already established imply, uniformly in $z$,
\begin{align*}
 f(x)&=x(\log x-1)+O(\log x),& f'(x)&=\log x+O(x^{-1}),\\
 g_a(x)&=\log(1+z^2)-\log(1+z/\sqrt a),&
 \sqrt a\,g_a'(x)&=\frac{2z}{1+z^2}-\frac{1}{\sqrt a+z}.
\end{align*}
As $\log x=\tfrac12L+\log z$, these yield the uniform expansion
\begin{equation}\label{R2E4:eq:Hscale}
 \sqrt a\,H(a,z\sqrt a)
 =\Lambda(z)+\frac{2}{zL}+O(L^{-2}),\qquad
 \Lambda(z)=\frac1z-\frac{2z}{(1+z^2)\log(1+z^2)}.
\end{equation}
For example,
$\sqrt a\,f'/f=z^{-1}(1+1/(\log x-1))+O(a^{-1/2})$;
expanding the reciprocal makes the $2/(zL)$ term explicit.

The sign of $\Lambda(z)$ is the sign of
$K_0(s)=(1+s)\log(1+s)-2s$, where $s=z^2$.
Since $K_0(0)=0$ and $K_0'(s)=\log(1+s)-1$, the function decreases up to $e-1$ and then strictly increases to infinity. It has exactly one positive root, which exceeds $e-1>1$. This proves existence and uniqueness of $c>1$ in \eqref{R2E4:eq:c}, and $\Lambda$ is negative below $c$ and positive above it.

For any fixed $0<z_-<c<z_+$, \eqref{R2E4:eq:Hscale} gives
$H(a,z_-\sqrt a)<0<H(a,z_+\sqrt a)$ for sufficiently large $a$.
Theorem~\ref{R2E4:thm:classification} therefore forces
$z_-<b_a/\sqrt a<z_+$. Letting the brackets approach $c$ proves
$z_a:=b_a/\sqrt a\to c$. This also proves the surjectivity in Theorem~\ref{R2E4:thm:motion}.

Differentiating $\Lambda$ and using \eqref{R2E4:eq:c} gives
\[
 \Lambda'(c)=\frac{c^2-1}{c^2(1+c^2)}>0.
\]
At $z=z_a$, \eqref{R2E4:eq:Hscale} first implies $z_a-c=O(L^{-1})$ by the mean value theorem. Expanding it once more at $c$ gives
\[
 0=\Lambda'(c)(z_a-c)+\frac{2}{cL}+O(L^{-2}),
\]
which is \eqref{R2E4:eq:basy}.

For $z$ in a compact neighborhood of $c$, the same uniform estimates give
\begin{equation}\label{R2E4:eq:Fscale}
 F_a(z\sqrt a)=\sqrt a\,J(z)\left(\frac12L+\log z-1\right)+O(L),
 \qquad J(z)=\frac{z}{\log(1+z^2)}.
\end{equation}
The defining equation for $c$ is exactly $J'(c)=0$.
Thus $J(z_a)=J(c)+O(L^{-2})$ and
$J(z_a)(\log z_a-1)=J(c)(\log c-1)+O(L^{-1})$.
Substitution into \eqref{R2E4:eq:Fscale} yields \eqref{R2E4:eq:masy}, since $L=O(\sqrt a/L)$ as $a\to\infty$.
\end{proof}

\begin{proof}[Proof of Corollary~\ref{R2E4:cor:power}]
Since $g_a(x)>0$, the stated power inequality is equivalent to
$F_a(x)\geq M(a)$. The infima and their equality cases are exactly those in Theorem~\ref{R2E4:thm:classification} and its analytic endpoint extension.
\end{proof}

\section{Scope, reproducibility, and remaining questions}
The main advance over Round 1 is structural, not an enlargement of a finite computation: one normalization proves the phase diagram for \emph{all} $a>0$. The trajectory results and both asymptotic formulas are new to this round. The known constant $a_*$, the elementary polygamma bounds, the derivative-ratio framework, and the $W$ sign-transfer argument are not counted as new. No theorem from another entrant is used.

The exact verification script checks the rational lower bound for $h$, the three derivative coefficients, all identities in Table~\ref{R2E4:tab:certificate}, the normalized cubic, endpoint algebra, and the asymptotic constant identities. Its optional numerical mode is only a diagnostic and supplies no part of the proof. In particular, the decimal approximations $a_*\approx5.321706$, $c\approx1.980291$ and $\kappa\approx6.671868$ are not certified enclosures and are not used.

The direct Round 1 route cannot prove the present theorem for all parameters: indeed,
\[
 A'(1)=4(a+1)(1-\gamma)+6(a+1)^2(z_2-1)
       -2(a+1)^3(\zeta(3)-1)
\]
is negative for all sufficiently large $a$. This is why monotonicity of the normalized cubic, rather than monotonicity of $A$, is essential.

This note makes no new classification claim on $(-1,1)$, where admissible domains and derivative zeros need separate treatment. A full parameter phase diagram on that interval, and analogous normalized-coefficient methods for other gamma quotients, remain natural directions. The current result is a special-function theorem of limited scope, not evidence of the resolution of a major general conjecture. The bibliography and targeted searches were checked on 20 October 2026; absence of a matching result in those searches is not a priority guarantee.



\chapter[Finite Gaussian--Newton certificates for higher-order q-log-concavity of Gaussian columns]{Finite Gaussian--Newton certificates for higher\mbox{-}order  $q$\mbox{-}log\mbox{-}concavity of Gaussian columns}
\label{source:R1E3}
\noindent\textit{Source development: Round 1, entry 03. Original mathematical arguments retained; notation is local to this chapter.}\par\medskip


\begin{quote}\small\noindent\textbf{Chapter abstract.} 
This is a partial-result research note, not a solution of the infinite
$q$-log-concavity conjecture. We give an explicit seven-term positive
Gaussian-binomial expansion proving fourth-order coefficientwise
$q$-log-concavity of the column $\bigl(\srcACqb{n+3}{3}\bigr)_{n\geq0}$.
A polynomial-degree reduction and a reproducible exact-integer calculation
certify sixth-order $q$-log-concavity of this column and fourth-order
$q$-log-concavity of the columns with lower indices $4,5,6$.
These statements hold for every column position, not just a finite tested
range. The unbounded iteration and unbounded lower-index problems remain
open in this note. The certificates were derived here; their priority has
not been established by an exhaustive literature review.
\end{quote}

\section{Target, conventions, and status}
Put $D_k(q)=\prod_{i=1}^k(1-q^i)$, with $D_0=1$, and
\[
 a^{(k)}_n(q)=\srcACqb{n+k}{k}=\frac{\prod_{i=1}^k(1-q^{n+i})}{D_k(q)}
 \quad(n\geq0),\qquad a^{(k)}_n=0\quad(n<0).
\]
Gaussian coefficients are polynomials with nonnegative integer
coefficients, as follows, for example, by their Pascal recurrence.
For any such sequence define
\[
 (\srcACLL a)_n=a_n^2-a_{n-1}a_{n+1}.
\]
Here ``$r$-fold $q$-log-concave'' means that every coefficient of
$(\srcACLL^s a)_n$ is nonnegative for all $n\geq0$ and $0\leq s\leq r$.
Evaluation at a positive real $q$ is a strictly weaker assertion.

McNamara and Sagan's Conjecture 5.3 in the March 2009 preprint
\cite{R1E3:MS} asks whether $a^{(k)}$ is infinitely $q$-log-concave for every
$k$. They establish the cases $k\leq2$; their Schur-positivity result
(Theorem 6.4) implies the first three iterations for all $k$ by principal
specialization. The ordinary specialization $q=1$ was subsequently
settled by Br\"and\'en and Chasse \cite{R1E3:BC}. Thus an ordinary numerical
log-concavity argument would not resolve the coefficientwise question.
The conjecture is explicitly open in the cited 2009 source. Targeted
searches on October 20, 2026 did not locate a resolution or the particular
certificates below; this is not a claim of exhaustively verified present-day
openness or novelty.

\begin{srcACtheorem}\label{R1E3:thm:main}
The column $a^{(3)}$ is $6$-fold $q$-log-concave. Each of
$a^{(4)},a^{(5)},a^{(6)}$ is $4$-fold $q$-log-concave.
The fourth step for $a^{(3)}$ has the explicit certificate in
Section~\ref{R1E3:sec:explicit}. The remaining stated steps are
computer-assisted, with the finite reduction proved in
Section~\ref{R1E3:sec:finite} and exact verification specified in
Section~\ref{R1E3:sec:computation}.
\end{srcACtheorem}

\section{Why finitely many certificates prove every position}\label{R1E3:sec:finite}
Work over the rational-function field $K=\mathbb Q(q)$, with $q$
an indeterminate. Define
\[
 F_{k,0}(t)=\frac{\prod_{i=1}^k(1-q^it)}{D_k(q)},\qquad
 F_{k,r+1}(t)=F_{k,r}(t)^2-F_{k,r}(t/q)F_{k,r}(qt).
\]
\begin{srcAClemma}\label{R1E3:lem:degree}
For $k\geq2$, put $e_r=2^r-1$ and $d_{k,r}=2^r(k-2)+2$.
There exists $G_{k,r}(t)\in K[t]$ of degree at most $d_{k,r}$ such that
\[
 F_{k,r}(t)=t^{e_r}G_{k,r}(t),\qquad
 (\srcACLL^r a^{(k)})_n=q^{e_r n}G_{k,r}(q^n)\quad(n\geq0).
\]
\end{srcAClemma}
\begin{proof}
The initial rational expression agrees with $a^{(k)}_n$ for $n\geq0$.
Its evaluations at $n=-1,-2$ are zero. Two consecutive zero entries
remain zero under $\srcACLL$: each new entry is a square of zero minus a
product having a zero factor. Consequently evaluations at all
nonnegative positions evolve exactly as the zero-extended sequence,
even though the rational expression need not vanish at more negative
positions.

Suppose $F=t^eG$, where $\deg G\leq d$. Then
\[
 F(t)^2-F(t/q)F(qt)
 =t^{2e}\{G(t)^2-G(t/q)G(qt)\}.
\]
The expression in braces has zero constant term. Its term of degree
$2d$ also cancels, since the two leading coefficients agree.
After division by $t$ it is therefore a polynomial of degree at most
$2d-2$. Starting with $e_0=0,d_0=k$ gives
$e_{r+1}=2e_r+1$ and $d_{r+1}=2d_r-2$, proving the formulas.
\end{proof}

For $j\geq0$ set
\[
 B_j(t)=\frac{\prod_{i=0}^{j-1}(1-q^{-i}t)}{D_j(q)}.
\]
Then $B_j(q^n)=\srcACqb{n}{j}$, with value zero when $j>n\geq0$,
and $B_j(q^j)=1$. The $B_j$ form a degree-graded basis of $K[t]$.

\begin{srcAClemma}[Positive interpolation certificate]\label{R1E3:lem:interpolation}
Let $G\in K[t]$ have degree at most $d$. Define successively
\[
 C_j=G(q^j)-\sum_{i=0}^{j-1}C_i\srcACqb{j}{i}
 \qquad(0\leq j\leq d).
\]
Then $G(t)=\sum_{j=0}^d C_jB_j(t)$. If every $C_j$ is a
Laurent polynomial with nonnegative coefficients, so is $G(q^n)$
for every $n\geq0$.
\end{srcAClemma}
\begin{proof}
The proposed expansion agrees with $G$ at the $d+1$ distinct
field elements $1,q,\ldots,q^d$ by triangularity. Their difference
has degree at most $d$ and hence vanishes identically.
The last assertion follows from positivity of Gaussian coefficients.
\end{proof}

Applied to $G_{k,r}$, this proves positivity of the entire $r$th
iterate once $d_{k,r}+1$ Laurent polynomials have been certified.
Their Laurent-polynomial nature is itself automatic from the recursion:
$G_{k,r}(q^j)=q^{-e_rj}(\srcACLL^ra^{(k)})_j$ is Laurent-polynomial,
and all subtracted terms are Laurent-polynomial. Any harmless common
power of $q$ can be used to make the computation polynomial.
Multiplication back by $q^{e_rn}$ preserves coefficient signs;
Lemma~\ref{R1E3:lem:degree} guarantees that the result is an ordinary polynomial.
No positivity assumption at a finite set of original positions is being
extrapolated without a certificate.

\section{An explicit fourth-step identity}\label{R1E3:sec:explicit}
Write $a_n=a^{(3)}_n$ and $D=D_3(q)$. Introduce
\begin{align*}
 A&=1+3q+q^2,\\
 U&=1+2q+3q^2+3q^3+3q^4+2q^5+q^6,\\
 V&=2+4q+10q^2+18q^3+26q^4+30q^5+26q^6\\
  &\hspace{18mm}+18q^7+10q^8+4q^9+2q^{10},\\
 W&=q+2q^2+7q^3+14q^4+23q^5+32q^6+42q^7\\
  &\hspace{18mm}+32q^8+23q^9+14q^{10}+7q^{11}+2q^{12}+q^{13},\\
 R(q,t)&=2A(1+q^{12}t^6)-4U(t+q^8t^5)
             +V(t^2+q^4t^4)-Wt^3.
\end{align*}
Direct polynomial multiplication gives
\begin{equation}\label{R1E3:eq:fourth}
 (\srcACLL^4a)_n=q^{15n-1}a_n^4\,\frac{R(q,q^n)}{D^2}.
\end{equation}
For transparency, this identity can be obtained from the following
successive factorizations, where $P(t)=\prod_{i=1}^3(1-q^it)$:
\begin{align*}
 F_{3,1}(t)&=\frac{t(1-q)^2(1+q+q^2)
 (1-qt)(1-q^2t)^2(1-q^3t)}{D^2},\\
 F_{3,2}(t)&=t^3\frac{P(t)^2}{D^2},\\
 F_{3,3}(t)&=\frac{t^7(1-q)^{14}(1+q)^4(1+q+q^2)^5}{D^8}\\
 &\quad\cdot(1-qt)^2(1-q^2t)^4(1-q^3t)^2
 \bigl(2-(1+q+q^3+q^4)t+2q^4t^2\bigr).
\end{align*}
Substituting the last expression into the defining operation for
$F_{3,4}$ gives \eqref{R1E3:eq:fourth}.

Here is the positive part of the certificate. For a finite integer vector
$v=(v_0,\ldots,v_m)$ write $P_v(q)=\sum_{i=0}^m v_iq^i$, and put
$[5]_q=1+q+q^2+q^3+q^4$. Define
\begin{align*}
 C_0={}&q,\\
 C_1={}&P_{(2,5,10,14,17,14,10,6,2)},\\
 C_2={}&qP_{(2,5,14,29,49,70,87,86,78,60,38,20,10,2)},\\
 C_3={}&q^4P_{(1,3,11,26,53,91,139,174,196,194,168,128,88,50,24,10,2)},\\
 C_4={}&2q^{10}(1+q^2)P_{(1,3,8,18,31,45,56,60,56,47,34,21,11,5,1)},\\
 C_5={}&2q^{18}(1+q^2)[5]_qP_{(2,4,6,8,9,6,5,4,1)},\\
 C_6={}&2q^{27}(1+q^3)(1+q^2)(1+3q+q^2)[5]_q.
\end{align*}
Every displayed coefficient is manifestly nonnegative. The exact identity is
\begin{equation}\label{R1E3:eq:certificate}
 \frac{R(q,t)}{D^2}=\sum_{j=0}^{6} C_jB_j(t).
\end{equation}
It can be checked by multiplication in $\mathbb Q(q)[t]$, or by
substituting $t=1,q,\ldots,q^6$ and using triangular interpolation.
Both sides have degree at most six, so these checks prove an identity,
not an approximation.

Combining \eqref{R1E3:eq:fourth} and \eqref{R1E3:eq:certificate} proves fourth-step
coefficientwise positivity for $n\geq1$, since $15n-1\geq0$.
At $n=0$ only $C_0$ contributes and the answer is $q^{-1}C_0=1$.
This also checks the left endpoint, which cannot simply be discarded.
The first three steps are covered by the known Schur-positive identities,
or can be checked directly. Thus the fourth-order assertion for the
entire third column follows without a large certificate.

\section{Exact higher-order verification}\label{R1E3:sec:computation}
The accompanying \texttt{verify.py} uses only Python's standard library
and arbitrary-precision integers. The input polynomial arrays are generated
by
\[
 \srcACqb{N}{j}=\srcACqb{N-1}{j}+q^{N-j}\srcACqb{N-1}{j-1}.
\]
For each pair $(k,r)$ it forms $a_n^{(k)}$ for
$0\leq n\leq d_{k,r}+r$, applies $\srcACLL$ exactly $r$ times while
dropping the unused rightmost entry each time, and applies
Lemma~\ref{R1E3:lem:interpolation} to the resulting $d_{k,r}+1$ values.
A common shift $q^{rk2^r}$ avoids negative array indices; the script
explicitly checks that no nonzero coefficient has been discarded.
Every coefficient of every resulting certificate is then checked to be
an integer at least zero.

\begin{center}
\begin{tabular}{rrrr}
\toprule
$k$ & $r$ & degree bound $d_{k,r}$ & certificate polynomials\\
\midrule
3&4&18&19\\
3&5&34&35\\
3&6&66&67\\
4&4&34&35\\
5&4&50&51\\
6&4&66&67\\
\bottomrule
\end{tabular}
\end{center}
These six finite tests all pass. Together with the first three iterations
from \cite{R1E3:MS}, they prove Theorem~\ref{R1E3:thm:main}.
Notice that the fourth-step calculation in this table uses the full
normalized polynomial, whereas the short certificate in
Section~\ref{R1E3:sec:explicit} first extracts the positive factor $a_n^4$.
Their sizes therefore differ.

For reproducibility, run \texttt{python verify.py}; the shorter
\texttt{python verify.py --quick} checks the full fourth-step certificate
for $k=3$. The program prints the number of nonzero integer coefficients
and a SHA-256 digest of the canonical coefficient arrays. The file
\texttt{verification.txt} records the run used for this note.
Polynomial multiplication is implemented by exact integer packing:
the packing base exceeds the bound
$\min(\deg f+1,\deg g+1)\max|f_i|\max|g_i|$ for every nonnegative
product. Hence no coefficient carries into the next digit. Signed
polynomials are split into positive and negative parts. Small products
are independently compared with ordinary schoolbook multiplication.
There is no floating-point root finding or probabilistic sign test.

\section{Limitations, failed routes, and the next question}
The result is a bounded-depth, bounded-column partial result. It is not
an infinite-log-concavity proof, and it does not address arbitrary $k$.
The exact-computation part depends on correctness of the supplied finite
integer algorithm, not on an external symbolic-algebra oracle.
The weakest non-computational issue is priority: absence of a matching
result in a targeted search is not proof of novelty.

The first attempted route expanded the iterates in powers of $t=q^n$.
Their coefficients alternate in sign, so this basis does not directly
prove coefficientwise positivity after specialization. Leading-term
asymptotics at fixed real $q$ would only prove eventual numerical
positivity and would miss the actual conjecture. A second attempted
route sought a uniform strong-factor log-concavity bound; such a bound
is unsuitable for polynomial-growth columns at $q=1$, where adjacent
ratios approach one. These obstacles motivated the Gaussian--Newton
basis, whose evaluation is coefficientwise nonnegative.

A concrete next problem is to explain the observed positivity of the
Gaussian--Newton coefficients structurally, instead of by computing
larger finite certificates. An invariance statement for the operator
\[
 G(t)\longmapsto\frac{G(t)^2-G(t/q)G(qt)}{t}
\]
would need hypotheses stronger than positivity in the Newton basis.
Indeed, at $q=1$ the Newton-positive polynomial
$1+\binom n2$ has values $1,1,2$ at $n=0,1,2$, so its first Tur\'an
difference at $n=1$ is negative. Thus closure of the entire positive
Newton cone is false. Any future closure proof must use additional
structure of the Gaussian columns, such as their forced negative-index
zeros. The small explicit certificate here and the saved verifier provide
starting points for that investigation; no subsequent round is claimed.



\chapter[A sharp three-step root factorization for Gaussian columns and its fourth-fold log-concavity consequence]{A sharp three-step root factorization for Gaussian columns and its fourth-fold log-concavity consequence}
\label{source:R7E3}
\noindent\textit{Source development: Round 7, entry 03. Original mathematical arguments retained; notation is local to this chapter.}\par\medskip

\begin{quote}\small\noindent\textbf{Chapter abstract.} 
For every Gaussian-binomial column, we factor its third log-concavity iterate into positive shifted $q$-integers, with two explicitly characterized fractional shifts and all other shifts integral. Equivalently, the first three normalized $q$-dilation Tur\'an iterates of a finite geometric-root polynomial have only positive real roots. This root assertion is sharp: a rational example at the fourth iterate has exactly four nonreal roots. The factorization proves fourth-fold log-concavity for every positive real specialization of $q$, uniformly in the column, and gives infinitely divisible reciprocal Hausdorff moment sequences. We also give an explicit positive polynomial multiplier for the coefficientwise fourth-fold inequality. These are not proofs of coefficientwise fourth-fold or infinite $q$-log-concavity. The starting rectangular identities are classical; the contribution here is their explicit factorization, root obstruction, and consequences. Historical priority beyond the sources checked is not asserted.
\end{quote}

\section{Question, scope, and main theorem}
Extend a sequence by zero at negative indices and put
\[
 (\srcBMLL a)_n=a_n^2-a_{n-1}a_{n+1}\quad(n\geq0).
\]
For fixed $k$, the Gaussian column is
\begin{equation}\label{R7E3:column}
 a_n=\genfrac{[}{]}{0pt}{}{n+k}{k}_q
 =\prod_{j=1}^k\frac{\srcBMbr{n+j}}{\srcBMbr j},\qquad
 \srcBMbr x=\frac{1-q^x}{1-q},
\end{equation}
with $\srcBMbr x=x$ when $q=1$. McNamara--Sagan's Conjecture 5.3 \cite{R7E3:MS} asks whether every iterate has nonnegative \emph{coefficients} as a polynomial in $q$. Their Theorem 6.4 proves the first three iterates by stronger Schur positivity; the fourth symmetric-function iterate is not even monomial-positive. We do not settle the coefficientwise conjecture.

Here is the precise additional structure. For $0\leq r\leq k+1$ define multisets and positive constants
\begin{equation}\label{R7E3:Mr}
 M_r=\{i+j-1:1\leq i\leq r,\ 1\leq j\leq k+1-r\},
 \qquad D_r=\prod_{c\in M_r}\srcBMbr c.
\end{equation}
Empty products are one. Multiset addition adds multiplicities; subtraction below removes specified occurrences. For $k\geq3$ set
\begin{equation}\label{R7E3:Cdef}
 C=2M_1+M_2+M_4-\{1,k\},\qquad
 R=\frac{\srcBMbr3\srcBMbr{k-2}}{\srcBMbr k}.
\end{equation}
Thus $C$ has $8k-16$ members, all positive integers.

\begin{srcBMtheorem}[Fractional-shift factorization]\label{R7E3:main}
Fix $k\geq3$ and $q>0$. There is a unique $\delta\in(0,1)$ such that, if $q\ne1$,
\begin{equation}\label{R7E3:delta}
 q^\delta+q^{k+1-\delta}
 =1+q^{k+1}-\frac{R}{R+1}(1-q)(1-q^k).
\end{equation}
At $q=1$ its definition is $\delta(k+1-\delta)=Rk/(R+1)$, where $R=3(k-2)/k$. For every integer $n\geq0$,
\begin{equation}\label{R7E3:factor}
 (\srcBMLL^3 a)_n=
 \frac{R+1}{D_3^2D_2}\,q^{7n}
 \prod_{c\in C}\srcBMbr{n+c}\,
 \srcBMbr{n+\delta}\srcBMbr{n+k+1-\delta}.
\end{equation}
The right side, viewed as a function of real $n\geq0$, is strictly log-concave. Consequently $(\srcBMLL^4a)_n>0$ for every $n\geq0$.
\end{srcBMtheorem}

The following polynomial formulation specifies exactly what root preservation is meant. It concerns an interpolation variable, not a sequence's ordinary generating function.
\begin{srcBMtheorem}[Sharp root depth]\label{R7E3:roots}
For $q>0$, $q\ne1$, define
\[
 F_k(t)=\prod_{j=1}^k(1-q^jt),\qquad
 \srcBMT f(t)=\frac{f(t)^2-f(t/q)f(qt)}{t}.
\]
For every $k\geq3$, $\srcBMT^r F_k$ has only positive real roots for $r=0,1,2,3$. At $r=3$ the roots, with multiplicity, are exactly
\[
 \{q^{-c}:c\in C\}\ \cup\ \{q^{-\delta},q^{-(k+1-\delta)}\}.
\]
This statement cannot be extended to $r=4$: $\mathcal T_{1/2}^4F_3$ has exactly four nonreal roots.
\end{srcBMtheorem}

\paragraph{Delta from earlier rounds.}
Entry 03's Round 1 \cite{R7E3:R1} gave coefficientwise finite-depth certificates for a few fixed columns ($k=3$ through depth six and $k=4,5,6$ through depth four). The present result is uniform in $k$, but its fourth-fold conclusion is only pointwise in $q>0$. It does not dominate those coefficientwise certificates. No theorem from Rounds 2--6 is reused. The distinct content is the explicit fractional-shift factorization and sharp root-depth obstruction; their consequences should not be counted as independent discoveries.

\section{Classical rectangular identities, with a direct proof}
Put
\begin{equation}\label{R7E3:S}
 S_r(n)=q^{\binom r2 n}\prod_{c\in M_r}\frac{\srcBMbr{n+c}}{\srcBMbr c}
 \quad(0\leq r\leq k+1).
\end{equation}
These are the principal specializations $s_{(n^r)}(1,q,\ldots,q^k)$ for integer $n\geq0$. We use the classical rectangular identity
\begin{equation}\label{R7E3:Q}
 S_r(n)^2-S_r(n-1)S_r(n+1)=S_{r-1}(n)S_{r+1}(n)
 \quad(1\leq r\leq k).
\end{equation}
It is the Dodgson/rectangular Schur identity used in \cite[Theorems 6.3--6.4]{R7E3:MS}; it is not new here. For completeness it can be verified without Schur functions. For real $n>0$, cancellation of the factors in \eqref{R7E3:S} gives, with $m=k+1$,
\begin{align*}
 \frac{S_r(n-1)S_r(n+1)}{S_r(n)^2}
 &=\frac{\srcBMbr n\srcBMbr{n+m}}{\srcBMbr{n+r}\srcBMbr{n+m-r}},\\
 \frac{S_{r-1}(n)S_{r+1}(n)}{S_r(n)^2}
 &=\frac{q^n\srcBMbr r\srcBMbr{m-r}}{\srcBMbr{n+r}\srcBMbr{n+m-r}}.
\end{align*}
Their sum is one, by expanding the $q$-integers. Continuity includes $n=0$. Every $S_r(-1)=0$ for $1\leq r\leq k$, so the relevant boundary agrees with zero extension.

In particular $a=S_1$ and
\begin{equation}\label{R7E3:firstthree}
 \srcBMLL a=S_2,\qquad \srcBMLL^2a=S_1S_3,\qquad
 \srcBMLL^3a=S_1^2S_2S_4+S_3(n-1)S_2(n)S_3(n+1).
\end{equation}
Indeed $\srcBMLL(fg)=f(n)^2\srcBMLL g+g(n-1)g(n+1)\srcBMLL f$, and \eqref{R7E3:Q} gives these equations. They are precisely the known first-three-iterate identities, reproduced to make the subsequent proof self-contained.

\section{Cancellation to two fractional shifts}
Write $P_E(n)=\prod_{c\in E}\srcBMbr{n+c}$ for a multiset $E$. The two terms of \eqref{R7E3:firstthree} are
\begin{equation}\label{R7E3:AB}
 q^{7n}\frac{P_U(n)}{D_A},\qquad
 q^{7n}\frac{P_V(n)}{D_B},
\end{equation}
where
\[
 U=2M_1+M_2+M_4,\quad
 V=(M_3-1)+M_2+(M_3+1),\quad
 D_A=D_1^2D_2D_4,\quad D_B=D_3^2D_2.
\]
Here $M_3\pm1$ shifts every member, rather than removing an occurrence. Direct multiset cancellation yields
\begin{equation}\label{R7E3:cancel}
 U=C+\{1,k\},\qquad V=C+\{0,k+1\},\qquad
 \frac{D_B}{D_A}=\frac{\srcBMbr3\srcBMbr{k-2}}{\srcBMbr k}=R.
\end{equation}
For clarity, the first two cancellations can be checked as a polynomial identity in an auxiliary variable $z$. The enumerator of $M_r$ is
$(z+\cdots+z^r)(1+\cdots+z^{k-r})$ for $r\leq k$, and that of $M_{k+1}$ is zero. Substitution gives
$2M_1(z)+M_4(z)-(z+z^{-1})M_3(z)=z+z^k-1-z^{k+1}$.
The last cancellation follows by cancelling the same consecutive factors in $D_3^2/(D_1^2D_4)$; it holds also at $k=3$, where $D_4=1$.
Thus \eqref{R7E3:AB} becomes
\begin{equation}\label{R7E3:H}
 (\srcBMLL^3a)_n=\frac{q^{7n}P_C(n)}{D_B}H(n),\qquad
 H(n)=R\srcBMbr{n+1}\srcBMbr{n+k}+\srcBMbr n\srcBMbr{n+k+1}.
\end{equation}
The elementary identity
$\srcBMbr{n+1}\srcBMbr{n+k}-\srcBMbr n\srcBMbr{n+k+1}=q^n\srcBMbr k$
gives, on setting $t=q^n$,
\[
 (1-q)^2H(n)=(R+1)(1-St+q^{k+1}t^2),\quad
 S=1+q^{k+1}-\frac{R}{R+1}(1-q)(1-q^k).
\]
Since $R>0$, $S$ lies strictly between $q+q^k$ and $1+q^{k+1}$. The function
$x\mapsto q^x+q^{k+1-x}$ is strictly decreasing on $[0,1]$ for either $0<q<1$ or $q>1$, with these endpoint values. This proves existence and uniqueness in \eqref{R7E3:delta} and factors $H$ as
\[
 H(n)=(R+1)\srcBMbr{n+\delta}\srcBMbr{n+k+1-\delta}.
\]
At $q=1$, directly $H(n)=(R+1)n(n+k+1)+Rk$; the increasing function $x(k+1-x)$ on $[0,1]$ has endpoint values $0,k$, proving the analogous factorization.

Every shift in \eqref{R7E3:factor} is positive. For such a shift $c$ and $x\geq0$,
\[
 \frac{d^2}{dx^2}\log\srcBMbr{x+c}
 =-\frac{(\log q)^2q^{x+c}}{(1-q^{x+c})^2}<0
 \quad(q\ne1),
\]
and the value at $q=1$ is $-1/(x+c)^2$. The geometric factor has zero logarithmic curvature. Strict log-concavity therefore proves the claimed inequality for $n\geq1$. At $n=0$, zero extension gives $(\srcBMLL^r a)_0=1$ for every $r$. This completes Theorem~\ref{R7E3:main}.

\paragraph{Small columns and $q=0$.}
For $k=0$, $a_n=1$; for $k=1$, $\srcBMLL a_n=q^n$ for $n\geq0$, so the next iterate is $(1,0,0,\ldots)$. For $k=2$, \eqref{R7E3:Q} gives $\srcBMLL a=q^n a$ and hence $\srcBMLL^r a=q^{(2^r-1)n}a$. These prove infinite log-concavity in the small cases, consistent with \cite{R7E3:MS}. At $q=0$ every column is the constant-one sequence, so fourth-fold nonnegativity also holds there. Thus the pointwise fourth-fold assertion covers every column and every $q\geq0$.

\section{Root locations and a sharp obstruction}
The numerator defining $\srcBMT$ has zero constant term, so $\srcBMT$ maps polynomials to polynomials. If a polynomial vanishes at $q^{-1}$ and $q^{-2}$, so does its image: at either point the shifted product contains the other root. Consequently the polynomial interpolation respects the zero-extension boundary at every iteration. Induction gives
\begin{equation}\label{R7E3:interp}
 (\srcBMLL^ra)_n=
 \frac{q^{(2^r-1)n}\srcBMT^r F_k(q^n)}{\left(\prod_{j=1}^k(1-q^j)\right)^{2^r}}.
\end{equation}
The first two identities in \eqref{R7E3:firstthree} show that, up to nonzero constants, $\srcBMT F_k$ and $\srcBMT^2F_k$ are respectively
\[
 \prod_{c\in M_2}(1-q^ct),\qquad
 \prod_{c\in M_1+M_3}(1-q^ct).
\]
Theorem~\ref{R7E3:main} similarly identifies $\srcBMT^3F_k$. The identities hold as polynomials, since they hold at the infinitely many distinct points $q^n$. This proves the positive-real-root assertion and gives degree $8k-14$ at the third iterate.

For sharpness take $q=1/2$, $k=3$. Exact polynomial arithmetic gives
\begin{equation}\label{R7E3:obstruct}
 \mathcal T_{1/2}^4F_3(t)=
 \frac{16679880978201}{79228162514264337593543950336}
 (t-2)^4(t-4)^4(t-8)^4 P(t),
\end{equation}
where
\begin{align*}
 P(t)={}&11t^6-434t^5+6097t^4-36737t^3\\
        &+97552t^2-111104t+45056.
\end{align*}
This displayed identity is checked by the supplied exact verifier. Its root count needs no numerical approximation. With
\[
 g(x)=11x^3-434x^2+5569x-22849
\]
we have $P(t)=t^3g(t+16/t)$. The discriminant of $g$ is $-278129439<0$, so it has one real root $\rho$ and two nonreal roots. Since $g(8)=-441<0$ and its leading coefficient is positive, $\rho>8$. The equation $t+16/t=\rho$ has two distinct positive solutions. Each nonreal root of $g$ gives two nonreal solutions, because a nonzero real $t$ would give real $t+16/t$. Thus $P$ has exactly two real and four nonreal roots, proving Theorem~\ref{R7E3:roots}. This is not a negative value of $\srcBMLL^4a$, nor a counterexample to the Gaussian-column conjecture.

\section{Two consequences of the explicit factorization}
\begin{srcBMcorollary}[Reciprocal moment sequences]\label{R7E3:moments}
For $k\geq3$, $0<q<1$, and every $\theta>0$, the sequence
$b_n=(q^{-7n}(\srcBMLL^3a)_n)^{-\theta}$ is the moment sequence of a probability measure on $[0,1]$. At $q=1$ the same holds without the geometric normalization.
\end{srcBMcorollary}
\begin{proof}
Let $E=C+\{\delta,k+1-\delta\}$. Since $(\srcBMLL^3a)_0=1$, the factorization yields
\[
 b_n=\prod_{c\in E}\left(\frac{1-q^c}{1-q^{n+c}}\right)^\theta.
\]
For $\alpha=q^c\in(0,1)$, the binomial series gives
\[
 \left(\frac{1-\alpha}{1-\alpha q^n}\right)^\theta
 =\sum_{j=0}^\infty\frac{(\theta)_j}{j!}(1-\alpha)^\theta\alpha^j(q^j)^n.
\]
The weights are nonnegative and sum to one at $n=0$. Each factor is therefore the moment sequence of a discrete probability measure. A finite product of moment sequences is the moment sequence of the product of independent variables with these distributions. Replacing $\theta$ by any positive multiple proves the power assertion. For $q=1$, each factor is $(c/(n+c))^\theta$, represented by
\[
 \frac{c^\theta}{\Gamma(\theta)}\int_0^\infty
 u^{\theta-1}e^{-cu}(e^{-u})^n\,du.
\]
The same product argument applies.
\end{proof}

\begin{srcBMproposition}[A coefficientwise clearing multiplier]\label{R7E3:multiplier}
With $D_A,D_B$ from \eqref{R7E3:AB}, for every integer $k\geq3$ and $n\geq1$ the polynomial
\[
 D_A^2D_B^2(\srcBMLL^4a)_n
\]
has nonnegative coefficients in $q$. The multiplier is independent of $n$.
\end{srcBMproposition}
\begin{proof}
For integers $a,b\geq0$ with $|a-b|\leq1$,
\[
 \srcBMbr{n+a}\srcBMbr{n+b}-\srcBMbr{n+a-1}\srcBMbr{n+b+1}
 =q^{n+a-1}\srcBMbr{b-a+1}
\]
is coefficientwise nonnegative for $n\geq1$. For equal-length multisets $A,B$ admitting a matching with paired differences at most one, telescope the product of these paired inequalities to obtain
\[
 P_A(n)P_B(n)-P_A(n-1)P_B(n+1)\ \succeq_q\ 0.
\]
All factors in the telescoping expansion have nonnegative coefficients; swapping $A,B$ gives the other cross inequality. This applies to $U,U$ and $V,V$ trivially, and to $U,V$ by \eqref{R7E3:cancel}: match common members of $C$, then $1$ with $0$ and $k$ with $k+1$. Expanding
\[
 \srcBMLL^4a_n=q^{14n}\left[
 \left(\frac{P_U(n)}{D_A}+\frac{P_V(n)}{D_B}\right)^2
 -\left(\frac{P_U(n-1)}{D_A}+\frac{P_V(n-1)}{D_B}\right)
  \left(\frac{P_U(n+1)}{D_A}+\frac{P_V(n+1)}{D_B}\right)
 \right]
\]
and multiplying by $D_A^2D_B^2$ expresses the result as four coefficientwise nonnegative differences times positive polynomials.
\end{proof}
Division by a positive-coefficient polynomial does \emph{not} preserve coefficientwise positivity. Proposition~\ref{R7E3:multiplier} therefore does not prove the unmultiplied coefficientwise conjecture. It is a uniform certificate for a weaker statement.

\section{Source audit, limitations, and verification}
The foundational Schur expansion is attributed to \cite{R7E3:MS}, not claimed as new. The October 20, 2026 source check used its arXiv v2 and the published paper, together with Br\"and\'en--Chasse \cite{R7E3:BC}. The latter already proves infinite log-concavity for ordinary Pascal columns ($q=1$). General infinite P\'olya-frequency preservation under $\srcBMLL$ is false: their Example 1.1 gives a negative fourth iterate for a rational generating function. A fixed Gaussian column nevertheless is PF, since its generating function is $\prod_{j=0}^k(1-q^jz)^{-1}$. The no-poles/entire-function preservation theorem and the polynomial-in-$n$ theorem in \cite{R7E3:BC} do not automatically cover this multiple-pole case when $q\ne1$. The checked sources did not yield a theorem implying the present root-location statement. We have not established exhaustive historical priority, including for the weaker positive-real specialization consequence.

The full coefficientwise infinite-log-concavity question is not answered here. Neither the fractional shifts (which depend on $q$) nor reciprocal moment positivity furnishes coefficientwise positivity. The depth-four example obstructs the simplest attempt to iterate the positive-root factorization, not every possible proof of further log-concavity. The multiplier certificate also leaves a genuine division obstruction.

The exact verifier checks the rectangular ratio identities, multiset cancellation, denominator ratio, third-iterate identity, normalized fourth-iterate factorization, cubic discriminant, and rational positive-specialization examples. Its finite tests are diagnostics, not substitutes for the uniform proofs above. Exploratory weighted-fiber searches and a two-parameter Gaussian Newton expansion did not provide additional theorems and are excluded from the claims.



\part{Permanent inequalities and operator stability}

\chapter[Determinant--Pairing Certificates and a Weak-Correlation Region of Chollet's Permanent Conjecture]{Determinant--Pairing Certificates and a Weak-Correlation Region of Chollet's Permanent Conjecture}
\label{source:R1E5}
\noindent\textit{Source development: Round 1, entry 05. Original mathematical arguments retained; notation is local to this chapter.}\par\medskip


\begin{quote}\small\noindent\textbf{Chapter abstract.} 
\noindent\textbf{Status: restricted theorem, not a solution of Chollet's conjecture; priority unverified.}
For a complex correlation matrix $C$, put $r=\max_i\sum_{j\ne i}|c_{ij}|^2$ and $S=\sum_{i<j}|c_{ij}|^4$.
We prove
\[
 \frac{\srcAEper(C\circ\overline C)}{(\srcAEper C)^2}
 \le \exp\!\left[-\frac{2S}{r^2}
 \left(\frac43\log(1+r/2)+\log(1-r)+r\right)\right],\quad 0<r<1.
\]
Consequently the self-conjugate Chollet inequality holds in every order when $r\le r_*=0.6307203397\ldots$, the positive zero of the expression in parentheses. The rational bound $r\le63/100$ is also certified. A determinant--pairing test covers additional matrices outside this uniform region. The argument combines a directly proved two-vector case of Lieb's inequality, a cycle expansion, and the classical matching-polytope theorem. An elementary variant avoiding the latter gives the weaker radius $0.5828116438\ldots$. Exact checks and an explicit dense, full-rank complex example are provided. No global result, optimal radius, or publication-level novelty is claimed.
\end{quote}

\section{The problem and the precise contribution}
For $A=(a_{ij})\in\srcAEC^{n\times n}$,
\[
 \srcAEper A=\sum_{\sigma\in S_n}\prod_{i=1}^n a_{i,\sigma(i)},
 \qquad (A\circ B)_{ij}=a_{ij}b_{ij}.
\]
Chollet's 1982 question \cite{R1E5:chollet} asks whether
\begin{equation}\label{R1E5:eq:chollet}
 \srcAEper(A\circ B)\le(\srcAEper A)(\srcAEper B)
 \qquad(A,B\text{ Hermitian positive semidefinite})
\end{equation}
holds in every order. The equivalent self-conjugate formulation replaces the left side by $\srcAEper(A\circ\overline A)$ and the right side by $(\srcAEper A)^2$.
The conjecture is explicitly open in Zhang's 2016 survey \cite{R1E5:zhang}. Rodtes proved the complex order-four case \cite{R1E5:rodtes}. A 2026 preprint claims the cases through order six \cite{R1E5:li}; another studies specified rank-four classes and expressly does not settle the unrestricted question \cite{R1E5:ma}. These recent claims are contextual, not dependencies of our proofs. A bounded literature search on 20 October 2026 found no unrestricted resolution, but this is not a certification of the current literature's completeness.

A correlation matrix is Hermitian positive semidefinite with all diagonal entries equal to one. Throughout, for such a matrix $C$ we use
\begin{equation}\label{R1E5:eq:parameters}
 x_{ii}=0,\qquad x_{ij}=|c_{ij}|^2\ (i\ne j),\qquad
 X=(x_{ij}),\qquad r=\max_i\sum_jx_{ij},\qquad S=\sum_{i<j}x_{ij}^2.
\end{equation}
Thus $C\circ\overline C=I+X$. A \emph{pairing} is any set of pairwise disjoint unordered pairs of indices; indices may remain unpaired. Write
\[
 L(X)=\max_{M\in\srcAEcalM}\sum_{\{i,j\}\in M}\log(1+x_{ij}).
\]
The actual research contribution is the quantitative all-order sufficient condition below and its certificate formulation. The permanent and polyhedral ingredients are classical; they are not presented as newly discovered theorems. A search did not locate this particular combination or constant, but priority remains unverified.

\begin{srcAEtheorem}\label{R1E5:thm:main}
Let $C$ be a complex correlation matrix, with parameters \eqref{R1E5:eq:parameters}. For $0<r<1$, set
\[
 F(r)=\frac43\log(1+r/2)+\log(1-r)+r.
\]
Then
\begin{equation}\label{R1E5:eq:main}
 \frac{\srcAEper(C\circ\overline C)}{(\srcAEper C)^2}
 \le \exp\left(-\frac{2F(r)}{r^2}S\right).
\end{equation}
The function $F$ has exactly one zero $r_*$ in $(0,1)$. It is positive on $(0,r_*)$, and
\[
 0.6307203397<r_*<0.6307203399.
\]
In particular $r\le r_*$ implies the self-conjugate Chollet inequality. Equality in that inequality within this region is possible only for $C=I$.
\end{srcAEtheorem}

\begin{srcAEcorollary}\label{R1E5:cor:spectral}
The conclusion holds if $\lambda_{\max}(C)\le1+r_*$, and in particular if $\lambda_{\max}(C)\le163/100$.
For two positive semidefinite matrices whose correlation normalizations both have $r\le r_*$, \eqref{R1E5:eq:chollet} holds.
\end{srcAEcorollary}

\section{A permanent lower bound from disjoint pairs}
We give the needed special case of Lieb's block inequality \cite{R1E5:lieb} directly, so that its applicability to complex matrices is explicit.

\begin{srcAElemma}\label{R1E5:lem:pairper}
If $C$ is a complex correlation matrix and $M$ is any pairing, then
\begin{equation}\label{R1E5:eq:pairper}
 \srcAEper C\ge\prod_{\{i,j\}\in M}(1+|c_{ij}|^2).
\end{equation}
\end{srcAElemma}
\begin{proof}
Use the polynomial inner product
$\langle p,q\rangle=\sum_\alpha\alpha!\,\overline{p_\alpha}q_\alpha$.
For a linear form $\ell_u(z)=\sum_k u_kz_k$, multiplication $M_u$ has adjoint
$a_u=\sum_k\overline{u_k}\,\partial_{z_k}$, and
$[a_u,M_v]=u^*v$.
Put $\alpha=u^*u$, $\beta=v^*v$, and $c=u^*v$.
Normal ordering gives, for every polynomial $g$,
\begin{align*}
 \|M_uM_vg\|^2={}&(\alpha\beta+|c|^2)\|g\|^2
 +\|a_ua_vg\|^2\\
 &+\beta\|a_ug\|^2+\alpha\|a_vg\|^2
 +2\operatorname{Re}\!\left(c\langle a_ug,a_vg\rangle\right).
\end{align*}
The last three terms are a nonnegative Hermitian quadratic form because
$\left(\begin{smallmatrix}\beta&c\\\overline c&\alpha\end{smallmatrix}\right)\succeq0$.
Consequently
$\|\ell_u\ell_vg\|^2\ge(\alpha\beta+|c|^2)\|g\|^2$.
The one-vector identity
$\|\ell_ug\|^2=\alpha\|g\|^2+\|a_ug\|^2$
handles singletons.

Take Gram vectors with $c_{ij}=\sum_k u_{ik}\overline{u_{jk}}$.
Expansion of the polynomial norm, or the complex Gaussian moment formula
$\mathbb E z^\alpha\overline z^\beta=\mathbf1_{\alpha=\beta}\alpha!$, gives
\[
 \left\|\prod_i\ell_{u_i}\right\|^2=\srcAEper C.
\]
Apply the two-vector inequality successively to the chosen pairs and the one-vector inequality to the remaining factors. All vector norms equal one, giving \eqref{R1E5:eq:pairper}. This also proves $\srcAEper C\ge1$.
\end{proof}

\section{An upper bound from cycles}
\begin{srcAElemma}\label{R1E5:lem:cycles}
Let $X$ be real, symmetric, entrywise nonnegative, and zero on the diagonal. If its maximum row sum satisfies $0<r<1$, then
\begin{equation}\label{R1E5:eq:cyclebound}
 \srcAEper(I+X)\le\frac1{\det(I-X)}
 \le\exp\left(\frac{2S}{r^2}[-\log(1-r)-r]\right),
 \qquad S=\sum_{i<j}x_{ij}^2.
\end{equation}
The first inequality is strict whenever $X\ne0$. If $X=0$, then $\srcAEper(I+X)=1/\det(I-X)=1$; the displayed expression with $r^{-2}$ is not needed.
\end{srcAElemma}
\begin{proof}
A permutation contributes a collection of vertex-disjoint directed simple cycles of length at least two; unused vertices contribute one. Give a cycle $\gamma$ the weight $w_\gamma=\prod_{(i,j)\in\gamma}x_{ij}$. Dropping the disjointness restriction and allowing arbitrary multiplicities yields
\[
 \srcAEper(I+X)\le\exp\left(\sum_{\gamma\ \mathrm{simple}}w_\gamma\right)
 \le\exp\left(\sum_{k=2}^{\infty}\frac{\srcAEtr X^k}{k}\right).
\]
In the second inequality each directed simple cycle of length $k$ occurs $k$ times in $\srcAEtr X^k$; all remaining closed walks have nonnegative weights. If $x_{ij}>0$, the exponential also contains two copies of the positive two-cycle $i\leftrightarrow j$, which no permutation collection can contain, so the first comparison is strict.

All eigenvalues of $X$ lie in $[-r,r]$. Since $\srcAEtr X=0$, the convergent logarithm series identifies the last exponential with $1/\det(I-X)$. Also, for every integer $k\ge2$,
\[
 \srcAEtr X^k\le\sum_j|\lambda_j(X)|^k
 \le r^{k-2}\srcAEtr X^2=2Sr^{k-2}.
\]
Summing proves \eqref{R1E5:eq:cyclebound}.
\end{proof}

The first consequence is a certificate that does not replace the spectrum of $X$ by its row-sum bound.
\begin{srcAEproposition}[Determinant--pairing certificate]\label{R1E5:prop:certificate}
For $C$ as above with $r<1$, if a pairing $M$ satisfies
\begin{equation}\label{R1E5:eq:certificate}
 \det(I-X)\prod_{\{i,j\}\in M}(1+x_{ij})^2\ge1,
\end{equation}
then $\srcAEper(C\circ\overline C)\le(\srcAEper C)^2$.
\end{srcAEproposition}
\begin{proof}
Lemma \ref{R1E5:lem:cycles} bounds the left side by $1/\det(I-X)$; Lemma \ref{R1E5:lem:pairper} bounds the right side below by the displayed product.
\end{proof}
For rational $X$, \eqref{R1E5:eq:certificate} is an exact rational certificate involving a determinant and a pairing, not a permanent calculation. A supplied pairing is sufficient; no assertion about numerical optimization of logarithmic weights is needed.

\section{A sharp universal pairing estimate}
Here we use a standard polyhedral theorem as an explicitly identified dependency. Edmonds's matching-polytope theorem \cite{R1E5:edmonds} says that the convex hull of pairing incidence vectors consists exactly of vectors $y_{ij}\ge0$ satisfying
\begin{equation}\label{R1E5:eq:edmonds}
 \sum_{j\ne i}y_{ij}\le1,\qquad
 \sum_{i<j,\ i,j\in T}y_{ij}\le\frac{|T|-1}{2}
 \quad\text{for every odd set $T$, $|T|\ge3$.}
\end{equation}
Only this classical theorem, not a new graph-theoretic conjecture, is used.

\begin{srcAElemma}\label{R1E5:lem:sharp-pairing}
For a symmetric nonnegative zero-diagonal $X$ with maximum row sum at most $r$, where $0<r<1$,
\begin{equation}\label{R1E5:eq:sharp-pairing}
 L(X)\ge\frac{4\log(1+r/2)}{3r^2}\sum_{i<j}x_{ij}^2.
\end{equation}
The coefficient is best possible for this assertion about arbitrary such $X$.
\end{srcAElemma}
\begin{proof}
Define $\phi(0)=0$ and, for $t>0$,
\[
 \phi(t)=\frac{t^2}{\log(1+t)}
 =\int_0^1 t(1+t)^s\,ds.
\]
The integrands have second derivative
$s(1+t)^{s-2}[2+(s+1)t]\ge0$, so $\phi$ is convex. It is also increasing and satisfies
$\sum_j\phi(t_j)\le\phi(\sum_jt_j)$ for $t_j\ge0$.
Put $D=3\phi(r/2)$. We have $\phi(r)\le D$, because
\[
 3\log(1+r)-4\log(1+r/2)\ge0\qquad(0\le r\le1);
\]
the expression vanishes at zero and its derivative is
$(1-r/2)/[(1+r)(1+r/2)]\ge0$.

Set $y_{ij}=\phi(x_{ij})/D$. The degree inequalities in \eqref{R1E5:eq:edmonds} follow from convexity and $\phi(r)\le D$. To prove an odd-set inequality, maximize
$\sum_{i<j\in T}\phi(z_{ij})$ over $z\ge0$ with all row sums at most $r$ on $T$. A convex function on this compact polytope has a maximizer at a vertex. The vertices, after scaling by $r$, consist of disjoint edges of weight one and disjoint odd cycles of weight one-half, with possible isolated indices. A proof of this elementary vertex fact is given below.

Each edge contributes $\phi(r)\le D$. An odd cycle of length $m\ge3$ contributes
$m\phi(r/2)\le(m-1)D/2$.
If there are $p$ edges, $q$ odd cycles using a total of $\ell$ indices, and $u$ isolated indices, then $|T|=2p+\ell+u$ and the total is at most
$(|T|-q-u)D/2$. As $|T|$ is odd, $q+u$ is odd and at least one. This proves the required bound $(|T|-1)D/2$.

By \eqref{R1E5:eq:edmonds}, $y$ is a convex combination of pairings. Therefore some pairing has logarithmic weight at least
\[
 \sum_{i<j}y_{ij}\log(1+x_{ij})=\frac1D\sum_{i<j}x_{ij}^2,
\]
which is \eqref{R1E5:eq:sharp-pairing}. For three indices with every off-diagonal entry $r/2$, a pairing has at most one edge, and equality holds.
\end{proof}

\paragraph{Proof of the vertex fact used above.}
Scale $r$ to one and consider the positive support of a vertex. An edge of weight one is an isolated component. In any remaining connected component, every positive edge has weight strictly between zero and one. Let $m$ be its number of edges, $v$ its number of vertices, and $t$ the number of tight degree constraints. Extremality requires $m\le t$: otherwise a nonzero perturbation supported on its edges preserves all tight sums, and sufficiently small perturbations of both signs preserve positivity and all slack inequalities.
If the component is a tree, every leaf is slack, since its sole incident edge has weight below one; hence $t\le v-2<m=v-1$, impossible. If it has more than one independent cycle, $m>v\ge t$, also impossible. The remaining case is unicyclic with $m=v=t$. No leaf can then occur, so it is a cycle. An even cycle admits alternating perturbations and is not extreme. On an odd cycle all degree equations are tight and force every edge weight to be one-half. This proves the fact.

\section{Proofs of the main conclusions}
\begin{proof}[Proof of Theorem \ref{R1E5:thm:main}]
Lemma \ref{R1E5:lem:pairper} and Lemma \ref{R1E5:lem:sharp-pairing} imply
\[
 \log\srcAEper C\ge L(X)\ge\frac{4\log(1+r/2)}{3r^2}S.
\]
Subtract twice this inequality from the logarithm of \eqref{R1E5:eq:cyclebound} to obtain \eqref{R1E5:eq:main}.
Moreover $F(0)=0$, $F'(0)=2/3$, $F(r)\to-\infty$ as $r\uparrow1$, and
\[
 F''(r)=-\frac4{3(2+r)^2}-\frac1{(1-r)^2}<0.
\]
Strict concavity gives exactly one positive zero and the stated sign. The numerical bracket and $F(63/100)>0$ are certified by rational logarithm bounds in the verifier; the theorem itself does not depend on decimal approximations. If $C\ne I$, Lemma \ref{R1E5:lem:cycles} is strict, so equality is impossible even at $r=r_*$. The case $r=0$ is $C=I$.
\end{proof}

\begin{proof}[Proof of Corollary \ref{R1E5:cor:spectral}]
For $\lambda=\lambda_{\max}(C)$, positive semidefiniteness gives
$C^2\preceq\lambda C$. Taking diagonal entries yields
\[
 \sum_{j\ne i}|c_{ij}|^2=(C^2)_{ii}-1\le\lambda-1.
\]
This proves the spectral statement.
For general $A\succeq0$, a zero diagonal entry forces a zero row and column, and both sides of \eqref{R1E5:eq:chollet} vanish. Otherwise write
$A=DCD$, where $D=\srcAEdiag(\sqrt{a_{ii}})$ and $C$ is a correlation matrix. If $h=\prod_i a_{ii}$, then
$\srcAEper A=h\srcAEper C$ and $\srcAEper(A\circ\overline A)=h^2\srcAEper(C\circ\overline C)$.
Finally Cauchy--Schwarz on the permutation monomials gives
\[
 |\srcAEper(A\circ B)|^2\le
 \srcAEper(A\circ\overline A)\srcAEper(B\circ\overline B)
 \le(\srcAEper A)^2(\srcAEper B)^2.
\]
The permanents are nonnegative by the polynomial Gram identity in Lemma \ref{R1E5:lem:pairper}; $A\circ B$ is itself a Gram matrix of tensor products. Taking square roots proves the assertion.
\end{proof}

\subsection*{An elementary fallback without the matching-polytope theorem}
There is an independent, weaker radius. Greedily choose a largest remaining entry $t=x_{ij}$ and remove its two indices. Every removed entry is at most $t$, and its total weight is at most $2r-t$. Thus the sum of removed squares is at most
\[
 t(2r-t)\le2r\log(1+t)\quad(0<r\le1),
\]
since $\log(1+t)\ge t-t^2/2\ge t-t^2/(2r)$.
Summing shows $L(X)\ge S/(2r)$. Consequently
\[
 \frac{\srcAEper(C\circ\overline C)}{(\srcAEper C)^2}
 \le\exp\left[-\frac{S}{r^2}\{3r+2\log(1-r)\}\right].
\]
This proves the conclusion up to the positive zero
$r_0=0.5828116438\ldots$ of $3r+2\log(1-r)$, in particular for $r\le1/2$.
This argument proves less, but isolates the gain supplied by the sharper packing step.

\section{An exact dense complex example}
Let
\[
 K=\begin{pmatrix}
0&1&1&1&1&1&1&1\\
-1&0&1&1&-1&1&-1&-1\\
-1&-1&0&1&1&-1&1&-1\\
-1&-1&-1&0&1&1&-1&1\\
-1&1&-1&-1&0&1&1&-1\\
-1&-1&1&-1&-1&0&1&1\\
-1&1&-1&1&-1&-1&0&1\\
-1&1&1&-1&1&-1&-1&0
\end{pmatrix}.
\]
Direct multiplication gives $K^T=-K$ and $KK^T=7I$. Therefore
$C_t=I+itK$ is a positive definite correlation matrix for $0<t<1/\sqrt7$, with eigenvalues $1\pm t\sqrt7$.
For $t=3/10$, one has $r=63/100$, just beyond the elementary radius and below $r_*$. Exact enumeration gives
\[
 \srcAEper C_{3/10}=\frac{479814841}{100000000},\qquad
 \srcAEper(C_{3/10}\circ\overline{C_{3/10}})
 =\frac{13690485269833393}{10000000000000000}.
\]
The inequality is far from equality; this is a scope check, not evidence of optimality.

The adaptive certificate reaches further. For $t=1/3$, $r=7/9>r_*$ and
$X=(J-I)/9$. A complete pairing has four pairs, and
\[
 \det(I-X)\prod_{\{i,j\}\in M}(1+x_{ij})^2
 =\frac29\left(\frac{10}{9}\right)^{15}>1,
\]
because $2\cdot10^{15}>9^{16}=1853020188851841$.
Proposition \ref{R1E5:prop:certificate} proves self-Chollet for this matrix without evaluating either permanent. Both examples are dense and full rank; no support restriction or rank-two reduction is being used.

\section{Limitations, unsuccessful routes, and next questions}
\paragraph{What has not been proved.}
Nothing here controls all correlation matrices, settles a new entire matrix order, or proves the global conjecture. The radius $r_*$ is a limit of the particular estimates, not a conjectured sharp boundary. The spectral corollary is sufficient, not necessary. The pairing lemma is sharp for unrestricted nonnegative $X$, but its worst example need not simultaneously make the cycle estimate sharp or arise from the extremal correlation matrices of interest.

\paragraph{A rigorous obstruction to extending the certificate blindly.}
Even the determinant--pairing certificate cannot cover the whole range $r<1$. For $C_t=I+itK$ above, its best left side is
\[
 (1-7t^2)(1+t^2)^{15},
\]
which tends to zero as $t\uparrow1/\sqrt7$. All pairings have at most four edges, and every edge has weight $t^2$, so no different pairing repairs this failure. These are admissible matrices throughout the interval. Failure of a sufficient certificate is not failure of Chollet's inequality.

\paragraph{Alternatives investigated.}
A direct expansion around $I$ leaves signed or complex cycle contributions in $\srcAEper C$; replacing them by absolute values does not provide a lower bound. A rank-one-plus-diagonal approach falls within known simple structured mechanisms and was not promoted as new. The elementary greedy pairing argument is valid but loses the interval $(r_0,r_*]$; the convex packing argument was developed to improve it. No counterexample to Chollet's conjecture was found or claimed.

\paragraph{Next avenues, not undertaken as another round.}
The promising remaining issue is to couple the lower and upper estimates instead of optimizing them separately. Larger principal blocks could replace pairs, and the determinant relaxation could retain the exclusion of repeated cycles. Another sharply posed task is to determine the largest universal $r$ for which some pairing always satisfies \eqref{R1E5:eq:certificate}. Those questions concern this method; they should not be confused with equivalent forms of the unrestricted conjecture. A fuller priority search and independent scrutiny of the packing reduction are prerequisites to a publishable novelty claim.

\paragraph{Reproducibility and assessment.}
The standalone file \texttt{verify.py} uses integer and rational arithmetic to verify $K$, both displayed permanent values, the second certificate, and rigorous logarithm brackets for $r_*$. Its deterministic finite matrix checks are sanity checks, not proofs of the general statement. The main theorem is analytic, conditional only on the stated classical matching-polytope theorem, not on a grid search. The appropriate assessment is a modest all-order sufficient-condition result with an explicit proof, not an Annals-level or Fields-level breakthrough.



\chapter[A strong permanent inequality for complex positive semidefinite cactus matrices]{A strong permanent inequality for complex positive semidefinite cactus matrices}
\label{source:R2E3}
\noindent\textit{Source development: Round 2, entry 03. Original mathematical arguments retained; notation is local to this chapter.}\par\medskip
\begin{quote}\small\textbf{Joint-credit note.}
The complex-cactus conclusion was independently obtained in entries 03
and 05 of Round 2. Both proof developments are retained, but the shared
principal theorem is counted once.\end{quote}


\begin{quote}\small\noindent\textbf{Chapter abstract.} 
For a Hermitian positive semidefinite matrix $A$ whose off-diagonal support is a cactus, we prove
\[
 \srcAIper(A\circ\overline A)\leq \srcAIper(A)\prod_i a_{ii}.
\]
Consequently Chollet's permanent inequality holds for every pair of such matrices, without requiring their support graphs to agree. The new local estimate allows arbitrary complex phases and nonuniform weights on a cycle: positivity of three-by-three path minors pays for the two directed cycle terms using only single-edge matchings. An elementary Gram-space decomposition converts arbitrary positive semidefinite matrices with cut vertices into one-point sums. We combine this with the previously established one-point-sum method of Pant and Singh, whose cactus result concerns graph Laplacians, and the classical order-three permanent-on-top theorem. The result is a structural restricted theorem, not a solution of Chollet's conjecture. No computational certificate is required for the proof.
\end{quote}

\section{Statement, scope, and provenance}
Write $A\succeq0$ for Hermitian positive semidefiniteness, $A\circ B$ for the entrywise product, and
\[
 P(A)=\srcAIper(A),\qquad \srcAIQ(A)=\srcAIper(A\circ\overline A),\qquad
 \srcAIh(A)=\prod_i a_{ii}.
\]
The permanent of the empty matrix and its diagonal product are both $1$.
The support graph of $A$ has edge $ij$ precisely when $i\ne j$ and $a_{ij}\ne0$.
A cactus is a graph in which every block is an edge or a simple cycle; disconnected graphs and isolated vertices are allowed. Equivalently, any two distinct simple cycles have at most one vertex in common.

Chollet's conjecture \cite{R2E3:Chollet} asserts that
\begin{equation}\label{R2E3:eq:chollet}
 P(A\circ B)\leq P(A)P(B)\qquad(A,B\succeq0).
\end{equation}
It remains open in general. The inequality $\srcAIQ(A)\leq\srcAIh(A)P(A)$ is stronger than its self-conjugate specialization, because $P(A)\geq\srcAIh(A)$ by Lieb's inequality \cite{R2E3:Lieb}.

\begin{srcAItheorem}\label{R2E3:thm:main}
If $A\succeq0$ and its support graph is a cactus, then
\begin{equation}\label{R2E3:eq:strong}
 \srcAIQ(A)\leq\srcAIh(A)P(A).
\end{equation}
Every principal submatrix of $A$ satisfies the same inequality.
If $A,B\succeq0$ have the same order and each has cactus support, then
\begin{equation}\label{R2E3:eq:mixed}
 0\leq P(A\circ B)
 \leq \sqrt{\srcAIh(A)\srcAIh(B)P(A)P(B)}
 \leq P(A)P(B).
\end{equation}
The two supports need not coincide.
\end{srcAItheorem}

\paragraph{What is and is not new.}
Pant and Singh's latest version \cite{R2E3:PS}, checked as version 2 of 18 September 2026, proves one-point-sum closure and establishes the strong inequality for cactus graph Laplacians. We do not claim those operations or the Laplacian subclass as new. Our additional assertion is the unrestricted complex, weighted PSD cactus class. Its essential new ingredient is Theorem~\ref{R2E3:thm:cycle}, followed by a support-level use of Gram-space splitting. The complex strong closure proof below is the same algebra as the real version of \cite[Theorem 4.3]{R2E3:PS}, included to identify exactly what is reused.

Round~1 entry~05's weak-correlation approach to Chollet motivated the change of topic, but none of its new bounds or matching-polytope certificates is used here. Our own Round~1 Gaussian-binomial results are not used. All six frozen Round~1 submissions were consulted; no result from another live Round~2 submission is used. Novelty beyond the sources checked remains subject to a wider literature review.

\section{The weighted cycle estimate}
A correlation matrix is a PSD matrix with all diagonal entries equal to $1$. Diagonal scaling reduces \eqref{R2E3:eq:strong} to this case when all diagonal entries are positive: if $A=DCD$ with $D$ positive diagonal, then $P(A)=\srcAIh(A)P(C)$ and $\srcAIQ(A)=\srcAIh(A)^2\srcAIQ(C)$. If $a_{ii}=0$, PSD forces the entire $i$th row to vanish, and both sides of \eqref{R2E3:eq:strong} are zero.

For $m\geq4$, label a cycle by $1,\ldots,m$ cyclically. For nonnegative edge weights $x_1,\ldots,x_m$, let
\[
 M(x)=\sum_{F\text{ a matching of }C_m}\prod_{i\in F}x_i.
\]
Zeros are permitted, so a matrix supported on a subgraph of the cycle is included.

\begin{srcAItheorem}\label{R2E3:thm:cycle}
Let $C$ be Hermitian, with diagonal $1$ and off-diagonal support contained in $C_m$, $m\geq4$. Put $x_i=|c_{i,i+1}|^2$ and $p=\sqrt{\prod_i x_i}$. Suppose
\begin{equation}\label{R2E3:eq:local}
 x_i+x_{i+1}\leq1\quad\text{for every }i.
\end{equation}
Then
\begin{equation}\label{R2E3:eq:cyclebound}
 P(C)-\srcAIQ(C)\geq m p^{4/m}-2p-2p^2
 \geq2p(1-p)\geq0.
\end{equation}
In particular, this holds for every PSD correlation matrix supported on $C_m$.
\end{srcAItheorem}
\begin{proof}
The only permutation cycles of length at least three allowed by $C_m$ are its two full orientations. If $z=\prod_i c_{i,i+1}$, then $|z|=p$, and the permanent expansions are
\begin{equation}\label{R2E3:eq:matching}
 P(C)=M(x)+2\operatorname{Re}z,\qquad
 \srcAIQ(C)=M(x_1^2,\ldots,x_m^2)+2p^2.
\end{equation}
By \eqref{R2E3:eq:local}, $0\leq x_i\leq1$. Every matching contributes nonnegatively to $M(x)-M(x^2)$; retaining only one-edge matchings gives
\begin{align*}
 P(C)-\srcAIQ(C)
 &\geq\sum_i x_i(1-x_i)-2p-2p^2\\
 &\geq\sum_i x_ix_{i+1}-2p-2p^2\\
 &\geq m\left(\prod_i x_ix_{i+1}\right)^{1/m}-2p-2p^2\\
 &=m p^{4/m}-2p-2p^2.
\end{align*}
Here the second line uses $1-x_i\geq x_{i+1}$, and the third is arithmetic--geometric mean. Since $p\leq1$ and $m\geq4$, we have $m p^{4/m}\geq4p$. This proves \eqref{R2E3:eq:cyclebound}, including $p=0$.

For a PSD correlation matrix, the three consecutive vertices $i,i+1,i+2$ induce a path (also when $m=4$). Their principal determinant is $1-x_i-x_{i+1}\geq0$, proving \eqref{R2E3:eq:local}.
\end{proof}

\begin{srcAIremark}
The theorem only needs the displayed local constraints, not global PSD. They are substantially weaker. For an actual cycle with all edges nonzero, multiplying $x_i+x_{i+1}\leq1$ yields $p\leq2^{-m/2}<1$, so \eqref{R2E3:eq:cyclebound} is strict. Complex phase cancellation is handled uniformly by $\operatorname{Re}z\geq-p$; no diagonal-unitary transformation to a nonnegative matrix is assumed.
\end{srcAIremark}

\section{Small blocks and the inherited closure mechanism}
We use two classical facts with their precise roles stated. Lieb's inequality \cite{R2E3:Lieb} gives, for a PSD matrix $T$ and index $v$,
\begin{equation}\label{R2E3:eq:lieb}
 P(T)\geq t_{vv}P(T-v)\geq0.
\end{equation}
Bapat and Sunder \cite{R2E3:BS} prove the permanent-on-top property through order three. The following consequence supplies triangle blocks; it is not a new result.

\begin{srcAIlemma}\label{R2E3:lem:three}
If $A\succeq0$ has order at most three, then $\srcAIQ(A)\leq\srcAIh(A)P(A)$.
\end{srcAIlemma}
\begin{proof}
It suffices to take a correlation matrix $C$, the zero-diagonal case being immediate. Form its Schur power matrix, indexed by permutations,
\[
 S_{\sigma,\tau}=\prod_i c_{\sigma(i),\tau(i)}.
\]
This is a Gram matrix of permuted tensor products and is PSD. Its diagonal entries are $1$; its row sums are $P(C)$. The order-three permanent-on-top theorem says $\lambda_{\max}(S)=P(C)$. Thus $S^2\preceq P(C)S$. On any diagonal entry, $(S^2)_{\sigma,\sigma}=\srcAIQ(C)$, giving $\srcAIQ(C)\leq P(C)$.
\end{proof}

For distinguished indices in
\[
 A_i=\begin{pmatrix}B_i&u_i\\u_i^*&d_i\end{pmatrix}\succeq0,
\]
their one-point sum is
\[
 A_1\star A_2=
 \begin{pmatrix}B_1&u_1&0\\u_1^*&d_1+d_2&u_2^*\\0&u_2&B_2\end{pmatrix}.
\]
It is PSD, being a sum of two embedded PSD matrices. The next lemma is the complex version of the strong closure calculation in \cite{R2E3:PS}; its proof uses no real-entry hypothesis.

\begin{srcAIlemma}[Strong one-point closure]\label{R2E3:lem:closure}
If $A_1,A_2,B_1,B_2$ satisfy \eqref{R2E3:eq:strong}, then so does $A_1\star A_2$.
\end{srcAIlemma}
\begin{proof}
Put $P_i=P(A_i)$, $R_i=P(B_i)$, and $h_i=\srcAIh(B_i)$. A permutation cycle cannot traverse both sides of the cut vertex. Expansion at that vertex gives
\begin{align*}
 P(A_1\star A_2)&=P_1R_2+R_1P_2,\\
 \srcAIQ(A_1\star A_2)&=\srcAIQ(A_1)\srcAIQ(B_2)+\srcAIQ(B_1)\srcAIQ(A_2)
                       +2d_1d_2\srcAIQ(B_1)\srcAIQ(B_2).
\end{align*}
The hypotheses bound the second line by
\[
 h_1h_2\bigl(d_1P_1R_2+d_2R_1P_2+2d_1d_2R_1R_2\bigr).
\]
Subtracting this from $\srcAIh(A_1\star A_2)P(A_1\star A_2)$ leaves
\[
 h_1h_2\bigl[d_2R_2(P_1-d_1R_1)+d_1R_1(P_2-d_2R_2)\bigr]\geq0
\]
by \eqref{R2E3:eq:lieb}. All quantities used as multiplicative bounds are nonnegative. No division is used, so singular summands are allowed.
\end{proof}

\section{From a support decomposition to PSD summands}
A matrix with a cut vertex is not generally the one-point sum of its two principal submatrices: that would count its cut-vertex diagonal twice. The following standard Gram-space argument gives the required split without assuming diagonal dominance or a Laplacian form.

\begin{srcAIlemma}[PSD splitting at a separator of size one]\label{R2E3:lem:split}
Suppose $A\succeq0$ is indexed by $U\cup\{v\}\cup W$, with $U,W$ disjoint and nonempty, and $A[U,W]=0$. Then $A=A_1\star A_2$ for PSD matrices $A_1,A_2$ indexed by $U\cup\{v\}$ and $W\cup\{v\}$ respectively. Their entries agree with the corresponding entries of $A$ except at $v,v$, where nonnegative entries $d_1,d_2$ sum to $a_{vv}$.
\end{srcAIlemma}
\begin{proof}
Realize $A$ as a Gram matrix of vectors $g_i$. The spans $H_U$ and $H_W$ are orthogonal because all cross inner products vanish. Decompose
\[
 g_v=u+w+r,\qquad u\in H_U,\quad w\in H_W,\quad
 r\perp(H_U+H_W).
\]
For $A_1$ use the vectors $g_i$ for $i\in U$, with $u$ at $v$; for $A_2$ use $g_j$ for $j\in W$, with $w+r$ at $v$. All required inner products are unchanged, and
$d_1=\|u\|^2$, $d_2=\|w+r\|^2$ satisfy $d_1+d_2=\|g_v\|^2$.
\end{proof}

\begin{proof}[Proof of Theorem~\ref{R2E3:thm:main}]
Prove \eqref{R2E3:eq:strong} by induction on the order. Disconnected support is handled by multiplicativity on direct sums. A connected cactus with no cut vertex is a single vertex, an edge, or a cycle. Lemma~\ref{R2E3:lem:three} covers the first two cases and triangles. For longer cycles, diagonal normalization and Theorem~\ref{R2E3:thm:cycle} apply.

Otherwise choose a cut vertex $v$ and partition the components of its deletion into nonempty sets $U,W$. Lemma~\ref{R2E3:lem:split} gives $A=A_1\star A_2$. Both summands have smaller order and cactus support, as do their distinguished deletions $B_1,B_2$. The induction hypothesis applies to all four. Lemma~\ref{R2E3:lem:closure} proves \eqref{R2E3:eq:strong} for $A$. Every principal submatrix has cactus support, proving the hereditary assertion.

For arbitrary same-order matrices, Cauchy--Schwarz on permutation monomials gives
\[
 |P(A\circ B)|^2\leq\srcAIQ(A)\srcAIQ(B).
\]
Apply \eqref{R2E3:eq:strong} twice. The Schur product theorem and nonnegativity of PSD permanents give $P(A\circ B)\geq0$; iterating \eqref{R2E3:eq:lieb} gives $P(A)\geq\srcAIh(A)$ and $P(B)\geq\srcAIh(B)$. These prove \eqref{R2E3:eq:mixed}.
\end{proof}

\section{Limits, exact checks, and further directions}
The theorem is independent of order and of a weak-correlation threshold. For instance, begin with a diagonal-one $3\times3$ all-ones block and attach arbitrarily many weighted complex cycle blocks through one-point sums. The normalized row squared-correlation sum can exceed the Round~1 entry~05 threshold $0.6308$, already because two entries of a triangle can be arbitrarily close to modulus $1$. Thus this is not a repackaging of that sufficient region. Conversely, dense matrices in the weak-correlation region need not have cactus support, so neither theorem subsumes the other.

The cycle estimate fails to directly address chorded cycles: a three-vertex path may acquire a chord, and additional long permutation cycles appear. A separator with two vertices also permits a permutation cycle to cross between both sides; the one-point identities no longer apply. These are genuine obstacles to extending the method, not missing finite cases.

The accompanying \texttt{verify.py} uses exact Gaussian-rational arithmetic to check the cycle expansions, their inequalities on a finite rational grid satisfying the local constraints, the one-point formulas on complex PSD examples, and every principal submatrix of selected cactus Gram matrices. These tests are diagnostics only; all-order validity follows from the proof. Triangle blocks depend on the cited classical theorem, not on numerical testing.

The literature check included the latest Pant--Singh version and Bapat--Sunder's primary paper. It did not locate the arbitrary weighted complex cactus statement, but this is not a certification of priority. The useful next questions are a direct local estimate for a theta block, a two-vertex-separator replacement for strong closure, and whether one cactus-supported factor is sufficient when the other PSD factor is unrestricted. None is proved here.



\chapter[Strong Self-Chollet Inequalities for Complex Cycle Blocks and Positive Semidefinite Matrices with Cut Vertices]{Strong Self-Chollet Inequalities for Complex Cycle Blocks 
and Positive Semidefinite Matrices with Cut Vertices}
\label{source:R2E5}
\noindent\textit{Source development: Round 2, entry 05. Original mathematical arguments retained; notation is local to this chapter.}\par\medskip
\begin{quote}\small\textbf{Joint-credit note.}
The complex-cactus conclusion was independently obtained in entries 03
and 05 of Round 2. Both proof developments are retained, but the shared
principal theorem is counted once.\end{quote}


\begin{quote}\small\noindent\textbf{Chapter abstract.} 
\noindent\textbf{Status: a structural partial result; neither the general Chollet conjecture nor historical priority is settled.}
We prove
\[
 \srcAKper(A\circ\overline A)\le \srcAKper(A)\prod_i a_{ii}
\]
for every complex Hermitian positive semidefinite matrix whose support is a cactus, with arbitrary edge weights and phases. More generally, the conclusion holds when every support block is either a cycle or can be made entrywise nonnegative by a diagonal unitary congruence. A quantitative estimate on one cycle controls its potentially negative phase contribution using only local three-by-three positivity. Gram projections then split an arbitrary positive semidefinite matrix at a cut vertex, allowing a known one-point-sum mechanism to apply. This strengthens the Round 1 conclusion on a different, unbounded-order class and covers full-rank examples with unbounded squared-correlation row sums. The Laplacian cactus theorem and one-point-sum framework are prior work of Pant--Singh; the proposed increment is the arbitrary complex weighted-cycle argument and its support-level extension. All analytic steps are written out, with exact finite checks supplied separately.
\end{quote}

\section{Target, source versions, and the precise increment}
For a Hermitian positive semidefinite matrix $A$, write
\[
 P(A)=\srcAKper A,\qquad Q(A)=\srcAKper(A\circ\overline A),
 \qquad h(A)=\prod_i a_{ii}.
\]
Empty matrices have $P=Q=h=1$. Chollet's 1982 conjecture \cite{R1E5:chollet} is
\begin{equation}\label{R2E5:eq:chollet}
 \srcAKper(A\circ B)\le P(A)P(B)\qquad(A,B\succeq0).
\end{equation}
It is equivalent to $Q(A)\le P(A)^2$; see \cite{R2E5:zhang}.
The stronger self-conjugate inequality
\begin{equation}\label{R2E5:eq:strong}
 Q(A)\le h(A)P(A)
\end{equation}
is false without structural restrictions: Drury's order-seven correlation
example has $P(A)=45$ and $Q(A)=6185/128>45$ \cite{R2E5:drury,R2E5:zhang}.
Thus the present result is a genuine sufficient class, not a proposed
universal strengthening.

The latest inspected version of Pant--Singh is version 2, dated
18 September 2026 \cite{R2E3:PS}. It proves complex one-point-sum closure for
the ordinary self-Chollet inequality, hereditary strong closure for real
matrices, and \eqref{R2E5:eq:strong} for cactus graph Laplacians. These results
and their cycle-cover method are explicitly credited, not claimed anew.
Our cycle matrices need not be Laplacians, real, or diagonally dominant.
The source calls the unrestricted conjecture open; current-source searches
on 20 October 2026 did not find a general resolution, but do not establish
complete bibliographic coverage.

Round 1, Entry 05 \cite{R2E5:R1}, proved the ordinary self-Chollet inequality
under the dimension-free sufficient condition
$\max_i\sum_{j\ne i}|c_{ij}|^2\le0.6307203397\ldots$
for correlation matrices, and supplied a determinant--pairing certificate
requiring that row sum to be below one. Neither estimate is used to prove
the present theorem. The standard Gram permanent identity and elementary
permutation Cauchy--Schwarz transfer are reused with attribution.
The other five Round 1 manuscripts were read; no theorem from them is a
dependency here.

\begin{srcAKdefinition}
The support of $A$ is the simple graph with edge $ij$ when $i\ne j$ and
$a_{ij}\ne0$. Its blocks are its maximal two-connected subgraphs and its
bridge edges. A connected graph is a \emph{cactus} when every block is an
edge or a simple cycle.
An induced matrix on a block is called \emph{phase-balanced} if a diagonal
unitary $U$ makes $U^*AU$ entrywise nonnegative on that block.
This condition concerns the off-diagonal phases and is unaffected by
changing nonnegative diagonal entries.
\end{srcAKdefinition}

\begin{srcAKtheorem}\label{R2E5:thm:main}
Let $A\succeq0$ be complex Hermitian. Suppose that each block of its support
is a simple cycle or is phase-balanced. Then every principal submatrix
$D$ of $A$ satisfies
\[
 Q(D)\le h(D)P(D).
\]
In particular this holds for every positive semidefinite matrix with
cactus support, with no restriction on its order or edge phases.
\end{srcAKtheorem}

\begin{srcAKcorollary}\label{R2E5:cor:pair}
If $A,B$ of the same order both satisfy the hypothesis of Theorem~\ref{R2E5:thm:main}, their
supports need not agree, and
\[
 \srcAKper(A\circ B)\le
 \sqrt{h(A)h(B)P(A)P(B)}\le P(A)P(B).
\]
If $A$ has positive diagonal, equality in $Q(A)\le P(A)^2$ is possible
only when $A$ is diagonal.
\end{srcAKcorollary}

\section{Elementary permanent facts}
The following standard facts are included to keep the proof independent
of unverified permanent conjectures.

\begin{srcAKlemma}\label{R2E5:lem:lieb}
For $A\succeq0$, $P(A)\ge h(A)\ge0$. More particularly,
\[
 P(A)\ge a_{ii}P(A(i)),
\]
where $A(i)$ deletes row and column $i$. If $A$ has positive diagonal and
some nonzero off-diagonal entry, then $P(A)>h(A)$.
\end{srcAKlemma}
\begin{proof}
For polynomials use
$\langle p,q\rangle=\sum_\alpha\alpha!\,\overline{p_\alpha}q_\alpha$.
Multiplication by $\ell_u(z)=\sum_k u_kz_k$ has adjoint
$a_u=\sum_k\overline{u_k}\partial_k$, so
\[
 \|\ell_u g\|^2=\|u\|^2\|g\|^2+\|a_ug\|^2.
\]
Choose Gram vectors with $a_{ij}=\sum_k u_{ik}\overline{u_{jk}}$.
Expansion gives $\|\prod_i\ell_{u_i}\|^2=P(A)$.
Applying the displayed identity once proves the deletion inequality;
iteration proves $P(A)\ge h(A)$. These are classical special cases of
Lieb's inequality \cite{R1E5:lieb}.

For strictness, choose $u_i$ with nonzero inner product with another
vector. Then $a_{u_i}\prod_{j\ne i}\ell_{u_j}$ is not zero: along the
line $z+t\overline{u_i}$, at least one factor is nonconstant in $t$, and
the product of its nonzero linear factors has positive degree in $t$.
The remaining factors are nonzero polynomials because their diagonal
entries are positive. Thus the first inequality is strict for that
deletion, and the remaining deletions give $P(A)>h(A)$.
\end{proof}

Positive diagonal congruence reduces \eqref{R2E5:eq:strong} to correlation
matrices: if $A=DCD$ and $h=\prod_i a_{ii}$, then
$P(A)=hP(C)$ and $Q(A)=h^2Q(C)$. A zero diagonal entry of a PSD matrix
forces its row and column to vanish, in which case $P=Q=h=0$.
A diagonal unitary congruence leaves $P,Q,h$ unchanged.

\begin{srcAKlemma}\label{R2E5:lem:balanced}
Every phase-balanced PSD matrix satisfies \eqref{R2E5:eq:strong}, as do all
its principal submatrices.
\end{srcAKlemma}
\begin{proof}
After the preceding reductions it is enough to consider an entrywise
nonnegative correlation matrix $C$. Each permanent monomial $w_\sigma$
belongs to $[0,1]$. Thus
$Q(C)=\sum_\sigma w_\sigma^2\le\sum_\sigma w_\sigma=P(C)$.
\end{proof}

\section{The cycle estimate with arbitrary complex phases}
\begin{srcAKtheorem}\label{R2E5:thm:cycle}
Let $C$ be a correlation matrix supported on a simple cycle of length
$m\ge4$. Number its edges cyclically, put
$x_i=|c_{i,i+1}|^2$, and let $q=(\prod_{i=1}^m x_i)^{1/2}$.
Then
\begin{equation}\label{R2E5:eq:cyclemargin}
 P(C)-Q(C)\ge
 \sum_i x_i(1-x_i)-2q(1+q)
 \ge 2q(1-q)\ge0.
\end{equation}
If all cycle edges are present, the last positive quantity is strictly
positive. The result remains valid if some edges vanish.
\end{srcAKtheorem}
\begin{proof}
Write $z=\prod_i c_{i,i+1}$, where subscripts are cyclic, so $|z|=q$.
Let
\[
 Z(x)=\sum_{M\ {\rm matching}}\prod_{i\in M}x_i
\]
be the weighted matching sum of the cycle. Its cycle-cover expansion gives
\[
 P(C)=Z(x)+2\Re z,\qquad Q(C)=Z(x^2)+2q^2.
\]
Every $x_i$ lies in $[0,1]$. Termwise comparison, retaining the one-edge
matchings, therefore yields
\[
 P(C)-Q(C)\ge Z(x)-Z(x^2)-2q-2q^2
 \ge\sum_i x_i(1-x_i)-2q-2q^2.
\]

Here positivity of the matrix supplies a local fact absent from an
arbitrary weighted cycle. Three consecutive vertices induce a path,
and the determinant of their principal correlation matrix is
$1-x_i-x_{i+1}$. Hence
\begin{equation}\label{R2E5:eq:local}
 x_i+x_{i+1}\le1\qquad\text{for every }i.
\end{equation}
Consequently
\[
 \sum_i x_i(1-x_i)\ge\sum_i x_ix_{i+1}
 \ge m\left(\prod_i x_ix_{i+1}\right)^{1/m}
 =m q^{4/m}\ge4q.
\]
The last step uses $m\ge4$ and $0\le q\le1$. This proves
\eqref{R2E5:eq:cyclemargin}. Summing \eqref{R2E5:eq:local} and applying
arithmetic--geometric mean gives $q\le2^{-m/2}<1$.
If all edges are nonzero then $q>0$, proving strictness.
\end{proof}

Triangles need a different argument because their three consecutive
vertices do not induce a path. Their strong inequality is classical,
following already from Bapat--Sunder's order-three result \cite{R2E5:BS}.
For completeness, here is a direct scalar proof.

\begin{srcAKlemma}\label{R2E5:lem:triangle}
Every correlation matrix of order three satisfies $Q(C)\le P(C)$.
\end{srcAKlemma}
\begin{proof}
Put $x=|c_{12}|^2$, $y=|c_{23}|^2$, $z=|c_{31}|^2$,
$s=x+y+z$, $t=xy+xz+yz$, $q=\sqrt{xyz}$ and
$p=\Re(c_{12}c_{23}c_{31})$.
Then $0\le x,y,z\le1$, $|p|\le q$, and $\det C\ge0$ gives
$2p\ge s-1$. Direct expansion yields
\begin{equation}\label{R2E5:eq:tri}
 P(C)-Q(C)=s-s^2+2t+2p-2q^2.
\end{equation}
If $q=0$, then $p=0$ and this is $\sum x(1-x)\ge0$.
Assume $q>0$.

If $s+2q\ge1$, feasibility also gives $s-1\le2q$, so
$|s-1|\le2q$. Using $2p\ge s-1$ and $t\ge3q^2$
(each of $xy,xz,yz$ is at least $xyz$) in \eqref{R2E5:eq:tri} gives
\[
 P(C)-Q(C)\ge-(s-1)^2+2t-2q^2\ge0.
\]

If $s+2q\le1$, set $u=q^{2/3}$. Arithmetic--geometric mean gives
$s\ge3u$ and $t\ge3u^2$. Since $3u+2u^{3/2}\le1$,
we have $u\le1/4$, $q\le u/2$ and $q^2=u^3\le u^2/4$.
The concavity of $s(1-s)$ on $[3u,1-2q]$ implies
\[
 s(1-s)+6u^2-2q-2q^2
 \ge\min\{3u-3u^2-2q-2q^2,\ 6u^2-6q^2\}\ge0.
\]
Indeed the first term is at least $2u-\tfrac72u^2>0$, and
the second is $6u^2(1-u)>0$. Now use $p\ge-q$ in
\eqref{R2E5:eq:tri}.
\end{proof}

An edge is phase-balanced. Every proper induced subgraph of a simple
cycle is a disjoint union of paths and isolated vertices, hence can be
made phase-balanced by choosing vertex phases successively along its
edges. Therefore all cycle-supported PSD matrices are hereditarily
strong, including their proper principal submatrices.

\section{From component matrices to an arbitrary cut vertex}
We give the relevant one-point-sum calculation from \cite{R2E3:PS} in complex
notation. Its proof uses no reality assumption, so the hereditary strong
closure in that source extends unchanged to Hermitian matrices with
$Q(A)=\srcAKper(A\circ\overline A)$.

\begin{srcAKlemma}[Strong gluing]\label{R2E5:lem:glue}
Let
\[
 A_1=\begin{pmatrix}B_1&u\\u^*&d_1\end{pmatrix}\succeq0,
 \quad
 A_2=\begin{pmatrix}d_2&v^*\\v&B_2\end{pmatrix}\succeq0,
 \quad
 A=\begin{pmatrix}
 B_1&u&0\\ u^*&d_1+d_2&v^*\\0&v&B_2
 \end{pmatrix}.
\]
If \eqref{R2E5:eq:strong} holds for $A_1,A_2,B_1,B_2$, it holds for $A$.
\end{srcAKlemma}
\begin{proof}
Set $P_i=P(A_i)$, $p_i=P(B_i)$, $Q_i=Q(A_i)$,
$q_i=Q(B_i)$, and $h_i=h(B_i)$.
A permutation cycle cannot cross between the two sides of the common
vertex. Splitting its diagonal contribution into $d_1+d_2$ gives
\[
 P(A)=P_1p_2+p_1P_2,\qquad
 Q(A)=Q_1q_2+q_1Q_2+2d_1d_2q_1q_2.
\]
The four assumed bounds imply
\[
 Q(A)\le h_1h_2
 \bigl(d_1P_1p_2+d_2p_1P_2+2d_1d_2p_1p_2\bigr).
\]
Subtract this from $h(A)P(A)=(d_1+d_2)h_1h_2(P_1p_2+p_1P_2)$.
The difference is
\[
 h_1h_2\{d_2p_2(P_1-d_1p_1)+d_1p_1(P_2-d_2p_2)\}\ge0
\]
by Lemma~\ref{R2E5:lem:lieb}. All quantities used are nonnegative.
\end{proof}

This does not yet justify applying the result to an arbitrary PSD matrix
with a cut vertex: its shared diagonal must actually be split into two
PSD pieces. The next standard Gram-projection observation supplies that
step, including singular matrices.

\begin{srcAKlemma}[PSD splitting at a cut vertex]\label{R2E5:lem:split}
Suppose
\[
 A=\begin{pmatrix}B_1&u&0\\u^*&a&v^*\\0&v&B_2\end{pmatrix}\succeq0.
\]
There exist $d_1,d_2\ge0$ with $d_1+d_2=a$ for which the matrices
$A_1,A_2$ in Lemma~\ref{R2E5:lem:glue} are PSD.
Only the shared diagonal changes; all off-diagonal entries are retained.
\end{srcAKlemma}
\begin{proof}
Realize $A$ by Gram vectors. Let $V_1,V_2$ be the spans of the vectors
belonging to $B_1,B_2$, and let $w$ be the vector at the cut vertex.
The zero corner block implies $V_1\perp V_2$.
Write $w=w_1+w_2+w_0$ for its orthogonal projections on
$V_1,V_2,(V_1+V_2)^\perp$, respectively.
Use $w_1+w_0$ as the distinguished vector for the first piece and $w_2$
for the second. Their squared norms are $d_1,d_2$, with sum $a$,
and their inner products with all vectors on their own sides agree with
the original ones. The resulting two Gram matrices are the desired PSD
pieces. No inverse or positive-definiteness assumption was used.
\end{proof}

\begin{proof}[Proof of Theorem~\ref{R2E5:thm:main}]
First prove \eqref{R2E5:eq:strong} by induction on the matrix order.
The one-vertex case and zero diagonal case are immediate.
For disconnected support, use multiplicativity of $P,Q,h$ on direct sums.
If a connected support has no cut vertex, it is a single block.
Its conclusion is Lemma~\ref{R2E5:lem:balanced}, Theorem~\ref{R2E5:thm:cycle},
or Lemma~\ref{R2E5:lem:triangle}.

Otherwise choose a cut vertex and partition the components left after
its deletion into two nonempty groups. Lemma~\ref{R2E5:lem:split} writes
$A=A_1\star A_2$ with smaller PSD matrices whose off-diagonal phases and
support blocks are unchanged. Each piece and each distinguished deletion
therefore satisfies the same structural hypothesis. Induction applies to
all four matrices, and Lemma~\ref{R2E5:lem:glue} completes the step.

Finally the hypothesis is hereditary under vertex deletion. Any block of
an induced subgraph lies within an original block. Inside a phase-balanced
block it remains phase-balanced; inside a cycle, a proper induced subgraph
is a union of paths. Thus the argument applies to every principal
submatrix.
\end{proof}

\begin{proof}[Proof of Corollary~\ref{R2E5:cor:pair}]
Permutation-monomial Cauchy--Schwarz and Theorem~\ref{R2E5:thm:main} give
\[
 |\srcAKper(A\circ B)|^2\le Q(A)Q(B)
 \le h(A)h(B)P(A)P(B).
\]
The Hadamard product is a Gram matrix of tensor products, so its permanent
is nonnegative. Lemma~\ref{R2E5:lem:lieb} gives the second bound.
If $A$ has positive diagonal and is not diagonal, the same lemma gives
$h(A)<P(A)$, hence $Q(A)\le h(A)P(A)<P(A)^2$.
\end{proof}

\section{Scope beyond a small-correlation condition}
There are full-rank matrices in Theorem~\ref{R2E5:thm:main} with arbitrarily
large squared-correlation row sums and genuinely complex cycle phase.
Let $B_k=I_k+J_k$ for $k\ge3$. Let $H_m$ have diagonal $3$, unit entries
on the edges of an $m$-cycle, and one directed edge entry $i$ instead of
$1$, with its reverse entry $-i$, where $m\ge4$.
The matrix $B_k$ is positive definite and phase-balanced.
The matrix $H_m$ is positive definite because its off-diagonal Hermitian
part has operator norm at most $2$. Its directed cycle product is
nonreal, a property invariant under diagonal unitary congruence.

Join a distinguished vertex of $B_k$ to a vertex of $H_m$ using the
one-point sum. The resulting $A_{k,m}$ is positive definite of order
$k+m-1$, with shared diagonal $5$, other clique diagonals $2$, and other
cycle diagonals $3$. Its correlation normalization $C_{k,m}$ has, at an
ordinary clique vertex,
\[
 \sum_{j\ne i}|(C_{k,m})_{ij}|^2
 =\frac{k-2}{4}+\frac1{10}.
\]
This tends to infinity. In particular, for $k\ge6$ neither the uniform
radius nor the row-sum hypothesis of the Round 1 determinant--pairing
certificate applies. The present theorem still gives the stronger
\eqref{R2E5:eq:strong}. This is a family distinction, not a claim that every
individual example escaped all earlier literature.

The exact verifier checks the instance $k=6,m=4$ independently by subset
dynamic programming and Ryser's formula, obtaining
\[
 P(A_{6,4})=103375,\quad Q(A_{6,4})=45444243,\quad h(A_{6,4})=4320,
\]
and hence $h(A_{6,4})P(A_{6,4})-Q(A_{6,4})=401135757$.
It also verifies the gluing
identities and all principal-submatrix inequalities for sampled smaller
examples. No numerical eigenvalue test is used to establish their PSD
property; their Gram or diagonal-dominance construction supplies it.

\section{Limits, unsuccessful routes, and remaining directions}
\paragraph{What is and is not new here.}
The proposed increment over Round 1 is a different all-order structural
theorem and the stronger right-hand side $h(A)P(A)$, not a new decimal
radius. Relative to Pant--Singh v2, the proposed increment is arbitrary
complex cycle weights and phases, together with the support-level PSD
splitting argument. Their Laplacian cactus theorem, cycle expansion,
one-point-sum identities, and strong gluing mechanism are not new claims.
The triangle inequality, Gram facts, and phase-balanced observation are
classical. Priority for the resulting extension remains unverified.

\paragraph{Why the hypotheses matter.}
For a four-cycle with diagonal $1$, three edge entries $1$, and the last
edge entry $-1$, the permanent is $5$ while the permanent of its entrywise
squared modulus is $9$. This Hermitian matrix is indefinite.
Thus the local PSD restriction in \eqref{R2E5:eq:local} cannot simply be
discarded even when all entry moduli are at most one.
Also, using the two unmodified principal blocks at a cut vertex would
double-count its diagonal. For example the PSD matrix
\[
 \begin{pmatrix}1&3/5&0\\3/5&1&4/5\\0&4/5&1\end{pmatrix}
\]
needs the shared diagonal split as $9/25+16/25$, not $1+1$.
Lemma~\ref{R2E5:lem:split} addresses this real obstruction.

\paragraph{Routes that did not close the global problem.}
A direct absolute-value cycle expansion still has the wrong direction
for lower-bounding a general complex permanent. Improving only the
Round 1 matching constant would not remove its determinant singularity.
The present argument instead controls a single frustrated cycle before
gluing. If two cycles lie in the same nonbalanced two-connected block,
their overlapping cycle covers need not reduce to one scalar phase;
no bound here handles that interaction. The theorem therefore does not
cover arbitrary complex bipartite supports, arbitrary diagonally dominant
Hermitian matrices, or all graph Laplacians.

\paragraph{Reproducibility and assessment.}
\texttt{verify.py} uses exact Gaussian-integer and rational arithmetic.
It checks the cycle identities and inequalities, the triangle bound,
one-point-sum formulas, sampled hereditary inequalities, and the displayed
obstructions and non-small-correlation family. Finite checks are
diagnostic; the proof is analytic and contains no computational premise.
The deterministic run checks $1954$ strong inequalities, including
$1664$ principal submatrices, and records $474$ agreements between the
two permanent algorithms.
This is a modest structural extension of known mechanisms, not an
Annals-level claim. Next directions are nonbalanced blocks with two
independent cycles, two-vertex separators, and a wider literature review
of strong self-Bapat--Sunder inequalities. No later round is claimed.



\chapter[A positive permanent-defect identity for two-hub bipartite matrices]{A positive permanent-defect identity for two-hub bipartite matrices}
\label{source:R3E3}
\noindent\textit{Source development: Round 3, entry 03. Original mathematical arguments retained; notation is local to this chapter.}\par\medskip

\begin{quote}\small\noindent\textbf{Chapter abstract.} 
We prove a strong self-Chollet inequality for every Hermitian positive semidefinite matrix whose off-diagonal support is contained in $K_{2,m}$, for arbitrary $m$, weights and complex phases. More generally, the normalized permanent defect has an explicit nonnegative decomposition under local column-norm constraints alone, without positive semidefiniteness. This gives an equality classification and permits two-hub bipartite blocks to be joined at cut vertices. A full-rank order-ten example lies outside the previous rounds' cactus and phase-balanced block classes and fails even the optimal matching version of their determinant certificate. The new ingredient is the local defect identity, not the inherited one-point-sum method. The general Chollet conjecture is not resolved. No unrestricted two-factor Bapat--Sunder bound is asserted; that stronger statement is false in general.
\end{quote}

\section{Target and precise advance}
For a Hermitian matrix $A$, put
\[
 P(A)=\srcAOper(A),\qquad \srcAOQ(A)=\srcAOper(A\circ\overline A),\qquad
 \srcAOdd(A)=\prod_i a_{ii}.
\]
The permanent of the empty matrix is $1$. Chollet's conjecture \cite{R1E5:chollet} asks whether $P(A\circ B)\leq P(A)P(B)$ for all same-order Hermitian PSD matrices. The latest checked Pant--Singh version, dated 18 September 2026, records the general problem as open \cite{R2E3:PS}. We establish the stronger restricted self-inequality
\begin{equation}\label{R3E3:strong}
 \srcAOQ(A)\leq\srcAOdd(A)P(A).
\end{equation}
The off-diagonal support of a Hermitian matrix has an edge $ij$ when $a_{ij}\ne0$. ``Supported on $K_{2,m}$'' permits deleted edges: there are two hubs with no edge between them and $m$ leaves with no edges between leaves.

\begin{srcAOtheorem}\label{R3E3:main}
Every Hermitian PSD matrix supported on $K_{2,m}$ satisfies \eqref{R3E3:strong}, as does every principal submatrix. If all diagonal entries are positive, equality holds exactly when the support is a matching and every nonzero edge satisfies $|a_{ij}|^2=a_{ii}a_{jj}$. Equivalently, after simultaneous permutation the matrix is a direct sum of positive scalar blocks and rank-one PSD blocks of order two.
\end{srcAOtheorem}

Round~2 entries~03 and~05 independently proved the same central complex PSD cactus theorem; these are one inherited result, not two advances \cite{R3E3:R203,R3E3:R205}. Entry~05 also included phase-balanced blocks, meaning blocks diagonally unitarily similar to entrywise nonnegative matrices. Here $K_{2,m}$ with $m\geq3$ has overlapping cycles in a single block. We allow arbitrary cycle phases, so its treatment is not furnished by those results. Section~\ref{R3E3:example} gives a fully explicit separation example.

Pant--Singh's bipartite-support theorem has a sign hypothesis (real symmetric $Z$-matrices); its nonnegative permanent-monomial argument does not cover the phase-unrestricted class here. Their one-point closure calculation is reused with attribution. A recent preprint of Li--Zhao--Luo \cite{R1E5:li} claims ordinary Chollet for arbitrary complex PSD matrices through order six. We do not audit that proof here; even if accepted, its stated result does not imply the stronger unbounded-order theorem below. Low-rank results about the entire PSD matrix must not be confused with the rank-at-most-two rectangular off-diagonal block: the matrices here can have arbitrarily large full rank. These distinctions establish an advance over the specific checked results, not historical priority over all literature.

\section{The defect identity}
Consider the Hermitian matrix
\begin{equation}\label{R3E3:C}
 C=\begin{pmatrix}I_2&B\\B^*&I_m\end{pmatrix},\qquad
 B=\begin{pmatrix}u_1&\cdots&u_m\\v_1&\cdots&v_m\end{pmatrix}.
\end{equation}
Set
\begin{align*}
 x_j&=|u_j|^2,&y_j&=|v_j|^2,& r&=\sum_j u_j\overline{v_j},\\
 X&=\sum_jx_j,&Y&=\sum_jy_j,& Z&=\sum_jx_jy_j,\\
 X_2&=\sum_jx_j^2,&Y_2&=\sum_jy_j^2,& Z_2&=\sum_jx_j^2y_j^2.
\end{align*}
Define the local slack sums
\begin{equation}\label{R3E3:slacks}
 a=\sum_jx_j(1-x_j-y_j)=X-X_2-Z,\qquad
 b=\sum_jy_j(1-x_j-y_j)=Y-Y_2-Z.
\end{equation}

\begin{srcAOtheorem}[Exact nonnegative decomposition]\label{R3E3:identity}
For every matrix \eqref{R3E3:C}, without any definiteness assumption,
\begin{align}\label{R3E3:defect}
 P(C)-\srcAOQ(C)={}&a+b+Z(X_2+Y_2)+bX_2+aY_2\notag\\
 &+Z(a+b)+ab+2Z_2+|r|^2.
\end{align}
Consequently, if $x_j+y_j\leq1$ for every $j$, then
\begin{equation}\label{R3E3:quantitative}
 P(C)-\srcAOQ(C)\geq 2Z_2+|r|^2\geq0.
\end{equation}
Under these local constraints, equality holds exactly when each column of $B$ is zero or has precisely one nonzero entry of modulus one. No restriction on the number of such columns per row is needed unless PSD is also imposed.
\end{srcAOtheorem}
\begin{proof}
A nonzero permutation monomial has at most two transpositions or one four-cycle. There are no longer cycles because there are only two hubs. The single transpositions contribute $X+Y$. Pairs of disjoint transpositions contribute $XY-Z$: the subtracted terms would use the same leaf twice. The two oriented four-cycles over leaves $j<k$ contribute
\[
 2\operatorname{Re}\sum_{j<k}u_j\overline{v_j}\overline{u_k}v_k
 =|r|^2-Z.
\]
Thus
\begin{equation}\label{R3E3:Pformula}
 P(C)=1+X+Y+XY-2Z+|r|^2.
\end{equation}
Apply the same expansion to $C\circ\overline C$, whose two cross rows are $(x_j)$ and $(y_j)$. Its row inner product is $Z$, so
\begin{equation}\label{R3E3:Qformula}
 \srcAOQ(C)=1+X_2+Y_2+X_2Y_2-2Z_2+Z^2.
\end{equation}
Subtract and substitute $X=X_2+Z+a$ and $Y=Y_2+Z+b$. The linear part $X+Y-X_2-Y_2-2Z$ is $a+b$, and
\[
 XY-X_2Y_2-Z^2
 =Z(X_2+Y_2)+bX_2+aY_2+Z(a+b)+ab.
\]
This proves \eqref{R3E3:defect}. Under the stated local conditions all its summands are nonnegative, proving \eqref{R3E3:quantitative}.

If the defect vanishes, its summands $a,b,2Z_2$ vanish. From $Z_2=0$, each $x_jy_j=0$. Then $a=b=0$ and $0\leq x_j,y_j\leq1$ imply that each $x_j,y_j$ belongs to $\{0,1\}$. The two cannot both equal one in a column. Conversely these conditions give $a=b=Z=Z_2=r=0$, so the defect vanishes.
\end{proof}

\begin{proof}[Proof of Theorem~\ref{R3E3:main}]
If some $a_{ii}=0$, PSD forces row and column $i$ to vanish, and both sides of \eqref{R3E3:strong} are zero. Otherwise normalize by positive diagonal congruence to \eqref{R3E3:C}. If $A=DCD$ then $P(A)=\srcAOdd(A)P(C)$ and $\srcAOQ(A)=\srcAOdd(A)^2\srcAOQ(C)$.

The principal submatrix on the two hubs and leaf $j$ has determinant $1-x_j-y_j\geq0$. Theorem~\ref{R3E3:identity} therefore applies. Every principal submatrix has the same structure, or only one hub or no hubs; a missing hub can be adjoined as an isolated diagonal-one index without changing either permanent.

For the equality claim, PSD also gives $BB^*\preceq I_2$, by the Schur complement of $I_m$. In particular $X,Y\leq1$. The equality characterization of Theorem~\ref{R3E3:identity} consequently allows at most one nonzero entry in each row of $B$. This is exactly the stated matching condition; diagonal rescaling gives the rank-one two-by-two blocks. The converse follows either from the identity or multiplicativity on direct sums.
\end{proof}

The local assumptions are strictly weaker than PSD: take $u_j=1$, $v_j=0$ for two leaves. They hold, but $X=2$ makes \eqref{R3E3:C} indefinite. Thus the identity is a local algebraic certificate, not merely a restatement of a global eigenvalue condition.

\section{Inherited closure and mixed consequences}
Call a PSD matrix hereditarily strong if \eqref{R3E3:strong} holds for every principal submatrix. The following proof records the inherited mechanism so that the new block class can be used without an implicit decomposition assumption.

\begin{srcAOlemma}[Strong closure and support splitting]\label{R3E3:closure}
Hereditary strongness is preserved by direct sums and one-point sums. Moreover, a PSD matrix with a cut vertex can be decomposed as a one-point sum of smaller PSD matrices, changing only the two copies of the cut-vertex diagonal.
\end{srcAOlemma}
\begin{proof}
For $A_i=\left(\begin{smallmatrix}B_i&w_i\\w_i^*&d_i\end{smallmatrix}\right)\succeq0$, let
\[
 A_1\star A_2=\begin{pmatrix}B_1&w_1&0\\w_1^*&d_1+d_2&w_2^*\\0&w_2&B_2\end{pmatrix}.
\]
Put $P_i=P(A_i)$, $R_i=P(B_i)$, $h_i=\srcAOdd(B_i)$. Permutation expansion at the common index gives
\begin{align*}
 P(A_1\star A_2)&=P_1R_2+R_1P_2,\\
 \srcAOQ(A_1\star A_2)&=\srcAOQ(A_1)\srcAOQ(B_2)+\srcAOQ(B_1)\srcAOQ(A_2)
 +2d_1d_2\srcAOQ(B_1)\srcAOQ(B_2).
\end{align*}
Strongness of the four matrices bounds the second line by
$h_1h_2(d_1P_1R_2+d_2R_1P_2+2d_1d_2R_1R_2)$.
Its difference from $\srcAOdd(A_1\star A_2)P(A_1\star A_2)$ is
\[
 h_1h_2\{d_2R_2(P_1-d_1R_1)+d_1R_1(P_2-d_2R_2)\}\geq0
\]
by Lieb's inequality $P_i\geq d_iR_i$ \cite{R1E5:lieb}. All factors are nonnegative. This is the complex form of the strong closure algebra of \cite[Theorem 4.3]{R2E3:PS}, also used in both Round~2 entries.

For a principal submatrix the same argument applies to the corresponding restricted summands. If it meets only one side and includes the common index, the other summand is a scalar block. Such blocks are hereditarily strong. Direct sums use multiplicativity.

For the final assertion, write the index set as $U\cup\{v\}\cup W$, with $A[U,W]=0$, and represent $A$ as a Gram matrix of vectors $g_i$. The spans of the $U$ and $W$ vectors are orthogonal. Decompose $g_v=u+w+s$ along those two spans and their orthogonal complement. Use $u$ as the cut-vertex vector on the $U$ side and $w+s$ on the $W$ side. This preserves all required off-diagonal entries and splits $a_{vv}$ into $\|u\|^2+\|w+s\|^2$. Both summands are PSD. This Gram splitting is the support-level argument already used in Round~2.
\end{proof}

\begin{srcAOcorollary}\label{R3E3:blocks}
If every nontrivial block of the support of a PSD matrix is a subgraph of some $K_{2,m}$, the matrix is hereditarily strong. More generally, arbitrary one-point sums of the new two-hub matrices and previously established hereditarily strong matrices are hereditarily strong.
\end{srcAOcorollary}
\begin{proof}
Induct on the order, splitting at a cut vertex by Lemma~\ref{R3E3:closure}. At a block with no cut vertex apply Theorem~\ref{R3E3:main}; isolated vertices and edges are included. The same reasoning applies to principal submatrices, whose blocks are subgraphs of original blocks. The second assertion is the first part of the lemma.
\end{proof}

In particular the previously established complex cactus and phase-balanced block classes \cite{R3E3:R203,R3E3:R205} can be combined with the new blocks. This application does not reprove or reclaim those prior block estimates. For any two same-order hereditarily strong PSD matrices, even with different supports, Cauchy--Schwarz on the permutation monomials gives
\[
 0\leq P(A\circ B)\leq\sqrt{\srcAOQ(A)\srcAOQ(B)}
 \leq\sqrt{\srcAOdd(A)\srcAOdd(B)P(A)P(B)}\leq P(A)P(B).
\]
The first inequality uses the Schur product theorem; the last uses Lieb's $P(A)\geq\srcAOdd(A)$. Thus the new class supplies genuine restricted cases of the original two-factor conjecture, not just its self-form.

\section{An exact separation example}\label{R3E3:example}
In \eqref{R3E3:C} take $m=8$ and
\begin{align*}
 u&=(2/3,2/3,1/30,1/30,1/30,1/30,1/30,1/30),\\
 v&=(2/3,-2/3,i/30,-i/30,1/30,-1/30,-1/30,1/30).
\end{align*}
The rows are orthogonal and each has squared norm $403/450<1$. Hence $BB^*=(403/450)I_2$, and $C$ is positive definite of full rank ten. Its support is the single block $K_{2,8}$, not a cactus. The oriented four-cycle through the first and third leaves has nonreal product. Cycle products are unchanged by diagonal unitary conjugation, so this block is not phase-balanced.

For comparison with Round~1 entry~05 \cite{R3E3:R105}, let $W$ have off-diagonal entries $|c_{ij}|^2$ and zero diagonal. That entry's determinant/matching sufficient certificate is
\begin{equation}\label{R3E3:oldcertificate}
 \det(I-W)\prod_{e\in F}(1+w_e)^2\geq1
\end{equation}
for some matching $F$, when the squared-correlation row sum is less than one. Here the two cross rows of $W$ both equal
$(4/9,4/9,1/900,1/900,1/900,1/900,1/900,1/900)$.
A matching uses at most two edges, and its maximal product is $(13/9)^2$, attained on the two heavy leaves. Direct calculation gives
\[
 \det(I-W)=\frac{42497}{202500},\qquad
 \max_F\det(I-W)\prod_{e\in F}(1+w_e)^2
 =\frac{1213756817}{1328602500}<1.
\]
Thus even its optimal matching certificate fails. The maximum row sum is $403/450$, also outside that entry's uniform sufficient threshold. The new theorem nevertheless applies and yields, exactly,
\[
 P(C)=\frac{94601}{33750},\qquad
 \srcAOQ(C)=\frac{63845870003}{32805000000}<P(C).
\]
This is a separation of sufficient methods, not a counterexample to any permanent inequality. Full rank also excludes an explanation solely by a small-rank theorem for the entire matrix.

\section{Verification, limitations, and next obstruction}
The accompanying exact verifier checks the polynomial identity symbolically; checks the two permanent expansions against an independent subset-DP permanent algorithm on Gaussian-rational matrices; checks the order-ten example, its determinant and optimal matching certificate; and checks equality and singular boundary examples. These are reproducibility diagnostics. All-order validity follows from the displayed proof, not finite testing.

The main structural limitation is an edge between the two hubs. Such an edge introduces triangular permutation cycles and removes the local constraint $x_j+y_j\leq1$. For a normalized matrix with hub entry $c$, the same enumeration instead gives
\[
 P=1+X+Y+XY-2Z+|c+r|^2,\qquad
 \srcAOQ=1+X_2+Y_2+X_2Y_2-2Z_2+(|c|^2+Z)^2.
\]
PSD constrains $|c-r|^2\leq(1-X)(1-Y)$ with $X,Y\leq1$, but the present nonnegative decomposition does not settle this case. Floating exploration is retained separately and is not a theorem. With three hubs, six-cycles and disjoint configurations produce further terms. Nor does the present proof show the Bapat--Sunder bound with one structured factor and an arbitrary second PSD factor.

Primary-source checking covered the latest Pant--Singh version and the cited classical sources. Searches did not locate this precise arbitrary-phase two-hub defect identity, but that absence does not establish priority. In particular, the complete full text of Sa-Nguansin--Rodtes's 2025 correlation-matrix paper was not accessible in this check \cite{R3E3:SR}; potential overlap with specialized correlation classes merits further review. The independently verifiable contribution is the identity, its all-$m$ consequences, and its strict separation from the specified prior-round certificates.



\chapter[One structured factor is enough: Bapat--Sunder bounds for long-cycle blocks]{One structured factor is enough: Bapat--Sunder bounds for long-cycle blocks}
\label{source:R3E4}
\noindent\textit{Source development: Round 3, entry 04. Original mathematical arguments retained; notation is local to this chapter.}\par\medskip


\begin{quote}\small\noindent\textbf{Chapter abstract.} 
For every Hermitian positive semidefinite matrix $A$ whose support is a cactus without cycles of lengths four or five, and every positive semidefinite matrix $B$ of the same order, we prove
\[
\srcAPper(A\circ B)\leq \srcAPper(A)\prod_i b_{ii}.
\]
The second factor is completely unrestricted. The new local result covers arbitrary complex phases and nonuniform weights on every cycle of length at least six. Its proof combines a weighted Gram-chain estimate with sharp path-product bounds; a separate rational certificate treats the six-cycle. We give quantitative local deficits, a computable criterion for shorter cycles, the full constant-modulus five-cycle case, and explicit open neighborhoods containing dense matrices. Cut-vertex splitting and the classical order-three permanent-on-top theorem supply the global result. This is a restricted positive case of a false universal inequality, not a solution of the general Chollet conjecture.
\end{quote}

\section{Statement, provenance, and the change from Round 2}
Write $P(A)=\srcAPper(A)$, $\srcAPh(A)=\prod_i a_{ii}$, and $A\circ B=(a_{ij}b_{ij})$. A \emph{correlation matrix} is Hermitian PSD with diagonal one. Say that a PSD matrix $A$ has property $\srcAPBS$ if
\begin{equation}\label{R3E4:eq:BS}
 P(A\circ B)\leq P(A)\srcAPh(B)\qquad\text{for every same-order PSD }B.
\end{equation}
The universal assertion \eqref{R3E4:eq:BS} is the disproved Bapat--Sunder proposal; see \cite{R2E3:Zhang,R3E4:Pioge}. A positive case of it also gives Chollet's inequality $P(A\circ B)\leq P(A)P(B)$, by the classical permanent Hadamard inequality $P(B)\geq\srcAPh(B)$ \cite{R2E3:Lieb}.

The support graph has an edge $ij$ precisely when $i\ne j$ and $a_{ij}\ne0$. A cactus has only edges and cycles as its nontrivial blocks. A Hermitian matrix is \emph{phase-balanced} if a diagonal unitary conjugation makes all its entries nonnegative real.

\begin{srcAPtheorem}\label{R3E4:thm:main}
Every Hermitian PSD matrix whose support is a cactus with no cycle of length four or five has property $\srcAPBS$. More generally, the conclusion holds if every block of its support is either a triangle, a cycle of length at least six, or a phase-balanced block. The corresponding property holds for every principal submatrix.
\end{srcAPtheorem}
The phase-balanced condition concerns the entries within that block, not a prescribed sign pattern. Edges and isolated vertices are included automatically. In particular, cycles of lengths four and five are permitted if phase-balanced.

\paragraph{What is new here.}
Round~2 entries~03 and~05 independently proved the same central cactus inequality
\begin{equation}\label{R3E4:eq:old}
 P(A\circ\overline A)\leq\srcAPh(A)P(A).
\end{equation}
Applying it to \emph{both} factors and using Cauchy--Schwarz gives a mixed inequality when both support graphs are cacti. It does not give \eqref{R3E4:eq:BS} for unrestricted $B$. The present result removes every restriction on the second factor, while retaining two unresolved short-cycle block lengths in the first factor. Thus the old and new domains are not identical, and the old theorem is not relabeled as new.

The Gram splitting argument and the block perspective are reused explicitly from Round~2/03, Lemma~4.1 and its proof of Theorem~1.1, and Round~2/05, Section~4 \cite{R3E4:R203,R3E4:R205}; Pant--Singh's public one-point-sum formulas precede both \cite{R2E3:PS}. The new ingredients are the weighted holonomy estimate, its coupling to matching losses, and the uniform long-cycle and six-cycle certificates. The triangle input is the classical theorem of Bapat--Sunder \cite{R2E3:BS}, not a new contribution.

\paragraph{Current source status.}
We checked Pant--Singh's latest arXiv version, v2 of 18 September 2026, and Pioge et al.'s v1/journal version \cite{R2E3:PS,R3E4:Pioge}. Pant--Singh already prove one-factor bounds for real $Z$-matrices with bipartite support; those have nonnegative permutation monomials. Our frustrated complex cycles need not have that property. A public preprint by Li--Zhao--Luo, posted 31 August 2026, claims the \emph{weaker} Chollet inequality through order six \cite{R1E5:li}. Its proof is not used or independently certified here. Neither that claim nor the frozen Round~2 results imply the arbitrary-order one-factor theorem above. The searches do not establish historical priority. The earlier fixed-parameter $a=6$ gamma result of this entrant has a public antecedent; it is unrelated to, and not credited as novelty in, the present paper.

\section{Gram-chain estimates}
All inner products in the proofs below are conjugate-linear in the first variable. The two estimates in the next lemma apply to arbitrary Gram vectors, with no support restriction.

\begin{srcAPlemma}[A chain controls its closing phase]\label{R3E4:lem:chain}
Let $u_1,\ldots,u_m$ be unit vectors. Put
\[
 w=\langle u_1,u_2\rangle\cdots\langle u_m,u_1\rangle,
 \qquad d_i=1-|\langle u_i,u_{i+1}\rangle|^2
\]
with cyclic indices. For every omitted edge $k$,
\begin{equation}\label{R3E4:eq:chain}
 |1-w|\leq\frac12\left[\left(\sum_{i\ne k}\sqrt{d_i}\right)^2+
                          \sum_{i\ne k}d_i\right].
\end{equation}
Consequently, for positive numbers $c_i$,
\begin{equation}\label{R3E4:eq:weightedchain}
 |1-w|\leq \frac12\left(\sum_{i\ne k}\frac1{c_i}+
                 \max_{i\ne k}\frac1{c_i}\right)\sum_i c_i d_i.
\end{equation}
In a gauge making the consecutive inner products of a path nonnegative, the same right-hand side of \eqref{R3E4:eq:chain}, using just that path, bounds $|1-\langle u_p,u_q\rangle|$ between its endpoints.
\end{srcAPlemma}
\begin{proof}
Relabel so that the omitted edge is $m1$, and change individual vector phases so that $r_i=\langle u_i,u_{i+1}\rangle\geq0$ for $i<m$. Zero overlaps cause no difficulty. Set $P_i=u_i u_i^*$ and
\[
 e_i=u_{i+1}-r_i u_i,\qquad e_i\perp u_i,\qquad\|e_i\|=\sqrt{d_i}.
\]
For unit vectors, $\|P_i-P_j\|=\sqrt{1-|\langle u_i,u_j\rangle|^2}$; this follows by restricting to their two-dimensional span. Thus the triangle inequality gives
\[
 \|(I-P_i)u_1\|\leq\|P_1-P_i\|\leq\sum_{j<i}\sqrt{d_j},
 \qquad |\langle u_1,e_i\rangle|\leq\sum_{j<i}\sqrt{d_i d_j}.
\]
Writing $R=\prod_{i<m}r_i$ and iterating the vector recurrence yields
\[
 u_m=Ru_1+\sum_{i<m}\left(\prod_{j=i+1}^{m-1}r_j\right)e_i,
 \qquad |\langle u_1,u_m\rangle-R|
       \leq\sum_{j<i<m}\sqrt{d_i d_j}.
\]
Since $w=R\overline{\langle u_1,u_m\rangle}$,
\[
 |1-w|\leq 1-R^2+R|\langle u_1,u_m\rangle-R|
 \leq\sum_{i<m}d_i+\sum_{j<i<m}\sqrt{d_i d_j},
\]
which is \eqref{R3E4:eq:chain}. The endpoint assertion uses $1-R\leq\sum d_i$ in place of $1-R^2\leq\sum d_i$. Finally,
\[
 \left(\sum_{i\ne k}\sqrt{d_i}\right)^2
 \leq\left(\sum_{i\ne k}c_i^{-1}\right)\sum_{i\ne k}c_i d_i,
 \qquad \sum_{i\ne k}d_i\leq\max_{i\ne k}c_i^{-1}\sum_{i\ne k}c_i d_i.
\]
Adding the omitted nonnegative term proves \eqref{R3E4:eq:weightedchain}.
\end{proof}

\begin{srcAPlemma}[Sharp product bound for a PSD path]\label{R3E4:lem:path}
If a correlation matrix is supported on a path with $\ell\geq1$ edges of magnitudes $r_1,\ldots,r_\ell$, then
\begin{equation}\label{R3E4:eq:pathprod}
 \prod_{i=1}^{\ell}r_i\leq2^{1-\ell}.
\end{equation}
\end{srcAPlemma}
\begin{proof}
First assume positive definiteness. The tridiagonal elimination pivots satisfy
\[
 p_1=1,\qquad p_{i+1}=1-r_i^2/p_i,\qquad 0<p_i\leq1.
\]
Hence
\[
 \prod_{i=1}^{\ell}r_i^2
 =p_1(1-p_{\ell+1})\prod_{i=2}^{\ell}p_i(1-p_i)
 \leq4^{1-\ell}.
\]
For a singular matrix apply this to $(A+\epsilon I)/(1+\epsilon)$ and pass to the limit. The bound is sharp: for $\ell\geq2$, take endpoint edge magnitudes $1/\sqrt2$, internal edge magnitudes $1/2$, and the terminal pivot zero; for $\ell=1$, take $r_1=1$.
\end{proof}

\section{Cycle certificates and the six-cycle refinement}
Let $A$ be a correlation matrix supported on a cycle $C_m$, $m\geq4$. Write
\[
 x_i=|a_{i,i+1}|^2,\qquad z=\prod_i a_{i,i+1},\qquad q=|z|=\sqrt{\prod_i x_i}.
\]
For a matching $S$ of cycle edges put $x_S=\prod_{i\in S}x_i$ and $M(x)=\sum_S x_S$. A permutation supported on the cycle consists of fixed points and a matching of transpositions, unless it is one of the two directed full cycles. Thus for any correlation matrix $B$, with $y_i=|b_{i,i+1}|^2$ and $w=\prod_i b_{i,i+1}$,
\begin{equation}\label{R3E4:eq:cycleexp}
 P(A)=M(x)+2\srcAPRea z,\quad P(A\circ B)=M(xy)+2\srcAPRea(zw).
\end{equation}
The entries of $B$ off the displayed cycle are arbitrary and disappear only because $A$ is zero there.

Call positive weights $c_i$ \emph{admissible} if
\begin{equation}\label{R3E4:eq:loss}
 M(x)-M(xy)\geq\sum_i c_i(1-y_i)\qquad(0\leq y_i\leq1).
\end{equation}
Three explicit choices, in increasing strength, are
\begin{align}
 c_i&=x_i,\label{R3E4:eq:singleton}\\
 c_i&=x_i\left(1+\frac12\sum_{j:\,j\notin\{i-1,i,i+1\}}x_j\right),\label{R3E4:eq:pairs}\\
 c_i&=\sum_{S\ni i}\frac{x_S}{|S|}.\label{R3E4:eq:fullweights}
\end{align}
Indeed, $1-\prod_{i\in S}y_i\geq |S|^{-1}\sum_{i\in S}(1-y_i)$, since the left side is at least each summand. Taking only singleton matchings, singleton and two-edge matchings, or all matchings proves the three assertions.

\begin{srcAPproposition}[A finite, second-factor-free certificate]\label{R3E4:prop:criterion}
Suppose all $x_i>0$ and $c_i$ are admissible. Define
\[
 \eta_k=q\left(\sum_{i\ne k}c_i^{-1}+\max_{i\ne k}c_i^{-1}\right),
 \qquad \eta=\min_k\eta_k.
\]
Then for every correlation matrix $B$,
\begin{equation}\label{R3E4:eq:gapcriterion}
 P(A)-P(A\circ B)\geq(1-\eta)\sum_i c_i(1-|b_{i,i+1}|^2).
\end{equation}
In particular $\eta\leq1$ certifies property $\srcAPBS$. If an edge of $A$ vanishes, the support is a forest and $\srcAPBS$ holds without the certificate.
\end{srcAPproposition}
\begin{proof}
Subtract \eqref{R3E4:eq:cycleexp}; the result is at least
$\sum_i c_i(1-y_i)-2q|1-w|$. Apply Lemma~\ref{R3E4:lem:chain} and minimize in $k$. For a forest every nonzero permanent monomial of $A$ is nonnegative, and every corresponding monomial of a correlation matrix has modulus at most one. Diagonal scaling gives \eqref{R3E4:eq:BS} for general $B$.
\end{proof}

\begin{srcAPtheorem}[All weighted cycles of length at least seven]\label{R3E4:thm:seven}
For $m\geq7$, every PSD matrix supported on $C_m$ has property $\srcAPBS$. In correlation normalization, with all cycle edges nonzero,
\begin{equation}\label{R3E4:eq:longgap}
 P(A)-P(A\circ B)\geq\delta_m\sum_i x_i(1-|b_{i,i+1}|^2),
 \qquad \delta_m=1-(m-1)2^{4-m}>0.
\end{equation}
\end{srcAPtheorem}
\begin{proof}
Choose $k$ with $r_k=\sqrt{x_k}$ minimal, and use $c_i=x_i$. For $i\ne k$,
\[
 \frac q{x_i}=\frac{r_k}{r_i}\prod_{j\ne k,i}r_j
 \leq\prod_{j\ne k,i}r_j.
\]
Deleting edges $k,i$ divides the vertex set into two arcs. Each nontrivial arc is an \emph{induced principal path} of $A$, so Lemma~\ref{R3E4:lem:path} applies separately; no arbitrary edge deletion is asserted to preserve PSD. If $i$ neighbors $k$, one arc is an isolated vertex and the other has $m-2$ edges, giving $q/x_i\leq2^{3-m}$. For the other $m-3$ choices, both arcs have edges and their lengths sum to $m-2$, giving $q/x_i\leq2^{4-m}$. Consequently
\[
 \eta_k\leq2\cdot2^{3-m}+(m-3)2^{4-m}+2^{4-m}
          =(m-1)2^{4-m}.
\]
This is $3/4$ at $m=7$ and strictly decreases thereafter, since the ratio of successive bounds is $m/[2(m-1)]<1$. Proposition~\ref{R3E4:prop:criterion} finishes the proof. Zero diagonals give zero rows; zero cycle edges were already handled.
\end{proof}

The singleton estimate above is insufficient at $m=6$. The next two elementary inequalities recover the missing case.

\begin{srcAPlemma}[Two rational path certificates]\label{R3E4:lem:shortpaths}
If $a,b,c$ are the edge magnitudes of a three-edge PSD correlation path, then
\begin{equation}\label{R3E4:eq:path3}
 \frac{abc}{1+(b^2+c^2)/2}\leq\frac5{27}.
\end{equation}
The same holds with $a^2+b^2$ in the denominator. If $a^2+b^2\leq1$ and $c^2+d^2\leq1$, then
\begin{equation}\label{R3E4:eq:two2}
 \frac{abcd}{1+(a^2+d^2)/2}\leq\frac5{27}.
\end{equation}
\end{srcAPlemma}
\begin{proof}
For \eqref{R3E4:eq:path3}, the Schur complement of the two endpoint diagonal entries gives $b^2\leq(1-a^2)(1-c^2)$. Put $s=b^2+c^2\leq1$ and $t=\sqrt{1-s}$. If $c^2=1$, then $b=0$ and the claim is immediate. Otherwise
\[
 a^2\leq\frac{1-s}{1-c^2},\qquad
 (abc)^2\leq t^2\frac{c^2(s-c^2)}{1-c^2}\leq t^2(1-t)^2.
\]
The last inequality is exactly $(c^2-(1-t))^2\geq0$ after multiplication by $1-c^2$. Therefore the desired quotient is at most $2t(1-t)/(3-t^2)$. Its comparison with $5/27$ is certified by
\begin{equation}\label{R3E4:eq:quad1}
 49t^2-54t+15=49(t-27/49)^2+6/49>0.
\end{equation}
Reversing the path gives the alternative denominator.
For \eqref{R3E4:eq:two2}, put $u=a^2$, $v=d^2$, $s=(u+v)/2$. Then
\[
 abcd\leq\sqrt{u(1-u)v(1-v)}\leq s(1-s).
\]
The second inequality follows by applying $uv\leq s^2$ and $(1-u)(1-v)\leq(1-s)^2$. Finally
\begin{equation}\label{R3E4:eq:quad2}
 27s^2-22s+5=27(s-11/27)^2+14/27>0
\end{equation}
implies $s(1-s)/(1+s)\leq5/27$.
\end{proof}

\begin{srcAPtheorem}[The six-cycle]\label{R3E4:thm:six}
Every PSD matrix supported on $C_6$ has property $\srcAPBS$. For correlation matrices with nonzero cycle edges, let $c_i$ be \eqref{R3E4:eq:pairs}. Then
\begin{equation}\label{R3E4:eq:sixgap}
 P(A)-P(A\circ B)\geq\frac1{108}\sum_i c_i(1-|b_{i,i+1}|^2)
 \geq\frac1{108}\sum_i x_i(1-|b_{i,i+1}|^2).
\end{equation}
\end{srcAPtheorem}
\begin{proof}
Index the edges $0,1,\ldots,5$, and choose $r_0$ minimal. Throughout, the factor $r_0/r_i$ in $q/x_i$ is at most one.

For $i=1,5$, the other four edges form an induced four-edge path. Thus $q/c_i\leq q/x_i\leq1/8$ by Lemma~\ref{R3E4:lem:path}.

For $i=2$, the surviving edges are the single edge $1$ and the three-edge induced path $3,4,5$. Here
\[
 \frac q{c_2}=\frac{r_0}{r_2}
 \frac{r_1r_3r_4r_5}{1+(x_0+x_4+x_5)/2}
 \leq\frac{r_3r_4r_5}{1+(r_4^2+r_5^2)/2}\leq\frac5{27},
\]
since $r_1\leq1$. Reversal gives the same estimate for $i=4$.

For $i=3$, the surviving principal paths have edges $1,2$ and $4,5$. Positivity gives $r_1^2+r_2^2\leq1$ and $r_4^2+r_5^2\leq1$. Therefore
\[
 \frac q{c_3}=\frac{r_0}{r_3}
 \frac{r_1r_2r_4r_5}{1+(x_0+x_1+x_5)/2}
 \leq\frac{r_1r_2r_4r_5}{1+(r_1^2+r_5^2)/2}\leq\frac5{27}.
\]
All reductions used induced principal submatrices. The five reciprocal terms have bounds $1/8,5/27,5/27,5/27,1/8$ after multiplication by $q$, and their maximum is at most $5/27$. Hence
\[
 \eta_0\leq\frac14+4\cdot\frac5{27}=\frac{107}{108}<1.
\]
Apply Proposition~\ref{R3E4:prop:criterion}. The zero-entry cases follow as before.
\end{proof}

\section{Short cycles, cut vertices, and equality}
\begin{srcAPproposition}[Constant-modulus short cycles]\label{R3E4:prop:uniform}
A PSD correlation matrix supported on $C_5$ with a common off-diagonal edge magnitude $r$ has property $\srcAPBS$, for every permitted total cycle phase. A similarly supported $C_4$ matrix has property $\srcAPBS$ whenever $r^2\leq2/7$.
\end{srcAPproposition}
\begin{proof}
For $C_5$, \eqref{R3E4:eq:pairs} gives $c_i=r^2(1+r^2)$, so the criterion is
$\eta=5r^3/(1+r^2)\leq1$. Gauge the five edge phases to a common phase $\phi/5$. The eigenvalues of the resulting cyclic matrix are
$1+2r\cos((\phi+2\pi j)/5)$. Some angle is within $\pi/5$ of $\pi$, so PSD implies
$r\leq[2\cos(\pi/5)]^{-1}=(\sqrt5-1)/2$. The displayed expression for $\eta$ is increasing in $r$ and at this endpoint equals $(3\sqrt5-5)/2<1$, the latter inequality following from $45<49$.
For $C_4$, $c_i=r^2(1+r^2/2)$, giving $\eta=4r^2/(1+r^2/2)\leq1$ precisely when $r^2\leq2/7$. The case $r=0$ is trivial.
\end{proof}

The constant-modulus hypothesis in Proposition~\ref{R3E4:prop:uniform} is imposed \emph{after} diagonal normalization. It is not automatically preserved when a cut-vertex diagonal is split. Thus this proposition is not used to assert an arbitrary-weight five-cycle block theorem.

\begin{srcAPlemma}[Classical small blocks and balanced blocks]\label{R3E4:lem:base}
Every PSD matrix of order at most three has property $\srcAPBS$. Every phase-balanced PSD matrix has property $\srcAPBS$.
\end{srcAPlemma}
\begin{proof}
For the first statement we use, explicitly, the order-three permanent-on-top theorem of Bapat--Sunder \cite{R2E3:BS}. Let $S(A)$ be the Schur power matrix, indexed by permutations, with entries
$S(A)_{\sigma,\tau}=\prod_i a_{\sigma(i),\tau(i)}$. It is PSD, and that theorem says $S(A)\preceq P(A)I$ when the order is at most three. If $B$ is a correlation matrix, $S(B)$ is PSD with diagonal one. Taking its Schur product with this inequality gives
\[
 S(A\circ B)=S(A)\circ S(B)\preceq P(A)I.
\]
The row sums of $S(A\circ B)$ equal $P(A\circ B)$, proving the assertion. Scaling handles general $B$.
For a phase-balanced matrix, conjugate to nonnegative entries. Every permutation monomial is then nonnegative, whereas every monomial of a correlation matrix has modulus at most one. The triangle inequality proves \eqref{R3E4:eq:BS}.
\end{proof}

\begin{srcAPlemma}[One-point closure]\label{R3E4:lem:closure}
Let $A=A_1\star A_2$ be the one-point sum of PSD matrices, obtained by identifying one index and adding its two diagonal entries. Let $A_i^-$ delete that index. If $A_1,A_2,A_1^-,A_2^-$ have property $\srcAPBS$, so does $A$.
\end{srcAPlemma}
\begin{proof}
Normalize $B$ to a correlation matrix. Write $P_i=P(A_i)$, $R_i=P(A_i^-)$ and use tildes for their Hadamard products with the corresponding principal submatrices of $B$. Splitting permutation cycles at the common vertex gives
\[
 P(A)=P_1R_2+R_1P_2,\qquad P(A\circ B)=\widetilde P_1\widetilde R_2+
 \widetilde R_1\widetilde P_2.
\]
These are Pant--Singh's one-point formulas \cite[Proposition~1]{R2E3:PS}, applied also after masking. Since PSD permanents are nonnegative, all four inequalities $\widetilde P_i\leq P_i$, $\widetilde R_i\leq R_i$ can be multiplied. This proves the result. Entries of $B$ across the two sides need not vanish.
\end{proof}

\begin{proof}[Proof of Theorem~\ref{R3E4:thm:main}]
Zero diagonal entries in $A$ or $B$ give zero rows and make the relevant inequality trivial. Normalize all remaining diagonal entries as needed. At a cut vertex $v$, let $U,W$ be a partition of the components of its deletion. In a Gram representation of $A$, the spans $H_U,H_W$ are orthogonal. Decompose its vector at $v$ as $g_v=u+w+r$, with $u\in H_U$, $w\in H_W$ and $r\perp H_U+H_W$. Replace $g_v$ by $u$ on the first side and by $w+r$ on the second. This gives PSD matrices $A_1,A_2$ with unchanged off-diagonal entries and diagonals at $v$ summing to $a_{vv}$, hence $A=A_1\star A_2$. This is the Gram splitting argument of the frozen Round~2 papers \cite{R3E4:R203,R3E4:R205}, restated to avoid double-counting the central diagonal.

Induct on order, using multiplicativity on disconnected supports and Lemma~\ref{R3E4:lem:closure} at a cut vertex. A block with at most three vertices is handled by Lemma~\ref{R3E4:lem:base}; long cycles are handled by Theorems~\ref{R3E4:thm:seven} and~\ref{R3E4:thm:six}; balanced blocks are handled by Lemma~\ref{R3E4:lem:base}. The decomposition changes only diagonal entries, so a balanced block remains balanced. Principal deletions of a cycle are forests; principal submatrices of balanced matrices remain balanced. Thus the same induction applies to all distinguished deletions and proves the hereditary assertion. The corollary $P(A\circ B)\leq P(A)P(B)$ follows from $P(B)\geq\srcAPh(B)$.
\end{proof}

\begin{srcAPcorollary}[Rigidity on long-cycle cacti]\label{R3E4:cor:equality}
Suppose $A$ has positive diagonal and cactus support with every cycle of length at least six, and $B$ is a correlation matrix. Equality in \eqref{R3E4:eq:BS} holds if and only if the principal submatrix of $B$ on each connected component of the support of $A$ has rank one.
\end{srcAPcorollary}
\begin{proof}
For a single long cycle the strictly positive gaps \eqref{R3E4:eq:longgap} and \eqref{R3E4:eq:sixgap} force $|b_{ij}|=1$ on every support edge. For a single edge the difference is $|a_{ij}|^2(1-|b_{ij}|^2)$ after normalization. Under cut-vertex splitting all differences in Lemma~\ref{R3E4:lem:closure} are nonnegative. A strict difference on either full side remains strict after multiplication by the permanent of the opposite distinguished deletion, which is positive because all its diagonal entries are positive. Positivity here also follows from the symmetric-tensor formula for a PSD permanent: a product of nonzero linear forms is nonzero. Each split cut-vertex diagonal is positive, because each connected side retains a nonzero incident edge and its summand is PSD. Thus the induction remains within positive-diagonal matrices and forces unit modulus on every support edge. Equality in Cauchy--Schwarz for the Gram vectors of $B$, propagated along a connected graph, makes them all collinear. Conversely, a rank-one correlation matrix masks $A$ by diagonal unitary conjugation on that component, which leaves its permanent unchanged. Direct sums complete the argument.
\end{proof}

\section{Dense perturbations: the sparse theorem has an interior}
The positive local deficits also allow chords. The following elementary stability statement avoids claiming a general chorded-cycle theorem.

\begin{srcAPtheorem}[An explicit full-dimensional neighborhood]\label{R3E4:thm:interior}
Let $A_0$ be a correlation matrix supported on a cycle of length $m\geq6$, with all edges nonzero. Put $x_{\min}=\min_i |(A_0)_{i,i+1}|^2$ and
\[
 \delta=\begin{cases}1/108,&m=6,\\1-(m-1)2^{4-m},&m\geq7.\end{cases}
\]
For another correlation matrix $A$ define
\[
 L(A,A_0)=\sum_{\sigma\in S_m}\left|\prod_i a_{i,\sigma(i)}-
                                   \prod_i (A_0)_{i,\sigma(i)}\right|.
\]
If $L(A,A_0)\leq2\delta x_{\min}/m^2$, then $A$ has property $\srcAPBS$. In particular this holds whenever
\begin{equation}\label{R3E4:eq:radius}
 \max_{ij}|a_{ij}-(A_0)_{ij}|\leq
 \frac{2\delta x_{\min}}{m^3m!}.
\end{equation}
The hypothesis that $A$ is a correlation matrix is part of the assertion. If $A_0$ is positive definite, the entire diagonal-one Hermitian ball of entrywise radius strictly less than
\[
 \min\left\{\frac{2\delta x_{\min}}{m^3m!},\frac{\lambda_{\min}(A_0)}m\right\}
\]
consists of positive definite matrices with property $\srcAPBS$. It is a nonempty relatively open neighborhood in the full diagonal-one Hermitian space and contains dense, non-cactus matrices.
\end{srcAPtheorem}
\begin{proof}
For a correlation matrix $B$ let $E=\sum_i |(A_0)_{i,i+1}|^2(1-|b_{i,i+1}|^2)$. The cycle theorems give $P(A_0)-P(A_0\circ B)\geq\delta E$. Gauge the Gram vectors of $B$ along the spanning path $1,\ldots,m$ to make its consecutive inner products nonnegative. The endpoint assertion of Lemma~\ref{R3E4:lem:chain}, followed by unweighted Cauchy--Schwarz, gives for every $i,j$
\[
 |1-b_{ij}|\leq\frac m2\sum_{\ell=1}^{m-1}(1-|b_{\ell,\ell+1}|^2)
 \leq\frac{mE}{2x_{\min}}.
\]
All individual entries have modulus at most one, so telescoping a product yields
\[
 \left|1-\prod_i b_{i,\sigma(i)}\right|\leq\frac{m^2E}{2x_{\min}}
 \quad(\sigma\in S_m).
\]
Permutation monomials are unchanged by the gauge. Subtracting the two permanent deficits and bounding termwise therefore proves
\[
 P(A)-P(A\circ B)\geq
 \left(\delta-\frac{m^2L(A,A_0)}{2x_{\min}}\right)E\geq0.
\]
Finally, telescoping products of correlation entries gives
$L(A,A_0)\leq m!m\max_{ij}|a_{ij}-(A_0)_{ij}|$, proving \eqref{R3E4:eq:radius}. If the entrywise perturbation is bounded by $\epsilon$, its operator norm is at most $m\epsilon$. Thus $m\epsilon<\lambda_{\min}(A_0)$ ensures positive definiteness. Arbitrarily small generic Hermitian perturbations fill every zero off-diagonal entry.
\end{proof}

For example, take $m=7$, edge magnitude $2/5$, and total cycle phase $\pi/2$. Strict diagonal dominance gives $A_0\succ0$, while its non-real cycle product prevents phase balance. Here $\delta=1/4$, $x_{\min}=4/25$, and the sufficient radius in \eqref{R3E4:eq:radius} is exactly $1/21609000$. This small explicit radius is not asserted optimal; its point is that the result is not confined to exact sparse zero patterns.

\section{Checks, limitations, and next problems}
The attached \texttt{verify.py} checks the two rational quadratic certificates, the infinite-tail bound by a symbolic ratio, matching expansions and loss weights in exact arithmetic, and Gaussian-rational PSD examples masked by unrestricted dense PSD matrices. It also verifies the one-point identity and the stated perturbation radius. A separate seeded numerical Gram-vector test checks the phase estimate as a diagnostic, not as proof. All universal assertions follow from the arguments above.

Arbitrary weighted frustrated four- and five-cycle blocks remain outside the unconditional theorem. Proposition~\ref{R3E4:prop:criterion} covers an explicit subset, including every constant-modulus five-cycle, but no numerical search is promoted to a proof of the missing cases. There is an exact obstruction to this particular certificate. Let $C=\left(\begin{smallmatrix}3/5&4/5\\-4/5&3/5\end{smallmatrix}\right)$ and $A=\left(\begin{smallmatrix}I_2&C\\srcAPC^*&I_2\end{smallmatrix}\right)$. Since $C$ is orthogonal, $A\succeq0$ and its support is a four-cycle, with alternating squared edge magnitudes $9/25,16/25$ and negative cycle product. Even the full matching weights \eqref{R3E4:eq:fullweights} give $\eta=2(32/59+3/11)=1058/649>1$. This is a failure of the sufficient certificate, not a counterexample to \eqref{R3E4:eq:BS}. A sharper cycle-sensitive Gram estimate is the natural next step. For larger separators, permutation cycles can cross between both sides, so the one-point identity no longer closes the argument.

The constant $1/108$ and the long-cycle factors are certified, not claimed sharp. The interior result gives only a small explicit neighborhood. No full permanent-on-top statement, no universal Bapat--Sunder statement, and no resolution of the general Chollet conjecture is claimed.



\chapter[Permanent-on-top certificates for complex cycle blocks and overlapping bipartite blocks]{Permanent-on-top certificates for complex cycle blocks 
and overlapping bipartite blocks}
\label{source:R3E5}
\noindent\textit{Source development: Round 3, entry 05. Original mathematical arguments retained; notation is local to this chapter.}\par\medskip



\begin{quote}\small\noindent\textbf{Chapter abstract.} 
We prove the permanent-on-top operator inequality for complex Hermitian
positive semidefinite matrices whose support blocks are simple cycles,
phase-balanced blocks, or bipartite blocks having no edge vertex-disjoint
from a simple cycle. The last family includes arbitrary complex matrices
supported on $K_{2,m}$, with overlapping and frustrated cycles.
The proof gives explicit positive operator certificates: the classical
Schur-minimum inequality pays for the long cycles, while even matchings
produce involution Laplacians. Odd cycles require a modulus-preserving
phase reflection that preserves positive semidefiniteness.
An operator-valued one-point identity and Gram-space splitting propagate
the certificates through cut vertices. Consequently the Bapat--Sunder
bound holds when just one factor belongs to this class and the other is
an unrestricted PSD matrix. This improves the two-structured-factor
cactus conclusion of Round 2. Neither the generally false
permanent-on-top conjecture nor the open general Chollet conjecture is
claimed to be resolved; historical novelty of the structural theorem
remains unverified.
\end{quote}

\section{Statement and precise scope}
Write $A\succeq0$ for complex Hermitian positive semidefiniteness and put
\[
 P(A)=\srcAQper(A),\qquad h(A)=\prod_i a_{ii},\qquad
 \srcAQSch(A)_{\sigma,\tau}=\prod_{i=1}^n a_{\sigma(i),\tau(i)}
 \quad(\sigma,\tau\in\srcAQG_n).
\]
The empty permanent and empty diagonal product are $1$.
The matrix $\srcAQSch(A)$ is the Schur power, indexed by all permutations.
The \emph{permanent-on-top property} is
\begin{equation}\label{R3E5:eq:POT}
 0\preceq\srcAQSch(A)\preceq P(A)\srcAQId_{n!}.
\end{equation}
Its unrestricted assertion is false \cite{R3E5:Tran,R2E3:Zhang}; our result is
only a sufficient structural condition.

The support graph has an edge $ij$ when $i\ne j$ and $a_{ij}\ne0$.
A block means a maximal two-connected subgraph, a bridge, or an isolated
vertex. A block is \emph{phase-balanced} when a diagonal unitary
congruence makes every entry of its principal matrix nonnegative.
Define $\srcAQClass$ to consist of PSD matrices for which every support block
is of at least one of the following kinds:
\begin{enumerate}
\item a phase-balanced block;
\item a simple cycle of any length at least three;
\item a bipartite block in which no edge is vertex-disjoint from any
simple cycle.
\end{enumerate}
Condition 3 concerns all simple cycles, not just induced cycles.
It allows many overlapping cycles: every subgraph of $K_{2,m}$ has this
property. No phase condition is imposed in cases 2 or 3.

\begin{srcAQtheorem}\label{R3E5:thm:main}
Every matrix $A\in\srcAQClass$ has the permanent-on-top property
\eqref{R3E5:eq:POT}. Every principal submatrix belongs to $\srcAQClass$ and has the
same property. For every PSD matrix $B$ of the same order, with no support,
rank, or phase restriction,
\begin{equation}\label{R3E5:eq:BS}
 0\le P(A\circ B)\le P(A)h(B)\le P(A)P(B).
\end{equation}
\end{srcAQtheorem}

The final inequality is Chollet's inequality \cite{R1E5:chollet}. Its general
case remains open in the latest structural source checked,
Pant--Singh v2 of 18 September 2026 \cite{R2E3:PS}.
A separate, non-peer-reviewed preprint posted 31 August 2026 claims
Chollet's inequality through complex order six \cite{R3E5:LZL}; that claim is
not a dependency here and does not subsume this unbounded-order operator
theorem. The stronger unrestricted Bapat--Sunder inequality fails
\cite{R2E5:drury}. Thus \eqref{R3E5:eq:BS} is explicitly restricted in its first factor.

\paragraph{Delta from earlier rounds.}
Round 2 entries 03 and 05 independently proved the same central complex
cactus inequality
$P(A\circ\overline A)\le h(A)P(A)$ and hence Chollet for \emph{two}
cactus-supported factors \cite{R3E5:R203,R3E5:R205}. Entry 05 also allowed
phase-balanced blocks. That shared result is counted once.
Here the conclusion is the full operator bound \eqref{R3E5:eq:POT},
the second factor in \eqref{R3E5:eq:BS} is unrestricted, and the first factor
may have noncactus blocks such as $K_{2,m}$. The cycle estimates of
Round 2 are not used in the new local proof.
The standard Gram cut-vertex splitting is reused, with both entries
credited, and the scalar one-point method of Pant--Singh is upgraded to
an operator identity. Phase-balanced matrices, the Schur-minimum bound,
and the implication from permanent-on-top to \eqref{R3E5:eq:BS} are classical,
not claimed as new. A search of the primary sources listed below did not
settle historical priority for the new structural combination.

\section{Schur powers and the classical lower bound}
Our permutation convention is
\[
 d_A(\pi)=\prod_{i=1}^n a_{\pi(i),i},\qquad
 L_\pi e_\tau=e_{\pi\tau}.
\]
Thus
\begin{equation}\label{R3E5:eq:regular}
 \srcAQSch(A)=\sum_{\pi\in\srcAQG_n}d_A(\pi)L_\pi,\qquad
 P(A)=\sum_{\pi\in\srcAQG_n}d_A(\pi).
\end{equation}
Indeed, the $(\sigma,\tau)$ entry of the sum is
$d_A(\sigma\tau^{-1})$, which is the defining entry of $\srcAQSch(A)$.
In particular every row sum of $\srcAQSch(A)$ is $P(A)$.
This regular-representation formulation also appears in
\cite[Proposition 2.1]{R3E5:Tran}.
Let $D$ be diagonal with $D_{\sigma,\sigma}=\srcAQsgn(\sigma)$. Then
\begin{equation}\label{R3E5:eq:sign}
 D L_\pi D=\srcAQsgn(\pi)L_\pi,\qquad L_\pi^*=L_{\pi^{-1}}.
\end{equation}

\begin{srcAQlemma}[Classical Schur-minimum operator inequality]\label{R3E5:lem:minimum}
For every $A\succeq0$,
\[
 \srcAQSch(A)\succeq \det(A)\srcAQId.
\]
\end{srcAQlemma}
\begin{proof}
For Gram vectors $u_i$ of $A$, the vectors
$v_\sigma=u_{\sigma(1)}\otimes\cdots\otimes u_{\sigma(n)}$
have Gram matrix $\srcAQSch(A)$, so $\srcAQSch(A)\succeq0$.
This suffices when $A$ is singular.
Otherwise take $A=LL^*$ with $L$ lower triangular and positive diagonal,
and let $u_i$ be row $i$ of $L$.
Project the tensor space orthogonally onto the span of
$e_{\pi(1)}\otimes\cdots\otimes e_{\pi(n)}$, $\pi\in\srcAQG_n$.
The coefficient of this basis tensor in $v_\sigma$ is
$\prod_i L_{\sigma(i),\pi(i)}$.
It can be nonzero only if $\pi(i)\le\sigma(i)$ for every $i$.
The sums of these integers agree, so all inequalities are equalities.
The projected vectors therefore have Gram matrix
$|\prod_i L_{ii}|^2\srcAQId=\det(A)\srcAQId$.
The Gram matrix of their orthogonal residuals is PSD, as required.
\end{proof}

This is a spectral lower bound, not merely a character-trace inequality;
see also \cite{R2E3:Zhang}. We give the proof to avoid relying on an
unjustified implication from immanants to operator order.

\begin{srcAQlemma}[Elementary invariances and the balanced case]\label{R3E5:lem:balanced}
Diagonal unitary congruence leaves both $\srcAQSch(A)$ and $P(A)$ unchanged.
If $A=E C E$ for positive diagonal $E$, then
$\srcAQSch(A)=h(A)\srcAQSch(C)$ and $P(A)=h(A)P(C)$ when $C$ has diagonal $1$.
Every phase-balanced PSD matrix satisfies \eqref{R3E5:eq:POT}.
\end{srcAQlemma}
\begin{proof}
Each product defining a Schur-power entry contains every diagonal
congruence factor once and every conjugate factor once.
This proves the invariances and scaling.
After a balanced gauge, $\srcAQSch(A)$ is entrywise nonnegative and Hermitian,
with row sum $P(A)$. For any vector $x$,
\[
 x^*\srcAQSch(A)x
 \le \sum_{\sigma,\tau}\srcAQSch(A)_{\sigma,\tau}|x_\sigma||x_\tau|
 \le P(A)\sum_\sigma|x_\sigma|^2.
\]
The last step uses $2|x_\sigma||x_\tau|\le
|x_\sigma|^2+|x_\tau|^2$. Positivity follows from the Gram construction.
\end{proof}

A zero diagonal entry of a PSD matrix forces its row and column to
vanish, since $|a_{ij}|^2\le a_{ii}a_{jj}$. In that case
$P(A)=0$ and $\srcAQSch(A)=0$. Thus local proofs may normalize to correlation
matrices without losing a singular or zero-diagonal case.

\section{Explicit cycle certificates}
Let $C$ be a PSD correlation matrix supported on a cycle of length
$m\ge3$, labeled cyclically. Write $c=(1\,2\,\cdots\,m)$,
\[
 z=d_C(c),\qquad
 x_M=\prod_{ij\in M}|c_{ij}|^2,\qquad
 \pi_M=\prod_{ij\in M}(i\,j),
\]
where $M$ ranges over matchings of the cycle, including the empty one.
A supported permutation is either a matching involution or one of
the two orientations of the full cycle. Consequently
\begin{align}
 \srcAQSch(C)&=\sum_M x_M L_{\pi_M}+zL_c+\overline z L_c^*,
       \label{R3E5:eq:cycleS}\\
 P(C)&=\sum_M x_M+2\srcAQRea z,\qquad
 \det(C)=\sum_M(-1)^{|M|}x_M+(-1)^{m-1}2\srcAQRea z.
       \label{R3E5:eq:cyclePD}
\end{align}
Zero edges are allowed in these identities.

\begin{srcAQproposition}[Even-cycle certificate]\label{R3E5:prop:even}
If $m$ is even, then
\begin{equation}\label{R3E5:eq:evencert}
 P(C)\srcAQId-\srcAQSch(C)
 =D\bigl(\srcAQSch(C)-\det(C)\srcAQId\bigr)D
 +2\sum_{\substack{M\ne\varnothing\\ |M|\ {\rm even}}}
       x_M(\srcAQId-L_{\pi_M})\succeq0.
\end{equation}
\end{srcAQproposition}
\begin{proof}
Use \eqref{R3E5:eq:sign}--\eqref{R3E5:eq:cyclePD}.
The full cycle has sign $-1$, so its coefficients agree on the two sides.
An odd matching also has sign $-1$.
For a nonempty even matching the second term changes coefficient
$x_M$ to $-x_M$.
The identity coefficient agrees because
$P(C)+\det(C)=2\sum_{|M|\ {\rm even}}x_M$.
The first summand is PSD by Lemma~\ref{R3E5:lem:minimum}.
Each $\pi_M$ is an involution, so
$\srcAQId-L_{\pi_M}\succeq0$.
\end{proof}

\begin{srcAQlemma}[PSD phase reflection on an odd cycle]\label{R3E5:lem:reflection}
Suppose $m$ is odd. There is a PSD correlation matrix $C'$ with the
same edge moduli and support as $C$ and oriented cycle product
\[
 z'=|\srcAQRea z|-i\srcAQIma z.
\]
All proper principal determinants of $C'$ equal those of $C$, and
$\det(C')-\det(C)=2(|\srcAQRea z|-\srcAQRea z)\ge0$.
\end{srcAQlemma}
\begin{proof}
If $z=0$, take $C'=C$. Otherwise $|z'|=|z|$, so changing the phase
of one edge, and conjugating its reversed entry, realizes $z'$.
Every proper induced subgraph of a cycle is a forest.
Any two Hermitian matrices on a forest with the same diagonal and edge
moduli are related by diagonal unitary congruence: choose a phase at
one root of each tree and propagate it along edges.
This proves equality of proper principal determinants.
Formula \eqref{R3E5:eq:cyclePD}, with $m$ odd, proves the full determinant
comparison. All principal minors of $C'$ are therefore nonnegative.
For completeness this implies PSD: every principal determinant of
$C'+\varepsilon\srcAQId$ is the sum of its principal minors weighted by
nonnegative powers of $\varepsilon$, including a positive leading
term. Its leading principal determinants are positive, so Sylvester's
criterion gives $C'+\varepsilon\srcAQId\succ0$; let $\varepsilon\downarrow0$.
\end{proof}

\begin{srcAQproposition}[Odd-cycle certificate]\label{R3E5:prop:odd}
For $C'$ from Lemma~\ref{R3E5:lem:reflection}, put
$t=z+z'=2\max(\srcAQRea z,0)$. Then
\begin{align}
 P(C)\srcAQId-\srcAQSch(C)
 &=D\bigl(\srcAQSch(C')-\det(C')\srcAQId\bigr)D \notag\\
 &\quad+2\sum_{\substack{M\ne\varnothing\\ |M|\ {\rm even}}}
          x_M(\srcAQId-L_{\pi_M})
       +t(2\srcAQId-L_c-L_c^*)\succeq0. \label{R3E5:eq:oddcert}
\end{align}
\end{srcAQproposition}
\begin{proof}
Matching coefficients agree just as in the even case.
The odd cycle has sign $+1$; its coefficient on the right is
$z'-t=-z$, and the reversed coefficient is its conjugate.
The identity coefficient agrees by
$P(C)+\det(C')=2\sum_{|M|\ {\rm even}}x_M+2t$.
The first term is PSD by the preceding lemma and
Lemma~\ref{R3E5:lem:minimum}. The matching terms are PSD, and
$2\srcAQId-L_c-L_c^*=(\srcAQId-L_c)^*(\srcAQId-L_c)\succeq0$, with $t\ge0$.
\end{proof}

Thus every complex PSD cycle matrix has permanent-on-top, including
singular matrices and triangles. The order-three case is classical;
the proof here treats all odd lengths uniformly.

\section{Overlapping bipartite cycles}
The even-cycle identity admits a useful extension beyond cacti.

\begin{srcAQproposition}\label{R3E5:prop:bipartite}
Let $C\succeq0$ have diagonal $1$ and bipartite support.
Suppose no edge is vertex-disjoint from any simple cycle.
Then \eqref{R3E5:eq:evencert} holds, with the sum taken over all matchings
of the support graph. In particular $C$ has permanent-on-top.
\end{srcAQproposition}
\begin{proof}
A supported permutation decomposes into fixed points, transpositions,
and oriented simple cycles of length at least three.
Every such long cycle has even length by bipartiteness.
If one occurs, the support hypothesis rules out any disjoint
transposition. It also rules out a second disjoint long cycle, since
that cycle contains an edge disjoint from the first.
Thus a supported permutation is either a matching involution or
exactly one even cycle with all other vertices fixed.
Every nonmatching permutation has sign $-1$, so its coefficient cancels
in $\srcAQSch(C)+D\srcAQSch(C)D$.
The only surviving coefficients are those of even matchings, with
nonnegative weights $x_M$.
Equivalently,
\[
 \srcAQSch(C)+D\srcAQSch(C)D=2\sum_{|M|\ {\rm even}}x_ML_{\pi_M},
 \qquad P(C)+\det(C)=2\sum_{|M|\ {\rm even}}x_M .
 \]
Rearranging gives \eqref{R3E5:eq:evencert}; its summands are PSD as before.
\end{proof}

Every simple cycle of $K_{2,m}$ has length four and uses both vertices
in the size-two part. Since every edge meets that part, there is no
edge disjoint from a cycle. Thus Proposition~\ref{R3E5:prop:bipartite}
allows arbitrary phases on all $2m$ edges, including cycle products
that are not real. For $m\ge3$ these support blocks are not cacti.

\section{Operator gluing at cut vertices}
For a matrix $F$ on a subset $J\subseteq[n]$, define its ambient
Schur operator by
\[
 \widehat{\srcAQSch}_n(F)=
 \sum_{\pi\in\srcAQG_J}d_F(\pi)L_\pi
 \quad\hbox{on }\srcAQC^{\srcAQG_n},
\]
where $\pi$ fixes $[n]\setminus J$.
It is not the Schur power of the zero-padded matrix.
The left action of $\srcAQG_J$ partitions $\srcAQG_n$ into right cosets
$\srcAQG_J\tau$, and on each coset this operator is a copy of $\srcAQSch(F)$.
In particular, if $F$ has permanent-on-top, then
\begin{equation}\label{R3E5:eq:embed}
 0\preceq\widehat{\srcAQSch}_n(F)\preceq P(F)\srcAQId.
\end{equation}
Operators on disjoint index sets commute.

Suppose $A_1$ and $A_2$ meet only in a distinguished index $v$.
Their one-point sum $A=A_1\star A_2$ retains all entries on the two
sides except that the two $v$-diagonals are added; cross entries are
zero. Let $B_i=A_i(v)$ delete $v$.

\begin{srcAQlemma}[Operator one-point identity]\label{R3E5:lem:glue}
On the full permutation space of the union of indices,
\begin{equation}\label{R3E5:eq:glue}
 \srcAQSch(A)=
 \widehat{\srcAQSch}(A_1)\widehat{\srcAQSch}(B_2)
 +\widehat{\srcAQSch}(B_1)\widehat{\srcAQSch}(A_2).
\end{equation}
Consequently, if $A_1,A_2,B_1,B_2$ are PSD and each has
permanent-on-top, then so does $A$.
\end{srcAQlemma}
\begin{proof}
No supported permutation cycle can use vertices from both sides of
$v$: a simple cycle doing so would have to pass through $v$ twice.
A cycle moving $v$ therefore belongs to exactly one side, giving
one term in \eqref{R3E5:eq:glue}. When $v$ is fixed, the two terms contribute
the two separate $v$-diagonals, whose sum is the required diagonal.
All remaining cycles lie wholly on their respective sides. This proves
the identity coefficient by coefficient in the group algebra.

Each product on the right consists of operators with disjoint supports.
Its two factors are commuting PSD operators, so their product is PSD.
If $0\preceq X\preceq a\srcAQId$ and $0\preceq Y\preceq b\srcAQId$ commute,
then simultaneous diagonalization gives $XY\preceq ab\srcAQId$.
Apply \eqref{R3E5:eq:embed} to bound the right side of \eqref{R3E5:eq:glue} by
\[
 [P(A_1)P(B_2)+P(B_1)P(A_2)]\srcAQId=P(A)\srcAQId.
 \]
The last equality follows from the same permutation partition, or by
applying \eqref{R3E5:eq:glue} to the all-ones vector.
\end{proof}

The scalar identity here is the familiar one-point formula used by
Pant--Singh \cite{R2E3:PS}; the group-algebra identity proves the stronger
operator propagation without trying to commute overlapping factors.

\begin{srcAQlemma}[Gram splitting, reused from Round 2]\label{R3E5:lem:split}
If $A\succeq0$ and its support has a cut vertex $v$ separating nonempty
sets $U,W$, then $A=A_1\star A_2$ for PSD matrices on
$U\cup\{v\}$ and $W\cup\{v\}$. The off-diagonal entries are unchanged.
\end{srcAQlemma}
\begin{proof}
Realize $A$ by Gram vectors, and let $H_U,H_W$ be the spans belonging
to $U,W$. The zero cross entries give $H_U\perp H_W$.
Decompose the vector at $v$ as $w=w_U+w_W+w_0$ by orthogonal projection
onto $H_U,H_W$, and their common orthogonal complement.
For $A_1$ use the old vectors on $U$ and the vector $w_U+w_0$;
for $A_2$ use the old vectors on $W$ and $w_W$.
The Gram entries agree off the $v$-diagonal, while their squared
norms at $v$ add to $\|w\|^2$. This is the required splitting.
\end{proof}

This standard decomposition appeared explicitly in Round 2 entry 03,
Lemma 4.1, and entry 05, Lemma 9 \cite{R3E5:R203,R3E5:R205}; it is not new here.

\begin{proof}[Proof of the operator assertion in Theorem~\ref{R3E5:thm:main}]
First note heredity of the class. Each block of an induced subgraph
lies within an original block. A principal submatrix of a balanced
block is balanced. An induced subgraph of a cycle is the whole cycle
or a forest. Bipartiteness and the absence of a disjoint edge--cycle
pair are both inherited by subgraphs.
Hence every principal submatrix of a member of $\srcAQClass$ remains in
$\srcAQClass$.

Proceed by induction on the matrix order. Zero diagonal entries have
already been handled. If the support has one block, apply
Lemma~\ref{R3E5:lem:balanced}, Propositions~\ref{R3E5:prop:even}--\ref{R3E5:prop:odd},
or Proposition~\ref{R3E5:prop:bipartite}, after diagonal normalization.
If there is a cut vertex, use Lemma~\ref{R3E5:lem:split}.
Both pieces still have the same off-diagonal phases and support on
their respective sides. Their blocks belong to the permitted types;
changing the common diagonal does not affect those types.
The two pieces and their vertex-deleted submatrices have smaller
order and are PSD, so induction and Lemma~\ref{R3E5:lem:glue} apply.
Finally, disconnected components give a product of commuting embedded
Schur operators on disjoint index sets; its scalar bound is the
product of their permanents, which is the permanent of the block
diagonal matrix. Singletons and the empty matrix start the induction.
\end{proof}

\section{One unrestricted factor, and a cycle equality criterion}
The Schur product of two PSD matrices is PSD, as is seen by tensoring
their Gram vectors. Also
\[
 \srcAQSch(A)\circ\srcAQSch(B)=\srcAQSch(A\circ B),\qquad
 \srcAQId\circ\srcAQSch(B)=h(B)\srcAQId.
\]
Taking the Schur product of \eqref{R3E5:eq:POT}'s upper-bound defect with
$\srcAQSch(B)$ proves
\[
 P(A)h(B)\srcAQId-\srcAQSch(A\circ B)\succeq0.
\]
Its quadratic form on the all-ones vector gives the middle inequality
of \eqref{R3E5:eq:BS}. Positivity follows from the Gram representation.
To recall a short proof of $P(B)\ge h(B)$, give polynomials the norm
$\|f\|^2=\sum_\alpha\alpha!|f_\alpha|^2$.
Multiplication by $\ell_u=\sum u_kz_k$ has adjoint
$a_u=\sum\overline u_k\partial_k$ and obeys
\[
 \|\ell_u f\|^2=\|u\|^2\|f\|^2+\|a_u f\|^2.
\]
For Gram vectors $u_i$ of $B$, expansion gives
$P(B)=\|\prod_i\ell_{u_i}\|^2$.
Iteration proves the desired bound. This classical deletion argument
is a special case of Lieb's inequality \cite{R1E5:lieb}; the same proof
appeared in Round 2 entry 05, Lemma 4.
This completes Theorem~\ref{R3E5:thm:main}.

For a sharper scalar consequence, retain a cycle correlation matrix
$C$ and let $B$ be any correlation matrix of the same order. Define
\[
 y_{ij}=|b_{ij}|^2,\quad w=d_B(c),\quad
 \ell_M=x_M\left(1-\prod_{ij\in M}y_{ij}\right)\ge0.
\]
Here $0\le y_{ij}\le1$ by the two-by-two principal minors.
If $m$ is even put $C'=C$ and $t=0$; if $m$ is odd use the phase
reflection and $t$ from Proposition~\ref{R3E5:prop:odd}. Direct expansion of
\eqref{R3E5:eq:cyclePD}, including the entrywise products, gives
\begin{equation}\label{R3E5:eq:scalar}
 P(C)-P(C\circ B)
 =\det(C'\circ B)-\det(C')
 +2\sum_{\substack{M\ne\varnothing\\ |M|\ {\rm even}}}\ell_M
 +2t(1-\srcAQRea w).
\end{equation}
All terms on the right are nonnegative.
For the determinant term this is Oppenheim's inequality, which follows
here from Lemma~\ref{R3E5:lem:minimum}: take its Schur product with $\srcAQSch(B)$,
then the quadratic form on the sign vector
$(\srcAQsgn\sigma)_\sigma$ to obtain
$\det(C'\circ B)\ge\det(C')h(B)=\det(C')$.
Also $|w|\le1$.

\begin{srcAQcorollary}\label{R3E5:cor:equality}
If $C$ is a correlation matrix with support exactly a cycle of length
$m\ge4$, and $B$ is any correlation matrix of that order, then
\[
 P(C)-P(C\circ B)\ge
 2\sum_{\substack{M\text{ a matching}\\|M|=2}}
 x_M\left(1-\prod_{ij\in M}|b_{ij}|^2\right).
\]
Equality $P(C\circ B)=P(C)$ holds if and only if $\operatorname{rank}B=1$.
\end{srcAQcorollary}
\begin{proof}
The bound follows from \eqref{R3E5:eq:scalar}. Every cycle edge has an edge
disjoint from it when $m\ge4$, and every corresponding $x_M$ is positive.
If the permanent defect vanishes, each such matching has
$\prod_{ij\in M}y_{ij}=1$, forcing every cycle-edge $y_{ij}=1$.
Unit Gram vectors of $B$ at adjacent vertices are therefore collinear
by equality in Cauchy--Schwarz. Connectivity makes them all collinear,
so $B$ has rank one. Conversely a rank-one correlation matrix has
$b_{ij}=u_i\overline u_j$ with $|u_i|=1$; thus $C\circ B$ is a diagonal
unitary congruence of $C$ and has the same permanent.
\end{proof}

\section{An explicit noncactus instance and reproducibility}
The matrix
\[
 A=\begin{pmatrix}
 4&0&1&1&1\\
 0&4&1&i&-1\\
 1&1&3&0&0\\
 1&-i&0&3&0\\
 1&-1&0&0&3
 \end{pmatrix}
\]
is strictly diagonally dominant with positive diagonal, hence positive
definite. Its support is $K_{2,3}$ and some four-cycle products are
nonreal, so it is neither a cactus matrix nor phase-balanced.
Let $V$ have rows $(1,t_j)$ for
$(t_1,\ldots,t_5)=(0,1,i,-1,-i)$ and set $B=\srcAQId+VV^*$.
This positive definite second factor has a nonbalanced dense block;
it is not required to belong to $\srcAQClass$.
Exact arithmetic gives
\[
 P(A)=660,\quad h(B)=162,\quad P(B)=534,\quad
 P(A\circ B)=75564,
\]
so the stronger bound in \eqref{R3E5:eq:BS} has defect $31356$.
Replacing the two hub diagonals by $m+1$, keeping leaf diagonals $3$,
and giving the two hub-to-leaf rows arbitrary unit phases provides
positive definite examples of unbounded order.

The accompanying \texttt{verify.py} uses Gaussian rational arithmetic,
not floating-point eigenvalue tests. It independently checks the group
algebra identities, phase reflection and all proper principal
determinants in sampled odd cycles, the operator gluing identity,
small full Schur-power PSD defects, and scalar inequalities against
unrestricted Gram matrices. Dynamic programming and Ryser's formula
cross-check permanents. These are finite diagnostics, not a substitute
for the proofs. Exact counts and sample ranges are recorded in
\texttt{verification.txt}. An initial bounded floating-point exploration
is retained separately in \texttt{explore.py}; it is not a certificate.

\paragraph{Limitations and remaining obstruction.}
The proof does not cover a general block containing several intersecting
odd cycles, nor a bipartite block with a cycle and a disjoint edge.
In the latter situation a supported permutation can combine an even
cycle with a transposition, acquiring even sign without being a
matching involution; the positive certificate then no longer follows.
An odd-cycle phase change is safe here because every proper principal
subgraph is a forest. With several cycles, changing one phase can alter
proper principal determinants, so the reflection step is not justified.
No general extension beyond these hypotheses, no resolution of
Chollet's conjecture, and no historical priority claim is made.



\chapter[Reflected PSD companions and permanent bounds for intersecting odd cycles]{Reflected PSD companions and permanent bounds for intersecting odd cycles}
\label{source:R4E3}
\noindent\textit{Source development: Round 4, entry 03. Original mathematical arguments retained; notation is local to this chapter.}\par\medskip

\begin{quote}\small\noindent\textbf{Chapter abstract.} 
We extend the preceding round's sign-conjugation certificate to nonbipartite supports by reflecting all off-diagonal entries. If no edge is vertex-disjoint from a simple cycle and the reflected matrix is positive semidefinite, the original Hermitian matrix satisfies an upper Schur-power bound and a Bapat--Sunder upper bound against every PSD second factor. The original matrix need not be PSD; this upper bound must therefore be distinguished from full permanent-on-top. For two-hub book matrices the hypothesis is an explicit Schur-complement disk condition, treating intersecting triangles of unrestricted phases. When both the original and companion are PSD, it is a lens condition; it includes every allowed hub entry when the two leaf rows are orthogonal. A quantitative spectral gap yields full-dimensional neighborhoods with permanent-on-top, and an exact diagonal-allocation obstruction explains why a stronger support-level closure claim is unjustified. The unrestricted book case and general Chollet conjecture remain unresolved by this argument.
\end{quote}

\section{Statement and attributed increment}
For a Hermitian matrix $A$ define
\[
 P(A)=\srcAUper(A),\quad h(A)=\prod_i a_{ii},\quad
 \srcAUSch(A)_{\sigma,\tau}=\prod_i a_{\sigma(i),\tau(i)},\quad
 \srcAURefl(A)=2\operatorname{diag}(A)-A.
\]
The Schur power is indexed by $\sigma,\tau\in S_n$. Its upper bound is
\begin{equation}\label{R4E3:upper}
 \srcAUSch(A)\preceq P(A)I.
\end{equation}
Full permanent-on-top additionally asserts $\srcAUSch(A)\succeq0$, automatic when $A\succeq0$ but not assumed below. The support graph has an edge $ij$ when $i\ne j$ and $a_{ij}\ne0$.

\begin{srcAUtheorem}[Reflected companion criterion]\label{R4E3:reflection}
Let $A$ be Hermitian with nonnegative diagonal. Suppose its support has no edge vertex-disjoint from any simple cycle, and $\srcAURefl(A)\succeq0$. Then \eqref{R4E3:upper} holds, as does
\begin{equation}\label{R4E3:BS}
 P(A\circ B)\leq P(A)h(B)\leq P(A)P(B)
 \qquad\text{for every }B\succeq0.
\end{equation}
Moreover $P(A)\geq h(A)\geq0$. Every principal submatrix satisfies the same conclusions. If $A\succeq0$ as well, then it has full permanent-on-top, and $P(A\circ B)\geq0$.
\end{srcAUtheorem}

The inequality with the final right side is Chollet's proposed bound \cite{R1E5:chollet}. Its general PSD case remains open in the latest checked Pant--Singh version, v2 of 18 September 2026 \cite{R2E3:PS}. Unrestricted permanent-on-top and unrestricted Bapat--Sunder are false \cite{R4E3:Tran,R2E5:drury}; the present hypotheses are essential to the assertion being made.

\paragraph{What is inherited and what changes.}
Round~3 entry~05 proved permanent-on-top for cycles, balanced blocks and bipartite blocks with the same edge--cycle intersection condition \cite[Proposition 4.1]{R4E3:R305}. Its sign-conjugation/Schur-minimum mechanism is the starting point here and is explicitly reused. The new observation is that a reflected PSD companion cancels odd cycles too. The resulting domain contains nonbalanced book blocks with arbitrarily many intersecting triangles, which were not in that theorem. Our own Round~3 entry~03 gave a scalar local identity for $K_{2,m}$, including indefinite matrices under different local constraints \cite{R4E3:R303}. Its shared PSD consequence is subsumed by entry~05 and is not counted again here. The present indefinite class uses a different companion condition and admits a full upper operator bound.

Entry~04's Round~3 dense-neighborhood theorem established scalar Bapat--Sunder stability \cite[Theorem 5.1]{R4E3:R304}. We reuse the perturbation perspective but obtain an operator spectral gap and neighborhoods with full permanent-on-top around the new book examples. These are incremental structural results, not a general-conjecture solution or a certified priority claim. All six frozen Round~3 papers and their assessment were read.

\section{A reflected operator certificate}
Use the convention
\[
 d_A(\pi)=\prod_i a_{\pi(i),i},\qquad L_\pi e_\tau=e_{\pi\tau},
 \qquad D_{\sigma,\sigma}=\srcAUsgn\sigma.
\]
Then
\begin{equation}\label{R4E3:regular}
 \srcAUSch(A)=\sum_\pi d_A(\pi)L_\pi,
 \qquad DL_\pi D=\srcAUsgn(\pi)L_\pi.
\end{equation}
For a matching $M$ put $x_M=\prod_{ij\in M}|a_{ij}|^2$ when the diagonal is one, and $\pi_M=\prod_{ij\in M}(ij)$.

\begin{srcAUlemma}[Classical Schur minimum]\label{R4E3:minimum}
If $H\succeq0$, then $\srcAUSch(H)\succeq\det(H)I$.
\end{srcAUlemma}
\begin{proof}
Schur powers are Gram matrices of permuted tensor products of Gram vectors. This proves positivity and suffices for singular $H$. Otherwise factor $H=LL^*$ with $L$ lower triangular. Project the tensor products of its rows onto the span of basis tensors whose indices are all distinct. The coefficient indexed by a permutation $\pi$ in the tensor indexed by $\sigma$ is $\prod_iL_{\sigma(i),\pi(i)}$. It vanishes unless $\pi(i)\leq\sigma(i)$ for all $i$, forcing $\pi=\sigma$. The projected Gram matrix is therefore $|\prod_iL_{ii}|^2I=\det(H)I$. The residual Gram matrix is PSD. This classical proof was also recorded in Round~3 entry~05, Lemma~2.1.
\end{proof}

\begin{srcAUproposition}[Exact identity]\label{R4E3:identity}
Let $C$ be Hermitian with diagonal one and with the support condition of Theorem~\ref{R4E3:reflection}. Without assuming either matrix PSD,
\begin{equation}\label{R4E3:certificate}
 P(C)I-\srcAUSch(C)
 =D\bigl(\srcAUSch(\srcAURefl(C))-\det(\srcAURefl(C))I\bigr)D
 +2\sum_{\substack{M\ne\varnothing\\ |M|\text{ even}}}
 x_M(I-L_{\pi_M}).
\end{equation}
\end{srcAUproposition}
\begin{proof}
A supported permutation either consists of a matching of transpositions and fixed points, or has exactly one cycle of length at least three with all other vertices fixed. Indeed any other nontrivial cycle would contain an edge disjoint from the first. For a matching, reflection preserves its monomial, and its sign is $(-1)^{|M|}$. For a long cycle of length $l$, reflection multiplies the monomial by $(-1)^l$, while its permutation sign is $(-1)^{l-1}$. Thus every long-cycle term cancels in
\[
 \srcAUSch(C)+D\srcAUSch(\srcAURefl(C))D
 =2\sum_{|M|\text{ even}}x_M L_{\pi_M}.
\]
The corresponding scalar identity is
$P(C)+\det(\srcAURefl(C))=2\sum_{|M|\text{ even}}x_M$.
Subtracting the operator identity from this scalar times $I$ gives \eqref{R4E3:certificate}. Every $\pi_M$ is an involution, so $I-L_{\pi_M}\succeq0$.
\end{proof}

\begin{proof}[Proof of Theorem~\ref{R4E3:reflection}]
If $a_{ii}=0$, the PSD companion has zero row $i$ and hence so does $A$; its permanent and Schur power are zero. Otherwise positive diagonal congruence reduces to diagonal one. Indeed $A=ECE$ gives $\srcAUSch(A)=h(A)\srcAUSch(C)$ and $P(A)=h(A)P(C)$, and reflection commutes with this congruence. Lemma~\ref{R4E3:minimum} and \eqref{R4E3:certificate} prove the upper bound. Its diagonal entries give $P(C)\geq1$.

For any $B\succeq0$, take the Schur product of the PSD upper defect with $\srcAUSch(B)\succeq0$. Since
\[
 \srcAUSch(A)\circ\srcAUSch(B)=\srcAUSch(A\circ B),\qquad
 I\circ\srcAUSch(B)=h(B)I,
\]
its quadratic form on the all-ones vector gives the first inequality in \eqref{R4E3:BS}. The classical permanent Hadamard inequality $P(B)\geq h(B)$, a consequence of Lieb's inequality \cite{R1E5:lieb}, gives the second because $P(A)\geq0$. Principal submatrices retain both hypotheses. Finally PSD of $A$ gives PSD of its Schur power and of $A\circ B$, proving the additional assertions.
\end{proof}

\section{Book matrices and the exact PSD lens}
Let
\begin{equation}\label{R4E3:book}
 C(c)=\begin{pmatrix}1&c&u\\\overline c&1&v\\u^*&v^*&I_m\end{pmatrix},
 \quad X=\sum_j|u_j|^2,\quad Y=\sum_j|v_j|^2,
 \quad r=\sum_j u_j\overline{v_j},\quad H=(1-X)(1-Y).
\end{equation}
Its support is contained in the book graph consisting of the two hubs, their edge, and $m$ independent leaves adjacent to the hubs. Every cycle contains both hubs, so no edge is disjoint from a cycle. A diagonal unitary changing the signs at the hubs carries $\srcAURefl(C(c))$ to $C(-c)$.

\begin{srcAUcorollary}\label{R4E3:lens}
If
\begin{equation}\label{R4E3:disk}
 X,Y\leq1,\qquad |c+r|^2\leq H,
\end{equation}
then $C(c)$ satisfies the upper Schur-power bound and \eqref{R4E3:BS}, whether or not $C(c)$ is PSD. It has full permanent-on-top if additionally $|c-r|^2\leq H$.
For a PSD $C(c)$ the sufficient condition is exactly
\begin{equation}\label{R4E3:lenscondition}
 |c|^2+|r|^2+2|\operatorname{Re}(c\overline r)|\leq H.
\end{equation}
In particular every PSD case with $\operatorname{Re}(c\overline r)\leq0$ is certified. If $r=0$, every PSD hub entry $c$ is certified.
\end{srcAUcorollary}
\begin{proof}
The Schur complement of $I_m$ gives
$C(c)\succeq0$ if and only if $X,Y\leq1$ and $|c-r|^2\leq H$.
Replace $c$ by $-c$ for the companion. Taking the maximum of the two squared absolute values gives \eqref{R4E3:lenscondition}. If the real part is nonpositive, $|c+r|^2\leq|c-r|^2$.
\end{proof}
The adjective ``exact'' concerns equivalence to simultaneous PSD, not necessity for permanent-on-top. No necessity claim for the lens is made.

\begin{srcAUcorollary}[Quantitative loss and equality]\label{R4E3:equality}
Suppose \eqref{R4E3:disk} holds, $m\geq2$, and all $u_j,v_j$ are nonzero. For every correlation matrix $B$,
\begin{equation}\label{R4E3:loss}
 P(C)-P(C\circ B)\geq
 2\sum_{j\ne k}|u_j|^2|v_k|^2
 \bigl(1-|b_{1,j+2}|^2|b_{2,k+2}|^2\bigr).
\end{equation}
Equality $P(C\circ B)=P(C)$ holds exactly when $B$ has rank one.
\end{srcAUcorollary}
\begin{proof}
The scalar identity in Proposition~\ref{R4E3:identity}, applied also after masking by $B$, gives
\[
 P(C)-P(C\circ B)=\det(\srcAURefl(C)\circ B)-\det(\srcAURefl(C))
 +2\sum_{|M|\text{ even},\ M\ne\varnothing}
 x_M\left(1-\prod_{ij\in M}|b_{ij}|^2\right).
\]
The determinant difference is nonnegative by Oppenheim's inequality. This also follows directly from Lemma~\ref{R4E3:minimum} by taking a Schur product with $\srcAUSch(B)$ and its quadratic form on the sign vector. In a book the nonempty even matchings have exactly two edges, giving \eqref{R4E3:loss}. If the permanent defect vanishes, all matching losses vanish. Each leaf edge has a disjoint leaf edge since $m\geq2$, so $|b_{1,j+2}|=|b_{2,j+2}|=1$ for every $j$. Equality in Gram-vector Cauchy--Schwarz and connectivity make every Gram vector collinear. Conversely a rank-one correlation mask is diagonal unitary conjugation and preserves the permanent. This determinant-plus-matching argument extends the cycle equality calculation in Round~3 entry~05; that underlying calculation is credited, not newly invented.
\end{proof}

\section{A spectral gap and dense permanent-on-top neighborhoods}
We now obtain an operator, rather than only scalar, stability statement. All matrices in this section are diagonal one. Let $\mathbf1$ and $\varepsilon=(\srcAUsgn\sigma)_\sigma$ be the constant and sign vectors in $\mathbb C^{S_n}$.

\begin{srcAUtheorem}\label{R4E3:gap}
For \eqref{R4E3:book} satisfying \eqref{R4E3:disk}, suppose $m\geq3$ and all leaf entries are nonzero. Put $n=m+2$, $N=n!/2$,
\[
 \begin{gathered}
 Z=\sum_j|u_j|^2|v_j|^2,\qquad w=\min_{j\ne k}|u_j|^2|v_k|^2,\\
 \delta=2(X+Y+|c|^2+|r|^2-Z),\qquad \gamma=\min\left\{\delta,\frac{4w}{N(N-1)}\right\}>0.
\end{gathered}
\]
On $\mathbf1^\perp$ one has
\begin{equation}\label{R4E3:spectralgap}
 P(C)I-\srcAUSch(C)\succeq\gamma I.
\end{equation}
Consequently, if a correlation matrix $A$ satisfies
\begin{equation}\label{R4E3:monomialball}
 \sum_{\pi\in S_n}|d_A(\pi)-d_C(\pi)|\leq\gamma/2,
\end{equation}
then $A$ has permanent-on-top. A sufficient entrywise condition is
$\max_{ij}|a_{ij}-c_{ij}|\leq\gamma/(2n\,n!)$.
If $C\succ0$, every diagonal-one Hermitian perturbation of radius strictly less than
\begin{equation}\label{R4E3:radius}
 \min\left\{\frac{\gamma}{2n\,n!},\frac{\lambda_{\min}(C)}{n}\right\}
\end{equation}
is positive definite and has permanent-on-top, including dense matrices.
\end{srcAUtheorem}
\begin{proof}
Let $g_{jk}=(1,j+2)(2,k+2)$ for $j\ne k$, and $L=\sum_{j\ne k}(I-L_{g_{jk}})$. The certificate gives $P(C)I-\srcAUSch(C)\succeq2wL$.
These generators generate $A_n$. Indeed, for three distinct leaves $j,k,l$, $g_{jk}g_{jl}$ is a three-cycle on hub $2$ and leaves $k,l$. Such three-cycles generate the alternating group on all vertices except hub $1$: a three-cycle not containing $2$ is a product of two containing $2$. For any three-cycle containing $1$, choose a leaf outside its support (possible since $m\geq3$) and a generator sending $1$ to that leaf. Conjugation carries the three-cycle to one fixing $1$. Hence all three-cycles lie in the generated group.

The Cayley graph of these involutions on each parity class is connected, has $N$ vertices, and diameter at most $N-1$. For a mean-zero function $f$ on one class, the identity
\[
 \sum_a|f(a)|^2=\frac1{2N}\sum_{a,b}|f(a)-f(b)|^2
\]
and Cauchy--Schwarz along a path of length at most $N-1$ show
$\|f\|^2\leq N(N-1)\langle f,Lf\rangle/2$.
Here each path energy is bounded by the total undirected edge energy, which equals $\langle f,Lf\rangle$.
Thus on $W=\operatorname{span}\{\mathbf1,\varepsilon\}^\perp$ the upper defect is at least $4w/[N(N-1)]$.

Every group-algebra operator preserves the constant and sign lines and their orthogonal complement. The eigenvalue of the upper defect on the sign line is $P(C)-\det C$. Direct permutation expansion gives
\[
 P(C)=1+X+Y+XY-2Z+|c+r|^2,\quad
 \det C=(1-X)(1-Y)-|c-r|^2,
\]
so this eigenvalue is $\delta$. Since the companion is PSD, $X,Y\leq1$, hence $Z\leq X$ and $\delta\geq2Y>0$. Combining the two orthogonal subspaces proves \eqref{R4E3:spectralgap}.

Write the sum in \eqref{R4E3:monomialball} as $L_0$. The regular expansion gives
$\|\srcAUSch(A)-\srcAUSch(C)\|\leq L_0$ and $|P(A)-P(C)|\leq L_0$.
Both upper defects annihilate $\mathbf1$. On its orthogonal complement their difference has norm at most $2L_0\leq\gamma$, so the new defect is PSD. Correlation entries have modulus at most one, and telescoping each product gives $L_0\leq n\,n!\max|a_{ij}-c_{ij}|$. Finally an entrywise perturbation of size $\eta$ has operator norm at most $n\eta$, proving positive definiteness under \eqref{R4E3:radius}. The constants are conservative and not claimed sharp.
\end{proof}

\section{Exact examples and a closure obstruction}
\paragraph{A new full-rank PSD block.}
Take $m=4$,
\[
 u=(1,1,1,1)/3,\quad v=(1,-1,i,-i)/3,\quad c=1/2.
\]
Then $r=0$, $X=Y=4/9$, and both $C(c)$ and $C(-c)$ are positive definite because their hub Schur complements have eigenvalues $5/9\pm1/2$. The single book block has intersecting triangles, and the triangle through the third leaf has nonreal product, so it is neither bipartite, a simple cycle, nor phase-balanced. It therefore lies outside the structural domain of Round~3 entry~05, regardless of its theorem's larger strength on the old domain. Here
\[
 P(C)=725/324,\quad \det C=19/324,\quad
 \delta=353/162,\quad w=1/81,\quad\gamma=1/2617110.
\]
One can check $C\succ I/100$ by its Schur complement: the least hub eigenvalue after subtracting $I/100$ and eliminating the leaves is
$49/100-400/891=3659/89100>0$.
Thus every diagonal-one Hermitian perturbation of entrywise size strictly below $1/22611830400$ is positive definite and permanent-on-top. This radius is small but explicit and full-dimensional; no general dense-matrix theorem is asserted.

\paragraph{Why the indefinite distinction is real.}
Take $m=3$, $u=v=(1,1,1)/2$ and $c=-3/4$. The companion $C(-c)$ is positive definite, with Schur complement $I_2/4$. But
\[
 P(C)=43/16,\quad \det C=-35/16.
\]
The sign vector is an eigenvector of $\srcAUSch(C)$ with negative eigenvalue $-35/16$. Thus Theorem~\ref{R4E3:reflection} gives a genuine upper operator bound for an indefinite Schur power, not full permanent-on-top.

\paragraph{Only explicit certified sums are automatically safe.}
The operator one-point identity of Round~3 entry~05 permits one-point sums of PSD matrices already known to be hereditarily permanent-on-top. Its proof uses products of commuting PSD embedded operators, and it applies to explicitly certified PSD book summands here. We do not extend it to indefinite summands, or claim that arbitrary PSD matrices whose blocks have book support meet the companion test after Gram splitting.

There is a concrete simultaneous-allocation obstruction. Let a root have diagonal $2/5$ and connect with entry $1/4$ to each of four leaves. The leaves have diagonal one, entries $1/2$ between the first pair and $-1/2$ between the second pair, and all other leaf entries zero. Both $A$ and $\srcAURefl(A)$ are positive definite: their required root Schur-complement costs are each $1/12+1/4=1/3<2/5$. But to split the two triangles into summands whose originals and companions are both PSD requires each root diagonal to be at least $\max(1/12,1/4)=1/4$. Their sum would be at least $1/2>2/5$. Hence the ordinary PSD Gram split does not guarantee a simultaneous companion split. This is an obstruction to that proof route, not a counterexample to permanent-on-top; the support is a cactus already covered in Round~3.

\section{Checks, current sources, and remaining gap}
The exact verifier independently checks the group-algebra certificate coefficient by coefficient, the book expansions, both displayed examples, the simultaneous-allocation obstruction, and the alternating-group generation in finite test orders. It checks upper bounds against exact PSD Gram masks, with no floating-point assertion. The universal certificate and spectral-gap proofs are written above and do not depend on those finite checks.

Bounded floating exploration found no violation of permanent-on-top in 1500 PSD five-by-five book examples, nor in 75 six-by-six unitary bipartite examples. First-compound tests on 240 bipartite examples of orders six through twelve also found no violation. These finite failures to find a counterexample are neither theorems nor evidence sufficient to remove the lens or support hypotheses. The unrestricted PSD book case with $\operatorname{Re}(c\overline r)>0$ and $|c+r|^2>H$ is the immediate unresolved target of this approach.

Primary-source checks on 20 October 2026 reconfirmed Pant--Singh v2 and Tran v5 (23 August 2022). Li--Zhao--Luo's latest posted v1 claims ordinary Chollet through order six \cite{R1E5:li}; it is not a premise and does not supply this stronger operator statement in unbounded order. The 2025 Sa-nguansin--Rodtes full text remained inaccessible \cite{R4E3:SR}; no exact priority comparison with all of its specialized classes is asserted. Searches for arrowhead, bipartite, orthogonal-vector and permanent-on-top results did not certify historical novelty. The precise verified increment is the reflected-companion certificate, its book/lens and indefinite consequences, and its operator gap. The classical lower bound and the originating Round~3 sign-certificate method retain their original credit.



\chapter[Polynomial-radius dense neighborhoods with permanent on top]{Polynomial-radius dense neighborhoods with permanent on top}
\label{source:R4E4}
\noindent\textit{Source development: Round 4, entry 04. Original mathematical arguments retained; notation is local to this chapter.}\par\medskip


\begin{quote}\small\noindent\textbf{Chapter abstract.} 
Let $n\ge5$, and let $A_0$ be a Hermitian matrix with diagonal one, cycle support, and every edge magnitude between $1/4$ and $1/3$. The phases are arbitrary. We prove that the permanent is the simple largest eigenvalue of its Schur power, with relative spectral gap at least $1/(128n^6)$. Every diagonal-one Hermitian matrix within entrywise distance $1/(8192n^7)$ of $A_0$ is positive definite and still has permanent on top. Thus full-dimensional neighborhoods containing dense, frustrated matrices have an explicit radius that decays only polynomially with the order. The proof quantitatively strengthens the attributed cycle operator certificate from Round 3: pairing even matchings produces an alternating-group walk, and an open-path expansion of permanental minors controls perturbations without a factorial loss. The unrestricted permanent-on-top assertion is false; no general permanent conjecture is resolved.
\end{quote}

\section{The statement and its exact increment}
For an $n\times n$ Hermitian matrix $A$, write
\[
 P(A)=\srcAVper(A),\qquad
 \srcAVSch(A)_{\sigma,\tau}=\prod_{i=1}^n a_{\sigma(i),\tau(i)}
 \quad(\sigma,\tau\in S_n),\qquad
 \srcAVDef(A)=P(A)I-\srcAVSch(A).
\]
Let $\Pi=J_{n!}/n!$, the orthogonal projection onto constant vectors. The rows of $\srcAVSch(A)$ sum to $P(A)$, so $\srcAVDef(A)\Pi=0$. For $A\succeq0$, the Schur power is PSD. The \emph{permanent-on-top} (POT) property is $\srcAVDef(A)\succeq0$. It fails for general complex PSD matrices, even at order five \cite{R4E4:Tran}. Nevertheless, it implies
\begin{equation}\label{R4E4:eq:BS}
 \srcAVper(A\circ B)\le P(A)\prod_i b_{ii}\le P(A)P(B)
 \qquad(B\succeq0),
\end{equation}
where the second inequality is classical \cite{R1E5:lieb}. The first universal assertion also fails \cite{R2E5:drury}; the weaker last comparison is Chollet's conjecture.

\begin{srcAVtheorem}[A polynomial-radius POT neighborhood]\label{R4E4:thm:main}
Suppose $n\ge5$ and $A_0$ is Hermitian with diagonal one and support exactly the cycle $C_n$. Assume
\[
 \frac14\le |(A_0)_{i,i+1}|\le\frac13
 \qquad(i\text{ cyclic}).
\]
There is no restriction on the edge phases. Then $A_0\succ0$ and
\begin{equation}\label{R4E4:eq:basegap}
 \srcAVDef(A_0)\succeq\frac{P(A_0)}{128n^6}(I-\Pi).
\end{equation}
Moreover, every diagonal-one Hermitian matrix $A$ satisfying
\begin{equation}\label{R4E4:eq:radius}
 \max_{i,j}|a_{ij}-(A_0)_{ij}|\le\frac1{8192n^7}
\end{equation}
is positive definite and obeys
\begin{equation}\label{R4E4:eq:robustgap}
 \srcAVDef(A)\succeq\frac{P(A_0)}{256n^6}(I-\Pi).
\end{equation}
In particular $A$ has POT, its permanent is a simple largest Schur-power eigenvalue, and \eqref{R4E4:eq:BS} holds for every PSD $B$.
\end{srcAVtheorem}
The ball is in the full real affine space of diagonal-one Hermitian matrices, not merely in the sparse subspace or inside an assumed PSD domain. Positivity of every member is part of the theorem. It contains dense matrices with arbitrary small chord phases. The constants and powers are sufficient, not claimed optimal.

\paragraph{Dependencies and the new contribution.}
Round 3 entry 05 proved POT for every complex PSD cycle by an explicit positive operator certificate \cite[Propositions 3.1 and 3.3]{R4E4:R305}. That qualitative sparse theorem, including its odd-cycle phase reflection, is an input here and is reproduced with proof in Section~\ref{R4E4:sec:inherited}. It is not counted again. Our Round 3 entry 04 proved quantitative scalar Bapat--Sunder bounds and dense neighborhoods with a conservative factorial denominator \cite[Theorem 5.1]{R4E4:R304}. That theorem did not establish POT in the dense neighborhoods.

The present increment is a \emph{relative operator spectral gap} and a \emph{polynomial-radius dense POT neighborhood}, valid uniformly over all phases and all nonuniform weights in the stated band. Two ingredients prevent exponential or factorial loss: the even-matching partition function retains a fixed fraction of the full permanent, and the total permanental-minor sensitivity is only of order $n$ in a diagonally dominant comparison matrix. The result does not subsume the arbitrary weights or singular boundary in entry 05's sparse theorem.

\paragraph{Source status.}
Pant--Singh's latest inspected version, v2 of 18 September 2026, states general Chollet as open and proves structured positive cases \cite{R2E3:PS}. The latest inspected Tran version is v5 of 23 August 2022, with its published counterexample \cite{R4E4:Tran}. Li--Zhao--Luo's public v1 preprint claims Chollet through order six \cite{R1E5:li}; it is not assumed or independently validated here. That weaker bounded-order claim does not imply our unbounded-order operator statement. Targeted primary-source searches on 20 October 2026 did not locate this precise quantitative neighborhood theorem; this is not a historical priority certificate. We seek a restricted structural advance, not a resolution of a currently open universal POT conjecture, since there is no such conjecture left to resolve.

\section{The inherited cycle certificate}\label{R4E4:sec:inherited}
Use the left regular representation $L_\pi e_\tau=e_{\pi\tau}$ and put
\[
 d_A(\pi)=\prod_i a_{\pi(i),i},\qquad
 \srcAVSch(A)=\sum_{\pi\in S_n}d_A(\pi)L_\pi.
\]
Let $D e_\sigma=\srcAVsgn(\sigma)e_\sigma$. Then
$DL_\pi D=\srcAVsgn(\pi)L_\pi$ and $L_\pi^*=L_{\pi^{-1}}$.

\begin{srcAVlemma}[Classical Schur minimum]\label{R4E4:lem:minimum}
For $A\succeq0$, $\srcAVSch(A)\succeq\det(A)I$.
\end{srcAVlemma}
\begin{proof}
A tensor Gram representation proves $\srcAVSch(A)\succeq0$, settling the singular case. Otherwise take a Cholesky factor $A=LL^*$ and use its rows as Gram vectors. In the tensor product representation, project onto basis tensors whose coordinate indices form a permutation. The coefficient associated with $\pi$ in the tensor for $\sigma$ is $\prod_iL_{\sigma(i),\pi(i)}$. It vanishes unless $\pi(i)\le\sigma(i)$ for every $i$, hence unless $\pi=\sigma$, because the sums of these indices agree. The projected vectors have Gram matrix $\det(A)I$. Their orthogonal residuals give the claimed PSD difference. This classical argument, also supplied in \cite[Lemma 2.1]{R4E4:R305}, is not new.
\end{proof}
For a cycle correlation matrix $C$, put $x_e=|c_e|^2$, $x_M=\prod_{e\in M}x_e$, and let $\pi_M$ be the product of the disjoint edge transpositions in a matching $M$. Set $c=(1\,2\,\cdots\,n)$, $z=d_C(c)$, and $Z=\sum_Mx_M$. Permutation enumeration gives
\begin{align}
 \srcAVSch(C)&=\sum_M x_ML_{\pi_M}+zL_c+\overline zL_c^*,\label{R4E4:eq:cycleS}\\
 P(C)&=Z+2\Re z,\qquad
 \det C=\sum_M(-1)^{|M|}x_M+(-1)^{n-1}2\Re z.\label{R4E4:eq:cyclePD}
\end{align}

\begin{srcAVlemma}[Cycle matching lower certificate, attributed]\label{R4E4:lem:cyclecert}
Every PSD correlation cycle satisfies
\begin{equation}\label{R4E4:eq:R}
 \srcAVDef(C)\succeq R:=2\sum_{\substack{M\text{ a matching}\\|M|\text{ even},\ M\ne\varnothing}}
 x_M(I-L_{\pi_M})\succeq0.
\end{equation}
\end{srcAVlemma}
\begin{proof}
For even $n$, the exact identity from \cite[Proposition 3.1]{R4E4:R305} is
\[
 \srcAVDef(C)=D\bigl(\srcAVSch(C)-\det(C)I\bigr)D+R.
\]
It follows immediately by comparing coefficients in \eqref{R4E4:eq:cycleS}--\eqref{R4E4:eq:cyclePD}: the full cycle and odd matchings have negative sign, and the remaining even matchings supply $R$.

For odd $n$, change one edge phase to obtain cycle product
$z'=|\Re z|-i\Im z$, preserving its modulus. Call the new matrix $C'$. Every proper principal subgraph is a forest, so its determinant is unchanged by edge phase changes. The full determinant increases by $2(|\Re z|-\Re z)\ge0$. Thus all principal minors of $C'$ are nonnegative, which implies $C'\succeq0$. One justification of the last implication is to apply Sylvester's criterion to $C'+\epsilon I$ and let $\epsilon\downarrow0$. If $z=0$, take $C'=C$.

Put $t=z+z'=2\max(\Re z,0)\ge0$. The attributed odd-cycle identity is
\[
 \srcAVDef(C)=D\bigl(\srcAVSch(C')-\det(C')I\bigr)D
       +R+t(2I-L_c-L_c^*).
\]
Matching coefficients agree as before; the full-cycle coefficient is $z'-t=-z$. Identity coefficients agree by \eqref{R4E4:eq:cyclePD}. Schur minimum makes the first term PSD. Each $\pi_M$ is an involution, while $2I-L_c-L_c^*=(I-L_c)^*(I-L_c)$. All remaining terms are PSD. This proves \eqref{R4E4:eq:R} and restates precisely the earlier input.
\end{proof}

\section{Even matchings retain a relative spectral gap}
Let $\mathcal G$ be the set of double transpositions $s=t_et_f$ over unordered pairs of disjoint edges of $C_n$, and put
\[
 k=|\mathcal G|=\frac{n(n-3)}2,\qquad
 Q=\sum_{s\in\mathcal G}(I-L_s),\qquad d=\frac{n(n-1)}2.
\]
Let $\varepsilon=(\srcAVsgn\sigma)_{\sigma\in S_n}$ and let $\Pi_{\rm ev}$ be projection onto $\operatorname{span}\{\mathbf1,\varepsilon\}$.

\begin{srcAVlemma}[A fully elementary alternating-group gap]\label{R4E4:lem:walk}
For every $n\ge5$,
\begin{equation}\label{R4E4:eq:walkgap}
 Q\succeq d^{-2}(I-\Pi_{\rm ev}).
\end{equation}
\end{srcAVlemma}
\begin{proof}
Every even permutation has an even-length expression of length at most $d$ in the adjacent transpositions $(1\,2),\ldots,(n-1\,n)$, by sorting inversions. Group this word into consecutive pairs. Equal transpositions cancel; disjoint ones give a member of $\mathcal G$. If two consecutive transpositions overlap, their union consists of three consecutive cycle vertices. Since $n\ge5$, a cycle edge $h$ disjoint from those three vertices exists, and
\[
 t_et_f=(t_et_h)(t_ht_f).
\]
Both right-hand factors belong to $\mathcal G$. Thus every element of $A_n$ has a word of length at most $d$ in $\mathcal G$. In particular these generators generate all of $A_n$, not a smaller subgroup.

For a unitary representation, telescoping a word $g=s_1\cdots s_\ell$ and applying Cauchy--Schwarz gives
\[
 \|(I-L_g)v\|^2\le\ell\sum_{j=1}^{\ell}\|(I-L_{s_j})v\|^2
 \le 2d^2\langle v,Qv\rangle.
\]
The telescoping prefixes are unitary; no commutation of overlapping generators is asserted. Average over $g\in A_n$. The left side becomes
$2\|(I-\Pi_{\rm ev})v\|^2$, because averaging the left action of $A_n$ projects onto the two parity cosets. This proves \eqref{R4E4:eq:walkgap}. No external interchange-process spectral-gap theorem is needed.
\end{proof}

\begin{srcAVlemma}[Pairing even matchings]\label{R4E4:lem:pairing}
For the matrix $A_0$ in Theorem~\ref{R4E4:thm:main}, with matching partition function $Z$, its operator $R$ in \eqref{R4E4:eq:R} satisfies
\begin{equation}\label{R4E4:eq:Rgap}
 R\succeq\frac{Z}{256k(10/9)^6}Q.
\end{equation}
\end{srcAVlemma}
\begin{proof}
Fix $s=t_et_f\in\mathcal G$. Let $H_s$ be the cycle subgraph obtained by retaining only edges incident to none of the four endpoints of $e,f$. At most six edges are removed. It is a forest, with isolated vertices allowed. For every even matching $N$ of $H_s$, both $N$ and $N\cup\{e,f\}$ are even matchings of the cycle. Set $U=L_{\pi_N}$, $T=L_s$, and $a=x_ex_f\le1$. These are commuting involutions, and
\begin{equation}\label{R4E4:eq:invol}
 (I-U)+a(I-UT)\succeq a(I-T).
\end{equation}
Indeed simultaneous eigenvalues $u,t\in\{1,-1\}$ reduce this to four scalar nonnegative inequalities. After multiplication by $2x_N$ and summation, the paired part $R_s$ of $R$ obeys
\[
 R_s\succeq2x_ex_f Z_{\rm even}(H_s)(I-L_s).
\]
Here the empty matching's zero operator is allowed. For a fixed $s$ a matching occurs in at most one pair, so $R_s$ is a subsum of $R$. Therefore $\sum_sR_s\preceq kR$; all summands in $R$ are PSD. This counting is essential: the same matching may be used for different $s$, and the factor $k$ pays for that reuse.

The forest correlation matrix with edge magnitudes $|(A_0)_e|$ is strictly diagonally dominant, since the sum at each vertex is at most $2/3$. Its determinant is nonnegative and equals $Z_{\rm even}(H_s)-Z_{\rm odd}(H_s)$. Hence
\[
 Z_{\rm even}(H_s)\ge\tfrac12 Z(H_s).
\]
If $F$ is the removed edge set, every matching splits into its parts in $H_s$ and in $F$. Dropping compatibility restrictions gives
\[
 Z\le Z(H_s)\prod_{e\in F}(1+x_e)
 \le Z(H_s)(10/9)^6.
\]
Finally $x_ex_f\ge(1/16)^2=1/256$. Combining these bounds and summing over $s$ proves \eqref{R4E4:eq:Rgap}.
\end{proof}

\begin{proof}[Proof of the base gap \eqref{R4E4:eq:basegap}]
Strict diagonal dominance gives $A_0\succ0$. Since $|z|\le3^{-n}$ and $Z\ge1+n/16$,
\[
 P(A_0)\ge1+\frac n{16}-\frac2{3^n}>\frac54,
 \qquad \frac Z{P(A_0)}\ge\frac12.
\]
The first bound follows already at $n=5$ and then increases. The second follows from $P\le Z+2\cdot3^{-n}\le2Z$. On the orthogonal complement of both constants and signs, Lemmas~\ref{R4E4:lem:cyclecert}--\ref{R4E4:lem:pairing} yield relative gap at least
\[
 \frac1{512k(10/9)^6d^2}\ge\frac1{128n^6},
\]
because $k\le n^2/2$, $d\le n^2/2$, and $(10/9)^6<2$.

The sign vector is an eigenvector of $\srcAVSch(A_0)$ with eigenvalue $\det(A_0)$. Hadamard's determinant inequality gives $\det(A_0)\le1$, so its defect eigenvalue is
\[
 P(A_0)-\det(A_0)\ge P(A_0)-1\ge\frac15P(A_0).
\]
For completeness, the determinant bound follows by Gram--Schmidt: the volume squared of unit Gram vectors is at most the product of their squared lengths. Constants, signs, and their common orthogonal complement are invariant for the Hermitian regular-representation operator. Combining the two bounds proves \eqref{R4E4:eq:basegap}; the constant-vector defect is zero.
\end{proof}

\begin{srcAVremark}[Sharper equal-weight specialization]
If every edge magnitude is exactly $1/4$, the same proof uses $(17/16)^6<3/2$ in place of $(10/9)^6<2$ and gives the relative base gap $1/(96n^6)$. Uniformity is unnecessary for the main result; the band version was chosen to include independently varying weights.
\end{srcAVremark}

\section{Relative perturbation control by open paths}
The next estimate is independent of the cycle operator certificate. It is the second quantitative ingredient.

\begin{srcAVlemma}[Open-path expansion of a permanental minor]\label{R4E4:lem:paths}
Let $M$ be a symmetric nonnegative matrix with diagonal one, and write $W=M-I$. Suppose every row sum of $W$ is at most $q<1$. Then
\begin{equation}\label{R4E4:eq:minorbound}
 \sum_{i\ne j}\srcAVper M(i\mid j)\le\frac{nq}{1-q}\srcAVper M
 \le\frac n{1-q}\srcAVper M.
\end{equation}
Here $M(i\mid j)$ deletes row $i$ and column $j$.
\end{srcAVlemma}
\begin{proof}
For $i\ne j$, a permutation monomial in the minor has a unique directed simple path from $j$ to $i$: $j$ has no incoming edge, $i$ has no outgoing edge, and all other vertices have equal in- and out-degree. The remaining vertices carry disjoint permutation cycles. Consequently
\begin{equation}\label{R4E4:eq:pathid}
 \srcAVper M(i\mid j)=\sum_{\substack{\gamma:j\to i\\\gamma\text{ simple}}}
       \left(\prod_{e\in\gamma}W_e\right)\srcAVper M[V\setminus V(\gamma)].
\end{equation}
Every principal permanent on the right is at most $\srcAVper M$: append fixed points of weight one to its permutations, and use nonnegativity. Simple paths are a subset of nontrivial walks. The total weight of walks of length $\ell$ from all starts is at most $nq^\ell$, by the row-sum bound. Summing over $\ell\ge1$ proves \eqref{R4E4:eq:minorbound}. This proof does not require $M$ to be PSD.
\end{proof}

\begin{srcAVproposition}[Relative coefficient sensitivity]\label{R4E4:prop:perturb}
Let $A_0$ be diagonal-one Hermitian, let $C=|A_0|$ entrywise, and suppose its off-diagonal row sums are at most $q_0$. For a diagonal-one Hermitian $A$ put
\[
 \epsilon=\max_{ij}|a_{ij}-(A_0)_{ij}|,
 \quad q=q_0+(n-1)\epsilon<1,
 \quad L=\sum_{\pi\in S_n}|d_A(\pi)-d_{A_0}(\pi)|.
\]
Then both $A_0,A$ are positive definite and
\begin{equation}\label{R4E4:eq:coefbound}
 L\le P(C)\left[\exp\left(\frac{n\epsilon}{1-q}\right)-1\right],
 \qquad \|\srcAVDef(A)-\srcAVDef(A_0)\|\le2L.
\end{equation}
\end{srcAVproposition}
\begin{proof}
The row sums of absolute off-diagonal entries are less than one, so positivity follows, for example, by bounding each cross term with $2|x_ix_j|\le|x_i|^2+|x_j|^2$ in the quadratic form. Let $E$ have zero diagonal and all off-diagonal entries $\epsilon$. Expanding each product and applying the triangle inequality gives
\[
 L\le\srcAVper(C+E)-\srcAVper C.
\]
For $M(t)=C+tE$, $0\le t\le1$, differentiation and Lemma~\ref{R4E4:lem:paths} imply
\[
 \frac{d}{dt}\srcAVper M(t)
 =\epsilon\sum_{i\ne j}\srcAVper M(t)(i\mid j)
 \le\frac{n\epsilon}{1-q}\srcAVper M(t).
\]
The permanent is positive because its diagonal monomial is one. Integrating the logarithmic derivative proves the first bound. For the second, $|P(A)-P(A_0)|\le L$ and
\[
 \|\srcAVSch(A)-\srcAVSch(A_0)\|
 \le\sum_\pi|d_A(\pi)-d_{A_0}(\pi)|\,\|L_\pi\|=L.
\]
Add these two bounds. The useful feature is relative control through $P(C)$, rather than $n!$ times a worst individual error.
\end{proof}

\begin{proof}[Proof of the neighborhood assertion]
For the cycle center, $q_0\le2/3$. Under \eqref{R4E4:eq:radius},
$q\le2/3+(n-1)\epsilon<3/4$, so Proposition~\ref{R4E4:prop:perturb} proves positive definiteness. Also
\[
 P(C)=Z+2|z|,\qquad P(A_0)\ge Z-2|z|,
 \qquad P(C)\le2P(A_0),
\]
since $Z\ge1\ge6\cdot3^{-n}$. Let $u=4n\epsilon\le1/2$. The elementary bound $e^u-1\le2u$ on this interval gives
\[
 L\le2P(A_0)(e^{4n\epsilon}-1)
 \le16n\epsilon P(A_0).
\]
For instance, the exponential bound follows from its power series and $1/j!\le1$ for $j\ge1$: $e^u-1\le u/(1-u)\le2u$.
Thus
\[
 \|\srcAVDef(A)-\srcAVDef(A_0)\|\le32n\epsilon P(A_0)
 \le\frac{P(A_0)}{256n^6}.
\]
Both defect operators annihilate constants, and are Hermitian, so their difference has no cross terms with that one-dimensional subspace. Apply the norm estimate on its orthogonal complement and use \eqref{R4E4:eq:basegap}. This proves \eqref{R4E4:eq:robustgap}, not a bound that erroneously subtracts a positive multiple of the identity on the constant vector.
\end{proof}

\section{Quantitative scalar consequences and limitations}
\begin{srcAVcorollary}[An unrestricted second factor and rigidity]\label{R4E4:cor:scalar}
Let $A$ satisfy Theorem~\ref{R4E4:thm:main}, and let $B$ be any correlation matrix of the same order. Then
\begin{equation}\label{R4E4:eq:scalar}
 P(A)-P(A\circ B)\ge
 \frac{P(A_0)}{256n^6}\left(1-\frac{P(B)}{n!}\right)\ge0.
\end{equation}
Equality $P(A\circ B)=P(A)$ holds if and only if $B$ has rank one. For a general PSD second factor, \eqref{R4E4:eq:BS} follows by diagonal scaling, with zero-diagonal cases understood directly.
\end{srcAVcorollary}
\begin{proof}
Take the Schur product of \eqref{R4E4:eq:robustgap} with the PSD matrix $\srcAVSch(B)$. The identities
\[
 \srcAVSch(A)\circ\srcAVSch(B)=\srcAVSch(A\circ B),\quad
 I\circ\srcAVSch(B)=I,\quad \Pi\circ\srcAVSch(B)=\srcAVSch(B)/n!
\]
then give \eqref{R4E4:eq:scalar} by evaluating on the constant vector and dividing by its squared norm. Every permutation monomial of a correlation matrix has modulus at most one, so $0\le P(B)\le n!$. If $P(B)=n!$, equality in the real-part bound for this sum forces every monomial to be one. Transposition monomials give $|b_{ij}|=1$ for all $i,j$, hence all unit Gram vectors are collinear and $B$ has rank one. Conversely such a mask is diagonal unitary conjugation and leaves the permanent invariant. A zero diagonal in $B$ forces its entire corresponding row to vanish; otherwise scale to diagonal one. The final comparison $P(B)\ge\prod b_{ii}$ is the classical permanent Hadamard inequality \cite{R1E5:lieb}.
\end{proof}

For a concrete dense family, take cycle entries $1/4$ except for one edge of phase $i$, and add the same positive rational number $\epsilon\le1/(8192n^7)$ to every noncycle off-diagonal pair. All entries are nonzero, the matrix is positive definite, and its cycle product remains nonreal, so it is not phase-balanced. The operator theorem applies at every order $n\ge5$, although the support is now complete. Exact rational versions, with $i$ as a Gaussian unit, are included in the verifier.

\paragraph{What has and has not improved.}
The previous sparse cycle theorem remains the essential qualitative input. We do not claim a new proof of all sparse cases or take credit for its phase reflection. Our earlier dense theorem had radius proportional to $1/(n^3n!)$ and only a scalar conclusion; here the radius is polynomial in $n$ and the conclusion is the full Schur-power operator inequality. The centers are restricted to a fixed band bounded away from zero and from loss of diagonal dominance. Therefore this is an asymptotic and operator-strength improvement on this common subfamily, not a uniform containment of every previous neighborhood.

The exponent seven and constant $8192$ are deliberately coarse. The word estimate for the alternating-group walk is the main quantitative loss. The method does not supply a comparable uniform radius when an edge tends to zero or the center approaches singularity. For $n=4$, the disjoint cycle-edge double swaps generate only the Klein four group, so our alternating-group gap proof genuinely fails. We make no assertion that the order-four POT property itself fails. The proof does not cover a fixed-size perturbation independent of $n$, arbitrary diagonally dominant matrices, or general Chollet.

\paragraph{Checks and research record.}
The companion exact verifier checks permutation-word constructions, the matching-pair budget, the forest determinant and partition comparisons, cycle group-algebra identities, the open-path minor identity, and dense Gaussian-rational examples. Small floating spectral tests are separately labeled diagnostics and are not proof premises. The all-order proof consists of the displayed inequalities and group argument. A first attempt based only on diagonal shifts gave an exponentially small gap; keeping the full even-matching partition function removed that loss. A termwise perturbation bound would retain a factorial denominator; open-path minor control removes it. No unbounded numerical search or external posting was undertaken.



\chapter[Quantitative permanent-on-top gaps and dense operator neighborhoods]{Quantitative permanent-on-top gaps 
and dense operator neighborhoods}
\label{source:R4E5}
\noindent\textit{Source development: Round 4, entry 05. Original mathematical arguments retained; notation is local to this chapter.}\par\medskip

\begin{quote}\small\noindent\textbf{Chapter abstract.} 
A repeated-coordinate tensor sector refines the classical Schur-minimum
inequality by an explicit positive transposition operator. Combined with
Round 3's permanent-on-top certificates, it gives a weighted interchange
Laplacian below the permanent defect for every connected positive definite
matrix in the previously proved structural class. We propagate strictly
positive edge weights through cut vertices and obtain a simple top
eigenvalue, a computable absolute spectral gap, and explicit dense
neighborhoods satisfying the full permanent-on-top operator inequality.
A classical theorem of Caputo--Liggett--Richthammer identifies the
comparison gap using only an order-$n$ weighted graph Laplacian; an
independent elementary, weaker estimate is also proved. The scalar
consequences include quantitative Bapat--Sunder bounds and a rank-one
equality criterion. A rational, phase-frustrated $K_{2,4}$ center has a
certified full-dimensional dense neighborhood. Positive definiteness
cannot simply be dropped: an exact connected singular example has a
three-dimensional top eigenspace. This is a quantitative extension of
previous structural results, not a resolution of general Chollet or of
the already false unrestricted permanent-on-top assertion. Historical
priority for the new refinement and combination remains unverified.
\end{quote}

\section{The improvement and main statement}
For a complex Hermitian matrix $A$ of order $n$, write
\begin{align*}
 P(A)&=\srcAWper(A),& h(A)&=\prod_i a_{ii},\\
 \srcAWSch(A)_{\sigma,\tau}&=\prod_{i=1}^n a_{\sigma(i),\tau(i)},&
 \srcAWDef(A)&=P(A)\srcAWId_{n!}-\srcAWSch(A)
 \quad(\sigma,\tau\in\srcAWG_n).
\end{align*}
The empty permanent and determinant are $1$.
For PSD $A$, $\srcAWSch(A)$ is PSD, being the Gram matrix of permuted
tensor products of Gram vectors of $A$.
The permanent-on-top property (POT) is $\srcAWDef(A)\succeq0$.
The general POT assertion is false \cite{R4E3:Tran}; the general Chollet
inequality $P(A\circ B)\le P(A)P(B)$ remains open in the latest checked
structural source \cite{R2E3:PS}. We prove sufficient conditions, not a
universal assertion.

The support graph has an edge $ij$ exactly when $i\ne j$ and
$a_{ij}\ne0$. A block is a maximal two-connected subgraph, a bridge,
or an isolated vertex. A block is \emph{phase-balanced} if a diagonal
unitary congruence makes its principal matrix entrywise nonnegative.
Let $\srcAWClass$ be the structural class from Round 3, Entry 05:
every support block is either
\begin{enumerate}
\item phase-balanced;
\item a simple cycle of length at least three;
\item bipartite, with no edge vertex-disjoint from any simple cycle.
\end{enumerate}
The last condition includes $K_{2,m}$ with arbitrary complex phases.
It concerns all simple cycles, not only induced cycles.
We use $A\succ0$ for positive definiteness.

Let $L_\pi e_\tau=e_{\pi\tau}$ be the left regular representation.
For positive edge weights $c=(c_{ij})$, define
\[
 K_c=\sum_{ij\in E}c_{ij}(\srcAWId-L_{(ij)}),\qquad
 \mathcal L_c=\sum_{ij\in E}c_{ij}(e_i-e_j)(e_i-e_j)^*,\qquad
 Q=\srcAWId_{n!}-\frac{\mathbf1\mathbf1^*}{n!}.
\]
The second-smallest eigenvalue of $\mathcal L_c$ is denoted
$\lambda_2(\mathcal L_c)$.

\begin{srcAWtheorem}\label{R4E5:thm:main}
Let $n\ge2$ and $A\succ0$ belong to $\srcAWClass$, with connected support.
The explicit recursion in Sections~\ref{R4E5:sec:local}--\ref{R4E5:sec:glue}
produces a strictly positive weight on every support edge such that
\begin{equation}\label{R4E5:eq:comparison}
 \srcAWDef(A)\succeq K_c\succeq gQ,
 \qquad g=\lambda_2(\mathcal L_c)>0.
\end{equation}
Consequently $P(A)$ is a simple largest eigenvalue of $\srcAWSch(A)$ and
its separation from every other eigenvalue is at least $g$.
For every PSD $B$ of order $n$,
\begin{align}
 P(A)h(B)-P(A\circ B)
 &\ge\sum_{ij\in E}c_{ij}
 \left(h(B)-|b_{ij}|^2\prod_{k\ne i,j}b_{kk}\right),\label{R4E5:eq:scalar}\\
 P(A)h(B)-P(A\circ B)
 &\ge g\left(h(B)-\frac{P(B)}{n!}\right).\label{R4E5:eq:scalarGap}
\end{align}
For a correlation matrix $B$, equality $P(A\circ B)=P(A)$ holds
if and only if $B$ has rank one.
\end{srcAWtheorem}

The first comparison in \eqref{R4E5:eq:comparison} is proved here. The sharp
second comparison uses the classical weighted interchange-gap theorem
\cite[Theorem 1.1]{R4E5:CLR}, not a new graph spectral theorem. For a fully
elementary alternative, if $T$ is a spanning tree of diameter $d$ and
$c_* = \min_{e\in T}c_e$, the explicit number
\begin{equation}\label{R4E5:eq:elementarygap}
 g_{\rm elem}=\frac{c_*}{2(n-1)^2(2d-1)}
\end{equation}
can replace $g$ everywhere; Section~\ref{R4E5:sec:gap} proves this directly.

\begin{srcAWtheorem}[Dense POT stability]\label{R4E5:thm:stable}
Suppose $A_0\succ0$ has $\srcAWDef(A_0)\succeq gQ$ for $g>0$.
For Hermitian $A$, put
\[
 d_A(\pi)=\prod_i a_{\pi(i),i},\qquad
 \ell(A,A_0)=\sum_{\pi\in\srcAWG_n}|d_A(\pi)-d_{A_0}(\pi)|.
\]
Then
\begin{equation}\label{R4E5:eq:stable}
 \srcAWDef(A)\succeq\bigl(g-2\ell(A,A_0)\bigr)Q.
\end{equation}
In particular, every PSD $A$ with $2\ell(A,A_0)\le g$ has POT,
without any support or phase condition on $A$.
If $A_0$ and $A$ are correlation matrices and
$\max_{ij}|a_{ij}-a^0_{ij}|\le\varepsilon$, it suffices that
\begin{equation}\label{R4E5:eq:roughradius}
 2n\,n!\,\varepsilon\le g.
\end{equation}
A sharper sufficient condition, also valid for noncorrelation $A$, is
\begin{equation}\label{R4E5:eq:polyradius}
 2\bigl[P(|A_0|+\varepsilon J_n)-P(|A_0|)\bigr]\le g,
\end{equation}
where $|A_0|$ is entrywise absolute value and $J_n$ is all ones.
\end{srcAWtheorem}

For an unconstrained Hermitian perturbation of a correlation center,
$\varepsilon<\lambda_{\min}(A_0)/n$ ensures positive definiteness.
Thus \eqref{R4E5:eq:polyradius}, together with this strict inequality,
certifies a full-dimensional open neighborhood, not only perturbations
preserving zeros. The quantity $g$ is an \emph{absolute} Schur-power gap;
the gap relative to the permanent is $g/P(A_0)$. No dimension-free
relative gap or dimension-free entrywise radius is asserted.

\paragraph{Precisely what is inherited.}
Round 3 Entry 05 proved POT on $\srcAWClass$, its operator one-point identity,
and the cycle sign-reflection certificates \cite{R4E5:R305}.
Round 3 Entry 04 independently proved quantitative scalar Bapat--Sunder
bounds for long cycles and explicit dense \emph{scalar} neighborhoods
\cite[Theorem 5.1]{R4E5:R304}; this supplies an important stability precedent.
Here the new-in-this-work steps are an explicit refinement of the
Schur-minimum operator, positive transposition-weight propagation, and
full \emph{operator} neighborhoods throughout the positive definite,
connected part of $\srcAWClass$. We reproduce the needed earlier identities
with proofs rather than count them again. The $K_{2,m}$ scalar conclusion
also independently overlaps Round 3 Entry 03; its explicit defect formula
is credited when used below \cite{R4E5:R303}. Historical priority for our
refinement and combination is not established by a search without a hit.

\section{A quantitative Schur-minimum refinement}\label{R4E5:sec:minimum}
The regular-representation formula and sign involution are
\begin{equation}\label{R4E5:eq:regular}
 \srcAWSch(A)=\sum_{\pi}d_A(\pi)L_\pi,\qquad
 D_{\sigma,\sigma}=\srcAWsgn(\sigma),\qquad
 DL_\pi D=\srcAWsgn(\pi)L_\pi.
\end{equation}
Indeed the $(\sigma,\tau)$ coefficient in the sum is
$d_A(\sigma\tau^{-1})=\prod_i a_{\sigma(i),\tau(i)}$.
Every row sum is $P(A)$, so $\srcAWDef(A)\mathbf1=0$.

\begin{srcAWlemma}[A two-row sector]\label{R4E5:lem:pair}
For every $A\succ0$ and distinct $i,j$,
\begin{equation}\label{R4E5:eq:pair}
 \srcAWSch(A)-\det(A)\srcAWId
 \succeq \beta_{ij}(A)(\srcAWId+L_{(ij)}),\qquad
 \beta_{ij}(A)=\frac{\det(A)|a_{ij}|^2}
 {a_{ii}a_{jj}-|a_{ij}|^2}.
\end{equation}
Consequently, for any $\theta_{ij}\ge0$ with
$\sum_{i<j}\theta_{ij}\le1$,
\begin{equation}\label{R4E5:eq:convex}
 D\bigl(\srcAWSch(A)-\det(A)\srcAWId\bigr)D
 \succeq\sum_{i<j}\theta_{ij}\beta_{ij}(A)
 (\srcAWId-L_{(ij)}).
\end{equation}
\end{srcAWlemma}
\begin{proof}
Relabel the selected pair as $1,2$ and write $A=LL^*$ with $L$ lower
triangular and positive diagonal. Use its rows $u_r$ as Gram vectors
and $v_\sigma=\bigotimes_k u_{\sigma(k)}$.
Project first onto coordinate words using each of $1,\ldots,n$ once.
A coefficient $\prod_k L_{\sigma(k),\pi(k)}$ can be nonzero only if
$\pi(k)\le\sigma(k)$ for all $k$. Equality of the sums forces
$\pi=\sigma$. This sector contributes exactly $\det(A)\srcAWId$ to the
Gram matrix. It also proves the classical Schur-minimum bound; for
singular PSD matrices that bound follows from $\srcAWSch(A)\succeq0$.

Take a second, orthogonal sector: coordinate $1$ occurs twice,
coordinate $2$ is absent, and every coordinate $3,\ldots,n$ occurs
once. A word in this sector with nonzero coefficient must assign
coordinate $n$ to row $n$, since all smaller rows have a zero there.
Descending induction assigns coordinate $r$ to row $r$ for every
$r\ge3$. Only rows $1,2$ remain, and both use coordinate $1$.
Each word thus receives exactly two permuted tensors, whose labels
differ by left multiplication by $(12)$, with the identical amplitude
\[
 c=L_{11}L_{21}\prod_{r=3}^nL_{rr}.
\]
Its Gram contribution over all such words is
$|c|^2(\srcAWId+L_{(12)})$. The remaining orthogonal sectors have a PSD
Gram matrix. Finally
\[
 |c|^2=\det(A)\frac{|L_{21}|^2}{L_{22}^2}
 =\det(A)\frac{|a_{21}|^2}{a_{11}a_{22}-|a_{21}|^2}.
\]
Simultaneous relabeling conjugates the regular operators and proves
\eqref{R4E5:eq:pair}. The denominators are positive by positive definiteness.
After sign conjugation each pair bound has the same left side $H\succeq0$.
Writing $H$ as the sum of the weighted copies plus
$(1-\sum\theta_{ij})H$ proves \eqref{R4E5:eq:convex}.
In particular, the separate pair bounds are \emph{not} added without
convex weights: their tensor sectors use different Cholesky orderings.
\end{proof}

\section{Local positive weights from the inherited certificates}\label{R4E5:sec:local}
Diagonal unitary congruence leaves both $P(A)$ and $\srcAWSch(A)$ unchanged.
For $A=E C E$ with $C$ a correlation matrix and $E$ positive diagonal,
these two quantities are $h(A)$ times their counterparts for $C$.
Thus it suffices to describe the local certificates at diagonal one.
For a matching $M$ let
$x_M=\prod_{ij\in M}|c_{ij}|^2$ and
$\pi_M=\prod_{ij\in M}(ij)$, including $x_\varnothing=1$.

\begin{srcAWlemma}[Local sign certificates, reused from Round 3]\label{R4E5:lem:local}
Let $C\succ0$ be a correlation matrix on a single permitted block.
If its support is bipartite and has no edge disjoint from a simple
cycle, then
\begin{equation}\label{R4E5:eq:even}
 \srcAWDef(C)=D\bigl(\srcAWSch(C)-\det(C)\srcAWId\bigr)D
 +2\sum_{\substack{M\ne\varnothing\\|M|\ {\rm even}}}
 x_M(\srcAWId-L_{\pi_M}).
\end{equation}
The same identity applies to an even cycle.
For an odd cycle with cyclic permutation $c$ and $z=d_C(c)$,
there exists $C'\succ0$ with the same diagonal and entry moduli such that
\begin{align}
 &z'=d_{C'}(c)=|\srcAWRea z|-i\srcAWIma z,\qquad
 \det(C')\ge\det(C),\qquad t=z+z'=2\max(\srcAWRea z,0),\label{R4E5:eq:reflection}\\
 &\srcAWDef(C)=D\bigl(\srcAWSch(C')-\det(C')\srcAWId\bigr)D
 +2\sum_{\substack{M\ne\varnothing\\|M|\ {\rm even}}}
 x_M(\srcAWId-L_{\pi_M})+t(2\srcAWId-L_c-L_c^*).\label{R4E5:eq:odd}
\end{align}
All displayed remainder terms are PSD.
\end{srcAWlemma}
\begin{proof}
In the bipartite case, a supported permutation is a matching involution
or consists of exactly one even simple cycle with every other vertex
fixed. A disjoint transposition or a second nontrivial cycle would
supply an edge disjoint from the first cycle. Every nonmatching term
therefore has negative sign; adding $\srcAWSch(C)$ and $D\srcAWSch(C)D$ leaves
exactly twice the even-matching terms. Their row sums similarly equal
$P(C)+\det(C)$. Subtraction proves \eqref{R4E5:eq:even}.

On a single cycle the only terms are matchings and its two full
orientations. The formulas are
\[
 P(C)=\sum_Mx_M+2\srcAWRea z,\qquad
 \det(C)=\sum_M(-1)^{|M|}x_M+(-1)^{m-1}2\srcAWRea z.
\]
For odd $m$, $|z'|=|z|$, so changing the phase of one edge realizes
$z'$ without changing any modulus. Every proper principal subgraph
is a forest: all its determinant terms are matchings and are unchanged.
The full determinant increases by $2(|\srcAWRea z|-\srcAWRea z)\ge0$.
The principal-minor criterion gives $C'\succ0$ since $C\succ0$.
Coefficient comparison using \eqref{R4E5:eq:regular} gives \eqref{R4E5:eq:odd}.
Every $L_{\pi_M}$ is a self-adjoint involution, and
$2\srcAWId-L_c-L_c^*=(\srcAWId-L_c)^*(\srcAWId-L_c)\succeq0$.
\end{proof}

These are exactly Round 3 Entry 05's Propositions 3.1, 3.3 and 4.1
\cite{R4E5:R305}. Their new quantitative use is the following.

\begin{srcAWproposition}[Explicit local recipe]\label{R4E5:prop:weights}
Suppose $A\succ0$ has one permitted support block with $q\ge1$ edges.
Assign weights as follows:
\begin{equation}\label{R4E5:eq:weights}
 c_{ij}=\begin{cases}
 h(A)|a_{ij}|^2/(a_{ii}a_{jj}),&\text{for a phase-balanced block},\\[2pt]
 \beta_{ij}(A)/q,&\text{for a cycle or a block of type 3}.
 \end{cases}
\end{equation}
Then every weight is positive and $\srcAWDef(A)\succeq K_c$.
More generally, in the second case $1/q$ may be replaced by any
positive edge weights $\theta_{ij}$ with sum at most $1$.
\end{srcAWproposition}
\begin{proof}
For the second case combine Lemmas~\ref{R4E5:lem:pair} and \ref{R4E5:lem:local}
and use diagonal normalization. In the odd case the reflected
$A'=EC'E$ has unchanged two-by-two principal determinants and
$\det(A')\ge\det(A)$, so $\beta_{ij}(A')\ge\beta_{ij}(A)$.
The convex weights in \eqref{R4E5:eq:convex} therefore suffice.
For a balanced matrix gauge to nonnegative entries. In
$\srcAWDef(A)=\sum_\pi d_A(\pi)(\srcAWId-L_\pi)$, retain every single
transposition, whose coefficient is the first line of
\eqref{R4E5:eq:weights}. Every other involution has a nonnegative
coefficient; each inverse pair contributes
$d_A(\pi)(2\srcAWId-L_\pi-L_\pi^*)\succeq0$.
The identity term vanishes. This proves simultaneous retention of
all the transpositions in the balanced case.
\end{proof}

\section{Positive-definite cut-vertex propagation}\label{R4E5:sec:glue}
For a matrix $F$ on $J\subseteq[n]$, let
\[
 \widehat\srcAWSch(F)=\sum_{\pi\in\srcAWG_J}d_F(\pi)L_\pi,
 \qquad \widehat\srcAWDef(F)=P(F)\srcAWId-\widehat\srcAWSch(F),
\]
where each permutation fixes the complement of $J$.
On every right coset $\srcAWG_J\tau$, $\widehat\srcAWSch(F)$ is a copy of
$\srcAWSch(F)$. Operator inequalities therefore embed unchanged.
Operators on disjoint label sets commute; no commutation of
operators on overlapping sets is used below.

Let $A=A_1\star A_2$ be the one-point sum at $v$: add the two
$v$-diagonal entries and keep all other entries on the two sides,
with zero cross entries. Put $B_i=A_i(v)$, $P_i=P(A_i)$ and $R_i=P(B_i)$.
The operator identity reused from \cite[Lemma 5.1]{R4E5:R305} is
\begin{equation}\label{R4E5:eq:glue}
 \srcAWSch(A)=\widehat\srcAWSch(A_1)\widehat\srcAWSch(B_2)
 +\widehat\srcAWSch(B_1)\widehat\srcAWSch(A_2),\qquad
 P(A)=P_1R_2+R_1P_2.
\end{equation}
Indeed a supported permutation cycle moving $v$ stays on one side;
a cycle visiting both sides would repeat $v$. When $v$ is fixed,
the two contributions add its diagonal entries. All other cycles
are supported on disjoint sets, which proves the identity term by term.

\begin{srcAWlemma}[Strict splitting and propagation]\label{R4E5:lem:glue}
If $A\succ0$ has a cut vertex $v$, it can be written
$A=A_1\star A_2$ with both pieces positive definite.
If $A_1,A_2,B_1,B_2$ have POT, then
\begin{equation}\label{R4E5:eq:propagate}
 \srcAWDef(A)\succeq
 \det(B_2)\widehat\srcAWDef(A_1)+
 \det(B_1)\widehat\srcAWDef(A_2).
\end{equation}
Thus a positive edge weight from piece $A_i$ is multiplied by
$\det(B_{3-i})$ when propagated to $A$.
\end{srcAWlemma}
\begin{proof}
Let $U,W$ be nonempty unions of the components after deleting $v$.
Write $H_U=A[U]$, $H_W=A[W]$, $b_U=A[U,v]$, $b_W=A[W,v]$.
Since cross entries vanish, the scalar Schur complement is
\[
 \delta=a_{vv}-b_U^*H_U^{-1}b_U-b_W^*H_W^{-1}b_W>0.
\]
Give the first piece diagonal $b_U^*H_U^{-1}b_U+\delta/2$ at $v$,
and the second $b_W^*H_W^{-1}b_W+\delta/2$.
Both Schur complements are $\delta/2>0$; all other entries are
unchanged. This is the positive-definite refinement of the earlier
Gram splitting, where the orthogonal residual at $v$ is split equally.
It is rational whenever $A$ is Gaussian rational.

Set $X_i=\widehat\srcAWSch(A_i)$ and $Y_i=\widehat\srcAWSch(B_i)$.
From \eqref{R4E5:eq:glue},
\[
 \srcAWDef(A)=\widehat\srcAWDef(A_1)Y_2+P_1(R_2\srcAWId-Y_2)
          +\widehat\srcAWDef(A_2)Y_1+P_2(R_1\srcAWId-Y_1).
\]
All terms are PSD: each product uses disjoint label sets, and the
assumed POT bounds apply. The Schur-minimum inequality from
Section~\ref{R4E5:sec:minimum} gives $Y_i\succeq\det(B_i)\srcAWId$.
The commuting products therefore imply \eqref{R4E5:eq:propagate}.
\end{proof}

For completeness, $\srcAWClass$ is hereditary under principal submatrices:
balanced blocks remain balanced, a proper induced subgraph of a cycle
is a forest, and type 3's two graph conditions are inherited by
subgraphs. Each new block lies in an old one. Changing a diagonal at
$v$ does not change a block's off-diagonal phases or type.
Induct on order, first splitting at cut vertices as above. Local POT
and weights follow from Proposition~\ref{R4E5:prop:weights}; the vertex-deleted
submatrices have POT by the same induction. For a disconnected matrix,
its Schur operator is the product of commuting embedded operators of
its components, bounded by the product of their permanents, which is
its permanent. Singletons start the induction.
Thus \eqref{R4E5:eq:propagate} constructs positive weights on every edge
of a connected $A\succ0$ in $\srcAWClass$, proving the first comparison
in Theorem~\ref{R4E5:thm:main}. Any cut-vertex splitting order is valid;
the resulting lower bound need not be optimal or unique.

\section{From edge weights to a spectral gap}\label{R4E5:sec:gap}
For a connected positive-weight support graph, $\ker K_c$ consists
of constants: a vector in the kernel is invariant under every edge
transposition, and these generate $\srcAWG_n$. The weighted interchange-gap
theorem of Caputo--Liggett--Richthammer \cite[Theorem 1.1]{R4E5:CLR} states
that the smallest nonzero eigenvalue of $K_c$ equals
$\lambda_2(\mathcal L_c)$. Their generator uses right transpositions;
inversion of permutations unitarily interchanges the left and right
conventions. Their rates $c_{ij}$ have exactly our normalization.
This proves the second comparison in \eqref{R4E5:eq:comparison}.

Here is the promised elementary substitute, which does not use that
theorem. On a spanning tree of diameter $d$, an arbitrary transposition
is a word of length at most $2d-1$ in edge transpositions, by following
the path and returning along it with the final edge omitted. Each tree
edge occurs at most twice. A permutation is a product of at most
$n-1$ arbitrary transpositions, obtained from its disjoint cycles.
It therefore has a word of length at most
$k=(n-1)(2d-1)$ in which any edge occurs at most $2(n-1)$ times.
For any such word $s_1\cdots s_r$, the telescoping identity and
Cauchy--Schwarz give
\[
 \|(\srcAWId-L_{s_1\cdots s_r})x\|^2
 \le r\sum_{j=1}^r\|(\srcAWId-L_{s_j})x\|^2
 \le 2(n-1)^2(2d-1)\sum_{e\in T}\|(\srcAWId-L_e)x\|^2.
\]
Here the terms in the telescoping sum are unitary prefixes applied
to $(\srcAWId-L_{s_j})x$, so their norms have the stated form.
For $x\perp\mathbf1$ average over all permutations:
\[
 \|x\|^2=\frac1{2n!}\sum_{\pi\in\srcAWG_n}\|x-L_\pi x\|^2
 \le 2(n-1)^2(2d-1)
       \sum_{e\in T}\langle x,(\srcAWId-L_e)x\rangle.
\]
Since $K_c\succeq c_*\sum_{e\in T}(\srcAWId-L_e)$, this proves
\eqref{R4E5:eq:elementarygap}.

\section{Scalar stability and dense operator perturbations}
\begin{proof}[Scalar conclusions of Theorem~\ref{R4E5:thm:main}]
The Schur product theorem and tensor Gram construction give
$\srcAWSch(A)\circ\srcAWSch(B)=\srcAWSch(A\circ B)$ and
$\srcAWId\circ\srcAWSch(B)=h(B)\srcAWId$.
Furthermore $L_\pi\circ\srcAWSch(B)=d_B(\pi)L_\pi$ by
\eqref{R4E5:eq:regular}. Taking the Schur product of the first comparison
in \eqref{R4E5:eq:comparison} with $\srcAWSch(B)\succeq0$, and then its
quadratic form on $\mathbf1/\sqrt{n!}$, proves \eqref{R4E5:eq:scalar}.
Every right-hand summand is nonnegative by the two-by-two principal
minor bound for $B$, including zero diagonal cases without division.
For a correlation matrix, equality forces $|b_{ij}|=1$ on every
positive-weight edge. Its unit Gram vectors are then pairwise
collinear along the connected support, so its rank is one.
Conversely rank one makes $A\circ B$ a diagonal unitary congruence
of $A$ and leaves its permanent unchanged.
Finally $Q\circ\srcAWSch(B)=h(B)\srcAWId-\srcAWSch(B)/n!$, giving
\eqref{R4E5:eq:scalarGap}. These bounds imply the restricted Bapat--Sunder
inequality. The further classical inequality $h(B)\le P(B)$ yields
Chollet on this first-factor domain, but is not a new assertion here.
\end{proof}

\begin{proof}[Theorem~\ref{R4E5:thm:stable}]
By \eqref{R4E5:eq:regular} and the unitarity of each $L_\pi$,
\[
 \|\srcAWSch(A)-\srcAWSch(A_0)\|\le\ell(A,A_0),\qquad
 |P(A)-P(A_0)|\le\ell(A,A_0).
\]
The Hermitian difference $\srcAWDef(A)-\srcAWDef(A_0)$ annihilates
$\mathbf1$, so it is bounded below by $-2\ell(A,A_0)Q$,
not merely by $-2\ell(A,A_0)\srcAWId$.
This proves \eqref{R4E5:eq:stable}. If both matrices are correlation,
each entry has modulus at most one. Telescoping each product bounds
its difference by $n\varepsilon$, proving \eqref{R4E5:eq:roughradius}.
In general, expansion after writing $a_{ij}=a^0_{ij}+e_{ij}$ gives
\[
 |d_A(\pi)-d_{A_0}(\pi)|
 \le\prod_i(|a^0_{\pi(i),i}|+\varepsilon)
       -\prod_i|a^0_{\pi(i),i}|.
\]
Summing gives \eqref{R4E5:eq:polyradius}. Finally a Hermitian matrix with
entrywise modulus bounded by $\varepsilon$ has operator norm at most
$n\varepsilon$, proving the stated positive-definiteness radius.
\end{proof}

\section{Examples and a singular boundary obstruction}
\subsection{An unbounded-order phase-frustrated family}
Let $m\ge4$ be divisible by four and choose
$\zeta_1,\ldots,\zeta_m$ by repeating $1,-1,i,-i$ equally often.
Let
\[
 A_m=\begin{pmatrix}\srcAWId_2&B_m\\B_m^*&\srcAWId_m\end{pmatrix},\qquad
 (B_m)_{1j}=\frac1{2\sqrt m},\quad
 (B_m)_{2j}=\frac{\zeta_j}{2\sqrt m}.
\]
The two rows of $B_m$ are orthogonal and have squared norm $1/4$.
Thus $A_m$ has eigenvalues $1/2,3/2$, each twice, and $1$ with
multiplicity $m-2$. In particular $\det(A_m)=9/16$ and
$\lambda_{\min}(A_m)=1/2$. Its support is $K_{2,m}$; nonreal
four-cycle products show that it is not phase-balanced.
All $2m$ edge moduli squared equal $1/(4m)$. Recipe
\eqref{R4E5:eq:weights} yields
\begin{equation}\label{R4E5:eq:familygap}
 \beta_e=\frac9{16(4m-1)},\quad
 c_e=\frac9{32m(4m-1)},\quad
 g_m=\frac9{16m(4m-1)},
\end{equation}
because the ordinary Laplacian of $K_{2,m}$ has nonzero eigenvalues
$2$ (multiplicity $m-1$), $m$ (multiplicity one), and $m+2$
(multiplicity one). These eigenvalues follow by separating leaf
vectors with zero sum, the difference of the two hub coordinates,
and vectors constant on each part.
The permanent is
\[
 P(A_m)=\frac{25}{16}-\frac1{8m}.
\]
To see this, the $K_{2,m}$ expansion of Round 3 Entry 03
\cite[Section 2]{R4E5:R303} is
$P=1+X+Y+XY-2Z+|r|^2$, where
$X=\sum|u_j|^2$, $Y=\sum|v_j|^2$,
$Z=\sum|u_jv_j|^2$, and $r=\sum u_j\overline v_j$.
Here $X=Y=1/4$, $Z=1/(16m)$, and $r=0$.
It also follows directly by enumerating matchings and four-cycles.
The certified absolute and relative gaps both decay as order $m^{-2}$
in this family, since its permanent stays bounded away from zero and
infinity. The crude entrywise radius has an additional $n\,n!$ loss.

\subsection{A rational dense neighborhood}
For $m=4$, set
\[
 A_0=\begin{pmatrix}\srcAWId_2&B\\B^*&\srcAWId_4\end{pmatrix},\qquad
 B=\tfrac14\begin{pmatrix}1&1&1&1\\1&-1&i&-i\end{pmatrix}.
\]
Then $n=6$, $P(A_0)=49/32$, $c_e=3/640$, and $g=3/320$.
For every Hermitian $A$ with
\begin{equation}\label{R4E5:eq:concreteradius}
 \max_{ij}|a_{ij}-a^0_{ij}|\le\frac1{5000},
\end{equation}
positive definiteness follows from $6/5000<1/2$.
The nonnegative polynomial in \eqref{R4E5:eq:polyradius}, computed by
finite permutation expansion, satisfies
\[
 P(|A_0|+tJ_6)-P(|A_0|)
 =\frac{257}{16}t+\frac{769}{8}t^2+366t^3
   +876t^4+1200t^5+720t^6.
\]
At $t=1/5000$ its value is less than $322/100000$, so twice that
value is less than $3/320$. Thus the \emph{entire} neighborhood
\eqref{R4E5:eq:concreteradius} has POT with a strictly positive gap.
The exact verifier checks the displayed coefficients and inequalities.
For example $A=A_0+(J_6-\srcAWId_6)/5000$ is fully dense.
Its triangle on the two hubs and the first leaf has positive cycle
product, whereas the triangle through the leaf with phase $i$ has
nonreal product; the latter alone excludes phase balance.
This dense conclusion is not obtained by merely preserving the
original $K_{2,4}$ zero pattern.

\subsection{Why positive definiteness was retained}
Consider the Gaussian-integer matrix
\[
 F=\begin{pmatrix}2\srcAWId_2&H\\H^*&2\srcAWId_2\end{pmatrix},\qquad
 H=(1+i)\begin{pmatrix}1&1\\1&-1\end{pmatrix}.
\]
Since $HH^*=4\srcAWId_2$, $F\succeq0$ has eigenvalues $4,4,0,0$.
It has connected $K_{2,2}$ support. Exact characteristic-polynomial
computation gives
\begin{equation}\label{R4E5:eq:singularspec}
 \det(t\srcAWId_{24}-\srcAWSch(F))=t^{12}(t-16)^3(t-32)^6(t-48)^3,
 \qquad P(F)=48.
\end{equation}
Thus its top eigenvalue has multiplicity three, so no inequality
$\srcAWDef(F)\succeq gQ$ with $g>0$ holds. This does not show that every
POT neighborhood of this singular matrix fails; it only rules out
extending our simple-top conclusion without additional hypotheses.
The small exact identity \eqref{R4E5:eq:singularspec} is independently
reproducible in the supplied verifier, not a floating-point assertion.

\section{Verification, status, and remaining limits}
The exact verifier uses Gaussian rational arithmetic and rational
Schur pivots. It checks the two-row bound on small full Schur powers,
convex weighting, local and gluing comparison certificates, spectral
gaps, polynomial radii, and unrestricted second-factor scalar tests.
The singular characteristic polynomial is checked exactly.
Dynamic programming and Ryser's formula cross-check permanents.
All checks are finite diagnostics; the general proofs above do not
infer operator positivity from numerical eigenvalues.
A separate bounded numerical exploration of general bipartite supports
found no counterexample, but also furnished no proof. It is not used.

We have not proved POT for all bipartite matrices or for every positive
definite PSD matrix. Even in $\srcAWClass$, our weights and radii are lower
bounds and need not be sharp; the determinant factor can degenerate
rapidly near singularity or with increasing order.
The main theorem is a quantitative extension of a restricted structural
result. The ordinary POT counterexamples remain counterexamples to
its universal version, and general Chollet remains unresolved here.
The latest Pant--Singh source checked is v2, dated 18 September 2026;
Tran's counterexample source is v5, dated 23 August 2022; the
Caputo--Liggett--Richthammer preprint's latest version is v4,
28 September 2009, with its published theorem used explicitly.
A posted 2026 claim of Chollet through order six \cite{R3E5:LZL} is recorded
as an unaudited, non-peer-reviewed claim and is not a dependency.
Even its asserted scalar conclusion would not subsume our unbounded-order
operator comparison. No historical priority or prize-level claim is made.

\begingroup\small

\endgroup

\chapter[A universal weak-correlation region with permanent on top]{A universal weak-correlation region with permanent on top}
\label{source:R5E4}
\noindent\textit{Source development: Round 5, entry 04. Original mathematical arguments retained; notation is local to this chapter.}\par\medskip

\begin{quote}\small\noindent\textbf{Chapter abstract.} 
For every $n\ge2$, every diagonal-one Hermitian matrix whose off-diagonal entries have modulus at most $1/[12(n-1)]$ has permanent on top, with no support or phase restrictions. More precisely, its Schur-power permanent defect lies between $1/6$ and $2$ times the weighted interchange Laplacian with edge weights $|a_{ij}|^2$. This gives the entire top eigenspace, quantitative inequalities against arbitrary PSD Hadamard factors, and spectral-gap estimates from an order-$n$ graph Laplacian. A more flexible criterion uses the maximum absolute off-diagonal row sum and the squared Frobenius mass. The proof controls complex long-cycle terms by transposition energy after cutting each cycle at a minimum-weight edge. Unlike preceding structural and perturbative results, the theorem applies directly to arbitrary dense supports, including at the completely degenerate identity center. Constants are sufficient, not optimal; unrestricted permanent-on-top is false, and general Chollet remains unresolved.
\end{quote}

\section{Main results and their scope}
For a Hermitian matrix $A$ of order $n$, set
\[
 P(A)=\srcBBper(A),\qquad \srcBBSch(A)_{\sigma,\tau}=\prod_{i=1}^n a_{\sigma(i),\tau(i)},\qquad
 \srcBBDef(A)=P(A)I-\srcBBSch(A),\quad \sigma,\tau\in S_n.
\]
Let $L_\pi e_\tau=e_{\pi\tau}$ be the left regular representation. For diagonal-one $A$ write
\[
 w_{ij}=|a_{ij}|\ (i\ne j),\quad w_{ii}=0,\quad
 H_A=\sum_{i<j}w_{ij}^2(I-L_{(ij)}),\quad
 \mathcal L_A=\sum_{i<j}w_{ij}^2(e_i-e_j)(e_i-e_j)^*.
\]
Each summand of $H_A$ is PSD. The support graph $G_A$ has edges $ij$ with $w_{ij}>0$.

\begin{srcBBtheorem}[A support-free, phase-free neighborhood]\label{R5E4:thm:main}
Let $n\ge2$. Every diagonal-one Hermitian matrix $A$ satisfying
\begin{equation}\label{R5E4:eq:ball}
 |a_{ij}|\le\frac1{12(n-1)}\qquad(i\ne j)
\end{equation}
is positive definite and obeys
\begin{equation}\label{R5E4:eq:sandwich}
 \frac16 H_A\preceq\srcBBDef(A)\preceq2H_A.
\end{equation}
In particular $\srcBBSch(A)\succeq0$ and $P(A)$ is its largest eigenvalue. Moreover
\begin{equation}\label{R5E4:eq:perbound}
 1\le P(A)\le\frac{4752}{4717}.
\end{equation}
If the connected components of $G_A$ have sizes $n_1,\ldots,n_r$, the multiplicity of this top eigenvalue is exactly $n!/\prod_jn_j!$. Its eigenspace consists of functions constant on each orbit of the left action of $S_{n_1}\times\cdots\times S_{n_r}$.
\end{srcBBtheorem}
Thus the whole entrywise ball around the identity is certified, not only its intersection with a presumed PSD domain. It contains arbitrary dense complex phases and all zero patterns. No lower bound on a nonzero entry is required. At the identity the defect is zero; the argument does not rely on perturbing a previously positive spectral gap.

\begin{srcBBcorollary}[Unrestricted masks and equality]\label{R5E4:cor:mask}
Under \eqref{R5E4:eq:ball}, every correlation matrix $B$ satisfies
\begin{equation}\label{R5E4:eq:mask}
 P(A)-P(A\circ B)\ge\frac16\sum_{i<j}|a_{ij}|^2(1-|b_{ij}|^2).
\end{equation}
Equality $P(A\circ B)=P(A)$ holds precisely when $B$ restricted to each connected component of $G_A$ has rank one. In particular it forces rank one when $G_A$ is connected. For arbitrary PSD $B$,
\[
 P(A\circ B)\le P(A)\prod_i b_{ii}\le P(A)P(B).
\]
\end{srcBBcorollary}

\paragraph{The precise increment.}
Round 3 entry 05 proved qualitative POT for cycle and specified overlapping bipartite blocks, with one-point closure \cite{R5E4:R305}. Round 4 entry 03 extended sign-reflection certificates to PSD companions, notably certain book matrices \cite{R5E4:R403}. Our Round 4 entry 04 proved polynomial relative gaps and radii around cycle centers with edge magnitudes in $[1/4,1/3]$ \cite{R5E4:R404}. Round 4 entry 05 proved a transposition comparison and dense stability for connected positive definite matrices in the inherited structural class \cite{R5E4:R405}. These contributions are credited, not repeated here.

The new step is a direct two-sided comparison for \emph{arbitrary supports and phases} in a universal weak-correlation region. The proof below uses neither a sparse center nor a reflected PSD companion, and it remains valid when edges vanish and the top eigenspace changes dimension. It is not a uniform enlargement of the earlier cycle neighborhoods: their large cycle entries lie outside \eqref{R5E4:eq:ball}. The familiar scalar consequences of an operator comparison, and the classical interchange-gap theorem used later, are not claimed as new machinery. All six frozen Round 4 papers and the assessment were read.

\paragraph{Current target status.}
Universal POT and universal Bapat--Sunder are false \cite{R4E3:Tran,R2E5:drury}; we do not label either as an open conjecture. General Chollet is stated open in the latest checked Pant--Singh version, v2 of 18 September 2026 \cite{R2E3:PS}. The motivating problem is to establish useful unrestricted-order positive classes for these permanent inequalities. The public through-order-six Chollet claim \cite{R1E5:li} is not assumed or independently validated here and does not give our stronger operator conclusion. Searches of current primary sources on 20 October 2026 did not locate the precise comparison above; this does not certify historical priority. Wanless's survey \cite{R5E4:Wanless} and the recent restricted-class discussion \cite{R3E4:Pioge} distinguish the relevant conjectures.

\section{A quadratic-form estimate for transposition words}\label{R5E4:sec:word}
The regular-representation identity is
\begin{equation}\label{R5E4:eq:regular}
 \srcBBSch(A)=\sum_{\pi\in S_n}d_A(\pi)L_\pi,
 \qquad d_A(\pi)=\prod_i a_{\pi(i),i}.
\end{equation}
It follows by evaluating the $(\sigma,\tau)$ entry at $\pi=\sigma\tau^{-1}$. Row sums are $P(A)$, so $\srcBBDef(A)\mathbf1=0$. Hermiticity gives $d_A(\pi^{-1})=\overline{d_A(\pi)}$.

\begin{srcBBlemma}[Word energy]\label{R5E4:lem:word}
Let $U_1,\ldots,U_s$ be self-adjoint unitary operators and put $E_j=\langle v,(I-U_j)v\rangle\ge0$. Then
\begin{equation}\label{R5E4:eq:word}
 \left|\langle v,(I-U_1\cdots U_s)v\rangle\right|
 \le\left(\sum_{j=1}^s\sqrt{E_j}\right)^2
 \le s\sum_{j=1}^s E_j.
\end{equation}
The assertion includes repeated operators and arbitrary complex $v$.
\end{srcBBlemma}
\begin{proof}
Let $V_{j-1}=U_1\cdots U_{j-1}$. Telescoping gives
$I-U_1\cdots U_s=\sum_jV_{j-1}(I-U_j)$.
Also $\|(I-U_j)v\|=\sqrt{2E_j}$. Telescoping the reversed word for $V_{j-1}^*$ yields
$\|(V_{j-1}^*-I)v\|\le\sum_{i<j}\sqrt{2E_i}$.
Consequently
\[
 |\langle v,V_{j-1}(I-U_j)v\rangle|
 \le E_j+2\sum_{i<j}\sqrt{E_iE_j}.
\]
Sum over $j$ and use Cauchy--Schwarz. No commutation is used. In particular the absolute value controls the imaginary as well as the real part of a long-cycle contribution.
\end{proof}

For each supported directed simple cycle $\gamma$ of length $l\ge3$, choose a minimum-weight edge, with a fixed tie-breaking rule, and delete it. The remaining path gives a word for $\gamma$ of length $l-1$ in its consecutive edge transpositions. Indeed, after rotating its labels so that the deleted edge closes the cycle,
\[
 (v_0\ v_1\ \cdots\ v_{l-1})
 =(v_0v_1)(v_1v_2)\cdots(v_{l-2}v_{l-1}).
\]
A two-cycle uses its sole transposition. Concatenating these words over the disjoint nontrivial cycles of a permutation gives a word with no repeated edge. Its length is
\[
 \ell(\pi)=\sum_{\gamma\subset\pi}(|\gamma|-1).
\]
The word construction, not an arbitrary shortest word, will be used throughout.

\section{Cycle masses and marked-edge estimates}\label{R5E4:sec:cycles}
Let $W=(w_{ij})$ be symmetric, nonnegative, and zero on its diagonal. Set
\[
 q=\max_i\sum_jw_{ij}<1,\qquad X=\sum_{i<j}w_{ij}^2.
\]
A cycle means a directed simple cycle modulo rotation; a two-cycle is counted once, and a cycle of length at least three and its reverse are distinct. Its weight and length cost are
\[
 a_\gamma=\prod_{ij\in\gamma}w_{ij},\qquad c_\gamma=|\gamma|-1.
\]
For two-cycles $a_\gamma=w_{ij}^2$. Put $T_l=\sum_{|\gamma|=l}a_\gamma$, and let
\[
 T_* =\sum_{l\ge2}T_l,\quad U_* =\sum_{l\ge2}(l-1)T_l,
 \quad T_{\rm long}=\sum_{l\ge3}T_l,\quad
 U_{\rm long}=\sum_{l\ge3}(l-1)T_l.
\]
All sums are finite, with $l\le n$.

\begin{srcBBlemma}[Scalar cycle majorants]\label{R5E4:lem:masses}
Define
\begin{equation}\label{R5E4:eq:parameters}
 t=X\left(1+\frac{2q}{3(1-q)}\right),\quad
 u=X\frac{1+q}{1-q},\quad
 b=\frac{2Xq}{3(1-q)},\quad v=\frac{2Xq}{1-q}.
\end{equation}
Then $T_*\le t$, $U_*\le u$, $T_{\rm long}\le b$, and $U_{\rm long}\le v$.
\end{srcBBlemma}
\begin{proof}
We have $T_2=X$. For $l\ge3$, closed-walk expansion and symmetry give
\[
 T_l\le\frac1l\operatorname{tr}(W^l)
 \le\frac1l\sum_j|\lambda_j(W)|^l
 \le\frac{q^{l-2}}l\operatorname{tr}(W^2)
 =\frac{2X}{l}q^{l-2}.
\]
Here $\|W\|\le q$, for example by the maximum absolute row and column sums. Bound $1/l\le1/3$ and $(l-1)/l\le1$, and sum the geometric series. Nonnegative closed-walk weights justify the first inequality even for odd $l$.
\end{proof}

A permutation is a collection of vertex-disjoint cycles with the other vertices fixed. In upper bounds we drop disjointness. For all subsets $\mathcal F$ of the finite cycle list, put $a(\mathcal F)=\prod_{\gamma\in\mathcal F}a_\gamma$, $C(\mathcal F)=\sum c_\gamma$, and $N(\mathcal F)$ equal to the number of long cycles in $\mathcal F$.
\begin{srcBBlemma}[Decorations]\label{R5E4:lem:decor}
The following upper bounds hold, also when the summation is restricted to vertex-disjoint collections or to any principal vertex set:
\begin{align}
 \sum_{\mathcal F}a(\mathcal F)&\le e^t,&
 \sum_{\mathcal F}C(\mathcal F)a(\mathcal F)&\le ue^t,\label{R5E4:eq:decor1}\\
 \sum_{N(\mathcal F)\ge1}(1+C(\mathcal F))a(\mathcal F)
 &\le e^t\{b(1+u)+v\}.\label{R5E4:eq:decor2}
\end{align}
\end{srcBBlemma}
\begin{proof}
The unrestricted partition function is $F=\prod_\gamma(1+a_\gamma)\le e^{T_*}$. Marking one cycle gives $\sum C a\le FU_*$. To bound the last sum replace its indicator by $N$. Marking one long cycle gives $\sum Na\le FT_{\rm long}$. Marking a long cycle and a length-cost cycle gives
\[
 \sum NCa\le F\{T_{\rm long}U_*+U_{\rm long}\}.
\]
Indeed the distinct marked cycles contribute at most the product of their weights times $F$, while coincident marks contribute at most $c_\gamma a_\gamma F$. These estimates prove the claims; every summand is nonnegative, so restrictions can only decrease them.
\end{proof}

\begin{srcBBlemma}[Cut at a lightest edge]\label{R5E4:lem:cut}
Fix an edge $e$ with $w_e>0$. For each $l\ge3$, sum over the length-$l$ directed cycles whose chosen word contains $e$. Then
\begin{equation}\label{R5E4:eq:cut}
 \sum_{\substack{|\gamma|=l\\ e\text{ in its word}}}
 \frac{a_\gamma}{w_e^2}\le2(l-1)q^{l-2}.
\end{equation}
\end{srcBBlemma}
\begin{proof}
The deleted minimum-weight edge $f$ is distinct from $e$, and $w_f\le w_e$. Hence
\[
 \frac{a_\gamma}{w_e^2}
 =\frac{w_f}{w_e}\prod_{h\in\gamma\setminus\{e,f\}}w_h
 \le\prod_{h\in\gamma\setminus\{e,f\}}w_h.
\]
Fix one of the two orientations of $e=ij$. Deleting $e,f$ leaves two paths starting at $i,j$, with lengths $k,l-2-k$ for some $0\le k\le l-2$, allowing a zero-length path. Their endpoints determine the missing edge $f$, so the encoding is injective once this orientation and $k$ are fixed. Dropping simplicity, disjointness, and the requirement that the endpoints are adjacent only enlarges the sum. The total weight of length-$k$ walks from a fixed vertex is at most $q^k$, by induction on $k$. Thus each $k$ contributes at most $q^{l-2}$. There are $l-1$ choices of $k$ and two orientations of $e$. This proves \eqref{R5E4:eq:cut}. This step is why arbitrarily small nonzero entries do not create small denominators in the final estimate.
\end{proof}

\section{The general two-sided operator comparison}\label{R5E4:sec:general}
In addition to \eqref{R5E4:eq:parameters}, define
\begin{align}
 s_1(q)&=\frac1{(1-q)^2}-1,
 &s_2(q)&=\frac{1+q}{(1-q)^3}-1,\label{R5E4:eq:sums}\\
 B(q,X)&=b(1+u)+v+2\{s_2(q)+u s_1(q)\},
 &\kappa(q,X)&=\frac{B(q,X)}{1-t}.\label{R5E4:eq:kappa}
\end{align}
The identity $s_j(q)=\sum_{l\ge3}(l-1)^j q^{l-2}$ for $j=1,2$ follows by differentiating a geometric series.

\begin{srcBBtheorem}[Row-sum and squared-mass criterion]\label{R5E4:thm:general}
Let $A$ be diagonal-one Hermitian. Form $q,X$ from its entrywise off-diagonal moduli as above. Suppose $q<1$, $t<1$, and $\kappa(q,X)<1$. Then $A\succ0$ and
\begin{equation}\label{R5E4:eq:general}
 (1-\kappa)H_A\preceq\srcBBDef(A)
 \preceq\left(\frac{1+X}{1-X}+\kappa\right)H_A.
\end{equation}
In particular the complete support/eigenspace and scalar equality conclusions of Section 1 hold, with lower coefficient $1-\kappa$. A convenient additional sufficient regime is
\begin{equation}\label{R5E4:eq:secondregime}
 q\le\frac1{32},\quad X\le\frac12
 \quad\Longrightarrow\quad
 \frac15H_A\preceq\srcBBDef(A)\preceq4H_A.
\end{equation}
\end{srcBBtheorem}
\begin{proof}
Strict diagonal dominance proves $A\succ0$: for every vector $z$,
$z^*Az\ge\sum_i(1-\sum_{j\ne i}w_{ij})|z_i|^2\ge(1-q)\|z\|^2$.
Separate the defect in \eqref{R5E4:eq:regular} as $D_{\rm mat}+D_{\rm long}$, according as the permutation is a matching of transpositions or contains a long cycle. For a matching $M$, $d_A(\pi_M)=\prod_{e\in M}w_e^2$ and $L_{\pi_M}$ is a self-adjoint involution. Its defect is PSD; retaining singleton matchings gives
\begin{equation}\label{R5E4:eq:matlower}
 D_{\rm mat}\succeq H_A.
\end{equation}
Write $E_e=\langle z,(I-L_e)z\rangle$. By Lemma~\ref{R5E4:lem:word},
\[
 \langle z,D_{\rm mat}z\rangle
 \le\sum_M |M|\left(\prod_{e\in M}w_e^2\right)\sum_{e\in M}E_e.
\]
For a fixed $e$, divide its coefficient by $w_e^2$, remove $e$, and allow arbitrary subsets of all remaining edges. Their partition function and first cardinality moment are at most $e^X$ and $Xe^X$. Therefore
\begin{equation}\label{R5E4:eq:matupper}
 D_{\rm mat}\preceq e^X(1+X)H_A
 \preceq\frac{1+X}{1-X}H_A.
\end{equation}
Here $X\le t<1$ and $e^X\le1/(1-X)$, by the power series.

For $D_{\rm long}$ use the chosen words of Section~\ref{R5E4:sec:word}. Its quadratic form is real, since inverse permutations occur together. The triangle inequality and Lemma~\ref{R5E4:lem:word} give
\begin{equation}\label{R5E4:eq:badenergy}
 |\langle z,D_{\rm long}z\rangle|
 \le\sum_e E_e\sum_{\substack{\pi\text{ has a long cycle}\\e\text{ in its word}}}
 \ell(\pi)|d_A(\pi)|.
\end{equation}
We bound the inner coefficient divided by $w_e^2$ in two cases.

If $e$ itself forms a two-cycle, the other cycles include a long cycle and have total length cost $C$, while $\ell(\pi)=1+C$. Lemma~\ref{R5E4:lem:decor} bounds this contribution by $e^t\{b(1+u)+v\}$.

Otherwise $e$ lies in the word of a unique long cycle $\gamma$ of length $l$. After deleting that cycle, the remaining cycles have cost $C$ and $\ell(\pi)=l-1+C$. Equations~\eqref{R5E4:eq:decor1} and \eqref{R5E4:eq:cut} bound this contribution, summed over all $l$, by
\[
 e^t\sum_{l\ge3}2(l-1)q^{l-2}(l-1+u)
 =2e^t(s_2+u s_1).
\]
These cases are disjoint and exhaustive, since the cycle words use each edge at most once. Consequently \eqref{R5E4:eq:badenergy} is at most $e^t B(q,X)\langle z,H_Az\rangle$. Since $e^t\le1/(1-t)$,
\begin{equation}\label{R5E4:eq:longbound}
 -\kappa H_A\preceq D_{\rm long}\preceq\kappa H_A.
\end{equation}
Combine \eqref{R5E4:eq:matlower}, \eqref{R5E4:eq:matupper}, and \eqref{R5E4:eq:longbound} to prove \eqref{R5E4:eq:general}. All estimates are monotone in $q,X$ on the indicated nonnegative domain. At $q=1/32,X=1/2$ the rational bound is
\[
 t=95/186,\qquad \kappa=69197/87451<4/5.
\]
The matching upper coefficient is $3$, so \eqref{R5E4:eq:secondregime} follows. This calculation uses a rational exponential majorant, not a numerical exponential evaluation.
\end{proof}

\begin{proof}[Proof of Theorem~\ref{R5E4:thm:main}]
The entry bound gives $q\le1/12$ and
\[
 X\le\frac{n}{288(n-1)}\le\frac1{144}.
\]
At these upper endpoints, direct rational simplification gives
\[
 t=\frac{35}{4752},\qquad
 \kappa=\frac{33853535}{41094504}<\frac56,
 \qquad \frac{145}{143}+\kappa<2.
\]
Monotonicity and Theorem~\ref{R5E4:thm:general} prove \eqref{R5E4:eq:sandwich}. A tensor Gram representation proves $\srcBBSch(A)\succeq0$, since $A\succ0$. Its diagonal entries are one, so the PSD defect gives $P(A)\ge1$. On the other hand, the permutation cycle expansion and Lemma~\ref{R5E4:lem:decor} give
$P(A)\le\srcBBper(I+W)\le e^t\le1/(1-t)\le4752/4717$.

The kernel of $H_A$ consists exactly of vectors fixed by every support-edge transposition, because its positive summands have that common kernel. Edge transpositions generate the full symmetric group of each connected component: along a path they generate a swap of its endpoints, and arbitrary swaps generate its symmetric group. Thus this kernel has the stated orbit description and dimension. Every supported permutation belongs to that component-preserving subgroup, so \eqref{R5E4:eq:regular} shows that $\srcBBDef(A)$ kills this entire kernel. Conversely \eqref{R5E4:eq:sandwich} implies $\ker\srcBBDef(A)\subseteq\ker H_A$. This proves the exact eigenspace claim, including isolated vertices and $A=I$.
\end{proof}

\begin{proof}[Proof of Corollary~\ref{R5E4:cor:mask}]
For PSD $B$, $\srcBBSch(B)$ is PSD and
\[
 \srcBBSch(A)\circ\srcBBSch(B)=\srcBBSch(A\circ B),\qquad
 I\circ\srcBBSch(B)=h(B)I,\quad h(B)=\prod_i b_{ii}.
\]
For correlation $B$, $L_{(ij)}\circ\srcBBSch(B)=|b_{ij}|^2L_{(ij)}$. Take the Schur product of the lower bound in \eqref{R5E4:eq:sandwich} with $\srcBBSch(B)$ and evaluate on $\mathbf1/\sqrt{n!}$ to obtain \eqref{R5E4:eq:mask}. Its summands are nonnegative. Equality forces $|b_{ij}|=1$ on every support edge, so the unit Gram vectors of $B$ are collinear on each connected component. Conversely, if this holds, $A\circ B$ is a diagonal unitary congruence of $A$ on each component, with zero cross-component entries, and has the same permanent. For general PSD $B$, a zero diagonal gives a zero row and a trivial inequality; otherwise positive diagonal scaling reduces to correlation form. The last inequality $P(B)\ge h(B)$ is the classical permanent Hadamard inequality \cite{R1E5:lieb}.
\end{proof}

\section{Spectral scale, dense examples, and limitations}
When $G_A$ is connected, let $g(A)$ be the gap below the simple eigenvalue $P(A)$ of $\srcBBSch(A)$. The weighted interchange-gap theorem of Caputo--Liggett--Richthammer \cite[Theorem 1.1]{R4E5:CLR}, already used in Round 4 entry 05, identifies the least nonzero eigenvalue of $H_A$ with $\lambda_2(\mathcal L_A)$. Their right-action convention and ours are unitarily equivalent by permutation inversion. Thus
\begin{equation}\label{R5E4:eq:gap}
 \frac16\lambda_2(\mathcal L_A)\le g(A)\le2\lambda_2(\mathcal L_A).
\end{equation}
This is a classical spectral input applied to our new comparison, not a new proof of Aldous' theorem. The upper estimate follows by testing on a least-energy nonconstant vector of $H_A$; the lower follows by minimization. More generally, the min--max principle compares every ordered eigenvalue of $\srcBBDef(A)$ with the corresponding one of $H_A$ by the same factors.

The positivity, exact kernel, and scalar conclusions do not depend on \cite{R4E5:CLR}. For an elementary explicit gap, choose a support spanning tree with diameter $d$ and minimum weight $a_* =\min_{e\in T}w_e^2$. Each permutation is a product of at most $n-1$ arbitrary transpositions. Replacing each by its tree-path word gives length at most $(n-1)(2d-1)$ and uses each edge at most $2(n-1)$ times. Telescoping and averaging squared differences over $S_n$, as in \cite[Section 5]{R5E4:R405}, give
\[
 H_A\succeq\frac{a_*}{2(n-1)^2(2d-1)}(I-\Pi),\qquad
 \Pi=\frac{\mathbf1\mathbf1^*}{n!}.
\]
Indeed $\|(I-L_\pi)z\|^2\le2(n-1)^2(2d-1)\sum_{e\in T}\|(I-L_e)z\|^2$; the average left side is $2\|(I-\Pi)z\|^2$. This reproduces the weaker elementary input explicitly rather than making \eqref{R5E4:eq:gap} a premise for POT.

\paragraph{A dense family with a determined gap order.}
For arbitrary Hermitian phases, set $|a_{ij}|=1/[12(n-1)]$ for every $i\ne j$. The weighted graph is complete, so its Laplacian is
\[
 \mathcal L_A=\frac{nI-J}{144(n-1)^2}.
\]
Hence uniformly over all these phases,
\begin{equation}\label{R5E4:eq:densegap}
 \frac{n}{864(n-1)^2}\le g(A)\le\frac{n}{72(n-1)^2}.
\end{equation}
Together with \eqref{R5E4:eq:perbound}, this determines both the absolute and relative Schur-power gap to order $1/n$. For example, take all upper-triangular phases equal to $1$ except $a_{13}=i/[12(n-1)]$. For $n\ge3$ a triangle has nonreal cycle product, excluding diagonal-unitary equivalence to an entrywise nonnegative matrix. For $n\ge5$ the support is complete and has edges disjoint from cycles, so neither the earlier sparse sign certificates nor the book support criterion explains this family. Gaussian-rational examples occur in every order.

\paragraph{Quantitative domain and genuine limitations.}
The criterion \eqref{R5E4:eq:general} does not require uniform entry bounds and can certify matrices outside \eqref{R5E4:eq:ball}. For example \eqref{R5E4:eq:secondregime} allows an entry as large as $1/32$ independently of order, provided the row and total squared-mass bounds hold. Nevertheless the total mass $X$ is an essential restriction of this proof. A large cycle with every edge $1/100$ has $q=1/50$ but $X=n/10000$, so eventually $t\ge1$ and our criterion says nothing, although the earlier cycle theorem proves POT. We have not proved POT under a dimension-independent row-sum bound alone, for all diagonally dominant matrices, or for all PSD matrices. No optimality is claimed for $1/12$, the scaling $1/n$, or the comparison factors.

The universal ball is much larger in entrywise radius than our preceding $n^{-7}$ cycle neighborhoods, but the centers and off-diagonal amplitudes differ; these are distinct positive domains, not a blanket containment statement. In particular a positive radius at the identity cannot be obtained by simply applying a positive-gap stability lemma, since its defect has full kernel. The edgewise coercive comparison resolves that degeneracy directly.

\paragraph{Failed routes and verification.}
Keeping only single-transposition positivity and bounding all other permutations by their number introduces factorial losses. Cutting a long cycle at an arbitrary edge introduces uncontrolled ratios when some entries vanish. The selected lightest-edge cut and weighted cycle-collection majorants avoid both failures. Dropping the global mass restriction would require controlling arbitrary remote cycle decorations relative to a positive operator, rather than by the scalar factor $e^t$; that is the remaining obstruction in this approach.

The supplied verifier checks exact rational endpoint constants, all chosen permutation words through bounded orders, marked-edge cycle coefficients, cycle-collection moments, the full coefficient budget in \eqref{R5E4:eq:badenergy}, Gaussian-rational Schur quadratic forms, disconnected support degeneracies, and unrestricted PSD-mask tests. Any floating spectra are explicitly separated as diagnostics. These computations audit the written proof; finite success is not used to infer an all-order inequality. No external posting or author contact occurred.



\chapter[Arbitrary-support permanent-on-top balls via cycle-to-edge comparison]{Arbitrary-support permanent-on-top balls 
via cycle-to-edge comparison}
\label{source:R5E5}
\noindent\textit{Source development: Round 5, entry 05. Original mathematical arguments retained; notation is local to this chapter.}\par\medskip

\begin{quote}\small\noindent\textbf{Chapter abstract.} 
We prove a permanent-on-top criterion with no support or phase restriction.
For a Hermitian diagonal-one matrix, let $q$ be the maximum absolute
off-diagonal row sum and let $s$ be the sum of squared edge magnitudes.
If $q\le1/32$ and $s\le1/4$, its permanent defect dominates one half
of the weighted transposition Laplacian. Consequently every such matrix
is positive definite and has permanent on top. In particular, the entire
diagonal-one entrywise ball of radius $1/[32(n-1)]$ about the identity
is certified, including arbitrary zeros and arbitrarily small nonzero
entries. The proof cuts each long permutation cycle at a least-weight
edge, compares its complex operator contribution with edge energies,
and bounds the remaining disjoint-cycle partition sums by a trace
estimate. No prior sparse-support POT theorem is a premise.
We also obtain the exact top-eigenspace multiplicity, quantitative
Bapat--Sunder bounds, and fully dense phase-unrestricted families with
relative top gap of order $1/n$ and neighborhood radius of order $1/n$.
These small-interaction domains do not contain all the earlier
large-edge cycle centers. Neither general Chollet nor a general
diagonal-dominance assertion is resolved; historical novelty remains
unverified.
\end{quote}

\section{Statement, scope, and the change from earlier rounds}
For Hermitian $A$ of order $n$, put
\[
 P(A)=\srcBCper(A),\qquad h(A)=\prod_i a_{ii},\qquad
 \srcBCSch(A)_{\sigma,\tau}=\prod_i a_{\sigma(i),\tau(i)},\qquad
 \srcBCDef(A)=P(A)\srcBCId_{n!}-\srcBCSch(A).
\]
The Schur power is indexed by $\sigma,\tau\in\srcBCG_n$.
For $A\succeq0$ it is PSD: take the Gram matrix of all permuted tensor
products of Gram vectors of $A$.
The permanent-on-top (POT) property means $\srcBCDef(A)\succeq0$.
The universal POT assertion is false \cite{R4E3:Tran}, while the weaker
Chollet inequality
\begin{equation}\label{R5E5:eq:Chollet}
 P(A\circ B)\le P(A)P(B)\qquad(A,B\succeq0)
\end{equation}
remains open in the latest structural source checked \cite{R2E3:PS}.
Our result gives a sufficient first-factor domain, not a universal
claim.

Use the left regular representation $L_\pi e_\tau=e_{\pi\tau}$.
For a diagonal-one Hermitian matrix define
\begin{equation}\label{R5E5:eq:parameters}
 r_{ij}=|a_{ij}|\ (i\ne j),\quad r_{ii}=0,\qquad
 q=\max_i\sum_{j\ne i}r_{ij},\quad
 s=\sum_{i<j}r_{ij}^2,
 \qquad K_A=\sum_{i<j}r_{ij}^2(\srcBCId-L_{(ij)}).
\end{equation}
All the summands of $K_A$ are PSD.

\begin{srcBCtheorem}[Arbitrary-support comparison]\label{R5E5:thm:main}
For every Hermitian diagonal-one matrix $A$ with
\begin{equation}\label{R5E5:eq:small}
 q\le\frac1{32},\qquad s\le\frac14,
\end{equation}
one has $A\succeq(31/32)\srcBCId_n$ and
\begin{equation}\label{R5E5:eq:main}
 0\preceq\srcBCSch(A)\preceq P(A)\srcBCId_{n!},\qquad
 \srcBCDef(A)\succeq\frac12K_A.
\end{equation}
There is no restriction on its support, complex phases, or positive
edge lower bounds. If its support components have orders
$n_1,\ldots,n_k$, the multiplicity of its top eigenvalue $P(A)$ is
exactly $n!/(n_1!\cdots n_k!)$. In particular it is simple precisely
when the support is connected, for $n\ge2$.
\end{srcBCtheorem}

\begin{srcBCcorollary}[Identity-centered balls]\label{R5E5:cor:ball}
For $n\ge2$, every Hermitian diagonal-one matrix satisfying
\begin{equation}\label{R5E5:eq:ball}
 \max_{i\ne j}|a_{ij}|\le\frac1{32(n-1)}
\end{equation}
is positive definite and has POT, with \eqref{R5E5:eq:main}.
In the full space of Hermitian matrices, every $A$ with
\begin{equation}\label{R5E5:eq:fullball}
 \max_{i,j}|a_{ij}-\delta_{ij}|\le\frac1{64(n-1)}
\end{equation}
is also positive definite and has POT.
\end{srcBCcorollary}

\begin{srcBCcorollary}[Dense families with relative gap]\label{R5E5:cor:dense}
Suppose $n\ge2$, $A$ is Hermitian with diagonal one, and
\begin{equation}\label{R5E5:eq:denseband}
 \frac1{64(n-1)}\le |a_{ij}|\le\frac1{32(n-1)}
 \qquad(i\ne j).
\end{equation}
With $Q=\srcBCId_{n!}-\mathbf1\mathbf1^*/n!$,
\begin{equation}\label{R5E5:eq:densegap}
 \srcBCDef(A)\succeq\frac{n}{8192(n-1)^2}Q
 \succeq\frac{P(A)}{16384n}Q.
\end{equation}
Thus any diagonal-one Hermitian center $A_0$ whose off-diagonal
magnitudes all equal $1/[48(n-1)]$, with arbitrary phases, has a whole
diagonal-one Hermitian entrywise neighborhood of radius $1/[192(n-1)]$ satisfying
\eqref{R5E5:eq:denseband} and \eqref{R5E5:eq:densegap}.
\end{srcBCcorollary}
The comparison with the weighted graph Laplacian in this corollary
uses the classical Caputo--Liggett--Richthammer theorem \cite{R4E5:CLR}.
The core operator inequality \eqref{R5E5:eq:main} does not depend on that
theorem. No uniform simple-top gap holds on the whole ball
\eqref{R5E5:eq:ball}, which contains the identity and disconnected supports.

\paragraph{Exact increment and attribution.}
Round 3 Entry 05 proved qualitative POT for cycles and other specified
support blocks \cite{R5E5:R305}. Round 4 Entry 03 added a reflected-companion
criterion, notably for some book blocks \cite{R5E5:R403}. Round 4 Entry 04
proved polynomial-radius dense neighborhoods around cycle centers
with edge magnitudes in $[1/4,1/3]$ \cite[Theorem 1.1]{R5E5:R404}.
Round 4 Entry 05 proved determinant-dependent transposition comparisons
on the connected positive-definite part of the earlier structural
class \cite[Theorem 1.1]{R5E5:R405}.
Here arbitrary supports, including several disjoint long cycles, are
allowed from the outset. The proof does not invoke any of those sparse
POT certificates, a reflected companion, or Schur-minimum refinements.
The new-in-this-work ingredients are the least-edge cycle comparison
and its global cycle-partition estimate. The transposition-operator
viewpoint, scalar consequences, and use of the classical interchange
gap are credited to the earlier treatment rather than counted again.

This changes the geometry of the known domain: zeros can appear or
disappear freely without a nonzero-edge lower bound. The radius
$O(1/n)$ and dense relative gap $O(1/n)$ have explicit constants,
but the centers now have small edge weights. Our domain does not
uniformly subsume the large-edge cycle neighborhoods of Entry 04.
Neither these exponents nor the constants are claimed sharp.

\section{Complex permutation defects controlled by edge energies}
For every Hermitian $A$,
\begin{equation}\label{R5E5:eq:regular}
 d_A(\pi)=\prod_i a_{\pi(i),i},\qquad
 \srcBCSch(A)=\sum_{\pi\in\srcBCG_n}d_A(\pi)L_\pi,
 \qquad \srcBCDef(A)=\sum_\pi d_A(\pi)(\srcBCId-L_\pi).
\end{equation}
Indeed the $(\sigma,\tau)$ entry of the first sum is
$d_A(\sigma\tau^{-1})$, the defining Schur entry.
In particular $\srcBCDef(A)\mathbf1=0$, and $P(A)$ is always an eigenvalue.
Hermitian symmetry pairs inverse coefficients and makes $P(A)$ real.

\begin{srcBClemma}[A word in involutions]\label{R5E5:lem:word}
Let $U_1,\ldots,U_b$ be self-adjoint unitary operators and
$E_j=\langle v,(\srcBCId-U_j)v\rangle$. Then $E_j\ge0$ and
\begin{equation}\label{R5E5:eq:word}
 \left|\langle v,(\srcBCId-U_1\cdots U_b)v\rangle\right|
 \le\left(\sum_{j=1}^b\sqrt{E_j}\right)^2
 \le b\sum_{j=1}^b E_j.
\end{equation}
\end{srcBClemma}
\begin{proof}
Put $V_j=U_1\cdots U_j$ and $V_0=\srcBCId$. Telescoping gives
$\srcBCId-V_b=\sum_jV_{j-1}(\srcBCId-U_j)$.
Separate each $V_{j-1}$ into $\srcBCId+(V_{j-1}-\srcBCId)$ in its quadratic
form. The first terms sum to $\sum_jE_j$.
A reversed telescoping of $V_{j-1}^*=U_{j-1}\cdots U_1$ shows
\[
 \|(V_{j-1}^*-\srcBCId)v\|\le\sum_{i<j}\|(\srcBCId-U_i)v\|.
\]
Since $\|(\srcBCId-U_i)v\|^2=2E_i$, the absolute value of the remaining
terms is at most $2\sum_{i<j}\sqrt{E_iE_j}$.
This proves the first bound; Cauchy--Schwarz proves the second.
The absolute value is important: the argument controls imaginary
cycle contributions, not only real coefficients.
\end{proof}

For each directed simple cycle $\gamma$ of length $\ell\ge3$,
choose an edge of minimum $r_{ij}$, with a fixed tie-breaking rule,
and delete it. The remaining path edges express the cycle permutation
as a product of $\ell-1$ adjacent transpositions along that path.
For a transposition cycle retain its one edge. Apply this independently
to every nontrivial cycle of a permutation $\pi$. Let $E_\pi$ be
the resulting edge set and
\[
 b(\pi)=|E_\pi|=\sum_{\gamma\text{ a nontrivial cycle of }\pi}
 (|\gamma|-1).
\]
The cycles are disjoint, so no edge is repeated in this word.
The orientation of the word is chosen to give $\pi$, not $\pi^{-1}$;
its edge set and length are independent of that convention.
For $\mathcal E_e(v)=\langle v,(\srcBCId-L_e)v\rangle$,
Lemma~\ref{R5E5:lem:word} yields
\begin{equation}\label{R5E5:eq:permword}
 |\langle v,(\srcBCId-L_\pi)v\rangle|
 \le b(\pi)\sum_{e\in E_\pi}\mathcal E_e(v).
\end{equation}

A permutation made only of transpositions and fixed points is a
matching involution. Its coefficient in \eqref{R5E5:eq:regular} is real
nonnegative, so its defect contribution is PSD.
Retain the single-transposition contributions and discard the other
matching contributions. Call a permutation \emph{long} if it has
at least one cycle of length at least three. Taking real parts and
using \eqref{R5E5:eq:permword} gives
\begin{equation}\label{R5E5:eq:charge}
 \langle v,\srcBCDef(A)v\rangle\ge
 \sum_e\left[r_e^2-
 \sum_{\substack{\pi\text{ long}\\e\in E_\pi}}
 |d_A(\pi)|b(\pi)\right]\mathcal E_e(v).
\end{equation}
Zero-weight edges never occur in the inner sum and may be omitted.
We next bound the bracketed loss without dividing by a small edge
unless a still smaller omitted edge compensates for it.

\section{Disjoint-cycle sums and least-edge charging}
Let $R=(r_{ij})$ be the symmetric nonnegative matrix from
\eqref{R5E5:eq:parameters}, with zero diagonal. A directed cycle is counted
modulo cyclic rotation. For length at least three both orientations
are distinct; a transposition is counted once. Write
\[
 w_\gamma=\prod_{ij\in\gamma}r_{ij},\qquad
 T=\sum_\gamma w_\gamma,\quad
 B=\sum_\gamma(|\gamma|-1)w_\gamma,
 \quad T_L=\sum_{|\gamma|\ge3}w_\gamma,\quad
 B_L=\sum_{|\gamma|\ge3}(|\gamma|-1)w_\gamma.
\]
For a transposition $ij$, $w_\gamma=r_{ij}^2$.
These sums are finite, over cycles of the support graph.

\begin{srcBClemma}[Cycle partition bounds]\label{R5E5:lem:partition}
If $q<1$, then
\begin{equation}\label{R5E5:eq:activity}
 T\le a:=\frac{s}{1-q},\qquad
 B\le b:=\frac{2s}{1-q},\qquad
 T_L\le c:=\frac{2sq}{3(1-q)},\qquad
 B_L\le d:=\frac{2sq}{1-q}.
\end{equation}
For any subset $J$ of vertices,
\begin{align}
 \sum_{\pi\in\srcBCG_J}|d_{A[J]}(\pi)|&\le e^T,\label{R5E5:eq:Z}\\
 \sum_{\pi\in\srcBCG_J}b(\pi)|d_{A[J]}(\pi)|&\le e^T B,\label{R5E5:eq:marked}\\
 \sum_{\substack{\pi\in\srcBCG_J\\\pi\text{ long}}}
 (1+b(\pi))|d_{A[J]}(\pi)|
 &\le e^T\{T_L(1+B)+B_L\}.\label{R5E5:eq:longmarked}
\end{align}
\end{srcBClemma}
\begin{proof}
For each $\ell\ge2$, simple cycles are among the closed walks counted
by the trace, and each such cycle has $\ell$ starting points. Thus
\[
 \sum_{|\gamma|=\ell}w_\gamma
 \le\frac{\operatorname{tr}(R^\ell)}{\ell}
 \le\frac{2s}{\ell}q^{\ell-2}.
\]
The second inequality follows by diagonalizing the real symmetric
matrix $R$: its eigenvalues have absolute value at most $q$ by the
row-sum bound, and their squares sum to $\operatorname{tr}(R^2)=2s$.
For odd powers, bounding each eigenvalue's contribution by its absolute
value still gives the stated upper bound.
Summing the geometric bounds, using $1/\ell\le1/2$ for $T$,
$(\ell-1)/\ell\le1$ for $B$, and $1/\ell\le1/3$ for $T_L$,
proves \eqref{R5E5:eq:activity}.

Every permutation is uniquely a vertex-disjoint family of nontrivial
cycles, with fixed points contributing one. Dropping disjointness
restrictions bounds the sum of their weights by
$\prod_\gamma(1+w_\gamma)\le e^T$, proving \eqref{R5E5:eq:Z}.
For \eqref{R5E5:eq:marked}, mark one cycle, pay its factor
$(|\gamma|-1)w_\gamma$, and bound the partition sum on the remaining
vertices by $e^T$. For \eqref{R5E5:eq:longmarked}, a long permutation is
bounded by the sum over its marked long cycles. For a fixed marked
cycle $\gamma$, its word length plus the constant one contributes
$1+(|\gamma|-1)$ and the other cycles contribute their total word
length. Bounds \eqref{R5E5:eq:Z}--\eqref{R5E5:eq:marked} on those other vertices
give
$w_\gamma e^T[1+(|\gamma|-1)+B]$.
Sum over long $\gamma$. Overcounting permutations with several long
cycles is deliberate and harmless because all these weights are
nonnegative.
\end{proof}

\begin{srcBClemma}[Least-edge deletion]\label{R5E5:lem:least}
Fix an edge $e$ with $r_e>0$. For every $\ell\ge3$,
\begin{equation}\label{R5E5:eq:least}
 \sum_{\substack{|\gamma|=\ell\\e\text{ retained in }\gamma}}
 \frac{w_\gamma}{r_e^2}
 \le 2(\ell-1)q^{\ell-2}.
\end{equation}
The sum uses the selected least-edge deletion in each cycle.
\end{srcBClemma}
\begin{proof}
Write $e=ij$ and fix first the orientation $i\to j$.
Let $f\ne e$ be the deleted least edge. Removing both $e$ and $f$
leaves a path starting at $j$ of some length $u\ge0$ and a path
ending at $i$ of length $\ell-2-u$. Reverse the latter, preserving
its weight by symmetry of $R$. The factor from the two removed
edges, after division by $r_e^2$, is $r_f/r_e\le1$.
The cycle is uniquely determined by these two paths and their
endpoints: $f$ joins the two free endpoints. Hence its contribution
is at most the product of the two path weights.
For fixed $u$, drop simplicity, disjointness, and the endpoint-edge
restrictions. The weights of all walks of length $u$ starting at $j$
sum to at most $q^u$; the other walk weights sum to at most
$q^{\ell-2-u}$, by induction from the row-sum bound.
There are $\ell-1$ choices of $u$. Sum over them and then over the
opposite orientation of $e$, giving \eqref{R5E5:eq:least}.
The empty walk has weight one, so the endpoint cases cause no
exception. Ties in the least edge are resolved once, not counted
several times.
\end{proof}

Define for $0\le q<1$
\begin{equation}\label{R5E5:eq:F12}
 F_1(q)=\frac{q(2-q)}{(1-q)^2}
       =\sum_{j\ge1}(j+1)q^j,\qquad
 F_2(q)=\frac{q(4-3q+q^2)}{(1-q)^3}
       =\sum_{j\ge1}(j+1)^2q^j.
\end{equation}
Use $a,b,c,d$ from \eqref{R5E5:eq:activity} and put
\begin{equation}\label{R5E5:eq:Gamma}
 \Gamma(q,s)=e^a\{2F_2(q)+2bF_1(q)+c(1+b)+d\}.
\end{equation}

\begin{srcBCtheorem}[General quantitative criterion]\label{R5E5:thm:criterion}
For any Hermitian diagonal-one $A$ with $q<1$,
\begin{equation}\label{R5E5:eq:general}
 \srcBCDef(A)\succeq(1-\Gamma(q,s))K_A.
\end{equation}
Thus $\Gamma(q,s)\le1$ suffices for POT, and $\Gamma(q,s)<1$
gives a strictly positive transposition comparison on every nonzero
support edge.
\end{srcBCtheorem}
\begin{proof}
Fix $e$ with $r_e>0$ in the loss sum of \eqref{R5E5:eq:charge}.
There are two disjoint possibilities.

If $e$ is a transposition cycle of $\pi$, its weight is exactly
$r_e^2$. Since $\pi$ is long, the permutation on the remaining
vertices is long. Its length contributes $1+b(\pi\setminus e)$.
By \eqref{R5E5:eq:longmarked} the resulting loss divided by $r_e^2$ is
at most $e^T[T_L(1+B)+B_L]$.

If $e$ is retained in a long cycle $\gamma$ of length $\ell$, the
other cycles occupy its complement. Their partition sum and length
sum are bounded by \eqref{R5E5:eq:Z}--\eqref{R5E5:eq:marked}, giving a loss
at most
\[
 e^T\sum_{\substack{\gamma\text{ long}\\e\text{ retained}}}
 w_\gamma\bigl[(|\gamma|-1)+B\bigr].
\]
Divide by $r_e^2$, apply \eqref{R5E5:eq:least}, and sum over lengths.
With $j=\ell-2$, the result is at most
$2e^T[F_2(q)+BF_1(q)]$.
Together the two cases bound the total loss for $e$ by
\[
 r_e^2e^T\{2F_2(q)+2BF_1(q)+T_L(1+B)+B_L\}
 \le r_e^2\Gamma(q,s).
\]
All functions used in the last substitution have nonnegative
coefficients in their series or in their displayed variables.
Insert this edgewise bound into \eqref{R5E5:eq:charge} for every $v$.
This proves the operator inequality.
Finally $q<1$ makes $A$ positive definite, since for any $x$,
$2r_{ij}|x_i||x_j|\le r_{ij}(|x_i|^2+|x_j|^2)$ gives
$x^*Ax\ge(1-q)\|x\|^2$. Its Schur power is therefore PSD,
which turns the upper bound into POT whenever $\Gamma\le1$.
\end{proof}

\section{Constants, eigenspaces, and dense neighborhoods}
\begin{proof}[The comparison in Theorem~\ref{R5E5:thm:main}]
All quantities in \eqref{R5E5:eq:Gamma} increase with $q,s$.
Under \eqref{R5E5:eq:small},
\[
 a\le\frac8{31},\quad b\le\frac{16}{31},\quad
 c\le\frac1{186},\quad d\le\frac1{62},\quad
 F_1(q)\le\frac{63}{961},\quad F_2(q)\le\frac{4001}{29791}.
\]
The elementary series bound $e^x\le(1-x)^{-1}$ for $0\le x<1$
gives $e^a\le31/23$. Direct rational substitution yields
\begin{equation}\label{R5E5:eq:constant}
 \Gamma(q,s)\le\frac{32224}{66309}<\frac12.
\end{equation}
Theorem~\ref{R5E5:thm:criterion} proves \eqref{R5E5:eq:main}.
The same quadratic-form bound gives $A\succeq(31/32)\srcBCId_n$.
\end{proof}

\begin{proof}[Top-eigenspace assertion]
A vector lies in $\ker K_A$ exactly when it is invariant under
every support-edge transposition, since each summand is PSD.
These transpositions generate
$H=\srcBCG_{n_1}\times\cdots\times\srcBCG_{n_k}$ on the support components:
a path of graph edges expresses any transposition within a component.
Thus $\ker K_A$ consists of vectors constant on each right coset
$H\tau$, and has dimension $n!/\prod n_i!$.
Every permutation with nonzero $d_A(\pi)$ preserves each support
component, so its left action fixes all these invariant vectors.
Equation~\eqref{R5E5:eq:regular} shows that $\srcBCDef(A)$ vanishes on
$\ker K_A$. Conversely \eqref{R5E5:eq:main} implies
$\ker\srcBCDef(A)\subseteq\ker K_A$. The kernels are equal, giving
the claimed exact multiplicity.
\end{proof}

\begin{proof}[Corollary~\ref{R5E5:cor:ball}]
Under \eqref{R5E5:eq:ball}, $q\le1/32$ and
\[
 s\le\binom n2\frac1{1024(n-1)^2}
   =\frac{n}{2048(n-1)}\le\frac1{1024}<\frac14.
\]
For \eqref{R5E5:eq:fullball}, write $\epsilon=1/[64(n-1)]$.
Its diagonal entries are positive. Set
$C_{ij}=a_{ij}/\sqrt{a_{ii}a_{jj}}$; then $C$ has diagonal one and
\[
 |c_{ij}|\le\frac{\epsilon}{1-\epsilon}
   =\frac1{64(n-1)-1}\le\frac1{32(n-1)}.
\]
The first assertion gives $C\succ0$ and POT. Positive diagonal
congruence gives $A\succ0$, $P(A)=h(A)P(C)$ and
$\srcBCSch(A)=h(A)\srcBCSch(C)$, so POT transfers to $A$.
\end{proof}

\begin{proof}[Corollary~\ref{R5E5:cor:dense}]
The classical weighted interchange-gap theorem \cite[Theorem 1.1]{R4E5:CLR}
identifies the gap of $\sum_e w_e(\srcBCId-L_e)$ with the second eigenvalue
of the order-$n$ weighted graph Laplacian
$\sum_e w_e(e_i-e_j)(e_i-e_j)^*$.
Permutation inversion unitarily interchanges the left and right
regular conventions, so its rate normalization applies unchanged.
For the complete graph with all weights one, this second eigenvalue
is $n$, since its Laplacian is $n\srcBCId-J_n$.
Under \eqref{R5E5:eq:denseband},
\[
 K_A\succeq\frac1{4096(n-1)^2}
 \sum_{i<j}(\srcBCId-L_{(ij)})\succeq
 \frac{n}{4096(n-1)^2}Q.
\]
Use Theorem~\ref{R5E5:thm:main} for the first inequality of
\eqref{R5E5:eq:densegap}. Also $P(A)\le P(|A|)\le e^T\le31/23<2$
by Lemma~\ref{R5E5:lem:partition}; and $n/(n-1)^2\ge1/n$.
These facts prove its relative-gap version.
The entrywise neighborhood around the stated centers has lower
magnitude $(1/48-1/192)/(n-1)=1/[64(n-1)]$ and upper magnitude
$(1/48+1/192)/(n-1)=5/[192(n-1)]\le1/[32(n-1)]$.
All its phases remain unrestricted. This proves the last assertion.
\end{proof}

More generally, Theorem~\ref{R5E5:thm:main} and the same classical input
bound the absolute top gap below by one half the second eigenvalue
of the weighted support Laplacian with weights $r_e^2$ when the
support is connected. This is the earlier spectral-comparison
principle \cite{R5E5:R405} applied to a new, arbitrary-support domain,
not a new interchange theorem.

\section{Unrestricted second factors and equality}
\begin{srcBCcorollary}\label{R5E5:cor:scalar}
Let $A$ satisfy Theorem~\ref{R5E5:thm:main}. For every PSD $B$,
\begin{equation}\label{R5E5:eq:scalar}
 P(A)h(B)-P(A\circ B)\ge\frac12\sum_{i<j}|a_{ij}|^2
 \left[h(B)-|b_{ij}|^2\prod_{t\ne i,j}b_{tt}\right]\ge0.
\end{equation}
If $B$ is a correlation matrix, equality $P(A\circ B)=P(A)$ holds
exactly when $|b_{ij}|=1$ for every support edge of $A$.
Equivalently its unit Gram vectors are collinear within each support
component; for connected support, this is exactly $\operatorname{rank}B=1$.
In particular the restricted Bapat--Sunder inequality and
\eqref{R5E5:eq:Chollet} hold when this $A$ is the first factor and the
second factor is arbitrary PSD.
\end{srcBCcorollary}
\begin{proof}
Take the Schur product of \eqref{R5E5:eq:main}'s defect comparison with
$\srcBCSch(B)\succeq0$. The Schur product theorem gives a PSD difference,
and \eqref{R5E5:eq:regular} gives
\[
 \srcBCSch(A)\circ\srcBCSch(B)=\srcBCSch(A\circ B),\qquad
 \srcBCId\circ\srcBCSch(B)=h(B)\srcBCId,\qquad
 L_\pi\circ\srcBCSch(B)=d_B(\pi)L_\pi.
\]
Evaluate on $\mathbf1/\sqrt{n!}$ to obtain \eqref{R5E5:eq:scalar}.
Each bracket is nonnegative by the two-by-two principal minor
inequality for $B$. This also handles zero diagonals without division.
For a correlation mask, equality forces the asserted unit magnitudes
on all nonzero edges. Unit Gram vectors are then collinear along
paths in each support component. Their edge entries form a diagonal
unitary gauge on that component, so $A\circ B$ is a blockwise
unitary diagonal congruence of $A$ and preserves its permanent.
This proves sufficiency too.

For completeness, the classical permanent Hadamard inequality
$P(B)\ge h(B)$ supplies the last comparison. One short proof uses
Gram vectors $u_i$ and the polynomial norm
$\|f\|^2=\sum_\alpha\alpha!|f_\alpha|^2$.
Multiplication by $\ell_u=\sum u_kz_k$ has adjoint
$a_u=\sum\overline u_k\partial_k$ and satisfies
$\|\ell_uf\|^2=\|u\|^2\|f\|^2+\|a_uf\|^2$.
Expanding the product gives $P(B)=\|\prod_i\ell_{u_i}\|^2$;
iteration proves $P(B)\ge\prod_i\|u_i\|^2=h(B)$.
This classical fact is not a new result here; see \cite{R1E5:lieb}.
\end{proof}

\section{Examples, exact diagnostics, and limitations}
For every $n\ge4$, set the diagonal to one, give every off-diagonal
entry magnitude $1/[48(n-1)]$, and assign conjugate phases arbitrarily.
Taking one oriented edge phase $i$ and all other upper-triangle
phases $1$ produces a nonreal triangle product, so the matrix is
not phase-balanced. Its support is complete, not a cycle or a
bipartite graph. At $n\ge5$, it also has an edge disjoint from a
triangle, so the earlier edge--cycle-intersection certificates do
not apply. Nevertheless Corollary~\ref{R5E5:cor:dense} certifies it and
its whole diagonal-one Hermitian entrywise ball of radius $1/[192(n-1)]$.
For example, at $n=6$ the center magnitude is $1/240$, the radius
is $1/960$, and the relative gap is at least $1/98304$.
These are rational examples when phases are Gaussian rational units.

The finite exact verifier independently reconstructs permutation
cycle words, least-edge deletion, cycle sums, the two types of edge
charges, and Gaussian-rational Schur-power comparisons at small orders.
It checks the rational constant \eqref{R5E5:eq:constant}, sparse and dense
examples with highly nonuniform weights, disconnected multiplicities,
and scalar inequalities against unrestricted Gram matrices.
Permanents are independently cross-checked by subset dynamic
programming and Ryser's formula. Floating-point spectra are not
proof premises. The all-order proof is the series and operator
argument above, not an extrapolation from those diagnostics.

\paragraph{Scope limitations.}
The separate row-sum and total-square assumptions both matter to
this proof: a small row sum alone can coexist with total mass
proportional to $n$, and our exponential partition bound then
loses control. We do not establish POT for every diagonally dominant
matrix, every bipartite matrix, or every matrix with a PSD reflected
companion. Nor do we prove optimality of the $1/n$ entrywise scale.
The full identity-centered ball contains matrices with every possible
support graph, but it does not have a uniform positive Schur-power
gap because the identity itself has multiplicity $n!$.
The lower edge bound in \eqref{R5E5:eq:denseband} is used only for its
uniform gap, never for POT on the full ball.

\paragraph{Source and priority status.}
Primary-source checks on 20 October 2026 reconfirmed Pant--Singh v2
of 18 September 2026, Tran v5 of 23 August 2022, and
Caputo--Liggett--Richthammer v4 of 28 September 2009 together with
its published theorem. The public Li--Zhao--Luo v1 claim of Chollet
through order six \cite{R1E5:li} is recorded but is neither audited nor
used; even its stated scalar conclusion would not subsume this
unbounded-order POT criterion. Searches for POT near the identity,
diagonal dominance, small entries, and Schur-power spectral gaps
did not establish historical priority. The proof is complete within
its scope; novelty remains unverified, and no general-conjecture
closure or prize-level claim is made.

\begingroup\small
\endgroup

\chapter[Dimension-independent Hadamard contraction for weakly interacting permanents]{Dimension-independent Hadamard contraction for weakly interacting permanents}
\label{source:R6E4}
\noindent\textit{Source development: Round 6, entry 04. Original mathematical arguments retained; notation is local to this chapter.}\par\medskip
\begin{quote}\small\textbf{Joint-credit note.}
The scalar row-local contraction theorem is shared by Round 6 entries
04 and 05 and is counted once. Their scalar result is distinct from
the full Schur-power operator theorem established in Round 7.\end{quote}

\begin{quote}\small\noindent\textbf{Chapter abstract.} 
We prove a quantitative Bapat--Sunder inequality under a dimension-independent absolute row-sum bound, without a total squared-mass restriction. If $A$ is Hermitian with diagonal one and off-diagonal row sums at most $1/12$, then for every correlation matrix $B$ the logarithmic permanent loss lies between $1/10$ and $19/10$ times $\sum_{i<j}|a_{ij}|^2(1-|b_{ij}|^2)$. Under the smaller row bound $1/16$ the factors improve to $3/8$ and $13/8$. Supports, phases, order, and the rank of the mask are unrestricted. Equality is characterized componentwise; consequences include extensive free-energy bounds, rank-one stability, and a graph-Laplacian comparison for the order-$n$ permanental cofactor matrix. The new step relative to the preceding round is connected-cluster cancellation of remote decorations, using classical cluster machinery and the two independently obtained preceding-round cycle-energy estimates. This does not prove a row-only permanent-on-top theorem for the full Schur power or resolve general Chollet. Historical priority of the stated application remains unestablished.
\end{quote}

\section{Main theorem and exact increment}
Write $P(A)=\srcBHper(A)$ and call a PSD Hermitian matrix with diagonal one a correlation matrix. For a diagonal-one Hermitian matrix $A$ define
\[
 w_{ij}=|a_{ij}|\ (i\ne j),\quad w_{ii}=0,\qquad
 q=\max_i\sum_j w_{ij},\qquad X=\sum_{i<j}w_{ij}^2.
\]
For a correlation matrix $B$ of the same order put
\[
 \srcBHE_A(B)=\sum_{i<j}w_{ij}^2(1-|b_{ij}|^2).
\]
Every summand is nonnegative.

\begin{srcBHtheorem}[A row-sum-only logarithmic comparison]\label{R6E4:thm:main}
Let $A$ be Hermitian of any finite order, with diagonal one. If $q\le1/12$, then $A$ is positive definite and, for every correlation matrix $B$,
\begin{equation}\label{R6E4:eq:main}
 \frac1{10}\srcBHE_A(B)\ \le\
 \log\frac{P(A)}{P(A\circ B)}
 \ \le\ \frac{19}{10}\srcBHE_A(B).
\end{equation}
If $q\le1/16$, the stronger bounds are
\begin{equation}\label{R6E4:eq:strong}
 \frac38\srcBHE_A(B)\ \le\
 \log\frac{P(A)}{P(A\circ B)}
 \ \le\ \frac{13}{8}\srcBHE_A(B).
\end{equation}
The logarithms are of positive real numbers. There is no bound on $X$, no support or phase restriction, and no lower bound on a nonzero entry.
\end{srcBHtheorem}

\begin{srcBHcorollary}[Bapat--Sunder, Chollet, and equality]\label{R6E4:cor:equality}
Under $q\le1/12$, every PSD $B$ satisfies
\begin{equation}\label{R6E4:eq:BS}
 P(A\circ B)\le P(A)\prod_i b_{ii}\le P(A)P(B).
\end{equation}
For a correlation mask, $P(A\circ B)=P(A)$ if and only if its restriction to each connected component of the support graph of $A$ has rank one. In particular connected support forces $\operatorname{rank}B=1$.
\end{srcBHcorollary}

For positive diagonal $A$, the same statements apply after replacing $a_{ij}$ by $a_{ij}/\sqrt{a_{ii}a_{jj}}$. Both permanents in the ratio acquire the same factor $\prod_i a_{ii}$, so the logarithmic conclusion transfers exactly.

\paragraph{What is new relative to the frozen pool.}
Round 5 entries 04 and 05 independently obtained the same essential word-energy, least-edge, and decorated-cycle argument \cite{R6E4:R504,R6E4:R505}. Both receive credit for that shared core. Their operator conclusions prove permanent-on-top (POT) on arbitrary supports, but their sufficient criteria retain a global squared mass $X$. Entry 04 also obtained a two-sided operator comparison and the explicit ball $|a_{ij}|\le1/[12(n-1)]$. We reproduce, and do not count again, the two local estimates needed from those papers.

The new step is to apply them to the \emph{connected expansion of the logarithm} rather than the full permutation expansion. Remote cycle collections no longer cost an exponential in $X$. This removes the global-mass hypothesis entirely for the scalar comparison and gives a multiplicative, extensive estimate. It includes the preceding explicit entrywise ball and arbitrarily large connected bounded-degree systems with fixed nonzero edge magnitudes. It does \emph{not} establish POT on this larger domain; a scalar inequality cannot be silently promoted to an operator inequality. All six frozen Round 5 papers and their assessment were read.

\paragraph{Precisely sourced target and status.}
Universal Bapat--Sunder and POT are false \cite{R2E5:drury,R4E3:Tran}. The relevant target is an unrestricted-order positive first-factor class for these inequalities; the weaker general Chollet conjecture remains open in the latest checked Pant--Singh version, v2 of 18 September 2026 \cite{R2E3:PS}. Pioge et al.\ \cite{R3E4:Pioge} study the associated restricted-class questions. We do not claim general closure.

The cluster expansion, tree-graph inequality, and rooted convergence criterion below are classical \cite{R6E4:KP,R6E4:FP}, not new lemmas claimed by this work. Zero-freeness under complex strict diagonal dominance is also already known, as explained by Barvinok with attribution to Brenner \cite[Theorem 1.2]{R6E4:Barvinok}. Cluster methods for permanents themselves have antecedents, including the recent random-matrix application \cite{R6E4:KL}. Our proposed increment is the deterministic, mask-uniform, two-sided logarithmic contraction, not the use of logarithms or polymer expansions in isolation. Current primary-source checks did not establish historical priority.

\section{The classical connected expansion, with proof}
We state the finite-polymer version explicitly to fix all normalization and repeated-polymer issues. A finite polymer list $\mathcal P$ has a symmetric incompatibility relation $\not\sim$, with every polymer incompatible with itself. With complex activities $z_\gamma$, its partition function is
\[
 Z(z)=\sum_{\substack{\mathcal F\subseteq\mathcal P\\
                \text{pairwise compatible}}}\prod_{\gamma\in\mathcal F}z_\gamma .
\]
For an ordered $k$-tuple $\boldsymbol\gamma=(\gamma_1,\ldots,\gamma_k)$, repetitions allowed, let $H(\boldsymbol\gamma)$ be the graph on $[k]$ joining incompatible pairs. Define
\[
 \phi(\boldsymbol\gamma)=
 \sum_{\substack{F\subseteq E(H(\boldsymbol\gamma))\\([k],F)\text{ connected}}}
       (-1)^{|F|},\qquad \phi(\gamma)=1.
\]
Thus two occurrences of the same polymer are joined. For a fixed root define the nonnegative series
\begin{equation}\label{R6E4:eq:rootseries}
 \srcBHLr_\gamma=\sum_{k\ge1}\frac1{(k-1)!}
 \sum_{\gamma_2,\ldots,\gamma_k\in\mathcal P}
 |\phi(\gamma,\gamma_2,\ldots,\gamma_k)|
 \prod_{j=2}^k |z_{\gamma_j}|.
\end{equation}
The $k=1$ term is one.

\begin{srcBHlemma}[Classical tree and convergence bounds]\label{R6E4:lem:cluster}
Suppose numbers $a_\gamma\ge0$ obey
\begin{equation}\label{R6E4:eq:KP}
 \sum_{\delta\not\sim\gamma}|z_\delta|e^{a_\delta}\le a_\gamma
 \qquad(\gamma\in\mathcal P).
\end{equation}
Then $\srcBHLr_\gamma\le e^{a_\gamma}$, and the absolutely convergent expansion
\begin{equation}\label{R6E4:eq:log}
 \log Z(z)=\sum_{k\ge1}\frac1{k!}
   \sum_{\gamma_1,\ldots,\gamma_k}
   \phi(\boldsymbol\gamma)\prod_{j=1}^kz_{\gamma_j}
\end{equation}
defines the logarithm continued from zero activities. The series is uniformly convergent under $z\mapsto tz$, $|t|\le1$, and its exponential equals $Z(tz)$ there. In particular $Z(tz)\ne0$ throughout that closed disk.
\end{srcBHlemma}
\begin{proof}
First, for any finite graph $H$ on $[k]$, the absolute alternating sum over connected spanning subgraphs is at most its number of spanning trees. Here is the classical Penrose partition argument. Root at vertex $1$. For a connected spanning subgraph $F$, use its breadth-first levels and choose for each nonroot vertex its smallest-labeled neighbor in the preceding level as parent. This gives a spanning tree $T$. Given $T$, let $R_H(T)$ consist of its edges, all edges of $H$ within one tree level, and all edges of $H$ between consecutive levels for which the preceding-level endpoint has label at least that of the prescribed parent. Exactly the graphs
\[
 T\subseteq F\subseteq R_H(T)
\]
produce $T$: allowed edges preserve distances and the least-parent rule, and every graph producing $T$ can have only those edges. These intervals partition the connected graphs. Their alternating sums equal
$(-1)^{k-1}(1-1)^{|E(R_H(T))\setminus E(T)|}$, of absolute value at most one. Summing proves the tree bound. The one-vertex case is immediate.

Replace $|\phi|$ in \eqref{R6E4:eq:rootseries} by this tree count. The resulting rooted labeled-tree series is bounded using
\[
 T_\gamma^{(0)}=1,\qquad
 T_\gamma^{(h+1)}
   =\exp\left(\sum_{\delta\not\sim\gamma}|z_\delta|T_\delta^{(h)}\right).
\]
Expansion of the exponential enumerates the unordered branches at the root, with their factorial denominators. Equivalently, the exponential generating function of rooted labeled trees of height at most $h$ is $T_\gamma^{(h)}$: distributing the remaining labels among the branches gives precisely the factorial $(k-1)!$ in \eqref{R6E4:eq:rootseries}. Repeated polymer labels are permitted throughout. Condition \eqref{R6E4:eq:KP} proves inductively $T_\gamma^{(h)}\le e^{a_\gamma}$. Every finite tree has finite height, so monotone convergence gives $\srcBHLr_\gamma\le e^{a_\gamma}$.

For completeness, the formal identity \eqref{R6E4:eq:log} follows by writing
\[
 Z(z)=\sum_{k\ge0}\frac1{k!}\sum_{\gamma_1,\ldots,\gamma_k}
 \prod_jz_{\gamma_j}\prod_{i<j}
       (1-\mathbf1_{\gamma_i\not\sim\gamma_j}).
\]
Expansion of the product sums all graphs. Decomposing a graph into its connected components gives the exponential of the connected-graph series. This is an identity of formal power series, including repeated labels.

The absolute sum of the connected series is at most
$\sum_\gamma |z_\gamma|\srcBHLr_\gamma$, which is finite. The same bound gives uniform convergence for $|t|\le1$. Its exponential agrees with $Z(tz)$ near zero by the formal identity, hence on $|t|<1$ by analytic continuation and on the boundary by continuity. This also proves nonvanishing and the branch assertion without assuming them.
\end{proof}
This proof is included for self-containment, not new credit. It is the standard Penrose tree bound and the Koteck\'y--Preiss rooted majorant; see \cite[Sections 2 and 4]{R6E4:FP}.

\section{Cycle activities and the inherited local energy estimates}
A polymer is a directed simple cycle of length $l\ge2$, considered modulo rotation. A two-cycle is counted once; for $l\ge3$ the two orientations are distinct. Its support is its vertex set, and incompatibility means nonempty intersection. Put
\[
 z_\gamma=\prod_{ij\in\gamma}a_{ij},\qquad
 m_\gamma=\prod_{ij\in\gamma}b_{ij},\qquad
 \rho_\gamma=|z_\gamma|=\prod_{ij\in\gamma}w_{ij}.
\]
For a two-cycle $e=ij$, $z_e=w_e^2$ and $m_e=|b_{ij}|^2$. The cycle decomposition of permutations gives exactly
\begin{equation}\label{R6E4:eq:polymerper}
 P(A)=Z(z),\qquad P(A\circ B)=Z(zm).
\end{equation}
Cycles of zero weight may be omitted.

\begin{srcBHlemma}[A local convergence bound]\label{R6E4:lem:local}
Let $a>0$ and $r=qe^a<1$. If
\begin{equation}\label{R6E4:eq:localcondition}
 \frac{r^2}{1-r}\le a,
\end{equation}
then the cycle system satisfies \eqref{R6E4:eq:KP} with $a_\gamma=a|\gamma|$, and
\begin{equation}\label{R6E4:eq:rootbound}
 \srcBHLr_\gamma\le e^{a|\gamma|},\qquad
 \srcBHLr_\gamma-1\le e^{a|\gamma|}-1.
\end{equation}
These bounds also hold when activities are multiplied by arbitrary factors of modulus at most one.
\end{srcBHlemma}
\begin{proof}
For a fixed vertex $i$, each directed cycle of length $l$ containing $i$ is one closed walk starting at $i$. With $W=(w_{ij})$, this gives
\[
 \sum_{\gamma\ni i}\rho_\gamma e^{a|\gamma|}
 \le\sum_{l\ge2}e^{al}(W^l)_{ii}
 \le\sum_{l\ge2}(qe^a)^l=\frac{r^2}{1-r}.
\]
For length two the first inequality is equality. The row sums of $W^l$ are at most $q^l$, by induction; all entries are nonnegative. Summing this estimate over the vertices of $\gamma$ bounds the activities incompatible with $\gamma$ by $a|\gamma|$. Lemma~\ref{R6E4:lem:cluster} applies. Decreasing activity moduli preserves every bound.
\end{proof}

For each long cycle choose a minimum-weight edge, with a fixed tie-breaking rule, and delete it. Write $F_\gamma$ for the remaining $l-1$ path edges. For a two-cycle $e$ put $F_e=\{e\}$. Put $d_e=1-|b_e|^2\ge0$.

\begin{srcBHlemma}[The shared Round 5 word estimate]\label{R6E4:lem:word}
For every correlation matrix $B$ and every cycle $\gamma$ of length $l$,
\begin{equation}\label{R6E4:eq:maskcycle}
 |1-m_\gamma|\le(l-1)\sum_{e\in F_\gamma}d_e.
\end{equation}
For a two-cycle the equality $1-m_e=d_e$ holds.
\end{srcBHlemma}
\begin{proof}
The word lemma independently proved in \cite[Lemma 2.1]{R6E4:R504} and \cite[Lemma 2.1]{R6E4:R505} states that for self-adjoint unitaries $U_1,\ldots,U_s$ and a unit vector $v$,
\[
 |\langle v,(I-U_1\cdots U_s)v\rangle|
 \le\left(\sum_{j=1}^s
        \sqrt{\langle v,(I-U_j)v\rangle}\right)^2
 \le s\sum_{j=1}^s\langle v,(I-U_j)v\rangle .
\]
To verify it, telescope $I-U_1\cdots U_s=\sum_jV_{j-1}(I-U_j)$, where $V_{j-1}=U_1\cdots U_{j-1}$. If $E_j=\langle v,(I-U_j)v\rangle$, then $\|(I-U_j)v\|=\sqrt{2E_j}$. A reversed-prefix telescope gives
$\|(V_{j-1}^*-I)v\|\le\sum_{i<j}\sqrt{2E_i}$. The $j$th quadratic-form term is therefore at most $E_j+2\sum_{i<j}\sqrt{E_iE_j}$ in absolute value. Summing proves the estimate, including its imaginary part.

Take unit Gram vectors $u_i$ for $B$, their tensor product $v=\bigotimes_i u_i$, and the unitary representation that permutes tensor positions. A swap $ij$ has matrix coefficient $|b_{ij}|^2$ at $v$. A cycle has coefficient $m_\gamma$ or its conjugate, depending on the convention, which has the same distance from one. A path obtained by deleting one cycle edge expresses that cycle as the product of its $l-1$ successive edge swaps. Apply the word estimate to this path.
\end{proof}

\begin{srcBHlemma}[The shared Round 5 least-edge charge]\label{R6E4:lem:cut}
For every positive-weight edge $e$ and $l\ge3$,
\begin{equation}\label{R6E4:eq:cut}
 \sum_{\substack{|\gamma|=l\\e\in F_\gamma}}
       \frac{\rho_\gamma}{w_e^2}
 \le2(l-1)q^{l-2}.
\end{equation}
\end{srcBHlemma}
\begin{proof}
This is \cite[Lemma 3.3]{R6E4:R504} and \cite[Lemma 3.2]{R6E4:R505}, reproduced here. If $f$ is the chosen deleted edge, then $f\ne e$ and $w_f\le w_e$. Consequently
\[
 \rho_\gamma/w_e^2\le
 \prod_{h\in\gamma\setminus\{e,f\}}w_h.
\]
Choose an orientation of $e=ij$. Removing $e,f$ leaves two paths from its endpoints of lengths $k$ and $l-2-k$, where $0\le k\le l-2$. Their free endpoints determine $f$, so these paths and the orientation encode the original cycle injectively. Dropping simplicity, disjointness, and the condition on the joining edge bounds their total weight by $q^kq^{l-2-k}$. There are $l-1$ length choices and two orientations. Summing proves \eqref{R6E4:eq:cut}. No division by a tiny edge remains in the bound.
\end{proof}

\section{Logarithmic comparison without remote mass}
Define
\[
 s_2(r)=\sum_{l\ge3}(l-1)^2r^{l-2}
       =\frac{1+r}{(1-r)^3}-1,
 \qquad
 \Theta(q,a)=e^{2a}-1+2e^{2a}s_2(qe^a).
\]

\begin{srcBHtheorem}[General scalar criterion]\label{R6E4:thm:general}
Suppose $A$ is Hermitian with diagonal one, $a>0$, $r=qe^a<1$, and \eqref{R6E4:eq:localcondition} holds. For every correlation $B$,
\begin{equation}\label{R6E4:eq:general}
 \left|\log\frac{P(A)}{P(A\circ B)}-\srcBHE_A(B)\right|
 \le\Theta(q,a)\srcBHE_A(B).
\end{equation}
In particular $\Theta(q,a)<1$ gives a strict Bapat--Sunder inequality whenever some positive-weight edge has $|b_{ij}|<1$.
\end{srcBHtheorem}
\begin{proof}
Since $q<1$, for every $v\in\srcBHC^n$,
\[
 v^*Av\ge\sum_i\left(1-\sum_{j\ne i}w_{ij}\right)|v_i|^2
 \ge(1-q)\|v\|^2.
\]
The same holds for $A\circ B$, since $|b_{ij}|\le1$. Thus both are positive definite. The permanent of a PSD Gram matrix is the squared norm of the product of its associated linear forms, in the polynomial inner product $\|f\|^2=\sum_\alpha\alpha!|f_\alpha|^2$; it is strictly positive for a positive definite matrix. Alternatively positivity follows directly from the real nonvanishing continuation below.

Lemmas~\ref{R6E4:lem:cluster} and \ref{R6E4:lem:local} give absolutely convergent expansions for both expressions in \eqref{R6E4:eq:polymerper}. For real $0\le t\le1$, the partition functions $Z(tz)$ and $Z(tzm)$ are real: reverse every cycle to conjugate its weight. They start at one and never vanish by Lemma~\ref{R6E4:lem:cluster}, so remain positive. The continued logarithms at $t=1$ are therefore the real logarithms in \eqref{R6E4:eq:general}, not values differing by an untracked multiple of $2\pi i$.

Subtract the two connected expansions. A tuple contributes
\[
 \frac{\phi(\boldsymbol\gamma)}{k!}
 \prod_j z_{\gamma_j}\left(1-\prod_jm_{\gamma_j}\right).
\]
Since $|m_\gamma|\le1$, a scalar telescope gives
\begin{equation}\label{R6E4:eq:scalartelescope}
 \left|1-\prod_jm_{\gamma_j}\right|
 \le\sum_j|1-m_{\gamma_j}|.
\end{equation}
The single-polymer two-cycle terms sum exactly to $\srcBHE_A(B)$. In all other terms, mark the index in \eqref{R6E4:eq:scalartelescope}. Symmetry of the ordered sums changes $1/k!$ into $1/(k-1)!$. Hence the absolute remainder after subtracting $\srcBHE_A(B)$ is at most
\begin{equation}\label{R6E4:eq:marked}
 \sum_{e}w_e^2d_e(\srcBHLr_e-1)
 +\sum_{\gamma:\,|\gamma|\ge3}\rho_\gamma
       |1-m_\gamma|\srcBHLr_\gamma .
\end{equation}
In the first sum the subtraction of one removes precisely the single-polymer two-cycle term; no longer cluster is omitted. Repetitions in clusters are allowed and already counted correctly by the marked factorial.

The first sum is at most $(e^{2a}-1)\srcBHE_A(B)$. For the second, use \eqref{R6E4:eq:rootbound} and \eqref{R6E4:eq:maskcycle}, then collect the coefficient of each $d_e$:
\begin{align*}
 \sum_{\gamma:\,|\gamma|\ge3}\rho_\gamma e^{a|\gamma|}
 (|\gamma|-1)\sum_{e\in F_\gamma}d_e
 &\le\sum_e w_e^2d_e
       \sum_{l\ge3}2(l-1)^2e^{al}q^{l-2}\\
 &=2e^{2a}s_2(qe^a)\srcBHE_A(B).
\end{align*}
The inequality is Lemma~\ref{R6E4:lem:cut}. The series is summable because $r<1$; all intermediate majorants are nonnegative. This proves \eqref{R6E4:eq:general}. Crucially only clusters joined to the marked cycle enter its rooted bound; no factor depending on the total number or squared mass of distant edges appears.
\end{proof}

\begin{proof}[Rational constants for Theorem~\ref{R6E4:thm:main}]
All displayed majorants increase in $q$. Use $e^x\le(1-x)^{-1}$ for $0\le x<1$. When $q\le1/12$, choose $a=1/64$. Then
\[
 r\le\frac{16}{189},\qquad
 \frac{r^2}{1-r}\le\frac{256}{32697}<\frac1{64},\qquad
 e^{2a}\le\frac{32}{31}.
\]
Thus the exact rational lower coefficient is
\[
 1-\Theta(q,a)\ge
 2-\frac{32}{31}\left(1+2s_2(16/189)\right)
 =\frac{18045878}{160509227}>\frac1{10}.
\]
When $q\le1/16$, choose $a=1/128$. Similarly
\[
 r\le\frac8{127},\qquad
 \frac{r^2}{1-r}\le\frac{64}{15113}<\frac1{128},\qquad
 e^{2a}\le\frac{64}{63},
\]
and
\[
 1-\Theta(q,a)\ge
 2-\frac{64}{63}\left(1+2s_2(8/127)\right)
 =\frac{41471090}{106165017}>\frac38.
\]
The upper coefficients in \eqref{R6E4:eq:general} are $1+\Theta=2-(1-\Theta)$. This proves both assertions with rational, not numerical, exponential estimates.
\end{proof}

\begin{proof}[Corollary~\ref{R6E4:cor:equality}]
The positive lower bound forces $d_e=0$ on every support edge in the equality case. Unit Gram vectors of $B$ are then collinear along each support component, which is exactly rank one on that principal block. Conversely such collinearity gives a diagonal unitary gauge on each component. The zero cross-component entries of $A$ are unchanged, so $A\circ B$ is a blockwise diagonal unitary congruence of $A$ and has the same permanent. This includes isolated vertices and $A=I$.

For an arbitrary PSD mask with positive diagonal, normalize by $b_{ij}/\sqrt{b_{ii}b_{jj}}$ and scale the permanent. A zero diagonal entry forces its row to vanish, and both $P(A\circ B)$ and $\prod_i b_{ii}$ are zero. The last inequality in \eqref{R6E4:eq:BS} is the classical permanent Hadamard inequality $P(B)\ge\prod_i b_{ii}$ \cite{R1E5:lieb}. One direct verification uses the same polynomial norm: multiplication by the linear form $\ell_u$ has adjoint $a_u=\sum\overline u_j\partial_j$ and
$\|\ell_uf\|^2=\|u\|^2\|f\|^2+\|a_uf\|^2$. Iteration gives the stated bound.
\end{proof}

\section{Extensive estimates, block cuts, and stability}
In this section $(c,C)$ denotes $(1/10,19/10)$ under $q\le1/12$, or $(3/8,13/8)$ under $q\le1/16$.

\begin{srcBHcorollary}[Free energy and block interaction]\label{R6E4:cor:blocks}
One has
\begin{equation}\label{R6E4:eq:free}
 cX\le\log P(A)\le CX.
\end{equation}
For every partition $V=V_1\sqcup\cdots\sqcup V_k$, let
$X_{\rm cut}=\sum_{i<j:\,i,j\text{ in different blocks}}w_{ij}^2$. Then
\begin{equation}\label{R6E4:eq:block}
 cX_{\rm cut}\le
 \log P(A)-\sum_{\nu=1}^k\log P(A[V_\nu])
 \le CX_{\rm cut}.
\end{equation}
\end{srcBHcorollary}
\begin{proof}
Use $B=I$ for \eqref{R6E4:eq:free}. For \eqref{R6E4:eq:block}, take the correlation matrix having entry one within each block and zero between blocks; it is the Gram matrix of one orthonormal vector per block. Then $A\circ B$ is block diagonal and its permanent factors. Its mask energy is exactly $X_{\rm cut}$.
\end{proof}
The qualitative block lower bound is classical Lieb supermultiplicativity, not new credit. The stated two-sided quantitative logarithmic estimate is a consequence of the new row-only comparison.

\begin{srcBHcorollary}[Quantitative rank-one stability]\label{R6E4:cor:stability}
Suppose the support graph of $A$ is connected, with at least two vertices, and let $\lambda_2>0$ be the second eigenvalue of its weighted graph Laplacian
\[
 \mathcal L_A=\sum_{i<j}w_{ij}^2(e_i-e_j)(e_i-e_j)^*.
\]
For every correlation matrix $B$,
\begin{equation}\label{R6E4:eq:stability}
 \log\frac{P(A)}{P(A\circ B)}
 \ge\frac{c\lambda_2}{2}
        \left(n-\frac{\srcBHtr(B^2)}n\right)
 \ge\frac{c\lambda_2}{2}\bigl(n-\lambda_{\max}(B)\bigr).
\end{equation}
\end{srcBHcorollary}
\begin{proof}
Let $B$ be the Gram matrix of unit vectors $u_i$, put $Q_i=u_iu_i^*$ and $\bar Q=n^{-1}\sum_iQ_i$. Then
\[
 \|Q_i-Q_j\|_F^2=2(1-|b_{ij}|^2),\qquad
 \sum_i\|Q_i-\bar Q\|_F^2=n-\srcBHtr(B^2)/n.
\]
Apply the ordinary graph-Laplacian Poincar\'e inequality to each real coordinate of these Hermitian matrices and sum. It gives
$2\srcBHE_A(B)\ge\lambda_2(n-\srcBHtr(B^2)/n)$.
Finally $\srcBHtr(B^2)\le\lambda_{\max}(B)\srcBHtr B=n\lambda_{\max}(B)$.
No interchange-process or Schur-power spectral theorem is invoked.
\end{proof}
For example a logarithmic loss at most $\varepsilon$ forces
$n-\lambda_{\max}(B)\le2\varepsilon/(c\lambda_2)$. The dependence on connectivity is necessary: distinct support components may carry independent rank-one Gram directions.

\begin{srcBHcorollary}[A quantitative permanental cofactor comparison]\label{R6E4:cor:cofactor}
Let
\[
 (F_A)_{ij}=a_{ij}\srcBHper(A(i\mid j)),
\]
where $A(i\mid j)$ deletes row $i$ and column $j$, and the empty permanent is one. With the same constants $c,C$,
\begin{equation}\label{R6E4:eq:cofactor}
 c\mathcal L_A\preceq I-\frac{F_A}{P(A)}
                  \preceq C\mathcal L_A .
\end{equation}
The eigenspace of $F_A$ at its largest eigenvalue $P(A)$ consists exactly of vectors constant on each support component. In particular its multiplicity is the number of components. For connected support of order at least two its gap below the top lies between $cP(A)\lambda_2$ and $CP(A)\lambda_2$.
\end{srcBHcorollary}
\begin{proof}
The qualitative implication from Bapat--Sunder to this cofactor top-eigenvalue property, and the rank-two mask perturbation used here, are due to Pioge et al.\ \cite[Lemma 1 and Theorem 1]{R3E4:Pioge}. We apply their mechanism to both sides of the quantitative estimate.

For arbitrary $x\in\srcBHC^n$ take the correlation mask
\[
 b_{ij}(s)=\frac{1+s\overline{x_i}x_j}
 {\sqrt{(1+s|x_i|^2)(1+s|x_j|^2)}}\qquad(s\ge0).
\]
It is the Gram matrix of the normalized vectors $(1,\sqrt{s}\,x_i)$. Direct expansion gives
\[
 b_{ij}(s)=1+s\left(\overline{x_i}x_j-\frac{|x_i|^2+|x_j|^2}{2}\right)+O(s^2),
 \quad
 \srcBHE_A(B(s))=s\,x^*\mathcal L_Ax+O(s^2).
\]
Hermiticity makes $F_A$ Hermitian, and permanent expansion gives every row and column sum equal to $P(A)$. Differentiating the permanent entry by entry therefore gives
\[
 \log\frac{P(A)}{P(A\circ B(s))}
 =s\,x^*\left(I-\frac{F_A}{P(A)}\right)x+O(s^2).
\]
Apply Theorem~\ref{R6E4:thm:main}, divide by $s$, and let $s\downarrow0$. This proves \eqref{R6E4:eq:cofactor} for every complex $x$. Its two PSD comparisons identify its kernel with $\ker\mathcal L_A$, which consists of the componentwise constants. The gap bounds follow by the Rayleigh principle. This is an order-$n$ cofactor statement, not an order-$n!$ Schur-power assertion.
\end{proof}

\paragraph{Connected systems with unbounded squared mass.}
Let $G=C_m\square K_2$, $m\ge3$, be the prism graph, of order $n=2m$ and degree three. Give every edge magnitude $1/48$, assign arbitrary conjugate phases, and set the diagonal to one. Then
\[
 q=1/16,\qquad X=\frac{n}{1536},\qquad
 \frac{n}{4096}\le\log P(A)\le\frac{13n}{12288}.
\]
The conclusions hold for every correlation mask even as $X\to\infty$. This is a connected example, not merely a direct sum of previously small instances. The support has overlapping cycles and arbitrary frustration. At sufficiently large order the global-mass criteria of both Round 5 entries fail. At $n=6144$, for instance, $X=4$: entry 04's required $t<1$ already fails, and entry 05's exponential charge is greater than one.

The domain also contains dense neighborhoods of large sparse systems. If $A_0$ is Hermitian with diagonal one and $q(A_0)\le1/32$, every Hermitian diagonal-one perturbation satisfying
\[
 \max_{i\ne j}|a_{ij}-(a_0)_{ij}|
 \le\frac1{32(n-1)}
\]
has $q(A)\le1/16$, and therefore satisfies \eqref{R6E4:eq:strong}. For a prism center with edge magnitude $1/96$, its squared mass is $n/6144$, unbounded with order. Arbitrarily small nonzero entries on all missing edges give fully dense examples. This is a direct geometric consequence of the new domain, not another proof mechanism.

\section{Limitations, source audit, and exact diagnostics}
This paper deliberately changes the conclusion from an operator inequality to a scalar logarithmic one. The classical cluster formula used here is commutative. Substituting permutation operators into it would require a different argument, because overlapping cycle operators need not commute. We have not removed the global-mass obstruction for POT, proved Bapat--Sunder for every diagonally dominant matrix, or resolved general Chollet. The sufficient constants are not claimed optimal. In particular scalar mask tests, however extensive, cannot certify an operator spectral statement.

An initial attempt to retain the full permutation expansion still incurred a factor exponential in $X$. Replacing remote sums by principal permanents directly did not give a valid uniform sign comparison, because complex long-cycle coefficients can have either sign. Neither attempt is a proof. Connected scalar logarithms remove precisely the remote-decoration obstruction; the repeated-polymer and branch checks in Section 2 are essential, rather than formal omissions.

Primary sources were rechecked on 20 October 2026. Pant--Singh's latest record is v2; Tran's is v5; Pioge et al.'s is v1. The original Koteck\'y--Preiss paper and Fern\'andez--Procacci's published treatment explicitly supply the classical ingredients reproduced here; the latter's latest preprint is v2 of 16 October 2006. Barvinok's latest v2 of 11 September 2018 already proves a much larger zero-free region, with Brenner's antecedent noted there. Koehler--Leung v1 of 1 April 2026 provides additional permanent/cluster-expansion context, not the present Hadamard comparison. The public through-order-six Chollet manuscript \cite{R1E5:li} is recorded as an unaudited claim, not a proof dependency. No literature search here certifies first discovery.

The verifier uses exact integers and rational intervals. It audits Penrose partitions on all graphs through a specified small order, connected-series coefficients against independently expanded partition polynomials, the least-edge and Gram-cycle estimates, the rational convergence and contraction constants, and full Gaussian-rational permanent ratios against rational enclosures of their logarithms. It includes exact anchored $LDL^*$ certificates for the cofactor comparisons, the rank-two derivative identity, disconnected equality, block masks, and the extensive prism construction. All-order validity comes from the proof, not from finite tests. No external posting, submission, or author contact occurred.

The next substantive problem in this line is an operator-valued connected expansion or another mechanism proving a row-only POT theorem. A second question is whether the scalar row threshold can be extended to the full strictly diagonally dominant range. Neither is asserted here.

\begingroup\small

\endgroup

\chapter[Row-local Bapat--Sunder bounds without a global mass restriction]{Row-local Bapat--Sunder bounds without a global mass restriction}
\label{source:R6E5}
\noindent\textit{Source development: Round 6, entry 05. Original mathematical arguments retained; notation is local to this chapter.}\par\medskip
\begin{quote}\small\textbf{Joint-credit note.}
The scalar row-local contraction theorem is shared by Round 6 entries
04 and 05 and is counted once. Their scalar result is distinct from
the full Schur-power operator theorem established in Round 7.\end{quote}

\begin{quote}\small\noindent\textbf{Chapter abstract.} 
For a Hermitian diagonal-one matrix $A$, suppose its maximum absolute
 off-diagonal row sum is at most $1/12$. For every correlation matrix
 $B$, let $E=\sum_{i<j}|a_{ij}|^2(1-|b_{ij}|^2)$. We prove
 $E/10\le\log(\srcBIper A/\srcBIper(A\circ B))\le19E/10$, in every order,
 with arbitrary support and phases and no bound on the total squared
 off-diagonal mass. A sharper estimate has relative error at most
 $11q$, where $q$ is that row sum. A rank-two mask variation yields a
 two-sided comparison between the order-$n$ permanental-cofactor defect
 and $\srcBIper(A)$ times the weighted support Laplacian. The new step is
 local connected-cluster accounting: classical hard-core expansion
 bounds replace the global exponential decoration loss in the shared
 Round 5 proofs of Entries 04 and 05. Their word-energy and least-edge
 lemmas are reused with joint credit. The conclusion is a restricted
 Bapat--Sunder theorem, not full Schur-power permanent-on-top; the two
 operator assertions are explicitly distinguished. Historical priority
 is unverified, and neither general Chollet nor all diagonal dominance
 is resolved.
\end{quote}

\section{The theorem and the precise increment}
Write $P(A)=\srcBIper(A)$. A correlation matrix is Hermitian PSD with all
 diagonal entries one. For a Hermitian diagonal-one matrix $A$, set
\begin{equation}\label{R6E5:eq:basic}
 r_{ij}=|a_{ij}|\quad(i\ne j),\quad r_{ii}=0,\qquad
 q=\max_i\sum_jr_{ij},\qquad
 s=\sum_{i<j}r_{ij}^2.
\end{equation}
For a correlation matrix $B$, define the nonnegative edge loss
\begin{equation}\label{R6E5:eq:E}
 E_A(B)=\sum_{i<j}r_{ij}^2(1-|b_{ij}|^2).
\end{equation}

\begin{srcBItheorem}[A dimension-independent local criterion]\label{R6E5:thm:main}
Let $A$ be Hermitian with diagonal one and $q\le1/12$. Then
 $A\succeq(11/12)\srcBIId$, and for every correlation matrix $B$,
\begin{equation}\label{R6E5:eq:main}
 \frac1{10}E_A(B)\ \le\
 \log\frac{P(A)}{P(A\circ B)}
 \ \le\ \frac{19}{10}E_A(B).
\end{equation}
Both permanents in this expression are positive. There is no
 restriction on $n$, support, phases, or $s$.
 If $q\le1/16$, the respective constants may be replaced by
 $1/3$ and $5/3$. More generally, throughout $q\le1/12$,
\begin{equation}\label{R6E5:eq:asymp}
 \left|\log\frac{P(A)}{P(A\circ B)}-E_A(B)\right|
 \le11q\,E_A(B).
\end{equation}
Equality $P(A\circ B)=P(A)$ holds exactly when $|b_{ij}|=1$ on
 every support edge of $A$, equivalently when the unit Gram vectors
 of $B$ are collinear within each support component.
\end{srcBItheorem}
The constants are sufficient, not optimal. In particular this theorem
 proves a restricted Bapat--Sunder inequality against an
 \emph{unrestricted} PSD second factor. The normalization and zero
 diagonal cases are addressed in Section~\ref{R6E5:sec:consequences}.
 Taking $B=\srcBIId$ also gives
 $|\log P(A)-s|\le11qs$; this is a consequence of the same theorem,
 not an additional discovery.

For row and column deletions let $A(i\mid j)$ denote the matrix with
 row $i$ and column $j$ removed, and put
\begin{equation}\label{R6E5:eq:cofactor}
 F_A=\bigl(a_{ij}\srcBIper A(i\mid j)\bigr)_{i,j=1}^n,
 \qquad \mathcal L_A=\sum_{i<j}r_{ij}^2(e_i-e_j)(e_i-e_j)^*.
\end{equation}
The empty permanent is one. For Hermitian PSD $A$, the matrix $F_A$
 is Hermitian PSD and all its row sums equal $P(A)$.

\begin{srcBIcorollary}[Permanental-cofactor spectral comparison]\label{R6E5:cor:cofactor}
Under the hypotheses of Theorem~\ref{R6E5:thm:main},
\begin{equation}\label{R6E5:eq:cofactorcomparison}
 \frac{P(A)}{10}\mathcal L_A
 \preceq P(A)\srcBIId_n-F_A
 \preceq\frac{19P(A)}{10}\mathcal L_A.
\end{equation}
The kernel is exactly the functions constant on support components.
 Hence the multiplicity of the top eigenvalue $P(A)$ of $F_A$ is
 the number of these components, and it is simple for connected support.
 For connected support with $n\ge2$, if $g_F$ is the gap below $P(A)$,
\begin{equation}\label{R6E5:eq:gap}
 \frac1{10}\lambda_2(\mathcal L_A)
 \le\frac{g_F}{P(A)}
 \le\frac{19}{10}\lambda_2(\mathcal L_A).
\end{equation}
The comparison factors $1/3,5/3$ apply when $q\le1/16$.
\end{srcBIcorollary}
This is an order-$n$ operator statement, \textbf{not} a bound on the
 $n!$-by-$n!$ Schur power. The qualitative implication from
 Bapat--Sunder to the largest eigenvalue of $F_A$ is a prior theorem
 of Pioge et al.\ \cite[Section 3]{R6E5:Pioge}. We apply the same
 infinitesimal-mask principle to obtain the quantitative comparison;
 the principle is not claimed new.

\paragraph{Target status and reuse credit.}
Universal Bapat--Sunder and permanent-on-top are false
 \cite{R2E5:drury,R4E3:Tran}. The restricted-class question remains meaningful;
 it is explicitly motivated by Pioge et al.\ \cite{R6E5:Pioge}.
 General Chollet, $P(A\circ B)\le P(A)P(B)$ for PSD factors, remains
 open in Pant--Singh v2, the latest structural source checked
 \cite{R2E3:PS}. Nothing here asserts its resolution.

Round 5 Entries 04 and 05 independently obtained the same essential
 complex-word/least-edge/cycle-decoration POT mechanism
 \cite[Lemmas 2.1 and 3.3, Theorem 4.1]{R6E5:R504}
 \cite[Lemmas 2.1 and 3.2, Theorem 3.3]{R6E5:R505}.
 This is one shared antecedent, credited to both. Their estimates
 control all remote cycles by a global factor exponential in $s$.
 Neither gave a dimension-independent row-sum-only result.
 We retain their local word and least-edge estimates, but replace
 unconnected decorations by the \emph{connected expansion of a scalar
 logarithm}. Only cycles linked to a marked local cycle must then be
 paid for. Classical Koteck\'y--Preiss/tree-graph machinery supplies
 the convergence estimate \cite{R6E4:KP,R6E4:FP}; it is not new machinery here.

The resulting domain allows arbitrarily large $s$, and the second
 factor is arbitrary. The tradeoff is explicit: our new result is
 scalar Bapat--Sunder and its order-$n$ cofactor consequence, not the
 previous full Schur-power comparison. It is not a blanket strengthening
 of the earlier POT theorem or of the large-edge cycle neighborhoods
 from Round 4. Barvinok's approximation theorem for complex diagonally
 dominant permanents \cite{R6E5:Barvinok} is a related classical analytic
 antecedent, not the mask-monotonicity inequality proved here.
 All six frozen Round 5 papers and their assessment were read.

\section{Two inherited local estimates, with proofs}
A directed simple cycle is counted modulo rotation. Its reverse is
 distinct for length at least three; a transposition is counted once.
 Let $V(\gamma)$ be its vertex set, $\ell_\gamma=|V(\gamma)|$, and
\[
 z_\gamma=\prod_{i\in V(\gamma)}a_{\gamma(i),i},\qquad
 w_\gamma=|z_\gamma|=\prod_{i\in V(\gamma)}r_{\gamma(i),i}.
\]
Only supported cycles of length at least two are used. For a
 transposition $ij$, $z_\gamma=w_\gamma=r_{ij}^2$.
 Delete a selected least-weight edge of each long cycle, using a fixed
 tie rule. Its remaining path gives a word for the cycle in
 $\ell_\gamma-1$ transpositions, with edge set $T_\gamma$.
 For a two-cycle, $T_\gamma$ is its single edge.

\begin{srcBIlemma}[Word loss for a correlation mask]\label{R6E5:lem:word}
For a correlation matrix $B$, put
 $\beta_\gamma=\prod_{i\in V(\gamma)}b_{\gamma(i),i}$ and
 $\epsilon_{ij}=1-|b_{ij}|^2$. Then $|\beta_\gamma|\le1$, and
\begin{equation}\label{R6E5:eq:wordmask}
 |1-\beta_\gamma|
 \le (\ell_\gamma-1)\sum_{e\in T_\gamma}\epsilon_e.
\end{equation}
For a transposition $ij$, $1-\beta_\gamma=\epsilon_{ij}$ exactly.
\end{srcBIlemma}
\begin{proof}
We recall the jointly attributed word inequality. For self-adjoint
 unitaries $U_1,\ldots,U_b$, let
 $E_j=\langle v,(\srcBIId-U_j)v\rangle\ge0$.
 Telescoping $\srcBIId-U_1\cdots U_b$ and then the reversed adjoint prefixes
 gives
\[
 |\langle v,(\srcBIId-U_1\cdots U_b)v\rangle|
 \le\sum_jE_j+2\sum_{i<j}\sqrt{E_iE_j}
 \le b\sum_jE_j.
\]
Indeed $\|(\srcBIId-U_j)v\|^2=2E_j$, and the adjoint prefix minus
 identity applied to $v$ has norm at most
 $\sum_{i<j}\sqrt{2E_i}$. This controls complex quadratic forms,
 not just their real parts.

Choose unit Gram vectors with $b_{ij}=\langle u_i,u_j\rangle$,
 using the convention linear in the second slot. Let
 $v=\bigotimes_i u_i$ and let $U_\pi$ permute tensor positions by
 $U_\pi(\bigotimes_i u_i)=\bigotimes_i u_{\pi^{-1}(i)}$.
 Then $\|v\|=1$ and
 $\langle v,U_\pi v\rangle=\prod_i b_{\pi(i),i}$.
 For an edge transposition the energy is $1-|b_{ij}|^2$.
 Apply the word inequality to the chosen path expression for $\gamma$.
 The modulus bound follows either from Cauchy--Schwarz or entrywise
 $|b_{ij}|\le1$.
\end{proof}

\begin{srcBIlemma}[Least-edge charging]\label{R6E5:lem:least}
For any edge $e$ with $r_e>0$ and every $\ell\ge3$,
\begin{equation}\label{R6E5:eq:least}
 \sum_{\substack{\ell_\gamma=\ell\\e\in T_\gamma}}
 \frac{w_\gamma}{r_e^2}\le2(\ell-1)q^{\ell-2}.
\end{equation}
Moreover, for each vertex $v$ and each $a\ge0$ with $t=qe^a<1$,
\begin{equation}\label{R6E5:eq:localmass}
 \sum_{\gamma\ni v}w_\gamma e^{a\ell_\gamma}
 \le\frac{t^2}{1-t}.
\end{equation}
\end{srcBIlemma}
\begin{proof}
For \eqref{R6E5:eq:least}, fix one orientation of $e=ij$ and let $f\ne e$
 be the deleted least edge. Removing $e,f$ leaves two paths, starting
 at $i,j$ after reversing one of them, of lengths $u$ and
 $\ell-2-u$. The factor removed after division by $r_e^2$ is
 $r_f/r_e\le1$. The two paths determine $f$ and the original directed
 cycle, so this encoding is injective. For each $u$, dropping
 simplicity, disjointness and endpoint restrictions bounds the total
 weight by $q^u q^{\ell-2-u}$. There are $\ell-1$ choices of $u$
 and two orientations. Empty paths have weight one; ties are counted
 only once. This is the shared Round 5 least-edge argument.

For \eqref{R6E5:eq:localmass}, write $R=(r_{ij})$. Every directed
 length-$\ell$ cycle containing $v$ gives a distinct closed walk
 starting at $v$, including the transposition convention at $\ell=2$.
 Thus its total weight is at most $(R^\ell)_{vv}\le q^\ell$,
 by the row-sum bound on nonnegative walk weights. Sum
 $\sum_{\ell\ge2}(qe^a)^\ell$.
\end{proof}

\section{Classical connected-cluster control, made explicit}\label{R6E5:sec:cluster}
We give the finite-polymer lemma needed below, including its proof.
 This is classical hard-core cluster expansion, with a simple
 tree-graph majorant of Koteck\'y--Preiss type \cite{R6E4:KP,R6E4:FP}.

Let $\mathcal P$ be a finite family with a symmetric reflexive
 incompatibility relation $\not\sim$. For complex activities
 $(z_\gamma)$ define
\[
 Z(z)=\sum_{\substack{\mathcal F\subseteq\mathcal P\\
                    \mathcal F\text{ pairwise compatible}}}
       \prod_{\gamma\in\mathcal F}z_\gamma.
\]
For a tuple $\boldsymbol\gamma=(\gamma_1,\ldots,\gamma_k)$, repetitions
 allowed, form its incompatibility graph on $[k]$. Put
\[
 \phi(\boldsymbol\gamma)=
 \sum_{\substack{H\text{ connected spanning}\\
          H\subseteq G(\boldsymbol\gamma)}}(-1)^{|E(H)|},
 \qquad \phi(\gamma)=1.
\]
For a disconnected incompatibility graph the sum is zero.

\begin{srcBIlemma}[Rooted cluster estimate]\label{R6E5:lem:cluster}
Suppose $w_\gamma\ge |z_\gamma|$ and $a_\gamma\ge0$ satisfy
\begin{equation}\label{R6E5:eq:KP}
 \sum_{\delta\not\sim\gamma}w_\delta e^{a_\delta}
 \le a_\gamma\qquad(\gamma\in\mathcal P).
\end{equation}
Then the series
\begin{equation}\label{R6E5:eq:clusterexp}
 C(z)=\sum_{k\ge1}\frac1{k!}
       \sum_{\gamma_1,\ldots,\gamma_k}
       \phi(\boldsymbol\gamma)\prod_{j=1}^kz_{\gamma_j}
\end{equation}
converges absolutely, $Z(z)=e^{C(z)}\ne0$, and for every root $\gamma$,
\begin{equation}\label{R6E5:eq:root}
 Q_\gamma:=\sum_{k\ge1}\frac1{(k-1)!}
 \sum_{\gamma_2,\ldots,\gamma_k}
 |\phi(\gamma,\gamma_2,\ldots,\gamma_k)|
 \prod_{j=2}^kw_{\gamma_j}\le e^{a_\gamma}.
\end{equation}
The root's $k=1$ term is one, so its higher-cluster sum is at most
 $e^{a_\gamma}-1$.
\end{srcBIlemma}
\begin{proof}
First, for a graph $G$ on $k$ vertices,
\begin{equation}\label{R6E5:eq:treegraph}
 |\phi(G)|\le\tau(G),
\end{equation}
where $\tau$ counts spanning trees, with $\tau=0$ for disconnected
 graphs and $\tau=1$ for a single vertex. Here is an elementary
 proof of this classical special tree-graph inequality. Inclusion--
 exclusion for proper colorings gives
 $\chi_G(x)=\sum_{F\subseteq E(G)}(-1)^{|F|}x^{c(F)}$;
 hence $[x]\chi_G=\phi(G)$. Deletion--contraction shows that
 $c(G)=(-1)^{k-1}[x]\chi_G$ obeys
 $c(G)=c(G-e)+c(G/e)$, with disconnected terms interpreted as zero.
 Parallel edges after contraction may be merged for $\chi$.
 Induction from edgeless graphs gives $0\le c(G)\le\tau(G)$:
 spanning-tree deletion--contraction has the same recurrence before
 merging parallel edges, and merging can only decrease its count.
 This proves \eqref{R6E5:eq:treegraph}.

Replace the cluster coefficient in \eqref{R6E5:eq:root} by the sum over
 labeled spanning trees and root each tree at its first vertex.
 If only trees of height at most $h$ are allowed, their root generating
 functions $T_\gamma^{(h)}$ satisfy
\[
 T_\gamma^{(0)}=1,\qquad
 T_\gamma^{(h+1)}=
 \exp\left(\sum_{\delta\not\sim\gamma}
                  w_\delta T_\delta^{(h)}\right).
\]
To justify the factorials, delete the root: the unordered collection
 of child subtrees gives the exponential series, with $1/m!$ for
 its $m$ children, and partitions of the other vertex labels give
 the factor $1/(k-1)!$ in \eqref{R6E5:eq:root}. Repeated polymer labels
 cause no exception. By induction using \eqref{R6E5:eq:KP},
 $T_\gamma^{(h)}\le e^{a_\gamma}$. Monotone convergence in $h$
 proves \eqref{R6E5:eq:root}. Multiplying by $w_\gamma$ and summing over
 the finite root set bounds the absolute series \eqref{R6E5:eq:clusterexp}.

Finally, write the compatible-family polynomial as an exponential
 generating sum over ordered tuples, inserting for every pair the
 factor $1-\mathbf1_{\gamma_i\not\sim\gamma_j}$. Repeated labels
 vanish by reflexivity. Expanding the pair factors gives all graphs;
 uniquely splitting a graph into connected components gives the
 formal identity $Z=\exp C$. This is the usual exponential formula,
 verified coefficient by coefficient. Substitute $\lambda z$.
 Absolute convergence on $|\lambda|\le1$, just proved, makes the
 identity analytic on $|\lambda|<1$, first near zero by the formal
 identity and then throughout by uniqueness of analytic continuation.
 Taking $\lambda\uparrow1$ by absolute convergence proves
 $Z(z)=e^{C(z)}$ at the endpoint as well.
\end{proof}

\section{Local logarithmic mask comparison}\label{R6E5:sec:proof}
Use the cycles of $A$ as polymers, incompatible exactly when their
 vertex sets intersect. The cycle expansion of a permutation gives
\begin{equation}\label{R6E5:eq:perZ}
 Z(z)=P(A),\qquad Z(z\beta)=P(A\circ B).
\end{equation}
Cycles in a compatible collection are vertex-disjoint. The
 logarithmic series, however, permits repetitions and arbitrary
 connected overlaps; those terms must not be discarded.

\begin{srcBIproposition}[General local estimate]\label{R6E5:prop:general}
Suppose $q<1$ and $a\ge0$ satisfy, with $t=qe^a$,
\begin{equation}\label{R6E5:eq:condition}
 t<1,\qquad \frac{t^2}{1-t}\le a.
\end{equation}
Let
\begin{equation}\label{R6E5:eq:eta}
 F_2(t)=\sum_{j\ge1}(j+1)^2t^j
       =\frac{t(4-3t+t^2)}{(1-t)^3},\qquad
 \eta(q,a)=e^{2a}\{1+2F_2(t)\}-1.
\end{equation}
For every correlation matrix $B$,
\begin{equation}\label{R6E5:eq:general}
 \left|\log\frac{P(A)}{P(A\circ B)}-E_A(B)\right|
 \le\eta(q,a)E_A(B).
\end{equation}
In particular $\eta<1$ suffices for the restricted Bapat--Sunder
 inequality and its stated equality characterization.
\end{srcBIproposition}
\begin{proof}
Give cycle $\gamma$ the parameter $a_\gamma=a\ell_\gamma$.
 Lemma~\ref{R6E5:lem:least} and a union bound on its vertices give
\[
 \sum_{\delta\not\sim\gamma}w_\delta e^{a\ell_\delta}
 \le\ell_\gamma\frac{t^2}{1-t}\le a\ell_\gamma.
\]
Thus Lemma~\ref{R6E5:lem:cluster} applies simultaneously to $z$ and $z\beta$.
 Both arrays are closed under conjugation by reversing every cycle.
 Reversal preserves incompatibility, so their absolutely convergent
 cluster sums are real. By \eqref{R6E5:eq:perZ} and $Z=e^C$, these sums
 are the ordinary real logarithms of the positive permanents.
 Alternatively their positivity follows from strict diagonal dominance:
 both $A$ and $A\circ B$ have row sum at most $q<1$ and are PD.

In the difference of the two cluster expansions a tuple contributes
\[
 \frac{\phi(\boldsymbol\gamma)}{k!}
 \left(\prod_jz_{\gamma_j}\right)
 \left(1-\prod_j\beta_{\gamma_j}\right).
\]
The singleton transposition terms are exactly $E_A(B)$. All other
 terms are bounded using
\[
 \left|1-\prod_j\beta_{\gamma_j}\right|
 \le\sum_j|1-\beta_{\gamma_j}|,
\]
which follows by telescoping since $|\beta_{\gamma_j}|\le1$.
 Mark the selected summand in this last sum. Symmetry of the tuple
 converts $1/k!$ to $1/(k-1)!$. Absolute convergence legitimizes
 all the regrouping. Singleton transposition terms have already been
 removed, whereas singleton long cycles have not. The remaining
 absolute error is consequently at most
\begin{equation}\label{R6E5:eq:errorbudget}
 \sum_{e}r_e^2\epsilon_e(Q_e-1)
 +\sum_{\gamma:\ell_\gamma\ge3}
 w_\gamma(\ell_\gamma-1)
 \sum_{e\in T_\gamma}\epsilon_e Q_\gamma.
\end{equation}
Here $Q_e$ denotes the rooted sum for the two-cycle $e$.
 This marking is the crucial local step: remote disconnected
 cycles have disappeared in the logarithm, and the length cost
 belongs to the marked cycle, not to the whole decorated permutation.

Use $Q_e-1\le e^{2a}-1$ and $Q_\gamma\le e^{a\ell_\gamma}$.
 For a fixed edge, Lemma~\ref{R6E5:lem:least} bounds its long-cycle
 charge divided by $r_e^2$ by
\[
 \sum_{\ell\ge3}2(\ell-1)^2 e^{a\ell}q^{\ell-2}
 =2e^{2a}F_2(qe^a).
\]
Inserting this in \eqref{R6E5:eq:errorbudget} proves
 \eqref{R6E5:eq:general}. No sum of all cycle activities or total mass
 $s$ was introduced. A zero edge is never divided by.
\end{proof}

\begin{proof}[Proof of Theorem~\ref{R6E5:thm:main}]
If $q=0$, then $A=\srcBIId$ and all assertions are immediate.
 For $0<q\le1/12$ choose $a=2q^2$ and let
\[
 u(q)=\frac{q}{1-2q^2},\qquad
 H(q)=\frac{1+2F_2(u(q))}{1-4q^2}-1.
\]
The elementary bound $e^x\le1/(1-x)$ for $0\le x<1$ gives
 $t\le u(q)$ and $\eta(q,2q^2)\le H(q)$. Moreover
\[
 \frac{u(q)^2}{(1-u(q))q^2}
 =\frac1{(1-2q^2)(1-q-2q^2)}
 \le\frac{5184}{4615}<2.
\]
Thus \eqref{R6E5:eq:condition} holds with strict room. Every coefficient
 of the Taylor series of $H(q)$ is nonnegative, its constant term
 is zero, and therefore both $H(q)$ and $H(q)/q$ increase on this
 interval. Exact rational substitution yields
\begin{equation}\label{R6E5:eq:constants}
 H(1/12)=\frac{8448929}{9611875}<\frac9{10},\qquad
 H(1/16)=\frac{64693927}{106165017}<\frac23,
 \qquad 12H(1/12)<11.
\end{equation}
Proposition~\ref{R6E5:prop:general} proves \eqref{R6E5:eq:main}, its
 $1/16$ version, and \eqref{R6E5:eq:asymp}. The sharper coefficient
 $\eta(q,2q^2)$ is $8q+O(q^2)$ as $q\downarrow0$;
 neither this upper-bound coefficient nor the displayed constants
 are claimed sharp.

Strict diagonal dominance gives
 $x^*Ax\ge(1-q)\|x\|^2$ using
 $2r_{ij}|x_i||x_j|\le r_{ij}(|x_i|^2+|x_j|^2)$.
 The same argument applies to $A\circ B$.
 The lower bound in \eqref{R6E5:eq:main} forces $E_A(B)=0$ in equality.
 Unit Gram vectors of $B$ are then collinear along each support
 edge and hence along each path in a component. Conversely those
 edge entries form a diagonal unitary gauge on each component,
 so $A\circ B$ is blockwise diagonal-unitarily congruent to $A$
 and has the same permanent. This proves the equality statement.
\end{proof}

\section{Scalar normalization and the \texorpdfstring{order-$n$}{order-n} operator}\label{R6E5:sec:consequences}
For arbitrary PSD $B$ with all $b_{ii}>0$, set
 $C_{ij}=b_{ij}/\sqrt{b_{ii}b_{jj}}$ and $h(B)=\prod_i b_{ii}$.
 Then $C$ is a correlation matrix and
 $P(A\circ B)=h(B)P(A\circ C)$. Thus Theorem~\ref{R6E5:thm:main} gives
\begin{equation}\label{R6E5:eq:generalB}
 e^{-19E_A(C)/10}
 \le\frac{P(A\circ B)}{P(A)h(B)}
 \le e^{-E_A(C)/10}\le1.
\end{equation}
If a diagonal of $B$ is zero, PSD forces its entire row and column
 to be zero; both $P(A\circ B)$ and $P(A)h(B)$ vanish. No normalized
 ratio is asserted in that case. The classical permanent Hadamard
 inequality $P(B)\ge h(B)$ then gives Chollet with this first factor.
 For completeness, with Gram vectors $u_i$ and polynomial norm
 $\|f\|^2=\sum_\alpha\alpha!|f_\alpha|^2$, multiplication by
 $\ell_u=\sum u_jz_j$ has adjoint $a_u=\sum\bar u_j\partial_j$ and
\[
 \|\ell_uf\|^2=\|u\|^2\|f\|^2+\|a_uf\|^2,
 \qquad P(B)=\left\|\prod_i\ell_{u_i}\right\|^2.
\]
Iteration proves $P(B)\ge\prod_i\|u_i\|^2=h(B)$; this is classical,
 not a contribution here \cite{R1E5:lieb}. Positive diagonal congruence
 also transfers all scalar claims to a first factor with positive
 diagonal whose correlation normalization satisfies the row-sum bound.

\begin{proof}[Proof of Corollary~\ref{R6E5:cor:cofactor}]
First recall why $F_A\succeq0$. The Schur power
 $\srcBISch(A)_{\sigma,\tau}=\prod_i a_{\sigma(i),\tau(i)}$ is a Gram
 matrix of permuted tensor products and is PSD. Its regular expansion
 is $\srcBISch(A)=\sum_\pi d_A(\pi)L_\pi$, where
 $d_A(\pi)=\prod_i a_{\pi(i),i}$ and $L_\pi e_\tau=e_{\pi\tau}$.
 The isometry
\[
 (Vx)(\sigma)=\frac{x_{\sigma(1)}}{\sqrt{(n-1)!}}
\]
intertwines $\srcBISch(A)V=VF_A$, since
 $\sum_{\pi:\pi(j)=i}d_A(\pi)=a_{ij}\srcBIper A(i\mid j)$.
 Thus $F_A=V^*\srcBISch(A)V\succeq0$. Permanent expansion gives both
 its row and column sums $P(A)$.

The following infinitesimal-mask implication is the quantitative use
 of the prior principle in \cite{R6E5:Pioge}. Fix $x\in\srcBIC^n$ and, for
 sufficiently small real $t\ge0$, form the unit vectors
\[
 u_i(t)=\bigl(\sqrt{1-t|x_i|^2},\sqrt t\,x_i\bigr),
 \qquad B(t)_{ij}=\langle u_i(t),u_j(t)\rangle.
\]
They give correlation matrices of rank at most two, with
\[
 B(t)_{ij}=1+t\left(\bar x_i x_j-
                   \frac{|x_i|^2+|x_j|^2}{2}\right)+O(t^2),
 \quad 1-|B(t)_{ij}|^2=t|x_i-x_j|^2+O(t^2).
\]
Differentiating the finite permanent polynomial and using its row
 and column sums yields
\[
 \left.\frac{d}{dt}P(A\circ B(t))\right|_{t=0}
 =x^*F_Ax-P(A)\|x\|^2.
\]
Consequently the derivatives of the middle term in \eqref{R6E5:eq:main}
 and of $E_A(B(t))$ at zero are, respectively,
 $x^*(P(A)\srcBIId-F_A)x/P(A)$ and $x^*\mathcal L_Ax$.
 Divide by $t>0$ and let $t\downarrow0$ to obtain
 \eqref{R6E5:eq:cofactorcomparison}. The stronger factors follow identically.

The two-sided comparison identifies the kernel with that of
 $\mathcal L_A$, namely componentwise constants. On connected support with $n\ge2$,
 minimize the two quadratic forms on $\mathbf1^\perp$; the lower
 comparison and testing on a least graph-eigenvalue vector prove
 \eqref{R6E5:eq:gap}. No interchange-gap theorem is needed for this
 order-$n$ statement. The same min--max argument compares every
 ordered eigenvalue of the two defects by the displayed factors.
\end{proof}

\section{An extensive, nonbalanced family and remaining limits}
For an even integer $m\ge6$, take the Cartesian torus
 $C_m\mathbin\square C_m$ on $n=m^2$ vertices. Set $a_{ii}=1$,
 put magnitude $1/48$ on every graph edge and zero elsewhere,
 and assign conjugate complex phases arbitrarily. Then
\begin{equation}\label{R6E5:eq:torus}
 q=\frac1{12},\qquad s=\frac{m^2}{1152}.
\end{equation}
Take one edge phase $i$ and all the others $1$; a square through that
 edge has nonreal cycle product, so this example is not phase-balanced.
 The graph is bipartite but has an edge vertex-disjoint from a square.
 It has no cut vertex: after deleting a vertex, choose a different
 column; its vertical cycle survives, and every remaining row is a
 cycle or a path connected to that column. Thus the earlier
 edge--cycle-intersection and cut-vertex certificates do not explain
 this nonbalanced single block.

The total mass tends to infinity. For $m\ge34$, even the condition
 $t_*<1$ in Round 5 Entry 04's general global-mass criterion fails:
 at this row sum its $t_*=(35/33)s>1$. The corresponding global
 exponential bound in Entry 05 likewise fails for all sufficiently
 large $m$. The new theorem applies uniformly, with any correlation
 mask, despite this extensive mass. No permanent of the large torus
 needs to be evaluated in order to certify the theorem's hypotheses.

The cofactor comparison gives an explicit relative spectral scale.
 The Fourier modes on the torus show
 $\lambda_2(\mathcal L_A)=\sin^2(\pi/m)/576$:
 each one-dimensional cycle has eigenvalues
 $4\sin^2(\pi k/m)$, and the two coordinate Laplacians add and
 are multiplied by $1/48^2$. Hence
\begin{equation}\label{R6E5:eq:torusgap}
 \frac{\sin^2(\pi/m)}{5760}
 \le\frac{g_F}{P(A)}
 \le\frac{19\sin^2(\pi/m)}{5760}.
\end{equation}
This is an order-$n$ cofactor gap of order $1/n$, not a Schur-power
 gap. The phases do not affect this bound. Such spectral consequences
 of a graph comparison are classical; the new input is the
 local-mass-free comparison itself.

\paragraph{Why no full POT conclusion is asserted.}
The scalar hard-core logarithm in Section~\ref{R6E5:sec:cluster} uses
 commuting activities. Substitution of noncommuting cycle operators
 cannot be justified by that scalar exponential formula. Even a
 tempting coefficientwise-positivity shortcut fails. If $X_1,\ldots,X_4$
 commute mutually but not with $Y$, put
 $M=\prod_{j=1}^4(1-X_j)-Y$. In the algebra with just these
 commutations, the coefficient of the word $YX_1X_2X_3X_4$ in
 $-\log M$ is $-1/30$, not nonnegative. Indeed its coefficient in
 $M^{-k}$, for positive integer $k$, is
 $\sum_{h=1}^k h^4=k^5/5+k^4/2+k^3/3-k/30$;
 differentiate the formal polynomial at $k=0$ to obtain the stated
 logarithmic coefficient. The finite exact verifier independently
 expands this coefficient. This only refutes that proof shortcut,
 not a row-sum-only POT theorem.

\paragraph{Scope, novelty, and reproducibility.}
We do not establish POT for $q\le1/12$, Bapat--Sunder for all
 diagonally dominant matrices, or general Chollet. A general row-sum-only
 POT comparison remains a separate target. Classical cluster expansion,
 the jointly credited Round 5 lemmas, and the known infinitesimal
 Bapat-to-cofactor implication are explicit inputs, not new discoveries.
 The proposed increment is their connected, local mask-loss combination,
 its uniform two-sided logarithmic estimate, and the resulting
 global-mass-free cofactor comparison.

Primary sources and their latest records were checked on 20 October
 2026: Pant--Singh v2 (18 September 2026), Pioge et al.\ v1
 (31 July 2025), Barvinok v2 (11 September 2018), and
 Fern\'andez--Procacci v2 (16 October 2006). The original
 Koteck\'y--Preiss publication is from 1986. Searches for Bapat--Sunder
 row-sum and diagonal-dominance criteria did not settle historical
 priority. The public through-order-six Chollet claim \cite{R1E5:li}
 is neither audited nor used and does not supply the present
 unrestricted-order, unrestricted-second-factor assertion.

The verifier uses exact Gaussian-rational permanents, rational
 logarithm intervals, all small incompatibility graphs for the
 tree-graph estimate, truncated cluster coefficients, least-edge
 charges, unrestricted Gram masks, cofactor operator comparisons,
 and the noncommutative failed-route coefficient. These are bounded
 diagnostics of a complete all-order proof, not numerical evidence
 substituted for one. Large torus checks verify the hypotheses and
 growth of $s$, not the exponentially large permanent. No external
 posting, submission or author contact occurred. No historical
 priority, optimal constants or prize-level claim is made.

\begingroup\small
\endgroup

\chapter[A row-local permanent-on-top theorem via heap derivatives]{A row-local permanent-on-top theorem via heap derivatives}
\label{source:R7E4}
\noindent\textit{Source development: Round 7, entry 04. Original mathematical arguments retained; notation is local to this chapter.}\par\medskip

\begin{quote}\small\noindent\textbf{Chapter abstract.} 
We prove full Schur-power permanent-on-top under a dimension-independent small absolute row-sum condition. If $A$ is Hermitian with diagonal one and $\max_i\sum_{j\ne i}|a_{ij}|\le1/64$, its $n!$-dimensional Schur power has largest eigenvalue $\srcBNper A$, with exact support-component multiplicity and a quantitative gap. More strongly, the Schur power of $I+t(A-I)$ divided by its permanent decreases in Loewner order for $0\le t\le1$. No total squared-mass bound is imposed. The preceding round's independent scalar connected-cluster theorems of entries 04 and 05 did not prove this operator conclusion. We bypass their noncommutative-logarithm obstruction by a classical Cartier--Foata heap identity for the right logarithmic derivative, then use an explicit positive Sylvester inverse. Classical heap and cluster machinery and the shared Round 5 word/least-edge estimates are credited. This is a restricted positive class, not a revival of the false universal POT conjecture; historical priority of the application remains unestablished.
\end{quote}

\section{The operator theorem and the exact increment}
For a Hermitian matrix $A=(a_{ij})$ with diagonal one put
\[
 w_{ij}=|a_{ij}|\ (i\ne j),\quad w_{ii}=0,\qquad
 q=\max_i\sum_jw_{ij},\qquad X=\sum_{i<j}w_{ij}^2.
\]
Its Schur power acts on $\srcBNC^{S_n}$ and is
\[
 \srcBNSch(A)_{\sigma,\tau}=\prod_{i=1}^n a_{\sigma(i),\tau(i)}.
\]
Let $L_\pi e_\tau=e_{\pi\tau}$ be the left regular representation. Then
\begin{equation}\label{R7E4:eq:regular}
 \srcBNSch(A)=\sum_{\pi\in S_n}d_A(\pi)L_\pi,
 \qquad d_A(\pi)=\prod_i a_{\pi(i),i}.
\end{equation}
Every row sum is $P(A):=\srcBNper A$. Write
\begin{equation}\label{R7E4:eq:H}
 H_A=\sum_{i<j}w_{ij}^2(\srcBNId-L_{(ij)})\succeq0,
 \qquad \srcBNK_A=\ker H_A.
\end{equation}
Zero-weight edges are omitted. The $n!$-dimensional identity is denoted by $\srcBNId$.

\begin{srcBNtheorem}[Row-local full POT and radial monotonicity]\label{R7E4:thm:main}
Suppose $q\le1/64$. For $0\le t\le1$ set
\[
 A_t=\srcBNId_n+t(A-\srcBNId_n),\quad S(t)=\srcBNSch(A_t),\quad
 P(t)=\srcBNper A_t,\quad T(t)=S(t)/P(t).
\]
These matrices are positive definite and $P(t)>0$. The normalized Schur power is Loewner decreasing:
\begin{equation}\label{R7E4:eq:monotone}
 0\prec T(t)\preceq T(s)\preceq\srcBNId\qquad(0\le s\le t\le1).
\end{equation}
For every $t>0$ its eigenspace at eigenvalue one is exactly $\srcBNK_A$. If the support components have orders $n_1,\ldots,n_k$, this eigenvalue has multiplicity
\begin{equation}\label{R7E4:eq:multiplicity}
 \dim\srcBNK_A=\frac{n!}{\prod_{j=1}^k n_j!}.
\end{equation}
In particular $\lambda_{\max}(\srcBNSch(A))=P(A)$, simple for connected support. There is no bound on $X$, no phase restriction, and no lower bound on a positive edge weight.
\end{srcBNtheorem}
The multiplicity in \eqref{R7E4:eq:multiplicity} belongs to the full regular representation. It is not the number of components, which is the multiplicity for the order-$n$ cofactor matrix in the preceding round.

The quantitative derivative statement underlying the theorem is as follows. For an operator $D$ put $\srcBNRea D=(D+D^*)/2$.
\begin{srcBNtheorem}[Accretive right logarithmic derivative]\label{R7E4:thm:derivative}
Under the same hypothesis, for $0<t\le1$ let
\begin{equation}\label{R7E4:eq:J}
 J(t)=\frac{tP'(t)}{P(t)}\srcBNId-tS'(t)S(t)^{-1}
       =-tT'(t)T(t)^{-1}.
\end{equation}
Then for every $v\in\srcBNC^{S_n}$,
\begin{equation}\label{R7E4:eq:numerical}
 \left|\langle v,(J(t)-2t^2H_A)v\rangle\right|
 \le\frac95t^2\langle v,H_Av\rangle.
\end{equation}
Consequently
\begin{equation}\label{R7E4:eq:accretive}
 \frac15t^2H_A\preceq\srcBNRea J(t)\preceq\frac{19}5t^2H_A.
\end{equation}
The right logarithmic derivative need not be Hermitian; no such assumption is made.
\end{srcBNtheorem}

\begin{srcBNcorollary}[A full-Schur-power gap]\label{R7E4:cor:gap}
If $A$ has a nonzero off-diagonal entry, let $\delta_A$ be the smallest positive eigenvalue of $H_A$, and let $\Pi_A$ project onto $\srcBNK_A$. Then
\begin{equation}\label{R7E4:eq:gap}
 T(t)\preceq\Pi_A+e^{-\delta_A t^2/10}(\srcBNId-\Pi_A).
\end{equation}
In particular the gap below the top eigenspace of $\srcBNSch(A)$ is at least
$P(A)(1-e^{-\delta_A/10})$.
\end{srcBNcorollary}

\paragraph{Why this is a new conclusion relative to the frozen pool.}
Round 5 entries 04 and 05 independently proved the same word-energy/least-edge/decorated-cycle core, obtaining full POT with a global squared-mass restriction \cite{R7E4:R504,R7E4:R505}. Round 6 entries 04 and 05 independently removed that restriction for a \emph{scalar} logarithmic Bapat--Sunder comparison, and hence the order-$n$ cofactor comparison \cite{R7E4:R604,R7E4:R605}. Both independent authors receive credit for each shared core, counted once. We reuse their local estimates and classical scalar cluster bounds, not count them again.

The new step is a noncommutative \emph{right logarithmic derivative}, not a noncommutative logarithm. Its heap expansion retains the multiplication order and is controlled by scalar connected coefficients only after taking an absolute majorant. A positive Sylvester inverse then proves radial Loewner monotonicity and full order-$n!$ POT. Our row constant is smaller than Round 6's scalar $1/12$, so this is not a blanket improvement of every preceding statement. It removes the global mass obstruction for the stronger operator conclusion.

\paragraph{Precisely sourced status.}
Universal POT and Bapat--Sunder are false \cite{R4E3:Tran,R2E5:drury}; the target is a meaningful unrestricted-order first-factor class, as motivated by Pioge et al.\ \cite{R6E5:Pioge}. General Chollet remains open in the latest checked Pant--Singh v2 \cite{R2E3:PS}, and is not resolved here. Cartier--Foata inversion and heap/pyramid theory are classical \cite{R7E4:CF,R7E4:Kratt}, as are the Penrose and Koteck\'y--Preiss estimates \cite{R6E4:KP,R6E4:FP}. We give the elementary identities and convergence argument needed for this application. We do not claim to have invented heaps, connected expansions, the Sylvester integral, or a new universal conjecture. Primary-source checks on 20 October 2026 did not establish historical priority of the stated application. All six frozen Round 6 papers and the assessment were read; no live opponent paper was inspected.

\section{Cycle words and the inherited local estimates}
A polymer is a supported directed simple cycle $\gamma$ of length $\ell_\gamma\ge2$, modulo rotation. A transposition occurs once; the two orientations of longer cycles are distinct. Two polymers are incompatible if their supports meet, including a polymer with itself. Define
\[
 z_\gamma=\prod_{i\in\gamma}a_{\gamma(i),i},\quad
 \rho_\gamma=|z_\gamma|,\quad U_\gamma=L_\gamma.
\]
Compatible polymers have disjoint supports, so their $U_\gamma$ commute. The cycle decomposition in \eqref{R7E4:eq:regular} gives
\begin{equation}\label{R7E4:eq:cyclepartition}
 S(t)=\sum_{\mathcal F\text{ compatible}}
       \prod_{\gamma\in\mathcal F}t^{\ell_\gamma}z_\gamma U_\gamma,
 \qquad
 P(t)=\sum_{\mathcal F\text{ compatible}}
       \prod_{\gamma\in\mathcal F}t^{\ell_\gamma}z_\gamma.
\end{equation}
All sets are finite here. For each long cycle delete a minimum-weight edge, with a fixed tie rule. The remaining path expresses the cycle as $\ell_\gamma-1$ edge transpositions; denote this list of edges by $F_\gamma$. For a two-cycle $e$, let $F_e=\{e\}$.

\begin{srcBNlemma}[Shared Round 5 word bound]\label{R7E4:lem:word}
For self-adjoint unitaries $V_1,\ldots,V_b$ and any vector $v$, writing $E_j=\langle v,(\srcBNId-V_j)v\rangle$, one has
\begin{equation}\label{R7E4:eq:word}
 |\langle v,(\srcBNId-V_1\cdots V_b)v\rangle|
 \le\left(\sum_j\sqrt{E_j}\right)^2\le b\sum_jE_j.
\end{equation}
Thus any ordered product of cycle operators, with total transposition-word length $C=\sum_j(\ell_{\gamma_j}-1)$, obeys
\begin{equation}\label{R7E4:eq:cycleword}
 |\langle v,(\srcBNId-U_{\gamma_1}\cdots U_{\gamma_k})v\rangle|
 \le C\sum_{j=1}^k\sum_{e\in F_{\gamma_j}}
       \langle v,(\srcBNId-L_e)v\rangle.
\end{equation}
\end{srcBNlemma}
\begin{proof}
This is \cite[Lemma 2.1]{R7E4:R504} and independently \cite[Lemma 2.1]{R7E4:R505}. Let $W_{j-1}=V_1\cdots V_{j-1}$. Telescope $\srcBNId-\prod V_j=\sum_jW_{j-1}(\srcBNId-V_j)$. Since $\|(\srcBNId-V_j)v\|=\sqrt{2E_j}$, the reverse-prefix telescope gives
$\|(W_{j-1}^*-\srcBNId)v\|\le\sum_{i<j}\sqrt{2E_i}$.
The absolute quadratic form of the $j$th summand is at most
$E_j+2\sum_{i<j}\sqrt{E_iE_j}$. Summation and Cauchy--Schwarz prove \eqref{R7E4:eq:word}, including the imaginary part. Concatenate the chosen transposition words to obtain \eqref{R7E4:eq:cycleword}. Repeated edges are counted with multiplicity.
\end{proof}

\begin{srcBNlemma}[Shared least-edge charge and local cycle mass]\label{R7E4:lem:local}
For a positive edge $e$ and $\ell\ge3$,
\begin{equation}\label{R7E4:eq:cut}
 \sum_{\substack{\ell_\gamma=\ell\\e\in F_\gamma}}
       \frac{\rho_\gamma}{w_e^2}\le2(\ell-1)q^{\ell-2}.
\end{equation}
For $b\ge0$, $r=qe^b<1$, and every vertex $i$,
\begin{equation}\label{R7E4:eq:localmass}
 \sum_{\gamma\ni i}\rho_\gamma e^{b\ell_\gamma}
 \le\frac{r^2}{1-r}.
\end{equation}
\end{srcBNlemma}
\begin{proof}
The first assertion is \cite[Lemma 3.3]{R7E4:R504} and independently \cite[Lemma 3.2]{R7E4:R505}. If $f\ne e$ is the deleted edge then $w_f\le w_e$, so
$\rho_\gamma/w_e^2\le\prod_{h\in\gamma\setminus\{e,f\}}w_h$.
Choose an orientation of $e$. Deleting $e,f$ leaves two endpoint paths of lengths $j$ and $\ell-2-j$. They determine the joining edge $f$ and the original directed cycle. Dropping simplicity and disjointness bounds their total weight by $q^jq^{\ell-2-j}$. Sum over the two orientations and $\ell-1$ choices of $j$. Tiny positive weights cause no problem.

For the second assertion, each directed length-$\ell$ cycle containing $i$ is a distinct closed walk starting there. With $W=(w_{ij})$, the total is at most $(W^\ell)_{ii}\le q^\ell$. This is also true at length two with its single-orientation convention. Sum the geometric series. This local mass estimate was already used by both Round 6 entries \cite{R7E4:R604,R7E4:R605}.
\end{proof}

\section{Ordered heaps, a derivative identity, and scalar majorants}
We separate three issues: a formal noncommutative identity, a scalar coefficient comparison, and absolute convergence. They must not be conflated.

\subsection{The formal identity retains the order}
For a finite polymer list with a symmetric reflexive incompatibility relation, a \emph{heap} is an equivalence class of finite words in the polymer letters, where adjacent compatible letters may be interchanged. Each occurrence is ordered before every later incompatible occurrence, with transitive closure; this poset determines the word class. Repeated occurrences of one letter are ordered. The minimal occurrences form a compatible set. A \emph{pyramid} here means a nonempty heap with exactly one minimal occurrence, called its root. Reversing all orders gives the equally common maximal-root convention.

Give letter $\gamma$ a positive integer degree $\ell_\gamma$, and let $D$ be the derivation multiplying a word by its total degree. For formal noncommuting letters $x_\gamma$, commuting only when compatible, put
\[
 M=\sum_{\mathcal F\text{ compatible}}(-1)^{|\mathcal F|}
                 \prod_{\gamma\in\mathcal F}x_\gamma.
\]
For a heap $h$ write $x(h)$ for its word class, $|h|$ for its number of occurrences, and $r(h)$ for its root when it is a pyramid.

\begin{srcBNlemma}[Cartier--Foata inversion and its derivative]\label{R7E4:lem:heap}
As formal series in the partially commutative algebra,
\begin{equation}\label{R7E4:eq:heapinverse}
 M^{-1}=\sum_{h\text{ heap}}x(h),
\end{equation}
\begin{equation}\label{R7E4:eq:heapderivative}
 -(DM)M^{-1}=\sum_{h\text{ pyramid}}\ell_{r(h)}x(h).
\end{equation}
In particular the derivative is on the left and the inverse on the right in \eqref{R7E4:eq:heapderivative}.
\end{srcBNlemma}
\begin{proof}
The inversion is classical \cite{R7E4:CF,R7E4:Kratt}. We give a coefficient proof that also supplies the required weighted derivative. For a fixed heap $h$, every decomposition into a compatible left factor and a remaining heap selects a subset $C$ of the minimal occurrences of $h$, and every such subset gives one decomposition. Thus its coefficient in $M\sum_hx(h)$ is
$\sum_{C\subseteq\min h}(-1)^{|C|}$, zero for every nonempty $h$, and one for the empty heap. The constant term of $M$ is one, so the formal inverse exists uniquely.

For the product $(DM)\sum_hx(h)$, the same coefficient is
\[
 \sum_{C\subseteq\min h}(-1)^{|C|}\sum_{\gamma\in C}\ell_\gamma
 =-\sum_{\gamma\in\min h}\ell_\gamma
            (1-1)^{|\min h|-1}.
\]
It is zero unless there is one minimal occurrence, in which case it is minus that occurrence's degree. This proves \eqref{R7E4:eq:heapderivative} without commuting incompatible letters or differentiating a noncommutative logarithm.
\end{proof}
Under $x_\gamma=-t^{\ell_\gamma}z_\gamma U_\gamma$, the derivation becomes $t\,d/dt$, and $M$ becomes $S(t)$. Therefore the formal derivative expansion is
\begin{equation}\label{R7E4:eq:operatorpyramids}
 tS'(t)S(t)^{-1}
 =\sum_{h\text{ pyramid}}(-1)^{|h|-1}\ell_{r(h)}
       t^{L(h)}z(h)U(h),
\end{equation}
where $L(h)=\sum_{\gamma\in h}\ell_\gamma$, $z(h)$ is the scalar product, and $U(h)$ is the ordered product. Compatible-letter commutation makes $U(h)$ well-defined. Substituting $U_\gamma=1$ gives the scalar analogue.

\subsection{Classical connected coefficients and an exact comparison}
For a scalar activity array $u=(u_\gamma)$ define
\[
 Z(u)=\sum_{\mathcal F\text{ compatible}}\prod_{\gamma\in\mathcal F}u_\gamma.
\]
For an ordered $k$-tuple $\boldsymbol\gamma$, repetitions allowed, let its incompatibility graph join incompatible pairs, and set
\[
 \phi(\boldsymbol\gamma)=
 \sum_{F\text{ connected spanning subgraph}}(-1)^{|F|},
 \qquad\phi(\gamma)=1.
\]
It vanishes for disconnected graphs. The following standard facts are reproduced to make the majorant auditable.

\begin{srcBNlemma}[Classical connected and rooted bounds]\label{R7E4:lem:cluster}
The formal identity is
\begin{equation}\label{R7E4:eq:scalarlog}
 \log Z(u)=\sum_{k\ge1}\frac1{k!}\sum_{\gamma_1,\ldots,\gamma_k}
                \phi(\boldsymbol\gamma)\prod_j u_{\gamma_j}.
\end{equation}
One has $(-1)^{k-1}\phi\ge0$. If nonnegative $u_\gamma$ and $a_\gamma$ satisfy
\begin{equation}\label{R7E4:eq:KP}
 \sum_{\delta\not\sim\gamma}u_\delta e^{a_\delta}\le a_\gamma,
\end{equation}
then
\begin{equation}\label{R7E4:eq:R}
 R_\gamma(u):=\sum_{k\ge1}\frac1{(k-1)!}
 \sum_{\gamma_2,\ldots,\gamma_k}|\phi(\gamma,\gamma_2,\ldots,\gamma_k)|
                \prod_{j=2}^ku_{\gamma_j}\le e^{a_\gamma}.
\end{equation}
The absolute connected series converges; the scalar heap series with weights $u$ converges and equals $1/Z(-u)$.
\end{srcBNlemma}
\begin{proof}
Expand the partition polynomial as
$\sum_{k\ge0}(1/k!)\sum_{\boldsymbol\gamma}\prod_ju_{\gamma_j}
\prod_{i<j}(1-\mathbf1_{\gamma_i\not\sim\gamma_j})$.
Repeated labels are excluded by reflexivity. Expansion into graphs and decomposition into connected components gives \eqref{R7E4:eq:scalarlog} coefficient by coefficient.

For a finite graph on $[k]$, root at $1$ and assign to each connected spanning subgraph its breadth-first tree with least-labeled preceding-level parents. The graphs assigned to a tree $T$ form the interval $T\subseteq F\subseteq R(T)$, where $R(T)$ additionally allows all same-level edges and consecutive-level edges whose preceding endpoint is not smaller than the prescribed parent. These edges preserve the distances and least-parent rule, proving the interval description. Its alternating sum is $(-1)^{k-1}(1-1)^{|R(T)\setminus T|}$. Hence the sign assertion holds and $|\phi|$ is at most the number of spanning trees.

Majorize \eqref{R7E4:eq:R} by its rooted labeled-tree series. Its height-$h$ generating functions satisfy
\[
 T_\gamma^{(0)}=1,\qquad T_\gamma^{(h+1)}=
 \exp\left(\sum_{\delta\not\sim\gamma}u_\delta T_\delta^{(h)}\right).
\]
The exponential enumerates unordered child branches; distributing occurrence labels yields precisely $(k-1)!$ in \eqref{R7E4:eq:R}, including repeated polymer types. Induction using \eqref{R7E4:eq:KP} gives $T_\gamma^{(h)}\le e^{a_\gamma}$, and increasing height includes every finite tree. This proves \eqref{R7E4:eq:R}. The absolute connected sum is bounded by $\sum_\gamma u_\gamma R_\gamma(u)$.

By the sign assertion, the formal series for $-\log Z(-u)$ has nonnegative coefficients and its total sum is finite. Its formal exponential therefore has finite total coefficient mass, equal to the exponential of that sum. Inversion \eqref{R7E4:eq:heapinverse} after scalar substitution identifies that exponential coefficientwise with the nonnegative heap series. Thus the latter converges and equals $1/Z(-u)$; the identity follows by absolutely convergent multiplication. This proves convergence without assuming zero-freeness or an operator logarithm.
\end{proof}
These are classical Penrose/Koteck\'y--Preiss estimates \cite{R6E4:KP,R6E4:FP}; their use in the preceding scalar permanent argument is jointly credited to \cite{R7E4:R604,R7E4:R605}.

Abelianize \eqref{R7E4:eq:heapderivative}, apply \eqref{R7E4:eq:scalarlog} to $Z(-u)$, and use the sign assertion. For every fixed multiset of polymer occurrences the coefficients in the following identity agree:
\begin{equation}\label{R7E4:eq:pyramidcluster}
 \sum_{h\text{ pyramid}}\ell_{r(h)}u(h)
 =\sum_{k\ge1}\frac1{k!}\sum_{\boldsymbol\gamma}
       |\phi(\boldsymbol\gamma)|
       \left(\sum_j\ell_{\gamma_j}\right)\prod_ju_{\gamma_j}.
\end{equation}
This equality is scalar and coefficientwise; no operator is reordered in using it. In particular any nonnegative factor depending only on the multiset can multiply both sides before summing. This observation is what permits the ordered word-energy bound to be charged by a scalar rooted series.

\subsection{The convergence margin used below}
Take
\begin{equation}\label{R7E4:eq:parameters}
 a=\frac1{64},\qquad b=\frac12,
 \qquad c=a+b=\frac{33}{64},\qquad
 r_t=tq e^c.
\end{equation}
For $q\le1/64$ one has $e^c<7/4$, and thus $r_t\le7/256$. The elementary estimate $e^x\le(1-x/m)^{-m}$ for $0\le x<m$ follows from $-\log(1-y)\ge y$. With $x=c,m=4$ it gives
\[
 e^c\le(256/223)^4<7/4.
\]
Consequently
\begin{equation}\label{R7E4:eq:massconstant}
 \frac{r_t^2}{1-r_t}\le\frac{49}{63744}<\frac1{64}=a.
\end{equation}
Apply Lemma~\ref{R7E4:lem:local} with the inflated activities
\[
 u_\gamma(t)=t^{\ell_\gamma}\rho_\gamma e^{b\ell_\gamma}.
\]
Summing its local mass over the vertices of a root polymer proves \eqref{R7E4:eq:KP} with $a_\gamma=a\ell_\gamma$. Hence
\begin{equation}\label{R7E4:eq:inflatedR}
 R_\gamma(u(t))\le e^{a\ell_\gamma},\qquad
 R_\gamma(u(t))-1\le e^{a\ell_\gamma}-1.
\end{equation}
The inflated scalar heap sum is finite. It dominates in operator norm the uninflated heap inverse in \eqref{R7E4:eq:heapinverse}, since $\|U(h)\|=1$. It also dominates the differentiated pyramid series, whose degree is at most a constant times $e^{bL(h)}$. Therefore the formal inverse and right derivative identities converge absolutely for $0\le t\le1$, uniformly by the $t=1$ bounds. Multiplying by the finite $S(t)$ proves that the inverse is its actual operator inverse. Thus \eqref{R7E4:eq:operatorpyramids} is an analytic identity here, not just a formal assertion.

\section{A dimension-free numerical-range estimate}
We now prove Theorem~\ref{R7E4:thm:derivative}. Subtract the scalar and operator versions of \eqref{R7E4:eq:operatorpyramids}:
\begin{equation}\label{R7E4:eq:Jheap}
 J(t)=\sum_{h\text{ pyramid}}(-1)^{|h|-1}\ell_{r(h)}
 t^{L(h)}z(h)(\srcBNId-U(h)).
\end{equation}
The single-occurrence two-cycles contribute exactly $2t^2H_A$. To control every other term, fix $v$, and write $E_e=\langle v,(\srcBNId-L_e)v\rangle\ge0$. For a multiset of cycles put
\[
 L=\sum_j\ell_{\gamma_j},\qquad C=\sum_j(\ell_{\gamma_j}-1),
 \qquad E(\boldsymbol\gamma)=\sum_j\sum_{e\in F_{\gamma_j}}E_e.
\]
Lemma~\ref{R7E4:lem:word} bounds each ordered product by $C E(\boldsymbol\gamma)$. Using the scalar multiset identity \eqref{R7E4:eq:pyramidcluster}, its absolute contribution is bounded by the connected sum with factor $LC E(\boldsymbol\gamma)$. Since $C\le L$ and
\begin{equation}\label{R7E4:eq:length}
 L^2\le4e^{L/2}\qquad(L\ge0),
\end{equation}
we can absorb the total length cost into the inflated activities. For \eqref{R7E4:eq:length}, maximize $x^2e^{-x}$ at $x=2$ and use $e>2$, with $x=L/2$.

Mark the occurrence supplying an edge in $E(\boldsymbol\gamma)$. Symmetry of ordered tuples replaces $1/k!$ by $1/(k-1)!$. The omitted singleton two-cycle requires its rooted term $1$ to be subtracted. Thus
\begin{align}\label{R7E4:eq:error}
 |\langle v,(J(t)-2t^2H_A)v\rangle|
 &\le4e^{2b}\sum_e t^2w_e^2E_e
                 (R_e(u(t))-1)\notag\\
 &\quad+4\sum_{\ell_\gamma\ge3}t^{\ell_\gamma}\rho_\gamma
       e^{b\ell_\gamma}R_\gamma(u(t))
                 \sum_{e\in F_\gamma}E_e.
\end{align}
In the first line $e^{2b}=e$; it comes from the marked two-cycle, not from distant components. All higher clusters with a two-cycle mark are retained, including repetitions. Long-cycle marks retain their singleton term. This is a bound on absolute quadratic forms, so the complex coefficients in \eqref{R7E4:eq:Jheap} have not been assigned a false sign.

By \eqref{R7E4:eq:inflatedR}, the first line is at most
$4e(e^{2a}-1)t^2\langle v,H_Av\rangle$.
Use the least-edge charge \eqref{R7E4:eq:cut} in the second line. Its coefficient at $t^2w_e^2E_e$ is at most
\[
 8\sum_{\ell\ge3}(\ell-1)e^{c\ell}(tq)^{\ell-2}
 =8e^{2c}\left(\frac1{(1-r_t)^2}-1\right).
\]
Consequently the error coefficient is bounded by
\begin{equation}\label{R7E4:eq:eta}
 \eta(tq)=4e(e^{2a}-1)+8e^{2c}
                      \left((1-tqe^c)^{-2}-1\right).
\end{equation}
It depends only on the local row bound, not $X$ or the order.

All constants can be checked rationally. Use $e<3$, $e^{1/32}\le32/31$, $e^c<7/4$, and $r_t\le7/256$. Then
\begin{equation}\label{R7E4:eq:constant}
 \eta(tq)\le\frac{12}{31}
 +\frac{49}{2}\left(\left(\frac{256}{249}\right)^2-1\right)
 =\frac{6857689}{3844062}<\frac95.
\end{equation}
Here $e<3$ follows immediately by bounding $1/k!\le2^{1-k}$ for $k\ge2$, strictly for some $k$. Equations \eqref{R7E4:eq:error}--\eqref{R7E4:eq:constant} prove \eqref{R7E4:eq:numerical}. Taking real parts gives \eqref{R7E4:eq:accretive}. This completes the derivative theorem.

\section{The positive Sylvester inverse and the spectrum}
\begin{proof}[Proof of Theorem~\ref{R7E4:thm:main} and Corollary~\ref{R7E4:cor:gap}]
For every $x\in\srcBNC^n$, the row bound gives
\[
 x^*A_tx\ge(1-tq)\|x\|^2>0.
\]
The Schur power is the principal submatrix of $A_t^{\otimes n}$ indexed by tuples of distinct indices, hence $S(t)\succ0$. Its constant vector has eigenvalue $P(t)$, so $P(t)>0$ and $T(t)\succ0$. Also $T(0)=\srcBNId$. These facts justify the inverses and differentiations as ordinary finite matrices.

The implication from accretivity to monotonicity requires a positive inverse, not multiplication of a Loewner inequality by $T$. Put $R=T(t)^{-1}\succ0$ and $Y=-T'(t)$. Equation \eqref{R7E4:eq:J} gives the exact Sylvester equation
\begin{equation}\label{R7E4:eq:sylvester}
 RY+YR=\frac2t\srcBNRea J(t).
\end{equation}
Its solution is
\begin{equation}\label{R7E4:eq:integral}
 Y=\frac2t\int_0^\infty e^{-sR}\srcBNRea J(t)e^{-sR}\,ds\succeq0.
\end{equation}
The integral converges because $R\succ0$. Differentiating $e^{-sR}(\srcBNRea J)e^{-sR}$ verifies \eqref{R7E4:eq:sylvester}; uniqueness follows in an eigenbasis of $R$, where each entry is multiplied by the positive number $r_i+r_j$. The integrand is PSD by \eqref{R7E4:eq:accretive}, proving $T'(t)\preceq0$ for $t>0$. Continuity at zero and integration prove \eqref{R7E4:eq:monotone}.

The supported transpositions generate $G=S_{V_1}\times\cdots\times S_{V_k}$, one symmetric group per connected support component. One direct verification along a path expresses the transposition of its endpoints by successive conjugations of its edge transpositions; thus every component permutation is generated. Since every summand of $H_A$ is PSD, its kernel consists exactly of $G$-invariant functions, constant on the orbits $G\sigma$. Their number is \eqref{R7E4:eq:multiplicity}. Every supported permutation in \eqref{R7E4:eq:regular} belongs to $G$, so $S(t)$ acts on this kernel by $P(t)$ and $T(t)$ by one. Hermiticity makes its orthogonal complement invariant as well.

If that complement is nonzero, $H_A\succeq\delta_A\srcBNId$ there. Let $\mu(t)$ be the largest eigenvalue of the restriction of $T(t)$ to this complement. The maximum of the differentiable Rayleigh quotients is locally Lipschitz, and its upper right derivative is the maximum of $\langle v,T'(t)v\rangle$ among unit maximizing eigenvectors. This elementary eigenvalue derivative formula follows by taking maximizing vectors at $t+h$ and a convergent subsequence as $h\downarrow0$; the reverse bound tests a maximizing vector at $t$. For each such vector, $T^{-1}v=\mu^{-1}v$, and \eqref{R7E4:eq:accretive} yields
\[
 -\frac{t}{\mu(t)}\langle v,T'(t)v\rangle
 =\langle v,\srcBNRea J(t)v\rangle\ge\frac15t^2\delta_A.
\]
Hence $D^+\log\mu(t)\le-\delta_At/5$. Integration, or the almost-everywhere derivative of this locally Lipschitz function, gives
$\mu(t)\le e^{-\delta_At^2/10}$, since $\mu(0)=1$.
This proves \eqref{R7E4:eq:gap}, strictness off the kernel, and the exact top multiplicity. If $A=\srcBNId_n$, the kernel is the whole space and $T(t)=\srcBNId$; no positive gap is asserted.
\end{proof}

\paragraph{Optional identification with the ordinary graph gap.}
For connected support of order at least two, the classical theorem of Caputo--Liggett--Richthammer \cite[Theorem 1.1]{R4E5:CLR} identifies $\delta_A$ with the second eigenvalue of the ordinary weighted Laplacian
$\mathcal L_A=\sum_{i<j}w_{ij}^2(e_i-e_j)(e_i-e_j)^*$.
Its edge-rate convention is exactly that of \eqref{R7E4:eq:H}. Thus the full Schur-power gap in Corollary~\ref{R7E4:cor:gap} is at least
$P(A)(1-e^{-\lambda_2(\mathcal L_A)/10})$.
This identification is an explicitly cited deep classical theorem, not part of our new proof and not needed for the operator assertions stated with $\delta_A$. For disconnected support, applying it on each nontrivial component and using the tensor-sum spectrum gives the minimum of their positive graph gaps.

\paragraph{Scalar consequences are not counted again.}
For a correlation $B$, Schur-product positivity gives
$\srcBNSch(A\circ B)=\srcBNSch(A)\circ\srcBNSch(B)\preceq P(A)\srcBNId$,
since $\srcBNSch(B)$ is PSD and has diagonal one. Its constant-vector eigenvalue yields $P(A\circ B)\le P(A)$. Normalization gives Bapat--Sunder for arbitrary PSD $B$, with zero diagonal entries handled by their zero rows. The permanent Hadamard inequality gives Chollet with this first factor. These standard POT implications, already available on the larger scalar row domain in Round 6, are not additional contributions. Positive diagonal congruence transfers the operator theorem to $A$ whose normalized entries $a_{ij}/\sqrt{a_{ii}a_{jj}}$ satisfy the row bound: both its Schur power and permanent scale by $\prod_i a_{ii}$.

\section{Extensive examples, failed routes, and scope}
Take the connected prism $C_m\square K_2$, $m\ge3$, of order $n=2m$ and degree three. Put magnitude $1/192$ on every edge, arbitrary conjugate complex phases there, and diagonal one. Then
\[
 q=1/64,\qquad X=n/24576.
\]
The full-POT theorem and radial monotonicity hold for every order, even as $X\to\infty$, on supports with overlapping cycles and arbitrary phase frustration. At $n=98304$, $X=4$: the necessary $t_*<1$ hypothesis of Round 5 entry 04's stated criterion fails, and the sufficient exponential charge of Round 5 entry 05 exceeds one. This demonstrates the removed hypothesis for a connected family, not just direct sums of small examples. Dense examples are obtained by starting at row sum at most $1/128$ and adding arbitrary Hermitian off-diagonal perturbations of magnitude at most $1/[128(n-1)]$.

\paragraph{What failed, and what the new identity changes.}
The scalar cluster logarithm cannot be substituted into a noncommutative algebra. Round 6 entry 05 gave the explicit negative coefficient $-1/30$ for $YX_1X_2X_3X_4$ in $-\log(\prod_{j=1}^4(1-X_j)-Y)$ when the $X_j$ commute with each other but not with $Y$ \cite[Section 6]{R7E4:R605}. We retain that obstruction. Our proof does not infer positivity of the operator logarithm from its coefficients. Instead \eqref{R7E4:eq:heapderivative} is a different identity, with root-degree weights and a fixed inverse order, justified directly by minimal-piece cancellation. The comparison to a scalar cluster sum occurs only after taking absolute word-energy bounds. Finally the Sylvester integral, not an unjustified rearrangement of an unsymmetrized product, gives Loewner order.

The new sufficient constant is not asserted optimal. We have not proved full POT for every diagonally dominant matrix, for the larger scalar threshold $1/12$, or for arbitrary PSD matrices. Universal POT is false, and this paper does not assert otherwise. Classical strict-diagonal-dominance zero-freeness of the permanent is already known \cite{R6E5:Barvinok}; zero-freeness alone gives neither our accretivity estimate nor normalized operator monotonicity.

\paragraph{Source and verification audit.}
Primary sources and latest records were rechecked on 20 October 2026: Tran v5 (23 August 2022), Pant--Singh v2 (18 September 2026), Pioge et al.\ v1 (31 July 2025), and Caputo--Liggett--Richthammer v4 (28 September 2009). Krattenthaler's author-hosted account explicitly distinguishes heap inversion in the partially commutative monoid from a logarithmic identity after abelianization. The Cartier--Foata and Koteck\'y--Preiss original publications and Fern\'andez--Procacci's treatment are established antecedents. Public through-order-six Chollet claims remain unaudited and are not dependencies. No literature search here proves historical priority.

The accompanying bounded verifier uses exact arithmetic for heap inverse and derivative coefficients, repeated-letter and multiset identities, rational constants, cycle charges, finite Schur-power positivity and radial derivatives. Numerical exploration, if any, is explicitly separated from certificates. The all-order assertions follow from the complete proof above, not finite enumeration. No external posting, submission, or author contact occurred. A next question is whether a sharper local derivative estimate reaches the full diagonal-dominance range; no such extension is claimed here.

\begingroup\small
\endgroup

\chapter[A dimension-independent Euclidean convexity region for the log-permanent]{A dimension-independent Euclidean convexity region for the log-permanent}
\label{source:R7E5}
\noindent\textit{Source development: Round 7, entry 05. Original mathematical arguments retained; notation is local to this chapter.}\par\medskip

\begin{quote}\small\noindent\textbf{Chapter abstract.} 
On the affine space of Hermitian matrices with fixed diagonal one, the
 logarithm of the permanent is uniformly strongly convex whenever every
 absolute off-diagonal row sum is at most $1/6$. Its Hessian lies between
 $1/20$ and $39/20$ times the Frobenius metric, independently of order,
 support, phases, and total squared off-diagonal mass. More precisely,
 the Hessian differs from that metric by at most $2q+O(q^2)$, and the
 leading coefficient $2$ is optimal uniformly in the order. We derive
 quantitative arithmetic Jensen inequalities and a random-phase-mask
 contraction theorem. A three-by-three strictly diagonally dominant
 example has negative log-permanent curvature, so the full such domain
 is not convexity-preserving. The new step beyond the jointly credited
 Round 6 scalar cluster theorem is a two-marked-cluster Hessian estimate,
 controlled by an ordinary vertex resolvent and closed-walk bounds.
 Classical cluster machinery is reproved and attributed. This is
 Euclidean, not geodesic, convexity; no full permanent-on-top assertion,
 unrestricted Bapat--Sunder extension, or historical priority is claimed.
\end{quote}

\section{A uniform Hessian theorem}
Let $\mathcal H_n^1=\{A=A^*\in\srcBOC^{n\times n}:a_{ii}=1\}$, viewed as
 a real affine space with tangent space
 $\mathcal H_n^0=\{H=H^*:h_{ii}=0\}$ and Frobenius norm
 $\|H\|_\srcBOF^2=\sum_{i,j}|h_{ij}|^2$. Put
\[
 q(A)=\max_i\sum_{j\ne i}|a_{ij}|,\qquad
 \mathcal D_n(q_0)=\{A\in\mathcal H_n^1:q(A)\le q_0\},\qquad
 f(A)=\log\srcBOper A.
\]
For $q(A)<1$, elementary diagonal dominance gives
 $A\succeq(1-q(A))\srcBOId$; its permanent is positive, so $f$ is real
 and smooth. The proof below also obtains positivity from a real
 cluster logarithm on the smaller domain used here.

\begin{srcBOtheorem}[Euclidean strong convexity, in every order]\label{R7E5:thm:main}
For every $A\in\mathcal D_n(1/6)$ and every $H\in\mathcal H_n^0$,
\begin{equation}\label{R7E5:eq:main}
 \frac1{20}\|H\|_\srcBOF^2
 \le D^2f(A)[H,H]\le\frac{39}{20}\|H\|_\srcBOF^2 .
\end{equation}
There is no total-mass restriction. If $q=q(A)\le1/6$, then more
 precisely
\begin{equation}\label{R7E5:eq:adaptive}
 \bigl|D^2f(A)[H,H]-\|H\|_\srcBOF^2\bigr|
 \le \mathcal K(q)\|H\|_\srcBOF^2,\qquad
 \mathcal K(q)=2q+O(q^2),
\end{equation}
where the explicit rational majorant $\mathcal K$ is given in
 \eqref{R7E5:eq:Krat}. It satisfies
 $\mathcal K(q)\le6q$ and $\mathcal K(1/6)<19/20$.
 For $q\le1/12$, the constants in \eqref{R7E5:eq:main} improve to
 $7/10,13/10$.

The coefficient $2$ in the small-$q$ statement is optimal uniformly
 in $n$: no coefficient $c<2$ can replace it in a bound
 $(c q+O(q^2))\|H\|_\srcBOF^2$ with a remainder uniform in the order.
\end{srcBOtheorem}
The row threshold and the coarse constants are sufficient, not claimed
 optimal. The asymptotic optimality concerns a different assertion,
 namely the leading uniform Hessian error. For $n=1$ the tangent
 space is zero and the statements are vacuous.

\begin{srcBOcorollary}[Arithmetic Jensen gap]\label{R7E5:cor:jensen}
Let $A_\nu\in\mathcal D_n(1/6)$, let $\mu_\nu>0$ sum to one, and set
 $\bar A=\sum_\nu\mu_\nu A_\nu$. For a finite family,
\begin{equation}\label{R7E5:eq:jensen}
 \frac1{40}\sum_\nu\mu_\nu\|A_\nu-\bar A\|_\srcBOF^2
 \le \sum_\nu\mu_\nu\log\srcBOper A_\nu-\log\srcBOper\bar A
 \le\frac{39}{40}\sum_\nu\mu_\nu\|A_\nu-\bar A\|_\srcBOF^2 .
\end{equation}
In particular
 $\srcBOper\bar A\le\prod_\nu(\srcBOper A_\nu)^{\mu_\nu}$, with equality
 exactly when all $A_\nu$ coincide. The same statements hold for
 probability distributions on this fixed-order compact domain.
\end{srcBOcorollary}
All diagonal entries are fixed. They cannot be allowed to vary
 arbitrarily: already on positive diagonal matrices $f(A)=\sum_i\log
 a_{ii}$ is concave in diagonal directions.

\paragraph{A sourced question and the precise increment.}
Bradley's public question \cite{R7E5:Bradley} asks for nontrivial matrix
 domains of convexity or concavity of the log-permanent. We address
 a quantitative, fixed-diagonal Euclidean version of that question.
 The existing posted discussion points to convexity after taking
 entrywise logarithms for positive entries, and to geodesic convexity.
 Neither is ordinary affine convexity of complex Hermitian entries.
 The latter distinction is substantive: the primary framework of
 Sra--Hosseini \cite[Theorem 8 and Lemma 24]{R7E5:SH} uses matrix geometric
 means. Applied to a positive tensor functional representing the
 permanent, it gives the known geodesic assertion, not
 \eqref{R7E5:eq:main}. We do not claim that a public discussion proves
 global novelty or that ours is the first response to that question.

Round 6 Entries 04 and 05 independently proved the same principal
 row-local logarithmic Hadamard contraction and order-$n$ cofactor
 comparison \cite{R7E5:R604,R7E5:R605}; their common core is jointly credited
 and counted once. Both used classical connected-cluster machinery.
 We reuse that scalar cycle representation and its rooted majorant,
 not rename their mask theorem. Here the perturbation $H$ is
 \emph{arbitrary}, may introduce previously zero edges, and need not
 be a correlation-mask variation. Two derivatives require two marked
 polymers, a path between them, and a Frobenius rather than a
 one-edge loss estimate. This produces a full Hessian theorem on a
 larger row ball, not a full Schur-power POT theorem. All six frozen
 Round 6 papers and their assessment were read.

\section{One and two marked classical clusters}
We give the classical finite-polymer facts needed for the proof,
 including the two-root estimate, to avoid a formal differentiation
 shortcut. The ingredients are tree majorants and connected expansions
 from Koteck\'y--Preiss and Fern\'andez--Procacci \cite{R6E4:KP,R6E4:FP};
 general cluster machinery is not a claimed invention.

Let $\mathcal P$ be a finite family of nonempty subsets of a finite
 vertex set, with labels; different labels may have identical subsets. Polymers are
 incompatible when their subsets intersect, including with themselves.
 Write $J_{\gamma\delta}$ for the incompatibility indicator. For
 activities $z_\gamma$ let
\[
 Z(z)=\sum_{\mathcal A\text{ pairwise compatible}}
             \prod_{\gamma\in\mathcal A}z_\gamma .
\]
For an ordered tuple $\boldsymbol\gamma=(\gamma_1,\ldots,\gamma_k)$,
 repetitions allowed, form its incompatibility graph on $[k]$ and put
\[
 \phi(\boldsymbol\gamma)
 =\sum_{G\text{ connected spanning subgraph}}(-1)^{|E(G)|},
 \qquad \phi(\gamma)=1.
\]
The sum is zero when the incompatibility graph is disconnected.
For $w_\gamma\ge|z_\gamma|$ and $a_\gamma\ge0$ assume
\begin{equation}\label{R7E5:eq:KP}
 \sum_\delta J_{\gamma\delta}w_\delta e^{a_\delta}\le a_\gamma.
\end{equation}
Define
\begin{align}
 Q_\gamma&=\sum_{k\ge1}\frac1{(k-1)!}
  \sum_{\gamma_2,\ldots,\gamma_k}
  |\phi(\gamma,\gamma_2,\ldots,\gamma_k)|
  \prod_{j=2}^kw_{\gamma_j},\label{R7E5:eq:Qone}\\
 Q_{\gamma\delta}^{(2)}
 &=\sum_{k\ge2}\frac1{(k-2)!}
  \sum_{\gamma_3,\ldots,\gamma_k}
  |\phi(\gamma,\delta,\gamma_3,\ldots,\gamma_k)|
  \prod_{j=3}^kw_{\gamma_j}.\label{R7E5:eq:Qtwo}
\end{align}
The two distinguished occurrences remain distinct even if
 $\gamma=\delta$ as polymer labels.

\begin{srcBOlemma}[Rooted and two-root bounds]\label{R7E5:lem:cluster}
Under \eqref{R7E5:eq:KP}, the connected series
\begin{equation}\label{R7E5:eq:cluster}
 C(z)=\sum_{k\ge1}\frac1{k!}\sum_{\gamma_1,\ldots,\gamma_k}
          \phi(\boldsymbol\gamma)\prod_jz_{\gamma_j}
\end{equation}
converges absolutely, $Z(z)=e^{C(z)}$, and $Q_\gamma\le e^{a_\gamma}$.
 Put $v_\gamma=w_\gamma e^{a_\gamma}$. Then
\begin{equation}\label{R7E5:eq:path}
 Q_{\gamma\delta}^{(2)}
 \le e^{a_\gamma+a_\delta}
 \sum_{r\ge0}\ \sum_{\eta_1,\ldots,\eta_r}
 J_{\gamma\eta_1}\cdots J_{\eta_r\delta}
 \prod_{j=1}^r v_{\eta_j},
\end{equation}
where the $r=0$ term is $J_{\gamma\delta}$.
 In particular, for nonnegative numbers $d_\gamma$, put
\begin{equation}\label{R7E5:eq:vertex}
 p_i=\sum_{\gamma\ni i}d_\gamma e^{a_\gamma},\qquad
 M_{ij}=\sum_{\gamma\ni i,j}v_\gamma .
\end{equation}
If $m=\max_i\sum_jM_{ij}<1$, then
\begin{equation}\label{R7E5:eq:twovertex}
 \sum_{\gamma,\delta}d_\gamma d_\delta Q_{\gamma\delta}^{(2)}
 \le p^{\mathsf T}(\srcBOId-M)^{-1}p
 \le\frac{\|p\|_2^2}{1-m}.
\end{equation}
\end{srcBOlemma}
\begin{proof}
The classical tree-graph inequality is
 $|\phi(G)|\le\tau(G)$, the number of spanning trees. One elementary
 proof uses the coefficient of $x$ in the chromatic polynomial:
 inclusion--exclusion gives $[x]\chi_G(x)=\phi(G)$.
 The signed coefficient $(-1)^{k-1}[x]\chi_G(x)$ is nonnegative
 by deletion--contraction and is at most the spanning-tree count
 by the same recurrence. Parallel edges created by contraction may
 be merged for the chromatic polynomial, which only decreases the
 corresponding tree count. Disconnected graphs contribute zero;
 the single-vertex case contributes one.

The generating function for labeled trees rooted at polymer $\gamma$,
 with no root weight, and with height at most $h$, satisfies
\[
 T_\gamma^{(0)}=1,\qquad
 T_\gamma^{(h+1)}
   =\exp\left(\sum_\delta J_{\gamma\delta}w_\delta
                         T_\delta^{(h)}\right).
\]
Indeed deleting the root leaves an unordered collection of branches;
 the exponential's branch factorial and the partitions of the remaining
 labels give the factor $1/(k-1)!$. Induction from \eqref{R7E5:eq:KP}
 gives $T_\gamma^{(h)}\le e^{a_\gamma}$, and monotone convergence
 gives the same bound for the unrestricted rooted-tree series.
 This proves the bound on $Q_\gamma$ and absolute convergence of
 \eqref{R7E5:eq:cluster}. Expanding
 $\prod_{i<j}(1-J_{\gamma_i\gamma_j})$ in the ordered-tuple
 representation of $Z$ and splitting graphs into connected components
 proves the formal identity $Z=e^C$. Absolute convergence under
 $z\mapsto\lambda z$, $|\lambda|\le1$, proves the analytic identity
 first near zero, then inside the disk, and at its boundary by continuity.

For \eqref{R7E5:eq:path}, replace $|\phi|$ by spanning trees and take the
 unique path between distinguished vertices 1 and 2. If it has $r$
 internal vertices, they are ordered from root 1 to root 2.
 Deleting path edges leaves a rooted branch-tree at each of these
 $r+2$ vertices. Its bound is $e^{a_\eta}$; internal path vertices
 also carry their activity weights, while the distinguished roots do
 not. There is \emph{no} extra $1/r!$: choosing and ordering the $r$
 internal labels contributes $(k-2)!/(k-2-r)!$, canceling the appropriate
 part of the denominator in \eqref{R7E5:eq:Qtwo}; distributing all remaining
 labels among the rooted branch-trees supplies their exponential
 factorials. This gives exactly \eqref{R7E5:eq:path}. Repeated polymer
 labels remain permitted in every path and branch.

Finally bound each incompatibility indicator by the number of shared
 vertices,
 $J_{\gamma\delta}\le\sum_i\mathbf1_{i\in\gamma}
 \mathbf1_{i\in\delta}$.
 After multiplying by $d_\gamma d_\delta$, the paths with $r$
 intermediate polymers sum to at most $p^{\mathsf T}M^r p$.
 The matrix $M$ is symmetric and nonnegative, with operator norm
 at most its maximum row sum $m$. Summing its Neumann series proves
 \eqref{R7E5:eq:twovertex}. This proves convergence of the two-root sums
 used here as well.
\end{proof}

\section{Cycle derivatives and a Frobenius estimate}
We now use \emph{all} directed simple cycles on $[n]$ of length at
 least two, modulo rotation, as the finite polymer family. Reverse
 cycles are distinct at lengths at least three; transpositions occur
 once. Zero-weight cycles are retained, since a perturbation may
 create edges. Put
\[
 R=(r_{ij}),\quad r_{ij}=|a_{ij}|\ (i\ne j),\quad r_{ii}=0,
 \qquad T=(|h_{ij}|),
\]
\[
 z_\gamma=\prod_{i\in V(\gamma)}a_{\gamma(i),i},
 \quad w_\gamma=|z_\gamma|,
 \quad d_\gamma=|D z_\gamma(A)[H]|,\quad
 b_\gamma=|D^2 z_\gamma(A)[H,H]|.
\]
Thus $R,T$ are real symmetric nonnegative matrices, $\|R\|_{\rm op}\le q$,
 and $\|T\|_\srcBOF=\|H\|_\srcBOF$. By permutation cycle decomposition,
 $Z(z)=\srcBOper A$.

\begin{srcBOlemma}[Local and differentiated walk bounds]\label{R7E5:lem:walks}
Let $a\ge0$ and $t=qe^a<1$. With $a_\gamma=a\ell_\gamma$,
 \eqref{R7E5:eq:KP} holds if $t^2/(1-t)\le a$. Moreover
\begin{align}
 \max_i\sum_jM_{ij}
 &\le m(t):=\frac{t^2(2-t)}{(1-t)^2},\label{R7E5:eq:m}\\
 \left\|\left(\sum_{\gamma\ni i}
                   d_\gamma e^{a\ell_\gamma}\right)_i\right\|_2
 &\le e^{2a}q\,\frac{2-t}{(1-t)^2}\|H\|_\srcBOF,\label{R7E5:eq:p}\\
 \sum_{\ell_\gamma=\ell} b_\gamma
 &\le(\ell-1)q^{\ell-2}\|H\|_\srcBOF^2\quad(\ell\ge3).
 \label{R7E5:eq:b}
\end{align}
For transpositions, $\sum_e b_e=\|H\|_\srcBOF^2$ exactly.
\end{srcBOlemma}
\begin{proof}
Every cycle of length $\ell$ through $i$ determines a distinct closed
 walk of that length starting at $i$. Hence
\[
 \sum_{\gamma\ni i,\ \ell_\gamma=\ell}w_\gamma
 \le(R^\ell)_{ii}\le q^\ell .
\]
Summing $e^{a\ell}q^\ell$ and then over vertices of a fixed polymer
 proves the stated KP condition. Multiplying the same bound by $\ell$
 proves \eqref{R7E5:eq:m}, since
 $\sum_{\ell\ge2}\ell t^\ell=t^2(2-t)/(1-t)^2$.

Differentiate a cycle once and take absolute values. The resulting
 products have one $T$ edge and all other edges from $R$.
 For a fixed starting vertex $i$ they are bounded by the corresponding
 differentiated closed walks:
\[
 \sum_{\gamma\ni i,\ \ell_\gamma=\ell}d_\gamma
 \le\sum_{j=0}^{\ell-1}(R^jTR^{\ell-1-j})_{ii}.
\]
The Euclidean norm of the diagonal of a matrix is at most its
 Frobenius norm, and
 $\|R^jTR^{\ell-1-j}\|_\srcBOF\le q^{\ell-1}\|H\|_\srcBOF$.
 The triangle inequality therefore gives the bound
 $\ell q^{\ell-1}\|H\|_\srcBOF$ for the left-hand vector.
 Sum over $\ell\ge2$ with the exponential weight to obtain
 \eqref{R7E5:eq:p}.

For two derivatives, a cycle has two ordered distinct differentiated
 edges. Bounding by closed walks and dividing by their $\ell$
 starting positions gives
\[
 \sum_{\ell_\gamma=\ell} b_\gamma
 \le\sum_{j=0}^{\ell-2}\srcBOtr(TR^jTR^{\ell-2-j})
 \le(\ell-1)q^{\ell-2}\|H\|_\srcBOF^2.
\]
The last step is the Frobenius Cauchy--Schwarz inequality and
 $\|R\|_{\rm op}\le q$; all walk weights are nonnegative.
 At $\ell=2$, a transposition activity is $|a_{ij}|^2$ and its
 second derivative is $2|h_{ij}|^2$, proving the exact assertion.
 No quotient by an edge weight is used anywhere, including at zero
 entries of $A$.
\end{proof}

\begin{srcBOproposition}[General Hessian comparison]\label{R7E5:prop:hessian}
Suppose $q>0$, $a>0$, $t=qe^a<1$,
\begin{equation}\label{R7E5:eq:conditions}
 \frac{t^2}{1-t}<a,\qquad m(t)<1.
\end{equation}
Then
\begin{equation}\label{R7E5:eq:general}
 |D^2f(A)[H,H]-\|H\|_\srcBOF^2|
 \le K(q,a)\|H\|_\srcBOF^2,
\end{equation}
where
\begin{equation}\label{R7E5:eq:K}
 K(q,a)=\frac{e^{2a}}{(1-t)^2}-1+
 \frac{e^{2a}t^2(2-t)^2}{(1-t)^4(1-m(t))}.
\end{equation}
\end{srcBOproposition}
\begin{proof}
The strict first condition provides a neighborhood of $A$ in which
 the connected series converges uniformly, with room to differentiate
 twice in any fixed direction $H$. More explicitly,
 $q(A+uH)\le q+|u|\max_i\sum_j|h_{ij}|$ even for complex $u$;
 the same fixed $a$ satisfies a strict KP bound for sufficiently small
 $|u|$. The all-cycle polynomial activities are holomorphic in $u$.
 Uniform convergence on a slightly larger circle justifies
 termwise differentiation by Cauchy's formula.

For real $u$, reversing every cycle conjugates the activities and
 preserves incompatibility. The convergent cluster sum is real, and
 its exponential is the permanent; it is the real logarithm used in
 $f$. This also proves positivity in this neighborhood.

Differentiate \eqref{R7E5:eq:cluster} twice. The terms with two derivatives
 on the same occurrence have one marked root and denominator
 $(k-1)!$. The terms with derivatives on two distinct occurrences
 have two ordered marked roots and denominator $(k-2)!$.
 Singleton transpositions give exactly $\|H\|_\srcBOF^2$.
 After subtracting them, the absolute remainder is at most
\begin{equation}\label{R7E5:eq:budget}
 \sum_e b_e(Q_e-1)
 +\sum_{\gamma:\ell_\gamma\ge3} b_\gamma Q_\gamma
 +\sum_{\gamma,\delta}d_\gamma d_\delta Q_{\gamma\delta}^{(2)}.
\end{equation}
This accounts for repeated labels, for two different differentiated
 occurrences with the same label, and for activity-zero roots.

The first two sums are at most
\[
 \left\{e^{2a}-1+
       \sum_{\ell\ge3}(\ell-1)e^{a\ell}q^{\ell-2}\right\}\|H\|_\srcBOF^2
 =\left\{\frac{e^{2a}}{(1-t)^2}-1\right\}\|H\|_\srcBOF^2.
\]
For the third, apply \eqref{R7E5:eq:twovertex}, \eqref{R7E5:eq:m} and
 \eqref{R7E5:eq:p}. Its bound is
\[
 \frac{e^{4a}q^2(2-t)^2}{(1-t)^4(1-m(t))}\|H\|_\srcBOF^2.
\]
Since $e^{4a}q^2=e^{2a}t^2$, this gives \eqref{R7E5:eq:K}.
Only local closed-walk row sums and a vertex-space resolvent appear;
 no total squared mass or minimum nonzero edge size enters.
\end{proof}

\section{Rational constants and the sharp first-order coefficient}
For $0<q\le1/6$, take $a=2q^2$ and define
\[
 u=\frac{q}{1-2q^2},\quad E=\frac1{1-4q^2},\quad
 M(u)=\frac{u^2(2-u)}{(1-u)^2},
\]
\begin{equation}\label{R7E5:eq:Krat}
 \mathcal K(q)=\frac{E}{(1-u)^2}-1+
       \frac{E u^2(2-u)^2}{(1-u)^4(1-M(u))}.
\end{equation}
The bound $e^x\le(1-x)^{-1}$, $0\le x<1$, gives $t\le u$,
 $e^{2a}\le E$, and $K(q,2q^2)\le\mathcal K(q)$.
Convergence follows uniformly, with strict room:
\[
 \frac{u^2}{(1-u)q^2}
 =\frac1{(1-2q^2)(1-q-2q^2)}
 \le\frac{162}{119}<2,\qquad
 M(u)\le M(3/17)=\frac{279}{3332}<1.
\]
At the two stated endpoints, exact substitution gives
\begin{equation}\label{R7E5:eq:constants}
 \mathcal K(1/6)=\frac{319789}{341936}<\frac{19}{20},
 \qquad
 \mathcal K(1/12)=\frac{178721191}{671304725}<\frac3{10}.
\end{equation}
The series of $\mathcal K$ has nonnegative coefficients and zero
 constant term. To see this without subtractive ambiguity, write
 $(2-u)/(1-u)^2=\sum_{j\ge0}(j+2)u^j$ and
 $M(u)=u^2\sum_{j\ge0}(j+2)u^j$; every denominator and composition
 in \eqref{R7E5:eq:Krat} then has nonnegative coefficients, while the
 displayed $-1$ removes only the constant term.
Thus $\mathcal K$ and $\mathcal K/q$ increase on this interval, and
 $6\mathcal K(1/6)<6$ proves $\mathcal K(q)\le6q$.
Expansion at zero yields
\begin{equation}\label{R7E5:eq:seriesK}
 \mathcal K(q)=2q+11q^2+O(q^3).
\end{equation}
At $q=0$, direct expansion gives
 $\srcBOper(\srcBOId+uH)=1+\tfrac12u^2\|H\|_\srcBOF^2+O(u^3)$, so the Hessian
 is exactly the Frobenius metric. Proposition~\ref{R7E5:prop:hessian}
 and these computations prove all positive bounds in
 Theorem~\ref{R7E5:thm:main}.

\begin{proof}[Sharpness of the leading coefficient]
For fixed $n\ge3$ and small $r>0$, let
 $A_n(r)=\srcBOId-r(J-\srcBOId)$ and $H_n=-(J-\srcBOId)$.
 Then $q=(n-1)r$, $A_n(r)\succ0$ for $q<1$, and
 $\|H_n\|_\srcBOF^2=n(n-1)$.
The only permutation contributions of degrees two and three are
 transpositions and oriented triangles, respectively, so
\[
 \srcBOper A_n(r)=1+\binom n2r^2-2\binom n3r^3+O(r^4).
\]
Consequently
\begin{equation}\label{R7E5:eq:sharp}
 \frac{D^2f(A_n(r))[H_n,H_n]}{\|H_n\|_\srcBOF^2}
 =1-2(n-2)r+O(r^2)
 =1-\frac{2(n-2)}{n-1}q+O(q^2).
\end{equation}
Given $c<2$, first choose a fixed $n$ with $2(n-2)/(n-1)>c$,
 then let $q\downarrow0$. This contradicts any proposed
 order-uniform estimate with leading coefficient $c$ and a uniform
 $O(q^2)$ remainder. No exchange of the two limits is assumed.
\end{proof}

\section{Consequences and a sharp distinction of domains}
\begin{proof}[Proof of Corollary~\ref{R7E5:cor:jensen}]
The set $\mathcal D_n(1/6)$ is convex. Integrating the Hessian bound
 along the segment from $B$ to $A$ gives
\[
 \frac1{40}\|A-B\|_\srcBOF^2
 \le f(A)-f(B)-Df(B)[A-B]
 \le\frac{39}{40}\|A-B\|_\srcBOF^2.
\]
Use $B=\bar A$, multiply by $\mu_\nu$, and sum; the linear terms
 cancel. The lower bound gives the equality characterization.
For distributions, the domain is compact at each fixed $n$,
 $f$ and its derivatives are continuous there, and all relevant
 functions are bounded. Integrating the same Taylor bounds proves
 the result directly, including countable and continuous distributions.
\end{proof}
The same calculation with \eqref{R7E5:eq:adaptive} gives the factors
 $(1-\mathcal K(q_0))/2$ and $(1+\mathcal K(q_0))/2$ on any
 $\mathcal D_n(q_0)$, $q_0\le1/6$.
Also the projected gradient in $\mathcal H_n^0$ is strongly monotone:
\[
 \langle\nabla_0 f(A)-\nabla_0 f(B),A-B\rangle_\srcBOF
 \ge\frac1{20}\|A-B\|_\srcBOF^2,
 \quad
 \|\nabla_0 f(A)-\nabla_0 f(B)\|_\srcBOF
 \le\frac{39}{20}\|A-B\|_\srcBOF .
\]
These are standard consequences of a Hessian comparison, not
 independent new mechanisms or efficient permanent algorithms.

\begin{srcBOcorollary}[Random-phase masks on the larger ball]\label{R7E5:cor:mask}
Suppose $q(A)\le1/6$. Let random variables $\zeta_i$ have modulus one
 and put $b_{ij}=\mathbb E(\overline{\zeta_i}\zeta_j)$.
Then, with $E_A(B)=\sum_{i<j}|a_{ij}|^2(1-|b_{ij}|^2)$,
\begin{equation}\label{R7E5:eq:mask}
 \frac1{20}E_A(B)\le
 \log\frac{\srcBOper A}{\srcBOper(A\circ B)}
 \le\frac{39}{20}E_A(B).
\end{equation}
The constants $1\pm\mathcal K(q(A))$ may replace those displayed.
\end{srcBOcorollary}
\begin{proof}
Let $D=\operatorname{diag}(\zeta_i)$. Each $D^*AD$ has the same
 row bound and the same permanent as $A$, and
 $\mathbb E(D^*AD)=A\circ B$. Its Frobenius variance is
\[
 \mathbb E\|D^*AD-A\circ B\|_\srcBOF^2
 =\sum_{i\ne j}|a_{ij}|^2(1-|b_{ij}|^2)=2E_A(B).
\]
Apply the distribution version of \eqref{R7E5:eq:jensen}.
\end{proof}
The mask class here is the convex hull of rank-one phase correlation
 matrices, not all correlation matrices. We do not infer an
 unrestricted-mask Bapat--Sunder theorem at $q=1/6$.
For comparison, the jointly attributed Round 6 theorem applies to
 \emph{every} correlation mask at $q\le1/12$. Neither conclusion
 implies full order-$n!$ Schur-power POT.

Fixed positive diagonals cause no problem. If $a_{ii}=d_i>0$ are
 fixed, put $D=\operatorname{diag}(d_i)$ and
 $\widehat A=D^{-1/2}AD^{-1/2}$.
 Its log-permanent differs by the constant $\sum_i\log d_i$.
 Thus the Hessian and Jensen claims hold with the weighted metric
 $\sum_{i,j}|h_{ij}|^2/(d_i d_j)$ when $q(\widehat A)\le1/6$.

\paragraph{Extensive and fully complex examples.}
On the Cartesian torus $C_m\square C_m$, $m\ge4$ even, put modulus $1/24$
 on every edge, diagonal one, and arbitrary conjugate phases.
 Then $q=1/6$ and
 $\sum_{i<j}|a_{ij}|^2=m^2/288\to\infty$.
 The Hessian theorem holds in \emph{all} Hermitian zero-diagonal
 directions, not only directions preserving that sparse support.
 A phase $i$ on one edge and $1$ elsewhere gives a nonbalanced example. The theorem
 also gives full-dimensional neighborhoods of any smaller-row-sum
 center: a uniform entrywise perturbation of size
 $\varepsilon/(n-1)$ increases the row bound by at most $\varepsilon$.
 These are illustrations of the domain, not extra discoveries.

\begin{srcBOproposition}[Strict diagonal dominance does not suffice]\label{R7E5:prop:failure}
The log-permanent is not Euclidean convex on all strictly diagonally
 dominant correlation matrices, already in order three.
\end{srcBOproposition}
\begin{proof}
For $0<r<1/2$ set $A(r)=\srcBOId_3-r(J_3-\srcBOId_3)$.
 Its eigenvalues are $1-2r,1+r,1+r$, and its off-diagonal row sum
 is $2r<1$. Directly
\[
 P(r)=\srcBOper A(r)=1+3r^2-2r^3,\qquad
 (\log P)''(r)=\frac{(6-12r)P(r)-(6r-6r^2)^2}{P(r)^2}.
\]
At $r=2/5$ these values are
\begin{equation}\label{R7E5:eq:failure}
 P(r)=\frac{169}{125},\qquad
 (6-12r)P(r)-(6r-6r^2)^2=-\frac{282}{625},
 \qquad
 (\log P)''(r)=-\frac{7050}{28561}<0.
\end{equation}
A sufficiently short segment around this positive definite point
 violates convexity. This rules out the full $q<1$ domain, not
 thresholds below $4/5$ or the geodesic convexity statement.
\end{proof}

\section{Scope, failed route, sources, and reproducibility}
The initial target this round was removal of the global-mass
 restriction for full Schur-power POT. That attempt did not produce
 a justified operator-valued connected expansion. The negative
 noncommutative logarithm coefficient already recorded in
 Round 6/05 remains an obstruction to a positivity shortcut.
 No scalar cluster argument here is substituted into noncommuting
 permutation operators. The proved pivot is the ordinary scalar
 Hessian, with arbitrary matrix directions and two marked clusters.

The old shared scalar mask result, its cofactor implication, the
 classical tree-graph bound and Koteck\'y--Preiss recursion retain
 their original credit. The proposed new content is the all-direction
 dimension-independent Hessian domain, the two-root/walk estimate
 yielding it, the optimal first-order coefficient, and the explicit
 boundary obstruction. Jensen, gradient and random-phase consequences
 are identified as consequences rather than separate discoveries.
No optimal row threshold, global convexity, unrestricted Bapat--Sunder,
 full POT, or resolution of general Chollet is asserted.

The public convexity question and current primary-source records were
 checked on 20 October 2026. Fern\'andez--Procacci's latest preprint
 is v2 (16 October 2006); Sra--Hosseini's is v3 (12 December 2014);
 Barvinok's diagonal-dominance approximation paper is v2
 (11 September 2018). The last already supplies a much larger
 nonvanishing/approximation context \cite{R6E5:Barvinok}, not the
 Euclidean Hessian theorem. Searches did not establish historical
 priority. A failed fetch of Watkins's 1974 primary PDF is not used
 to assert a new theorem or a complete literature audit.

The accompanying verifier uses exact Gaussian-rational arithmetic.
 It checks full real-coordinate Hessian matrices, two marked
 connected-series coefficients with repeated labels, rooted-tree
 path factorials, local walk inequalities, rational constants,
 arithmetic Jensen gaps with rational logarithm intervals, and the
 negative-curvature certificate. Tests allowing a perturbation to
 create zero edges are included. The all-order proofs above do not
 depend on extrapolation from finite checks. No external posting,
 submission, or author contact occurred, and internal checking is
 not external peer review.

\begingroup\small
\endgroup

\part{Symmetric lattice points and simplex bundles}

\chapter[Sharp Ewald-Point Counts in Sequential Simplex-Bundle Towers]{Sharp Ewald-Point Counts in Sequential Simplex-Bundle Towers}
\label{source:R1E6}
\noindent\textit{Source development: Round 1, entry 06. Original mathematical arguments retained; notation is local to this chapter.}\par\medskip


\begin{quote}\small\noindent\textbf{Chapter abstract.} 
\textbf{Status: a proved restricted-family optimization, not a solution of the global open problem.}
Crespo, Pelayo and Santos ask for the smallest possible number of symmetric lattice points in a smooth reflexive polytope of each dimension. Their simplex-bundle construction gives an exponential upper bound with rate $9^{1/3}$. We prove that their dimension-by-dimension counts are optimal among all sequentially marked simplex-bundle towers starting from a segment, allowing arbitrary positive fiber dimensions and every nonnegative integral twist compatible with this construction. The lower bound follows from three explicit rational invariant polyhedra for a family of $2\times2$ transfer matrices. A finite certificate and an elementary trinomial-coefficient estimate make the argument independent of computational searches. The construction and its counting recurrence are prior work; the proposed contribution is the sharp lower bound for the full restricted class. Its novelty has not been independently established.
\end{quote}

\section{The sourced problem and the precise scope}
For a full-dimensional lattice polytope $P\subset\srcAFR^D$ containing $0$ in its interior, put
\[
 \srcAFE(P)=P\cap(-P)\cap\srcAFZ^D.
\]
A polytope is \emph{smooth} if the primitive facet normals at every vertex form a lattice basis. It is \emph{reflexive} if each facet is at lattice distance one from the origin. We call a smooth reflexive polytope \emph{monotone}.

Crespo--Pelayo--Santos \cite[Question 6.3]{R1E6:CPS} ask for the minimum of $|\srcAFE(P)|$ over monotone $D$-polytopes, and for its minimizers. The question appears in the March 2024 preprint and the 2026 published article. Their Definition 6.4 and Lemmas 6.5--6.6 give small simplex bundles and a counting recurrence; Corollary 6.7 gives the upper bounds reproduced below. They also report exact global minima through dimension seven. We do not prove a new global bound, nor claim that a literature search certifies present-day openness.

The following specifies the class optimized here; in particular, the choice of marked facet is part of the restriction.
\begin{srcAFdefinition}\label{R1E6:def:class}
Start with $B_0=[-1,1]$ and its marked facet $x=1$, represented by $\ell_0(x)=x$. Given a marked monotone polytope $(B,\ell)$, where $\ell$ is an integral linear form and $\ell\leq1$ defines the marked facet, choose integers $n\geq1$ and $0\leq q\leq n$. Set
\begin{equation}\label{R1E6:eq:bundle}
 P=\left\{(x,y)\in B\times\srcAFR^n:
 y_j\geq-1\ (1\leq j\leq n),\quad
 \sum_{j=1}^n y_j+q\ell(x)\leq1\right\}.
\end{equation}
Mark the \emph{new} facet with defining form
$\ell'(x,y)=\sum_jy_j+q\ell(x)$.
Let $\mathcal T_D$ be the $D$-dimensional polytopes obtainable by finitely many such operations, including zero operations when $D=1$.
\end{srcAFdefinition}
The case $q=n$ is precisely the small-bundle operation of \cite[Definition 6.4]{R1E6:CPS}. Smaller $q$ are included here. The class does not allow an arbitrary old facet to be re-marked, arbitrary twists involving several facet forms, or arbitrary starting polytopes.

\begin{srcAFtheorem}\label{R1E6:thm:main}
Every member of $\mathcal T_D$ is monotone. Its minimum Ewald-set size is
\begin{equation}\label{R1E6:eq:min}
 \min_{P\in\mathcal T_D}|\srcAFE(P)|=
 \begin{cases}
 7,&D=2,\\
 3\cdot9^k,&D=3k+1,\quad k\geq0,\\
 13\cdot9^k,&D=3k+3,\quad k\geq0,\\
 59\cdot9^{k-1},&D=3k+2,\quad k\geq1.
 \end{cases}
\end{equation}
Consequently the $D$th root of this restricted minimum tends to $9^{1/3}$.
\end{srcAFtheorem}
The upper bounds in this theorem are the existing construction, not new examples. The lower bounds show that varying the fiber dimensions, their order, or the permitted twist magnitudes cannot improve those counts within Definition~\ref{R1E6:def:class}.

\section{Geometry and exact counting}
\begin{srcAFlemma}\label{R1E6:lem:geometry}
Operation \eqref{R1E6:eq:bundle} takes a marked monotone polytope to another marked monotone polytope.
\end{srcAFlemma}
\begin{proof}
After setting $z_j=y_j+1$, the fiber is a simplex
$z_j\geq0$, $\sum_jz_j\leq n+1-q\ell(x)$.
Its last right-hand side is at least $n+1-q\geq1$ on $B$.
Thus its face structure never degenerates. Each vertex of $P$ lies over a vertex $v$ of $B$ and has either all $y_j=-1$, or all but one equal to $-1$ and the remaining coordinate equal to $n-q\ell(v)$. These are integral. To justify that this list is exhaustive, at any point over a positive-dimensional minimal face of $B$, the active fiber equations can be maintained along a sufficiently short segment in that face, so the point cannot be a vertex.

At each listed vertex, the active normal matrix is block triangular. Its base block is unimodular by smoothness of $B$. Its fiber block consists either of $-I_n$, or of $n-1$ rows of $-I_n$ together with the all-ones row, and has determinant $\pm1$. All inequalities in \eqref{R1E6:eq:bundle}, as well as the base inequalities, define facets and have primitive normals. They have right-hand side one, and the origin satisfies them strictly. This proves smoothness and reflexivity.
\end{proof}
For a marked polytope define
\[
 a=|\{x\in\srcAFE(B):\ell(x)=1\}|,\qquad
 b=|\{x\in\srcAFE(B):\ell(x)=0\}|.
\]
The only possible values of $\ell$ on $\srcAFE(B)$ are $-1,0,1$, and negation interchanges the outer two slices. Hence $|\srcAFE(B)|=2a+b$.

For integers $n\geq0$ and $j\in\srcAFZ$ let
\[
 c_n(j)=[z^j](z^{-1}+1+z)^n,\qquad T_n=c_n(0),\qquad A_n=c_n(1).
\]
Coefficients outside $[-n,n]$ are zero. The sequence $c_n(j)$ is symmetric and nonincreasing for $j\geq0$. This follows by induction from convolution with $(1,1,1)$; for example, for $j\geq1$ the difference of adjacent coefficients after convolution is $c_n(j-1)-c_n(j+2)\geq0$, and the difference at $j=0$ is $c_n(1)-c_n(2)\geq0$.

\begin{srcAFlemma}[Transfer matrices]\label{R1E6:lem:transfer}
The state vector $v=(a,b)^t$ changes under \eqref{R1E6:eq:bundle} by
\begin{equation}\label{R1E6:eq:generalM}
 v'=M_{n,q}v,\qquad
 M_{n,q}=\begin{pmatrix}
 c_n(q-1)+c_n(q+1)&A_n\\
 2c_n(q)&T_n
 \end{pmatrix}.
\end{equation}
Moreover $M_{n,q}\geq M_n$ entrywise, where
\begin{equation}\label{R1E6:eq:Mn}
 M_n:=M_{n,n}=\begin{pmatrix}n&A_n\\2&T_n\end{pmatrix}.
\end{equation}
\end{srcAFlemma}
\begin{proof}
An Ewald point over $x\in\srcAFE(B)$ has each $y_j\in\{-1,0,1\}$ and satisfies
$|\sum_jy_j+q\ell(x)|\leq1$.
If $\ell(x)=0$, the numbers on the new marked and central slices are $A_n$ and $T_n$. Combining the two old outer slices gives the first column of \eqref{R1E6:eq:generalM}.

For $1\leq q\leq n$, unimodality gives $c_n(q-1)\geq c_n(n-1)=n$, and $c_n(q)\geq c_n(n)=1$. For $q=0$, use $A_n\geq n$, by counting words with one $+1$ and all other letters zero, and $T_n\geq1$. The last column is unchanged. This proves the entrywise inequality. Formula \eqref{R1E6:eq:Mn} is also \cite[Lemma 6.5]{R1E6:CPS}, expressed in this notation.
\end{proof}
All matrices are nonnegative. Therefore replacing every $q$ in a tower by the corresponding $n$ can only decrease both final state coordinates: for nonnegative matrices, entrywise order is preserved under products. It suffices to minimize
\begin{equation}\label{R1E6:eq:word}
 (2,1)M_{n_s}\cdots M_{n_1}\binom11,
 \qquad n_i\geq1,\quad \sum_i n_i=D-1.
\end{equation}

\section{Two uniform estimates for the transfer matrices}
\begin{srcAFlemma}\label{R1E6:lem:cone}
Every $M_n$ preserves the cone
$K=\{(a,b)^t:0\leq a\leq b\}$.
\end{srcAFlemma}
\begin{proof}
Put $d_n=T_n-A_n$. We first prove
\begin{equation}\label{R1E6:eq:centralgap}
 d_n\geq n-2\quad(n\geq2).
\end{equation}
Expanding two additional letters gives
$d_{n+2}=T_n-c_n(3)$, so
$d_{n+2}-d_n=A_n-c_n(3)$.
The reflection principle for walks with steps $-1,0,1$ identifies $A_n-c_n(3)$ with the number of length-$n$ walks from $0$ to $1$ that never go below zero. There are at least $n$ such walks, obtained by choosing the position of their single $+1$ step. For completeness, the usual reflection bijection reflects the portion of a bad walk through its first visit to $-1$, identifying bad walks with unrestricted walks from $-2$ to $1$, counted by $c_n(3)$.
Starting with $d_2=d_3=1$, induction in steps of two proves \eqref{R1E6:eq:centralgap}.

If $M_n(a,b)^t=(a',b')^t$, then
$b'-a'=(2-n)a+d_nb$.
For $n\geq2$, \eqref{R1E6:eq:centralgap} and $a\leq b$ make this nonnegative. For $n=1$ it equals $a$. Nonnegativity of both coordinates is automatic.
\end{proof}

\begin{srcAFlemma}\label{R1E6:lem:growth}
For $n\geq2$ and $v\in K$,
\begin{equation}\label{R1E6:eq:growth}
 M_{n+3}v\geq9M_nv
\end{equation}
coordinatewise.
\end{srcAFlemma}
\begin{proof}
We establish
\begin{equation}\label{R1E6:eq:coeffgrowth}
 T_{n+3}-9T_n\geq16,\qquad
 A_{n+3}-9A_n\geq8n-3\quad(n\geq2).
\end{equation}
The coefficient recursion and unimodality give
\[
 3A_n-2T_n=T_{n-1}-A_{n-1}+3c_{n-1}(2)\geq0\quad(n\geq1).
\]
Consequently $T_{n+1}=T_n+2A_n\geq\frac73T_n$ for $n\geq1$.
Also, for $m\geq1$, direct convolution gives
\[
 T_m-A_m-c_m(2)=-2c_{m-1}(2)-c_{m-1}(3)\leq0.
\]
It follows, for $n\geq2$, that
\begin{align*}
 2c_n(2)-A_n
 &=A_{n-1}+c_{n-1}(2)+2c_{n-1}(3)-T_{n-1}\geq0,\\
 A_{n+1}&=T_n+A_n+c_n(2)\geq\tfrac52 A_n.
\end{align*}
Here the last inequality also uses $T_n\geq A_n$.
Since $A_3=6$, these inequalities imply $A_n\geq2n$ for $n\geq3$.

For $n\geq3$ we obtain
\[
 T_{n+3}-9T_n\geq\tfrac{100}{27}T_n\geq\tfrac{700}{27}>16,
 \qquad
 A_{n+3}-9A_n\geq\tfrac{53}{8}A_n\geq\tfrac{53}{4}n>8n-3.
\]
For $n=2$, the first difference is $51-9\cdot3=24$, while the same estimate for the second is at least $53/4>13$.
This proves \eqref{R1E6:eq:coeffgrowth}.

Finally the coordinates of $(M_{n+3}-9M_n)(a,b)^t$ are
\[
 (3-8n)a+(A_{n+3}-9A_n)b,
 \qquad -16a+(T_{n+3}-9T_n)b.
\]
They are nonnegative by $0\leq a\leq b$ and \eqref{R1E6:eq:coeffgrowth}.
\end{proof}

\section{The rational invariant-polyhedra certificate}
All upward closures below mean addition of the nonnegative quadrant. Define
\begin{align*}
 S_0&=\operatorname{conv}\{(1,1)\}+\srcAFR_{\geq0}^2,\\
 S_1&=\operatorname{conv}\{(2,23/9),(20/9,7/3),(83/36,83/36)\}
       +\srcAFR_{\geq0}^2,\\
 S_2&=\operatorname{conv}\{(4,5),(19/4,19/4)\}+\srcAFR_{\geq0}^2.
\end{align*}
Their equivalent inequality descriptions are
\begin{align}
 S_0:&\quad a\geq1,\quad b\geq1;\label{R1E6:eq:S0}\\
 S_1:&\quad a\geq2,\quad b\geq83/36,\quad a+b\geq41/9,
                  \quad a+3b\geq83/9;\label{R1E6:eq:S1}\\
 S_2:&\quad a\geq4,\quad b\geq19/4,\quad a+3b\geq19.\label{R1E6:eq:S2}
\end{align}
These descriptions can also be verified by intersecting adjacent boundary lines.

\begin{srcAFproposition}\label{R1E6:prop:invariant}
For $r\in\{0,1,2\}$, $n\geq1$, and $v\in S_r\cap K$,
\begin{equation}\label{R1E6:eq:invariant}
 M_nv\in9^{\lfloor(r+n)/3\rfloor}S_{(r+n)\bmod3}.
\end{equation}
\end{srcAFproposition}
\begin{proof}
First take $n\in\{1,2,3,4\}$. The matrices are
\[
 M_1=\begin{pmatrix}1&1\\2&1\end{pmatrix},\quad
 M_2=\begin{pmatrix}2&2\\2&3\end{pmatrix},\quad
 M_3=\begin{pmatrix}3&6\\2&7\end{pmatrix},\quad
 M_4=\begin{pmatrix}4&16\\2&19\end{pmatrix}.
\]
All entries are nonnegative, so checking \eqref{R1E6:eq:invariant} at the listed vertices of $S_r$ suffices for the entire upward closure. Table~\ref{R1E6:tab:certificate} prints a complete rational certificate. To define its entries precisely, apply $M_n/9^{\lfloor(r+n)/3\rfloor}$ to each listed vertex of $S_r$, evaluate the left-hand side minus right-hand side of each target inequality in \eqref{R1E6:eq:S0}--\eqref{R1E6:eq:S2} in the printed order, and take the minimum for each inequality separately. Each tuple in the table consists of these minima. All are nonnegative.

\begin{table}[ht]
\centering
\begin{tabular}{cccl}
\toprule
$n$&$r$&Target remainder&Minimum slack tuple\\
\midrule
1&0&1&$(0,25/36,4/9,16/9)$\\
1&1&2&$(5/9,65/36,47/9)$\\
1&2&0&$(0,4/9)$\\
2&0&2&$(0,1/4,0)$\\
2&1&0&$(1/81,22/81)$\\
2&2&1&$(0,1/4,0,4/9)$\\
3&0&0&$(0,0)$\\
3&1&1&$(8/27,0,4/81,0)$\\
3&2&2&$(2/3,0,0)$\\
4&0&1&$(2/9,1/36,0,0)$\\
4&1&2&$(91/81,17/27,733/324)$\\
4&2&0&$(14/81,25/108)$\\
\bottomrule
\end{tabular}
\caption{Exact finite certificate for Proposition~\ref{R1E6:prop:invariant}.}\label{R1E6:tab:certificate}
\end{table}

For $n\geq5$, apply Lemma~\ref{R1E6:lem:growth} with $n-3\geq2$ and induct on $n$. It gives $M_nv\geq9M_{n-3}v$. The target remainder is unchanged, its required power of nine increases by one, and each $S_r$ is upward closed. This proves \eqref{R1E6:eq:invariant} for every $n$.
\end{proof}

\begin{proof}[Proof of Theorem~\ref{R1E6:thm:main}]
Smoothness and reflexivity follow from Lemma~\ref{R1E6:lem:geometry}. By Lemma~\ref{R1E6:lem:transfer}, it suffices for the lower bound to use maximal twists. Their initial state $(1,1)^t$ belongs to $S_0\cap K$. Lemma~\ref{R1E6:lem:cone} and Proposition~\ref{R1E6:prop:invariant}, applied successively along any word of total fiber dimension $d=D-1$, show that its final state lies in
\[
 9^{\lfloor d/3\rfloor}S_{d\bmod3}\cap K.
\]
The minima of $2a+b$ on $S_0,S_1,S_2$ are respectively
\[
 3,\qquad 59/9,\qquad 13.
\]
For $S_1$ this also follows directly by adding $a\geq2$ to $a+b\geq41/9$; the other cases follow from the displayed vertices. These give the claimed lower bounds except in dimension two. There the only possible fiber dimension word is $(1)$, and $q=1$ gives seven points while $q=0$ gives nine.

For attainment, observe
\[
 M_3\binom11=9\binom11,\qquad
 M_2\binom11=\binom45,\qquad
 M_2^2\binom11=\binom{18}{23}.
\]
Use $k$ fibers of dimension three for $D=3k+1$; use $k$ fibers of dimension three followed by one of dimension two for $D=3k+3$; use $k-1$ fibers of dimension three followed by two of dimension two for $D=3k+2$, $k\geq1$. All twists are maximal. These are the constructions underlying \cite[Lemma 6.6 and Corollary 6.7]{R1E6:CPS}.
The limit follows at once from the three residue-class formulas.
\end{proof}

\section{Computational audit, failed avenue, and remaining gap}
\paragraph{Independent exact checks.}
The accompanying \texttt{verify.py} uses only the Python standard library and rational/integer arithmetic. It checks every entry of Table~\ref{R1E6:tab:certificate}; verifies the coefficient inequalities for $2\leq n\leq100$ as a diagnostic rather than a proof; checks the transfer formula by direct enumeration for all towers of total added dimension at most five and all permitted twists; and computes the exact Pareto state sets through dimension 100. In that dynamic program a state is discarded only when another state of the same dimension is no larger in either coordinate, a valid pruning rule because every later matrix has nonnegative entries. All minima agree with \eqref{R1E6:eq:min}. The proof above, not the finite search, establishes the theorem in arbitrary dimension.

\paragraph{The unsuccessful improvement attempt.}
The initial approach was to mix fiber dimensions and twists to beat the published upper bound. Exact state searches found no improvement. A scalar estimate for the one-step total is too weak to recover the sharp residue-$2$ count: it loses the correlation between the outer and central slices. Keeping both state coordinates led instead to the invariant polyhedra and a sharp obstruction. Thus this note documents a proved limit of that construction method rather than a new global extremal construction.

\paragraph{What remains open here.}
We do not show that the right-hand side of \eqref{R1E6:eq:min} equals the minimum among all monotone polytopes in any new dimension. Nor do we optimize arbitrary repeated facet choices even within small bundles. Re-marking an old facet requires joint slice information not present in $(a,b)$, so the proof cannot simply be iterated with such choices. General linear twists also fall outside the entrywise comparison. These are concrete next avenues: develop a joint-slice transfer system for branching facet choices, or find a base polytope outside the tower class whose slice statistics beat the three invariant regions after dimension normalization.

\paragraph{Novelty and evaluation.}
The inspected primary source provides the constructions and their transfer recurrence but not this restricted-family lower bound. Targeted searches for Ewald-point minimum, small fiber bundle, and trinomial formulations found no matching optimization theorem. That is limited evidence, not a priority guarantee. The result is a modest, self-contained candidate contribution in discrete geometry; it is not an Annals-level claim, a proof of Ewald's conjecture, or a resolution of Question 6.3. The principal audit points are the scope of Definition~\ref{R1E6:def:class}, the transfer matrix, and the printed invariant-region certificate.



\chapter[Branching Simplex Bundles and Improved Upper Bounds for Ewald-Point Minima]{Branching Simplex Bundles and Improved Upper Bounds for Ewald-Point Minima}
\label{source:R2E6}
\noindent\textit{Source development: Round 2, entry 06. Original mathematical arguments retained; notation is local to this chapter.}\par\medskip


\begin{quote}\small\noindent\textbf{Chapter abstract.} 
\textbf{Status: an explicit construction improving a cited global upper bound, not a determination of the minimum.}
Let $\srcALemin(D)$ denote the minimum number of symmetric lattice points in a smooth reflexive $D$-polytope. Crespo, Pelayo and Santos give an upper-bound construction with exponential rate $9^{1/3}$. We construct smooth reflexive polytopes proving
\[
 \lim_{D\to\infty}\srcALemin(D)^{1/D}
 \leq\bigl(19195+4\sqrt{2297138}\bigr)^{1/14}
 =2.0627969707\ldots<9^{1/3}.
\]
We also exhibit an eight-dimensional example with $507$ symmetric lattice points, improving their displayed upper bound $531$. The construction branches at distinct simplex facets and repeatedly re-marks an old facet. A two-state weighted coefficient formula yields a repeatable fourteen-dimensional operation with an explicit positive integral transfer matrix. All geometric and counting arguments are proved below; finite enumeration is an independent audit. These bounds improve the specified 2026 publication, but no exhaustive current-record or historical-priority claim is made. Ewald's conjecture and the exact global minimum remain unresolved here.
\end{quote}

\section{Target, provenance, and main result}
For a full-dimensional lattice polytope $P\subset\srcALR^D$ with $0$ in its interior, set
\[
 \srcALE(P)=P\cap(-P)\cap\srcALZ^D.
\]
A lattice polytope is \emph{monotone} if it is smooth and reflexive: at each vertex its primitive facet normals form a lattice basis, and every facet has a primitive integral inequality with right-hand side one. Put
\[
 \srcALemin(D)=\min\{|\srcALE(P)|:P\text{ is a monotone }D\text{-polytope}\}.
\]
This is a minimum of a nonempty set of positive integers, because the cube is an example.

Question 6.3 of Crespo--Pelayo--Santos \cite{R1E6:CPS} asks for $\srcALemin(D)$ and its minimizers. Their small-simplex-bundle construction gives $\lim_D\srcALemin(D)^{1/D}\leq9^{1/3}$, and their Table 5 gives $\srcALemin(8)\leq531$. These statements remain in the article published on 9 April 2026, inspected on 20 October 2026. We improve those particular bounds; we do not certify the absence of every later unpublished or published improvement.

\begin{srcALtheorem}\label{R2E6:thm:main}
There is a monotone $8$-polytope $P$ with $|\srcALE(P)|=507$. There are monotone $(14k+1)$-polytopes $P_k$, $k\geq0$, with
\begin{equation}\label{R2E6:eq:recurrence}
 e_k:=|\srcALE(P_k)|,\qquad
 e_0=3,\quad e_1=96463,\quad
 e_{k+2}=38390e_{k+1}-331693817e_k.
\end{equation}
Writing $\rho=19195+4\sqrt{2297138}$, these imply
\begin{equation}\label{R2E6:eq:rate}
 \lim_{D\to\infty}\srcALemin(D)^{1/D}\leq\rho^{1/14}<2.063<9^{1/3}.
\end{equation}
More explicitly, if $D=14k+1+r$ with $0\leq r\leq13$, then
$\srcALemin(D)\leq3^r e_k$.
\end{srcALtheorem}

\paragraph{Delta from Round 1.}
Entry 06 of Round 1 \cite{R2E6:R1} optimized sequential towers that always retain the newest sum facet as the next attachment site. It proved sharpness of the cited construction \emph{within that restricted class}. Here old facets may be selected and different new facets may receive independent branches. This produces a genuinely smaller upper bound for the unrestricted minimum, not a stronger lower bound for the old class. The present proof reuses the elementary bundle geometry and slice-counting viewpoint, with their provenance acknowledged, but does not assume the Round 1 optimization theorem. No result of another Round 1 entrant is a mathematical dependency.

\section{A geometric attachment and its facet states}
Let $B\subset\srcALR^d$ be monotone, and let $\ell(x)\leq1$ be any facet inequality with primitive integral $\ell$. For an integer $n\geq1$ define
\begin{equation}\label{R2E6:eq:attach}
 S_n(B,\ell)=\left\{(x,y):x\in B,\ y_j\geq-1\ (1\leq j\leq n),\
 \sum_{j=1}^n y_j+n\ell(x)\leq1\right\}.
\end{equation}
This is the small-bundle construction of \cite[Definition 6.4]{R1E6:CPS}; we will choose its attachment facets differently.

\begin{srcALlemma}\label{R2E6:lem:geometry}
The polytope \eqref{R2E6:eq:attach} is monotone. Every old facet remains a facet. Its new primitive facet forms are
\begin{equation}\label{R2E6:eq:states}
 h_0=\sum_j y_j+n\ell(x),\qquad h_j=-y_j\quad(1\leq j\leq n).
\end{equation}
Consequently any of these facets can serve as an attachment site at any later stage.
\end{srcALlemma}
\begin{proof}
Writing $z_j=y_j+1$, the fiber over $x$ becomes
$z_j\geq0$, $\sum_j z_j\leq n+1-n\ell(x)$.
Its last right-hand side is at least one everywhere on $B$. Thus it never degenerates, and all old and new facet inequalities define facets.
Its vertices lie over vertices $v$ of $B$, with either every $y_j=-1$, or all but one equal to $-1$ and the remaining one equal to $n-n\ell(v)$. They are integral. There are no additional vertices: over a positive-dimensional minimal base face, the active fiber equations can be maintained along a sufficiently short nontrivial segment in that face.
At every listed vertex, the active normal matrix is block triangular. The base block is unimodular. The fiber block is either $-I_n$, or $n-1$ of its rows together with the all-ones row, and has determinant $\pm1$. The new normals are primitive because they contain a coefficient $1$ or $-1$ in a new coordinate. Every inequality has right-hand side one, and the origin is strictly interior. This proves the assertions.
\end{proof}

For a symmetric lattice point, all facet heights in \eqref{R2E6:eq:states} lie in $\{-1,0,1\}$, and
\begin{equation}\label{R2E6:eq:conservation}
 h_0+h_1+\cdots+h_n=n\ell(x).
\end{equation}
Conversely, any such assignment satisfying \eqref{R2E6:eq:conservation} corresponds to exactly one fiber lattice point, by $y_j=-h_j$. This elementary conservation equation is useful when several new facets receive branches.

A branch uses new variables and depends on the preceding polytope only through its attachment height $t\in\{-1,0,1\}$. Suppose its numbers of completions are $u$ for $t=1$ and $t=-1$, and $v$ for $t=0$. Negation of every variable explains equality of the two outer counts. Call $(u,v)$ its \emph{conditional counts}; an empty branch has $(1,1)$.

\begin{srcALproposition}[Weighted facet composition]\label{R2E6:prop:composition}
Attach an $n$-simplex, and place disjoint branches of conditional counts $(u_i,v_i)$ on its $n+1$ new facets. The combined branch has conditional counts
\begin{equation}\label{R2E6:eq:composition}
 u=[z^n]\prod_{i=0}^n\bigl(v_i+u_i(z+z^{-1})\bigr),\qquad
 v=[z^0]\prod_{i=0}^n\bigl(v_i+u_i(z+z^{-1})\bigr).
\end{equation}
\end{srcALproposition}
\begin{proof}
For an attachment height $t$, equation \eqref{R2E6:eq:conservation} requires the sum of the new heights to be $nt$. Assigning a height $h_i$ contributes weight $v_i$ if it is zero, and $u_i$ if it is $\pm1$. The branch variables are disjoint, so these weights multiply. Coefficient extraction counts exactly the permitted assignments. Symmetry equates the coefficients of $z^n$ and $z^{-n}$.
\end{proof}
With no further branches, the conditional counts of an $n$-simplex are $n+1$ and $[z^0](z^{-1}+1+z)^{n+1}$. In particular,
\begin{equation}\label{R2E6:eq:leaves}
 n=2:\ (u,v)=(3,7),\qquad n=3:\ (u,v)=(4,19).
\end{equation}

\section{An eight-dimensional improvement}
Use coordinates $(x,y_1,y_2,y_3,z_1,z_2,w_1,w_2)$ and impose
\begin{align}
 -1\leq x\leq1,&\qquad y_i,z_j,w_j\geq-1,\notag\\
 y_1+y_2+y_3+3x&\leq1,\label{R2E6:eq:P8}\\
 z_1+z_2-2y_1&\leq1,\notag\\
 w_1+w_2-2y_2&\leq1.\notag
\end{align}
The first extension is a $3$-simplex. The next two are $2$-simplices attached at the distinct primitive facets $-y_1\leq1$ and $-y_2\leq1$. Lemma~\ref{R2E6:lem:geometry} proves smoothness and reflexivity, without a floating-point vertex test.

Put $Q(z)=z^{-1}+1+z$. By \eqref{R2E6:eq:leaves} and Proposition~\ref{R2E6:prop:composition}, the conditional counts over $x=\pm1$ and $x=0$ are
\[
 u=[z^3]Q(z)^2(3z^{-1}+7+3z)^2=60,\qquad
 v=[z^0]Q(z)^2(3z^{-1}+7+3z)^2=387.
\]
A full coefficient certificate is
\begin{equation}\label{R2E6:eq:P8poly}
 Q(z)^2(3z^{-1}+7+3z)^2
 =387+320(z+z^{-1})+178(z^2+z^{-2})
   +60(z^3+z^{-3})+9(z^4+z^{-4}).
\end{equation}
Hence $|\srcALE(P)|=2u+v=507$, proving the first assertion of Theorem~\ref{R2E6:thm:main}. This is an upper bound on $\srcALemin(8)$, not a proof that $507$ is its value.

\section{Re-marking as a repeatable transfer operation}
For a marked base facet $\ell\leq1$, write
\[
 a=|\{x\in\srcALE(B):\ell(x)=1\}|,\qquad
 b=|\{x\in\srcALE(B):\ell(x)=0\}|.
\]
The negative slice also has size $a$, so $|\srcALE(B)|=2a+b$.

\begin{srcALlemma}[A two-simplex bridge]\label{R2E6:lem:bridge}
Attach a $2$-simplex at $\ell$, place a branch with conditional counts $(u,v)$ on its new sum facet $h_0=y_1+y_2+2\ell$, and then mark the other facet $h_1=-y_1$. The new marked state satisfies
\begin{equation}\label{R2E6:eq:bridge}
 \binom{a'}{b'}=
 K(u,v)\binom ab,
 \qquad
 K(u,v)=\begin{pmatrix}u+v&u+v\\2u&2u+v\end{pmatrix}.
\end{equation}
\end{srcALlemma}
\begin{proof}
The unmarked remaining facet $h_2=-y_2$ has weight polynomial $Q(z)$ and the branch at $h_0$ has polynomial $v+u(z+z^{-1})$. Their product has coefficients
\[
 c(0)=2u+v,\qquad c(1)=c(-1)=u+v,\qquad c(2)=c(-2)=u.
\]
For a fixed old height $t$ and new marked height $s$, the other heights sum to $2t-s$. Thus an old positive point contributes $u+v$ points to the positive marked slice and $u$ to the central slice, and an old negative point contributes $u$ to the central slice but none to the positive slice. An old central point contributes $u+v$ and $2u+v$ to the positive and central slices, respectively. These are exactly \eqref{R2E6:eq:bridge}.
\end{proof}
This lemma also identifies why the Round 1 lower-bound mechanism does not extend unchanged. Taking $(u,v)=(3,7)$ and initial state $(1,1)$ gives
\[
 K(3,7)\binom11=\binom{20}{19}.
\]
Thus there is a monotone $5$-polytope with a facet whose positive Ewald slice is larger than its central slice. The invariant cone $a\leq b$ used for the sequential class fails after re-marking; the earlier theorem is not contradicted.

As a particularly small exponential improvement, the leaf branch $(u,v)=(4,19)$ gives
\[
 K(4,19)=\begin{pmatrix}23&23\\8&27\end{pmatrix},
 \qquad\lambda_{\max}=25+2\sqrt{47}.
\]
It adds five dimensions. Repetition already yields rate
$(25+2\sqrt{47})^{1/5}=2.0776269582\ldots<9^{1/3}$.
The next section gives a stronger construction using the same principle.

\section{The fourteen-dimensional operation and its rate}
At the bridge's facet $h_0$, attach a $3$-simplex with coordinates $z_1,z_2,z_3$. At each of its three coordinate facets $-z_j\leq1$, attach one further $3$-simplex. Do not attach anything to its sum facet. This branch adds twelve dimensions. By Proposition~\ref{R2E6:prop:composition} its conditional counts are
\begin{equation}\label{R2E6:eq:uv}
 u=[z^3]Q(z)(4z^{-1}+19+4z)^3=976,\qquad
 v=[z^0]Q(z)(4z^{-1}+19+4z)^3=17731.
\end{equation}
For independent verification, the full polynomial is
\begin{align}\label{R2E6:eq:branchpoly}
 Q(z)(4z^{-1}+19+4z)^3
 ={}&17731+14119(z+z^{-1})+5500(z^2+z^{-2})\notag\\
 &+976(z^3+z^{-3})+64(z^4+z^{-4}).
\end{align}
Together with the two bridge coordinates, the operation adds fourteen dimensions. Re-marking $-y_1$ gives the fixed matrix
\begin{equation}\label{R2E6:eq:K}
 K=\begin{pmatrix}18707&18707\\1952&19683\end{pmatrix}.
\end{equation}

For clarity, the complete list of \emph{new} inequalities at a base facet $\ell\leq1$ is
\begin{align}\label{R2E6:eq:explicitgadget}
 y_1,y_2,z_j,w_{j,l}&\geq-1
 &&(1\leq j,l\leq3),\notag\\
 h_0:=y_1+y_2+2\ell&\leq1,\notag\\
 z_1+z_2+z_3+3h_0&\leq1,\\
 w_{j,1}+w_{j,2}+w_{j,3}-3z_j&\leq1
 &&(1\leq j\leq3).\notag
\end{align}
All base inequalities are retained. Every step has a strictly positive simplex fiber, and the output facet $-y_1\leq1$ remains an actual primitive facet by Lemma~\ref{R2E6:lem:geometry}. Thus the operation can be repeated indefinitely, not just applied once.

\begin{proof}[Proof of the remaining assertions of Theorem~\ref{R2E6:thm:main}]
Starting with $[-1,1]$ and its facet $x=1$, repeat \eqref{R2E6:eq:explicitgadget} $k$ times, always using the stated output facet for the next input. The resulting polytope $P_k$ is monotone and has dimension $14k+1$, with
\[
 e_k=(2,1)K^k\binom11.
\]
Direct calculation gives
\[
 \operatorname{tr}K=38390,\qquad \det K=331693817,
 \qquad (2,1)K\binom11=96463.
\]
The Cayley--Hamilton identity proves \eqref{R2E6:eq:recurrence}.
The two eigenvalues are
\[
 \rho=19195+4\sqrt{2297138},\qquad
 \sigma=19195-4\sqrt{2297138},\qquad 0<\sigma<\rho.
\]
In fact
\[
 e_k=A\rho^k+(3-A)\sigma^k,
 \qquad A=\frac{96463-3\sigma}{\rho-\sigma}>0,
\]
because $\sigma<19195$ and $96463>3\cdot19195$. Therefore
$e_k^{1/(14k+1)}\to\rho^{1/14}$.

Cartesian products preserve smoothness and reflexivity, and their Ewald sets are Cartesian products. Multiplying $P_k$ by $r$ copies of $[-1,1]$ gives the bound $\srcALemin(14k+1+r)\leq3^r e_k$ for $0\leq r\leq13$. Hence the upper limit of the dimension roots is at most $\rho^{1/14}$. The limit exists: $\srcALemin(m+n)\leq\srcALemin(m)\srcALemin(n)$, so $\log\srcALemin(D)$ is subadditive. The elementary subadditivity lemma gives convergence of $D^{-1}\log\srcALemin(D)$ to its infimum. Alternatively, the established upper-limit bound suffices for the improvement claimed here.

Finally $16\cdot2297138=36754208<6063^2=36759969$, so $\rho<25258$. The exact integer inequality
\[
 25258^3=16113759473512<22876792454961=9^{14}
\]
proves $\rho^{1/14}<9^{1/3}$. The stronger displayed bound $\rho^{1/14}<2.063$ follows from the exact rational inequality $25258\cdot1000^{14}<2063^{14}$. Also $(2063/1000)^3<9$, proving the full chain in \eqref{R2E6:eq:rate}.
\end{proof}

\section{Audit, comparison, and limitations}
The new family need not improve the earlier bound in every small dimension. For example, $e_1=96463$ in dimension $15$ exceeds its earlier upper bound $85293$. At dimension $57$, however,
\[
 e_4=1864620809452020987
 <1951230258860988573,
\]
where the right side is the value from \cite[Corollary 6.7]{R1E6:CPS}. The strictly smaller exponential rate guarantees an eventual improvement after filling dimension residues by bounded products. Separately, \eqref{R2E6:eq:P8} improves dimension eight immediately.

\paragraph{Reproducibility.}
The standard-library program \texttt{verify.py} expands \eqref{R2E6:eq:P8poly} and \eqref{R2E6:eq:branchpoly} exactly, checks all transfer entries, the characteristic polynomial, and every integer comparison used above. It independently enumerates symmetric lattice points from the explicit geometric inequalities in dimensions $8$ and $15$, by generating new coordinates in $\{-1,0,1\}$ and retaining exactly those satisfying each new absolute facet inequality. It counts every pair of old and output facet heights for the fourteen-dimensional operation, independently of the weighted-polynomial formula. Its output is recorded in \texttt{verification.txt}. These tests audit finite identities; the geometric induction and spectral argument prove the infinite family.

\paragraph{What failed and what changed.}
The attempted extension of the Round 1 invariant cone failed at the explicit state $(20,19)$. Rather than overlooking this obstruction, we used it to construct lower Ewald counts. A rooted-branch search using exact conditional counts found the small pattern in \eqref{R2E6:eq:uv}; numerical spectral radii were used only to select candidates, never as proof. The final theorem needs neither exhaustive search nor a claim that this branch is optimal.

\paragraph{Remaining questions and novelty.}
We do not determine $\srcALemin(8)$ or its minimizers, prove a matching exponential lower bound, optimize all branches, or settle Ewald's basis conjecture. The method may admit stronger finite branches or a structural optimization theorem. The 2026 primary source and targeted searches concerning Ewald-point minima, branching bundles, and the reported counts did not reveal this construction. Absence of a search hit is not a priority certificate. The defensible conclusion is an explicit, proved improvement over the bounds in the cited paper, subject to external checking and a broader literature review. This is a more substantial advance than Round 1's restricted obstruction, but no exceptional-journal claim is made.



\chapter[An Exponential Barrier for Arbitrary-Facet Simplex-Bundle Towers]{An Exponential Barrier for Arbitrary-Facet Simplex-Bundle Towers}
\label{source:R3E6}
\noindent\textit{Source development: Round 3, entry 06. Original mathematical arguments retained; notation is local to this chapter.}\par\medskip


\begin{quote}\small\noindent\textbf{Chapter abstract.} 
\textbf{Status: a lower bound for a broad construction class, not for all smooth reflexive polytopes.}
Starting with a segment, allow simplex bundles of arbitrary fiber dimension, attached at any surviving facet, with arbitrary branching, repeated use of old facets, and integer twists from zero to the fiber dimension. Every resulting $D$-polytope has at least
$3(203/100)^{D-1}$ symmetric lattice points.
More generally, over any fixed monotone seed, the symmetric lattice points projecting to the seed's origin already grow at rate at least $203/100$ per added dimension.
A two-piece rational potential is multiplicative from below under parallel attachments and expands uniformly under simplex composition. Unimodality handles all fiber dimensions at least nine; 216 explicitly reproducible rational corner checks cover the remaining dimensions.
Together with the attributed construction from Round 2, this brackets the asymptotic minimum for the segment-generated class between $2.03$ and $2.0627969707\ldots$.
The original Crespo--Pelayo--Santos definition already permits arbitrary facets: we do not introduce that permission or relax their global minimum problem.
\end{quote}

\section{The question, the class, and the result}
For a polytope $P\subset\srcARR^D$, write
$\srcARE(P)=P\cap(-P)\cap\srcARZ^D$.
A \emph{monotone} polytope means a smooth reflexive lattice polytope.
Question 6.3 of Crespo--Pelayo--Santos \cite{R1E6:CPS} asks for the minimum of $|\srcARE(P)|$ among all monotone polytopes in each dimension.
Their Definition 6.4 permits a small simplex bundle at any facet of a monotone base; it does not exclude old facets.
Their Lemma 6.5 tracks a particular new sum facet, leading through Lemma 6.6 to the upper bounds in Corollary 6.7 and Table 5.
We checked these distinctions in the published 2026 version.

Here is the exact construction class studied in this paper. Given a monotone $B\subset\srcARR^d$, a primitive facet form $\ell(x)\leq1$, and integers $n\geq1$, $0\leq q\leq n$, define
\begin{equation}\label{R3E6:eq:attachment}
 S_{n,q}(B,\ell)=
 \left\{(x,y):x\in B,\quad y_j\geq-1,\quad
 \sum_{j=1}^n y_j+q\ell(x)\leq1\right\}.
\end{equation}
The choice $q=n$ is precisely the cited small-bundle operation.
A \emph{facet-simplex tower over $B_0$} is obtained by finitely many operations \eqref{R3E6:eq:attachment}, always choosing any facet of the current polytope. A previously chosen facet may be chosen again, and no restriction is placed on the attachment order or on the positive integers $n$.
Unimodular changes of coordinates do not affect membership or the conclusions.
Let $\mathcal T$ denote the class with initial seed $B_0=[-1,1]$.

\newpage
\begin{srcARtheorem}\label{R3E6:thm:main}
Put $\srcARrate=203/100$.
\begin{enumerate}
\item Every $D$-dimensional $P\in\mathcal T$ satisfies
\begin{equation}\label{R3E6:eq:main}
 |\srcARE(P)|\geq3\srcARrate^{D-1}.
\end{equation}
\item Let $B_0$ be any monotone seed of dimension $d$, and let $P$ be a tower over it of dimension $D$. Under the natural projection $\pi:P\to B_0$,
\begin{equation}\label{R3E6:eq:seed}
 |\srcARE(P)\cap\pi^{-1}(0)|\geq\srcARrate^{D-d}.
\end{equation}
In particular $|\srcARE(P)|\geq\srcARrate^{D-d}+|\srcARE(B_0)|-1$.
\end{enumerate}
\end{srcARtheorem}
The lower bound is on a construction class, not on the unrestricted minimum in Question 6.3. A lower bound for a subclass cannot be transferred to that unrestricted minimum.

\begin{srcARcorollary}\label{R3E6:cor:window}
Let $M(D)=\min\{|\srcARE(P)|:P\in\mathcal T,\ \dim P=D\}$. Its dimension-root limit exists, and
\[
 2.03\ \leq\lim_{D\to\infty}M(D)^{1/D}
 \ \leq\bigl(19195+4\sqrt{2297138}\bigr)^{1/14}
 =2.0627969707\ldots.
\]
The upper bound is the construction of Round 2, Entry 06 \cite[Theorem 1.1]{R3E6:R2}, not a new result of this paper.
\end{srcARcorollary}

\paragraph{Precise advance and dependencies.}
Round 1, Entry 06 \cite{R3E6:R1}, proved sharp lower bounds only for towers that always re-marked the newest sum facet.
Round 2 introduced improved examples using other facet choices, but no exponential lower bound for that larger class.
The present theorem allows all attachment histories in \eqref{R3E6:eq:attachment} and establishes a strict barrier above $2$ for them.
The geometric attachment argument and conditional-count formula are reused and reproved; the two-piece potential, its preservation under arbitrary parallel branching, and the all-dimension lower bound are the proposed increment.
All frozen Round 2 submissions were read. No result from another entrant is a mathematical dependency, and the upper bound in Corollary~\ref{R3E6:cor:window} is explicitly inherited.

\section{Geometry and the historical facet decomposition}
\begin{srcARlemma}\label{R3E6:lem:geometry}
Operation \eqref{R3E6:eq:attachment} preserves monotonicity. All old inequalities remain facets. Its new primitive facet forms are
\begin{equation}\label{R3E6:eq:heights}
 h_0=\sum_j y_j+q\ell(x),\qquad h_j=-y_j\quad(1\leq j\leq n).
\end{equation}
\end{srcARlemma}
\begin{proof}
After writing $z_j=y_j+1$, the fiber is
$z_j\geq0$, $\sum_jz_j\leq n+1-q\ell(x)$.
The last bound is at least $n+1-q\geq1$ on $B$.
Consequently all fibers are full-dimensional simplices, and all old and new inequalities are actual facets.
Vertices lie over base vertices $v$, with either all $y_j=-1$, or one coordinate equal to $n-q\ell(v)$ and the others $-1$.
They are integral.
There are no other vertices: along a short nontrivial segment in a positive-dimensional minimal base face, any active fiber equalities can be preserved, because the fiber does not degenerate.
At a listed vertex the active-normal matrix is block triangular. Its base block is unimodular, and the fiber block consists of $-I_n$, or $n-1$ of its rows and the all-ones row, with determinant $\pm1$.
New normals are primitive because of a coefficient $\pm1$ in a new coordinate.
All right-hand sides are one and the origin is strictly interior.
\end{proof}

For a symmetric lattice point, every facet height is in $\{-1,0,1\}$.
The new heights obey
\begin{equation}\label{R3E6:eq:conservation}
 h_0+\cdots+h_n=q\,t,\qquad t=\ell(x).
\end{equation}
Conversely, such an assignment corresponds to exactly one fiber point, by $y_j=-h_j$.

It is important to organize a tower by the \emph{birth of its facets}, not by whichever facet was most recently marked.
Create a node for each facet when it appears.
An attachment uses exactly one existing facet, its parent node, and creates $n+1$ new facet nodes.
All attachments directly at the same node use disjoint new coordinates.
Later attachments depend on an old coordinate only through their single parent facet height.
Thus the descendants of distinct newly created facet nodes have disjoint variables, and each historical node may have any number of direct attachments.
This is a rooted decomposition of the construction history, even when the chronological sequence repeatedly returns to an old facet.

For any collection of descendants of one facet, let $u$ be the number of completions when the parent height is $1$ or $-1$, and $v$ when it is zero.
Negating all descendant coordinates equates the outer counts.
An empty collection has $(u,v)=(1,1)$.
There are two recursive operations:
\begin{align}
 \text{parallel:}\qquad &(u,v)=(u_1u_2,v_1v_2),\label{R3E6:eq:parallel}\\
 \text{simplex:}\qquad
 &u=[z^q]\prod_{i=0}^n\bigl(v_i+u_i(z+z^{-1})\bigr),\quad
 v=[z^0]\prod_{i=0}^n\bigl(v_i+u_i(z+z^{-1})\bigr).\label{R3E6:eq:compose}
\end{align}
The first operation combines disjoint attachments sharing their parent.
The second follows by weighting the height assignments in \eqref{R3E6:eq:conservation}.
The added dimensions are, respectively, $d_1+d_2$ and $n+\sum_i d_i$.

\begin{srcARlemma}\label{R3E6:lem:unimodal}
Every historical branch has $0<u\leq v$.
For nonnegative $r_i\leq1$, the coefficients of
$\prod_i(1+r_i(z+z^{-1}))$ are symmetric and nonincreasing away from the central coefficient.
\end{srcARlemma}
\begin{proof}
The second statement follows by induction. If a symmetric coefficient sequence $c_k$ is nonincreasing for $k\geq0$, multiplication by $1+r(z+z^{-1})$ gives coefficients $d_k$ with
\[
 d_k-d_{k+1}
 =r(c_{k-1}-c_{k+2})+(1-r)(c_k-c_{k+1})\geq0.
\]
For $k=0$, use $c_{-1}=c_1\geq c_2$.
For branch counts, start at $(1,1)$. Products preserve $0<u\leq v$.
In \eqref{R3E6:eq:compose}, normalize each factor by $v_i$ and apply the second statement. Positivity follows because all input counts are positive and $0\leq q\leq n<n+1$.
\end{proof}

\section{A two-piece potential}
For $0\leq u\leq v$, $v>0$, define
\begin{equation}\label{R3E6:eq:potential}
 \srcARF(u,v)=
 \min\left\{\frac{96v+399u}{400},\,
                 \frac{91v+209u}{300}\right\}.
\end{equation}
Equivalently, $\srcARF(u,v)=v f(u/v)$, where
\[
 f(r)=\min\left\{\frac6{25}+\frac{399}{400}r,\,
                    \frac{91}{300}+\frac{209}{300}r\right\}.
\]
The two lines meet at $t=4/19$, and
\begin{equation}\label{R3E6:eq:basic}
 f(0)=\frac6{25},\quad f(t)=\frac9{20},\quad f(1)=1,\qquad
 \frac6{25}\leq f(r)\leq1,\qquad 3f(r)\leq1+2r.
\end{equation}
These inequalities follow by checking the three endpoints $0,t,1$ of its linearity intervals.
The potential is homogeneous of degree one and increasing in both arguments.

\begin{srcARlemma}[Parallel preservation]\label{R3E6:lem:parallel}
For $0\leq u_i\leq v_i$ with $v_i>0$,
\[
 \srcARF(u_1u_2,v_1v_2)\geq\srcARF(u_1,v_1)\srcARF(u_2,v_2).
\]
\end{srcARlemma}
\begin{proof}
Normalize to $v_1=v_2=1$.
For either defining line $L$, the expression
$L(rs)-f(r)f(s)$ is affine in each variable separately on each rectangle cut out by $0,t,1$.
It is therefore enough to check its corners.
By symmetry the six pairs below are sufficient; both columns are nonnegative.
\[
\begin{array}{c|cc}
(r,s)&L_1(rs)-f(r)f(s)&L_2(rs)-f(r)f(s)\\ \hline
(0,0)&114/625&1843/7500\\
(0,t)&33/250&293/1500\\
(0,1)&0&19/300\\
(t,t)&621/7600&1001/7600\\
(t,1)&0&0\\
(1,1)&19/80&0
\end{array}
\]
Taking the minimum of the two lines proves the claim.
\end{proof}

\begin{srcARlemma}[Uniform simplex expansion]\label{R3E6:lem:expansion}
For every $n\geq1$, every $0\leq q\leq n$, and all pairs $0\leq u_i\leq v_i$ with $v_i>0$, let $(u,v)$ be given by \eqref{R3E6:eq:compose}. Then
\begin{equation}\label{R3E6:eq:expand}
 \srcARF(u,v)\geq\srcARrate^n\prod_{i=0}^n\srcARF(u_i,v_i).
\end{equation}
\end{srcARlemma}
\begin{proof}
Homogeneity permits $v_i=1$, so write $r_i=u_i$.
Unimodality gives $[z^q]\geq[z^n]$; monotonicity of $\srcARF$ therefore reduces the proof to $q=n$.

First suppose $n\geq9$. The central coefficient $v$ is maximal among $2n+3$ coefficients whose sum is $\prod_i(1+2r_i)$. By \eqref{R3E6:eq:basic},
\[
 \srcARF(u,v)\geq\frac6{25}v
 \geq\frac{6\cdot3^{n+1}}{25(2n+3)}\prod_i f(r_i).
\]
The scalar factor is at least $\srcARrate^n$ for all $n\geq9$.
At $n=9$ this is the exact positive integer certificate
\[
 6\cdot3^{10}\cdot100^9-25\cdot21\cdot203^9
 =46950774614029217521425>0.
\]
The ratio of consecutive left factors divided by $\srcARrate$ exceeds one, since
$300(2n+3)-203(2n+5)=194n-115>0$.

It remains to treat $1\leq n\leq8$.
For either line $L_j$, the difference
\[
 vL_j(u/v)-\srcARrate^n\prod_i f(r_i)
\]
is affine in each $r_i$ separately on each box cut out by $0,t,1$:
both coefficients $u,v$ are multiaffine, and so is the product of the selected input lines.
Hence its minimum is attained at a corner with each $r_i\in\{0,t,1\}$.
Permutation symmetry reduces these corners to triples $(a,b,c)$ of nonnegative integers with $a+b+c=n+1$.
Here $a,b,c$ count the coordinates equal to $0,t,1$.

Define integers
\begin{equation}\label{R3E6:eq:certificatecoeff}
 U=[z^n](19+4z+4z^{-1})^b(1+z+z^{-1})^c,\qquad
 V=[z^0](19+4z+4z^{-1})^b(1+z+z^{-1})^c.
\end{equation}
At this corner $u=U/19^b$, $v=V/19^b$, and the two relevant ratios are
\begin{equation}\label{R3E6:eq:corner-ratios}
 C_1=\frac{(96V+399U)/400}
 {19^b(6/25)^a(9/20)^b},\qquad
 C_2=\frac{(91V+209U)/300}
 {19^b(6/25)^a(9/20)^b}.
\end{equation}
Table~\ref{R3E6:tab:corners} gives their minima over all triples.
Every entry in row $n$ is at least $(203/100)^n$, by exact integer cross multiplication.
These are 216 corners, hence 432 rational inequalities, not a numerical sample.
Appendix~\ref{R3E6:sec:certificate} contains a complete short verifier of the table and all comparisons.
This proves \eqref{R3E6:eq:expand}.
\end{proof}

\begin{table}[ht]
\centering
\begin{tabular}{c rr cc}
\toprule
$n$&$\min C_1$&$\min C_2$&
attaining $(a,b,c)$ for $C_1$&for $C_2$\\
\midrule
1&$543/200$&$691/300$&$(0,0,2)$&$(0,0,2)$\\
2&$25/6$&$316/75$&$(2,0,1)$&$(0,0,3)$\\
3&$4648000/555579$&$171/20$&$(0,3,1)$&$(0,0,4)$\\
4&$6891/400$&$2843/150$&$(0,0,5)$&$(0,0,5)$\\
5&$1593/40$&$939/20$&$(0,0,6)$&$(0,0,6)$\\
6&$40521/400$&$18613/150$&$(0,0,7)$&$(0,0,7)$\\
7&$13683/50$&$102409/300$&$(0,0,8)$&$(0,0,8)$\\
8&$60987/80$&$28753/30$&$(0,0,9)$&$(0,0,9)$\\
\bottomrule
\end{tabular}
\caption{Exact finite certificate for the small fiber dimensions.}\label{R3E6:tab:corners}
\end{table}

\section{Induction over arbitrary attachment histories}
\begin{srcARproposition}\label{R3E6:prop:branch}
Any historical branch using a total of $m$ new coordinates, with conditional counts $(u,v)$, satisfies
\[
 \srcARF(u,v)\geq\srcARrate^m.
\]
\end{srcARproposition}
\begin{proof}
An empty branch has $(u,v)=(1,1)$ and $\srcARF(1,1)=1$.
At one facet, combine its direct attachments by Lemma~\ref{R3E6:lem:parallel}.
For an attachment of fiber dimension $n$, combine its $n+1$ descendant branches by Lemma~\ref{R3E6:lem:expansion}; the exponent increases by $n$.
Induction on the finite historical decomposition proves the result.
It makes no assumption that the chronological construction always proceeds down a single path.
\end{proof}

\begin{proof}[Proof of Theorem~\ref{R3E6:thm:main}]
For a segment seed, all attachments descend from its two facets $x=1$ and $-x=1$.
The conditional counts of each collection are even in the parent height.
Consequently the two root collections combine by \eqref{R3E6:eq:parallel}, despite their opposite orientation, to a pair $(u,v)$ using $D-1$ coordinates.
There are exactly $u,v,u$ completions over the seed points $-1,0,1$.
By \eqref{R3E6:eq:basic} and Proposition~\ref{R3E6:prop:branch},
\[
 |\srcARE(P)|=2u+v\geq3\srcARF(u,v)\geq3\srcARrate^{D-1}.
\]

For a general seed, associate a collection of descendants with each original facet.
At the seed origin every facet height is zero, so the number of completions is the product of the corresponding central counts $v_i$.
Since $f(r)\leq1$, we have $v_i\geq\srcARF(u_i,v_i)$.
Proposition~\ref{R3E6:prop:branch} gives \eqref{R3E6:eq:seed} after multiplication.
Every other seed Ewald point has facet heights in $\{-1,0,1\}$, and every required completion count is a positive integer by Lemma~\ref{R3E6:lem:unimodal}.
Thus each of the $|\srcARE(B_0)|-1$ other seed points contributes at least one further point.
\end{proof}

\begin{proof}[Proof of Corollary~\ref{R3E6:cor:window}]
The class $\mathcal T$ contains cubes, so $M(D)$ is a minimum of a nonempty set of positive integers.
It is closed under products. Indeed, after constructing one factor, an operation $S_{1,0}$ creates an independent segment, and the second factor's sequence can be performed using only the facets and coordinates of that segment and its descendants.
Thus $M(d+e)\leq M(d)M(e)$.
Subadditivity of $\log M(D)$ proves existence of its dimension-root limit, and \eqref{R3E6:eq:main} gives its lower bound.

For the upper bound, the family in \cite[Theorem 1.1]{R3E6:R2} uses only $q=n$ instances of \eqref{R3E6:eq:attachment}, hence belongs to $\mathcal T$.
Its repeatable fourteen-dimensional operation has transfer matrix
$\left(\begin{smallmatrix}18707&18707\\1952&19683\end{smallmatrix}\right)$
and dominant eigenvalue $19195+4\sqrt{2297138}$.
We invoke that attributed construction only for the upper bound, not as a premise of our lower-bound proof.
\end{proof}

\section{Scope, obstruction, and reproducibility}
\paragraph{What the barrier excludes.}
No family generated from a fixed seed by \eqref{R3E6:eq:attachment} can have Ewald-point growth rate below $2.03$.
This includes arbitrarily complicated branching and repeated use of old facets, as well as nonmaximal twists.
For a finite collection of fixed seeds the same asymptotic barrier holds uniformly, by taking the largest seed dimension.
It does not extend to arbitrary growing-dimensional seeds, arbitrary fiber polytopes, or linear twists depending on several base facet heights.
Those operations are not encoded by our one-parent historical decomposition.
In particular the general bundle description in the original source is broader than the class treated here.
The theorem does not claim to settle Question 6.3, prove a global exponential lower bound, or optimize the class's rate.

\paragraph{Failed routes and the useful replacement.}
The earlier slice inequality $a\leq b$ for an arbitrarily re-marked \emph{global} facet is false: Round 2 exhibited $(a,b)=(20,19)$.
Our new inequality $u\leq v$ concerns conditional completions of a \emph{descendant branch}, not those global slices; these are different objects.
Simple monomial potentials and a single linear potential did not produce the desired barrier in the exploratory calculations.
The two-piece potential makes parallel reuse compatible with composition and reduces all real input ratios to a small exact corner set.
No candidate search is a premise of the theorem.

\paragraph{Verification and novelty status.}
The standalone \texttt{verify.py} reproduces every finite certificate, writes all 216 corner rows to \texttt{corner\_certificates.tsv}, and checks the large-fiber integer inequality.
As independent diagnostics it compares historical-tree counts with direct geometric enumeration on 60 deterministically selected towers in dimensions two through nine, using every twist $0\leq q\leq n$ for $1\leq n\leq4$.
It also recovers the inherited 507-point example.
Those sample counts are not the proof of the theorem; the finite corner certificate and the all-dimension argument are.

Primary-source and targeted literature checks on 20 October 2026 did not locate this arbitrary-facet exponential lower bound.
That is limited evidence, not a priority certificate.
The result is a proposed structural obstruction for a precisely defined family.
Further useful work includes improving its lower potential, closing the remaining rate interval, or identifying genuinely broader operations that evade the barrier.
No later-round result is claimed.

\clearpage
\section*{Supplementary material}
\section{A complete rational certificate for Table~\ref{R3E6:tab:corners}}\label{R3E6:sec:certificate}
This standard-library Python program checks every corner, both defining lines, and the exact table minima.
The list polynomial is indexed from its lowest exponent; multiplication by $[4,19,4]$ or $[1,1,1]$ therefore implements \eqref{R3E6:eq:certificatecoeff}.
Together with the six printed parallel checks and the large-dimension proof, this is the entire finite computational component.
{\footnotesize
\begin{verbatim}
from fractions import Fraction as F
def mul(p, q):
    r = [0] * (len(p) + len(q) - 1)
    for i, x in enumerate(p):
        for j, y in enumerate(q): r[i+j] += x*y
    return r
table = [(F(543,200), F(691,300)),
         (F(25,6), F(316,75)),
         (F(4648000,555579), F(171,20)),
         (F(6891,400), F(2843,150)),
         (F(1593,40), F(939,20)),
         (F(40521,400), F(18613,150)),
         (F(13683,50), F(102409,300)),
         (F(60987,80), F(28753,30))]
count = 0
for n in range(1, 9):
    values = [[], []]
    for a in range(n+2):
        for b in range(n+2-a):
            c = n+1-a-b
            p = [1]
            for _ in range(b): p = mul(p, [4,19,4])
            for _ in range(c): p = mul(p, [1,1,1])
            V = p[b+c]
            U = p[b+c+n] if n <= b+c else 0
            den = 19**b * F(6,25)**a * F(9,20)**b
            pair = (F(96*V+399*U,400)/den,
                    F(91*V+209*U,300)/den)
            for j in (0,1):
                assert pair[j] >= F(203,100)**n
                values[j].append(pair[j])
            count += 1
    assert tuple(map(min, values)) == table[n-1]
assert count == 216
\end{verbatim}
}



\part{Integer fibers and corrected concavity}

\chapter[Barycentric Log-Concavity for Disk-Network Fibers and Rank-Two Unimodular Matrices]{Barycentric Log-Concavity for Disk-Network Fibers and Rank-Two Unimodular Matrices}
\label{source:R4E6}
\noindent\textit{Source development: Round 4, entry 06. Original mathematical arguments retained; notation is local to this chapter.}\par\medskip


\begin{quote}\small\noindent\textbf{Chapter abstract.} 
We prove the finite-support barycentric lattice-point inequality asked in Question 1.7 of Leake--Mohammadi Yekta for two restricted classes: all unimodular matrices of rank at most two, and arbitrary-rank directed networks embedded in a disk when the variable balances are confined to boundary terminals. In particular, all outerplanar incidence matrices are covered. The result permits finite coordinate boxes and positive monomial weights; finite nonnegative fibers follow by truncation. The proof combines established four-functions marginalization with a planar face-potential representation and an elementary integer-energy balancing argument. The rank-two case also has an independent difference-constraint realization and a concave piecewise-affine logarithmic interpolation. We do not resolve the question for general unimodular matrices or for infinite-support distributions. The underlying $L$-log-concavity machinery is known; novelty of this particular application has not been established by a comprehensive priority search.
\end{quote}

\section{Target, statements, and attribution}
For an integer matrix $M\in\srcAXZ^{m\times n}$, integer bounds $l\leq u$, and weights $w_j>0$, put
\begin{equation}\label{R4E6:eq:fiber}
 F_{M,l,u,w}(b)=\sum_{\substack{x\in\srcAXZ^n,\ Mx=b\\l\leq x\leq u}}w^x,
 \qquad w^x=\prod_{j=1}^n w_j^{x_j}.
\end{equation}
All bounds in \eqref{R4E6:eq:fiber} are finite unless stated otherwise. We call a matrix of rank $r$ \emph{unimodular} when every nonsingular $r\times r$ submatrix has determinant $\pm1$.

Leake--Mohammadi Yekta \cite[Theorems 1.4 and 3.1]{R4E6:LMY} establish midpoint log-concavity for unimodular fibers, including weighted box versions. Their Question 1.7 asks for arbitrary barycentric combinations with integer mean, and explicitly distinguishes this from concavity along lattice lines. We use the version dated 30 September 2026, checked on 20 October 2026. For a nonnegative function $F$ on a lattice, the finite-support conclusion considered here is
\begin{equation}\label{R4E6:eq:jensen}
 F(\bar b)\ \geq\ \prod_{i=1}^{s}F(b^{(i)})^{\mu_i},
 \quad \mu_i>0,\quad \sum_i\mu_i=1,\quad
 \bar b=\sum_i\mu_i b^{(i)}\text{ an integer vector}.
\end{equation}
We state results for positive input fibers; a zero input makes the assertion immediate.

A \emph{disk network} means a finite connected directed multigraph with an embedding in a closed disk, a set of distinct terminals $t_1,\ldots,t_k$ in counterclockwise order on its boundary, and all other vertices and all edge interiors in the open disk. Parallel edges and bridges are allowed. Directions have no bearing on embeddability. We fix the divergence convention ``outgoing minus incoming.'' This explicit embedding hypothesis avoids any implicit assertion about an arbitrary prescribed planar embedding with repeated vertices in a facial walk.

\begin{srcAXtheorem}[Disk-network fibers]\label{R4E6:thm:network}
Let $G$ be a disk network, with finite integer edge bounds and positive edge weights. Fix the divergences at all nonterminal vertices. If $b^{(1)},\ldots,b^{(s)}$ are feasible integer divergence vectors with these fixed nonterminal entries, and $\mu_i>0$ sum to one with integer mean $\bar b$, then the corresponding weighted flow sums satisfy \eqref{R4E6:eq:jensen}.
\end{srcAXtheorem}

\begin{srcAXtheorem}[Rank at most two]\label{R4E6:thm:rank2}
For every unimodular integer matrix $M$ of rank at most two, every finite integer coordinate box, and every positive monomial weight, the function \eqref{R4E6:eq:fiber} satisfies \eqref{R4E6:eq:jensen} for all finite collections of feasible integer margins with integer mean. The number of rows need not equal the rank.
\end{srcAXtheorem}

\begin{srcAXcorollary}[Nonnegative fibers and outerplanar matrices]\label{R4E6:cor:unbounded}
The conclusions remain valid for unweighted nonnegative fibers whenever the fibers in question are finite. In particular, the finite-support case of \cite[Question 1.7]{R4E6:LMY} holds for matrices of rank at most two and for incidence matrices of outerplanar directed multigraphs, including their reduced incidence matrices. The disk-network conclusion also applies with arbitrary fixed interior divergences. Disjoint unions are covered componentwise.
\end{srcAXcorollary}

\paragraph{What is, and is not, being claimed.}
The four-functions theorem is due to Ahlswede--Daykin \cite{R4E6:AD}. Marginal closure for $L$-log-concavity is already explicit in Speyer \cite[Lemma 5.3 and Theorem 5.1]{R4E6:Speyer}, whose current version is v2, 30 April 2026; see also his discussion of Murota's discrete convexity. We claim neither this closure principle nor the distinction between line concavity and barycentric concavity as new. Below we give all details of the shifted inequality, the balancing step, and the network and rank-two reductions. The contribution of this note is this restricted-class application to the stated question. It is a partial result, with priority of the application unestablished, not a solution of the general question. Earlier rounds by this entrant concerned symmetric lattice points of smooth reflexive polytopes \cite{R4E6:R2,R4E6:R3}; none of their bounds is strengthened here, and none is counted again as a new result.

\section{Difference constraints and finite barycentric balancing}
Throughout, $\vee$ and $\wedge$ denote coordinatewise maximum and minimum. A difference-constraint set consists of integer vectors $z$ satisfying finitely many conditions
\begin{equation}\label{R4E6:eq:difference}
 a_{ij}\leq z_i-z_j\leq b_{ij},\qquad a_{ij},b_{ij}\in\srcAXZ.
\end{equation}
Coordinates are partitioned into prescribed coordinates $p$ and summed coordinates $q$. One or more prescribed coordinates can be fixed at zero. Write
\begin{equation}\label{R4E6:eq:partition}
 Z(p)=\sum_{q:\,(p,q)\text{ satisfies }\eqref{R4E6:eq:difference}}
       \exp(c\cdot(p,q)),
\end{equation}
where $c$ is any fixed real vector. Assume that the summation set is finite for every $p$ under consideration, and in the lemmas below for every prescribed vector produced by an operation. In our applications this holds for all $p$.

We use the following classical form of the four-functions theorem: if nonnegative finitely supported functions $A,B,C,D$ on $\srcAXZ^d$ satisfy
\[
 A(u)B(v)\leq C(u\vee v)D(u\wedge v)\quad\hbox{for every }u,v,
\]
then $(\sum A)(\sum B)\leq(\sum C)(\sum D)$ \cite[Theorem 1 and its distributive-lattice corollary]{R4E6:AD}. Restriction to a sufficiently large finite product of chains gives precisely this formulation. This is the only correlation theorem used below.

\begin{srcAXlemma}[Shifted marginal inequality]\label{R4E6:lem:shift}
For every integer $\alpha\geq0$, the partition function \eqref{R4E6:eq:partition} obeys
\begin{equation}\label{R4E6:eq:shift}
 Z(p)Z(r)\leq
 Z\bigl((p-\alpha\srcAXone)\vee r\bigr)
 Z\bigl(p\wedge(r+\alpha\srcAXone)\bigr).
\end{equation}
Coordinates prescribed to be zero in both inputs remain zero in both outputs.
\end{srcAXlemma}
\begin{proof}
The set defined by \eqref{R4E6:eq:difference}, before fixing any coordinates, is invariant under common translations and closed under both coordinatewise maximum and minimum. Indeed, each upper inequality $z_i\leq z_j+b_{ij}$ is preserved by either operation, and lower inequalities are reversed upper inequalities. Hence, if $z,w$ are feasible, so are
\[
 z'=(z-\alpha\srcAXone)\vee w,\qquad
 w'=z\wedge(w+\alpha\srcAXone).
\]
For each coordinate, $z'_i+w'_i=z_i+w_i$. Thus exponential weights have exactly the same product before and after the operation. An anchored coordinate has outputs $\max(-\alpha,0)=0$ and $\min(0,\alpha)=0$.

Extend the feasible weight to $W(p,q)$ by assigning zero to infeasible vectors. Set $p'=(p-\alpha\srcAXone)\vee r$ and $r'=p\wedge(r+\alpha\srcAXone)$. On the summed-coordinate lattice take
\[
 A(u)=W(p,u+\alpha\srcAXone),\quad B(v)=W(r,v),\quad
 C(t)=W(p',t),\quad D(t)=W(r',t+\alpha\srcAXone).
\]
The pointwise product assertion just proved gives the hypothesis of four-functions, since the second output summed coordinate is $(u\wedge v)+\alpha\srcAXone$. Summing, and translating the summation index in $A,D$, gives \eqref{R4E6:eq:shift}.
\end{proof}

\begin{srcAXlemma}[Integer-energy balancing]\label{R4E6:lem:balancing}
Let $Z:\srcAXZ^d\to\srcAXR_{\geq0}$ satisfy \eqref{R4E6:eq:shift} for $\alpha=1$. Then $Z$ satisfies finite-support barycentric log-concavity \eqref{R4E6:eq:jensen}.
\end{srcAXlemma}
\begin{proof}
First take $N$ vectors $p_1,\ldots,p_N$ with integer average $\bar p$ and positive $Z(p_i)$. If a coordinate has range at least two, choose an ordered pair $x,y$ among the vectors with $x_j-y_j\geq2$ in that coordinate, and replace it by
\[
 x'=(x-\srcAXone)\vee y,\qquad y'=x\wedge(y+\srcAXone).
\]
The sum of the vectors is unchanged and their product of $Z$-values cannot decrease. In one coordinate put $d=x_i-y_i$. If $d\leq1$, the two output values are $y_i,x_i$, merely interchanged. If $d\geq2$, they are $x_i-1,y_i+1$, and their squared sum drops by $2(d-1)>0$. Consequently
\[
 \sum_{a=1}^N\|p_a-\bar p\|_2^2
\]
is a nonnegative integer that strictly decreases at each operation. All output $Z$-values remain positive by \eqref{R4E6:eq:shift}. The process terminates with coordinate ranges at most one. An integer average of integers drawn from $\{h,h+1\}$ can be an integer only when they are all equal. Thus all final vectors equal $\bar p$, proving $\prod_a Z(p_a)\leq Z(\bar p)^N$.

This proves the assertion for rational probabilities by repetitions. For real probabilities on fixed $p_1,\ldots,p_s$ and fixed integer mean $\bar p$, consider
\[
 \mathcal P=\{\nu\in\srcAXR_{\geq0}^s:\ \sum\nu_i=1,\ \sum\nu_i p_i=\bar p\}.
\]
This is a bounded polytope defined over $\srcAXQ$. Its vertices are rational, by solving the rational linear systems of active constraints, and rational convex combinations of these vertices are dense in $\mathcal P$. Approximate the given probability vector by such rational points. The rational result first ensures $Z(\bar p)>0$, and continuity of $\sum_i\nu_i\log Z(p_i)$ yields the real-weight assertion. Zero input values were excluded only to avoid writing $\log0$; their original inequalities are trivial.
\end{proof}

\begin{srcAXremark}
The operations in Lemma \ref{R4E6:lem:balancing} do not in general preserve a prescribed set of input points or keep them on an input line. That is precisely why midpoint log-concavity alone is insufficient for this proof. The stronger shifted min/max inequality supplies the intermediate points and their positive weights.
\end{srcAXremark}

\section{Planar duality with boundary balances}
We prove Theorem \ref{R4E6:thm:network}. The case of zero or one terminal is immediate: conservation fixes every divergence entry. Suppose $k\geq2$.

\subsection{Reduction to zero interior divergence}
Choose an integer feasible reference flow $x^0$ in one input fiber, and write $x=x^0+y$. The new bounds are $l-x^0\leq y\leq u-x^0$. The new divergences $\beta=b-\srcAXdivg x^0$ vanish at nonterminals and sum to zero on terminals. The weight becomes $w^{x^0}w^y$. This constant factor cancels from the desired inequality because the probabilities sum to one. We may therefore assume zero interior divergence, signed finite bounds, and
\[
 \sum_{i=1}^k\beta_i=0.
\]

\subsection{An auxiliary boundary cycle}
Adjoin counterclockwise boundary arcs $e_i:t_i\to t_{i+1}$, with indices cyclic. They may be drawn on the disk boundary, disjoint from all old edge interiors. For $k=2$ they are the two distinct arcs of the boundary circle. Prescribe their flows to be
\begin{equation}\label{R4E6:eq:cumulative}
 p_i=-\sum_{j=1}^{i}\beta_j\quad(1\leq i<k),\qquad p_k=0.
\end{equation}
The added divergence at $t_i$ is $p_i-p_{i-1}=-\beta_i$, taking $p_0=p_k=0$. Thus old flows of divergence $\beta$ are in bijection with circulations on the augmented plane graph having these prescribed boundary flows. The transformation between terminal balances summing to zero and $(p_1,\ldots,p_{k-1})\in\srcAXZ^{k-1}$ is an integer linear bijection: $\beta_i=p_{i-1}-p_i$. In particular it preserves integer barycenters. Boundary flows are auxiliary, signed, unweighted, and are not subject to the original capacities.

Each new boundary arc has a different bounded face on its left. Here is a proof that includes bridges. Retain a spanning tree of the original connected graph and all $k$ boundary arcs. Euler's formula gives exactly $k$ bounded faces. A bounded face must contain a boundary arc: a face bounded only by tree edges would require a cycle in the tree. There are exactly $k$ such arcs, each incident to one bounded face and the exterior face. Hence each bounded face contains exactly one of them. Reinsert the other original edges one at a time. They split faces, never merge them, so the bounded faces incident to distinct boundary arcs remain distinct.

\subsection{Integral face potentials}
For the augmented graph let $L(e),R(e)$ be the faces on the left and right of an oriented edge. Circulations have the unique representation
\begin{equation}\label{R4E6:eq:dual}
 y_e=z_{L(e)}-z_{R(e)},\qquad z_{\rm ext}=0,
\end{equation}
with integer face potentials. For completeness, the divergence of these differences at every vertex is zero by telescoping the incident face potentials in cyclic order. Conversely, any integer circulation decomposes into signed simple directed cycles: orient an edge against its displayed direction when its current flow is negative, find a directed cycle in the nonzero circulation support, subtract the minimum positive flow on that cycle, and repeat. The sum of absolute edge flows strictly decreases. The indicator of the bounded faces inside a simple cycle, with the sign determined by its orientation, represents that cycle in \eqref{R4E6:eq:dual}. Adding these indicators constructs integral potentials for the circulation. Uniqueness follows because zero edge differences force equal potentials on adjacent faces, and the dual adjacency graph is connected; a curve from a face to the exterior gives an adjacency path. Bridges cause no exception: their two sides belong to the same face, and their circulation flow is zero.

Let $f_i$ be the bounded face left of boundary arc $e_i$. By the preceding distinctness result and the counterclockwise orientation,
\[
 z_{f_i}=p_i,\qquad z_{\rm ext}=0.
\]
These are genuinely distinct prescribed coordinates; in particular $z_{f_k}=0$ is a second anchor, not an equation identifying two variable boundary coordinates. The remaining bounded-face potentials are free variables $q$.

For every original edge, its finite bound becomes
\[
 l_e\leq z_{L(e)}-z_{R(e)}\leq u_e,
\]
and the monomial weight becomes
\[
 \exp\left(\sum_{e\text{ original}}(\log w_e)
                   (z_{L(e)}-z_{R(e)})\right).
\]
This is exactly the difference-constraint partition function of Section 2. For each fixed $p$ its summed-coordinate support is finite. Indeed, the augmented dual is connected, original-edge potential differences have finite bounds, and differences across each boundary edge are prescribed finite values. Along a dual path from the exterior these bounds bound each potential.

Lemma \ref{R4E6:lem:shift} and Lemma \ref{R4E6:lem:balancing} now give barycentric log-concavity in $p$. The integer linear transformation \eqref{R4E6:eq:cumulative}, followed by undoing the reference flow, gives \eqref{R4E6:eq:jensen} in the original balances. This proves Theorem \ref{R4E6:thm:network}.

\paragraph{Scope of the embedding assumption.}
No variation of interior balances was permitted in the argument. No assertion is made for arbitrary planar networks when all their vertex balances vary. A connected outerplanar multigraph admits a disk embedding with all vertices on the boundary (parallel edges can be placed alongside their underlying edge); it therefore is covered with every vertex terminal. Loops, if present, can first be removed and their independent finite weighted sums factored out. For disconnected graphs the flow sum is a product over connected components, so applying the theorem componentwise proves the asserted disjoint-union version.

\section{Rank two: normalization and an independent realization}
The following proof of Theorem \ref{R4E6:thm:rank2} uses only Section 2, not the planar argument.

\subsection{The relative lattice and three column directions}
Let $M$ have rank $r\in\{1,2\}$. Choose $r$ independent rows, let $R$ denote projection to them, and choose $r$ columns such that the square matrix $B=RM_{J}$ is nonsingular. By unimodularity $\det B=\pm1$. Put
\[
 N=B^{-1}RM.
\]
Projection $R$ is injective on the real column span $V$ of $M$. Moreover,
\begin{equation}\label{R4E6:eq:lattice}
 V\cap\srcAXZ^m\longrightarrow\srcAXZ^r,\qquad b\longmapsto B^{-1}Rb
\end{equation}
is a lattice isomorphism, with inverse $t\mapsto M_Jt$: this follows from injectivity and $B^{-1}\in\srcAXZ^{r\times r}$. The original and normalized fibers are therefore identical for $b\in V$. Every barycenter of feasible margins lies in $V$, so redundant rows cause no loss of information or unmentioned lattice rescaling.

When $r=2$, the chosen basis columns of $N$ are $e_1,e_2$. Cramer's rule and the maximal-minor hypothesis show that every entry of every column of $N$ is in $\{0,\pm1\}$. Thus the only possible nonzero directions up to sign are
\[
 e_1,\quad e_2,\quad e_1+e_2,\quad e_1-e_2.
\]
The last two cannot both occur, since their determinant has absolute value two. A sign change of the second row, if necessary, followed by individual column sign changes, leaves precisely the three directions $e_1,e_2,e_1+e_2$.

An individual column sign change is legitimate even in the presence of opposite columns: replace its variable $x_j$ by $-x_j$, its bounds $[l_j,u_j]$ by $[-u_j,-l_j]$, and its weight $w_j$ by $w_j^{-1}$. This is an exact bijection with exactly the same term weight. Zero columns contribute the positive constant $\prod_{j:N_j=0}\sum_{a=l_j}^{u_j}w_j^a$, which cancels from Jensen's inequality. No pointedness assumption or deletion of opposite columns is used for finite boxes.

\subsection{Three paths of difference constraints}
Partition the remaining variables into groups $A,B,C$ according to the three directions. The margin equations are
\[
 a=\sum_{j\in A}x_j+\sum_{j\in C}x_j,\qquad
 b=\sum_{j\in B}x_j+\sum_{j\in C}x_j.
\]
Introduce a common potential $t=\sum_{j\in C}x_j$. Realize the $A$ variables, in any fixed order, as successive differences along a path from potential $a$ to potential $t$; realize $B$ along a path from $b$ to $t$; realize $C$ along a path from $t$ to an anchor of potential zero. All internal vertices of these paths are distinct. The sum of differences along each path is its endpoint difference. Consequently assignments of the original variables satisfying the two margin equations are in bijection with assignments of all internal potentials, including $t$, with prescribed potentials $(a,b,0)$. Bounds become difference bounds and monomial weights are exponentials of linear forms in potentials. If a group is empty, replace its path by the difference equality of its two endpoints. Finite edge bounds and connection to the anchor ensure that every fiber is finite. Lemmas \ref{R4E6:lem:shift} and \ref{R4E6:lem:balancing} apply directly.

For $r=1$, normalized nonzero columns are $\pm1$. Reflect the negative ones in the same way and use a single path from the normalized margin to the zero anchor, with one edge for each variable. For $r=0$ the only feasible margin is zero, so the conclusion is immediate. These arguments, together with \eqref{R4E6:eq:lattice}, prove Theorem \ref{R4E6:thm:rank2}.

\subsection{All three rhombi and concavity on the full support}
This subsection supplies an independent geometric description of the rank-two conclusion; no piecewise-linear assertion is needed in the preceding proofs. Denote the normalized function, after dividing out zero-column factors, by $f(a,b)$. The shifted inequality gives the following three inequalities, valid also with zero values:
\begin{align}
 f(a,b)f(a+1,b+1)&\geq f(a+1,b)f(a,b+1),\label{R4E6:eq:r1}\\
 f(a,b)f(a+1,b)&\geq f(a,b-1)f(a+1,b+1),\label{R4E6:eq:r2}\\
 f(a,b)f(a,b+1)&\geq f(a-1,b)f(a+1,b+1).\label{R4E6:eq:r3}
\end{align}
For \eqref{R4E6:eq:r1}, use $p=(a+1,b),r=(a,b+1),\alpha=0$ in \eqref{R4E6:eq:shift}; the output pair is $(a+1,b+1),(a,b)$. For \eqref{R4E6:eq:r2}, use $p=(a+1,b+1),r=(a,b-1),\alpha=1$; the outputs are $(a,b),(a+1,b)$. For \eqref{R4E6:eq:r3}, use $p=(a+1,b+1),r=(a-1,b),\alpha=1$; the outputs are $(a,b),(a,b+1)$. The anchor coordinate is zero throughout.

Write $[L_A,U_A]$, $[L_B,U_B]$, $[L_C,U_C]$ for the intervals of possible group sums, taking $[0,0]$ for an empty group. Each such sum takes every integer in its interval. The convex support is the polygon
\[
 P=[L_A,U_A]e_1+[L_B,U_B]e_2+[L_C,U_C](e_1+e_2).
\]
Its integer points are exactly the positive support of $f$. In fact, for integer $(a,b)$ the allowed values of $t$ form
\[
 [L_C,U_C]\cap[a-U_A,a-L_A]\cap[b-U_B,b-L_B],
\]
whose endpoints are integers; a nonempty real intersection thus contains an integer. The polygon is bounded by lines $a=h$, $b=h$, and $a-b=h$ with integer $h$. Hence it is a union of faces of the lattice triangulation obtained by dividing each unit square along its $(1,1)$ diagonal.

Suppose first that $P$ is two-dimensional. Interpolate $\log f$ affinely on each triangle of this triangulation contained in $P$. Across a diagonal, horizontal, or vertical interior edge, respectively, the condition that the directional slope does not increase on crossing is exactly the logarithm of \eqref{R4E6:eq:r1}, \eqref{R4E6:eq:r2}, or \eqref{R4E6:eq:r3}: in each case the endpoints of the common edge sum to the two opposite vertices. For clarity, the corresponding shared edges are
\[
 [(a,b),(a+1,b+1)],\quad[(a,b),(a+1,b)],\quad[(a,b),(a,b+1)].
\]
Along a line segment in the interior of $P$ that avoids triangulation vertices and is not parallel to a triangulation edge, the affine slope is therefore nonincreasing at every crossing. The restriction to that segment is concave. Any segment in $P$ is a limit of such generic interior segments; the interpolation is continuous on the finite triangulation, so the concavity inequality passes to the limit. Thus this interpolation is concave on the \emph{entire} convex support, including its boundary.

If $P$ is one-dimensional, all group intervals with positive length have the same one of the three directions. The shifted inequality with $\alpha=1$ applied to points two unit steps apart in that direction gives ordinary consecutive log-concavity. Linear interpolation on the interval is concave. A zero-dimensional support is trivial. This verifies the full-support assertion also for degenerate boxes.

\section{Finite nonnegative fibers, limitations, and checks}
\begin{proof}[Proof of Corollary \ref{R4E6:cor:unbounded}]
Intersect each nonnegative fiber with the common box $[0,K]^n$. The relevant boxed result applies for every sufficiently large $K$ for which all the finitely many input fibers have a point. If the input fibers and the barycenter fiber are finite, one common $K$ contains them all, so the desired inequality follows immediately. Equivalently one can take the increasing limit in the boxed inequality. The same reasoning permits positive monomial weights on finite fibers.

For clarity, a nonempty finite integer fiber of $Mx=b$, $x\geq0$, forces $\ker(M)\cap\srcAXR_{\geq0}^n=\{0\}$. Otherwise the rational polyhedral recession cone contains a nonzero rational vector, and scaling gives a nonzero nonnegative integer kernel vector that can be repeatedly added to a lattice point. Thus every feasible real nonnegative fiber is bounded. This explains the finiteness requirement and in particular excludes zero columns or opposite columns from a nonempty finite unbounded fiber. They were \emph{not} excluded in the finite-box theorem. Outerplanar and disjoint-union conclusions follow as described after Theorem \ref{R4E6:thm:network}; dropping one incidence row in each component is an integer bijection between full divergence vectors with component sums zero and reduced vectors.
\end{proof}

\paragraph{Unresolved scope.}
We have not proved barycentric log-concavity for arbitrary unimodular matrices of rank three or higher, for arbitrary varying balances at interior vertices of a planar graph, or for infinite-support distributions. Nor have we derived a new general counting algorithm or a capacity approximation factor. A boundary cycle cannot encode arbitrary interior imbalances as independently prescribed face potentials; that is the specific obstruction to extending this proof directly.

\paragraph{Verification versus proof.}
The accompanying \texttt{verify.py} uses exact integer and rational arithmetic. It checks the shifted inequality and barycentric products on finite difference-constraint examples, compares explicit planar primal flow counts with dual-potential counts (including fixed nonzero interior divergence and a bridge), checks a rank-two system with opposite and zero columns, verifies the three rhombus forms, and records concrete noncollinear barycentric examples. These finite checks are diagnostics, not replacements for any universal argument. The manuscript's proof is independent of their success. Earlier floating-point searches found no counterexample in several small rank-three and rank-four samples; they neither support a general theorem nor enter the proof. An initial search for improved smooth-reflexive bundle constructions produced no improved exponential bound and was abandoned in favor of this target.



\chapter[A sharp quarter correction for two-dimensional unimodular fibers]{A sharp quarter correction for two-dimensional unimodular fibers}
\label{source:R5E3}
\noindent\textit{Source development: Round 5, entry 03. Original mathematical arguments retained; notation is local to this chapter.}\par\medskip

\begin{quote}\small\noindent\textbf{Chapter abstract.} 
For a unimodular integer matrix of nullity at most two, we prove that the square root of its nonempty finite integer-fiber count minus $1/4$ is barycentrically concave. Integer coordinate bounds may vary along with the margins. This implies the unweighted log-concavity inequality in Question 1.7 of Leake--Mohammadi Yekta on this arbitrary-rank class, including countable distributions with a finite first moment. The correction $1/4$ is optimal: a fixed rank-four totally unimodular matrix has fibers with counts $3k^2+3k+1$, which violate the uncorrected square-root inequality. The proof normalizes a two-dimensional kernel to three possible normal directions and combines Pick's formula with classical planar Brunn--Minkowski after adding half a fixed hexagon. This is a restricted result, not a resolution for general unimodular matrices; arbitrary monomial weights are not covered.
\end{quote}

\section{The theorem, target, and exact increment}
An integer matrix $M\in\srcBAZ^{m\times n}$ of rank $r$ is called \emph{unimodular} if every nonsingular $r\times r$ submatrix has determinant $\pm1$. Its nullity is $n-r$, not the number of redundant rows. For finite integer vectors $l\le u$ put
\[
 F(b,l,u)=\#\{x\in\srcBAZ^n:Mx=b,\ l\le x\le u\}.
\]
We also use the nonnegative-fiber count $F_+(b)=\#\{x\in\srcBAZ_{\ge0}^n:Mx=b\}$ when finite.

\begin{srcBAtheorem}[Sharp corrected square-root concavity]\label{R5E3:main}
Let $M$ be unimodular with nullity at most two. Let $(b_i,l_i,u_i)$, $1\le i\le s$, be integer parameters with $0<F_i:=F(b_i,l_i,u_i)<\infty$. Let $\mu_i>0$ sum to one, and suppose the componentwise mean $(\bar b,\bar l,\bar u)$ is integral. Then $\bar F:=F(\bar b,\bar l,\bar u)>0$ and
\begin{equation}\label{R5E3:sqrt}
 \sqrt{\bar F-\tfrac14}\ \ge\ \sum_{i=1}^s\mu_i\sqrt{F_i-\tfrac14}.
\end{equation}
The same assertion holds for nonempty finite nonnegative fibers, with only the margins averaged. In either setting,
\begin{equation}\label{R5E3:jensen}
 \bar F\ge\prod_{i=1}^s F_i^{\mu_i}.
\end{equation}
The universal correction $1/4$ in \eqref{R5E3:sqrt} cannot be decreased, even for a fixed rank-four totally unimodular matrix and nonnegative fibers.
\end{srcBAtheorem}

\begin{srcBAcorollary}[Countable distributions]\label{R5E3:countable}
The two inequalities remain true for countably supported probability distributions on the indicated integer parameters, provided their first absolute moment is finite and their mean is integral. For nonnegative fibers, assume all support fibers are nonempty and finite. The products in \eqref{R5E3:jensen} mean $\exp(\sum_i\mu_i\log F_i)$; this sum is finite as a consequence of the proof.
\end{srcBAcorollary}
The hypotheses are on the entire matrix nullity, not merely the dimension of selected fibers after incidental bounds become active. In the boxed statement both capacities and margins may vary. In particular, fixing the bounds is allowed. No positive nonconstant monomial weight is asserted.

\paragraph{Precisely sourced target.}
Leake--Mohammadi Yekta \cite[Question 1.7]{R4E6:LMY} ask for \eqref{R5E3:jensen} for nonnegative fibers of every unimodular matrix and probability distributions with integer mean. Their Theorem 1.4 proves midpoint log-concavity, which they distinguish from the full barycentric question. The latest source record inspected on 20 October 2026 is v1, 30 September 2026. Theorem~\ref{R5E3:main} answers that question only for nullity at most two; Corollary~\ref{R5E3:countable} includes its countably supported version under the standard integrability meaning of a finite mean. The general question remains unresolved here.

\paragraph{Prior-round credit and scope comparison.}
Round 4 entry 06 proved barycentric log-concavity for rank-at-most-two matrices and for disk networks with variable boundary balances \cite[Theorems 1--2]{R5E3:R406}. Rank and nullity are different restrictions. Its rank-two normalization motivates the complementary-minor normalization of our kernel below, and this intellectual antecedent is credited. Its four-functions and $L$-log-concavity proofs are not used as new results. Our proof uses planar area instead and supplies the stronger corrected power inequality, variable boxes, and a countable-distribution conclusion.

The rank-four example in Section~\ref{R5E3:sharp} is a theta network with underlying graph $K_{2,3}$ and all balances variable. This graph is not outerplanar, so the earlier all-boundary network result does not directly include its full parameter family. Conversely, the present theorem does not cover the earlier arbitrary-nullity rank-two and disk-network families. Our own preceding rounds concerned permanent inequalities; none of those results is repackaged as a contribution here. All six frozen Round 4 papers and the assessment were read.

\section{A planar counting identity}
Set
\begin{equation}\label{R5E3:hex}
 \begin{aligned}
 H&=\{(x,y)\in\srcBAR^2:|x|\le1,\ |y|\le1,\ |x+y|\le1\}\\
 &=\operatorname{conv}\{(1,0),(1,-1),(0,-1),(-1,0),(-1,1),(0,1)\}.
 \end{aligned}
\end{equation}
It has area $3$. Its support function is $1$ at each of the six primitive vectors $\pm e_1,\pm e_2,\pm(e_1+e_2)$.

\begin{srcBAlemma}[Half-hexagon case of Chu--Karshon]\label{R5E3:pick}
Let $P$ be a nonempty bounded intersection of half-planes whose normals belong to those six vectors and whose constants are integers. Lower-dimensional intersections are permitted. Then $P$ is a lattice polygon, segment, or point, and
\begin{equation}\label{R5E3:areaidentity}
 \#(P\cap\srcBAZ^2)=\srcBAarea(P+\tfrac12H)+\tfrac14.
\end{equation}
\end{srcBAlemma}
\begin{proof}
Any two independent allowed normals have determinant $\pm1$. Every vertex is therefore integral. In dimension two let $B(P)$ be its boundary lattice-point count, equivalently the sum of the lattice lengths of its edges. Classical Pick's formula gives
\[
 \#(P\cap\srcBAZ^2)=\srcBAarea(P)+\tfrac12 B(P)+1.
\]
The planar Minkowski-area expansion is
\begin{equation}\label{R5E3:steiner}
 \srcBAarea(P+tH)=\srcBAarea(P)+tB(P)+3t^2\qquad(t\ge0).
\end{equation}
Here is a direct justification of its normalization. Refine the normal fans of $P,H$ to their common ordered list. Vertices of $P+tH$ are sums of the corresponding vertices, linear in $t$; the shoelace formula is quadratic. The constant and quadratic coefficients are $\srcBAarea(P)$ and $\srcBAarea(H)=3$. Its derivative at zero is the sum, over the edges of $P$, of Euclidean edge length times the support function of $H$ at the outward unit normal. If the primitive normal is $a$, an edge of lattice length $\ell$ has Euclidean length $\ell\|a\|$, whereas that support value is $1/\|a\|$. The contribution is $\ell$. This derivative can also be obtained by the thin strips along the edges; the remaining corner regions have area $O(t^2)$. This proves \eqref{R5E3:steiner}. Set $t=1/2$ and apply Pick.

If $P$ is a segment, its direction is perpendicular to an allowed primitive normal. Write its lattice length as $\ell\ge0$. The width of $H$ in that normal direction is $2/\|a\|$, so adding $P$ to $tH$ sweeps area $2\ell t$ in addition to $3t^2$. At $t=1/2$ the area is $\ell+3/4$, while the segment has $\ell+1$ lattice points. A point gives area $3/4$ and count one. These calculations prove \eqref{R5E3:areaidentity} in every permitted dimension.
\end{proof}
\paragraph{Direct prior theorem.}
Lemma~\ref{R5E3:pick} is also a direct special case of Chu--Karshon \cite[Theorem 5.8]{R5E3:CK}, not a new counting identity. Indeed, $Q=P+H/2$ has exactly six facets, including when $P$ is a segment or point: each edge of $H/2$ contributes positive edge length. Their primitive outward normals are the six allowed vectors; consecutive normals have determinant one. Writing $h_P(a)=\max_{x\in P}\langle a,x\rangle$, its facet constants are $h_P(a)+1/2$, hence half-integers. For lattice $z$, the six strict inequalities defining $\operatorname{int}Q$ are equivalent to $\langle a,z\rangle\le h_P(a)$, which define membership in $P$. Thus $I(Q)=\#(P\cap\srcBAZ^2)$ and $V(Q)=6$. Their formula $\srcBAarea(Q)=I(Q)-V(Q)/8+1/2$ gives \eqref{R5E3:areaidentity}. The independent classical proof above is retained for self-containment.

\begin{srcBAlemma}[Concave extensions for three-normal families]\label{R5E3:extension}
Let $P(\theta)$ be nonempty bounded polygons, segments or points given by a fixed list of the allowed normals with affine right-hand sides in real parameters $\theta$. Let $\mathcal D$ be their convex feasibility domain. Then
\[
 G(\theta)=\sqrt{\srcBAarea(P(\theta)+\tfrac12H)},\qquad
 g(\theta)=\log\bigl(G(\theta)^2+\tfrac14\bigr)
\]
are concave on $\mathcal D$, with $G\ge\sqrt3/2$ and $g\ge0$. At parameters with integral right-hand sides,
\[
 G(\theta)=\sqrt{\#(P(\theta)\cap\srcBAZ^2)-\tfrac14},\qquad
 g(\theta)=\log\#(P(\theta)\cap\srcBAZ^2).
\]
\end{srcBAlemma}
\begin{proof}
For $\theta=(1-t)\theta_0+t\theta_1$, linearity of the inequalities gives
\[
 P(\theta)\supseteq(1-t)P(\theta_0)+tP(\theta_1).
\]
After adding $H/2$, the same containment holds for the thickened sets, since $(1-t)H/2+tH/2=H/2$. Classical planar Brunn--Minkowski \cite{R5E3:Gardner} and monotonicity of area give the concavity of $G$. Each thickened set contains a translate of $H/2$, so $G^2\ge3/4$.

For completeness, the planar Brunn--Minkowski step has a short slice proof. For two convex bodies of area one, let $h_0,h_1$ be their vertical slice lengths and choose horizontal coordinates $x(s),y(s)$ enclosing area $s$ to their left, $0<s<1$. Almost everywhere $x'=1/h_0$, $y'=1/h_1$. The mixed body at $z(s)=(1-t)x(s)+ty(s)$ has slice length at least $(1-t)h_0+t h_1$. Integrating in $z$ and using
\[
 ((1-t)h_0+t h_1)((1-t)/h_0+t/h_1)\ge1
\]
proves that its area is at least one. Rescale bodies of arbitrary positive areas and adjust $t$ to obtain the square-root form. Our thickened bodies always have positive area, including when $P$ is degenerate.

Finally $\phi(s)=\log(s^2+1/4)$ is increasing, and
\[
 \phi''(s)=\frac{2(1/4-s^2)}{(s^2+1/4)^2}<0
 \qquad(s\ge\sqrt3/2).
\]
Thus $g=\phi\circ G$ is concave and nonnegative. Its integer values, and those of $G$, follow from Lemma~\ref{R5E3:pick}.
\end{proof}
The added hexagon is an auxiliary geometric device. No additional lattice points are inserted into the actual counted fiber. General discrete Brunn--Minkowski inequalities often enlarge the set being counted \cite{R5E3:IYZ}; that is a different conclusion.

\section{Why a nullity-two unimodular fiber has these normals}
\begin{srcBAlemma}[Complementary-minor normalization]\label{R5E3:kernel}
For a unimodular integer matrix $M$ of nullity two, there is an integer affine parametrization
\[
 x=x^0(b)+Kt,\qquad t\in\srcBAR^2,
\]
of every feasible real fiber, such that $x^0$ is linear and integral on feasible integer margins, the map identifies the fiber lattice with $\srcBAZ^2$, and, after an integer unimodular change of $t$, every row of $K$ belongs to
\[
 \{0,\pm e_1,\pm e_2,\pm(e_1+e_2)\}.
\]
\end{srcBAlemma}
\begin{proof}
Choose $r$ independent rows and $r$ columns with square submatrix $B$ of determinant $\pm1$. Project a margin to those rows by $R$. After multiplying by $B^{-1}$ and reordering columns, the equations become
\[
 [I_r\ N]x=B^{-1}Rb=:a,
 \quad x^0(b)=(a,0,0),\quad K=\binom{-N}{I_2}.
\]
Projection to the selected rows is injective on the column span of $M$, so this system is equivalent to the original for every feasible margin. For integer margins, $a$ is integral. The final two coordinates of $x$ are exactly $t$, proving the lattice identification.

Every maximal minor of $[I_r\ N]$ is $0$ or $\pm1$. Replacing one identity column by a column of $N$ shows its entries belong to $\{0,\pm1\}$; replacing two shows every two-by-two minor of $N$ is $0$ or $\pm1$. Together with the last two identity rows this proves that every two-by-two minor of $K$ is $0$ or $\pm1$. Equivalently these are the complementary maximal minors. This reasoning also covers $r=1$; there are then no two distinct rows of $N$ to consider. For $r=0$, $K=I_2$ directly.

Up to sign, the possible nonzero rows are $e_1,e_2,e_1+e_2,e_1-e_2$. The last two cannot both occur because their determinant has absolute value two. A sign change of the second coordinate if necessary leaves only the stated three directions. This is the dual use of the elementary rank-two direction argument in Round 4 entry 06, Section 4; no new general classification of regular matroids is claimed.
\end{proof}

\begin{proof}[Proof of Theorem~\ref{R5E3:main}, except sharpness]
For nullity two apply Lemma~\ref{R5E3:kernel}. Bounds on $x_j$ become
\[
 l_j-x_j^0(b)\le K_jt\le u_j-x_j^0(b).
\]
Zero rows impose only feasibility conditions on parameters. All other normals have the form required in Lemma~\ref{R5E3:extension}. Its domain includes the componentwise means because a mean of feasible real vectors satisfies the mean equations and bounds. At integer mean parameters, the resulting bounded polygon is integral and nonempty, so its integer count is positive. Concavity of $G$ gives \eqref{R5E3:sqrt}, and concavity of $g$ gives \eqref{R5E3:jensen}.

For nonnegative fibers, use only the lower bounds. A nonempty finite integer fiber forces $\ker(M)\cap\srcBAR_{\ge0}^n=\{0\}$. Otherwise that rational polyhedral cone has a nonzero rational vector, which can be scaled to a nonnegative integer kernel vector and repeatedly added to a lattice point. Hence every feasible real nonnegative fiber is bounded. The same planar argument applies.

For nullity one the corresponding parametrization has one free integer coordinate and nonzero kernel entries $\pm1$, by the one-column minor argument. A nonempty bounded real fiber is an interval $[L(\theta),U(\theta)]$ with integral endpoints at integral parameters. Its length $w=U-L\ge0$ is concave: $U$ is a minimum of affine functions and $L$ a maximum. Thus $\sqrt{w+3/4}$ and $\log(w+1)$ are concave, nonnegative functions with the required integer values. Nullity zero gives a single feasible point, integral for integer parameters, and the constant functions $\sqrt3/2$ and zero. This proves all finite-support assertions.
\end{proof}

\section{Countable support without a hidden limiting assumption}
\begin{proof}[Proof of Corollary~\ref{R5E3:countable}]
In the boxed setting the real feasible parameter domain consists of triples $(b,l,u)$ admitting $Mx=b$, $l\le x\le u$. It is a closed convex polyhedral set: it is the linear projection of a polyhedron, equivalently the set obtained by finite elimination of the variables $x$. Every fiber is bounded. In the nonnegative setting, the domain is the finitely generated cone $M\srcBAR_{\ge0}^n$, also closed and convex, and every feasible fiber is bounded by the preceding recession argument. Restrict to the column span when needed.

Let $\Theta$ have the given countable distribution. Its expectation and every conditional tail expectation belong to this domain. For example this follows by taking expectations in each of its finitely many defining affine inequalities; the finite first moment makes all of those expectations well-defined.

Choose any finite set of support atoms, of total mass $1-\varepsilon$. If $\varepsilon>0$, replace all remaining atoms by their conditional mean $\theta_*$. Finite concavity, valid for real parameters by Lemma~\ref{R5E3:extension} and the interval argument, gives
\[
 G(E\Theta)\ge\sum_{i\text{ retained}}\mu_iG(\theta_i)+\varepsilon G(\theta_*),
 \qquad
 g(E\Theta)\ge\sum_{i\text{ retained}}\mu_i g(\theta_i)+\varepsilon g(\theta_*).
\]
All terms are nonnegative, so discard the last terms. Exhaust the countable support by finite sets. Monotone convergence of these sums proves both asserted inequalities and proves their right sides finite. This avoids assuming uniform integrability from a mere family of finite truncations with changing means. No assertion is made for conditionally convergent coordinate sums presented as a ``mean.''
\end{proof}

\section{The optimal quarter and an explicit new class}\label{R5E3:sharp}
Consider the fixed matrix
\begin{equation}\label{R5E3:exampleM}
 M=\begin{pmatrix}
 1&0&1&0&0&0\\
 0&1&0&1&0&0\\
 1&1&0&0&-1&0\\
 0&0&0&0&1&1
 \end{pmatrix}.
\end{equation}
Write its variables as $(x,y,s,t,e,f)$ and its margin as $(a,b,c,d)$. Its nonnegative fibers correspond exactly to
\begin{equation}\label{R5E3:rectangle}
 0\le x\le a,\quad 0\le y\le b,\quad c\le x+y\le c+d,
 \quad(s,t,e,f)=(a-x,b-y,x+y-c,c+d-x-y).
\end{equation}
It has rank four, for example using columns $s,t,e,f$. Multiplying rows three and four by $-1$ makes it a reduced directed incidence matrix: the six edges form the three two-edge paths between vertices $0$ and $3$, through $1,2,4$. Thus its underlying graph is $K_{2,3}$.

For clarity, reduced directed incidence matrices are totally unimodular: in any square submatrix, a column with at most one nonzero entry permits determinant expansion and induction; if every column has two entries, their opposite signs make the sum of the rows zero, giving determinant zero. This proves total unimodularity here without relying on numerical minors.

At margins
\[
 b_k=(2k,2k,k,2k),\qquad k\in\srcBAZ_{\ge0},
\]
translation by $(k,k)$ identifies the polygon \eqref{R5E3:rectangle} with $kH$. Pick's formula, or direct row counting, gives
\begin{equation}\label{R5E3:hexcounts}
 F_+(b_k)=3k^2+3k+1,\qquad
 \sqrt{F_+(b_k)-\tfrac14}=\sqrt3(k+\tfrac12).
\end{equation}
Thus the corrected inequality is an equality along this whole ray. A genuinely noncollinear example uses the three margins $(4,3,2,3)$, $(3,4,2,3)$ and $(2,2,2,3)$. Their equal-weight mean is $(3,3,2,3)$; direct counting in \eqref{R5E3:rectangle} gives respective counts $14,14,6$ and $12$, so $12^3>14\cdot14\cdot6$.

\begin{srcBAproposition}[Sharpness and failure of the uncorrected power]\label{R5E3:optimal}
No constant $c<1/4$ can replace $1/4$ in Theorem~\ref{R5E3:main}. In particular ordinary square-root concavity of unimodular fiber counts is false, even at nullity two.
\end{srcBAproposition}
\begin{proof}
For $c<1/4$, the function $q_c(k)=\sqrt{3k^2+3k+1-c}$ satisfies
\[
 q_c''(k)=\frac{3-12c}{4(3k^2+3k+1-c)^{3/2}}>0\qquad(k\ge0).
\]
Hence $q_c(1)<(q_c(0)+q_c(2))/2$. These are actual fibers of the same matrix, with $b_1=(b_0+b_2)/2$. At $c=0$ their counts are $1,7,19$, and
\[
 \sqrt7<\frac{1+\sqrt{19}}2.
\]
This does not contradict log-concavity: $7^2\ge1\cdot19$. It only rules out removing the optimal correction in the stronger power inequality.
\end{proof}

The graph $K_{2,3}$ is not outerplanar. Indeed, if all five vertices had a disk embedding on its boundary, the three internally disjoint paths between its two degree-three vertices could not all have their middle vertices on that boundary: the union of any two paths separates the third path onto one side, and the three paths bound two bounded regions in a plane embedding, making one middle vertex non-exterior. Equivalently this is the elementary forbidden $K_{2,3}$ obstruction to outerplanarity. Our statement varies all four independent balances, not merely a selected boundary-terminal subset.

\section{Limitations, prior art, and reproducibility}
The theorem is a positive answer to the sourced barycentric question on an arbitrary-rank, low-nullity class, with a sharp stronger power statement. It does not address nullity three in general; in particular the initial $K_4$ search target has nullity three and remains outside the proof. It does not assert that every unimodular polytope of actual dimension two arising inside a higher-nullity ambient fiber is covered without further reduction. Arbitrary positive monomial weights are not covered by the unweighted area identity.

Pick's theorem, planar Brunn--Minkowski, and the elementary normal-direction classification are classical ingredients. The original matrix-fiber barycentric target and the distinction from midpoint concavity belong to \cite{R4E6:LMY}. The preceding tournament's rank-two and disk-network applications belong to entry 06 \cite{R5E3:R406}. Our proposed increment is the complementary low-nullity application, the sharp $1/4$ power correction, simultaneous capacity variation, and the countable extension. Targeted primary-source checks on 20 October 2026 identify lattice Brunn--Minkowski \cite{R5E3:IYZ} as related context and Chu--Karshon \cite[Theorem 5.8]{R5E3:CK} as a direct prior generalization of Lemma~\ref{R5E3:pick}. The counting identity is not a new contribution. Historical priority for the matrix-fiber application and sharp power statement has not been established. No claim of exhaustive novelty verification is made.

The exact verifier uses rational polygon intersections, convex hulls and shoelace areas to test \eqref{R5E3:areaidentity}, including segments and points. It verifies complementary minors, all minors of \eqref{R5E3:exampleM}, finite noncollinear barycentric count products, the sharpness family, and the gap in the uncorrected power inequality. Rational square-root comparisons are performed by exact squaring when required; no floating eigenvalue or numerical concavity test is a proof premise.

Before finding the area argument, bounded numerical convex-envelope searches checked 120 $K_4$ boxed systems and 48 additional $K_5$ or $K_{3,3}$ boxed systems, some with positive weights. They found no counterexample. These failures to find a violation did not prove any new class; the larger cases remain unproved here. The scripts and logs are separate from the exact verifier. No external posting, submission, or author contact occurred.

\begingroup\small

\endgroup

\chapter[Sharp Power Concavity for the A\_3 Partition Extension and Simple Rank-Three Unimodular Fibers]{Sharp Power Concavity for the $A_3$ Partition Extension and Simple Rank-Three Unimodular Fibers}
\label{source:R5E6}
\noindent\textit{Source development: Round 5, entry 06. Original mathematical arguments retained; notation is local to this chapter.}\par\medskip


\begin{quote}\small\noindent\textbf{Chapter abstract.} 
The classical type-$A_3$ Kostant partition function has seven cubic polynomial chambers. We prove that their continuous extension to the nonnegative orthant has concave power $4/13$, and that no larger positive power is concave for this particular extension. This supplies a power-mean strengthening of finite-support barycentric log-concavity for the complete four-vertex acyclic flow network. Cube-root concavity fails even at integer midpoints: the relevant counts are $1,13,51$. Combining the new complete-network case with the attributed outerplanar result from Round 4, we prove the barycentric inequality for every simple rank-three unimodular system, where simple means nonzero, pairwise nonproportional columns. A small explicit classification certificate makes the reduction self-contained. Repeated columns, arbitrary weighted or boxed rank-three fibers, and general higher rank remain outside the new theorem. Sharpness is asserted for the canonical continuous extension, not for the optimal exponent of a discrete interpolation problem; historical priority of these applications remains unestablished.
\end{quote}

\section{Statements, target, and scope}
Let $K(a,b,c)$ be the number of nonnegative integer solutions of
\begin{equation}\label{R5E6:eq:A3}
 (a,b,c)=x_1(1,0,0)+x_2(0,1,0)+x_3(0,0,1)
       +u(1,1,0)+v(0,1,1)+z(1,1,1).
\end{equation}
This is the type-$A_3$ Kostant partition function in simple-root coordinates. Its support is exactly $\srcBDZ_{\geq0}^3$. Write $\mathcal K$ for its standard chamber-polynomial extension, whose formulas are derived in Section 2.

\begin{srcBDtheorem}[Sharp power for the chamber extension]\label{R5E6:thm:power}
The continuous function $\mathcal K^{4/13}$ is concave on all of $\srcBDR_{\geq0}^3$. For every $p>4/13$, $\mathcal K^p$ is not concave on that orthant. Thus $4/13$ is the largest positive concavity exponent of this specified continuous extension.
\end{srcBDtheorem}

\begin{srcBDcorollary}[Power-mean and barycentric inequalities]\label{R5E6:cor:baryA3}
For $a^{(i)}\in\srcBDZ_{\geq0}^3$, real $\mu_i>0$ with $\sum_i\mu_i=1$, and integer mean $\bar a=\sum_i\mu_i a^{(i)}$,
\begin{equation}\label{R5E6:eq:powermean}
 K(\bar a)^{4/13}\geq\sum_i\mu_i K(a^{(i)})^{4/13},
 \qquad K(\bar a)\geq\prod_iK(a^{(i)})^{\mu_i}.
\end{equation}
Only finitely many input points are considered. The optimal exponent for the discrete inequality is not determined here. In particular it is not $1/3$: at the midpoint of $(0,0,0)$ and $(4,6,4)$ the counts are $1,13,51$ and
\begin{equation}\label{R5E6:eq:cubefail}
 2\sqrt[3]{13}<1+\sqrt[3]{51}.
\end{equation}
\end{srcBDcorollary}

Call an integer matrix $M$ of rank $r$ \emph{unimodular} if every nonsingular $r\times r$ submatrix has determinant $\pm1$. Call it \emph{simple} if its columns are nonzero and no two columns are proportional over $\srcBDR$. Put
\[
 N_M(b)=\#\{x\in\srcBDZ_{\geq0}^n:Mx=b\}.
\]
\begin{srcBDtheorem}[Simple rank-three systems]\label{R5E6:thm:rank3}
Let $M\in\srcBDZ^{m\times n}$ be simple and unimodular with rank three; redundant rows are allowed. For any finite collection of margins $b^{(i)}$ with nonempty finite fibers, and positive real probabilities with integer mean $\bar b$, one has
\begin{equation}\label{R5E6:eq:jensen}
 N_M(\bar b)\geq\prod_iN_M(b^{(i)})^{\mu_i}.
\end{equation}
The central fiber is also finite and nonempty. Equivalently, all finite nonnegative flow fibers of simple directed graphs on four vertices satisfy the finite-support barycentric inequality. The power-mean strengthening \eqref{R5E6:eq:powermean} is claimed for the complete acyclic graph, not for every graph in this theorem.
\end{srcBDtheorem}

\paragraph{Precisely sourced target.}
Leake--Mohammadi Yekta \cite[Question 1.7]{R4E6:LMY} ask for \eqref{R5E6:eq:jensen} for arbitrary unimodular matrices. Their Theorems 1.4 and 3.1 establish midpoint/line log-concavity, including weighted box versions, and they explicitly distinguish the barycentric question. The latest posted version checked on 20 October 2026 is v1, 30 September 2026. Our result is a restricted rank-three answer, not a resolution of that question. Theorem \ref{R5E6:thm:power} is stronger than the needed logarithmic conclusion on the complete four-vertex system.

\paragraph{Prior work and exact increment.}
The chamber formulas are classical, not a discovery here: Billey--Guillemin--Rassart \cite[Section 1.4]{R5E6:BGR} describe the seven $A_3$ chambers and credit De Loera--Sturmfels' earlier computations. We rederive the formulas to make the proof independent of tables. Khare--Matherne--St. Dizier \cite[Definition 1.2 and Theorem 1.6]{R5E6:KMS} prove root-direction coefficient log-concavity and normalized generating-function log-concavity for much broader flow classes. The latter concerns the generating function evaluated at positive variables, not a real interpolation of its lattice-indexed coefficients, and is not our power-concavity assertion.

Round 4 entry 06 \cite[Theorem 1.1 and Corollary 1.3]{R5E6:R4} proved the barycentric inequality for disk networks with only boundary balances varying, hence for outerplanar graphs, and separately for rank at most two. Its marginal-closure input was explicitly attributed to Speyer \cite[Lemma 5.3]{R4E6:Speyer} and Ahlswede--Daykin \cite{R5E6:AD}. Here the genuinely additional case is $K_4$ with \emph{all} balances varying: $K_4$ is not outerplanar. The Hessian and interface analysis, the exact exponent for the specified extension, and the rank-three application are the proposed increments. The underlying small unimodular classification is not claimed as a new structural classification. Priority of the application and power result is unestablished; a targeted search is not a historical novelty certificate. All frozen Round 4 papers and the assessment were read; no other entrant's result is a proof dependency.

\section{The classical seven-chamber extension}
Eliminating $x_1,x_2,x_3$ from \eqref{R5E6:eq:A3} gives
\begin{equation}\label{R5E6:eq:count}
 K(a,b,c)=\#\{(u,v,z)\in\srcBDZ_{\geq0}^3:
 u+z\leq a,\ v+z\leq c,\ u+v+z\leq b\}.
\end{equation}
The expression is symmetric in $a,c$. It suffices to give the formulas when $0\leq a\leq c$:
\begin{equation}\label{R5E6:eq:pieces}
 6\mathcal K(a,b,c)=
 \begin{cases}
 (b+1)(b+2)(b+3),&0\leq b\leq a,\\
 (a+1)(a+2)(3b-2a+3),&a\leq b\leq c,\\
 (a+1)(a+2)(3c-a+3)-d(d+1)(d+2),
       &c\leq b\leq a+c,\ d=a+c-b,\\
 (a+1)(a+2)(3c-a+3),&b\geq a+c.
 \end{cases}
\end{equation}
Reflect the formulas across $a=c$ on the other half of the orthant. The first chamber is shared by both halves, giving seven full-dimensional chambers in total.

\begin{srcBDlemma}\label{R5E6:lem:formula}
These polynomials agree on their common boundaries, are positive throughout the orthant, and equal $K$ at every nonnegative integer point.
\end{srcBDlemma}
\begin{proof}
If $b\leq a,c$, the only effective constraint in \eqref{R5E6:eq:count} is $u+v+z\leq b$, giving $\binom{b+3}{3}$. If $a\leq b\leq c$, group by $t=u+z$. For each $0\leq t\leq a$ there are $t+1$ choices of $(u,z)$ and $b-t+1$ choices of $v$. Summing gives the second polynomial.

If $b\geq a+c$, its constraint is redundant. Grouping by $z$ gives
\[
 \sum_{z=0}^{a}(a-z+1)(c-z+1)
 =\frac{(a+1)(a+2)(3c-a+3)}6.
\]
In the remaining case $a\leq c\leq b\leq a+c$, put $d=a+c-b$. A point counted by this last sum violates the remaining constraint exactly when, on writing $u'=a-z-u$ and $v'=c-z-v$, one has
\[
 z+u'+v'\leq d-1.
\]
Because $0\leq d\leq a\leq c$, these three nonnegative variables automatically obey the original upper constraints. The violating count is $\binom{d+2}{3}$, including zero when $d=0$. This proves the third formula.

Boundary agreement follows by substitution at $b=a$, $b=c$, $b=a+c$, and $a=c$. All product pieces are positive. In the third piece, $H(t)=t(t+1)(t+2)$ is increasing for $t\geq0$, so its numerator is at least
\[
 (a+1)(a+2)(3c-a+3)-H(a)
 =(a+1)(a+2)(3c-2a+3)>0.
\]
The same conclusions hold after reflection.
\end{proof}

The formulas also specify uniquely the polynomial on each full-dimensional chamber from its integer values: a polynomial zero on the integer points of an open rational cone is zero, for example by restricting to arbitrarily large translated integer boxes in that cone. Thus ``the chamber extension'' denotes a fixed classical interpolation, not an interpolation chosen to optimize a concavity exponent.

\section{Power concavity inside the chambers}
We prove the exponent $p=4/13$. Multiplication by a positive constant does not affect power concavity. In each product chamber, the numerator is a product of three positive affine functions. Such a product $G=L_1L_2L_3$ satisfies $\srcBDHess(G^p)\preceq0$ whenever $0<p\leq1/3$. Indeed, in any direction, put $t_i=D L_i/L_i$. Then
\[
 \frac{D^2G^p}{G^p}
 =p^2\Bigl(\sum_i t_i\Bigr)^2-p\sum_i t_i^2
 \leq p(3p-1)\sum_i t_i^2\leq0.
\]
The affine functions are positive on the entire chamber being used, including when one coefficient in their formulas is negative.

The only remaining chamber is
\[
 0\leq d\leq a\leq c,\qquad d=a+c-b.
\]
The change from $(a,b,c)$ to $(a,c,d)$ is invertible and linear. Put
\begin{gather}\label{R5E6:eq:central}
 A=(a+1)(a+2),\quad H=d(d+1)(d+2),\quad
 F=A(3c-a+3)-H=6\mathcal K,\\
 h=1-\frac a3-\frac{H}{3A},\qquad
 \tau=c+h=\frac{F}{3A}>0.\notag
\end{gather}
Primes on $A$ denote $a$-derivatives, and primes on $H$ denote $d$-derivatives.

\begin{srcBDlemma}[Central-chamber Hessian]\label{R5E6:lem:hessian}
$F^{4/13}$ is concave on the central chamber.
\end{srcBDlemma}
\begin{proof}
For $0<p<1$, negative semidefiniteness of $\srcBDHess(F^p)$ is equivalent to positive semidefiniteness of
\[
 Q=-\srcBDHess\log F-p(\nabla\log F)(\nabla\log F)^{\mathsf T},
\]
since $-\srcBDHess(F^p)=pF^pQ$. As $\log F=\log(3A)+\log(c+h)$, the $c,c$ entry of $Q$ is $(1-p)/\tau^2>0$. Its Schur complement on the $(a,d)$ coordinates is
\begin{equation}\label{R5E6:eq:schur}
 S=\begin{pmatrix}J/A^2&0\\0&0\end{pmatrix}
       -\frac{\srcBDHess_{a,d}h}{\tau},
 \qquad J=(A'^2-AA'')-\frac{p}{1-p}A'^2.
\end{equation}
For example, the cross entries before taking the complement are
$Q_{i,c}=(1-p)h_i/\tau^2-p(\log A)_i/\tau$; their outer product cancels the corresponding terms containing $h_i h_j$ and $(\log A)_i h_j$, leaving \eqref{R5E6:eq:schur}.

At $p=4/13$, write $C=2a^2+6a+9$. Then
\[
 J=\frac C9>0,\quad
 h_{aa}=-\frac{H(2A'^2-AA'')}{3A^3}\leq0,\quad
 h_{dd}=-\frac{H''}{3A}<0,\quad
 h_{ad}=\frac{H'A'}{3A^2}.
\]
Here $A''=2$, $H''=6(d+1)$, and $2A'^2-AA''>0$. Both diagonal entries of $S$ are positive. Direct substitution in its determinant gives
\begin{equation}\label{R5E6:eq:detS}
 \det S=\frac{\mathcal N}{27A^4\tau^2},\qquad
 \mathcal N=2(d+1)AC(3c-a+3)-A'^2\bigl((d+1)^4+8(d+1)^2+3\bigr).
\end{equation}
One way to verify the simplification is to use $H=t^3-t$, $H'=3t^2-1$, $H''=6t$, where $t=d+1$.

Set $u=a+1$ and $t=d+1$, so $1\leq t\leq u$. Because $c\geq a$,
\[
 3c-a+3\geq2a+3=A'=2u+1,
 \quad A=u(u+1),\quad C=2u^2+2u+5.
\]
The function $\phi(t)=t^3+8t+3/t$ increases for $t\geq1$, since
$\phi'(t)=3t^2+8-3/t^2>0$. Hence
\begin{align*}
 \mathcal N
 &\geq (2u+1)t\left[2u(u+1)(2u^2+2u+5)
                    -(2u+1)\phi(t)\right]\\
 &\geq\frac{(2u+1)t}{u}
       (u-1)(2u^4+9u^3+7u^2+9u+3)\geq0.
\end{align*}
The final equality underlying this bound is the elementary factorization
\[
 2u^2(u+1)(2u^2+2u+5)-(2u+1)(u^4+8u^2+3)
 =(u-1)(2u^4+9u^3+7u^2+9u+3).
\]
Thus $S\succeq0$ and $Q\succeq0$, proving the lemma, including the boundary by continuity. In the strict interior $a>0$ the last lower bound is positive.
\end{proof}

\section{Interfaces, global concavity, and sharpness}
Piecewise concavity does not by itself give global concavity; we check every interface. In the following table, $q=\mathcal K$, $n$ is the gradient of the displayed linear form, and the interface is crossed in the direction of increasing that form. The table gives
$\nabla q_{\rm after}-\nabla q_{\rm before}=-\kappa n$.
\begin{center}
\begin{tabular}{lll}
\toprule
Interface and region & Increasing linear form & $\kappa$\\
\midrule
$b=a$, $a<c$ & $b-a$ & $(3a+5)/6$\\
$b=c$, $a<c$ & $b-c$ & $(3a+4)/6$\\
$b=a+c$ & $b-a-c$ & $1/3$\\
$a=c$, $b>a$ & $a-c$ & $1/6$\\
\bottomrule
\end{tabular}
\end{center}
Also include the first two rows with $a,c$ interchanged. The formulas follow by differentiating \eqref{R5E6:eq:pieces} and substituting the facet equation. Across $a=c$ with $b<a$, the formula is unchanged, so the jump is zero. This exhausts all full-dimensional shared facets of the seven chambers.

Every $\kappa$ is nonnegative, and $q$ is continuous and positive. Thus the gradient jump of $q^p$ is $-p q^{p-1}\kappa n$, also in the concave direction. On a generic line segment in the open orthant, avoiding intersections of two chamber walls, the directional derivative of $q^{4/13}$ is nonincreasing within each chamber by Section 3 and at every crossing by the table. The restriction is concave. Arbitrary segments in the closed orthant are limits of such generic interior segments. Continuity then proves concavity on the whole orthant, not only chamber interiors. This proves the positive assertion of Theorem \ref{R5E6:thm:power}.

\begin{proof}[Sharpness for the specified continuous extension]
Use the central polynomial $F$ of \eqref{R5E6:eq:central} in the coordinates $(a,c,d)$. For any $p>0$, negative semidefiniteness of $\srcBDHess(F^p)$ requires
\[
 R_p=(1-p)(\nabla F)(\nabla F)^{\mathsf T}-F\srcBDHess F\succeq0.
\]
At $(a,c,d)=(0,0,0)$ one has $F=6$, $\nabla F=(7,6,-2)^{\mathsf T}$, and
\[
 \srcBDHess F=\begin{pmatrix}0&9&0\\9&0&0\\0&0&-6\end{pmatrix},
 \qquad \det R_p=11664(4-13p).
\]
For $p>4/13$ this determinant is negative. By continuity it remains negative at sufficiently small points $(a,c,d)=(2\epsilon,3\epsilon,\epsilon)$ in the strict central chamber. A small open ball about such a point remains in that chamber, so the Hessian test for concavity applies there and fails. This proves the sharpness assertion. It does not prove an optimal exponent for the values on the discrete lattice or for other continuous interpolations.
\end{proof}

\begin{proof}[Corollary \ref{R5E6:cor:baryA3} and the cube-root obstruction]
Apply ordinary finite Jensen to the concave function $\mathcal K^{4/13}$ and then weighted arithmetic--geometric mean to its positive values. This proves \eqref{R5E6:eq:powermean}. Formula \eqref{R5E6:eq:pieces} gives the counts $1,13,51$. To verify \eqref{R5E6:eq:cubefail} without decimal approximations, put $s=\sqrt[3]{13}>47/20$, since $47^3<13\cdot20^3$. The polynomial $103-12s^2+6s$ decreases for $s>1/4$, so
\[
 (2s-1)^3=103-12s^2+6s
 <103-12(47/20)^2+6(47/20)=5083/100<51.
\]
Taking cube roots proves the strict inequality.
\end{proof}

\section{A complete small unimodular reduction}
We next prove Theorem \ref{R5E6:thm:rank3}. The structural classification in this section is elementary finite linear algebra; its role is a reduction, not a claimed new matroid theorem.

Choose three independent rows of $M$, with row projection $R$, and three columns $J$ such that $B=RM_J$ is nonsingular. Then $\det B=\pm1$ and
\[
 N=B^{-1}RM
\]
contains the three standard basis columns. The map $b\mapsto B^{-1}Rb$ is a lattice isomorphism from $\operatorname{span}_{\srcBDR}(M)\cap\srcBDZ^m$ to $\srcBDZ^3$, with inverse $t\mapsto M_Jt$. Indeed $R$ is injective on the column span, and $B^{-1}$ is integral. It preserves fibers and integer means, so no ambient full-row-rank assumption is hidden here.

Every entry of $N$ belongs to $\{0,\pm1\}$ by Cramer's rule. Ignoring column signs temporarily, choose the representative whose first nonzero coordinate is $+1$. Apart from the three basis directions, the ten candidates are
\begin{equation}\label{R5E6:eq:candidates}
\begin{array}{c|ccccc}
 i&1&2&3&4&5\\\hline
 v_i&(0,1,-1)&(0,1,1)&(1,-1,-1)&(1,-1,0)&(1,-1,1)\\[3pt]
 i&6&7&8&9&10\\\hline
 v_i&(1,0,-1)&(1,0,1)&(1,1,-1)&(1,1,0)&(1,1,1)
\end{array}
\end{equation}
Simplicity ensures that each listed direction occurs at most once.

\begin{srcBDlemma}[Rank-three direction certificate]\label{R5E6:lem:classification}
Every simple unimodular system of rank three is, after an integral invertible change of coordinates, a subset of the six directions
\begin{equation}\label{R5E6:eq:roots}
 e_1,e_2,e_3,e_1-e_2,e_1-e_3,e_2-e_3,
\end{equation}
with arbitrary individual signs.
\end{srcBDlemma}
\begin{proof}
Here is a complete, small certificate. A pair $v_i,v_j$ can coexist with all three basis columns only for the following 24 pairs; checking their determinants with the basis columns gives this list:
\[
\begin{split}
&14,15,16,17,18,19,\quad23,24,26,27,29,2\!:\!10,\\
&34,36,45,46,47,57,68,69,79,7\!:\!10,89,9\!:\!10.
\end{split}
\]
Here $ij$ denotes the unordered pair $\{i,j\}$, and a colon separates indices involving 10. Among triples whose three pairs appear in the list, exactly four have determinant of absolute value two:
\[
 \{1,4,7\},\quad\{1,6,9\},\quad\{2,4,6\},\quad\{2,7,9\}.
\]
The only four-element cliques in the displayed pair graph are
\[
 \{1,4,5,7\},\quad\{1,6,8,9\},\quad
 \{2,3,4,6\},\quad\{2,7,9,10\};
\]
each contains one of those forbidden triples. Thus at most three candidates can be added. The 16 permissible triples and a change-of-basis certificate for each appear in Table \ref{R5E6:tab:systems}. Every permissible pair and singleton is contained in one of them, as the table also verifies. These observations cover all $2^{10}=1024$ candidate subsets, not a sampled list.

For a row of the table let $S$ consist of the standard basis and its three listed candidates, and let $T$ be the matrix with the three displayed basis vectors as columns. Direct three-by-three determinants give $|\det T|=1$, and the six vectors $T^{-1}S$, after normalizing their individual signs, are exactly \eqref{R5E6:eq:roots}. All input directions are contained in some such $S$, proving the lemma. These are ten-vector and three-by-three arithmetic checks, also fully reproduced by the accompanying verifier.
\end{proof}

\begin{table}[ht]
\centering
\begin{tabular}{cc@{\qquad}cc}
\toprule
Added indices & Columns of $T$ & Added indices & Columns of $T$\\
\midrule
$2,3,4$&$-e_1,-e_2,-v_2$&$1,6,8$&$-e_1,-e_3,v_1$\\
$1,4,5$&$-e_1,-e_2,-v_1$&$2,6,9$&$-e_1,e_2,-e_3$\\
$2,3,6$&$-e_1,-e_3,-v_2$&$1,7,9$&$-e_1,e_2,e_3$\\
$1,4,6$&$-e_1,-e_2,-e_3$&$1,8,9$&$-e_1,e_2,v_1$\\
$3,4,6$&$-e_1,-e_2,-v_6$&$6,8,9$&$-e_1,e_2,-v_6$\\
$2,4,7$&$-e_1,-e_2,e_3$&$2,7,10$&$-e_1,e_3,v_2$\\
$1,5,7$&$-e_1,e_3,-v_1$&$2,9,10$&$-e_1,e_2,v_2$\\
$4,5,7$&$-e_1,-e_2,-v_7$&$7,9,10$&$-e_1,e_2,-v_7$\\
\bottomrule
\end{tabular}
\caption{All maximal candidate sets and their unimodular root-system certificates.}\label{R5E6:tab:systems}
\end{table}

\paragraph{Signs are orientations, not variable reflections.}
We have normalized signs only to classify directions. In the original nonnegative fiber no variable is replaced by its negative. Rather, each signed vector in \eqref{R5E6:eq:roots} is the column of a reduced incidence matrix on vertices $\{1,2,3,4\}$: $e_i$ represents $i\to4$, $e_i-e_j$ represents $i\to j$, and a negative sign reverses that edge. Simplicity gives a simple directed graph, and rank three makes its underlying graph connected.

If this graph had a directed cycle, adding any nonnegative integer multiple of its circulation to one feasible flow would make that fiber infinite. The nonempty finite input hypothesis therefore forces an acyclic orientation. Conversely an acyclic directed graph has bounded nonnegative fibers: in a topological order, each prefix cut has fixed nonnegative total flow, and every edge crosses such a cut. This bounds every coordinate.

If all six edges occur, an acyclic orientation of $K_4$ is transitive. In a topological ordering its positive roots become exactly \eqref{R5E6:eq:A3} after the integral change from vertex balances $(\beta_1,\ldots,\beta_4)$ to prefix sums
\[
 a=\beta_1,\qquad b=\beta_1+\beta_2,\qquad
 c=\beta_1+\beta_2+\beta_3,
 \qquad\sum_i\beta_i=0.
\]
Thus Corollary \ref{R5E6:cor:baryA3} applies. If an edge is absent, the underlying graph is outerplanar: $K_4$ minus any edge has an embedding as a quadrilateral with one diagonal, and every proper subgraph lies in such a graph. The attributed outerplanar barycentric result, proved again in the next section, applies. This completes the reduction for Theorem \ref{R5E6:thm:rank3}.

\section{The inherited outerplanar case, with proof}
We include the exact earlier input to keep the rank-three theorem standalone. No new credit is claimed for this section.

\begin{srcBDlemma}[Outerplanar fibers, Round 4]\label{R5E6:lem:outer}
For an outerplanar directed graph, the counts of integer flows in a fixed finite integer box satisfy finite-support barycentric log-concavity in all vertex balances. The same holds for finite nonnegative fibers.
\end{srcBDlemma}
\begin{proof}
An isolated vertex contributes one at balance zero and zero otherwise, so it may be removed. It is enough to treat a connected loopless graph embedded in a disk with all $k\geq2$ vertices on its boundary; disconnected components multiply. Add a counterclockwise boundary cycle, allowing two arcs if $k=2$. For vertex divergences $\beta$ summing to zero prescribe its edge flows as
\[
 p_i=-\sum_{j=1}^i\beta_j\quad(i<k),\qquad p_k=0.
\]
The augmented flow is a circulation, and $\beta\leftrightarrow(p_1,\ldots,p_{k-1})$ is an integral linear bijection. Distinct new boundary arcs have distinct bounded faces on their left: retain a spanning tree of the old graph; after adding the boundary cycle Euler's formula gives $k$ bounded faces, each needing a boundary arc since the tree has no cycle. Reinstating other edges splits, but never merges, faces.

An integral circulation is uniquely the difference $z_{L(e)}-z_{R(e)}$ of integral face potentials with exterior potential zero. To see existence, decompose it into signed directed simple cycles and add their signed interior-face indicators. Uniqueness follows from connectivity of face adjacency. Bridges have the same face on both sides and circulation zero, so cause no exception. The distinct boundary faces therefore have prescribed potentials $p_i$, including a second zero anchor at $p_k$. All other potentials are summed. The original edge box gives integer difference bounds between potentials. For fixed prescribed potentials there are finitely many allowed assignments, by paths in the connected dual from the exterior.

Let $Z(p)$ denote the number of assignments in any such difference-bound system. The unconstrained potential set is translation-invariant and closed under coordinate maximum and minimum. Thus for two feasible full vectors $z,w$ and integer $\alpha\geq0$,
\[
 z'=(z-\alpha\srcBDone)\vee w,\qquad
 w'=z\wedge(w+\alpha\srcBDone)
\]
are feasible and preserve their sum; zero anchors remain zero. The classical Ahlswede--Daykin four-functions theorem \cite{R5E6:AD}, applied to the summed coordinates with the first set shifted by $-\alpha\srcBDone$, gives
\begin{equation}\label{R5E6:eq:shift}
 Z(p)Z(r)\leq Z((p-\alpha\srcBDone)\vee r)\,
                    Z(p\wedge(r+\alpha\srcBDone)).
\end{equation}
Explicitly, for feasibility indicators $W$, use functions
$A(u)=W(p,u+\alpha\srcBDone)$, $B(v)=W(r,v)$,
$C(t)=W((p-\alpha\srcBDone)\vee r,t)$ and
$D(t)=W(p\wedge(r+\alpha\srcBDone),t+\alpha\srcBDone)$.
Their pointwise four-functions inequality is exactly feasibility of $z',w'$. The theorem compares their finite sums and yields \eqref{R5E6:eq:shift}. This is the marginal-closure mechanism already attributed in Round 4 to Speyer \cite{R4E6:Speyer} and four-functions.

For $N$ input integer potential vectors with integer average, if some coordinate range is at least two, choose ordered input vectors $x,y$ differing by at least two there and use \eqref{R5E6:eq:shift} with $\alpha=1$. Their sum is unchanged and the product of their positive counts cannot decrease. In each coordinate with difference $d=x_i-y_i\leq1$ the outputs simply interchange the values. When $d\geq2$ they become $x_i-1,y_i+1$, decreasing the squared sum by $2(d-1)$. Hence the integer energy about the fixed mean strictly decreases. Eventually all coordinate ranges are at most one; the integer average then forces every vector to equal that mean. This proves Jensen for equal weights and hence for rational weights by repetition. For real finite weights, rational points are dense in the probability polytope defined by the fixed integer mean, since its equations are rational. Continuity of the finite weighted sum of logarithms gives the result.

The integral transformation from balances to $p$ proves the boxed assertion. For finite nonnegative input and central fibers, a sufficiently large common box contains them all. Alternatively, increasing the common boxes and taking limits proves the assertion directly and forces a positive central count. In the acyclic graphs used above all fibers are bounded, as already proved.
\end{proof}

\begin{proof}[Completion of Theorem \ref{R5E6:thm:rank3}]
The relative lattice change, finite direction classification, and inherited signs identify the fibers with those of a connected acyclic simple graph on four vertices. In the six-edge case use Corollary \ref{R5E6:cor:baryA3}; otherwise use Lemma \ref{R5E6:lem:outer}. All coordinate changes preserve integer means and the exact counts. These inequalities give positivity of the central count, and acyclicity gives its finiteness.
\end{proof}

\section{Limits, verification, and further questions}
Theorem \ref{R5E6:thm:rank3} excludes repeated proportional columns. Repeated columns correspond to parallel graph edges and change the cubic counting formulas; the present proof does not remove them. Arbitrary positive monomial weights or coordinate upper bounds on the complete $K_4$ system are also not covered by the new result. The earlier rank-two and outerplanar weighted-box theorems remain valid, but must not be silently promoted to weighted or boxed rank three. Nothing here resolves the full higher-rank question or asserts a new capacity approximation algorithm.

The continuous exponent $4/13$ is exact for the displayed chamber-polynomial extension. Discrete midpoint failure of exponent $1/3$ does not establish $4/13$ as the optimal discrete exponent. Determining that exponent, extending beyond a simple rank-three configuration, and controlling weighted or bounded $K_4$ fibers are natural remaining questions, not conclusions.

The exact verifier independently checks the seven formulas against integer enumeration, symbolic Schur-complement and determinant identities, all interface derivatives, the sharpness determinant and positive factorization, the $1,13,51$ obstruction, all 1024 candidate subsets and every printed basis certificate. It also checks noncollinear barycentric examples with integer arithmetic. Floating convex-hull searches for weighted bounded $K_4$ counts found no counterexample in 300 samples; these failed searches are preserved separately and are not premises of any theorem. The universal argument is the explicit chamber, Hessian, interface, and finite-classification proof above. No external submission or author contact occurred.



\chapter[Optimal-order corrections for power concavity of integer flows]{Optimal-order corrections for power concavity of integer flows}
\label{source:R6E3}
\noindent\textit{Source development: Round 6, entry 03. Original mathematical arguments retained; notation is local to this chapter.}\par\medskip

\begin{quote}\small\noindent\textbf{Chapter abstract.} 
The sharp constant correction for two-dimensional unimodular fibers does not extend to dimension three. For the fixed type-$A_3$ flow matrix, we prove that no sublinear additive correction, in particular no constant, can make the cube root of its integer-fiber count midpoint-concave. An explicit piecewise affine correction does work: subtracting it gives exactly the flow-polytope volume at a fixed half-root-sum shift. Classical Brunn--Minkowski yields global barycentric cube-root concavity, and the correction's coefficient is sharp from below. A product construction gives, in every dimension $d\ge3$, a fixed totally unimodular system for which corrections of order $o(\|b\|^{d-2})$ fail, whereas an explicit degree-$(d-2)$ correction succeeds. The exponent $1/d$ is optimal for the corrected counts. Classical chamber formulas and volume methods are credited; the proposal is the quantitative correction and obstruction, not a new proof of the already established $A_3$ logarithmic inequality or a resolution for general unimodular matrices.
\end{quote}

\section{The correction and the main assertions}
Let $K(a,b,c)$ count the nonnegative integer solutions of
\begin{equation}\label{R6E3:matrix}
 \begin{pmatrix}1&0&0&1&0&1\\0&1&0&1&1&1\\0&0&1&0&1&1\end{pmatrix}
 (x_1,x_2,x_3,u,v,z)^{\mathsf T}=(a,b,c)^{\mathsf T}.
\end{equation}
This rank-three matrix has nullity three and is totally unimodular. Indeed its columns are interval indicators, and this property persists after selecting any subset of rows. For each square submatrix, consecutive row differences turn it into a reduced incidence matrix: adjoining the negative sum row gives at most one $+1$ and one $-1$ in each column. Row differencing has determinant one. Incidence minors are $0,\pm1$ by expansion at a column with at most one nonzero, or by zero row sum if every column has two. This proves the claim for every minor. Its fibers are finite for $a,b,c\ge0$ and are parametrized by
\begin{equation}\label{R6E3:poly}
 P(a,b,c)=\{(u,v,z)\in\srcBGR_{\ge0}^3:u+z\le a,\ v+z\le c,\ u+v+z\le b\}.
\end{equation}
Throughout, volume is ordinary Lebesgue volume in these three free coordinates, whose integer lattice is exactly the fiber lattice.

Define a function $R$ \emph{on integer nonnegative margins}. It is symmetric in $a,c$; when $a\le c$ set
\begin{equation}\label{R6E3:correction}
 R(a,b,c)=\begin{cases}
 -(b+2)/6,&b<a,\\
 -(3b-2a+3)/24,&a\le b<c,\\
 (5a+c-4b+1)/24,&c\le b\le a+c,\\
 -(3c-a+3)/24,&b>a+c.
 \end{cases}
\end{equation}
These half-open cases specify all walls without ambiguity, including $a=b=c=0$. They are not claimed to define the appropriate real interpolation between integer points. The correction can have either sign, and $|R(a,b,c)|\le(a+b+c+3)/4$.

\begin{srcBGtheorem}[A sharp linear-order replacement]\label{R6E3:main}
Put $G(a,b,c)=K(a,b,c)-R(a,b,c)$. Then $G\ge23/24>0$, and for any finite family of nonnegative integer margins $\lambda_i$, positive probabilities $\mu_i$, and integral mean $\bar\lambda$,
\begin{equation}\label{R6E3:root}
 G(\bar\lambda)^{1/3}\ge\sum_i\mu_iG(\lambda_i)^{1/3}.
\end{equation}
The exponent $1/3$ cannot be increased. Conversely, if $E:\srcBGZ_{\ge0}^3\to\srcBGR$ satisfies
\begin{equation}\label{R6E3:sublinear}
 E(\lambda)=o(1+\|\lambda\|_1)\qquad(\|\lambda\|_1\to\infty),
\end{equation}
then $(K-E)^{1/3}$ fails the midpoint inequality at arbitrarily large positive counts. In particular no fixed constant can replace $R$ in \eqref{R6E3:root}. Moreover, for every fixed $t<1$, $(K-tR)^{1/3}$ fails that inequality. Thus the coefficient $1$ of this specific correction is sharp from below.
\end{srcBGtheorem}
The negative assertions use points where the corrected quantities are positive, so no ambiguity about roots or small counts is involved. No assertion about coefficients $t>1$ is made.

\paragraph{Target, credit and exact increment.}
Leake--Mohammadi Yekta \cite[Question 1.7]{R4E6:LMY} ask for barycentric log-concavity of nonnegative integer fibers for general unimodular matrices. Their v1 of 30 September 2026 remains the latest posted version checked on 20 October 2026. The unrestricted question is not resolved here.
Round 5 entry 06 already proved the $A_3$ logarithmic inequality and the sharp exponent $4/13$ for its specified \emph{uncorrected continuous chamber extension} \cite[Theorem 1.1 and Corollary 1.2]{R6E3:R506}. It also exhibited a failure of uncorrected cube-root concavity. Those results are not reclaimed. The present increment is a corrected power-$1/3$ inequality and a sharp obstruction to every sublinear correction. It neither establishes the optimal uncorrected discrete exponent nor proves new weighted or boxed $A_3$ logarithmic inequalities.

Our Round 5 entry 03 proved the sharp constant $1/4$ correction at nullity at most two \cite[Theorem 1.1]{R6E3:R503}. Its half-hexagon identity is a direct special case of Chu--Karshon \cite[Theorem 5.8]{R5E3:CK}, not a new lemma. The present result explains a genuine dimensional obstruction to extending that constant-correction pattern. All six frozen Round 5 papers and the assessment were read; no live Round 6 competitor work was inspected.

\section{Classical count and volume formulas}
We recall the count formulas, credited to the classical $A_3$ chamber description \cite[Section 1.4]{R6E3:BGR} and reproduced in entry 06 \cite[Lemma 2.1]{R6E3:R506}. For $0\le a\le c$, and $d=a+c-b$ in the third case,
\begin{equation}\label{R6E3:counts}
6K(a,b,c)=\begin{cases}
(b+1)(b+2)(b+3),&b\le a,\\
(a+1)(a+2)(3b-2a+3),&a\le b\le c,\\
(a+1)(a+2)(3c-a+3)-d(d+1)(d+2),&c\le b\le a+c,\\
(a+1)(a+2)(3c-a+3),&b\ge a+c.
\end{cases}
\end{equation}
For self-containment, if $b\le a,c$, only $u+v+z\le b$ matters, giving $\binom{b+3}{3}$. If $a\le b\le c$, sum $(s+1)(b-s+1)$ over $s=u+z=0,\ldots,a$. If $b\ge a+c$, sum $(a-z+1)(c-z+1)$ over $z=0,\ldots,a$. In the remaining region, subtract from this last count the violating triples: with $u'=a-z-u$ and $v'=c-z-v$, they satisfy $z+u'+v'\le d-1$ and number $\binom{d+2}{3}$. Substitution verifies agreement on every common boundary.

Write $V(A,B,C)=\srcBGvol P(A,B,C)$ for real nonnegative parameters. The corresponding homogeneous formulas, for $A\le C$, are
\begin{equation}\label{R6E3:volumes}
6V(A,B,C)=\begin{cases}
B^3,&B\le A,\\
A^2(3B-2A),&A\le B\le C,\\
A^2(3C-A)-(A+C-B)^3,&C\le B\le A+C,\\
A^2(3C-A),&B\ge A+C.
\end{cases}
\end{equation}
They follow by the same slicing integrals. The first case is a simplex of volume $B^3/6$; the second is $\int_0^A s(B-s)\,ds$; the last is $\int_0^A(A-z)(C-z)\,dz$. In the third subtract a simplex of side $A+C-B$. Thus these formulas are independently normalized, not an unexplained use of leading Ehrhart coefficients.

\begin{srcBGlemma}[The fixed half-root-sum shift]\label{R6E3:shift}
For every nonnegative integer margin,
\begin{equation}\label{R6E3:identity}
 K(a,b,c)-R(a,b,c)=V(a+3/2,b+2,c+3/2).
\end{equation}
The shift $\rho=(3/2,2,3/2)$ is half the sum of the six columns in \eqref{R6E3:matrix}.
\end{srcBGlemma}
\begin{proof}
Suppose $a\le c$. The four integer regions in \eqref{R6E3:correction} place $(A,B,C)=(a+3/2,b+2,c+3/2)$ in the four respective regions of \eqref{R6E3:volumes}. In particular, $b<a$ implies $b\le a-1$, so $B<A$; $a\le b<c$ implies $A<B<C$; $c\le b\le a+c$ implies $C<B<A+C$; and $b>a+c$ implies $B\ge A+C$. Equality cases at $A=C$ are harmless. This uses integrality and is the reason for the half-open convention.

In the first region, $(b+1)(b+2)(b+3)=(b+2)^3-(b+2)$. In the second and fourth use $(a+1)(a+2)=(a+3/2)^2-1/4$. In the third also use
\[
 d(d+1)(d+2)=(d+1)^3-(d+1),\qquad A+C-B=d+1.
\]
Subtracting \eqref{R6E3:volumes} from \eqref{R6E3:counts} gives exactly the four expressions in \eqref{R6E3:correction}. Reflection proves the other half of the orthant.
\end{proof}
The shift-volume viewpoint belongs to the established theory of vector partitions, Euler--Maclaurin formulas and box-spline deconvolution; see \cite{R6E3:McS}. We do not claim the general count/volume mechanism or the classical chamber polynomials as new. The elementary identity above is stated explicitly to make the correction usable and its chamber conventions auditable; historical priority for this particular formula is not asserted.

\section{Global power concavity without an interpolation shortcut}
For real margins $\lambda_i\ge0$ and probabilities $\mu_i$,
\[
 P\left(\sum_i\mu_i(\lambda_i+\rho)\right)
 \supseteq\sum_i\mu_i P(\lambda_i+\rho).
\]
Every defining inequality in \eqref{R6E3:poly} is affine in the parameters, so this inclusion follows by adding inequalities. Classical three-dimensional Brunn--Minkowski therefore implies
\begin{equation}\label{R6E3:BM}
 V\left(\sum_i\mu_i\lambda_i+\rho\right)^{1/3}
 \ge\sum_i\mu_i V(\lambda_i+\rho)^{1/3}.
\end{equation}
All these bodies have positive volume: they contain $P(\rho)$, whose volume from \eqref{R6E3:volumes} is $23/24$. Lemma~\ref{R6E3:shift} gives \eqref{R6E3:root} and positivity. This proof uses a genuinely concave real volume extension, not the piecewise affine formula for $R$ at noninteger margins. Brunn--Minkowski is a classical input, not a discovery here \cite{R5E3:Gardner}.

The same proof permits countably supported distributions with finite first absolute moment. Indeed $P(a+3/2,b+2,c+3/2)$ lies in a box of side at most $a+b+c+2$, so its volume's cube root grows at most linearly. Jensen for this nonnegative concave extension follows by replacing each omitted tail with its real conditional mean and then dropping its nonnegative contribution, exactly the fixed-mean method in \cite[Corollary 1.2]{R6E3:R503}. Thus the extension is integrable and the finite inequalities pass to the limit. This argument is reused with credit.

The exponent is best possible. Along $\lambda=(m,m,m)$, \eqref{R6E3:identity} is a polynomial in $m$ with leading term $m^3/6$. If $p>1/3$, the proposed midpoint inequality for $G^p$ at $\lambda=(m,m,m)$ and $3\lambda$ would, after division by $m^{3p}$ and passage to infinity, require
\[
 2^{3p}\ge\frac{1+3^{3p}}2,
\]
contrary to strict convexity of $x^{3p}$. This proves the positive and exponent assertions of Theorem~\ref{R6E3:main}.

\section{Why every constant or sublinear correction fails}
For fixed integer $x\ge0$ and integer $z\ge0$, put
\[
 \lambda(x,z)=(z+x,z+2x,z+x),\qquad h=z+2x+2.
\]
All these margins belong to the third integer region in \eqref{R6E3:correction}. Direct substitution in the attributed formula \eqref{R6E3:counts} gives
\begin{align}
 K(\lambda(x,z))&=\frac{h^3}{6}-\left(x^2+x+\frac16\right)h
                  +x^3+\frac32x^2+\frac12x,\label{R6E3:hcount}\\
 R(\lambda(x,z))&=\frac h{12}-\frac{2x+1}{8}.\label{R6E3:hcorrection}
\end{align}
Consider the actual integer midpoint triple
\begin{equation}\label{R6E3:triple}
 \lambda_-=(k,k,k),\quad
 \lambda_0=(3k+1,3k+2,3k+1),\quad
 \lambda_+=(5k+2,5k+4,5k+2).
\end{equation}
It corresponds to $(x,z)=(0,k),(1,3k),(2,5k)$ and has $\lambda_0=(\lambda_-+\lambda_+)/2$. Its counts are respectively
\begin{equation}\label{R6E3:threecounts}
 \frac{k^3+6k^2+11k+6}{6},\quad
 \frac{9k^3+36k^2+35k+10}{2},\quad
 \frac{125k^3+450k^2+355k+84}{6}.
\end{equation}
In particular at $k=1$ they are $4,45,169$. The asymptotic argument below, not this single diagnostic, rules out every sublinear correction.

For a fixed real $t$, write
\[
 e_t(x)=-\left(x+\tfrac12\right)^2+\frac{1-t}{12}.
\]
Equations \eqref{R6E3:hcount}--\eqref{R6E3:hcorrection} and the Taylor expansion $(1+s)^\alpha=1+\alpha s+O(s^2)$ at $s=0$ give, for fixed $x,t$,
\begin{equation}\label{R6E3:expansion}
 (K(\lambda(x,z))-tR(\lambda(x,z)))^{1/3}
 =6^{-1/3}\left(h+\frac{2e_t(x)}h+O(h^{-2})\right).
\end{equation}
The remainder statement is rigorous because the argument of the Taylor expansion is $1+O(h^{-2})$, with the next polynomial term $O(h^{-3})$. All radicands are positive for sufficiently large $h$.

The three $h$ values are $k+2,3k+4,5k+6$, whose middle is their arithmetic mean. Their leading terms cancel exactly. Further,
\begin{equation}\label{R6E3:ecoef}
 \frac{e_t(0)}2+\frac{e_t(2)}{10}-\frac{e_t(1)}3=\frac{1-t}{45}.
\end{equation}
If $\Delta_t(k)$ denotes the average of the two endpoint cube roots minus the middle cube root, \eqref{R6E3:expansion} proves
\begin{equation}\label{R6E3:limit}
 \lim_{k\to\infty} k\Delta_t(k)=\frac{2(1-t)}{45\,6^{1/3}}.
\end{equation}
For every $t<1$ the limit is strictly positive, proving failure at every sufficiently large $k$.

Now let $E$ satisfy \eqref{R6E3:sublinear}. Along each sequence \eqref{R6E3:triple}, $E=o(k)$ while $K$ has a positive cubic leading term. By the mean value theorem,
\[
 (K-E)^{1/3}-K^{1/3}=O(|E|k^{-2})=o(k^{-1}).
\]
Thus $k$ times its midpoint defect has the same strictly positive limit as \eqref{R6E3:limit} at $t=0$. This proves the sublinear obstruction and completes Theorem~\ref{R6E3:main}. No effective cutoff for an unspecified little-$o$ function is asserted.

\section{An optimal-order obstruction in every higher dimension}
The following product consequence is included to show that the obstruction is not confined to a single dimension. Tensorization itself is elementary and is not claimed as a new mechanism.

For an integer $d\ge3$, take the direct sum of the matrix in \eqref{R6E3:matrix} and $d-3$ copies of the row $(1\ 1)$. The resulting fixed matrix $M_d$ is totally unimodular, has rank $d$, $2d$ columns and nullity $d$. Its nonnegative-fiber count at parameters
$\theta=(a,b,c,t_1,\ldots,t_{d-3})\in\srcBGZ_{\ge0}^d$ is
\begin{equation}\label{R6E3:product}
 F_d(\theta)=K(a,b,c)\prod_{j=1}^{d-3}(t_j+1).
\end{equation}
Define the explicit correction
\[
 C_d(\theta)=R(a,b,c)\prod_{j=1}^{d-3}(t_j+1).
\]
It is piecewise polynomial of degree at most $d-2$ and is $O((1+\|\theta\|_1)^{d-2})$.

\begin{srcBGtheorem}[Sharp correction order]\label{R6E3:dimension}
For the fixed systems $M_d$, $(F_d-C_d)^{1/d}$ satisfies the finite barycentric inequality at every integral mean of nonnegative integer parameters. Its exponent $1/d$ is optimal. No correction $E_d=o((1+\|\theta\|_1)^{d-2})$ can make $(F_d-E_d)^{1/d}$ midpoint-concave. For every fixed $t<1$, replacing $C_d$ by $tC_d$ also fails.
\end{srcBGtheorem}
\begin{proof}
The corrected count is exactly the volume of the $d$-dimensional convex body
\[
 P(a+3/2,b+2,c+3/2)\ \times\prod_{j=1}^{d-3}[0,t_j+1].
\]
These bodies obey the same Minkowski containment under parameter averaging, giving the positive result by classical $d$-dimensional Brunn--Minkowski. Along a ray with all margin coordinates growing proportionally, the volume has positive leading order $k^d$. The midpoint argument from Section 3 excludes exponents above $1/d$.

For the negative result, use \eqref{R6E3:triple} and set every additional margin $t_j=b+1$ at each point. This prescription is affine, so the middle full parameter remains the arithmetic mean. Each factor $t_j+1$ equals $h=b+2$. Hence
\[
 F_d-tC_d=(K-tR)h^{d-3}
 =\tfrac16h^d+e_t(x)h^{d-2}+O(h^{d-3}),
\]
where the last expression holds with fixed $x,t$. Taking the $d$th root gives
\[
 (F_d-tC_d)^{1/d}=6^{-1/d}\left(h+\frac{6e_t(x)}{d h}+O(h^{-2})\right).
\]
The same cancellation and \eqref{R6E3:ecoef} show that its endpoint-average-minus-middle defect $\Delta_{d,t}(k)$ satisfies
\begin{equation}\label{R6E3:dlimit}
 \lim_{k\to\infty}k\Delta_{d,t}(k)
 =\frac{2(1-t)}{15d\,6^{1/d}}.
\end{equation}
This is positive for $t<1$. Finally, along these parameters $\|\theta\|_1$ is of order $k$ and $F_d$ of order $k^d$. A correction $E_d=o(k^{d-2})$ changes its $d$th root by
$O(|E_d|k^{1-d})=o(k^{-1})$. It cannot remove the positive limit at $t=0$.
\end{proof}
This is a necessary-order obstruction for explicit fixed matrices, not a theorem that every nullity-$d$ unimodular system admits such a correction. Its direct-sum form also does not give new uncorrected logarithmic classes beyond products of the previously covered $A_3$ and interval cases.

\section{Limits, source audit and exact verification}
The correction in \eqref{R6E3:correction} has both signs. Therefore the corrected inequality is not automatically a stronger inequality between the \emph{uncorrected} counts in every numerical instance. It is a sharp power-concavity statement for explicitly corrected counts. The uncorrected logarithmic inequality for $A_3$ remains credited to entry 06; its continuous exponent $4/13$ concerns a different function and is not contradicted by our corrected exponent $1/3$.

The count formulas, homogeneous volume formulas, Brunn--Minkowski, half-root-sum shifts and general count/volume relations all have classical antecedents. Billey--Guillemin--Rassart v1 and the current Leake--Mohammadi Yekta record were checked on 20 October 2026. McSwiggen's published 2021 paper was inspected for its box-spline and shifted-volume context, not treated as a new theorem of this work. Targeted searches did not establish historical priority for our exact correction inequality or asymptotic obstruction. No exhaustive novelty claim, general-conjecture closure or external peer-review claim is made.

The exact verifier independently enumerates integer fibers, computes rational slice volumes, checks every half-open chamber including all walls, and verifies all symbolic identities and asymptotic coefficients. It uses rational isolating intervals for roots to certify finite failures for constants and underscaled corrections in several dimensions. The all-order conclusions follow from the displayed proofs, not these finite tests. Initial weighted theta-fiber searches found no counterexample in a bounded experiment, but yielded no proof; arbitrary monomial weights, general nullity three and weighted or boxed $K_4$ logarithmic inequalities remain unproved here. No external posting, submission or author contact occurred.

\begingroup\small
\endgroup

\chapter[Series Invariance and a Dimension-Three Obstruction for Unimodular Fiber Concavity]{Series Invariance and a Dimension-Three Obstruction for Unimodular Fiber Concavity}
\label{source:R6E6}
\noindent\textit{Source development: Round 6, entry 06. Original mathematical arguments retained; notation is local to this chapter.}\par\medskip

\begin{quote}\small\noindent\textbf{Chapter abstract.} 
We prove that finite-support barycentric log-concavity of nonnegative integer fibers is invariant under coherent series extension of a pointed integer matrix. A cone-monotone concave extension transfers the inequality through the convex maximum lower bounds created by the new variables. Combining this closure with the attributed low-rank results from Rounds 4--5 gives a new arbitrary-rank class: unimodular matrices of nullity three whose nonzero coordinate functionals on the kernel contain no negatively proportional pair. Coherently oriented subdivisions of the complete four-vertex acyclic network additionally inherit the earlier $4/13$ power inequality, with every internal balance allowed to vary. Countable integrable distributions are included. In the opposite direction, the quarter-corrected square-root theorem of Round 5 entry 03 has no constant-corrected dimension-matched-root analogue in any nullity $d\ge3$. Fixed totally unimodular matrices have explicit ray counts $a_d(k+1/2)^d+b_d(k+1/2)^{d-2}$ with $a_d,b_d>0$, making every constant-corrected $d$th root eventually strictly convex. An order-$k^{d-2}$ correction is necessary and attained on these rays. These are scoped extensions and obstructions, not a resolution of the general barycentric question.
\end{quote}

\section{Results, definitions, and precise increment}
For an integer matrix $M$ let
\[
 F_M(b)=\#\{x\in\srcBJZ_{\ge0}^n:Mx=b\},\qquad C_M=M\srcBJR_{\ge0}^n.
\]
We call $M$ \emph{pointed} when $\ker M\cap\srcBJR_{\ge0}^n=\{0\}$. All its nonnegative real fibers are then bounded. Conversely one nonempty finite integer fiber forces pointedness: a nonzero rational ray in the rational recession cone would give an integer vector that can be repeatedly added. Throughout, matrices have no zero columns unless the context explicitly removes them.

Write $\mathsf{BL}(M)$ for the following property: for finitely many integer margins $b_i$ with positive finite counts, positive real weights $\mu_i$ summing to one, and integer mean $\bar b$, the central count is positive and
\begin{equation}\label{R6E6:eq:BL}
 F_M(\bar b)\ge\prod_iF_M(b_i)^{\mu_i}.
\end{equation}
A matrix of rank $r$ is \emph{unimodular} if every nonsingular $r\times r$ submatrix has determinant $\pm1$; redundant rows are allowed. Its nullity is $n-r$.

A \emph{coherent series extension} replaces $x_e$ by a chain of $k_e\ge1$ nonnegative coordinates, with equations
\begin{equation}\label{R6E6:eq:series}
 M(x_{e,1})_e=b,\qquad x_{e,j+1}-x_{e,j}=\kappa_{e,j}.
\end{equation}
All components of $(b,\kappa)$ may vary. Denote the resulting integer matrix by $M^{\rm ser}$.

\begin{srcBJtheorem}[Series invariance]\label{R6E6:thm:series}
For every pointed integer matrix $M$ and every choice of the lengths $k_e$,
\[
 \mathsf{BL}(M)\quad\Longleftrightarrow\quad\mathsf{BL}(M^{\rm ser}).
\]
If $M$ is unimodular, so is $M^{\rm ser}$, and the nullity is unchanged. Each positive conclusion extends to countably supported distributions with finite first absolute moment and integer mean.
\end{srcBJtheorem}

For $V=\ker M$, let $\lambda_j:V\to\srcBJR$ be its $j$th coordinate functional. Call the kernel \emph{unmixed} if no two nonzero $\lambda_i,\lambda_j$ are negatively proportional. This condition is intrinsic, independent of a choice of coordinates on $V$. Equal or positively proportional functionals are allowed; zero functionals impose deterministic coordinate constraints and cause no difficulty.

\begin{srcBJtheorem}[Unmixed nullity three]\label{R6E6:thm:unmixed}
Let $M$ be unimodular of nullity three with unmixed kernel. Its nonempty finite nonnegative integer fibers satisfy \eqref{R6E6:eq:BL}, including for countably supported distributions with finite first absolute moment and integer mean. Rank and number of variables are unrestricted.
\end{srcBJtheorem}

\begin{srcBJcorollary}[Subdivided complete four-vertex networks]\label{R6E6:cor:network}
Start with an acyclic orientation of $K_4$ and replace each directed edge by any nonempty directed path following its orientation. For nonnegative integer flows on the resulting network, all vertex balances, including every new internal balance, may vary. Its fiber count $F$ satisfies
\begin{equation}\label{R6E6:eq:power}
 F(\bar b)^{4/13}\ge\sum_i\mu_iF(b_i)^{4/13},
 \qquad F(\bar b)\ge\prod_iF(b_i)^{\mu_i}
\end{equation}
for finite or countably integrable distributions with integer mean and nonempty fibers. The network has nullity three and unbounded possible rank.
\end{srcBJcorollary}
Coherence matters: an oppositely directed edge in a path may create an upper bound rather than a lower bound; it is not covered by the series statement. Theorem~\ref{R6E6:thm:unmixed} excludes exactly such opposite kernel directions, not all repeated directions. General finite boxes introduce opposite rows in a slack-variable kernel and remain outside the new theorem.

\begin{srcBJtheorem}[No constant correction in dimensions at least three]\label{R6E6:thm:negative}
For every $d\ge3$ there is a fixed rank-$d$, nullity-$d$ totally unimodular matrix $D_d$ and an integer ray of margins $b_k$ such that, for every real constant $c$, the positive sequence
\[
 \bigl(F_{D_d}(b_k)-c\bigr)^{1/d}
\]
is eventually strictly convex. In particular the midpoint concavity inequality fails for every sufficiently large integer center $k$. Thus no constant $c_d$, even one depending on the fixed matrix, repairs dimension-matched root concavity for all its nonnegative fibers.
\end{srcBJtheorem}
This does not contradict logarithmic concavity, which is strictly weaker. A sharper necessary correction-growth statement appears in Section~\ref{R6E6:sec:obstruction}.

\paragraph{Target and source status.}
Leake--Mohammadi Yekta \cite[Question 1.7]{R4E6:LMY} ask for \eqref{R6E6:eq:BL} for arbitrary unimodular matrices. Their Theorem 1.4 and Theorem 3.1 establish line/midpoint statements, which the question explicitly distinguishes from barycentric concavity. The arXiv record rechecked on 20 October 2026 lists v1, 30 September 2026. The present results cover a restricted nullity-three class, not all nullity three or all rank three.

\paragraph{Attribution and exact delta.}
Round 4 entry 06 \cite{R6E6:R406} proved rank-at-most-two and outerplanar all-balance cases, using the known four-functions/marginal-closure machinery of Ahlswede--Daykin and Speyer \cite{R5E6:AD,R4E6:Speyer}. Round 5 entry 06 \cite{R6E6:R506} proved the $A_3$ chamber-extension exponent $4/13$ and the simple rank-three case. Those core results and its small direction certificate are reproduced in the appendices and not counted again. The new positive increment is series invariance and the arbitrary-rank unmixed-nullity-three reduction.

Round 5 entry 03 \cite{R6E6:R503} proved sharp $\sqrt{F-1/4}$ concavity in nullity at most two, with varying boxes. Its half-hexagon identity is a direct case of Chu--Karshon \cite[Theorem 5.8]{R5E3:CK}, not a newly discovered identity. Our negative result establishes a sharp dimension boundary for that \emph{constant-correction strategy}, not a failure of its logarithmic target. All six frozen Round 5 papers and the assessment were read. The unrelated shared permanent results of entries 04 and 05 are not proof inputs. Primary-source searches did not establish historical priority of our closure or obstruction; no first-discovery claim is certified.

\section{A cone-monotone concave extension}
The following elementary envelope lemma makes the closure argument exact, rather than assuming that integer line concavity has a concave extension.

\begin{srcBJlemma}[Extension from the full barycentric property]\label{R6E6:lem:envelope}
If a pointed integer matrix $M$ has $\mathsf{BL}(M)$, there is a finite nonnegative concave function $g:C_M\to\srcBJR$ with
\[
 g(b)=\log F_M(b)\quad(b\in C_M\cap\srcBJZ^m),\qquad
 g(b+z)\ge g(b)\quad(z\in C_M).
\]
In particular all integer points of $C_M$ are feasible margins.
\end{srcBJlemma}
\begin{proof}
Let $S=M\srcBJZ_{\ge0}^n$. Its convex hull is $C_M$: if $b=Mt$ with $t\ge0$, independently rounding each coordinate of $t$ to its floor or ceiling expresses $b$ as a finite convex combination of points of $S$. Put
\[
 g(b)=\sup\left\{\sum_i\mu_i\log F_M(b_i):
 b_i\in S,\ \mu_i>0,\ \sum_i\mu_i=1,\ \sum_i\mu_i b_i=b\right\}.
\]
The defining collections are nonempty, and concatenating two collections proves concavity once finiteness is checked. At an integer $b$, any such collection and $\mathsf{BL}(M)$ prove both feasibility and $g(b)=\log F_M(b)$, using the one-point collection for equality.

Integer translation in the cone cannot decrease the count: adding $v\in\srcBJZ_{\ge0}^n$ injects a fiber into the fiber at $b+Mv$. If $z=Mt$ with $t\ge0$, the same floor/ceiling rounding represents $z$ as the expectation of finitely many such $Mv$. Translate every member of a defining collection by that distribution. The mean increases by $z$ and no average log count decreases. Taking suprema proves cone monotonicity, even before ruling out infinity.

For any fixed $b=Mt$, choose $B=M\lceil t\rceil\in S$. Since $B-b\in C_M$, the same translation argument shows that every candidate average for $g(b)$ is at most $g(B)=\log F_M(B)<\infty$. Thus $g$ is finite, and it is nonnegative because positive counts are integers. This proves the lemma without a hidden compactness or growth assumption on the envelope.
\end{proof}

\begin{srcBJlemma}[Convex lower-bound transfer]\label{R6E6:lem:transfer}
Suppose $g$ is as in Lemma~\ref{R6E6:lem:envelope}, $B(\theta)$ is linear, and each coordinate of $L(\theta)$ is convex. Then
\begin{equation}\label{R6E6:eq:composition}
 \widetilde g(\theta)=g(B(\theta)-ML(\theta))
\end{equation}
is concave on the domain where its argument belongs to $C_M$, provided that domain is convex. The same assertion holds for any other finite concave, cone-monotone function in place of $g$.
\end{srcBJlemma}
\begin{proof}
For a finite convex combination let $\delta=\sum_i\mu_iL(\theta_i)-L(\bar\theta)\ge0$. Then
\[
 B(\bar\theta)-ML(\bar\theta)
 =\sum_i\mu_i[B(\theta_i)-ML(\theta_i)]+M\delta.
\]
Cone monotonicity followed by ordinary concavity gives the result. In fact the domain is automatically convex when $B$ is linear and $L$ is coordinatewise convex, by this identity and convexity of the cone.
\end{proof}

\begin{proof}[Proof of Theorem~\ref{R6E6:thm:series}]
Set $s_{e,1}=0$ and $s_{e,j}=\sum_{h<j}\kappa_{e,h}$. Equations~\eqref{R6E6:eq:series} force $x_{e,j}=x_{e,1}+s_{e,j}$, so their nonnegativity is equivalent to
\[
 x_{e,1}\ge L_e(\kappa):=\max_j(-s_{e,j}).
\]
These lower bounds are integer at integer parameters and convex in real parameters. Shifting $x_{e,1}=y_e+L_e$ gives the exact identity
\begin{equation}\label{R6E6:eq:seriescount}
 F_{M^{\rm ser}}(b,\kappa)=F_M(b-ML(\kappa)).
\end{equation}
The forward implication follows from Lemmas~\ref{R6E6:lem:envelope}--\ref{R6E6:lem:transfer}; the reverse sets every $\kappa$ to zero. Pointedness is preserved because a homogeneous nonnegative solution is constant on each chain and reduces to a nonnegative kernel vector of $M$.

If $s=\sum_e(k_e-1)$, the rank rises from $r$ to $r+s$, and the number of columns rises from $n$ to $n+s$. For unimodularity, a nonsingular maximal minor must select all $s$ difference rows and $r$ independent old rows; the old row space has rank only $r$. In each chain it must select at least $k_e-1$ columns, since otherwise its difference rows are dependent on the selected columns. Precisely $r$ chains have all $k_e$ columns selected. Eliminating the selected difference rows within each path has determinants $\pm1$ and leaves the maximal old minor on those $r$ columns. Thus every nonsingular maximal minor is $\pm1$.

For countable distributions, the real feasible domain is the closed polyhedral cone $M^{\rm ser}\srcBJR_{\ge0}^{n+s}$. Every conditional tail mean belongs to it when the first absolute moment is finite, by taking expectations in a finite half-space description. Replace the tail of any finite retained set of atoms by its real conditional mean. Finite concavity of the nonnegative function \eqref{R6E6:eq:composition} bounds the retained sum of log counts by its value at the full mean. Discard the nonnegative tail term and exhaust the support. Monotone convergence proves the countable assertion, including finiteness of the sum of log counts. This is the same standard conditional-tail argument used in Round 5 entry 03, not a new limiting principle.
\end{proof}

\section{Compressing an unmixed three-dimensional kernel}
\begin{srcBJlemma}[Integral kernel normalization]\label{R6E6:lem:kernel}
A unimodular matrix of nullity three admits an integral parametrization
\[
 x=x^0(b)+Kt,\quad t\in\srcBJR^3,\qquad
 K=\binom{-N}{I_3},
\]
after reordering coordinates, where $x^0$ is linear and integral on feasible integer margins. All $3\times3$ minors of $K$ are $0,\pm1$, and $t\in\srcBJZ^3$ exactly when $x$ is integral.
\end{srcBJlemma}
\begin{proof}
Choose independent rows with projection $R$ and basis columns $J$ such that $B=RM_J$ has determinant $\pm1$. The equations become $[I_r\ N]x=B^{-1}Rb$, with $x^0=(B^{-1}Rb,0,0,0)$. Projection is injective on the column span, so no equation is lost there. Entries and minors of $N$ of sizes at most three are $0,\pm1$, by replacing the corresponding identity columns in maximal minors of $[I_r\ N]$. Expanding a $3\times3$ minor of $[-N;I_3]$ on its identity rows gives precisely such a minor. The last three coordinates of $x$ are $t$, proving the lattice statement.
\end{proof}

\begin{proof}[Proof of Theorem~\ref{R6E6:thm:unmixed}]
If there are no nonempty finite fibers the statement has no inputs. Otherwise $M$ is pointed. By Lemma~\ref{R6E6:lem:kernel}, each nonzero row of $K$ is a primitive vector in $\{0,\pm1\}^3$. Thus positive proportionality means equality. Remove zero rows and retain one representative from each group of equal rows, obtaining $K_0$ with $q$ distinct nonproportional rows. Unmixedness ensures that opposite rows have not survived. The inherited direction certificate in Appendix~\ref{R6E6:app:directions}, applied to $K_0^{\mathsf T}$, gives $q\le6$. The original identity rows survive in some groups, so $K_0$ has a three-row unimodular basis.

Choose that basis, make it $I_3$ by an integer change of $t$, and reorder rows to write
\[
 K_0=\binom{N_0}{I_3},\qquad M_0=[I_{q-3}\ {-N_0}].
\]
Then $\ker M_0=\operatorname{im}K_0$. Complementary minors show that $M_0$ is unimodular: replacing $j$ identity columns reduces a maximal minor to a $j\times j$ minor of $N_0$, which is a $3\times3$ minor of $K_0$ completed by identity rows. Also $M_0$ is pointed: $K_0t\ge0$ would reproduce a nonnegative $Kt$ by duplicating rows and adding zero rows, contradicting pointedness of $M$. In particular $q>3$.

If $q\le5$, $M_0$ has rank at most two, so the attributed core theorem in Appendix~\ref{R6E6:app:cores} applies. If $q=6$, the direction certificate identifies the six rows of $K_0$, up to individual signs, with the six reduced-incidence columns of $K_4$. Any four of these directions span $\srcBJR^3$: a four-edge simple graph on four vertices is connected, since a disconnected such graph has at most three edges. Hence a vector $K_0t$ supported on at most two coordinates is zero. This proves that no column of $M_0$ is zero and no pair is proportional. It is therefore a simple rank-three unimodular matrix, again covered by Appendix~\ref{R6E6:app:cores}. In either case $\mathsf{BL}(M_0)$ holds.

Here is the exact transfer, including the parameter dependence that a mere count of directions would miss. For each group choose representative $j(g)$ and put $c_g(b)=x^0_{j(g)}(b)$. Write $y_g=x_{j(g)}$. All coordinates in that group satisfy
\[
 x_j=y_g+x_j^0(b)-c_g(b),\qquad
 y_g\ge L_g(b):=\max_{j\text{ in group }g}\{c_g(b)-x_j^0(b)\}.
\]
The lower bounds are convex and integral on integer margins. The affine constraints on $y$ are $M_0y=M_0c(b)$. They identify the lattice exactly, since the three chosen representative coordinates give the integer parameters $t$. Zero rows of $K$ merely impose $x_j^0(b)\ge0$. Consequently, on these linear feasibility conditions,
\begin{equation}\label{R6E6:eq:compresscount}
 F_M(b)=F_{M_0}\bigl(M_0c(b)-M_0L(b)\bigr).
\end{equation}
Lemmas~\ref{R6E6:lem:envelope}--\ref{R6E6:lem:transfer} prove finite concavity of the log count. The domain is the original closed feasible cone, with the stated zero-row conditions; the conditional-tail argument gives the countable conclusion. All coordinate maps are integral on the relative lattice, including when $M$ has redundant rows.
\end{proof}

\subsection{The explicit power statement and graph interpretation}
For the transitive $K_4$, use simple-root margins $(a,b,c)$, where the six columns are
\begin{equation}\label{R6E6:eq:A3}
 A_3=[e_1,e_2,e_3,e_1+e_2,e_2+e_3,e_1+e_2+e_3].
\end{equation}
Appendix~\ref{R6E6:app:A3} reproves the previously established concavity of the continuous chamber function $\mathcal K^{4/13}$ on $\srcBJR_{\ge0}^3$. It is also coordinatewise nondecreasing. In the central chamber $a\le c$, $d=a+c-b\in[0,a]$, direct derivatives give
\[
 6\mathcal K_b=3d^2+6d+2>0,\qquad
 6\mathcal K_a\ge3a+5>0,\qquad
 6\mathcal K_c\ge3a+4>0.
\]
The product chambers have nonnegative partial derivatives by their displayed formulas, and reflection handles $a\ge c$. Continuity transfers monotonicity across walls. Since $C_{A_3}=\srcBJR_{\ge0}^3$, this is the required cone monotonicity.

Thus Lemma~\ref{R6E6:lem:transfer} applies to $\mathcal K^{4/13}$, not merely to log counts, for every coherent series extension of $A_3$. It is positive, so the same tail argument proves countable power Jensen; weighted arithmetic--geometric mean then gives its logarithmic consequence.

For completeness, an actual edge subdivision has the claimed serial form even with variable internal balances. On the path $v_0\to\cdots\to v_k$, use divergence outflow minus inflow. At its internal vertex $v_j$, the balance is $\beta_{v_j}=x_{j+1}-x_j$. If $S_e$ is the sum of these internal balances, the terminal incoming flow is $x_1+S_e$. At an original vertex $v$, the contracted margin is therefore
\[
 B_v=\beta_v+\sum_{e:\,\operatorname{head}(e)=v}S_e.
\]
This is an integral linear change of the margin data, with inverse obtained by the same formula and the retained internal balances. It gives exactly \eqref{R6E6:eq:series}. Every acyclic orientation of $K_4$ is transitive, and changing the topological vertex order is harmless. This proves Corollary~\ref{R6E6:cor:network}. It does not assert the same for arbitrarily reversed paths, weighted counts, or upper capacities.

\section{The obstruction to constant root corrections}\label{R6E6:sec:obstruction}
Let $D$ be the reduced incidence matrix of the transitive four-vertex tournament, with columns ordered as $(12),(23),(34),(13),(24),(14)$:
\[
 D=\begin{pmatrix}1&0&0&1&0&1\\-1&1&0&0&1&0\\0&-1&1&-1&0&0\end{pmatrix}.
\]
At margin $k(3,1,-1)$, cumulative balances give simple-root coordinates $(3k,4k,3k)$. The classical chamber count, or grouping by its three free coordinates, gives
\begin{align}\label{R6E6:eq:ray}
 N(k)&=\frac{(3k+1)(3k+2)(6k+3)-(2k)(2k+1)(2k+2)}6\\
 &=\frac{23}{3}(k+\tfrac12)^3+\frac1{12}(k+\tfrac12).\notag
\end{align}
The first values are $1,26,120,329,699$. These are exact counts of one fixed matrix, not values of a chosen interpolation at noninteger parameters.

For $d\ge3$ take
\[
 D_d= D\oplus\underbrace{[1\ 1]\oplus\cdots\oplus[1\ 1]}_{d-3\ \text{blocks}},
 \qquad b_k=(3k,k,-k,2k,\ldots,2k).
\]
This matrix has rank and nullity both $d$. Incidence matrices are totally unimodular: determinant expansion handles a column with at most one nonzero in a square submatrix; otherwise each column has two opposite entries and the row sum is zero. This applies to $D$ and $[1\ 1]$, and direct sums preserve the property. All these acyclic nonnegative fibers are finite. Independence of the blocks gives, with $t=k+1/2$,
\begin{equation}\label{R6E6:eq:Nd}
 N_d(k)=N(k)(2k+1)^{d-3}=a_dt^d+b_dt^{d-2},\quad
 a_d=2^{d-3}\frac{23}{3},\quad b_d=2^{d-3}\frac1{12}.
\end{equation}

\begin{proof}[Proof of Theorem~\ref{R6E6:thm:negative}]
For real $t$ put $Q(t)=a_dt^d+b_dt^{d-2}-c$. For sufficiently large $t$, it is positive. The second derivative of $Q^{1/d}$ has the sign of
\begin{align}\label{R6E6:eq:curvature}
 dQQ''-(d-1)(Q')^2
 &=2d a_db_dt^{2d-4}-2(d-2)b_d^2t^{2d-6}\\
 &\quad-dc\{a_dd(d-1)t^{d-2}+b_d(d-2)(d-3)t^{d-4}\}.\notag
\end{align}
Its leading term is strictly positive: $2d-4>d-2$ for $d\ge3$, and all other displayed powers are lower. Thus $Q^{1/d}$ is strictly convex on a real terminal ray. Its values at every three consecutive sufficiently large integers violate midpoint concavity. This proves the all-constant, all-dimension assertion with explicit fixed matrices.
\end{proof}

\begin{srcBJproposition}[The necessary correction scale is sharp on the ray]\label{R6E6:prop:scale}
Fix $d\ge3$. Suppose a real correction sequence $R(k)=o(k^{d-1})$ makes $(N_d(k)-R(k))^{1/d}$ well-defined and concave for all sufficiently large integers. Then necessarily
\begin{equation}\label{R6E6:eq:necessary}
 R(k)\ge b_d(k+\tfrac12)^{d-2}
\end{equation}
for every sufficiently large $k$. Equality in \eqref{R6E6:eq:necessary} makes that root exactly affine, so the lower bound is attained on this family. In particular a correction of order $o(k^{d-2})$ cannot suffice.
\end{srcBJproposition}
\begin{proof}
The assumed size of $R$ and \eqref{R6E6:eq:Nd} give
\[
 f(k):=(N_d(k)-R(k))^{1/d}
 =a_d^{1/d}(k+\tfrac12)+o(1).
\]
For an eventually concave sequence the increments $f(k+1)-f(k)$ decrease to its asymptotic slope $a_d^{1/d}$, and hence are at least that slope. Thus $f(k)-a_d^{1/d}(k+1/2)$ is nondecreasing and tends to zero, so is nonpositive on its terminal domain. Raising to the $d$th power proves \eqref{R6E6:eq:necessary}. Equality follows immediately from \eqref{R6E6:eq:Nd}.
\end{proof}
This shows why a dimension-two constant correction cannot simply be carried one dimension higher. It does not rule out geometric corrections depending on edge lengths, lower-dimensional volumes, or the particular parameter vector. Nor does it refute the logarithmic inequality proved for these same matrices.

\section{Limitations and verification}
The principal positive restriction is the unmixed kernel condition. Opposite coordinate functionals create upper and lower bounds on the same free direction, and the single convex-lower-bound argument no longer applies. No assertion is made for all nullity-three matrices, arbitrary monomial weights, or arbitrary variable boxes. The series-invariance theorem is conditional on a proved base barycentric property; it is not a reformulation that settles the unrestricted question automatically.

The appendices reproduce exactly the earlier core inputs required for standalone use. Their power exponent and finite direction classification are not additional Round 6 discoveries. The new negative ray uses the same classical $A_3$ counting formulas but a different conclusion: no constant correction works, even asymptotically, in every dimension at least three. Historical priority of that obstruction and the series application remains unestablished after targeted primary-source checks.

The exact verifier checks the serial elimination formula against direct enumeration, maximal minors, the mixed/opposite-row diagnostic, arbitrary-rank examples, the classical chamber count and ray polynomial, the symbolic derivative identity, all printed finite direction certificates, and noncollinear barycentric tests. The universal conclusions come from the proofs, not finite tests. An exploratory attempt to extend the planar half-hexagon argument by a three-dimensional fixed thickening did not produce an identity; the obstruction above shows that a mere constant correction cannot be its conclusion. No external posting, submission, or author contact occurred.

\section*{Supplementary material}
\section{Attributed \texorpdfstring{$A_3$}{A3} power-concavity input, with proof}\label{R6E6:app:A3}
This is Round 5 entry 06, Theorem 1.1, reproduced rather than newly claimed. Its lattice count is $K(a,b,c)=\#\{x\ge0:A_3x=(a,b,c)\}$, with integer variables. The function $\mathcal K$ below extends those counts.

Eliminating $x_1,x_2,x_3$ from \eqref{R6E6:eq:A3} gives
\begin{equation}\label{R6E6:eq:count}
 K(a,b,c)=\#\{(u,v,z)\in\srcBJZ_{\geq0}^3:
 u+z\leq a,\ v+z\leq c,\ u+v+z\leq b\}.
\end{equation}
The expression is symmetric in $a,c$. It suffices to give the formulas when $0\leq a\leq c$:
\begin{equation}\label{R6E6:eq:pieces}
 6\mathcal K(a,b,c)=
 \begin{cases}
 (b+1)(b+2)(b+3),&0\leq b\leq a,\\
 (a+1)(a+2)(3b-2a+3),&a\leq b\leq c,\\
 (a+1)(a+2)(3c-a+3)-d(d+1)(d+2),
       &c\leq b\leq a+c,\ d=a+c-b,\\
 (a+1)(a+2)(3c-a+3),&b\geq a+c.
 \end{cases}
\end{equation}
Reflect the formulas across $a=c$ on the other half of the orthant. The first chamber is shared by both halves, giving seven full-dimensional chambers in total.

\begin{srcBJlemma}\label{R6E6:lem:formula}
These polynomials agree on their common boundaries, are positive throughout the orthant, and equal $K$ at every nonnegative integer point.
\end{srcBJlemma}
\begin{proof}
If $b\leq a,c$, the only effective constraint in \eqref{R6E6:eq:count} is $u+v+z\leq b$, giving $\binom{b+3}{3}$. If $a\leq b\leq c$, group by $t=u+z$. For each $0\leq t\leq a$ there are $t+1$ choices of $(u,z)$ and $b-t+1$ choices of $v$. Summing gives the second polynomial.

If $b\geq a+c$, its constraint is redundant. Grouping by $z$ gives
\[
 \sum_{z=0}^{a}(a-z+1)(c-z+1)
 =\frac{(a+1)(a+2)(3c-a+3)}6.
\]
In the remaining case $a\leq c\leq b\leq a+c$, put $d=a+c-b$. A point counted by this last sum violates the remaining constraint exactly when, on writing $u'=a-z-u$ and $v'=c-z-v$, one has
\[
 z+u'+v'\leq d-1.
\]
Because $0\leq d\leq a\leq c$, these three nonnegative variables automatically obey the original upper constraints. The violating count is $\binom{d+2}{3}$, including zero when $d=0$. This proves the third formula.

Boundary agreement follows by substitution at $b=a$, $b=c$, $b=a+c$, and $a=c$. All product pieces are positive. In the third piece, $H(t)=t(t+1)(t+2)$ is increasing for $t\geq0$, so its numerator is at least
\[
 (a+1)(a+2)(3c-a+3)-H(a)
 =(a+1)(a+2)(3c-2a+3)>0.
\]
The same conclusions hold after reflection.
\end{proof}

The formulas also specify uniquely the polynomial on each full-dimensional chamber from its integer values: a polynomial zero on the integer points of an open rational cone is zero, for example by restricting to arbitrarily large translated integer boxes in that cone. Thus ``the chamber extension'' denotes a fixed classical interpolation, not an interpolation chosen to optimize a concavity exponent.

\subsection{Power concavity inside the chambers}
We prove the exponent $p=4/13$. Multiplication by a positive constant does not affect power concavity. In each product chamber, the numerator is a product of three positive affine functions. Such a product $G=L_1L_2L_3$ satisfies $\srcBJHess(G^p)\preceq0$ whenever $0<p\leq1/3$. Indeed, in any direction, put $t_i=D L_i/L_i$. Then
\[
 \frac{D^2G^p}{G^p}
 =p^2\Bigl(\sum_i t_i\Bigr)^2-p\sum_i t_i^2
 \leq p(3p-1)\sum_i t_i^2\leq0.
\]
The affine functions are positive on the entire chamber being used, including when one coefficient in their formulas is negative.

The only remaining chamber is
\[
 0\leq d\leq a\leq c,\qquad d=a+c-b.
\]
The change from $(a,b,c)$ to $(a,c,d)$ is invertible and linear. Put
\begin{gather}\label{R6E6:eq:central}
 A=(a+1)(a+2),\quad H=d(d+1)(d+2),\quad
 F=A(3c-a+3)-H=6\mathcal K,\\
 h=1-\frac a3-\frac{H}{3A},\qquad
 \tau=c+h=\frac{F}{3A}>0.\notag
\end{gather}
Primes on $A$ denote $a$-derivatives, and primes on $H$ denote $d$-derivatives.

\begin{srcBJlemma}[Central-chamber Hessian]\label{R6E6:lem:hessian}
$F^{4/13}$ is concave on the central chamber.
\end{srcBJlemma}
\begin{proof}
For $0<p<1$, negative semidefiniteness of $\srcBJHess(F^p)$ is equivalent to positive semidefiniteness of
\[
 Q=-\srcBJHess\log F-p(\nabla\log F)(\nabla\log F)^{\mathsf T},
\]
since $-\srcBJHess(F^p)=pF^pQ$. As $\log F=\log(3A)+\log(c+h)$, the $c,c$ entry of $Q$ is $(1-p)/\tau^2>0$. Its Schur complement on the $(a,d)$ coordinates is
\begin{equation}\label{R6E6:eq:schur}
 S=\begin{pmatrix}J/A^2&0\\0&0\end{pmatrix}
       -\frac{\srcBJHess_{a,d}h}{\tau},
 \qquad J=(A'^2-AA'')-\frac{p}{1-p}A'^2.
\end{equation}
For example, the cross entries before taking the complement are
$Q_{i,c}=(1-p)h_i/\tau^2-p(\log A)_i/\tau$; their outer product cancels the corresponding terms containing $h_i h_j$ and $(\log A)_i h_j$, leaving \eqref{R6E6:eq:schur}.

At $p=4/13$, write $C=2a^2+6a+9$. Then
\[
 J=\frac C9>0,\quad
 h_{aa}=-\frac{H(2A'^2-AA'')}{3A^3}\leq0,\quad
 h_{dd}=-\frac{H''}{3A}<0,\quad
 h_{ad}=\frac{H'A'}{3A^2}.
\]
Here $A''=2$, $H''=6(d+1)$, and $2A'^2-AA''>0$. Both diagonal entries of $S$ are positive. Direct substitution in its determinant gives
\begin{equation}\label{R6E6:eq:detS}
 \det S=\frac{\mathcal N}{27A^4\tau^2},\qquad
 \mathcal N=2(d+1)AC(3c-a+3)-A'^2\bigl((d+1)^4+8(d+1)^2+3\bigr).
\end{equation}
One way to verify the simplification is to use $H=t^3-t$, $H'=3t^2-1$, $H''=6t$, where $t=d+1$.

Set $u=a+1$ and $t=d+1$, so $1\leq t\leq u$. Because $c\geq a$,
\[
 3c-a+3\geq2a+3=A'=2u+1,
 \quad A=u(u+1),\quad C=2u^2+2u+5.
\]
The function $\phi(t)=t^3+8t+3/t$ increases for $t\geq1$, since
$\phi'(t)=3t^2+8-3/t^2>0$. Hence
\begin{align*}
 \mathcal N
 &\geq (2u+1)t\left[2u(u+1)(2u^2+2u+5)
                    -(2u+1)\phi(t)\right]\\
 &\geq\frac{(2u+1)t}{u}
       (u-1)(2u^4+9u^3+7u^2+9u+3)\geq0.
\end{align*}
The final equality underlying this bound is the elementary factorization
\[
 2u^2(u+1)(2u^2+2u+5)-(2u+1)(u^4+8u^2+3)
 =(u-1)(2u^4+9u^3+7u^2+9u+3).
\]
Thus $S\succeq0$ and $Q\succeq0$, proving the lemma, including the boundary by continuity. In the strict interior $a>0$ the last lower bound is positive.
\end{proof}

\subsection{Interfaces and global concavity}
Piecewise concavity does not by itself give global concavity; we check every interface. In the following table, $q=\mathcal K$, $n$ is the gradient of the displayed linear form, and the interface is crossed in the direction of increasing that form. The table gives
$\nabla q_{\rm after}-\nabla q_{\rm before}=-\kappa n$.
\begin{center}
\begin{tabular}{lll}
\toprule
Interface and region & Increasing linear form & $\kappa$\\
\midrule
$b=a$, $a<c$ & $b-a$ & $(3a+5)/6$\\
$b=c$, $a<c$ & $b-c$ & $(3a+4)/6$\\
$b=a+c$ & $b-a-c$ & $1/3$\\
$a=c$, $b>a$ & $a-c$ & $1/6$\\
\bottomrule
\end{tabular}
\end{center}
Also include the first two rows with $a,c$ interchanged. The formulas follow by differentiating \eqref{R6E6:eq:pieces} and substituting the facet equation. Across $a=c$ with $b<a$, the formula is unchanged, so the jump is zero. This exhausts all full-dimensional shared facets of the seven chambers.

Every $\kappa$ is nonnegative, and $q$ is continuous and positive. Thus the gradient jump of $q^p$ is $-p q^{p-1}\kappa n$, also in the concave direction. On a generic line segment in the open orthant, avoiding intersections of two chamber walls, the directional derivative of $q^{4/13}$ is nonincreasing within each chamber by the preceding chamber calculation and at every crossing by the table. The restriction is concave. Arbitrary segments in the closed orthant are limits of such generic interior segments. Continuity then proves concavity on the whole orthant, not only chamber interiors. This proves the inherited $4/13$ power-concavity assertion used in Corollary \ref{R6E6:cor:network}.

\section{Attributed direction certificate, reproduced in full}\label{R6E6:app:directions}
We reproduce the finite certificate of Round 5 entry 06, Lemma 5.1. The structural classification in this section is elementary finite linear algebra; its role is a reduction, not a claimed new matroid theorem.

Choose three independent rows of $M$, with row projection $R$, and three columns $J$ such that $B=RM_J$ is nonsingular. Then $\det B=\pm1$ and
\[
 N=B^{-1}RM
\]
contains the three standard basis columns. The map $b\mapsto B^{-1}Rb$ is a lattice isomorphism from $\operatorname{span}_{\srcBJR}(M)\cap\srcBJZ^m$ to $\srcBJZ^3$, with inverse $t\mapsto M_Jt$. Indeed $R$ is injective on the column span, and $B^{-1}$ is integral. It preserves fibers and integer means, so no ambient full-row-rank assumption is hidden here.

Every entry of $N$ belongs to $\{0,\pm1\}$ by Cramer's rule. Ignoring column signs temporarily, choose the representative whose first nonzero coordinate is $+1$. Apart from the three basis directions, the ten candidates are
\begin{equation}\label{R6E6:eq:candidates}
\begin{array}{c|ccccc}
 i&1&2&3&4&5\\\hline
 v_i&(0,1,-1)&(0,1,1)&(1,-1,-1)&(1,-1,0)&(1,-1,1)\\[3pt]
 i&6&7&8&9&10\\\hline
 v_i&(1,0,-1)&(1,0,1)&(1,1,-1)&(1,1,0)&(1,1,1)
\end{array}
\end{equation}
Simplicity ensures that each listed direction occurs at most once.

\begin{srcBJlemma}[Rank-three direction certificate]\label{R6E6:lem:classification}
Every simple unimodular system of rank three is, after an integral invertible change of coordinates, a subset of the six directions
\begin{equation}\label{R6E6:eq:roots}
 e_1,e_2,e_3,e_1-e_2,e_1-e_3,e_2-e_3,
\end{equation}
with arbitrary individual signs.
\end{srcBJlemma}
\begin{proof}
Here is a complete, small certificate. A pair $v_i,v_j$ can coexist with all three basis columns only for the following 24 pairs; checking their determinants with the basis columns gives this list:
\[
\begin{split}
&14,15,16,17,18,19,\quad23,24,26,27,29,2\!:\!10,\\
&34,36,45,46,47,57,68,69,79,7\!:\!10,89,9\!:\!10.
\end{split}
\]
Here $ij$ denotes the unordered pair $\{i,j\}$, and a colon separates indices involving 10. Among triples whose three pairs appear in the list, exactly four have determinant of absolute value two:
\[
 \{1,4,7\},\quad\{1,6,9\},\quad\{2,4,6\},\quad\{2,7,9\}.
\]
The only four-element cliques in the displayed pair graph are
\[
 \{1,4,5,7\},\quad\{1,6,8,9\},\quad
 \{2,3,4,6\},\quad\{2,7,9,10\};
\]
each contains one of those forbidden triples. Thus at most three candidates can be added. The 16 permissible triples and a change-of-basis certificate for each appear in Table \ref{R6E6:tab:systems}. Every permissible pair and singleton is contained in one of them, as the table also verifies. These observations cover all $2^{10}=1024$ candidate subsets, not a sampled list.

For a row of the table let $S$ consist of the standard basis and its three listed candidates, and let $T$ be the matrix with the three displayed basis vectors as columns. Direct three-by-three determinants give $|\det T|=1$, and the six vectors $T^{-1}S$, after normalizing their individual signs, are exactly \eqref{R6E6:eq:roots}. All input directions are contained in some such $S$, proving the lemma. These are ten-vector and three-by-three arithmetic checks, also fully reproduced by the accompanying verifier.
\end{proof}

\begin{table}[ht]
\centering
\begin{tabular}{cc@{\qquad}cc}
\toprule
Added indices & Columns of $T$ & Added indices & Columns of $T$\\
\midrule
$2,3,4$&$-e_1,-e_2,-v_2$&$1,6,8$&$-e_1,-e_3,v_1$\\
$1,4,5$&$-e_1,-e_2,-v_1$&$2,6,9$&$-e_1,e_2,-e_3$\\
$2,3,6$&$-e_1,-e_3,-v_2$&$1,7,9$&$-e_1,e_2,e_3$\\
$1,4,6$&$-e_1,-e_2,-e_3$&$1,8,9$&$-e_1,e_2,v_1$\\
$3,4,6$&$-e_1,-e_2,-v_6$&$6,8,9$&$-e_1,e_2,-v_6$\\
$2,4,7$&$-e_1,-e_2,e_3$&$2,7,10$&$-e_1,e_3,v_2$\\
$1,5,7$&$-e_1,e_3,-v_1$&$2,9,10$&$-e_1,e_2,v_2$\\
$4,5,7$&$-e_1,-e_2,-v_7$&$7,9,10$&$-e_1,e_2,-v_7$\\
\bottomrule
\end{tabular}
\caption{All maximal candidate sets and their unimodular root-system certificates.}\label{R6E6:tab:systems}
\end{table}

\section{Attributed low-rank core theorem, with proof}\label{R6E6:app:cores}
\begin{srcBJlemma}\label{R6E6:lem:cores} Every pointed unimodular matrix of rank at most two, and every pointed simple rank-three unimodular matrix, has $\mathsf{BL}$.\end{srcBJlemma}
\paragraph{Signs are orientations, not variable reflections.}
First consider a simple rank-three matrix, normalized using Appendix~\ref{R6E6:app:directions}. We have normalized signs only to classify directions. In the original nonnegative fiber no variable is replaced by its negative. Rather, each signed vector in \eqref{R6E6:eq:roots} is the column of a reduced incidence matrix on vertices $\{1,2,3,4\}$: $e_i$ represents $i\to4$, $e_i-e_j$ represents $i\to j$, and a negative sign reverses that edge. Simplicity gives a simple directed graph, and rank three makes its underlying graph connected.

If this graph had a directed cycle, adding any nonnegative integer multiple of its circulation to one feasible flow would make that fiber infinite. The nonempty finite input hypothesis therefore forces an acyclic orientation. Conversely an acyclic directed graph has bounded nonnegative fibers: in a topological order, each prefix cut has fixed nonnegative total flow, and every edge crosses such a cut. This bounds every coordinate.

If all six edges occur, an acyclic orientation of $K_4$ is transitive. In a topological ordering its positive roots become exactly \eqref{R6E6:eq:A3} after the integral change from vertex balances $(\beta_1,\ldots,\beta_4)$ to prefix sums
\[
 a=\beta_1,\qquad b=\beta_1+\beta_2,\qquad
 c=\beta_1+\beta_2+\beta_3,
 \qquad\sum_i\beta_i=0.
\]
Thus the power-concavity assertion in Appendix~\ref{R6E6:app:A3} applies. If an edge is absent, the underlying graph is outerplanar: $K_4$ minus any edge has an embedding as a quadrilateral with one diagonal, and every proper subgraph lies in such a graph. The attributed outerplanar barycentric result, proved again in the next section, applies. This completes the reduction for Lemma~\ref{R6E6:lem:cores}.

\subsection{The inherited outerplanar case, with proof}
We include the exact earlier input to keep the rank-three theorem standalone. No new credit is claimed for this section.

\begin{srcBJlemma}[Outerplanar fibers, Round 4]\label{R6E6:lem:outer}
For an outerplanar directed graph, the counts of integer flows in a fixed finite integer box satisfy finite-support barycentric log-concavity in all vertex balances. The same holds for finite nonnegative fibers.
\end{srcBJlemma}
\begin{proof}
An isolated vertex contributes one at balance zero and zero otherwise, so it may be removed. It is enough to treat a connected loopless graph embedded in a disk with all $k\geq2$ vertices on its boundary; disconnected components multiply. Add a counterclockwise boundary cycle, allowing two arcs if $k=2$. For vertex divergences $\beta$ summing to zero prescribe its edge flows as
\[
 p_i=-\sum_{j=1}^i\beta_j\quad(i<k),\qquad p_k=0.
\]
The augmented flow is a circulation, and $\beta\leftrightarrow(p_1,\ldots,p_{k-1})$ is an integral linear bijection. Distinct new boundary arcs have distinct bounded faces on their left: retain a spanning tree of the old graph; after adding the boundary cycle Euler's formula gives $k$ bounded faces, each needing a boundary arc since the tree has no cycle. Reinstating other edges splits, but never merges, faces.

An integral circulation is uniquely the difference $z_{L(e)}-z_{R(e)}$ of integral face potentials with exterior potential zero. To see existence, decompose it into signed directed simple cycles and add their signed interior-face indicators. Uniqueness follows from connectivity of face adjacency. Bridges have the same face on both sides and circulation zero, so cause no exception. The distinct boundary faces therefore have prescribed potentials $p_i$, including a second zero anchor at $p_k$. All other potentials are summed. The original edge box gives integer difference bounds between potentials. For fixed prescribed potentials there are finitely many allowed assignments, by paths in the connected dual from the exterior.

Let $Z(p)$ denote the number of assignments in any such difference-bound system. The unconstrained potential set is translation-invariant and closed under coordinate maximum and minimum. Thus for two feasible full vectors $z,w$ and integer $\alpha\geq0$,
\[
 z'=(z-\alpha\srcBJone)\vee w,\qquad
 w'=z\wedge(w+\alpha\srcBJone)
\]
are feasible and preserve their sum; zero anchors remain zero. The classical Ahlswede--Daykin four-functions theorem \cite{R5E6:AD}, applied to the summed coordinates with the first set shifted by $-\alpha\srcBJone$, gives
\begin{equation}\label{R6E6:eq:shift}
 Z(p)Z(r)\leq Z((p-\alpha\srcBJone)\vee r)\,
                    Z(p\wedge(r+\alpha\srcBJone)).
\end{equation}
Explicitly, for feasibility indicators $W$, use functions
$A(u)=W(p,u+\alpha\srcBJone)$, $B(v)=W(r,v)$,
$C(t)=W((p-\alpha\srcBJone)\vee r,t)$ and
$D(t)=W(p\wedge(r+\alpha\srcBJone),t+\alpha\srcBJone)$.
Their pointwise four-functions inequality is exactly feasibility of $z',w'$. The theorem compares their finite sums and yields \eqref{R6E6:eq:shift}. This is the marginal-closure mechanism already attributed in Round 4 to Speyer \cite{R4E6:Speyer} and four-functions.

For $N$ input integer potential vectors with integer average, if some coordinate range is at least two, choose ordered input vectors $x,y$ differing by at least two there and use \eqref{R6E6:eq:shift} with $\alpha=1$. Their sum is unchanged and the product of their positive counts cannot decrease. In each coordinate with difference $d=x_i-y_i\leq1$ the outputs simply interchange the values. When $d\geq2$ they become $x_i-1,y_i+1$, decreasing the squared sum by $2(d-1)$. Hence the integer energy about the fixed mean strictly decreases. Eventually all coordinate ranges are at most one; the integer average then forces every vector to equal that mean. This proves Jensen for equal weights and hence for rational weights by repetition. For real finite weights, rational points are dense in the probability polytope defined by the fixed integer mean, since its equations are rational. Continuity of the finite weighted sum of logarithms gives the result.

The integral transformation from balances to $p$ proves the boxed assertion. For finite nonnegative input and central fibers, a sufficiently large common box contains them all. Alternatively, increasing the common boxes and taking limits proves the assertion directly and forces a positive central count. In the acyclic graphs used above all fibers are bounded, as already proved.
\end{proof}

\begin{proof}[Completion of Lemma~\ref{R6E6:lem:cores}]
For rank at most two, normalize a nonsingular maximal minor to an identity as in Appendix~\ref{R6E6:app:directions}. In rank two all column entries and two-by-two minors are $0,\pm1$. Up to signs, only $e_1,e_2,e_1+e_2,e_1-e_2$ can occur, and the last two cannot coexist. Changing one coordinate sign leaves three root directions, which are the edges of a triangle after an integral basis change. Repeated columns become parallel edges; pointedness makes their orientation acyclic. This is an outerplanar multigraph, so Lemma~\ref{R6E6:lem:outer} applies. Rank one is a two-vertex parallel-edge network; rank zero has no variables under pointedness and is trivial.

For a simple rank-three matrix, the relative lattice change, finite direction classification, and inherited signs identify the fibers with those of a connected acyclic simple graph on four vertices. In the six-edge case use the power-concavity assertion in Appendix~\ref{R6E6:app:A3}; otherwise use Lemma \ref{R6E6:lem:outer}. All coordinate changes preserve integer means and the exact counts. These inequalities give positivity of the central count, and acyclicity gives its finiteness.
\end{proof}



\chapter[Canonical Half-Facet Corrections for Unimodular Fibers]{Canonical Half-Facet Corrections for Unimodular Fibers}
\label{source:R7E6}
\noindent\textit{Source development: Round 7, entry 06. Original mathematical arguments retained; notation is local to this chapter.}\par\medskip
\begin{quote}\small\textbf{Integrated source-status qualification.}
Chu--Karshon, Remark 1.1, attributes a halfway-lattice Euler--Maclaurin
analogue to the 2025 thesis of Natalia Juliana Lozano Duarte, supervised
by Alexander Caviedes Castro. See the detailed source qualification in
Section~\ref{syn:geom:framework}. No priority claim for that general
mechanism or its parity phenomenon is made here.\end{quote}

\begin{quote}\small\noindent\textbf{Chapter abstract.} 
For a fixed primitive integral normal system, relaxing every integral facet bound by one half cancels the entire codimension-one term in the classical Berline--Vergne local Euler--Maclaurin formula. We make its uniform consequence explicit, including chamber walls and degenerate original polytopes: the error is bounded by a constant times the diameter to power $d-2$. The half-facet shift is uniquely forced, modulo translations, by second-order accuracy throughout a fixed-normal family. For unimodular normal systems the correction is piecewise polynomial without residue classes, and all odd-codimension local terms vanish. In dimension three this yields an explicit edge-only formula, valid also at nonsimple vertices. Applied to arbitrary-rank unimodular fibers, including varying finite boxes and opposite kernel directions, this gives an explicit piecewise-affine correction restoring cube-root concavity. Finally, affinely Cartesian-indecomposable type-$A$ range polytopes give sharp-order obstructions in every dimension: their uncorrected dimension-matched root is strictly convex, and every asymptotically small concavity correction must dominate an explicit positive polynomial of degree $d-2$. The general error estimate and corrected concavity are applications of classical machinery, not a solution of the uncorrected barycentric log-concavity problem.
\end{quote}

\section{Scope and main statements}
Let $A$ have $m$ primitive nonzero integral rows $a_i\in\srcBPZ^d$, and assume
\begin{equation}\label{R7E6:eq:recession}
 \{v\in\srcBPR^d:Av\le0\}=\{0\}.
\end{equation}
For $h\in\srcBPZ^m$ with $P(h)\ne\varnothing$, put
\[
 P(h)=\{x:Ax\le h\},\qquad Q(h)=\{x:Ax\le h+\tfrac12\srcBPone\},
 \qquad N(h)=\#(P(h)\cap\srcBPZ^d).
\]
Volume in $\srcBPR^d$ is normalized so that $\srcBPZ^d$ has covolume one. On a face it is normalized by the lattice in its parallel linear space. A row system is \emph{unimodular} if every nonsingular $d$-row minor has determinant $\pm1$.

The local coefficients $\nu(Q,F)$ below are the constant terms of the Berline--Vergne transverse-cone functions, using the standard Euclidean scalar product. Their defining classical formula and the properties needed here are stated in Section~\ref{R7E6:sec:EM}. In particular, the correction is specified geometrically, not defined tautologically as count minus volume:
\begin{equation}\label{R7E6:eq:correction}
 \srcBPcR_A(h)=\sum_{\substack{\varnothing\ne F\text{ face of }Q(h)\\\dim F\le d-2}}
       \nu(Q(h),F)\srcBPvol_{\srcBPZ}(F).
\end{equation}

\begin{srcBPtheorem}[Uniform half-facet application of local Euler--Maclaurin]\label{R7E6:thm:uniform}
For $d\ge2$ and fixed $A$ as above, there is an effectively finite constant $C_A$ such that, for every integral $h$ with $P(h)\ne\varnothing$,
\begin{equation}\label{R7E6:eq:uniform}
 N(h)=\srcBPvol Q(h)+\srcBPcR_A(h),\qquad
 |\srcBPcR_A(h)|\le C_A\max\{1,\srcBPdiam Q(h)\}^{d-2}.
\end{equation}
This includes all chamber walls and all lower-dimensional $P(h)$. In dimension one, $N(h)=\srcBPvol Q(h)$ exactly. If $A$ is unimodular, $\srcBPcR_A$ is given on the integral parameters by finitely many rational polynomials of degree at most $d-2$ on a finite affine polyhedral stratification, without a residue-class subdivision. Moreover, only faces of even codimension contribute to \eqref{R7E6:eq:correction}.
\end{srcBPtheorem}

\begin{srcBPtheorem}[Explicit three-dimensional correction]\label{R7E6:thm:edges}
For a unimodular row system $A$ in dimension three,
\begin{equation}\label{R7E6:eq:edges}
 \boxed{\displaystyle
 \srcBPcR_A(h)=\frac1{24}\sum_{E\text{ edge of }Q(h)}
 \ell_{\srcBPZ}(E)\,(a_E\cdot b_E)
 \left(\frac1{\|a_E\|^2}+\frac1{\|b_E\|^2}\right).}
\end{equation}
Here $a_E,b_E$ are the primitive outward normals of the two actual facets containing $E$, and $\ell_{\srcBPZ}(E)$ is its lattice length. Redundant inequalities are not counted as facets. Formula~\eqref{R7E6:eq:edges} remains valid when vertices are nonsimple or $P(h)$ is degenerate.
\end{srcBPtheorem}

For example, for $P(k)=\{x\ge0:\sum_{i=1}^3x_i\le k\}$, the only nonzero terms in \eqref{R7E6:eq:edges} are the three edges where the sum facet meets a coordinate facet. Each has lattice length $k+2$ and coefficient $-1/18$, giving
\[
 \binom{k+3}{3}=\frac{(k+2)^3}{6}-\frac{k+2}{6}.
\]
The correction need not be positive.

\paragraph{Target and attribution.}
Question~1.7 of Leake--Mohammadi Yekta, arXiv:2609.39917v1~\cite{R4E6:LMY}, asks for an uncorrected barycentric inequality for unimodular integer-fiber counts, beyond their midpoint results. Its latest posted version checked on 20 October 2026 was v1. We do \emph{not} settle that question. Round 6 entry 03~\cite{R7E6:R603} obtained an explicit variable shifted-volume correction for the $A_3$ fiber. Round 6 entries 03 and 06~\cite{R7E6:R603,R7E6:R606} independently gave overlapping dimension-matched-root obstructions; their common obstruction is one contribution, not two. The present extension is to all unimodular normal systems and all unimodular nullity-three fibers, with opposite directions and boxes included, together with the canonical-shift characterization and an indecomposable sharpness family.

The uniform counting estimate is explicitly a corollary of the classical local Euler--Maclaurin formula~\cite{R7E6:BV}. Piecewise-polynomial vector-partition counting and Todd/Euler--Maclaurin methods are classical~\cite{R7E6:SV,R7E6:DLS}. Corrected concavity follows from the classical Brunn--Minkowski theorem. Our proposed increment is the precise uniform package, the parity reduction and explicit edge specialization, the uniqueness statement, and the sharp obstruction application below. Historical priority for these special consequences is not established by the literature search.

\section{The classical local formula and its uniform consequence}\label{R7E6:sec:EM}
We use the following classical input, with exact normalization. For a rational polytope $Q$ in a Euclidean rational vector space with lattice $\Lambda$, and a nonempty face $F$, let $\mathfrak t(Q,F)$ be its transverse supporting cone in the quotient by $\srcBPlin(F-F)$, with quotient lattice. Berline and Vergne associate to each such rational affine cone $C$ an analytic germ $\mu(C)(\xi)$ at zero. Set
\[
 \nu(Q,F)=\mu(\mathfrak t(Q,F))(0).
\]
Their local Euler--Maclaurin formula, arXiv v3 Theorem 26~\cite{R7E6:BV}, at the constant polynomial gives
\begin{equation}\label{R7E6:eq:BV}
 \#(Q\cap\Lambda)=\sum_{\varnothing\ne F\text{ face of }Q}
                 \nu(Q,F)\srcBPvol_{\Lambda}(F).
\end{equation}
The top-dimensional coefficient is $1$. The cone functions are invariant under lattice translations (v3 Theorem 20(a)), and their coefficients are computable from the rational cone data (v3 Theorem 21). In the primitive one-dimensional coordinate, with $s\in\srcBPR$,
\begin{equation}\label{R7E6:eq:halfline}
 \mu(s+\srcBPR_{\ge0})(z)=\frac{e^{(\lceil s\rceil-s)z}}{1-e^z}+\frac1z,
 \qquad \mu(s+\srcBPR_{\ge0})(0)=\frac12-(\lceil s\rceil-s).
\end{equation}
These are known results, not theorems claimed as new here.

\begin{proof}[Proof of the identity and bound in Theorem~\ref{R7E6:thm:uniform}]
Every point of $P(h)$ satisfies all inequalities of $Q(h)$ strictly, so $Q(h)$ has nonempty interior even if $P(h)$ does not. Condition~\eqref{R7E6:eq:recession} makes it bounded. Since $a_i\cdot z$ is integral for $z\in\srcBPZ^d$,
\[
 Q(h)\cap\srcBPZ^d=P(h)\cap\srcBPZ^d.
\]
Every actual facet of $Q(h)$ lies in a plane $a_i\cdot x=h_i+1/2$ with primitive $a_i$. Its one-dimensional transverse cone therefore has a half-integral endpoint. Its coefficient in \eqref{R7E6:eq:BV} vanishes by \eqref{R7E6:eq:halfline}. This proves the identity, and proves the dimension-one assertion.

For uniformity, every vertex is
\begin{equation}\label{R7E6:eq:vertex}
 v_I(h)=A_I^{-1}(h_I+\tfrac12\srcBPone)
\end{equation}
for some nonsingular $d$-row submatrix $A_I$. A common denominator dividing twice the least common multiple of the nonzero determinants works for every vertex and every $h\in\srcBPZ^m$. Thus these vertices have only finitely many classes modulo $\srcBPZ^d$. There are finitely many possible face directions and supporting cones, since their defining rows come from a fixed finite list. Passing to any of the finitely many quotient lattices preserves finiteness of the apex classes. Translation invariance of $\mu$ now shows that every local constant appearing in \eqref{R7E6:eq:correction} belongs to a finite, computable list. Let $B_A$ bound their absolute values.

A $k$-face has Euclidean $k$-volume at most $(\srcBPdiam Q)^k$: project onto an orthonormal coordinate frame in its affine hull and enclose it in a box with each side at most the diameter. The covolume of its parallel integer lattice is at least one, since its squared covolume is a positive integer Gram determinant. Thus the same bound applies to normalized volume. There are at most $2^m$ faces, so $C_A=2^m\max(1,B_A)$ suffices. Also \eqref{R7E6:eq:vertex} gives $\srcBPdiam Q\le C'_A(1+\|h\|_\infty)$; hence the assertion is uniformly of order at most $d-2$ in the parameter size as well.
\end{proof}

\begin{srcBPlemma}[No residue dependence]\label{R7E6:lem:polynomial}
If $A$ is unimodular, the correction in \eqref{R7E6:eq:correction} is piecewise polynomial as asserted in Theorem~\ref{R7E6:thm:uniform}.
\end{srcBPlemma}
\begin{proof}
For each basis $I$, $A_I^{-1}h_I\in\srcBPZ^d$ for integral $h$. Consequently
\[
 v_I(h)\equiv \tfrac12A_I^{-1}\srcBPone\pmod{\srcBPZ^d}
\]
is independent of $h$, and $2v_I(h)\in\srcBPZ^d$. Divide parameter space by all affine feasibility and equality tests for the candidate vertices \eqref{R7E6:eq:vertex}. On every relatively open stratum the actual vertex and face incidences, including coincident candidates, are fixed. For each face choose a fixed incident basis by a fixed ordering. Its transverse cone and its apex modulo the quotient lattice are then fixed, so its local coefficient is constant on the integral parameters of the stratum.

A pulling triangulation with a fixed vertex order triangulates every face in that combinatorial type. In a fixed rational coordinate frame of the face direction, each simplex volume is an absolute determinant of affine vertex differences, divided by a fixed lattice normalization. Its sign is fixed on the stratum, since a nonzero simplex in the pulling triangulation cannot reverse orientation without degenerating. Thus each face volume is a rational polynomial of degree at most its dimension. The same argument is made separately on every lower-dimensional stratum. Summing proves the claim. No extra congruence classes are needed, precisely because the displayed vertex classes are fixed, not merely bounded in denominator.
\end{proof}

\section{Parity of local terms and the edge formula}
We record the parity argument since it explains why no vertex correction has been omitted in dimension three. In the following lemma every facet plane is measured using a primitive functional of the relevant lattice.

\begin{srcBPlemma}[Half-integral parity]\label{R7E6:lem:parity}
Let $C=s+c$ be a full-dimensional pointed rational affine cone of dimension $q$, with $2s\in\Lambda$. Suppose every facet plane of $C$ has half-integral, nonintegral offset. Then
\begin{equation}\label{R7E6:eq:parity}
 \mu(C)(-\xi)=(-1)^q\mu(C)(\xi).
\end{equation}
In particular $\mu(C)(0)=0$ when $q$ is odd. For a half-integral polytope all of whose facet offsets are primitive half-integers, every odd-codimension local coefficient in \eqref{R7E6:eq:BV} vanishes.
\end{srcBPlemma}
\begin{proof}
Write $S(C)(\xi)=\sum_{x\in C\cap\Lambda}e^{\langle x,\xi\rangle}$ as its meromorphic continuation. Cone reciprocity gives
\[
 S(s+c)(-\xi)=(-1)^q S(-s+c^\circ)(\xi).
\]
This standard reciprocity follows for simplicial cones by pairing complementary half-open fundamental parallelepipeds, and for rational cones by a triangulation and inclusion--exclusion; see~\cite{R7E6:BR}. No lattice point lies on the boundary of $C$, because a boundary point would satisfy a primitive integral functional equal to a noninteger. Since $2s\in\Lambda$, translation by $-2s$ therefore yields, for $\sigma_C(\xi)=e^{-\langle s,\xi\rangle}S(C)(\xi)$,
\begin{equation}\label{R7E6:eq:sigma}
 \sigma_C(-\xi)=(-1)^q\sigma_C(\xi).
\end{equation}

The defining local recursion (Berline--Vergne v3 Theorem 19(d)) is
\begin{equation}\label{R7E6:eq:recursion}
 \sigma_C(\xi)=\sum_{G\text{ face of }C}
 \mu(\mathfrak t(C,G))(\xi) I(\srcBPlin_+(G-s))(\xi),
\end{equation}
where $I$ is the lattice-normalized exponential cone integral; the zero-dimensional integral is $1$. Quotient germs are pulled back by orthogonal projection. The term for $G=\{s\}$ is $\mu(C)$. A cone integral of dimension $k$ has parity $(-1)^k$ by homogeneity.

Induct on $q$. For a positive-dimensional face $G$, the projected apex is still half-integral. A primitive facet functional containing $G$ descends to a primitive functional on the quotient lattice: its image on that quotient is its original image $\srcBPZ$. Its offset remains half-integral. Thus the induction hypothesis gives parity $(-1)^{q-\dim G}$ for the transverse germ. Every nonvertex term in \eqref{R7E6:eq:recursion} has parity $(-1)^q$. Subtracting them from \eqref{R7E6:eq:sigma} proves \eqref{R7E6:eq:parity}; dimension zero starts with $\mu=1$.

For a polytope face, choose a half-integral incident vertex. Its projection is the transverse apex, so the same argument applies. In a unimodular $Q(h)$ all vertices are half-integral by \eqref{R7E6:eq:vertex}. This completes the parity assertion in Theorem~\ref{R7E6:thm:uniform}.
\end{proof}

\begin{proof}[Proof of Theorem~\ref{R7E6:thm:edges}]
Only codimension two remains in dimension three, by Lemma~\ref{R7E6:lem:parity}. At an edge the two facet normals $a,b$ are a basis of the dual quotient lattice. Indeed, extend these two independent rows to a nonsingular three-row submatrix of $A$; its determinant is $\pm1$. In the two-dimensional quotient, let $v_1,v_2$ be the primitive inward edge generators of the transverse cone, with $a,b$ the negative dual basis, in either ordering. Put
\[
 x=\langle\xi,v_1\rangle,\quad y=\langle\xi,v_2\rangle,\quad
 c_1=\frac{v_1\cdot v_2}{\|v_1\|^2},\quad
 c_2=\frac{v_1\cdot v_2}{\|v_2\|^2}.
\]
The apex coordinates in this basis are half-integers, since both facet offsets are half-integers. Therefore
\[
 \sigma_C(\xi)=\frac1{4\sinh(x/2)\sinh(y/2)}
 =\frac1{xy}-\frac1{24}\left(\frac{x}{y}+\frac{y}{x}\right)+\text{higher degrees}.
\]
The full cone integral is $1/(xy)$. The ray integrals are $-1/x$ and $-1/y$. Orthogonal projection gives transverse coordinates $y-c_1x$ and $x-c_2y$, respectively. Formula~\eqref{R7E6:eq:halfline} at a half-integer is
\[
 \mu_{1/2}(z)=-\frac1{2\sinh(z/2)}+\frac1z=\frac z{24}+O(z^3).
\]
Subtracting the full cone and the two ray terms in \eqref{R7E6:eq:recursion}, the constant term is
\[
 \mu(C)(0)=-\frac{c_1+c_2}{24}.
\]
The inverse-Gram relation for dual bases gives
\[
 c_1+c_2=-(a\cdot b)\left(\|a\|^{-2}+\|b\|^{-2}\right).
\]
The quotient metric is the one induced from $\srcBPR^3$; the covectors $a,b$ annihilate the edge direction, so their inner products are their usual ambient inner products. Multiplying this coefficient by the normalized edge length proves \eqref{R7E6:eq:edges}. No simplicity assumption at the endpoints has been used.
\end{proof}

\section{Why the half shift is canonical}
The following characterization concerns a \emph{fixed normal list throughout its whole bounded full-dimensional family}, not one chamber or one dilation ray. Distinct primitive rows exclude repeated positive parallel normals. Condition~\eqref{R7E6:eq:recession} says that they positively span $\srcBPR^d$; in particular all of them can be simultaneously realized as facets. For example, their distinct unit-normal tangent planes to the unit ball define such a polytope, and a small rational perturbation preserves all its facets.

\begin{srcBPtheorem}[Uniqueness modulo translations]\label{R7E6:thm:unique}
Let $d\ge2$ and let the rows of $A$ be distinct primitive integral vectors satisfying \eqref{R7E6:eq:recession}. For a fixed $\delta\in\srcBPR^m$, the following are equivalent:
\begin{enumerate}
\item For every integral $k$ such that $P(k)$ is a full-dimensional lattice polytope with every listed inequality a facet,
\[
 N(tk)-\srcBPvol P(tk+\delta)=o(t^{d-1})\qquad(t\longrightarrow\infty,\ t\in\srcBPZ).
\]
\item There exists $v\in\srcBPR^d$ with $\delta=\tfrac12\srcBPone+Av$.
\end{enumerate}
In case (2), $P(h+\delta)=Q(h)+v$, so the uniform second-order bound of Theorem~\ref{R7E6:thm:uniform} holds for all admissible integral $h$.
\end{srcBPtheorem}
\begin{proof}
Let $S_i$ be the Euclidean area of the facet with normal $a_i$ in $P(k)$, and put $w_i=S_i/\|a_i\|$. The classical second coefficient of the Ehrhart polynomial is half the normalized boundary volume~\cite{R7E6:BR}; hence
\[
 N(tk)=t^d\srcBPvol P(k)+\tfrac12t^{d-1}\sum_iw_i+O(t^{d-2}).
\]
The first variation of volume under parallel facet motions is the sum of facet area times signed displacement. This follows by taking the volume of the thin facet prisms; overlaps near their boundaries are $o(\varepsilon)$ under an $\varepsilon$ perturbation. Thus
\[
 \srcBPvol P(tk+\delta)=t^d\srcBPvol P(k)+t^{d-1}\sum_i\delta_iw_i+o(t^{d-1}).
\]
Assertion (1) forces $(\delta-\tfrac12\srcBPone)\cdot w=0$.

Every rational support vector with all facets active can be multiplied by an integer to give an integral support vector and integral vertices. The preceding relation is homogeneous in the facet areas and thus holds before that multiplication. Rational support vectors are dense in the open set where all facets are active, and facet areas are continuous there. The relation therefore holds for every real polytope with exactly these normals.

Minkowski's existence theorem for polytopes (for example~\cite[Theorem 2]{R7E6:Klain}) says that all positive vectors $w$ with $A^Tw=0$ occur: assign unit normals $a_i/\|a_i\|$ and areas $w_i\|a_i\|$. Their vector sum is zero and the normals span the space. A strictly positive such vector exists, and the positive vectors in $\ker A^T$ span that kernel (perturb a fixed positive one by small multiples of an arbitrary kernel vector). Consequently $\delta-\tfrac12\srcBPone$ annihilates $\ker A^T$ and lies in $\operatorname{im}A$, proving (2). Conversely (2) gives the asserted translation identity, and Theorem~\ref{R7E6:thm:uniform} implies (1).
\end{proof}

\section{Application to unimodular fibers, including boxes}
Let $M\in\srcBPZ^{p\times n}$ have rank $r$ and nullity $d=n-r\ge1$. We use the same convention as in earlier rounds: every nonsingular $r\times r$ submatrix has determinant $\pm1$; redundant rows are allowed. Margins lie in $\operatorname{im}_{\srcBPR}M\cap\srcBPZ^p$, which equals $M\srcBPZ^n$ under this convention. All volumes in a fiber are normalized by $\ker M\cap\srcBPZ^n$.

For finite integral bounds $\ell,u$ and a real-nonempty fiber, write
\[
 F(b,\ell,u)=\#\{x\in\srcBPZ^n:Mx=b,\ \ell\le x\le u\}.
\]
Alternatively allow the nonnegative fiber $F_M(b)$, provided $M$ is pointed, meaning $\ker M\cap\srcBPR_{\ge0}^n=\{0\}$. This is exactly the boundedness condition for nonempty nonnegative real fibers.

\begin{srcBPcorollary}[Uniform corrections in every nullity]\label{R7E6:cor:fibers}
For each such fixed $M$ there is an explicit geometric correction $\srcBPcR$ with the following properties, on the admissible integral parameters:
\begin{enumerate}
\item $F-\srcBPcR$ is the normalized volume of the fiber with every nonconstant lower bound decreased by $1/2$ and every nonconstant upper bound increased by $1/2$.
\item $\srcBPcR$ is piecewise polynomial of degree at most $d-2$ for $d\ge2$, with the uniform intrinsic bound \eqref{R7E6:eq:uniform}. It is zero for $d=1$. In nullity three it is the explicit piecewise-affine edge sum \eqref{R7E6:eq:edges} in any integral kernel coordinates.
\item For finitely many admissible integral parameter vectors $\theta_i$, weights $\lambda_i\ge0$ with sum one, and integral mean $\bar\theta=\sum_i\lambda_i\theta_i$,
\begin{equation}\label{R7E6:eq:correctedBM}
 \big(F(\bar\theta)-\srcBPcR(\bar\theta)\big)^{1/d}
 \ge\sum_i\lambda_i\big(F(\theta_i)-\srcBPcR(\theta_i)\big)^{1/d}.
\end{equation}
\end{enumerate}
There is no restriction on rank, opposite kernel directions, or varying box sizes. These are corrected inequalities, not assertions about the uncorrected count.
\end{srcBPcorollary}
\begin{proof}
Choose $r$ independent rows and $r$ basis columns. After permuting columns and applying the integral inverse basis matrix, the effective equations become $x_B+Nx_J=\beta$. The selected equations imply all redundant equations for compatible margins. Put
\[
 x=x_0(b)+Kt,\qquad x_0=(\beta,0),\qquad K=\begin{pmatrix}-N\\I_d\end{pmatrix}.
\]
Then $t\mapsto Kt$ identifies $\srcBPZ^d$ with the full kernel lattice. Complementary maximal minors show that every nonzero $d$-row minor of $K$ is $\pm1$: selecting the identity rows reduces such a minor to a square minor of $N$, which is also a maximal minor of $[I_r\ N]$ after adjoining complementary identity columns. In particular each nonzero row of $K$ is primitive, since it extends to a row basis of determinant $\pm1$.

For a box the normal system is $A=(K^T,-K^T)^T$, with parameter $h=(u-x_0,x_0-\ell)$ and zero rows removed. More plainly, its rows are the rows of $K$ and their negatives. It is unimodular: a maximal minor using an opposite pair vanishes and every other maximal minor equals, up to sign, one of $K$. A zero row imposes only a deterministic inequality on $b,\ell,u$. Such inequalities are retained in the admissible parameter domain, but omitted from the polytope description. Since the original parameters are integral and feasible, relaxing those deterministic bounds would not affect the fiber either. The box makes the recession cone trivial. For a pointed nonnegative fiber use $A=-K$ and $h=x_0$ after the same zero-row convention; pointedness gives \eqref{R7E6:eq:recession}. Theorems~\ref{R7E6:thm:uniform} and~\ref{R7E6:thm:edges} now apply, including at degenerate original fibers.

For real admissible parameters let $Q_\theta$ be the relaxed polytope in these fixed coordinates. Its right-hand sides depend affinely on $\theta$, so
\[
 \sum_i\lambda_iQ_{\theta_i}\subseteq Q_{\bar\theta}.
\]
The classical Brunn--Minkowski inequality~\cite{R7E6:Gardner} and monotonicity of volume prove \eqref{R7E6:eq:correctedBM}. The deterministic feasibility inequalities are affine, so remain valid at the mean. Unimodularity ensures an integer point in every nonempty bounded integral fiber: its vertices in the $t$-coordinates are integral. This also confirms that the central count is nonzero.
\end{proof}

For a pointed nonnegative fiber, translating the relaxed variable by $\srcBPone/2$ identifies its volume with the ordinary nonnegative fiber volume at margin $b+M\srcBPone/2$. This is the familiar half-column-sum shift underlying the $A_3$ formula of Round 6 entry 03. Corollary~\ref{R7E6:cor:fibers} supplies the uniform lower-order error, including walls, and the explicit all-nullity-three formula; the shift itself is not claimed as new.

\section{An indecomposable sharpness family in every dimension}
Let $d\ge2$, let $z_0=0$, and define
\begin{equation}\label{R7E6:eq:range}
 H_d(t)=\{z\in\srcBPR^d: z_i-z_j\le t\text{ for all }0\le i\ne j\le d\},
 \qquad t\ge0.
\end{equation}
Equivalently the range of $0,z_1,\ldots,z_d$ is at most $t$. Its normal system consists of the type-$A$ roots $\pm e_i$ and $e_i-e_j$. This classical configuration is totally unimodular: its transpose is a reduced directed incidence matrix, whose minors are $0,\pm1$ by expanding a column with at most one nonzero entry, or, if none exists, using the dependence of the rows. We claim no novelty for this configuration or for the elementary counting identity below.

\begin{srcBPtheorem}[Exact, indecomposable sharp-order obstruction]\label{R7E6:thm:range}
For integers $t\ge0$, put $s=t+1/2$. Then
\begin{align}
 n_d(t):=\#(H_d(t)\cap\srcBPZ^d)&=(t+1)^{d+1}-t^{d+1},\label{R7E6:eq:rangecount}\\
 \srcBPvol H_d(t+1/2)&=(d+1)s^d,\label{R7E6:eq:rangevol}\\
 R_d(t):=n_d(t)-(d+1)s^d
 &=\sum_{\substack{3\le j\le d+1\\j\text{ odd}}}
       \frac{\binom{d+1}{j}}{2^{j-1}}s^{d+1-j}.\label{R7E6:eq:ranger}
\end{align}
The polynomial $R_d$ is positive, has degree $d-2$, and has leading coefficient $(d+1)d(d-1)/24$. The real extension $n_d(t)^{1/d}$ is strictly convex on $t\ge0$.

Suppose a real correction $E(t)=o(t^{d-1})$ is defined on all sufficiently large integers, $n_d(t)-E(t)>0$, and $(n_d(t)-E(t))^{1/d}$ is eventually concave as a sequence. Then
\begin{equation}\label{R7E6:eq:necessary}
 E(t)\ge R_d(t)
\end{equation}
throughout that terminal concavity domain. Equality is attained by $E=R_d$, which gives an affine corrected root. For $t>0$, $H_d(t)$ is not affinely equivalent to a Cartesian product of two positive-dimensional polytopes. These counts are nonnegative fibers of a fixed totally unimodular matrix of nullity $d$.
\end{srcBPtheorem}
\begin{proof}
For an integral point $z$, subtract the minimum of $0,z_1,\ldots,z_d$ from all $d+1$ entries. This gives a unique vector in $\{0,\ldots,t\}^{d+1}$ with minimum zero. The inverse map is $z_i=u_i-u_0$. There are $(t+1)^{d+1}-t^{d+1}$ such vectors, proving \eqref{R7E6:eq:rangecount}.

Partition $H_d(t)$ according to which of its $d+1$ entries realizes the minimum; ties have measure zero. If the minimum is entry zero the region is $[0,t]^d$. If it is entry $j>0$, use the $d$ differences $z_i-z_j$ for $i\ne j$, including $i=0$. These range independently over $[0,t]$ and determine $z$ by an integral linear map of determinant $\pm1$. Each region therefore has volume $t^d$. All primitive inequalities have right-hand side $t$, so half-facet relaxation is exactly $H_d(t+1/2)$. This proves \eqref{R7E6:eq:rangevol}. Expanding $(s+1/2)^{d+1}-(s-1/2)^{d+1}$ gives \eqref{R7E6:eq:ranger} and its leading coefficient.

For real $t\ge0$,
\[
 n_d(t)=(d+1)\int_{-1/2}^{1/2}(s+u)^d\,du.
\]
Thus $n_d(t)^{1/d}$ is $(d+1)^{1/d}$ times the $L^d$ norm of the nonnegative function $u\mapsto s+u$. The triangle inequality gives convexity, and strict convexity of the $L^d$ norm gives strictness: for two different values of $s$, these affine functions are not proportional almost everywhere. Hence, in particular, every consecutive integer midpoint violates concavity of this dimension-matched root.

Here is the asymptotic-slope argument, also used in Round 6 entry 06~\cite{R7E6:R606}, reproduced to make the stronger conclusion self-contained. Write $\alpha=(d+1)^{1/d}$ and $f(t)=(n_d(t)-E(t))^{1/d}$. Equations~\eqref{R7E6:eq:ranger} and $E=o(t^{d-1})$ imply
\[
 f(t)=\alpha(t+1/2)+o(1).
\]
On a terminal domain of sequence concavity, the differences $f(t+1)-f(t)$ decrease, and their limit must be $\alpha$ by this asymptotic. They are consequently at least $\alpha$. Thus $f(t)-\alpha(t+1/2)$ is nondecreasing and tends to zero, so it is nonpositive throughout that domain. Raising to power $d$ proves \eqref{R7E6:eq:necessary}. The exact correction in \eqref{R7E6:eq:ranger} attains equality. This strengthens an order obstruction to domination of the entire explicit polynomial for this family.

Every listed root gives an actual facet when $t>0$: place the two relevant entries at the extremes and all other entries strictly between; the convention $z_0=0$ is accommodated by translation of the entries. If a full-dimensional polytope is affinely a Cartesian product, its facet normals partition into two nonempty classes lying in complementary linear subspaces. The three directions $e_i,e_j,e_i-e_j$ form a minimal linear dependence for every $i\ne j$; a minimal dependence cannot cross a direct-sum partition. Thus all $e_i$ lie in one class, and then so do all root directions. This contradicts a nontrivial product decomposition.

Finally introduce nonnegative variables $y_i=t-z_i$, $x_{0i}=t+z_i$ for $1\le i\le d$, and $x_{ij}=t-z_i+z_j$ for distinct $i,j\in\{1,\ldots,d\}$. They satisfy
\begin{equation}\label{R7E6:eq:rangefiber}
 x_{0i}+y_i=2t,\qquad x_{ij}-y_i+y_j=t.
\end{equation}
Conversely these equations and nonnegativity recover exactly \eqref{R7E6:eq:range}. With the $x$ variables first, the equation matrix is $[I_{d^2}\ N]$, where the rows of $N$ are $e_i$ and $-e_i+e_j$. The latter matrix is totally unimodular by the incidence argument, and adjoining identity columns preserves total unimodularity. There are $d^2+d$ variables and rank $d^2$, so the nullity is $d$. The fibers are bounded, as is already clear from $0\le y_i\le2t$.
\end{proof}

For $d=2$, $R_2=1/4$, consistent with the sharp quarter correction of Round 5 entry 03~\cite{R7E6:R503}. For $d=3$, $R_3=t+1/2$; the first counts are $1,15,65,175,369$. In every $d\ge3$, no constant correction works on this fixed family. Unlike the earlier direct-product witnesses, these polytopes are affinely Cartesian-indecomposable. Moreover Theorem~\ref{R7E6:thm:unique} shows that no other fixed facet shift can uniformly improve the error order in this normal system: every shift with second-order accuracy is a translation of the half shift, whose nonzero error \eqref{R7E6:eq:ranger} has exact degree $d-2$.

\section{Verification, limitations, and next questions}
The proofs are structural and do not depend on the diagnostic search. The accompanying \texttt{verify.py} uses exact rational arithmetic to reconstruct three-dimensional vertices, actual facets, edges, lattice lengths, and volumes, and compares \eqref{R7E6:eq:edges} against direct lattice-point enumeration. It includes degenerate original polytopes, redundant inequalities, nonsimple relaxed vertices, translated bounds, and $A_3$ fibers. It also checks the range-family formulas, the explicit fiber representation, and the relevant symbolic coefficients. The log records the exact finite scope; no finite sample is presented as proof of a universal claim.

The principal counting engine, its effective computability, first-order Ehrhart coefficient, Minkowski existence, and Brunn--Minkowski inequality are classical inputs. The edge formula depends on unimodularity of the transverse two-dimensional lattice cone; dropping that hypothesis introduces arithmetic terms. General primitive normal systems still satisfy the uniform bound, but may have residue dependence and odd-codimension terms. Corrections may be negative and are not claimed to be monotone. The corrected inequality does not imply the uncorrected logarithmic inequality asked in~\cite{R4E6:LMY}. Weighted counts and infinite-support barycentric distributions are not asserted here. For a concrete arithmetic negative control, $P=\{(x,y,z)\ge0:x+y+2z\le1\}$ has three lattice points, while its half relaxation has volume $343/96$ and the naive edge expression in \eqref{R7E6:eq:edges} is $-49/96$. Their sum is $49/16\ne3$: primitive normals alone do not justify the unimodular edge formula.

An initial attempt to prove uncorrected barycentric log-concavity for all boxed nullity-three systems produced no proof or counterexample. Exact small-box tests and exploratory weighted tests were inconclusive; they are not used above. The most direct next problem is whether the explicit edge correction can be combined with the cubic chamber volume to control the uncorrected logarithm across every chamber. The parity lemma also suggests computing the codimension-four term explicitly, but no such formula or general uncorrected inequality is claimed.



\chapter{Source coverage and verification boundary}\label{chap:audit}

All forty-two supplied \texttt{verify.py} programs were inspected and
executed successfully in the audit environment. The arithmetic,
tournament, special-function, structural-permanent, operator-permanent,
and geometry groups contain respectively $7,7,4,7,8,9$ programs.
Some runs required the locally installed SymPy and mpmath dependencies;
initial missing-package errors were resolved before the successful runs.
Where a program writes results, execution used a working copy or an
audit directory. The original mathematical source texts are retained.

The checks have different logical roles. Some verify exact identities
or a finite certificate whose completeness is proved in the text;
others exercise bounded examples of an independently proved uniform
statement. In particular the complete finite Gaussian coefficient
certificates were run, rather than only their quick mode. None of the
random or bounded matrix/tournament tests is used as a substitute for
an all-order proof. Additional exact checks include the optimized
operator constants, Gaussian pair constants, and a prime-power Paley
construction. The fixed-seed Ewald grafting proof, optimized operator
radius, and tournament-pair sampling normalization received independent
cross-checks. No substantive proof gap was identified within this audit;
that statement is not an independent peer-review certification.

\section*{Recoverable verification supplement}
The \texttt{.tex} file is self-contained for compilation with pdfLaTeX;
it needs no external figures, bibliography file, or source inclusions.
After its \verb|\end{document}| it contains a base64-encoded ZIP in
comment lines. This carries the supplied verification programs and
supporting text/data, plus the integration audit notes and additional
checkers. It is ignored by TeX. The following Python standard-library
code recovers it from the same TeX file; run recovered programs only
after inspecting them and installing their stated dependencies.

\begin{quote}\small
\begin{verbatim}
from pathlib import Path
import base64, io, zipfile
p = Path("Keio_Mathematics_Tournament_Integrated_2026.tex")
s = p.read_text()
a = "% BEGIN INTEGRATED VERIFICATION ZIP BASE64\n"
b = "% END INTEGRATED VERIFICATION ZIP BASE64"
block = s.split(a, 1)[1].split(b, 1)[0]
data = base64.b64decode("".join(
    line[2:] for line in block.splitlines()))
with zipfile.ZipFile(io.BytesIO(data)) as z:
    z.extractall("integrated_verification")
\end{verbatim}
\end{quote}

The source identifiers below use \texttt{RrEe} for Round $r$, entry $e$.
The map contains each of the forty-two research texts exactly once.
Duplicated principal conclusions retain both proof developments but
are counted only once in the integrated contribution discussion.

\begin{longtable}{@{}p{.10\textwidth}p{.77\textwidth}r@{}}\toprule Source & Mathematical development & Ch.\\\midrule\endhead

R1E1 & An elementary proof of a cubic power-sum conjecture and a sharp fourth-power strengthening & \ref{source:R1E1}\\

R1E2 & A Deficiency Budget for Strong Seymour Vertices under Tournament Perturbations & \ref{source:R1E2}\\

R1E3 & Finite Gaussian--Newton certificates for higher-order q-log-concavity of Gaussian columns & \ref{source:R1E3}\\

R1E4 & A unique-minimum theorem for a logarithmic gamma quotient & \ref{source:R1E4}\\

R1E5 & Determinant--Pairing Certificates and a Weak-Correlation Region of Chollet's Permanent Conjecture & \ref{source:R1E5}\\

R1E6 & Sharp Ewald-Point Counts in Sequential Simplex-Bundle Towers & \ref{source:R1E6}\\

R2E1 & Three p-adic branches of power-sum block defects and an endpoint obstruction to naive lifting & \ref{source:R2E1}\\

R2E2 & Locally Resilient Strong Second Neighborhoods in Doubly Regular Tournaments & \ref{source:R2E2}\\

R2E3 & A strong permanent inequality for complex positive semidefinite cactus matrices & \ref{source:R2E3}\\

R2E4 & A global parameter classification and moving minima for a logarithmic gamma quotient & \ref{source:R2E4}\\

R2E5 & Strong Self-Chollet Inequalities for Complex Cycle Blocks and Positive Semidefinite Matrices with Cut Vertices & \ref{source:R2E5}\\

R2E6 & Branching Simplex Bundles and Improved Upper Bounds for Ewald-Point Minima & \ref{source:R2E6}\\

R3E1 & Endpoint stability of power-sum valuations at regular and locally nonirregular primes & \ref{source:R3E1}\\

R3E2 & Arbitrarily Sparse Strong Seymour Sets in Regular Quasirandom Tournaments & \ref{source:R3E2}\\

R3E3 & A positive permanent-defect identity for two-hub bipartite matrices & \ref{source:R3E3}\\

R3E4 & One structured factor is enough: Bapat--Sunder bounds for long-cycle blocks & \ref{source:R3E4}\\

R3E5 & Permanent-on-top certificates for complex cycle blocks and overlapping bipartite blocks & \ref{source:R3E5}\\

R3E6 & An Exponential Barrier for Arbitrary-Facet Simplex-Bundle Towers & \ref{source:R3E6}\\

R4E1 & A Wilson-quotient dichotomy for exponent valuation trees & \ref{source:R4E1}\\

R4E2 & The Sharp Spectral Scale for\ Strong Seymour Matching Deficiency & \ref{source:R4E2}\\

R4E3 & Reflected PSD companions and permanent bounds for intersecting odd cycles & \ref{source:R4E3}\\

R4E4 & Polynomial-radius dense neighborhoods with permanent on top & \ref{source:R4E4}\\

R4E5 & Quantitative permanent-on-top gaps and dense operator neighborhoods & \ref{source:R4E5}\\

R4E6 & Barycentric Log-Concavity for Disk-Network Fibers and Rank-Two Unimodular Matrices & \ref{source:R4E6}\\

R5E1 & Finite boundary exceptions and mean valuations of power sums & \ref{source:R5E1}\\

R5E2 & A Sharp Three-Quarter Law and Spectral Tail Bounds for Strong Seymour Matchings & \ref{source:R5E2}\\

R5E3 & A sharp quarter correction for two-dimensional unimodular fibers & \ref{source:R5E3}\\

R5E4 & A universal weak-correlation region with permanent on top & \ref{source:R5E4}\\

R5E5 & Arbitrary-support permanent-on-top balls via cycle-to-edge comparison & \ref{source:R5E5}\\

R5E6 & Sharp Power Concavity for the A\_3 Partition Extension and Simple Rank-Three Unimodular Fibers & \ref{source:R5E6}\\

R6E1 & Distinct-shift correlations and maximal-spike formulas for power-sum valuations & \ref{source:R6E1}\\

R6E2 & A Sharp Four-Fifths Law for Prime Regular Tournaments & \ref{source:R6E2}\\

R6E3 & Optimal-order corrections for power concavity of integer flows & \ref{source:R6E3}\\

R6E4 & Dimension-independent Hadamard contraction for weakly interacting permanents & \ref{source:R6E4}\\

R6E5 & Row-local Bapat--Sunder bounds without a global mass restriction & \ref{source:R6E5}\\

R6E6 & Series Invariance and a Dimension-Three Obstruction for Unimodular Fiber Concavity & \ref{source:R6E6}\\

R7E1 & Gaussian square sums at composite side lengths: exact valuations and effective moments & \ref{source:R7E1}\\

R7E2 & Prime Regularization Preserving Strong Seymour Failures and Quasirandom Tournaments of Vanishing Strong Density & \ref{source:R7E2}\\

R7E3 & A sharp three-step root factorization for Gaussian columns and its fourth-fold log-concavity consequence & \ref{source:R7E3}\\

R7E4 & A row-local permanent-on-top theorem via heap derivatives & \ref{source:R7E4}\\

R7E5 & A dimension-independent Euclidean convexity region for the log-permanent & \ref{source:R7E5}\\

R7E6 & Canonical Half-Facet Corrections for Unimodular Fibers & \ref{source:R7E6}\\

\bottomrule\end{longtable}

\begin{thebibliography}{999}
\small\raggedright
\bibitem{R1E1:SM}
J. Sondow and K. MacMillan, \emph{Reducing the Erd\H{o}s--Moser equation $1^n+2^n+\cdots+k^n=(k+1)^n$ modulo $k$ and $k^2$}, Integers \textbf{11} (2011), A34, 8 pp. Published 27 May 2011.\newline
\url{https://math.colgate.edu/~integers/l34/l34.pdf}.

\bibitem{R1E1:SMpre}
J. Sondow and K. MacMillan, preprint version of \cite{R1E1:SM}, 9 November 2010, arXiv:1011.2154v1.\newline
\url{https://arxiv.org/abs/1011.2154v1}.

\bibitem{R1E1:M}
R. Me\v{s}trovi\'{c}, \emph{On a congruence modulo $n^3$ involving two consecutive sums of powers}, Journal of Integer Sequences \textbf{17} (2014), Article 14.8.4, 20 pp.\newline
\url{https://cs.uwaterloo.ca/journals/JIS/VOL17/Mestrovic/mes4.html}.

\bibitem{R1E2:BLP}
Y. Bai, B. Li and B. Park, \emph{Towards a strengthening of the second neighborhood conjecture}, arXiv:2607.18047v1, 20 July 2026. Conjectures 1.2--1.3 and Theorems 1.4--1.5.\newline
\url{https://arxiv.org/html/2607.18047v1}. Status update: version 2, 24 July 2026, Remark 3.1, \url{https://arxiv.org/html/2607.18047v2}.

\bibitem{R1E2:GN}
K. Gunderson and J. D. Nir, \emph{The oriented chromatic number of random graphs of bounded degree}, arXiv:2202.12323, first submitted 24 February 2022. Section 3 gives the double-regularity parameters and Paley examples.\newline
\url{https://arxiv.org/abs/2202.12323}.

\bibitem{R1E3:MS}
P. R. W. McNamara and B. E. Sagan,
\emph{Infinite log-concavity: developments and conjectures},
Advances in Applied Mathematics \textbf{44} (2010), 1--15.
Preprint arXiv:0808.1065v2, March 24, 2009.
\href{https://arxiv.org/abs/0808.1065}{arxiv.org/abs/0808.1065}.
\href{https://doi.org/10.1016/j.aam.2009.03.001}{doi:10.1016/j.aam.2009.03.001}.

\bibitem{R1E3:BC}
P. Br\"and\'en and M. Chasse,
\emph{Infinite log-concavity for polynomial P\'olya frequency sequences},
preprint arXiv:1405.6378 (2014).
\href{https://arxiv.org/abs/1405.6378}{arxiv.org/abs/1405.6378}.

\bibitem{R1E4:BZ}
T. Bulboac\u{a} and H. M. Zayed,
\emph{Monotonic nature of the Gamma function},
Journal of Inequalities and Applications \textbf{2026}, article 27 (2026),
published 22 January 2026.
\href{https://doi.org/10.1186/s13660-025-03425-0}{doi:10.1186/s13660-025-03425-0}.

\bibitem{R1E4:KS}
P. A. Kup\'an and R. Sz\'asz,
\emph{Monotonicity theorems and inequalities for the gamma function},
Mathematical Inequalities \& Applications \textbf{17} (2014), 149--159.
\href{https://doi.org/10.7153/mia-17-11}{doi:10.7153/mia-17-11}.
\href{https://files.ele-math.com/articles/mia-17-11.pdf}{Publisher PDF}.

\bibitem{R1E4:YT}
Z.-H. Yang and J. Tian,
\emph{Monotonicity and sharp inequalities related to gamma function},
Journal of Mathematical Inequalities \textbf{12} (2018), 1--22.
\href{https://doi.org/10.7153/jmi-2018-12-01}{doi:10.7153/jmi-2018-12-01}.
\href{https://files.ele-math.com/articles/jmi-12-01.pdf}{Publisher PDF}.

\bibitem{R1E5:chollet}
J. Chollet, \emph{Is there a permanental analogue to Oppenheim's inequality?}, Amer. Math. Monthly \textbf{89} (1982), 57--58. \href{https://doi.org/10.1080/00029890.1982.11995380}{doi:10.1080/00029890.1982.11995380}.

\bibitem{R1E5:lieb}
E. H. Lieb, \emph{Proofs of some conjectures on permanents}, J. Math. Mech. \textbf{16} (1966), 127--134. \href{https://doi.org/10.1512/iumj.1967.16.16008}{doi:10.1512/iumj.1967.16.16008}.

\bibitem{R1E5:edmonds}
J. Edmonds, \emph{Maximum matching and a polyhedron with 0,1-vertices}, J. Res. Nat. Bur. Standards Sect. B \textbf{69B} (1965), 125--130. \href{https://nvlpubs.nist.gov/nistpubs/jres/69B/jresv69Bn1-2p125_A1b.pdf}{Original NIST paper}; \href{https://doi.org/10.6028/jres.069B.013}{doi:10.6028/jres.069B.013}.

\bibitem{R1E5:zhang}
F. Zhang, \emph{An update on a few permanent conjectures}, Spec. Matrices \textbf{4} (2016), 305--316. \href{https://arxiv.org/abs/1608.02844}{arXiv:1608.02844}; \href{https://doi.org/10.1515/spma-2016-0030}{doi:10.1515/spma-2016-0030}.

\bibitem{R1E5:rodtes}
K. Rodtes, \emph{Chollet's permanent conjecture for $4\times4$ matrices}, Linear Multilinear Algebra \textbf{72} (2024), 2633--2638. \href{https://arxiv.org/abs/2211.02839}{arXiv:2211.02839}; \href{https://doi.org/10.1080/03081087.2023.2279150}{doi:10.1080/03081087.2023.2279150}.

\bibitem{R1E5:li}
Q. Li, Q. Zhao, and W. Luo, \emph{Chollet's permanent conjecture through order six via border induction}, preprint, 31 August 2026. \url{https://www.preprints.org/manuscript/202608.2196}. Cited as a claim, not used in a proof.

\bibitem{R1E5:ma}
Y. Ma, \emph{Quantitative Chollet inequalities for matrices of rank at most four: Spectral bounds and an order-nine tight-frame construction}, preprint, 28 September 2026. \url{https://arxiv.org/abs/2609.37461}. Context only; not used in a proof.

\bibitem{R1E6:CPS}
L. Crespo, \'A. Pelayo, and F. Santos,
\emph{Ewald's Conjecture and integer points in algebraic and symplectic toric geometry},
Proceedings of the London Mathematical Society \textbf{132} (2026), no.~4, e70144.
\href{https://doi.org/10.1112/plms.70144}{doi:10.1112/plms.70144}.
Preprint version 3, 20 March 2024:
\href{https://arxiv.org/abs/2310.10366v3}{arXiv:2310.10366v3}.
Full text consulted:
\url{https://arxiv.org/html/2310.10366v3}.

\bibitem{R2E1:ST}
J. Sondow and E. Tsukerman, \emph{The $p$-adic order of power sums, the Erd\H{o}s--Moser equation, and Bernoulli numbers}, preprint, 2014, arXiv:1401.0322.\newline
\url{https://arxiv.org/abs/1401.0322}.

\bibitem{R2E1:AMSV}
T. Amdeberhan, V. H. Moll, V. Sharma and C. Vignat, \emph{The valuations of power sums}, Scientia, Series A: Mathematical Sciences \textbf{35} (2025), 115--129.\newline
\url{https://doi.org/10.71712/fwa0-kj92}.

\bibitem{R2E1:R1}
KEIO (AI-assisted research experiment), \emph{An elementary proof of a cubic power-sum conjecture and a sharp fourth-power strengthening}, Round 1, Entry 01, 20 October 2026, Theorem 4 and Proposition 6. Internal tournament manuscript, not an externally published or peer-reviewed source.

\bibitem{R2E2:BLP}
Y. Bai, B. Li, and B. Park, \emph{Towards a strengthening of the second neighborhood conjecture}, arXiv:2607.18047, version 2 (24 July 2026), especially Remark 3.1. \url{https://arxiv.org/html/2607.18047v2}.

\bibitem{R2E2:SV}
B. Sudakov and V. H. Vu, \emph{Local resilience of graphs}, Random Structures \& Algorithms \textbf{33} (2008), 409--433. DOI: \href{https://doi.org/10.1002/rsa.20235}{10.1002/rsa.20235}. \url{https://arxiv.org/abs/0706.4104}.

\bibitem{R2E2:R1}
KEIO, \emph{A Deficiency Budget for Strong Seymour Vertices under Tournament Perturbations}, competition entry 02, Round 1 (2026), internal manuscript. Cyclic-link regularity reused and proved independently above; its open-status framing is corrected here.

\bibitem{R2E3:Chollet}
J. Chollet, \emph{Is there a permanental analogue to Oppenheim's inequality?}, American Mathematical Monthly \textbf{89} (1982), 57--58. \href{https://doi.org/10.1080/00029890.1982.11995380}{doi:10.1080/00029890.1982.11995380}.

\bibitem{R2E3:Lieb}
E. H. Lieb, \emph{Proofs of some conjectures on permanents}, Journal of Mathematics and Mechanics \textbf{16} (1966), 127--134. \href{https://doi.org/10.1512/iumj.1967.16.16008}{doi:10.1512/iumj.1967.16.16008}.

\bibitem{R2E3:BS}
R. B. Bapat and V. S. Sunder, \emph{An extremal property of the permanent and the determinant}, Linear Algebra and its Applications \textbf{76} (1986), 153--163. \href{https://doi.org/10.1016/0024-3795(86)90220-X}{doi:10.1016/0024-3795(86)90220-X}.

\bibitem{R2E3:PS}
P. Pant and R. Singh, \emph{Structural Classes for Chollet's Permanent Conjecture}, arXiv:2604.24192v2, 18 September 2026. \url{https://arxiv.org/abs/2604.24192v2}.

\bibitem{R2E3:Zhang}
F. Zhang, \emph{An update on a few permanent conjectures}, Special Matrices \textbf{4} (2016), 305--316. \url{https://arxiv.org/abs/1608.02844}.

\bibitem{R2E4:KS}
P. A. Kup\'an and R. Sz\'asz,
\emph{Monotonicity theorems and inequalities for the gamma function},
Mathematical Inequalities \& Applications \textbf{17} (2014), 149--159.
\href{https://doi.org/10.7153/mia-17-11}{doi:10.7153/mia-17-11}.
\href{https://files.ele-math.com/articles/mia-17-11.pdf}{Primary PDF}.

\bibitem{R2E4:YT}
Z.-H. Yang and J. Tian,
\emph{Monotonicity and sharp inequalities related to gamma function},
Journal of Mathematical Inequalities \textbf{12} (2018), 1--22.
\href{https://doi.org/10.7153/jmi-2018-12-01}{doi:10.7153/jmi-2018-12-01}.
\href{https://files.ele-math.com/articles/jmi-12-01.pdf}{Primary PDF}.

\bibitem{R2E4:BZ}
T. Bulboac\u{a} and H. M. Zayed,
\emph{Monotonic nature of the Gamma function},
Journal of Inequalities and Applications \textbf{2026}, article 27 (22 January 2026).
\href{https://doi.org/10.1186/s13660-025-03425-0}{doi:10.1186/s13660-025-03425-0}.

\bibitem{R2E4:SP}
D. S. P. Salazar and Pudim AI,
\emph{A Laplace-Positivity Calculus for Special-Function Inequalities and Source Problems},
public manuscript, version v036, manuscript date 23 June 2026, APP-0029 and APP-0034.
\href{https://github.com/pudim-project/pudim-ai-demo-zetalaw/blob/main/theory/versions/v036/THEORY_v036.tex}{Versioned primary source};
\href{https://github.com/pudim-project/pudim-ai-demo-zetalaw/blob/main/APPLICATIONS.md}{application ledger}.

\bibitem{R2E4:R1}
KEIO (AI-assisted research experiment),
\emph{A unique-minimum theorem for a logarithmic gamma quotient},
Round 1, Entry 04, unpublished research note, 20 October 2026.
The present paper reproduces all ingredients needed for its right-half-line results.

\bibitem{R2E5:BS}
R. B. Bapat and V. S. Sunder,
\emph{An extremal property of the permanent and the determinant},
Linear Algebra Appl. \textbf{76} (1986), 153--163.
\href{https://doi.org/10.1016/0024-3795(86)90220-X}{doi:10.1016/0024-3795(86)90220-X}.

\bibitem{R2E5:drury}
S. W. Drury, \emph{A counterexample to a question of Bapat and Sunder},
Electron. J. Linear Algebra \textbf{31} (2016), 69--70.
\href{https://doi.org/10.13001/1081-3810.3178}{doi:10.13001/1081-3810.3178}.

\bibitem{R2E5:zhang}
F. Zhang, \emph{An update on a few permanent conjectures},
Spec. Matrices \textbf{4} (2016), 305--316.
\href{https://arxiv.org/abs/1608.02844}{arXiv:1608.02844}.

\bibitem{R2E5:R1}
KEIO (AI-assisted research experiment), Entry 05, Round 1,
\emph{Determinant--Pairing Certificates and a Weak-Correlation Region of
Chollet's Permanent Conjecture}, 20 October 2026.
Internal tournament research note; corrected frozen version.
Lemma 3 and Corollary 2 supply the standard Gram and transfer ingredients
reused here; the Round 1 main theorem is not a dependency.

\bibitem{R2E6:R1}
KEIO (AI-assisted research experiment),
\emph{Sharp Ewald-Point Counts in Sequential Simplex-Bundle Towers},
Round 1, Entry 06, unpublished research-experiment note, 20 October 2026.
Theorem 1.2 proves restricted-family optimality; Lemma 2.1 gives the elementary attachment proof. These earlier results are not claimed anew here.

\bibitem{R3E1:AMSV}
T. Amdeberhan, V. H. Moll, V. Sharma and C. Vignat,
\emph{The valuations of power sums}, Scientia, Series A: Mathematical Sciences \textbf{35} (2025), 115--129. DOI: \href{https://doi.org/10.71712/fwa0-kj92}{10.71712/fwa0-kj92}. Current journal PDF, Conjecture 6.1, p.~128; received 10 October 2024, revised 23 November 2024.
\url{https://revistas.usm.cl/scientia/article/download/62/VOL35_2025_9/137}.

\bibitem{R3E1:BM}
I. N. Baoulina and P. Moree, \emph{Forbidden integer ratios of consecutive power sums}, in \emph{From Arithmetic to Zeta-Functions}, Springer (2016), 1--30. DOI: \href{https://doi.org/10.1007/978-3-319-28203-9_1}{10.1007/978-3-319-28203-9\_1}. Preprint arXiv:1510.06064v1, 20 October 2015 (only version in inspected record). \url{https://arxiv.org/abs/1510.06064}.

\bibitem{R3E1:R2}
KEIO (AI-assisted research experiment), \emph{Three $p$-adic branches of power-sum block defects and an endpoint obstruction to naive lifting}, Round 2, Entry 01, unpublished tournament note, 20 October 2026. Context only, not a proof dependency.

\bibitem{R3E2:G}
AustinBGibbons, \emph{Strong Seymour vertices and the Nine-Mass Conjecture}, public repository, notes dated 22--23 July 2026; inspected 20 October 2026. In particular, \emph{Exactly nine strong Seymour vertices at every odd order at least 15}, and the ladder/substitution discussion. Unreviewed public mathematical notes, not assumed as a theorem dependency. \url{https://github.com/AustinBGibbons/ssnc}.

\bibitem{R3E2:CG}
F. R. K. Chung and R. L. Graham, \emph{Quasi-random tournaments}, Journal of Graph Theory \textbf{15} (1991), 173--198. \href{https://doi.org/10.1002/jgt.3190150206}{doi:10.1002/jgt.3190150206}.

\bibitem{R3E2:KS}
S. Kalyanasundaram and A. Shapira, \emph{A note on even cycles and quasi-random tournaments}, Journal of Graph Theory \textbf{73} (2013), 260--266. \url{https://arxiv.org/abs/1108.0011}.

\bibitem{R3E2:R2}
KEIO (AI-assisted research experiment), \emph{Locally Resilient Strong Second Neighborhoods in Doubly Regular Tournaments}, Round 2, Entry 02, 20 October 2026. Internal tournament manuscript, especially Section 2. The DRT link and skew identity are reused, not counted as new.

\bibitem{R3E3:SR}
S. Sa-nguansin and K. Rodtes, \emph{Permanents of correlation matrices and the Chollet permanental conjecture}, Linear and Multilinear Algebra \textbf{73} (2025), 4084--4096. \href{https://doi.org/10.1080/03081087.2025.2588584}{doi:10.1080/03081087.2025.2588584}. Bibliographic identification only; full text not available during this review.

\bibitem{R3E3:R105}
KEIO (AI-assisted research experiment), Round~1 entry~05, \emph{Determinant--Pairing Certificates and a Weak-Correlation Region of Chollet's Permanent Conjecture}, internal tournament manuscript, 20 October 2026. Determinant and matching certificate.

\bibitem{R3E3:R203}
KEIO (AI-assisted research experiment), Round~2 entry~03, \emph{A strong permanent inequality for complex positive semidefinite cactus matrices}, internal tournament manuscript, 20 October 2026. Theorem~1.1 and Lemmas~3.2, 4.1.

\bibitem{R3E3:R205}
KEIO (AI-assisted research experiment), Round~2 entry~05, \emph{Strong Self-Chollet Inequalities for Complex Cycle Blocks}, internal tournament manuscript, 20 October 2026. Complex PSD cactus and phase-balanced block theorem. The common cactus result is counted once.

\bibitem{R3E4:Pioge}
L. Pioge, K. K. Pietrasz, B. Seron, L. Novo and N. J. Cerf, \emph{A logical implication between two conjectures on matrix permanents}, Linear Algebra and its Applications \textbf{725} (2025), 309--318; arXiv:2508.00111v1. \url{https://arxiv.org/abs/2508.00111}.

\bibitem{R3E4:R203}
KEIO, Round~2 entry~03, \emph{A strong permanent inequality for complex positive semidefinite cactus matrices}, frozen tournament manuscript, 20 October 2026; Theorem~1.1, Section~3, Lemma~4.1, and Section~5. Internal source: \texttt{round2/03\_high/paper.tex}.

\bibitem{R3E4:R205}
KEIO, Round~2 entry~05, \emph{Strong Self-Chollet Inequalities for Complex Cycle Blocks and Positive Semidefinite Matrices with Cut Vertices}, frozen tournament manuscript, 20 October 2026; Section~4 and the main theorem. Internal source: \texttt{round2/05\_max/paper.tex}.

\bibitem{R3E5:Tran}
H. A. Tran, \emph{A simple counterexample for the permanent-on-top conjecture},
Math. Inequal. Appl. \textbf{25} (2022), 1--16.
\href{https://doi.org/10.7153/mia-2022-25-01}{doi:10.7153/mia-2022-25-01}.
\url{https://files.ele-math.com/articles/mia-25-01.pdf}.
Preprint \url{https://arxiv.org/abs/2101.03428}, latest version 5,
23 August 2022, checked 20 October 2026.

\bibitem{R3E5:LZL}
Q. Li, Q. Zhao, and W. Luo,
\emph{Chollet's Permanent Conjecture Through Order Six via Border Induction},
Preprints.org, version 1, posted 31 August 2026.
\url{https://www.preprints.org/manuscript/202608.2196/v1}.
Latest posted version checked 20 October 2026; not peer reviewed.
The order-six claim is recorded, not assumed.

\bibitem{R3E5:R203}
KEIO (AI-assisted research experiment), Round 2, Entry 03,
\emph{A strong permanent inequality for complex positive semidefinite
cactus matrices}, 20 October 2026. Frozen tournament note.
Theorem 1.1 and Lemma 4.1 (Gram splitting).

\bibitem{R3E5:R205}
KEIO (AI-assisted research experiment), Round 2, Entry 05,
\emph{Strong Self-Chollet Inequalities for Complex Cycle Blocks
and Positive Semidefinite Matrices with Cut Vertices},
20 October 2026. Frozen tournament note.
Theorem 2, Lemma 4 (classical deletion inequality), and Lemma 9
(Gram splitting). The cactus theorem independently overlaps Entry 03.

\bibitem{R3E6:R1}
KEIO (AI-assisted research experiment),
\emph{Sharp Ewald-Point Counts in Sequential Simplex-Bundle Towers},
Round 1, Entry 06, internal research note, 20 October 2026.

\bibitem{R3E6:R2}
KEIO (AI-assisted research experiment),
\emph{Branching Simplex Bundles and Improved Upper Bounds for Ewald-Point Minima},
Round 2, Entry 06, internal research note, 20 October 2026.

\bibitem{R4E1:Sondow}
J. Sondow, \emph{Lerch Quotients, Lerch Primes, Fermat--Wilson Quotients, and the Wieferich--non--Wilson Primes 2, 3, 14771}, arXiv:1110.3113. \url{https://arxiv.org/abs/1110.3113}. Cited for classical context only; our required congruence is proved in full.

\bibitem{R4E2:G}
AustinBGibbons, \emph{Strong Seymour vertices and the Nine-Mass Conjecture}, public repository, notes dated 22--23 July 2026. Current README and \emph{The Nine-Mass Program} checked 20 October 2026. Unreviewed claims; not theorem dependencies. \url{https://github.com/AustinBGibbons/ssnc}.

\bibitem{R4E2:R2}
KEIO, \emph{Locally Resilient Strong Second Neighborhoods in Doubly Regular Tournaments}, Round 2, Entry 02, frozen tournament manuscript, 20 October 2026, Section 2. Paley/DRT input reused.

\bibitem{R4E2:R3}
KEIO, \emph{Arbitrarily Sparse Strong Seymour Sets in Regular Quasirandom Tournaments}, Round 3, Entry 02, frozen tournament manuscript, 20 October 2026. Theorem 1, Lemma 7, and final section. Spectral obstruction at $o(n)$ scale and the question of the $O(\sqrt n)$ scale.

\bibitem{R4E2:LM}
L. Li and T. Molla, \emph{Cyclic triangle factors in regular tournaments}, Electronic Journal of Combinatorics \textbf{26} (2019), P4.24. \href{https://doi.org/10.37236/7759}{doi:10.37236/7759}. Different packing problem, cited to delimit scope.

\bibitem{R4E3:Tran}
H. A. Tran, \emph{A simple counterexample for the permanent-on-top conjecture}, Math. Inequal. Appl. \textbf{25} (2022), 1--16. \href{https://doi.org/10.7153/mia-2022-25-01}{doi:10.7153/mia-2022-25-01}. \url{https://arxiv.org/abs/2101.03428v5}.

\bibitem{R4E3:SR}
S. Sa-nguansin and K. Rodtes, \emph{Permanents of correlation matrices and the Chollet permanental conjecture}, Linear and Multilinear Algebra \textbf{73} (2025), 4084--4096. \href{https://doi.org/10.1080/03081087.2025.2588584}{doi:10.1080/03081087.2025.2588584}. Full text not accessed.

\bibitem{R4E3:R303}
KEIO (AI-assisted research experiment), Round~3 entry~03, \emph{A positive permanent-defect identity for two-hub bipartite matrices}, frozen tournament manuscript, 20 October 2026. Theorem~2.1.

\bibitem{R4E3:R304}
KEIO (AI-assisted research experiment), Round~3 entry~04, \emph{One structured factor is enough: Bapat--Sunder bounds for long-cycle blocks}, frozen tournament manuscript, 20 October 2026. Theorem~5.1, dense scalar stability.

\bibitem{R4E3:R305}
KEIO (AI-assisted research experiment), Round~3 entry~05, \emph{Permanent-on-top certificates for complex cycle blocks and overlapping bipartite blocks}, frozen tournament manuscript, 20 October 2026. Lemma~2.1; Propositions~3.1, 3.3, 4.1; Lemma~5.1; Corollary~6.1. The originating sign-conjugation certificates and operator one-point method are reused and extended, not reclaimed.

\bibitem{R4E4:Tran}
H. A. Tran, \emph{A simple counterexample for the permanent-on-top conjecture}, Math. Inequal. Appl. \textbf{25} (2022), 1--16. \href{https://doi.org/10.7153/mia-2022-25-01}{doi:10.7153/mia-2022-25-01}. \url{https://arxiv.org/abs/2101.03428}; latest inspected v5, 23 August 2022.

\bibitem{R4E4:R304}
KEIO, Round 3 entry 04, \emph{One structured factor is enough: Bapat--Sunder bounds for long-cycle blocks}, frozen tournament manuscript, 20 October 2026; especially Theorem 5.1. Internal source: \texttt{round3/04\_xhigh/paper.tex}.

\bibitem{R4E4:R305}
KEIO, Round 3 entry 05, \emph{Permanent-on-top certificates for complex cycle blocks and overlapping bipartite blocks}, frozen tournament manuscript, 20 October 2026; Lemma 2.1, Propositions 3.1 and 3.3. Internal source: \texttt{round3/05\_max/paper.tex}. These earlier cycle certificates are reused, not claimed as new.

\bibitem{R4E5:CLR}
P. Caputo, T. M. Liggett, and T. Richthammer,
\emph{Proof of Aldous' spectral gap conjecture},
J. Amer. Math. Soc. \textbf{23} (2010), 831--851, Theorem 1.1.
\href{https://doi.org/10.1090/S0894-0347-10-00659-4}{doi:10.1090/S0894-0347-10-00659-4}.
\url{https://arxiv.org/abs/0906.1238v4}.
Latest version, 28 September 2009, checked 20 October 2026.

\bibitem{R4E5:R303}
KEIO (AI-assisted research experiment), Round 3, Entry 03,
\emph{A positive permanent-defect identity for two-hub bipartite matrices},
20 October 2026. Frozen tournament note, Section 2.
Its scalar $K_{2,m}$ result independently overlaps part of Entry 05;
its exact identity and indefinite extension have separate credit.

\bibitem{R4E5:R304}
KEIO (AI-assisted research experiment), Round 3, Entry 04,
\emph{One structured factor is enough: Bapat--Sunder bounds for long-cycle blocks},
20 October 2026. Frozen tournament note, Theorem 5.1.

\bibitem{R4E5:R305}
KEIO (AI-assisted research experiment), Round 3, Entry 05,
\emph{Permanent-on-top certificates for complex cycle blocks and overlapping bipartite blocks},
20 October 2026. Frozen tournament note, Theorem 1.1,
Propositions 3.1, 3.3, 4.1, and Lemmas 5.1--5.2.
Its Gram splitting in turn credits Round 2 Entry 03, Lemma 4.1,
and Round 2 Entry 05, Lemma 9.

\bibitem{R4E6:LMY}
J. Leake and M. Mohammadi Yekta,
\emph{Log-concavity and Approximate Counting for Totally Unimodular Polytopes},
arXiv:2609.39917v1, 30 September 2026.
\url{https://arxiv.org/html/2609.39917v1}.

\bibitem{R4E6:AD}
R. Ahlswede and D. E. Daykin,
\emph{An inequality for the weights of two families of sets, their unions and intersections},
Z. Wahrscheinlichkeitstheorie verw. Gebiete \textbf{43} (1978), 183--185.
\href{https://doi.org/10.1007/BF00536201}{doi:10.1007/BF00536201}.
Original scan: \url{https://noah.nrw/ubbihs/download/pdf/5106477}.

\bibitem{R4E6:Speyer}
D. E. Speyer,
\emph{$L$-log-concavity and a proof of the conjecture of Lam, Postnikov and Pylyavskyy},
arXiv:2601.05007v2, 30 April 2026.
\url{https://arxiv.org/abs/2601.05007v2}.

\bibitem{R4E6:R2}
KEIO, entrant 06, Round 2 research note (2026),
\emph{Branching simplex bundles and improved upper bounds for Ewald-point minima}.
Internal tournament manuscript; cited only to distinguish the earlier research direction.

\bibitem{R4E6:R3}
KEIO, entrant 06, Round 3 research note (2026),
\emph{An exponential barrier for arbitrary-facet simplex-bundle towers}.
Internal tournament manuscript; no result from it is reused in this proof.

\bibitem{R5E1:Evertse}
J.-H. Evertse, \emph{The $p$-adic Subspace Theorem}, author-hosted lecture notes, May 2011, Theorem 1.1 and following ineffectivity remark. The theorem is attributed there to H. P. Schlickewei (1977). \url{https://pub.math.leidenuniv.nl/~evertsejh/dio2011-padicsubspace.pdf}. Checked 20 October 2026.

\bibitem{R5E1:Schlickewei}
H. P. Schlickewei, \emph{The $p$-adic Thue--Siegel--Roth--Schmidt theorem}, Archiv der Mathematik \textbf{29} (1977), 267--270. Original theorem attribution; the precise formulation used is the accessible source \cite{R5E1:Evertse}.

\bibitem{R5E1:BD}
G. Bayarmagnai and S. Delger, \emph{On the $p$-adic valuations of sums of powers of integers}, Journal of Integer Sequences \textbf{25} (2022), Article 22.8.1, Theorem 1. Published 7 October 2022; checked 20 October 2026. \url{https://cs.uwaterloo.ca/journals/JIS/VOL25/Bayarmagnai/bayar10.pdf}.

\bibitem{R5E1:R4}
KEIO (AI-assisted research experiment), \emph{A Wilson-quotient dichotomy for exponent valuation trees}, Round 4, Entry 01, frozen tournament manuscript, 20 October 2026. Theorems 1 and 5 and Corollary 4. Its elementary classification is reused, not counted as a new contribution.

\bibitem{R5E2:G}
AustinBGibbons, \emph{Strong Seymour vertices and the Nine-Mass Conjecture}, public repository and \emph{The Nine-Mass Program}, notes dated 22--23 July 2026, current version rechecked 20 October 2026. Unreviewed public claims, cited for problem status and the module/substitution viewpoint, not assumed theorems. \url{https://github.com/AustinBGibbons/ssnc}.

\bibitem{R5E2:R3}
KEIO, \emph{Arbitrarily Sparse Strong Seymour Sets in Regular Quasirandom Tournaments}, Round 3, Entry 02, frozen tournament manuscript, 20 October 2026. Theorem 1 and its spectral-count limitations.

\bibitem{R5E2:R4}
KEIO, \emph{The Sharp Spectral Scale for Strong Seymour Matching Deficiency}, Round 4, Entry 02, frozen tournament manuscript, 20 October 2026. Theorem 1 and Section 2 (Hall compression); Theorem 2 (single prescribed-deficiency construction); final section (open count direction).

\bibitem{R5E3:Gardner}
R. J. Gardner, \emph{The Brunn--Minkowski inequality}, Bull. Amer. Math. Soc. \textbf{39} (2002), 355--405. \href{https://doi.org/10.1090/S0273-0979-02-00941-2}{doi:10.1090/S0273-0979-02-00941-2}.

\bibitem{R5E3:CK}
C. Chu and Y. Karshon, \emph{A Pick-type theorem for halfway-lattice polygons}, arXiv:2610.00696v1, 30 September 2026. Theorem 5.8 directly implies Lemma~\ref{R5E3:pick}; an independent proof is supplied here. Latest record and full HTML checked 20 October 2026. \url{https://arxiv.org/html/2610.00696v1}.

\bibitem{R5E3:IYZ}
D. Iglesias, J. Yepes Nicol\'as and A. Zvavitch, \emph{Brunn--Minkowski type inequalities for the lattice point enumerator}, arXiv:1911.12874v2, 7 April 2020, Theorem 2.1; latest record checked 20 October 2026. \url{https://arxiv.org/html/1911.12874v2}. Its enlarged-set inequality is cited for context, not used as an unaltered-fiber theorem.

\bibitem{R5E3:R406}
KEIO (AI-assisted research experiment), Round 4 entry 06, \emph{Barycentric log-concavity for disk-network fibers and rank-two unimodular matrices}, frozen tournament manuscript, 20 October 2026. Theorems 1--2 and Section 4. The prior low-rank direction argument and the motivating restricted-class target retain entry 06's credit.

\bibitem{R5E4:Wanless}
I. M. Wanless, \emph{Lieb's permanental dominance conjecture}, arXiv:2202.01867v1, 3 February 2022; current record checked 20 October 2026. \url{https://arxiv.org/html/2202.01867}.

\bibitem{R5E4:R305}
KEIO, Round 3 entry 05, \emph{Permanent-on-top certificates for complex cycle blocks and overlapping bipartite blocks}, frozen tournament manuscript, 20 October 2026, Theorem 1.1 and Propositions 3.1, 3.3, 4.1.

\bibitem{R5E4:R403}
KEIO, Round 4 entry 03, \emph{Reflected PSD companions and permanent bounds for intersecting odd cycles}, frozen tournament manuscript, 20 October 2026, Theorem 1.1 and Corollary 3.1.

\bibitem{R5E4:R404}
KEIO, Round 4 entry 04, \emph{Polynomial-radius dense neighborhoods with permanent on top}, frozen tournament manuscript, 20 October 2026, Theorem 1.1.

\bibitem{R5E4:R405}
KEIO, Round 4 entry 05, \emph{Quantitative permanent-on-top gaps and dense operator neighborhoods}, frozen tournament manuscript, 20 October 2026, Theorems 1.1--1.2 and Section 5. The transposition-comparison viewpoint and classical interchange-gap application are credited; their sparse-domain theorem is not a premise of our direct cycle-energy proof.

\bibitem{R5E5:R305}
KEIO (AI-assisted research experiment), Round 3, Entry 05,
\emph{Permanent-on-top certificates for complex cycle blocks and overlapping
bipartite blocks}, 20 October 2026, frozen tournament note, Theorem 1.1.

\bibitem{R5E5:R403}
KEIO (AI-assisted research experiment), Round 4, Entry 03,
\emph{Reflected PSD companions and permanent bounds for intersecting odd cycles},
20 October 2026, frozen tournament note, Theorem 1.1 and Proposition 2.2.

\bibitem{R5E5:R404}
KEIO (AI-assisted research experiment), Round 4, Entry 04,
\emph{Polynomial-radius dense neighborhoods with permanent on top},
20 October 2026, frozen tournament note, Theorem 1.1.
Its large-edge cycle-centered neighborhoods are not uniformly subsumed here.

\bibitem{R5E5:R405}
KEIO (AI-assisted research experiment), Round 4, Entry 05,
\emph{Quantitative permanent-on-top gaps and dense operator neighborhoods},
20 October 2026, frozen tournament note, Theorem 1.1 and Sections 5--6.
Transposition-gap interpretation and scalar consequences are reused with credit.

\bibitem{R5E6:BGR}
S. Billey, V. Guillemin, and E. Rassart,
\emph{A vector partition function for the multiplicities of $\mathrm{sl}_k\mathbb C$},
arXiv:math/0307227v1, 16 July 2003, Section 1.4.
\url{https://arxiv.org/html/math/0307227v1}.

\bibitem{R5E6:KMS}
A. Khare, J. P. Matherne, and A. St. Dizier,
\emph{Log-concavity of characters of parabolic Verma modules, and of restricted Kostant partition functions},
arXiv:2504.01623v2, 21 October 2025.
\url{https://arxiv.org/abs/2504.01623v2}.

\bibitem{R5E6:AD}
R. Ahlswede and D. E. Daykin,
\emph{An inequality for the weights of two families of sets, their unions and intersections},
Z. Wahrscheinlichkeitstheorie verw. Gebiete \textbf{43} (1978), 183--185.
\href{https://doi.org/10.1007/BF00536201}{doi:10.1007/BF00536201}.

\bibitem{R5E6:R4}
KEIO, Round 4 entry 06,
\emph{Barycentric log-concavity for disk-network fibers and rank-two unimodular matrices},
frozen internal tournament manuscript, 20 October 2026.
Theorem 1.1, Corollary 1.3 and Lemmas 2.1--2.2 are reused with attribution; the $K_4$ all-balance case is not in their scope.

\bibitem{R6E1:R4}
KEIO (AI-assisted research experiment), \emph{A Wilson-quotient dichotomy for exponent valuation trees}, Round 4, Entry 01, frozen tournament manuscript, 20 October 2026. Root classification reused with proof.

\bibitem{R6E1:R5}
KEIO (AI-assisted research experiment), \emph{Finite boundary exceptions and mean valuations of power sums}, Round 5, Entry 01, frozen tournament manuscript, 20 October 2026. Finite-exception theorem, classical arithmetic input, and degree-one largest-term identity reused with attribution.

\bibitem{R6E2:G}
AustinBGibbons, \emph{Strong Seymour vertices and the Nine-Mass Conjecture}; \emph{The Nine-Mass Program}, public notes dated 23 July 2026, current text checked 20 October 2026. Unreviewed public problem and reduction claims, not assumed proof inputs. \url{https://github.com/AustinBGibbons/ssnc/blob/main/notes/05_nine_mass_program.md}.

\bibitem{R6E2:R4}
KEIO, Round 4 Entry 02, \emph{The Sharp Spectral Scale for Strong Seymour Matching Deficiency}, frozen tournament manuscript, 20 October 2026. Theorem 1 and Section 2: Hall compression and deficiency threshold, reused with proof and attribution.

\bibitem{R6E2:R5}
KEIO, Round 5 Entry 02, \emph{A Sharp Three-Quarter Law and Spectral Tail Bounds for Strong Seymour Matchings}, frozen tournament manuscript, 20 October 2026. Lemma 6: near-row packing; Corollary 7: all-strong threshold without homogeneous triangles; Theorems 1--2: earlier density regime and nonprime attaining construction.

\bibitem{R6E3:BGR}
S. Billey, V. Guillemin and E. Rassart, \emph{A vector partition function for the multiplicities of $\mathfrak{sl}_k\mathbb C$}, J. Algebra \textbf{278} (2004), 251--293; arXiv:math/0307227v1, 16 July 2003, Section 1.4. Latest arXiv record checked 20 October 2026. \url{https://arxiv.org/abs/math/0307227}.

\bibitem{R6E3:McS}
C. McSwiggen, \emph{Box splines, tensor product multiplicities and the volume function}, Algebraic Combinatorics \textbf{4} (2021), 435--464. Published version checked 20 October 2026. \href{https://doi.org/10.5802/alco.164}{doi:10.5802/alco.164}. Classical shift/deconvolution context, not a proof dependency.

\bibitem{R6E3:R503}
KEIO (AI-assisted research experiment), Round 5 entry 03, \emph{A sharp quarter correction for two-dimensional unimodular fibers}, frozen tournament manuscript, 20 October 2026. Theorem 1.1 and Corollary 1.2. The half-hexagon identity is explicitly credited there to Chu--Karshon.

\bibitem{R6E3:R506}
KEIO (AI-assisted research experiment), Round 5 entry 06, \emph{Sharp power concavity for the $A_3$ partition extension and simple rank-three unimodular fibers}, frozen tournament manuscript, 20 October 2026. Theorem 1.1, Corollary 1.2 and Lemma 2.1. The uncorrected power and logarithmic results retain that entry's credit.

\bibitem{R6E4:KP}
R. Koteck\'y and D. Preiss, \emph{Cluster expansion for abstract polymer models}, Commun. Math. Phys. \textbf{103} (1986), 491--498. \href{https://doi.org/10.1007/BF01211762}{doi:10.1007/BF01211762}.

\bibitem{R6E4:FP}
R. Fern\'andez and A. Procacci, \emph{Cluster expansion for abstract polymer models. New bounds from an old approach}, Commun. Math. Phys. \textbf{274} (2007), 123--140, Sections 2 and 4. \href{https://doi.org/10.1007/s00220-007-0279-2}{doi:10.1007/s00220-007-0279-2}. \url{https://arxiv.org/abs/math-ph/0605041v2}; latest version 16 October 2006. The Mayer, Penrose and rooted-tree ingredients are classical.

\bibitem{R6E4:Barvinok}
A. Barvinok, \emph{Computing permanents of complex diagonally dominant matrices and tensors}, arXiv:1801.04191v2, 11 September 2018, Theorems 1.2--1.3 and the Brenner attribution following Theorem 1.2. \url{https://arxiv.org/html/1801.04191}.

\bibitem{R6E4:KL}
F. Koehler and P. K. Leung, \emph{Approximating the Permanent of a Random Matrix with Polynomially Small Mean: Zeros and Universality}, arXiv:2604.01367v1, 1 April 2026, Sections 2.3 and 3.2. \url{https://arxiv.org/html/2604.01367v1}. Related context, not a theorem dependency.

\bibitem{R6E4:R504}
KEIO (AI-assisted research experiment), Round 5 entry 04, \emph{A universal weak-correlation region with permanent on top}, frozen tournament manuscript, 20 October 2026. Lemmas 2.1 and 3.3, Theorems 1.1 and 4.1. The word and least-edge core is shared independently with entry 05.

\bibitem{R6E4:R505}
KEIO (AI-assisted research experiment), Round 5 entry 05, \emph{Arbitrary-support permanent-on-top balls via cycle-to-edge comparison}, frozen tournament manuscript, 20 October 2026. Lemmas 2.1 and 3.2, Theorems 1.1 and 3.3. Independent credit for the same local core; its global exponential criterion is not the new result here.

\bibitem{R6E5:Pioge}
L. Pioge, K. K. Pietrasz, B. Seron, L. Novo and N. J. Cerf,
 \emph{A logical implication between two conjectures on matrix permanents},
 Linear Algebra Appl. \textbf{725} (2025), 309--318.
 \href{https://doi.org/10.1016/j.laa.2025.07.011}{doi:10.1016/j.laa.2025.07.011}.
 \url{https://arxiv.org/html/2508.00111v1}, Section 3.

\bibitem{R6E5:Barvinok}
A. Barvinok, \emph{Computing permanents of complex diagonally dominant matrices and tensors},
 arXiv:1801.04191v2, 11 September 2018.
 \url{https://arxiv.org/abs/1801.04191v2}. Related approximation result,
 not a mask-monotonicity theorem used as a dependency.

\bibitem{R6E5:R504}
KEIO, Round 5 Entry 04,
 \emph{A universal weak-correlation region with permanent on top},
 frozen tournament manuscript, 20 October 2026.
 Lemmas 2.1, 3.3 and Theorem 4.1. Independent joint credit with Entry 05
 for the shared word/least-edge/decorated-cycle mechanism.

\bibitem{R6E5:R505}
KEIO, Round 5 Entry 05,
 \emph{Arbitrary-support permanent-on-top balls via cycle-to-edge comparison},
 frozen tournament manuscript, 20 October 2026.
 Lemmas 2.1, 3.2 and Theorem 3.3. Its global mass loss is removed here
 for scalar masks, not for full Schur-power POT.

\bibitem{R6E6:R406}
KEIO, Round 4 entry 06, \emph{Barycentric log-concavity for disk-network fibers and rank-two unimodular matrices}, frozen tournament manuscript, 20 October 2026. Theorem 1.1, Corollary 1.3 and the rank-two theorem; their core cases are reproduced in Appendix~\ref{R6E6:app:cores}.

\bibitem{R6E6:R506}
KEIO, Round 5 entry 06, \emph{Sharp Power Concavity for the $A_3$ Partition Extension and Simple Rank-Three Unimodular Fibers}, frozen tournament manuscript, 20 October 2026. Theorems 1.1, 1.3 and Lemma 5.1. Analytic and finite-classification proofs are reproduced with credit.

\bibitem{R6E6:R503}
KEIO, Round 5 entry 03, \emph{A sharp quarter correction for two-dimensional unimodular fibers}, frozen tournament manuscript, 20 October 2026. Theorem 1.1 and Corollary 1.2. Its low-nullity power theorem motivates the dimension-three obstruction; its half-hexagon identity is attributed to Chu--Karshon, not newly claimed.

\bibitem{R7E1:KZ}
N. Kalinin and F. S. Zottor, \emph{A $p$-adic ($p\equiv3\pmod4$) depth-5 supercongruence for Gaussian $p$-th power sums over a square}, arXiv:2602.00206v3, 28 July 2026. Latest version checked 20 October 2026. Sections 2--5, including Table 1. \url{https://arxiv.org/html/2602.00206v3}.

\bibitem{R7E1:AGO}
P. Fortuny Ayuso, J. M. Grau and A. M. Oller-Marc\'en, \emph{A von Staudt-type result for $\sum_{z\in\mathbb Z_N[i]}z^k$}, Monatshefte f\"ur Mathematik \textbf{178} (2015), 345--359. \href{https://doi.org/10.1007/s00605-015-0736-5}{doi:10.1007/s00605-015-0736-5}. Author's copy: \url{https://zaguan.unizar.es/record/36750/files/texto_completo.pdf}. Related finite-ring context, not an assumed higher-depth theorem.

\bibitem{R7E1:GO}
J. M. Grau and A. M. Oller-Marc\'en, \emph{Power sums over finite commutative unital rings}, Finite Fields and Their Applications \textbf{48} (2017), 10--19. \href{https://doi.org/10.1016/j.ffa.2017.07.003}{doi:10.1016/j.ffa.2017.07.003}. The finite-field special calculation needed here is reproduced in full.

\bibitem{R7E1:R6}
KEIO (AI-assisted research experiment), \emph{Distinct-shift correlations and maximal-spike formulas for power-sum valuations}, Round 6, Entry 01, frozen tournament manuscript, 20 October 2026. Lemma 3 and Theorem 5: one-root polynomial layer sums and separated-root reduction, credited as methodological antecedents. Its exponential roots and Subspace-Theorem estimates are not premises here.

\bibitem{R7E2:BLP}
Y. Bai, B. Li and B. Park, \emph{Towards a strengthening of the second neighborhood conjecture}, arXiv:2607.18047v2, 24 July 2026, especially Remark 3.1. Version checked 20 October 2026. \url{https://arxiv.org/abs/2607.18047}.

\bibitem{R7E2:G}
AustinBGibbons, \emph{Strong Seymour vertices and the Nine-Mass Conjecture}, public mathematical repository; ladder and substitution notes and \emph{The Nine-Mass Program}, 22--23 July 2026. Unreviewed; headline lower bound labeled open. Checked 20 October 2026. \url{https://github.com/AustinBGibbons/ssnc}.

\bibitem{R7E2:R3}
KEIO (AI-assisted research experiment), \emph{Arbitrarily Sparse Strong Seymour Sets in Regular Quasirandom Tournaments}, Round 3, Entry 02, frozen manuscript, 20 October 2026. Constant-eleven ladder, exact link separation, Paley substitution spectrum and fixed-pattern argument reused with attribution.

\bibitem{R7E2:R6}
KEIO (AI-assisted research experiment), \emph{A Sharp Four-Fifths Law for Prime Regular Tournaments}, Round 6, Entry 02, frozen manuscript, 20 October 2026. Earlier prime construction, module norm inequality and Paley row-separation strategy; not the all-fiber coded surgery proved here.

\bibitem{R7E3:MS}
P. R. W. McNamara and B. E. Sagan, \emph{Infinite log-concavity: developments and conjectures}, Advances in Applied Mathematics \textbf{44} (2010), 1--15. \href{https://doi.org/10.1016/j.aam.2009.03.001}{doi:10.1016/j.aam.2009.03.001}. Preprint \href{https://arxiv.org/abs/0808.1065}{arXiv:0808.1065v2}.

\bibitem{R7E3:BC}
P. Br\"and\'en and M. Chasse, \emph{Infinite log-concavity for polynomial P\'olya frequency sequences}, Proceedings of the AMS \textbf{143} (2015), 5147--5158. \href{https://doi.org/10.1090/proc/12654}{doi:10.1090/proc/12654}. \href{https://arxiv.org/abs/1405.6378}{arXiv:1405.6378v1}.

\bibitem{R7E3:R1}
KEIO, Entry 03, Round 1, \emph{Finite Gaussian--Newton certificates for higher-order $q$-log-concavity of Gaussian columns}, frozen internal competition manuscript, 2026. Cited only for cumulative attribution; no proof here depends on it.

\bibitem{R7E4:CF}
P. Cartier and D. Foata, \emph{Probl\`emes combinatoires de commutation et r\'earrangements}, Lecture Notes in Mathematics \textbf{85}, Springer, 1969. Electronic reissue with introductory material: \url{https://mat.univie.ac.at/~slc/books/cartfoa.pdf}. Classical partially commutative inversion, not a new result here.

\bibitem{R7E4:Kratt}
C. Krattenthaler, \emph{The theory of heaps and the Cartier--Foata monoid}, author-hosted account, Theorems 4.1 and 4.4. \url{https://www.mat.univie.ac.at/~kratt/artikel/heaps.pdf}. Heap inversion is noncommutative; the displayed logarithm theorem is explicitly abelianized.

\bibitem{R7E4:R504}
KEIO, Round 5 entry 04, \emph{A universal weak-correlation region with permanent on top}, frozen tournament manuscript, 20 October 2026. Lemmas 2.1 and 3.3, Theorem 4.1. Shared independent local core with entry 05.

\bibitem{R7E4:R505}
KEIO, Round 5 entry 05, \emph{Arbitrary-support permanent-on-top balls via cycle-to-edge comparison}, frozen tournament manuscript, 20 October 2026. Lemmas 2.1 and 3.2, Theorem 3.3. Independent credit for that same core.

\bibitem{R7E4:R604}
KEIO, Round 6 entry 04, \emph{Dimension-independent Hadamard contraction for weakly interacting permanents}, frozen tournament manuscript, 20 October 2026. Theorems 1.1 and 4.1, Lemmas 2.1 and 3.1. Shared independent scalar connected-cluster theorem with entry 05; not full POT.

\bibitem{R7E4:R605}
KEIO, Round 6 entry 05, \emph{Row-local Bapat--Sunder bounds without a global mass restriction}, frozen tournament manuscript, 20 October 2026. Theorem 1.1, Proposition 4.1 and Section 6. Independent credit for the scalar core and explicit negative-logarithm-coefficient obstruction.

\bibitem{R7E5:Bradley}
B. Bradley,
 \emph{When is the log-permanent concave?},
 MathOverflow question 358038, 16 April 2020, edited 5 October 2020.
 Current question and discussion checked 20 October 2026.
 \url{https://mathoverflow.net/questions/358038/when-is-the-log-permanent-concave}.
 A public motivating question, not a certified current novelty claim.

\bibitem{R7E5:SH}
S. Sra and R. Hosseini,
 \emph{Conic geometric optimization on the manifold of positive definite matrices},
 SIAM J. Optim. \textbf{25} (2015), 713--739.
 \href{https://doi.org/10.1137/140978168}{doi:10.1137/140978168}.
 \url{https://arxiv.org/html/1312.1039v3},
 12 December 2014, Theorem 8 and Lemma 24.

\bibitem{R7E5:R604}
KEIO, Round 6 Entry 04,
 \emph{Dimension-independent Hadamard contraction for weakly interacting permanents},
 frozen tournament manuscript, 20 October 2026.
 Theorem 1.1, Lemma 2.1, and Theorem 4.1.
 Independent joint credit with Entry 05 for the shared scalar cluster core.

\bibitem{R7E5:R605}
KEIO, Round 6 Entry 05,
 \emph{Row-local Bapat--Sunder bounds without a global mass restriction},
 frozen tournament manuscript, 20 October 2026.
 Theorem 1.1, Lemma 3.1, and Proposition 4.1.
 Independent joint credit with Entry 04; the negative noncommutative
 logarithm coefficient is in Section 6.

\bibitem{R7E6:BV}
N.~Berline and M.~Vergne, \emph{Local Euler--Maclaurin formula for polytopes}, Moscow Math. J. \textbf{7} (2007), 355--386; arXiv:math/0507256v3 (28 July 2006). Theorem numbers in this paper refer to v3: 19(d), 20(a), 21, and 26. \url{https://arxiv.org/pdf/math/0507256}.

\bibitem{R7E6:SV}
A.~Szenes and M.~Vergne, \emph{Residue formulae for vector partitions and Euler--MacLaurin sums}, Adv. Appl. Math. \textbf{30} (2003), 295--342. \url{https://arxiv.org/abs/math/0202253}.

\bibitem{R7E6:DLS}
J.~A.~De Loera and B.~Sturmfels, \emph{Algebraic unimodular counting}, Math. Program. \textbf{96} (2003), 183--203. \url{https://escholarship.org/uc/item/666982kh}.

\bibitem{R7E6:BR}
M.~Beck and S.~Robins, \emph{Computing the Continuous Discretely}, 2nd ed., Springer, 2015. Classical cone reciprocity and Ehrhart coefficients. \url{https://matthbeck.github.io/ccd.html}.

\bibitem{R7E6:Klain}
D.~A.~Klain, \emph{The Minkowski problem for polytopes}, Adv. Math. \textbf{185} (2004), 270--288. \url{https://faculty.uml.edu/dklain/Klain-Minkowski.pdf}.

\bibitem{R7E6:Gardner}
R.~J.~Gardner, \emph{The Brunn--Minkowski inequality}, Bull. Amer. Math. Soc. \textbf{39} (2002), 355--405. \url{https://doi.org/10.1090/S0273-0979-02-00941-2}.

\bibitem{R7E6:R503}
Entry 03, Round 5, \emph{A sharp quarter correction for two-dimensional unimodular fibers} (20 October 2026), Theorem 1.1 and Proposition 5.1. Internal predecessor, not a public priority claim.

\bibitem{R7E6:R603}
Entry 03, Round 6, \emph{Optimal-order corrections for power concavity of integer flows} (20 October 2026), Theorems 1.1 and 5.1 (explicit shifted-volume $A_3$ correction and higher-dimensional correction-order obstruction). Internal predecessor.

\bibitem{R7E6:R606}
Entry 06, Round 6, \emph{Series Invariance and a Dimension-Three Obstruction for Unimodular Fiber Concavity} (20 October 2026), Theorem 1.4 and Proposition 4.1 (correction-order theorem and asymptotic-slope argument). Internal predecessor.
\end{thebibliography}
\end{document}

% BEGIN INTEGRATED VERIFICATION ZIP BASE64
% UEsDBBQAAAAIANKpVV26WQS0+EkAAFzoAAAeAAAAb3JpZ2luYWxfY2hlY2tzL2ludmVudG9yeS5q
% c29utX37d9vIlebv81dgc3BOS2tCEqmXH/Gc47bb6Z64u512T3p3wtgLkiAFCwRoAJTMuPW/731X
% FQhKspLMSSayRAL1uHXrPr773b/9RxR9gf9G0R9WaXvxh6fRH9pqXZfpMivbD0fDD0fnh3W1LmdH
% w0P4V1FdH67SVVYftNnnPwz4i23eFhl+80UZZUWG30zrTbSqq2oewX/SaLqe5NNoVV1nddKsl9G0
% Kj9m03ZdZ2P4v7ScwWeai7ReRXN4eXuR0Eejpq2zctFeZGVeLvRtk02bNfC24ej85In8Lp3AR9Np
% i4P4LYsW+VWGT8yKeQKvatO8zGZuPPDAKG2arG7zqoxW60mRNxfwgbSJXtrA4PnRZBO9q8pZdR3h
% EH9Mpz/mRZGWUV5Go6Ph8GmUz6P43Ydy73J/PM4+rfOrvctHw/335Xi8WlazL5fvRzfxAN9Xbn3u
% cLWPv1nty2dX749vYph+HWVXGa1evsyieDUeL/NZdAmPub7IYFj0nM/7z8djWMgPXz4+H968//zx
% fRkfRK/hy/iFRXYaR9dZVE1w5jRdWJ6qXMCSwqvWxbqBz70/iSPYV/jddZZewv9cbFYVfLbJmoPo
% V/hO9jldrgp44145uBys9p/vDUeDk9HgfD+GzaquG3hy2sI6r0tY77pJi6jOVkU6JQnAxYOXwEjy
% JpqnRZPxU2kXmmjdZL6sTPKyWubwhOzzKi0b3Bdc8lWa17D1z3AOeQ3rvqizWQ7faaK0xlnBh9J6
% dhCNxyCP7WT+5V2btuvm6Y3MOm1hX5v1pMlnOb7HbTuMCoZyBX/GZR1Ek3UbXeRNW9X5FMYByw8/
% tRsaRt7ih2FlYFrVCjazobdEFynIWVnBXDP4ZQa/VVGa4BdhMhfpumlRGou8zeqUBKvOrvLs+iD6
% qYqaqljTaEQqv6tn4/H3X6qbJkl+rGCoEUhLyp+A6S9pxLS68AUc1WVZXZfRoqgmMOYJntMGZzYt
% UpCe2cEf4HjcDL7miI8+wPfy9fK2Ux69yub5FLZhuom+Xc8WWUty+45kLHqXbZbw+OivuNDTrMET
% znL2q701egt/XNcTmlqzdbJPRqNRz8n+Ns0H0ZuctuRtWl96amSGZ/K/1sUGDubojAWTzxHsIrwQ
% PrCo09UFbBmITgR/gZnBiYI1jqp1m5RZvriYVDUsHezbJIuWaTuVjZzlNbwE9UM9hdUt24qWHo7H
% tAUdA4OYBc+4qCqQyN8yli9428omC/9sVtkUBX2aNiB1KDzhAWrWRcsnBUa+yMu0gElNc3y/+xgf
% Z2/+DRxiEK1Z3ohQw3LgZ1FyRipdDaz/NKPDZA/Hw7HMG/og/i98lRZvvZrB2TmIfgAN94rOMGsT
% +Pu8rpYo3XEZJ7KQswqU6AZGv1gXaR058aL1g3PewtGN4lnMa4gbWGc0Pvx1HUeoeGr6o6pLeCnt
% VRsVWdq08LZklozgo816ekEbiMJ1EP1Y1RnMuB7wFNfLQYT/jGDd7FMDXYF1qfvqbxiMUOQ5z+j4
% wEuXFb5z9ui4FsXKswBphSdXZQaPgw9OYPeXGUpUW7WwmDN3MCZ8MODFMG1P9cGmLUHb0Zd0vSb5
% KoWRtiJ2sCiJ9yT4MCiBApXRdd5eiOCpVCZtnaflAvR0DuOfwTcyeMW6xBXn14KUoKJsUXZlFXI4
% r8ue0ZJSpXs4LeosncHxwa/UWWZr+QxkNoJBL7KkqvFUV6s2X/LwcOlIwYLqwVfD4cInenpepD9B
% 6Rdhb+QY0j08ZWVXwhYW8ECUgD79ehB9R4swz0tcNdjR6WUDm79aVXXLuhxfjF9dN3QaWDz48nmA
% Wjz+cAGycotSfM0j+RPo+wb2I0l+yq5bmMkUl22O84IdRy2Jz8lquNUn1ecvyQ2t4XhcZtcF7HMU
% f4rtT0W1sJ9hD6fpFa4xSJC+BHa2WC979Ofx8elpj/78FUUW9ygicYNt4B3Ajcjg7F3ggplS6t5N
% ecmrPS5hkAkMLnGDcoqIFB/bX3j/rWBLc9SSV1mZwA24BCOwyVERjktbq57LH7UY6gYxB1nSplU2
% Z4ltr/OmdyAyVl6YCNZyki+KvfH40+RL+ej45svxDf2q3v/wpURD6dPRTXwwLl/AsIoNjwKO3gIF
% Hi6V9bRVSwQPKwwKfpVPiozPYAJ3QYZGFRgMUzjI+OFxyTsOop5/tpHvGCje1DxSfIU/1zvn1rAu
% KMhOzssZHs0oPhmcDs5wQnz4yP5Zkrl0URUzz7qU1/JmVCXv+UewVWCicqhAluAcj8satIvYbuuS
% TAw8UGTN6OK439N4YF1mcCnAcsFSLRtYuSXcHOOSTCfSQTkda3lqcET4GsvqXC0ztf3MIAOdMC4f
% anR9/ck/+fD5jqP/AhXqp3WWLGHhluDdwIDhUlrScqewJIsUBn6xBA9okS6XafRpXbUoxFvn9vRs
% NNx1bnG96GDh9Qn/SYtNiz6V79R0xuH5OFXD/g18lXZB9gaX7Nt1ManSKRhpX9Ib2s7/STdiUKXl
% GrU22VQ4G5Y/vFvW1RpNoDbj81rNQWLHY1JZf8JJ7n0GR+iQfrO39/n96NHZ/iH87mwfnAcwoOK9
% ZDgYj0GntJt9FNgfQEZp29lKggnD/PNpC1ZFmS34l5MMxgCmC+jJVQVnL4on+Ii94WAET8WR23dU
% y0TpHETgGnwEMALwEzgDXR9aVRDVRVaimIA4vkDNCGIA37FdBNelQoMadNUc7QN8gu4gXo/xnjcT
% WqX46I8pzDz79Di4+OEmagL/GFQOC4QY7jg+U5meQpLJkKL1zgrdxOV6mbHHAu+gC3GGehOnlmUz
% uYfBhJqztqqv4JNwea/JOMProJwW6xlv93xdFH2uVIoTIBEa0GfADgGBblmM1MUv07quruFvzRQk
% bOa51b7A6eVCDgp9sYxelCDMTVKAZgIfBsyOS7A4qvXi4iB6q4d+km0q2T22YnHkaEygXbzTSvj6
% 0376YZl+vuWsv0LRAPkBcypJ3rJ3Gr30FRhfFr/BLJKXVV1nfC1Ev2QLuUtfgioGU/KbBp2gZVri
% RriwQ1cnjIaj8767HEyGCtVs2Xb3C29zOgZkypMM913pbhju8n7mtOy6hHsCBX52cDMuX5MmA8t1
% VYBin3rzAoO1zj9H8ct4AFsOR7J+DkZL+vlDrgEKNG6i/Ob36Ycv+ceb39+P+KjG7zSEkf/xo/vr
% CSoD9Z3GcEf/bVyCSM5h3l/GY9iWvZfj8TSvQV+hmU9G08v9my979Ef48f3ohr4Bl/R4DKdpPP5f
% +I95+7dEHjN6d/Oldh+bt3v8h5NjVlfDR/XhaP+R/COp9x+BgQbu0kW7L//7dzjwYJLPoqM/1n8c
% HuAw/z4uYRMbsNRhS+AUkKRK/OnjGhRYpivu2/N4K5NZrr4qhWPI/6lpCvWH//386ODs+Oh8dHR8
% /OQcfjmr2iYWc1a13D+yutJLAOYMEkBaGR4M+gwG5EV1+NrWaIG85+z4cHh0RG4eqzuWaFQfL0At
% ezIvERmyDkAQ0IeL0tksl4eSOODpBP+qAW+E73p1OGo6BDwOU0KeP8Q+TbHRyEx7XYGHOW2rmtxl
% nOGbPJt803hrOEDB3EzJJhPr0Wl6VttTHhg7VqhXW1BJejZghkHU8ipFfwpOy1WVkyrF54Cw45XA
% 1y/+QkJmdTrLMZh2dHD6ePR4ODw7OX6sW6Q+ijgnpBc8kxiOLgYAUJ8mYGJd2unSsBs6MGQGkxb/
% qdIoD9vsA/W7ZAwDDBB5LpSoU3Wk/qmY0NkHeGGd3qIW35HP+N11WsySt3Q1v0QXlET7HZ8JtO/f
% 5TTF5Ft4MMzwVzQWt72X87OTYa/G6+i5VMXEKbxkni5zkB9am/wfqbNtt/0ZWU7/dkJV9xKetqoG
% oJyLdFPRtr0DeajQtb00I6iBlS/wCMARbMgjgIt4gt7wPGo2S7AgYEAkNnAa2FqhxYBhLOGevoAx
% z2El8OyaQMJXsxScsBlsE8kxHRQwfRtZtgkvG/nztfgmZhGCYFWlrDP4wTAUPuDkJqB1E8VP3n8Z
% Hh6Dw+OCUxTrYQPbXptMNon9g2MJHG5ViUuXGHzCAEtjewuLvkzrS4y2hqNtaZPRFakpCMTRI/jm
% AkVtgI+prvEPaT3JYbNrz3ib57ikNhaxkEhVguFmdiG5YTUMDN1CUigg07QpNHnxs9yiiU1ErhMv
% 0rzCUTQ8OI532EE1hZmXqhtwy8C+qFMx8VXq5lE8AjGF8TZwy8FkymaOMXdRiqhMNWLhTAbTDE4H
% wcfF9PNcXrRsYPkpSHSZcRLDDDmwBECOZxIZxhVYtzpu9u1V/wfSo8FtDRiBWE7XcLljEIn1D5gD
% 0XVVXz7T+Am7EugE1PlkraF0OhOkBMJl5dOCSs43S0gtH0Ro8/dGerzpwLL+UxYd+kBNg3+4RX39
% jDrg5XpJPjwI1I/wVPBZ4ew2XtAaA9m/0LyGT8Ge49gVruBLsJjoknlhL9tOVp2N+sy4d2A7WfQl
% wzVVH5iDDRj1pSBpQx+E04ZOOfyUNlVJv89a3DgeyAzEHwz6FiVOJoCnolrBRHD3c/TFMUQW8bLI
% ZUYh4JafS7NilcluUYEpo6ygQ1YU+QJlw7yaq7xhGcvAAgDDEX6Cs1xiSKKJvoPpbKKjkwEfQ/nn
% qVzMIEZ4Vi9Qj2JmivMhLhjtdr3ZchJArF7nGZhOvU7Dr9f+gxqORvHtCodHXRc9OxyVA5NpmoO9
% Zu6L+c6NOM9opKVgb2bBNQtOGIZt+B2ehMMj+aQnIMmk7mlhD6Lve1JMHB9BMwmO1TyvlxkfHdKw
% +WdyS6umhUcVuURvajSUustEylyWeRS1GKOltMV2usDCu0HCgD0rlzd4RrohlCVxismoS+dzcsA4
% 4ErbSsqjTfBSTTTGALb1UnxSFARQm9clunboOuIoSPmT6+uZS3M6aODgoiGJaT4S9vWEItm42It0
% pfFsJ78u4KVj4ohnmeWUYMCARY3yusrQeqOYELykxuXGm419Ht7ar9U1o7vz47/SzRKv4gRmPQWx
% TUtUzKiyXW58UlTTS4zNw9o2lh6HnS5nHPWoJk6BtxU43hzqmrd9SfLzs9HZvYypEsUt4XOAT4ZR
% YJZK4h4Sl+c8nl2N38JiVqDb8+QqLSRP6bLIuwyv3UnOnf7nG3CbJPEdY3aOnR9dMRcqYc/y1Ycv
% 5WB1s5fuP6cUe7qfpJxrl1wO5dYbyZWfx2pWYEoLf3OCyXbfQhqIT3b1YbX3LTw8Gd3sY8hnPP6S
% DAdH47Fk+VGUMBXyDxBF3kC610iwySfE7Wf7yEkBxW1UFAaYe4qP4HHxMB5EioyAzUXtjJ+MQSCK
% 1UUqQIJkyPiBaBWzJ6N7M0/RbxIb2Hw4nDXOwi3Rvg86iFOeFp75yST6nw+r2LO5QH8s6jUZByYn
% fBXovFO1B+LV+8xBFrLnx4/wreVemQz39x/56yizpFHJ1B4N95+XSSbBswaOGr56WqzFICWnM9iM
% 58NYtmhvBW9gzARCIl5i7rdWNw6FBZV0VavNjBqHlhWXPRnGkhwVdWJHDleVVPT1Bdi5MCfYTVi1
% KaUv05ivs5oCm/AQkqV2TUYxvGxFUTa8snUcoOhpXPAKOzio5eFyzhCxkNCQzBGEoa2en8ZkYKOb
% DHfwGuMeBfg4beofLb3UHqC77pH4f4MqAeYEtk9ekElqKX9Kgv/kJcBJc4F+fsWZ4V8k0+msqW2/
% 7/T4fAeeh08iayS0NHATaQQgiSi8cr9xLj7Iw5vfT97XzjR1Q4luig2nnzk2zMnUlt8nKSE/NVwe
% Pk5+BoltPoE5VaIY8yEJoAX4ZU429zyBDsMhhsNmS7xoKbGd2ojFYFvhuUVDM8KIEzojU0k2O4/T
% y01bZh2spQEIWpGLryDuLo4okWyvZP8wggsmjJi2ZIVuenPegyCpDl83eBH6BfBrjNPCDYhTw+j/
% 40SX5/Ak1kA03/eakwY1d/iY4QXiXIo1BsfpSo+LPd1AHbzIGAKAOYHmhDX+HsfmX40C5MCDh3tA
% Nqjtqs5Mdxa385FuJ6YGi2oRlTf7sZ/H9uw7fs26diZR73UnWJggyAqiVnICFnOGiltQfJbDeHjW
% Glhp03UjVqju0zXcLLvxHuTVkfAp6uPa+bSFmZBsM/JHEkE1ZXVd1Q/QIHfmyF/oQV1Z4NuLhcLG
% j8cWBbMYQJOBLs/UbQadsG7Mod4GDp0eHfWoEA5ff08nCfPm/Q/XUPaLWIFB83kyy9MFe9EMLuD8
% OY/DWQkuUI0R6hdbEeoX+5QP0r8j5g/cuw95lH74kuc3vRFkF6DvXS4OHXt4wTTn6BNqEV0gdr1A
% /ZDvmOtFiDElmQ6hohq8uVI8CHKdg1Us2lZjDikHSFyQxjYKLBxJe5RVqZHea9TBIIIYY+Dw7NPI
% S2PRdQWvw1ATW0UoXpiYQwjWKt00FjtowZszBBYHetGUacQvrkq8ZHMEwCQZaDWn8LuB3T+BG5I0
% q3SK+WUcveS+8aChPmn6IlDbAoKaj3TJFFbVdCGhwkBFJmwvgKHRUJRNgtvsiEg0iqynq7xaN2Fw
% w32frNplBjs3w6V6S/FoDETCvC4GIp1yFkSJo4EE1ngjMLc3KQIxQdabvmA45Rl04VW4kqpMwFhx
% cXHKFvDTSeg9nffwFBO5+WF8yo+FUYQCJTVz0SNKMj5AH90ncy+3osv6dhweXLslQ1Eoa5w6qM7X
% ZvdHRydHox3KSa3v/zxi10NGwWkk/xX68HEZv/6QMg74lqR7Skl3MvIxU/35P4ecaw/y655xul6R
% JhBYKQf0/Iw47DvqnWdROvGix/JbuHMCdwfz9EF+XRPvFoqClZUQlOBHEvcyFTHJY5F1u4YTgZep
% XoE+WpKenf8DZ8FnMA+AEoPe+c7BEvftZ3s9WHNg5i/IyiFjEQ1+Dg5n5hbw0bpImyCxg0krQjox
% SM7NKG02yxWYbvmUVAMl2jyItvuzjkZNIg3tcwKBpczdBE5lizA3bbbiE+u5o4zAZ83BB98PLes6
% NJZ1IpPC1mmyEV9JFWPi+dw+JuF2+ICLwnm6xAeR3IojMAgB6dEeIIKav12cFiz1IgMPm527ef45
% ELQ4fX4Wc36RwcEc3lO0JgEVJllac5jKgoCgNNfNtM5XYExqFJ9hOpppQ7wNbgv6wA9QX3dBEcz3
% KuaJppZ/UKnIBW34Um7pl3RxfouhJfLMxiWBqPWae+dfcz8G19xLuOYUyr0NTjh5cv514ITgGhEY
% opzjZxamI39Fos86Nc8cLglGuR1HRaWWtbA8FDP67atss12mmVhmPnKOV/Se1iSfuH7jEZfX2Rxk
% vajdRMUPZFkxwlmXo9gMFISl8RA2BCn6wSPUl3EgEV4pq6p5csqmC8h9lrG/h3jKIKeGWLrIjF8w
% 69o0iABhuAJErWwpAgFfMEsRDzIsK7+L8kRV0XAJRdVastBeRLMM80meZcdWaNdWNEPygKw63OeP
% mSgwQo83oItbUv/3Mek+R1TGogZd9tnPSspd2LXNphdpmTdLui5Xq2Ijd5qrWOJ7TW8qb7/IJJ7l
% 83lWUwLUwJOKcsW7nwRAIl0OIiARITmZDnXZwE6AxZT4wBxQomKG4vVg9qCajYZuK7t2J4brM0z5
% dTKAaokmCZmhnYQgXRYKKuNQVNdTYNnGYZJoWPJSk5Ait5LTsWsBJK0oHOARbze+plT3g3sz8JHx
% 26BwUNiwSBnq/DaDvfp6ZXw3BOJbiuCiyAjKIWKUA2/lD0u5VP+bUvPfMuAPtYoPmviRTMyeGP7j
% 03vG8M0KCdK8XDjE4sxBATEVPKCAGvCG9vHtIjXcNBQ/HsMhKvdegV0JWh6RqT6u8sFQiPhVbPic
% g2g3EkPx5TT8hHPNwYRFGhwiogcAYV/YDcloFIfu7pAiX3748gqWvhLM542tBT78laLKPgn0fPhk
% +OT00YkElkajJ+fD48cKQsdvDE/wK89HB0dno/MnZ0/OjxTk9UcZrkYH1GbMwK2ZsHqjE+XAGmRN
% cdCYViA+PTqPd27CwJMLhYCA4iTwr48giU+Ph3EPeMDSV2nrMs+C/sD0Q9aaQUYhOIq8JhLEBtFC
% QDp+Cq8SNJopeWdaIrRPwZAFKzzaUNaXwfjwzBRNPILNZ1kZLAMGDFMnCv7JsIvAwUYCoMZnVjiL
% rJJVY1UsAAmXNvYsZYp6PlPFkxEm11IhaYBnACdmlrcaVdSKvaXmmqQ+DPNejLr2cs2c0iwrH+Gu
% Cdc6A91PaWdnGCVmGLERmzd04R+wygm8cosRsAYV7aDnmfPTcN/AgawKBeZ/vRb9t2AxRobFGES/
% wCUJm6R5M0IPJy9qRKNpwnxLuQ4fjx73Yu7vADiAdpMBeFiHs6dezctureKdu+vKP2tWosVamu8G
% dmdDAJNiKzxIooSQsYRD78pOMENLL2g/+BnoJxwd03odnXZwNxwXW2RcAyaJOJCpcpqvKD7NOC8x
% KPqClBiqyqcXhD7EA0QxLoQWNZzncLgC3/7j8/dxPVvwQZPQFo+BncJbcCrTdJVOcnaMHWhSRHlg
% H3bDdaeAMJQ70eg+hhI0FoaB2D5t1UtFVdFm04zOORiGfd4me+8SVmOs7QUW2ITmoqKi3In52tN2
% fDca4TvNbmIIklYsBCJYapJRCK3lzVBaCskFgs+Q1/oHzqtvJ/ZGp0e3Jfb8mEpCUSHM5LvilTon
% 0QnL888OhijQL5azDIyNCyy5+xE8RbBOL1LY2iT5a74AIybaAzV6uv+UTfxgoHBpXTHEN/4ptvJR
% 3Aje7/gnzOVScvqnmGu97Rh4aVuK34mp4dWrwskesLebg0Z1+AjErSCAjyqNp1kh5nq+VCTpMxfv
% QT8qxzJzNCQDC9aq7Cyi4l2XyawyAy6EOXHyGbyMYoPzdqOSCxZtAioIsYw54lxlXweC0gGbe9Z1
% QrEmuiMG5JKiR0qfBb8ku0o5lc6TK/IAJk71/fYIQ5ZQqne2lpCogdj5Ywnv4/Y08HnLJisIGytQ
% J3DTsrQOFSwfDngke0RbsMaUbQMG0VOWH0VP0mtOLlVS+dphn7VGFzNjjKtf84IrVV6RhqDQgA2Z
% wsmzLFtxLE+VGmUBAtzYPK+bNuHwVA+GzHAzX68x7lP8L24d4l3fwbrADDp1/++yVpP/mvX/yzpt
% clhIkMvbEABgHfeGkwJjHfNchJjqyeJLkhf1r8tzsH9pv2sus+sknX0EyxMLn9lQhLXHWOmAThw7
% D34yX9KWaiRxsboalqm3Ik1BoctFzcED/p4WyTdo6n6LNgEn+TDEhkbkJ7c6/vHT6xcvTISG04Do
% OgE7HixO1AheoT3WE6AgCDuBrDYeDUnqoamxzNAzayT8gEEeAwvgO64wpkE+7BYMwopixSzj6w1T
% jrjGAw8lNE9zHJpMGkOqmpLzC869JAEWSmheeiBpFcd+QiinYwY5nbA6rmagki8RqTU81koYzH5e
% 0CmietsVEqbgVtO+0qcJibASuBG6lmgQXiGSBz/Nen44XMW6245i4FezrozvQBZbgOWWs8WDzoEw
% 8kwY6q5RGs2NgpWFwF4TC3LYCBtOIW2rvTGQRIK6E6ziIp2pPL9Ni2zDeMi8XbdaMp2JrtW0iNR7
% l13MIPO0WHULKSyNy2aRgTrQliFUOM7gm8aH3Wj8CERI0w64oKiyPHkuwXpBFKlFdkErrluXVUk8
% AQYjuWnp2RjfFVeDNhTuMpCrfiwFgUp4XZzPR6Wg/togzpNtHbrG0217bRCAI2oG7jzTS4GxKrIN
% EuUnn7JFvVuEzCx85Rms4hszG6MSboQHqOV7ACvMqXXpXAEgsuZjfAW52RfrSUg40Q+lOB+djm4z
% 2lKHuvISECGeQzTCV4XKd4EuHI8TiGf85w9fRoMlwi5JX1isMV5SAtUF0UOYQn8w3UuTObfAw2/6
% 0QM/Rh7CCBgbK7BZKq9PSPfw1ZVS2EtgVArK6F0OPDNydl1xjy5qT3paAWLbe0sJAMJ2TLLoY0WL
% B2LpAxdU7XB8WXABWWlhpK6RZtAF1gnNNxpPtoRFYogszj+I3cYMB3A1NIy09fObdj8ogomjjhiQ
% cqWHfprRpTwdI5TGnXkDOsLPJhizWIAaYlBeD9qCH3xL6knOvgZCyMWkuIjBIWAXGHobfQueKEbL
% WTA4luYxpDwTJaHIL8OuG2EWSrqM5QFa4274w8+lS8RmM0EMU8KoRAfoKVpxPZNgV6LA65Aj+Sxn
% W9bc45PeOr7XX6sXHCCrJ41mZ4mGQlZMQRmYhgszqBDjKvMNlNtf9m1sdEno6Il50wF6BbnE6Ntd
% 6K5JgO4iM4KRsm6p1X8rNoEcdXFYhvLhgtuvR2FJ8pITgLpGzlulMpMfyKnpcBRZPJbwU9MLjARa
% ro9sYIWJIG0fZd5bkZNnVOfHPqsro/PhPi2CMyV8lX9mcXK+bZBftHONargZSDn6mmPAUzzStXiA
% SE5XE5aHRGLgStDMoFIWPBQNEWLm4wooGKg01EcT2y3ExiwYfl5t38t1ayYdph/b3HMZvg57xUmr
% jY/b3bLpfAyWSbRWSaeiQBya1A/K7aiH3an1HqB87gIvvO3OfIsnycgGPBXjkAuUvQePnsApnQtv
% SxGdnR2d3RV/2toKcww7Js1W/n9c3oUWDHSXXMrEtMWVICKm8Jye+5MLu7fudDBjceqYDEDYAMsd
% 2E0NXvPtuNQqW+8Nais3rdarCvdHnz7RwY9LGTddmL7dRUff3wa64Os1rq9BNRumIhKtwubMdiLG
% ltoXgqfhuRmX76YX69q4bbxdCeCiBRW6ycG/vsgRA4T2hsFCYZm5shLDMHoEHFzyIPp5poNXKCIG
% q1lhJA6ZLtvF8X1xrPA+l09kzS3W3bh8UdqsqXbKt0ecyY5L6sFWnVrBtHuKpBLjsu0yKGn4MjTz
% tsgpgot9zFACH0tCZFAYDpG7CnNcWEZBjj+HwgWm4EIieYOnMzSI3r57ZXk1Vl7sKzWK7k2cCSKW
% 07iUu91HTVi65QDrTLowIYTnocaDXdnSLCFwiA3PrLzdyoMhiM/G1rMafM984JFG6dVe2KoJGJf/
% krDc3eADEKfvvGz3t2ldI0aNIPXjscbsNslrzLfej4ZhNDx9cj/UwVbJdxpN6irtpORJXvjusdrW
% nRkyBBq8U8YADrR1GAMs0TwRlAXeeqbGOsQBWFXXpsT7iGmAEiMncIiZYQ6XZAuMNVE0x8DS10RZ
% i7UymreWgLiSwBH5gFAIEB+LZYCCsRxE35Epxhc6ZR881EPAeTku4+O90RFxs+y///IqGd7szuaD
% Uun6tsyDWWrUcFnB4uOJbjKs2SWR3YXPUEQVRogqIbPJMATEvKEWVMdIJ46XTDzH1SmDjvFaRXoY
% rIe0BUC6PTz3qzybenahocMoUY1rs+LM95VE0ri2iB0RNCrh1BaysZKRZ+MtRRdFTGAyq/ELKi6e
% x34Q/TfcJdVM6ilSQesURQ/vhM4MIzjP4Gyc+eGiAHHq7NyqxjyxwJCmlcuELNWCtBfgNVktMhdX
% JgRVy7mIrHOSaDFEGQ6EhxRO52UmlrQHJNbb0mhL+AwlisgWKgZY2vYaiRfi0cHRscRCexApwizo
% yGMZopMkDNGB24TxOXLZUYRCBMXK2NwZpTP0FF2rWUV6IUesIF3NfJPiV5gMgAKCRfpZwgIdfIKy
% t3y1Zv23ABKOQ24IJCOp+Zu30UM8PuqrfTTOBAlu3A+RcMyxN7Ozdpq2gd1AIawOFFGsQbVd/ViO
% RkbZqNgiqe2aAYHj6+EXQDk0oaNHRBboW4WFnEoWyySPPqLGIXU4iPPM0BiRWBppcIASlnkfHjbh
% C/OZY0+Q0Fs3mM2hZYyNEUR1sg3JOFbEIUHJ1QTsBmNJVfUH9dRslecNXSg8tZRjL6M5w2adp7md
% FvVTGQgU0TWloDhzGXtWDGq2HNRXbcCjLgnHMv0I9qtnOG3hJRyP2PdbdtOgn46BqVC36BgcCumf
% 4IM5uUcPg+i3vGhASozjcQaLABfncqPHivPXbnlb8OWbHmDR+fmuWsSSclia5AJvCuUzp8zVr1gn
% b/T+KdL7f1k9Gt6k78s4SpuQ980GE8MfkZ0oawLmcTf0Bd4/JZOk4iN4ko7IkhycCgcjpdCggO2P
% WNA/lEybmzaRxV2j88X0c3VkwW1XXO8V4IViDo4WXncEPE+lgq9DimDgcQ5eIVyjyQWpIjIokWkv
% Qg9SayS1eAcRARLdekFQPYQKK7rPkvEohE6vuDlz4i+3rCDTWSvxAZfpq9IpjQegE9hWClWko2ND
% 0+9LcDem5ZsACHN6MOQDOquyRjOQfEAoExyoYeT4cBLMoyKJyUHHpILj5QgiVTEuKdgyCRmjWLp8
% /SlRPpkiY/2FFy37PM1WrYJTBSbzgHN7D2wC7ggz0b3TXOQ70DeZlBZ3cAo/atbBdTPY9oGOnjzZ
% WU/cw3aP7gEnoEtMNIdkABhUyKf99ARk93vpcaRlXGRD4uaYKsAvRC5QcgmM8wbWvZlviH+Xfvfn
% 9yP88oj4BfbqR8P9hMv+wrJz7+N/PCkT4pdAWFIf8kHWTMy/rSL+Tvm+D1Vwc3rKf0OGO5udC8o3
% a0UvIaoOK+dUPZZdDABXSBogpHcjSAkpiKEHdoA9B+Dl3qJhCF9Q2nuj+tFov7x5NDw5ovVjinDM
% JySI2GJgARpHSrTpRTS8w/LMS/5jHANj0qsipRucCoZTyeO7gQ+Eio4VtVpjBXrT6DKl9dQhMRDz
% MNBXUxDeSx2a0aCY4raWQhqu1RSYSHcrXaJMF1WADy47j1gZsPY01s2M3+z8Ez0Z+TM91/2/AKJ0
% co9c+C8cj5PIE3EBlkaN4dK7nKFAq31OMV7Ym0Yc3soCgNsApaOjk/5IsgeL28ZNNKC2jQg2D7sC
% EIVTVboIrxjbDaNzOLgodItBblzY4ui61jgwbHwnFBz1RILNfq9trST7le+IWg7CniEucSffYx2U
% 863HeGYO16680BCj48KEol9nkC5QPJTiA3ev4zQErqe8GVmuhTgQv/JMCGe9xy4xkImmBrOHTzIp
% VlisuayUFojSQ11HibGdllevqksX1Mf1UKorQ9mYy8WTZ/eJ0Vl5Q+1qZkK3T4kvIkXzJU9beVA0
% K7i3Fb7wG4ZmJ5VECGwxFOlAok65BViLAQJvKFVUYDmavZ0AORb759VORengRKkgkIPASsAAamKO
% FWXCA1qDm0VCuFUFaJoGYa5SrUH4Ar84I0ypkaruLv3A8eVysxMW+gTHKf6Ab5UE/Eypl1MPSsoU
% v2nlEOsSo9zioHCDA2/NacMpo8bW4M5w8XZ9BBnZuPQP5GM6uU/6/q1jiBceYk5Jll0apnCFCc5f
% rbZ028molzuF6r1KbirGu4Kzj198OCImtnRbGdDrTE2RT8Bqx/qzuGtfKDsWoGRAGl0canh4opC8
% w2Nl0Zccd+0VF/ZQ2fqVCJpaZiwLWhWoX5DEkpIuaqzScWXku/gvhnMMBNoFOIeHe8PR4/L9GTIE
% fSfuicgohlh7FwUz51Zsi0oI83z0rMfDJ6PyPTZWQw5ZWltfFVtOkXsc5KCr0F7sbqkAme84b90U
% 9sDP2pl+Y5xc6dPgkojRKs+yKbp9VKHriACKjYtdkhXsNz/wtQTRooV1sn6skzEKPTEqP/J5/FSb
% w3Xye0qcyuwPBQUYULkmCzhaq+g6LS5NuTBNZcq1gqmRENi6coyRuoKpK7ryG5cZ+CRVuj34BiiZ
% pkebbEXhglZ0lLyihg63FahEzrB7iO10V1L+L74m3xovHAAvAc+qxqGnfRHrqpbjo5OjvlSSR3g3
% rRjnSX1nygZBHMz7zjD6pgOiCHLB49ILPmoPmG6OmXH8itXTcWNelJAuzA82LkW4fNKlPoQCXa0a
% IjSEDN3lXM4zLl3FNectQs0k8UCfVqAssxDMoad+XKp1VnbZg6T+0Ms8CrUya0XOERvziYdUCMgG
% +jLGDLRmdhKzHmEZxqVTnx0QDmwt5tNDndmB1JDYjMtQI4kHq3DkXlusD49BXpkTCgXPIPNQCoOq
% kuRNvlhkLZibv+TTC1BdyyVmqSkIq71cxiUZTnXeSHGOxz6AKoLgOhhes23uMCyBl1eiCLqAtIt5
% DbRRgCGmtMMC75xEkWD0KWa9NZzOjFJkoIXmVS8aT9pTpeUlXj1jdxYMFUXckpKXGghM01P6jOo4
% uYkjpqXkJKR2skIk+taVsm1pHDg2ER/rgLl86tZEIrQh8xsEc5XNnjrzzgk/rj25nQpAlaEwA4Q/
% ApQJEkXCRvjAqJCWImhRZOjVkMvK1RKmrF8NGNW1+jDIOWfQhWa2FGl1L13vO8Y4HKne0fwcAv5c
% 4ZCa9RMt5vqXeND3oDIAyZ2SezmN3lSL5KW1IMNxvgJXJvkJDDXkhHiNmVJjKcb62wTptzmxShv5
% 4w6A+ej89PzkLlSWcGgbdsobmCSsFTLjpD5tLhkk/heJRkTDg3NqHYJHESOqqAfAmon+b3bZpoqL
% 3+KnRwgSOt1rNxfX3WRO5824PLnYD8+hQS3o70ZcV/J6gbu1nGQz6beUsltonha1F0BNKhAwtnWJ
% E7wUJApuInUIYFh00WyxSVKcALRwjZEmqnfTrpA2en7qFfLPBixvmqO1qLjdyZPqs1wJdoMgmoDo
% n+QWsfJ3HyFPifRGYlrkFdXrUuP/PS0xfTY8Crc5YqolWO84Y421CypEZonJ5MQhCDAjT+16HUw+
% 7LAgoI0Ej/diIwvu1/drsxwQaBQOdgNRcUtBgK/vlcwFhuDA/khtboNlFc2t/GDSCHxAJyBJ53Mi
% 8vUY3siIWFWFLNJvlh0Xw0/6j4loo/CqjuoVVAIBWe9GO0nofWiho1mrVhcav+n2HlymaF2joZJL
% 0+FnIfwJ3RWTQqLE0bzdru5cTDHEUVT0A8joVIXIHRG+Xrv9WzL6Jw/K6B+fne/O6I++LqN/Irgq
% h+ftaxRrdlc3YYOekIFCpFeL5B16CrF2MA0w6L6gTIZ1baBxmbPc6zOI2rhHrMKjJDhlSoIzV5Gb
% DYKivUF4zZt1qFavTj24MgKRJqS4kW2IgpZ2lnaQrOrYsvsc6G3C+Gnigm9C20llvr1B5Wcus2/w
% hG7KWRNlCaaWvSyu5VabRjQdGOhUcVd4idUXRUioYFTNdLwSafPAmTz8Rc3h19ZjYPip6po+g05w
% zFtXt14DRursyBxoQwhHpBBpGc2/JV1wene2X9r22r1q2UoW22UGi+14EIwfgYgCto7849NeEA9F
% 0G5N6897MAHqmjoS7G5WYZfQEB7BhWKcg2TZ9XfrCSOKf1Xy1qDdAcXEepekVpNMUgjyWavRDNeL
% SzrleqE/EZDgkDj6Y2f3dLEIDDfwIAnSImDF38QOgpViB8j+8Vq00UUtVA5L1Agbfzx4UJiLDqn/
% 6KJEp5r6EhfCAaBhQuykeZXP1j6SxqEb4p+pXwHShP8fmAsLq8W8nP8jNrwr5U6pwDaj1imgucKm
% DSH+gQlNbCOpHsmBEjimbYAEJR/geZFVTHFGToIQEQM8pMAh4Y3MtMpCQhe9zMCG/39pXXFPASOI
% ai64kJ3D6qRr/OahStXLk+5KlRG7OUoRz8OlhWBuvbwEb10cMKQ6R9M0l/6mAWAh9M2IcgxE5ooz
% KbZJrDTDqNkDm9+d3o8TQVAH1N4l+cs6peqjNyljmQyL8Csun8f8tgOPsO0pHT8ePt7hKSGQTTN8
% +h7ecSHtgUutaXsr8Cn/zLVb9g3YbiQ23ZD8NUoTqV8Co5h1mUT+7wRAcHlYCFvQOrs/Sy8WxQJI
% +9DFc0y+MwjJRyjsLR4N98tk72TxfpTgz/HAdVKAczgvKgxDHi7G45p+jl0oi2LZgr6A83HZuxwu
% tl/HW06VP5LTMjlzwAz3NRjENMuL4/LwBMaAP8Yd+IRXjsEpA1cT2kptMypUsdXJSiulKWjz1KHm
% t5ESjCKwrjakL73I/W3btFghyY1P2xzTb2zxqL2Eh0uJ8Zzye2vvLwOpPYRFsm1ZKOnCo9ECISiy
% 39ishfgetSvVmrWJDeBYus7k1B8O+4biuY7E7qp2JdlJlF1jiVHQfNkL+l1UywqNG4wDcVWhZy+D
% GgIddtkI0hiVaRKSh1AedKCscETHyfvkPcRCqpaGD+Su2dVVQo8wr0m3s0So/kIzKwoAGgE15xZa
% g+XuAbrwPh0eaBc+iQ70GpUp/YAfw/MMbY4UbKu+89PdDR28r2sRhmgXWJ4S2YO4Z4gL0VxvJQuZ
% OpXvW8kGlnDZgK+0UVNHYwVcFcBbs8zLdSNpSsTfO2uY3Cvx9FGP8LiCWArdAMuUes1uiBNg4W55
% jnM0rliq0Kbj69JC0KEn4wW/7h/2qkpNUQcRKynT8WloYL4UlQpiBhp/kTUiTiJYaBc68bbelklA
% Yk/pe4gOo/gK1WlTSxuq1giDGMwgI3Ek7rvADUnj48v3o0fHl4+GsdLMXeUOw4i99GoJvvEec1Ol
% Lm86awdjgCDKp46UXiJwtmDFI/TD0vKVUcoc5DOfwYJZb/Pp5TeNcVfy4M0Yl8DVt/W6RJgnqAdQ
% K5c5N6ynlsqECizmtlgXoCAXHDjrqQHWtsBbquKOENEzrySjG9Mj/USM5hIt/HqVcZ8mDC6qgdmS
% kMWYOkbfH8Rw9vjx49t7LBCUYWTtmu7I2m9TkmiTTjIqtI7crrzh4d+GDLz8e9yboxejqKysqAKD
% FlJJah0PlKjE2lIOHE5BMFVbuUTSEx6A4kzreOKIOtJK3+rtXKVHDs3UV35uMP49tb7tAYCaSfta
% LI8NUzGdyEzuU9IrxMuJHMK8vk9nKag9BXrJ5aqXYYJJOc2fqbXrJ+Y6+ThMdi1x8aiqMGQHWGvL
% U4WGWs4y2GPhzjacnPBrR6/rapKVOXUg0iS/RrAFH6AlMx49BreroeC3n4OW0DOf9+maC3sJHisg
% PTzHkuJOeDtoa7Bsjbp6ORPIy2YJylEwCleZS28Rlks9M6lQsYbr6AzYnnBoTnGI/mUgl6ZHWuF1
% 97DiFs4ics0x2mJS7W52KyspuQye3ZE6U3DEoBeF5YJED8VFnN6Ni3D1s3pmt0Y5gctL0BFXuRDt
% wx+4JZFLMG8nwc6Pn9zOsrSNQzBpvo8mORiX3e5XgcpjRTdgUNcna67ZPR/g/uw6IPTNxr6Jv0bm
% bzkyIcALCx3RsfuEaOvh4fFIdFTD/0ZsNwUMurpNCC2pSi2jWxEHFGq08HR56iAofVe0uTbJ+jwu
% dyKt+jFWPd3bWA+Oy2BlHdRrwthdhVDhJXFsl0Q6qaR5rp4fGpHl3oPjZ0cUK26aIONI9IPE1gn2
% K/55XBpC2NNTa1heUjKFNmKTRky+1iGPVnQOPAa2cCBCnDXSL5y1nENp+cSZ7hYhHZcTuQUF49n0
% 7VbBsosk2pJh0PgYImOlJBGeCwQ2yCVA4WluTNAQBaWdy7c//+q3KMaCuWyJUULHnu513GHwAWxq
% 4m4vV3lMVCn+VTYu++AXrJeYuo8VJ+MrAq1GhBu5NG5AVgBjUaQcieePDw9LPhNh8SNLTvgxrsZo
% hKQ1oRs9ZYsbzgupRUkdCuSPwgc0ce5hMy65RRErKWbBIdXdOI6FrsItSZnIbz2J5xPK5M4eru1W
% 3gRR3uPynwrx341n4OjcWzKbQiwDLkb84sNxHL01sftOESOKamC6BAE3kAfgwRte9/utT07OdrGK
% e4ClzSpL+PV/rjhq4aTfWlOhFmoI48g9nPzOS+DSTWi7tlCwea1VZhjmcCgYCZ/4SXqEc6clU91Z
% ipoWKz45JIZNC/a7yI8jHKRP5u6rvK5hPthr5sFtUvwSWQwOU3bAgUIpYza/Ff6xlcvzDRMGD0jM
% I+XQH9iDcGtJng/pmibiEnqPcfx0FIiTmELO1NRCSWNMxuKIom0TDwfD48HpcCtrU0qhEA4q0Rwj
% 5Sj6yvx93IZgMjzI60kQv8iC5ejlXRTUHoMYSHA9/6/ZNC1lfqhFs3wUt6HRC2RACFvtDEThpNq4
% DZDfkNgD+c5xfUdCgkIfvcvlZbyEM8lOE3mkUTseRL8o2Ya8YeDdeXbTIxtNZeEDnhlHCEIjER1P
% rCvGuIZg9H0qQ9wba3b4TuunfIo+J1cp8pwrUX33UDlWk96ubcTHhXmNOpNokuE7NCX6rLefFVs4
% GCZ0cIrGM3n/iULo038PUuL0IUiJ4Xlv4YE+kcUGNWBNwRZpNJLOrhgjNck2lau8YnjfqqoKL2dv
% efYgOpncFSGiLd2KUe4ILh447gC/JM/FdF2sPKFEmqgP6UscxvylbV2rQP0d6AwfMXHS38TBM3co
% f+iQA+ndXo21uZKohd91grsJWaWNKQAydunCQGqnfCZd/9j/OdiiZhDyQLpEevS53ZXeueMbU+4+
% 72rjBOi20vNOz4HDXlgowEOVdxOYreB0Ncq5CyZAZBdWX0Ws7X4e3Azw3Ugf10Lk2Ftr2tgRb+yJ
% YG+spOEqC8ssTW0oS/82s4JG89S3eEaxbs2lBharl1h1mdSmB7rxlYrn7G5Mxis54klzkc/byAsH
% CigDfVMwOZsVRkGsZJ8ah/c0r9juw3o22pVTIM/IYBjefXsLGfts53B1YUMASdyLARlsNSoqNj64
% 3miFKEF/EP0sTUOnIQ8EEoNogt140f0hSa8mrjyaYTdbrjWK92Ak+3urZPX+S3L1YbU327/Zd3gL
% USEgLpZg68JP0KpRqId2xyMfzFoZBPvmKm644zsRhzqEDf2k5KGzDjDCLSiSlip3QeP4XTk+z/d/
% wlmJABRH4FtX6cohC/deD/qB4XkjxE/1JHdzj9s6JFEdQuAGr35II7q0nqHsDKS2m8IHvJzUZ8ZL
% 5TpYiFjrLJWtJmBQ7zBZS+n9Uepo8a8DOLcC3KGBGWkEV0Ju1GjiRZPYAYNTM2optG6r+XwL1SLm
% lQdsIfRGoFJ2AzSobZW/V/6jH47eOPsq9MZrdBlewzZcNITdQLF8S2gsbWNxW+uKx0fnO+jSlcMa
% g3m4mVtmAbcSwEzHkr2nDmrAlbZxrBguH4kH+Zlsg5pbEbFjmrsVpqGeCGXglW6RAvDWPkGDeYQe
% Qmcxis/KZDg8iwO+OsNAnJSHp10MhN9OIW9c/TED3rwmHnZbirvBBnmAXjUTm+j4tSoKn0MH8Sux
% EgGbBD+mr7eIW7PTlSyaTM0bK2IUhDQPnOfTuLOQpFYY+hwL2cTp6iYOvsa/P407tWk7UA8Y+ZFh
% MJqEq83iU9ifx7HQw1HgKyx7EjFU/gcHHS6bjOTDMDLoFEtXiJ2L09MdkOwrMmaZN4PCS67biAIk
% UryyBLTGzKkVslrsZsEQ6Q4eRSWpuAoa33XRRbLdM7kmlPnYbw0hhUZPtd6ESovcrksP1CWlXTRN
% uQW50ESRtJW4ix9ju2L8AartTjDGz3xKjMuttjS1mUqeyU3XHgc8MFDSYzadnhzvCGl1oDdfi/rw
% mIuERaNy3IicdBcqBm2o5ofOKKyjaYwQ4EG5EZUuTKkzSbYNb0BC7cJU8HlHhDKlNPAl21HT9SQE
% ifSDQqSHmcE/fGynVX1EUvXhLRMtAMaGnuKAaW2JHKL12arEZsD5Or5QpDBeStheYAJpMSdPluxg
% MixAlZiBsgU4EJYG4VT0g0rT7RDZwOP+09FTRxbXqzNXp86RdqIdo9TwgZKg2dEuaDZedx1t00VG
% wUfd8m2ECMewSMgYAeILudPX1d54/PsE/vv+yywZ3ezHpNMk9NXp7yEWWLwHn4QPenvUrKfTLJuJ
% rnPMa8PDWdxlQPIWKJPmoY2/Ceq8BnxjspVcQCBVWkEjGkZ0MUEqVwF5eXZvpB3OLsWEoPHBWSDx
% S7SC0S/K4TMVFiRZPERor+4LLXmAUrsbLvJKRSTxS7AMQkCGNmdBlNwHQSUFl3zpwbJIx3at+uj8
% 5I6+N7dU5Hprxb9InUgH4zXcQQ2HGc+p9GB2hAIk7tbZG/EGQVqX8LEvSPBcZneLeYNFoSdzG0BW
% hiPm9iuDenTnM3qtMZhi3QmHS48i7cEWBmV4pCweT+hnBqKoB5z/8eONg5fsDZPfJ/oPpNX4b1pB
% 60KIcQwsIuSexPDsMx6OgEZcE12woY4PH+uL8ccDcMsUzMAsIgO/b1krlX4GKk8bbn4eduT4zkB2
% DR5glIyspq5BRCqMq4Dq/VnUW8ctcSoiMwYFo4WHkkPU6m3Xi1Ob2BK6JfH6t7tqdYs4GxTGZ66Y
% VsLc7tjXOeaNHdS7jdy6lqUkdEqpbhLysqBjJUFQl9j7G/sgcXwCZtJlIdXvugpCXXIsQegLViIK
% REcjtUmMp5A1MyyQuBJmPegBxTPF3CS7+mvoyhHVgMf+ErnOY93MZ3+szajfKUrhlz3+S8L0Z3eD
% U36ByTKKuS87rcg51inKm9xRJls6cDg8f7wTdXsrkAR10oABKdi/qG16kCT9WKu8q5Kks+oObQSP
% AX0kuJXvnt9Po/xmXXzgO6CQmPMP9Nkete6JXhx2G/tQZ5/hE/xw7NkprDzGZZc2XiPp1g1LW/KE
% RRmBbu8uCEPLBEntU0egq23HNqtr3A1eMZzPcPgp1mweAnlycYxkfZkCQqqZScFRtTmtqHZVh8f0
% Be1Noe/UNYnpGsbrjNkGibUhUi8CUdfKQ+oEGgqbMFg7ta4mSqfK2frU0zWwYliKWXvRxXGpqBPu
% eZj5TXV8ClOnxPgq06TJheyRpqPIgiLrciv1QhPIazTkZ6qsGCGFcBpEWeCMsuUybQT5v27Uu+cq
% Abb3zJvWaEwH6wtPCc96q2V8lIc0ndZFiooWfKbad+yIcBxqg4fmBaEDtrxeLTgumUxNQRwDBbDs
% RpAUhbNTDDoyLv8ZgqOze0FBqLIs+qEk6VfXPY2cVckIj5+d+axmZBf04aAkXf15cjYajm6ntkjb
% +4MbCJ/u4BrmpbPvTLE0nk0rRue0uuCyvoZn63OgUNw7VwpSr2YCFcQUe9BYmwgFdHjgEYTYzfGt
% 3ErPjF2N5gi58RV1EGK971WcNCjfdFEKCTWedSW9aELoBEV9ma6vDymBPrDk9Tnf7HASTXSSJKfG
% zlRS4mu7xuHpLmIPze1y8pAR4IKI8Ms4FJ9DrXNZV0ihgKKtQOvrpjFVk8NPYFoG58v7hJZPjb4z
% M7szz7W4sHeiWqxwfabZR/J1qMuhD/YSMI/idbx2YF6/Y/KRSgoDFHw6JAwHVuIV0e29tHIQEp+6
% pzKEK1zJ6p1ZeW61onRL5qolBurD3pKL9yidVP9mmtRnXgnX0809wblcFMSEa4MelsLEqsWaIpkY
% l9SN1lADGBy7Ag0hLZ7HibexcpT0w2zv8tHwcLT/fvZo4v5BIQeJF+OHBvC3/0RjYpleCnsd0391
% ZxHPYiQixJErjzgT5xGXlxwzCjHJnXyp7/JCARjzBFO6aVIxvK3SkGUWwcLpJqhHa6ZwKczcqW+6
% 8YQd9Wh+DzlfmWnQ8+v1+b8FonL2IIjK6PjxDjKPEHyBC9il9eCrnRP9jRPlMzwbnvczYF5X9Kmu
% AitlVwyAdR5j/fSQO39c67QJQrvllN3q2g3EEQt61AThi9qcDuYu61R5oLyYnWi2oeRA1UoCrbZE
% hFglFdp4jkgukc3LsDQWLmxQuNbwnmJzS5LGQWYU7MM0QpzW6+TawqxSDzaFr8+8YyyU3ftkXS5R
% cSSiUAR2Ikx8BkLCnatqKrroqKeoA0PscBj6SfdOlE8GSlWLdsyFe3OLZYQaG/UgFhw/gg+16FEJ
% eteSLvNL931z0V3SA0dbtG79++weQJOvRLZYPuWhOOLzu2Epf0rXcHoQpsoloxpF045SCKKcZdp9
% lprnMozLg3/QBDRDrjn27UzL8flOJt4/fSj3fnI8IoMJoUh+SoY3e+mjSb7/vgwc506bkWs/HY1S
% Y7gKJwSYSsDUqTcZbb2B1GRYIG54EMWtuXC4Id3j8vmbR7VzSU/G42U+i97Ab7DC3PvXG7oOyWEk
% jzyq8f8r6y9dKcEAcydvkmGBJYJ5rbAh5XS6roWZlvhF4lW3fYBwibwf+YNIOeTl51kMJgO2D0Xm
% GFiWFotsUqcwohCTov4iLFsi3XXVtBVOEMzlM68C91LUhO+bwZtHQ/jvCP6LzbeI3t/IT97jQrxo
% pbdI67UFZ0DYzApMXfDWdsPrr+C2M9zNOWE1UnrBcezzmprZiV4Bq6hOBxfHQB1dcIGaFTdfZi5U
% ItwwDO0w8Vd6CONp4ZBG/DPTqry/RGKVgWAe4CXSUNfBe8giJyJShQyAnFP0IsyriKNVM9269gsX
% W47LSLAnwzSbBalsV88i6RY76PyOnlbseKLc8cagIGpauXA1JOylY9CeY8LwQCU+vPPu+b0QLwGq
% RSnq3lpn0y4xyWsGABjL419cSywfEoPa/a8I+7rwnvGKb9wesvHHO8gfqfNzqnk41291EALjZuQr
% YSnCht1phglrR40gxc8KkLtqwZExelcSNWKoWGJOEK0pAgCaF9TDE2L6bBmrXHAKWG7dF1TYMNMy
% bizfd429cr89rGAIYXHVAkDvvHR4Hey46Bk3B8QkQ0+j7o6OfyW+iK3zNH4n/j6WmN7AUwp6DWy3
% U8HfLbLh6JyyrbuMKy+hOjq9eD/C+8QhWyx2S38z4o+To4tHJ6Jft9Ey9qXfx2MQmL3v938/vIiF
% xpxaNjj8hY3I68gWmH53I3JQYYCqK+idUguVq07meu4+nhz3ALnuhDpAbW7z3bAwBK6cAhte1tIS
% wotZkK3kJ2Px4sqWSvfhwp9OSVEyPceStFRQXplLnhsOlTEuljvkzjDcknmjIkhFaHSOJdBjWGVX
% jEFw511LzfZ5LXAZL+vSZG1baFc5lcNtNpLbkTEPbOF4/hWcJEwhzMkvPG3C2P4Pa5EXmYUnxSeq
% 8IgTFQbQXiRz7J0Thuc81b7lN56dHY1upSPQVyaTXGu66N1kqonflBNFd15v0X60UneNrYQcmBVt
% elQCn+JEInyNXLUY+veVaJDDnEvWHH06fIQg27FylcK48kuJ+hSUDgXDvRDflRVhTa1W0H4xRg0e
% ySyX7M2va7DuvkEqOB59I+3q2Yw0djO2N7xSNzqe0nBAZooEq3x4tecRRXn88kPa+afEjm0tyJhX
% WirZeVttKX0eJLLXwPcKX5N15Eacq93i4SngYOCd1nEIDMcEjutbKwalygMVNpFl5hWLOEsOYcUw
% FrQ1vgeZmlX1fK4mlyWjKQNGYEC+ZfuI8r0111JYofnk0RjQh6BM/rw71CriO0piVnImt33f46kl
% mQnWUVtXSk9odIJToQqXuulcmY01yvHMeu2ZJUdXSq7kTjb5YEsHchn7ABqBO3sL2eeeepVIzLHZ
% cPffbGbYQykse4Ceuw+RigvL7Mp8o0pDgoKLLF3BZVKT6321zc59fHL++FZigjuzTHcAYLhscAsG
% YxfT1wBcYvDj0s8fcnGKP2JVQ5RbIpi5Bc6UXCAu/1cc4BH9/D8qgJ4+MbHkh7VMgx177S5k8I+g
% eFwySwFaaJGuhNmFr0vVnAEEAeb9w6N270XyA1ilHmK4y4yAyZSG9EH0psquSzVC6bCKK93S5GOq
% mO9BFBkBXb4krrBd9Q3BznHEbzuEqHRlZJhsRUZn+cwDaZAf67MHSKqTNNRkQ7yGwsCG0a8lsxRc
% ZYmBj4J6DqaStvDmS9QSYFskr6u0FWE3dhJVZTU5rD6YyZ0HgUOtmx1K8t0Ge+u1lEhBTHDmg/vo
% daQxdkJfOKtsOQ1MGB+6LLHHduODAHcQP9mghLpLY/NXeDtryE2bFCiCGYkSnJt5W/npvx7Ucn4P
% xpUdWuO79bSAjSQTDdMfGsx1WVoBqCV2ULYrNZ4Mhz2aTQuqGJvLdYCwAh1KqKlwOFgk1e9yRd1E
% Ik9C59v9qNz9rhpAE6YYrkJDga0FhFDcyk0U4mXOkLOyRVXJFqwPxCN0yOFIoXjHT+hnh8jwWI3I
% 9hp0UVlzh3exFl4KoyOJ9hVLH5ilw2U15qXWsTINbaPpDJwVM9xONp5b+enRz3ufEMZjmD0CU2j7
% Soc/RsIrD467ZVCJo06VOTXRenT6j7tip/9CPo9yO7lB3VbAKUqY6gPhNOPSh566uvIX4ndMNC1g
% CTxdI/yRERCIm/H7nljKPxDoaLquQUvBsZVaLrmMidOGiT8YSCHMH3xMvEhKB2fjWS2UqUKJFI1j
% sJcz0/qhrmdlhKChZVpfeneB7q1qsoGuD2JAZpLhsNCgOISC/NDGJxTdTLBhl4QVD+AhL28BFhJ4
% RFoipR6juac64QmmQ1hVgssBTjVKva0Vla71dyEy7wInFATtQkSOxw1AxW3b2vXhZW3n9wG6vDTG
% gu+xHuA1qLM2eung8btQLT2W4HCnD6vQfIwa5A6cUhPHC3EjKssEppU+u3S3fWxO45LKpY2xO0k6
% MmCZm1buTsDPEWu3OkmuvCGrsd4jSf6a1QtqrYG//W5doDnwYwofXNd5qQFiUgMUrOZO2aQqApZq
% vfsDSkxB718j6VfE3dCMBc16cmqJhtCGcHQ7b6wjsLS/kMSeU8egF5YZnbFKaVhmiZLZUnEHLxon
% 7/hKgcGSxwo+x2yguQNC6xSKy51spAxU6oFSzJSkU8PuSBs7TmTy5lEuZ8PpJA8MEWxt0ykGob6F
% VuniOZMuOtXks3WmnRUGLswwQ+yeqwDhnWPyvCuKKxOYpFMYxJakIhg9Ww0NKUYBy14PMFeQz9j/
% peKgUlI6rjRRo7c+HZ7meDtFS75AIDSGSuWUJkC75RjqRaBBjip0EKRy+sqEku0yIU4hkwhul+Yg
% fKWU0DB9E6aOhnBG8SW80jmBQYgdqZ+K8SJbZE5WOTBAYRORk4DWeLtQ56mY6T7f6o4Ed9500St+
% J9BOiSf7h36UzVaBGvoqI9ydQQzir6aUmneMFKfCh9IQtuzmu7IdfTm1HPU5WDCSYSrHLiC1vbuo
% GH9tdsP8HhwC/bfAZM4fApM5Hh6d7e55c9LX86bb6eb8aUQsQpXSczhnFEldvhHKbpD960PGx7An
% iyxqeJ/dGZfQCgWlECVi3SAYgsHCC7JxMH5fEv4DFYe5q2bhG8iiy7qyyq8wp+DXMe7wY4wEXd0Z
% kZnA5PPwMJqcDQL1LhhpOTw5VN3gv9fWUUuvd+VTYCBXls6TYQqAxofKpNPLdJE1HkWSd0d59Xgh
% J/12Pz9yfDWVTVH3UGNZg3s1kRss0awzZt6T4uYNV036hAxbV3+GmffMw8foclg7Qi8PEGR8PaDM
% du53C9bRxz+kFoYLtic9uQih1ua0RQmmZCLgBj8kySwP7a0RVdGHtGPYFAhvMiYS8HB+RkXzcT1b
% 0Mj7wwH3p5VhFfYff///UEsDBBQAAAAIANKpVV33P47+4AIAAGEGAAAoAAAAb3JpZ2luYWxfY2hl
% Y2tzL3JvdW5kMDEvMDFfbG93L3ZlcmlmeS5weYVU7U7bMBT9n6e4MFWNg0Obj1JW5klIjBfg58Qk
% L3Gpl9Tx7ITB2+/aSdMUmJaqip37fc6xP50tOmsWP6VaCPUM+rXdNSoLzs/Pv73wooVS8ifV2FYW
% UOxEUVnYNgY018JctuLlBlQD7U6aMtbctK9oKSr+JOwlZgi2ptnDnrc7kHvdmBaKZv8zCEqxhYdQ
% Uago7JuyqztLNgHgY0TbGQW224e6+RP+oqCOLr7yL5AKDFdPIkwwAVxAQgjMDk59cm3kXthQnWb9
% rvvejxlSqi4SAnILvK5DPSu9Qzl1kKoNdRRdrgh6ksc+/xaRaYwNq38XGFqofPZqphlbYrAzq2P6
% nCbLhKZDFmfk0/HQOJjccwfMo5ZFnKZrEvPI79x69DE7y+IsUlGo4oREPAp5xHG1WKSjD7dWIBXh
% XYzeZJausTUK4TyLeSmLOVWUkwDbx8GHb4Azi6LdQP55uQThdTGIQbts5ZwERdOpli2DtwBkhxEQ
% Bv1ltUEJqFaqTowjT/DQcULzyL0Qbepex/HfobOaWt2DEmBIVX7y8Y45lDRihuYDaNpvTvxu2VF0
% Kk69w3vBaXIa5fA+gO1iryhO4EIjjb8R/uj2JOoNA67vnoJ8oED3JExjnEadihCUzSHB3RB64uiJ
% uGDJgcL7pjPtLsbuhAFZCkS/fd3Mqfej8/+zWU0ByEdJ4hlgx3PwEZsJvZ5SJPtDX9EqwlNxxsAB
% ViHPyn95o4yTUaYa8OLCwh9yH+ZeZ1/ZCkRtBWTkQ+D7Prwgwmqx0GSiipEP76PHo3Df1bW7xmpZ
% 8FY2auPuDMjSJR4JzaWxYAsjhBLlzQju747Xcvsq1dPgMgL8zGvLsI8kpXlKk/VVTqjDI89octiP
% 1qvr5XrYrofd2NbDjhuthLVQ4Ghy67oTFvl1FUgwTOw23zfpI2Pharmi+B9Nobelj/FV5FfZI5n5
% GoylyPb7Ov1l0F+4Dax/rLDYP5OQ4C9QSwMEFAAAAAgA0qlVXZ7vQ0ylBQAAvg4AACsAAABvcmln
% aW5hbF9jaGVja3Mvcm91bmQwMS8wMl9tZWRpdW0vdmVyaWZ5LnB5vVdNb+M2EL37V3C32EpM6O/d
% FnXCBWygPRdFb4Yg0BZjM7AoLUUZdtP8986Q1Edsp91TdUgicTjz5s2bIfPDh3FdmfFG6bHUR1Ke
% 7b7Q88HHjx9/PYmtJUpbuZNmvC1ysBG2MEociJWVrUbkd2dNKit0JkxGDmpjhDmTQh/OI3AxeDJF
% TpSVxhbFoSIqLwtjSeNLFbpipDRFVm+ttzXgCbd4wz/c2yC8PVeFHgwy+URKcZDnuKSLAYGn+sZf
% 1J36VJKnwhAFkNHNTsZTVtJXZ2KkrY0m65dnZ/Pc2ZSUqCcSPw8VBQfwufr2eumnpImPmwu73Su9
% i0X2zMhBPllGjNrtbUAiqkrttMz4iw+Le0S9y6W2cc1IJaUOlvhglKMLCdnJDJ2u6+THvsPmAYje
% EjwsgD9tla7lGxNcGoksi4/0eqsusJItPgKRG1zNt/UxYRcAm6dnwuur1cDtn6YHKHz7TRwq2S9A
% VedxR0glbUyp46Hu8YC8UuoZr6wpgO8lOwVcuMiX61Pi3SJX3PkZ1Rr0FN/FSyCx8+mdIQv4F5EA
% iDj7bv+Q49LfL6c3WmlrvWSu0L4sHGy1R+gBHqVRT+d0IyoZLwNG7YyW9IEIHuv7KR2PPz+AGHbw
% NsS3WaMWaAyiP33mfN7/Ig6HGD0cIRzs6oSy9Kjx/cTO+KXfSrFXq6Zs1iti8Bk7cySOkg+cxCf/
% ek7opaXDDnY/ulXOxXDaRe2aoi+UtigPvYo0ZnSIy8OG3ht5LjvZQ7iL2t3ahDLyhFzXuvHg1r3T
% Sxc9Ufki4jiDVyIzZat/KaJbWPE1pnekvbpA4nspsspLEWrNDZ+0rNXsyIwEpYD2wNpFuSqQc7SC
% fNoFfBkZmRdH2W9qkHJw9rZTV9CebgDUNKAJ4wCKcs+nrS22wIJkzaddUcC8Ol2XFwMFplbAVDvR
% 1Bbg83XyXaJYNZ3adWv8tl1X39euN1TUYBmJspQ6i9vOHLa9u+r3Lu33GEqocUAfeXY/vzNufQvl
% hbW6aRcCZxAJ3XJM2nqmbSXpNQnJZaQtyvIuzu6NB7H3QY6cn5z/Rh3vqeW/YmBTNNkA5ke+hZ/j
% sbif3e0xh6vdHsYW9qonJd8XwKk5OpygHFZ0/She++FbP+SRO0G5xU1R64zrIdbFbadDx8VNQNfB
% oOHE4GIuoW/6lTvPX8F1NpyFssF2YPdRLG6Y9yH0R3xrwnKl41yUKCG2opSFk8gtNkc/mAR1G4n3
% Ef4SCaPsPpdWbaNFJN19CcRKilIaP5IdYeHAAz0OS2H37clSRQw27UVdWXWU0WKdsMjff/Dv17a/
% SpaLk5MBcrOO52xOWfwz/Ey6blvy5lL0cHEmdY2f7SQ0Lhy8x96h216DLq8jOFtp17xb4NyGodYg
% s52DBiIMyrdjCe1kdXVWOTDM3rht4IZNSDbcDeP1hE0T6IpSCstvbcInTHGXINv4hFyybIO+/lJl
% DCHRNU2AJZdQfzD6sq77NUma2fISFSaTJlqULIJUVV7naV5kTvSpMNsqWjQMsGgLxMMHF+A1KE7v
% uL/NxrPJ7KfpZDbpznJ371xPp2z6C5tP2ef5/1lXV0d3qW83fZlMLii2HBIYoTi9xcR1zOyuZNhD
% DgClUPo3mxwbHhd7u33W3PqYP6VhsRJ5eZCtLJLBjcr6UXhVr9AzN2sVSgEZvV+360xC0Rru5yEx
% YLkH9TbhYPsO42z2Jbjp8QL/XsXqKwc1+G2Kea84+HWd4xwBUqrAB84jZB4uheyIM3cLkK7J6YRc
% HvBISGE0IeQo4S03oLMokwfpujFazIAff+jAlR34AraUdmwVtYVwRkrHFEaO/K0gbRBEC4er/1me
% 0oOqbLRobGDSIb5o4X752VYazB7/qxtldV7C/dXhZhgFRg2IZABTPU21yGWach6lKU7hNI0WYRwP
% /gFQSwMEFAAAAAgA0qlVXQ358wF0AgAArwYAADAAAABvcmlnaW5hbF9jaGVja3Mvcm91bmQwMS8w
% M19oaWdoL3ZlcmlmaWNhdGlvbi50eHStVctu2zAQvPsrCJ8TY8kll2SAFuih96D5AIOPZSzUlgxJ
% DpoW/fdujLRIEKFNg5xEjqglZ2dGvP50c3Olajcd9+meq5r4jvvLmceDKjzOXetKmvlC6Xh5HLp+
% Vm04jfPuspMl8kJ1lfu5m+8vVOqr+vJnvln9WKf9flsGblKjE3Ta9kPf822auzteX6l5PEnh9ZNt
% JkF1fMCGw2Hot9Oua/MZMwJWvh2Zt3k49fUBDIJ9lQHKUyp/53HYygn5lsdn254LoJZVowytPKdd
% Mo5ksqaUTQQHpWiIZBlAJ2zZG2IoNTMGQ4VSo4o1VUiBLelIMWZCMi6uf64+79NxOneuDH2drhRs
% QMe30kf3kr4N8JI+2v+gjzqGR/7uGX822WQgXQp5NJVJhh5tJWswO2tS5gbUGriSBSnkmAMBYWLv
% rYFl/hb8W/mTX5BfuwX9iV7fgEje/DYAPWtAM626liPYwITGNB9Ea8BonfHa11J89s5HF6ptPnEI
% pVqOyRUWc0S31ADjNyb493SAHPYvDrCvcIDxhIsJiM2wCO5sw1q4gqQh1IiOIlTWhiL6Ai1gcs3p
% gJg552QLAnkotuhlB5j45gQ4/ZI/moUEOHjk7/7N3wKG5T+AyJ+qz1riDFSjbgC1QNA1t6xJmDtn
% k8ueTCsWnTOSAEpiEF9QS2KW+OMGvXnPBGCwfwkAvSIAIqRb5K99Ida6ibKtuQo2cZJ4Z5E3Fcl+
% FMvHXLU2toBp2Iioaq4etNG2tiX+OmwMwOr6fLfwt1Tmp5fJpI7706TmHavjONzJhw+HHo/DXhoz
% 9GrPh0NSPM0p77tpp6SLqlcfP4itVr8AUEsDBBQAAAAIANKpVV0/Zxlc3AkAAFsYAAApAAAAb3Jp
% Z2luYWxfY2hlY2tzL3JvdW5kMDEvMDNfaGlnaC92ZXJpZnkucHmtWG2P2zYS/u5fwSYoJNm0Yvlt
% vU5VwN0EyQFFcEjb9IPhGpRMr7iWKEWUN3YP999vhqRsyXZSoNddWJapmYfDeR+9/O7VXpWvIiFf
% cflMimOV5HLUefHixdsDiytKVMXkhpWbfiqikpVH8sxLsRW8JPmWbIUUFSef+x/4lyqXJOZlBQ9j
% VnHldz7kZJvmrBLykRS5kIDGDxUvJUuJQRNcUbLJv0gg2/AN2bCKUZKXFrgviZBbXnIZc7/zhm/Z
% Pq3mdpcjqRJOlDiQgolSoaAVQAgJvwte+hU/+KTf/7wX8W5Otvm+rJK+qnhBduHIxxN2RFbkZUUS
% phKQp/75pHJZ36ujqm8rkfFOZ8O3pCpF5jJv3iHw9yURKScpl7DyY0BAWYQt+8EqDAdzwvwiL1xP
% E5a82peSMIPBNhuX0Ygq8SjDwGLFIfPjvDi6Xm85WHUzdnAHFKEjr292MFBb0JCgBzwrl/uMl3By
% oAHNLMWqFyJm99DcVEsce2brIleiEs98XZT5Zh9XKIcVQGyJzCs4xBH2QjvUvxDcYoFkmvYlecvi
% hOT7qthXJM75FgwvuKyIUOSHkGRCuiD0Y5UoTx+Fma/I8zV/JCoVmnUf7teGFo5uidprBgxorT4u
% ngfGFmJTJaGLyL0779WrmV5dhOB6/rbMs3V0BMd0I8fxn8Af3YNf5XZNs1InFVWVcsfTKtYKhh1P
% qxrup/8LLrqEQ5cP3UX3J++Cu2sP3LMOEFxyNm27vJAJUZdiLnoaanUhiUBJSiYfufUvJPeopvVW
% 1k2yffo3PUOwUEKgu+K2owI/oB2810RE3ySMLKEGZSEciM1Xr0kURksRzVe1XOgbEHvhQAcf/orw
% 1xxdM7zp7JqTp4qbs2n4Ilyi4x3ooGl82I5l5kn/8tGJN7rJG6Got3mjMy/KqHMBfK5lLWhUePR6
% PaNR5nn0hHI6xddQsq+gFIACntX0J8w8rmA9EXk9kM44wyPbKyWYXFcsSrn7gZIYNPYhl9z6x7tw
% uVwGq9XqlKDk2ckC+qFX5zi9U/4lXJ51gNRPZ2q0oKSAj8bVX0p7G+6mrUak14KzfvAEKRc9E77l
% nIBNglWLRBsc1lFJ75YScvTyaaUZf8AYMEue2QGUcK1d+6cVJPtPXq9GQb4WNRzQZ0XBJZjj/ORd
% vQaPWyp/Z5QsGZinEXNxqLe6TANfLwGsoRStVHq8WSN6T1AlDt3jV0uE4ul2DXmkci0iP7CsSLkC
% K4NmQLPwoUM6gu9+QEe0P6FDuB/QAR3D525FUwHcxqCjgeed/YJpkSxeW96o9YgwpTgWXpSszkYe
% GPekJut5F84ZDK2KIB2SjFUJsRU8zrPo7J4UrNDWzrsr9RVtArRb2+2siGqfgaHDEHcA333ybhEV
% YVgs5/N+UNfPX6F/qb7kKF7KM6ic0CrlUhFWcth3w9FXsJ6iLNDrgFsCPpBCx6UylqbERrLybZ1b
% gkP2RqvlaHUZgfcNA5TNwBw3z8xUFS4xPIy2oSTgzWIpQGYMFHEODljsgddj8rC+yNq6WizxpN45
% BFWKUY/otQEB39z+QztYPetjhLjd6ckixEXj2/xQpNCoVGvTEq4bPat7SmYXLjWc2GKNsZGILZYR
% 1Sp8XdVjZ5KYpXjAuEFy6MZXxeMBw4lCMJHlkE5oAPdjGtzp64BOMaTqI5hdDdmYDu/p+J7eDejs
% js6m9G5GpwM6mtHhABmBLfAuOTFKA4jZKZ2M6H1Ag9E9bAUb3U/hA9/TGQ2GMzqb0QlAjC3Q+BLI
% nA39AiC18BoamGd0BP40oZMpSgPX8R0dgawB7gtigy2HHsB+HbGNuupOPAonHoMmZvQerhO4r2Fm
% 34Y5QdUiDjSnUYO+0fCANLyzLmTdBwJLl6BYtxRoshhN9uD9ZZANRxYpL8WjgCknXNh2/lbsNSNv
% cTPspA0K2QwK2QoK+c2geEk+k5//GBO2liQkn/9wg4n08BesQcoiD+snsoSpLclzwH9a+ZZrsYUp
% DWpuCVlG/KmzEolySKRKbLiCeekZEld/wx9LzqHVB8sTlZPg3kx5MIrtcRbg/s1e4L5x7vdh3TN+
% tQ+Y6loPlKig91ozD1C3KVgB67d3TrXq8x4yZ4gUtQGAhjbuz7RlokLjNrWzGG5qvjz6Hlxs0pXe
% ZXIxTAsEDiDhAw4lrlPnFCIwYYvq6FDLCspMQTdxwuOdHlg3QhUpO8Kk+hHUlR5lngmYh2tGGIeN
% WnGwBVVUNreDhxjXfU1+C03t1f9DasvJp1CHCuaQGUb5aIBXuA+wJA8N0e8h5htg0jlmCADAM8Qr
% 3MPK3RnvTdgIo34dQ30bRAN7f8vCdw0Df9RmM6G5WGC4NYPVrATD3hQ07V0wfaSatUH+G2Yj2YbQ
% a7Memurcyd4G+USH3RP3JzrujZHpmsMQ/E5HSN2E/Ee91frTx/Ck5zf0jfY78KeP145UlBBarvPv
% xS+/zBtOpPgzl30I16z58oVCMPZ1MNYvPkSl80PDz3BYOu/jO3Sb7lUS/lruue0C7XsWd0dLa9JN
% 6A673dLrurv+0OsNO6f0G5Zd+2gHPp/L9AiNXsLKDDo5hX1LBjnkZ7YvsaHRHK+tCuAQEa+L9VXl
% 3fTK3o7uWsl3t1rurpIvEgaNLke/6Wk0OmVrZPg7CbfZ0/117r0+iVfX/Jsz0qYlX77dKl4ZbcPW
% XdnXOus0xh1D8gPMuYmZE/pmxethduq0pp5b7Wg9wCP1cm544aDE+ZeMS85Mt3llNacFlYSa2zDP
% L2KDxm31PHgoKqo90fqOT1HR0q4VD+MpwRKM8SD5I8PZldTvGc9vmxwK7kkbifqhnrKSOgO/ZyrB
% t5Uxk7mE+Eihwy7ZETJ3s8sGAnzFp+OCy2dR5hJ7cpN/N+IRpqHQvin0VcKGk6mLbwr9zT4rlPtA
% FS8YnDMvVeg6FP7njuf5XMb5BhpLz0/4waC03wf+x9k58x11SmdeUsck/3WU7+XGmW+o03yj6sxx
% Bny4mPsdONSfvMzXEO38kZfrhnKABUcT/By+s92MLiux1+xsLgChA2qBrGGH2gLOHBMEyKV9Y218
% Yq6/qGPUAnLrc/630wE3Xa8ly/h6HYbOep0xIddrx/hjY8i0M+Y3GnNNUbHykVcwhLojSPs6VB37
% htfBk6ij8oHk2YavpqJwnejrFK5jvTLR1yliWFOofYq458gEp9Ll1+x4DiAFK1WIjuLD+fMKHcr1
% XluIsJky60V1Gvv1z7OnmpTecCJDQBUMq+sdPyqTjFuJuc3rvE1ZoXQZiHO5gZn5e3+0db53LwXs
% a7mvoVpVheMb/9Y7fFIAvW5bYNR81m/XoY5A22J6QphaM0bAgpDgBMQYzqSS/BiSge94nf8BUEsD
% BBQAAAAIANKpVV2o1BW5BwEAAF0BAAAxAAAAb3JpZ2luYWxfY2hlY2tzL3JvdW5kMDEvMDRfeGhp
% Z2gvdmVyaWZpY2F0aW9uLnR4dBWPTU/DMAyG7/0VPgKCyk4cJ5nEgcsQlxVt4ozSNhuRSouaMLF/
% T3q13+f9eH85nXYQpgniXxgK5Nt3v0xpgDTGuaSSYoYwj7CGkpY5TJDTZYYhriWd0xBKzG1z6A5P
% x7fX7th9nGBM4TIvuaQhP8K8FPjNcYT+BuUrws+6LOddEz4f4BlMqxVZFGKLiq0oK+TROSaxbC06
% 5euBDRujvG/6ilBLShRuEBoS1izGOE3o2ZNy2jki8kYjOcPN/q6/r5BqfXVwLEqs057YGWSvNQkK
% VuVWgnx1IVf9mj3kOCx18hjXdK2zrxFCgS0eW4OyidAr1o4rxFZLbYeqJovx3kn9eGEmq5p/UEsD
% BBQAAAAIANKpVV2mMxEPnwYAACoNAAAqAAAAb3JpZ2luYWxfY2hlY2tzL3JvdW5kMDEvMDRfeGhp
% Z2gvdmVyaWZ5LnB5lVdLb9s4EL7rV3C7B5MSJYuknmlUoEDaRYEiTVIs0CAoCtqWHHUtyStKrZxf
% v0NKtpxuDruXWEPOfPPgvPL7b8tetctVWS/z+gfaH7rHphbWq1ev3g1y3aF13nZlUa5llytUNC3a
% y33eel0+eOjzobo5oDb/uy/bfPMauW7dV3kLzDtUKrQp5bZuVFeuUVPvDh5gWmW1b9oOyXa7l63K
% j7Q6VPsDkgopyxookhT1KEPKg/NVs1N4MSCJ+gWxOjgeHGbdwi8e7MGRxMb6r3Wlb+CA29IeXGl9
% MQD5sJf1BitvUxYFvqUDsa/cW3uir4Am1j1w3tpXllQq17Z4CqzalcXhKKa8XbPFk7qlUUdA0r1a
% 3pIs82fBSdsXZxK811zcfqb9RYmJ4wtweKpfKTzQ3mHExYHd27ZwMAO3HMYJkL0Dp1KT0fHMgQMn
% IAb6LnuuEHO70xCd3RHw3JAdRKkDfOtOPZZFlzFfqwkdrBGTSGuxA1CTAMmTmExGuHzUGwitLhKT
% NS6Im/Mkhb8iTLVFGGhQzvzYZcBO/uXy3TNHR0NejuaRkZEJ1ebaDeky8qLA/SzgTgLCMP6O3u52
% 6A6tm7yAlC7zulOQijnaN6rsyh+5SXD/Ul5mCUWrA+oec5WjvN7sm7Lu0Krp643yjgohDhCR5I2P
% QCsC73UgTuToOJCW9SXKJickjchrdB9l9zP9b4cjF3INHoAH9qADD+8AueeyAGg3iJLR7Y/ZYOPI
% Ce0BcpLrKhCmFiAqM2QBNdy0mC0xcyAXp8T4qPM+yzT4EhsRE1ODYHNiXWcB16pjh/mhr78iJ4gT
% pr9CJwh5rL8C8JfxxNjoxmmqRQYnDtNIG5uIyHo34SSAk06IYSLEiJiKETt0QxGkE2IcBpPXiYEG
% wNjXfgsm+OzXWtbrfAehsvFHl2n/Ibu1lya9CXHu4YYtIcmXeDpyr4/Xo6e/JM//R9RfkSkuAH83
% cdriHF13T5U9WAhfU7fXj9ATGtMHHR3K/SBIaMoZD+A78QWnPGR+xCiLeAJXIU/8mEZhwL8SajB+
% hQh5GNE05jTwE8oTHoQ0iNIAgEQYhpSlkT+KvjuJJvRBh5KyWIScCp4GPuURSwIqfB6kNI6iOKFJ
% DNlNRQDZRhmLI43y1dLlcUMhc+mGmiJSqKzNiFAXlvnNVF9BN/tQd/k2b/HatI7S1FVJ15o7NzMC
% xgkeESBb0a+P4BpTbXtj3xxrGX6IqzWYuB4l5G6H11BxGt+gT5gW1Pqd7MqmhlE0D6V1U6tO6qrv
% Vb7R/Pu2aQpF0c+ye0R1g4pdA2L11h0rft1UMKdK1dSeBTZnyjuiWjAS+K75mbe634LTcI0ZhdT4
% rstJ2Pb0SYxx37WyVtbbHAdzfZ4w3mjpKBWU+b5/un526BoFggYJMdycMRpq5qd+DxM5m22o7XpU
% Ws9KIafAFGdkYLOKSTgziH4aChFTxmMf0pGYPjYyXOp7ISBNT5LCnc+gDRqAJDUmjXBpPNo38cf2
% zO8w9oJAEMQxVIBPLllyHgJwGobcOb8rDAnZKiJy6Z8zQzaDp+N1NEZoFKXpGMJYuz85zCGcKSh8
% c4KALjjGSByDJUZ5RnkMfZaNzwBg5AgSGiNO0Cb18qr5IVe7fEoj09gV5A2sUG1T6bkCb68zEjIA
% 8qt7rGBJ2sqqkgjMKHMYMoaiT5wesrM9aOR54ugAy1DhZ645gC75ZLrvklvvfc1vQHDhL892l8js
% LpHeXSjUPPFabWb+CZMXVp/3vgvNzaA7eH4FcwDhf4LvYexxBcsgLOaC2AcHP3EYzPbBPmhj2Jkx
% 7GgMVPdBDxpjUQzfhB7+g0XMxbF9UuUA8cR1Jo02sYRolaNN+xbCjhc3bz9/vtBNAl5AL7NjGCHU
% 5QZmP4x82Gl1krfHVqHKbf1s5fUgzFZZnBZW7227hXZSdzeaajHxzPG3v+rmZ/0NuBQmD+wryjL0
% sDjbhxdfoUFOy261ryR0G9h2KwtVXrX3NnuVRdDVNnmBir7Gw0AuYK/u+rYGBoiZcRGOoZ8sq+lF
% 4UmnNzWPaiFoZV1WeUVZb/Qn3slqtZFoGC6qcewDNAVoirDWWuAF8xhfEDpTYkHGftzJNgOup7yT
% mBO38vJ+l7egTLjzKbeP5xaaIn796dq9+/DHp7tPf34++wcAOmzddGPXHfeqsfNeLGZR+c1G2QKM
% gR7dYmMCDUIyM6zOrrWDz2/f4xU5Y9Bh1EzkFy6oL5gCGwh1W/6QZumTsNmdSZ4Fy2jhE8Q/UEsD
% BBQAAAAIANKpVV37hZqcqgEAAMkCAAAvAAAAb3JpZ2luYWxfY2hlY2tzL3JvdW5kMDEvMDVfbWF4
% L3ZlcmlmaWNhdGlvbi50eHSFUsFu2zAMvfsreEwGt5Es25Hb5RB4wQa0hxTpYacFgs0kQm0pk5Su
% +ftRdrIVQ4EJhiU9UXyPj1ovN5s7cCpoa1QHnd0Dmqaz/uTQQ4Mu6N0Z2G0p2DxjQlRz+Axu+4n+
% 78AqWY95FqWYccZAe9DG6xYhHBCOzr5iCw73xHIN/S9lITPJeZkLOVCygfIPeKXMGeCbagIY63rV
% 3VjXotNmD0Rugg6acirTkhz8eVLdAFyuPkSZ/gV/QWPNDh2pwHuot7GEeHS0nsJfEVrcaaMDUiLV
% nYNuVNedkzUsIJ9Xkucy57HqccA9PNEJF2XFcllkZSUFWST+RlxH8jRb/8g+jM0Ey2jIofyioMVF
% c73lM0GCArpeG2XCzVHpod7RONJGOheQ/cM147KgnIxLKQtOgpP1eO8OJhOWAp+mMMlSEHHOUyji
% XKYwn04vzGUOHrGlPo5+N7Y/dvgGX53qIxR3PoXBfw8ZNd7Z0/4AVfJ8QI9wsbA5YPNCLXEIXhFy
% viApNTCAim/F7oC++HDI57GhcWcd9rfJ8vERVt+X9TPU31b1wwaiutWX5DdQSwMEFAAAAAgA0qlV
% XV+gCU/fCwAAtSQAACgAAABvcmlnaW5hbF9jaGVja3Mvcm91bmQwMS8wNV9tYXgvdmVyaWZ5LnB5
% tRprc9s28rt+BerMNSRFyaRkO7ITZZpz7VzGnSZN8qFzHpUDkZBEmyJpkLQlZfLfbxcA35TtTHtp
% xyaB3cW+sC/6xU+HWcIP5354yMJ7Em/TVRSOewcHBxcb6qbEXTH3NiGLiJN0xcjnKAs9YpPzVRQE
% LCWcJYxyd0XCKGXDXu9zFp7lRMg94/5iO4y3vY9hsCWfxPLLhCQpDT3KPRL4c075lvgJELrLfM68
% IXkXBIQmCeOpH4UJyRJG/DBlS8YTEvHeggNbuDMkl37op4ywDV3HAUsI5UzxayI/JOZRtACkhWB9
% yULGaYDPEWfrYe9rZVGAIh9+SGIaMz5M2ea1wIuQ9zAKByxgaxamyLDHYhZ6LHSR996Ft45CLwHJ
% 3AA4910guKapu/LD5SCOgm0axaw8F3QLUkRrUohC/HUc8ZRcqgWQn1xKGBCQp1EUFDDA2zpLqcBT
% ZLLQrYEEPHNcCpqQ+8DKKt9awAER92lgosieBOBgDjxKgnwWb73efy8+fyRTolkmsfTex98v8MUW
% Lz2PLQj1PI2aZK6f9Qj84yzNeEg0em3N+nP4YRJ6beOjPVMY6yzYh2EgxgARDERAXEs+9tWilVNx
% o/BGox00TCIIKLAErMDwuKTruERyZyQKOoz4etSkKiCFPAIWf0hotAENwRecJc3A4DR0lIsWJMDI
% n7fgxHhzQG5KIrgN5L0CLzzaJA8+WIeJu0Y5PK9Z6rtD9BEkE4LSAxYCWfGaRin41pSgbcQC3ss1
% TW7Rb8GMS6bZ5M0bEiouBA4wCyhgwGIJsfwSpQotxI8enCRbV8+pYt7sx8R//kJy9PYtQP5M7DZE
% /Qz0I/WGRvFn1zczvYajJED3wUczRy6h4EwtHOCxw7mfOi5EqVTTdfIvMjrrIiWdQxIb2CWdXL/I
% k3g2BYZe9QqxLv3AY7jrhxROyw1PAf36+lLb6EJZG6Gs6GEm3uAB3+msy7g0TB7AR6bkUlMsPaLt
% GOBCtkk1zW+a8wain44qyZVpkt+jkNW0FWOsw9WG5aWIl5rVgP5pSm7qoBQoo73iGeoLfpliqQ6T
% SzSQTyX3/j3E56nAqCIpBGMqIfY57E3fNknT9xJBT0h8WMfOKdw+QUGKBRRuZ2QATmII9m5ntZgg
% BRHmn7MkdWLqc4jzDobTzE3B7mc1426kJn8pgrJWGkPEqSi4Zxq6boUbfyESGK7uMZANVtUqRsLg
% LK7dz0TcA11cBGBgma40fWD3SgoJal7A/kVkuPBLQvc0yBgkB5AKFSq5Qxy9tycG+F2qBAHEQfuD
% QGwS5i1Z/ZCCoxtdb2OgX9j9jbRxV9yJyVspQHfMacgGHMDtMeEsU+8LXnoNNVcRqk4gGc5DLQYQ
% 4Q9BtBR5gAOetpGhI5lOrEpKEHkbQgyUDkGUZFCwZAl4D6JCqJiSEaFQHK00bQNUD7VN39b1Ihls
% RGxQDiVLJFh7SyyxsEMXKNHkGi7ujJ14iSN5GXe1VFLc9cYNEcxX04iA7wOLhiB0qI2M234ldkry
% YKKdPIBxHuFxI4POE03sAmuAJSgDpqHZg91OrwXXIurir74gIULZjrydEouwAKpBCTMQmwpSGeCy
% VD/PwzHUG3gTq6ax+/xwJI91TeK1tge8xtORQQ/HfbcPhx0Zc3j0+ioAvH/meQ2CY4P3QSumepir
% ogIKRQdyXOrH8BD4IRTW2tIkmaIcZXhvv30vjEWDeAVEdmiz5RAKxXWiVQwmrqlJXNxmYbaGOjdl
% Wta4qHOWYsYK/ARSGBLUW9tw2dDudhdemkHlreFLHQ+YvcbVmUqm8D5cslRAmqKsgPiFCX0HHNZd
% 4Bs9c6V8puAdUZVw6Aku5iIk8L2iNhqG/soPUMD/u8YwsSIdUEs7zDyhzqpKB02VPkOt4iIW9Uuu
% PyyIQYlmwVcb6Vn2SP+OHfwQOikNYuqd0ldnrVq1QNy2QEf9VYoq5NwBn3eCdUVK8K6bnQG/nj4g
% QWtF/6Pd6mXEk5bS95V5oqd0sEWggRNxj2HCz9nmEL2nqmnSji1LRURRgsKGVYjusSVnrAywxw3P
% c8qtScPlMqwpNThpiK2ahBmMIH6AMrpW9SbBkV7Pl/f/NEHUwTdMplLKga+fkR+h3zZfo+STdCFr
% fK9XmGg5vBDZWpM9XCp5TRE107tiVjfsfR3W3df9iAAP8Ds/1jIT0NpBwFUe7MoQpxpWbHzrZ1DU
% Wmf8qrO9B6zBMe2Em7fpBWwBZSvA5eDNpLMvE93rnbSAjLz8BWGzPEOH9rlewfrLFaJo8gIbaJS+
% NAfEMENSgsOWLcziX18CUZQNfjwCiMQVSYrAArYvD5bLSONREiND2NAkJRmA18U8oqZ+WY9JhUyl
% mHsBoJj5YfFlRCkScQxBKNUWB5/efflyRr6J3e9qkCBD1SAPVcT3WJj6qY8zstADSdhdRgOxcACl
% 0xVGgx65tky7+t8Mgur1wDbl8kD9GBTrza3Kem0rX2/uVOBrWzX4yk65Xt+B9Vk1UMN9WzAOFTbT
% Kr3YpFo10yDQrmQfgcYaXGGj58/ag5F2C6436WA4uZJdo3ElG8ZmLQ1YeIz2ChOnj483spS1WoHv
% 2Ry8IOdTzbbIh/4YEK70Q9saEvLvKF2V4yk1Eo3WcZYytHs+t0yDrSIy3xI5p0LHgBfP58ytTRmL
% UgiehwJrLkYccDgOA5si4bxwbFy15zgtOTrULT1+Dg2DPEJeDMz5llTErmucMpdoOBJ5ZDI3l8zc
% PQEFZ0tAFNvBLhV/36l0YFYSeuIv11QUMpWRrFZYqVLVNMZvmM7t5w7j0nzyBXEHlSpOhYdG+Mmw
% 85Lq2gcmBVK5KZcubezfYYTJqj4eo3EVLjrJnXBlCWzmHlyF1Y5enU7so8mR3dyWqPb45NQ6mhyP
% Tk4n4/H4dNyignNXYwSd7V0eBOVY7ARAbasOXOv5ALzoh2WAVPHxCuddyS17IGV0eE3OnTFcG9yK
% owTC4T3DkYz8qkChTd+mOMsPtgd6jSKZHpjATSxmzrZlGBNw0IPX5A+1cac2xE4N9Y/DT3+NalBK
% VgW2UcPD9rW61GzzVP+BS+Qx9JzqhBIJS5L6QA1RnkdPiSAHXeXYZs/4S5a/+PlmAdoDTQouDAUI
% slatV4NDE49ASuMSYwv8NuzjfcBv1Q1SW6cAe4IU7Mnx2BpZ9mQyObbBBduucO7Yh+OqYgaKtToz
% YKPKe90BJPzZgdIFpNBK+uFRlDpzTt1blhbNDWZ9s6h90JOtVyMLnP+V9CDb0s3q8mmxvMfXkWDN
% 3TuuA56mX9sz8mYvjLhTeKVqtAKI4dwSjB5PRhPbPjkaT2qMFsvdjFYGI4GVkxbRozoysWrM1YzE
% 80EZDsaKYVmiTLIl1rDUIFDgjoF0ysXTxpWVRKcneN8t+Z0vgapIfOAD172H7MjZEo7sxHuSmVJL
% ghlLMFPq6KDuITLCOEtO107+8bKznRxZoxMbvHlfSxlWGjK0Q6OjTATp/W3lPcNmWKXaVlvW2at1
% tIJH+mxf892ADJtNIyhAHC4+NbaBn8J/InM+83tV2dgVfZ1SDMRIM1cSFjNne4cMKEleTor0Wi6Y
% +WdP0QbO9UYP6JHmGOgBOjdGvvKuKbb8uNQ8QPPMdsIQc5Hefrs8L4tU/z2WnNT30+KDk3iC/8UH
% VPGNp9Uw/V12uPiKsRH1NxSCerUi3Ohdnwc4eaOqCIx43dacc0ZvWzteewD6grwDBt6jC/eJ96Gz
% iBiSD1CDuxHnLBAxpEFB9mn+TtbZQOHTl1/VF+Ei5iTQqnEIT2u6BJKZB3Hnz2Gv/l1k+vgX6Rr0
% E/Vvteh+rOam3e1vjBEdHONO/ZZBP64mlxaGKvamZbVXSJaXHM4T5UZhp3+66GkwWy9kkGf7x3v0
% hDEPDCpbdWzOAraRjpQnA5OI1j0hI0hQPMqWK9JIaF9XDO6dKlXV3+hgr5dQWNnW/gqGqj9uUX8G
% A/WsHAwUf5CC2Qluh+OEdM0cB5X10nHW1A8d56W8JZ2VTWWnNRut7FXb8cpydxKsyvjut9/IxZ/v
% zr+S8/9cnF99IajGi1+B3f8BUEsDBBQAAAAIANKpVV1yClWdMgIAAEcEAAAxAAAAb3JpZ2luYWxf
% Y2hlY2tzL3JvdW5kMDEvMDZfdWx0cmEvdmVyaWZpY2F0aW9uLnR4dK1TS2+bQBC+8yu+YyJRw+wu
% 7O6hB9chFVJtRw5Vq96oWduoDkSwVvrzO6zjKI9rJQvvfI+ZYXa4LVdlVWAzr8r1av4Ni2JTlbfl
% Yl4VESFFGosskXmsEhtTnlgGCRkH+RnWARKsY0UkgoMSFad8ZmliKBaCnxyKCxWkMkhT/ieYROgA
% s3pCWJvIeGIVq8RUeip2Bgg2pCU9ubSUiRSKca4XEnC/lJpo0bvdrt22rvOouwb+qR09mrbed/3o
% 2+2Ix3ocHROHoT/tD+g+U5rOopt2cFuP4qk+Np+OtWepg+tOD26ofdt3qPeDcyN2/YD6eIS0OXz/
% 5IYR/Q5103DKpn1w3RjEHg9cDlmMttseT03b7YMttDNyuXJZrO559khQ/JwvKiz5Rpbflxyv76py
% Wf4qV19xW34pNvix3twwfjffFNUa9xVf0j1PX+LqGtOENa4ono48VAZFOPNkGJfhnCGzEw5xDRHl
% IArUOdQQSoZQTlqDTNKZDQIZWVCaPSueIeLrIaMvWDASQWltX3RBqiISsIy/SC8w92pzI1/h5ywK
% SqZEr/XBkkWUwWTCvrFcKH4nrUnpt9w5o4Y0OrX2nS9Y84gMdK5z+d56ofntM6ukkB/4kF2kkMpI
% Y+mjP6TQPKwUVnAHKpfcorHaSJ2aTOdkRaZ4uOZj8f//m7qd8wq6vzUv+vbgtn8uH8MMqx67tmu9
% g3e8t+PpN38s/uSfN94fHB6Hnje97djz6IaZd39n0T9QSwMEFAAAAAgA0qlVXeIok9d7BwAAqBEA
% ACoAAABvcmlnaW5hbF9jaGVja3Mvcm91bmQwMS8wNl91bHRyYS92ZXJpZnkucHmFWHtv4kgS/9+f
% ondWq9jQENvkyQ4rZQOskJYkynB3q0No1MZN4oxpk3Y7kBvNd7+q7vYDyN2iKNjV9e7qX1Xz80+n
% RS5Po0SccvFGNu/qORM959OnT6MdWyrCijhRJFuRDdtw2VV89yvJFRMxk3EnTSLJ5DvJRPreJY+F
% 6JcKyBuXyeq9u3nvOs7smZNErBKRKN6JkzUXeZIJlhImnwp4UyTJgYEo4DNmyFgzE8VzlZOIp9kW
% mLmTiJhvOPwDmThhTyLLVbLMCfhDGJF8I7O4WCZRyokxRyRTxtaSS5WskiVTvIvhOSuZrclKQpDA
% APbXm0wqMrYEwnIyNjygRqosSyseY0U5zvTmr6/DyZQMSOD7zu39aDyG5/l3v0+CHwtnlUkiMDDJ
% xBN3A0pKgTY59/oOgU+WxiCiReedYKFp+q3LNhip+/2lT/Ji7QJj94kr96WTU+J7BJXrrLkdUOxT
% Enielt7/INtL7UNHUCLagffDc5yYr8jSBcKL9UVyVUhhnRELYw6NGd41UzLZocDr4C4T3EolK/KK
% G4ikfuXBK0QlmlpdF229dgKvrR/ACarNg9+UuGFLUy3N9zxrFLKQvrtTSt6sOUYjUP3WVD2d+wv4
% azHIq34OFq0Ikg3fNTkwZMf55+hxNrkdfcGd0lrm7tgFZ/Cft6AVKURS2KPX2j949OERSJe0ZylX
% Pdq7QJJ5aAqfIfnc8gXXVL/rb+RamNjylC2/uZJCTCSqsynJYED8OpVlAhludAQJ3GMMPmIMkbFy
% kLA2vp0FGAC89Fp28dqWTC14ZgSNx5Y1uIa9GP31MLqdjYaQte9aBsrZ9/oYnq8TdW5zcWayFFzo
% xJWsgWE9N4sXFfPlHlfYUHjWWAn3TAUmm35j2aoP6JXeyDDEh3o5/Ej6wEZvz4bfoFvlVzS8bEga
% S00+ayWE+qBHas5K9aHNj83AEZu1dl3GElxau5c92M3wrMFqDQalM7AHgX+FDD8cZyMTodyT8eRu
% MhuRx5vZ5P7u5k9yi9U/ntzezEYn3jFClbCEC7Je6Hl1lSlAba4oyZcMYHZAXNkW3i89Sq5bLVe/
% nJ72aih6Y2nBczxsC6cJSm+ovTyLc7no74GXOeXm8Je449G3fYQzqkugNKfJesdOtXs0Mt8NaFwn
% Ar1RxSbloHrjAoH+J9m4LaPOa/CyPAfotyIDUp6BuaBy8QETqso98huc3mrVbAMI0BN60n3JkAuM
% 5kpSzY04d7ALIbVtAiDSYp4xYZCxY0ATzYhOeLDexjLuXLcMK/IEF8csQcli1Fy1RKdXbftr7Ymo
% PMDPlN5B4qouAFhd7cxhMliautN5spi/LH4b3JkHrTxphOkdNqcQkmHL9jbjK2jWCTZ6bO5qm+T7
% LX+DpmDhWWbF0zMRA2jBXShp52cyTCSH0eWJZ2sODi4JFzBmmEkA2l+GKnVX1xQcbnDyWGaFUIl4
% AjRcFlJysYRBAWFaZVsuczeuETo+QOj3hEMbd00aeJo3+uDRCYv3cvq3SW8yKZakek6yDnXEB4wN
% f8wmeW2UMx0n1pn5CgOc4u42k2VMIpNrlupTCi56FGcKz5T4mslveMqBoN+xAwRVsYAF9AhV1a40
% tBVt16deS2jmAlnt4j4YNIrgMCbLX55yrW4IHV0/vOAD+loSXNEJOi/N427cN8f9tbXT5nZoDle8
% NsbVEs7/sKZ5qkUwWw41bBAN/CoLWp+dCl30x4eTRWEc5UwNho2AoHbwaLAod3GmK5LWLjGpSehO
% nw1EooLuABc+D4KjrB3k5pknT88Kwvs/yjAE1LcPrxiJW4rDhu4vRwfL/t6YANjsOMtnvvzGcXb1
% NX7F9Q5eNHoIFkazZBsB8N0GalGr0APC/pH5qK4OpPa7A5Z6uXiER8d1j2GV7FTbqWTK0NpY6AaP
% VicWU0ZblsK9hynAIN7EFcKeJIc+h87DFpPvVssPGzo5sfpPAG5YHIP+6iZEGLQPgDVyTiHoZQpX
% LoAh1KJhL7ew9gA3IJWR4UNlF+wh15rtEqiOjgFJjAWuCNtEPWOHTMBQBq2GAZ5BjRYCdBtcA1fX
% RcrcYQ1sQ0xLeDRRXlbrv/Q+njl7Ld39hzCdVt2/5j8eZoNKoFcL2LXz63Lt3Kw5et80niA8BdSF
% KXpRtgpolKO7LzDdkFPo0De3MzKFmWf6jym83z/MJtPJvyd3f5Dx5PfRI/nX/eMQ6A83j6PZPfky
% gzHoy0nZdWAmA9VYigcVHZQduby0bcydsRpq/gbldVHQiG6RxQQzB/Q+mHlu6O8fzDwuCB6cXmu+
% gkQQpNs2HIFydikZcrit2q4EN1noblzWPkdwt/6KUxZe3JwjN62OPeyyIvayR4A9Ip8t9aCnWXOV
% j1rvQRxGcBBpoklLyV/Ke+WdD5wa6BmrXKHf+PsgZesoZiTvh60crnrtHO54XnPeCVtw98ESLIs9
% 1vdO/KLGpaqntwPyGYrfx7DwBYTsrvcPpjkU1orplqa8dqkeX27gVHL9+4mGgXJa6ULeyh8m8KcN
% wO0IphlVKIsc+hcQmQFEYEMpf3HB8/9fUEsDBBQAAAAIANKpVV0OqXraMAEAANoBAAAiAAAAb3Jp
% Z2luYWxfY2hlY2tzL3JvdW5kMDEvUkVBRE1FLnR4dD1Ry3LbQAy7+yvwAbZmck2PbY5NZpz+AC1B
% ESf7UJZcW/37Unaml31hCQLg28qCnz33JK5X4rf4whzn0fCn9lYks/gzzrWXCU/D4V03NBqljQsC
% akqDBCYFnNRrU0kQM5rtpQNeJH6usrINzg1qMI8CSbVwOJz51WnOCUZ3LR/2jFRvR+Rg6/mIRT+W
% I7bHlmU7oidvMuDci2smcp2YoFM0U/+LXq5sOiun4fBL55ktgCBP8ykWjnsrr2sY/IFSHYKxho2a
% UgB3rlNInWvLUkbiwjIuWdrnLlUmcHcjFwsJoz+M//eKqAK3HQiL3h/BJM3qNuC1XplC4P7GzRnZ
% Jqxkizyvyhuk8a5oZPNvA2d2k0viHnnYNuRuHhcXLbCYTyjc+dp9POIxjkt3rWU4/ANQSwMEFAAA
% AAgA0qlVXZ37xMQsBAAACAgAACYAAABvcmlnaW5hbF9jaGVja3Mvcm91bmQwMS9TSEEyNTZTVU1T
% LnR4dG1VW3IbOAz7zyl6gi4liqR4mowoitvMNG2mSbvp7RfOx9pN1h8eP2QYBAGouLVRPn3vCopj
% 6T1zkHSqMGJaU23yWkZnzpY8yLRP12xNeK9Pn6jdf/3+z19P6+n8+PyUdac4syZ3mtl7Ntt9N3Hv
% PqLt4aOP1VqbUkdkCCWvcOPR8LlSD3sH+XJe76wG8aYtjWQeO37wvGsubUeora4qgI81WsnY5wQd
% OSlJOZ0i/oP8dX481O/PT7/vbExQyNn58KgCO6gwd/rWlGbiR3abHK2tSd3aPizSNuceEnIaIPv9
% 48mHn4+3s29SNl8cAyTX0pkQwkn7OVaTdouzhlZUqDRRDk5y6zIH1+H4iHoZn1UIam2XoM0+W0G1
% VhbYj1s7tKeOhCKeOyZoamj2KjbKDJZb1KsCHrnXOnufcWa+LQIesJW2Nvk4Eus0LHCLboNOPdK8
% 90HZM0QHUPn+y8PfX27mJwmXfmo0KDE9szdJUK+lzr6cq0RAq4/Jc8/ucy/rvLD/2ZfHe8zL9Hpy
% 9FW+h+4GZRePtcYIPu2kbUsROCmLnXunVFiIJoygMHVQb1fM6+xFQpOmrjIsIJ2bg8gKnTFsYjib
% RNOGMXIBunKiQ41FuTbssN9hPuz18vD92+eX15e7o3PQGEehBH6fhaxANh4bb1hcd07p3RSrxxSD
% MD6iVOvyyOqAHvev73TtNq3gOwMw1j+xCDZmi+lQrrXqwrMTdZhadox2ipRAZbVKGPwD6EVYyLYQ
% 4zJz6GWB5FguUYcBDroAsssAYztsGH5j6VZ7sYxhyho3oFdlx7Q+hds4KA04vs1uIpMGgoOfWdWG
% 2U8VEgoH4M1Frt4IDTSQkfegt9IqTI88NlJdTDolOzYYsFRg+FjRsk94t4hgWLDtyqcbRpiGpmBg
% y/3jer0NbMwMd+GI9IRzJd0iJzrkDF0RMCpLMZfbVmEclt1no9Yy9K0GbiHfygrkUEAbO8pq6AKN
% occF2ZKcqgYuCB0CsHzzzA0vUDmjJw1FcWV5FfW0hDcrUX4I6NC+eu0obvD3Qjfx4NF5uzrO4Z+S
% Da5GkdmBpbT+hLyVlDeCWtgXinSK7e2nYBkvNB+jGhPVje8RO1oogFy9n0xq1g95yKX/9f7n15cf
% 60ZUaBajlmnB8xuUj43Ya7inbCSqgbJcHIHiW5vJuC42RFUNT1/tA+hbCQbMGbChbgyes5AGXBzo
% GMXrWKSoftjCnNaimKW90Lo9ju9S0xvQq7AYF/cA0rk2oqW5caWsN5vhDgnXS4UbEl9H/aAaES3d
% HYdrE/o434P+UQTNbAaSOwxSeSz4e46CyVHQbZ6BgZxaR7y0d5w62lvseQRFxoR4refn8/z8eL69
% 3IiLysPfVxPydTA71nxwD15sXoY2NyrUazjyWnlJtOLOXBfvXsqznf+Bhbz/AlBLAwQUAAAACADS
% qVVdUgQOyAcBAABvAQAALQAAAG9yaWdpbmFsX2NoZWNrcy9yb3VuZDAxL1NUQVRVU19DT1JSRUNU
% SU9OLnR4dE2QsXLCQAxEe75CP2D7bBjIkC4z6ZgJQ5VWnGX7gLtzJJ0hfx+5S72r3X3qXLevWld1
% 7gifSfnXdeALMyWt8kwJBsYY0gjPoFPP+Ew1fGCoTqE6I98B+Tss3d4d6vbN7Q5LBxeKJmzrFsaw
% kMB2Xy3ESi/wuSQlphfG+WGKZrhmnUCUszX4nG7ktTBJDWc7KXxFDTmBTpSZIpSEw2AW6mv44jCG
% hA+Qco1BxHwCs90SL9S/W5pBrFaYcSZuUIREonGZKy+hX0MutqhvjS0lYrezAj9hGlfphEqiILmw
% p5VgbQA/kb8LMP2UwJZty5hwMKp6M6nOcmwa5FdY6sxjM2l8NP+fs/kDUEsDBBQAAAAIANKpVV3c
% jYnJQgEAAPABAAAvAAAAb3JpZ2luYWxfY2hlY2tzL3JvdW5kMDIvMDFfbG93L3ZlcmlmaWNhdGlv
% bi50eHRdkU1uwyAQhfecYnZJJCcyGP/Ekhdd9AbdVU1EzCRGtQEZXKWtcveOmyppKyENfOi9YR6P
% Z9VGOKHFUfWgJm1iDZ+LH7D3yoxhUYOUZQILPEe0GvUdE3xT/aSicXbfdti+fuOcp3SjzYht3Idp
% uF8V2fbCDg2H1g2eZIceAc9q8D2GGp6X1CYjLS2egFwlMJNS/Cecyz+I02b7S/fCfFMmtikgdOYY
% Ua9H5yKMGIyeqBPjUDABsmIZZFIwCTwVFcuvpYCiEnnBSpCCizRjFXBBEcjiZks5eGdshDgqG8w8
% f006yWQFOSPL2V+mKZXZEipS5qQUEqwbB9WbD9TgXf9u3WAo+sHpnJLn36FmNdDIdJYXaKRaqp1Y
% 89VOsIe+h2uW4FUIqDfw1GFAUCOCNupkXYimDQl1obd1ePtaTwEcN+wLUEsDBBQAAAAIANKpVV0u
% 5pt9cAUAAAoMAAAoAAAAb3JpZ2luYWxfY2hlY2tzL3JvdW5kMDIvMDFfbG93L3ZlcmlmeS5weX1W
% 227bOBB911dwWwQmZTqWZDtp7GWBDbB5LBxg3ww1oCQ6kleiGF1cx0H/fYfUzZduDbSRRsMzM4dz
% hvz8x7Qui2mQyKmQe6TeqziXM+vTp09/H3hYoYJXSS55Os3yqE55gXgdJdUKrY0jKisuI15EKE2C
% ghfvKJfp+y2strZFnqFtARiwvkRJpvKiQk+tAfESPTc+Ga/i7nOYZ4FlRWKLAlHIvE7TBH8jSwvB
% 75FtnrFLfPOyzQuUoURCfvJVYJd+G7utn/G95UoJGeFJWWf4GWtcnI1duiPEftzsfAOwGwAyQqbP
% 2oMQA1KIqi4kemyS2St8oKrFP7BnfGi8ki2SeYUOy9bddWz7wXwxqxJZ4eNJVkfGgxIsqz1zeuOP
% OEkFOt4oxpzlfszc1XE6Zar/3kLvT9MyyIdbWWcC9icvyKS1RELmWSKNrUkd/omwelF5+o4lVfSx
% zWet2XSIb2M59ppq1hvpj6E4ZduSemS13rj+ZHjveY8G2mYUVsO3ocaQtWxLGk1cQiPNtxy7k8i3
% ASjqHdebCKKFbZTwtLx1k7rY4zU9tNhHptPtcwhNDmIvilJEeE2WR3a0D+MzmGMDA437hkOakW6X
% woE4W+U/cHhKG5244HqTdSnAap2FXt7m4fwuCayz6CNqnOuE0mRb4SfYi5rGLSpnauL2uLvTxo5P
% yH1jTUJPlFMgE+90v06n8LjrffiY4cmbKcyza12OIjeK2L0TL0sBUjtDiglr62pT5Zb1yAYNuh5w
% r0q2uaeuS90ZdeHhgXoz6j3QmetbYV7LqmQfo1chgdr0RfGkKEdLh47EoQIliujEtOdpbebKSxiL
% 8N/GGCWF7lQQ7GD9aYkDz1QqILRvaW6U5kaVy54sOZA1p67nnnVjyUC6Enpct6Iiq0C/Qz9OPB9e
% BzclwoSnDMsbDBuhyUAw11DAmDt4gdoDo3cIiIFXhxgnbWkBlmEuq0TWYpCD4WVzSYLGatcgkZYC
% XfAGynB7DMFm43IcDMph16JePbFNqBtBDJ259vslNaino+HOKBIYmBppY4AmxCQ0Dr5C4SafX7jr
% DoJ5OgkGRtpWAk5roLNjTX4Vlx48TTF4heClQ/QpPg1Y4gB8VCJiG8e3UyHxE1l1NhgRbGJUVZtu
% Hj7MfNbbL4NuWh0qchLQZ6xbTLGikpb0uiAzeUwf6NlnijImByzmwe31on8xu1vxVMWc9dLukwL9
% 9n46C66z2DgU+pQqLU44lJStn1zHpQ8P99QgNf+PtTT9oZ317zWvGMwqTeeFhC+imcOFy94bvMhY
% /zUjoXlqYl2va3nQsZjGaKjiFAxUv597d01+JeuzNtaH5Z/Mc87L6TlpJPxAzj9fbgoH4vWRvtP7
% cjUrlc31RJxw+zc+/Rn/qxKuh9BlDZ3Uu8HU3TS6ZqIC4FWhz+NRc4Vqpd3cnJYj2sTqnQLm6nuP
% AuYCuAl0uODYPYIru19J6LL/l/5sPjtV/qpudPGsB9Udofd9OICiAIXKGHpVRJMih/lViDKJah2U
% mCkbnzIGe9KsjenZ0XUFCTSoHCyogqVlojuhA2za/o7Ov9DZHG4Mcwck4Hhf/BaaU9PPZochWWJ9
% Rg1520QmlUB5UFZF3VwdqxxFRa5UIl+R+srul5C+EhyKQbqYEvEKKba4tdaLM7q8OV0YvhaaMNdb
% nBC28K3LUbW4GFULYhUwnj6SZTtXFs1cSaj5LtpLhcDgqRul9/rZQevl7MNdzuls6dHFcv5zYHAB
% DHpzOEyKjKfJEWrROetrCXQOQC2gH2A9HbE5x/y7B9Ppuzfqt+CvNEVNwyKlo0W36J9YwCDnhUBR
% wl9lXlZJWFJzXFWx6LtSAWfbW0D6D1BLAwQUAAAACADSqVVdyj53HqcDAACEBgAAMgAAAG9yaWdp
% bmFsX2NoZWNrcy9yb3VuZDAyLzAyX21lZGl1bS92ZXJpZmljYXRpb24udHh0bVRdbxs3EHz3r9hH
% G5DOulNi1SpaIG0ctE3Quk0e+lSA4q3uGPNIluRJvvz6zvIk20jzIODEj9nZmeH+5UfXUkPscpxo
% 1dCBo9kbrbLx7uJnPwzKtVsKU+69W8+7UxUm+pF0z/ohVZ8TDt49mkwpqzymLa0WpKwllRJHgUkU
% 5Luly2ZFf+jsdxypWTU3V9UFbiqdT1jbiyXdK8sTZT9GpwawosjdaFU0ebrWoOPd0rHp+p2PgI04
% kzkmAksyLnMH6Pfv//n0g/t1+VsFvLdfTE6eD9c/KbNcfsDvXsUHOjS0vlmincyPL6ptibE20Wlj
% UFn3xnXUMjQx7PREwSeTzYEXFJXruKbcRz92fV1LubvHXo1J9il6j8YmbY1eWuMe6BdRRY8ZErkN
% +ZDNMA41Xdavy+rVglxdn9cbqFWvytrteW1Nl6/Wm9ffXUmlZkOJuYWsWoXA7bLlLjITt7DCOJjh
% NEOYDIRms6hv60Vzs96S9cfz0ZSjR3Mh+oB+pwXatCyOiX9ncZ53o2iTlCU+eQb5NdQ3X0paFsUE
% wMMCjVjlkoLZ+lmaACnMzOckDLVmJlpX1c1T+S1tbhqUG/wB7amo0wIdadNKHmbudVNCJkjC00ir
% czv/L5XUECyAiv476/UDx+0Z53YxG/UNq+uK3o248qSKBVkQeZn878mJy1aZARWKT8pWF28k/8hs
% j3hKlxLSwTjBR1w49d62tLfeI7pjAvdzdhFxpNoAGT0LNA5BXbnnLUeRqiJB76IK/RPrgnLyRbWf
% lS5RTZznl6HGTtItMEHlPuHd3aPrvVhq+Li9uFO6BzP2EJ16hVu0l9aP6CGzkxDgtHEtB3bFBvzF
% eZREOIyAp2oW+IxiUjmQQFu62hmnso9GiQVKP9DOP1b0EZTkdRMnUS7zzLefgNqxh3oR8mVlLKUx
% wFSIfISsMyGU/Iinq/k0el4StBOqxsgIKJSOfsAbw18wPzQVYbCFMZfYSp4PEi3AgLNVO9jdzuKn
% cQdeeSwYMnCkqoj39h39+WYr8zEYiVY+IoHwQBTejca2VeoXJMf8mFFJdCqJBC12aoc7CJeR4nJI
% FN+jE7jTcTrviDkyNmWRXn3VXJSvKM+jTL4USqcV/e4JryYW9yAxI0aY8Mi0w9nIe9wpcwHdfFCf
% +G86qii5RNV7Kf4MPOvc7rMPYShltHdZEpZ65vxcdEuYE/ma/x3VPO7lJp7JoGJn5hzjhZgQRB3U
% HUwqQnV2CpLF/wBQSwMEFAAAAAgA0qlVXdgLkjuUCAAAmBUAACsAAABvcmlnaW5hbF9jaGVja3Mv
% cm91bmQwMi8wMl9tZWRpdW0vdmVyaWZ5LnB5jVhrb+O2Ev2eX8HdYivJoRM/0hRrLRdIevfiAhsU
% bbLIF0MQaIlJGEuUq4ebB/Lf78xQlOXHbmsgsR7DeR6eGfqnd6dNVZ4utDlVZs1Wz/VDYaZH79+/
% //Ikk5olRQ7vZF2UWmZMNqmuQ1bV0qSyTNkfJM4yvShl+cwKkz2fwNIjna+KsmaPVWE4y2X9cHRX
% FjnTtSrrosgq1go47bowlRUpQTNK2vfXdHd0lKo7tpKZevZXweyIwaf6S7zqgf6wYndFyTTTBpfe
% K3/MV8EbiZSqbkrD5q+PJPO4kVkFTN8x/3GoA1AAj6u/3nb1rILI2gX3kwdt7v0Lnqm7mpf6/qFu
% vZBVpe+NSsWrNYnysrn3G14pZVoh/KDyNVmCsFTqX8yb6Oe+JvcBv6wcrJ9BfkytTaO2RPDViUxT
% fx3sLzVFjcudYwzsokPufr6Odl1zn56IaEKXvW9lz3r77L8yq1Q/xVWT+y7s2g8CCrfphYuJCwKb
% T/jTia4hnU+tF1fiYv4Uhexa4HKb/kwZ/yIIguHV8PXJJvcn9mWtAGWuIKyW5b2icK+ZriDOuoB3
% OpFZ9sykYYoQXCnIYsqMgmwvivKk7zhauQqGvRJf8evWT8J6vJCVAkeso0ZYt0ImhW+Ox8Hp6VnI
% FnA9xOuJwwSgnJkPZ0JMwY2UgT8Uz00gxIJyc4NOX9jyEe44obO/Hdo8mIBPesVqlfskfjHXUcDe
% CeZre/sYBbuS5DDI/UxvhZDDMYl8FfP5CCGjhXhkCgrKxnjbKbbPhuPdvWOCaHerwKNNJP3He34j
% UL6i10KMtvbG4+FlB5bOl9HgK8QC37Ry2V8JerEW23FBbWxFVUqwSx5QumrNXIo5oq6yoK0ofMBi
% qu7FfBQNbMG78Bq+5qWiHdqq2YtxbTPYRFR7uruEu04Mb05KlRdr1d/B4DIonl3C9qPN3WxegS+w
% 5liMQ7xa41UfxJc8l08+vOIQpGyyWoy6rZYpBFOc6RxCd5ndgLevBpUsd1MqQQ5dmwzk8Gyw/CxG
% FJXf3geDweSzMK25XD+pdMfWzub4ka1Fa2vRt7QY+IuNpbOBGThrlcxXmYrbMmBd5YqX5p7nRapa
% 61/EHEhp3SOkLXbZ0DLWKwph9Un10NzdZcr/ErSM3sdB2FoT86gPCdTwZdYvpS3ZJ3CJorB1w9tt
% bLfqTuRqpUzqo6/kPqrA73diYjGMnmFrtO5PgBi39PwDQlozNm9lUQCtZUWyVGXHbNTvH2RT1chW
% UEmjc+j3/8ObpIHdp18UMGxNq9mIs1VZrJGDJbsD2VqxYlXrvMmp/W9z5ZXomt4oCojm7f2G7WHn
% DvHt8HX0FiCjWmK2CFJVLRYh+1sbo0rxe2FUyBKg5arlkG7fugEAgbTJMyw+Hg+rLbq52WPbK17t
% 8M7NDbWjm+1M4+rbvdXXvD7QUW9vScMtRKRSQg0VeNSD480Ggrdd7ycw7mmDypYKlVhKh/nh9rav
% KTqeI0du6Gfzqgf020OKs1qKlkPQBMAc0+tAtDNjkPgnrMqMSkP3YVsdirPfBu1zbM84mWDxLOty
% a7WlZCsV9Bfa9wJNtFvIjg2XkL/PlhsuATFCIG76YH/1ijJVpTcz3EO0erMR9wjQTR4DlmOINLaR
% ehQB9+BpFcO0AFIq9WYUPfcce1bezPr3ZjdQUjQGBlmUBwKanvtt6f8U89cxP+O/vPHXafs94mM+
% pe8z+D+hZxN4Yj2mXSXmv/IpH4Mc/8inUZhksAtViRwTLnsQf6F5ClfMnIjjjUxXbict+fL4Begh
% XB6LF1p7YRucxUC82SfLftfmv+ELZZpclbJWvrOwM8Kuioo325K/7QAfoXfSrFJ6N4cFx+NZtL+H
% qOH/Cf181l/gzHaTTFt1SEZ/bCQNT/uhtNDBecstDLahURVNmUDdvUuph1d6+Icsl2w9Ydcqx6vp
% yZiz/7zAugKaPMyNVV02CYLA4w5WS0QGKlcG/gAcztSbOy9o4yABtC3sGcafjCbn49FkNJoEYakq
% aNOVcP7ESVGWytqZ7cOLezTJxhVQcuaoGwzPI+7VAGC6fOuKaTAxACqA1PhjtKnOhbBHKBOEW7Ot
% c2d+2EzkULbTOIIdg+MJmPw45pPz6b8xmmZibzwJ80xsDxFb6MOOiF2ezPkjnmYBB7rP8WuCX9Hs
% MM4sSqYHKNr18x8PE+EOXbkJcp9Ke0C9PABUmJTDrPhbzJ/23yC3IrVfwlko+NSNTdGhg1pHkJ8E
% OPo9CdoHFaiDaXtjD+wHh3gdIwXBmTu/FMi7z6DgB9LjWc9UZwun8M7UfqDB9/zF9vXifCW6A+Of
% cVilWCEbPl1Ae99T0WHYbogOtK+eoU6A/nozQpCHDQAr2XUB0gptgMqaxrJMYEthLVyhuQdEUJj7
% mDanNzvoKfcg3lanE0QlmIW34Ki/G2Cv2AfdPLk3GbTDDAhuT6mUfGAYfz1sgg+46fxmuKYrQM15
% dKi9buG1dzB0z4X49XzSdlnquuNJd3Dtw3kLRxsHHce2FcjwWB53m7slMi8SVAv02LbWLtNg/EBR
% xpMu63Dsh2M9cm27mmZOYstWFUuek0wnrDGpRiZVEIrGn6kSPNfWD2XR3D+wc8/9knBDPsH0upAZ
% CqWbUbcbckMGpM7upM6aUtEPDO4nhQT9gbRAyzppT1TTSbgz5x4YcylZvUGX5id4SBcYvOgGJExz
% Tb+5dan+ZTTqMZidTunEQrHABCuh7X/bfXyNj492+e6HkygNl996J9Z/YL/vDoVdWO4AckCbk4H+
% 4HitP+htQ8vGBO2BINb2ol1k1bLb0dSyum793PVqUr414Tk3EHOybuAJdMNVhsKsASop+0XnLTva
% Q4/MdP3MkkzqvMXXqtSm9vHXz5O0yVeV3wbAtUmVqQWc344gaXFsZK7iWAgvjnFsiGOgKBofjv4P
% UEsDBBQAAAAIANKpVV3mtaNpbgEAAMECAAAwAAAAb3JpZ2luYWxfY2hlY2tzL3JvdW5kMDIvMDNf
% aGlnaC92ZXJpZmljYXRpb24udHh0nZLBTsMwEETv/Yr9AFJUmqAQBFLVVqhSBZXSA1fjbNJFjm1s
% h7R8Pes2lcoNuDly3szOrJd7IQNUJBptfCDp70FooQ58BOuMqcF31ipCD2GHIJRKjKvQxS/jsB2P
% 5EEqhMZRBe1DWsAky0AZKVRSo/D0xpd4NOmRml2AT5TBOA+bWVn+hLMCsun0n/AtO+d5+gfatFbh
% Hk4quLdCezIaqEIdKMTEnGc8visgzQcdKTwOuNGYWEM6gEXXCs3QFfiPTjhuSVfg0FPVCXWhV8AU
% zq6bcgHvjA9qnnTTKeHA7tgCgiOhGx4r/oIVBDNE4HOWHCcu4i5gcpPzmkhLsmzlu7dWMCvPQ4b+
% AjxbpycB/ysFbTicoi/mnemHgFUijXOoRIiF+a4F0rFBRZJiSTJ0PjYW/Xin1yk8Tk5yLe1ZiTOx
% jqJwgJ7Cjp9fXaPjms5sfHPGhWGG2XoNy9fZfAuL1ezp+aXcrubl8Wq5GH0DUEsDBBQAAAAIANKp
% VV0y2HzEQQkAAEoXAAApAAAAb3JpZ2luYWxfY2hlY2tzL3JvdW5kMDIvMDNfaGlnaC92ZXJpZnku
% cHmFWG1v2zgS/p5fwTboRlLo1HaS27u4NM52m6LAtltc+iGoYRiURdty9RZSdpP09r/fDElJlO3s
% BW0ikvPG4cwzQ56+ertV8m0YZ29FtiPFU7nOs8uT169ff3jki5JEMV9luSrjBVmsxeKHIstcknIt
% yH/6ZAEUW4WjXIp0QLK8JHGpCE+STi4jIUkh83x5cfJViyWq5FnEZUSSOJRcPpE8S54uyJecLJOc
% l3G26hR5nJVkkacFl7HKMxAmBdkqEV2gUSdLmadkKUFxjItxWuSyJLd2gnBFbg1NXApZ5nlS04Ap
% 0XZRmtWUl2t3wcrdZosWTyK3c9jkWpycnJKPfKtUzDMiOeriCcm2aSikIlIUUiiRlSJCCwoeS3VB
% PsF4ZZZTHmcktuOLk+/M69KuPyB/fvnAvB58nkRiSXgUeZyG/g1wlFsJu5l2Z+ch/KJ82sOv3kwT
% ptvkkDBAwg4SBkhI7Rzw2bmuYV7k2cbjbV6q+fRylst0fznQhmgxNd2qQKoTAj+SwT70F8bGDjZK
% +I1kaKWkO9/QGGnyRDMXQnojy52xRGQw0oN/1x73vuSZMJPIIcXCS7n6YZnwJ2Y4cRHG5XyRb7PS
% 85ulJayy7MZqrczDH8W+199o7gbNlTxbCS9zhFspENFa7XC4+a3nEwhfMprGs+lm9op9b1Mb4XiG
% iuLWLR2tLP9v7927je83RlrjVMs/QGyj4cE6ya7gaDr19PE8+xAy2vpnbX3+c6ZH8IHj0cyvDjvC
% 4Hal6CFaBv+6hituPGBOwjcGFDLOFnHBYSv0rpYwndqNNd67c+TcGb3PQuboz5rpuyaZu86eHZnS
% zPkWcneOnhxB+NzUkXVo6M1LR1mtW1v1sVQHgoezg2+qM2EHU9WORbQS85XkqZdR/FZUPJaSs65V
% NGJ2X41JdEP1GWjyxpwdm36fBdkA9WB2DFANex60t3ZsZ+A0Yg+osRqM9bQpcO6tfBqZfAI4W6wB
% PpWXVjlZOT7eV5D6GNgmqOPfeuYUcNiQQKSm/n4qAKh7TToAp84GM/K8+Lznv0l9TJIDffZYP3uP
% NFV1TKht6ulwfITtOQFkBNzhN1Cb/S2eFgmeTBx5dn9lXvKEdWsPpo3CK/q7ExipYq53BhoqLGPF
% DDiukN+WCM/IufYhdwvBS5YeAgPPnjxkA9vP9V/rgdnw6ogDbiDUoLxtRUvMI5veejt65Te4iaJm
% A1L0mXFN+xQUQ6ftAq/X2TlMe1Q82mg69GvwWBt2xKw2m1JQMYkaQthtjq14MP/2FvmCIB2yoh8E
% fXJKhO4TRp87nY+fScnlSpQt7lM0iAzZlXqQpVf0fUp2QsbLGIpl+ETUwxYqfbYCqM0ysYLSuhNE
% xZFQx2wAUQH8H7LeP4Ki36JY8WKe5XNoOEpmQ63z2Zvugp3jqZme7u/zWumOiCHr6vh2pmC7sIt9
% Vh1P56xXT+rIPGd6fqBBtPSWZzqCCUYwSdmv9K8b8ktT/EWSfMGTzlJwFYdAYtz5U8SrdQmOgn4E
% Ooivo7u7s1biay1ucoDHIMTKWKgqRdxIP0iRXvflHGnlhhJwTjXjP/cyYccTxabeLQIAUvpvLink
% DYz7QTPzu38EFWYtQRZa92DngOmw5FZYiSBLNsxG+6Aq0wwthM/BCAZIZ0DfTPr76Wira1NakRDy
% 8ZlBu4PfR1PGFGrGPBt25/3gGZuq7lHqh4r8eHACs7XiZVUgfMisnJeDsQ49aKYT8WjiBOKr4JnC
% XrmJGAjJq4uLfzUxaYJwwZWoY0+H2ga6c6hcY3sOp3A9UNi2b2O11mESgZ5YkYSrErcU5uV6gPW8
% 2JbE3AhGnRSAUHUwo6j+NXZmtFhuO8JBqD/G/mBiooOfh52eP9Apzju9Y+WTv9wTwNLEBoUNDk0Z
% P7IEdmFBn1OjxD+foprZMR3hyzpCrQORd4Z/NjNTxPfm6Njob2X0xGZzkqutFHNz23oplzOqs3nq
% XUL5oJDPkHTeNb30Z41pI+b2M1AHqEmOjMKVw4M9BkF8mJjQhNmup9fE1tgRlTqiUvj7pk97R/Pb
% iuk3YiasDqDBqMvcBtMqR88Pxu7S2C6luFRLKjgtQlihMmRTTIR7v2pn79EI1EFHXTruOjjzwOlD
% SBWnCpgeDNf/YYk4jUIIl05vpv938eyawcmR9JwgEhQ8kOG55AFY6cCACwFI98ADFZ4rHjyEkPcR
% D6IwUDhXc6x7dN1nto3v+tR8jR2ZIYRGxNa9YN0PPBBhVIMgrb0WK3HavXxAjd1CC1UxAkUYgN0d
% IJfcP9d/YCLs6CX/YBPGkokfmM3g5dOYwirZQ6fNau/8HdOkJ3+LW3D/s88BwJTyDNCKmnZB6NJc
% qXGArEGwCvS+3r3XqNXGMehl5270JTwUyd59dHBQP/d7ZPe+eN+K5sOuO9vruv1BwUwEHsQQ3LCZ
% 7T4KjfTG0/fHw+genFkRBIUbk0ed+ktvFdoPcEHtrNp0aMpD6ARkvBDHmg5NbjwIFQIdXAMU8BiM
% vvSdSw7MAlrgAwf1ugATiDw93+Ie1CMHVq4dWLn+W4S6rgSMDJygEpBl5vZP9kxheUq4JMUayhna
% CUKgDiIrlKwyt70WfF93dIm0e55cOdZdOdZdoXUAe4eWXe1jZ9sabe3kik4Ass/Kn47mKlqvjAXK
% mnBK7kqZQ2+8yKUUCTdvXuBe7I/NI1wPhPAfHYcADmsFfy60hG9sOoXOaL9MXfoH14FLi3pKazSe
% /db4VXb1neLWPEAYIgg2LGK0Gel/gR1vcBXGrRblyD3dkPs6QTavmIU2G+OyO3QD+AzV8yR+Bqfh
% a4eBg6i1f7BT38gfiyRexGX1Smkj9uaMyi49G/bc+D4l7+PlUkiR1eRqW+ATILg7055WPBV4cSnh
% mJQoL9pYgQhl2wgBREpEXpX3WNAccDC7pchgo3gMZ+Q+GFXNweHz1OxIsdYy0vhRRBpPRlBarev0
% pAMmZoyA0nU9XE3D9aZCkpEf2Dpj0F13nPpr7FcIY89Dc4NBAg4iicsn8jMu19AWvuTNCoXhrOfz
% DHw6nzN2Np/jC+l8fnbjyt5/gVYD2AdPnvAxWj8sa7EJNrB4Qs2rs32QvjirgaZ+PRgc3pYG+y3X
% oAE315zRH3+QD/ejyTfy/tPo45c/7759mtzpDX14D5r+B1BLAwQUAAAACADSqVVd3k4GF+sBAABX
% BAAAMQAAAG9yaWdpbmFsX2NoZWNrcy9yb3VuZDAyLzA0X3hoaWdoL3ZlcmlmaWNhdGlvbi50eHSN
% k81uE1EMhfd5iruEIi7+t2+lLlggxKZFrdiCpmlajRgyIZlK9O3xtCWoYRGWI597Ph/b8/n91dVp
% 6YahrH51y6lsu6kf191QuvVN2T38uB6HflmWq+3U3/bLblrtyu24LZfjfdapfFhP24cCUhfnF+dv
% Lz99vLi8+HKVXpthTKsxizd9d7ced1O/3JVxPTycLrpvJ+WsaGVCB0NxIHEjN2wQIWgu7qUsU4S1
% BVBDBhAmQmoKruEtlFVbKd+7zaZLoVVzDPMgTSNAac1S7RRkYbzoZqCUcj17ZjkUMbQpkzRGi6y8
% 2/3cTq+612cFqrA5kaB6ZFceSepyGptpzBxlu9r1N/fdcDKMd/ngK838FBqioM8sq/DMQjKCOSYo
% mrCY2iFLmzshOGevNnd2DOYVnBqIAs2w2MPEkTFdGmUwChHVA5gSYDTNabMxGP0HS3OcZtoeWQh7
% mIeGUsxzcgLO7IdTVJPEkAmRSpgiHaeFuqQT6RPtD65V5kBoLMEm2Xow0AGu/SM5houqjK7u8BwO
% 9kBkqS7NlD2Cm+b3y3hYWV4K9Hg6t9a0RZNHXIXVm3iiWS6zCjG4sKfEmA9os+KF4BitVZQ8bc7t
% 7Wn0HC4EzIjzN6l5n8yWV3PA+yt5UhzdHeaxpzxP3Re/AVBLAwQUAAAACADSqVVdmnGR2sIJAAAf
% FAAAKgAAAG9yaWdpbmFsX2NoZWNrcy9yb3VuZDAyLzA0X3hoaWdoL3ZlcmlmeS5weY1YbXObSBL+
% zq+YSz4I0IBgkJAch1Rlz1tXu5tK7GS955RLmxqjESIRLwHkSPfr7+kBDPZ5r9ZlwUxPv093Tw8v
% /zE71NXsLs1nKr9n5anZFXlgvHjx4uejjBtWFvtTXmSp3LNYVU26TWPZqPqcFWWTFjnA+SFTVRrX
% TFaKvf/wOyuroti64GCkWVlUDatPWXlismZ1D5FVUsqqVsaRS77ju5If+H1Uu8C8K/a1OTkyyXZs
% V7IDu59YRhMdp77xRyTsnSNsGv8WmcJu7GaK59S3ZibGttlOLGPngdlH2WpoBlxY9tEZQc6471nG
% u+hom+F0YR9BL2zzOA0sehKH7/viR2S+c3xaaSx6+1qgBVx/BjGzVgGLRiFGdgCyj1Ho2Qfbnk/F
% /IwGwVT4cwwOzgqv6SI0ZF0rcoobyzxWe9MkUc7Os9z6cFebR36AAs5HSPFs8zCdQ6WQBr5l22Lq
% L9rxFAZYVhR5PT+wt4Uzt0ncInztDXJqOH2fbk/Qd+e17nPMR+6Y88VjZiOi2t2k2y1e+yIx/enR
% ct5xPI4wbWZqZ0Fu5ztbtFzeR7dYn9vmDvj8qBHMObYAxuymgX2ExznUORIIP5/AR7uFYEcJsja+
% l3F067SMWixO2mtuS00bEjagK4IG01AD/bXxkl3efl3fZmuWIizrWOWbNE9YXKgt4jdVecPuVdwU
% FUtzRNgWg2anHq0XW3b/Z+Yal9GtwW5vF77Hw8XZkvsLQU+xDFd8vjgD1PPWnN36i3Dh8XlIoPlc
% zD3uL0XocYz0+srzAPIFngt6YL5ec2ItlgQSPvdX8+UZF+FqBQHLlbfki0VAYlsBYgUBi3AO9LNQ
% YGEZ+BAmFnMIC4JWzFyLWYb0PJuTHBJGYjtpoAOjBeJ/uRLcD1YLwZdkjPDARnTGeB4ZA6MhZU78
% Fwta9lfLkUnaiOeMWhtl8UNVdXQbcJ9763NWx3KvMF+QvhwGGeTzrzynDVC6gqCqmO+tVwbboJrc
% A4QMVsdS5huzj8EcgTd9GO8QNKVl/+HkFDi70kHwGewhgDvaDr1nilJjOQgshIeOVPZT1C+16YdS
% 5E9/s3dgRfUuqg+ZeW/bmU2D2v0lb1SiKjO2KLm/6dD5xuPHZnTBZ1l6OaPFSuaJMoNHGnb5/1Mn
% mO+86f3jItBlfetNMLVJpZnZepPmQBBt/gkqLl0S9yLkfm/GbyJPq1EVP0gR0k3Ptc4ADhpptT36
% f+MZV5FJuSh1QZTW+UXUpqZEhknjJur8eoUtsS+cK7ubX2B+/jm6si+Mpztxgxz/7JjOH+CBZH1/
% K9YYSQz8tX556yc1qCfsfYKqZ8q+4EAVVIRnCT4PBE5HEDzPuY+mz1TShP3IqL9TDzsfzbSPLGJy
% MbvShC/Zz/mmLFKUkt/laQ+Hb9L7tEa9ZZDM8iLfIJByBEx8YukGJSdtUlW7Rsm/84pv1L6RXNaN
% rHji80TwJODXHr/2+bXgp/FRWbLvDNyJgGkClvgsESwJ2LXHrn12LdgJZ+hfZMYNtB75C7WWPPas
% u0buGVEEI4q3ZZVmyo86LnaJmvywZd/73Qvs6n94d5TOWC2QP9kT8EBo2dVIvpa7TaJrz058+zQ1
% aSCm1z6mln2yTzPRwnAS2wTFmujWcISFRpJEmjARHXYS9EtQEXtCztYDc5vMkoSfOOqt5VYqK+7V
% B9PqTbCeixZN6OqTxTyRu679mfir0HqMLDorsece6MQs7CLrXZHIKm12WRqzTEkqoqhKFXqxLpBO
% zESKZwBs0rpJc93H1Qiwe4Uy1YZlbeFw41f+OJQu2ZU/4Q+40e/VQT1nFRqgS3REV+i72jy4AgC5
% RiCHhh34klD6dNDtZFzkCFFIZ5U61GrDtlWRsY/FATnhcxiyUSU0hBX7E1CqQ852qlKuUY0aul4j
% oYtyZfrUS3yl0Apsuxu2tffrUHvnlvWmMsMzHEk4tR6sGoEccAr4fEV4wvdxTA1ovaDczlvO+cAZ
% JzfkTWkZTdnrygzQcA6kgdNDUGHAeXWmOb95iICl3SNMfX+E8dpfDZZq6UErBZ0WvcSSWlHSmp/p
% yg/SBQ9CUmGpG9yX7J+jlqbeF6VidZrkba8+tPbOQ3SQacc3ke8aj84QbzgyavcSdOZwOLfNl6Ob
% Lzh+fIBZ/GC54PBFh3RtWqTSW7oONEWD4B2iId6p+BvCQW4bVTFigTA+NNS1UWe2Uds0pwmuFtgz
% tdeR4Box/zGO35j9mPAn0ftOZncbGfmzGO6P0bL609jG8d03tDSxjE2H1lWbdsbjzpoxLgcX4tNT
% PmTHVlJHaXacHBPLDl1M8LZ7bJ0Lv0bxbMzxrw+ZX/+OBl1+fVK4UaH7YPKwSXX/So6Do5Sz1emF
% Aos93KZHtXFkFLbXNI6TCJvM+rSLT65xFUZXrVTJQ5z8YXQxmt6E0c0wNTTXL3vqUHrn3YRUq83R
% /UIgZFEXcFca6rlGusDx8OgyZTlXYQ+l646Y6YtVJ4YiNnrSiD3pxJ60YmiIQ46OekXt+AqJjnYZ
% na03n1MrzUWwHm9hfyML9HXOHtn2uC8btPn/TcJFONN2Tm9oQNejEbrOrIfUeja3bsJB7Oq5bIIi
% eYO6/fbTp1fEkCldZavO8brVaLODsm10f9dC27Ir0Kk01Yl5cxddgqHv5VXU39Ddt1UCZ+bNpYab
% 1nmL4MrN5ovs1syJ43SfAeR+wqECxEeTGvmgvjRIw4l1Dtw66mj1i6hhhJFu6XNA7T4wwB2g+0yQ
% lZlsdvTlIDNY5maluylr3LCNttGJzMz9j2qkKSwnc9Vhj32amYEzQIXdw9HOx1HmooxsqIaYe52l
% LH7V14NslGGOzjAcwAb7JstSRpi3SQxwTDlNb2S3wbodeP/hvfPxl399+Pjh+hM2oUTDJ2H9CQew
% TPIC1SyuWZHvT68mA5H8YrNowjMXVbAy22Yv8Cw+QaUd4HHcAbUqowU9pzWLLk1bFqMnIP+Zkq5R
% GvRv86jHLImyZxrWc5aUsO0468D0lYPeRFGp5lDlLIFjNmkis0zqy/4sKR3NagABe19EtD1bc+K7
% Xv/nY9d3aZTJo9ku4qoegF39vWqgI9H92KV7BS13qfU68l7tUjsSAFNwfhkOWIGjUVuRpZvI3BdT
% oM8Ij6VkIqA4TV+RqHRDULWv1StSSU87Q0Z0ZaXuI9/QYqQkObeThTuf8EmI3wo/39OP7jm8+7ES
% 3mRNKknZGS5xR7qLhi3AtMvzuzck71wLvQNJv/mjnZd8RRt8N4DuqB0g0Kzz1mhp1rtQWriRE5Yc
% jtRK1enmIPc27ba0/hQD4SNKJ46nOoBmbWAA1OcAhihV1EBYxn8BUEsDBBQAAAAIANKpVV0rETd9
% zAEAAAwDAAAvAAAAb3JpZ2luYWxfY2hlY2tzL3JvdW5kMDIvMDVfbWF4L3ZlcmlmaWNhdGlvbi50
% eHRtktGO0zAQRd/zFfMIUrM0bdKyK/VhpQKvhX4AmrWn6YBjZ8f20v4940RFFSBFiuKcGd97Z/Yc
% R4dXsvD8/d1m0b5/qgCCWJLd4wIOu2a5Xm+7BXzdtV3btqt2vYDzrl2vlvo+1Hq8bJp1t+22cx17
% lCsYx6+Z6jeSRBeIrxmFbG2CCDlMHDxI+AUxD7um+dAsq+rTBU0Cy9j7EBMbMCH7FIuY/QG0T8wR
% vl0jCdCEYi9EAxUG2m2rnNFjhc6kV3EqKtjHhN6QIh8LEIbRqZwvggMkYfS9o3uoWU3Y1eixCwZd
% fTjuAb0FNeCTNk38RvWA0rMHcybzU6u67lakBlUpsKUCX2FElhtgWcikeiQZcpojMBJirG9tmo1S
% 7C2d2HMiCC8xSTYTqV35xGSVUmjgi47rBV2RralOev9vu9wcPNVjYJ9qzfsfcdO14xkj1beOc0B3
% XVYl3lHYGx7R6dheBtT89N8cwiRssylUZN9nhwImpz/jHx2nvzyot+B7iOROmpr/kXtMZRakSSvN
% 0zgeu7aqDs/H4xOgc7fBR12CkkpUIcYQWbIP1ec5tDlM0G2726W4AB+0UC2EcAJ90pkgey5rpX70
% KwgND9VvUEsDBBQAAAAIANKpVV280K725A4AAMstAAAoAAAAb3JpZ2luYWxfY2hlY2tzL3JvdW5k
% MDIvMDVfbWF4L3ZlcmlmeS5webUaaXPbNva7fgXqzHZIh3IkOUebRpnpbOvdfFm7TT9Zq2hAEpJo
% 85B5OJK9+e/7DoAEDznJ7GwmiQji4eFdeBf47IcXVZG/8KP0hUrvxe5QbrP0fHRycvL7XgalCCO5
% SbOijAIRbFVwW4h1los/syoNxcwTv6dlfhCTV2ej0RUtFUUp01DmoYgjP5cwmaXx4Uz8PUt2sdqL
% IpCxzAshcyV2MoKnbC2itFQbhc/56CKHbaMsLX4RaSbWcSbLKN2MdxkAwaqo3CYKqQEq0ipReQQI
% xdXH30SpipLxVoUKz0Z/bZWQqYwPCL3Ls2ztAcZSlNsImIjSqFS0RhSq9BDgXhUwqfBflqvkDIUw
% Gq3zLBFrQ5SIkl2Wl8JQKWQhLjRMlQZllsU1TK7CKlA8CZvlrckgS0DmkpDC7ipPqpJHvCCDV7IE
% JjV8UsU8kYN4ESG//pNGo9G1mAtn4omJO7rExyk9ftBvp+4oTzcwYHBnNpm9nk5mABGAIkFqc/H4
% ZTQahWotShnHByeVifKETHB6PnXfjgT8YeAFzi1hCQ/PNqrU4BNXPNdrNDIZho70hK8R5KqscpDZ
% YrIESB9+YI/FlAfTpV6Uqo0j2yvGkmDHsoEqKn8INW+IKHzXNSxFiSqO0XFKdIgxE3JKhHjNFBFX
% T03M7kGW3nSJ7NOYZnkyGwITp6diViOGZ70ArBBspnT2RXsNm5IDRuCJfUEK1XR8ZQWxzmsuzRo/
% lumtk7bhF4trOtkrOIxoYxsFEMuBVzWlURpEOxmjXKM0jAJVdDHKRbRc3DCSG0Si4fhNZL/RWEMV
% q1IRyjay1naLmwYl0xWrFMTsimgNb3+Yi2hpmC3uKnAJYVcJi4XDytm7ZLeIb0/4ss9MHjzgWBrS
% tjUKWRRwmAWcE4c4hL+gxflcTBq+WmS5bU5YYWbpZHl0lZY0+IZVuKu3T+HoMQgNwx2MF9don85U
% vHsnUvMeUc/FJY1wi0QWt80uGpZx4p8IgBHkzI/KFZ1tx20m1zg/h80BEaBGQOL5ukHALiIFd12p
% +mVHUfaGGi/6ZCLtR83AjQt+OxTGen7obUJS2JeaXvGfel0PCigFQBQDOgY98rRDMGx4Zqu2omB6
% XJ9kUkJ+AMUf0wNyQkbQ0KoRsQbKDDwrLLmmUVElRa23lF7tcnUPrya1vmxlwZHv6muTywPA34pP
% wrkV79+jU6jVsMVlMEtAnwh3PXmDvpsASNXAw6bcOi64wGkNs8uKqIzuEYUPYcshPD/qZW5Lu9Fx
% 7SKXIFotfj2qxY0iq/dRcaHY87d0QZID2SPJ5uQgnpZlOhAgiFPbdF3xNzFrk6MR4S74aO2gdYNU
% 0rMn2gBaN7iHbSIEa1mITBVszQ5R7RT8B9HzQgJrWjB4VpvzbOzGgm4I1l5mhzbVsr6GaorUJ79d
% iXvInapC/IkQQlHeJje5UglgLE5c22vtjK9q+aQmsBVlnqWbb+FhiFne6q4F0XjgAVBN1US8m8My
% +A/dLByJXW+WJuCXZzTvTC1kcPF6jCG52kjI6qJUwZ4xGJYy3BtGPXHnESqjtSJckcDqY11nvv8E
% DgCFTMXHYFvlYhfdZ6XO1aISk1LxoPJMvyeXVcDLKsYk9XCGyeOAm3jixHzFVZqwow/PfK6TEBgs
% MN4hjI8+ZeFcOHvQKMgbH6Zwlp6MbwP+xt5b7+svbpf41zYfcrHoIvTcZGkfy1C8s72hpQYypzZo
% 23Pq1zI9OIRbh4GuiG4hhQLXmLpuP0AMbnU0REUDSNs4j23d31m7GQ4yPqoLY4zmox+jfKNQSmn1
% yBPkojCEvxAhuyMUPQzaUeqvHDhhWwbvlDj33TxHJ4Lg2hAn0bT32HQOxjA8cUDOHqKdU3niswuW
% c90hFOHIYO7ZfCp+1ltzUbIqczgum1g5ML7VhNyzSUL9cYZQUMU5YygcZ7BH/529aSf1JJS9hPTc
% XdrsahGY5PgQAC0q3GAKjIVsBuXg/F9ZqjohnEBaDikV7+finN+ggeuUuTnCnsBFSIiiMhT8jkbz
% thVonQgNBYNR2pwoo2N9dgEMl2qt4GMdGgzVAmpWJLxBXs+AeKewB0RFR/qFiZ94/rHAsN5MB4zv
% CQtPrXR6eGHPly1t+YRt4Rh6bc/CKTBKKaR61a6StA4TWQbbFTC3Ym1CvOGA0opJ4Gw/NKGl3hUL
% 9EzjANdcYAAK4ipE312l9FqFtSSLYz7bpAeTb02lsQGB51nZKfRnFW22mLQ2WdY9BCkEpEP8jTlV
% tG5ny9GABzpmeK1sHnYgOo3WcfB2UM+GzL4vrX1YruRtb4b6MNUuRPU7YBE3bt+O9piQcZmsE78e
% iBbc6RyATzGKmaqOk8a9HUiI1L7jfkKetTywDnlSDCiCpmobBGlIDYHUsEtq2E5Hyayez/WqY4kl
% td2M9Q/XHrbTekmZiJ3XcNZGmQ+aOhrmYJKHVsjAD3a6TSxb1rTsS7SdYWIs7x1bdEYzkMmDEZ2G
% vhuErg850aSXspk8uHYCSPCQ+kEGLCLkAlIv7iuetHxoo1E0uMWFY9mczZvrdfVq1emnMLCB24PJ
% ciBkDrjHTguBktt9hP9PORrTmr07BLxHd/lc7Nv6qJcOKIXKMMotEgcwU5NgDHu4R/YKb2SALlSv
% wP1Ou/s9sVWOqb9pSblCPBN3nyCyy/ssIq8LGZ2IcnbNYPz6cORZVp51+S3Qmg09PWEYOk9P2exT
% fjoFCqif1kNGHTcAnL5GoO58jkK8wDp7Rogajp4JLgeY1DGSOl5DbYVSSDC6YMs4iCVkViG2lYGr
% ROabKH1roXCuxn+4L7afiIIZ4MpLJ3ed6Vg/uQ33aJjgxB5UuNpIrLUuKHRiTQSKM2cYuelZSGcp
% n5oct5x0QZ1hWLcW08uWlOzTJuIskPEYW+7oaUAseOokVvFj5lzfFJzURRb27leUoIDHwT4s/Jhe
% rJRkMnZ7LzrSEWM/jfmIxFOwgB9d9vgtHD7i8BlkCJVvo/I1jtQTifGqWKoSoM9zQZMCAtIEOyX/
% azUX1Fm/1IFvCGFyHGHSQ4hdEMwEAZt5vDG9lyOzoArfDrs6/ARwYlFNfisEbeKKXDNqT84nqMP5
% xFwNmBWoiGF1ExyVGvO6zys9I3CSqWcmfL/OB3dwJu/g38qKWyamgWe5m7WmjMZyWFG0V2kvl8OC
% or3oYIdI631wLD6GgDmcse7G0yX9m7AwJ0v+24uIuymeKGz65/i0m/Wi4B2+LxCiwKe7mT6VIY5C
% fCoYooW7DvRBKzBCrcB3ZmNQxHBoNMo2DW4FpSIc5Pzw1TQDfOXP35oL+18/3lxpdLJaK5oa3fZa
% XzVeLJ0hicjxLoGtVYW1+/ElHN1AhfUZHChJJ1yR6puAdoE5dAliCkNTcJL9bGWhqKuLmINtBtQ4
% i0tPfOALqUuXfz8Ab8dwPuFQmmqJLrl4RNc6rXqfS3z+n+pJJgtbRE2j29PE9vPtnttaDuUxLNi1
% TCJQwy14Tk2n7j/N6vuCG058p7Vo7RZKH/etFgNV0ojrEtOPpHmHRw3ef6A3W2oL20X+4pzg2/bd
% OC/wKtvErSvLKHU03c/ER4iGsdioFOrG2Or91ValilL6cVRsIeDLUvhZucVL2JiarPrmllFBLQTW
% yR1YT3yOym1WlXh4qgRTIMwXwKIhQ6pwK33bDPVoCDksXkif1ZaQtm4A3nSCgWU+L7u9wm7P4juK
% IasqH2raTPoFWkevQxWp1fY4hvjpbpCNxjRNdP9zoGx8qkdt+XH7Fgqd2ss+4WGUqwDT4VYfrV0c
% FdEGfp7Mi8XXGi+AA1faHwI4NZqBfpyVyu3ocMFhY1r7/Uh2ljw7tnYQQZ4VxbjJ1gwx2Gtrmd6r
% 46Z3PhuwvcF+4FA3u1endnz+V1VnMlP9eck/wB8Lsy0sxa9RAtVizjpULz0xPT/O2usOZ7VX+lb7
% Peblm2u6prS3/JirqX0mrsz1GDgs/mLFr0rqVJimFTgakDSUHWn5VqgIHNi9BFdHH8JoJNOJeP5i
% LN5w2YGkkvMyV2+/6PsNKKGjdRSwZUSFuUcqxGcVx+yT0iw1m9Wed+G8QfuDQ9H9HdMDhKjFlIrn
% l273+tz0LB+sbuUD9inZpXjiod1DzLPPTzUl3+Olhn2n7plbjgaFxYI+nrYS7NlBiwG9/vSdbrjJ
% Pv6fB4Di+dhsxudg0PwtR21zMnRqMQG9wWtUt5XNGb3jpTTpeaoVPtEvYDye2glcQDfOnRoCMZ/z
% z8sG1MpCg16+F4CcqsKC+SqLrzos1to417ydD/J2yVwgNx80h9/GEGfL38dVEu1BY4ayMRfZT/GI
% jOjc67VntvmuBt+zpkgTu6rk797A+IstXtWKII7uKoUeoaTcZw1FQ+kJbLhAjgLZz+14yh4Bzhch
% 4XZR01ab6DRT1gURdc9uzNXE6KvXD1NdGrrt7mK9IbVrUC+TfvdxOjk/f/OK2hMsEAjt57OJ/swC
% HfTJb1Gxi+UBmP11BVKEk6urIp5fnwiR5ZCRzR+Zii+euJo/7uDnj/njHfxs54+IG5+uxn/w4HQ3
% vvsyhAc8CiiSxTpGseKXkHwvPg6yPFcxu110V8Dd/FGzaZC1ajO2c9b/Txib2Sp8GbYMmP9y4UFO
% uO+DMZA0DggQtEVZJ1Aw4wluxYIoXw0DmXt+RGNB/2xXpWiNOpBlPifByDaHHi7Z6HhVpamtFgtH
% u5drZAJM7BxZ1h7nAkM4j3D2GlheLt0j9XGlvzmI1bqkOuWSMJzbGHr4f4ZArqf1TaO+xaFr9unr
% erpPjT0mwobpQnLcTtOeNmFqb+CgKhD/ilsTTXMF13lYWRFwnZhr/LyMvhKoSivOkb/ipXrdtF5q
% PqrQXzHgytpYd3FU9vT0TPzK30DUl5HIBvoSlezKAxdPe/omly/egCK1B1RBVMYH9iDaVS202swQ
% dXnd1ak+u/9O+98m07eo5gxT0JYJXc8WGTAQOgxgRYPmfD4i6Je34tH+vrV9jGHLq18/fnyLWUud
% GKGc6ZvgogoCpYC5s/aiC7Z0/d00SqGhuOCPkSV/mWz6yFUa4bc8TVl4Rl0MSHlWKyRrtUKNnqxW
% WLyuVifMDleyo/8CUEsDBBQAAAAIANKpVV27WEsrzwIAAJkFAAAxAAAAb3JpZ2luYWxfY2hlY2tz
% L3JvdW5kMDIvMDZfdWx0cmEvdmVyaWZpY2F0aW9uLnR4dI1UTW+jMBC98yvmmEgksiF8OKs9dNUe
% VupK1ba3KqpcGBKrYFhj2lSr/vcdm5CkVaXuBYbxm+fnN4Ov9rKwIIdS2TU8GqmLndJb6FXT1biH
% x0GXNfYh/G4pgiiEK23NK7A0uEZZwe1DBEWrqVq1WtYUD9r2a5jFIWTzCRN/jlmFwMU8uMnhWg4G
% taUlrCpVKIp7mOG+azWF4Xl+vob72YJKxTyE2YL2SZmPSBvPch/yEOLIZxlFeeaiU+4EPBZ7uk1w
% 94L1My5K1aDuR7GjJZ8KnISkq0mJyNJJSpIwNmnhK87FQQ3Pspgf9BzzZ/gTy8i8CX4YVW5xLKDD
% zxw6OlTF8/lx3Tvu1rnbxpOlFJxDRsM9B6Gj2GFyCrIzDO0+qfRkecYySriXly0S559Ic0d8mUOp
% DNIEbbFt0BpVAOqhQSNds9eQsAxM21roa1UgWfZ3wddkeghs7TpDVO7zzTE9o7E0codZzCKfcFXf
% gEdQyQIthbo1DfWlRIumUVr6RnjWOPVscUpsPPlCmEhXaexxM6W7wcIO1XZHo9YO9vQ5Pw7r/cy3
% cuF6dzTDp5zTh5a5Bn9AsBEwGnZIvEd8oODn6xuv8L0xPInTM2si/h/WZGnuvaH3W3BnCB+ew1wr
% YsEgjnkq4pxnAT4wsDvTDtsd4MOKzu8G03nmxoXlnKBc0MTRwQRbkXjBY+FkEyZiOROrJGIRE3m2
% CX5J8n8PXfuCpgdJ9wh1ZjD0QxUIcmsQj3s9fY/Y0h/awS6TjBrQdNKovtU9vCi7IyhCoSyW7lLp
% rRmKsaM3F7e3wUVdA/orTWmLWzRQOKfol5UWe6haA70bBgs0CEgg/DPImm4mN5uewT3WIImnUpq2
% gaKWqumh2GHxhOUSfupxYVHJRtWvXijZqJqhWdSqUTY4VMhyuvIINfTolXf0N1Q9bQyd7NAsqbEh
% Nc5Ou5FKSxfwMvgHUEsDBBQAAAAIANKpVV0ewASCfgwAAGkjAAAqAAAAb3JpZ2luYWxfY2hlY2tz
% L3JvdW5kMDIvMDZfdWx0cmEvdmVyaWZ5LnB5nVptk9u2Ef6uX4Gep1PSpmRS1HuiziSNP2QmTTOJ
% p180ioYSIQl3FEnz5Sz15v57dxcACYiS7TbjnCgA+4LdZxe7oN785X1dFu+3In3P02eWX6pjloa9
% h4eHD+doV3msrKI0joq4n4htERUXFtWxqFi2Z9WRs9+zOo3ZkH34HCVxP89EWrFdlpZVUe8qAQ+D
% Xu/3Ol1oxuyZF2J/GeSX3kcgF+lepKLi7MCzE68KsWMgjZU531VFlLBdEolTyaKCs7zInnkMFCyP
% cl4MKn4eMPbxKMoeTB2K6ARzMc85/Emr5MJ2R757KlFLUTAlJkoOHHZBQvg5T8QOtpJnyaXKcg66
% /pqxfZJFlUgPajMFL+ukYqJkdSnFR/CvLHmB2xugoXr7IjvBrpOEyz0zccqzomL/AONUvJDzIL6o
% sixpZkHrGIwkZ09RddQTh10sB/d1urNICg4UvNfrxXzPTqCXyJOLk3vsk7voMfivYEst1XFpZJ8V
% THgsIsMNQIlT6ajFevrRY1uc/tSdJp4rwd6xxzV7twQ2b9lWSuJVXaQsFrvKKVylUp59BsGgT2ro
% 8+IvWPDaKLNBUUWUHriTGpJwZbOjwmO5a4oplIDycpIw2ewyvt+D+8DXpWN+0ZIl4cvjgpmzq2hb
% Oo/uWm68VSVgfZbw1ObkWXb4pv86PNxXpXqcVU4Ellb6SQjR+shlyyU9ba1Nl/XJOYPBL6Ts2YMH
% 0Pc/Ipd8tNFjDt4+iTRKwRPZZ20AAObPgIMDL9iPED+iLBlPBK5DjHrsswDEcYxxcOKzKAm4FDM8
% HiCqSUvwyioRJXF2SQ94ILOBpDWtSdlSbcPcWJQkDo4SHWwvNYnV0rzgzyKrS0gy4oBsAo8FDVCe
% DKCAdwIDLGLPotXTGv4ha9/GqwBGKT9XjiMk+ls+TwDkANEpOQjk4Hrs1yzlrs1jj2QlzSw6KFD+
% 8a0JVMgjpiAfPzwastaobfbxs5nIxXNWEQ3tyArNW7rb+lzh+M4q8lN94kVUZYXSbwUh/VZJ72tr
% wAjp8bju0CvHtmz+2jiQvNAlUEKWBs379w1Rr7MW3WnwablLLa3YQEuirv1gjf9DLMBZARj/TSXz
% Jgg+FtEOocQ/1VEiKsEBbs+YvXf4hAcBnV2M0j1keTg+ChZVVbQ7njCCm1jAWNts8BzZbJySJ3vD
% xvh1EAsgwDhCJNtTaVacogR3snICDyDn9OFjbS/SWtGqvlzWXcVJW2OJ3y5s9NxlWRFjqPPNPtrx
% ivT1mDATrjRkVecJd/qwd+cR3Sjca0jZe8O0o6VIKynekFJ4khgCUEmbtncFJcwPSINyY3kmJwkK
% Nm3WEKX888Y0kR0p54ZOr7EjQInELAxCPHZ22femm64lDKIcSwlIwO+kswBtqful0LsVcxdUFOCJ
% xN1JFRwQfYhlQ7Vu7N7RTHrv4rYUb6Ak4iXXcMbKiavwg/KFygqj1tqWWVJX3A4PgxceElldQe0D
% FREWdlD6gSOxMsMqo2RggohBqZaWewgbqGIKcR40DHbgUxGDZHQYHSOq6HHQpMyHdA/YLaBoi6pl
% auwC99sA/Z6jaYFtc6j1DkdMpveMiRzoGG1V63oNDwcoE/D8vbhgZfSO5CxBs7hZFzQ6mw66uHcT
% AZ6F7yh8ER3m8WiuXN+mh1psZWmh+MQNKL4U0qapbxw1Bs5vi9c7lJLQPBXxqJAHRfQ7ykh2zFjR
% 6dnZzAS415ryXoKF7adGHsKOZKNgfbnOzd+Wh0qVJBsF3aYqs4e/SExKX1HKsWsyrI9kOe9Ave9B
% 6EBZt4W9Y8FEDIJ7kGh5UUtyXfFLq1TlzQlk+fyVVCmhThjFMAIN4MR0vw5QY3en6OwoNjfyLNSc
% 4pkjf6H7kyMhRyeqlhZC8UjGuCmGyk5iJs8Qaw3YABOBUR2vTL2xUGtwL5ms3VtCMBUAF+mTawnl
% Cv5Qa2RPoWcGdY4JRmt4Czi4zEZLx8dv2L8KcUCjY8qGhlxgVY/NcRsveQGFPdmUTnusa+AgxYTL
% Y4MREvAzkyJMy0A+ed8PBoCXEwR0xVMOldQeGloo8oFVmdHBYXCqqHc3u+2yogMmz0pSpL8XW6yj
% ikONVRTaGGlk634rFhqk6UQGuPmWEMAF+ZeOBEMGeVLLyV2JzK+LMFvdEtpwHiMeSt0tN43YthDx
% gW/kIejUGDRWL+o4NeTFZ4/RBxZtQ9gorKMPGtOsJI8NpFnohp/wFLM4ybzbbOVsBedTy0SkupkH
% hEAAyDsdmS8XoG+U7o54sJcC/X5mW6h+EwSUutbx2IcUMirzJw/SHJ+wn+8HC+wxfPobWM09FoFn
% vLyRVyUrB1g4ocemrisfnBGQzF133Topx/qeLg4+QebH/qUTKk6+SqGfylf+mqJFyzBaBdze/uEX
% Hu3ZH5uX9PVGsbJgL5rwVe0nn5nXDo0WQ9BWfqHNhrTZKW02fMVpq9VFJst7VxOrcAaE4RBrnenM
% YxN4mKtEIx1gavDJkjuSRp6T4BEIDm3Bmv6+8GA6DdFJowC5jMc+Sp9OQI3R2jWB8duM/RLVBafr
% u5YDc8Bi0AOnkHeMcXfxYN6NqKDIZ0ZIGLw/fuYJZITm+AWHKM1vibzJWq43I07moX/TbaKKPDg/
% wARQa8aiAC9DUqJyV9WEYDPoGamlI6hGxVZUeKGpWEEliKitPQjQFFri7YUy1i7J8JCVAYmE4Kto
% 0AAeo1wiPZDFrER7A3R4IHuTI0zUq/r53hmdJVATHMFbR+B7HNLt3f3SeXTVeuCp6UMkHQP8M0SE
% YI7JrnNjq8hKCgzoKHtWDLCrh3IOuom6vXoxsz6g7ZTXFRUhjl2OKrZoFeCppfTpu6e/4i7WV9ds
% BqV/RelblL5J2UkZ9/Rc3krTnTzyowSU81JD1L08v7qQO25z1JmkKXyMxG1L0sAgoFAacygrzBWa
% deuibzAUImh4oFrtUHb6N2Yg0crvndY/+N+WXyU3ucq42liyGbXq1LzTpCxwcWbsT4m2TLAumRno
% xtTdntUmnSmsoVOnDCZLyICUQunrq5lXfpqpSDc6Wl1CgqqQSDpKmlB7KLJMi8S15vmuFGkTjnsl
% WJVS8iS1BDWNgsce9PN3Vk5r1+pSA5bKwu27W8WZnRGVG+2WR2fEvdhXnKddFKkJw/0tjq7mEGMw
% L2Hau3HREZrvDzq0eqQDrBDfItjg0msteAXjBl96vkXYfDKahAbGYHELMkKZx/r5CrJKizabi4wF
% VQcQ+UuzG0qt/QCiPZhN/am8pmM+fMck3jN6bXuVLxcF88kslF/b2ZbKYKWygLXmtRsJqN7SUNYC
% IUx+C/zt3bvXLByRQh5Th6THsrpqv7q6cLp1JGsNb5/5yPpWjGhtvi1K9Or/J04aZN6OlCd2fRIY
% 57TpiSfK1I5yNX1QFp+PMUjQ42rnUEzssKl9AhCu6Oh6AhyuVOcqu1E991bPsb4cC5oxf21KV0yX
% kAPDuU+BQZxgIAwm83AWTLvL376F477PRthGycUTfB4O59MgnMnowXYNzgA8uFXlou7ZroJ9aL53
% 4friB6sJYiE3Kh8Do3+mEaxsrxoZGrcPmNViTK9xVhA6FN5QYE/9WQAbxKJ1Cjae+yPwzjwI53de
% yQUzIBz6M38+Gg/9oT+fTdfXb6P46hFUHZIsZScmx8gNaCka6Lz5uH5VOLeTWBAC3RyNPqLTcTyc
% h/YGgf/fu+MtXUAnHiAqGIb+cDybTWADs/H0is1ojS3xnVW61saNeWZYYORVzagVpHzjQ41bZPXh
% yPhmhCvJG9aif8rKl9qSkgAIOacuoG4HC0aHgvOGx9Ny6A8eOnkAaX4aTxmWTFEhSnzRSG8fqb4W
% WJaZPxlYsN9++OOPBxWn2liTFsCE/sl0PAKHg0Um/iSUkFfjcwCMSTscg7VwRUhmngRBCItG0xAs
% CfTSBeS74XA2nUznw9F4NJ8EN3iwwPd9tf57NlSSg5G1VI2GkndDEqpG4wdWcrADxofM1NigYgqE
% DggH6SIVLFNCfks4/nTBbp4UG5m8/laqeyFGDKGPx59JYJezYEl02sYRaBHOBlPYKvpBPv5Juv05
% Hph6j6aGsmiYyWw8oWdzFdRiwKrd3ls27uxR+f6HJFEvmIV6EW1sUzZhlu7GO4gGA2aLCiMLuthS
% P+ZQPw7R76zZz+rHJP19dBLQ/eF+IQbEqT718c13ZfPTvy2J9TUBkNQll/dUUBvuS+vHJrIrVKKx
% WRTpAcHe6wl8L5hGJ77ZIIgeNhu8edlsHmTilNcwvf8CUEsDBBQAAAAIANKpVV1uLxaBMAEAANoB
% AAAiAAAAb3JpZ2luYWxfY2hlY2tzL3JvdW5kMDIvUkVBRE1FLnR4dD1RS3bbQAzb+xQ4gK1Flumy
% zbLJe04vQEtQxJf5KEOOrd6+lJ3XzfwwBAHwbWXBz557Etcr8Vt8YY7zaPhTeyuSWfwZ59rLhKfh
% 8K4bGo3SxgUBNaVBApMCTuq1qSSIGc320gEvEj9XWdkG5wY1mEeBpFo4HM786jTnBKO7lg97Rqq3
% I3Kw9XzEoh/LEdtjy7Id0ZM3GXDuxTUTuU5M0Cmaqf9FL1c2nZXTcPil88wWQJCn+RQLx72V1zUM
% /kCpDsFYw0ZNKYA71ymkzrVlKSNxYRmXLO1zlyoTuLuRi4WE0R/G/3tFVIHbDoRF749gkmZ1G/Ba
% r0whcH/j5oxsE1ayRZ5X5Q3SeFc0svm3gTO7ySVxjzxsG3I3j4uLFljMJxTufO0+HvEYx6W71jIc
% /gFQSwMEFAAAAAgA0qlVXfAyhTSqBAAAHQkAACYAAABvcmlnaW5hbF9jaGVja3Mvcm91bmQwMi9T
% SEEyNTZTVU1TLnR4dG1WW3IcNg789ylyAocACBA8jQoEgY2r4sRlK1nn9tuTqsRjeX/0oRk1gUY/
% NEmdJ/uafrVGrmAv7sokP4unZ6kOcgoKsTCNsfuWTw+vSfbTT4Nefv39vz9/ik/1+f2n2+92HDeh
% rdRte4zSg78JqTXkHLpFZ1b0XHjY9hrVe04/d8126dFvIF/r6zsnV56rdE8mj8XCpKp+e2tkiVfL
% ciHBJ8tGe0YF1cVLSzzoX8g/6/OH/uv9p7/esU47su1s8eEVuT12Sh9WyxpTfetxQOPVnunRsRK/
% IBM5PMf3kB8yXj/8/tv716+v7+awiBYXFSvvvCz3dJSoUvBeETyEjQdYJ7u6aEvx3sxaBkaAzC8f
% 63744+MTq14HdyBPshwzNScYuHlOr4PljxjutlafFbQwJpf52KK8sFlP/RH1QezZGClSRJPlDBoX
% J7DApHLVCcOtIgfPfufoMciKKk9eo+RD/Iz6jdsesu8NvHuGzpLYeMdoz8tJqXd1zLy1No3iOTON
% oauWVAfLsn5AfaZXsrlGTM1FxXbGONk0OHHCKtZLJGqkYLMcKx01j/JNJ4qlH/KSl18+/OeXJ3Ih
% ny3ENqDceSWsxsXAj7WTKkCi4/T7OO08+AyE09jXeNY4185bzAe10CEJ9DTvHUYTRhiJ718h3UXz
% iKwsj9NahE8eOpwnzkolnRT7G+Y3YsH4aR6uwMOWrLsx1RlnB5+jq28M2GT7EAJBZzPMFRG6R9Sd
% /AbzmdYhY+fGzRguGqp7QWu+gwgKwptz8ChwsBhcLmI3ZphvTeKx5Cag58vXN7wS3FNr8hiG27jP
% 6mXiWrk5qrIO4qbvhtXwrMDeM64VNC1Cw/QH0AexDPbSG/o+sTBg3ol8Gjk2I3C6D+h9AIn1Xq1L
% Ox2EQntQJel6Av3G7DXYSqOQQbtgFwjCaoEDHGufBwacZdSIlYVRq7HZwWqhI2leewv6TO2CnaBZ
% w7enIwvqjgXOKLMMKdQnLpyCjINQ/fJNMcFXkUOJUNMLbH35GF+fiIV0ZImOoofyl0OYcIJaB9iA
% 0O6Z065obAmPhAzKIm3vvJhE8g3k31FgU2cMeAtJPOJgSAprOwkEqD2HcfgEzTjOQSmcGI4dYvLl
% zfov5DdScRwvG2bIVmQAfCXoBA/G9LkADKN251qhcDWKaMQyKNr5JLbn7yGfKfXQyQg/bZ1isCAh
% +DuHt1yk2GF5hFZiB5hEHcc7oB0nroAFS4BsL3/8+vo5nkg15ZwX0xhX3bR1otFHCBWciwxdGHu3
% zTi8rkyU5F0lYnM9wsLPD6APWicN+Bq3cfdEFqflNIjnUT0LUt34EZv9wtA4FM6Xa6I6vFBL8ncb
% /AP6pFbclVMPYiOOMjcZRup5UN4ERSJ1UVyIloN+wW9tIge9mSdtnC7egj5Ty9uJ78ZWc8/Fe2LI
% Pm0XHBqNuBAyIs9A4+xRqeNaSHuDeoQC7BVfvtSXLx/rt9cncpnwHwHa66IcUpDzu5GPqOm4O+wu
% 10O10VmKYjgXBYc3WjeaE7TY/T+wD3oHMgPFiv8p7qMDF+Iqo4+tRjp1x8HG5Qg2xCK2Kg2ocJNf
% +AxR3N/BvuXif1BLAwQUAAAACADSqVVdV4AvjiQBAADOAQAAMwAAAG9yaWdpbmFsX2NoZWNrcy9y
% b3VuZDAyL2Fzc2Vzc21lbnQvdmVyaWZpY2F0aW9uLnR4dE1RS07DMBDdc4p3gCgE8dl0hcRHbEoX
% XMC1J2REYltjO21vz0zDgk2U2O+fgxOKFUIuwM0zCp+RJaWx9Hg9O1+h13LBSsIjkxSFSosYk2C4
% 64b7bnjohsdueOrxEQNl0ocKzul0+y0uT3AtcIVPS56pUtihTgSOlURa1gOMTX03rClzwf7zSwlN
% MaHHO6WFLAL/k6fYFhJXOUW4GFDFxTKSwAnXSfHskV0pJvCs8luJw8ub5VcJUV/j+Yn8j74fLxaL
% ld/qlKTslH0dxgVNry66jSUombxltt5W+ipiPqUsCu+xTyjtWKqLlVfCt8tgS2zjBZw0HEcUn7J+
% Ca1MJ+OsNNdLBzrrKmqFTFplu+6uFjr3oufbX/B/vYUQk05Lssn3N79QSwMEFAAAAAgA0qlVXXHd
% TWfzAAAAZAEAAC8AAABvcmlnaW5hbF9jaGVja3Mvcm91bmQwMy8wMV9sb3cvdmVyaWZpY2F0aW9u
% LnR4dE2QvWrEMBCEez/Fdm6cI/69s8HFFQmkyQuEwwhpHS+RJaGfYAj37pGMj6QbdnY+afbNWvwM
% klmw6EgEdEAKWBDkUYCxtKIb4Kc+D/BRV7cC2j6qpomqS7P2ElX5XEbZXW737JqSwHVQPuVyrnGe
% iRMqP+1Ulw9QVnXTFpALssj95ML65136uooWKmE0/Q/1Vd8lY+MyCBQTPX4+4Wa0ig/s6P6evRwr
% gBtbjUQwY30u1BjB72Nq8s1kYJ60ggos8wta8AtTcACfjlPE+iiI74tldpUyAXlstyD/cmCYcyhO
% 8EoqHusxZRZBaR+BKa/1fMp+AVBLAwQUAAAACADSqVVdd8ffKbYEAACdCgAAKAAAAG9yaWdpbmFs
% X2NoZWNrcy9yb3VuZDAzLzAxX2xvdy92ZXJpZnkucHl9Vttu2zgQfddXsC0CS7JsSw7aYtUyQAO0
% QB8aIM+GN2DEcUxborgklToo8u87pC5WbO/qSZoZnhnOmYs+vFs0Ri8ehVyAfCbqxW5reR28f//+
% +4EVlnDBnmRtrCgM2dSagOSqFtISY9mjKIV9+eJeJWeak1I8aqZfSC3LlzlCBBtdV2SjEUjU0hBR
% qVpb8qMTEGbIfWtTMbvt1UVdPQYBhw15BC3rpixF+CvKA4LPLV3dh1m09h8unj0RkmgmnyDMkl/T
% rLPztnOmFMYbzkxThfehww330yzZRVF8u9qtPcDuCLCPosW9s4giD6LBNlqS2zaYZxUeEtXhH+h9
% eGitxIbI2pJD3plnaRz/5TX+lAjZKCZG2aNByRdL00H4eytKIOxKUZrmdkqzL2yxoGrQd8B2HBTi
% HuayqUAzW+to5r85yLoS0kvaoJUWFYQYdXdM3WQEySKsLEN1xX0G+DEDywSpDVUczz9GPg3BLT2S
% kH1Ko0AZulL+nBqfy9Jl5FLRO1wHRd1IS/9Miho2G1EIkPaBNVxYM8nTZMKFhsI+IDVjaV9eb2SH
% omw48AehNTw1JdMPcFC1REBn8BoMcvpH5Svtg9Pj4NQsS9rwkMTblV4jjzfp+ngLZXzsN8vXYCRD
% 2tyXfHPPT32J7an03F8tkTYCpQEiZ1lfEojlE72/CtH7qAI07USD5CSso6l7fBZX/5+DtSuZopZW
% yAaOpWYMYDd56L2DphTf9/gyNI9xN1thbpLrdT6OBytxmQM1s+z0isvYzJaDKZSDemzu7zccMVP0
% K3u/7TEDaI9Yvao9GWICo5PTvVHslWOU7oLwVcZmEPKUnsc8aAtD/8gcJwiWw+sgPWmCLJHTZfSW
% hh2V02zG38jQy47SLD/LvPdEu4EjE+7i5u3IOQUo8sKs+JoWZv4ENuTJfZhG0bS4wKIzTB2PU57G
% hlJ4G39SuBsgjrBQmfAkfnTF31Ge5iM8DxWbG0hClcjEJDyBpDjmtyu9C/3rCm5cLvIr/ZT6cjdf
% sXDeeHbBNRfrbATQtLPvYibdc0dxJpm4OVN8ID8lBzflMbzyhbCiaCpsDwvEbsFNYCGFfCKq/g1Y
% 701lSFXzpqyJ+jsEpHl+Bol6581rz5S2tqykbp8gYrjDtKF5dLpIsuTOTc8r1J3D4wUNDT3QzIHc
% dSCXzY98+RO+jcGnGkUeqxW1HN4hg62dV53H33F6Pny7GdIOm5Mh7On+QCR21hZXNvzTsJKUwLjL
% 7DMrG+a3uUtCzTlRORGWSABuiAHFtGdDA7MVkjS/MGGbceqyj13tOtJbxo9JuItDl9rFYtkPNCeZ
% 9ZJgPDIv3QLXEy64yc9+iuI6NYI34CehNwTe7jCTTxJcJ+02UYnfJ8Pw7ZvM1a5+jXrYbw6g9e+O
% +5cIM/cNZbLAFQykH+REMaFzB+7K1Tth5CQoIgx5AtkICVjaG+Zm2RY0zAPPNb3+nFxjS30Onl0i
% XFHu4lj+RzG67HRhfu9jgAOrFP56KAclKYI50PzI6SR5TibIH3rFfmKyj3DWR6g0cNH+zGWTKOiJ
% wvUx5KQsnR/8lSy2UOwNXhyN+Jz8wM7EhHRSpsH/RrmuVbquN3OE+xdQSwMEFAAAAAgA0qlVXRQx
% 2ECgBAAACwkAADIAAABvcmlnaW5hbF9jaGVja3Mvcm91bmQwMy8wMl9tZWRpdW0vdmVyaWZpY2F0
% aW9uLnR4dIVWYW/bNhD97l9xHzdAdiwrrmMbG1A0KZa1S7LUG7ovA2iJkrhQJEdScfTv945y3KzF
% FgRwbFP37t3du0ff3/52c1nQ1c3u/o/5gn6/ur9+f/3u7e769mZyKaLcLOa3ZbR76RfzxZvJO9t1
% wlQbckNsrSnoUXpVDzM30I9UtrJ8CLO/gjWTqycVN/OMhNYkQpA+KmsCOX5fzehOeGkiKVNJJ/Fi
% oh7ISy8Mhb4sZQh1r/UwmwBIlJE6EVuJF1UKTZUSjbEBH8JmMqV72fRa+DyfgkyUTxSkBMNK1qpU
% 0pQDSKIET9/l2fzbv+9nwHgLmlFycue0khU56WtEJSzOXrbKNIEehVYV2lJx0EdRVdKTlqaJbchn
% s3MUDCaNDJGsx9k6R/8u5uStjVRyE8AJ4YHDr0TZktXVNJ0+5xi7iKR7G1uKrbd905IAESc8Qsl5
% iXY+Cm4ogPuojCTMBG3xzNgJcKfa+g5dSXnQloGMPIyJZGroTzwYI1XT7q1vra0SBL6c7vs4tYA8
% EfpXzddfJkaibzr8xzNTh/kg4kl1ffclUiBbSp4Sg9KLXuULhrsTII1ue1v1ZcQ4iUyxzlyRPeTF
% pigoRG9Nk715MU5mtw0P8kAGRVL4u4eYquV8tc2LJUSEjEHokIEOBFYq5prPZwy9zEfo1Qvo/OJV
% 7Is3wF5f/Bf2ImGvgb1KtFerE/b5q9h5flFsi/wV4vl8OaIvX6AvXmeeL1fLbbHKX0E/LzKX54l8
% vshP+IvX8S+W6+35+v/YY8g2qKgeJZVDqVU5rRXshDrZ8T+W3Jhv9BJePmWOunGQxygTLQe6vN8R
% L0EnI3IlwSoTZQOUD/Thz90P7nr6M1RZqaNYd9IEiC5RUXE4RjwKr4SJ09DvgxMllqdMu/S8eAcF
% LR+BQ1JtWkXmciQdtmMTaqu1PTAToYfkTXrIcASte4tjbQVvB2HNJNLqflz8Xas82IyLWGLdKhVK
% j/1mrxpNlP3nmCwr5nOKNgqdOtHv0UIqYaaRizC8/emT78cyvvbU0yFKE0dfCvkyy1dZvs4W+Zad
% oLa9B0xvYljP6MZ+8WxYEKLgvkFSrQxae6LoLWaqYkDK4wkes152sO07rHx99sLyng0MkJvJrpWj
% s/JEouycZlMkKApdgQGX6iu7PJrYKVGphepQXED2xFHSwauYLHzo9pZ75JjCOA4+P0Z+Y1KRQWb0
% SYAFf8ZEJY+8Z7Ecq4QbYzKYtgpJKbZPz47gQbHTsuS4cb1JoQgJUnh2eI+aIFhohILVqIN6vgPR
% o8v39OvbzeSTeoKBNpLwRQbwzimN+HhQEGdS46kpHALLd32c4d5DGnbP511p5Om+cFUdrXNdhpYF
% J8fhg+kxETQRWXyhlTKmIWB/VWAsvOdHQioGpXu+iNGOJxky4sLQR8nP18htcFNzgR/FTn6mg/Am
% 3ZLYXzyFerVtEg7W3jkcnXUqBG5yowfXInTPVxl37hrDUYKFCFGn26QWSvcQxEEEHjR3Lo3Byz5B
% 8EwfjD0YOlj/wN/srj7/8v6Mbz7BPy1qL0bhA2ULnfI0eV74jYEW8vpUHLXXFmsPm/kHUEsDBBQA
% AAAIANKpVV3bbKVRVwsAAIsdAAArAAAAb3JpZ2luYWxfY2hlY2tzL3JvdW5kMDMvMDJfbWVkaXVt
% L3ZlcmlmeS5weaUZ23LayPKdr9BmK0EDYxuBLzHKPBBvXHsqIWvs3apNqVQqBWRbBiRFFwxh+ffT
% 3TMjiUuyrnMeLObS09Pd03f/+stJkaUnX8PoJIgWRrLKH+Oo13j16tWHpT/Oja9xEU2CieEXkzA3
% 4nsjfwyMYJnMwjHM0+ChmPmpMfMnkyA1/Ghi3PizYGVkxdcsD/MiD+MoO258jo37WeznYfRwlMRh
% lBtZEozz1J8ZcWrc3P1m5EGWZ7aR5YDDTwENEWLMwq+pn66MOJqtjpGqRjhP4jQ3nrI4atyn8dwI
% 8yDN43iWGWprHM+BHZ+uliApIEVIuX9Ls0bj7sOH34SztniX9/gFf7vh6y7+8ktudWDSq0+s2qSx
% 7nB56gx29Pi8Nj6F8Sk/4+d4Bg/IicLV4XK6NbnYuI33g7sP3nDw59XvYt2w+o4JuBg3L/g5fN/y
% M/he8lP4Wh3eYS5vdAGmx62fwvQABqj6KcwpwHT4W1wBquCnyy/hCxKAL1CHMGc/gzkHdABz/jMY
% 4BNhLgDmlPg6I6rOif63Ev5S0vP2IAxuEs0XCHNZ3rUPqWGszk+ALoA2d9O4/eOPPzsCZWRZRDER
% 0iPC6xKzushjozEJ7qU1mAPWbxiRmAURDG0/y0ARjeh1l+zAn81M3EnjZyaEGR1Z7OSka9yDvsOS
% EUbGgDVoGuIEVPQhMCPEaChMoRHFOQE6oQvLCPtUwYasrwDNJw3FjF+EYYZy+uQySe3cz8ePYHrm
% gC/xgk9i4Cxd+1ZkQW5KZJIHxo4+Ha2XG+QlfIiCiVhvGoZk+MEseBYEkkIkZYG3ZGBQwcQcOIX7
% 5pa2jPAethTleIAWaXQMXsJcMJrXofRt6AvwIj13Fq66srYiCjsN8iKNjD/TIgBcanbtz7JACrSo
% UfYJDivac5PB3QpcbcPGgkl2eFGn5RjcyjzDEyTC8WMwnnp1QfIhhw1wYqLzIpk29LPi8hBUAn8/
% sSOJZGt3XUg2OIl4uFHA+tB6sbNdHkadIy4+kQrS/m01xFfaOqqYk97bTPkYkIT34dgHbyxIoMja
% QDjIVCbllOFJdJ2ufVNjFvwjY/ZIgD8rHWXDGIp12J+FWW4mfpgqDCGnCeKpHJ6W90Y+YV7pucXT
% tkWqtdSmthLLtmUPjoNlHkQTk8hj8oVdptST5HDTR55J75bMHjgrOS62gEZ9fDq1YZcHVghU7qxY
% Xe2HlabTLUNQzGM/SZAac8lXjJQ8AAFK7UeGRQTkosI9M/UI5AXwKGKS9NJj0fqNtBPCnAbzeBGQ
% DJkN5wSObNrSIgC8S3BSK8BOIgCMuRAWUNZxBb0A+TncutHyGJWcAXT97fvKSAZ8WBrMQOpK4idB
% 6sWzSWUMKZhQn3AshOj0K5NEn/q3mbKfedUvJjhd8KttB4ZPDOCfwFOyXWfXJSVwUaOImUpxQAiS
% gcU7ca4JHwK6vxEd4txDZpXItGQpOlTOAY+fatKJwt6LKZTuYjWGBMmcoly0MNZm2H5ir6f71JjT
% NgUHttmNB1ON8G909wrTZbvbCuXyl9qy1anWEx8e0595oE3eI3qFlId8GkYT9VA4FKK5bNJrxXEu
% 8AL7TqxL4/1njQLcCzptKTwwU+OzkGkNgH75F9AhxNcLyjh09oEa0KEwTPLt4nMYQ6MtjH8Tcyjf
% TgP/+Jk1AQCsLZFY/aJYfSnxwOcPxbHF4SnlJ2dacSzKerqUfaDiv4hoi4cvFUWNP6UCyB6/45/B
% bJUaQC5uJvjo2Tdwxa3wdbKrYhZPkIFSS592r0kYagyQEDI4jZH1256eJlpPs2nwjHmRRud08HAo
% xBO9gWHhtExX5NqRtXulipzu7jV6vcadV0vF9NBOdKD4KBRBjX2K+/splVzU0TQr5uZHoNLJ3dZH
% yKbgdzc2JQxzu+TI2uYSXktdyAl7vSIx9UnQ+X4t6KM83lDSRhh74BBOtwI7Am1lciodwBAldWO5
% RZh6EF2IBeaIXyN/czo1YvaUBtcQ8oTj7gloTrKI4WAq1k+taXtRSWuEb1cGQ+2qNvWgqhf7Ax0X
% CdU/oIQaFR2+hnC7YbtRRleH6nWBcj7gafRwmP7yle33cgSbV3J0Dah/NT5Eebp6DuFlIFhmyOck
% iEAqKyOZFZiILPw09LEqBWkl/jjIOIrcgNAAsFEePAQpVLF7ImpND2mRXtaPRzr05MKzvnfCk5Mp
% TPBLKgODXwROpeJcOeFr3IeP0iCvwntGWL8LBwRxjCWtXD7C+mwXdk5O5Pt9MZsJ57t+rq3n3X8r
% PLKAZA/0HkoUILlFCGBwiEVWr2U+ulXdAlmAM20hlveS89b3QygUghfQ9J3087CCIgEvksgUjNK1
% AVrr45HiE9LW7zqNwum+EP4XAeD5K9Bt5xkRoNI/SyTPde5eKoIdXnpUcoBbmUCevsN8r8b5wMU4
% FypXNMalICrmQQpJnkmnybePIVHc2H+9DLKLZh4I5E96n8x981e9NGCVIMxuKzhCqDvWwp+/GGu1
% usY7YUxb09a8tbWlLD8u0nHgRWEUQBQFLl+9evWfaBLgkwVooXM0zLIblRRfIdFSp2wqKH1sUcWQ
% 2Bn61Hh1jJ0jYyScdbdst2DrpQezt9TdUd0aWu3C+IIaQPLXUvCW6uxgj0fD9uTJDTzScxA+PObw
% IAgn/yz5i6USwSK+rRrxEsrDm733TWvlzkiXO6PD5Y6tri134c7dImhUK4JGh4ugWwCqiqDRVhE0
% qooguywebsv9MfhQcO3AuGvnce7PRKdRqboir6+htPlRGi/ZpUOcvm0wRyaxtMVzWXvuWTNBaMsI
% +dWOwqq7Su+cxJks7yuYK+k+tmKarBaLZFJHQn0UWe3he1xB2O1VZeWVYwJyyMJe9xCsTEGyPI2j
% B+EsdyOzpFxj+2FAd7c6AhIb7F+WkXLdTJv9lDfjFGr3Zl/KsCkBvXFcRHmzXzu60a2gMDKpKgHi
% ZCPU7Ha651an2+n0mJ2CQ1lvbMkGVvkqDC358HC93t/pjOAZbI7AuSwIJnCncMjMYR35OtriWkKz
% vfTFsmoC0HjAsKAI67gtq4N0Ok3c8GAnHIdg5GGQNV2hgBuyY5NpVUzrCe+pbCdAdSvK1gc2k5it
% n0+pzbLqVAxLhlEXDrSDZIPDPrhzoGQGpiVs1hbWfnCrmh60g/UabjpQsvHmqunKtkKZ79Pl4sdV
% H0Ef7E0NHETiHq3xZ0NwSux374TapKYEmeFxEUEOa7bMqpUkfT57cyvE5/px6dulRn9mbYv29qRD
% HBD1UAtuCQQMJH0nrC5xOsmkFg32VWhwSH/KQqGiB1AIUh6rDXrUgtCUUtoJeiSVwKt3QECT1k3A
% 8pA/Zs1+0zo+Pu2A5Gd++hBkuaeM7tLiTeRg+6g0iYw3wypulTxn3hZWq9vcKCqSNJ4U4xyv1un4
% nJPe6h5/T/03QDavsZFtgbNnpAwjIYu9OVQ/tboIosd1qeSY6e/UA3WN/z+EXJOxuStk1poDxA8z
% eti7EqpfIvN54ABc7zalV5pS3KplGHuFEe5XrhRnh8ojvceoUbdH2xUnNIq88SOemQinahYXe1Ko
% KhoyjXp7m4goXJRTCG4K1YE0UZ7UQVi2ZNVd9fKhODmRchFioYcUjxQuQA0WY5fThSsNqGaHCit7
% J2TOdc1aStbqvxJo4HN/aWokAKn3qZ23pZ46gq+b82Z/zptTGWiuGW9Gcjhgu5Fokh3TyOzAVumy
% V14cBXsgFp6G4s2L4nTuZd8KPw0m+g7IH1twJ7oMcP5+OvYkc5kE0JyWEMCVp7mqQLd43Wg/IJNJ
% T6WglGlKlUCT3E5Md4MK/mfISFKoFU38z+TxpJgnUOpDIRmS/YsuaBoohedF/jzwPCGanofB2POQ
% GgzKjf8CUEsDBBQAAAAIANKpVV2AMNHckAEAADsCAAAwAAAAb3JpZ2luYWxfY2hlY2tzL3JvdW5k
% MDMvMDNfaGlnaC92ZXJpZmljYXRpb24udHh0TVJNr9MwELz3V+wREGntfDVpyaGvfYKKSuQpSHDd
% OJtXC8cOtgP03z8nVSV83NmdnZ1xcxtao6QAVCr6g1Ziqwg66kl4kB1pL/1tB/WhaVZn3dFIei7C
% Z5yck6gji14ajQrc1Dry0amGkeyAeu6ifyNqF3C3gzIHgY4cDBVfr4s75TfbkY08aeinIMCi/gV1
% cwqDOIyKdvD09KFKWbJJMwbn+CNooy2FbeImgs6Fo67KNGd8kyTbjO3hpcqTIs2KLWNhLokLlrHl
% 7aGOXqq44CxPGC/L7f/gqiP/7hz9eF+lcRqgmMXZPGJGL4ewb0AvrlK/wmhNNwlf8bzcFHwPgqyX
% vRToqeIxDxrygm83PHDnC8cnfpd5MSLw0O8JVbAURvSe7GwMj7P9/egHhroDF5ZNCi04b2XIojWT
% 7tA+sjgoBXGWBzlSCzne/Q8iQ2+w2PSAC6X/a6Lr1ILRFI1GhkjcNDwoLhd4/nk4fofjl+fj12ap
% ruEk8VUb56VwgJaC4x78lZYPYua4ZgtMv169AVBLAwQUAAAACADSqVVdo1E2SBYJAABwFQAAKQAA
% AG9yaWdpbmFsX2NoZWNrcy9yb3VuZDAzLzAzX2hpZ2gvdmVyaWZ5LnB5rRhrc9pI8ju/YpLUrSUY
% ESRjxzGe1GEnvktdbhfHrlowpaMGMYTBehA9HIMv//26Z0ZIYPv2tupUBZpHd09Pv1tvXr0tsvTt
% VMZvRXxPVut8kcSHjdevX3964EFOZpJ/i5Msl0FG5klK8oUg+Y/EWRRTMhNzASByJuJc5us2GShs
% ckha5HodDdbtxq8JmYcJz2X8zVklMs4JzzKR5hLAeDwjcUJSEUoeB4LAElIXD6swSXmepGsyTZI7
% kgWpXOXtxtciPjUMknuRyvm6vVojp415mkRkngK/QDcjMlolaU4uzQIcSS4NTBEHeZKEW5gwLSYB
% DxZC78sceKvvr9JkVgR5w0yzdbRaI71s1bhlVod27B757ddPzHJh2ACBED6bWZxO7dNU5EUKh487
% fmsKf5SPXRy5voKLivAJXBPhHIRrIhw1a4Bm1joaN0jipcV3UKlCU7txkkZ7u82mp0k0PQWyEWli
% xVuQ8fhWKXdCZExSHn8TsOk/s6SQVyK1+vZpg8ATs1DEMFOTv26laf2axEIvIkYqAivi2Z1Bwkcy
% XGhPZT4JkiLOLbvamsMui0vmQLzbrYzdbsfI3bLOXUXcUImTXB374cPyF9dW9tYfS3+89F+x211o
% TRx1l1HUjIGjJef/ds/OlrZdMWmYyxq1CQIbK/huhGR2cDYeW0ozGxtMRXG/UdwnP7SkYYDzvm+X
% ao74nbAKer8j68Jueb0+MxpslJKQO5JQ7EuflbJDiAU1J4i4iAQ4mLDGQNyvyU2JlG52gQCLekhy
% gRJhm14fXjBhygo3dl0CfcX3jOe8xvcDC2WGmlgpAdDCtnvrvbV7WBuyrIisB7s3UoO1pjz01Cxv
% 5jr+IHMI88yywdioHd6cqh3wMdzbyJX1QNdwzMYz+wDxAoy+krG1PQB1r1SZSkq1EytBTEvr4Gzo
% DD1n05uykTOCgVodMLc1bI1aw+bI8ZqbljKFVGNc4Z7XGnnw3wQM2Pdam6ZGTBdgl60pzC0FY7em
% TRhwAMQ12LJhMm0ppDpVoxK4EB3SEU3pgF5RoNZQOoJINk1CGVhGSUOPjjy6oRuPwmXpV5at2hom
% sw6GHhl5YK4bj3AyJV8PjGoYMLJpcVAYMtOaPnvTr//nO2p6Oo1AHG5DvgDPtgbOlQN0bMY6CmCV
% QraxDq7NPQkPQ+eep5JPQ7GfuU7JoH99DbdSolEEM8wk1tbzOltTiCpHc+n7PefJhJjV9r29mFSw
% sXVpWcsWwtl/6dJ3NoW51yxXDoGivR/ZIoi7dSr3mkqF1aUnis4SBFTScTt/TOiJabCt5+7AGVmD
% BFEVZ8x92WnKp8+2wes5Ujo6MmYNMBZiYP5eLl2pKWiTMeBITeD9wWigfOIWc2tqnh98jmdiJWJU
% KPkbL7IMagoH4heokYckK6aZyJ2PAzw54jFCVWo+JY/xTxLwTGQkYm67fbJnDzxahaI0Bq1Ejx5i
% HPchseLUpYcdPT9uVDoyQNRyauMONeBq7FSTiopCeGFSO6uxJ+neHynUSH/I2IixS6vbOaTdo459
% 5moxMxPx/gcl1eEGSOt997gDjB2+OzKwV7h6fHjSPTp51+nASYfeSeeoox77bKAI/I5R5p88T+UD
% pEcYfzUKKzOlrYqBV2xJRJgJ0nkxQ5lqRT7NcX3bLwsRKO4iGfM4ZxbGjTV4aMd2frfbsGXtSKgO
% y+p8db3u+3fU63hwE43xhuhCWTwswOxyeS8gb+fBAipeGDzISG4Ucpt8gqJ1XW0uoIw8Yx4Rs28i
% aytaU5HlIDW3Sus17z2pBRMFCNQtHFC39YC1ittaw8veCUh3z+ObAmn5it2d7tKyNDG7aSG5O9/4
% dQBSYTWpNBEYq8ma0BQhZP/4PT0x9ZbC25Gg67lgJMcn7jvqgkkca1GeuUaYl0mROsE6CFHbzgcS
% Cj7XI+w43O2SVy51tOgAhWE2Ln99KKLGnk/78AeFrw0D1x931aDrw+aOvgEbtl/tpo7f0plInVzE
% 0DZA8gAp3pHB9ccyHpyS8/MmAx96Cz5EPnsU6s04FRBtNPsmhNSD1IA9Dn72wDUer+CFMe4R/n7u
% QaE1fgazZI81gQN8ssplBOS3FmSaE/aIkgcAFLacywDMnj3i5Cd49k4k+17wEPJdGcrekH4YkjAJ
% gGq5Z+JgAsYKwfCQkh8yX5DVAlcpWFMQFjM8XEKsncsYOiZo3EoDDpKwiOIMgt8tvYVgBfWnekOg
% c201uqWwpt4Y/FzbRLF6hg1U/xWX17MMUajFV4Ln7LBmxxCLA9ClRkMcxPV797jsPll+LintxzoT
% 4rCx2+aYN0rtGdy6CHlayR8NnFfroP7tFhqLA3LBniDnBmAKkr5rb+8JfFAldeRwbI2VsHw6VqJx
% bZ/epIWw6ZM+RT+WTjueyUJ0b3rJIWDafiWqP5Mn6rkC0z1edFQOLNcZYnhwRvYHVi84a1gW5IQr
% EKe6XuNPJP8dR/iya5grnoM7mHTd065Y7iGVrSIySCgQkafQV854ulfdIfFAriAS9el1rQM2HV8V
% da/9qrG61g3ZMpHQ7NLzvea3F6nBud27MH1ZK3Lcl3qz/9LAXmgemOFFQcoHNo4d12+phsmAUn3C
% s0dELx8RqSMegLqPr6Xvqz5mb42e6+N32ogLLb2FSMVM5iDViRKGOcsYl7bBbr3O6YKH05317dAc
% cF7iaoSjshQ6qsDVsCRzVCevoJ0KpszJl2kyFbEsMvU15F8eOYPc8a3gIIdcQHgTPFiQIE2yTHsl
% kRm6sjYc6Ody8+2o8tYb5aWgfP+07D2gi1T2f6NbY11vLypxe08qcI+iodxAq2gS3gXb2lRPfQep
% dxo8q6Vv9+wMcS/smnqv2Vjuq98AYXrX3z7kL67fu2GV3V+A3W9JzKp64xlbMtyezprsRn9SwG8T
% JfCAod/e2L2yEISYC+1X5dI39hmbNQeAoy+3X7Rj9nlUOz8rv8RiPcKaMMA0NAe1oKeX3xwhppoP
% iSD+yrPhtpNJzCMxmTB2MJlEXMaTyYEWVdXr9uq9XW9b2PeqxNh7YuH1kqD/5Qv5NOxf3JCLv3+6
% +Me1Or9NPtY+kvJU4Icn9SkTe84EqwjMZcm8DZz+B1BLAwQUAAAACADSqVVd/loFGe8BAAA5AwAA
% MQAAAG9yaWdpbmFsX2NoZWNrcy9yb3VuZDAzLzA0X3hoaWdoL3ZlcmlmaWNhdGlvbi50eHRdkkGP
% 2jAQhe/8ijmClECgwLKtOETZLEVFUDUc2uPEcZIRjh3ZzgL/vk5IKrbHyHnfvHnz4t9hdIYo/nXe
% v++j8Bwno59hknyFVNkSNFpSEgXU6L6SUwKMa0s5MbTcAMoMDN18dmeCA7e4nQcvs3mwmfaQTma5
% 4IapmmQBtVZZw6z5BkLJohcaiw7qXlPVOOKX2fJBvlepEsQg40xzNO0fnaGB3kjKla4gpw8+eLgh
% s8BlViuS1oNcNbp/sqXmplQi8zq6xowaA/PZYr4OXoMgGKgPhEZ58e1VPRNUaqx27l0kYJUDOutN
% 7sIgLu1zMjB+ZLHazNbL18lAzhshoELLSreKf+VUlBZqJe5SVYTCuFkalM64NrD0Vt76s6V/SqGM
% gUJTBhlhIZWxxAyMB2W33mryWTzEU6M0zr/xXKrScECdktWo77DTWLkZ5mIeAfX78AwKrB/equ3a
% e/E2n8lKcr9LG8ghLVlyzbiSOzz2M1poh0SwmlAWzgnJurHT0ej9cArP++MO3vbh7nhKzvsogdPx
% 8AfGUtm2Lyqf9PMWi1UAhvPMeWrd+qxEksBKzi6uURXeqGqqsSiNr0sz2frr6Zz7bRtH4eEA0fc4
% +pFAy4rfphA/cvm/zzlJchd8Tha1OzSv0ZXPTcZUffDp6C9QSwMEFAAAAAgA0qlVXTSD9+Y/DAAA
% Wh0AACoAAABvcmlnaW5hbF9jaGVja3Mvcm91bmQwMy8wNF94aGlnaC92ZXJpZnkucHmdWVFz2kgS
% fudXzGYr65EYYSTATuwoddiJc6nLJltrl+tuOWprgAHLEZJWEga8tf/9vp6RhAR4H44H0Iy6v+nu
% 6e7pHn784XSVpaeTIDpV0RNLtvlDHPVar169+riR05xNH9T0e8bmQZrllyxTaqZmbB7GMg+iBZsF
% chHFWR5MM6Y2SRhMgzzcgiyRqczVrNO6e1BsFQVPKs1kyJI0jucZk6liQcQSmai0k6tNh32NQTWJ
% VxHBx0keLINnLBFHLMhYuoo6rV/VH6sgVRn7RYvIeoLdbpe/bAWT0Yx9XeGxw+62STDFOuAAhCJm
% yeZqDYmmcTTLOqRYa57GSzZPoR4WyFiwTOI0ZzfFBJMZuzE0Qa7SPI7DimYaL2EpLVhmSJYyfyjf
% zgEQp4EMW8VEtl0mW8LLklbrR/ZJrrIskNApytUCFhEsVQl0UhjP2GQLi8DOLJ6XSpaEndZMzZmc
% zbgUE+sCbPkqjRiXo+64PcGXkCOXntyxpUmz1eQIqVOROjvS5So8QmoTqUOkNpGKYg6LFHPdgh12
% feRyj1toTkMQxemyTqCRtOAaib5av/m8K7rWJfv29aPPXTxqVrjHUkawDh9aFy2GT+SHKsJIDwKA
% +373ogAGr56OV7n/m36axyk2KPtOzpbKaKEA7b57FxVg9En8kq1kSOM10Q93NPR5KiDrlI873Dpk
% +YF4tPj7948/uRdPPu3ekwD66HFsNYgTn/YgEU+7afDyyCH2ziTIf58iMnJuWa+9i8TnTqJNnGgT
% lxykNS2BX5GY6cIumGlpc6ZKhjBpZcwnv27gS5llcHf2BFhYtY7whPXMhqRBNA0SGfKhCBBPVmn7
% 0Wg4CsbQbGcYem+GQTXUIA9yBvarGm/pgpoaD0T/HCRcpmKSWgZEP5cviN2ALVK55Pc1rGy1XHAC
% 1I65scTWwG7EtoLFUmZSr3RfLGCei/ABSKb+KAz1XPMnTYZXF8/a3M9CNoz9bEw93U5DeIWgBOkP
% hJotVOZ/jSNVIJoJ/c0AOuIefH5sY89dqz1CMHjQjwiH/mj0m17497q3jY9MVSIGx/1S71Aw9jkJ
% RdH26POg7VqvoybJ49gIhufLIUbEo41ZzhqFC3eRYagB3/swGpeTjBcglCHa9TG89TBsyNMff/AD
% 61Dyhl2Hxq5It9MHnDxZpVm58fkqgcWDQ5hdHAY/udZBEO/lB8oOrb0oJhWjOOcVij5yzIgXJrSA
% fUSFMmziWc2dEt89dKfE9mVdocQo/DNcuHRuMrCG2sCiu8VuzcK39LyzD+XJjWUViRTnyloFi4c8
% I7xaKt00rDwiZJu7bVpq0wjnxoYhqec0yQPBA4f0F6UdrFPviCGMHearMHxJjEtKYKMbTmEQVebZ
% U6ruzpT9IcXtBXwTxchKVW9wai/9o4Y6pbVurUa+N68usDyI2z4xH2RPyns5P/n47+H1Hbv++Ovd
% 55vP18O7j7cnVisXmZ8lHZz1E1QK/ATH/omgVOvfpStltYo4AQkqJDgO77+1czt3Bn07b7sDh8Y8
% d/D+V11YILl656L/1rJs22vXp8/0rEXJ+QDUO7czO3M8z87awMSQZw1M1xXe+SGm29fTe6CzIJum
% qMEiCbWb8ooctI7Xf4m6IYjIiHhw1lo3TLQ+OWIW7jo5iQf3c9aWs6ZfWhfPfO3ot/S6LukNzvO+
% 1e7bN3xAavg+Zrrnwu2+sco9+2V4e3vBJjHKtLRQGuGA0e23WzYFSjBHyZgjD1NQZ8HG0bmbqVz6
% wDoFVgfionyjUpY4HXKtFWrjYIaTM8i3VGYmcbiN4iWKP+1UCgXvFlVzhBKSwekWWBBroMy7BFKQ
% Z2WxuKtyUSWDc0nemNNKVCp3WhosDBsVjFvmkaRh1aR7cdLO8pSDHrFovDx8yEwoJIkJbYce2mUu
% rgUpuMqUSzxUYxTkDqjtQ5BDCFfsQA72F5I4QNb7t2wIvjwRRaGrQ0YkcRbkaBkOAihDXR0G8y1f
% wk1s3nOW1innSyQgPe5jbDnLUw4fosmGV/NzIrvhnqUpjbP04D5NR9GukatQZdM4oRan2OzskoVx
% tCicI8tlqhsg3bKw3mnfeI/WKJiymZoiA2REoZ1Oe1Dqc9LhjzTnAwvCnHotcrK6XgM7te3eKRJw
% aqfWUdXB4vCeXUMaUMqtadofvOu/3YVIX2stzi2CpYjxygkwua1UzoJV5nu2CSb8ICMAkp9DErtq
% afj5TpyCxddg7ln3bbfb3TNj6cxzbGMZULqfVNEsiUEo4DurtHiVP6AJeojDmWnmDD5zT0vwzi5f
% YGuxbZ4YvLXatH8IhiLsB2/EGfLje7d1fU1G/VnmabDhqAhr6a4nBpaoT/QxMRYjZ39O7HONx5UM
% 19f29XXnzqdV1FbBpZrKG00REt+dfB3XFY0niM+VaTPzWMd5tpojAQXIJPVsxLjJP4M3p6SU9p/N
% phE1m5P9QKFojARcc6H06TyCJpsNHeebDVzEEsiS5ZhGZ8WoZ9N7TCK9nvZQtJj0Mo3VfO7Xz+zR
% ZoPD+aD828X4k2NWJ3fUyeGJ5NBAezYi2Opgd8wCtSSaae44naH5ZdgPcVYmYSTRIrteoTNX6KbJ
% ZR5kOHfKCapTmcxzFa1MdiVrYyOq/K83IzPJNdKWwhKF1scb/zLlo5efIvQVtfGafEM1C/Kp6MEw
% x+oebUm/WflU5Y1uSSpArssfYWKUfuCZqUqUzP164VNYnipD52c+kvZkv29Cr2ONLVORTwPK19ui
% wJ4GYhuUZFOQHXXdal/COMvQYgWzxk0PL/dFm35gvLNhS3EmzsWbwqJDv+iGrJ3F9KWAbgyKunFs
% UdZ5yYDP5gyrGt7nRdW2HzY8eEn93/NCNBeou+0N3zXDYmBj73yfDIqc4tnPC+qv9Wy1RKSk7j5H
% N3BOJSjgxrstuUeXhhTU1UUpUSoQMerpUBe/7jkufmwtR9/B1mKDsYweDzAev6Q2fa78osWtZh4m
% xhpXpp+DrC/1TmUBsNO27L2N0vbDZEeYSP/QKtXbbbVrV81d20lhX1XTu4fu2Po77daLg/uX5mau
% zWauF+LqcDPp8yN7ttf6agMapDhX1gFOcF3vlUe3voKSYfBMF5Il3/O69JL1wtqPLTKa/3exZfxk
% XfpJw5A7jHd+Ihv9ynv/rHlPNPWbrVlxp3VmHGjToJ2pMJf60KX6tknpOvrqYFff7HWxhUyJdCDW
% e19D2f9veijPc6T8jG5CBUSLIIRMJ0GeSqTmT3BZ3SBn5kQvTjY1YwuZmNy+9HWOKNL6t0g5ui6g
% JpcuQGWtFCcEiU2MJjKU0RQoONnhIKES5ohg2n6d1tAtUs2ZdTn0iueeuYTpF5cwI329KIrvrkCS
% tehuNlyhUJc4XHC25GpDMgxdOqT11Xd91sNCx25l3hzeymCqtefTZ0Um32usMV1evgxd87DP2TvO
% 2dOc7QFY8KUvpWpjMfQMmtXSSapX5KQ+cpLOTCYbITu9kItIhSoHXV77ZQZpfXD92nWgK8Igy7nh
% GaD0vvzg1d97YoRTbVdD7RKN71fPcODy+YNntavn2vzQO4C4rkFo6VyxW/iqLteZRR+7TvzBaxCP
% 4JRj0LQbNC8CDg4Bh3uA+iwkyEOx3+00L/qrl1P6m+OhGFdxUzSjARpZnQFlEZT6fstEUBk1QE1W
% uY68AvK/0c2Xb8O7z18/sQ+fh5++fru9+3x9y759/fIfRldfphu1wFHUQ9Gq+GcjSlpptPCjpAM5
% Z/GyM1NzuQrz3zHLva535nY9NAfrGEHkO/qPI34SRPMT61JfaqNtqSqHwp+FW/d0ErrW/va6Vv3Y
% xSodk915hvzu8wg9jNV2H+3jb+rZWCO/9ugfhPvRheiO2yi4kTzN/GBHe39K+oUBSkeDye+F3ASZ
% 71rgo1vd+nF939GXpdY/7jt31fSMEGSayi0fLeWGdzvCdehyFDsuqkONrjeOnJY7SdbVyV8x/d3Z
% T6++v/wnxSwjqXAWqFzxmfjePDLoFoDzqOgzZ5nVofPC3CDNsmKAHrbOo9LUJ630DQ4AGi+ND5D2
% +kkYZwCLVf33gME7TzluE1U7Cjan8NZ5EQF/6vm/yj8l6cxxpg+SKmr9p+UlfH8TLCFndfXwp175
% otNb/EXuX3n/8MsXdv3Pj9f/umWE/fFDhxV/f+5fEBVXOvVymP7NrP7xZHISPykC/x9QSwMEFAAA
% AAgA0qlVXYdGzDxFBQAAsw4AADIAAABvcmlnaW5hbF9jaGVja3Mvcm91bmQwMy8wNV9tYXgvZXhh
% Y3RfYXJpdGhtZXRpYy5weaVXS2/bRhC+81dMBTQlXUqVHPRSWD41TnOJiyYnC4qxJFfSSuRS2F3a
% ktEf35l98KFHcqggQPuY+eY9sxqNRh8OLDfwkTVaCybHihlRS1ZCY0QpjOAaFG80L2Cl6gr+qRtZ
% wG0KH6RRR5j+PomirxuuOTAlzKbiRuQp7LmqmOTSpPBRsSoFhkx/f/lzvBcvtQFVI7hEZKY45PVe
% ePTIbJih1RuXwEnALxpeuBKr42R/nMAno6EuCygEW8taoyhtkSsmJBRKIClBRkIa5LZmlEeoK2EM
% Lybwmb96/d8PIVALBOgERaPRKLLmrhQ6B4E0iGpfKwMP/gCYhgdP08jc1HXZ0iheNDl3lzW6gpla
% hbuqKaPoCeYQT1OYJtEjLWd2+cmfzpIoigq+AlYUMUshS/6IAD+Km0ah4MV0Cb9Chj/o18XMbWZL
% zyT5OmZDjjGztGPWUekmuwTtBBJElgQtjKi4vqbHjdUDxk6RG6tI2l1Z5dqraZCe13J7quS5jrJW
% 1e0lMri5gdsWGNeeYa9qdLyJD3rI4+IRo+tTOOieg/MfcFjTHc9j4MlKJnexHNIvFk+wwiA/Ux4p
% JtccKZYXjlpNhczFnpXkVyELkXN9isgWYrnYOpAtgXg6dyL6Jx510zqLac0x2TD9YwuDX3TVfA7T
% jtnpVHKJXBjsnuzglcA6XV7l8uZguT8X+1a8xEx2JHZb7HG/eKIkiGdwdwcynBP0HB7tjkRUTO86
% KZ7WYdJHIDGRTDJhnnOsZBMn3eWK7ucoHIEQmgitzU8dAH0w+bD5NLw9bB0cojSkF5SJxqn2zhuw
% TWzfCSH66UyI9cLBeH3h35bvjAo1RUJyA1Wf36W+6oIZaRA1DBRej9tysUFQRwz8tTiQJTYJOl09
% kIuAqQ32fbTF7nRT6TZu0h7tFX/Bo2kbr36wsK5O47VW7Ij0O/gG8Q7u76ny2jBsiA1vLdE3i91e
% bvHcEdhQow1rs4kT7DOzlmZfaxxPLwSRYfeNLc47z5YMoiuuR5esRNd69/td625yWSuHlzjlbHsd
% xMJ6Dn1PKofKIZxBZsbYha2l/dRN4Ge4HarjgUgKLXsSfGxIS7tOYUjgY0My+iliaUOG6OKZ07Rv
% E2QUxv9fiIVWMglf8k2jwM5p7QY5Tk8h1/DGVe3PbfJrPGxKmvrHCc3LCwn3Hd//oOhCA/NhwKR1
% MwM3eOTQM8rORfwQH7CNJCnQYoZRsdgHi12/utaFC9qz5aXM7cv2crPFbklf3zTb64KSzd9Nl/0A
% F3DXr6teBB4Y5s2QdFiD/pjJY2yxfUM5ddEOJx4WmUyS81ZzUdTVZicugA4xr4k+l+wT1rWrjMJF
% 3crbcd7tshBQ+wLxuxRsstMw+A0Kl9jketwM+91XhZa4XMY8r+KX05np5zYWCWFanQ6pS51jSIwU
% jmTZm9jHTQqvCWbO04miRGcT5sWlT+PW4dVwzEse82JNbwN6R9I7c/65lvyk7VqSpD+QJdzP4b07
% oVTyb4muWFIgJhLHZVPR09FL6rmemmMsKCTUQGSXu8GbvkqQjFi9/bRsx0DQGoQGUrwDb2+wtmYo
% AztZzDIdeh5VGr28eiezC2H+Ti7JhBTYUoqLy4xnXWPZ908xdE7Qt1/D7tlCXirsw7r/fAytsMb/
% CM/WOHyA0eMWf8IDl1FsBk+0hbjyAnKvDrKF0VN3gT++OWUDjIwwMkdyCSrrQ2UeQ6ZQhYbKME0t
% Yebu8i59ELSiyfh/e27e1ibzsb0EWF0HrM4AaepRFiFaWG7DrL1yi6HwbWEQOPxTyShMWfQfUEsD
% BBQAAAAIANKpVV3Ouh2YbAIAANwDAAAvAAAAb3JpZ2luYWxfY2hlY2tzL3JvdW5kMDMvMDVfbWF4
% L3ZlcmlmaWNhdGlvbi50eHRVU8tu2zAQvPMr9pgAlqKHJbkOcnAcowUSFEacQ28FTa4kNhSpklQS
% t8i/dyU7SXMRIO4Od2Z2eG8HIyGfwcYEd4CkWAK+cBHgKx+8V9xEjgdlDdfwhE7VSky/bIcol1Ak
% eZZkZZpkSQzfLdTaUtk0UW+VCdDx4NQLBPTBx+yhRY9QK6MCglS8MdYHJTxIC8YGoC6+18q30Cof
% rKNRmgpPqMMhZuuD0AjWSXQe8jhOk0u4/fk3m3WvUFsH3VUWx9UlNHpA+dYXWmeHpoU0jdnmpddK
% qHBE5a+jzq7XuITt2ep8Bu3ZNX23p+8KLFyfwxWclWUyg7TMZlDk8xlURVHOzz9uu+Y9D9GOXEQH
% Emsk764gT/OiZG/89tqKRw/PKpA2dChV4GS2RI2jmSBapPoSSoa/h0m0MdiQk08IvbM93dw7ZYTq
% qSgxoOuU4SYQJJ0XFasHrWEn2sGR8c/Uvd3dXATbX1Cb6oYOaA2k9H1OtjhCGnKnj7hucO84jHM4
% +Q4CXThuGkFJNEEFheOstGCKdPZoxlO42V7cH/xIjghxMx7xxiF2OFFbFF+YlTIS0+b6ltP2Hdaa
% HCLR1DDPmTUY1RQ3mvrJSE6hXLdWk0HvrNMyZe8cR+AxZCeGhwkjlf81nQrbdUPgn+xNF2xiEe25
% 5kZQTAhKpBo+NGQOVaipYI6bx4juh2kZ483KUDQJMHJeME+55OTSJOu08XE2pV20FH4Yc6yOwP/c
% qxLmqTp8YO94r7mgN/Y2u2J/0NlofBvTizs5c6pmbHV3B5sfq/UDrL9t1rc72K52u80N22jee5Lj
% UVgjx9Y4qyr2D1BLAwQUAAAACADSqVVdOw0HTawSAAB2OgAAKAAAAG9yaWdpbmFsX2NoZWNrcy9y
% b3VuZDAzLzA1X21heC92ZXJpZnkucHmtW/tz28bx/51/xVmZtkAD0nzIzneUKDOOY7dpM7EnSme+
% Yw6HcwCOJCS8jAf1cP2/d3fvgTsAlJy0msQigbu9vX1+du/01bPnbV09D5P8uciPrLxvDkW+mpyd
% nf1atHnMVgF7kzfVPZu/uGDijkcN2yV50ggWJ3yfF3WTRHXA8qJhzUEwnvP0Hh6xsiqK3WwyeZen
% 9+w9Uf1LzeqG5zGvYpYmYcWBalKzthbxjL1KUxYVWZmKO5a3WSiqmvFKsL/xtq4Tnk8q3iQFUK9Z
% JcpK1CJvRMzCe1bypKpnTLK7NNz+BecnzSETwE7ASlFlPIdJweRvFc8C4DRmV9GhraZlcgTuDyIt
% 9aKVQKZYkssdbztCs/J+NvmlYLu0AH7y/bQskryxVmJFhfyLKol4ykSyB6nytBWsEXVTsyKK2grW
% qsQMZTzZVUXGdhUsApurWZKVRdWwt+oB4zV7K8eAxKumKFIzBoQFSiOh1HJ7bSO/yQkVbBDnydG/
% 0jf5pkkyoZ/DvN02AtEBefm2v2M90psw+PkQsHcB+wnEF8cBq9sQVC/2AdEENqIiv0ZjqLIlfAEb
% iNuoCZj+QCTClOc3+CzJo6TkacAOkv9tXKoP1X0tKvhcx1tiJ2B70ll0H6VCUiG5b+s2Cyb+ZFLl
% e3apNum9mK+W8+XLxXw59ye0txpefvo8mUxisWMNT9N770bcA6OXC/+CyMlha3i6gbHy22wvGjlu
% 7rOvWa7mx8kxiYXHAxaqyUeYQgLAhygCL/R9ehPDG5IGPKEHvIadNSymL5Vo2ipnb73jer4JWOwH
% 9HlBn9VqNZgRLVapxdQkDlPYXxlIicME/KTZE6DKDCwjbzyu5oClaT9iIk0yZTfsFpTMquKWkQuA
% OX/LbovqpqYXRQt2zd5f/QiqEllSixkaLJLLYVepyIG8VCh8Xa+9t94d8ER7uAOWfJ/twBfu0Itg
% iQ19w7XgO99IYeSomHf0Gd/e0Fie74WXK85J1eSfIEhx13heQkOTbijYUu6zZMfCdbJZ32zYs0v2
% Adj4pciFb4jAe0kHIg6+6chbQv0wHA7EbtyxIawBul/Te7QW/TGgV85YuUPwEA8+dczIx8pk0H1x
% IvzXjejvEexvgft0WbE2fQmbdl9Ks85Bra1wXuxgbWXEavYIA5qJ6yeZkDIBOtcoDAgJnvqmooK3
% U+SvN74/GcxC3pXgtWlDCJO2DMEATT/JY1EK+CdvLn+rWqEYKGGiihvaFJV7legSIJH5REnJotAx
% rwfjyC7saFL4I0PFmTWb/fj++a84rEsojO8rITL4VJ/59j5KcAe1kQOPeQZ5z44aatR6LaV0F7B7
% 5TLwCaX+kMC+wMHDyt8M9CJfmGFIVi9WJ/vcK91VpgsUA4RLrwShs++BuetN38rQp8uejk4ZQyKN
% QU3x2Z8g9Yq0FmyhuMBkXtTCg5j+0eWlaSHLAx8fgZNTPOjwl+RHSMrCbOcjxhqKfHKgiR1JILkT
% lHx5002hadI4kyEbH/VKHDJz1qY4EzIbhX1K3IpKrHMDJD4v7lIDRho5TtuaHuk4ZCxS9cKIR+x2
% SZSg2dhx+o3mnxWInOq2xOwLYMTO8AHinqhI2yyfNsUUoyoYM6gdrbAL0wHDGK5jdYApUPuNlCGk
% NMncjh2TOmk8SN6IfAJ2C8jl0FgShI1d455y1/lhgbXSp4/SkfNGYuupyDYS0hBLEv76M/MW7Lvv
% WOITWuMqyDwbDXWlrWL7R+1M2ivR/behq0OUZDvQS6D5dVPnAAAA9ji+DfvuLBQUgWrZZryJDpBF
% +84HUQTMvSRzBxEmT9i8JrOVONhQU/FqfEF7PePmzx5ZjD1/zpb9FU9Z5U9WBKwEQNg6OQrjawgl
% ih1itPJgaNVfZIdpUjeetgS/b45tftoayUro/X9rkQmMAkikVuvCHxQZyLV8zKbsU/J54hoVjtDc
% gTuR9cB/814ip8BUkxvTlKHJP2XaCSRTZeDgoqdsnBgGPq8/W0xJ+NkZ9qPEMYrapl8LSzsBW5zw
% AcLlW4Xyjd2oYZF5vpZZA5OFNAWU1ph98s4ZKrFLhUX0AeF5bzn5AgoKqGsgQIe19yCR6PQBgECH
% USFOri/GgCho4KEn+xBmYpDUME19DzRykssF7MESaQgDCHmoIkBOcZCJy3roYzyQpOxRUmfyOQ1R
% DxxSGFVc5XYjFQIbaHtc1v3cblRgaTFU+ihK9Pii2kbAQwKxAhMmwI8ijre0ucu3UKODYMCjiqqW
% 37pA8hpQAVbYAtL6PUSMoi2nPN2LsOJ2SkResJ0QpRxEAIX6uyuWYPRJmnsTVzR+01WitoReeaLn
% oSopLnThhqhEpC8n7skXGcaFE5GxrwgeESQoERCgGrxbqhpRsPBM2lo2gyI+qz0qVm4dEhixsQik
% CUcaHc0IVsBwXwHU3rJyzuKRORL+xhzxvlMVOqUoJ36RmgKOEDQML8C8Jq2Y93FjkgAG8M4/cVfG
% CiQW5Cf0pNJVTBQc1kpNeqgV9SoVO4zKn8oLqquO/gk+PyvtGlSHEwNjDQHzVHUhLRR3pQyeAqdc
% QVbfx0Bj6t5i5enVKgkouuWQWdiw30HW0kRpxzq6GHRjG6ABA3YSvCFghKnwBhG4mzoe48ZbgpY1
% 6pn30sJwIrDqTcem2DrvVn+AuFiOBOqg/8TaTaOwNc31+0Vaoys62nCD5vm9NnAqdI13dzkmH2JN
% i+4jwpFah0fLRydEUqnN41RVARP5I6O1B1IJSoUFV6rxTBaDhK6yma0mVZmC7KfS38oDr4X2RqwR
% zhwvl15zyRRXXh4w6Q/0QA5VNHdtmvYisw76zAr6Wl6J0EuBLcigf+Ggnzp5EKfxPqkoj2scYbc0
% TZQOiMBIzyHuxQ7TUFQtg3oE7MTl6UnlyVlaTyhAoIAWGDv2GA/ssacn8bHlKfhpDiYA+wPsDC5Q
% YhdBr29z1esj8NL0Ag9tZZBQafCzXRp6bgr33VaDhcVqLEfWzaka3OR/R1MUqspSzqjVF8Ud2s22
% gwcHEd0YXrHdpbmXzAe4YabzWtBLUNg5IfEVVKyusavkldTJQalfy/SChRlMvRvpWgWy8dh1AwDw
% +RvdKVDor3tbqw5LWDRNkZkl78BTkM2xdUex1R9dd5Aia3/8OQjkxBvJ+tCT1RlHcQvW9v7qx+dA
% 4TnIOcnajDVVAiGTka7Q6BTkpjb6FlvuGCfg680ltkwhDRRQFID1bmEYHQI58A4HIiLAX7DfjN95
% 2C6kQtM3jfL12qvy/QwXSUDX0yWGWVhl8MwWMMpv2xkoLqGkurVDixImrELMH01cGjB+8aX9CFXT
% qeSkvvW7AVz3K3liLL9r9I3gVIoalu6UU8CkMoT/eWh5Bv4O1W+3j+hYwpx9d0kz6RekkAOWF/pL
% GdqGAcqc7mBh2PoPvOTN9KqF2r4irl4fijQFmzdGYXc0Dd0pLmUdUfBKYoDh7gGt/73Iir3IRdHW
% KO4MmwWI7uVEO/EWLKn0QR+9BuTD+LGAIfGp44cRGJp37UgJQdNDbSRqbeGERMHl9zzJVTfbjkzd
% 6JKGByO9Uokeeqg2oDVdjall3Eyintn5JCMc/FQxkhZAFXc5V7aUVPjIPBggTkn2j6NOQpzEuesx
% X7E3WIWYFiZqOhcCNGjIsqzIiyzBsz/ErmCULH/GUGB1j5SylDbH3kaVRNgJJd8pGjQaNGeQWQ6K
% rgUU3TGTZj1zyMToTib1hbITN9J6yPs4WOkpdI4T9I+Srmm0eMqq4tBpANkgDGd8P0YHFPf1JX3q
% d4ZucOHlSE9I69eepzH0F+BuTLE9V1F2kqDvwZqwJXtb/c5Fv7Mh55nd6S2j413aZv01UdabltP6
% M76Xy5s9GqoqfqnAIZEvhCDwfLIKY191w8MkxaLfQakyXgkw6Vj6UK2xYnesJVuVqkVpV3s9pxs6
% zsdhP9NJMo+eUMhTCmRFtomJq/V1d/DT9TU/Ul/zeKIVV5j2sB2FT+6KD/eBrz+qpkU4fN2rc4By
% 4B726KY6lM23/omO4T5tKZD1swXUVBy/UgLUR/w4ADBXaDoPgEpuAjDaVGUCLK4oI3D1IdTH7a2r
% kxvFT9h/DgZOR5tq2hyX70oDHjBrnKbhDgpDPWhBrKlRdzDINTggxltci97f99+HuNN2vbjYmDEP
% AxpzpLG+mC66QbcDQiESMu/vIL8+3MrepLEQOnQMnEcPQa83dXePHmwPuTeYWw1Buu6QW+yMOsZG
% l2DwuPJ2aFKWOSGbpXOk1mPDcsLoBBTuPYf8FXU9NAIAhKnoVgc9m5tnoQOgTd2LgEne9DHdRMrW
% SX1NT8H+dQU2CpwiZfY3y8zL1MYRX2QYDg1opZsxntIXFfSYYNbywo03XchyB8+jpgv88NZbBewF
% 3bg4xw8a/2IbmlrXeDVENqS9TJ5/zTulWA3fpXuqzyHuyPMAXGBuNSyQ8rWkfK3WeGe9Xai3kve1
% dw3ustTtfvnQOsomGl27XE0fBdYZJA9tMBDbK1U5mqtLnpJZheXD6Ztr+hLM1GDMo6hkPwO+njlE
% rgCyXDDrLtGMDe99QcKpkjt5t2vmzv/tICCnDu/Ksbgg0CMwR0EIOrADxKFC3hfLi6NIscXtkHpN
% ea6oYryftprNFvNv2T+3n5ZB9pk0mV0uZ7NvvsWgCvBIjWsOVdHuD2yxQGoTCadeqxt2yrgAXUHZ
% eABpI9xO9nSTrK0EgKok43tIzSA+lfJreWXuQVRFPTNGlHdGBPJeLHppBKQDlM2Ql700IuK9NPE/
% UBKyk9Vh3rsuAShH8jHEqg4T7+SR+Oh7aevkg9OFf5L84r8i/5ObY3UD1aMhPUSJRc/TZzFhUaQe
% QbzRUmXsR53bdDNlx5V2+N1AB6Es+Af9giWEGznnT+CIjy49bCrQ2np2f8V+jT3p115jpejgigEW
% xufOzlZD1Y13tHoVyr+A5WkGLpJCdYt7n0KyYNTwo5qW7tjJ8CO9zq1NujCP5h4diiQS3omQv/Gf
% MnVcH5dXDZD1upVFTkutOlpqvDRxparI+MOiRyXGwT6BPt7tjSTMlpK5gjjZdjg9xKgMAeWCKjoK
% OfLFz7xMeQRhGaF4lLYxIniwYfn6dLCZ2ynLOMtaSY28LWDrJX0YNt2dfiPd1umg9u/yK/ekczJu
% pz0fWVlGPW6zj0ww9tvrpD9urbp0cpViZM8iNI5Oea+qMGnoira+lU06Rx+VJ7ToN/k+Fery93Vb
% N5TmMK1iUV5I4Xbao74gKBDx8RJ39GiqeNGvOHqBZiWt1DXmpxWH1weNxugu4eR3evxpva6GKepp
% 3apJSuq/FHkEkQIiSQhFBewBsEOYFpG6ksv2UDwmOeZnrRXr9MfylMzBdv/XcxSJQp8wd/8x7ax6
% 2jktFEK4sEedDwJGUnepxyIVjbxwL+nTrL4FcLcoA9++7gNZOU9eVnt2qeluvsRKuONg2Ujb5csc
% TAMzW2143z9OGvQmYskpE5Tq39yVKbhWw3gWkoOZ4qPrj3bhEWEYpS/jz2UiIm0C8jMmle6AVAZH
% J7FsLFDgvH83OgTNZjWcYsXbcx1vz61hsqrRAGRj39jkvX4Isd2LCyGlLnqzTi42rkrw2HfYSjhl
% mpFrmrpzsHRsMjrI1q9D1vmimMHLQOrjAo9n1efzjT/pr4skR9bGxzBx5az/FftHgbdfbrFM0NGA
% jOpcGxWUM6k4ihyrFygptNbDZN/nG1UGumDy94AxmDFkCx46QlFc/QCRDPJ7rNqtkhPbJDnyCz5+
% xEihMwFrc2n3e97uxWgqh5W+sXQ+ckY0l/XAfAB/VqcPgcYOgv63eIvviVN1i9+98q4KYDwpotpW
% fccLeOPnhs4Sw770pg9fes3IQWuE78eRyH6ILAYR+ffgi/0g/tFWp6FrLaT9PsT4gDHMGArUogJi
% pahUCx+FgEGMClTV8K0hFEJlW7Ajz6FyluZEoVA3T85Vf2OEVxw38AB86IrkvC8SC1qa6Phu0/3p
% y4nJY32Wc//CvSx4Lf9+RPdSPkxOQvLQAuNK0sj71MhPSU3KGB1Gi/k9NhFAE1CDIGQIZOahbGsu
% 9ab3s4mNDla+OU9Ymz/5GriHEYEqNmDY3ab7q59mY13PH0d2oXO6GvZOV+VlNuso9GAOQr/gQFTP
% hrL95UvgePESIs2LFajomxcvXp7re/v1UzJfLVYvXtrtGJOpKSqvPmvJXrD33ivi8gf497369xUr
% 2A9A6Eyv54/Scs5h1ZEGzqHVOzneiHvrMrH8OzX7T6WI6u7sE4z7fME+WX/W9tntKb36+Wf25v9f
% vf6Nvf77m9f/vGLvX11dvfnRGbQ7e5PyEi8/y8M1SDyf3K4bm8p23MVstftMBywQPLbbnGdiu0Xh
% nW232Lvbbs8kk7KRN/kPUEsDBBQAAAAIANKpVV0wSYD4qgIAAFYFAAAxAAAAb3JpZ2luYWxfY2hl
% Y2tzL3JvdW5kMDMvMDZfdWx0cmEvdmVyaWZpY2F0aW9uLnR4dHVUy27bMBC8+yv2GCN2xLdEBQkQ
% oJcCPQRNb0UPjEzbRCXKJek8/r5LmY4aNDFgU+RyPDuzu/pukhu96eEwJuuTM30Lg/MXF1o9XXKt
% j8tKELK60PTpkpG85YQsF/cmmL63/foQxs2xS9CNwdsAsTfd79guAC7ICsiyBUpFpZhcAW0Er2pJ
% SAmKimqMc14xiVumeUXnKMVQXnTmy2f5+htIMVrViuQLhJwe5zsFOp3Q0xb/ppmOHtxw6O3LuhuH
% wxhd1p7lusHA+IT5oyiIr8NgU3BdETXJ8Te0BYkKWGZVmua8rsGkZBDud/BzEow6lit4e/w1IVkL
% TFZqBZwqdOA9ik1X6YziBcVbEEo0hJBK4qfOymqK/P+x8nd4UfACbWowz1w+YNl8tPejjOWMlQUr
% 0TKJBRF4prn+iDS7MANVASpMmkhWWGmj6Ke09YyuC7pGWq4aXsmpskwQ/ZnNzYxuCrpBwUQ39VRq
% 1tSSI/gjrJ6xGrF3WHNG1bnaYPwGBGdgtlvnLeD3z9H02Cs2tnB/9/Cw+GbCzq637hE75tFEu+7w
% By8muzvPQC6elqSuhaKCMM1ojbYIJt+Bnc/Dk5uwGy3SdQ5nMLer8GtKJdwCge0YwN/e6ML91e9t
% cMlu4EsD9sXkfsbOJPU1pGDtlP3GBYszubNjbuRXMDuMXC0UgY1NNuSGjwn7O07oDaTxGYWvEDdY
% HzGdCAzSPozH3R7OxD+eERSR0obORbtp0U+aZ3xVxgxXVvZs3k/zwMv5uVP5fM7zKkpclLgocTHH
% c1tP1sUE1h8HG0w2IRYDVMMkvsOwBnHxY2+xHKcIljNYFGZ2fsyaI9q0t4CvrXEL5nDosaqoHyxO
% /+vJiKtF1ttmc9FE7AF0uwA6G5LDMiH1qU+KxejlPxRXi79QSwMEFAAAAAgA0qlVXQu9C0/nCwAA
% KSAAACoAAABvcmlnaW5hbF9jaGVja3Mvcm91bmQwMy8wNl91bHRyYS92ZXJpZnkucHmlWVFv2zgS
% fvevmHNxOCmRXdmOkzitC3RvU+Bwh70i24cDkkCQZTpRK0uKRDv2Bvnv9w1JSZTsdA84o3Etcjic
% Gc58M0O9+8v7TVm8X8Tpe5FuKd/Lxyyd9Pr9/vUujCRFopDxKo5CKUoK0yXF6VLkAl+ppAeRrYUs
% 4oiWcfiQZqWMo5JWWUE32Qa0k2Gv9+802ZN8FPRVsaZSgktYLCmJF0VY7CkuqRBPm7gQyyHRzSat
% hKCtKOLVfpjvh71vYLCK01gKirIiFQVFjyL6AZEKQXmRZauWpEO6wS7ZGluXIhGRFEtabGRvKaQo
% 1uDDklrii124zhOh2Vm6eJRmkkKzwwLcmC5OH4Zsod6qyNa0KmCnOEtLitd5Vkj6YgYoLOmLpoHc
% hcyypKaJsjVMHqp1wXMsH4NC5EkYiTXs6vF+y00k9eI8lI8wVrX0Kx71RKF0rMa1xr3ezedv1zSn
% L87Yn3g08n2398vN9ed/qrEzjMzc3g0enC+O73qk5jxMjVy314OFVjiZVJRO4V71CJ9CyE2RMvm5
% R+OpS6cgnsxmHp2BN51Q4Sm66vPFmY08mvDcqZJiZp5AWe2w6nLHmTjVthVRnkkYIw4TZ+PR1tCH
% ZQlTkk8f57Thr63yyi19Ir/LEJKcY9ct5IC8+LXRMrflNSIbQoirCVnkSpL1JpFxnuyd3KMnI0i2
% kTDirX8PcicRqZOzvvzjyaUBjVxFxbEQe7RD2JBIN2tRwD1BelWLwBTfPdq3KZ4sCrPbbQz+3+/p
% dA52J7S3tcW0EVVHR7ANkw2MmXoUerTwKDL84LUh/SGKDM4D54ODL8wvOns/wklFtEGUmbEPlJ6O
% SGYyTIbs79YBhBBmgb+I5nNK8WOkZnO2yei+1j1gteCnD8JZWCoxmW3UW+WXOJ1799jS6GdL4Wz8
% z6zcKjfPb5Vs98yVTk5ooSc39iT+0nuKV5AebqSHRFIK8jurAHXZmoMVQvF6EwYgCHEO8B08++p5
% YZ+J3ABRnC1IRJLQ+xYb1o9HoaH2+g0Itsrfrv/z9frv365/5QhV7ODG07MJ78Hhiu2r8DJ+jPnx
% 1KNzNTkZQbqLqTV3dn526fvQaYrPxUxRjS5GSuaG6vySuZ6ZLcaXvOFoalOMprMJUyiC2WTWYXDm
% T8cWh9Hl+eiQxeT8EoNTQ+KPzypssATxZ5cXHl1WglxMJ0zCFDUuMBgHNtw7xj/yIk6l079RuBom
% DYJcKTwAHGxPgQQb9z3k9BD021PEOx5Zhr5r+/cK4MiuXR13e06BppnXlmjPj9TcyB7U56wXMmpg
% vjVkguYdfc2gxHMMR+T5EPiRIj2GSSxjJKhUIJVlJqeCZlGI8EfOS8phHTq8hm6amKnlKlz61NHJ
% mp/AVxUN4mHEWMiQbRv2a1iESSKSgclPVS4ukbl+lFf9Jnh3BtF+muecG4hhhbZmA8fXocMBMoBA
% O1fJtXcPw0YBoXugCB+2Zqb09RvwUHqs+nC1l92rRy/7V/eKXjQtToWHzMPo/hX66IDOnlmq2356
% J8M7ubiT0Z1cxUUpA4WTd7IUUZYu9VNfH2SECogzhG8b8PeYy4zdAGbJszJWVQKXI+uQMlQ7BOtS
% uV+bqkRbt2XXtAFFRNvMMp5hAzF/y1LhEX/f17NJFulDeIuAmYcNc0b0cScHMcmiQ4IDCjtkSnky
% SYHn8bc4oNAJCmRvJqyDJXwMwzDn4tPp38n+8Dvc3lmHuVPKwiNrrUcnmp/ruod89MmcVvHZVfE7
% q+hwEjiimOVkeodbZGS4mCq6kAHSoyuQZPTxMDXqXTY9YauGxceG4PimzRGDxGvOkxfPG0Yww4H9
% jLw6qDQPBU9VprlNuV65PxIk6fwlfUV8mI11fJgHxMcHCqUM8Zg+0Est0WsbSE0czGk8OlfjXL86
% QbCKExEE7lCBQhquhdM3rtAq5GW57YOoQAUdSLGTTi1m/y41LsCOwcUXj7g9W4XPiKcXJcFrFU2q
% XHxhaDPj4QqNhWih7BV9/fz771X4vyNmk36az1R78AgGCfcTuq1R2ehvzDZM9txVREA2wCRiXIBz
% FHOjtOB2SOPzIixFoDAG58bl6YT9ZuQTf/n8e4bjGE/xPOaiFFW8GrRtavGoyt53nEedcXo6cQdY
% wr+m7nw0O0sHo9GUPvkaPqDE0OY04S3JYXPMuEzm0lXt2IxNdS6bnZEWTbFrw9q/wuJBDFbxAhDG
% sg0ifMGiUjxU2eGq71liu2+tRne50c2TZb4ravSgRpHmlHpRAn3oC/KK/C1birrW5YNjH40eOd+U
% aO0KNINInpuSnVY+Ihjxh5IM7pbQihnUlS7XGkHATWcQOGgjV3aqwuPQ5gxYve/Vy5Rrld1F3MZ4
% kAErSlZxTrp29VsoDBh7Qvw+xsmyEArvu5u1EcIuuWuI49WBvZN/ZL7K4nPV/R1gvSJSSdzIcgTn
% oVIEnaIlIzmTDY3qh7B7ULjXizf3R0C6owHgOlq+QVXpcTK3ekbNvc2YQ1MUZDcsNUiCPOTmIb81
% VKf0BDitHtv2bbpQLORyKdwem2+k0ezdVqpgCLLlb3FgXyFucU+UaNxGdPZo2YYV6pis96Ywqp22
% RTlcY3qYjsfWkXa9yxPEpvyWPYuijjZ9Y3T9HMJtqlYWqzw4cB5yW0vq2oJRU0XaQBZCYKtoU8C7
% IvE/RR7gVnlsFeuWp6k4MfP8X3tCbakC1UHthPaCSfCf4wyax/v2Gl1YqzWaCHUBf+N3h7IV1e2p
% LDeG6KKEDmmloKctEvDtGmrntLln6F56PPFXc1Aoh+GkyEuq2LQUvbUYNrIuK6LDI4+Q1+Kluuer
% KnDjmMpCqi1mM4lchHKeWoVVKp6DxlLtmnJXA5imaGOIknmzdriRXuga1NM15h9x7ijddp0Krrp/
% WJSO1J3KAXDpzqNW5xC2UI/xet55z4XDE7Y3zI5XX42GVQG6w7K9+4a3NORvuiDfn7AzYeeUbxmX
% wq3054faanpJY9M6L4CJFQPdGxMFcJaPvqPHwOfSk4MP4rHVtct9UEMZvgp9/RllWbHkns1Eqekr
% O0pUdnC0p7AFw6aJwMmxXTlOlIJvlrS2SqgvrbOuS3GjUKcYPyaKsefyp7vx59T494Drj5jLm+/a
% grG9oev9KaNa9O9sb1t8HZJ21q7FbCd4130LSE7tOO8gSc3sEDZcC2A2y1i+UYVUQMCwd5i0TZQt
% VQ83fxMzSmQBhRcvKNI4nHJztSGVOfNO8GsjS9Vg1ZDy2t20Ygqug9EVy+tfscj8+4CYL1xbO3Cn
% wTe4kHz7p8TIgVzr1nmw5X35T6GrwqEEyd1gyXEUQ3cChEYjeghWHdxSlEEn9N1jOdnchJkXGPvA
% emVR3YW9I35nEqeI65hff+g3MuPqVQdjQahfolQvORgWIviWLHQJruNexA+PDNOtlG88Rc0ZH+eO
% ma/qjszw5SXf9LBQg/3okODMIshQPCjNQTrmdMMdJuSXj0W2eXhE0yC2cbYpqZQit1sZw1G5vHID
% 55xFurxotaPsAprScpipf2E3JP+ojfbrZWWvKyYCVnLFwh2k7iTqE6DwATND034W6C3m5oWMM/bH
% 5yO/uiIshUgDvt5TeaAU0phyHe6aHOq3r53quLuTmuRONkjQby7bZYHqrkGwc9+KesnH9sYxqnlu
% wvicx4yOitFf6bKefgZaCc3EQqiPZlk7Mjg3wQJDfj3F9kSg84UcDtlsMugy6qT4pw4DHGPaaU8a
% 1LNotdp8wpq/CaD2Uj1lHO8QPTtJpj6sYbg00G3xM7y0yzU3cfZR4sFpBjxqpDP+Z8W3db+16r90
% bPR6J18OFvOgHmgconP/8sL293XRqgo4Z2yex83z2H1ldLL0/ZOrmvrVqQ09/9dlTf/cp/brWfWu
% lSOfFSytNqSEk1Zw0NwCWKy+6fgSO1FEcSmWfPlQZgUi2rFUdA9vIUrZvINbGgF4sX2Cff2js6O6
% jC+Pv0X+0FwUEc434Xt8mZHYCgCHPj51ixFz18P2DQJGpX4QrIF9QdCv3m4cvvdQE8fTQOviHia6
% YigDZJnX6IcvzRWuHX2lz+L9F1BLAwQUAAAACADSqVVdabMnYdwBAAAYAwAAIgAAAG9yaWdpbmFs
% X2NoZWNrcy9yb3VuZDAzL1JFQURNRS50eHRFUsuO2zAMvOcr+AGJUaDtJTm2vXW3RbIoeisYmY6J
% 1cMrUUn89x3Zwe7FtjzUkDPDX5NE+lZD9Wx6FXpiGyXg2xV6STVHDhJtT8dUY0+fu81J75SlCGc3
% EqCsUoiBcSTp1VJW9sSlSCntakc/GJUTT5I7kztpoWK4wD5F6TZHeatSTHoqYqbxUvbk021LAWw1
% bGnUy7il+/oKfN9S9Za5o2ONpkEopF48aY9majPVeJWsg0rfbb7rMEgGAHI/7PAQ11pZmiDwQDEZ
% MbkEGcl7AAvXDqMOKQeOTugs0Y2B82sblXuSpobPBSM4W4W/ayXcIrk3ABKtrsZ4DWqlo+d0FY8B
% 2z+5m8BbT5NIhp9XlRtxlmUiJ9keAo5SC5+9NMshu1CoxXAw1kgF+WDCxpeXeNgQx7maptht/iwu
% OG6nPeUaaZptTJEWe+ZummnIKayCWpAz9ZrhT8pzR6cEZx9G5oKObxUgnebwez6s5Z++IA7I8SW9
% 488V+MOGyatT8zM0xd2UUxrQgC8xlbZd2IuF5CuN4mH46ts/zmpjEFS0AbErGp2vPaK5jdjUR5++
% o5/8In9pUN/WD50/duqASGOLz1Mr6bW8u0I3sD9KM9LGsmMz3StfwKLlg37zH1BLAwQUAAAACADS
% qVVd3ikUmhcFAADdCQAAJgAAAG9yaWdpbmFsX2NoZWNrcy9yb3VuZDAzL1NIQTI1NlNVTVMudHh0
% bVZbchw3DPzXKXICh+ALwGlUAAjEqopsl7ROlNunJ6mK1nJ+Rprd2R6yX6BUZVutNduxddOYlpPU
% 194xalFn2Yt5yFTTY7qy6rgwOw0yW7/80ujx969//vrNvuXLp2+nHs7UThYjaRpg+sT/tAXo2nqY
% Sq9Zspa3TQBrk3ltHdxaufTpHyBv+fZgvsjHCT88qvEe2Y5aW5uthRZb9og6O3HPyRWtnTUmrbVb
% nEb/Qf6RL0/116dvfz3EOOk8FfsbRVGiMaOdYs65szYJSfal+4ytOtm4B5/ts1/Ie/0I+RR2e/r6
% 5dPt7fbAHJSlc1vz2lhCyMqtWL/MCdTFc655Dh9lx8eG/ayzTGats0cDcn98zvP0/fmO1YpDrtNY
% VqPte5ENtzAee5y+9lyb2gh8Be3YRXk3q7akaaoy/4x6EdtyCxdUX+Z8QIRnSlJzrn42FmnKFzsn
% thvt3kgh67JT++TQeY/6zq2aVZDK8OLVA2CsdWZ5QREpvK1RpyP7NNYOP1E368ZnQmcr+Qn1nt7j
% AVF0aFymatT2mOJ0hlpi055CTHPM6MPd0/Uiw0dX7np5F+Dj8fPTb5/vyBU73mXYsIMI9OxDicPc
% dQcWe5w9iji15WzYUfo6cFcXXeVH5SPmRe2S3qUfwJzsCbevZJm7d9tivmOFLttLtuQ8CosFYxUb
% j1wvoDvMd2LF5rQBoYP5hI2smdBfGkQM6jElznUB4QhM4cVz1uzchcKi7w+YP7jWmHJeAYSJeC2G
% 9bl2nsUExXL2FdWtVCnDu8PS4KISjJa2NgA9H98+8JpI/hkG0/YFg6bP3L7amqNjrw07JzSPxqYY
% 6It9Nk+pS07lGWo/gV7EztkHvEcS2hPpnEm9BkcPdUO3nPIYWjR5ufiAgLKGo3nObk1G3YG+M2vw
% NzoIafQEgbHQDWMgHpNiMWyKIgMonIdWmIoAlg0Gt0PRFkEfQe+pbakh3fpVgSGEBljaYf1WVlOp
% oGnvo2mhxWolasDawcd4VkT8qpr1+Gxvd8TCPqeu7oLxFahdc4cl8che4LINNsvg5u3EQXcIzzwN
% hsYPwfYHyItW7wTlYdJETAGNNkVNrYmKQTkerDgmlowldhGtsXys0RzEWzDp+Q/ynVRMAlSxIIrN
% rr8cglIMKR2JZj52xE9nomqxkQfUQ8OMGO1sw8LrR8h7SicKhU/BR5nGrXvuMzljtc1bS0yvNmkJ
% uR0lOBuUjStWGYF4vy823yxuj/bydPv8nLen+Ke+1tV+6KWNuI+zRHPInuw9lVArGzWhsjAZRIFK
% ZatMBC05faBmAL4fv/9+e7E7xcwhJiyOKzU8PGNtGBUSJY2dnbDfwD1MgNyuUUrDYJVrkOwp5yfQ
% S7MUxkJhf8zw4LmqOqZLh9sOwbJMCBeENAIdRBB+w17HMHHLYeh70HfV5hUpjG60DAKubaGe5+o6
% CAbAnDqzS+2+GktCMyUMIxcYHF3mmO8fQe91w2Do1YqHYmJFjDa3OCbuPx1CSMCVCgzuBi+GiV6F
% 4wM3u/Xu5HfY8fXlS748Rr7c/n1Hvn66vf7xoE0n+joPehbhAd5CA2Eg2xjwdGC4aQ+qjqBVrpXD
% EA+9bIgHB15hr6/5+vqcX273iRvJ55yG8B4Q0Tw4mgtdu9GJvaA1GQWP+TkwdQdC7SEYm/i+44zz
% P7CXguozaMJuNWTAUmfhzNCVMGvT+6alp+E0ZbjiuIBywnxQ8evk0vaW9QPsR7r/BlBLAwQUAAAA
% CADSqVVdzYRqa4gBAACVAgAAMwAAAG9yaWdpbmFsX2NoZWNrcy9yb3VuZDAzL2Fzc2Vzc21lbnQv
% dmVyaWZpY2F0aW9uLnR4dFVSy27cMAy89yt46yW7TYIml6I5pQF6C9D8AFeibSI2pVCU1/v3oeT0
% dREEaTjDGfJlIsioJAZKGAHnGQpvMFS/LCi1BOVsBVCiIxQFbCJWWEl5YNJyhB8bBuuAYU5oLCNE
% xlFSMQ4FzqTkD6V9VC4TRS8R08v1/ecCWVmMItze3B9CUiH1pzQqLnDGAiyRMvkhNl+ANlPXcnhT
% o41CtUb38y8KsEb2fsNE4dWBu9KdK3VLKbsHS10kDZ1mR3x1xFtFMTa3sNIXJyt0EOJxOiWdUop7
% zTeQBKWeijXwSjBi7q0OqTrbmW3iFhIalOByv83eNLMHjBzAe7CqJ9dJAiG5zlslCdRpaMszB252
% lYopd7+WQB6+3x7/6bY4zg4t6SQ4dyvPvx7/79gdzT6m0KU+RrFgpD8qR3jiVv38+PTxHdKSeaZ4
% 5fKeqLZb42YpmUKP+2WiQnvCvhheI8nabEgbVSafodLKdL7yUHTxt31b9jbA0/c1om3CWnqCklaa
% 7eKp6kqX46d3UEsDBBQAAAAIANKpVV0RSsEYLgEAAPEBAAAvAAAAb3JpZ2luYWxfY2hlY2tzL3Jv
% dW5kMDQvMDFfbG93L3ZlcmlmaWNhdGlvbi50eHRlUdtqwzAMffdX6APS4JucpNCH7TfGBl7iFoNj
% p7Yzur+fkm5jFzDi6Eg6uvghBHA3O1aY07QGm8GW4nL1KRZYNjy17Lqm6l2s4CeynnA5gkTmS5pd
% ze8wpnjJq4vjFkCtlWTVelK+rv7Nhs9AJwc0zN1Gt2z6NhyW7GcHlLHaveMRes7ZmFbqdU55poGI
% E71mr8RNllpReVjLPVkIpPTlhJBTqpBd8dPqyn2VBPiSG8gn0bRt20h+ZE+iAbM/ocRm+m+rpNko
% aTrOtzjXvUQCekCNG6O0xK8CMfDOqL8QpRG8w10dcTBa3qV+Okr0nUIt1S5Oi2kj6FrD3gINDop2
% GnAf75f7zB7/XwBsdnD20VcHk7eXmEr1Y2kg0jVsBBvCwd2WFLe/W+hG55Z9AFBLAwQUAAAACADS
% qVVdPXOR9eUDAABbCgAAKAAAAG9yaWdpbmFsX2NoZWNrcy9yb3VuZDA0LzAxX2xvdy92ZXJpZnku
% cHmNVkuP2joU3vMrPLpCOMQBQsuiTL3grrppF910JMRFhpjBauK4tjN3RqP57/fYcUICmfZmgSyf
% x3fOdx7mr7t5ZfT8IOScyyekXuy5lB9GJ10WqGD2jEShSm3RiR1tqQXLR6OMn5DSouAGf4vWSHNb
% aYm2Cp1KjRQSEmkmHzn+QL7FKVlGSJwQy3OsxplXyS4qSyKkxWo6na2iOI2inXf+WGBFCqKj9QjB
% V5QZBRWsQcNfBERsKlAs/8WMYJWk0bQgoBp5DOYwcmEsroFSoqIo3qo43UVjpzWfK4+lyxICIN8D
% 1obiNlPvFDBBNUm9lNHFPSorS7eLnb9wUPqSTkq+g37tyX2ZeBSWJtgnxCAhcOUSAb/R1EW+IYmL
% bKxaExZTbzVtFC8SYzgUovVF6aKVQUwzphSXGWY9jkBQF0xXEofIjmUlraGvk19VaQWXFokMfgWc
% zWS9IBNhyoJb/QKa8lFXXB6DwDKRI/6rEk8sv9zy5yNXVpSS5YnvCwTiirmbWsEDOq6KKmf11QGu
% MgYQYJxXJqi+taT6NgpNli66nP4YrlCHwg11jYEZEOjkiZf/uSvUNdMbSvEPsB6r+8DZdpCyXUzT
% 1tbBHDodseyEPtAxH6/EjYrtNpUaUPKDQQ8xtIme2kFxSAM34wREJf58qM+16XQDCYbp6rZU92uy
% H2yMXvow6kDb4jZn2U3oY6+izfdEF85ejn3ZEM8NR+mNlqGXsZcEIl92Jr5lLF5GPq3l6D1WzFjd
% 1YBPEHCNBnfuzGSGDAij0XtE/Lble4TUVtAqsuLtpVs6hobVAy3S655zN5FPVzyFXXj+Hb/nd/iF
% qalY3qOvvy/77DmogX5Qmmc01KghC8v5vJ41n9n2nPiBgq46J8NdFaoQQqLO6WXIbtfMDaUPdLWA
% XfyF4nOskmUU8O/RV/oQzj39Q5X/rNdC8V6MnofiwsMX8tU9SVeV1Nq7+dPTM0TlNUBKCiEbmD5O
% JjQ/2isk+X8xerP2cJNE4N5REvuEaA3Xj4BLWqfWUDTkgx0Mro2nYJA8RJ8pTmOvHqoQOcFoaIau
% XoWbLdpjanXT0U14xT3S1BEp4SntzlJ32Jui67rk+u6qJZuYhp6lNjAYdPizMtnk0JzP0LluGCF0
% HVCcMlLunM0mUfuS/eQvxC0H7tKpcWbC8sLAcxxcniavoPW2Rq9e8S1YBzxFV35hwJNuRFZxU+OW
% aPWPJpB7SmazGVku1j0zv1xWcB1t0/Wu5/Dv2yQR0xydhITAoPfYoyyNFUdDkARcJt3ft4Q/q1K6
% p0+B75NLEXbnfi9Zwfd7Sif7fcGE3O8n6/rPxug/UEsDBBQAAAAIANKpVV32g0LGdAEAAC4FAAAy
% AAAAb3JpZ2luYWxfY2hlY2tzL3JvdW5kMDQvMDJfbWVkaXVtL3ZlcmlmaWNhdGlvbi50eHTNlE1u
% wyAQhfc5BfKqSaPIBtux23VPUVWI4nGCaoMDOHIU5e4FR43SnxQvs0Fi5jHMfHqa4wyhiKu202CM
% UJJ2TLMWLGjKt8A/TPSEkiKO19nyrJTG6p5bqCgMrO0a8IpXl0PoOJ5OJf0jvF5+3TXsejD+TQW1
% 4AIkP3jJRdAyy7dCbm7lt6xpqFG95uDi5HvcMr0B69vAl4TVTJoatI+W+OqfgQrphmuhEswCbRSr
% xmZ/SLioQFqqYe9qsGYsnl80HLQVtXDzSKVb+q566atEZqftQ7pwo88fcRFHo/60/A0HZ0kIDg7A
% wTfgZLfgkD/hkDgP0iEkTCcvp9HJF272+WOK/6MT5+ROvJNkKQnyydMwn4QkU+3jxw/4h8QFuRcD
% FWkRJFTmEwiVxVQL+fGvPeTOt3E9dVqpmu4Ndb9vpDJWcF86ehkYt6gWUlhA572GlGwOz6iX4twC
% cguwb6xBTANydfZQISFRxzrQKwvDKpqdZp9QSwMEFAAAAAgA0qlVXYUfwAy6CAAAShQAACsAAABv
% cmlnaW5hbF9jaGVja3Mvcm91bmQwNC8wMl9tZWRpdW0vdmVyaWZ5LnB5lVhtb9s4Ev6eX8G2SCzK
% kmvZbfdglQF6wXW7uGJbNPupPseQJVpmqxeHlBvJ2exvvxmKenOyu3cCEovkcN75zFAvnr08KPly
% I7KXPPtB9lWxy7P52fPnz/+ZH7KIR4SXQViQSARxlqtChGpCfs3JNsmDQmQx4SKGjUFy4IoEkpO9
% zPMt/OepUFxNgNHZVuYp2UpgI/JMEZHuc1mQ92aCBIq8r2nCPEn4kCritwd+ZgbfVJ45JA2KnUNU
% pc7gb6J4IXl4kAp2JSIVhTWbwkPPIr4lG1Eoq6SLMwLP3U4knJT1AJ8NKy/c0ieV4ElENhOgXic8
% i4udRV3PJ6XLNprNXoqUW5nhI3lxkBnJLtmMBFlEgiSxsnNBtrkkgoiMyCCLuTVzUNGJULeygL1j
% j9Y67YOEVy2zQCkOljUSNMPs/BVjc72sbg/gVMXuhS3Os1MRnpPRh75OS3VILe/tW8sS42/0HPjh
% jm+4w3CipzwyutJqbROxt945pVMNNbPeLcvV5WVFLzw9jcMbBjIqH96r+r080zxCCFwhDyHY60jD
% JmYzW45nfsSszPXoy5czPf2ZNX7w37HPy8XKTwP1naHyGTi/1oBNffKVJUIVlo4kkpAb8nkZrOAH
% aQNK6XIRu95K7/iZwdL4q3/NvvokTplxR6mNLtHon6n/RUvCfzdxqrflSZQfip4glHCh6Sj1yScG
% SWbVREuQtVjB5FV/ciHH3opqXr/oec3FMHA/uVcDj0KOWZ+ogz9X9c8vlDJmRW489hxg5USuBF/p
% TZ3mXeJiDNh0uNyTuNAEF+wPYx6GJpChpWOr137XARwyuF4gTeCUnVzh6BWeHVIug4Jb17RTAgl2
% /VwExemilnO9hATc0XO07Zqu/twSnKtO1W9X8RHbmuATAVKrfMYCfUb03FUjr6KDPTxRXK9UjTXG
% 9XhU0X591EMAuMICx0f99NDkXEq2PCV0e3Td2elxT4MSUnRvBRvlAAdK37JYK4tpiBPMBA2wM2LL
% 6crOfHButM9FVqhmItwh72b4kXlv7NgvQKDacqkMB56IWGwSzrTLevD2mzzwzoEly3hZABp0p77n
% aPQtqLUUq8spdX7NM965EZbAUEVwdrGRPPjeLlU107/j+XbacQvBCyKCDFIMvUouSHOWEULMWa4o
% hYXGMnj9Q5/djotxc8fMIQg0ziD+R2Z1BBdu905P8L3LwQb5jtQ378c+R4zW8rgaM68XLbACJ7px
% hePO4DqIhqoZVYMRspz1HW4EMfZx0XhBo+uxJULPlisXBcNbza5NjUb+/5rrTQ7T0xxGPSB3P+rc
% xYnWyjajcdbYAXPzuR33mbQqAbkNsB/XIl6YXqLIoVYFKc8KBwtXIDJoMz6AvgSg06lPDIdaErkZ
% dBe7TS53eR6REFiLrQghmEM8aQxZnGZKlrf1q2zq1wB16p3QIDTQ3Kt2dIyjyuwF19X7Rabq2vEO
% asRNXYN8sgPtTxH5eoGzv+uMv4BtfQ+1NQJJEPyvtNk1Xro4OQybNA4U0JLtIRzgu6QiW4HtBwZD
% pIeUQOeFASzCHTZmm4p8yPehzLeF6/47kPsaevi2X+jQCChx+0BIdv/gS/4DfmpFo2/svlw0yndm
% IYeHP0WcCFjjvulwQ40nGA+YQGkP/i3TvZ2FO+DgYbfJ3geA3C2vmv/tsB6U7Hayz/fI1Bqi/rCW
% gP6g+kkxaQuK0QTsXdSS0YpHpHAKt0iDGGV2oLYL/Les51dMD+Cce/7tJNhjcKx6aQCmuF1LOgFT
% rM3RtmtS/09jjgxkTWLIppMSaMQeGwzX1RNV1cRHRIJGb515qMORPiEBHwwYthyVX1vGSt80nI/c
% Zrgy6J8NyTCmg6x4VOkHKbKo/dK0MCIUPAsrFrl4SpBgAFw9iva0XO14+J0UO04+Y7NJPthQsPDs
% iKKqsQgO0Z2AK88BysphI0KSHpJC7BNEGbhQTDSjL70D0zSF8gdmuYdK626xRgndeRDX68z88jAE
% hS+dzclO6QbVarg0yNNwobbUqdfmwhedvM9YrybKp5gcB0zapSMsBf0lzfmInL/2C6LxKOrHgH//
% ghEJbO4ZVF08u6rg0brneOk06NOf7V6dR9mFSLdWUA9Czmrwc/RUEUjIU8VMj9y1P+3bY14Ag2so
% UlymHG6rBV/rNqstaE69HuoEWEMiAxNITdYvZY+ZmqIDdma5TNcbDRbbkb7U3ccP9n32QMf3P03t
% +GFE61tQCgXNMicpDOAS3CsMUVd1Xjvea6934DK4KEXjYZGS/f4a7gczuD/Nx94TWKF6F0/bXCGe
% OM0Fi8bSVf4dCnOVWzyiMMG/u5xqXLi71A5STkF7ujUP3PvZe0tf7GxlF+PMttS4oM7dYzCCC/6G
% weIY6WBP4Spq2zMYWgVYhK8F7tavTzEwiqFIhsy0ejCCa7hda2DhDcT1nsJByGTtlkWfy9/t08Fr
% uqoX5F2SuPqe7ipxhI4FWB3SvflagbUEMQne4GQECSCMQzbmE0r/40mbCnEXsVeON505s0EyvPZm
% pqnqGf8eFJ47r6gbu2/gWuDiOHa8N/Ry2pK+IBm5ZB5+BLHjmzl+SEkEVxoGwyQA/SOCuUwOUKxk
% rSPRFBue5HcEm3sgllzt4HY76alkeNr2vCdMn4TYzigug8BZC6wBsoHy03SPk1NrLL3Djsd4fDDs
% 5C2x4vGM1nExEQFs3h/werRqXQfwg85bWt7sJ8ejjjV77YH74Hf6Zq4n5tN/zGFm1XnUsGlK9PBT
% BTUYpwD92f0ozNM9DPCT0nofSGhVAVXWIVYTNVrotHBGLQfABjAZvMxh0YgZAslIfxRb/1DrXiaM
% FqN/6XYYujhRQGw0e5JnSeWTQyZqfDJKtd/WfvCoro8Qu0nBy8moLjCQmNAp4geySQRZqax6nyN0
% r8hmULJyMNwa1WImSDlyRncjOrmTIP6vto5H/8kA3OAYrdfYt6/XjI3Wa8S59Xq0qPHu7L9QSwME
% FAAAAAgA0qlVXZqwCbtbAQAATAIAADAAAABvcmlnaW5hbF9jaGVja3Mvcm91bmQwNC8wM19oaWdo
% L3ZlcmlmaWNhdGlvbi50eHSVUctOw0AMvPMV5p6SZNO0UJRD1fISHCrCgauTuMHqZjfaB4+/x22E
% xBFulj0zHo+faa+pDdRB72wcZ6h7ahwCd2QCByYPaDqgT2wD7Oot3DkcYEB/8CtQSqrguBXUbl3X
% Z7dR65lDc4DG2sMKdtVSlWmh5gl0FKr8aqp7HAas8lQt8mWeZwnILvf1wZ7AYcfRH2cyzC+LbJ5l
% k/aD6WjPhgOBddyzQX1cMC/SfJGAoR4DvxN47g0Q92TeUUeqZkV5ArR2GNGwNZMT6U2yWxsbTbMg
% rv1ovZwsEB+bUxxgqnIl6zpysMgqc56qv7EWP6ziX7TlD02V6jev5iHqgIZs9L8vYeyt5ACotW3x
% pGYbH1xsj/VKoD54kGNVkqfzawg2CFql5SS7fnqCm9f15gU29zebx/rUvYDbKeWjuhEBbr14HEdN
% g/xJorYBHI0aW0ogvBGMztq9vzj7BlBLAwQUAAAACADSqVVdraVJAEkKAACOFwAAKQAAAG9yaWdp
% bmFsX2NoZWNrcy9yb3VuZDA0LzAzX2hpZ2gvdmVyaWZ5LnB5lVhtc9s4Dv6eX8G204uk0LGlONmt
% FXbWaZvtze723OQ+JPV4PLRE24plSdFL0jiX/34ASMl2nNzLF5sEQQAEHgCk3r1pV0XenkRJWyV3
% LHso52lytPf27dsvP2VQsou0SsIuCyM5S9KijILikA00EytKmYQyD1kcTXKZP7A0iR8O9/pxzGRR
% qLyM0qRgVaFYlJRqpvKifZ6DUCS3f5dVUUQyYblEgowLn6mfWZzqecFkrlihMglzdYgG7U3zdMmm
% tQQWLbM0L1ktEnSyc80TlaA8TeOGJ1P5siq1YCOmSoItljivxoEM5kqvL2U5r5emoCDNIxnrpSCN
% Y7VtQ6huK7X3Q1gd3rH9f3z7IiwXRnuhmjIZhpbkE7uXq7LKwcxhZ3QwgR8uhy6O3BHxJWpmyYar
% hWy8JevVoppsSyGpuGdiaz3LKt7R46AeEuKgHm5ooNbQOlp6kCY3G8qf6U7SfPls1XE8LcLxtHmB
% jBWoL7e5SjqjUxLPSuWplTQMw+EPNk1zdgPoABAkMwWLIyJFWyTaDBG0+nZvjyUiVgkM99hvTcis
% b2migIKMuQqspSwWyMsigcPDSVSOAwByadlInAJdJLUdEC0gyqQQP+B/xyIUg1uStCS5Hz/e/M21
% gT9k/WE0Gt6M3ogfPdxOIUkKjoEwS7y25l/u6emNbaP62j+AxHqMXAYtQaqmU6vPM1SMUsk8cgon
% u1RSLRXmhAUsyEBxB7V90Ada7b1NDRSbaJYgsyFbLdd2nKJaWhmwf8xg23Ono4cz237ujOjA5WZJ
% G5urKWYCBqYJ6sr494apGDIfEboykvhq+wB5el8HnMN4e7FvAh+qUgfeRAjZM2TdzGlrbXcfjNsK
% h4bm2q/c+MPe9RQkfTAfqzs6Y2+9HMfgK/TWSIjoFWdpSNRuffM63/uuEJ090heFKimj8mEb2f6F
% WDvWHwgNfT/MBfriwvaVKKsMzlRDFA5CloPy/+YhYEZAS7H2hz8xkwuYwBoeX2x5Aojx3DjUuHOy
% 6UWWw2pdATwbAYA7NQAwqXQ3ICECeLmV8IzDjMPEZCRu6DWnQFX1hEtgMRKQPuBhbgujr2HygOkd
% +3sSqkwl6FY6vUxwdFc0uAJXHDbSBsI4EqO2dkidEP/BhxupASXf3YASWqeDO0nThVXxOx6gy/uC
% yh9GuLIP0NoX4AHY7VHliEY67/tQRKGAisDvwx9MBFXqwGyf76bNEDSOKMavpRy4CpTMsTyJlU9V
% Y27krtbn6K8PMYY+Xxbro1wJ9NdSZhY2Bl7Ztn+9TboDUl5DEXoS2rCKMhQBlYKim3Pdr0jvhGC0
% WpEU020c+p/YL8kA5q/QYVtXtgO/1zAdCPfg6uD64Mq5bnnOanWgpYCigOfIH4qvLaJhIyUaJpSh
% bfCZ41/xa57zr3y14hDRJqazXIKBvFAqpPqQ3hdiOLTOLaiNRHzf5cc2h7nn1JRj/ottc7gPnFvR
% +yPexQkOPX5EnDWfBztf7n/sTJjeube73Hu1aT1zHBoLwOL0fwMYOdM9isJhxrtBMf442zMFsliM
% oeEGi0LXrLo64UWCFOJZ1pZ0yZIzQX7TCNfO87/TxuGwvrU8t1bmfJIbd9C4XujzMxubHJuLc8u1
% zelfSKS5I850LpFxddJ/x3wSHarW32HlVAycuT5dIAtVWGgyXRagSpPs5Vq2y08ah790UlYJQsON
% iekR/0AxvjnwnJridnY76xKjzNid3u056/2uawQ0BO/V7YGA3cjX8lAL1xMAlnsEcO+LjYqE7A3G
% DcTFTrIjl3Hb1alwyWfX9SDMP6KHGMsyHoZio5f5Zk+WCQHluqNbI/AIK8ybKfV1JOFFea1o0AI+
% z7Ewmyk5A1v/53ZrRYzbIPQpVgfCxfoP98AWngL6cIYX86LH/uiSOsmm0Z1qBQ9BrA7Z5RJ6Ogia
% JVFZhapgM1iE90wewZPn82eW4kVsmclE9wt0d0LFHtMbA91fZ+QLKblVw5kKZwrKhAUXOHuX+aVL
% Ft6D8copRFc3UdyMK/b7xH6xQhBxwbWOrYpP2o1JkOyAEBcxhKCii+CCu8eIDn171H3AMFNr3gjr
% q47PcnjgWdP9C31tgbyY5WmVtWQ8U/AyrIVE4GiMhaKX5eDyM/sd6gLFEwL1SPKe8CoDYQDWQf/y
% ct/WqQlbltCpdXbqJHOxfnbskdP17zYJ3GptjDtcT2hoFuyRH2hHeDaBbys1/P8lLwxYcyF+0JGu
% hLgWUJC6kPAa3kLATixRH/iR17U/dppNGy4F+IMdv3jHxEMGNztgRk+auJQCksK/F8so0d2xwuDo
% Bnk3XIx2MNTVlMUWBQB180YsQOY30bxorRO73fb8mVwuJcknfbzr3Letb471DSBir4+rjYFDHR1D
% ITrx9FH1ZkH+PHF/cV10aS7DqCoErbWhop04mzqbEmHY9GbY7f561Ol2OvYpEk46nbVu8O0HqGod
% G4IIHPzXD65Nppwcf8AJrKCPtzG6huZ5Fcct8MWCrjQ9aFyPgyeOVUg8hjDQh3ikP5hChPKH+whS
% z1j4qP+falQCCGvUeRqGRwgqMA5bPH4FeFZv/w9UGeCAh+FINZpQ8jGSiDAgsGEQ7NdwNSA4UcF9
% nrprt+CFeRpBFVQMYjOLEhk3voGnG9x7sS7CVZ+paKaSOxlXSjusqY/Gh/nTVsJSARiXqigpZZ/1
% 0SPTR+GhceDtvmd8UAUVc7R7tfH4xuVm8QKdYC7EogeFrIySShFxJuIILFEgGT9tzLAS4g/HKdBc
% HC2QtsCROyILDmWGjwlLWzejpw501EQ8qif/tlLgCvrwYw0V1cr7eRQrRgtkSyZofJilsH+qPz6Q
% 4TM0HDVok2/N+WfDn/rC85NeHnZ9nlv8+IAkVN7Dn0O8s93a2ojazFt746kFlyBkBCSs8y7BXNfd
% EF605gEnoJ03rx3aAlJqeHxOq0msWiU4ucjSIqLvbHCDpuhC9B6Tpx4AJ1Q5e2xUPonkDejZRAPo
% SwN6RBEa3rGLFE6UxVXByvsUrJVTlskoL6A1B/MqB3QVZQEXsjioYokNhXpG/EDWE46BNnlA7MED
% DWCblMVh88o63r0ou9i4N/qyT08rGPf1BVH3g67pB7imaXgrNzR4gHnI79GTbJ36IOpoNOziUncE
% Q6EbkOktdBKoFJ6jFTiOBy+Ucj2z2/CGKZ1S12yKcy3byMEGb+KqpYnhuoujlFHz7YHWsXqi3aee
% s5Q/DU175C8S3+frDwyUNu/Y5UN8B8mq8n1we44fUCFY3LRqiE8YJTOCRRBlMn7u9+Z9DrjCi90Q
% HhnD3kKj2bxR/0LCCF8KUNSepS/eq5uqtH8ZwaughLd7CvjYqDORnOGH4g00sXQCl7aKPsb2DGjc
% tgs33nbXZ2VaArfXPm6wCMk0HidyqcZjIfbH46WMkvF4H2/8+u7vry8a/mYJ8zcR3Bja//NP9uWq
% /+mf7NPXL5/+uCQ9h+xcF9SNb+Z0H43VEgo0p2TOVRbLQHFWzhVIS9NpcQgG/htQSwMEFAAAAAgA
% 0qlVXc5z/oN8AwAAEAYAADEAAABvcmlnaW5hbF9jaGVja3Mvcm91bmQwNC8wNF94aGlnaC92ZXJp
% ZmljYXRpb24udHh0hVRLb+M2EL7nV8yxBRyb70eCHAzHSY0GyaL2oT3SFu0IK1MuRWW3WPS/95Pi
% BEGAtjxYkjnzPWaGXP4+X2xosfxts7pbLeab5Zrmj7d0t3pcbZZ0u5rfPz6tN6vF+uLLfL2+ohxK
% 3abQ0K5NXQmpdBM6hNMsh6ruOwq5Ls/HWOrdhNKNpjrVpUZ0Vx8Sbds+VdfkDZVQN1TV4ZDaDrHd
% 9IwemoaE5IxRfImJTjEf+zIyAroMiNOpoxxH8tzvSgTeMXyvj/2RvrW5oiamQ3nu6Ie+Ij8hc0Vc
% TsjigS93RUL/fQ0cRV2/PeS2PxGyYib1JkEwB8Sye67T4fIU6kzbvjrE0l3Tvs2xK1TFAl11Gsxf
% QwzMpFjRKeRSD1JpH3alzWNCn1GCl7bpx41YH2J6CU0fu8mrGS7eeLmg+B2Jo/NZW1W0+2vXRBpF
% XobmELc50C6CZF/vQhkgtm15HosUUkwF1Tu0Y/2xleO+iUN96PQcuojypeqj8rEjbyrsuwjOziqq
% mLpI7SkmFAEsyGrhBX8PJs94pS3o7bfQfH3t7YgnP+C9otyHvuvqkC7fhwccx1MTPzT1+sy7a+Me
% /urBD5yVPm/f+g9CR30aepDr0dp9DsfLY+i+0p99GIQh9CWCdUAomCoa0K/O0OJh9oXRDTmtjXRM
% SKG4FY7hzXDDsJw3jo3Lz5TzymMStVLs/9fIZD4zcamcstwqbxTTzHptvLXeaGaUYYYJw7V2jimJ
% b43AmZAKyxknrZcWC/vIE0YwroxlWlquFTfCCXiQmnsFTA0vowL7WYEAMDdOgMEqI4TQ1joPSKsk
% EwJ25fCEaw4SL5jSDJtOM4nQGWdcI5lD/KAcihRimOHIdJ5rMahzwkOwcJJ5yEK4tVxJ7fGrJIqo
% jRvFuc/iDJPMMuWdFgj2nrkB0FjgOWEFfEvYZRZRKB7Dl1FeKo4MywEuuId37Eo709Yp45lXFgXl
% 3I9ZChhGSwF4VMjDqTLKuqHoAq1VmqP46A1n3iqASeNRJof4IcKKi4u7h6f5ZvV4//EmpKfHhz/o
% p9RiQnPb7n8+DztGbULttov5BbOZY/M6jbgdb9gUerSCGKsmb4cYQU37DXfPeHZudLxk7wcHw/Qf
% WBhWjIIQ9t+w+BQ+MMHOvYJezB8eaPHLcvHrmgaC5e2Ulk04dcjjU8moG67UCifmH1BLAwQUAAAA
% CADSqVVdgWyVbO0RAADHLgAAKgAAAG9yaWdpbmFsX2NoZWNrcy9yb3VuZDA0LzA0X3hoaWdoL3Zl
% cmlmeS5wea06a2/buLLf/Su4W3SPpMiO5aTd1goDuE2yDU6aGm2xuDc+vgVtMbZcWVIk2Y6TzX+/
% M0PqZTntLnACRJbI4XA4nDf54pfDVZocTvzwUIZrFm+zeRQetX799dd30Sr0pMfkvZhmbCqTzL/1
% pyKTTKw8P0ttdhsFQbQBkMmWBWIiA3i9DSKR+eGMeb6YhVGa+dO007qOWCpFMp2zKGHhaikTwBQw
% 6c9gUhGsJPNTJlicyKWfShbdsmwO8wRBO0o8mUBHFN12Wp/l3cpPZMqGRCY7ctn1ajnc4uhVCrNH
% YbAFshIWxZkfhTDFfnpgea3bJFqy2wQWB5Ap85dxlGTsQjcwkbILBeNnsPYoCgqYWCbLVSZonM2m
% 0RK4p77UgKXI5jnsLaCLEl8EqmsKLJP1GT15K1ZB5vnTTMFk/lJWprr9NoWdABpaaSagidcaDbPV
% AgzIIc+4T81+i8FfzB36RVbcMz9k92k/tvg9NSYyWyUhi9VAID82hD0x+3lHtooDaYjRZOSPx4TC
% RxSJCGfSCGRoCNM0aWwGbakR2r690BMLHvhpZijY0DRdAUhsMVqMOT5s/HTr8yhUsUj8bGvEZl/3
% pqulEQP0aQzj9lERmyY1L8pm/8CxdZf5sqeWt50G8pv0ZhIILZCPgNEPMkxlZhi+jePMl6HGV5km
% NMeEBDZ0OgchIhRVHlbw+M3BzMeh6ffTU/83RyHHzxLEOTkJTcKX/8EIkHojjDKjGMlE6Gk8Rk5q
% jq82H1BL5G5AaZCThPkF++zP5hkT3kJMZZixdCPilKUkXS4sZC2TFBUENM4HxSF9k/egimkKmhZ3
% qvsam+6Gj8bUtJn7AehouDVwU0/hceA8Iy5tx9TU5EK5D6Tf4ESJt6/lCN+5+lGytOmIOJahZ/hm
% dWM2o36/7QA/4sQPM+Nf5/8zeP+VvT///PXy4vL94Ov5Fza4PmMXl9eXX8/Z2eXgj+tPX75evv/y
% L7MlYOFoCQyna781Les1O2G9SuvvtvNaN18YR3avGOEcXBivsLN9YfTsI8t6ZZ5iy7HZwjWH5ZoB
% qpsv+DsPLSNsH5mHhz3XUx8OfijGF9Par5ye9d3yLM88rNB2yqnT6b2xQvhWXJBxyrH1jfO2h82/
% m3VkQJ15QPNYAHqCyziugRzDKNXj4AIrPUc91cVpgt6r13pazejh4MsXsCRCm98pGLpMhOgtZiI+
% TITnr8DSg67PlxJssc1C/gr44mdgIlnqz0I2Qb/jsrevWSb8oGa3YXNamynvuktxj0LOH5+anH2r
% +Uo6z+v6Tx1Zykel5co1ap/uIzR4qJQ/kpXMUhRCeC7GGtwm41P1AIXpA7ahDNPMMKzjp56fLiJk
% kmoDJE9VxqIe4GQmrwgEAQQigQEZ7xY2PcZpq26otLilFvlVqwobAW5wJYveDddmwq3MvwGzyfU8
% tPg01/Z8ZvAUOPeDHxuoY72xvRk5+LtHfTmfNOfVnYoHYg9fJoBqluZqnTNeEOOhr24vZZDK+sT4
% N+ehvM+MOVE83zXJapr5nqnF+K+CBrOBFYgCrEjUDlXzsWnnLfOCzhKBZjAa9hkSg7tMpKmvtISM
% VhlXbrHYzxr39YA+wtGM8GLPGjNhN4/Jb5BUpebJrl0hEZgecMfNpQt0ytDvth6lEGtlG4Vjrvtb
% L9g5RYXkLdq30QocxyTNkhWFNn2K3rJNxLxoNQEnoVwOLFsmGD/+eQxqH2VscNxphfzYlXJHTd0s
% +7sqOpvxkeJ9RsqZ/SPllLuaKXH4uJVKGfLHnZ14crPIi5QjRACzpVwgtuZhF7534ig2zFZ9z2al
% lN6prZvZsVnV1TviCcAi7j4+OsLzjDuTps314a7wNrhHRAbnx9r63mrzC5LGHjfTJwbOvW4nmMjQ
% 4nY6b8BTknHGPZNgb2GX/eVqScED4p5l85Q96r1/cmHUMcRkk1kSrWIdJRyTQY7BIO/xb0c/MMPu
% MuXVmMqVaz5aqgAJkSxTZAgucKkMkrJBCR+hzzmmvfNf9kj7yfsd7ZMN956PhCWU2aIeHbhk/HHZ
% VwEz7H85cmlWSXhybzgGoZusQylKamh9mKxgGYSjW8LL9ZMbCz8Bg2lgOPxf8Q/jXf9AMzQdBM4l
% 7VtyCghRytoa07YpbIDCKXMDd1ta9Q/88ZnYdR9tOULtu/DvAsJfjiFwoSntD7sWCYlHODBEr4u+
% BM1Oc9+XJ/xDSd6N1LswWo5LSBy6IyOlMt1Ef2dIw2reyFN+E5HVhNfSCb5gl5AIo/Zh+BxDctdG
% BY5C/PQkJGFL2F14B41aJYkMp7JT4vYWwN/+aNxk8VMzIv5Q92XgbLU7yLcCvLW3QLepzYExsVGG
% TdU+KduFbi/wrcK1n0Iq69W3yoUFYCRXAiqjVoDXKYI8PFPOtSIzBWzOXqQFJj/hTt2LQlKx5tfA
% OBcYxRGX63W559QJKIn4mqz2ePdiuk4il9FaGoCs6a3De5AtY2FvdJZob5BUpAzAx0jp4heOBI0b
% Q6EPTTFg6E8SKb43AAAZh95Rdwz020C/58CzvbG8bgNWC5fnnDb7cHobWQH/9qLWDdtiAdpdGcXd
% 4jeyfRM1pJedcGbcyIObyLTIuDkHdfNG+vfDYWVaURPNQn0aIV6uS/vDvGXMl389gmF6au3hyJKC
% WGVMVYpbftfgVRPoMkYr+UeMX+Vi+D3YNuu+ateQ8hWiHLUhPRrXacfOrNKpaTKc9so8EBb+WpnZ
% prfMPK2ExPE0n7gS1k1OeGmeCfmkXEzFfdTd9GMMLjr3g20crEekLuJAe1UxLykm65AShaBnENdn
% PpWpVH2JBkAQ5ofrKFhRR1lbS23l8p0e+esX7A+xSlNfoL/O5AwceZmPddhHkSVo4KFR5gDtdCoC
% 6dkosLDmNZYLqCwYyjTttG7OP3/iRtfumu6n63MO/ririjoYv1BpSWflhgCVOZjAA5J3B9+csYKE
% yKIJ2S4g2yXkchU0IS2EbCOkhZC2boMpdFtXjwY5XRiiPtimgZoOXCjgz3bwZ0QGiIRa1yTtVbAQ
% BC2NgPChF7WclTW5B458apTlHjhy6cG+r9UvHnRVTybGQI8PScQGCgxNFPi8nAJgeyvPH2qz7NSb
% bKo4laoQ83xkPiCJyE4O+qC7wOnYpkVAa16FW+yvcC1+c8x6HoPLwvREsTS2Daw2YNyCAzoTP1Ml
% TFALs7Z0GFMs/psXN9bvejF/7PaB8KdWg+qCgpCKAzUrBtPaawTz4g64jyVoZNMo1JbXf841FGtu
% QlBJR245gvwF7F64oTeChjExJPQ6oN8GfNu4T6aNTF7bir8VBsIaQ6/KFS8eqXJh2yHnFargF7eP
% mAV+KqhKS8zVV57dxyCUIDB57TUGadV1Yz+c+rEIjIHth15alkhHA/AetV3H/oo/oU+lk34YJYgA
% 68CV8WIMIUlpDst638Ak0Zn8wjV+sbdf/ML9sdZbeXubFqtD+Xp8+mk5RKOq7OO6IfB+Y+b+mqtd
% GYywAD3GUKYq2aN4zNd7JXYpvstvlOpA3nrGX3fteC5SyY021tc0GQM+GuHe0+zfahnLnqbWHjqr
% gkmb5I+5cYbmd83PDg+NI50fwX4rKTku6X/gsLCuyqBOwrajAIy1RZSCVBwevrLzLwe/yrELrjNx
% V4sGf4C3Bc1PhvWhpsiDVmu6Nzd8o+mnraPyzAhauyaYCPXTtV/hBzKt4rcHvMZfxdp4brpTwQv5
% KKAvdwspLgzUbXoZeyroLkaoOVD7GaBijiEn6zgVpYvP9W3ItQLSr7JjrUpwp0Yq27jWtrGo1qkQ
% IlbWCrBra1Ugh/FKnSs5yte5LGIJhsdySe3MDhBhUSaRs1UgEqbSeBHM5CQRZbbymVdOo4xALCee
% 6JOhagaEzUORXSOJh30/jhBjftmwwEQr5OeqShLbeR1I2obU5zS1McB63rOmAtQSA3LNIvVNTKLg
% EqCqURwtdnSpjDK+2AZAgPCZ7mfUb9UcQ3Nbt1dLNSGsp4ZqEPMR2vD+uOaO6qnFAweiQAjHu8x6
% ALpPuk1XMohHoKG4Kq4EpWywQWDcAVpxalHqV3b/nSIp0jzS8VDFBkeb+hriOrnTSaFrsek+wC5N
% mmvSewC93IBYycAFmnb7gcKsnSwn/geaMCk0Yd90iAwmOuWoHt06RRmICBY4kRK7Uiog1uCahEu0
% wMPNKnJSeJepqHMwmKccg9Yhldk5v1SGFHXFJtmrk5ggOC0R+EViVV/mrkggxn6i54BfG6lrokQx
% xd7PjQmrOHIoowdxbHfPZLhAA3bRButn1oe0myPykg6AcYAp+qaUGe2kOFOsQlKtGGuRh5HnqSNZ
% ZYHa2gJVTRWkK5Mom5M/F1RjAaCIMhToSuQtHqBjDoT2PyXlrhZh8Awpz3h+1wmPLlbHMmxj+QaP
% N8GmUD0UMiUPD3BZul1CBgSpD9iuMJQzCCTWZFEpG+q0on11zqPCl51xp2uF9PoR9Mo4Iw/L+ULJ
% hXMAPmdhvjxC39uIM8fPHT4NhzwP7D6in5vOZR73UPAW0ylx/cyHQn5dSqYRfXqOsH1coCvjvo/2
% 6Gdn19XIVDv3Ck6VDAcBfvA8SHWzKAPD3P1Z/PKTmFszcb/3UPcGsnkKdoOi+43Eo+79Yfuacz+P
% TRWcBezTZP8fsbExrHoVgebZ2HlYv9GTWR9H6/FoQ0aXFrPZ5WKZL2x+c5quK98QFUN/pBh6r74h
% MFdkLDCPg0cdjjh+wAHOjcoqxR0d7FwYuIiPKmwDWvencguIt037zHwu5tGU3J2oOwlnFk15wi8A
% wro7NJz2nWkNh7s2IIoLI6B0rdREWjXzsbLqZ75WZ8LKNiL4rg6DSZ2PSnU+IySAcAlxGh1pYGiT
% QuwPZkEdMdtsAwaDLrcEYByCLVuFWFABZUagPxKxpE0BxSaSvuFoOux89jwZ4PiRVZyoQ8jl+ViI
% xAOuyZY5PZfBKhK/pAi2FeRnQimJWolItlRXOToEbCrsGnT/u/lAN08IYAaVEsALJAXHtUOTI9Mt
% ovmf5QZde23u4Fehf3dP7D/YGwh1x//QDuBNnj22IM87qMjRVPM1L+xskaIUS1Idvcqi6smLHmDu
% 1vZ0+y+82KS6oWxoUgMDqh6GQVgiMg/UmzPO69IUdkm6FAcCTNOVyngCm/ecNg67hfEYwEYPB+Vn
% mfbnuQdEYsNBHvlMI3saFnFcF0+vq+nT1e5Ng2nUr+QuQSYoLJmGGMlMI4w9yswEejFghbAbRqvP
% eqpydcCRBTmkYkgO+Heur1ydcOe1uoRiDbvVnp51hcZo2K3cS8kXEkqRUKp5IUAAbDxaqOSWf6LP
% drpdMtgISVLSo2zUf3nUduDH8s2Xr9o9aHEwlIIvaH/ed+PfO0BL4S0WFEjK7k17q3b+3t7mNymw
% nlkpG/9ZqYn8WSKbTzjV9d8pDc+t+E4CO3xXiMG7itoiJXnpdPcih0jsSaIXQu95x8B+V1nMcFCi
% HrxrnCQUNx3RGcwn7eG7asKV75DeF6sGDaJJIwDrKR92LaOJy8xtcG6ri5Ou0L6A3biyh13T3LmG
% pLxNXu5uF7eSqq5DnYa7+TVXVAR/6tNJX8V8K8f0pulFqJrJ7lYCHZiKGSEsAQwZXVciybNpYqo8
% lvT3cyeJlc3H8KmvKehdHQ67jLNHGvREJ+xqTf8JL64+Db5eXv9RvbDGPl1f/S/DG4PqlqwJI7Jk
% q0Rb3yIF4xJv8UZrGJf6QMrwyn79bIWlIj1THsYdkSRia4xGyqneq9o5Ke3h6+4PMshKZhLH+h7h
% 3qtLpnutjlEqmc4XnDkKZWoY1/a1aXvZNpZck2A+V8kD1/HF4oPpqCAbcI76EFzh6ePYrjVjkw19
% 44qoV0wpLM6ySic59OsUmYfX5XUpEWNnAJlIMOtAfAixWDo3hhY0yq00rs32l4OhNfTNDgQ+hnk4
% 3NUPwHDqQCBVu82njgqrIRVJjI13bmSyxqBHBkr2kIJHePQ7b2cAoNMqgFAVIQpC+GN1BoJEUZX3
% Uxln7JJE5jxJokRXkpX0qYvWq1CshR8ICHr6P75kzdLvPmiol2uWujReRGQ4Y76m/4SDqyv2/sP5
% +39/YbjA87MOOw9EjFe6H+u3ndt0Cbrf6d0+sRRvsHikZ/8PUEsDBBQAAAAIANKpVV0o+9avhAUA
% AEcPAAAyAAAAb3JpZ2luYWxfY2hlY2tzL3JvdW5kMDQvMDVfbWF4L2V4YWN0X2FyaXRobWV0aWMu
% cHmlV0tz2zYQvvNXbDXTlHQpV3LbS8fyqXGai91pcrKqeEASkiCTIAcAZcmTH99dPChSj+RQjWaE
% x+6378VqNBq937HcwAfWai2YHCtmRC1ZCa0RpTCCa1C81byApaor+KduZQE3KbyXRu1h8vt1FH1e
% c82BKWHWFTciT6HhqmKSS5PCB8WqFBgy/f3pz3EjtrUBVSO4RGSmOOR1Izx6ZNbM0OqNS+Ak4CcN
% W67Ecn/d7K/ho9FQlwUUgq1krVGUtsgVExIKJZCUICMhDXJbM8o91JUwhhfX8MBfvf6/DSFQCwQ4
% CIpGo1FkzV0qdA4CaRBVUysD9/4AmIZ7T9PK3NR12dEoXrQ5d5c1uoKZWoW7qi2j6AlmEE9SmCTR
% Iy2ndvnRn06TKIoKvgRWFDFLIUv+iAA/iptWoeD5ZAE/Q4Y/6Nf51G2mC88k+SpmQ44xs7RjdqDS
% bXYO2gkkiCwJWhhRcX1JjyurB4ydIldWkfRwZZXrriZBel7LzbGSpzrKWlU358jg6gpuOmBce4ZG
% 1eh4E+/0kMfFI0bXp7DTPQfn3+Gwpjuex8CTlUy+xHJIP58/wRKD/Ex5pJhccaRYnDnqNBUyFw0r
% ya9CFiLn+hiRzcVivnEgGwLxdO5E9E886rpzFtOaY7Jh+scWBr/oqtkMJgdmp1PJJXJhsHuyg1cC
% 62Rxkcubg+X+XDSdeImZ7EjstmhwP3+iJIincHsLMpwT9Awe7Y5EVEy/HKR4WodJH4HERHKdCfOc
% YyWbODlcLul+hsIRCKGJ0Nr8dACgDyYfNp+Wd4edg0OUhvSCMtE41d55AzaJ7TshRD+cCLFe2Bmv
% L3zt+E6oUFMkJDdQ9fld6qsumJEGUcNA4fW4KxcbBLXHwF+KA1lik+CgqwdyETC1wb6Pttidbivd
% xU3ao0bxLR5Nunj1g4V1dRyvlWJ7pH+BLxC/wN0dVV4XhjWx4a0l+mKxu8sNnjsCG2q0YWXWcYJ9
% ZtrRNLXG52lLEBl239jivPNsySC64nJ0yUp0rXe/33XuJpd1cniJr5xtr4NYWM+h70nlUDmEM8jM
% GLuwtbSfugn8CDdDdTwQSaFlT4KPDWlp1ykMCXxsSEY/RSxtyBBdPHN67bsEGYXn/y/EQiuZhE/5
% ulVg32ntHnJ8PYVcwRtXtT+3ya/xsC3p1d9f03t5JuG+4fvvFF2vgQkUpA2T2I93KcSChop79B0h
% 7AghhCOFf0+qa/RQw7KsGVkwbmpkhi0rWz96sMKNBhYFJw7rihwFi6XImeGjszr51MBCcu8YbvDI
% WZxRxczj+3iHrQ01osV00ddW1a+uneLCyl2cq6a+P7zcbP6yoK9v5N11QQXg7yaLftIVcNuv9V5W
% 3DPM5SHpsC/4Yyb3scX2Te44bC/4CmPhyyQ5bX9nRV1swOIM6BDzkuhTyb6IXAvNKFzUQb0dpx04
% CwG1U5HfYaYRDD1Qv0Dhio1cj5thD/6s0BJXX1h7Vbw9fsf9LIGFS5hWJ8xjmzr7kBgp7MmyN9HE
% bQqvCWbO05GiRGcTZuvSp3XrMMns85LHvFjRvEKzLc2+s4da8qOnwJIk/SFBwt0MfnUnlEp+vjkU
% cArEROK4bCsaZ72knuupYceCQkJNTR5yN3jTVwmSEau3n5bd0xS0BqGBFD+AdzdYW1OUgd01ZpkO
% fZgqjabB3sn0TJi/kUsyIQU2lOLiPONJJ1v0/VMMnRP07dewG6XIS4Ud9vsjbWjP1J6erXE4FNLA
% jT9h6GYUm8HYOBcXpjI3CZEtjMbvOf745pQNMDLCyBzJOaisD5V5DJlCFZo8wzS1hJm7yw/pg6AV
% vdb/9x3Iu9pkPrbnAKvLgNUJIL3ElEWIFpab8P5fuMVQ+LYwCBz+0WUUpiz6D1BLAwQUAAAACADS
% qVVdwuFxq+wCAAB7BQAALwAAAG9yaWdpbmFsX2NoZWNrcy9yb3VuZDA0LzA1X21heC92ZXJpZmlj
% YXRpb24udHh0ZZTNduI4EIX3PEU9AAb/YWxmNdOd1Sw6p9Pr9BFyYXQiSx5JJjBPP1cGHJJhwQFZ
% VXXvV1X+O1+WNLDrhWETdlQ26yIn6pwYjtSJYUcFDlIiJ1o1+h1l602apovnewgNVl+M7ZXQJC0f
% DkoqHPsltdw5Zk/palXtaJMhz5LyzXadVUvaVs26XlJR4Xe9xVeWpymO83TxdBbyU1ptvScRZglV
% XmdN3mzKpsqKHHKaddZsiiz+fPwsvrEL6qC4Jce9UEaZjsTeWz0GvprL621RlEVdb+u62pZNk62b
% bbWpvqRavCB01MKRH1gGN/Yk77l3lL5mcJZVrwUM5a8wU9avxR8fWPF/tVj8qTXx5E0eWb55GoT3
% 3K4WRPBN9P15/fPi2T3EiUiwjzjjJSqI9kILI+HIq37UAdfs6InbjpN3Vt0RyW0/CKe8NVNQVhOO
% zInP8cleGREUHpE9kD2x02IYIpZBKEd7O5r2WmpD6J/xTKNx7INTMnCbHKDeOvpLDOLuYqoB9awn
% ocJdKED0DZTQCTh/1UQZ3UjMVpOHfk+FExPt7K07WtveaUsReMqQl3MG+w5kfuw94gLSXdscjpDw
% tWlT7eoe6dU5aRVEewBB2Ss/oL0OP/4xsEXDU8kKJg/j3MPwbhNn3+lFHkeXxKo9pmImuI07o0yr
% TqodP69GNB1Gt58aQYCr+rutCFJbOd3vB83nRF6k5g9pzz9+/Q/mJi6neUusYeJ/UEyFC0nheWY9
% g5BH4SCeER2UfGSuAD2ooG5BwHttOskxJJiTgPnxA1JjE01Lg7OD6ESUFJkkFp5EHI0v4mKfHieI
% PGMY50GC2lgVJE58myplbh7uUmDvX3YWrRLd1KfPKa5GnzDHfk6Pd0S5KvNm8SIhbEcHdAcbf21c
% zGNs9I93lLFYMjThfBSjn1TAWVwNh0cn1pFkHPMV3kqDtm7q2e+9gkfo5lU4Q7FHWRyABh20xRXT
% Aa1Cp/kjaMKGq7HiOK39f1BLAwQUAAAACADSqVVdnLfvImITAACCNQAAKAAAAG9yaWdpbmFsX2No
% ZWNrcy9yb3VuZDA0LzA1X21heC92ZXJpZnkucHmlW21z27aW/q5fgTRzx6QEyZKsuKkUZubevfVO
% Z2duOu2d/RCtqqFEyIJFkTRf9OJM/vs+BwDf6TTd7aQ2CQIHB+f1OQD89s1tlsS3GxnciuDEomu6
% D4O73g8//PBbmAUem7GfgzS+svG7OdtQi/CYuLjblHnSfQzCJJXbZMGCkAXZUcRy6/pMyEfQcv1M
% JKNe71dFkiWpG3hu7DFfbmIXFMPAv47Yf7pZkkg3GMZuKsMAo91YpvujAF2GESwS8dENRJD24jBL
% ZSAS9BBMBnvMloIb64z+zE3TWG4yomHjo+ZxXdIaRZjsQQYY0tvuxfaQsCSLojBOGUYxD/yHKYtF
% 5LtbwdK9wOSufyUuojgMd0QzcsHMKBWXEcmnt4vDI9vFmAeTJkweFbUH08DchD3oPpgzTsPQL/ps
% wyMErhaccLXCLNVvesDRxYpM3x3IhbF0ff0phkyIpP74m3rTX1J5FHk7SO7WW6gLM+uvTXnkPa3P
% nH3i7BfOXM/jkMmGs0A8ckUOzG3D4AktYXyc8h7Df1uIw8u2KWcb3w0OYD+WwVZGrs/ZXi9m7UXm
% Ib4mIsZz4q0VA5w9xu7RELpufYFvoQzSdZId7V4vDh6ZYxZlvRvPpuPp/WQ8Hds9tZYEH7987fU8
% sWOp6/tX6yCunB2ciT3XJFWvJVpX6KrfRo8i1f3GNhuwgx6ewE6F9cJZbIbGIs3igL0sx6s+OH5Z
% TvBb9/XkSXqq89l0zkBdyYcaSULW2bYXzEOzkhReVT83wfJhXNUZHqwMc3Dm2Vw9T9SzmQmsumYO
% F8SWS+vButj8wbraNtuFMbvwK1liHJ5X6h0P9O6u1Bg4CEZ9Us/09aD6usGjsHwRgLShTf9FxKy4
% pJYlVV9Z9j1w05vJHXOXcrU8rNgbh30Gy/8KA2EXRPA9YjJRrfN8hZ/rnzHwUE6rlraMVtwlog6j
% X5waFoZ7GJ+Fp3IO3azFjWc1EP/KDk3uBxPeWq3hJl+Mg8XMSXMIKJmo9XrGZEbjpnd7xnzWp++Y
% VS8YhJ5oufAvy7xxvaRnTf5pZdu91hglos9V84EEtK1EXmEq5hM8wiILKhzScrlhzrZtGB37yMZN
% 3hpmki9iMIH2VfjNXReNZmIRFzNfwJ7xePpesfkLybiMAfiqyF1g7/RF86Gd+Oafv97+Rp3KWM/c
% x1iII56SG7u6xguWodnYuxAA39RFsFxqocJNav7yIsFfzDexrd1GPRcfQGWVh4V9Vq4uIh/xZQKR
% VkK0VXMo267Pn4dHaDmBBlfLlH42bTQf3GvpIlWJJtIjEvNiQh7GJmvkhlherIAz8T1sBjZFJom8
% Tbr6Es2Nt/NIZUmVtN1UWFFkf63EEBXbFZv4QAT4E1pFGVk6xpd2Tz6UZhEi7BN5nbIFyYSfCBKC
% aXjSDRetJ7XSmtMvFdPL55Vxg0+dbqDFsVZ+1bYGav6/24Lru/EagQBRIr1C4pcmeetCq5GV1Yz5
% 2G6GhsBu6T8oJpGPEHGdrjWc2P0+MqIVwXY+RhQ5OuwnslszSROEIjKtgv6aMlT2SFpyeWDmSiiZ
% Ls30io5SaLcFrRpSVckzD2MJ2XkfPxHCmgwZO1+1Iqfs6KMYBj6I9mvhPcKLG/FtacEK7VeGvyaK
% ahZ70lnMzLQRqQuB5F5kLHyhYprO4maQyew6Odfyeh7MVGjzqhHWsExt/cutpoR/eHNBkf6NV8OL
% UdIZgHmfJmR8MLlgKzznwYU1Gc6EwOx1uVTn+ILM+2Dt0drP2RagvsLPyQpZ5bbSQD8VD+qb+Tlu
% ZDaIK+dDG3UhqVuSqBBa2EIFAGHwmO8ScgaQX8fiMYPjwF/OOV6qxpQCz5JhFfFEcIVjzucRwPIx
% sSrBJFXpvx77xKvZnzgk+NDIwiqekAoIcnieVbxz6wyX/WZqf4WkEVWaVkxrnk9UJPvivTpRM4h5
% wicJ565JQ00iqnavhLpmbCIWE5vr1GzzxBiWLhbWgAq52580qgTQHtFHQG9rOOVTm7datJLXtZjV
% bimgKuF661QqtBbs5m5N+OaNf6p5kwI0NTxj8IYr47Uu2XKzoAkCjmJNJcilNeXvbG7d8R/xc8bv
% 7dW8Zh7oWHJDoxqaJP6rokI5G6Cpaqm1/hs3IZes6UOri7eSRsB13CAva1uZ8qFqNVjEXGigbW7d
% rrBgG/Ku3Elbo2JxVFonRon1PG7rKA59WKR06ThPNiyU5wZtIxW2ib3iI4GxDllvaiNho+sCVYI5
% u8utAgqs03k+gLDti4Lt1bKHVkavL6YmalMyCHOX+b7ZtkjP4ZCG/04qGx5lII/ZUW9t5Fgz/08F
% XgWsXlNRvX+6hw4IZYn5Q1EQIFz2LfOgULWGT6IOn/Dxa69DUIQCFNmR3k+xQABSeABQn9YnP5/1
% zKr3Uqz6hUHUo/U3tVE16TpyIMMBeOCvxPmmJLTYMfwE3FmVHgt3LDyJGHQiiXKfvLuUvsG5QqxV
% UzWHvGX/FCnQCRRGe05sEwvXS/fDnYwTyClyg4Do0WCVjmWalBtVRxo6MsHFRTDxnC/jrwvV21lq
% Mz3vpS8Yacp0sT8EpQuKvFwWpTiTME6FZ4FFsliL0qu2Sj2cMoJFKbbaWMqJ5h5BCgLrhBObDqMs
% 8pRB6I7EuzMugl5Si6oFKQ8icb4k83Gp3nI19LG2lCJEc+XAxEdnZlO+TKOVPJ/UBplpmdMPxAhH
% /ZarwaSLwFONgGwTkIbAU5WAWvLRvVj0wPVDkpb2X8+JMoD8Ye6lVmg99u2DNe1bgULSeJj2iRje
% jIn5IYJfV1JR+eSOzxoZpFIh3zfisoJkTj2lTvjdd2bUujs6ak/MUiTtRSsvlq7uFJCRwKnn5ABi
% wQ4H5xUX/W7P9zx+OHQ7tJIb+TNqu8tQcWvQK7Dir5/+rT65sUzg6Y1g+ug0/ZrSrPNqkn121O5X
% kZrsIQU9JJWuCus7ss+3Vnw48LyqufBHu6xIaZ+tkm+eV91iEb7aq6BdbRV/kkhskab94aMbdYik
% tLQpr9vayVl2WNL4L2AxbUdVNPZNROZ04LFOTPanRmMssDtBFAXOv+NM2NVsYSRYFBwJ0jEoBSLM
% EuVbQz24IUZd3Ea+BLbi8sSzrMwUXuY0t8GyTG2ALbxz6xMlL9pUUgs8OZDFaaX+N23PmeOehm2C
% A+qkht56me557upJ3QZmklvvbIJXIr3M9an/czZ8Pue+nn/4qEO+O3G6Jl24U6drCjNmOZys6H/H
% es4GOcXbKcwIA8EF/uHTuf6picQnZvsPj9Pax5I+6AxygkQTq6nB9wl3p9zLOBatS3sE8I6YW4mt
% s0psxdp1NFzCCQ5/Q4mC34jg9GtsV0pWCANOc0fQFZ8m9oov8560sFXFL/iar53iwAGL5FPicVwh
% hhp8o5l2cvOa8uWYT0D2rjLrW/abALxJ0jjbpmp/jRI5goB/Ze4OaENbZwpMwplaNKETAeRzBexL
% 4+uoV+LzTSb9tMkcTQxexnZpHKof5Fxm+UmZCRD9z9PiFSOVr5Vdtw7QoXfun0zBfVIF9yQvuL/W
% goXaajFdpkVNft4uLTmY8qfB1EbqzvqnZmwgnH75OC4j6HlbZu7/byCZ8fO2I3ZABxI62GbpEOJN
% ATiV6BNtwHEYuYRgPUZFwDCMCHCDu+5MZeo3Y6WIDls+s2v7geWnM6/UkpHr6Op7ETh636nLwBvF
% 58aplp0lAhI75A8ncvt7a2MPiW6+2112UonUIUs59y3VT+/+bEz1hrFljNjwZVpuL+dJgHYvclCm
% FL6yOyu+0iAamzQtjWrOPzqKuY/OuKmqLEDQUeqiUE8O5A01AkCcdalkRpVwEuwfbuQSWwKtUKUU
% uY42jt5FMtKK9ihLgL4+8V+4NZyoGDAm31+9shHRe61+tedGcI4+QdCUl/JvsxVXJ3ym4QkNdj1a
% VtXjQGvVow1QPwzDQDCzkivbEpliNS39t1neUHhFpt7oDUTnc6/ThTY6YJMlOM44PzqvclY7c3kR
% cTikI3x15l5XRc6iMvude5QYcCzKsKMz6yd/TBnkwYpD+8mtNWXJcwxXpuOCdC+TygUBloEgDRxV
% eT86zky/G6UeB9NOIVB7AVjK080n/lIvomEGxgh+yc2gmlBIjohcwDzql0qDhClndi0DoGDT3fIz
% Yt0HYVgPQ2JVBlF07Ny0Cv1rEB6Ba9fFoZblbpJ1vslnooRqWYRZ6iwfLCQr1CyDSSmEb+zJV05y
% Ixo8qQDAb9i9AirPxWzmhCefs1bz8Ev7iOf5eXlYDZxLX3G+jNQRl1wtqH0woS81OmDt+bl6qOMs
% L4Nr8zSGTnVBfPXK3miQNGGD6xirnNWrJL0Pr/bbHNLtT3xybzfBjYrSjvNgzX7id8bg6hVVvvcC
% Qoh1RfpynC8PwBn3s7FJlo8Ovd9NzQRv2X+LWO6ueq/eFAHZkdH9FNoBzxKCAG7AgjfDMPYAEXZ+
% 6BJCGKq0bzYUtY/4PlANKamB8+9byP++Euw6o3SpfN/XTjRwzgs8q4MI86xi37Bol/SsTx/k4/ok
% KC5QlQskBCzEycv4eEUbx3nTWDWqpnYKKTtRNzMuJ6NJzbrHoXCmbvm/Fb9XQ1f9e8J2xcJPHHwq
% cy3ZnTez05KyZb7W/ik/V3uqytJu+A6k65C3aL33Qb5fObE8raoxVW83JfKCwIp6MNGhsaiPtVno
% y1GawxusRtsmeZNS+ISyMhW8+uhFx57yAOuN8zn/oE3jqWUasm0aFI2cb4QkregLGazwaLXvJuQa
% iHzTdz8qH3qwfrz/ib+nh7t79f7+R/VrMh2P1edpro7cz2hOJyeqxeTQct6NMWLhh0mi0Av1Q0zp
% p/3+oX0b4cd6rqVRH8DCdMonY/pPu/y0rz48qpf7fvqhsj+a7OVOrWqpF6vlOEj5X5BfJXQYckBl
% xDdVPTS1rjspUmGeoqSmQ59U7aS9cSrnw399Xs9StPVi9TWPIN/HJ4toDW+bdZVecUdhqemMKalx
% /YycNsufZ5QjVzba3hjw8Jb9UlY6ZnsVSqRNXbDFYATA1WK3k1tJHbZ7mj4ZmbF/RxfTxPYiFoAQ
% wHrlxb9NEvpZKpiKuOp+IDvvBRJgFgs8x8L1F4aUROUlXE+jYNp4lyi4ALFi+NkRi/OzhG2u6h6f
% biYWdTfUanS9B3DM0DLSGMJWXeW0mirdbBOQ14mi9reuKerlZRHMwyG3/NPsfV/N3qHvOZVLI3lG
% bakUhYU4lz21hl7vXdlSoqCyyNSv1xDCvT0/9Z16Th+ki6zZVqlcTsMsT5S68KDiDSxyLMi2Pzib
% /mahZDJwNlXj03IqvaYWO4vANCyjlfarYUBhdBPG+zD02BaUJKwMsCTH0oZOqd2aJYIwYMVGqYGh
% BJFHNZSX24+5zDrKvPd5fv97vJEpXV4d/sP0o2utG6mAPQoWMjdYkY9JaR9AMW6yOmdPWUI8m1j4
% trhfSnzntlpUp9XVcraNIa2hYim/fosx5EKjrirz3Ssl5n2JlkxloDik6uCDql51lFFVZ7Ny02up
% 1m95taBrNS0vowsqPFPr5r+mfFaqdH5jzqf5jUmFj26Exke8x64nswQvaY3Ar8V9sIo5VNSacMiY
% boolbDwa3WP8DWc3oyegKevoRhZY5TQyL2sN2Z+1rZUkyRhJfwUb1FAb8h/a4CD9WBxdqY6ainhl
% 1jHUaSg/e0e0olPaY+ZXbiblmwMa8tNdjOC1U4E/KV5b8JxSTvvG46V9tKNrp2t99IYuOM47j3tB
% 5Tov7jEQ+L92InWCt1nhOyVW18XdrLO0m5WFnVVsxDX6TF8v25uH5C+O2gKk0zhCwm/0K2VeS8HH
% jiuWqtR7yYs7aYq7F7u7ym7UGJ+LOpveZu/N1Sonvw+wqB5fU+XojKvn17C69vE1YPaSriFXv6xq
% g3JQd0YgTRZpnJ9gttDTtBoJ0jg/bdSwC6O15JvinlV3gSDKw4fpbK5nqxq1auFJEzwScdrCV1DZ
% QDo8cmrTtzVQCKjd3Hs62wCWU5B+9h5vK7trCTUshNXWQeVb9i9xThHVzZ0LKZI5m860eAicJKQt
% dXCs7/XriDFEl2MYyK0hAlxCF/dRwKkNizI4cHCz9TNPeTwADJHY+i4APpAG1iR1uJc5zNnWqvFO
% lWyjXBNDktY2Wh6GEggY+kCF01DHhB/o2sDtwQi6RT4Ow/T7pPvqmSaNbl+Yqm0SiO5NgqdmFEI/
% FStUoFiop8EExSRx2a9vDmApFDtqwBTKjXp/ejWhcclmOuMzlCYIvTUokMejhm6r0KK0mTqKMGik
% bULqigGBypx2XuXf8OmsljF+b/bIQYvw5mz8x2TK2eT+jzvO7qZ/3HM2e//H3aJyAXr2fpTvwFG+
% yaNpOnaqf9thmSmrV6QW9cPtRe3YZVHfTlk0Y3Z1Bf8T/L24NWP+ZCYitXijyjnmQR8SmKsQ5k8u
% 8u3huaa0u/lymr/zvjL25fC1BhF2Nz/7bpQUm8Fw3S/19Q3T8Xx0t6sPu/l9C6Q0Zzv19zztv0fi
% alfbpb8C2rsEu1Bm6L/iIWgVnoRPG7G0pNFNAxlEfqgh/nojIxcKS8UovaT0Rw6JQANsvLltI8pB
% +qJDoqbPjKR6CKHrdeAexXrtODfrNelzvb6Za732/hdQSwMEFAAAAAgA0qlVXZ895zP2AgAArAYA
% ADEAAABvcmlnaW5hbF9jaGVja3Mvcm91bmQwNC8wNl91bHRyYS92ZXJpZmljYXRpb24udHh0rVTN
% bts4EL77KeboAFIlUpRsuZuD0RaLAsVu0aRoTwVoaWxzQ5FakkqcPH2Hku00WWwRFD0YoDmjGX4/
% M9d7h5j2MuzhDtVuH7CFjT3ACgSH3noV1C1CJ91OGf8apNZgXYuO0nqpnE/oqt/LyzxhCYeP66ur
% 2fX/lmRiKQD/HaROpwgoE3CHLt1Id98g/XEQnOo1+qkWgLGmsVorg/IUg/k8ZQmk7CKB8ZTHA0+A
% XVxAh9LAPI+X9HXvbDs0AS6hLjKWwx+XMLX5VtCdyOsqY7ykxNND0KDb3cNGamkaZXbgA/aecvmE
% C9zedpvBw9a6jtAr0+ihjYkP6Gzqh763LhDewbSECW6lHtCvJjSfpLlJw50lUK6TWj3IoKxJgNiM
% 6SaAs3cJ2H4kHrNYEgj90BkP0rSUuNXYHAl9LHu+dacGW7UhKlfAi18U8Wc1K/6bRORnEflRRPaj
% iJO0T1QsFqxY5mVZZpyVFefPBF3kVbUsRFnWGV8uRF29TNkve0QNvSZZY7ZT1sHGqXaHKyhycSpB
% 71AkG2y1vSPOluxI6cSMp15rYvbs+Sk7a4kp8EPnAQ+yCfp+Iu91JGUUeKsOlP3YeHzgWdy3yt+k
% BkkAdwPN0A1ajmr2NhBuJbUnSegpvybzi6oz8Vun9snkRr2P51HtJ9r/V/6SsQVf1LysWS1qVlRV
% non6mQt4JcRC8CVni6qsl3lVsLzgWSGKl9nh0ziWpBo13siN0kQsQZuzjCcsE/F3sYLTDjr5FOZF
% AvzJAhrZSI6VxWz2LhoAWiV3xvqgGgg2EMUretbjO4KT/9DkWaeiB8QYm2h28YuJZxJlWecU6+Rh
% 9GSnaKfE/HIkLN6lo/Mmb0aLUOA8y6MVx73TRKw+booYPy430i3qPeKmEahiJ79X2/Aswisawtn6
% wwd493X95hrevl//+dffV9fv31xNTvhsyE3OxwkIMmBHKGiTOYzU3o6uJ2v26F4FPCTkmthgiy5a
% dutsB2GPHiGgD/7V7DtQSwMEFAAAAAgA0qlVXcG2JqExDAAA7R4AACoAAABvcmlnaW5hbF9jaGVj
% a3Mvcm91bmQwNC8wNl91bHRyYS92ZXJpZnkucHmlWW1vGzcS/q5fwTYotCtTsl5cBJXDALm0Lgo0
% vsZJ0QN0OmN3RVmM9y3krrzy4f77zQy5b5KdXHsBIq+4wyE588wzM9SLb85Lo89DlZ7LdM/yQ7HL
% 0sXg22+//akKooJtVaoKyTYquEszU6jIsG2mWbGT7CYr0w27YEmQlibSKi84+43mM1ME6SbQGxar
% UAf6MBl83EkjWaAlK6QpDGdpVrBI6kJtVRTAGMu2pLVM1V5qE8T4LdMyMRN2nbFtnAWFSu9Ynqm0
% mAxuynRZ75bBBLU9TPIDe22fUaXK0klRFXiUwVZnCYuyOJYRjhumkjzTBdvIbVDGxUZFhZXZ6qAv
% ceUGWGDYeysD9tBFlsWNTK6zTRnB6aMsAUPS0qb/7fZBFbtbLfM4iGQiU7dcEhS7rpbBINrJ6N4w
% 0d2aByf2BwMYYQ9S3e0K7wHUV/5ywOCflkWpU5ruee+9B5+NRiwgLz1weFApe1S5m+Jz9I0uxHtv
% 5judZqe2hdx4FWcHmBDnu6Cv2ivKPJZeElReMKb3nIU+rRDwsF6g4gdQT/Pqf26eSj2QO7Oan5zm
% dkKHv6X9eFecxUEoY7cVnFXhnKtlswSOHY7G6nFaDN95Uz7jc78vgf+qnB9y0R7+wO0GTwQDYyQi
% YVWtR1erw5q9EuxqcicLD1RM/ZF9PuAzZx5tmn9BnfXwakgrww7l5zKIVaGkGa7ZmWAzmpFrcLpT
% xobLIY9l6l35fJhnBoT3ErCj71RqLuGoMcv0Rmq5YXmgtHFOFHRy9tubDx+GzsAhxGIE6JPaq4yz
% ybVA1fD1kpkyMcI6DR49nT2sPq3JnPCIxgQp+voJv+ggvZOenbyartGLqE9tWZAePPPdNYkaFEXF
% /rLG03WWyi6+3Irn58cz7KZlKvXdARZB1PWAiZsESEb+aDRvEVIZh7GoxVjoNxaIgzSSCC9kNhXU
% AKuMiJXBWLOjlywUfXNd1kgA7Brir+YgWby5dSwgKBARLN0dWdOYQuZGTOn5YadiydB6BpFkfP/1
% rAUp+lGksio8z1P8k7W64qRNpmUiNXCm17iDH05enOCudQ1Y7LWYnxoJA5HTodrZ7si4n+6pAejg
% j7CVg00KFLo8jiuzUmsOn5/WfHai9pkwei3AfjgRA47mrpuZaGqLCNECo7ORzoKCtnK8ZnfWq1bb
% sVjf90KEjQB58UzYOHXSgB6xCtcjF0qDHm/AeP3itehAZdClAwtMzHCFDj5Bnso0UUK9UA15XLzL
% lgUk3lgaALQli1glqrilLGkE+dKiyo0YSDXglqsmVrvyXQ93sGhnrqoW0LX4FsnHG89eif0rMaP3
% ewK8vz6ebv+slt31rFAqpkA9EnOz3Igmpmgx4oIvJlLPquKLDsUfR+6gEwAUvHRA0JtCSVHKp1EZ
% IihHC/ZasKdDuiX7HgDT2mNuwfpk9boQghvWsiZGYj1SM4EQi36+Kh0rQ+SehK1ZzQjvyMC9Sfsv
% Tpo/PclRRAlBNNpDDH0jSviER7V2HPTp2CVemwhKIJC5/0SybdzbidYW+dZXWkXMoRlhnz6dB1M+
% pIw5tqUQbKaQd1KPW4/XSurMd+SHZYcqeCjq4UvHzm2C6ALHbmMI3s1SrCMhbQf1QkMEwDCRQTrs
% 4qCZ4qKdiSF/BkhDqCjs3v+1QDFLGYsndNWntZzFWtKgzeNcevB7lOFOOBi8YB93Wkpg82JnMAFS
% xS2p0Ie4J3diHf7BlsnsYglT3rAd1L7FQ8bk5k7CtL8xSgBv8c8lS+VdQPVIlEERgpiQbB/EJdgf
% q33gh+xBbiaDK9GtaN/7g5ZMrH0cjMYzCGVun5uH/qDDVyAqwO9ZBfiHj8UaM3YFuKYvJHG1wsLT
% BzC5utmb8wX/nr/0eeUP+tXmkEwzRtO4KhtCNsyqYS3YEu0XRCnO2uJovOA/dArYsPfmZSdOtiIO
% knATMMjCS8s/lJCBg47ZaUtnGsGfsxlU1TNgEPdMgzT05Tl2Bg/Hs66a5+e0q4xn/aUHx0Wt3gEx
% lOaorD0Ti4GDsAVgLQZGSQiIUVxuEMaPUmdjU+bUEYXYYEJYOzwtm1IWQQlmvB/POVuMsTaFrjEB
% FMUQlNBTaVVRHlMFYH/LPv4+AeBLbfvPLKcSWoIWiOQywVYtoIXdd2iR1F1qXYp4R3aGsAeteygM
% WVOAW9ebyeBaeN4Mq3/sNeBzRp9jOwR/xjN6NQWqfec4GSw4H4U85MFZv5O69hl7AdGcYTUPcee2
% OHBFuHmGfBd+0+N8naXf+X6vI9LcCNTdDGxkId4B46/0evQO6H9l1mP6S98Vfj9GC0yhbgtPzZ8A
% BjSPKgHrQTOYaVfWuKnUV4bG21tDUJJ619YSERTGA4KCATuPrXGn7nPeGnpaWx0b2weQnXNocpE3
% IODhcY6PM9gaFFbQOAXxESEBlQEaAbOPkKaOX2m5dQnsayw2qg3NSwgQOgY8ooQ9Q81eYdLptMi4
% kHGrut86arDAY64FaMcXLnOHaUfR9VcUXT+haO4U1XZZhUnLmA8GqZLqmsY6qzA9FniezOc1b89P
% WH3eofdZ57mpIF6wG2v3Zm9ABVoFIWb3BQXmhb1JCuwtRn2J5HrjST9PuNQAHxcuWcya4Yu1S5nO
% z09kjgZN+AAp5L23wNMBpNACDswtUERrMDJPmHDgMbRRfZgJsFBivBoQqfDCBHY6BmZIsKijz14v
% ASPzNfQiJEavHTt1XSMErDNoGBl5krI3pXmgldLecGEMIiTq3fg1Rd/UM2q1RCCcpgIfpwUSEW94
% 9LxDnYa209jAMmjD292E226ZD28a+WazWxVKfZJ6/8dJwOw/KnPP7nSQ75ausmJTDsAvNXv79g8G
% 3xM8s2GANigJLizLBykVWFplGnTEMtiy7yfsJtgoFJ2OXyvMVVB6segQQfeuYAQyIQu1gtIIBb6f
% wMR6wgpxveZOeDWlL04Wvs3XE3alKjjBBq887ySUnSwoMC/ZHUPDcDazGyvsbmAEdDKTEcqdKizP
% tlBlMZKF6ZrWpzGDdySsyCi/YRZU1OQFDFYrIUnHB0xt46yE9aClSgM8+MNOyphho9W1iFvuErNq
% Y0AXaeQDVLyYDOzqtw+iKbYurQVwaAYWWfDZSz77AW9YSCOMz6k+SE44GZQ93EbAm4WY2hT4LMfI
% XAaFuHDBhPU+3mD534jpstfnNSmmr+VUBcWkLALHrhHQ6jhaeQqy+XcX67FuOqKWSS/8kyoKd4Fa
% oKebtuW80zmu366WgKP1sxTvWgCwzyqHqK3NNnLcVFuca78ecgbnUaugMSWmXogQKM+WjBwPJU0B
% iEPM1EGyUQZMBnxhMPZaAcO8fMbzOc8XmHRBD4nboNhCS94V/XyrbFw5oMZyW5xr6tpMkEiSnww2
% sBHx7/+Qe3PrGDypu7XIRX4GqR3AAqQm3nuuHv5TXvzszJ3n4LOz6CuO0076Mzr8c+PwL00BtFES
% cnchVX0XQhlR+8uTi4b/H1e5+AsYelF7YqO2W6iHkXCQEDbosQi2f7B1cLBFVACNOxmi7EYLOOKv
% YBD9DPCl5FQXjRjl1t0QHyjQpC07OsYxy+tNvnLydbb6g8gqj6np6DHVcsgbzPOh65wduiw5cnut
% 7jRCD2/JzK3XLMHYmzhuOz0rfk5bw2tqm1nBdtTyYBmZkhW3xO3tpuzFwtPp0OrkQ8xb41RCNtP3
% LCoT6Gio12ij6iQp/ompfP6SsuONtGGjoJACmGB6h7N4iOEEXF4fpa26XKvj41ETWBXvBUr8PY/Z
% cINzEhp7SQPSL+QHGUDpANkK8FiCaUuwZHNrMRmc3EpusUgWV7BbaYtKfuCPJ73M0Z1fBdV+Rb+5
% PDbXPf0bwObSrxeKZcXLgzhgMYUlIj9AQTXGqpDe7iu+P4jH5u1j7y0sUVYjEBDlYbSv+opPfzp4
% 7nKpd0l8QRff8DCnH5rg43HduQWDysy5CtwXBqGy/TXzZudzPju/wP8+oJ4M0dxI8SHijdsd9O6H
% rK0/6lJehloG93VY0ngN/n+m9ofg9hdgKCcK8M/SXXfcc2rTnAtt/NaVrb+sQ2jI78/w+m7fRNWb
% X39lP/3jzduP7Mdf3vx8/fcPH395+6GODSfze/NjsCkAZXjfa6+VwAB7Ci6WB7nUk0JWthFQKZAW
% 9s70+2pBvzrTL84T0PpfUEsDBBQAAAAIANKpVV2rLfDX3AEAABgDAAAiAAAAb3JpZ2luYWxfY2hl
% Y2tzL3JvdW5kMDQvUkVBRE1FLnR4dEVSy3LbMAy8+yv2A2xND+nFPra9tUnHznR668AUZGHCh0KC
% tvX3BSVPcpFELbjA7uJl4ohvNVRPKlfGL9KRg327gtdUc6TAUfc4php7PHWbk9yRuTBlN8KgLFxA
% hlEE96IpC3lQKVxKu9rhB1nlRBPnTvkOKShqF8inyN3myO+Vi3KPwqoSL2UPn25bBGOrYYtRLuMW
% 9/UV6L5F9Zqpw7FGlcAIqWcP6a2Z6Iwar5xlEO67zXcZBs4GGLkfdvZg11ppmkzgATEpCC6ZjOS9
% AQvXzkYdUg4UHePM0Y2B8lsblXpwU0PnYiM4XYV/aIXdAt8bYBK1rsZ4CaKlw3O6srcB2z++K5u3
% HhNzNj+vwjdQ5mUix1kfAo5cC509N8tNdkGoRe2gJBHF8rEJG19e4iG1OM5VJcVu82dxwVE77ZFr
% xDTrmCIWe+ZumjHkFFZBLcgZvWTzJ+W5wymZsw8jc7GO79VAnObwez6s5V+eLA6T40v6wJ+r4Q8b
% Ji9O1M+mKe6mnNJgDegSU2nbZXuxkHzFyN4MX337R1l0DGwVbUDbFYnO196iuY22qY8+fYef9Mp/
% MYhv62edP3fqYJHGFp9HK+mlfLiCm7E/SrOlbctum+ne6GIsUj7pN/8BUEsDBBQAAAAIANKpVV00
% maOs3gQAAHoJAAAmAAAAb3JpZ2luYWxfY2hlY2tzL3JvdW5kMDQvU0hBMjU2U1VNUy50eHRtVkuS
% VDkM3HMKTsBYsixbp+nQxx46ggYCipnm9pOPBV10z6I2VfXyyan8uFdN2YGP5mmndde2RltHk30a
% vjqkOqOnLyKNYqMp0fw0Xmnt/ftGD5++/PvXV/+6v334WuddZXDfg5boFmnhk8Vkrs57T83VqNZa
% LXtvmSU5XLjWLNq88+h5BXnbz+8kANDFqYbxsZ60dFSfNb3UucIxFZEwR8NsIoNH09JBNjHKC+Q/
% +9vj+fnh6893o+Pxs3XJMhujX4/LcpdT0qdZYvhNK1rr4XmOxeA93edqDb/Yn5CP6bfHL58/3J5v
% 78TVh1bP2v1sy8iT3DIkfVxsnhUrrLNvwpsGtzi2uB+tI5M7A5kfnnY9/ni6YxXHnG2Y51zevM3u
% x8/gtTqQFvPa5N28xTx25hk2VqV6Sai7sr1FvYjtow5R2xVJZ4QCu5I309wqSwHg5Md2iYIkrr5d
% iueYNrC4Ne9RX7gVrLm0wuYe1TyjvK8ooGA2VWZZfQRoT2MwZSm1cZS91LMX0xvUe3oNzDbsiNaq
% odRBro7OvCWEmNrkGtnHVIipOkeAJwgBXOUoOxd4f/j4+PfHO3Kn4cHRSOYGAWVrnzUDhM/SFLe5
% FNNWpuQe4Dj35RFtbsq7m73GvKg1iR0WGI6bieWkvitJyInycKNjhmmFj5pDjXaOexLrCtLe7jBf
% iPUU3g5pL1auhHEhSGq59m5jzrZ1N6GApdq2cS4Pzpw4h0brtV5j3tMaArUeYhkdNCoFl2+DY6Ei
% amvJIDhFdJqTDAdsG2sM1b6wsrkBLQ/Pr3hdu4upN91Gi2MgBNZBkhzyMX2fc8CI/sJciXOPc9FE
% 1kwtmtUb0IvYtjczt9aWZZsbUq2lq8SgdC0uvqjAZiRzQmhdETM7ZIgtlRN2B/rC7LaKAvWXbydi
% RiEcE+9zhvElubM3mJ8TGlhjMjyMbSkJLNG79Neg99Tq2u3UkvKQSw0bxl9zbh7LOzK3xGGmjqQQ
% DNinBl5C2xjZNqUS2OPhyZ/viB2Uw7CTkciShXiCik4cZMCVAF5wJ0LrWmQ0xkto2kXTPLSRbcGv
% IC9al9VkxA8lgh3BSUUYGmGNWdARmSodknCKjRogiB+6hhYDzhvi7TfkC6k2rM04ZIm5ijw8rqCC
% WR1H9oEQvKRKdNig+BQ4ZCgW0AschfwJeU8pkmLspt1HQIfoLJ5gFgDkC57UkfguqhpTIHKE09Lt
% OLxNeIJ+I+9nz9uDf3u8fXzat8f8FV+8wZ/1IDk5kI+EQuitcM5AGrhstCAywqBBp51XhZylVBaO
% MrvG1ocfn27f/D6/rVoZUnE2H+vMrge+6Sq/SMHKBypiownJdSg8VtLQv3ADdbQbvwG9doY+4XGg
% 1cXempntMTHUVTeGLijnbOiWmajXiX+11XZDMeG1Y6Of7kBftnYmGn7phLQgnKsE88yCQ3vAsZB/
% d8EdYeJuAGoIvYwgX3AXFWZo+hr0fm8ltHwrZVc12MCxl7zCZA3cNQSJHXRw78gGO1MfMvLqGjc8
% 13xBZP79+/7+/Wl/vt2RiwbFdcAClxK4kh1JH0HNutJ0VEOcBackDVKEVkO20JAj7VrgwRr/B/ai
% d+NEkTMgsutG0yaEet2TUGwIXfzEMZHDPJHiyB5FIGV063hfdIr4A/Y1F/8BUEsDBBQAAAAIANKp
% VV2/LI1O2gAAAEwBAAAzAAAAb3JpZ2luYWxfY2hlY2tzL3JvdW5kMDQvYXNzZXNzbWVudC92ZXJp
% ZmljYXRpb24udHh0bVBLTsQwDN1zineAUVUhYMMKCSHYoLmCcdxikSZR7Jb29ngYiRWbLPL+PlOX
% 4uhCCZQzTHdwXVoWF7Re62SPf4DsxI6kNJdqrgzjrs0t5H0tsJVZzKY152PAW0nSJJ6wT+KkWSJi
% TRp8rpt0BNJVbLw9jXen8f40Pgx4r2HzYU7FdRPM1DDVtSRoiZRN5TtcOJMuNuBVzWtXpowSjtkP
% UFBld+klPptEylWF2BkkB0t3nVTSgKeYdX5+MWxqa2w8wJ/CX7+QWSxZLt0v17h2D+vo0tFoln9E
% Nz9QSwMEFAAAAAgA0qlVXZ46xr2UAgAAWwgAAC8AAABvcmlnaW5hbF9jaGVja3Mvcm91bmQwNS8w
% MV9sb3cvdmVyaWZpY2F0aW9uLnR4dKVWXW+bMBR9z6+4PyCpv3AI3XioumrqS/bQSps0TRM1bmsN
% bAamXffrd01IAglVp+YN+55zz73HF8zVt4vLW/h0ffF5/eXm9vry5gNYB/fGGq/B68ZDVbsn3UDW
% vJSVd94oUM4+6fpBW6XB1T3Y6qYBdw9XP6uzWZWKOXyF0uVQpWy+zVfqzMI65ZTSc2CEz8FnmMen
% 3XOZ/UkZUuWeSqepkrDljrtZBLKcAajM5ibPkHHnWptn9QvoP0pX3jiLTZTOPkD5MZUhzfcfqBbv
% 1eS0mlglJDztFGMiTlFk7D+8YXRgTlhs3WHiTXsSJAzogrAo2iQQ76r3bYuYEGOLWEy47E3i7xJN
% 9qJ8WpQnRA4lEyJ4dIIkHzi7mpaM4mDtXpMLEq1O0hy0yV4RjQ+8xb7jk0TFcACTV040HpkrGEko
% PeU8afzm1Cdjc8PcL7eD37xY/6jD16ebh3XKzjtyyIUIU7YlWU/heI/jQxw6eoyUPXJF5Bgrj7HL
% HivJKG9MlhPYpAdHiBYH8OQYH9Nta1QQMaolZiSmI4YY9jj24qDHDimm3IiImMi5q5qTMKVDOE8m
% Mm+rxnEVdGQgw40RQQZCNF10dIyUU0UvD46lQwq2K1skjJKwHpI4jlQkpnj78mmE9bLDFsRmc3ZR
% FDjpmfJ4ITa63sx7FZ5zvPbcs64XTVuCKnDL3BuVdQj8WIURn7U49IvaOQ8m19YbbzTGohAqwtA3
% fuF1XYKx+nebFX182QHwHs5b5RfW1SWG/u4yRyNznbH+2TQa1KNWvwKbBvYeUegnXeAd3lqPUSFm
% t4+6r1dlBdy0d02V4c2O267WWAu+z6DwZc3DrwDe8MZWrZ/jb4IPAVfiqqsFySo40nWtz2b/AFBL
% AwQUAAAACADSqVVdDR5ve08HAAA2EgAAKAAAAG9yaWdpbmFsX2NoZWNrcy9yb3VuZDA1LzAxX2xv
% dy92ZXJpZnkucHmVWG1v2zgS/p5fwesisGTLL3KaLuqEAZLd7HW/eIF1gB5gGAYtUbEakWIkyrUv
% 6H/fGZKSZcfu4fwhkDjD4bw8fGaUX/41rMpiuErlkMsNUTu9zuXVxYcPHx63LNIkSWWqOYlT9izz
% UqdROSC/57wkMtck4oVOkx1hZFatSsUi3n9a87zggkSVzpNkAHYukiIXJCnAWprLkqRC5YUmf7gF
% KxZMr2tJAoK8SFl2cRHzhGxY5m0D5U8uCPxYWcKhZGtetpStSm/r35CCjszK93WacbK9VJSOJmQ7
% HFIFwh4NjbTguiokKazhshJoe8VKXgayOWFNw3HL1lNRcSvAn8hjqrrddbNQUjDjqfy7twIbIPZJ
% khdkRVJJjGX/EhYb9RSOndR+4OklnNtI1106bl5cpOvbcIS/gHT2RSAlS7jekYipjg0m4mm2zPJn
% rx3I6Ia8utBtMK+3ckLWkA4QdKlqJ2Vt7agiFbz0pN94OVcmIoURFUw+c28cyF7oYywsyzx1GRuF
% uK2QSu1BngbXPmj6C2v7GRIViKBw/rlcegWotD2pE8oCT/VDvyv2aWV4RpaW2rMHhRCs35urXrgw
% eR4OlT2qyHNwIPi7Rg2m4p56XoMtYxoOhh3wcGlTkVeazkcL84zHFfuQwuBv0N4DIaZ9z8TDIB4w
% gnGAIb+Lnt8HffTMWTUO9GjcrZWOS9zYoQ7FzpcBU4rL2GMH6QGBi7GSnvMoyiupS/rWgdN50YcM
% kigD62mSRsxcu85kFHQquMt9TA1JYy51qlNuBRkrnnmp+5oXAmLmrxXL9lJV5HEV6b7MCwHr/21Z
% LHdSrzkiUuVQ8+9pyUm05tHLsTjjG545P0H0w7gNYAOcdB7/c//bE/n9z/t/T/+aPf352+wG2KVm
% Hg1+gWK+Ac5h5U4onaO9KJcbDk7LiBOolFWWvCxJnpDHpRp0/KaMBrrzq+A6+DUIwyC8CkJ4+ByM
% r4Lx5+AKVka/LvalxSzR0yBrdL7Sk1BqVXxT0vleH92QbTSNR6M2nMwO6hipObY39g0x3YC1Ggwb
% /2ATXEJ5aTyY1HDatGGEP56B1ldKw7ZKeKBiCzP/CX4WvdYWUdGav9HPvmUDPAIOAwg0QuXucHe8
% d9tWPeko+qZ+BOQrEgGBl6/w4oouOJNkSiFJwOJvjTVkhk3pY/JGPmhrhK2mb6KCF8G28MS2qPGj
% 47cp18YOkNGprHgjwZtQUkcV4cf9lgRrd0Mejyso2hW8fldAXRcQEdSQlzosmKuAvg2vT66Dqz20
% IcxFLefhx8WxibpYJ6/zQZ1c/PqOOm8mjzWOxKHNpAGYfl8pOJLJOI0ZlGYFh8es2BG+jbiyLZ2J
% XD4TcUuvTb0e2+mfIfU+4J+ItmCh2jyIuZ0GCWaXy0rwAg7ykjI4TvCsRxM0hlV+CBK4F19o0/im
% vfBMroF3v9xNT0lwfJj1o+60/+Df0nH3Sw+eh3s3YePp1J8nzCb/vxA7Qc36KoPJiDgWnYC6QTl0
% 7zwv4lRCsJBCGbfeMSMCFrWzZLoyMXbKgHxPYVhqKYfQGowsJqsdAcLFFkpymcFkBtEPGiI0Awlm
% 2ZsDAwVXwUfgxE+LoHmzzx+Dz0H4Kbj6tGjlf0+kqApUujiqDS259uZqcYSrg2HocEstf0W5GzxW
% /nsdh+LV5StOdbMBi2Pv9f0xLwf8Gp4wtIV7vep2X468WrxTdLVqkZx/0vMt2tjCROc2dOn2fISz
% 05GhPA22h+DflmfygL+dmTlBBbMCOVEGPSkkx7Lv9uxO5+aQvna7SDK74PV9YPWtcTk4YHq7zR2P
% ffHd7qaTnB4aWtdj1gwH0AB4NiEcbFcIaIQwN5cH+IcXsGwuRJUxG6mMOVIVcF62c8aiXChWcJJq
% ez0Y3IEEuldqlEiGBOWai+qzGAeWgifpdnBxAtwtaE+xrWC/AQITAB+1B4sdqUU574eLW4qKE3hr
% +NWs9yBj9mmfqRXXjNpFcDyDgAtwvyi4+0Basw0nNhXwTu7Q8nFr+kb74QlGvWno9nzfQnO949th
% Y/nWC28zLsF33yQanIQlCE5Mvh23lQ2gDeD37daBDnUXuCIotc9m2bYys3CGnWEWsUqYl3M64MLh
% faebHiRg37WwGRy3h4MdrU4hznaKuh+IVj8QP+sHmIE7/NR0Bh5c8L3T49VPxuV3ffv/aDUtb0yp
% oXVg1XuhfZmc4Dfb2Fv+2GFsCiPUj4mZv+h+7poFwrcTVioqMZy2RA8o6hy28nUbbuH4CGr1Z9/a
% 4LXgZRpXnGLxD76UzUdWipfCsF1yR9fGuJkSkvKwDM6KSZlRF2iLUrd+9hKcBIE7lbaN/o9qHnzd
% NGVxnzf3WebYzB5gLrnC57j9lfISbNBDa2QARCVK+LyrK/X2AmV529SpdqafgCjtoB4BR9b/hCH1
% P2FSGCFIBKZiZPMcU6cqHZj/3TBDmZU2JAOb7f9ycNzn6BVAabmUTPDlktLOcilYKpfLzsR8dF78
% A1BLAwQUAAAACADSqVVdXmEV+zECAADbDAAAMgAAAG9yaWdpbmFsX2NoZWNrcy9yb3VuZDA1LzAy
% X21lZGl1bS92ZXJpZmljYXRpb24udHh0tZfbjtsgEIbv8xTIVz2kK2OMD+191XeoKsRikkXB4AJu
% Eq323Yu9q9jZNjZRnBtLA79nhi8z4/C8AiDiB8ocYbpuDLdWaEUMVVtO2BNnOxt9BUlWxmW+ftPW
% jeTd6k9vA/DcP/1O45d6TW8Zb8GTtfVWerJU53KQ7pTeK7KhQvKKGK1d5zw7bb9m94NKSaiqSEOF
% GTJDg66mB1G3NampY09CbV9dEcet49V5xM7Xu11nWn7at85o74DpVjkiNqTTD36GYyltavLoVd16
% lH6yv437kH/8nEf/kZDG6D/0UR6JUGQvVKX3/q0NlZb34pf1vzwhDAaaommgMA4jmsXBREchbyKK
% 0DRRCJdEWgYjzfEM0iIMKYyLYKajmDcxxfkM03JBpii8TGEyU6cosE5hEV6o46A3US1nKhXB+zV/
% ckY1O6Oa4WWaH+JwqNlSlYovMM2G7ofwTu0/BRVCtEz/J/kVpQqXKtU5rOUdsaIzrMU5VgyXwYrj
% NBzrKOhNWOHFj1Ux4oqX4ppkKHy0xnNkkywJnQN5Ej4I0im0/fnC2KpWyun56nFcNWC73/UiW1zk
% wbMA4WSGLS4C/7IWeRk+DkbSu7F9Gwgex3UTYYDrn7/624HX6Q2xjrq2O2r0XSjhOKgE3SptnWAW
% aCWP34DP7Ys2FTegowSYpP7isRGMOn/9WIOGsp2nsQaGb1tJjXBH4EkC23DmDJWgT8sCajjocuMV
% EMq/1XDz4PjhIVq9rP4CUEsDBBQAAAAIANKpVV00yakaRwgAACoSAAArAAAAb3JpZ2luYWxfY2hl
% Y2tzL3JvdW5kMDUvMDJfbWVkaXVtL3ZlcmlmeS5weY1YW1PjuBJ+z6/wzhbESpxAwuwtHlFFbS2c
% eeAwW+xbTnCJWEkEjuyV7ODAcH776W7JTsLMoZaqGTtS39T99UX+8YeTypqTe6VPpN4ExbZc5fqs
% 8+HDhz9qMS+D+7zSqUyD+UrOH20cVFptpLEiCwqT54tA2UDpoBCFNMNS1kNg7CxMvg4WBthVrmF/
% XeSmDC79QiBscOlo5nmWyUOqVP5dSberSmnKPM/avXm+BjsF0Xf82lqUq+jB5jqyW9uBf0MrSyPn
% lbFAlqm1KsPxKfyxTioXwb0qbVizSSd4WqlMBjW8Bfe8Ph7U8VbJLA3uh0CTZFIvy1XIBqO4HvD7
% DjEXIpPbsEBuYS0YFxRHHzk/C4ROA5FlYXGkgkVuAoU+MUIvZTiO0MKhsn+bElj7I8Y6wZ/8RfXU
% UfGWeBQV7LUTGFlWRgdTW63D0adPYaj6NTsqGJHXSP4ne8tasJkzEtTNV0ovw0wuykikD2huIZTh
% L6+xkRt4NIf/y1QSz58qW/KXenK6U4DMgVrAD52XLsLKvMYQm0pyilGIXCxeID74pcisBElOLlGh
% 4KDm9D4s8gIlhgwXUckWZVIwwMJpPWNEjhq3jUawdeKko520LTMggPXpdtZQoRUT/G/q1mecftSz
% /shZOxRFIXUaum00AGQgM8meeF/j8dAT4MB04RHyrqXPHAQOl7IMt6wx/Rmz4d+5lgHwkX+I4Jnx
% nVGEFVTxzLwk0g17fBs7G3kde6uag3tuDmj0O43DDwJG8t4GbeLO48Axh8wpTTUHKEYG9S/5uGf6
% 4/gLb8AdC34ar4V9LDiCr4AU6AQ3PAMTQnLDl6mYMRZ/3lsi8jvcuEMewWBf4SL3EK53hn4Gd4En
% EW9f4EzHRLfbvnmNH3lYQJ6cnHxs8+zxnC9PTsZtoj1Tks4hgoApzh3/M/KD6OFGZJW0IWZakw2W
% T2cdokp2OYMiKQjXvM2aG0qZ2OuFMhBegwJ83rC4ldaA6po1MYgIJ9dDKFprUD1xYLnjcPgtuNnk
% VYE2xJqfdr7NXTTiiS8HIwyf4lwA1K0MlrFjbNSRxx2PjnT/Cf2s+/yp9ZPmfNkrMGIXfHo66+n/
% o6sUBpDZRGdDRA8t0L9M1cwVmA2uORumDzO2Dzi/qmaTC4Smk9ihDBJmHoJDGO189S6AblGSB2xe
% mbl0r0o/7iKjoisULHW1lkaUEB9SQQbf8qtD11xNR5OZt2cV1YeMt7v0taUs9gusAWRNnIG303DV
% x312hPG9ZU11UD9wQQI2/Aq8GJPpTQw2zGcZOfXG5bA/U0OjXElofXVLGjdRTetov2fD87/PVEeo
% UKYKnbeL6DVStHCceEeo6OEQhVSDoy1vgwUiojaeeDYPnBBDdX6+Zccj5LlosRtfQEWi1zomK7CK
% 8ZF/3+I7BJ2HeoAZO6bKQYVDU+Hw4jFn68OcTXcnvWAt4VrUIYkGCkzPzo7K454O1cA9by2vneVt
% xXbUUMebA+6dkPXx19ZzgSbgfFKlltZixbRtir4DSA/GCVTUUmkq0x4tAFXEZpxXJb+Ybmax0paj
% W+5gheojQUhLtVyhpjdY2os/kXzlF1QmYeRB7NQFDEwy5S8+itcQVAzkd1ABfX6HtW/FUtW+o0hT
% iWbsGAz8rhIta2hjrsi+KXTUsjlX7B0jfAhgMHPtggxg4PdGEdV1l4MDfDTrQGIOC+z+CBheRa58
% U38c8489F+Q7jC07/m/ojvUVH1solfv464+RzxtG7J+AfzkYx/tAIHD/GHyGCRhTVOqSyhiiVK2r
% 9aA9ZYQID2AwLhWUgUDAYde4kpsUhuVhJyilRYfuF3BGvtOfOIynrqZZGGhlGqKfqORMJxPMh1GE
% PqEVSLHRmM36U1eyr/YKcUNbOJJd+b6aYSOEugwutTT9tVvOKoJzg1FErAdpjDPFXpfHLWg4IsXe
% 3YDycAJ5dTMUagJinlI0306kaE7r+h0xDw06hAyjwzqnwOAOIbiswJswuuR6OaAQegF4n4AxH4ZG
% zFE1B6dv3V0FUAV75UpiaAYUtrmwEoOxAGE8RNOcAwBmGrSActxpKgY2x5rzvaHYG7ubL4BvEEK7
% ZVg/dJo/NYPw7wgmgNOyjzVQUFMOf6Ul2OifOW1CnJ9OPKMQPfgZjoEG/wOZrNcb9woiLM756Gfk
% bmxzTO1NIVU40XGY6biJlnwZaa6jR50/6WQhYCJPE9d9wdRI4rUu+Rc4JYGsSxDnibvc8X3oR242
% bf484pMmlk5g0sB658oIBR9sok/fSHOBdLmYqEWCPJ76H3i9CZTDB87bb8Tr3KwTurjyRfdl+dqj
% y9dL8cr6L+Tb1260o0ngIrsR99k2UTrx0XAP1tyolKZe6t3ku4PZHyvGp1SJmnH6IP5Ene6oz4Dk
% I3CMXPHSwJMSnRtY9u6OKCtKB2bgSYOSp30zoE4QgD/4ZUiNt2d7ZV/3QtsvGdK7ScJjBejOQzCE
% 9QCsiL4xa3//HrszUaELCo+yXXp+g9vz091m+M0uQva7GAbQreEGCCPM4Q2ETgwveOZp+Es0YlE4
% GrnHb/Q4GzWLY7foH2fwGP98Rps//foLLM9mmA+2yuAu23Uoh25RwBJ+CUjIpR7p3Yl7Rt3GsO6k
% edvDUpc+cSS2FGUFFN1LpaHrQbaJpc4tlHpLtSemCkOlPsC8CuYZeEgtoBphfYrgGjZ/hJyJwLxl
% lQmjyi11PItdzojMfWSBtmEkflXZyPTwm0oXymphFPQt/NAxTKt1YUN31EhRW+JjKMw59Kiw6042
% RMpu1H3qsuETKJTvsfa7/9Fd1oG0ShIt1jJJOO8mCeI+SboTh//O/wBQSwMEFAAAAAgA0qlVXVyM
% HzmbAQAAZwIAADAAAABvcmlnaW5hbF9jaGVja3Mvcm91bmQwNS8wM19oaWdoL3ZlcmlmaWNhdGlv
% bi50eHRVkk9r4zAQxe/9FO/YUjkb/4nTeOlCCCkL28NCUuhpQbUntbAsuZLcOCz73XfiuIWCLpLm
% vXnzk7aDLANqqQ9RTYN8tQbSVDj0WkM6klAVmaCCIl8gye6WMNZQ24UTOqtPXO+/451coAGl7U3w
% +BsXuFvEAikr8lwg4f08E8gKZItUYFFgmQjkvI3/4fd6t7vajjHC0UbKdH2Ak0FZIzV7OkdloCpy
% 1oYxm7avUIbeeqmnXMv4i0vtiCafF+lOJQ/gVInO2arn+7KmsmHV6iLa2LbT1HIR16IhZ0ijVcY6
% P7brjWpZqKX7xqNH+yc4qnpTSRM40xFM7WzAhml+cdzXFCRayV2HApJJJvPVh+X+SXD6ijoyZ7Lw
% aojepVPyRdNEUMDX0nUfPGOxFPFKjGE4QWm15vGl+7zPBK8c0Q/EySXBbtKP7BhkcU7Jo0+PeUL2
% cN3cRPF9ep00t/HNnwQH69Dcz2ezdH7xWD8+Yvu83uyx+bnd/NqNpzM8KKMCIZDn1r7vJnZMnXB0
% KgQyZ2T8JTw/HzO3Bz+7+g9QSwMEFAAAAAgA0qlVXc9nI/FaCgAA+RYAACkAAABvcmlnaW5hbF9j
% aGVja3Mvcm91bmQwNS8wM19oaWdoL3ZlcmlmeS5weYVYbW+byhL+7l+xp1VVwItj7KRV7W4kN21V
% qae9aZNIV0KcCOPFIbaBLNixU/W/32d2AZO4Obcvht2dnZ155nV5+dfRulBH0yQ9kumG5bvyJkuH
% nRcvXnzahlHJZkk4T7OiTKKCxZli5Y1kP7N1OmMnLF0vl0m5c8v7jOYzJVc9dlGG6SxUM7ZMpipU
% O5aly10P/DqxylYsVuCaZGnBklWeqZJ9riZYWLAfhiYppSqzbNnQRNkKAoZ6H2e5ymbrqDS0q7C8
% qcnm0YzTIk+KO1WtK0hDLA3FTz3qqHQuzKt10h96b+3O9//8/HYhfMvjfZtbrnn0uad/XXp4eoAV
% jILOl4bWM8sVldsw8CoGQWcmYxaprCiskE95ZI+ULNcqZdbU7wduiB/bsSLfo3cvsF3M1+80X9No
% PjeA3MqzJC0Le9RhuSiglpxZhSzrabvDkpgtZWrl9nvh1aflHaYZhMsYC9jLMuEHeJBVtywBBU2y
% +5tkKfX2zD71GFCqhM98dxBw/HoB34Jzf5T18iy3bNqV9cI8l+nM2tKwOjLr1G/Vqf4Im7tm4I9o
% YKa0aqGSoZFsL/6wlv6H1bcbbuG0sIr1CjySAMg4uW8lXc9+ZTYFhJ5e8g6X+oHWNyF94RpzaZkV
% /Dn6YQ0MyHm23FlTcnIN8ktmLbN7CZT4GkrS04x31XhXjbfdegZvEHeJDaDFJCgKvi6EYTpWN4Xw
% seQSxQ4PkBR4FMF4o42iZeS3JGXb9y0j8hubD7QF4U9Cey60HUd8Vg1ug/FMliJ0Zu7UiUAHQNMM
% sSzLUZSlZZKuJWa34ocFSbBXE9LrbcBBZI93WAqrGbeiifSS4RaSFzrb7p2zey+UFtfK+Z3NFYn8
% kOSWloRjq22PNrVzWFu+s/d21M68MZhHQKZsgf4cds1mcoBl8V4A6vdiXezduDIr7Yfl9fyuNU8c
% MW9OTWYScJRta5+LtvnHsXgkWQdpqkB2YlPkJ+vcFkK/xCboCOTz2mW/Z6lsyAmwsIfTshUZM1PY
% +HjoaUlhUZL1HPwuhcbHt7ZdGAMW5zu8TPFilAKWmrS970uwl1DH06WNPR4/hqAxr5Tg57xa4yT3
% S3Z2I6OFzuyEhpLsbh3OFFwu4lqllVRyuWOl8I4GvSY8SXlv9EH0O0wuW3MDzA0c5GIL6ciiWD1H
% JFIaO8cPZTKbP1rxqhVkA3IOuSwkWJB9HzPZR/K5XfFLDvg9JvIqIs36T5F/ro/UhY0WLAJraHNK
% Bfg5IbAp0Fq2KJ0Qhiid6f+xwhMzCKFfzm1s/dAdOqVTNq4ccyNKR/skosya0KkT4fs/kFBbrm2y
% lw6ySTBOBe2b2OMZwtWr9FjsFUy17HmyyUqRym1pWclTDBY8tcmgE0JpEdicvLYKcr2RJYX25MeJ
% uFn+SyxGE2zkE1+PA1G/cJqGZO4M5A+CRnrCEQ9V0WmL0fW4kRZeVslyRHSa8PaAUJMgNYnS0Yxv
% gwbMWafJ4XNKmbqgoC6KZVKU1h+yqTsgi8PQXTgPTw7cRK+Df3ET5rIQv36Po7DAS52ndVHX9dP0
% JeSxBVcyl2EphlopE3ci7E67ETL806xj8NyMzAn+Bu4aCDPozeENNIG2ouuZk5tsWoXzRscAmgFF
% YRC/No0bVVn3Rm7DOfVWqOIxPFh7Y531ElmM2C9yIc3W/o1gT+UqL3cNemO2gQ/LrUnPBftlpPrN
% zicXF6918vhC5+BgOUfdo3OQPLBFsVWyhUFkwe4TatGIQoWQABzDuZxVoCCfRJR+5IwSCeF5vYfe
% 6+OPRhAJMA7RUU15PLUF+rdeEa7ypTSiw3qgudPz0U2WRNLyYTV+zE8Ce5zreWoEK64oVCBHgoER
% qYpZd25uO7t2PFMJg2HJ7rrepTvr4dWdpnigddpst4tpJMo1yfNwdPSUbBxHVRGJ7HGVFeLoFOqy
% Cf/Az8SxE4foGPGYmkfkelj8KO6cO+fMzZ3cmbhGSGfgfNjnlo+noq8x/+jg9dghyoZu0qaMI8e5
% OxVx6IAknjqGjOQ24HeF99SB0NXDrvm6ZErHC2wXZUrJCP2mqzKkBjp5mc2hqITRl7VHVSzbToIM
% EmW4LKQyRKZVCZBCI693sXuZzG9K42pJvDOuVgdTwSQJgxvEv/nJce0meZioQ+8Y2mNjbKopD1Rq
% boO9jWiP/TTPvGlbfvtquM/BB5ZvrGscYHt09JS6ZXXHGZ4K0s3aIqhbvTcJMf4XY9woKStzTHGr
% ihDBKolqmAw2fwJfJ8OFVCnKqs6FJnkBw1miAAnuKOYWs7/wNDed/qN7jb737OtPHUtvNfA0XUhE
% dQtBnc6/C78VlXSm/dR6SmPNvqHckcqJELcH9gBNVyfw78j8Bym64vAVHNyQu9OgXY2/B13fhzIB
% 96FfENRV5fn+WnUHdYfNaD0X/uJpbSUa8o+F7pCodQA/I0ZT+dGQUDH3fZXdo0rtWRDPqpRn9zTx
% DQkcHUK94SuU5F/hpZh+T87ATAy5l1dsnSYr2HwZms0lxCnAaRWaSzRCMkQpnCG3hggjIsmRhKXa
% SNyVt8kKEYeeM1OUeRm7Elq40RNpaAnKqVNvdIV4EX6IlmXaRpUSJC3xK2qu9IarujDpOLsyJfrQ
% VIe2BZS6Y6plKJ61C1xRVXYh2kh/I3jehqA9sEVbrsoBcO9qiYnz97DraHkUj2cZZZYV4o8+bpjI
% qhDV6O/tc5QaizXWcAlf5DNioGOVmD+OVDThZajjlMIBXsvpf7/23nqMkVetuLVrm7+ecXEjUJ0q
% F+2APWkCVqlnwDu2+cLArFGOnr+HLv6A8LcDhKM2vsqgOzYiamwrFvtgae9uKvnTiPnrD+oNPaNd
% lZKRvgbOgpvfBR86i1YyFgJjZwHPRle5L5Wov65Ha9aAFqiWHh5j2jrhE1/DG78UBVO1LqW+vNCV
% Tz+dbRU42zpwoAPfUMQ/1Crxh3rtG7oOEyKbdlNpzjUS1c3lG7t1zzAHP1bpJaMPYdZbm71nltfV
% A++dbR8NEN0xq8enx73GBt67U+8N7bykikMd4b5wr0I1T9J2Q4c2r65GUoFJThXlGPbSPbVFdjNv
% Az7gVR8dg6axTgjboH3mUXdWaa3baa4LSV4EjVzYJXzvmOPfm8B8lap5DHlfn3hCiHoDvegNqNLq
% QhsX5pJL5OF0KWHCHYpfgfo5XeukuVpbC1sM/rHcBRU73HfDlC1wud03PEw3PFWbQ83IS3axW02z
% JWow5EPjo0v9iH7BT3hcvy3q94x9sgZkOzHcY33svDWeNtQ+VmeYS8oB9GFTJdsRfTlgv0ys/K7z
% zOUVBz4zSdkW0LMi2bqbUCVaPdOqc7qvqLxu3FGnufeOa2we2bReN8gy9xTQNY2bkej1RcVJYwHA
% RqYpa76esOPPABCqDI2O/wxMvIh+rzfsN/kNBeX6Og1X8vpaiNfX16swSa+vX9NnzOayNm5alXGV
% ChspJn//zT79d3J2yc6+fDr7eqH59tjnJE1KyUpZlGSDvErO+nPGvUrKUqaUk9FcFih98IgsLnoQ
% 539QSwMEFAAAAAgA0qlVXQbLPKBIAgAA4AMAADEAAABvcmlnaW5hbF9jaGVja3Mvcm91bmQwNS8w
% NF94aGlnaC92ZXJpZmljYXRpb24udHh0dZNNc+IwDIbv/Aod25kmkJAPaIdDBtI2swx0Sg67R+MI
% 8OLYWdtZyr9fBejXdnrJQZYevXql5D+zaQnPWVksF9kcpvlzWdwX06zMV5AtZnBfLIoyh1mRPSyW
% q7KYrnpP2Wp1C/jCuANUVaOFcsC1so4pZ2/opdEKlRNMQs1+a9OFgakKovEYKlGjskIrz3ImhdoC
% 3yHfW//CjZJwGBKubiQ6hAZN3TrmqAAO2lTEV5PQ90d3l8yzjnP+yynFQ4Vme6RObKu0dYK/w2Pg
% Ry6RdFkLa92qyt5BnAQUMHuk0mqLwFv39tZVaCmRnxTUmsS7N1wQpe9KK+Q0qsOKQrjZCC4oFdYt
% EWl8oSCI4YBiu6OUfi2spdn7TqjjpSmzeBpu6Pvpa4PhGDatlJchV3zXGvjTsor6CO5ttKlP/ZkR
% lvy/WJPcQTAE2zaNNq6/R6NQgjZrQTahcYKkkUz7occZz7oUw8wRHgyrPbJoT7KR+knhBHYzcNlW
% 3cpopXuPlgzvr6+4Sli6BUWGnZyoz7dwEBa/VB3hINwOtlKvmZTH87s76G47e793P1/SWS4evKdl
% sSg/niAsF/NfcKW0g8ZovaEvkqNor3tqEt1CLdTVzHvsJ9cTb+inQ/SC9OYUDR+92TVZnI7OMb22
% aP6S0i1rJgN/MBhEcZAOBkE6JFT8CRX7Ufw/yYv9MPwOlaRRGARpEBEq+YQK/GT0BUXB5FtV4zQc
% D+Jg3Mvmc8hPP+30MZ/+WEFnez7zIZessVQT+qMRWDpGul+/9w9QSwMEFAAAAAgA0qlVXVDHSTwQ
% EQAAaisAACoAAABvcmlnaW5hbF9jaGVja3Mvcm91bmQwNS8wNF94aGlnaC92ZXJpZnkucHmlWutu
% 27gS/u+n0G7RXcmWHV/iprXCAm6TtEGD1NhkT3siGIUs0bYc3SJKiZ0g735mSEqibLcniy1QR6KG
% w+FwLt+M9Oq3g5ylBzM/OqDRvZZssmUcDRq///776dpxM23uR35GNSf3/Ixp8VzLlnAXBO049Wiq
% JWkcz00cjLQ4yfw4cgJtHsRO5kcLzfOdRRSzzHdZpzHhrLWBxjIn8pzU0wJ/ljrpRmP5fO67lHW0
% yzycbDSfVbyAVIvoPSzl0jTz5z5lmoPrxSkNO41PTs6Y70RtP8roAqiWNEjgj5P62TKksDJyczwn
% yainzdM4hLkwQqMsdaLsT6b9FeewxKEGSyD3tNMYBwEs+aA9wBZNzd24AW3TiKaLDdzFFGX1YX57
% lnsLmplcRDcOE1iTwQYzykBTTkq1GQ3ihw7qssFXnqegUdgW0/wwidNMO5MDmsO0M0EDyk6zOA5K
% GthNmGcOn4frh3BU4k5McOMgoHW2Hp07eZB5vpsJmtDJlsXDOSwZp74TiEeZH1JlqfkPF/QBMjSu
% rsd/XRN1SDcaDWCNZ+7pa2aMGhr8c0iP/53HqbbW/Ehbs5HTJGs+mNIsT2F7YiIqKWZUd8yZMZKP
% sjwJYMSe2f50ypn4yAQOZ0H1gEa6YxgGn40HxvTIpOXCgc8yXVBGhmH55opQywFGpmOvpgR/TLy1
% 6msJdhQOT4cppSR66Ed8wAydNb8QhNwAmJ7IZRmlEWE00w0rzjNiT8vdbwmeGHIG/vPn4jHOHrlx
% BO6R0/KpC2ysFfHLgYelH1BtpUVxVs7Cn47jefrKsNyOkyQ04tcrksBG1ZVwddd43xuBgAWh2LsL
% e1LOBZ43qj2iveuu+U2KXTIipF/oyOZaa7qGWO+V9jX0M/DHgDos4yo1tQfwPAgba/A3WDiJwTPB
% ymh7llLntsPn3RLUdaUqEMu8pRsSOOHMc7Tbkf7Ndu3b6RR+9dtWz3gtqKZTkwuAz8ytR4bcWgK2
% ToCg1RtNW649goupumexBaRC8xB/gcTgR7jaOkJ4arR7sFth+E7qQEyhKdPvzO9STRn53tR7rX7z
% 7kAfwFX7DiSxcjF6ZxzwEU45I/3md5XMui9GShrWIz1x22z22z2L9YnCptkctIW3fSAz5J8brXtY
% Wmf9Vt5kvdrhZmZuzsx784P5AWdn4L1JCoeh/3n6ffzxWvtrfH3+9XJ8oX08/ev6/Oz84/j69Eob
% X55oZ+eX59en2sn5+NPl16vr849XfxoNhyjbP9N7Zq9vmPzv4SFsZbb9eCAf9+FcHMYgqmmO3Z0S
% cqYPhubh0bBv8Mjp2O2eGB28HQ6G+KzXfXc47B4q84Dk+Ewfmm/EHOB7OISFB0ZLPOsXpDO5xDt4
% /FZSz4oV3rzrvTsy3x4dDnsGsDs0h+USXOkon4GEKB2I2DsyGmgVUWUVfXPY7RkjOesMwlH/7dum
% HoGRGMekUEeh6Mn46mqkUZ5IS18A78cMmEE0p+skjiCVQDSGGL2KMSUxLvPhu3eQPUMaMQjsbeY6
% AeZTd0ndW9aB02i80s5lxitSoJrz9Cjm0Tag6yobqzkxEmmQicxhdBo3RO+aXcP6enkKFgdX3OIx
% 3KjRWkcVtVDJEFh7eAWewylZPtulbJeU7YoyzINdyiZStpGyiZSmHIMl5FhXzgbtrSCE1yabfKKU
% A1SFGSbb4p9xMZqZ3NeM9RUunIJvjRPhj9xUuFAz3c1upiNCRWtjT76D5KmP5eSIYEAZi6RxUwTY
% iJBuIQRovmQfOuy2MjowquPjSMkmCSmIiwlp/ID041FCUL+JyUWHUQhy25EtMnBpXOL9+9UfPRk3
% 8R/KhvuBv6ZQZGLqYNpNtHCc0Jn5mYADumH8LJcgRNLHZpEwH4m6s60kOYYk+chFfjTHdoIgAP7X
% GD+isQswihlK+E+aC+CELMcXFzWUBKadxvliqQmISv3FMpNpCZBTngL0eaVOgIWSlDLwQzm/AHTo
% PQWuRFQHGBGxExN+iC4lYGGngYIJvZCuBfABtkaeni3xuHiwG0veSRV9I7bd5RiBkJVGA0a1Xkv3
% W6uW31wZr492T3AHKkUyJT8SW9f9pt8aNCPjda/XHpp6H27h5qg9MH427ZQ80RH3CnTkR5uCM0xN
% /NubQmrl0yhOUwFoib0gyD+XJ5wgmXoeFUSrLNhlpMRVFmqP2JTP5rHJZcqCNWhS4ZxFnBEBakr2
% NY/gk3HeCCkL6AnXZgUkqxkyoCMpSbgBoHXidEgI4EsSprR7qpA705EIoSGfBxNLHiXhHd8r6lhH
% rlzja6MKKY8GhF7hvOgUPNo9Cr8w4Dj86Z4TVLYxw2KmWraJa5zadFrXyI7cXAwuGiQxzqPJfys8
% Wtp3i/RqZt2S8F8YvR1NCeRCmf/mMgE+VdOfpWNBSal6IBJAPoxIv9N5a2lP6grPMoEWHom0sh6r
% FZgiK37Dei91RPVoakuorDwKhSaFMRpsAFFHeeSDOkJeA0I5aWqhDwkUvRmQIdyCqhwsBzfALsDg
% wR90tNMoD6lgLUq7oirGkhgArR/mYRtJYcUoDtFLMFigMUGIAmYedWMUw6vtXSkpUQWA+OMowCI4
% Djit8JMOj60PPJj9gHou9ddgwhCEPCWGnOld4x/ECj5srn7l1pXHPiC6UYISgLGuBD5q7YESEdLj
% u0Zi4/UAcxzO7u4h7EtCqLYIIhDAVg8CRnWbzaNqwjcwfazpvsEPXJIHNUN8azRELf6DnyCEYDeH
% NMXBEo/HjFV3/BCqWznRdRhMhGJuK0YPzLdcBziMEitPSr3vPRRLiCKj2i9DJ2rjm4i4It6K01lj
% SK6NQ0lQceP8n607gtVq6CQYTyBAGtZ3Hq/WnXsnyCG+yuOpqoEvZKuOsWQUyI7FsX05LlyaoVKU
% boJ+ZqBqt8dkLfgRDRW8qvIS1wF3811s40AdC54i6kTwrYyuZUqexRlmar9Mv53SNgM1V0ZQ6xXW
% iA+RCTaNfq7ZoLJenJA5frCTmgouNlSLanoi1RPAl0YLJ1eukBHeA8Ei1RdFqs8r0T0djLp3oHc/
% ZLsNAKFqO5i2yENWK+Xfk0EllvaztDjCU7EDk3IGB2sI+XzSNbcFwb0yCOtvPqwHCO0coWrTQa6S
% 0s9oWFrO9QWn/gVdISgwvngZ53JGQ8lCUgl9qNW+c0O8PiYZv/j7mORi5ALykxiCq/t9ljKoWco2
% 82B6LGpugLRBu28cBGp4KJJZTdPc0dQzKHmWSkeeYs8FXyUCFUzdcMeV+jsjwT7XAn/NfLRXHkSt
% MIa6MBPXM8dTbv8NBittugDheEZ4/RvxJSTdAz1eYtwXPwdQVhBHC+JEG/kYTGI/wCpVgPYtNdAi
% F80tbwFuo1IlLWydXBgqjcK7fqI/86sakbJPvpIb2uhwKEXlcuo/AB5qE83t/196RkdusIdKWlyp
% hmPRJspEh0PsF5xENnqUGaU2joletowMQVfLhS0yaPzK7gubxz0fo2a/cy7fd8Ak3+We1XboAlt6
% znaHq5XDVKULZvyMA9en5PPFUjFA4XRqei9aoTJF8zavwmJr54b5xTC2gawSK56FqfDAInA3A+Ba
% uf0zPElvqScQIYxXRKran5UGvjxHDmW3Flb3pmDoClEqKFJummfG3lBiRuodSJh7AN65kVKhWhB2
% DzqdIwmgz/Ig0K7cZZ5qd7njYSZ3UTUhA1yq7bxvkSVxRxsfnCD0BnlSGvCg02lMPo+vTq+gIh1i
% VaO3xZ+uOYTfgXmIQ/hHNFtC55b+GMvDIV3peydk0JU40xoDxL35p/hWHR6NEUYCfNRPTCVavhgD
% PxLRFRH7wqzf51B4ataQ8UsA8aNsAu1Fw8d94zeir+BPjW4sQfCjNZYwWNSHtXbJ2DwRfZiULvLA
% SX94TubstqESJl6f7E8ShoXyPCUjXyqIJ5QC2lE9YbLg9wDiVg0fJfewqTUhvI71WC0gTQgRfTG+
% 14mNrdmuIOBOAGcMa9tFwZ6YmTzgTPKtLaLuPGEmujUUkhPeduTsGg1uu4D0b2ka0UAifV4V3JZl
% wG5r5mgH9us2lCdmn+dFiL9DkRSxMSlUC4onNRtGJZcSSXHI1qmomjlpNoHvRHa4Jybe86ay0o5G
% 8kt+hnAE/7DGGKu1xG/kpnISRumeyqZsJ62aq9YR/Ecy43Vv0H4D7tvE0RU4gBzttYc7teel0q+J
% IlJ2PBxB6YiOR9UcuSD2dusjjR945wOipOx9wJUp252KQYIZyzYVXuJDofBKAO9O9F0mzSiqNVmw
% GS3kMWdcIj8BozXvLozdBpF3pxgs/lveVduS/mk0T3jDtG/Ab9S+u+AGXXSdxJtFOIZu0csR0Uck
% ne313jRhxe70PVnecbvgd5g0l3cWN+wi0b3SPmFPBHQBgGOeaSxP+Dtd2PIijXPR0/IjllB8F8Eb
% KnE689Gp7p0UIrpLO9X7zaetztqzlcVeLCIGEgg5xVtKfKK2pvG+k8SJbrwIS6dlb67qy0EoaWyB
% rnT/u9DU4JIVqT2VDuWtwJLwHe0LWh9CJpgCp1e8X8W7lbjzi+gaOn4E6ZO/+61eOjP/UTQOKoWU
% pNVG0zjOSDkutGPhvskTPpLqtfF6uvUSuK5e8R5ij4qLTXELRvGd6R+VINXScBXfU30mluc7nNV1
% OKuY8s0V4xzHwhyjFrREyxMsgoj6ofy8QGdC+RyGcEa1eU4QgJPZt+BMN1oR0JGRiLl7Eo6Yrkbx
% yvb/jsCSP6UO/xwiD/FLCVhAY4BgUghtEYcv/WFH0GDoZ1oepZRlKdRZIvhJjxFeMAP+kIkEcLkx
% b6amDYgF0cshH4H7G1OOdBHCwPOuOUBaTjGtzCyMPbo/uN6D4eBCtr5qIZXx+nC6D97g2xp4fNyX
% KUcRaw95yf5DCDuox9O1YW5kt9ncFLHunuEnE/C7Kjr9L4NXW6f5IZTgiuh9lO9X1eJ4AqKhVDJk
% mmLyavpP159MBJqYlF2syWQrQgMWgphHJs3+EBJqe4LgoHLMJdsTv/U3/WGbD5VyKUEd+fCrF4Vv
% vvp7AgupeA/PE0vCIq0IGRWpFXiCZr5VDvDQ/6zNEauLYL6F2Nu8w119lVQ0199A+aG60HNh9gdi
% VCYF+a2VC663rxpRhCu68w7O419zoYu15QtMCgIFULDyNp0b5B622eEUb9txRLXqqaw9/oPfYPEe
% PfjlPb4V93wGhhtR7qSSfjOC4fmcpjQSbwX4O3Th2p6fyo+hOo2IDK0aKOsbljQXA22A+wd/140n
% wlE3Eahb+NlLS42iNgfG/8qoDf7twqT+8UBNAeVmH3yQb1uPG9FQXQTxDPxxI55nDzG3Ja5hyfns
% 4uv4+vzyU3vy9fzyWv3OQ/t6efFf/HogE1/2wS+FshEsG6Zn6UZ+GSQ+F4PgnGzw07UoKYMdx9D2
% oTk030yrQKeewktQsVVC3JLFFYmSziNNY6brl+alYXrZJqFEvhQyrM/7n/OPH+pAxMNFFZwnRDBG
% VzZcmSVwbRHdwy8Cequmh18XHCAcrzF6cdkouvm1EIPc+tZnG4R2iiVN9QbWX9cYfLaFoHvBZCV2
% W5l2QudkwuVuAmu6QYL2Vfk4iB9QaQEIHyw61F/cOwFb6jCr/fngDVqjtfQXy300/ebnNtAZaihd
% OMlP2Bl2ryIrQgmAzOh5hO/N9BO+HnkCgUadweLZ5KN9WAIGUQQ5Gs/A0e7BD3CpJ/gZdd4tnsEy
% 6dqlSaadc7s8TdM4FcqX9i6+KgX9g+dY1aelHAunPMiUbxE1dusD3PHUoIffF4hPpj5+Pv345UpD
% zzw96WingZMwEOep/pVkm386Oer058/4gU0ceTy+/Q9QSwMEFAAAAAgA0qlVXfbl5C6EBQAARw8A
% ADIAAABvcmlnaW5hbF9jaGVja3Mvcm91bmQwNS8wNV9tYXgvZXhhY3RfYXJpdGhtZXRpYy5weaVX
% S2/bRhC+81dMBTQlXcqVXORSWD41TnOxgyYnq4qxJFfSyuSS2F3KktEf35l9UKQeyaGCAO1j5pv3
% 7Gg0Gn3YsdzAR9ZqLZgcK2ZELVkJrRGlMIJrULzVvIClqiv4u25lATcpfJBG7WHy/jqKvq655sCU
% MOuKG5Gn0HBVMcmlSeGjYlUKDJk+f/lz3IhtbUDVCC4RmSkOed0Ijx6ZNTO0euMSOAn4RcOWK7Hc
% Xzf7a/hkNNRlAYVgK1lrFKUtcsWEhEIJJCXISEiD3NaMcg91JYzhxTU88Fev//shBGqBAAdB0Wg0
% iqy5S4XOQSANompqZeDeHwDTcO9pWpmbui47GsWLNufuskZXMFOrcFe1ZRQ9wQziSQqTJHqk5dQu
% P/nTaRJFUcGXwIoiZilkyR8R4Edx0yoUPJ8s4FfI8Af9Op+6zXThmSRfxWzIMWaWdswOVLrNzkE7
% gQSRJUELIyquL+lxZfWAsVPkyiqSHq6sct3VJEjPa7k5VvJUR1mr6uYcGVxdwU0HjGvP0KgaHW/i
% nR7yuHjE6PoUdrrn4PwHHNZ0x/MYeLKSyZdYDunn8ydYYpCfKY8UkyuOFIszR52mQuaiYSX5VchC
% 5FwfI7K5WMw3DmRDIJ7OnYj+iUddd85iWnNMNkz/2MLgF101m8HkwOx0KrlELgx2T3bwSmCdLC5y
% eXOw3J+LphMvMZMdid0WDe7nT5QE8RRub0GGc4KewaPdkYiK6ZeDFE/rMOkjkJhIrjNhnnOsZBMn
% h8sl3c9QOAIhNBFam58OAPTB5MPm0/LusHNwiNKQXlAmGqfaO2/AJrF9J4TopxMh1gs74/WFfzu+
% EyrUFAnJDVR9fpf6qgtmpEHUMFB4Pe7KxQZB7THwl+JAltgkOOjqgVwETG2w76MtdqfbSndxk/ao
% UXyLR5MuXv1gYV0dx2ul2B7pX+AbxC9wd0eV14VhTWx4a4m+WezucoPnjsCGGm1YmXWcYJ+ZdjRN
% rfF52hJEht03tjjvPFsyiK64HF2yEl3r3e93nbvJZZ0cXuIrZ9vrIBbWc+h7UjlUDuEMMjPGLmwt
% 7aduAj/DzVAdD0RSaNmT4GNDWtp1CkMCHxuS0U8RSxsyRBfPnF77LkFG4fn/C7HQSibhS75uFdh3
% WruHHF9PIVfwxlXtz23yazxsS3r199f0Xp5JuO/4/gdF12tgAgVpwyT2410KsaCh4h59Rwg7Qgjh
% SOGfk+oaPdSwLGtGFoybGplhy8rWjx6scKOBRcGJw7oiR8FiKXJm+OisTj41sJDcO4YbPHIWZ1Qx
% 8/g+3mFrQ41oMV30tVX1q2unuLByF+eqqe8PLzebvyzo6xt5d11QAfi7yaKfdAXc9mu9lxX3DHN5
% SDrsC/6YyX1ssX2TOw7bC77CWPgySU7b31lRFxuwOAM6xLwk+lSyLyLXQjMKF3VQb8dpB85CQO1U
% 5HeYaQRDD9RvULhiI9fjZtiDvyq0xNUX1l4Vb4/fcT9LYOESptUJ89imzj4kRgp7suxNNHGbwmuC
% mfN0pCjR2YTZuvRp3TpMMvu85DEvVjSv0GxLs+/soZb86CmwJEl/SJBwN4Pf3Qmlkp9vDgWcAjGR
% OC7bisZZL6nnemrYsaCQUFOTh9wN3vRVgmTE6u2nZfc0Ba1BaCDFD+DdDdbWFGVgd41ZpkMfpkqj
% abB3Mj0T5u/kkkxIgQ2luDjPeNLJFn3/FEPnBH37NexGKfJSYYf9/kgb2jO1p2drHA6FNHDjTxi6
% GcVmMDbOxYWpzE1CZAuj8XuOP745ZQOMjDAyR3IOKutDZR5DplCFJs8wTS1h5u7yQ/ogaEWv9f99
% B/KuNpmP7TnA6jJgdQJILzFlEaKF5Sa8/xduMRS+LQwCh390GYUpi/4DUEsDBBQAAAAIANKpVV0M
% KLmcTAIAABkEAAAvAAAAb3JpZ2luYWxfY2hlY2tzL3JvdW5kMDUvMDVfbWF4L3ZlcmlmaWNhdGlv
% bi50eHRNU11v2zAMfM+v4B9w0thp0GzDgA3b+rYVxYA9MxLtCJFFTR9pvV8/SnbS+smGeeTd8fiH
% Q0ygKCTTG9IQMBl2aEGdMAz0Abq2bXeb/b67O8An2G7a1TMlNE5qSQ/UkKMwTKCY+t4oQy4JZnf3
% cL9gPlfM6ou1QK+oZNaJ1DmCxxhJr1cgz7YFUJOy1PQ4GjvBhULMUkNhRCctBerRRWEWKwD2AJpc
% JPAnjAILOSahLqQcmeF05HBi1mBcTOgUxduYhUMdJm/mYtIER85Ow9+MOmRvl+puf632KO4UWwCl
% bMRwJt3UFjNyrn/Y3eoLJw0vHPTVHhYpmDhAopjm+vZWn9g3wlr0eFQEY7bJeCtmJkNvtX0WCx8x
% x2jQNbc9Pf36vRHpguU4k1Q8CmMTr2a1BwAjbqUitSiozkkHbcTjmUJUaDEsu1lVr7aCcpo8uYKF
% b0+b5ylSeLcUHALRKG8V0XUyxxLG1JRgLB4nbtILNy9oz9KOxGP7JgtkhmOnuK4pNdrgUEUVrc0R
% bZX3jhVsH0Ai6s4NO4Kl2yQrifX/vttJx5i951CyUKjmVJ2CXmRyMP/qV5w3yTpbsWAx58rqcAdQ
% OKfJ05zx+RYqZrB8LNexUH4fgBKY7IIsOBhVxj8GHJvCrUSrTJDZF4Kv6PF6BqvvFn0JSyRpqeMH
% uF8fusPqh3Em0bKcYovjmIyKwM5OHyGdhI1YJAkrGwnMPRgR4iSrInqd6HW9+snQW5aZbmg8m3JF
% JWQXtJmkURCylX2djkEIiS2BINe7/A9QSwMEFAAAAAgA0qlVXcTY5cQFEQAAsC4AACgAAABvcmln
% aW5hbF9jaGVja3Mvcm91bmQwNS8wNV9tYXgvdmVyaWZ5LnB5rVr/b9u4Ff89f4V6hyGiTTu2nBY9
% Oyyw4pbtgGI93AIMqGYYtETHsmVKkWQ7Ttb/fe+RlETJcltsK4pYIvn45X39vEf9/OZmn2c3y0je
% CHlw0lOxTuTk6qeffvoj2cvQeevcOH+RRXZyRm+nzhLbROiIZx4UThjxR5nkRRTk1PldUTp5wWXI
% s9CJo2XGs9Pw6urvibOKE15E8nGQJpEsnB0vsujZETBvJHInyRy534ksCnjsiOgRNsLjvciHzsM6
% yh14iUJeiPxqFcmoEE4kcZUAKGEtJ82SZAU7S7ZbIVJYhDoyKRygLBI9ZZpF8FSccKFiLa54HA+S
% LBSZw7NHWFgWQ+evfJ/nEZfQFBXrnYBDObB2JNewrwKOfIRmhxew4+W+iBI5RB5drbJk56wy4AY0
% wfBdmmSFc28aHJ4793oMzJEVSRJXY4JkB0znio46qch2+0K/aQLg0bocu4Lp4AQ81l0ZHBun1J1/
% qDfdU0Q7UbbDlKtFAAKDlXWvktrCPqAe6X6hn+lvlIchzfdLKsUjxYlyGiRyQ2WS7TwaAJvDfVDQ
% Kwf+LWMutxT4KoMo5TFdUzzAIkz1b3bKRUbTPFyoJeljxnfk6iqTj0zv1n07euuNvHfjkTciV2qT
% OXv9enUVipVTgHxO7pZKNiZT3edv50w/DR9FAX0j0pd6dA4SFu4LfSbTTBT7TDov/mjee6Yv/hh+
% 9KAwOkQhjjqSqTrAgakTQgue0T0SMguZOik8z3gO+wf9VkPNrPfuAealIaH4NMYnPTdIY+tyMy9n
% vl/sU9jRvftMVqBwz6BDzot6fMHHLDnO8QV+8ZXP9SJspH6xZ6OGcfko3FhIl8OyJFo53BFxLpyR
% WQn/bZkUz4XrbhXdtqbLqKIkig6Y52/mb9gXQv+eSEEqcujcopJjKzBagoXuRdXL/WxOkZjhH4qv
% VR8uF1nL9cflgvXmzArcj3B5xr6cL6ElYYRjBqqF4Jc0RuGCcb3gplwQeUOmijSeM9Be1zxrDXYP
% eroYRjXmK3f1pWrN+mxscyZjTB9puswE39qqkGnBg6pXck+ZMQBoKdUnZay2BmjX3gpUhxlha0W/
% jsClpgL+gGf89febP3C4cghcYhN/zIRAJ5VfE3sXKZxdbwRcSZrkwk3pU2UEWglT/wkOOm+JC8+V
% EmL0t4DG3JVUVEeJwXW6eqgEw4joholZChPRFHm2oZG9D7OSmS04BTHwPTWT5UJIloPJklmyL5g/
% v+pQH7Ofqc1+1Y3U51oTwDSzDYuqhuM6igWYDbr9kgr/DMGjuRsyC4Y8RQbj84bhGeyVcPWAfBhP
% YYPlQH2mgDQYDv36jCKETQNXWszOwZeK0FU9xGbHIgVf7gY0M0esFmXMK6fw1aS9gOi9bdkukm7N
% HtgK3YoTi/luGXInmrqZH6Bo4a8b9cfkT3rUfE7VRNhHW13EHOfAAn/bH0/n/cCfwsPcPqTexwHJ
% 4Q90kg5hHchgTIzyHSGSguaVZ2tJOUDCWiegt89aPNFKB9GFaTa2dLScCcM+zj5VQyuVhxeqNVhT
% gBprmtIIcTRLKWwRiCm+NvqRCjqItnbUVHzrljsG9UNUAGRxy+NKRZadabflqzwqQQjThu885Iov
% FgSo7I3GLR+K49MUx9sgwT3kPoiwNVabB/ZBqCL9ND3rPbJ7d0zOmttCJtNjj52p2Hx+Rgi6fLQt
% B0R6pC0JX2Dn06Lmo+nb8Wd3x1M33++AjsIovo8L2PII9B8aQeuV3+71vC4nAku1AmhUGmLKAYMI
% gEK5+0TzKlrnN+548ERmS+b1yucAn3tPN+6kpxoQGpgWfNUyHrOnnuvBq26EDc1WHrbdDia9p/5T
% r+6Z2CfkdEkDGtLVmK68EkCgRHm8eOS7Hbe3Z49lrROUUYbfje35Xa+38vpeb9lbjfsBHKC/JP1Q
% 7YUDK9I1z0XO/M/UHYwBRQHsc0cUrJkCYJnQt4hubuFHvb+l4wk2DMYePpUWn/Md2qmkOWwpF+ye
% Ay6BeCtP+rFirkKJoNKzDFARyhCFs6iFI8n8vKUrQEgT4qM5+3yhv2FetgIAUpAtKzmAQgEWHQbr
% JAqE64+pR+Hsc0LH70Y9SdoQRp9ThW8kQ/itJ5+A2xhND0o/20TID0Wi4kEkfeAzsnlEJTrP6aGH
% lkjHo17PnSDBhrGxhnjvWmDlhWmQa+1Zy5HQQzeseZlx+EF+KXj7AiIwPZlpPzTptC5pAPwCpzr0
% Dg0vWkJ5gDGqHTSQafst9fDpTp1n4mmck+vXW9JU/qxELGK1ioIIUM0iWEMWJuBHBICjadOzApqy
% lwKLAJMAm7hgFHrPQcEarlq1PjD0H+5RhyVwUqlKsCQORyUfGaZ/1ONM1AVR9b5L8vDpG3NbEMNr
% rPPpxxa6RG/Y/nDHOFcc/3jHlkv19PDpjgWBboTHMNStxlEY1KlTeOWpB4ZbJ53fO097HmYYi0vE
% KQTzlRqf2V2ntSF61uqmrTkWPC8wvatQQfOI0w44xibnUVBBAKSbqhl9V1CDjeZMtajksG7V3Oof
% bzJfQECcw9/xHGNHE1QI0fQeVuyetGJ3AzjYC5q17iBQuDHK8qmHD17TPA3rFekAgZbhf5EMimMy
% OPIYczgB/I+V7hr2v4CiKB8z2/FsK0L9HCfy0XpXI7UteexVTEt/a474dab74q4+TZuabEgBjjO8
% UeUDNTdaWEIpB1XaIFQ9hxeiSgRsaP+GbTS8UDpiCx8R/PEc8QPeM+UHSO9ShUmi+ZkuVq5Iu7EQ
% 3Nixd6zThpRhPhCkNKz3fDyyCsACCEDFOR4BApxYBVqrsSiGPjsaIcBTb2nvncuTWxqqjX1bDKjF
% 1mcqOltbbKHmaSdMExWkbucV7X+ISY1GNHIOHWWMPlyi9MW8eUSjvPk+TVWiY6tHWaF60cqi/M0O
% BBbvAehiYmurs5ESsvOOjRGZPOiYodlyxz7qtgZMr7h2x9yHT4hrPpL+x0/EGtpwbBAYcFGswVUz
% DxTXtJPLr+lEkynUxTpA2GUv8cw6PYq1X8N2YOIde+5d2vEZQWwIwIsAivsIGO6bw3H+fpNSbb89
% WjXeYTy2HJJhFzqe4pQKldUO9GSKZ49xsuQxhGpVbS1KtjXEAjbFXF0HocYFZVgCcyNZuBFjG9I3
% Vk5H7Vhh0F90jv5K0AHzZhwUFjXFlrB2miu+i+KTc4DYD1pW10vEcwrpIKrhNSlLhOu9VaoxBY4O
% 72ZqWM18xfct55Mr51NcKKogLTarMkSaqiG5eTZYXzzuY54BeE6/vx1waWQWyZC9ptPIuNe05V5T
% 8rXKuzXeVnlwamXOG/p0RoQpmw9T+3btaK4Q4ufOZE1ZMWz7gEXZijPuQTl0ELQpTv6omNWcO1Wv
% o0trQmx5piddOaUnVTuNUpdndJnpidRj1Q7EZrIY0kZJj0erDqH5URXPkZ+1TVNVfj0eh1Ehdrlr
% udBYsFpMdXWMXKp+liWElhdWHFZYG8NO9UbdIx2R8+rm5rsTVgBJW9Qb9mVqZt3o0mf1Zq/RFqQS
% OBrKIo9eIMCVZhGLVaEqdU11nqlRZWlDV9pwaL01nJC9YrkK28nXWZGEidZp7AKgAs0DVbFpVexw
% YAsdMGwbpknqNjkkc/a6qfmEU1rFZWAF4BtY4N9M5mY9eMC5+nonMq/nUycq6xUqKOI2G9xSQzS/
% Vvs4XiQpWk4CToTu9nERpTEkL8WJPWR70c5Xclb6m5mqCuNDKJhSdm1FqtJEwAOUQQZQyKsuKN67
% GryYgxFw2M4Pg26LG93Jm9pDCHC1tBViitLViLwRSvHs1c3YoAyRzu+fH26UXYDn0DEWOcizKNdO
% 1yiqzalazMCBZGeV2qubiw7FVGN7rLbg7Vkg1LwcqEuYUECwsM395kbN0A55GiEUSTpQl42Q3gfC
% 3q2GKlV1s5K+zlErX4E7l7hzf0Jv5/UJFzRjZY2k3u93/HyDHxp9p2fIGaKtXSNFMN55p9JAtQ13
% BuYotabGomlh2y3TaqHyAwOESY1/jsev505L3Wj5eHnYKh/Nz13Xz87f4OjAXS4VOnOSlfPSc38b
% fAJDqG6Wlye8olV8h4glssfT8GymMLStSWLNVJyXNNdrACK6cIIOuLzvC8PStHSLKZDoe0DoVOUR
% opJHhEukEx+3PbbsiHIlmunIHVsWud3S9bpjpSagxYpPqBgz0IxxSr10CpEXFtrhIOu6niLSnE3a
% TirlxjVdNWzQXOcBiRV7DgjoysIXAruBB4yhZy3t8h4Wcz0qb26g61KhT10mM7wfdg/Ni8k/eecF
% gGbJb2mFV/NMP1sOYol3sUrG5dXi0oj+h5GoTtjynKW8t3aXZKCYtqwXUWqrCjkdnrvnKhrdXq4N
% 81T35u6S+oVTIcZybawemsssFRXmpFMNfywunHtMOM8Hpjb+wVxEWsq2lxkoUxYFmOH9FeQy2PF8
% iyUhTOLAbR2E85FD4uloZ1iWKJivS5J+pK6bTH3yEl9r0byoiyplg8id/ypJ+F/EDBEj5XbMw0gy
% gDjkmDLMyUEGVNYFLicXi0DgVxVV3Zkz44rI/6E6/cJ+QwNQkofsCqvGClt/7qz0lp8/qGLr+56L
% xWWshNjl31IlG6XYEg3q7C5fgCUv1OHK8PZU1nMB0Vil3OBiqlzmhAyGTzzPu6Xv3k1Gv2iIEdjp
% JxpA4V7/M8lyGA9E0SoCdasAhs5Cp9c0oNd3zvjGu27Q/SEKHkn8DAlTVuMMrZIyEI4HQPrBIq2C
% dXkpN7HrembrYMPUG1VsvDMVeu+W3FkcaBCM6btbPRz+lyVwQ981+vY9GahZf/EIY5r+0sB+ObAq
% rXdt+D2M0Av2eh75oIaqNnmW65tSjL6XUMIG4BRGkD7rGKNDR9OyW4yb0F9svrGWPcww2phtt7kA
% Rp8lR3/a8fnLd6zmBy1HYYPyssS4Ixjo9Tak4ZQoxrczSuM16nACniR8IbOlbUpL25Sa6OZhLbCq
% ViNhZw/vOXiRTJhqWEQdfkgiCJmPpgP4kBT58BJQsPULxXvHmsGkpXAdhagaaizPi0hlXcPH72jw
% myYOf3wPf7051rcIvL+xgkQDWyztahIWb+/Y7bSZLi3PVFDpiwYzg1W2h1DDMdJIQOLrZZKtkySs
% v+8zKvizc4+JSAUf9Qd+mz06jyTLRKz8BmBIWGLqcMuPJPIRMjSYCjLOJBdmNvxuUfWXHyAeeHai
% +gu/MMoEWsKu/kjwH5jOGaMYXnWo/W2n27+97PZvLxTFm069DGdtIFpehbpjcBCusvsIvNRoVNXg
% yqvhODagpPowCmsDVW0OawRGhUDJJqTjVv2WnG2dfPNiMLfS3260q0MlRAcNQCEJzkkj65SJLMPS
% oBIV6tUABYzZZ+2hdB0JokEZtYAsK5j97aOpJHRGuu7Q8L4lueKkEjIfQnErr+FWqldehy+TJHaB
% hOgrcVUiaBB98/KzrRcdRsVb2/bx9tpOQFtwvq1tFzcNmPvipttbmLXSDK801V9FHC1VnTE+gTXl
% wHcpFJo0NwfUSfFaALO8Kh5FRS7i1bCuzqmU0/fnVF+bl78UA7f1hoYxn9vxSGvXbfetlKgsDOhV
% lPLefxMsdZ1cV2w66wI2SPmX/DMoq7llVf2Q9YJBhEMrsG7FiR64unA0X4+Zr1xNUZJM9Wyr61cY
% Nn0XfnWcVyD62gBEq+u/QOqO2WEu4AxhPnVemxYwUGYxHU5WTcrre/1F9dn33E4i49NMCal2hPpL
% 60hXtDmsMCzE87A54fnH3tYH3fgB9l6WqX4ueAacyR2Mg3vDmSvQ+cVC8p1YLBi7XizQuBeL66k2
% 8qv/AFBLAwQUAAAACADSqVVdrQ9o3p0DAABvCwAAMQAAAG9yaWdpbmFsX2NoZWNrcy9yb3VuZDA1
% LzA2X3VsdHJhL3ZlcmlmaWNhdGlvbi50eHSVVk2P2zYQvftXTG92ItuivvyR3QILN0CR5rCocws2
% ACXRKyISpYiUEzfIf+9Qoiw71mrti0dDzjzOIx/N+TevRAw+MKHKA9jBGtgPGinYs5LveEQVz8Xo
% 8aCSXIA7I86MeLA9ZI8HIPg5s0ePD9vtGmiagrciLkSIp2CXl1mVUgmYZcPdPVArtCL9QYJZnQJc
% KFbuaMTWEN5TC+hdBO/gKy0KChg/d+At+POgJzjqD3bmbh/y24tw0hNIG9TwzqFnkUF/ZPjnlXF3
% p2F2P5noLCgaYE77gweYn4e3zNcgD1mYpzyCbZRUJZ5aVqQsY83RAcPTP0ABf2DKzGREWiE0nUYJ
% zUJWQgIxSmSPAtkzacGH+818hVuIYmoEFDOsJeOCIqaoMlZSlZctWJFLrhOnWCwOA48RnatDnZ/l
% Ile54JEeOM7ouoqEtwgyoWUhmJRnC93DlJAg8N6MifumgCl4k7Z8rUsJxCKu5ZN6nbIWN02RWqka
% sTPw7aU7JzaKFnzkvsFAHuOEXMNPsoaxbSEGrjKxwDm62nPRq2fMrNf5Nrp+5+rZoHFtE7w4utpb
% Nl6LtDq6GojYR5dMfp1cP2I7HsgqlEzJd7CnKY8hxB2NYxaD5P+hKH5irs7DX7uu3/Hqwknwa7RB
% CeCGhCmDgvIS+X7Wy3i4pLa+sYGxC2OXxq60dRDNWM/YwNiFsW0c0VzGrgl0TaBnFvI6v070jQ3M
% goEBWnS2AVyagVU98DR6RCrTqKOGKgo57ogAVXKU/JFmxyjoivRO6tejT6MH3OhdXiFmyr9Vp+n+
% GcLyZEcucNra9JUNqeSI0gZO4BNKNa0yIVHM4/G0UUXNze7U1H6jPib6AtTCPwc0Vd0I2OGRvvqC
% V+BaPffU11tecEN59iBd9za8aXvbXqZbS+Km8obI+q+jmdt/eRS/c/Wugvud67BW6pt1/eEOa88x
% t2gY78WzvaxucSvaINflTWiDMg5uQrtQXc/GNX9lVyvllTu7ugLtnOvLxS1ug7uUHTGvVdNe0Djj
% Utb/ys2rhY2DYiKWQKEodSsTty8zfmX0B8/wrZYHqVjWdgEEX3sP2LcKG5PvjD8nqv1fh+9cJXVD
% 9IzNSsaoQAc+25b39MXVLIizBJELJJJywSi2Jg8fP8Lm7/ebf7ag0d//NRt9SljXKhVlnu8sTMJF
% EiYZhLp3xtpiTp9FLhWPsBFiUtEw5RIj4DnNQ91f4DJ0j03MbPQ/UEsDBBQAAAAIANKpVV29jZ8t
% lw4AAKwnAAAqAAAAb3JpZ2luYWxfY2hlY2tzL3JvdW5kMDUvMDZfdWx0cmEvdmVyaWZ5LnB5nVpZ
% c9vIEX7Hr5iVKwkgkRRBUtRh01UqRSrXJtlKLGVfWLRqCIxErEmAwkGRdvzf83XPABjwUHbj2hXA
% QV/T9zTw7qfTIktPp1F8quKVWG7yWRL3naOjo9u1DPKWmCZFHKpQBDMVfM3EU5KKz7QkzoSK83Qj
% usOO43wu4iuDLFYqjZ42neVGfNT3USDzKIk7+Tp3PquXIkpVJv6pgftCgtb9ZvHPTUf8kohY5a9J
% +rUlnuYJsOJnIdMony1UHgUtAeYp4JOFyORiOcfjjvPvOAKXTM5FkMSBXEX5RkSZWKbJCmJHsVjK
% pUo7uVq/F09RHOVK5CrLMxBWIozkc5xkIJ61hFoHapk7+UzhFtSDKJ9vcDuTBSBWSvjd3qCdFdNM
% 5SKYyyyr9iYCleb6l+qQ9pynFFIGyXyuAgLIRLRYJmkubqC8XKX6+VMqm0/vzIJ+DFnTPEnm1eMg
% WcBSzBHyYothEeQadiHzWQlGD5zyfi5zGG1R/s42C1hGZiJzHCdUT9qwj99UmrjYdepdOQL/sDlF
% 0J1AxoGa60diNBLdFhafIGWS6sWSDO3rMYRtg9yV8JuWCAytVOVFGousWLj8m/751R151DeyEyz7
% rNxFFFfo4kT4XgOwqAGlaANvB2JVQwS7ENETSJwA6ARPPozElJ8090DaKuZyaxPAlHDo4KqiheeB
% GAm4pSwhpkRT1iBxsQCEO2UhxLG+61V3fS2YmpeowQ6qrFBlhdrHnyn21sNV2mQy9bvxSTcWbkN+
% LFuC0L9QmGVgTRtPiE17BADQDSteIfOyHYng/iSG7EG2U9D66akYwgLLNIpz92g7u1whAuFujVRy
% 5JXQOo0ctSo/7+gs9Eg5AZCu1xJHnF6OyHEfy/XHR8+hNKBoa12HPEdOA04XOqpc7UP+OQikaqlk
% Puo3Y6Ph8cfA5vho+hAvM5JhdjKC65eyX9/fXwk5n0M0/RiyMr4w+JlAaumyUVrTVkA3/rCD3TvG
% OVsCOMuWKEABG8k6CO8pEoZ7JGHMAHZZwuFzYKw339gfjFM7L/7bninIKC+93++DGmHw+5yOgcOm
% W7308RME2iLU7mT5k3YooR2FQjVVT5RZ7ZRlPIpWOpSk3e/yijQUQMc/YPwIGs1lrJIiGz2khSqj
% PqJ8jIym3KmC3hXi+QkrpvTJdNMSMcwh5y3xVS6X0GCOwqAMVyt9uoxHe2E6npaipOJpPyBPW4Mi
% udq3aOnCLiV9rw46m2rWCaOnpy3iLbGm9MgCQUPxfmaWo9XbvDoyO4C3vTcURkdmc6SUWiEvfku8
% 9Fri+1Tr0G1jAf91ERRsUbLeGdsFxKYj6EZ+COBtFgmgv/Q1iYBIdJlC27dJDGoSwR4SQH8ZGCnI
% XSxJmE7W+cxpQc5dLPU9LcvJAVKl67wQ4HfjHYTY3UtuSOSklmv6oSeb9AYWvcEfpDf9uE3Nt6j5
% B6h1K/QPb2D3vDdUHjStRjGyRatBp6mwQ64QNO0o3yDa97a19sctq4k71/vyjfNJHCpJd3hyvScf
% tcUnZ4ZnPi+doi/FCq4Md41aIQuq9QCeOddIuNfLJedbDs1r+AYLW/1qEa+fidfy+LgHWsSTcNpI
% yKBKbJYer+G5EyRJGmYmP1NS95znVIYRaiBz+YfM02jtjg2HO45+ziSURjT2xHNmCi2pjBnF3BOs
% fu45/zI9QcXcrZgci/K280DKuMOKoeBE4RpYYzhRb+Lcm13LZ/dnbOQa0pMPAKVmOWvxPkKPvAGK
% c96Jm7mSqaDO+mucvMYiVHGyoE42STOT0Shxy5hqM/S3SqKQm3+RzZNX8axiVP8AhFLjDu2nIuZO
% mfI6evXqkCHus0DOuay7fhvbvP7Sw3buO079oFYoZ8mxa6lky2AeW8zYiQDudI4+sdT4ifB6JZCF
% 29Kwbcvc+N8Y8RNpiAEnGm78v+AMOQuK2DUAye9ADr5A3hG1xG+7LU2v7mh6puTMZ+x8LnnIGPYe
% R5OWoOtvkwkY8DJF5aTR+7VFA963QQ2yZyMc74p+pwFS5m8MNCaxidadhtniXp1M2F9Cl2RvEwV9
% PGk2V7odigJxH8yKlM5Py7laKN1hCYVmcAPz/oS4p67KqrpGpzMTy7D3QStXWeL4uO8dIBIaIlum
% qhPMYeahDq4tX7CY9oB8A/VxK0ainaDP1nX10qb6s24Rlo2sivrlQ2qwuAHJS8/5xVCqUyft82Zv
% ytSawBM2jME4Ph7gelGT0CLhoHGf5TJ9rKLQmPsNqRDvIY7alHP4alCxdQ5tBPb1l4G4+9Lz3gNU
% p+i7U7d/Tck44XSzkHGRBWm0zO2cQwOCvs4MX3odZ9ulbEE74OxqE1TiXZA3Di8tP/6FSsY+Bwzo
% GIN0FczkYor+bQYx0mglaZ6AI/zPo5vTyxbPQPQpB9xUSlLGfGZSKYlLvqmP3I9z9ZQbCxVG+UVV
% 46rVEw2gy3zb2Kdc0rBFbaeispBhkkbPM+IC0u0maUK51Ld93J7X2JeGer/hybbUbWGTxwluFoFH
% rimRHDmhU8nYEwsARtfKmyE3zDXTC/xmDB0F78R1DiofR9Q/AOMvdHvxHo4Ae4dRhjPiBv4TKsoD
% SRbZkxuViSjvvMlZc4Mycp0A/Jpzw+aa9Eq19X6RC+EENJYiMy+SOMmTmIZLm/oJZSMwIkMnUE9E
% NZyOL10+vuBviL8/Ske90zGjIdnrhuWjqoxvQ9Q177xFLVq7h+RskEzhfgOHqj/8tIs0P77kHpRv
% u7otHVKi/9xoLfbLcWC98wC0Yd1yNB7aJvlcRaPvD4cDsgWdF32yynLbDNlMpssYFBtBNTqyBliG
% HjCdZJrlaaFbCnQ7rt4ad7focPqUrXE7YN0NvEmpunFzCrDy6iGURXFC2uQaBkJnfoU9OCfv/6A3
% 0Ovih8MnOGpEu9SC+lwkzVDQHcB0va5X5/jtJyVdQ2RUA5x1L8Da76JV+wAJtpIU7SFD0EC8M5/d
% tOyz7MGmICKnoMEkeAxxy22DX+uqWx8KtAZ9rzxnk6qbM7V34nPymolVRoPSYhFnYiZXivN2JhfK
% ttt7aHRZmKkt2lmV0kxaY3Xs478rx90JjzHG3JAE496ETs10oV9+syehcqYBgdDdh9DdQjgBQq+J
% 4GsEv0KoBooyhu5XzQHFiiZtsVrn7rru4XkRTf1H0eVBnsgLdCpu2wZhquxoPDJGfCZprkL3e8mm
% 9r2y4+PDFFvBmmLxLDPeAOGH5/xKOZh5rWp0iwtgVwIpi5ZvKwebq9itgThZ+HqUT09+1QtVKbzB
% gyiEB2U0e/geURG6EpofkiuzVKbaKWBDrJWcRxQG4wl3sguZfa0HuzSHL2dxYcjNhN4CU9ZkbWje
% MZP4M6WoDx9EZAYk2SbLFR2KboGpafw6jijwJjUZ5uFVU+D53JVIT+TOx8iOwEH/RJM5xtAr+lhW
% D+tdzYiOsNaUhzfZkculQtNhuBgFm9cE7kIuXai0pWE9Vuz37hUFGVTo06HsSvSQlvr4NfzhLGWU
% kmt8Nza1fEKrFDsgE62YUu+Hg5YHVlThY4Wpmzmf0pyePJyZ69Bcz831wlwvzUGHE6VJmANzHZrr
% ublemiuMYpD6BrhvgAeG4aD+zchnfNU4Q8N8aAie11efc8+FWbg0nKrSaXY5Es19O1MZPmrTkQ70
% udAKiIYtjVuBNIVVv+EZr/qNEzEh5Ncd5BVVEjMaR7hUnrTH71YoTEgIvlPVi4aUVKXISrVBhrV+
% B5bqaXXiBPPopVDZ4A9ub/B/ba+SuOZai3vWkPiiFrq/IzebryLGoYfElUFhOccc3b1oH+cMZSlI
% r76whEYIz1nIdbRAVYMwqz2RQXmuZLCiSkk3r169Vx2FEzsNGpI65Q21USFnHWT9mlEJu/MGgF8t
% Cv1qMXtvBJpuTHrLom88tQ0jpHST9A8liE6EPJNRT1NxuUG3C/NMkZbYdETKUOHfVuOEn+2gBofg
% 0wgixCatMaqt4xLxGlt4Sgoga1UTXK11XQxj9bxdCg8UOQXPUPAARW8Ebp0VZSVTNF7LovG6UzQc
% q1WxEpnxK+9KuCSBIqfmm565WaGce2VmMR76FrRfQ2vaw13o/mHaw8O0gWZnxTdhV+NhQw5OA3uh
% SZhKgLNtMNXfs6/BHsAt7ueNfXFC3qsFpluLySmqwb+3s3ff5KxtsMZGLvZB2MyGB2D2KU/XjW21
% sPFqoMtdoF4T6PwQkK2yH06aJDm3b0gut574j/ju8rzb9M3WrL+e3FcFjJDslEO/bcfXkz/OGy3x
% wNnPelpmBx2FYa4nK+7xQ+NlLRelXGe0Q23SahwdaJEW8llV4ccvs7g91Wi61j2Mr4B9ItzXFk3F
% HsY6uic0MT4lsbZauL7nNd/yc/RrqXj9hy2+kWAkWNE7L8KmMos4QRklHT1Uxw86HD5ghaZNdA8l
% 7Ks966r2fPPq7xcaVrAyGuflZsbXg0cZLiKcdynRmo9KcChAK4hDjhaYjjhmOhGWtcPsmk9f78Rt
% /XEKdEyfm0SxnrSnSrWXCYhgv+mGx1BRYL58IYOJbEF1JyimqgNCDzh6q91PYrgg0mkieRIJtwDm
% G5p2Viz5exKbOni/FNhuHsHTHOZOfj6PstxtzqDP7PMIjvj6AAqnWV9tv0G333Rokj8c8zHSqLnp
% 8kX+gRZcI1PDpr0/KxZZdXKgj1NoYl3zMq09L/zWcMWqIYI7rMWfTIFnJKJpv8ZN4jyKC8ULC8Wv
% ZjS/NTn6NuLO1wXZmJAmNCT4OOJDnWvW17uS1q+jM/2ZAYe4X7HUYMiQtM22wcJptdz1nk2GvW3s
% 3h/ANioKfUoVxyBGsD8h4/jj38xCmUPMi4o9fSjVkJ59aGqYfOdzCuo7WAWIYnbG9qviaWbZNr9G
% 9J0UYutZpdok4DvutgaTL/3WkXnF0mACSvbPTv3xyfXf/y5uPt3e/O1eEPPbv1rPHmiIUb594ADS
% 0ZRzoJVf1TW/P8tyOUWwAEI8z5Op/UkbUf4vUEsDBBQAAAAIANKpVV1GcxGu3AEAABgDAAAiAAAA
% b3JpZ2luYWxfY2hlY2tzL3JvdW5kMDUvUkVBRE1FLnR4dEVSy3LbMAy8+yv2A2xND83FPra9tUnH
% znR668AUZGHCh0KCtvT3BWVPcpFELbjA7uJl4ohvNVRPKlfGL9KRg327gtdUc6TAUfc4php7PHWb
% k8zIXJiyG2FQFi4gwyiCe9GUhTyoFC6lXe3wg6xyoolzpzxDCoraBfIpcrc58nvlotyjsKrES9nD
% p9sWwdhq2GKUy7jFfH8FmreoXjN1ONaoEhgh9ewhvTUTXVDjlbMMwn23+S7DwNkAI/fDzh7sWitN
% kwk8ICYFwSWTkbw3YOXa2ahDyoGiY5w5ujFQfmujUg9uauhcbASnd+EfWmG3wHMDTKLWuzFegmjp
% 8Jyu7G3A9o9nZfPWY2LO5udV+AbKvE7kOOtDwJFrobPnZrnJLgi1qB2UJKJYPjZh48trPKQWx7mq
% pNht/qwuOGqnPXKNmBYdU8Rqz9JNC4acwl1QC3JBL9n8SXnpcErm7MPIXKzjezUQpyX8Xg738i9f
% LQ6T40v6wJ+r4Q8bJi9O1C+mKe6mnNJgDegSU2nbZXuxkjxhZG+G3337R1l0DGwVbUDbFYnO196i
% uY22qY8+fYef9Mp/MYhv62edP3fqYJHGFp9HK+mlfLiCm7E/SrOlbctum+ne6GIsUj7pN/8BUEsD
% BBQAAAAIANKpVV2m174t3gQAAHoJAAAmAAAAb3JpZ2luYWxfY2hlY2tzL3JvdW5kMDUvU0hBMjU2
% U1VNUy50eHRtVltuXTcM/M8qsoJUJEVJXI1BSmRjIE6D5KZ1dt85LpDc2P0w/LjHI3FePId9T22x
% 6KzOR106R9hSrs3jzJjEnOltWqRM2zVoLCWS3ZdQvH/f6OHTX//88cW/5NcPX069o9HNK48Ojr36
% 3rxO+lrV91ouwd6Cm1IPpz5ODi5peLbNTubNX0He8vmd7pO8Fsk6fdU63s+qrD6tB7GsrbQilyUZ
% DbM5GhHhqvjWlbh+Qv6dXx/rx4cvP95tYAkRnohxYXT2WcYh4qdFnT5xFKC67Wt0Ozt2TcHok3fL
% 9jvk4/bb41+fP9yeb+/OsBktaOdZw8+WVk413XC/PUj7KT2KA2nxdXb0E+EbjzrGUwEyPzzlefz+
% dMdqRV87u+xUU4uz9yjuJDSbt1XVk9siYtCwe41SDc3W8DNTG73eol7EikJ24MphEds9oALEd9IJ
% H3jDreYxJeOhrG503dXajJ5tUtv3qL+4zQBOCx1qYg1GqkWlkVkwQzH4JJvVpUGDxHMraM6Gay5J
% lnqLek9vgxgQIWSNpjDFXjxgG2F/wYMrs0+FSNHtmp3SQyGzhc52eABcHj4+/vnxjlxocbI7+W6H
% ThkdYdE1xkjXWk28tcLVoBUpvJ0+hw63S+HsLzT8hnlR69WrTsu18uC/GTeaVg2GjGSHSysRk2Az
% n+RzsQR1QkKurNi8w/xFbC/pUHT1DlPDQkvP2Qxr0WE7MD+ffSXJ+mlko05Wa7189C1Ku73CvKeV
% DU2wOR0yyGxgSjAYS05dl8cmDjTeCFwKFTU9U08TRSiKsymg+8PzK17H4j7T+IwBY9tes3eJDLXF
% +L0I3rOaIGpAP8f0x4cvMkyGAdob0IvYbkhYK10oErpU8M0nKk7tokQPwEut5TZTbj3wIa/r2lsC
% gti4A/3FLAafk0JmlEVItoR1YdpCQ0mOVUfcBKlicD/QmCXIV08oQDvWfg16T+30YlOESJb0jI2w
% r2Eress8HKVTFaI3KVqHOtozm6BwLjtm2FUI+vDkz3fEbuIN6kHR4Oa4jWuOshyCvkpGW+LMgcYi
% fMJjF1qzg7AQFPr0eAX54tc1tooiSBj5YlfxFx1U3WBLiIwvq96xKnB5pC9BjiLAC53dbP+E/EUq
% rNwH/C/XkJaRtFcsnknIKqFlNY1Iq3JBMp7VDDZY1hoExm+/Q95TikDKUGAVAvNf0k+N4IMS5dUN
% 3eCuy31sa4wTuqHjUL+QdhjlT+R89n178K+Pt49PeXvc17XRscPq2nHAsVrbt3gXyHHgdLpSMjGI
% CfQRNLFgq+FIQdHPiHWBj4fvn25f/b5iJvxPGBiSYdh5EAvt28LZbcCgIMFzIB8Cn859dquO5kqH
% t7GJ34BemmHo3BNztzOPG1YL2kQLC3YehZMbPsLq6jKvzlwHbbEaF6HND0zhd6B3mxFZzwpDD6oj
% 8OMckIl3gxDeWwYo3dgUyK6NONf1zzTF+IIS0/0G9F63WGcgTQ2rwY0Jm0z6xStujSuFd4spGmh3
% mBELEmEZhJUEH05B7N+/92/f8tu3p/x8uyM3Vs+RnDuaX1JNvzoG7xYn0LuOboUFNnaaMlwGf8yX
% nVmE6w5w9T+wF71TZeItqGCb7YyXi73n7HBAp446YbyPyIDt8I6xLMY8rRsh3WBtqwr/Bvuai38B
% UEsDBBQAAAAIANKpVV2PjmE/+gAAAHUBAAAzAAAAb3JpZ2luYWxfY2hlY2tzL3JvdW5kMDUvYXNz
% ZXNzbWVudC92ZXJpZmljYXRpb24udHh0PVBBTgMxDLzzCj9gWQItXDghIQS3SrzATbytRZpEtrN0
% f48jEJco8oxnxnNAoWIghAkwZ1C+wtL906TWRQFLclCw/KPaW8tMCeiK0WAl4YVJ1IEYSXVsbzO8
% /NEPr2+/KkvtctvwRICqzrsM35W1u/AG8Uzxi9IMHyVRI38cTWTI2a2wJzaFWC8tk9EQk3A/hYcp
% 7Kawn8LTM5TqCY5qWIxXghO2YenGXMCnY0tjbaQzfJ796hT2d+ExViHX7WXgtUSaw64JVwEeEdg2
% QDPhYzeuxZkiFG0EfWe1Khwxu/VKeTDLaMVIig8bkXh1K9O3EwwiiY2m0nzzA1BLAwQUAAAACADS
% qVVdyAQ7UMsBAADPAwAALwAAAG9yaWdpbmFsX2NoZWNrcy9yb3VuZDA2LzAxX2xvdy92ZXJpZmlj
% YXRpb24udHh0jZPLbtswEEX3+or5ADUiKcuxAmgR9IUAhVPA2RRFFrQ0tgfhQyWpwP77krLjqnEB
% lytyODwzc3X1eS/bAE4GskaqgkzALTroSG6N9YFafwNfyFBAaHfYvnjoLBgbAH2Qa0V+B9IfdB9s
% zIV2CHaz8TfZp4f7r8vH1dPDR3hcfvsBfVNBusHgm58sB/6cw7Kp6zoH+YpObrHhs3pRjBFFmkJT
% FuIqpjphqgnmthbFGDli6qK63s0/SYLN+JQUT4L9D0vkUJ5wtxPcG+dyLN/TC4K2HarIE+mpYHHl
% sCHnY+W4j3IcQx5ba7qGC8YZfwtquW/44gq4vACXjJf8Pbiq57yY/cHO+BVudcHlouIF/4tas3k9
% hk5YMc++369WdyCVAhxdODVd1ktyEJ9Gm5kA1KEJFAj9HcSP0KMj2yXDSdOiUqN9TwaNCYsFY5mm
% PXapFOlBAzpnHaztYLqYUc5jH5kns1X4wdlo596qg7GapDoniaRBJh2Fncbk7pRYjLfSHc7VOB9Z
% vwbpTqyjOtOWUzmWLW2U5/grWd0P4dh17+wregg7hNWw9r1sEZ52aB3qPAkQ7bPFOGUOcQC59mkf
% HQfnRnDfYp9QUbXfUEsDBBQAAAAIANKpVV0WFvblGwgAALMUAAAoAAAAb3JpZ2luYWxfY2hlY2tz
% L3JvdW5kMDYvMDFfbG93L3ZlcmlmeS5weaVY62/iRhD/zl+xvVMUP5aAyUXVQfakS9qUND0aQaWq
% Qla0wJoY/Kp3TUmi+987u14/Ypt+qSMFszs77/ntDB9/GGQ8Haz8aMCiA0pexHMcXfY+fPjw85Gu
% Bdr4dBvFXPhrjlEUC5SkcexxFHvIj5jnsbXwDwzR1BfPIQMytM5E7Hn8Alj0vDQOkZcCIz+OOPLD
% JE4FutMLiHJ0l9OEVDwX2+s4XGHkAU2c+jTAyOd/pwJLyZuceh0HAXvP8zbOIsHSXu+W6FfD7G2Y
% hw6JccSJOe4heCjnDIiP6suR0BU3juYEHchQrfzz7AcMHc8SQoZjdBwMSAKbNnHUbspElkbo0FN8
% 1zRJ2EbyxlNzXGxOke8BYzJELOAMhX5kaAWASp0T1A+MBO8xFzQVWq9Xcmc4QDRB38jSGRhO/9V0
% 1Y4XpygDV6OURlsGRHt9RD7fLqQS0cZ4zY9YPAsN6T4jwzvT+rbcuYrBrmKQmaZZt+bVsgylSd+p
% Hd/jzNTrsL/vZ2a/osoXJPvMbeq3B+7KzISmNORgqOB4w7XOU+LYIT0ayif9FzBYnffxUXJgURay
% lApmCJ5vvMhlwZe+7YxdEwNfmgWCDHML1pwsZU4YhmMDQ9HPgCGot+Gl3ThrM4YA7X4gfiFatChy
% xz+SxLKm6vVXiI50TSGsCL0tVPClSMVM4I1k9uqDMsrqXMaxcs+jaeLHXPt78qut/G3pjNjgqT3K
% T6wrRuucUT1kU/yIYfk+dzSUo4DT8yLDyXCCvhLDKOvHSCBotmMOBvLlLCnTKq2n1RwoqsQK4400
% 30hhtVzcEuWE+B9jhRVTi2IgzFVeSWaBz4VRcIRY2MvEdlzzDKhAesmI2sTob2FFSpB8JM+vuC/P
% nCXwVa9XB/Ky/R/iia5w7UKa+w4YHsBBK46jEiLmxBnVwOCPNGOVX3jNBXDEsuY18SsO2sNSSQ2Z
% xscFaiQGBxnl3twio0lh2PyaOJ9HOp5ZZGhVktSPhHGeAzFkJ+AdDQawxrYsrSPzBbrzI18wtH5m
% 6z1Hm1hBNYOSXYFTnkHOS5iIuA7P52aZCInUfjnCl/gK/+iOm27X+elgxyRQCQlOOmKjiUYFkWWA
% 700dKGtU0UuBvMq8q1rWyecJy797UuLHcij9zF28BAVc8x2xln2vJNpOJQ2k9xWiwknz/Znb5XlC
% fagxuDwEjaD8NywSvvAZP3cLqC8UFUpTYzl0pR6O+p/k7/oTTLaKFfCg2zBHMtloJo5rBSySCAO2
% 2P7ZZY5AlS/0rgnsRiVtg6GCUQ0ApAGyLcKPiAbJMyVwo6hoIAp5EakESmmAkj7dQEZkkDoTpBEQ
% NNrChc4RU1kH5ZGpxLtoMZfJ6nEjNIsUXyoYVpIsI7SFqQo692MO5G4nk5SFkktrT/mPKBGde1qs
% QuVjgcHHCjo9hZz9HGRLgjU+vkNXr8LX+jMny1wvdSisQ3jbCp2H+S0xlxDfhHbnFLATct9i15D4
% aTg84R0tN1eUkPkyPHt0Oyll1rPUj2W81zRasyBQYdWI8T7zi2ch75IbuGWHZTZ26jqrXyRXQ+dU
% LPlMRnMGXc7CJqfDxmfdIZGP1EZ2EEe80pcrXtUP3nCz2wEPHfZVKA848XBNZuMH+wTZCnrKDRlZ
% j5aUHtLEgN4RLt0i0MbIegBb7FkV8geog/+8zk9EUzali/7MuteZu6oyd/WO0U3Oy7wmSrtTcQ/9
% I9tAg330wyxELE3lfSUPNKL+Ef3kp9BVozhifdlVoCQOXqI4hB6iuFiyBIkYinabMoa4f7x4h5b7
% eh782MiCBbnRd/DJ1HGGnbnjkRqwzBSoTCCBPMvaT26IjMcN9tounXaEsoz2VEZ72hXtRhTKdt0x
% +zcgEbw9sqbwYtf2pvU2qeZ67kfboO3MpverO1jdwNgBV1xiB14+49ElHn3Gl9I59av5T9LZ4J3V
% mixKZhBJ2YT8SYiTzyO6WXQ+VeoCwkNNub0T4NMZlQPRnVPZbEHnqm/fUI0xkmsxnBze+0YrNNaO
% PpBGEKSi7SQQeRKEfUrrbVSzC7l2rlp7WaFtd39Y0/skW0Ky1h7YIa712ZoxokXIAkn6pU2qVzoO
% gAcKsi+dZJBctZFbhnWgsoqmL52QDhrAUHvVbk7KDkd3N1e6l7k63dPk00u996i1Nubysg3AMwJJ
% NOzLUu26Q+hhW5+wZO4soX9w661D6x6ewSg1a/PKe2bv/Kf7r7/Mfl/8cX+Lfp/99hdYf4Vk38sE
% J2+Cf8eg1NsMPugB2p4tI2+gBXwN/NAX5O3+e3d/XCvBGZhDRmDXcBLKSdSyPlWW50AT8mXfcQFq
% gHQclgWRL9twQv2yob/XboUVE5Tkq5Md6TsAcwuHLEY1AG2VqZRhNws1V2NnO9cyPCF4kUYbqRYs
% gV7heNeEQDkdXw/HB6KKDSih9W0UhsrnkEgFdy5QlkbsXLNdyRWnXYtTVQqKRprdJFHgX4P5ycIB
% 6J8sRvIC8FqqfyHD8akq9mDKsrWyJ0v9hryzx3aaZcf/zmiqMZ0n/p6dHCBOpWJ+CoZSBi04eUt0
% LoKF8Ob5KYf8uzMWeUhNWOMMhpWNWhyVizL73m6KNNWj4uPXxWIMTX+g2/f6jFhL6D17wbKzZzKB
% bi9gcoRKNseFwm+w/32M3hRJQ8IsBhXzWTMOk0zkXSTUrZwZxDNDi2zFE7qGyfmZxdCZYjlpQY1t
% GXSdGIF0uFrlu/wJscQsdlyzRP2gJ/WESD49RTRkT0+EnD89hdSPnp7Ox2o27v0LUEsDBBQAAAAI
% ANKpVV3KNA8h6QEAAD8IAAAyAAAAb3JpZ2luYWxfY2hlY2tzL3JvdW5kMDYvMDJfbWVkaXVtL3Zl
% cmlmaWNhdGlvbi50eHTllsGOmzAQhu/7FBZnWgXIQtJVe+sT9Fi1lmsPxIqxkT1sE6323TsmCgnt
% gtpIPaz2guSZ8fzz+0OGpzvGEunazkMI2lkudyD3IfnAqmqT3acxbYRvgMseXV3zzpmjda0WJhat
% hwI4iLYzEANfac3Y0/CkTBdrivS8tLTMT12HNYmCf3TaUzyrxvDeup+W10IbUHFHOWY6r1vgHoTc
% iR/aaDxeBl4Xq7GuFQcOSiNX0HiA2D6fJBuhGkBunFBnG2MOvRa2MfBSFuVO24Z75zD8ZkYYc4rz
% oWwYHX0PY0FA72ivdL3FONFmO5WFgKBo4lpLDVYeY83FuRGWjKOP4YtR6zC6S760JP/OeQWeKS0a
% 6wJqGR4Y7oCFvj71RNZ9+pitvufrGHceWqYDox5M20e3B8Vq52MqADtDfZ8MWs/pn2TLakK2KF4m
% m6/nyF69Gktky2qRbPk/yF6b+Vey+TZ/5WSzfIq2nEFbbObQXr3cS2hJ6Ga297eyrcpFtDXdbbNs
% bW/MK4ebl9MbOStmruSynKObbau/wktSS3jz1QLe8ma81dvBS89vwze4I5s1Dyiwj2eQfD4IiazW
% ViNca5NsRw1akk0HDQ9kQkI6zHQZd+jHhFVMhGPboaPNTDpL59VLpD8FmuD57hdQSwMEFAAAAAgA
% 0qlVXT7yxN2oCwAAbRoAACsAAABvcmlnaW5hbF9jaGVja3Mvcm91bmQwNi8wMl9tZWRpdW0vdmVy
% aWZ5LnB5tVlbb+JIFn7nV9RMq4MLCppLkllBKlJmNd1vq57u1b6wDDK4AAdje1yG2OnO/Pb9zikb
% nHR6NKvVRgrYVeecOvdL8eaHdwebvVuG8TsTH0Va5tskHrd+/PHHXwp/lYtlcogDE4hVEts8O6zy
% MInFamtWOzsVcSLWYRzmRuTG5iK0wo+FH0W9JAtMJtIsSdZ9kGqts2Qv1pnP6FaE+zTJcvG+WhC+
% Fe8dzCqJIvMcKjC/H4zbxUlZniTRaW+V7MG5z/Ctam3v51uV+XGQ7NW9TWJlS9vCf9+aPDOrQ2YB
% HYX7MPdGA/zJVmDWYhnm1ivkpCUetmFkRIEnsdTFRa+YlqGJArHsA2YRmXiTbz3ZG06Lnl62GDn1
% I1N6KWH71oJHkb691HoMdQSkDy99G4p1kolQhLEAbxvjjRQx2g/t71kO1O5Qypb4VX8JO+Hb9CXw
% UKXyqSUykx+yWMzsYe8Nb248L+wW8m0qGbwg8F/lS9RUzh2TOG61DeONF5l1rvzgnthN/TDTX56m
% mTniqxb+n9nBkPxBaHP9pZgMzgcQsgjXeImTnBaIwtMUJjoYzabyCEtO1+Q4+r0fWQNKji5DEWFR
% aH7up0lKFD1Ji3RISTTZGOBwVswlg9OJZX0ieJ046sQnb5sIAFiflfMairiY0MfMrc81vxTz7tBx
% 2/fT1MSB57aJAdAgZKY9qXRN4pEmoMBgXXnIn3L6qEGwvzG5V8qa9UcKjX8ksRHAY/0wwKPUZ6bY
% V+iIR1lR4rOxp8up41EX04qrWvAKW8Mbq51a4c8MxvReGm3i5Kk8OAv3ZsFx7d3R+W/Ee1Aw/moL
% xlZboqaETyGdIrL3SXCAPZEUch8JIN7QWVYYShlReYIjMjHE3qd5KVZRYpFI7GGJQBTJWuRboJAE
% wD+ElnwTr5vMT7dKPIT5FkSh2iSMTdBviVgj9sDcdO/bnSb/jykKKfNYPWidJXZ+H7MWkQTAoCaU
% 3wilkFOb+1muY1PkHlvPgVD0MQmomoh4rEhFinbmyEgXJtB5EiQalJgMbTjXpuXaATx6uejRp3yZ
% NHiLifSGtXuQQzuLxybcbJdJZrV3N3uc0yYeivnt7aO8GMLNrREsC+2yQPCWCydBFQjWTJ5RKeYM
% XENd/OGwWu64B30CvvijEnFafX/F3gMzzE+tmlne/EE7epOm3ymvUBBGVTAs+wrRhPjmA9lWXT08
% ZTJSr+JV54ZwNDjJIo382EuVNSYg5X/UdYKdZvFGu9ze/8RfHgNN97qRTZuJtdMZ09nOSPtO5/KG
% liZ7x4XehzioN1R76RzIz3Py1kaaRoEg43zS8FoPx/etv0fO8ursqjJZuWSVlRuZ/hNJ/Ub8XNVQ
% VyRRRxFRxooDjJlvjQBBI6Jk5Uckf4gtxXFKe9vDxjANe1ivw1Vo4hwlpgSLSR6uxOqQJ+t1X/yd
% AtcRM1Q/cVhmUIiOSZghplD2EJaIIXGHAjOJEHPO9z/OwvkFMS9fqRsoB+LnE/yn3pfwqUdKuAPS
% 9+DhIKRS+udoJVBFTz9/B0nejCeURsKYMxp0ErCm5TRK/IBOHzT0+TTd+AFSp9Wz+TTZ1QXgG6p1
% wnNYzQNqFeuZd692jqF7giJW+W1Hb8QuR9/H2T2ib0fRRwkabcijiYlBRpd1QiW+59WhtFSdMgGP
% HBnTJWJix3mEXGgLa8KHKig5BS12xdr+O1PqyN8vA1/c7yZw6MIjdczud7PBfK7q5+GctHve6DY3
% OOKIr74fBN5Lxp1+IRwCwT3u+LHSb10cvVAx+Lk6JrumOvMshMojs2gYLo++NVsNR4YDWpZp51Qc
% 0Sd9ZJl8Zs1s8r9bjHj2vjTEh0RSqpf6aC7sVCjl08Uz2eT/17p5VNvWPbFlT4vd86LLo03O+oc0
% 8HPjef+9jPLUdj1yzXPGnuC0xzknyPNR3/EIeEOVyj+qT6pyHnVCYueEHKrKq0i5foi0d8ddMvrw
% X7IsyTyvXY8ZlPnF2keyDtroeavu5DR8NIvCXzlPv1JQqJO46qQv2og7DUV34tbruQS1Hgdxhh9T
% 433VuW844imVnpoI/0zgSk7uZledsOvPtSPztW7fadXzu1a+vXLULFthqEZEyAQh5TniarraEq0q
% wFrcja7KFXzsaJk/7n0Ud6SPYYpVdbSz4WTexddkWLWm1WhSNRTlKykNJM4prTrzr1LnToMao3J6
% h16VH4ua8ZcpiA+SUxaRul8kHvdcOr9jI5Cf0ZmVlekUJzSpU8EE3SE+d4rUOOYG/sX26LR9OZcv
% iJ5cZnJG6l4yGn3uHEqgvbg3lO/ejU6THY1zx2Zngx6eSXPbeCe/04hW2JyNnAEKGKDVGCUcNAaM
% 2k4NQ8kuvZUVFo4E5ooiCD0BYv/U/76msyO8PZx+dgnXHXKE5o7dK1gA3YiuGsqjayiPcsotIZE8
% i/LZNZVfNfF0AayzSGRQjgCGkGDuS22bygRPz32oObTXDDE78PtJPdWM9Kh72fHqBhYDGjpX18F/
% pS/yn6YVGg7O6Dd6NCKPaCiJPeuN+IVmGvRR6IrEOlximsEYssySnYmn1ICji8JSnJssRkeWJ0hu
% aNDQd219K3yisE32ycbEJjlYtFmbQ+Rn6L0yY6g96n8nhXxoGsB5aMgm2Ox1lVNCedaqH5fe4IbF
% v9jsn4l6c/Wtiz2b7j7UaejzN+r+oMZOx3Zf96tFw8pNNbKrkrc7JuxzJuCCTbzTTQfPLS+BOa0N
% cPS51nxwc3HF8WdKeDyDKiTt58PotL5S4W1k8TQHMPXSSIgNnWKE4/p8o8dXAzcq2SQDkEceSvFd
% hchCLRohMp9NJlSBh9ymVouI+OFIzrsucb806EClyqGkDCe/zflcCpCm0Xhbul/ZxclDTIHxOhNP
% rXMKcZKRkV4NTprn9bmDpy04UXUUXnTQI0Ge3/Z8KSZ14D6/GXiSDZ87E7mFdaFL5odZd/psAlDG
% OUQRD+WOZfhEfJrsgpCqtU5VrGN1mkUY+pNUTHPh6rxuxqhytucB0l+GUZiXzhGsTnMFVqH3BZWK
% RWA2iDhNduDSIckkC6fPBc8OddfcP/rRwVjMoAxyap5OQHnUAOEznO4WWZLAfA0JFVzcrS4YCNyT
% EhS6kwTw7O4cV4XWjUuzSm/nU0i5hOjUSjdDtWzuIEjnpr1VyRx+Q0BRRwM15Vmpq16HKCCcjG5/
% 3p/vYIPQ38SJxbBop25APA+SKcw8+G10SetJZvaUCF1AHpENXVXDFjg0BQ+9tt+W9WUi+lnKJJVx
% qvoTnGPgJ3U5GHKuifWoE3TrWmfPIJcq6FUpP9eA6Nmqx85vLp/NbBBbv/e4FHdsJ+/GHc92cwn0
% YTNjAez2qgOov1V0qOlGdyhv9dWkAaOvATS+vKWv4fDazRgkRlUh/mWycF0KjNnLJMKUjbY3TVAR
% wLn5/eDDKUOM734uCp0iLN9dSlbktViWhM63YEi75qToOIljs0EKPsKfhb9GbQEuoXRLgKd+THfS
% qBzVFTYN+CXditu0RVVT27TvuLFeuxBlm7JlEpVoD0fXnaLTueyhfBSdovu3wWDYKWSnM1LuRgl/
% 75GoxteS4K579DK6GlRvhHoG9CpIAI0ZqIKCIQf0NG7SvEK17nh89lD2cO6oU/Su5s4PiDm+ZyQm
% uXEjbViS4yOWPHyz2IFHEH26E6SesItzOtfw7lL2KdIcltcsiygxq9sqtlauvBGMK0DuedYbzm/h
% kLXT6llzfhgNRtfDwWjQra7NUy5Ol2N1/ZMajn5So+sxzXpIWYcIXklpDBU0xTvZqE5G7ktF1NAv
% 3C3MgkSJk32IiZBzBgsv1YmN+oG0yD+NLGzu5wer2+6nluqnlEbEIlhTYFD34a6EMoO4XxnFkfzi
% hxZXght3Q41fbPptV19RjOknkX5w2KfovVhGFWLyQtpCHZsmGPG8thOuT5Bt1X5oy/5DBs7+DLXb
% /neMMxBziwW1S4uF1u3FgtLEYtGeuHTR+g9QSwMEFAAAAAgA0qlVXdXcdW7QAQAAGAMAADAAAABv
% cmlnaW5hbF9jaGVja3Mvcm91bmQwNi8wM19oaWdoL3ZlcmlmaWNhdGlvbi50eHRdksFu2zAMhu95
% Cj6AlMSxnTQbegiCFAPWw4DsMGDoQZHpmKgsuZLsLm8/ym6SbgcLhkB+P/+fWmXbDZCtsEM+bDQX
% QNu36FXECmo6oQ+gbAV8Qc4qA8GQRjk4w1VAqYciYYAfu+PxK3g8cxlo19sY4Pe6XArYrB4EFGUp
% IF8XL/CujAGtAveURTbL8035eYK70qQhozoZBN2gfp1UZjsGPOSwy6El624D2teF7Y2heLmX1a73
% EC7tyfHY01iL0FAdrw608x51UgwCzmjZuYFGds5crGtJGQGK+7vo4gjAuiZNac4kGhuPCJ0nm+Ia
% 8f9nMjv8UZpdORclBWdGd6DRR2ISx/yFXVSdY8Q1tkIUpcjWWxgo1fOQ/QllInCF1Wq4WVwVgCO+
% c4ElB5QV1mznMz9wCB6qx3w+5z3Ex6XIFiuxXWRLsU3nckKtP0isEKKyUd6TgVqR6T3+QxW3ygCS
% KUJmgtEi/fJ3xRYFhOiJwWNWkmzXc0ac1xm9bFHZ6wqwmiySxbdemSnBD0UO93S5v4zU7gdlru/h
% +RkOv3b7n7D/dth/P463c3giSxxeRepsXeD1BQh91xls0/pig/DuKUa00FuOzgcmdzxCHeazv1BL
% AwQUAAAACADSqVVduPSW30gKAADrFgAAKQAAAG9yaWdpbmFsX2NoZWNrcy9yb3VuZDA2LzAzX2hp
% Z2gvdmVyaWZ5LnB5lVhtb9s4Ev7uX8Hb4hBKpmzLTrJNHBVwgi7usNuiSYvDAYZR0BIVa6236MW1
% c7j/fs+Qki077h0OSRSRHA6HM8+86d1fhnVZDJdROlTphuW7apWlk94vv/zycSv9ij1ldRqwaxZE
% 8jnNyiryywH7oqnYhPXZ113yZTdlacbCOJNVlD47eRalFcsLlUSlGoBTLyyyhIUF+EVZWrIoybOi
% Yr81E0yW7NHQRJUqqiyL9zR+lkA0qfeJvMiC2q8MaSHTgHYYuic96jWjcpfkO2Jb5r1eoEL2O5di
% KXzrtscKVdVFyso64S4Ls4K9siglbs+KJ1HaEPZdSy/Wh0XpvLazm8Osb2ajkNX9Tf/1zlta5kTQ
% JXUsDwdL4XvmAN8Sidzqlx7tXN7J20asR86X4GfjOdbPiSWuWyq/QyU1ldRUE3vpjG15ROvJvo9D
% WUAvznL6s52+Q/ucwOaBng8wr7n89w1EQtd8+j/v5zzSzc5cCQv7a4jx5dE19qJc2XSZSxsqMjTd
% zVowPU2C5VkZVdFG8a3V7ieZtuKRjyxDssniOlHfyzjyFZ+Je/FA13jH/p5W6rmQlWJ+HOW5CnCM
% X8HaMWaKrCydUhkkT1mQqRLgr1hdYnElk6UqWCWXsRr0mObpJTLnj6I5wJp+00oyI33Jb3fe6HDF
% 0eFWnN/b9sTh9843Cy/WcOJ8s/n+YvfODNPjfmfigSas/hGN89DX253jyYajNRwbDymyrPoOv1XF
% Rsb8RQSi9GWsPHcEwhEp5sV75C/WVJYlfJS9fPBGPfbZexmkUCK0lRW23mHbwXD4MghUmiXkuVnR
% Y3HmjaaryHPv7jj/PFhG1fdYpc/Vilv9wHGt4TDAtX+sImg4zvru3Soi7CYex2gVYX2MIXSVgPud
% 9/kWDBPMqLhUt+CL90asOAOFfXT83VkZ74g1cH1K3QF+nBkdWOJREzcjo69nmX9fUmwsORRWq1IE
% OroU3vyNKk3QeKGgYWgXBxsXxXy0wG8fL2N6gUUcvLsL/Ipm2W2X3c7yaNHrEVgDlSs8EHKVjtg+
% pKqGBtyApPLXpcDRflwHiM5MbVSxYyTeM6D6Q8bxoFeoZ4KzNx8J/bOY0jwMTHJrByfZm/DLTeBz
% EQQKHC0rb0I3D702yLYQCT3vJAr22MY78rpH3ji8JSgs0NDXQ2tv0UceWk4bZjxvwxDp2eYDwDhu
% AkVDKJclb+msO6zL/rJP3DQNgtH58LT1Rk2MkhpPzG0ik2+G4zYQIfQ0UxMyn9bYPNou+p5rghXp
% iLgKUFqMNOd5/i1pkkjyAjrnF2P35lcQ7o0W75hq0IlAE0aIH6W+YqGTnoyZ1pTT2DOiPXBjRJ0v
% s69fp40gxujlhWjkEhd0LPNlqTBJ7yKM63LlfStqQvA79g8ZRwHFuFWWZM8qVVldNhHRxC8mKzDH
% +XtBdF4/wlK5ikISmw4oBxotOrB10TJ/5H+KiXGBPw950xXutbU4QlCbQGaIjNpCMx0iQ4/CIGn4
% HlYwFpC2pFx0T/nCapb8oyVkA8vhlC7urc72/luqPYDC4bXnnUkKrekmk1+vuqY7aMZscozejMtp
% 81ycan0Go+CKGWychaxaKVh8CwXOJshOVRFtBfsRobaqKxKqTkjJkqWq+pEVa6LPUFUNep+8Mh98
% 0hv4fO5qp6UnIsa8eaMfPdqPFwur51eNUxddS1xq/6XZ7EdJC92qq3F38nYi03S+rtDO0l23dK1W
% Pw3UtqJKjxN3QVstBNyKW9pfHC2dNfUrcpL9HjBbc7j7pFU+NHch/EpckKqMBo2XpOthWsdxVO3O
% qxwF6jKDNTuu04UwQb+ACbeanZ8VhcntDGGrivI4onw+6G3Fq1iJivReaoYlv9iidFyx6sLqhR7n
% r/2trpPoP1VKY/u1P7a3VFy90qxZQ2QbXvdWobcCKIfXDt+CZNsH26cGTLDHtWXZq/4WFEcLFBdt
% TY902GgK62qbQ3QeOqtwUNYIgithTsZZHsxdePzV2VpDd3xyDAXPN2wKh69A6nDigMz83jrLVHkQ
% /ZSfpesR7joVnfaW9+FtFTqVfXqQNfAzFYY4yrUcpY85sChR2sdRuOPKyLMVI6TC/n40xpEjZz90
% qVoyklxeaVbrI9utLxq3VoE352tS9bW9ttd918XjGkYS/MbW8xOzMLmiVTpUcHd8ZdYur0bNIq2+
% v8S+hfYvvgViXi2RE8hfo5zP+UisLWwVE5v+jwV2IAA2Qli3bw1q7vKv7S3xun19/bfl5EYrlOOO
% riORfXxcaYYIdwqZ6b2H3Dp98Py3S0HTHvTWIdRw2nmQEs73GbTyk0aCOommlQgM3MVZStLUBqfe
% a0cQM3vWRvSZ3tNMPKDKdfis/2DiuHUgpRViUoCJo6XGmtMVcnwpui2DnnAOAowvFy3AkL866Fyv
% nc0GVRYpW8e79VpsNqIoWmOuQ7EJRRFaVic8gbIuWNmGG1OG6QzZJu9DdEEIagPPCoV5vKP6U8YC
% QROta5VVmgGcIfIjyjQUm6pVoRRr8GLYn5YDJ8GP6lTkiDxWfL3vgoB2sdZQnBCgCY6ovPS7Reqi
% ET0v9XNs7XU03xdzdm7SuQZ3c4BrLTxvfikur5DYbxY9FM8ozDtV8v/cjiJhuq/iP4xazZrPEFRT
% O1GZxRq+zAdVBOWggLlFOgjMJwdTBLFWCLaJiB56r5fKIQ6gSH25+Wmu+EMWz4otkX2jFFhQaYl+
% qS2rD0fCekj9K8oMK5nuP3yg3Ep9FRsRBz3a0GTb4JBtJ+KmzbadDIPlOfV+qIApjuLfjXD18IZe
% RtaCUqopwLxWZejMqI+mngKlO+/q1znwtp9oAo6XUxsBD4RjOpOOBQxXtCTsjc2IM9qXrlkED8SB
% udBbSAi67aHMpZG4MIpr+04HcFQnitRSBN5kMHgvGLQlXITXG0RycUPP0Rkz6fpDK8wBgXCa6gav
% +BuNSFFGIUf68N8iDup7D3T/5NbnwXi977HSspJp5XQKhlBGcV2oE6C0lCU7K+55HH4GTDPUNKmS
% hWi+ObQfF5ay2Pnw+gIholUtKvdIsiJ6zgoq4zXU2wYUdXmRPnvmAxm/HuFiOd0saNu7qlIIOZ16
% EGIRRHPosKpJVdhPBVnQrhvwfO+A2jozBdWiggUT+tSWz/9cdCBXvukHND2KdJnu+PavE728pWXi
% Yd1CjXCxWvVYomTqGbm2w+Ep4VnjPx2HnLxcTDfJoTe1iaWm0i9ARCLqxDtu4zcJQcIg5LTF37T9
% zaZp8Utzlzj5QHdfwvH0su4QDQuUJ7fGDLp91B8y4EXv2JNKZJTqRv2lllTWDgkFTjsyPR2TABp9
% dvJjGSVIB7I8wt2gt6+/6+SDZ6QYnZOije5GmA+ee9Vg3UyIi5KQVpns40RpXlftxwOH1NUmNRU4
% DewaSSknNSJBvuXu0C3tkXkG/G06/eMP9vGfs4dv7OFvHx9+/6opB+w3aAbxvPMtGlbPAYWEkiS1
% Uz+KCHhOWZ3CLYqSulaIFZYDFEf/AVBLAwQUAAAACADSqVVdJjl42X4CAABQBAAAMQAAAG9yaWdp
% bmFsX2NoZWNrcy9yb3VuZDA2LzA0X3hoaWdoL3ZlcmlmaWNhdGlvbi50eHSFU8tu2zAQvPsr9pgU
% dizKethuHCBI0xaoUQRND7nS0lomQpMOSSVxv75DSXXaXHrxg8vZmZ0d/rCtqakgNsEdKcmWxK+y
% CuT44GzdVmqjmWolG2N9UJUf3Qdpaulq0mrjJDDW6ONH4mfGb+k9u6CsodazJ2UCN+z89LNDTxz7
% C/puaautDMo05A9cBSc1WUfaNhTYB38xuru+vycnIwC1CjBwBr+kp5WYinRMEt9FNiZ3uRLFVMwX
% Y8ArqS9XaV5MZ2mxKOPJC7urlZgnWT4v51NRJHmySNPyCk2S/5AUPYlI5x0L0Gl5IimyqciFmJ04
% MpGVIlkk6FuIIk9EeTWbznsKfg1svHrmaOz+ADe3UmlPGxt2MGjHTgWuqdF2E3XEf05JOhvW8Edg
% xwQQtuXP+853bJz1TAcJyzt3lwQVC2qcPOz8mMSizPM4moHNXE98u+lKQCiH+lzMCgqOuVuUe5ba
% 953/gkBMXGSNfKhwXFKWDAGpLG+3qlIoRKpkIVKssWaHYUKLOeP6K93WcdFIEw8ae4bOSaqOFfyQ
% uBONn6EJ2kaPXumLk/tJX1eGn1qpAWcwwWnSLH2YcN0wVTvpmnhe5oQIBkD3SOHgVE+2bbXGaXDq
% 9USGfZ7MncTsAfVIVUwvppKho0qHWaWpdjYOVtkt/iOt60/rD+9uixxKawxqoll0YLdvQ8dAvt0P
% UjYY/LGrSRNv9e3Ur/4eXhbemj/ZP6RCHyHVPE7Ci6XBimOn933C3oAHp/z+H4FLyvv9QOrbYsyq
% EBle0sMqG12v13T7cH3zk26+3t58Q8DQfXSr5cGj4wt0UFB7pjN8Rmx8+Odw8mKWZeQZ5DD8N1BL
% AwQUAAAACADSqVVd+f0GhsQVAAC8PwAAKgAAAG9yaWdpbmFsX2NoZWNrcy9yb3VuZDA2LzA0X3ho
% aWdoL3ZlcmlmeS5weZ07aVPbyLbf+RU9TM1cSbRBMoQkJqZqkgBzaxIuCZma1HP5uWRZxrK1GEkm
% QIb/fs/SWlqSk8lNVbCW092nz760fv7pYJOlB9MgPvDjO7F+yBdJfLizu7t7du96uXA3syDPxDxJ
% xcdkE8/EsfDjPH0Q9tG+uCJokeVuPHPTmQiDaerCuyQOH/Z3ds6DOMh9kfsZzOCmvlinSTIXs8C9
% iZMsD7xMijiBNWIRxHM/9WMPoBPhhqFI0pmfZvs7l4mYh4mbB/GNWCdBnIsgE5vMn8EQGPgg3Czz
% 0zxI4n1EemeeJpGYp4A6PMpEEK2TNBfn6gFAi3OGAczSPEnCEsZLIqCCS+OkWPtptMnLuzSZbbyc
% R0ZuvigGzWHeJA3cUAJet6mCSIEeuATDfKQ7fuMlYejruM38ubsJ81ng5RJubjc+g+ZB5BcwgM18
% 4gH9Aemdnf87+/gfKf5zeSaGwrClsE0pDAd/d2Ay4c5mhivF1ByI1M83Kex6ZI/3pvBHwqWDl86Y
% QGP/xnAruJ5LMD23eJ9tps2peHIcODV5vWgTdqxn4Xo0k4XLyeIhLK8e2moNzw19nCBvTJAztlZO
% YF4SL+u4NlGNkzTqNwEs2jnNUgIiLx+NzL81BzsC/rnAiiFSk+5Q0O9RtgBgoN7RBlEM7k2CKRaI
% sx2eEWTFjUEtjPdqzhhGhX4M93Q7W8P9CNk2tgxHvHol4uIF4NdYPXKzFSIAQnTjF9A8Lf4LECEA
% 2Z8GOYuEYVYv5/h+KOIB0guUBoSpeIdzL6uJ63OqoaiLBi1/eiqWpvhVODoM/lup9cXfgpFbmi0Y
% 2NYKt4XCQteSaAiXOBLu3o+C8Wg5NjV6wuse8IhlKgB+Gu9BqyQsMKiDjUbvR+l45I1pR161Iya4
% iTvxxE9DsRxriCF02gmdInSgVvZnN36G1CnWC4MsN+rGwSgIKEVfqQAQOwa19mdGLIk6UpwprJGu
% 4tUQSFlM+ClVXMl8H8DzZJYArb7aT1KgTuCbL4sg9OlNxYANAOGT/XWyrrEctwXLGUQo3J4fbyI/
% dXPfOGuzuODuqpu5d7DGspSijfDDDGwQPljWHvSc1kCAuBOnQ2GDVszgEkWJlMiP26sUe99H+bgz
% T3hb7nrtx3hfZzZyCUFNkmrm0HSeTfLU99u0dmdLVLQRi8akLuzjHZ1Y36JVi044MUhsgeISUMYn
% y/JJwEiH/p0fIgYgx1Z8wves4faJ+AA/ZN8NeGTW+PxBZ/IH5HDoz/MGl+/I5cG6m3GLr7zS3Vi8
% Evaguhuqy814zwEEGjReg08Gh8r47o1A54wNrbQpVoIpysk3Y+RBMXXPMTv4WqKpjJcsCB/EM/8e
% pdwfrBQX/CYDngjyE0Blm4htCw0b5Rsgh5GBIwQFM+4k4w1ImKY57lpTE6FPkpGWar+FzY7TJPMn
% FN8Yivk3qbteoMUP0NHeuSFer90gxV+7lKC4vtpxTXDQIRf240TMNyEKA20E5fisZqfjOel8NKgJ
% cyTPTPYA1fw4B1u0moVBwCdNNH5vjNDl4yZNNuuMpK8MNQw0azoLLwDid+0Jiyfaq0GXxl/gqriV
% bhXnZUcIMLoYg8yX6nIB1GGy7gF9uqcGWQNJnqa+u2oBIKbGBYggaKeO8XoRKPkx4LVl3dTdJJHq
% JkOsGbd9YPEGuKXe0ItMJwqHl8KdZgbMbaIlR17yeLPlYEDYtAUgyowyw2wTaEJSSTJJ8gY0+jRu
% QX2Ed592urRsJb9jxDp8++81i9YKD9oGBSye0niwdAPA5W8lzKvOQX4Iw5BOPLTH48hwO91okclM
% JLFM7aZc+JValz2OsZSBuXUOHJWg9wFijhbBd3AtOBqGhnEDxPhEKH4ixwUP/hYf6cHH7TJBu73P
% Uxfx/ij+H6bojMYaK5LcZDR537Jogm2g3xRghRLOY3TSxEaKMIJMPZrJiOGnjVxl6lq6qMxh+Xyd
% ArAx37367fpaXLEBRZrnAcVFA/GVRzypkVJ8JS1/qqxXDzIKesn6DxDl+k8C7WCFz66pLPUimNAQ
% Azw3OCUl45DsvSlmFdnajWNMEd0QRsecLgIRxfRBYPoK4ds6wQid5N4NCBaS0VzY+5g1KilasbAW
% noN3DYJC0cBQxbsrU8U/kR+REWdTjKiiu0IDh6kAjqmpY5EzgNSBKtLb09PgV4fEju+WvzrMYVAF
% 5YDNcmqgSHPKIkwBUESlRBtvONCuVsetqrj9V9HD3/Id5zjI2Rb6lbRc4/bxGRtdbQL2EtetmOQa
% AGlh3OI1xti0Pt41VwIduu6wXYhab0h7v9Yll/C57kCm2jzmH5ChNRmgHqlbnFoxWIlbmNyAuvnz
% ORh8Kf5QWP0sMFMbQhAFDIrcUASIe5A/iKt/Da+EAaPElfmvfQL2MKQ6N2wTsr0/9lSg1Aoc8E1N
% PijuiC1jPYrHKmkgi7E2WYltU4PtsZdbWR6kVzSqt+JxvdX2kU0n0hU20Q4QiSEudBDXyeUxjbxw
% k4Ge6cFTCio1VMUO49g+to94PiJm4AVAL9wiBXN6PJVjAaVC5UWNLBGwAGCPpHhWC56IwCPKUZbN
% vDZSIWGgPdKiJaVfHKlFzSQp9jg5JaXw4IL8oJgmSWjAHvexuMPz9msUe2SeI4C3SAIPgvzeoez1
% pSP78nAMueLL52YzLakjtq4JTbTn6KH/daMg0EJ6hhQFirR2zooIfj9Avz9u6qgbPxjVjmvE0RJe
% mF7CZrdEDaBDsPi54bTDIUIFRg8IyBqKR1hJD9hGKKoAAmzbI1UvX+dueuOj2dJU8o9G6rtdpSrj
% hmRt4zajQE2V9oyCYBJkfe27+XDVZZAoweR8hCOWljytiORE0xkEv+Mx/S7H47YWghccau6ttPj1
% f4+PtIP1wuwMBCsSA2AXhQsyAHUfH08K/Ws5+yIsirODsqiJmwEdY0aMVvq8ml6X0/0s/l3Z9vBB
% eAvfW4k/08yHZCiIYO0IXrhcjL0BT5xhtKNigiKSpWruD2Vb27KhWvLVYOcd1mwxkbDlj2U99Uzs
% AjIxlm0gAIdZF98MBVUkfvHjmU1d9PxtKTTOX0X5nZlMJW8xyRsyuNA6LcCrxWx+GoDIFC4P4rw6
% 758g1MROQf0ZxHUsZ09cyIcQTckdKEYIXgPDL5ASFTkWoR4sib2EPNM9C1WJ+0Za1eQMZy81Dwyn
% lwLFD1U8liZfyAoaYIqk0zcl/h4f8a9j0+8L++jZi+cvpHNsP7Nf9vvPTVNqVOLBx2pQ/wVeHEr6
% OXKOnjv2SxvmOnaOn9kODK7KSrfSlZ4kCnJzAtGp+XfA7BYxdiuuKI6k4K85DkytlDcFT6rcIEy+
% +Di834Nt9y2iBG2+b7UnU8BDUWCC8/LDU3Dfpempc5o1DoS4pP9A3A6/3j5BvD386sJP+mr4NYXf
% MPHcEK8LROkZTH46/Eq/T6dfvaddxupWis/sGIigcKEigiScTdCuf7ZoQ0CWQ7jq3Zpm9f7GjSJ3
% Umwd4D4fMIhFFLhlSLVpnvBUUbE5/FRLX3jLkB35cRbc+Si8aJXE3A0gB5om2NuJFyDxSLybMJki
% XfAODKIwWNRLgvECU2yMZWYhxRBPTrwHL+QycikgtQinL2PNU5UuvrvM3KznUMwE6JL7qvWpjOuR
% Mxh3uK2HwA8hub3Glom5h0OVviGWE7ahXqPKrhfdMEKQ3sgIAO9f0J56JlbfmvEGvzGLIr678ie/
% YZk25xZZlqcJgKVDNK5dYaNj27YV7xE4c/gtcv+wTyltMZzD2v6RbVqRe495rVS2br1wM5+NwDMJ
% Gm/0+MeWz+DvoTzCR/jDentJcST2ZToKxt+qIQdaH0VcYvxEUaLxVtpVvK8FBvWGC9bZue/FGI/0
% wJKid3oBtJT6Owe2ojV8ltS92FMWRIN9ZqoyGq6H2yzHXXLEJy/L+PZOUodNlYZhXo7KsVhBM8dY
% YDmstsC0403cS2ScbVbdMzB+TEC4wPvL8QmwEsITAqyhYY+x+Ax4wF+ulMMOkWVOQUglkJfybU2q
% XpdSVSZoZ3d++kDdGlDcC3eTZYEb99AIosjc+RjTCCCqyG43Lnok7BiK/jOONNYQ3tOOWHCoxUoN
% O8mNO1k9H+seY1TCMgyOAFHDPRxtG6Le2yiQ+irVwNYYWx7qc8Kz8k7S6/G4rVOHTZX6E/eJ++2S
% Onhsbpf918z0TfRoYEOPJGZjyjvm/EZSHf4xWBt/or34s2rxbWlBjtvaNN5pFOleF8oFotFHArfH
% aILyWvaf7ajuddRu9JZq0NXppe52nMktjV6ImyI27ROg2sTb5HqgAg9wFgSjtMWjzKizcYAxhe4A
% GlnwYcOOg/xTBUeZ1Bo/i3+v5UUB8bob4q+qXFVf2e82UnVjxd31y5GPZXz463TlM1TJxlMQ1NQr
% ohuLwpG7FvRfI39M4UEq3+pz3dI27g0scPzVcESYdBUCqolSZQq7C8VbhKbMZhaY5MwaTZJz2Aew
% MYK9VOl5rIfVRYgDseIJFeybYV6xOggOlfXa4YFWWMaVNMf84CEW6kQB3rUJv0i6MnCdvxDUIRzY
% YKT8tswEQL7RHCiJCGgMCsqgcu5hsSjhjVtW2DU7+5KBWPlr6jmC17PHbTzQt3e3C5IooCJmEKvQ
% KDTlyn8Yhm40nbliOQBh8SnVlvRjdguCWh9ABjjjeA+u8GIPYqetkkMkxJGDQlJGofTLfSNFLavf
% Gh5BboY5Ph05eT0CqkEWPaZfY0Wh1HjcLGOEHXhD5gXTXHRR9sZd94n5rKJ4YscAcPQLuDrkN3Bn
% wf/2rGFCYk3NBhSeXjHJa3U8Q15YF6ZW40ASbG1wECqvwEwbIabCOL3ZRRWyj60CBK5iAFHBnbAd
% VITe3kNTyyL4K+yp8LK3Fl70QWfIHHeuQ30/XEOJ8LfmBkCY/aP1EaX7I6VmJ8q2d3dG+B2psDo+
% x52R6In6EKF/Ly7grscQQexDPBJCIuxj2ow4P4Gmu1neQ/1XVMBXNO8TRDJgCe8F2STOPHZrpeui
% fQJWVcBVlA0hEy47JnymD4vVBZzYZNQ3gbRvcSK+JOmKD/v5FE2tkwwQgyypzHjSsl9SZq+nysPR
% SZpX9WZf0dzTEHMOUkmYmc3afG8RyF6Y0GNUUSMFfh4YaVENDdHKQfqH8poDl/sW6tDBubpoKhIv
% gpm7XUYIEc3Ag+k9TmCc1+4wE82tXI8pYCNhsgfDixA0W0382E9vHgzwyxI8r549KZUqvKbSprfW
% W9NqShos9321Y0+A9v/SLHdUoUJFtgIT6TUOoV2q3Ajvq6Nql5WLXsfgzylToEgfbiEQL7j6nkI+
% jPaKnKFA8wfCuXWkrf2+tnakrx3V1yap4+on4mSBAYwlgfCWzlRPq8mLnZpDOxuUksqzFTqrO5pO
% UaURrRKLZ53RMYAE/8IgtHj9nmdaZ3URIIxVB1Np0QRu6EjshMoI5WnBUjEL8/Du7TvrhE7RYs4L
% gUhYU0U39hYJ5i9TMAgrPliLR2zBUhaq+YEzmXPjEUkl8dcZK4l67ErP3teOR+nZ/Id63RR7FB/A
% g9VLwQUfdTYqLkphrORar87WxAO0VrYW6UgUtsLRZst2Dfo+dSc5MUWxpSerMae3HzDfXQEtpHOA
% SJpl5ZGL3jVVutqmSOeXLZWoRJsPLl5KjlM7VGTnGyFpK/ehDAZ4xM3+qyL5UVw7vzSbA84vC2oM
% +cTsucrwO7KmNnI83TvaHqhb/M0SyLbaxmMtXWBDUb56p1I5bEWcwN2SUC3uGPFe9S7gu5pIU3Rz
% VR5QXKoGJwY7563VfiTfDIGixQIsOoUgefsQO0F8GLsgHhaacNkSQcPbV0V4ALmy1FYAqzaNwaO1
% kcCHNOkbKvNymXYR3CyqTb/ZskIbGaHj/6aFfyN++zEMt9kyICAe6tryFvdiqjrNm6pPIJK5yBdg
% 0nKwctMNFm5gvVUv/5KgtwA7OqW4Q4KRmwY5fVtQ2Mh7rtvcU7Ev6PWlEVhUnTzqOeY2Ns/8NNBq
% YRB6zZpFgaD7kPR3zk/PlAEq6yL3qHXynhx6GQvQw73iGl/1TWk3WFJgScen8VripIWAy1kDXO0B
% ofGSgBtYlMrB+GhmQ63G2NMExhXFWIzlY81t3Fe4Ki/XZ/OJhw4nkQtcvNcrIvztx7C0r3iJFnbr
% gcq+fPljdREuCmO0yIA/DZuNL48iCCy4VEVkNhyqMfTdOovkQTrgVSOcKs5u1xNqjA2O2y5rmoKs
% c+b1aHBySAxaY7+4o1JBj9atcn/5fmtGphZSvuOkoH1nTjT9DqE7SkzTjhoTpb1dQemJEobW4j+L
% f6vOomT1T2JfVfIlRRMu+JMYlrpD/Cjq4S8gdpprUpGmhQ1EQnS6iTyGWfjQ7/oF2TFRqVYjVaW2
% pYMtheIGzeQo+OVo3GWXizru9pFLHLnlWPX/hC9t3Ah+6ZvDobGEn3++/fbHE6/ppBmSedDmsPMt
% DpfavzfcFmlp8MRj1IvapzF1TNZuikevKqPRIarbjsRwpRBIQk1Smqhd9qkjgDwv72uR3mh0WTs2
% s1THMMbaoYxmDQqye6Qc1Rmr8RU+wCK2TGVP6p+FL81gvESyXJAiyHIbyhW/DbKyyy9UGeJBfAny
% RdH5pMN5qJQCfLIUqJmlirImwlrV8UpWSc2EHmsNsW7z0tjoIR7zEGUkK6omVemnuRGhxPvVIYk3
% /HSJ93GbbMd6pN1OIR1unO10GjNHHQloy7tWBqLj/+wTqzIQDXgqyyk9LMUQjT38iBHCIjfnMpDS
% kuJwRZn0laqEOWJrFJn3p/rxyrrHQH+TqeZ8vRleMpKFpvio8VGd0gELPKsLCosGpKRlrFaKDm6l
% PLebBlnU6I/gqRjN89OnDqNDeSSfy0NHHtrP+7WPW2Kq1kTaiZ/Mb34XkwUzf7staJ2Ia5uKM/oQ
% Se8z4KRWtBdIdcGd78isO9v21GoqI5AwtAY58zFCoZxqRz/Q842PlRqZ1lkj3vRvVCrlnNANZ1JO
% 89DQP/luqfwYrfYtWj2c6focTdNbnLX6Nm3ceW5qqX0UVn0Atmx8ALZsn2nRPgLjkABy3RneHhIS
% dOAGiFAN5ZMnfPJLHr2wrH5r2s84HgJz6Tw7PG69ZR23SqAj+2UXkFOHcg4tmK3ff/GidpSWZL7b
% RFSnUCrtIrXRFJvOW8Ek6iwVKHp1iCoeHjtHR1J8Hh7tlhXBIC7Db5uD1PLjYUPT/uYX3QNlbVKf
% T0UGU2B77XNt3XZcd332faIqx+Vn2fixdiZUGz47KL7EzvZF/dPubA27T1387pvK02Qn93eLE8SN
% Y2GMhf4FF0PqB5OLj7q2NG3ZCrVTl2KTlfmqb/u3d+/E2eff3nwSb34/e/PHtUBWapSZ756F7hq/
% Uf+Cn7LTh9wG/MWNIo3ATnzVmdLL7cH+4fwJtAI2yxV9UJnJJHYjfzJB6dqdTJCzk8nuQLF4579Q
% SwMEFAAAAAgA0qlVXdXAw3aEBQAARw8AADIAAABvcmlnaW5hbF9jaGVja3Mvcm91bmQwNi8wNV9t
% YXgvZXhhY3RfYXJpdGhtZXRpYy5weaVXS2/bRhC+81dMBTQlXcqVXLSHwvKpcZqLXTQ5WVWMJbmS
% ViaXxO5Sloz8+M7sgyL1SA4VBGgfM9+8Z0ej0ej9juUGPrBWa8HkWDEjaslKaI0ohRFcg+Kt5gUs
% VV3BP3UrC7hJ4b00ag+T366j6POaaw5MCbOuuBF5Cg1XFZNcmhQ+KFalwJDp709/jhuxrQ2oGsEl
% IjPFIa8b4dEjs2aGVm9cAicBP2nYciWW++tmfw0fjYa6LKAQbCVrjaK0Ra6YkFAogaQEGQlpkNua
% Ue6hroQxvLiGB/7q9f99CIFaIMBBUDQajSJr7lKhcxBIg6iaWhm49wfANNx7mlbmpq7Ljkbxos25
% u6zRFczUKtxVbRlFTzCDeJLCJIkeaTm1y4/+dJpEUVTwJbCiiFkKWfJHBPhR3LQKBc8nC/gZMvxB
% v86nbjNdeCbJVzEbcoyZpR2zA5Vus3PQTiBBZEnQwoiK60t6XFk9YOwUubKKpIcrq1x3NQnS81pu
% jpU81VHWqro5RwZXV3DTAePaMzSqRsebeKeHPC4eMbo+hZ3uOTj/Doc13fE8Bp6sZPIllkP6+fwJ
% lhjkZ8ojxeSKI8XizFGnqZC5aFhJfhWyEDnXx4hsLhbzjQPZEIincyeif+JR152zmNYckw3TP7Yw
% +EVXzWYwOTA7nUoukQuD3ZMdvBJYJ4uLXN4cLPfnounES8xkR2K3RYP7+RMlQTyF21uQ4ZygZ/Bo
% dySiYvrlIMXTOkz6CCQmkutMmOccK9nEyeFySfczFI5ACE2E1uanAwB9MPmw+bS8O+wcHKI0pBeU
% icap9s4bsEls3wkh+uFEiPXCznh94WvHd0KFmiIhuYGqz+9SX3XBjDSIGgYKr8ddudggqD0G/lIc
% yBKbBAddPZCLgKkN9n20xe50W+kubtIeNYpv8WjSxasfLKyr43itFNsj/Qt8gfgF7u6o8rowrIkN
% by3RF4vdXW7w3BHYUKMNK7OOE+wz046mqTU+T1uCyLD7xhbnnWdLBtEVl6NLVqJrvfv9rnM3uayT
% w0t85Wx7HcTCeg59TyqHyiGcQWbG2IWtpf3UTeBHuBmq44FICi17EnxsSEu7TmFI4GNDMvopYmlD
% hujimdNr3yXIKDz/fyEWWskkfMrXrQL7Tmv3kOPrKeQK3riq/blNfo2HbUmv/v6a3sszCfcN33+n
% 6HoNTKAgbZjEfrxLIRY0VNyj7whhRwghHCn8e1Jdo4calmXNyIJxUyMzbFnZ+tGDFW40sCg4cVhX
% 5ChYLEXODB+d1cmnBhaSe8dwg0fO4owqZh7fxztsbagRLaaLvraqfnXtFBdW7uJcNfX94eVm85cF
% fX0j764LKgB/N1n0k66A236t97LinmEuD0mHfcEfM7mPLbZvcsdhe8FXGAtfJslp+zsr6mIDFmdA
% h5iXRJ9K9kXkWmhG4aIO6u047cBZCKidivwOM41g6IH6BQpXbOR63Ax78GeFlrj6wtqr4u3xO+5n
% CSxcwrQ6YR7b1NmHxEhhT5a9iSZuU3hNMHOejhQlOpswW5c+rVuHSWaflzzmxYrmFZptafadPdSS
% Hz0FliTpDwkS7mbwqzuhVPLzzaGAUyAmEsdlW9E46yX1XE8NOxYUEmpq8pC7wZu+SpCMWL39tOye
% pqA1CA2k+AG8u8HamqIM7K4xy3Tow1RpNA32TqZnwvyNXJIJKbChFBfnGU862aLvn2LonKBvv4bd
% KEVeKuyw3x9pQ3um9vRsjcOhkAZu/AlDN6PYDMbGubgwlblJiGxhNH7P8cc3p2yAkRFG5kjOQWV9
% qMxjyBSq0OQZpqklzNxdfkgfBK3otf6/70De1SbzsT0HWF0GrE4A6SWmLEK0sNyE9//CLYbCt4VB
% 4PCPLqMwZdF/UEsDBBQAAAAIANKpVV3hUi4QLAMAAF8GAAAvAAAAb3JpZ2luYWxfY2hlY2tzL3Jv
% dW5kMDYvMDVfbWF4L3ZlcmlmaWNhdGlvbi50eHR1VMuS2zYQvOsr5mhXhSJBSlzSTg6pcpI9pVyx
% fwALjiR4QYAGQD389WmQWi21TnRQkVLPTE9Pzzy+E7ko39Nv1Gw2TVu2eVsL0Txs6Vdqc1GsHhOg
% ToB6U7dVWz7g11rU20I8AFPm1UqU9HhLIwpRNQ9i0ywSCbH621nl+n6MMuojU4jSZ8btyT2F6EcV
% tbOkHO92Wmm28QNlIq8KesdnqeL79eqr82OgPtEgi++qpu+pGqpSmB6qcgESxYQSRXEPK7d52TQL
% YLWZgWL7JmHZtG+g9QzdFO0baFXm7Wr1uzE0saVOy711IWoVaJAhcLdeUfpst0R/yTEELW3mZWpa
% GoIM0ut46LWiXoZn0jayP0oT5ih0TF/UYfTZ4E7sc+V2qOI89OoHz8gG7XQH1XTUfA0qCyJlxoBM
% +a3AUuGfIsSGrvxvBaIbqB9N1INB0Au0arbiFdoPhs+kLspwNrE/Od/Rkxtt95K4WqKdRdmTDkxe
% 2ucMr8TfR2n+k4m2O2115KB7aeb0HXt9nE30/x3Ek8sC/u5ee3EDQ/GrapAjODuFiaJtU9gBc5mz
% 2gkR9ZMGqUu293I4UPTMd13RVCyyDSkoTi45XAYXDxw4kLQd7Y17mmmHpc+v8dBwpqZ/gCdEY581
% 8OtoMdPotYrczbFLg6gDq+cpA1aRQLbjgW0Sgj59zv+5BPaETnuZZCa5B+0eTwtXLEP4PM9Dw4ae
% 1eg9W/WzsMmByzDL+2kC0wYvPTU4c7GuT+kmonN8VUJhwzLEjLs9z2YBQPr9tUJZQk3jVIozDguT
% naR5Jsijj5jBvfAQzt5fk53UBjHejZFvVntlpdhHjRcZGSnKusR1SCl+sHc3uLU8KX5dGoojjD2X
% TFeEFusKmhkCjgz+Sa40a/be+Wziudz/q+yYFKZvE4OOvHPLQr385rAKLyPCTaFELEtZpoKfv3ya
% 74KScNbqDyMHKESBwaELH0is62a7+nNalLvb46y5fCQY8rZHLyeBIP3YT1cAJ8q7I/Jpi2MF66wj
% n9e417QzDj3bPc4OLhIowJRnQpCHMX4h1ugeV2pML9irm0thfs80TmfvX1BLAwQUAAAACADSqVVd
% Q7SgH2cVAADsOgAAKAAAAG9yaWdpbmFsX2NoZWNrcy9yb3VuZDA2LzA1X21heC92ZXJpZnkucHmd
% O/1v4zayv+ev0LZ4jSTTiWUn2V07XODdJXk9YLE9tAcUWJ2eIcl0LFtf0YftON3//WZISqI+sk1v
% 0cYSyRlyhsP5pH58d1nm2aUXxJcs3mvpc7FJ4tnZDz/88GtSxivtRruPi+xZm1zPtXUQBwXT2NH1
% C20VuI9xkheBnxMtL9x45WYrLQy8zIXhSRw+X5yd/Z9b5nngxuPMLYIkdkOtLIIwKAKWa27GNLco
% ssArC7bSglggXrpZUGwiBogvUsDxJdHiMmJZ4AN0mDyKbi3JNBY8wpLdsGRaAOi0NEuSNfxlUZCz
% C+3vsP4C2uPVWZGVse/iLOyYunEOS8k1j4XJQXPLVVBoXpLsdoylQfxItDgpYDWSWD+J/bDkEBfI
% lbN1lkTaOoOVcixBlCZZoT3IBs3NtQcxBqCzIknCeoyfRMBmzghgWcqyqCzqtyxZlX4hICO32FRA
% a8CbZIEbiq4MyEHcovNX/iZ6iiBiVTvgXi99pJ9lorfL2mqk/pX8Qv5B3NWK5KVHYvZIEFFOgO4t
% iZMsmhJfLo6cafDPC914R9IsiP0gdUOyIUjJcpWK3+w5ZxlJ89WST0keMzcyzs6y+JGK1eo3k+vp
% ZHpjTaYT44wvMqcv385WbK0Vbhg+6zsSU8uYiy5751DxdPHICuibGKOYD85BIJh+IkdjnrGizGLt
% ZE8c80hOtgU/fMwq2AcrHHQw5nz1K8pp0g/GYk85pdCJtEKLsXBzWDyINh8qkT7oe0BLVgbBJwuf
% OGrAGxd6XE9u23oADQGlWwMWuQYB3aJQw449MhjnYEvQauF4gGu6SzLmhvRfWcnkOlMquaq79bJS
% ShsWQzsfKFh2HsQrlrIY16Td/fPyVxzCRcyNscl9zBiL4Ck/F2DBWsMp5xVqIIzSyULSkgLFKg/S
% M75UWPgO5hUrdKltF2UKW/CgHzm5RyTuxB9PnM7kwImGX3x1nUWW0QkH7nAnZLHuwpwGLMvVWJgz
% bSKnwX87GrNjoes7Drdr4LKMcFCDA4Ko2FvnHf1qkC9JzIwaHjp3qCKwFcQqLoK4ZHWva2eZQxCa
% 4h/C3+vezq5l2ciq5mwWKCdx7QBXQOnX/iz4z6dSHuVAMRU8GK1hOGXYTLmtZkQGGXMOGzoUjqsu
% n8WR1X2JL4RhLYTVur7WrUAGtVT+wM5QQdbcA8HYqbufZWL7f/bQIuhPtQBMzSfzaVHSp0vdGkt5
% XNPS1K/GM7MclWZpYE9pmOZMRahbo6m55n1T0zXGlsAPyp3PkOtHAooryql1U80lxPT4aVJJ7/HW
% avgfJmQT0AbeupQYjEqix5uAjMOEQ5yofhxbMP1xZIlFhwnQkpeRrp9MU5+awHLj8kE+dGRV4EVt
% MDGqIxiEAC9AeTeC6w/KmwmknsyTBJBrglWHyQihBf0gM2hHi1xvUy35/qBbxJoaBqUP+oerqw8f
% px/JxxvL+vD+2rh90D8Sa2K8BnYjwG6ubj7OPk7fw9Ab6+Z6Yr1HyCmZtQCtqdmZ8tay6mOrHD/o
% nFjKMXiCKXYAMJmAzhqUDmWW8tZC26yBlAghuXX5ey1lt9SyzKcaTCq62pEIE9D/Y+DZnmXgB/iM
% Q7MsS7IxR6F6KFLrodUq9POfQUCALI2ekw6h5PxW+3hpTQbG3/TG34jx08tZe7g11ZQZ+txEIMsC
% mLMftbsgYz66BnEMv2w1zsE/icENGcP5BtuZbqCPrdeBH6D2JlqxYbEGfTkrtBMrXK2AvchhY6KL
% s3QTLAvXC1ltUGu8ekzY6pHlcrfgAMXA38p2oeHh7TljMX2ZfOMvh00QMt7XbHF8LCgO+uNlK1Qj
% 4UeD40asXFXigG9/vASDI7b1CFX/IF6OuFoSKiN8B7mN64F8eTBUnBfOnqW/Yf6uPjI4Y6xK6I0i
% n4zRMMgLXXXE9MoaE9iZxRdq3d7i1EwxH+BiUttZFGBCc3homYbIzZUD8aVjExi1j8LWEW4cGfdk
% wRFF/EA1Qn/6tP3Jctp2Iolpa+fa6hzWc+GmaO11HfSYafIFIz4Aqqxn1zZhFx4QMRaYOrbmnKIK
% Fa6lgQJRojCPPW+T27HZQEWH4j/niVxPTfocprLxzRnR6vH/ddiGrfFXWF24JeUaHIf9VACFBYdC
% V17jhHZsotBCk1sqeKjH+FOv5ZYCws4EXP+w4wZimiLYg7MdgyClIEYehjTPY3FccS5N2KHzFj/l
% 0bRjB+kUIsybuRMFfwmc6eqA4iIUX8kjOpwjoyVCA1KMewKIDC7Mc/V8IWp018ATcH6qX7aOM8eZ
% /kCx91oeX7PeCqkjOCMXHpWhnpInUu1DUhbU/uqY+pfKqIrD35H7VFmWOBjP7QFP9hwwjAPHmANK
% OxiB3wLRiV69SF/nSJ7bxhT6pRWFaA3MrtALOpJKTuD/HwhEhxABVfylQsEAwxd39Fr4Bamg4A4o
% WJzQB6a/nA1LHyjy29tIZfEKVMO2e0aiVw45tOL0AGR8uuv7iuhKxs98dbBfYru2jqJMW1sPWPh2
% 99AACdUsgodqQx3V6acTTFK7udgpGHuKAFtOAWjuLMQLskQ4kWlyYBm1f3FGyLK7BfheaRI+Nwwc
% 9hewp2GaQCJECR+JmIXcGa8pHQ4v57KlZDSvRISEHBe+Pkj9qO8AjOwMSZiv7jNKBbXRmXPM6C2L
% 5kqlTHFAxcFqsyGKS5lb0F1HN6UbOnDQuxoa0w7p5pXIIeXhCVDLnyShr2whHlagPd2QOnmgI/Ht
% GZFwrgQmDmjeB931cgAxVBhgHqhEmONJPxw6M9jWvBO2SAUZJ/GJZYnMEjWOjTyXGudCpRqlFoYd
% oXIb1aBWglw2WR/FFxIBOE8lnVcS0wnW1AMqZxJKAGi5tYa6sINa6I/WA42uB9okk3CMQlnkbpPM
% bWLsH7UHcMvAVWVPpYsWArNPKK7aV4N+JZoSs4fPmvcMqswvs4z7slvt63KLBg1E8e/LHbzp2/HO
% uBiKnrsSmm1yJdTriTPGNHMcg+IEv1KYZAzJxQvOPMzm8DPTZZNMu6DxAItEKWDockhNRgDJCaYg
% QKBU+pTNM6RF4XIWFIo9AYe52VduSIICrYJJj4OqH31MbicrJ1KcuDzJ0I/iPYZI3vjPfsiWYLs3
% EDVnjVOM6tEHoqZ1VofjNH3pi+woaCjF2PpganfsmYZu5K1cLZjrme2jjYW/egCM/h8xynG4+61j
% H+l0VQdzT30bFNXcGfn2HB5aGRixjj2Cwx/oNDqyjsj2EErLrBInENwCSVrMzV1mLLiZds4G0gxT
% Evf03D7vWZvaXQ4HPL9UqEUls6nvc64p+j6gT7Fv4hBjlKa93gMV8tBjp9PNo4UdvSZ38oCCUrvJ
% PjmQzp4PuA/pxs3Rxf+FgN0gE4P8g+gTAjwlELvNyDVq1Sv44e/XxJphw9ia4pPke+5GKHIxgTgu
% yxl9cMEdB58lfhaPde6kSh8uwIpy84NkLXsJw3bLkJKLux5V0wUcI3GH9Xs4VVn8eOFvksBnum0R
% CP/JFZz1qw/CCe4ZJkEKjx4QElPQAv8MTspkjhj78QaSzEGk42oDJ5GRE4JzOPO9SXk8PIF/beAT
% FSpGWaXYGIPsh/Nap4ULP/BCeSb3BEyVPZls35+J/ETkHiFASHVQrCgCVWb16ZbK4BxXXCevqxyr
% lBGXZORJaBo384ICCx1LzG/DdqPJa+3wHrZVr9iFaYHxFLw00mvp7vvsu/suU6n7lm8PSlc/UfpV
% UzOvoAOTQ+PCwjBc4hxz2dz7aSXkOBF7qYjRX5XeM1LstVSIayx4apo/ePzJtoXpqJKanvh9UwZc
% 8Hfj0Y3u1YlpdIQ2Xs90w3R0spAmBokZY4KH54L++dudcNB9lJPzKuvHMdyLsFBKhGlOMRf3oIsy
% QLVWsWj4b+KYHsrM1uHZ1t6KB0+Ymgzsy9jPtE5pqT5P+mnCz0fqyYf7TzLg4ykuEMjU3Hgk9YxF
% N73JB7SQAUk/G+b9LQy6pZsAglprhA0ERPOJ3BOOgfzcgrm/BJmfGAjD50cwaOF5RPNeHYmuYZiM
% 72VKzrzn47F1EyitquvWL/zVTlzgV8EUOE17ODJNSeIe9UnljnEuUGtRB97Cn4yEP3EIQCVhQWIM
% r5WHJd0JxcTL7NCgIA+ogyprCVPR2oSKoGHjZo+Mvnxb8MQjJtiqQ4nGJeUlO9ARfqGkpakw8QvP
% o7Wv7qm2TFi371g0OBrCEDYCWo8WiIWy9DxiXW68BpClQu7fLOtKeInEqDKtiIJAhbWHX4jnGbD3
% uh5ikMVSCLSmZ71sSbVpITuKLRnzzT8k2aqfJeE+Bydtzrls63vwMxzKX3j9jzeIpY0OrXj6E531
% HRJWkTMX+2frjCMUbwIjUzBeZjbDgAj+Wo5TEVQva9jifseJUjg3SMMtfTLNcCiM4sM1P0xyiDAO
% brhD0Qr2GEW02IYLECT4nHGCsKBgUZ0TVQOtWxBwsWEwMTxMe+FNyNy8GKO7KfZLoqzmq2yhX6Vb
% EhEw1nXB+nitXcprxcCvv+i3dH2WCGypa5eOvXd6WwFbv39HRXak7HSU72jQzqyuK5ehbbLQjEVE
% mO8WmWtXlqZASfR86c//LX1DftmRNhaq1fFZnFYerclnoh/BO4UFiNMru/iz7OpjgGGfpTOkj3GM
% SubnyvIXWXBchkHMXCxKe8RXyts4y1EGi8/QI2q+In93ClLdBRWbCSPPH+t20LwyKKmUd15LS8Yi
% mrNCCalc8MwXvi9toahBwKiGW0WySqjN4zAWGQ7Y+ugC/k/2EMNCF+oygFfrmwIJdnZKthTbLtIk
% 1Y1FnNMXNYHHIqWk+45+/YYTjWmcLxBoJCoJcd7io+8LRj5hPQYIP0oqTzIeH4jZjsbrclL1n0T2
% rirzoo97xMwEaUvxEbW70bLxp3ZdHy9ItI+t4ujh7mQsx6suLU+27/jByVaOPZwDcTwWf+773MFJ
% RuOhp7xYROlW1isw2qoO59vuTqj+CTjA3NhVGKSVbOGQGpfSlHsvOL4ebg0MB7b1j3Tbg6vjhLVr
% aO3IoX2U7shnMk5N7qW9hqMNYdsxe1TvVXQuU9w5gBIxjroY0eUCpXQHBLTO21nPLBeHZJxDHLqq
% pUFLUszCo0+DlY0syCGEPx8ArQGKJNWiMiyCNAz8ygGrS061QCk531ScnFaSoGZvo7h+A0mpfSbX
% znkWoBjwlmKhiHidJ015by6fX7vEEQ+kLr5jfvgJPikpNRWuJHmnwpEOgFeRNAgm2KThHG/LCSLF
% G7FW4TbGdXQ7P0ldcSK/CYPZT46083gngpFvk/TlGQCQnepsdN2D3/xNmY15iv2ylgJlp9uZPZEG
% vZcy0+RukzUWsoX/DmLIA4LzXFsHGfge4BoyYAGIHxcPkfM80v8qnK40xepvUvP0dCjXmpjQEa4t
% bxxVz9g1xetcb9dJf9kNkPsxnpqrv9Xqi9LG08YlSeUuaGFZsHexBKmYFTxPveRbfUKaNXgQIVPl
% YMGECDJ8sGqo1aqTUW5TCN0odasVESRUGBUEzZr5yPpV2jBcFlmt2qpMAaK6sKlrNKqgSlF9Cw4V
% zfE/GqR9BUbAfAYQyqMhJQkwyN43hvt9fVhdF82DyA1FhNMsvpPs5vY3TrB4LHZqz5ZrNwjLjFUe
% e1WR5JbahbH6QdlCkcNdZGXcvY3ALcWhdy/tiGE1QI2oTIUDqDECKevi4IEj2OM5tFe502Ozia2s
% OuJT0Bn1gmVluLPgl2+tlR5EuJL2I5VGBfI7wU/DI/BfSQ+j/VDuF/2V0vh0h5eEZZkX3o1b0f66
% 9i2pYHbJ8+SgPynmkDFkK2WA6JurLjdeDnNfJQlB5JLRBgqyf6cv+oQYc6xofHu9pqxeVjlQwWeQ
% ue6Ju2rKysFPlrH43T44tC52IpUHY1QlcgXzF6LW+qLLNSwCvDcFqqBueXvBtqnX/m4Mbir2DW8b
% svXQsPVQsfVyVy0Iu+Vje4h6rhEN5pNF7hrieYp5eDKTp79D6zBl779Dl1xAL/shKsvtmVGvbEzz
% is+wUWfYjSzjTSge9J1pXpNrY8SfrsDkiKcZmRljvFNXUyZumH1paQ+8fZ+Nsa6YeHmRleIiumJu
% 59rYupxNNJ0rK+Oi7dG1VZGGqoitxhnwmNVF3MZy+0BEAC/glJy/egMZfFeOjC9KhcYKb5xEWAwU
% mblzclOpRHYsWJyjNoRZMWHUutTFjZt9Q6wJmV2RmytHuZNGIzNa8AtmPJY0XrNU0ff9vmhAx+xp
% YEajba+dz3aBpkwU5ogucnqRgcMh8ukNQDx4lxfHqHbxkWIiLm6vWbkzN4chYKEgnl3g09ZRI9sq
% t4R3sfj1PkrBtIkABzgBANDycvWtdUXzilx9MBY5PDVwWP0xpz1pfaJqYSTHt5hY1vW0dbk6+kRn
% V5Uf86DPrslsZpj5J6ufDpRbDHFDVuba5jlNwA/MGf9MQ3sME0+Yz1wVZfUalZD/f3HoCK9XRuQ8
% xt+YnD/h7xMBPxJ+87oGL1zbE4Q5ISaypHvZ/YhlwcsOLojuQRMzlxlIKawuwQSD/PQhAZGlVxMI
% qCu+LDwqaoNK6h6isszM8H4z/pqe6VXXjF/N499TgdxEAIgTa5husl4Okzl7+fZ66r7BKndD7sS6
% 4Qrnx/jDZKKVMTjxEIXy+xp8F9SkPT+wtQMTuUFcnU/UPwVVPzrRjYVym3nRvqm5eM35GbrBOSUf
% Oh60osWn3VQrVuhotwbrJQlegBFVWP6NRQvIo91qXqfKWNUC2sWEdoayWzNbdFqqcm+vo1dJ5Avs
% VkTtfjH0lUCnpsqjVWFOFEV5sHP6S+mV16jz+neXbunN/E9ySjDmqg/I1ToWoAduCGySnMUUcNrz
% G6fXjZeoqI1qzje6ZRgB6gyF8WpWoXtJY7gcI+55vHmOw4HaB+2to1+9rmjzSpQpZn8FW+fEIB/J
% NbkhHxRuts9E2za26s7zP5FMox3Or4K8vtzV5IN2LIshgmjuKTFegrFth4gbANUv4Z8dNG/gDsEG
% KMuWNySumhsSjnSylqon7AwaTjavEv2AHjX1DG93t24JVMnG3l0BOfysuY3dFeum7/ssu5IsG/Js
% VHfu3/H/hmH/s0otRSW+ulAqPTvCix8y7pLfxUlP25gLdOvzl/38/eqbpr3svrW+TFif34duCijh
% 7AAPVvlce2lr7TFX5fOL2boNef4gvoZUV4efeC54Gqf5WFLehAM7W0b8jl6O97r24vPO1IXJLgp2
% vDjveLPaOkzAFsSP4zSBJlmI0ABDBkEzUb70hBcM6iqzJL4iLSWXzkChLJexG7HlktLz5RKN1HJ5
% PhfG6uw/UEsDBBQAAAAIANKpVV29jZ8tlw4AAKwnAAAvAAAAb3JpZ2luYWxfY2hlY2tzL3JvdW5k
% MDYvMDZfdWx0cmEvY29yZV9jaGVja3MucHmdWllz28gRfsevmJUrCSCRFEFS1GHTVSpFKtcm2Uos
% ZV9YtGoIjESsSYDCQZF2/N/zdc8AGPBQduPaFcBBX9P3NPDup9MiS0+nUXyq4pVYbvJZEvedo6Oj
% 27UM8paYJkUcqlAEMxV8zcRTkorPtCTOhIrzdCO6w47jfC7iK4MsViqNnjad5UZ81PdRIPMoiTv5
% Onc+q5ciSlUm/qmB+0KC1v1m8c9NR/ySiFjlr0n6tSWe5gmw4mch0yifLVQeBS0B5ingk4XI5GI5
% x+OO8+84ApdMzkWQxIFcRflGRJlYpskKYkexWMqlSju5Wr8XT1Ec5UrkKsszEFYijORznGQgnrWE
% WgdqmTv5TOEW1IMon29wO5MFIFZK+N3eoJ0V00zlIpjLLKv2JgKV5vqX6pD2nKcUUgbJfK4CAshE
% tFgmaS5uoLxcpfr5UyqbT+/Mgn4MWdM8SebV4yBZwFLMEfJii2ER5Bp2IfNZCUYPnPJ+LnMYbVH+
% zjYLWEZmInMcJ1RP2rCP31SauNh16l05Av+wOUXQnUDGgZrrR2I0Et0WFp8gZZLqxZIM7esxhG2D
% 3JXwm5YIDK1U5UUai6xYuPyb/vnVHXnUN7ITLPus3EUUV+jiRPheA7CoAaVoA28HYlVDBLsQ0RNI
% nADoBE8+jMSUnzT3QNoq5nJrE8CUcOjgqqKF54EYCbilLCGmRFPWIHGxAIQ7ZSHEsb7rVXd9LZia
% l6jBDqqsUGWF2sefKfbWw1XaZDL1u/FJNxZuQ34sW4LQv1CYZWBNG0+ITXsEANANK14h87IdieD+
% JIbsQbZT0PrpqRjCAss0inP3aDu7XCEC4W6NVHLkldA6jRy1Kj/v6Cz0SDkBkK7XEkecXo7IcR/L
% 9cdHz6E0oGhrXYc8R04DThc6qlztQ/45CKRqqWQ+6jdjo+Hxx8Dm+Gj6EC8zkmF2MoLrl7Jf399f
% CTmfQzT9GLIyvjD4mUBq6bJRWtNWQDf+sIPdO8Y5WwI4y5YoQAEbyToI7ykShnskYcwAdlnC4XNg
% rDff2B+MUzsv/tueKcgoL73f74MaYfD7nI6Bw6ZbvfTxEwTaItTuZPmTdiihHYVCNVVPlFntlGU8
% ilY6lKTd7/KKNBRAxz9g/AgazWWskiIbPaSFKqM+onyMjKbcqYLeFeL5CSum9Ml00xIxzCHnLfFV
% LpfQYI7CoAxXK326jEd7YTqelqKk4mk/IE9bgyK52rdo6cIuJX2vDjqbatYJo6enLeItsab0yAJB
% Q/F+Zpaj1du8OjI7gLe9NxRGR2ZzpJRaIS9+S7z0WuL7VOvQbWMB/3URFGxRst4Z2wXEpiPoRn4I
% 4G0WCaC/9DWJgEh0mULbt0kMahLBHhJAfxkYKchdLEmYTtb5zGlBzl0s9T0ty8kBUqXrvBDgd+Md
% hNjdS25I5KSWa/qhJ5v0Bha9wR+kN/24Tc23qPkHqHUr9A9vYPe8N1QeNK1GMbJFq0GnqbBDrhA0
% 7SjfINr3trX2xy2riTvX+/KN80kcKkl3eHK9Jx+1xSdnhmc+L52iL8UKrgx3jVohC6r1AJ4510i4
% 18sl51sOzWv4Bgtb/WoRr5+J1/L4uAdaxJNw2kjIoEpslh6v4bkTJEkaZiY/U1L3nOdUhhFqIHP5
% h8zTaO2ODYc7jn7OJJRGNPbEc2YKLamMGcXcE6x+7jn/Mj1BxdytmByL8rbzQMq4w4qh4EThGlhj
% OFFv4tybXctn92ds5BrSkw8ApWY5a/E+Qo+8AYpz3ombuZKpoM76a5y8xiJUcbKgTjZJM5PRKHHL
% mGoz9LdKopCbf5HNk1fxrGJU/wCEUuMO7aci5k6Z8jp69eqQIe6zQM65rLt+G9u8/tLDdu47Tv2g
% VihnybFrqWTLYB5bzNiJAO50jj6x1PiJ8HolkIXb0rBty9z43xjxE2mIAScabvy/4Aw5C4rYNQDJ
% 70AOvkDeEbXEb7stTa/uaHqm5Mxn7HwuecgY9h5Hk5ag62+TCRjwMkXlpNH7tUUD3rdBDbJnIxzv
% in6nAVLmbww0JrGJ1p2G2eJenUzYX0KXZG8TBX08aTZXuh2KAnEfzIqUzk/LuVoo3WEJhWZwA/P+
% hLinrsqqukanMxPLsPdBK1dZ4vi47x0gEhoiW6aqE8xh5qEOri1fsJj2gHwD9XErRqKdoM/WdfXS
% pvqzbhGWjayK+uVDarC4AclLz/nFUKpTJ+3zZm/K1JrAEzaMwTg+HuB6UZPQIuGgcZ/lMn2sotCY
% +w2pEO8hjtqUc/hqULF1Dm0E9vWXgbj70vPeA1Sn6LtTt39NyTjhdLOQcZEFabTM7ZxDA4K+zgxf
% eh1n26VsQTvg7GoTVOJdkDcOLy0//oVKxj4HDOgYg3QVzORiiv5tBjHSaCVpnoAj/M+jm9PLFs9A
% 9CkH3FRKUsZ8ZlIpiUu+qY/cj3P1lBsLFUb5RVXjqtUTDaDLfNvYp1zSsEVtp6KykGGSRs8z4gLS
% 7SZpQrnUt33cntfYl4Z6v+HJttRtYZPHCW4WgUeuKZEcOaFTydgTCwBG18qbITfMNdML/GYMHQXv
% xHUOKh9H1D8A4y90e/EejgB7h1GGM+IG/hMqygNJFtmTG5WJKO+8yVlzgzJynQD8mnPD5pr0SrX1
% fpEL4QQ0liIzL5I4yZOYhkub+gllIzAiQydQT0Q1nI4vXT6+4G+Ivz9KR73TMaMh2euG5aOqjG9D
% 1DXvvEUtWruH5GyQTOF+A4eqP/y0izQ/vuQelG+7ui0dUqL/3Ggt9stxYL3zALRh3XI0Htom+VxF
% o+8PhwOyBZ0XfbLKctsM2UymyxgUG0E1OrIGWIYeMJ1kmuVpoVsKdDuu3hp3t+hw+pStcTtg3Q28
% Sam6cXMKsPLqIZRFcULa5BoGQmd+hT04J+//oDfQ6+KHwyc4akS71IL6XCTNUNAdwHS9rlfn+O0n
% JV1DZFQDnHUvwNrvolX7AAm2khTtIUPQQLwzn9207LPswaYgIqegwSR4DHHLbYNf66pbHwq0Bn2v
% PGeTqpsztXfic/KaiVVGg9JiEWdiJleK83YmF8q223todFmYqS3aWZXSTFpjdezjvyvH3QmPMcbc
% kATj3oROzXShX36zJ6FypgGB0N2H0N1COAFCr4ngawS/QqgGijKG7lfNAcWKJm2xWufuuu7heRFN
% /UfR5UGeyAt0Km7bBmGq7Gg8MkZ8JmmuQvd7yab2vbLj48MUW8GaYvEsM94A4Yfn/Eo5mHmtanSL
% C2BXAimLlm8rB5ur2K2BOFn4epRPT37VC1UpvMGDKIQHZTR7+B5REboSmh+SK7NUptopYEOslZxH
% FAbjCXeyC5l9rQe7NIcvZ3FhyM2E3gJT1mRtaN4xk/gzpagPH0RkBiTZJssVHYpugalp/DqOKPAm
% NRnm4VVT4PnclUhP5M7HyI7AQf9EkznG0Cv6WFYP613NiI6w1pSHN9mRy6VC02G4GAWb1wTuQi5d
% qLSlYT1W7PfuFQUZVOjToexK9JCW+vg1/OEsZZSSa3w3NrV8QqsUOyATrZhS74eDlgdWVOFjhamb
% OZ/SnJ48nJnr0FzPzfXCXC/NQYcTpUmYA3Mdmuu5uV6aK4xikPoGuG+AB4bhoP7NyGd81ThDw3xo
% CJ7XV59zz4VZuDScqtJpdjkSzX07Uxk+atORDvS50AqIhi2NW4E0hVW/4Rmv+o0TMSHk1x3kFVUS
% MxpHuFSetMfvVihMSAi+U9WLhpRUpchKtUGGtX4HluppdeIE8+ilUNngD25v8H9tr5K45lqLe9aQ
% +KIWur8jN5uvIsahh8SVQWE5xxzdvWgf5wxlKUivvrCERgjPWch1tEBVgzCrPZFBea5ksKJKSTev
% Xr1XHYUTOw0akjrlDbVRIWcdZP2aUQm78waAXy0K/Woxe28Emm5Mesuibzy1DSOkdJP0DyWIToQ8
% k1FPU3G5QbcL80yRlth0RMpQ4d9W44Sf7aAGh+DTCCLEJq0xqq3jEvEaW3hKCiBrVRNcrXVdDGP1
% vF0KDxQ5Bc9Q8ABFbwRunRVlJVM0Xsui8bpTNByrVbESmfEr70q4JIEip+abnrlZoZx7ZWYxHvoW
% tF9Da9rDXej+YdrDw7SBZmfFN2FX42FDDk4De6FJmEqAs20w1d+zr8EewC3u5419cULeqwWmW4vJ
% KarBv7ezd9/krG2wxkYu9kHYzIYHYPYpT9eNbbWw8Wqgy12gXhPo/BCQrbIfTpokObdvSC63nviP
% +O7yvNv0zdasv57cVwWMkOyUQ79tx9eTP84bLfHA2c96WmYHHYVhricr7vFD42UtF6VcZ7RDbdJq
% HB1okRbyWVXhxy+zuD3VaLrWPYyvgH0i3NcWTcUexjq6JzQxPiWxtlq4vuc13/Jz9GupeP2HLb6R
% YCRY0TsvwqYyizhBGSUdPVTHDzocPmCFpk10DyXsqz3rqvZ88+rvFxpWsDIa5+VmxteDRxkuIpx3
% KdGaj0pwKEAriEOOFpiOOGY6EZa1w+yaT1/vxG39cQp0TJ+bRLGetKdKtZcJiGC/6YbHUFFgvnwh
% g4lsQXUnKKaqA0IPOHqr3U9iuCDSaSJ5Egm3AOYbmnZWLPl7Eps6eL8U2G4ewdMc5k5+Po+y3G3O
% oM/s8wiO+PoACqdZX22/QbffdGiSPxzzMdKouenyRf6BFlwjU8OmvT8rFll1cqCPU2hiXfMyrT0v
% /NZwxaohgjusxZ9MgWckomm/xk3iPIoLxQsLxa9mNL81Ofo24s7XBdmYkCY0JPg44kOda9bXu5LW
% r6Mz/ZkBh7hfsdRgyJC0zbbBwmm13PWeTYa9bezeH8A2Kgp9ShXHIEawPyHj+OPfzEKZQ8yLij19
% KNWQnn1oaph853MK6jtYBYhidsb2q+JpZtk2v0b0nRRi61ml2iTgO+62BpMv/daRecXSYAJK9s9O
% /fHJ9d//Lm4+3d787V4Q89u/Ws8eaIhRvn3gANLRlHOglV/VNb8/y3I5RbAAQjzPk6n9SRtR/i9Q
% SwMEFAAAAAgA0qlVXVf5CJ/nBQAA5RUAADEAAABvcmlnaW5hbF9jaGVja3Mvcm91bmQwNi8wNl91
% bHRyYS92ZXJpZmljYXRpb24udHh0tVjbbts4EH33V8y+Ja1si5YsXxoXyKZdLHaLbtB0n4ouQFF0
% TEQSVVJy4hb99x1KlK+qYxX1i2lKM4dnyKHmkLPZDK7zXImwyHkETCoObMHZg4a5kgl8kEUawRB4
% mqsVuAHMZrPO/sMp8CfKclhyJeaC0VzItHO7yhcyBa9HBj3iw90quV0Bwb89t3N7fXc3BRrH4E+I
% h6MWaQ5zqZIiphrQy4WrGVAndJj5Q4Je6QIizbmaU8anEM6oA/SKwSt4oFlGAe37A3gJw37QYMya
% jQd9rwn55YE5aTCkFWp4NaA7lkGzZfj6RLurbTO3ORi2Y8SORE6bjY9EvmteRz4FvUpCGQsGd2xR
% KFy1JIt5wqulA46rv4IMfkOXnvVgJkNo3GULmoRcwQIiTJElJsiSawf+mt30JziFmExVAkUcuSQi
% pYiZFglXNJeqBsukFsaxi2TxMYgI0UW+Kv0TmcpcpoKZB+s3hle2EDWCXlCVpVzrnYFm0CUkCPwX
% F8R7kUEX/MuavslLDcQhnjMk5TiqTG4aY2gqr5Kdw9Ade33iYtLCEGO/QUMR4Qs9hW9kCheugxg4
% yqUDg3XX9DzslW/sW3/Td7E73HTN26DqutZ4tO6a3rjq1UiTddcAEXfdJZfft7YfcQc+6CLUPNev
% YEljEUGIMxpF+DHQ4ismxTf0NX7465b8B35JnATfOzeYAjghYcwho0JhvJ/MMD4OadqhbQPbjmw7
% tu3EtANEs61v28C2I9vWdsTEcuFZQ88a+nYgf9MvHYe2DeyAgQUabdoKcGwfTMoHnzu3GEqXbULD
% LAoFzkgK+J3ElF+HuYko2JD0t/ibp5871zjRc1kgZiy+FNvuwx2E8daMHODU3MyWDakWiFIbXsJH
% TNW4SFKNyXxx0a2yoozN3WRT/R/z49JsgDLxdwEtq5aAGzzSxC94Bq7O5wZ+jfSCFvTco+F67fC6
% 9W77cbhlSrSidyzY4fNodvcfLsV+rP5JcPuxHs+VcmedvrjHc29gd9FxvB+u7SG7UVu0o7GOW6Ed
% TeOgFdpB1jVMXPUpOzlTntmzkxPQdmP9MblRO7jDtCO2WlXygkaJ0Lr8KldVC4VDztNIA4VMGSkT
% 1ZUZ/yX0SSRYq/VK5zypVQDBau8D/1KgMHnk4n6R1991eBT5ohRE9yhWEk5T7MAn1/E//+eZKMhg
% DKlMMZBYpJyiNLl+9w5u/nx78/cdGPS3b3qdjwu+kUqZknLuoBMOsuCaQ2i0M3KLBL1Ppc4FQyHE
% dU7DWGi0gPtYhkZf4DB0iSKm1+mg5ob3/NGK8aAW6EaKl5NtapVCWYPKZIrAijMjUjT4EBVZXGoU
% jZJa0fQBPPjKlYQHrlIeg5KPGgageITIRgyVD/6gsebngf6oijbIgwp5eAbSLaBbkh5WyOMzkG4B
% /TzpYcMa+r+G9M9CtyRt1zA4A+kW0C1J2zWcnIF0C+jnSQcNa/iLNuLPQrckbddwdAbSLaBbkrZr
% SNwzsG6DvaY9BS9w7QE9pjlWLI4Fk6Z6zhXWXaYkHqjJ2FSsubgvqhMyFrWJW5fSbllKkUCOj70a
% y+8Tr5vJR15dJVAlNPrVVRod+JMhLvJ4BTIrT/98izKI6hjPTM2vrhBYXGBpraoti6lIyhstFA48
% qmHHrp0cPOwqgUWXd8sLEDzSd7Hs0hSD01wJnKzq9N/bOi0nIpUYspybWm4kDscImbltQK+Eoop4
% cuDa25slqZJySqqLDo/UBBRd7Y1hL0GQd2rWsRtBJuNVKhNB4z4r1JLmheJ9JpVd3vqmw9A1dx3R
% zHN6vZ5D3EqGJDQtNENxkxshskQrQzwSOovpqpwcPp/jFJs7nJ1LE4onXmSPisvOHGiMDsnZa5NS
% nyz5E9pPgZW6z3UrT5ynHF7PYHCS3085kb6/60Ymp/nteo2807z2IxsNXHKqZ4Ovi95GNn7459/3
% b3B/7snH3w8FImrdzF61OeYGaoFbJl/gDlEcF5Jxp1zWIhWY6hppPCqRoyquxCdm1/9QSwMEFAAA
% AAgA0qlVXUuFns/bEQAAaDMAACoAAABvcmlnaW5hbF9jaGVja3Mvcm91bmQwNi8wNl91bHRyYS92
% ZXJpZnkucHm1O2tv27iW3/0reFMsIMmS40frW2THA3SmLe4iTXemzX7yeg3ZomM2suRKsmN3Mf99
% z4OUSFnJzL4GmFoiDw8Pz/scKq/+dn0oi+uVyq5ldhT7c7XNs0nv6urqwyleV2KVH7JEJmK9levH
% UmzyQnzBITEVMquKsxhOB+LLIbvRS8VRFmpzHuzPg17vNx6aiL74et79dv5nkeUik9VTXjwKwLRJ
% 87hS2YOIC1Vtd7JS64FY54Vc8naARahSVFvZi6uqUKtDBaTw/m94JyULWLFXMJ7m6zhNz0RjWcVZ
% Eqd5JkUh90WeHNZqpVJVAVkfVaYqKSpZVqVIgKC8QqA0XkvcShwyBajLOBXrNC8PhRSAS6wPxTGu
% 8A3Q5ZtygDzqbYp8B/unqVxXKs9KoXb7vKjEr0BjJQue3xSxO/tRD/A0EFNUeZ7W0+t8B/KIaUko
% mPqKYXdxtTVgOMGjexhN1cpM/AavPf1cAO35rn47ZPuzeSnPO+BvXIqy19sXKqu8q9lsJt41nEZR
% 1JLHjQzrjegFLLgKQY6Hcju7Lw7S79EWA/h3iVR5ZVV4SI+3XG5UKpdLf/AEsl5m8U56V66sr3zf
% N5T8e4a0fJZPtbppOjp2BAWa6YMOvtCPNx1Ox8Px1O/1eonciERWEy8OxSoUa/+mJ+C/QoIwM+HF
% 8+FCBMJbzUeLYD0fL0QkVvADz6OFT6Dmv0jEMMbAwzbwsAXcB+CxDTxi4JEGrknLK6LMJas87LxT
% wLp8CgWYQSZ+qD1DmqUqO068e70wqYAJdM7gnimJyxIUS8Sr0ksqHxgnRvYW1WGfSo//rSn3vGjk
% B4Gn+t98H4l3znQ/nz8SSY9ID3D8QXoTX6gNDPxtJr4t4GCLl2EUwThYg7+Ad/QX8I5aeKP/I3pH
% /0/02nzwxfU1ihAXfLMX+DSknKFa/JV8kMVyn8NT6d02uo3eOzvsZBFX7NKMF9/n6bnK92DU/VtR
% iZ9nYhiKQ4kuWJK7B8dXyRPDlwP0cIivUg9bVK///INekaIifwrFMU4P0qimtT9pWf4EK1i94Lkx
% DuAAetw4O9P4jcNbmGSkP4mhO2Mbh6w819jWeQZx5CDrQaJ4DvgXQMROZR5hDXl88AAImgP4jAsG
% wLuIVJWVR2A8TCA4Mbdwag6QVGAZyxF5p9YE21BI0gNeIqjt1z2WZiozpKT0/VBMLF7c27wr52rR
% qAFia05/afYWj5Pq5mU2PeciiAQFmLWLqQe1fhjiIvRT92oOGr9dBMyqORIIr0zy9lKZXf22aYbw
% zYIC74l+kWVEW/qkrM+oHvJIDwEn3TMbqQziJPE0rp7FIzNvaW5by4hWUgGmmvQJzkzUHJEIgwSk
% uItPz032Rx0MWGifXB7S6kJzKoTUCYAXMB3W+RpBXHCt8l8yU4freBzNX5d1TBQxTvNMs4YntBvK
% 8mIXp+qHTJaPsshk6t014Qw8yL9k6/SQAP8KmYBXibMqQtPBnKqQKVjDUUbwgywSWxhNwRuF7I5S
% Sk+QBfFeFrU7guVLUDJta3eD+0FRyI3ng9el+S84OidDtkBBNjc8D/matf5La/WmkBJlrdqe924A
% C0ty4YpUh0y6RrWwgy7aNSIim5rQzC+41/wmZA9jLfQX7XD9yyBBNbAMEuwQ1/8ygAevdkwlW+gv
% EJLKwV0MedvJm68cb2ExQCcopyGq2eCHLPJSnykUI5677ZiaNOr4LcQDI1ppwotzEEszh3OYQOeL
% ZII51DOk1pDAuViIV6723SIGUGREEnn6mMTAzcI3GP8Sttv5RuMZ2Zy+C5AVs4Z1K9+dvuVZww72
% MhM3uUKPNfhOgmos/jT0n8u0MLk9+g3oLaL11GWcpwmKC7xwF+9xcYjITQYAAWUPtujE/ociP+zL
% JlgT9ajPC1uMOng1LLu1+PVihCZ8g3i/l1nifWsciUxL6UJqFulIcWzOjEiNATG5zjoeGoAnhDPG
% 4GjYo80XfmtbHbF5A15l+E4cmPMYhWzM8jujNpcgNZq1cd3spKVBmeZPUGnOamGcABJmF5G7wNrR
% csAA6IYqwnuhSISqDnOEj7jtN4Ryrkdm7yR/jJaJO0WXdQOfMuRj+HrrvFAPkI+kHfhIoQgIos83
% KG4heQSV0hsAndEno0qf6qPBoeodmAKcqnf5wz6vhXZmH41gXkH9nEHlCFFP3A1FvIEKGdRRQEg4
% QFlebcG+KYhAqb6D4JjGhVjFpSqhSETjGRAWqJXZn2nzJnuiGYadEcRF9LjXExxACLQJHZaTvr9w
% 0ref9dLg3vLTWZ6Z/S5jSpP/OWGFFvCOiPP289xgsYJYIe/oePLMeAyI76PrWH4DfnrR58Z7r6nF
% 0JmGEjIObyFjZu9z05Up0rwdyCCCGV78VAesjajPhhyaXqICjW+iVmeaCxYP9D16VhRt/fcn6fUU
% jvPaynZsY3NOQUn1wtpt/D/cbew7rh95UBsATNc84Udt3rqeK3LI3jLwDprvaR6KLWbhUKSNxE8/
% QXWeDVaqWsLSh2rr+f0CqvXr64K3fNqqFFIoFaW5+FmMGobvFBqvl+b9rQLwsZv8JUFQoNwy13sD
% EqyckhccPNFmQMyBcxObMOvdnMGswR94OxlnyzU2xVC/4MfoltYTTHt5HJLYeQO90PbVzs61m/hE
% XpmqVapzcaFANlJ6eYBoYc/m8E9B06VYnXXFm5xjSF7KQW0mwGASsPca2P4WWT8dWoZQruNUjiaY
% TJBIRpMAVzQ6lm7R1CcBi7M5ShC8DvRiwDlpFhS0ACsprQFtyD6nFhmrHTHPFiFuCFED0LTT95pT
% RaxKKd6R9oPCfiiKvPC8q0MG2UOeHqkz8ESd1N0+LlSZZ1eh6JAZquoHdGTeKEStHIIee8ii5hHf
% IKPN5AOl9ku0M0pHiJDXN7A2AnD+31+ENPzGHjaYXJCpAaEZBLHQtBf80YNSd6MeDgUbKAZusNFy
% Iwt8XsXFeY09TIWBn46+5HYw2FoPef09FPUJKM0A1rtHGqhK7kpPawauSSA4qjUkU6w+QNA4FG/s
% PLTK97CDgxic7v5sdTNeia/xTkaJKrijXEOTRu/zUtELxFpZKdMhTg8l9gBkAaUceDM4FaF2cm6r
% 6GwIbafcQKDJsTyHzP4HSLnFP4nvVpeTjtNHZRg6ykCPTcJ/y2epcydcBuga9cemC7gAHG9Gm9BW
% 2OGsDhgEfeHY7ywPXjjHfyH41YGv8DsTWIx5fxrv6r6YKXJJBz6C15ehoD61i/sOD3jX7n3Vqy97
% X7yAaGVe3M2HmAz4zYgpVcAM9JmsyehujvZxs2h3ztLDLlvmRULpLRfNEu8/ZGL4Y5C5K1Eof7KM
% Mxu/4+B383p56NCwaDesbCvu29ymNh8kPFise2BpEwy+zuxKz7LevY6a5ghO+HY94CpfXVkWKk4b
% u3njX4olGTWduuwB1S9hYGT2mPdYNiiKFjcIxfh/jSLfbCCCUooJFugHwNQEFidjUxNEXUUBg/j+
% 4gKfDsV11dhmDO9H2T3vfNOVLoEA6pOd+pf7k3xCjaLjVLTbIU1XRsNWq26gly3HRYUOa4VXMhxJ
% 4DFqN/jbe8/pwS0MjfZ2r7zFphgoNv6gx79skd2FjL77QLpbXUcBcPJo2hyNxcwq/K2NulExjmcL
% R+2dQ5BVv3be3ZggmdodLgrBhie2GVv1n0PAH52YLf48S+efHlR76w5cs5r0lxZyl3TGwcgh2od0
% evjS0kZOuP47hWpbXs+UE3Y5cncX1LHt5DttKdaThqEXJ+zmB9uwCed809EJ2GRGF/7VIpLR4VVm
% MMF0E7vSOmWfj9pRhYvsJsnqRAxW+71VGjr0OwVEvb3JRbs3JT218rmLjfX18m/vvn6trQgEdwP5
% 7lUjRnj7jgN1ogQDzUvYse0VJh4ABZ73CkUv2NC5RzVG7E1mgGMIagbc6+ya324CCwd5S4plJbJQ
% ZkyHPetIeIx6HvbkEsd015uV8RryBKTB3cUc7Ko1DICWNAGvNs6I6i3iNMJYjK8xMQWvr0eTqF1i
% lIMrv5fvKaeVtMxOjinJnk8XUA9+wJpQVxxYWxizyc7eM81FDEwOYivUOxO+wz3sMckTilnh3YOB
% tGSJn6Ngpk2dK5SGPNHtRhJS62adxmpHH05AKiQTPGDvlXjPCTAQrWQZ4WcRUEyrnU5DQ4EfRAhM
% ns7k9YoMko5VnMYZyAw9l6pwtyoHF1WcB713E0pmZw1fRk4lRuXQSD8udDn+bsLOqPUVhL7+MYPu
% tasu6YftDxO4HP3RpCQ2CnPJdWim4+hHf3Rhq43EuAVFQEjQoX/s/8DMeoXlZkce9tbU49wH4Zs5
% J3MaYdunnThNdaLzCC4xpkWtVRFmkRfLHrGkbG7y9Z4L3RyVG3WSrf46bICQvE/DThI8AA7dy5Ct
% hnW9IAIb343PmKX0tw0Tzc42DM9+wj0iFAmN0R60NV0l8iqmQG426O+OspOBby44YS4rm6yuxoCt
% uz4pB+gmvAQ/TNUVsqJgvseKG4pP3V81MEOwWTPTOtehZitUsBuHgZamYV5HenTTe0HV8IudDiiC
% VAW5IcZ0iH6E9MkOPB35GVbS0yEURzia3xXVkE5C5Hd/xtB5FW8LnpXz1N82CoLDC5NBG3EiU2mj
% sJZsRzKvKapcTaBtmjtmUpX/PrVWBjMnNf3WB27x06Ldd8DEijYii4LoBVr/bLauxYbwobEoPC8R
% Hmrj6shUWYP6rbs+rVezxg8Gte62YsDboc7n0N2qeJXKyHjlyHhlduUmxWIv/xHYj166oo8QrZuJ
% fMVXGdjN2WFKh/fOGDsAEPuSqhDlDp938UnBA0orL8pB7719ezGf65Ybu3v4d9HOQ6jAd6AuQXSb
% zLTaEAbSKPpkEjJKalqMqHE1sTpbpgPW0TyZQLDBlS2bfKHdMr1c0Wq2vK/v7qlP/kzXRS+4ZBCk
% iPOROb95hBrPD947eTW7I39w33sM8RIZEpsYo9iK2F6edyuY9q4ewULWMAMTIOgiPnNn6rSPsUvm
% TYJH8CQB/Y7hdwq/EEQiMQ4eA2+sZ8c0619P/YZujQIxRvD6hbgUp954gtwPvMe+NUitAB9y7wuF
% d9biScfda6mo4O7m49EOpxPjCDVlaAtA1KA8gCyANQjNq2vjgcPCIARZ/cvv7WQqNaoMJT5/hob9
% GJWtVSLRisgYwJ29m7RSVCxwKDW6qvNItJfJyFgm8qwxPmq+qp11IuCCOdTvIK44DqogQJA+SBGf
% ocDbRWMU05qPzjeLM8QT/A6lWKI2G+930gsCowUgSWsGpMFlXVnFXBaDmHE9bKd3offoNauD3hPn
% grE1PW07sQiRrAMP8OCT2boh+zmn2QfUZhf6nTSrXvvuhXOtf3z0iA+hlYTglsgwKF2DGteFknI6
% 8CIga2STGzmqN7WcAJBj7pkvtI8MCB769HEqb9J860ffTVmkaVCnH6hP3Wwy4wMGFWtGf1W/2Ux+
% hR+ugx+AGjbebFQmdZuAr6RRp7nAgWpF144hf6gMWxxVLPAuKYGIS7XPoJP/nqbDIaOmolvWURss
% spFoITrGqC+7oDrJsN6N+FvQLN9BZn1df9t+3ZxDoJHiFYPkNCGZTcLBYBCOhmhyGvc93rnF2aFc
% F2pPl/pHSd/rw07lPo3PVAtBpIWKSmJPnHFyKzAuVgrsHaoeUyh9xDCEhRB+e99QssaPAWJs5eie
% gIjF0zZP8RoEPwegew94QK9Q3+vpT1g5tG3Vhhvy0WjI0dF8gs+lAocInIJ/hqa8eIWHtj78h9oP
% 9B58ePUfY7Do0eQtWGmFT9d/H7Nw8dOEEkhL6EuAE/alR+B0/j4O8A3wEynGEuMJhhu8S6zJYbty
% 6TP2oz/WybHjPgUvMwFXMgnqLcGrIC1vaYb2aebIASG0dlkJJDDHWOf+o/ElshYmW3MbMn4WQ0pk
% LGymU2bCeewQOOk7m0QsGLPG7s2UmCdVjWxBB47yBKpzA8EY/wKAlobiijQJEynxMw7XuP36jxre
% ffokvvzrv31+L6bi1398+PX2q8AtPrwftP6QQMP/oj+WTlT8kOV4nVZCZbOHegeMB7aEs2/J+OPM
% /N1I2PrDkacCa/bM/KkIaPd/AVBLAwQUAAAACADSqVVd+Q+EYOsBAAA7AwAAIgAAAG9yaWdpbmFs
% X2NoZWNrcy9yb3VuZDA2L1JFQURNRS50eHRFU8mS2kAMvfMV+gBw5ZDMAY5JbskkBVOp3KZEW8Yq
% evFIasB/H7UhMxdvei29Rf41UYavNdWIxheCn2gjJX8OCi+lSsZE2bawLzX38NStDnwDISWUMIKX
% hEkBvYYZqGcrwhgBVUm1He3gOzpywomkM7oBK6j5AYwlU7fa01slNepByYzzSbcQy3UNybvVtIaR
% T+MabvdbwtsaajTBDvY1GyeCVHqKwL0PY5uh5gsJD0x9t/rGw0DiBW8eh41fKLRRViYXuINcDBBC
% cRklRi8svTZOdSiSMAeCI+UwJpRzo4o9UFODR3UKwe7C37WCnwK6tYJLtHo3JnJi0w6ey4WiE2zf
% 6Gbk3kaYiMT9vDBdAYUWRoHEHgL2VBWPkZrlLlshVTV/MeQM6vk4w9ZPlnjQPI5jNS65W/1ZXAjY
% 3rYgNcM021gyLPbM3TTDICXdBbUgZ+hZ3J8icweH4s4+jBT1iW/Vi3CY0+95d4d/+uxxuJyo5b3+
% XL3+sGGKHNji7JryZpJSBh+Ap1y0bZfvxdLkC4wU3fC7b68obGMiRzSCi1UL7Ok/LBSh1zBSOOuC
% 8KGcQ6y9p3cdfZkfVPoOfuAL/YWBY9tQx32s3c5Tzy3hCA3Ss74bB1cn8ICKL4T/D7684Ywn78L6
% 0X71D1BLAwQUAAAACADSqVVdcHtKRBUFAADUCQAAJgAAAG9yaWdpbmFsX2NoZWNrcy9yb3VuZDA2
% L1NIQTI1NlNVTVMudHh0bVZbchw5Dvz3KXwCLwiCBHEaBUgAa8WMxg6pPau5/Wbpw2rL86OIVnVl
% E/kCVfkcPd4a7a1M1oZlY5vV2jlLLacEVQrv6n3QdqZqIjk9Dd/6/Jnaw5/f/vef7/49n798j/o0
% hqycIbtbyp45W7quVWd2tRNHJVarOVIa2dYcvjV4R5siOi0/QN7y9dOZtrNvElLts1lGNV3eqs3g
% Ji5OnZbuxW7WzOo4S2IwTduA/Qn5dz4/1j9fvv/zaRDHrE3ahWZ3lZ1DiKLHqY1hS5y39aObF5eN
% yjcyxjJh6tN/hXw8fnv89teX2+vt0/Be3lefnXhwUaoNH6dv7bv5Dj4kmzJiBMZqMaekZ9OomcOI
% gcwPTxmPP57uWFWdxXuNjnNRrzWpQDRG8FXedg1usaGJ+1rMaxro4RpuJbIo2++oF7F1eobOoXSG
% cT+LcJbhrKRlPhcb01igcU2haCNih5xZSQ6+KO5R37mFLK2lgF44YCV+Q41WMYw0eknfETxaGo8z
% Rp9BcVpb8J81EhX7DfWeXgO5Mj2AP3xFHaOZFCC9e/bWWVdbZGdXJnCjkqEo/IAzGa9LuP7w9fG/
% X+/IFZJxclwmpLV2WNbqlyJ+bNMEYEarNaZN+Ha0CWZFXVQp8GR9xLyo1YgEdXzw9dg8dW9MCbvv
% ZTrbgu3KatZcNGrgQ+1QzNM3ZjPv75jvxPLpiqDiGKevlUybews2Khk9eicodyw6hcBerk0mwzZL
% jtY6bX/AvKe1I6cVtkHWdKAC56ANrslD1syzcW4b2qxrSCfNiXzx5urwL01Ay8PrB14bn+XKQGUU
% i1CiW2YiDqiPAL84rcKwvRQz6VizjVNcbY3G0LN+A33zLHy5snl0HTLVekMsXTEveNNFKpdYZX1J
% wa9KNqgzoz9SGx7dgb4zC3kGDUXZAQRKj8uRxvCsoXN8uiOtJOsaCH/XCJwANWAIxaZTH0HvqT2C
% kJk6xF9xYqIZlrhOBGoGysaUUawNkrL1SJSBbHA/Yb4i5gt7PDz5671hHboaIf+5a3eUVJMG/shp
% HmS2uS2Bw1DqbnOge1GWECuQSIdJPkBetFLEEcjutGI0UzPaqiUHmcyB9kIJo4KVV4chCs9Oxyyi
% nc7Eiz8h30nF/3E0mFAxfp9ahzv3Sx2YyHQr+oT6OqXMAc4XDxjPz/FNY0v8CvlLx+K1hf4xdP8p
% tHQj1O5EGxKYE889vA6iBlMcL7vwmhGBstpD10/kfPVze/Dnx9vXp7w9nreUydVKAnOLoqdh+U0N
% +cC+arMAguk7wiBYG7QSJEFO25I7E/a9Fth8+PHn7dnvFEPiwZUvxbCjM3yzPBXFNAPQzEQED6cQ
% tqddFYNNM5ab4sSNG/8G+qaZBXpkxYZtEAu0F9pwE1pb0BKBrV6FtZaTumGdFTXsIuaNeqRz2rgD
% fVctHU2MPY0Y8tXK2KX4fUiITXUi0M9H0L2YoOu6doZALjRDN8xHfj6C3uvGqOOjczOFhhsWmJkP
% qIYrwNhwNB6loV+wuEdbAU8uQiHwRPxO+R32+facD+drnj9erlNvkwkJAgbrJwtexaUD71hgtzva
% pebsHVuAUL+7Yy9CMcSvgT/Enj5/9peXfHl5yr9u95eZCXddzR1a0RIWc2xxpoU9ySOwh2fmPon0
% IoltYuUPVB2JLRil67/Avl1oCvcZQWN3hB6QuCOh+GVowgnt2j6w3Dwu8+xeqDGJQ+6eKKBjtn6B
% /cjy/wFQSwMEFAAAAAgA0qlVXT5denzsAAAAYQEAADMAAABvcmlnaW5hbF9jaGVja3Mvcm91bmQw
% Ni9hc3Nlc3NtZW50L3ZlcmlmaWNhdGlvbi50eHRFkEtOxTAMReeswguo+gLvswcmiC2Y1G0tUidy
% nL6ye5wKiUkGVnzO9f1EJTGYW0pQNOcZlHam5wB0YDTYSXlmUh9rkwooE+xcGyaIK8XvCjFvJZHR
% BHNWQOdUPmBDaTUqF6sjvMtEhfxx00SGnPw3tomtr7sivA7hbQjXIdyG8Dglhrp0ari7emahzbed
% 9ZGhtq9qKMY7wYLFvc0XWMCnPUeNudAYbpdwr6vf57CVstLWweF6CY/uTFgKywLZYdqicRYP06QT
% skSCJ9vq1P/o0VFsPcJOyX7OmHQYqXgbhc6SenfQ5K+2aXz5BVBLAwQUAAAACADSqVVdORuZFH8C
% AABSBgAALwAAAG9yaWdpbmFsX2NoZWNrcy9yb3VuZDA3LzAxX2xvdy92ZXJpZmljYXRpb24udHh0
% lVXLbtswELzrK/bYAo5Bik8F8CFImyBAkBRNDimKFKBlJiYikYIo5dGv79J2Wsl2ANcHw1rNzM4s
% 5dXXu5PTW/hycXJ+dX1ze3F6cwy9d8+2jaaCMviy6qMLPkJrqzcIHrqlhZfWdZ310LQhPEzAhw6i
% qZvK+cdp9j3gZWujW/Q2QjNjE7ic6QnUYQHNL30MP0VRYFEIjdVc5XICTGt1v59J5ZDKSS6QijXJ
% coJ0ruUHxJwPiVToAuFEIT2nSYAXBd2lqkTlY7sizxmSBSGSFKjKlVaSYBPKGSVE3Gf/5gfXV5c/
% 1hb8KrU1OLynGZ1OOZgO7mY5IQRVz1pTdjjZT6xAbU7I5wkMaoIlt+MqJUKr3TKTTL2X7zMA+4o3
% oHK16/BAMEhtfReHLSlaHyrkW9cUAfkIQBLiw5zFITk1znvXPBdqXzlXJGf7sio8BbWe12FZi+0o
% bCus4Gn+w0qhP4wrDj5WXhS79hXfkzVNm8phRTL2PxHpln8k6+3IbDQEoVaVfQlVSsgPiqjTMyrG
% 50bXf5RRPrWLk3KDOzghGz01EyhGLdKhjjLztBwUyn87ucGdZqpq06Zc2vIpTrMz01dLM7ctbjn7
% 8OBKh01xEyz6lQC251xnc+dD7XAVLmz64U0XWpiH3i8QQPEjMx/a2lTut8V9Eaq3DX4oJHW27o1b
% xh5V1j92S3g2VW8SYOMowQjJnF/YxuIXmjk3fYzOeLC+r21rNnK0kNnA84uLFpqjsp+jgzIsbWt9
% af+qCpqtBwqVecO0LknjoC3eY+lEs4j2qyP72gSfYO9EzbOrAKvOrsREBl8K5tGCi9BHbGUioDUT
% 3+qmC50r12+DafYHUEsDBBQAAAAIANKpVV3jAq5cNQoAAHYbAAAoAAAAb3JpZ2luYWxfY2hlY2tz
% L3JvdW5kMDcvMDFfbG93L3ZlcmlmeS5wea1Ze2/bOBL/P59C16IXPejYstNuapfBFrlNUyDNLprF
% Yg+CG9AWE9GWKEEPx3bQ734zFCXLltzeAZcCjUQO58V5/EZ5/Y9+kaX9mZB9LldGssmDWI5OXr16
% 9duazXNizOJC+tw35jF/fBRzwWVuMOkbLBV5EPFczA1fsCcZZ/CYTYw/FAcjy4GIpb4RilnK0o0R
% y3BzBmxPHtM4Mh5TYC5imRkiSuI0N671gsEy47qkiVgeVNvzOJoRI5xH5dY8DkO+z+AK9Mx5qtkX
% cp7HcVjvhmnxMGfzgJ9cUU1pWicnPn80Vom5Jok1PjHgRzwaMs6N9TjleZFKwx3Y9vtqR2RCol1z
% DieurYoGGZzJIuIpy+MUWPXUis9lHAmp1xSPNWWzzFxbkxUdqIXnQITcWL9JKB2M1/0+TSYrh7pq
% r+JeapmnIjKZ1rI8FnIJK5euug7Umnk9dzpmZ0mcmFaTByt5REVoMjLTTOIip961ObCmtlmycvDX
% zOq55eHHODUEWRtCGlybx2sVtEvWu7fqyIJs9o/MrH0ifXQzBg084SymDl3bm6a+yljY1TfEfB/1
% JtmchZy6mh2jzNH6R2xtDojWvrTlByaAPswTIFXxs9ctyQzk/lpHjHkXS24pRWY8lWa0C5UIr02f
% vDbdPZ/3siIyr00MXDNyXLK0LFudX1pKrSUqlTL5xIGj1b9GomNy783FTuiiIRTsh9siKHvaFF76
% Bd2j5C/IEsmWWoNFbwmyDrVwycIZAptuFT4lcbgxpVYj5Tp0JiLST7W/FzuW0nEbdw85F1GMQTSH
% 3Juyt7CsyZYqFSVZWM3AWrw5L+2kePspJ3iabHc0PKyo3DEogVQi+gHV8JBXb48s44dcqn0sCaUy
% LMs41BIWhuYa3pW1KrZEVJIqabLUSRODUCqig93hcVYpr1iBRg0eXm/dlLd322ARiFDXxFcsLDDI
% 8brIuk4VvKL6hjaloBVPM+4rSsgIyuy1s+mqGrGP1RpKXlSXvEa9s5P4+aDa9VwgfVNrBAy0OiSq
% FRr8RBsT1alEb+AgMmwrd1OGZV2/UwoOnuAirSLWEzI3U3B6FaPQaOKUei4Zwr/e+dRLy42jd4KM
% vHHquNPSh0xm1CvZ2Gs7se1FXyeaJKmFC6mlS+FB6VHWqfC+pMjvUK7uRZdN6SDNapcoXFQeSOM4
% L92bkK+1e4dNzo0rYMCdauf7PBUr6i1sZd9iepi9rqqoSmXtOtUyBoTtkj1okn9tJDtIpOAIM4AS
% ADm+pwI8Wv0+7AY1uS+eRE57WxVNJbFSTymMAZVYb5KamjlUHbD3eHRZi6Iqe7UFZyxJuPSxyje8
% ChulP598kQK0MO9IVexeG58BAeEZxD4cYZHxiRVZJpjEnpqLJBRzVqIXaMV8nUDZlLnQS3AGiilE
% EQTimfY8NHQR8czckuc6qbbeYGo/w3+9refikzsles2dOnptMK05gKsAxWyhrO/cvqJwD4NdYSuR
% wnK/AUP8Lf/pjle01GHVLJj4s6WVcltrsgRUMqy3a1CiuwDUzF0ms2Yw3DW0ws3ZsU0FjMiGluYo
% gAKun6QcYAG0F4dujtS6KI7Ayw85E6GZkCUABJ8APEvz+t7ui+jhJbikavW7YZozxw+sb8sePpgB
% 9ER4sfrJt4Doe03VnbGwvKct1E0MvckXKBd904WWsAv9vGkQ3MGXKrC2JaWNAEBVhRx6m/2lI8Ny
% y9oLwq1tK1V3R5d41IcKYyvly20L8moJzVOZYZtqDW3Ry92yAH/oglFIU3soSbE2nv7298erP41/
% ff746e73+z8/X92PjUIKLMgsBKgt52GRKaid8hCBvJEH3HgG/J9zKI5QgB6Jwp8ZiyAT5NPZ6a7X
% JKiBNyJvyS/EdYk7Ii48vJ/uB8de1Xn39iA6soR6qs4uS8xyWFURT+iLOagFx2pqluxHfMr9Qo0T
% 1Nu1u8RqHujkvzuHdRGwWAldzASuA1sIdnFjAKULsNhA42yQ3Rvui7/yTq9ZEQYM4NneoFXzz06n
% 1VxQ+awB3c7J+cC12lA8bZB0gHBtBDiojDUnhfYFZn+gy96oRQ1KwpQITR6CotHtywGxod+cZRz6
% 4y5REnKbqSgwR8Q7JxfEHRL3HRmeTy1ivm0s4fsvxIPnIVh0ga8QMt5woJ5Geke9IJmid98TbwRP
% BwF1iwJvs32bg1qt/8FJkt466eSG7pDGJMiqVL+xJj4Pc0YxjxKEdiPVAm4hBGwdBFaLJdzqsTi7
% mbaooZ1AP8IjQ4gwJa4OqN3gc9MKqv0wpbTi441L+mnXDcs4jVgotjDwo736uo+FYW0QYJCMKiRy
% AzDEHSqMo1QtUwDniNYpFSiVJ6GCS3JD1Bmi+Fltc9BNRbNUuK2wb/SwIgm5uQ+IE7sonY1vyKgC
% iVYnkyYibCJcwPe7K9seOTvPCxZSOFFVoJ8eCVa0VveGFN1E4Nhg1W0z/szAqQAtq5yWZUY7JqBN
% y4bVgiTW0cN8naeMAlWw+hFZ5ZbSQooiHXX06IkqGo7r3QRwyOqDO/yvaSmqXPRU1OAnkB+a2Jir
% 9OFBJzEkQwkJMuFz/KjwlAdGeTuI6mA6ni+PpEMdpmrGGE3Hxy4SwvEDHb0f/8zTFSYFeoBG9P8R
% 2tpE0UC2NaatOmt3wqOMGxQQZJ477jDOF4+PapK+IUAymJLqg9LPGvRolyHIw+oqUI3u+CygsiS9
% eTFTnycDnnI550dvBv3dXaHbJqTQKepydmFhUDXtnh4zByvrC3sDUGXUQMTZd5DSDpMSfT2efgVB
% UGIhzgqeGQl9Sb4T45a+3MIv6BJG8u1ibLwAFw2rDoLiHShIqXnxdnR+/m5AGnN+q6SiNqrq7n3D
% w9JrKQCnm/QF9NURea+684XqyecX0GABCeZCFrtCXg7bs3a5mfg43e5mwSLDj0QzAHCW04XaXctR
% WtrtTajxHV+q3ja6l3YHCIFYg2rkXJt+yY4kzcgDTA1qQNM8PznaT4aDVkP5YVHWn4uhILecU05p
% SR5019O/KGJ3XZAdRdjKsgOTx2iBt4RYdOhf4Ms9+teG+l4PiZMtcNBU1dAI2QYQpcD0Fvnm7KSz
% JmsHgiGm1udyoHJEaXV5AGC6FGuFtpJL8cIdHGXUnBLAfLLUI0tQjivL9jeEyiNuR92qNVbsHapZ
% 14yRaydTZUgnyypvlb5tr+qiU4blvjcFP6gxVTrv5ijj97vbf1cJLemLxITmDAapJXXPzs7x08Df
% FGJuAPkNCZJh/A3KcFe5ihc+rdK+KcPQ42ooItBEPulZOAM+kAffu+evA4gsj+PeVl5cdNzxlu63
% pWO+PYKAKJW26Tr6vjtrfQZQNOxVn1NahV1Pr398vIeBFSi1T0qy5gwqWcTJHP/qg6KvzkTOowxm
% 4MqbL0jwHXynaCrvafZ3saHaoZgDEGYwFLMnbojMKDLoOiyDRIH/N1GSx/inMDUNo2yI1YcH5Pvw
% QOnpw0PEhHx4OB2r6fvkP1BLAwQUAAAACADSqVVd9iGM5VICAADABwAAMgAAAG9yaWdpbmFsX2No
% ZWNrcy9yb3VuZDA3LzAyX21lZGl1bS92ZXJpZmljYXRpb24udHh0zZXNbtswEITveQrC59QQ9W/n
% FLQFeiwQ9EzQ1NpaRCIFkqodFH33LmnZdYEkVooccrEJabU78w1FfT1I5Zm06NsePComdcMcQMNU
% C+rRsUE6B82a5au6ZqZpmLENWHfHeJlz9k12HVNgPW5RSQ9uefMFLVDLwWIPTBntvB2VR6PX7Ndi
% WKwZX92yRU+LtKCFpkVehdWGVjX9W/hJA6AR0OzAhQfyWNnLA11CLxrYWYDQIKXLcPBgtezEINGK
% Sy1UUVRJQjXK9IOFYEQYK1UHQlnj3NEilcX5qBsYgH60F71pRqqyIFUrN9ihfxLU0YfiVZH8vu4y
% 4yeX1eSyzqrJZfacy7Tg/HmX+TWXPKlrPsNmLJprk5dVPsMnj/7+iTPj5zyfdboq/jvPrObvH2iZ
% hUQ/n1syiC9F3NtyJIXBaxu8Bn6PMdNiMp0l0ctfEhZCFxoZpdGrElXnRyRRH93ocBvuW2O8EwNY
% 4a3UrpMB7Lm9HA+kU9onqgvSY/HUIWztKAG1iGrFHr0m9QSTaKGPagMp34Kx0IvO7GnOxow6PpxP
% AWjYC9QacNdujI25R5fa+JDI4p4YDOwkj9hscUNpMgxUKZbujtEEB3SEwCW3y9xuGXVjNA37sf/U
% S69a1DvaTaP2brmYhz6b0GfpCX2ezkBfZGXxFvRlkVxHn1whn5Yvky9fIf8RuRcn7uctX5QzuFc8
% z97CnRfJO4DPkpfBVx8I/Pf7h4c1C48dOzQod9o46u9YONo6oMlL9kNjODolKSISWxdnoqZvMgFc
% ejgsb/4AUEsDBBQAAAAIANKpVV2zZHtv0g0AAH4gAAArAAAAb3JpZ2luYWxfY2hlY2tzL3JvdW5k
% MDcvMDJfbWVkaXVtL3ZlcmlmeS5wea1ZbW/bOBL+7l+h3UUjKaHdKG23t1YYIOg2WNxhFyly98nn
% GrJNx0z0Fr04ctL2t98zQ0qWk/R2D7gCjSVyOBxyZp55SP30w+u6LF7PdfpapRsn31brLH0z+PHH
% Hz820aJyVjrVlXKWOrpOs7LSi9JZZYWTFzpRTqGu6zgq9ENU6SwdOVdVlC6jYunEel5ExdbJ0ng7
% Gvwr1RtVlFHslFVUqUSlVelEhYKWbKOWjk6dPMpVMapUI5w0q9CyUkWBrlWRJU61VqVyKlVW5Ygs
% G3DrIotjtaCJS0cneVZUzlKtojqulnpRGRmYXlRZFncSiyzBUtnc0ogkUbVue3V5V9iRqyLa131h
% Gwb2/abE8wAzOnNdlV7jjwfO/VrHymnw5FSyORg24VareOlUI8jMYpVeV2vPHwZhM5SVGZxHsdp6
% OY2OyhLWOvmrt1K+cbCTThTHXv5qyTvO21RE6bXyTgQbilFHge8PnE/yUR/qVznL6Z1cIHL/68C5
% lJOyTrzg9NTz9NHSf5X7O42f/Kejcn86cOYyVU3lzblz3p8agzXWbN3kfGJDvfkwgN5dA9bqPWL+
% YSDm+DsXVkJ4wXCO36++lD/DdLW8VqUsVeXh5ZkhtJGL7UJOtIDlZrw+ovEwkcUbsaUBDzr3ICjw
% fxKMp0f0Ow6mrKDd18tJMz072x4E1HartrKq81h5JZyplp4HTb7vh1YY/e0K2cSQ/46i5dJDl985
% i9bJXVjPm8N8gJSo6iJ1LsXc+BeWxHrhJWSK7dt5oyFvJP5T/wbCS7Da169PfNPX7PoSeIcVR/VS
% V9456U0lmXHuh0lU3kpSnSLIOhvTVyfdM0VUxNG4yOoUuy6llw55LvYb1uxFZ2f6IPBNBB58Y6Xy
% 2HhHRLwlaZ2oAqmMSZ/7LeVtp8abXaP2x9aEOfLRO59oOOMG0/wgbcPN1Mxr84IAZrZYq8WtWeRP
% zu/I3rSKAEjpNYxbrM3OAEqcJFvWyDy9Wjm///ux+UqtWV0NF3FWGoD5ddYIUnK/VgCe7fDswUlU
% hPR+ALYB19LrWpfAGYooArIiwySYLyWI2ehq66hmEddLVZIWBTzbEnwBtbBnqUpySPRmRESPnA9k
% PdYLgClUtFhHcx2TpqUuLHCNvu+8J343m7rMEqxfkuhnEm38EJBaVHAPkrJxVAygDOzuF2pD472L
% CK3in0WtTD6wLWopq2yZSShhDdRh4IuaWc55kB69HAzpr/8Uw7iLlQw5pcgCTDlOlb5ez7OilHDq
% w5Sazzn1Hg4CYyCbT328hgffPzDrYi0ksa+jmbJwK3XwzYxi6VTdy0744JtdWmh/v6Dvng3lpx4Y
% tHsgjU6B9Iftwjb3E8Tv8vbk0AI+wqKsinoBBBac2Mh22QJ5eCG7nA8LaZA6QYodBSGyL2EXJ+zi
% cznxCAuoGXBwc2ihgNOGq8ol0gSo9MW7wC5AREPkBcB+CSScrNDXKHWxPB8tsnzr+SEiEL1LA7fh
% Ul0XSsnJ8fTQROB3APiZbg4OA8uHyVEjPIvNPr29fl1ww7xteFUQVrOWWmz+FKxbB51PaoTMxqD1
% C3ANXS/DtV1ku9geZOMfaf1MMb8Jzycb81jbnUDXkQzalw297MG8VWiA/jBheEyixjMD/FP59rA4
% 2kdahrnPrSPIl/vQa+fS06c7bxxCsdfCPKHOBYQM8qz0HMCTR7oQwKUc0aareOswXIL0lGBCpb4G
% zwEUoY6rAtPv4VwB5LnXVapgbCmP/5vrTYXtEyevjQRxYpxG2APH/yC3+DsGbDyJiEGb16arFxu9
% lGzrAnUYtNjVBmrc2saws5s99JPzocjKcmjWDqpIpuYFCSyRA9ECkBZdI8fLintXdYydAutThY5i
% /QCp6yLK11wddEoYz/WlbgDWxGGLLCO6Cm+X6IhVlaXOb/DusMzqYoHp0KWX0EdobiacMVeVxztM
% 93LK3Bdw3aTTR0mzzNSSqaRIREG0ybcuuNundH1S83FyN0UJNyjZBVoz9c/OvDua8oDABQOiTaTj
% aB4rGSMMPEYXAhWgKdlXwDaON57QroyKRzdOdE+UrP238cm0CwNKQjN6TMVap7ViaxdVDSTSqQnF
% WZ0ijuQ/eHuOQzgsydK/S2uqwQqa3WoyuMAqvkhK4XBf0Rd5QY2s7gultFe/KvxW7YFpQRj6O1MO
% pC1CzdQWUupD2NDZARDp7c9w8M3uFYvuZXFr8EtOcrJ4CVJgovgSSHt2dkcRTM5GGfe4ahsZUxlp
% B/zP5FQefkWDqe2p3HFbcO+kDMYkdiC/2fX2uubcZfaZm7sFHkmPRxlLni2oxTDeq33M6lSE/Vjn
% VLR18tHN3XEu3MQdJ8JN8XyI37k7ngu3IPxCZs6YNLvjPrBiRNRwDswMMGL8Dl+F2wLZjHBvtoCB
% eqUXlHjuuMME4e7Sf2YtXBBAMD5Asm+1cHW6VLnCn7SaGRI561M1I4YV7LHRr4PB1cePv8oJTjgn
% 4o14L/72VTye0K/4RQTHeHnTfwl6LwPn8ViYYe/Q1T7/3Ht+i+e34p34mQbxCPNmtR0L87r38v7r
% lKlJjHKnCm9N8Qd6QZW+NPSgpPgks6fhJTOANlbfo4ziBHm8M3LwlLx762FgziMU1k3LWLeyAak5
% H5Fj0qXHs/mC/05bKOGKfzmmasulGIz1fLI1z5s9oU9jykbbEXYDuHJ3PVs/vGz1fOr6KcjbU9UD
% a3xg+Oit/HwamlIKWYyg6JGTaajiWHarG7QIhtaxkfCOhXEy9pr3hpjVk90JBORR7/i0ilHdfni/
% HJ0c3ojHbmO/PFJLxQqqnYKbo1ZF63lIBsd/IuoLylSP5W7+h3EImu8a0umeWja4EVfiD67/WBgt
% 0CLD1akkDw06aGevjxgsvUMuR7tKd4XK55H4l8fNV3/wBF6k/KO7Mrjyh/TzB3BmhyfwmKDpDfl+
% sVDy+dBcWFk0oVKukdYLHOmylQM/Nr5Apc/qHJ3zLbOkjhoxn0L5/ih7V0hUA2AtTbqKwVnvxJZ9
% fPdK5i3xvOtOzQTafL3gh/TI1HI7MCMDQSzID/llLsB7fI615rRoadZi54XCFI9aNiCUC1tY69Nk
% zMOHgahbTXhetDkUvaAggoLmmYJ5pyDAc7S7zKAMeuij/S6VPo42UVyrEuehjuha93y09xEJaOhK
% A3XXRH3XIpfByXuy5Df2nuygKbyVJ+8O12Ei14e3OCLRc3ssOyykBAfYOyiF57KlKOHeSetDe9K6
% xRrmUakMCDChuSVCQ3X/8NbfEQqmPb8hEu0parPbtDVl9oWckB4KVe/DJKIj1wYKXhB+uue3NLxH
% /Yvs/sk1S2KvWfhont3T0AtCIhKQoJtEF3WqkzqRf2SpCrkYEgsBHTn+TjZix6sr+SeLvvLDq1yy
% 6MG3c4wytKpHBK9y1ADDyY7D77CyLpt7pKxjX80e+2qesa/Gsq9nm8ZqGolNPorC77PgDnZyw8Sk
% ZAbUkjTp/RWW1ma7bHc1VqvKKPofCRx3/VUSR2cPTEbvTy3ajRxbm46kEd9dhZCRe63/d1LY7stf
% 44Tt5vHuMfrZllN5AlwI+qJW8SlOzLYHSKEWlbzK+9qHrR+H7YieEjPiTN4OATvD4IxX2SaK6aSF
% 2Ra6+KPssRdNOvVshzCivJJ+ZhHD7L0L+0vxC9voQgCpeSSB1D1+u3bHa+HeuuPbluUaxlsoYp6k
% grgn0M5yXyaOaKMcxA+fKGe5KmYV4q2M+VTtjnlGLurmn9udQTGCiCkPs7rcMXoxu05niurezDJg
% UOcV3QbAKrtwF6ferFDJLM7uMeUcs2A0QBfVgxl3qu5nOu1u0Zh1d/vRMwdVDpTcPY9jJ3day+kb
% DRVPkCwi0eDjcWg/2NAXHraNj8tOn7CbrzyYhywc4jC+WGu+b8UOlCP3qykpdDPnmdp+YT5DGXVM
% 5Lisq5i/JtExHflxBzahK63Maf2SyoRDvJ9oPR3ObUS1X3K8vwnwvqP+KwrWaUCnePiuuxZJep8/
% 3oHGHR8H9uLj6fWe0c/l8VQmoVHDR6PODJBOo5U/lkyCX8TJG/EmEG9xBnhPBogAbyc/43hw/GZK
% k+yVvN1yojxnkpmLeUe2n13g7L596OmTcpQPLee1F/s337nlybtbHqvL41vJz5eTm6lvbw+0/4V+
% bnzf/94spvqbGG43dt1D3ED8cmInesYX+mOP+PBBAlg0jmWYxzXELyp0tU4UQtBc1iiEpjn0OTmZ
% vhy7wnhEuNlyaZ0cuqKvXrjPYnXkilVcl2vJN+gDikFQnb4nBUfIxCNnvvOFRw58j1/yJhrYi6j4
% cv/eOGwVta6EiB/aJf3KnwjsZ9ZuHNyCRdCXx9GyTvKSh+xbtzv9dvatTaQFiCwRvOuseUrYuGTb
% sS+Z9GF3rdbLas7FPzOK6s2jSyh/jY3mjTc4xOdx4xLrhf1j/b5rOAtmXQq44+5RuOabyqy/VRBo
% t1i4+6slsG1X00KcWy6yvAW4XjDxYkfOVULLzU/h1C4EGOJ6VwjtFSy2qPfBvP2kvckWJrMIuegq
% kr+GI0jrKlutRqhbOC1kEd3ROmWO1RQ4mSiooMtKDCraj/GokkWUZ6ZusA01+QywSXe6awfLSD1T
% dMoROcYV7r3rI4Gd1bhzFJ2ExUpotlye7NLp8vzqakwsxvq5/+2fti1WsHrk7L7q5/DPqt2M3af8
% kesPUJlnszRK1GwmpTubEaTPZlRiCNoH/wFQSwMEFAAAAAgA0qlVXVI8t3pLAQAAHgIAADAAAABv
% cmlnaW5hbF9jaGVja3Mvcm91bmQwNy8wM19oaWdoL3ZlcmlmaWNhdGlvbi50eHRVkctuwjAQRfd8
% xSxbKUGEh1qKWCBUCakskPgCYw8wwoyDPSmkX99JCFTd2p4z515vFtvtB5wrL5RQwLADhxzOxEZC
% BGvYovdGKHCCvZ6c5qOs3+9nk0Fv087izViBVJ93wZOFVcuIaMXwofImwiUnFjxgBFK0kBCmbpbY
% YYncHHccOVJ0OQlGIwhl8HUjY/xjts6gmMDLJTu9QmkoPkiOUulNjU4lqyjHJ2Kv1BDpp40wA1vt
% VFJf20hNSF2cD9/ei+F0PJp2rIQ30UdHk+5Svga5BpVJ6v6NGk59YgiS2qzNQuDAf8cdp73oPBIk
% idSynhxiGA+63LH1U4ANvjozpBKtxu6807+ubcD9nixpIVf9tbYHdPk9KViMQnqtSzNdYX3liA/A
% 8wJ2oWJnYtP/Yr2G5epz+bWFBj0D9KZMWl/RHw21ARvYpd4vUEsDBBQAAAAIANKpVV2yOLbChQcA
% AFwRAAApAAAAb3JpZ2luYWxfY2hlY2tzL3JvdW5kMDcvMDNfaGlnaC92ZXJpZnkucHmNWG1v4zYS
% /u5fwaYoqvdYku04CbRA4STI4doim3T7xTACWqIcJhJli5TXTu/+e4ekXr257vmDTQ05w+HMM89Q
% /vGH84qX52vKzgnbo+1RvBQsHJ2dnd0ecCxQQvGGFVzQmKO0KJF4IWhbFnuSoIpRkOSIJoQJKijh
% 1+hB6aMQ2ejpmD8cvdHvBUqzAgvKNu62oEwgQbjgCJcEVZwkHrqjjAqCSsw2RMt7mzqIFULuWKTc
% k26N0rLIUVxkGYkFLRhHNN8WpUCLomKClHo+LfFw9q4WIMzR51Et5Md8e5QS3kgEzcmIC1yKSA49
% +WWYo50jnHfnEHEPVNZFxo2fd0igd3T42RwlJEW/GW9OaV6NEHxKIqqSNf4Y1H51fRU7iijTxzR8
% p7R9U0lf+9I3O3BLU9vkLzQVRuwk5tWpzb9e7eRqr9WdvTQQexDDnBvmf7XyjhowE+06XV7lxt6y
% 6Kkrr/V2EOPEuHXSivV0PCUFmRGblpUr3djJpfZts6NW/8NInX0dgRQC9VBkRxAJU0lwLfFkdp4T
% GgvjL4M65lVsgU+GS3UspMiJpfXUg3Oq8zQm1v9ogn5Xn1V5lFqpi621dolzAhkHscfEizE2o2jc
% z99Hu7m+3O9kK2mht5mH+TM5bEuAzY/otyoTlBOJd5ZAwCUEcYagYoqcMixkKnT5HB2EswwRWXXe
% SO7w1uUodKZ+HdsvkcSab9q9n0D/TPRB/4w0cqQoNB23WzSQ+3r1IvriNrBaAgBXZj86X6JoYQ+n
% 1Un+HMjHgFt/qCjRtvD2OKuIhEgUza03158pZTi5sTA/6WgnUWNG68uT3zqcbpg8/rLxNjAdQ5/Y
% cZvxRB1upePS6J6i8woly3hlR9KilaulUN4bItp9l6Hz5gYr81rLl28rN/L7Z4G8GNJIFNUr4pXe
% Sm4EyTUS8z/yR8+aiKYo/iGC+G5L4Dvj7OGXp6crlNdQUCHoAyDGLCZZhjVfqcxHoeN5njMdn0kQ
% 3WYkB4zg8tjjWvSVihfEEGFxkQAbr49IRDvLYg54FWdVAoSr2PoepQCpovwAVP64BtVjBHwRmhZ8
% QyjMc/mr03EfPVqG7+4sYepf6w2GNgxbgQHJN626yp4i325k7uO58Sjn5LpWfQgTT5/euHfbpU+W
% sDuzYLitTHmAckiXTVnITx7Bc/vU7gDFCCE39PZldxAjB6pt/G4+bv9kOcy6on4oB3raqUGCVeki
% 3SBoDHFXJQ9dCpytMlyinQuryYaUvTSqBP+LJWRLmBRmR+ht+baChsigueKMvkNyAcvAHZD0PSl5
% xZH4WriScurU0ncFHp1h1Q6W3HusyQbCBAXTew5kPfWfmxL6Bh/zXnDvoroZ+C7Q7aslPmpevtlF
% 82Y81Pje+pPkhlC5d9Efxh10lXbN4p/I7yMaU46EkXKjZpIM5+sEo/gKQB6DcfP6JugWBB8saC3p
% MgFRXSgwUqXSd/E+MlTF7Gug7U8rZt9VjHkun6GxBq06gBUQQxLt0aJzRmvGoKWVTOv+3LgJQde6
% CcxvUa/67935zdiy5m5jVbYnyp/fSVl8q5GQTUkIBFw0dD0ZApx2MK3BLl5ombg1OtEWtpS0Bg2u
% a2r+FBk7581EW0xLDffmZsm3GT4CutOiKoHLajOOrBkKx49x5sbVGkqpswZKcUklczLokvvoBObX
% /wdKJ5BQKXzuhBPzqoXaQ+T7wDrWzJ2EEzmY2rPx5YUcTdxwdhGqYWhfXkyngRwGrg+fMay1J9Px
% dDaKgcah5sXAudns4nI+B0PzYOw7F5dBMPdnwdSfBLNJGF5ML8PpJLycjsNwZo5OqevObWxahgB+
% BlfkYNIM5mrwoChpI/0/gIfKfxgE9nQ6u7QOLmw5uRydEu+DK49jbTxerblxcITtz85rzu2g0Qu7
% sXEOMOsGF3M/uASnFc+1+nM5N5n4oyESHyT0Ylmgz3CXF9xwuVcUjvwCheB7q8fdygEkT0HUYnFA
% jNdI46h/DNQd4OykU3NygFcP9AIvBgrmwMlAuABvDpS9hzcVAghXfqmjy52BrFknVihvcg/mdhV0
% aYKIvA7VbX7Yo2VBQpEKwrh8R1nDyRNo931C/6x5/LOmb3iSX2Et+p8M7vf7I46kFXM1oFzWrw1/
% 2luuVDy8lSVv4KXrryTtMfutpj024EZpC+hB1nGvrIbGoiVespUsGbDHwKC8KsEVJuMEjU0LZHCN
% PPUqI8zAptv3++Rutt9/GiutvQoV/oYN8XIMlzd/mGWVtraxclFSlek2y2BpMq55rr24w3tnlcOl
% D+iUyt6ss6nyvShImtKYAlV9hXue4juSuAlNU1Kq/Dc0pq9uTZJV8mu8cp3xQU6alMgJ3E30MyWn
% 1t1Ujg9QMBhC5mA7OEmC9CdSecSmzufadPUzKNQS2x/eitquJK8/3cq1i09XnlKX3K5tP+1F7oMk
% xp+ice9WrXlgYMZ0dvBylWXPsYy0evH84P4Vf5wGHV8Uw3YUplWf6WqQRX6TD30nqw3/8uuvaHF/
% u/j3E5J7XANU8ZYD2/zkBfCCToCWE36GfkJG798CV/1/AN79DVBLAwQUAAAACADSqVVd9jOscagC
% AACHBAAAMQAAAG9yaWdpbmFsX2NoZWNrcy9yb3VuZDA3LzA0X3hoaWdoL3ZlcmlmaWNhdGlvbi50
% eHRdU8lu2zAQvecr5tgUVuRFlpfEByNNG6AGGiQ55EpTI4koRaokZdd/30dZrdMcLEskZ968hc+2
% MwUtiE1wJxpna+LfQgbax3Uu6MBOlUqKoKy52mpNjQg144E1TcJ7dnHLU+eZlAlcsfPpV4cecfmW
% jKVSW5w3VdJanCDfsgwO1XxQBRvJN1dP25eXAViiKggTPCYJYpMv54t8uUpnyywb59O7VTofkbYR
% vAH6Jlul+WyRZXeTNM/OffY21BilxuQBDCpt90In8TTJuOSUoFIoTaV19GsTC0f0thmqaxYtqsHb
% c+pUVYekQMkBDA68pkWOTWmbFt97pVU4UWML1n5E2Wyej/t6fMym2WpG7cmJRhX+3PrIsR1G8hhf
% OGo6HZTnkIKzgSbYiYIENF3TfJpPIAaXUF9hET0BrLsCQpLjlkU8r8Ue2CTgIWyhEjMmmgNIElR0
% Z1wHIZXB6ec8Hc/J7n1wXe/PGvZENk0Xen6UaFu9R90kk3Q2vqWLBP9GJFmz/BnZcCscptGnM5w8
% Sc3J0bqC8GMHLpMlFcqdKfbbCMYqGxzXLHxIuKgYLYWr2NNRwUJE5kRxeZCv7BC/F1l3DkE6gqGM
% 2evDyUgLML6JznslTOL6vCIjThQKf33woGChGja+z2uone2qmqaD78t8GOeRXaMCukBVWVuHkXdf
% dp8/oD39eB1RONrEq3hNhJRQWR2gy+gv6EUzIEeHJlMYZJBA2lk+GjAoVFmyi3dg4FgofwlDJJwA
% 6ByUVsMSyJ7TJ2PjPSkxrXX/b06vz1AXk4TnC9neDxoYx7CAc4iK4K11yjcUrWw4fFD3ffbMZrWc
% jbMeJt6a7W5HD2/b+1e6f3y4//5CsfnVg8Y1AIejgGkBstMnPGO9Nfp0jYt0M8+myA7o4n78AVBL
% AwQUAAAACADSqVVdeuH+fCkVAAD9OQAAKgAAAG9yaWdpbmFsX2NoZWNrcy9yb3VuZDA3LzA0X3ho
% aWdoL3ZlcmlmeS5web0723LbuJLv/gpkUichKci62U4iGaljO/bJ7CQZJ5rKpFalcdEULFGiSJqk
% ZFte70ftJ+yXbXcDvEvJzJmqTaUsXoDuRt+7AT5/1lrFUeva9VvSX7PwIZkFfm/vp59+Or+3nYRd
% Byt/Iids4tpTP4gT14nZTRCxL6+Y9JPogbUP9tkwsf2JHU2Y515HNjwMfO9hf+98LeF6aSczCX9c
% x/aYHccyStzAZ6tYxsz1EzmVUcwA4kUE+ODNPvvNXbr+lNle4EvmAj4vgOn+dH/vwvXdRDIHYdwA
% wARgaMokAzQwx2sG0URGLIyC4GaADx+YHUnmBwlzk5hFMvRsRy6B+n1c5t5NFCzZjUYOJC3DIEoy
% aoBidqHGAOYoCQIvG+MES+CbTfM4C2W0XCXZXRRMVk6iZjqB58ky/Im8sVdeMnHTMTcr3ylB96LV
% lWM7M6neIxvTVzdAWxC5tseBPbeRhhCBEJBMNeYL3ak3ibuU6XMg8+bKAanCavb2/vP8y6+c/frp
% nAlmtHnb5MzowM8e0MfsycSw+bXZB6Ylqwh4MWqPG9fwh9ujDl51xjQwXl3XBzazgc1s4HLl1Qda
% OLCJAy0cyPUzQKCftdVkJ/Dnhl2ey2meIgL0SwJ0Jx/hWDQGf/QgP4iW3SoQGqWQ5dBWy41xH5v9
% PQb/NgIZRZeo/PeguOw+7sNz5NGG35v0TsPc7KX0xmAYSWxoKLa4AOYeHZgD5sBlr6euZQQ3rzhe
% RkJGraMDNZpMhV0Y3cMj3u32TMs6YMcwuvg2sqKW0WlGpgAYB2/4Ue/VwYF5bNMYmRDGLu91zMZr
% S0b436DhlmU0u2azYxaB4XiYcPT68NXR6ze89/rgoH3UNY8vjDf8sDSy24Sxb3E1hyYDRWPdBjyA
% gZ1s5HN2NpPOgqzS9WcyAvOZsKkXXIMbcPAOFBg4Bhb7AArtgm9IaPDEjcFCH2CwvLeXoSf3Cd4t
% /5ayj8NSFZLAm1wl4ptldBpd67Zl9OCqeWua6m3cFfDi1mzRQ8vqNTuaNHkfGt/UY5O9FazTsBv2
% H91WF/72WkecvIWdOR4cD77IT8Dk9rVC6OmK0ffhlRfcgSA7jU1jY20AEPyF36OMzKm9XNp6VDbe
% 6lpxt8hYWtDbDvG0MuutIj6MwGkaLy9PhkNkkJPketYnCWrxtbT4jt+0DjnzAvS+S8AiDt60SEuO
% O6BpL80a0OsAHU1FYk2cmouN5IWWcCsQCmffBILa+2fmtIxPwDDlRsKZayw4gxgQL7QluDdsIUQn
% s0K1tHPhuXFiFN2qAU5tKo2FybtaqPZkLpIVqIUhiYIZR/2CYLRayggCgnFuInjC9vbt7IVW8Rpl
% +BCpcyG8hRL+AAMKFBbMGTTUAH0wFET3hZKOupsDfCLD5XMkA6j7AT6Qlg+xQE6uwMVUMdp+LFAS
% NaLQNwRBIvDmBQpjkc0ZChoCtvyi9Pxu5nqSDXPgmvHDFwiJ1lBF88cQHCPQ0BRlKodmCQhgHNbQ
% pczy46IrLMMxOsfHC/I6ykHafuAbdxCvues7mg0Qkj/Ie9cJppENmuOQf/CkHYN3CMLAC6aUSFCQ
% R9MMbpTX0AtxJAshGUj2MbQjPF94UiEBHoaRnAj0+H7my0lsSsv8giRIpvmruVnjI5A8QrAjdzxW
% F/PxuE8o4Oq/BKzVpTmQ6UxEe8CCVSJG4wzx1XbEwJSJO8HMRozmdQoRMyojAtXqh5LUWF/8NyET
% 7XEGby4gmTJyqHwhH4RnL68nNpv3DU04n5u5jIHQfTtEfqavgXUImBY1L8pXWSJM0CIFXO7Srsq0
% NFiDrC4tk1JZ11Juk0P2H4ytXK9Ly9T0zKQdXtmriZukQXgZTKQXFwWxzCdiPMtFsdsfLXN/lEEB
% YygAOj7GBZ2bVb3xHTEauULURLs0a8uAR6XJ+HrBDXA1ZtXl9Ukd3TFwQ+DVfAw34hrySe22Fqkj
% TP8pPqRyRr7C2tPQfX4feq7jJgwCS8Qoh/xy1GofvoxBED4wRGW6a9kEe2TBdZxEK8pvOcOcfeZO
% IXwoG1XxsrBuXCNoKP7O8bfCh8M6Hw41H7aQfKRJJjlTVivCh8heuhN95wTyBioFF1ycftLOBI8A
% fkE0CnAuqSX5DNTfATiWaSTBGEkVFD1dvmwUxe947u1KZjq1QyM46MSyqg6TWAczty58VPss5phV
% 96OjUi73YhgqqSwgQW3t1+xK050zFAYa4J0ti5Y/ic3CIl1/LUaPYEWdp3yZyPaYHj+V115a+C+N
% TgX7xI0gKohHHQHIVdC8O5yna6fc1qBmk3YiFuZTCQphT6lXIMtscgJRqLNAoJX3iPGMx1OfeKa4
% UWcTMBv5cWa+XfQxoEFGuJK1USoXWSMkYNVo0VRzxvuQMi2xBICpI7Xes8aMVjxuCsBtras0P876
% 6yI8J0iBIC3rMhN04gjTcF5HzcNZiiFPFfezTvnlBKDcudk0SOW3MRFs+P+Pi2BtAjMFbXWg2bld
% nJn/BtOB+jrXAbpV5zwutMp6dGC7eQ95PCU3OC2lWEdAhQtSjXFVIiknSgOhfutslSsSJVI8pQEl
% P6fEl44rM8qTN8mPJagXnILIOYjTR8pJxUEEr4yZaY4bosy8CFx+GctFHccPrLuuNJTFk0Jg4qFC
% BTdWfF0JgN8vGbb/S1O40QozidF6PK6PDmdCVS+qeKm910JSTtPAHNwKZ291iKmxp8zFO+LihWFf
% x0Y4M61dan9n8qzfA6syt7C9qrX09Mc+g8SE4gXtoykVR1SJnErHaGCZiOfsIwjdDb0HTMivH7DN
% k7iQhDez8oJhqtaSvoymDwA5TmJs7FH9ul9XxWwNdY34sNM9zMCfnYkP5AJmdVEhUqVJDddyi9N2
% SZUUfza2PlhnFs1WTCo+KhTPN6p6Ro+KvkhGsWzR+CYWKmvKlvrskSyf0g3zCfUvWIbw6tr13ORB
% 5yGcPeZ++UlFOXhWsven1PzjUgWvibiTiBiKd+yJQQKXyqOVlWOQeWA7I3kAkmqCfirKHlADmR6k
% 0CBbZaow3bOvgVIqPbD6uqFUUCYJZH2YMmIvQGnGF6nsXbVnF35w57Nmp9VrF7NGyBQSTBiTme2z
% GMpWP/EeKEhENiopc7WWPGc/++zjH4/NxRMngIBrKnWaiipWpBw1DAR+9TgD3/rHgs3+OMiggP7l
% YsHO0wIyURiPMIsgdHUZ20up4IPGUM77MScoV3Lt4HFg0/am8jqyGQ4Gpwp1MS40mUXBajrTWSXx
% rZ8lyEr8o2ZnPNIFHMQtuQyTh1K5Uk8ri0XL2a6E8nBXQvmXkskzSiUzwrLc8ayQOZ6leWNI3S6V
% NA7Y2vZWEhxeGwxVyU0Ybd7hXd7jup9XSx2PCivz778Txii8KO9HWLP4tTtVSdewM1e5MxtqOcfi
% sI/odSZx1zhTmQQEQcwjsvl6vXfaG2t6YOIOX0wcwRigeLewFOnAGUMxCBvyi1LnVTERpjQ7vNeu
% d++g1LZdHyyU6rWikfUrpZvS44KyC7LLQdEwUi/BHOzlokORoY0R13t4mTZxwIMZIb/NW+pK/8LR
% LdboVS1EhoaoH6r3Rn4SHugGuR9Td0YNyus0pYd5uA9Vm4qK3XqDAd4o8PGdHRo+x3pZwQ+LVZxv
% moMQSOQhwsE/HG8H5VWkq3xwPBlnjRqkyssX1eV+qcTB18Oa6aRoubdFKxPsqKLyFnaRjOGo0x/D
% Wh9c6U0Y3LbH3Gzg0AJV5fZG5E+F2voxXrVpp2GGmhQLdEpxofj1c/J7/FWBonei9/rAMvxmwUX8
% Lh5l/8IA4Pu4y5Ra5yuTv1N1m9y92m6hanvOTtgU5AgaCg4eSoAHJidTCb48YLQMvQOwdslbgrk5
% 0vMIYp4m/D5qc6BuTNsBnTb8s97lpN6KpX1vYKD/vZxvUd/ENMfbe2ruMzE3q8rqF7ItbX+3x+ku
% RF7+E4cnu5JfhAmCIgaVtShL0Cmxgpe4GQQWUCYbnqOW4o/hgpr9wxvDMmp2VekVRbMACa3n4Kpb
% HvdxhCV+H8nyvAW1C1OYxWah2weOUs415vRTSUUhj5kJpHLRAK1t4FUfLsvgQx2d6vzd0nT1QAX7
% MEW7GG3OiIfMln4BQQWIDi1lS/8uGkCiJo0Ur8eiyOxtmXMo9DrUPJAbmVdDdKptGu+t6PW31q3V
% hW4tVaQoK8OWxW+vcbRSjjwuIUyBrFsoasvqlvTSgNcmp5hC+qkmbY+dxQB0LLoWUW3dWnjRzR1N
% yoNSQkqK36QEipqCkHUSw550ygSRRdnGANJRBedJ725R67+JXiLFwO7cZJY7jziNREs7wf3lE36q
% KVet/5PS3uxoRDu7NBBTncWYn8If7HFXsw9fW1nJV4zrPkJtGHuUzF8Deoef8km58z0a4T6x2qE+
% URkWd0yuHpzqB5M6PkW/uTWK4nMVRXWQBttwRCePw6OR4ZBrw64rJJaS4Qb/X1gSpORLN3FtxNUv
% 7obpJQhBW/Inqtts7iByx5r03mAQu5hpXMENne1AVOk+0Lne3sRN4Hv24d0Ha4BCjtz7NFVPp0M2
% 4MyCCNWIRtPpjvIO0Amo6GcBPLkwNhBHcRN5M0L7Ifo2RF9wp3gBF3h/Mt6WlBa9dyg+o/IsxoPM
% MwBM7GrTpgzgedvmxoKH5q6NJfCQ3N+SD8y/PwT/fdZSwNMX+lprFGr3Z63dJKLPKKIFLJZ3WkhV
% uj2ykBFE4qvCeRrQ32svcBbpyQe9rKIqVJh6B++keHT7C71BcVpO1jS4fOGnT3ulKhsSagRhCpHv
% yrnb+a3nkAmjEsKc7GBG4X1JR3HA9phPaGHYM6Gu5uM856RFKNL7e5W64NR82ymL42dxOupDSjKo
% 6/NodFIoqYiGnwsW9/O45J46qStbyKsTMOiETvlM3Dir3MWFDaa8LdtrtxudtuU3aI45eCc63dft
% QiIXzuwY0wvjEMsKo6l+2vwQ/mIFBo/wRyn9J1FwK+9qufiuVC9ny0YoXcScEWzTdSCZJgJMXk0j
% D81SMTp/JiCykQWVR/ZA3O3+pixymFBkD00z3OOu+UwYc/ipjv+k9WIz+KR3ychCNunW10f73l2u
% 8LzVFOS3mkjWbbfecRYH5BZA9dgxHXrY/5GyxstCI/A7W81pMxIy0LqRr4Q6sqTJNgcbQUe+jJU5
% 0Nq+sTZCrAbxsiE2VUu4MOIliK+ctGpd+8Tf7emTU7NVdDWxE9uAZzwphc9PgIdfi2RfG3QQ8WQf
% VCNYovyDaADX4p11TVM+goNVYtf08utyBLL/bPwhhQWnozfeSmVRFp7NgXsvHsO+7nXwsFIm4nzz
% aQBlJqh9XmzS6DCttmKFjAph8fg0WDvwN5NtPqyo2r9+Ot/l0H3UOHTAG/5xFNI2tVvoLBOaUYj6
% t3bwV7Frw7FiCckZuVvKEAXgUpDnIxj7lIhBlmYOvqrHa6fwTMv/ksLR10JQuqSgtKfOcYxGih7K
% sGPurlNhpFsqsbqPK8wahjAV6P+r88rtsn9F9rKpTpFqd0k7yMVmIB4kxUbgOvBW1Cl0goksWB6P
% t4m8XInP+er7g+oy3R72c7nGEEzHoxX+NbclyWCQQ2UA9PY9cCsLQ1c5T6r3NHiKXXK0rNIhhb/g
% d5XD+Kgdxl7lEMWmX9t8C0WhUzLI8afNvUr2svgz/Hyv8iI6H6mvObCuXSlXZsK9H+n6bmWOB+9H
% M5qGOY2+Lk5b+bEEdxPLtLSjThGhB51X4Trlmjr/pGYUznTRKSoocdULc3AqHvHZ0yAJJoEY4fV4
% oF7uR3IZrCGEwTOzcq4KR1dcuMBn+2EQGvVyk9xIztm6n58XGEELQq8x3taZnKujzJBH9U/3kb0g
% M8KspQW3ZfLnOTmKQ+lIXVSemub2VG+ozhYVHoSlkZhmkeeBwAih+bKYOw9LoWbIhyFHz8O/4p/3
% HOxf55lZ7QZpfbmVVTzlLRzKXZYBrD3wXWd7I6tbOr5DbTXMhQoDKloK4U6Ucy1z4EG9KfDMXr1H
% h4BGGEl7WD501U/HHNfFOeRf+SVfw0pP9TpFLcoOJhRgh9/ZgPLT7FlUOgvKmV6yn5tDNoV8M2Y3
% K89j/jN2+etv6vjqfBXjkSnmNycu1EPYUMOTt6BikYzjUkctC04FjjfEluIgLXKz3HDCL6Gm4EMO
% SWZaOGwhdD1sXn5lgjUhIhmJ+UeX/fYSfjlbhdj6w3P5WSlXyC3+BoVDYP5X3rzcTdVwKHTHALTT
% HAy/5rdfIaZmb7/yLRL6LDJEQ9611jy7/Qrs+DrkwA5AXkf6PoObYXtvbsOAh7D06d5ptMKDjTHu
% Rn2R7D+aRqd1aCb/+z/vVbbbeaNvm/j2bzDtMz+0Tq1TDlTyZtf6DvP+vBwAVO+1dck/8yYB3w0T
% qyuwve2NsGGbX7YFvt/e7vrT9AAcVNfLdkbIIPMqjS1WRu4ApHRpDsgHpWOes3fFsiNehfh5RZ8l
% dwEZGZ0XJ5HF7kbiY0422gQDVRuy6ojdAztSEiv4ogPe5r+B0BWPfuxL0AmZe7s8R3564MBstfK7
% rmlZXSGO9v490+8e7LL9v2mkkCSm206q7QdsL25C+6Kzv837lw9vfp9reXGLS8j5V+DhUOQlcM40
% bApSGg0emS7WolDlFQLjptB50DHxfanVVEO4Q2j+lg1+cvZDXE6TtgxLLMctfdTUJ/YvewWe3vab
% EWWM4P5xPx1+wsClnf0sMsTZznT3YNuJgscigrQ1+z5NDPL2GzbpKsRQQAL1b8YufkhmOw4kBe4a
% 2MpTcvIdR6AJmfqYmeQTU30U9iGQd5A+Ack3NzIif7jl44VSK2CHuTHa6nYCxeLyy66pCMg3PpGT
% GXPUp2UHacdZ3ifIPsixCvuCpJGjI/6adw74UZe/ed1rH4xrPSz/H92i5lAnW/Qsv9XKW/T4xQu9
% 4J03XTDWmsp8w51gn3cPDl8d1V72LNoge9Olj4Py8j9rDRBlerIQBzVm4umEfI34Kl4y3P9dyqSi
% c2X7JMDER/1diErxIOvWfEpE8TM0oyTHL/jFIcs/L+xrbUu/RIQqXmEFtS1rwInn7fjaMMbj69nX
% hq304754AJlS9pVPk6yCxXgALoLZcu3SBwX7Gkvhe65B8Vz5oLQNOyinsoOCjpRo/fCBnX87OfuN
% nb0/P/tlyJDjFcs79+wwhhXfgT9R3/AZifo0EstiE+y8zMVm0t/v3TyxWAKpdDIIQurVlQ/yuroS
% 4uXVFcrg6uplX8li7/8AUEsDBBQAAAAIANKpVV1k/qSjnwUAAIwPAAAyAAAAb3JpZ2luYWxfY2hl
% Y2tzL3JvdW5kMDcvMDVfbWF4L2V4YWN0X2FyaXRobWV0aWMucHmlV0tv20YQvvNXTAU0JV3KlVy0
% BQorlzZOcrGLJierirEkV9LK5JLYXcqS0R/fmX1QpB7JoYIA7WPmm/fsaDQavdux3MB71motmBwr
% ZkQtWQmtEaUwgmtQvNW8gKWqK/i7bmUBNym8k0btYfLLdRR9XnPNgSlh1hU3Ik+h4apikkuTwnvF
% qhQYMv316c9xI7a1AVUjuERkpjjkdSM8emTWzNDqlUvgJOAHDVuuxHJ/3eyv4aPRUJcFFIKtZK1R
% lLbIFRMSCiWQlCAjIQ1yWzPKPdSVMIYX13DPX7z+vw0hUAsEOAiK/nA6vaBFqCrkayZXQq7ArPt2
% Oo94dR3wrz2/jEajyFOgh1EbDaJqamXgzh8A03DnaVqZm7ouOxrFizbn7rJGfzJTq3BXtWUUPcIM
% 4kkKkyR6oOXULj/602kSRVHBl8CKImYpZMnvEeBHcdMqFDyfLOBHyPAHgzOfus104ZkkX8VsyDFm
% lnbMDlS6zc5BO4EEkSVBCyMqri/pcWX1gLFT5Moqkh6urHLd1SRIz2u5OVbyVEdZq+rmHBlcXcFN
% B4xrz9CoGh1v4p0e8rh4xOj6FHa65+D8GxzWdMfzEHiyksnnWA7p5/NHWGKQnygZFWYcR4rFmaNO
% UyFz0bCS/CpkIXKujxHZXCzmGweyIRBP505E/8SjrjtnMa05JhvWUGxh8Iuums1gcmB2OpVcIhcG
% uyc7eCWwThYXubw52DOeiqYTLzGTHYndFg3u54+UBPEUbm9BhnOCnsGD3ZGIiunngxRP6zDpI5CY
% SK4zYZ5yrFoTJ4fLJd3PUDgCITQRWpsfDwD0weTDDtby7rBzcIjSkF5QJhqn2htvwCaxzSuE6LsT
% IdYLO+P1hX87vhMq1BQJyQ1UfX6X+qoLZqRB1DBQeD3uysUGQe0x8JfiQJbYJDjo6oFcBExt8PFA
% W+xOt5Xu4ibtUaP4Fo8mXbz6wcK6Oo7XSrE90j/DF4if4e1bqrwuDNSYOd5aoi8Wu7vc4LkjsKFG
% G1ZmHSfYZ6YdTVNrfOO2BJFh940tzhvPlgyiKy5Hl6xE13r3+13nbnJZJ4eX+FTa9jqIhfUc+p5U
% DpVDOIPMjLELW0v7qZvA93AzVMcDkRRa9iT42JCWdp3CkMDHhmT0U8TShgzRxROnkaFLkFGYIT4g
% FlrJJHzK160C+9hrNw3gE0wP6CtXtT+3ya/xsC3pSd1f03t5JuG+4vtvFF2vgQkUpA2T2I93KcSC
% JpM79B0h7AghhCOFf06qa3Rfw7KsGVkwbmpkhi0rWz+/sMLNFxYFxxbrihwFi6XImeGjszr51MBC
% cu8YbvDIWZxRxczju3iHrQ01osV00ddW1S+uneLCyl2cq6a+P7zcbP68oK9v5N11QQXg7yaLftIV
% cNuv9V5W3DHM5SHpsC/4Yyb3scX2Te44bM/4CmPhyyQ5bX9nRV1swOIM6BDzkuhTyb6IXAvNKFzU
% Qb0dpx04CwG1U5HfYaYRDD1QP0Hhio1cj5thD/6s0BJXX1h7Vbw9fsf9LIGFS5hWJ8xjmzr7kBgp
% 7MmyV9HEbQovCWbO45GiRGcTZuvSp3XrMMns85LHvFjRvEIDMg3Qs/ta8qOnwJIk/SFBwtsZ/OxO
% KJX8fHMo4BSIicRx2VY0znpJPddTw44FhYSamjzkbvCmrxIkI1ZvPy27pyloDUIDKX4A726wtqYo
% A7trzDId+jBVGk2DvZPpmTB/JZdkQgpsKMXFecaTTrbo+6cYOifo269hN0qRlwo77PdH2tCeqT09
% WeNwKKSBG3/C0M0oNoOxcS4uTGVuEiJbGI3fc/zxzSkbYGSEkTmSc1BZHyrzGDKFKjR5hmlqCTN3
% lx/SB0Ereq3/7zuQd7XJfGzPAVaXAasTQHqJKYsQLSw34f2/cIuh8G1hEDj8t8woTFn0H1BLAwQU
% AAAACADSqVVdZNmmIIADAADzBgAALwAAAG9yaWdpbmFsX2NoZWNrcy9yb3VuZDA3LzA1X21heC92
% ZXJpZmljYXRpb24udHh0dVXLcts4ELzrK+YYV5kWH3ra8cFluZJa78EV5bBXiByRiEGAAUA52q/f
% BqiXnY2rXCyQ6Hl094yeP2Xj2RXdU5Et54vluJhky2JGnylbjvP0jp7xPcvDhWy+mOdZtsT9eVak
% k3k+xbVinKWjZ/ou9spYKg1vt7KUrL0j31jT1w1VXFtmmt1Sek35NWUZHovwxLlIw5s0L0Zrb2Xp
% 1T5Zraj4VRD/Em2n+JZ+3k/G02vq2N5nsyWqwUGZmsre7oTvLd8n83SajvPFdJbdjEZP/zw8fqfH
% r0+Pz2t6eVivn1YjCn9TImGlb1r2sqS/WDvWVIvO0WZPVnhptFAh9HCLpPaMFMpFOFihUvUO72jH
% 1vWOKmm59Eko5ovonZNC0w/2JCsQIL3kgMzSYgHkvlTsEBIhKqlr+petSd5Y1o1HnO2WbcAIz9Xh
% bkya5XTIguIclZZRJtBbY1u2ah/DEFf1kGoyyYhY9/gWIymxYYWnfzOJNcaThxJD5Akig+LSn+Oj
% +a88dCE1/+yFOvaA7tPjdWVKXBR6iFqxlVBB7hjkGMdV8ibUK5VsvYQTUMUh3fSI74RvxhsrdNnQ
% Fi+MlYh34ZyP9C2XAdqA4JgGFJoWQeRGor59UlvRNbEx2pheV0M+mgWQh8bAJEgS5DKtkPpdbfRJ
% gxVYqxU6pGbI3UcnuKtBAMTZ9koRmFdJaYyFekCemFLmje247xCDXtar3zsPtjt0tD+hAn3oXe5C
% B9EKUZ6P4CyHnlJX3LEOIWj1Mv62d0h1LlmE6WrDxA1ET94jLtRNgjeP7sCbj1BCMs11VDM5TRe5
% i8EcJLwoMwBnObykjY5e/OhlozXywxfH2fF9d3I3TGU0x97HJ49eWOrsj3eeGEiFhaBo0jXCMbXC
% vbpI61DiDjMMi/EfgK8J8tIAPTh9TyVOA+1pSueFgBYw7zWHaFG3SnQDRXCJF+C4bLh8PbRUXCBb
% 8cPYcOF/1mO8vQh5/mCsSrbBvXHsB765Ou0E1wjb0VZa5xMgQSv+ha6jmkln1F6bNs7VpaP+FmjD
% +cHQv+W5pTyN1dC5GgzIEH56dTNaXSwKfMKKjxNWnTb9cHVxN0xjGPTzr8DstIx2oaSS3c3oAWU4
% bBNh3w9lH1SJMp739R261RL8DJkrQ3FyrYFNfANllDowgZOx3N6MnhR2OxI6hlKVu6XiJk2L0X9Q
% SwMEFAAAAAgA0qlVXRtdBhXJGwAAWlQAACgAAABvcmlnaW5hbF9jaGVja3Mvcm91bmQwNy8wNV9t
% YXgvdmVyaWZ5LnB5xTxrc9tGkt/5K2aV2l2AGkoEqIdNh6lzYtmJs6WkHFftVXhcFkiORFAgAAHg
% S17/9+vueWDwoOzcpupUtkQMpnt6enr6NT385i/nmzw7n4XxuYi3LD0UyyQedE5OTj4km3jBrtlN
% XGQH1r8cshm2iAVnYh/MC7YIg/s4yYtwnrNNHsb3rFgKlhdBvAiyBYvCWRZkh7NO53UUsXweREGW
% syATLIwLcS+ynLO3GSAKkxg+Jhl7F2zyPAziXhZgYxCxNAiz/Iyx24TdRQm0xvcdEd4DpUG0EZxF
% yX2QhcVyTfB5sE4jpEPsiyxIk4jQsBBGZWmWJHfwW6zDXADGH2AuRd4JNouwYHdhHBZA10LERViE
% AsiJk4KmE0RRL8kWIsOnBMAB9jUNKWDibCmiVMhpdeaZAGRiAfOTHJoGpuNZegDAj5kQ5/NokxeA
% b1OEEQ3GgkWQytFWCfAmOpSo5CJcKfaxdTBfhrHIDq+o+4+C+HWeBsWyshzI5VjsgLxMnOFidu6y
% ZM3uNLtZuE6TrDD8Z0HO3so+MGxWJElk+syTNUhHoJYJZrveFOYpSxabeSEh10iFAroDvEkWBhEn
% eNkhA8nAEWSXD/Qk3xThWuh2GOFuOsflEZkiexPPKyRF2WY6B1YI+b7ObN3N6TD4+Z2zXzj7iQOb
% QXTzzQwWV9xzGjNH8uIVLne29uFBTYizWRTED5wQ4JSni1TOfZodcpHB53wxpXF5x+10MhC6kZqR
% c92/9Pv+ldf3+26HJpLDy0+fOwtxxwoQqIPzIA44Mrwaee5QfsrH0DqBnvLp7F4Usl/fZaeykTCg
% KAjnibM9QGai2GQxexr3J6zL9hw+efSJupJAO3HZbzx2QMCckI1GbOUS6jvYOCsUWVidewGdJ9QU
% VpoIG/LG2efukLiSbAqg9a3jufSIMHuE2edDetcdARH4Rg0NbYQlE0HkPCkkQQ7MLCTRQFLfBnjr
% 4KxcObTInEDBpDCsWhJos9GkiKNcJP1WsvwkBN2VihhZwt78ev4Bu1DvIMam4B425xo+5SeuTQbR
% m0oy7lBMkBD9Mt+sHYdkB6ZEPMiSHXIhoIcnYmKyA06/dfquRALbBCTRCUDGYA1HyEE1M7NI+B4W
% WK70DldaIoc2wv4Upg7IYK7WCj/qZkSr1msJemUNmli21YYg6cdBdv8X1LDtZtMwjpHJDeTEMeDM
% 3JGD4A4D9pihGjjbCVDsIkQ1qfvdiNyTLXIj3OIOfIJt4NZlr/MN+yBQHxUZqJMAlSvodVBsYNzS
% ZAeyAPIA2vhMChxigi3/yNltjXVOOg4nLITdxb5lkYhBOJiIciH3qfNYe/to3hIW/VPbY7cA6ul9
% Jtc9rezw8R529YHADgiWqr7rTVSjU3JijAKHmkChNhwLudynIt6sBRhaGGg8vD31Jgpc9wOleKj2
% e6R+vRB64iDj8HQ1YaewzZGytq2ehvEWibut7vYUtnWVUs9l59QM+ErKbw3NDyWjPEBHEyqpRVoe
% UHH2aDum4xVCU2tvNakrOECwDmPn4RQ+yNWDH7U/FRWtcxH79MhcSs1lTeg/mwlNZAWQ/9lkHhqb
% AKfyM8ikHHMDQz1CP8djPeZ3H7uPUk5uoN3T7Rdl+xraN91N1/F7G/fc8eB3t+vbg9yYVoD0YHo3
% XQMAjeeOen3RhQ9r2OJEEWgIdBuL3Kly92fnLU7xCiYzQlsz8F5ev3jJ2eDCezm4cmGHwXt49vtu
% C5jnazjv+sW173kvEde1N+hfXPuXEnoA3arAV936qC8vYbBr6Hjdf3n1AuGujqyn17dX85Hs4wNg
% 6vfdV2ASRpLHryTfkVNBqRSe4a1F3QZG9xj4GdhX9oKWgFrW+K4OgIvNvh3BtB6pk34uJ2kglJ00
% zjesylZk4GzPBYGSmxpuRU8vFwMPbP6g7aUaUO1mR8mUGgWm4uKSq8de5YVepSvwwGCVcEZ+FSUu
% 0zX2RfboR+KSeUCRIzT+NUrIYOBLsm5RliW6JBbN7dm9NVuBlOm4D/gmXCqvMdAIzz1/gjsf/5v9
% UXl/MTHvoFm/JSMCg3Kt0jfQ1XMrqKCfr4eGz5zRL/Wy2CVT2N+bfErkkU2SmAir77ZgJsCtRmlh
% INS+6bLWXegXvN1IusrR35gxscdNCW9GJSrscb+vIC1/3Rj06g+1btVo9L9iHY/wb13hnxn1IdeU
% vuHse9OspOchHw8HpKRxaUnEJlKeo8jZs+9AeZfe60Ne8RnNXlgHqyRDkf8YHCLoPE/E3V04D8lf
% hBFPlYVNM/SvT352vHPQTyN4VW40zk5gg7489/uvGHYA/WT3QHVFXQbnXv+kiq5tWIgAs2RzD4Gf
% QNeVXQ0BF/47wyjSWQepkxfgSeH8ryfoToELpENscHmiQ5ysIUJjK1HkwypqCJU3/+rzzb88+O/L
% aBhi4HAboAIAH+n97zcffgGuOxBd4T+38/6XW9wWzi+6ARX7SrtRpSNTbFJYSWyHaOUgfb+98jXQ
% 96PeCrrNEZLwVhBU+oHKS19pn6jquCnA0i9F2wq9yLA24iAwqcrptlT8QBPWFgohA9pDoZEkiRzT
% vdvmXawgDqF4hrOlQhoDGJr0AEzHAmOeMfF8AlqUffsti6mZvI/qyOsgt42S7FsaJQj9qMvZLCxk
% lO2UFoDcVtgoMYpDXIDeEFWf0A4Kh5UdC6AoJA4N/9130PVvFddG/wDN1OXfTNK2cnEGJCeNvkf6
% c8lO9QYenQB8blrMpf7wu+tWGA2de96k5DV6tce4rZYTuW24mpKvbic/HM2Hqv+mJ0OLLeVEkpeG
% 8FuTqB9+d5vBttsqIfNNBptvkwmbati5OJ4tPB21GHIyZPmHdZ+AwEZ1RtT9ADteXoSZmEtV2ENo
% HRRQ6qgeOGOmBPRZ6i+D6I7EDSI6Cr0AtBawfwfKF3WxDt11zK5m73cllvMUnInU66beuZN2YZ93
% fv3x9W83v+G2AIXj9NAOu5TlcfrKSDjkGlySO3yBH1R4RVlCYDTEmGmQ5WL0NoD4jMOcFimm3+Sz
% dkRhBJ1FecUyHI/cBlq2aTNtUm3qtER6cT3OKl+F6MbXN9ZeZlnAgnSduIdKCVZX0ypDy7dOFt+f
% YXYNrQV0vYBJ+xeqf32bylkT12HElcv+ygYUwgzVWLVIVW0ujEpI6eJY82USzoUjF8E1Wq0EWU0w
% DB5Rdk1vT2SgwZWZLjJL9EiCspdmawOGPnMrovKIoqwsqZQYnYCriQzIMhi9R1pqI7aG60sYhlJ7
% tJ7F/9t6tizYZZWHyz/O9mWV7UvD9sLgKmpsV0xbclboUExmRQoxXcrsskm8nZyc/GTpBHArk+x8
% nshUL/D67k5kmEBXWfeYCQgfDpQOstCeYTa6qnLp8b9MUte5BZ9PNlK+DMdxsmRHydoot/iZ0YYM
% m0siLRkaJJmAyycGZj5HmFXTniHMSsPgQCWM4pLOOo7Hek8YJPO5nTbNJpL8WZCH6JuO/xTJkdhO
% R5TF5Zig+QXVHPwN8SOYE93+U9kei3vnJ1B9E5M7pfScZKqD/cDgvKpnUNF0SLX8nUpt3IVZTtul
% nt174mqFYGw5rKscJ0kLeWYz1TSjB5yHpOddjTngCM6qSSfqbC95sithDBxm05+Hw59cZlY05VYK
% VM9B+YhJVogFTuiB0inV1hWPMMvitrorauIrLqctHUgecUptzqVY/mUEPiVyd4UfoyoimN9ZkOIe
% cySnXbJ9xP5xuJ901afD5DzF2NkAvzNgya6iOTHMeacVwIi9UxqgJoy4E9+5DTdYNVc07Dsul1Aq
% DKUlpjIXoBM4iCYu0UAI8bIm7fMgF7ZnXV0qpcVxwermmgBhIp5lsnWb30yyWhrcb/FHQfNNYRTb
% UjT6oO+FSrXqgr2i1II8Elg2gRT3Ha+HKRe3e4sGjDDBX8c7Va3H4NDU+f0aGHg0L2VrAwwkK4gP
% TlAu9O8kY9qM/AUbmnqyXe+08MlyDA2vcjYHGcVDWcQDGy86sCeRJUws7oX2BlsQ6ONj41Pqg0wg
% RDxuAnkseuIekaSrZyXJ/3MkqeYEaZkn29piISEYU+bsXRUwoqMF0tl6E/aY43e1XJiYa6WPCXh9
% 7Ko8W6u1aJ7VLdxJBXYDOkENT4OeHhkUiHqnY6c/cXhtVoyvRuyoO3BEZHVQykbIZJ3tKvV8mfUa
% NBy2+si2o9fQCdrMYL6BRlKr2yL3VX0+tNwyjPJMvGnlI9w2DtR0B/Kecvw0OiDpKu53ZcNqcmQJ
% jpocmxPNhTIHkC2kaQ1G+rNLVO27+zKLQRS1IlDb+W4TReRV9MqtYXY0Lfi5FMNff3vD5jBkeBfi
% ARwmzvxWjMeQLcBix7nUPBKPWACSxTFNcW0t6HOKvrFIOuTTa9XQ8jpGlhUjB0Mh5RUB/xYae6gL
% e1kCHmVl1pSIw5oYXXYSz5N1Cnp0hgUhB3afBekSnN1dWCwhAM03sxxC7idRBKyAqeWoa88gAP6p
% rCoAfz+GaYF7AtqNtK+aeS4E+tmf+p/pcbcMI8E+ZhtRMibeo19HHf/NPq20xJM7hZiksoAn7PIZ
% +4StfVamj51PIuwSvUnhoZrEBmJtbHorWgGC5kRseN6t8KrGQGB6PwrzwrFrVkyiBqVNWW1KJSEZ
% QpRrugVDhK4lxEuZEHnTzaym1m5r2kIyAoCEOj8VVY8UhkKGVBJkk5rSi2XsVl/LSi8gU7t6DsT3
% 3S7Ng/ohfkTSduCsXqGAWgC4AD1vKGes0SKJJXC6xJwhjDoeTo7FK4qVNY58mWeKsApPhjiiTO1h
% sKMf2L/KDOAfWJUi2OjzVOr4N1bQrAuCpViP5t6qGvvkqxGTKZfSLSmDF4C60+rdLINNjnny+s7u
% kUjTeLKWzvaRYD+PY7QrMIY8cV6GUxl5wG+wMKLQm5omos+cKaziEKus3Kq8tewBXCZABubAt9cK
% k0KAfQyvMC0JPDKPq8lkKMf7t9o1MzsUQKI10onkjMqxHq1L0Flsc2QidclX1iR8TUlCmYGlZ5Uu
% 3vNDM7MqlacsyJvaJyCO0Z9UvDUGBcpBi3IP/3ify5CV/F2gaayyiKhbJyXVTwjsOMqbl0lKerom
% i+pxz3N1kzegDvWjMITueXzQVxAXfQPxUmPsK2TXbis4DH5ZDg7IXtikeC8VZn/gutbmArHA2Q1V
% eRad9eDHMR6P0gd/UsovnhZevqJjRZSIp/LNU+X04g2s+ytoqx1dHDm+4FLo1vXThkXemvlZV1Rs
% 2FSxKkwi+Q6NpKtMjrJolY0DI9FmaR6IyKmR+EOn8hjDbtNHAE+ViBtfWexBlgHQEAzP01qxxRw+
% yBjivoWFqY4skIeTU/32TYWdVV4iYJWRGonasPgEa7uGrk2v2dL3TUSaznG5AdUjN+d4iJ0a3irL
% RadsnD3Y2ZR5Y6rHJ9QkIxewTgtTImZDgo6iMxxZ4elogYHwUKQQxY4eWmaULqVCbijiNiuGqcN0
% eURS5NRkdQ9xhx4Mb9rEBJUq8ioF19XU0zoPrbmnb9gb8HsgEt+E+VIsGHqdOZ7Yg3JbYIAuJwlv
% omAG7gH5ATsR3i8LclMlwFkDL8TM6FTJVJkKcsgw9Ce4/WU0pqtmiX4UYHsKY284aSHYLBTyQj7o
% KArHbJk3nmW4rYx/wEN8/0jOYpcY+lXCz9Av9RjRSE911dny8/xU/dapfmG6QGHbbP3W2SofI05i
% WjU76Q5ra1zHni4yJ7mtp2KUdyOHJzlUSyk9dSSuzbFJYkHRzDlQLMOasiSgJN2qoj9pRliUYEbF
% UI+kNMFbkeWbXIVpPehZlizgwaOFHKwX7v/uQBWRpdkGeDFF7wo4/Fjmz1V9BGgVb9KNq0f00HGo
% Ooz35G/KCirjzE+a/Uu9LALMUcRijyXV0oIsqg6KRE1WaUGpppIlNITx5BEX32N8oqjBhgnrAcgr
% i8CeJhDcGcsCfnHYSac5ZMMjoteySh2gwR+KYEXjexkrAD85C6J0iSW7EIty9h4z5+qYdLGiJNPz
% h2ZWwDhEENhCmpgVnkFAy8q0hGWZEprTXXkgBJKBpUNpgIJvKW0IUVqCtrntHTRqFpwtJftIHc4x
% SUozdPE2iXzlla9w1i02AofVRNdiLozDngfwkIgtkgg+nKdKhXfj+aRq9eYQt9MSI4+2k8Y05Hs8
% uyGmHLFB1EtxCZlIz3zbCC+Ji/gX6KCE0Hv4MF5MugYBPBBdC5u5bpPfjSM0BLZFDrEjKX1OR9Wy
% 9LpYTo06MZ44hu0X0r8cUMN7kjhwX3l/wukv6PNxn/6CK0XHRODT8oEk4qM66/Vc8JVU5V9dWteW
% tE7bY0q8VfNHZCw3tdhWyGPDLo7DqsVV3B8yU+aWm4q4j/ACtiH8LsvhLFK1mFH1Mpir0xwCJqpV
% LLt+JBW2o+d/EpOgXxM9AjYJVseKePOHkswwVUWuxWa3Be6YrETBQSWs0SOWtTfqEHdiX7DAIb92
% scb2ArSs+TEx+Irl/YoVxC7i+S5moYVcaKBbMlEvuGlQlZPEJWwR0PJP+n3bFAC9NFVcZQtnFl4L
% Vq+BlgoQHvIq5LKYos6PAMpJQmBsC+sfWu+aT39b8egxHfZw6r9CNZEr0eyuv7RcYKObjn7scuXn
% xz3/SOau7kG0nBzMv7zWVXEg0tUCgHfRZlhjPucL/p7vWnPhxqgr1x3+Gv+LMlcn7p8orOb0sSSb
% V2IPNDJaHuR7cB+fP+hD0TifwZjzpeUmWtmWqssob4mtgwJl7DX/vn5zSZ1QHNRJjikcBcXAgbOu
% y+X1B2JGEunX3e/digJ5LTNV88McHGRTi4EdIvsMIa4IJL7e5o0kgRGxqIXjadooHdzmFBUN2SEU
% YLfxGWyYe5qmRNMuiB6+dJ5+XaNKVXWVmahaIkoehX7glYNQdQ4q/9RqifjH2mnJK7okeezI+x/y
% MKI++Wd4aSDHEWWXlJKOufnQPHEyAl5btkoWQnmM86YylPWS5AQQcvo1ad3Flk6Kjmlsjk7vHMP4
% +ZhSGX+N2s/t5LikNOEjH39Qx4Yf5d8J95u07nZ8seCzGQAihzpHT/3mQ+iLiQNQL4AdJAkChoXV
% 4E2O8VyPcep3saN/dCNLdlvpBDrakvmDRhBKfU+aCxdxRw9Iqep/nIWFWBsRb1FDskL/2xF77Haj
% Mgjb7Y6WWzSPLxcLrRAAT9SNuoDL8bsRWIEuCvIxTMB5TP5HmKhCkOMAWG1Lh6dDCxYPVto6UxIs
% 114JXcYdjVbuV97DbReDaUVW1RCmHEhpUtmKVdD8Q0tiAW+PUymGcVfjzhfO/4/tC7wEoketjY4y
% D76CfADWYuV9e9m3Rc7+9HvcJzo1Vj8q0F1bNnJdjBo2A4TRgLckW7JgLqYgUm1lsW0M6Xl/gCX6
% 70dumONWuOMf5U5J2Kk626pwqFJN3qJFG1Je4vtKka/Y9yiBWIHCZHRMrJzQPEpyseihOWucgKOd
% i5J7eQyGV0EKMJD5yL+o3gDcmwJE2GV7vOBmpV8STkeTJRrvfM8Jj1uPO3vLkPeipKPPZJw9TPLc
% 2euILEqYqX54IgWxglfnb9WH+krLQWzeFkEYEQoJTR0Qg/PWesKqn6fuUzXrAtOIklOEl0dnIs5F
% PAV2TvEcw6hHqWxVOAlxLQ7u8YH8c2V500fdhACL2cd14/8wGnmuzmRW7rqU+2ktgritngZBdjzQ
% u0mRyAO8Ia9g1G13fOKB7eJuA3AD8WqhkuBdV/oV5no8gfT0LZz2YerSDaJAHCpFQhazShyEIQis
% eUV8vaxIEHbHgSs45RGJonKvebWTisimCIfXNPXWy069Csug8L4SRVRPoGqugWRe9KVJI/LggxxE
% VwnC267uXM+1Wt+U8Z6kjd0Hac5mB2buu5kvVqEva8kwbWNnc3lWdSRN+Rx6ndbRljR0zdL1Z8vt
% FSvKZZope2l/nQAoO5c/jVtubrVbzvKW2JON+bhUz1b1pZnlINazUqZn0MeWaPXFFFQNGGA/6wsY
% ZnhgR43UY2algWOR3R+0cKgLEnTT1evJb5aYqTr+lrtCX1eoh1WdTUuw3dbN2nMTrwYSq9rsVlYJ
% Oc5vRXFiWk0sbtFa1bf5igczXrmmoiV9KwvSJH+eUfrpeWrzs153q/hLOwV/A7RdgFtDb+594je5
% 9NIlVoOSGib7Ju2d0VytJx3fsNcMvwwpEkyCJ9ksLNg93ptkqrr1IG9UggMjogPDsrE0zZJ9uMZi
% N3Stz+rTQa7aDMeA0a0x2fJ2lOBgLyU8Mpf9xQLguO0+ePzQw6vTcjp6ClR2a0x5vgyyNBZ5Tobr
% DowZltGZA5hMJmP7rTkfn3uXLp7CZOZoAg/7ug60oL96Sn99fbRWs3EDgq5sQjCQGJ6DXnpwuwj8
% MOnKA+eHxpVOjEknDW6PhxfyxjDQzBUusLg9v6sebOv4DfuhdncWj77Zr+7f/05fxOSzX6c+yc8V
% fBqwoFAHOjmtNcd76Y0F97spxGMURsjao04jRXMFXQYTrt7r22EA0UOfw8E8F69A4hk1Cv+L6m7O
% 1A2zq7YLY4HM69KtfiwEhGBFHlL0Mv7V3x7UVhIsL5noOlg8hOxmsEBt61PXXEowSebkzQj1FV3w
% H4HwMmLPuuBc8z3LOfv8EtMa5kjhT5vsUiPq1xBh9Y37x3A1aqjV3Ry/m8kvPLjg9GUWnt3b65nX
% Hr5GAO9UtVxzddesYrGsG0D4XQwvueerbirUSpBpzlXPA9RuF2+lgAhmPfiPz93KVzZYIFIeX/j8
% SuNrqYmWZbYGiu62KNDr/mWf+y8ur7zK1fxY3Afy2yg0GpYXWTgvokPvzRulqq2Vr96o/83qOtgP
% sDu6MkP2OLo4v6Sv+hoBD86BB/QtcyWxox4SdE4EnWn1J/YFFiZvxVQhqiTu1pSWu+BX/AX3rvgA
% WHFhpea2SOScMtBUvNqoUFGJa9/aAjdUsmvf3bYiW42wus23FPbwQzMxQ/nQsUPR0F/X0INDP+dA
% T+6kGdHenKHprt6K2oInZBcsgIbDjOB2KL9cYWvT9bmaT5WXo27oHH68lYfwdOY93lkn8pbgYHbv
% RtWrr7trEm5kBoCc0RcDAvfp9aeLz021ibdMKz05BJ16q1y16Fk5HPTC+1ayI4wKIvmioZKMHPRA
% iLGMIVkHmKKzFBBzqE7IfO+YQDpkXtjV4oRApgKxCLJC7lPztXhq2a2vzaHnasW07NJa22gO7Srn
% rDLMtHPPUpfVLnipm+f1KFUS2+4GqFOk5h6xt+T/xDf//fqHj+yHH29++Pk39uvr3367eWNdAnoQ
% JNxK4OSX5GG2i6DvTj5ZX6I3vFp8ZuwTfPxc3fb/CLJ7kReMriw0bhUMmd+vX1eF9ZK25dI9q+J6
% Y91fgn5g+WmGC/O9HBLuhSwnJ3aX39lxZQ5z9KaoYS+/M7MqPRt0wUi3lVHcK32TQJKxSGQpWpZs
% xZHvsKyMdXdyE0H8Jxaq9icfsk9VaeuREA7PBnfIzw5YtOk0DtZiOsXtcDKdosBOpye6PBmlt/O/
% UEsDBBQAAAAIANKpVV0N9VmXcgIAAEEEAAAxAAAAb3JpZ2luYWxfY2hlY2tzL3JvdW5kMDcvMDZf
% dWx0cmEvdmVyaWZpY2F0aW9uLnR4dF1TXW/bMAx8z6/gYwvYzofTFE2Rh67LNmBFUbR9HqBYtCPM
% FjVKTuMN+++jbLdNB/jBNqW7I4+3PaoigDaqsuSDKXwCD13YkwUflNWKNdRmx4q7a/CIGi5nq8Vs
% scomN3UNjbHEHqiEsEeYL2DXQQ5MFEAKjZITKrA5rmF5sYSHm6enyWes0CKrgODI2ADEpjJW1Wso
% qLVhM0+AsVZH4TpQ3Ta4mU8XCaCuUE4wYxEM2fjzWkisN42rEQ7Ioh49OEaPNmSTR1GR+s4HbBLw
% 7e71NbCyvhZ+HYl0K03akIA0C/pdW4XUYOAOCuXRr2GxHPXf5FCaHUrX8Qb241MVIzbCCi8m7OFR
% +tCwAvkhALP8RDbI82eWZFmWXPz9kQvufDXgfqLjIK1EHmDIOfImYJTpMAqe/kYm+IlssZYpvwwa
% 3poYh91X1pDPBuAvRkaZesEObQPOYIEvxmOqytJYEzqx9VeLtui7HO88jzOKio09KDZKDkSfe4+S
% 0ZlhBB+NWQ8I92TT1pqGdFsrBouVoB3iQRkKvZmdv0Jt8mU+vVolYFU89r/b6fJKqnDWWhEra6W8
% jN5JEfW5WK1shamjugvk5LIVQB7UR4W9X2C0+BH7PdObRZYJVdjMsiw/f29727vJEe3UswhsqTHC
% WxDK2AojUH4Emi/Px54/JqKHSQfu+8GauJ/xUtKT5vkpq5a20q+smnehJTGgLHYnHmhZARsLfbrS
% MV1OGX4lv7uD22/b2+9P/XcG21o5LyFaZpdX4rEMXvts8ryXeMAubqjUTnIvEXESpLjFYxhI4hXi
% 7tWqkCWMCW+UbX3BxgXJGVEpgP8AUEsDBBQAAAAIANKpVV1K662NWhAAAPQsAAAqAAAAb3JpZ2lu
% YWxfY2hlY2tzL3JvdW5kMDcvMDZfdWx0cmEvdmVyaWZ5LnB5rRr9b9s29nf/FeyKXiSbcmw5TQEn
% LOA2S7dlLYZuKDYIPoOWaFuJLKmS7NgZ9r/feyQlUZLT7Q6Xpo5EPj7yfX/QL1+c7/LsfBnG5yLe
% k/RYbJJ40vvuu+++P3C/IKswDgtBgpCv4yQvQj8nqyQjn5NdHJA3RMRFdiSjyyH5Ra4kE5IXPA54
% FpAoXGYcZpM4Og57nxKyihJehPGapEkYFyTMyS4XwZD8thFky+Nd7mdhWpzlgCLb+cUu4xFJsyRZ
% 5YRngoRxIFIBH3ExxAP2VlmyJasMjhkmcU7CbZpkBbnVA4Tn5FbBAAVZkSRRBeMnW6CYy3UU9whg
% P41vF/tN2G26KJLFgzgqgC0vNuXc2g+oREbJCnZNspBHlET+VoFmwAncXgF/lm9qpgi3ohzfJnFS
% JHHo93q9QKxIkBQWp0t7SjIBTIhJvttah/5RMv5Aj8AI8hSmEsaWK/Ldsrmi2KWRsA7ON9b4WZLn
% zVUW98bz/tJz5w6HD3gaz6l+GuHYSM1S/TTGsbGaVUgDIY9O/RqpokbthhMKEFQhyaxDDXYYxrut
% yDjw8Pz8MAQhJ1uUUJKp04owMsEdhcA52L1yY5HhgriwsuQxt6c9Aj+ced4trFNsQCbA5Fy+wYN+
% z+dXJGaRiC1uX5FkV7Bba2zL9Qh4L8F4vBZWrNHiT8picSgsK5RAYQ10T2ObhCvCvXDu3c9t+imJ
% hV2tg5kUVR9HK3JurVET4gW7nwKGFHl9P2f4QfFVnc+Bjwp8L2fhV8712b6awYM9NA+GSAH4YX7e
% gjMJGIypSWm5iSSnMfrMBiFu4LB9X+8l12hK8eQ9P+J5Tj6IZCvAfaidUIqLBXqbxcLKRbSiM+MM
% ODCcMaXX6pMrqXLcfwaCkyBLnoucefMGbT8iiGnyljovinxm23TSIlZqMLNmQEjNmx9hixtWK3hj
% BYjspomjOrQ80ZCn6Lks60d6Q63aGLRd+JSXj9JAbbtiidjzaMcLoViyof5G+A/st2wnjEPPmOLP
% FfnKpL4PQIOpa6j9BrRjD04w9IE9ORBhNzkEx3qPcPWR2/JXTEdX9NX7EeQ6779X0m0bycS2z2/a
% qjGxOwzjUWQp57C3r9lSCZMuS081o19t0Kby0EMeBNa+RvKF5eA7BYxpABAPKBW8EBTrF/stu6iA
% fQgZIjNI2MPR5YZ73O6Lfa7WnCClZhP3RYGqhZom4hNMDKnURaEdmbBmLb0Kg1yf4V7tRPfNBV+k
% 4yiZwthX0MAO3/CogMm+nhDcNMiV3EQ8BR2PIb7uRHsJxBfC46NWPIwWX9ChfPFgNfhu8FFq7MEc
% s18wa0Thn91RbIPoe/rQsS7EMJ7OQQHtZ86Ex5USRUKuNG8rIwkpDtfb7pMIOMQaTlIxXFOv1reM
% mKH9gk/choVy1fdd8So/DU4alOeAjuMZmPsXDPE0vZ9vKNQXbz9v6BRScF6KqkaKuzT5iUYOuQXo
% dh2HJTYERUE4PnzSemgsh8bzLhoQQgqJkrWnj3bXG0EWwNRGV+R3+od8fuzKdrNnIyTZOr4doYId
% GbxD5kIOb9nIJiLKBRl3Vz3qVX+oVX+Uq37/1ipYsNm/YJvHaRnWx2rsevP43Jonduj/4Rz7v3dm
% Gjie4BwKg3y71m+jjkFqNwKPFDjC6mzP0uy0gV/7EasMoxMBDW0Z01LkzviECJQiDxhf5hZGErS6
% /Ygq91QaIW4D7vXU6GBcj4OXvawdXJJlQma90kquiAjW3TgIhrXPbQpxGhKkjtUqIwKjnT7HI3R5
% gOBf+Pcxt5/3TO70pNUb/hnBGHPbYVfaLJ3JZCZldUS8ImsGubaF/5F5KWacVD2BJZSP7hyFlWq7
% PJyfr+sImDaP+6B8wkM3UMn8C1whYIKTrouNFJcWAQgAc5tKGRDuXEI3XUMiVisQRZnKU/fC7lsY
% lNUIt1WIxreljPit5aU4B0ydoS8xarlWjlLRibykIb2nCpRKUAPjEms1UAZFoMykIVPuxECbqpx6
% yw/ducH4RHCcG/q3iwuGjrCO6k+no/rGVpiepFBU5WX11RntRhIsU53pKQ3CTSy3f7DNOoGxsXHk
% WvB7+xSKW0ueGZRQG2XNc4qnpHKaqklaT9bI8nAN++4yUYYBZSRaJs9lBCcCfIPRvZY3e+YYtNqd
% fqHacqVu6BoSquZiESfZlkfgarRJC+bp00GVFDJ2b7dT/8DuiDmwG/m7GGgczr4+t5mKi/nAQ5cl
% 0JIFxtbuJt09kC0hFD1zdXwp+gUINsnyKo3aUlWlzWzpZWfoA65IIXJIy0a9E/XImKKewzJU3+k/
% Kgm2Nn1oOUAE/+kZ8LgLbuiYWZR63kwVUDXtP5nVxVw6ZMsZQ8Y17kZlSeWAjU1RyDEt7y0HSvVJ
% 8oJnBau6ChYwKYvXTLUfrDejS3fkXrY4N2MNjZnA9AdWFmjAcQmUZqg4Z6oxZHSEaNn8abd+ZKoc
% EL3l8KyBZhZFREmYJCtSbCBUu2R5JBOpvEQdBbstWXiYntGWStCzX2a//qoxviS3O8Am1634NoyO
% FLjpR7ug7jYlWbgGUUS5TEtikGG4BT0GZkaQVgV1oVFpUlFr0htDyDHd04wuWnbHPgyrUs0r5n1d
% WwIHIAwUZUHWa2lIzJhVgHL2+xdOAR/ybHvGLvp5vz+Rb+DXctMrFpBZnXaKkOCjB77X6aekYdFM
% ke23k7btqcqnFYC0iG7EWsTSc7WYONUe/xXIvmKgytJxCDli5iSv8iuD4yWnYReRYzPv7JUlmWoc
% oqHvbbO+HBniaNR2oOdDbLopQGdCL5SvWZyu6fKIQzHtNVe9NvsKdYjbME+3slT+Nch1bMvr2Cbx
% GWHR0IlNaXAlUS/Jr7tlfswLsc3Rorn0MEkWoIsROmhLDdjyI2QifI+6GsAouBOl6mEsvu54FBah
% qbYGp16bjHoHydX0ErwzN+jzLqdz1KomC1yDhh9qN/AOaPgWu106+Sa7Oyx8nvHvjAN8k4nfH9Io
% 9EN0tpWu5mK9hR3y8zTisZDyWSZQOlFSwE4gJIjTkrOPYbEhwQ4xwJjmYcUnDDAUPyAZbXChYrVK
% uHIzkymppplIBS+Y2VvasCjMC0uvsgeeLK9hJz2C+9WPA/ef8UDb6mfQCEfpE8WGcPlYk0xr/aGS
% fINla00d8WXv54zKLVpOdjZRZi/9SSaEoa852YEFP5Uc9FQgG0GxMaL4KB/wER5UiMOBMVXv+FcR
% 29I2I9/DJPzGLpO8Bb2hUst+UN0qdUQ+MRIB2abriubypGiUU+8vmMFpJR4K6kqXRn4fhOjWZLKr
% zgKUA+ndnZbg2Ou90Mp2g6drxpWT3+OzL593A/W2PBUf1Hb1BBzHLwsyNCXDvpIoWHxmlnMLW7v0
% Um65vOZclb4wPOkvHbfP+WCCBYma9n053c44bq3XCOj7zkUfCKnhmRw1UAJIibFzfghe8lAyRGfJ
% k4j1zdFleXM0gWgBTkxe2IDZDGsMk5aCg/atwqXIlFMUMgnha1BDtHVlym3cRgwi8PvniA6HQ/r6
% r39PQMP5pKXenyCHCIujU2xQt6XHIPKyhpMHkcUiUpsoiaI5JWma5GEhlDU9iSyR1wmKhjswgtoG
% ahNAjQeLkL/qUcHUsPhXWcNH5nnSOqicBV9BPf2kRuDdKQfkEAyMaP1vrNsVHyvLZx8H3kfwLt50
% Tgle47xREIaloWbXC1TW2r873ex9Y9vY6Wll+Jf2c/1fY5tmdtze8H9Kk7+V2YPd62TguWT+DUIo
% nJCsoDghNUgezduiO9kwhEQL3m0IC++YhqyqokPjqsk2F2vQ/4eja3PXjPOHUTujaQXmOzh48XdA
% 3Rr/wLx0gPEbyKHFSMGmVBOHCdBhRO+MBRHwb+80d5kY1e4BzrEJAWbwLZie2Yjyyrpa71pX1nd2
% QzRmzrEJQaWcw6hxnQPYIPDKQSdKOlNGYGgGhc3JQIByk1iumUQp+XSHFXBhXzN5gLaNYN591867
% DT0uGv0R3byReyhSqEFW2Z0JG5Kz/yacaL/6LjmoJGElMuXhar9W+rpz6dy0F0Qfp0qTKh1VVZqc
% mZLJiDTThtUKMlXTF+Mu2DaIyS72N3jeoDJIeYeO3wDgxW6rXCmXCIzQXkiYilbXVH9FIFYdLQUf
% 09eGfhuZPd6oemb59tA3rLjCN7e7/YXXhqbpBuifiTdRwk4QEpD/hcbcLNkyUOXEc5tw81NVXZbj
% dbzbz7D3OMgwMSxb6/Ay0VOunBrjVBOJ5FwZSEsnIwdNFbgNoYBzNM9JGgpfPIa5cCQkhETIqL/u
% ROzLzFAtb8VO1RMoc035dYx8Aza64dFqMNtrRYaAinV+zrdCl4tKwJuurC4aslLX3h86ZphvwlXB
% LJdCqTfGYhvfRdAoxpcDVXTIOftUQ3LeiE2wkzedACs1MnwxmfVbTSQg2XPQReAM9jBUr05TprKU
% ZiE8bRrGpyR2dnG4BdZEPEP8W6jloM5TuftUMkuigENvAYbE2E7B27wNpFX/MNnGvoMioEpyjdDT
% SXfHc33HfMshxWtwJn7BwFebvGhREIs1lwRg+z9LGm0Coz0Qc4TpNgmItVMVLX51RxxSmBCB3W4P
% vMQv0silgBqwLcN73RPViaEuCMtmlIPMw+BbnU6lkYp9Ut/M/mFg+BX6ptU2NFpCF62+n6J0dKWw
% t+yw0/DWMbeghVEkBCdaiUc4HHbkQaL24Ml2sKOpn69Z0QGXpxgwWNXFxLzCTE+e5t0rQojxbr9w
% jifiVWBDvCwGOAXz9/P/rqt7giyXIS3HweFgv22xy9A5pB/JcTsAks+D7pxeqOTBlDTKLpsVwB/s
% tMmHpkuWitAuqnHQSZPoWCSpqNIN+e0y0DWJHPIFqDfQS1oBc4fDS6i52Wg4hPLyjCqsLWf5C+DD
% awuMYWptCJ5xH3J1bwTaixVNEEqtBv+CnQs0em2q0rvyKInF8LTWjl+b/TF5FfVnOsUbj+0SCaep
% TV3FAycFL6FTOUNymFDgba6EsF+5EMD+6nUYDIi9ABwlQDW3G6ZJCuJvr0A9VpdTsMZxJQs1FscF
% PLeW3LkfwNHQeUFF2Ttx8+O/VQWHr/J43A8ZI3LLtp+z1suWaZ3uyuLPZ5nQ+diATXWKa2wUYrvO
% sk9eKZ3UMVA9SROgCwafu230TPuyKjdKa+UwtCHXyjUGt2MGkZO9dInTUcr5SeVmGHoRA5VaKSk0
% ++mt6yJvGsz/SUXTXub2g+lc50lSIy3csN2rn0Vrscx46JNAZOFemRNeAWBg3YNnFiblVOWkHD14
% khWbZA02EREohkKj6akjeqOqmzXusm9YefXaL29dnfJ6tt93e89/levDeKyvcnERvcG7Edetbne5
% Hhq7zDFufG+aOhJtcuZYAHQO6Abyr+va5+5F0wlt8v/91ri8XgccDBCduLKB0Op8yPi2dlnIOwE8
% P5rfrpVydfQFTMrD7KTK/fwzef/D9+/vfpWTQ/J9xNMc0pNXQ3cFCSOkAEGOsdu4iHLk1ZTdwPPb
% RuS63Q2Lze8Z57sUFHArqnZlIr0gBMuI+1CcFI2vDetvCuMlUy/ErxLGkGUuFoydLRZ4O7ZYnE3V
% LVnvP1BLAwQUAAAACADSqVVdtOOepOsBAAA7AwAAIgAAAG9yaWdpbmFsX2NoZWNrcy9yb3VuZDA3
% L1JFQURNRS50eHRFU8mS2kAMvfMV+gBw5ZClijkmuSWTFEylcpsSbRmr6MUjqQH/fdSGMBdvei29
% Rf41UYavNdWIxmeCn2gjJX8OCi+lSsZE2bawKzX38KVb7fkKQkooYQQvCZMCeg0zUM9WhDECqpJq
% O9rBd3TkhBNJZ3QFVlDzAxhLpm61o7dKatSDkhnno24hlssakneraQ0jH8c1XG+3hNc11GiCHexq
% Nk4EqfQUgXsfxjZDzWcSHpj6bvWNh4HEC948Dhu/UGijrEwu8AlyMUAIxWWUGL2w9No41aFIwhwI
% DpTDmFBOjSr2QE0NHtQpBLsJf2gFPwV0bQWXaPVmTOTEph08lzNFJ9i+0dXIvY0wEYn7eWa6AAot
% jAKJ3QXsqCoeIjXLXbZCqmr+YsgZ1PNxhq2fLPGgeRyHalxyt/qzuBCwvW1BaoZptrFkWOyZu2mG
% QUq6CWpBztCzuD9F5g72xZ29GynqE9+qF2E/p9/z0w3+4aPH4XKilkf9uXr9bsMUObDF2TXlzSSl
% DD4Aj7lo2y7fi6XJJxgpuuE3315R2MZEjmgEF6sW2Of/sFCEXsNI4aQLwodyDrH2nt5l9GW+U+k7
% +IEv9BcGjm1DHfe+dk+eem4JR2iQnvVhHFycwB0qvhD+P/jyhhMevQvre/vVP1BLAwQUAAAACADS
% qVVdbkVYG90EAAB6CQAAJgAAAG9yaWdpbmFsX2NoZWNrcy9yb3VuZDA3L1NIQTI1NlNVTVMudHh0
% bZZRclU5DET/WQUrYGTZsqzVpGRZGlJFgILHTNj99M0HPJL5SFKVd2+XfNzdepJcEntTrRjqw1ZM
% GbSPnr0Oz90spwpXW9XbiBV48pDZtj3WovfvqT18+vLvX1/9a3778PXUO8/Eq82HW5rm9OxLtc1m
% bfCg7qvaoWikI+fsyjOHUh9UGKMPeSV5y+d3K/pgXha9bSbx0Smm7cO0elTNMLZSV/y/dK9NeFCa
% Kk4xI0t/Sf6T3x7r54evP9+JRZ0tNrMf34zHp+wWY86isaSf3Mf6mZulnfQ1q3rSKE5uewX/KfkY
% fnv88vnD7fn2zjnwxNkWTRaRd9HKHUdqV9EcnZ1kruIxTob21VlbrcoxmiQtKPPDU57HH0/3VEeY
% Bi5n9WRnxW8BgMGT5axc3LyfdXrwxHkODmSURzQ3Bw/Xt6ovYH1uxXgpNohx3OU9NHQEyTghi8dy
% EzqB+Wo64/O15lEagDDqXvWOrZy2CWinDq3RGiWRmA+dITgE56xGDqGpIUWxi5WMdQwT0fNG9R7v
% hj35DMHI03rPFrz9SO845Vj9glE+c4FtU2oDxqNmMkbM0Tpf9uoPHx///ngHlyeOfYZ24lGCH661
% p9lyAhJmnacz7NlkJuHBXl2zjsCy4Jv7teaFtmRLi9pzZI+584RuatpittNXSVlDIqLIWThX2tpL
% uztPWKbl3Zy/wfJl0WZLcLGONBDl7mEbpyU4vo+5uwMAKIvCgjrX5r7VTGKds15p3mONnb5xSD/7
% NBxvEA5eBsaDj/ZLZMuVqR7AjIeYG7hoKkATX9Lj4fkV19n3aitkUjbbNA1ZgkN7l2kHGp07eqe2
% g4MCQOIOUEIDecNLOt+IXmBHR7dIOOJAHg6q6CSzYFz5qoDFLNLSvavn8JYgxQhOYAB0k9yJ/iar
% jQZIknZeeA1DwQSuyNyq05tVIFQUoFKe0hxoEqhggwnz8hvRe7TISzuqgGABImW4RKmSVNrZ9qy8
% Eki9GngGSq2xbwyEpkEdxBUyeXjy5zuwiSTug9rKY5Ph9ILislR4qVsZyrLCDsCrBEq156JWsqby
% Isv+SvKlCgTFiodHmRsNaplgiUbBH7yFoKK0AAgh6Q2A5+EV64o3+gKAfkn+hgodCguVg+6c2CWC
% kbFNOnbDwg4QqznY0LxXDWSntaItUsLsEPY/Je+RAro7SpOOwAa2EdrmQo5tg2ZsojAI8V4IWVYS
% zHuufAG+WN+t/VLOZ4/bg397vH18yttjvIyNovCNsj8+DR3KMXYmb+QZ1Y0NoXgC42LSmduuJKPD
% q7A/BdG+duJ8+PHp9s3vbixGR37owASpaXyJBZaolRO7bCwHJ8etO3xq5yUsfKSh2sMy9hvRlyis
% YTq5sLLnFVbD/kNuccGBXau6XCajI871mfTJp8OxTQtbvWXnO9Hft9Y3HUtxJVkja3eR5ckIArKJ
% zsNdeRV23ED1+tRDNEGpM2FBHK7Xon/cm02xo5gt8LUjVkvsw3F9xWjW58bGrFmi0EbcDFuTTvmG
% b48fw2nfv/fv3/P796f8fLuDi0JeMxRbBesb3YXLgyv4OLyAtqk9sHRUBvxdc3KPkYk6QzZswen8
% P7IX3p194vuFsNGygy7rOfZAC1ztaLHQiiv8KgdshdYGxrURgk4+B4luf8i+ZvEfUEsDBBQAAAAI
% ANKpVV1Cn+fkyQAAADsBAAAzAAAAb3JpZ2luYWxfY2hlY2tzL3JvdW5kMDcvYXNzZXNzbWVudC92
% ZXJpZmljYXRpb24udHh0PY9BTsQwDEX3nMIHKL0DEkJiN1cwyS9jTXAix0mH2+OW0WyyyLf/e76w
% QZ0MnIlLoVR/WoGDmtW69SUSYz2jLnfCnZPThMkmsE6smab0EfkvpSvSDf89l/ePvtKnZjTEE4gM
% ZylHPLJ4f5IybdWegGaiSRqXAPdRYm4Xv5JW6uOrO6vLBH1zi62heaW3WLxp3fVVtA1fqKfacIpF
% VzXxEOMJjio7FBT5EHOYnpQp2JcAeBz3+GyAPRIKN60T5aiBedyd2KXq+vIHUEsDBBQAAAAIAMyq
% VV0V0NzNZRUAAG8vAAAnAAAAaW50ZWdyYXRpb25fY2hlY2tzL2FyaXRobWV0aWNfYXVkaXQudHh0
% hVrrctvI0f2Pp5h/IWsBkiBF3VxKlSJRa8Wy5FCydzepWDUEhiRWIIDFALpks3mHPNT3Xt/p7gEI
% apX1D1sSCfT0dJ8+fZk5nV/evf84u7s8U6efzy/v1On1ubr96fru/ez28lbNZ2c383PvdGWy6lgN
% yzyvhlVel/e6TKr1xlRJ5J3ryhyrcahuoipfmFKNR+N97zbKC3ys01RZ82gyleaRThW9nOkNxN2P
% wvvRAUTWWXx9PcRfaf40LHRhykFlntUySY0deDdlskoyvMl/qydTGpXllTJxUpl4oE7bBWxdFGli
% YvVoymSZQJGizFel3shbXml0rBZmmUOCeTZRXSV5NlB3a+N001WlozXpZn1V6moNEdVaZ+oqWZS6
% fPFl9dpiXc87n11dfpnNT/9yNfM6u3rKy4fh1jr39iWDIJtY2pR3kTwaVegV7WQNnZMsNoXBf1mV
% vqgo3xTYZQxNVBgWla/CIMmitdroEkawA3VZ4aGs0vjD06rOaJ+xamREL0Nb6aq2sMVmwxprFpqa
% ysCAGUwWkA/le9q/Sg1+873Oc5l5UlWpM5vCsbHaJM8mDjY57Q7mMDDfBnIz2BDmjhLS6BEb1llk
% vCgv8zTllc2zJoGWno1VmmySChu40guTWrgB+1GwzTGb6njg3VZ4TKd5ZlSc6FWWW1hPfTq/4Ned
% ZeColdLkA2y3VB0zR2sTPdgh/HI9+wGwvf18dYcfV6d3cJK6u1EAtLqdfZldq4v5zd/x49Ppp9n8
% Fh55pl2pOeFQHdCuoAO2utSbBD5h5U2lvvTq/snjfdH7/j7rFXW/P1AXOVTIXpxhPYun8iW0h+ZZ
% VAGlWWZWsDJcnmSVWUFju06WlVXVfcKCi9wm/H2RA1pWPdwncAUQzUZXYnRLUgHluI7oPSjyHX72
% v/6Kp38jOADvK0B6uTQRCYPSVa60V7KDgev3WpeiQalTgBhGcwDIyxdlyhL7uOnBtL0fIRXfYBWS
% 3YfrlpUpvRy+w55WbCjZrXK71VFVY4nvdW0tIKD+ZcrcimmKIE4eE5ss4DfZt0fvy3rJBjt6RBhA
% 1e3aG/3crP2OF9sA6QoqbNRaW4AZwbupU+114oa0oEfXyYoCNgZdiMnoQ52uDGIXSGLYL0ylJeKx
% er4EXDeLJCMt1rSjQpcEeQmR2BDo2EE5BapDyL7Yj3fSWNwr8vQlg5VgibWBHliJH27fOvDVoq4a
% irKKVXzKtkEHgjXwjsfhKNggDXdoSFtrSucHzbFI3o/S2pKGSwRe/mTZRIDQC0Lhji2QBg3qJHix
% swp6W7ycpgm9aylYIzDQsSKEk5UCLWh/39P9vrfMyxbCWgKavjgZtbLML4BBUr2QdWvan4pz7JOY
% usR3CWJWZ16dsddhmLJBJ95RUaqTzUBdG+DUgIK3MRTBwbx90wCYrErc1zE5rEy08ajTmoW+w87w
% dsLoEAnAThyzGPgYfJEUmgzpcdxhE6vSINmoT7/z4yohkDJC1Y99Ylo4H9amIC0NEbGJvTaMnY0g
% aQ7Kj2sTYGuWLF1THK81nmFkwAQSB/Aj7zhFCksR+9gehSY5eZmUFmZMgLIS3G4Y1NopG3zoWiAC
% Wzotv36AngN2fRPyMBiFjsRtywqlaKigU4mEpArgaPvtyoB7qhLpXUyEnaaWnAtZBklmAeQgrWE3
% sEcTCORL2kHHt0Dh7Pd5gkAKlboM0dqQwaqQ+9YMwKp/7P2jF38dfxebtNL9XvFd2O+FQfH112D9
% Wz/oxe6Lr+N/DntFEOIX30NyBeIWJ+V3DOhE7NTL/LIPVovxecia8qtkn90UEMHNvAMimiQiomxp
% g2hYbQzFYxxb9Y/Fd60Ksv4/v44Hauay37FXnEz87OTQr04mx/Je6O/K2RtO/K5xwuFkgNem7rVp
% 89pw79WLk2E42n3zaHg4gsk/zW9uLlxB95ef1Pzm8/W5Nw+P1RmlSSORkBlgyiJXK4vgLc0SsBLm
% CcJ/EzVh62QFZPBCI4BjlcexgwMwio8bxlTFyVjePJkgKa9qV0PRRxxzfz6ZwuV1Wa0Dft9boOJ6
% CGJDuFEcqMwd17la6YIeFfKsCG66hh5ifqYpIESXBVjbeq7GAJYpUViigCQ+Vk7uwdc9pFB49OuU
% 4qwXjv29sX/Q5zKqST9rsHdJkRKrW+w4fwqCjzr6mKCKydRZnv0MSTWwFI4b/s0zn3ltVsb/91+L
% x3N8PPDeg7NyBAyCBwLxG4KBklbGeQe1XkQvU70GB83Hu7640HW61lQ8p6hSKfqjepFE5FwEF6KH
% qKwlOOEQ63vvTWYNSmNYBQv+S7tUhZxKKcyqkR/6Oi3WjrORMAswI5GKLVIdGfKTcIVIbLkaHjto
% 3Uc1JfT8iymzvEbK6CiiU5BrxjWOFUGC2/FeU/+xDUoka5bBagE4OWUIpiUkvhokqvKFrUogG1I9
% Kl+QKF+akkMAQpa1CcmUBM0J8i3MNBqj4rdEYZUYmXLfO7gDxWOLbar6kSDhUfpRdPMptPu5tluP
% TXY9Ji3D1iY63iQogRYJZzQgkOiWsfIBfG4ofxZ15SMFJr8gL4Dhk0296TieqhyfI8kVOWQnVn/r
% N9loqbI/n4wlWXQLF5WdhFxR1JWW3qZrG4mg0qwgugwA0Y1pV/KoptbRQ5Y/of1YUXbSZUotFExY
% p5WVaozzdMOWTdQ0ynlteEgq0c5E5rkCSttWC+Sa1pT/epMDfzL2Jwf9LlS8LbaEUqDiIgEyLcp/
% bpLiMi8KQoaIR5WNYkJ2RZZflvmGo9rtjYO9rUVQd5mUa3SO/61lKaWJ66jeMuTwvVchCv/oavgD
% 8iCU+6XOKw5Km6LPhd8ChG0E23DaRPuz5LZJqgiPsieVSAl2R6mcHLje8mP1lAdckHdjdGsIcik2
% h2UH3q5HLzqLAHMwM2Gv2Wybp6lbI3w7bZxluPZA8n9pex7xEFMBpwXDMWqxuIVzYBYhvhfj6DP2
% XFEBSJpNQQzeYczpIKTXY2OjMllQ+sioNVmj41+tuQKmVjah8F8FHqp06nQJUxkRZoSOwrB18mYC
% 0CCOiv9nQgCkc9nJTbCEvmeRZiJqHLKa1i0qiYI0XxB95Cl3/UghiFqaBsA0CGC4hJtuzjew2ooB
% MN0FwG29gCkg+87V0RlX2g3tgjWgEsWV1DGPhkjZer3QX9yPvz74g8EAtUe3cyFHIMqrNrM12b31
% kADEi1G+ZDqrgoW2LoSYo1AQJpn4AKbdUDPK4LFkDa4rIbHgfC26C0dLUJO6ee+h/zqvCC9b7sfL
% mJmYSiC3nlZc0hIbe7ExheM1Bobm2YFrZRrKHbRGRISugMnAFsmD6YwgsBB3hshIdUbA43IZcBNa
% IJjTDinjkwYtPts6zbaGykskT/pOimeukqQxgp6dDplBwZANGnF+p5+L6ipfLn0aN8iMAaENzZON
% FrBzMHG/2QxGxKRNE9YwrydZSaYVnVTamrDtydiKBLr9XdAhLl2221b7qX6REsx/3bPC1UkOHvKk
% IMKqPlDxDL1TZ/UoB1lGrqxjvnmE+tTCNx0cmVje8XhEtckfuRtxvCqNNMoqQ8M8abLbMYx87mZq
% mroXapFSD3bjJoxyl4w6jomTbnphv9PgEGhlZSWDoCcO0oznVVQFOfhuYxsZI4tT7qcHO2aTitbN
% utzGGVOi4CsapRkiAyZw/aBrP6VxFpulWpDWzFs4I3pu6ODGN031xJuoN19/5c9/O8l714j8p3VC
% BGfdhtywcgGD42uPuLDM02Zw4US4hEXAONgCY6f3lR4PjvO5xAYCY9f0wm1+O/Wjbhc1vFRYvmpa
% IsoaTC4VcbiDuO/oqCG4nSV9r5PksGDQTsyifO0cScCiGHHB4RJU0xkwPrxtiaSwfo/6tNLotO/T
% b1BjxYHc7wsduUZRUjSi3HDrjCo2qTyuZezuqAT7NjBEZyZQNVMfGuBaSUbcL7TtqmcRkboUNmkn
% WijYKh4NWGlUmzkIT7p4GNkMQiGQJ2WUFyla3O6pV/qjMvXVyJHGQ8W/rwduetPMr9CsmwDFwwr9
% NsOaUmszJmmDsxFFnuWn8qU3l9J+Lwj233Hade5oCS2iHEg5V8a+iEZq9ji4tQt1auN/vJvNr0+v
% 1Kf55cfT+U/B7c3n+dlMnb2fnX249cKB+iuMyNaXRorN+CEhSKi2pfLVnEDYVPY7PRRTEDcAStGy
% dZqrB5bygJ5eXbrhigpD1RuPwhAxdTrZG9DT66oq7PFwiP2sB1GerrCvARYa/qcZyQzTyR79GxTx
% kt74wSyUqxUuz0FHKFmyg8ljYp7QUvMfe/wHiz9PiDThXz5DKAYHw5THhuPJKAjGexSaO50i9Txu
% pCaFMzVzJIjalyaZ0jkCTwhDIUynzbZEpQxdy5gl2fqYxLxtNjeeAZOAzDa27aslOUc13EogE6Lj
% +SdEtb0qxf/zWlPLg+oFVUOm5DABvh9DNGhyExt0Kmvy4pfB+4H6iGxHv6pbdOQbNJlnA/UlWQEF
% Pu+ojT4ZYXPnjyVleHDLTW2i1amaTMmh4ykcGobTIAjHR746v7lU4WhwEB6E4+HySY+Ch5+Pxjve
% Lg1IodJ2UNvNIEqH1okc8owvNUOgMEtzHQ/3x8MvN1eT6T0tc380DCcHfwyCUEBwFFp0QtFaYHDH
% 428qLEqajZlHbm+K3Lqs8zMfA6XtqQUFW9shDqUPJDkLtFqrUnLP3JCpuXQchONDfg/JzaCYX7vh
% qaHhHlaEcCz0/u7uk9obTd6RpIpKKuIQUfPVeZKQTdVi60+2i9N91OdEUIQc2V1tOxOX7gC3QHnj
% jn+aSpWVZQJjkqNBh3lKX0gQJV6KE+r8yB7pi/AOswvnbdfVQH7LXNzNNRE2n7T23iZ7aL9tUd1p
% lOaW7lVX28KOZDQtLoBuYmI7GX09uVyQ2HaMyxwtLS6+2MYCScnQ7aQIkgTJHx+uXDlc7pTzSWYL
% qa1IgMTEn7hZpuM/z5tgheQhqbT6gLwDn7OzLzSdJ+Bposy/5xUSMSqz8sfk8Xi8PxoPRqPxaP9x
% 4ivY+681LETHqzthoMvn5HGQl6uhXtjh9qW3n1lXm3TYlfz2Y6DJV6LeDJYpB4vqoSbijmTNU7CX
% vuPQfff1EhFAYdH3GbXtV9Lj09+HB+7R93cfr/q8vy9yjossPmmGDg2JufiTw14ZLlzSkUq50lnT
% kFFbQSDSlLVDwarrW+FmDtoWHFSYPt5P+HBvgpqDtHq8n/Lf037fAYmn/mjdQ58L1BXlZFaLjwWC
% JRFn2/WcRhF3IyvOAVjOse3GBPxeOy6xagOguTmhcqFG/SiHFOoPfJKnpLVEu2XjnLnDKcOVgZLK
% oG2E/vAQV3oVEkLVC9MDHdZzhy+N0TYPtOlBAkOjW3Jxx+8A1XtI/IMA2WAGVyB3CfW7ecjrnhmG
% q2GMMoBBSC9K47RzY3cQXdSLAafy1FASQwn2OMjS4X+MLPHzehgnOb0bFLRM292+ldtdMheCDJln
% KtOOsOiU/JVOcR5JmUqiLgi4nAxbJkbVRMOM2o1Lmg7OYYmh96QdiRqGIG3FULMrFKdtU/m5OlAm
% WfNph+xoUECDMh5zVJak7NbkLca7DbpkeiZitTRVtJY1OPrcWOvsZj6fnd1d3lzfqqH62+fTq8uL
% y7NT+WA++9vny/nsXF1e84H55fXd7Pv56R0+4eNy6pvI2JXOGB3bbpKPnDvdLLOecTczYMySyton
% jjRKe94ngnn5aHYvWphSE4DwQVI6DwbuRoPkRYKpG/nIju2xpwJ1RvXVsHvCQFXQqqylrbRvzevx
% mswVuuc9zcyLx2c0HmgDSkYC9F1nkMBSvm+OytwEjnbKk0sutFzMNW0j6RPnLJ+OFt2ptRS8cthP
% PPO6lHc1OBabH0jK2g62c86Av+8h5CSncwTfoQkW5erpV2f3QkYgItfKSIbuTPzluFBEzJ6lSu3y
% Gp13TbfMujDElM2oIKKQ2K4oLmZRd4603rpu0i1FmApfG4gbfSaQnaFM50yBKpDmbRf82ya61Wfb
% 68zOPhPqv8zmCI/Z/Fb1iPvNMx9jYpejvle8gDWy/3GjaRQ2N5r4NtLLoKDSaM8/Go2kG2gOv+Ti
% Cgpf/xBf7UKYv3qnJmN/hO96mf9AmYgO46yrEU2MN8cTFxN0YPJuezzWnDgJ4RVIOINvKD1+S+lT
% 3jhxT+cwyt9CwkevQixAI2x3CjjdOc5R3VzHagitulNxqoN+d5HH6TNI89W3lJ68pXQ49id7053j
% MyZkmPPQP5qMnc7ci7xTR/74aH87endPkpij7YFFe8zQIk+CrL0H1VokGDdHGt/Sfe8t3ad7/h40
% bMYsLQ4Oxv7RdF+9PkygLQEeHRaTOhiSSCP6Pjzcc9bG76E/JaQJ93WHqTJU5lN+4m4e4MnanpQK
% MuOVwrm9gvL6uEAOAr6x7+lb+x4Lzhn8DCYZk7ZU/w7xQw/QREZGe67DTWiP+/wlxDRjZr42BLKW
% uSY/M94btTP2nURK749YuNzXq7hdZzQ8JbBFG4qT7QNyZyNgu7oHvrXt/Te3PeKQhqCME+vOrsaH
% zAvNYJdmNZFJ00bpyT4rzaTpqXZqKKe5tF9n0Zbqmm+AAvqmM7p1W4SUxhLbGaobHm8V+9ZGD95k
% v73DznlnO4wUbcJwfzuUbLTcP3x7Bumpndf32UbdCdl28gDhCO1um9awPYSgtERL3jw4DTsjzFfk
% 6zJK1zWIKUs3rQIS1M5AHQ6u8hW1SP/7XugOz4XEc+3RmDOhcJ93en5+SdXZ6VXnXq4bup1mjpkv
% Ss3GCDoO7VyjlBNEs20NCHCqAZyXL6QD3nbjRcAFddM4tDeJdi/FtTfhpJ3ydi/48AkWyH18eEi3
% rmAGusoVocKD8SibarrIhLb7suMcoWWv8VHQge7vaxAaAf94Mib/8M1cOfuWvs7rFX7WP+lN/MO+
% j/+P8P+Ufz/w9w777hYkz/eLhvEYDXwss12hNF5nvpvS9VXqEVrbWnf0ZV/AezkZm/mRrrrmq2O1
% vf7L3myvHu7WvVIsO0puKoLuNatO+UjXe84vhnfmR0FB9+asXJPVW+egIaIG8Bl5LopgbOkHpnIF
% GT0qXe+k8yR0M56O6b6Du7qT5QqNckmdw5CPr7iHWOTPNDov+URryZApTXs9TQ6dUNuz5SZBsIea
% hfrQpvdzO+/ApB0g9cWDJcGAgErae9C+yotCqgrUsTW3T27MYuJ32yqTHnCXpDlT0RSJ+n5ShgcG
% 7sBXmiqpH5fwsAR11lxs5at0fCU0Vos6wQt8nY1vr3hP6zw1QdMdbjc+FPu5a+dRcwlHMmQjFAH0
% jI7y/wFQSwMEFAAAAAgAJKtVXfpPPHB8IgAAzVAAACUAAABpbnRlZ3JhdGlvbl9jaGVja3MvZ2Vv
% bWV0cnlfYXVkaXQudHh0nXzrcuPGkuZ/PUVF7I9DnSZIgZSolth0hC5st06rJZlSt4/H69YBAZCE
% BQI0LpLomY2Yd9h5wnmS+TKzqlBQXxy7YYdaIoFCVVbml19eCj9Orz9M72a/qLvZydl7dfLx/OJO
% /fd//pca+Or67O76dDpTg73BaOekjpIqjlSZ10UYq0WSxuWxypIsVusgq8uwSDZV2VVZrvIiWSZZ
% kMpFKuYbezuzPK9UlBRxWOXF9lj1n/LiodwEYdzH3UEVrvrhaz88mr/2Dxah36/wpCxYx1l1v+ff
% 7x3u3G6zahWXSakWRbCkL46VcxEN11/G+Tquiu19aS7uVfHzzk0RPybx03eu38gVvU20UJ1DtQmW
% cTlWYb7eYBGRqp6SMJbVPcbFok7Tfp1Fzm8LSCJST0EBkSxLFSyquMDtBS03ybPd3s5JmqosflJp
% MI9TXJFFap7M0yRfFsFmtVUP8bZU8xiiU5j8Mc3suLezc/1pOju5vFQnt7fT29sP06u7natcLYIK
% Al4HWCN+JCH+WAYbPL9USYTVJYsE08FQuOCLXVJ5pgK1KfJ84aXxY5yqIg4izLun7la0txl2TGGd
% NEqBqRaxKuvNJo1JcHgURIWHiEAqGQgXzMsqqeoqLmUYSLSMi0fSmTDfxOqPOkgxYBiQPGTQuCxp
% rkHaUxcZZF5gJXUaFDQupoxRafpmj3DBhqaDcfP0EZp1GyRV/rdS3SSlTGOOGY7pnqRQP9VxSU9S
% fu9QQWcu4+Ah/u///L8f8lWwXmO56pf4oQpwe50VcRqQctdZss4jmoGaB8U2xOSKJFRpvvTCPAuD
% x6TaYh0yMPbmavqzOp+efzy7u7i+ulUXV+ru3VTd/nKFf24vbnf8nnp78c/puXc7nZ6rj1cX2Mvb
% k8uLu1/U2+uZmv58cnmu7q5/xqe4Mi8g2BjrXCTPmMwmL5MqeYy9KIHYSzyRdjyHxEk2ZYxLTklE
% 9msVsVLJGOsfJnvdHaXUh/syZmPprF/5u/2hejPBZ6ed6NV6l37/j2nndPc/XlzWoztveF8huvlW
% QUcXFTREBdBiY+DeAsYLm45LSCoKsgoKn6ai8CVMTclEeVjF13Z5hpAl5ljGqTMubXT1lJvLvYLQ
% 4ivDySMxr9OeOsszYIsIJszrrNJa9Rhn6impVmrvTf1m8jhWUc6Kiucu6jKmZ+EnXxEU86QqsNdq
% meZzwiwev0xh7qXqPK2ScKXCIFOr4DFWwQ/zXX5sGUMLsgoLcHeM92QFE6TFlEAZFZTb9abKodZq
% Data12sFpItFDNUKsg3SMlerPI2Aai8UQNthKeulwbVlAZWgBzySShPMU9t5kq1gsxUbPuDnEcv5
% ddDbG3YVfo4Gh0ejo8O9w16v9xuNg4fHz0FYybQTYNgGu0hmzTBCAgtgn4WMp2ZDNR3BTrHfeSxf
% h2ku0pSp2IdiHfQh2VVJBkQjTJ+CNDJC6O0Meup2OruY3qqbj5eXp+R7rt+qs5Or66uLs5NLdXY9
% m03ZrMQwwhxLIzUKV0ECbXi+j3+fbO/jVyV+YWk+BkUSAADUHJoQwRMRdmMCrDmX9/FkHTzf/95J
% cb1HN+121Uf6NMnwaW0/pe2FJhECFWsA1p+8penCI4ygXXnM03ptd5CQJkoWC5ocBEfWERPQsQDm
% QcmwB4iraJuCSnXm3cvux13GO2BcFi8DsnEgflnFa9GcGMOF/Ok6KMgdkAl6Hy576qqZEwALWhXb
% 6fB2recJyWhahykAGlprvsUIUPyIHhslAG+aYZqT15jWaVwQLAZhGtQFnkZYC30Nc8wjCWlA7Dap
% CD1iTX66q+Z1pRdf1ph1lNPcoE4wJZZ7XTHK03PpVlLmnK6d09MJq0qtgbh4Du9UqtkI6gVVLEhQ
% cZpgX2QINtPZIX0LO8wz9nWNYxVHiidDd0V88zx/eIjjDaNK7urFM1zTzhA7fP3x6u7k9HKq4Fi9
% KyggIbJWOQD16Qexs6/su2JcIhsPjH262yheIpb5MCbnYprk5gBN8LOAmDgoE5oRHB5AgojCJk+3
% qzgq4FXULN6kACEVB4CefJ1UZD5VkKSEw0lVyiChg31rbPXYTGeRFCVUYQ4fCV8Mh8Hou4bzKwVx
% UlCMKsFGxj11XuSbjYFfdyH8PMxsTW71MRbJMsayLLFl7M7hDkn8ZPlA8nnCn2i3TZL6YpcP1HQI
% l0v7K+6kWNY8wWab/0agmHoZmBWN1mx1G3yS9QbwW2f6e8hoGWdxwYrtuOvezn5PndypD9e3d2q4
% 5wCLupvOPrDTNhpw9242nRIArA3VIUeaAAlKtSSpBC5DKILswatWRRyDuhUZsaj8SRuy6A+tdtQX
% zQ55UGyIMGBSjRDYKoQo5uFLuOMRngD3imdjQfbSkgw3TOuItBhyKpOleEQICs98M/EH4rhw4YkA
% IJkcJNgiD0OP9KwiMoYNyliN30wGe0T1KnZ4tJlvJhBTHIH/dknjGCD4C+PB6Or4GW5/SWtnCeh9
% pv+tH2GvZSCJhMWD8IzBQDdMb0hLnrRXK4EZLHWN4AQZa0gaq/0ztoRaMB3r0ZAHwE0y4A24dKWF
% /MzbvnPhzMTQ/z4pJUcW94CFagXrI6hbxeFDHInbJqSzrESobbAkn1OpmQ+gnMGdJrwZdAFWjWiG
% MUk2oKtk+xi8ug45K1fJgr7WfoOGjuqw6r+NH0gtoDWbMcUXROMXJADayrv4n6rabnL1L5hb9C92
% zI3CQwH+9b/5i567snITh+DUfDHTi6B8kKuxCnxXNQulyWjMNZaIqAdfEZpgCERBBB6QOH5jDCcA
% IuNKa4HxnZ2bk5vpzDv9xeNfdAh5cnXOzPjmEu6crG1nZ+aTHVBsWdDy9vz+3ui+TuEy+kzsOVC7
% S6o0Pla3q6DYCGfwbnLAi3hlZjq3wr1ohbcJ77d3ivGASnf5E2glIs0Y8JUda3rTkBAPLg58gPU4
% UIfkj88ng7Ea/v3o84P8NXx45Y+V3/5kOFYHR/ik8+D5u/bTgeo8/DDxd8fEfsTAkoy9DSvc0jon
% rDArF4TypE29nTNRt2OMVAHjWeXLNbZvBYgs+xpMCfz6gGSYGeMaQ/s8LsbNeB/usz/G5FqK7VNS
% El14xsJSD8FqSVMg/cIsKIzNbDwWmI+se4e650/VaowrkkzYnl6GJ8sAHQ7CB8XQP+ZZJxnNurdz
% S5HdMYe1IAKaPDOrfQq25Kk8UJgHIh5nQcb2uiHU1koEq34zmTOOWRqub+Fg9OzmFiNRbLplN5RE
% 2khxL/ZUVscihkHUBVMw7AbvvgoLSjz0dj7pMBafF3WmbhBHG81wltVtCSRKgmWWI8oLyfFBXZcr
% lU38vb2uxmTsF/moihQOGwDDKbTB3yACwYKibRasAS2YNvZybUdp8ACj9XbOsUebYxvh9tnDNhtc
% xMTamH0PGHrADnwgCRlH2ZhBytNg1MQtaybIAR6FBWRiADAzWOps0LbAwXcs8BRzCFcEcdrIlBiZ
% qPTFmmlBpD5ucJ86tYDdstkPbGiNPdrFe3GyXFXYhQ3irIR+CegZQgIO9g7BhxJGogIIHgjl8Pfh
% 1PMiIl4IfxnAQVXq/eTXX/3XCGq6/PO37q/+0cGg6x+NXg9/+w33r/KJf4TPXu2XfxRVZzA4OvSH
% r2GzovNGaKv8c8fv+/u7kxeRkmOtEpohUtNq3khah7aQBLE1klmjxWP1xIuNo5aGkafi+J5scV4k
% WA1x6yJ5HjdiIW+dgQECxMcwoG0ab6EF76BYaUXxn6P0JlkAfkTMmZ9gnzY2zkY2DzQ8iIg8Ej8a
% a1tgHdbceg1impSM7trAa95mEVaepVvJhG2IRhrKR9uG/Tc5G94fIrMWfHWAbdSSGZC4Mvb/iCU2
% q4IipkSUi3ExIRBJNSfNxfg5wiVs4PloHmDSVRAzCwWMMi9SRgYozXdhwBGzpCdgX/gcNxACqAgK
% ODZ2ry0VHztWTwZ6PnnNN59P/IMGor2VaPomACsfOztmAeEB/GuMuw4OHcGbTdFhtebaCYUvGjEQ
% jetoX7N6Z4dA602+g9AEisTZSwKAYRsAht8BgJNMTZ83AGnBmNOgKEh8ZOMnBq29t2wJf+GGh3/v
% DPAoAN7u5845uVAXsHSAZPI+xtlFjU/QaSZ23mPJjXjMXv6Mi9xjp6jHaj0n2jViHPgjYAcoOgVa
% lV6Qw74dI3/J58ByHVPVJl5y5JamcdoXgKLZWl/OaeMmOKNsFPl11g5JR22SOIzJY3spNjdwZ0UX
% mcGdKGuswIRxwYLSzEqvxXHw2Q+Toxa+IAgtiLKxXFpCY6UYc+6mLihJEwkzEXzgMF1n7FjgFgKc
% FBeuscGm/bDJ1WmLbDJrgFogYSsU4WzWmDJ6HpZLTxPS8sffUwswlSwaqqU1QEuLr6RIR1MDOPEs
% f0o5ZGmKDwQQ5zHPk3Zu1Ns3sUeQbWVzv4EKFFNAZbSYy/7+cNAyQqNWCAKXsZdJmodZG7Gz0d4L
% rNhSKnCTSprOyqBU5NGPutqt/MGSz95M9hklqMgAmQJTPXaFxkNa+5cUoMcAweJHYEo5VR2TzzEA
% 8zCmZqFJWIacUycFnrN7xxKyZU16YqKS2aAPkDAaULBx7Mz227Cx/z3e4OTNLxGIn9m8OS3wPCkf
% vKu4otqLekuWK3Yxo1j67ilXH5v4+gO5wzB2gMQ6OcqCrznE0b6120qZevP8uZXAd3eP5wEn9ADR
% 8DxKoR1sKaSucFKQDS60mSOT4iiSvGi+5ieOOVehP1N5jYs2+B3TD6EVY5fYYokItsXLhxRcN1Gk
% iYb7hGhkSfU6Y3XROSsLTxxExhyK98H2bY2ldJOQwo6DJNP+3qFNIOnuc0kHcnJMCDgLb1FnOqXG
% EbXcbJLQ8C95QdimFV0RTi+h2+zewekBlDDHMi41pATpZhVMfPxVJSnblDMPCAPG+mbij03uSLJY
% kpfADpGAG/CzAZbsN1YLI4pN2lvAjW+nGpBJQ5lcB3x3reHca4BW596K5E9N3IWE6bqc0QW6ifw2
% 9m4gmQ7AFZTFgV5wx/W8LkVB8EzDCIyMnfhNC7EJm3QCQ2uM0azxC/w0St4nbYjU+/0eTId+Y+0Q
% pPydUAK7AJ3d0lMlhdJT100tli1DT0HHBHg6RWE2pwayBoVPSrcukmRaBYi4Fcm8loTU16FzMOoe
% jAZWTduoOQQ0NtJyv7EUGREoGGWfN6yhQ2KgiIXZPMTSrT0ma2N9vHO8i2bHiniRcqqEDatvk07a
% wkjQGfkt8SsGMZi8aJx1LVqSQfD47Cx1XNak+HrqRlujKZJLCMrzilxGxtlRQxx0ionjCNy/Mp6Z
% cJeypQ7uHvT3hvcriOorbE2nTH6qg4ISZWc2a8p4B2j1zh3368CsoHADsjrvCqjqvkCv7hdaZrCA
% FXGsOMR66/n9/d0W/Npk7BhoYPc6zZeSg5MrjIcaO2lmrbSc0vZ0Jtum1sc6EMbj1GMSkFSrQL2/
% H3SHjbu0yBm6SV2Pk7ombWsqPTozwiyBWE3yzIlWrnzQFVjhTRI+wFyhwIhtH0AmHhIqOQa66A7U
% +V83E/qgc/PqXX+w+wqTcwFXO1ZBGRPkNkNBDCa3MnZDRJJV58fPg1csWlNkIJKgfvxhwmIf7vYH
% QLrJXtdXK4yecuzpElFK6Vvu4uYQoKnDBwxOua8Gm3T6nbot8KsOTs3HAjpYF6+b4YVYuGU3ruNj
% G4NHoqD0hePWZURhO9+ClO7+60PuB1hvNKovxUd9ibBjdeg3+VGuIfchOg6/VAV0w/qPEJwRpzR/
% D0daD/oE6lzG9iiRT/4bXw/2jrRm6UKAlJu+lKGFDK4arXDJkvyA0Yuk1LQQ5oNYgYFUkrJnqxpM
% 8H1QlCvcgEg4L+K1Oui97hkrNqxdUQ2UCnmu+kHNhKDokq1k3mwFuglYuZAl3gloMztUnbi37GFZ
% ZUlSFB13ag67No4sk1SK3bq1ZOt6rVYpznlwPv+dQntBsVEbxf4i73vD4WGbOdK4J0NKr1ViElNb
% YqPVSwiqeSTXZL4DcbJxZGxJVucwZkT8a1ziNWW7WIe+ar/vD8etMpMFvTbCOfhVz2PpX8jnTZpS
% Nyj4XX/YPQDvKWXGBKxDW/4lHPuLShGttsUzLbC5aEehG0exJXdDiProgoro8FMgVs1JvrG6DVe1
% O1j/HZVZQKlAgSgidtARBlckj5KT+b1eb2AlEVUxdeLJ90ej/c6+5w83u23qjM2u6mIuOsgJXYJE
% 2gwOkHXqiMsXun3JoVQGuRrpuvkSf2+w71ELUlxpkToyGxscf+TeM+kS4a4rt1hTEs2EhWESQWao
% A3+Oibm1URMXYN1BuA05VwclLfK1QHoc0XI46UY5mff7lHHjx5XuzlmotbomeskhuJXBl1raeEDb
% fkWZWJ2NI5oG3lHFdliuKAvXMJEFF+K+zS//EpWFkN1u1zfbY1uma6XBy+16nkMyoldc0ygW3Opn
% t1tDY0IUVoLgRuH1TjKjC6jjz9nNLjSs64/2WwYomF6qjg8HOHjdIna7FpoBIFqtmxkJPfuqqb2g
% elzHDTZUU0xC6Q5sWZWTrR8RcePAVm7ixLyUJaxlnv4465/HCJIhNDiAW5jGesGdgZXl2QSfozYJ
% HH2bBF6LCnjXRdSigBLzvkRVaPWFhrK3lBdx62Y6Ac2uDDqaGdDzjA42qK9mog3vvdnkUyd4NewP
% uvNXg25Iv1GiXdK+DX42Bnz6gSue2GzNwJtR2wWq6o0PHEpS0OKxU0WITN6qBbQFgDgpNDEgTxB5
% g9bIxlKaKVkL9PuRW58Lq5qreiIls3aObhQ3aTAj7TgVO3D+moqvu2MbBOnGEaMmYxNHcl7HaRBb
% UF5/qzoPXfy321Ud4mNdKjTSD58+OaA/8GOffpBw/f4DVkntQu18P6YCjfDyNuSarLStBwDxtrbw
% a/Fopkw7a5454Tn3Igl2SPNXgI1gBjOXJjCYBhAP2FjpRpcXIXFVxulirHQHnpGPZxpr6A4LbdTr
% UbpaRr0KTe8DtxuZfg/XOVeaPi3yVHT6r9CLikWHra4FprJ97gA0c+MGAZdyHuz77ECJP3aHhwd6
% kz0JW+TysXo9JMQR3jhuINFhTP1m/4lr7nPTjAeXJv2oWudF6CNlaLTnyMVcMlb7++1wts2Dhe5a
% JGSKto6h0llSrpkQMlI5hiRdXrp9jT0fZXl0mqlq2jXKvo62sCtgldwozA0FYPvkZ2WXuKXXM3rX
% sFJnKaJZzUWcJHJnNLb5LqEh7QJnyfMNHvNE02Xd10RNanA6ECc2wwFr8YnpVmPsqI2x36ao0i5x
% oWvk2oEEyobXmn9eO5BE8PuSjzZI3ADv6aXTQiBpNNv4aFrjLEWlQSmVrTmHhSlTRvxOexrV3j3b
% SuxE1E3NhRt2nG6nOltzmk0HgENTjtE9fVSZ12SICMSGfQ1xZ91mm2o0oEC0yXO5kb4T4lMi2/Qt
% sML3v2DT8KSaM/wwGbpF2gxqSg1O0mSX2BY1kQYTj/m2lYqkHZdkaNSgXRgnKTcAUaWU+LEmn9S5
% /Oxx/cqTmlOrkkubRDDW9C/2jfnrZIPTweTgCSe2qBVqrP6gJjAu00qAwLkwobRPlBz9YzIa2+7P
% BtG1r5LeLJOQ9qrc03ojV1KMp91rTa0gVAFqyp9jLlZVknxzaELQJAocf0VtG1ySMflaEIZ5TP1o
% INBSrGnqx6I9nOOlOFUZlUy5s4NSkpKkMEk+Z4NMWpoTMGK2OmXPIEb+wORseoiL8Zgm0ddEs24y
% q90SKi1ZNi0F7aWUQatkIKV+Iu8lNF0XXL/0Oj31o4CdETrsmcjIY5Cy8LgTsknIMOTZBhebe40j
% qSidXv6192o6vhFlhwS34nu6bHFfqpxoo+Q2ukosmpIdTcNLV722BSwjJ8+Yri106PVJTNs1FbBQ
% ToW4mZKhzxzDXOhyN6uBfccHNEGB1Gkmw26v1+uC1TuJFx1rpp6YI2kpcaqyaaOn3J0Bj5LS5SRv
% 0shjavlVs+uPV+dwqGfvpmfvb1ms0/OeOqmrFQJCcXakBU1nzr3ItT+DY7C+gTflnjvletVzZZ0r
% FdiaOKEJI9RsX4ICzohw4lcn+7lKJ3nfUqgPSc0FPIrG0mAjTq5J2JM8yC3XxNjJWjhc+FtpyZEu
% Ga+T6CXhJL932PZ7h9/xe2e2+fodRQVS+X8ZY3wn6SLt5mE7L90671JwU/BjbFJRkhTRph6T5Xg5
% RzZwS4RkprjteKmmst6gmg4D3kwQB0j+Io8iqE7TimWmpvUFeg9StHIaceylQ26MdVfRYDjEoyHb
% tNP2NUXEBAk+dY6DgiodQ/C0c8HmtEmw2N5na4AAqn7bkRvooRiK2Ku4IUY15k2eDickp+BoUm/n
% bR5yoBrQkT6eEaDPK6gkyqS+PN7pBLvqpy8biDmXlCzUjZNzLY+ZhdwID2nqgSVmWi6YCKYNrp+0
% gRUP/4nyE6H2iKwJXVXmpomZhUanvrjkv9OZ76pbQn1DdiSHfizN8sF9ov6UJmTtxjnwXnHqhAqO
% 5PqNyfj9AcYLd0lnG2061kQBF6xNn4DjRZrqIGeTdB90xj1XFbt5o0vO4YkNZb8ABZrR2cm1tl0/
% yHYHBKQdZUXBKx8WcZSc0Ns8HHiRRDVWplXp3z5H2pwoiUlC4fxJV3ixtKaVcv4Du1AmLTMx538g
% l2hXXUdEbuyUjqkvIQFocbaLayDw44aIyD6YZh7u7bRFVJLZgMvZ5mrpdqfkXXNcwDlvE8q5MXIx
% +lgcLeSr5+FIbF2C0t/FHIBaz7Yzj9MXLV3QPXymKNpAjjSYNFTDlGLM2QjZHjnoVjb3QU1Ag/6o
% gadEXM36OsmamirgT6H6/4Y4PV8sSqdlkJDUTKuZJPCWzudowDcNYvT000/SOlaaxGkjPW7w5t4R
% DT8J936TCMumKz/QemmPW0hNMrbHvCTKlfyGyddSM416AWlQjXhXTaOlTWgcCyRSfG5Mi87x2dyJ
% lIh0d7w+m0c8FoLUSmOBc2xPAeqCmz7XI1TTsYChJ1lBJvZku3SfVuoXW8EX9tQ5vvR+pE5cW5Hh
% KhlZjgSiWofktMVHWJgZgZcHAresVpJwbY5XyaeQyYJhRIP8sbEBYXq62UHKdaKAAR/vMl1N1oSk
% X6HbtDuYan2JvYMm9qntgatv3CbR1IDAgeAqmZhtiGDbRhkwwM7J57vdlg9KLFqcqEcXhXQvAbBZ
% 51XMLVjhcpeKK+14/Fhd3UedanfSqV75u587EX56lfzbVZ8m/At91x/sfo6cJlWJn8SZu1CpG9Ve
% EAB4mhpeI93qTs8zKnFzbUJnCsCh2IMQ2DiumQtlNqKxefEmtGH1mXCjip7lq7xDU6fn1NnvdanP
% UdusR+KccWriCeeZJlBy++ZajReN59YtGBRXaDzibKht5GpFGO0zTXxehpvu7ZF4cwKOmue0/pl7
% Gy3ChL7T87Z/sG+oO5bnD56HctpN7wZik4pyVJKk6lOKnven72iQraGG277T+KJrqc0nOvtpgdg0
% eJiimDlzp1FnzP14J0NX0k6bMrmdf9/jQOHg/3wejumc17oVCNo+TW4E4ZCHesRM+Mx1Xk2VPeqd
% qup1s5W2RTrPHAxykypVkadjna9i3kXcJpoMer2R7sBypu0wA77E35cOpTY89XZuCINJ14iZHfOh
% HN1XAvWMJdSmjBvIfus9AA3GIwQJCqKzgCt7ZpNe53DAXucp2DYHwk486UMB9ocxIza9EUH3KkAV
% tRJRcFrkAZkovB/lu4kx0img2cWHk9kv6vb64+xseiuh1fRcXV58mna/eJuE3/vWefwudw22bfSE
% EsvPfGRRzgGRV+a2lryC4m7dsONGL4gD/n8mj8eD0d5Rb3h05B8++l3SjNt4A60lR0dT4ZPuq6ra
% lMf9flA8J4+9vFj2V9U67bu30mWtVwqY4IDq+eVD6cCB24TWappSHVDBFQVx7a663W67C67bMn1p
% rbAZU2l0c1uDJQ+6rCl3UHKU27fhXiNIrk9yK8dXW3N6tvPA7+2zizBDNI/qqZ/jORWywMzonHcR
% L0qFCD47fE3H6xCj8x9H/Aef8T5F0Ir7sc1AnWUW0+5+/dSxqSLStm6aPfyQA7WeqBd01VP/6KlD
% 6vHeO4TEhgcHGGD4ejTWG01G0N872DscHIweh9/Y10206LsX0lWfzMkBK4Ej4r/MQi3tGttvB3sU
% HbUycLbpz+SEW1eDMDXne6kFs+GketXO5SPN1Iw8GBn1wzyHIzMhhSrolMFPL3mPfhOJkv7ynj+U
% oxR8/k8XKQ/NfupdfD1wdvH1kP7w5Y+jofON/OHLSWq3d6ULPkENUl61lROmuhJTWLRxqc5SJ+XE
% Rv293t7e6Gj0/2Gjza3f3vLmKrrG6bE5NnNrZRFketLZNWn6F8z0pe7rNXH+a0XsgQJZOvUDC/Ip
% 6UP9soy4FF+D77hxi3nql4YQ4Ot8WevDUBqWTWDCr7BICqq7IZQ952Y/hsgzGDFgu8QvAJycpmLa
% gX89+E0CrmNa+xW9OCYJ1D9q/MwC2OOfAViHDNVV71f5IxT6IUm8m/phlSYP+WPLNjlcoZlbI+2p
% D73bsKeEiHQJ1OLiMSlljicpnCfHvS+m2AVU83tqEK8QgHPSP4Jn+UcAt8Nz+1X7OyhtKv1sdAyX
% 3jFDohBjAPzSDWR2wW9ddm6sLle5yWJyT2mx1fNTH2eXfYZBhM8Vn5Elg9M2UVFLP3MmmE24khCq
% bpXkyY7sWSPq12nZgA5h5FFS8TSdWU6NMd3yG3covtLv4Wmib6Ds+zjeSJLZPSxlHM5XrdZXHUf5
% dPq9HO1hZqODXdek91Sn9dohfe1waC8byeL3+j4fnj8j1Mgx0A1Y2ZZ+uQVLyLHTfNYP9OMszyh0
% omKCLaLY3AyJNUiX8bwI4G+4sLElkhdKDYGckEmyGiQfDGGt/t5wBBynLQX8Yz60t6BKm9pkyW+K
% HGp3iTgL0hEPcYtP/OGAnMRghPXEh3v+/v538cN5lntZlCd8EX3r+4P+Jl2XPR6txQRGvaG4/y9e
% eGKtXY7mytn58TePpDj8gbPpOiFBTpuGptIDnafTBQYdlj+D1EkQy2Dh1ltbh2W44Yl5m2lvYi21
% 78LCHlIPoD1oQ6dw5PRPW9fEzw9cP+9TGnl6M7s+/3h2cXrB71Hoq+vZxY8XVyeXt/zWK3sQxrxS
% SpkXUUljDPVrxFJJCzKSTG+z1S8Pe6KOR8QtiXQafPP1XSY9//x8z+ehQY8XJI34GQ5ctyTl5jRu
% KRWqJ6qcqLvbT27TD9mcftVJHlGro8YgLWj7qgJKbUEFtxU1bcbZYwLfvhYHzO8S4AKN7vVoSv/N
% O7U4A3vzy92766ubk7t3k/+Ht6Hxyjf85HvABr8cIuFDDVAKCvSatgGSgDM380XXvMuliRqaPoKy
% DimYBkZKhFgKomO+DGzm3RZcMq+zec2hcaRfGMClMzzYraSMGeTix4Sa2Kgl3gTAtYng+U0EShp8
% ad91g2KpQ9SGatNrHa7OpzdT/Li6a960pWbTTxfTn9XHm/OTu+nOV16G0EZf88YH7ighW5DXE+mD
% bPbcgTkVZYrwnq172XhuA1cyx5Z3tceXSsAXb2YDNElLX8guyzzfaT+0BwXJg8J/8ouUFvatEqYp
% oqlaBK/8V5255+9Oglfzrj38LxU+pwKNyC/RmS0KlHVRUHMZXf3EzVQaFO8utYmmEicFqbJeLoFu
% 0i0lZZGH+/iHia9W202uPR536cQxJhdFVI49JfkmX7yIgwzAeU2d7Rk2r7sYUy05Ke1RD2kvalXn
% yi9eBtjb+R9QSwMEFAAAAAgAZatVXXVVwtnkEwAAUysAACUAAABpbnRlZ3JhdGlvbl9jaGVja3Mv
% b3BlcmF0b3JfYXVkaXQudHh0lVprc9vGkv2OXzH5tGIuwfdDj/hWKZJsqyJLKlHee3dTsWoIDMmx
% QIDGAJLo9f73Pd09AEFb3qqkKhYJAjP9PH26Bze3F3en9zd3qquubt6F+Pbh9Pri+l6dfjy/vA9O
% y9gWKtaFOVaDvrqJimxucjXoDSbBZbopC+WibIMf86xM494oDOXDtK1MWuTWONUbPbys7HKldBqr
% 3vhhrV/UgcGFQm30xuSu1Qlucru0qU7UWqeli3K7KRzf/2Ryu7CRLmyWquqHZ5MblWaFMhDOxJ0g
% 6HfU9cW/wsvrcPZf1/fvL2aXM3V3Mft4dR8EEC/XRZY/uG1arIyzrlOYF7W0TxBOqyhbbxJTGAiT
% 6zU+5Par4Y0d7ZktFB5Sd9NubxSsjN4onS/LNZRj+ex6k2e0kIVcscV1eipc5MbAJM94NjdulSWx
% WuTZWvW7k5EqsqDfHQw76n5l8SApqrJNYdf2q+jp97TpCsJAQYU9skWbVU7Ns2Ix1jpa2dTk24Dk
% oJ+0wnpQFOZKwk1u8anYqijRdt1Rs0LPE6MSPTeJOw6UgjWOs81xrmNbOgUH4gMckGb5WidsgkWZ
% JOEsWpV5uMlgc7XOsE2W2ojW1YX6on57o0iVE6z33X86inJT2Ce6daERNrnDrf1ej802PDqiz688
% t4Rq5mVzEMYmKfSvxadBd9DrIUR+uPN0J8DoeG8PvwU+nfB6sUmzNcILNyByOzvl19o9YqVjDoIs
% hVOfrTMNG4SIHYQkeZvvnVNwqyLXqVuYPDfxj2Kxn8l9ZD5VBR/JgTi9rWLMHSv9BrIODttq/mbc
% xZ/ozaG/cPZm2h13gqtPA9LvjO0xv2qpRZYk2bOTLQbjSfdgMFbm06ClflP8BN0YtX7rk+ptfmww
% b/027Q78F90K+2KyybATvIVYX96Q/45JO0e/DCf97uhwPOxhSZbmJDCFpl+OjvrT6fBw0h33eoeT
% PnbGHeJGfBo0lhvtL4fFRj8u1+8ND6eTMZQYj45G9XrkPV7tHhZ0RV5GRZkjLPWSfGDT2GwM/kmL
% ZKtgfuTHEyKVzI0b8C1StD6URdpRJrYD+q14zhDflFsMMZCNHenYLOx3Cv3qGkUPHi+TwqZLhEbq
% Cp0WrrNDkt21zVY5Q8hRGJaIAqksgAcGALIFIrhNorcQsd6edwmebbGCHxG0uOo6b/2nNlSMEmAu
% Nt6HgPekD4AC2m8yi61P1O3pbIaYIkvV8RtymC4QwGWigS2ITwh9cKoy9Xuryx9bZP58lal/qIN+
% iA8twpcYln4445t///Psr1b329m3X9rB84qwFje92SUlJ2Rbwd1711oK0PXlN46ATnBZkFBzIJTb
% AShnI6vOl0ZdKhi9sfIgI1LU6SVBYL6UyEUCMjJETvj4rIGZSW50DBh60jZhaLOpupuQDLR2ApQG
% YjlAoc4DQuI4W2ubineBuRA9ISCD165v7iFrmZKZaWWAHLRmpI2SzJW5gZUHHfXhFJUF/1yenV5J
% eVRvL6/PL6/fzYKAlTlWdxCKBYi2EWSKTF5I/TLYOYH8QCE4N1yiTMIUWR7X9aQdrHVBmL4MN9rm
% al7GS1O0CUJSXIHRajACvgPP4ABH9cajLJcBqqvzDP7K4mAXpeo6A0DiMXbAgiKQqo/xqKkem+bY
% 6K2jCFKVOEFuSmf8A1ylIQjSci7VCku4yuxZXqyyJce51FUulKhntSztwGUAWOycENKaJRZ5Qqqv
% 7KIgv6CgZg6O8Avu8BeamfTJkLwSGJss2RKuYy9fxKzIzb4PDaxXeyElyK1cT7WyTC0lSbINXDl3
% sH+s5ls8ifuUJXyBRUN5jn7SSYL8ZyePG04GroQUW2eo8MY9EhawPXMDO7HubdL8CXTjmRkP5XW2
% CGK7QHyTVuxpeQrIA9XJmwY5hQUQ7ioqixBIQoQFSGILWqFN+brRSzZ/OyCXkrF1bh0hiMR8nUUS
% 8eYFblIOT5f0o4R7jqvkpE7wWnzMPK6JjAKOKtIpmW9ulI5jWIaSOSsLr2UgWgKb2HfE7vAL6mS0
% F1gFeEpH3XplVaUsoIJ9iD+CyIHcbITDYDdnSVwQqI0PB52UCMx3SUmG9fjhFKKVkPcz4aQwnsAZ
% hJ6wpBTO/1w6ykssCcuty6LmXeByeaI3G1pvnmTRI7t9TNhFZqRPzQBwK53LIpDrRfLZYOk1J3xi
% tCskEAGSRtwlIQnTIsRwoRujZkhxqAvQMsnmlOJUqDzq6bZ6BHgY5JVOIvhQFvO+puwCTpo0MhLl
% 5EbkW8Ex+PP8d1IahYYHaxgFoVsAJpXYDhSO6H/lBq8tadlt6Ab5BR0olMsF4M5i10AAl1gxBXsF
% TwytgBoAF2SNEchsa1HZ1+YtsWmTLMj0k9r0k33TIyZQESNGmjRFhJm4C5pMII4K6iiRK08AdOBK
% 3JlkS2RJsVqDJwifY4xlHkD5Qbm8Q6q69EgKQWuThylg1aNm0+qvm/gObEUzgQdWrUlxpt+AfBOA
% IPuaQ+alDbhuZYg7ZG60okxAiUvIHDmxCYfPbrsGfcwF7D1aQMCOJwFesNrztAaLrVIJfsFovvCL
% alL7+hFZShQXaLappCeST4LyEZkGmmIpFa3zQUIsLFO3NkNAGGRZ0vG+kz4LVC5LqSKDBU4gqanC
% xKmDIdhvH/+0IJ5OxdGgJt1he9wdtk4QSJK4UcTIDCzeqlW2xk4pODrpJPkObIyKJl8LmHQcK25i
% fUMjvVRH/WEASEykVLQy0aOJqSn6Vd2QdXAv91hVJVDhwfmH1odP/xP2/1ew01EFge+pDpNd1GYL
% Zm+BkISI3BVQMCLvl2gFTzz8U2zCtyZJPOA8kelRgAyV0IWqFsuiqARqUmDxUr7H5ByETOx9Z5d1
% p8g8WLxVw22H1TlFvIE2cj/XsPpOM8twSfUWVtULShs9d1kCEsuJHhKtXm5ZDJgWbkw6CuQuzoxU
% 0txIQMH4VIQKS1ysYspOpJgJTCGSEk4HMm7oDI8HuHKgYENZVcGyI1KBEMQTFPisMEvQUFQebxOx
% Jkjh0BElmP7Sw15HkYHCOkWhl/YYKbakbKPYl6bOVqL4UL7KzHP6XdPLMjB/1XWtVrNtAucwcRBp
% O+rckg8gjgnJX7FaUzNBDMsD/Vbdy1pZmcRUTS3VMhvvURRHow0S/R2FIqBi6etESdC51i80MSDw
% bFAjPzXhterteatdwVQ1/bC83gnxsca1uWW+jdjSyA2wsVQenmtnJR41lYbMEZ2QAYZjCvU6j2Cu
% VvW9+4MOPc+eGENEYLKsTZkeBA73V7l6ms8t6Axgr9EAKq49ugBvKGU1rENUnjiC5akQpf/49fQn
% ocg+lKRdoXEZ9fj5Y6N+CJyCiKCACCh3mYcbjiMpwUoqS9OkVCJ18tgoPxXdDDSY8bZA9fnOBVJh
% 3ubZHJgGJmvyHDhetXEoFwxrgLgaoet+km/lkiJFDBkw7B9ND4+6w1H/aDjhlfvTw+mg3z9C4z/t
% ox+fDsa+xUZJh6dsjmIuedxEiYGE2sK+QCcm5Ojwwn/2QDIz+pYbCdagCl/zEq04sxA8CfxM3GNW
% IOUAGFZLYxDzLIbSXSdwAa1AbTJBXBwgM6TvpS7YVLx1+DKsW4UQ0AgIKEVfpq/qDjpm63Cz0sAK
% qlAuqHrRCVTkJPO01dmkmh6gM4+E9+s6wBg+PTrzQq/H9BnTXcqSGgfLDdIT2HyLECzYOjRm8a0J
% gnEIMS/Cu4/X6j8v7i7foom8v7y5DoLfDdxMuhhwfQ6F2y04dcqdUO7HoJusEEqcbOuezxAVTGQa
% Cge4DXMg2NtPSYkF11EitMdQ6HHxQIWMNdA9sXOvN/XqmqeH7iSo2/Oa9Y7q6SuIv5VVk60QGnVd
% rm8hGEmQa+qhydlphviPxICpQY7lj22EEqpDRbgRRivUwqbuRCb8OAb1Hp0AWBhCq0u3tAPkhPTm
% MQUZWmiQRbPSTxY/0FDAe+dCgz65csM4ICPkLU1q6BbZC5HGQzSa2mbPKeTNGT5B9RjeXRGjq+ni
% D9ILdt5QvMVIY8qQIivRxhNZeiCl4NtTUoIn2tU0GSYuWX4Cna06oGDBj4iPXouIRj56qObiD5rw
% 6SEv007xUijg1dWVOnt/cfbHjEc8F+dtNTzq9PrUJ2ZoEOqpqSQHtchVB0MVGyarOvcdv90NEXxV
% pNpZrfNOlzxcCutwaTb46OaqEZdaJBnPLypfd7wqPNX/URH2LcnIFMuhDKC1i9tq2jka935Qh5G0
% u+Rmrstd7W56woWem/xd29/e9bMiTrlmecb/v2kv/n16dv+9gUfDzmj4g0SjSXswHLxm4fGkLwUI
% eAfygiv90VT5Xo5ARQxfr/RTCyb7uSLi/01z9iedybT3k/DwndVuekYZXR96OBCl3ZS6W/WTDcOz
% TJO/adK2GnT6r0RsfzCtMSmkXofbieZ2LN5ooKjaZtJb+67m6vzqV4U7ilqk1830qolGncHkR3Fy
% 359J9fdl37XrPbu7oT9PD9J6WkAEtNA56VAtls1lXk02bhbREJ1Mr5rCFATeNPOQ5pg1mf594w47
% h+Mf1ZlOGixcSNwakAP6gvhujyc9Jt74Nhy0R0fDumepgWDcZitJEDBnRWey162dqMOJt3DtIXZM
% 04fVchH1NxTrqIltHpYJG26cDbWrc6cGsydLVfwoMrWNXvf2q9k87UzHox+Drz0a9aViEcWvSF+B
% Bo0sNFaEj2FN+l5DSvSuzCKpMQ+jDFhARdPU03luKLqoPGY/iWoZpNAIyPoQLVYoZMuV78sPT9RX
% VLZQGsamWUhYJ/N+6ohyHsbyjJDX8dMtbichNfVv1dwGKe7JHbNfJ9SP59u+WfY9cO7zxTeC9KyI
% kVXnszB88yBkk5sna575QJUCD9XdqZ8cuBbPFpxPSugmXlzpe/Pv42CCLYGf3IPQbzQ4peOXyvtt
% xTeiducptRt87kBxxX6YZy909srUNBB/h6QCzVEa56uADSJTRHSE+TDx9MLvJiMy3J5nGdX0EZjU
% zce7s4uZOr0+Vxfnl/c3d5enV+rq8sPl/SwIbnOw+1wOfVAT8DekKgvpfQ+jOFap2q+KYuOOu12d
% v9inTpYvu+ixu72j3qTTHwwP1cGZ3pRFFoZXdomyUYThnY1WaGnXaxrgPo1aP11kAJjt9IajAVa5
% R4+Cu8c/v3tNEb5ZdXuT3riHZDh4SycSMPvXMLzNUX+jyGKFQSuQcR4Yo45BfoWEytSHxlD1aFGo
% hkRZDs7xRHFtpVZSnGUZ7gZfgufhU3esQCvSwxHZHZi4+9JvfhkwZzQvK03zAIT+a+faztCGzOkQ
% b9QqGd9s0ssJ6hZRHYYzhMyKD8ivbBj+90qTjcus9ln9JsH+ySI3M1LHt+JKZBUfQTI13nL20WMI
% obWlFskPYiindr0urgAg1zxBbtrM6S2hoqZ5Y8gzi/qY2D9SzQCkK2r7AWjQWCPhedfeItx5+Oe1
% 8/K88CLzxKzbLPXOlLxa6ZqC1Q04mn46Zf6Kfk+Gw2xZ0RvEmqYcsVhHXoyoxpd1xkuXsarGGwsy
% CQM/Q+Yy2fqkrecaht+QSAs6BCPbNcc6dJhrquM5mYzRCQ6uZgkd/4Ir1h4gyIaQPon92egaJVkq
% rtB7Hgxwh8Ad1hxNRVWHyIj7vbiMLoOoavmqUbflQyVB893MuprVMZxZP7bl84u2NMp1nIX0gsz2
% u5cq+FWOXT8q49baqTKhFiIcoLuUeDcpn3jF/vB1KKNxmjSS6Xm6x8ea9SGfos4d0IQNgnqMtbEm
% MpVXvj95/w/H0xIemIauXO8cDSOPO+r0HPCIdhbwOLu/+3h2//EOH7vq/cVsdnl6rX6/uzx/d+Hf
% KXr1bB870KkWn/3Nc0vUmlJ7Py/rSS2/UUPMFf9nEUdPNYeRKQPqRGMH1s35g/1Fwo2yvxjwcOME
% KQ9fcBrxjHyxCOuBRfUik4+wZ6MfqxVB4WgeW1BABI1XZ8hNcmio/kXEW/3x8PnN+cPnTwdgg4OW
% ev/wWTW+UqE21LTulghqpuAnaLWf6hNR6ZTIq5rfOpJSh8pHHgpxmQa1Lqi4iXVefD+J1ghOLaMa
% 3L8m49Cz/na4A3QY66V0uiKjlwBGsjy5wTM8M2rwE1KVcusZmt5CveqdgoNE/fJGfW6pu4cEEXF7
% cNqSlzKaHNmblms62TeUaKQOoSE+lq4PShpPB3iaGnkiSCEftHE47D3aD/EwWbJ0/pylOflFvTyY
% dvu9FlnigR38J2327Rv89u3bw9tPA/WPA16jxdf+/Pjw+S/54a9OsDvJrKIPZfHRKJihfUem8L/z
% 2ExeDMFK/+x5ZYSkOB6TBXIYtCWzEw+UaGO7+1iA+wWo6R0VG9OUvOi96Qk00+wTYby2RT3i8KdK
% VUycNOd29MMBvVx1XwcGLhsj07KVJjJezS53EdiwXUe9x5cnIim4J2BIRH2WelRL7dQBM21qmP00
% sdWuDtA55mn2TuRCzjaCoiGOHJgcNwRp+m7nNvW9Yxgk+EUcoD5NLpmkgN6u58T/UWdy+wIl6SiH
% F2O191IYyxLA1YRWgOmBYYgGSvvoxL3p/lsL4WdT52mlGKpZRG+ZNZxCnVKVWruTmZTmWbk/dxQe
% 7ieejt+v4dmslgipG3q5e7FDdH9cUR1R84y0TX1CVdbkuhxL0wk4Nwq0M3e8MHohW+yCW6KvmhW5
% LyWdZ3J70t4FJXu8cqKE+kCIDJfb/VDxMwl6XSjmlbrV77uOCQCzm4BUL0w1+/wTTvp6WGfTnzjO
% 9wbURVV+96t4euPfW6GKIJPc3dy+8SamllfOKP98fZazheD/AFBLAwQUAAAACADIqlVdUVeMilYR
% AAA7IwAAJAAAAGludGVncmF0aW9uX2NoZWNrcy9zcGVjaWFsX2F1ZGl0LnR4dI1aW1fbSBJ+16/o
% fdnYiSXrbhviOcchkDBLgMVkZrIP4bSllq1FN3QxOL9+v+qWZEPCnH1IYkl9qa7LV19VZ3l9enK+
% uNDPvl6e3J5fXbLF5Uf2afF1uTxfXLIPN4vLk89s8fXj+a22aMK4zssjNi7zvB7XeVPeVYUIYp4Y
% 7COvxRGzTdvXLVO3LUPTzs4vTpfs5nTxkZ1fsrOvFxfajXXqHDGamfFUZPWdad2ZE6zXZKFpjU3n
% bhOvN+OCF6I0avFEE9y/m+DePb2cYf/NDPuXMyZ/I9TkF0ItkoRFGM6CvCxFVeRZGGdrthVlHO2M
% YseiOBEVexSlYHFGKqpFyFYiyvFCPImgqeM8M7SrMl7HGU9YhcWCbkaW1yzNwziKRQgtXv1xerO4
% uGC/f/346cvp5a12mbOU1xuBv+IAk9e8YI+8YnEI2eUsbMowQMnIqypeZ3ipNmEFjBdVhnaLAWue
% ppzxuc9wjCbBvk1VSwFWcmCak+BYm7NMPOJYJT5jaJ7IExyxG1vjSSl4uMNraCOkoUWzSuKACV4m
% sSj1kNMiKc+aKijjApvwe1IXSVhB4SxIeJwqgT7xBuLyjBV5FdfxVuhYHPpRXhb/4LRtNYLiRRTF
% QYwTP8YVaVk8NBhQx6IaaeSeJD7sJV8wnoUkX4x30Jie5mTmZ5PUyUsoFboL46qOs6A2SNdSOgjL
% a6XTBsaPsyjO4lr8JEeQZ/+FtRsyfAVFPkLryVba8fSvxcktgzHPz85PFjLUbk6XXy9ul/Lk1a6q
% Rcqud/Umz6TAK7hfAs2VDcwKNbWfaGPoItmxhAf3+L7cpde7Y8ajWpTq6PC5GiPIEeihzg+d+zEv
% 78eFXOsuFAW0Q8eqN6UgIdJVDttpylIVK2GKx7jeqPNff7v9fHV5vbj9bDDILPUuhWFxWuRlzSIe
% Jzi7cmRNZNu4zDPalScj6VYyROC1ZMZ+NJRDuKBUC8mzkJehnsSrkpc7xAsPuokC4HO9WC7ZgBb7
% L5lM1x+aOLgfGtrHuCoSviNPF1uR6VBHygJyAbmjGDFrphc59CEDo97osGCJD23k1Dup9htt/wh5
% CsQP7Nd7pn4pHmsIf7BwH7dZJtacvBYblEcaG9yPyuF84Izc4RE2fzYH0oxmjkWzfogyJzOJtSgP
% Pao61li14bbnM5+v7JnpmUFgmTPfFaZpcSdaTWxfmEG4Es7U9gOfR37ohDw0+VS4vjXzZ7OV7/i2
% N4MszsiDFI73UgpnZM2mr2wr7JW9Mn0rCPyJY4fCx8+J44a+azsrz7U5MM30o8j0ghXeBL4nxNQ3
% fYeLycS1Tbmtj239yYttZ/5oYluvbBvZUehFq5npToXv2HY0meJgpjNzPXtiTcIgmKwm3mTmTUM3
% mnAxnQahK2bcCwQ0MfOwrSt1/vNp7dHEd17ZdhbZAofz3MgJAxGaUPg0nDmePzNDYdn+zJkEZjR1
% uBd51tRxVmK14m7gmP7EDNzAwrae3NazXmzrmiNn+uppQ4+Hk5UFQ5l+OLMi0wwDc2qFq2hl+djQ
% 81zurSa+HQWu43k2lOxzaGMSOBaMgm19ue3PSrZH/sx7ZVtrEvjCsiKcI4q80HS54DDcCofhAaw6
% g1pnq9CybDcw7ciJfN8PLRFOTMu23DCSqAVABbqxUKwBH2Py4bLIExXeK8qfBIMCEZQRTBwR1AA1
% OKKlBVAZ3dqzYELsyMjLFGQgRRSUTfOIZRJ2ulDRC+AfJZJUAMxCFmxEcF9pAS9LFcmVKHh5IEtl
% sDMJMSoNBTktTLmpagJk34rwZ6fAyG1xhuRQANRBYwcWlDNkWpEbUI5Cln2mf6OlIi8W6idIEbtV
% X0yUjEROTJGWkVnqMbbNU7AFMDAW8CwQSavn+7kzMgxj5JnHL2TVkPMAp+sm4eWB3MfQIJBfZPSm
% nVJvYmDuz4BYM8vTBg+j+yGQEOn/mJJjC7MKRQHjAWRqc/MxC5oVzoNR0DLJm9Wabk+mlj1zndkx
% jPJUEzmgXZHD6se8z/VM5npSiyQuQEb5glIa9nVN7YUCe47wnB4wqJAmkLwi1JV4LzP1octBHUHS
% SAqXzS3lK7yEogy2VK4iEl4gCzDL8KY4AbI8nEk7w7LkK2Oy85hs1vq38kS4VQEHk1SDUuxjGde1
% yFiTxXDxtCdi2vXN1dWZ4tjs8ur2dNnS5J4MtSnneXjFRLUSsYYKUjoEPEzPoMKwCSS7pFDRVGC2
% oWh/Lwf3uj18Z0slb3nSqKXwQbdkVEpCKw3TyKxH4aaoYaXto1aH98DwL8kPUnctj1upyH1OlR70
% JF9DcwHfwrP6MLsRab7lqwRYQqm5YqYUzhoxmRYR9YeOH4IGbGWK1aqHsh649lD3Fb+TdBIxCKTI
% utBXWFHipJU8Xg2HooAXTwXCI64p3K8xEAiC1IsVQNBTtuGS+KwFM9/z9/PpsaRR0p6ihFqwALZ4
% +g2usmNkezqMDPVL2JX8ECNUGCznOv/uvCvswdOQf7ffFRb9eFeY+GekmBUhg2Rf2iFOk09C2ApO
% OWJrkTUgqwmRbGgLyEOHlUaFg3OyUF3Gyuzaec3CXFTSAlikN2gJsouAgPtleykXb0ANO9rbxyHc
% UyOxuHKAP+frQfRmvH4z1COocd0ol4aKqggKkcLAqR8a+E+ZVyTz3pEUr78HFc7kMatNnsiswN/O
% Bz9EzQcwoaxDhuOBo3dvbPWKGDgOCQUAmIIduQPtBneI0yYlw/NkR3iCSgWnHzGAmqJ4MDHkk+/g
% HRoQcC30orc1NCbDpmrdHkPBvzGvZbviGc6VQGEueaoE5jPI0qH+Jo5qtUivPvJGpQ+oPRTwfwql
% DwpWdsAYU8P5X+4p3R3gKKTUxwTrTA4D9aXqEUw4LQDJX1xllYffLMX2W/FoMBw9DlU4PPQFD+E4
% EkZvMNaWP/tw1SKUeuzhvaWKCOAKqk+MpFFrgdHkXaxzG67g5UxxaJyg6KBayqOUgf3yFbxSgZF0
% MnzTn+UKptKNtBF5K0cENBSvwHlKzTRLGlOiS6teOs3PcH5QeBnaYm8KoOWO0AOWkom0oIqUtNHW
% uG0ipXh9viBsfXn6J1t+u7z9fLo8XzJqZhz9VEe1meeu2mUE/XEl+wNncmu+FhTGpArLKuoRQEPo
% iEiEPFJPQIhXim2MurrtDESyFyDZjVKLoX2tiPlk8KOuLpK1PyPqQ+uPD2osxUcScHKgHQQ6IuGO
% QCaENCX1F1QL4IiQd9Q/2IcPSDz7hwk9SJcnfIBKLINdFRQ4Sds6QJThVHCTLwM+nAPr757Y2YCP
% nobkjuskX8lKVcE+KrATa6TxiCICh05QHvO32E6C137YuI9qvsopnt4qlyabndjESfhbQzuVZKAP
% c5RbkOHtO1S0w3lBvyz5W3fYR5HUAF+Gx+/2eDBlN2+H764G9OgAYdiXN29o/HC+HzseuDQKEL2J
% YSg1Qh+yOTOxc7aF7IXoIw+f52d3nOv46wmTz+6enhQ1rtrDZaCXhrbc8LJgvyPCQQPa3oM8dRir
% dkrXeOhhrALKw9zv5+rECGFV7ncU44DIKf1FcZnCmzdv3kgevRnUw7k1RuYdWO+scT0caYRteAbB
% qTepTiQYPytE4z63tgmV8EIKdhgc8AKbgr9k5dwcWSN7BB9BvQ1gw8o1/03l74ffzNFBywWB+Eip
% AvhdNXGNfL/TemSRLnKQlHre09JRWRKkFEcdmSBcBRRUMXiD9hnDw7yMog7YIp7GSDxASSo+CFFD
% laRuJhDaGT28n1MtVuZg0BC75WPJbqRJIO7IwR5FJMj3TLxqD2i9QERtKU39DKgiRq27zzy7Fwll
% rbykwyTw/IoQ/YQ3wWan6x8QELXMUyrTKT38gReiTCkBHix6LI9SCoCHLKtkAcQrajkRcelUpO2l
% Bw3Ny9pg1LbsVQrRPrLVXTbHH/w7yN4h68p/rSFiRVGC/SI/KbmC/ypQ2cacbLNXW1NtoNBHCVZK
% MiIBsEW8F49dXH3ST64u/zj96/z2G0lz4Cz92iov5G2bVHs5+WTxB00mGlSyC3imozzz4R9z8Edi
% MIngkte3Fggp96uipIK2tuJnY2sc1VNR5wQ/lKJr3iEreFa8htc0YIciUX2eq8G/v2dDRcVUWlaJ
% X5MMhaqZGtknkHH6hHPVzJS1Wgq+3QnT0rlus7ZIldWBpmBv31dSZqHjUpLtKoq23btPQlRTnH9Z
% 3Hxjy6uvNyen7OTz6cm/lmyw788PtS/BJU9BiHR9ydccOfi8Uy/gQe+ZOlBJSLzb55h9sq1G2iLc
% wrVQ6RjsC1KTwVyX9rFMoKelW94xRP0r3h6ZU3NqWKbvbW1D29R1UR2Nx7x8irdGXq7HfFWN+yG/
% +L6p02R8sIb2p1hBHX1qgzTZ1KR0CuCRDzN6sNTDTH6xDx8cQzvp8fJk37f1jJZ3veww/7KOgedq
% PV4FedKkGUrP6zJX/gTs9Ay3ywX37+f2MXmJLDN8fJAlYqw8QFPtddWJXQabptT78OvDVlbHqm9K
% QdOW4MAISXK0A+wh+FfMvCsPhOw3KMzVoyZT1KwLNjjNh5IQJ9P1kw11PV/zCLn0AbO6xk/OolIt
% tNMOwheKCOAdX5bMch3pFh7cwrPcie5Z3rTzDcs1PcN3JtOt9Zpv9ENe842DNV73DevQN+xDD3hA
% TrjmFV0PKDO+6OV31xxtP599UkDNrs+0IMkl/LY4RKMjVJxi3zTe9ywUvkvU7ixQoV4jnNN4A7BQ
% lzpIcKkQsm8QH+paf94D2OyKXLYcKOT/1YAKIZh/8OrHiH2BS9R5Fkv631a2yn2eXXt0XiLpnNaJ
% NJKh3N8wnR/O+KcM9raHXzFrIs3qUrS7M8T7bG9BeQ1miEToRFoNcI1xS32rMU6jWxPdsowijLQw
% j48s05hYnnPwhT2z40Hsmr+uKul07a1TX/izAVjKUMKlrDHkRRZ1p6iuojbEhyZZ5TwADP6Hqr0D
% zbGM122c1fJuihjvXkW/y1ogIQR4piCS4pmKCHRHHem3JyN1QVZtRGjb7HeeNWAhNGavuCTO7o2q
% KOEmojzU2xhKsqypP64sx/dN3bQ93XRc+vt1p7cPnR6wtxTyBOBtCWUfb+rovulQFikJG8SWbmv2
% LRLGq/u2K5vtC26Vn1rVtNV/W5Z3CqN7xYcGmZT4Arx1oFujFkZ3cBd1YZU1qeg4YCioabii4IiR
% Bw87GQA2sb9/o14Xp5s5RFd3nQRTLnnCf/BSWuC6CeOULc5HbMEuOCr5QOjXe3Q84UnQJI0611JV
% cfpZa9uf7blUd6eAMzDOFMAG56Gqg21NB7alK06wRPJGaPFW/MVqaE4caTLfmr5uO3vjrkEym5W0
% akFC6lAlAU37xGM9FGkuOyEJf+ynlfzRUFObSpQAJ2od/L+rjIk2jCUM7Mat7NWYhB+jvr26+XZH
% v2XxKt2I+gq/CjzEpFRWRBmEKqNHjEYplFUyK4mQ+sPU+NGCpkxUY50aZsB817JnPkivvJmo+FZd
% LHcFNHw9L3shqGjloVKiclKY8fpaN017RvUfvXJde9q9ddzure97s74fGJdaU1TU9EsVAuv7Gmff
% yOLyol9t4zoIKqxsSjK3O4gCTbZL9pfWUfxEvTe4uHJ+HZ4r0bK9Ymwht7solQ0ieS+uychR18ot
% weXd9WbIpB7IQaVL5eWLCk/ZhC6FZVPW0K7VXTvcIAWJlZPo/wLQqj2/RLa6zPcrk8dL6WWAwmub
% JNRWgv6jgCjrdjQYbky9EYrbmrj7wf093WxLq8qlqFXft7KXJ1fXp+zqjF2c357eLG6/3rTcU91x
% IxVIGgJZUiq9nv2/B1nSttuLpw1Sp7TTY14m4WMcCi0DZCdUMMsb+baMgkN2tLfrVLdnov5OiyjK
% 14D8sKN0AbygWi9uD0dVcJwctKLa2wGqpFGS532RitzR5VJD+x9QSwMEFAAAAAgAUatVXYAmc0Mc
% EwAAHSsAACcAAABpbnRlZ3JhdGlvbl9jaGVja3Mvc3RydWN0dXJhbF9hdWRpdC50eHSVWm1X20jS
% /e5f0ftp8WDZlvwCxmHPIYQM2WQIC8wzyczZcNpS227Qi6MXwHPmxz+3qrslmTDZmZyEYFnql6pb
% t25V6/rm6ufTm5+vTj6Iy7Orn04uzi5uxMnPb97h58Ubcf354ub87PrddeeNLNWRCIbB1POHXuB3
% Pj6mKj8SZVblt0WZV2FZ5TLudK7DbKM6Hv3pXCkZCZniX7xSi1zqUMbxVsgq0qWKRLlWolAPKhWy
% KPQqxaWi2mxijV9u1CdRYOxQFWKh4uyx37nB7WGWhnFV6Cy1l/mKylMeLNLFJpZbPC7zVZWotCx4
% enync1GU2EPUWW83GS4UqugLDLkVUSbSrBQYpdTLrVjrosxyWqrY5Bq/lVuR5Zgn2cSqVKkqCpEt
% acyOeioxNW6MsZ9cwgKqLy4y3K5Xmq6bHYiljpV4lIUI1zJdqajf6fh9ceWfTYwBU0lrvR36t8OD
% QZ5VaTT0B8PJbSKfBhu5UXm/VE8dIcRPUqciV0UVl0ciwmryBPOkJe7SuU5XMGe89E7XWYylwpjL
% pQ41hsYzK9isJ4qvlcxhBIwVZnmuYlniushhyAXNK3Lx6ljktz+IYzHsT0fDg2A4Gs0O+v1+z3gy
% kptSP6j27DTaBl4RiSzDNS3D2BI2LFWfvj1dq/BeRdjtY+Y9qLAkg+aKJx/INNVrbVcCl6rcrGUu
% wm0Iw6mnTZZiExrwwlhAiLuDVwTw6bAkt8zFMpchDQPTu7UMsijyClgDQMGtMGBZzGmcAh6WOSZc
% Sfh4nehQ6FTBPjF5nEbOZaSxqyrVXytlJoi0XPHw9DRG593B42SNkDCNadcKPxg/K7nBmAbnDL62
% S3A1y1XCI7xLNllewpQYF+EDOwGFAn8JmMsqjsXlxxtCIQyB9QABvGkaUGFfkbeUZFMa6rXcyNLz
% rmEe4Ea8Qwg03sWIsLEFgUiyqIo1bzWpipInWyhGugNJiijKXdhiRUuhSwEPKd4UfAZXyEWRxVWp
% PB6vas0HnAfAeXA2+lOcB4Ph6HatV+sW0Ef8zOR7z7wQG6dZkmBzIe4GEdVR8p+9ky5hek3/X9KP
% Jewo84XGbfnWxvWTuLx+I0JYsSpoMHgQ21YFQA/jZzlsCbo4/fgz6PHm/N21+HhxesaWA4zBZuL1
% x5tzzFkCX7DsJs+yZcGupZ1gkCIDEkBdEfhss5aF8haAH6jLcB68LxZxFt4Xu/FiIqCOK4SCTJn9
% 5ojEO0mbde70VLRSLZTTOE+3ev/pdk/v+91Xxz7CI0uEhOtypRCGOSyHVUUVrQKoWc/FyU+e9+NP
% mDAHgc0FjADCihWNxR5li2H1mH+teOd0BTNmWN0qrmiROqJgRQyZDdMlAgvIMpc0UCAiX0QBPPPa
% 7+JH0BVEJXPjgqp0KwOX67KkxxfEwOU6M6H3Yy4Tppw8u1Mc7obJiVwQHEAhqIedgAXoDZ4wlsXC
% UwtxrBQ3LFQoq4I3V5pEUC1i5YXYaGmCds1x4oK+vxPpCG7j4BEcbCLBRCVCBS4FYBQ5m8YxuLJm
% 8oijW4lMPIJ9aCAX5z2xqLBKJCmXPwycGG5IcTBJWAo4nawsHRlLHSPyQD06RdiN+7SuPw+70Qth
% 1+DO+gurZsJeVwsKmqSKJQfH5bG//2n/8/6nz17w6/4f+R9fAlobffUf+irY/xzg5+cAXwf7v34J
% 2E7X1QJLLyum+U/HuOvXfdkTn48/028LsdUqjopneXyTFdrucAmz0ziRopDl6xiIVxlv++JDRowb
% gomStM31iIG7/e3tHeHfsC+2Bs1BQ5HlM9gawDPokxFW8Kn3mW7mDGCJthB1YigzWt5dpimryvTe
% AyUzgmCoxdYjDFq40QD1+m22sTH+IpBoAe+DXtLGBgi7gN3AwhFFAWFtziZKs9SjB2Ledh1y1g3G
% d/Waw5gEFudjGhSgRt4AxwO0DWfplHzMd/By3gKOHu3Q0J8H6eOesoYn4uRdZpAEiYy9l/I/Ixvg
% hDYyiQOEozYqpRUjUnKVV40XO50JA3f8PeCOb5/+HLlEDh5kFtIucacOPQQPGA7WKGiRpephGXic
% LH8qq3C9RbIM148y/53GeZBEeWXP8aJHvCjcGMyBvYaP46wojC4CM2MyzQgt9JNniHsg/nV8YH9n
% TGLob73vNELLaiYNEKkpmUMQ55jAYQLis0lfp+PB6YRgEmItJHxLeIO0KJJcDJyD/LaOpQyxsLSg
% JLn9Z7HLIrVdKJN4tLfaaI5aOAAtBc2dcsaFUObMGgmpBCmiiqzh1N8F4SrWv6smoPCrsUqjxhmW
% NJysyswqKABkg0dU/mDkCEKAHExJOXNgLilnCMoZlLBNJqgFONCBAcP7NHuMaVcFb39uKV+k6lEU
% 25QWAKNUtArnSPirQNxrOA1VDJA5ZWROvofMP1PsOt1giaQ76smOaJ2k/7CO5IhMOxcfVJJIeC05
% ovw0t5zB8Qw/maj7emSZhEB3muVQ+kbF5Ecu3PvPmRzwrnIBra6TKjEGbHJnXSr1qEgAxedertjo
% aclzmlohfSCGoRKCyNNCmoq2Nmp7gsS2+Y5FDoZcxmaenknZ7EyzF8JxTEUWAbaVrbHQLEdYLfAJ
% Q4OjGI9mXF5Ne0aZLFhOg4a9DbNyjX/Hir2GsZE6ksrsq+c49xmiejvKxmozEtnafMvWrFkFUZgW
% S5XzaM/UGkeQKSRMnsBXxX1DywQIfqQATYy/KSUoLzCav1kEPbnQEQ3Otmn2xwrv5OKzoSEmaJ6h
% x5GVpSTj01rz8VdUIyCSVeEEKnkYW6dUBNgfAPbj7yiJ8ctKosF9ceRKnTYaCLGMXY2JY5UWgxfh
% a58glUxYazwaknaO4xpJMQlQa0wufbhApwom85x1WCEPLHStTR3/aAJclt3zWmzEmNKA2wg9WxQ8
% c6xbM5UJ3AogueqtYJmNWKFkzOVOiLn0TKnAhAFYF27eTSHkOJQuMlVZVXgtqssWptHiMvSLYKkt
% 5nEdl9KDzv66NNoTENhQ8c7Z6t3Fm7O37y7e3ZyJE1NQmx5MUcPFVNwW9019TsNSw+Ot1a00WA6L
% aAQ5BxbcckL54BGOosWBYl3usNpGkeYIWXwYwe3o0QkiqtdoTfTrFT50Id2zBScGJJ8LFMUEcqzy
% m8KBxR2raFsrbIg1eFcW4ButQiz/J/nkZUsWbrHVkET8XBEvCChJwh2bd41ugatUyELWNLqe/+mc
% wCKms2UVPCakoiQvjAYyqqeoQmSsgnT/1iwXYpQiBubotMLtMcvvB02P7Tak4CgG+dBXw0k/zlZ/
% 6eZADUd/5+a/PvLo74yMm8d/5+a/voyxW0aHSh+GDUlJGFNT7yimLFnFETOh5oaLuN4ml6gdrsgd
% nP9NwkS5hVhFAJf/LDqkAPA01YHW/yHBiKu2y8835x8vLk9uzo8HtK5ig5J8UIQ5McUgPPTD2eLQ
% nyxDf/B8B5stao/0FmMC7K+p2MPMiEsodcBNfKTg+hY9WKIV929tq0uY7pWC+ul3eL+OR6gOtSq/
% EEEwGeIjStJILOOMiWpHLLMVTf4ATXTMJK1Ma++YN0/bR5xsK6mfRwq0MPUqr2ZioV73PvzDOmOb
% dsGglZLtkD3hT31O6LattdPTqm8irfFcZbDXl+RxZ6cWbXXMk4Yn4VKPCxvbKSpYmwrDlhhGxluY
% 1JbeJn2yToWkBD2ADy52tGOkYpA5dc28l/98D8IN7VH2/CAXKkbWxNUjc9NRW520LqMoJkcXlrYX
% hY40qZaYRzBd5jOXVQjUXt2/XFTIhaWgvEiPb6G6nqh/DdI+IsLcu+4+6/DYxXInkqnTyOi9q679
% 2DCuYX+nwU0h7fJ4k5VsZ54GEf87Udfix8z6S7cpgjx3T9oqNLJHz9B5s16qFci5OxWom94dSoAZ
% 8gzlhvg6PPYHwch2gY0qocxwJH7/4fWXSK5WwM7vc6iCJ8p0bruDZqsGeo2KMJmF12/33up0GLRi
% MBbpGNCQEO6LlVy2tublfOrBIOS2Ct9tktarYyEfUGzLRayaXrXtrLTc6VpNdZX5jWZuIpYqkmUJ
% k+TUdEjVysgWXl7dGeNGL0VELWIsNj0nkF0HidH1RkHf7J30zrtQvpRib+/ctdu73vlv/3d799+u
% q7pv9+J/HN91SbLsvb6Nu3V3wgQnFYWWrRoJQPJGpiZV25Yp27K1AxrFJXxqGRtdJZC6s8diR4M5
% Cc36kNHswMUyewNv59z9gLhvnV+AaM6c1Cf9IVfMxGUGkxoVs9NAap/H8FCPlA+UDNe0wNRAyhIB
% dcYI0gL01zpeaLbUHIo46UPNlRHUjn5SdFyQ6CY9tNo57J0fUX2Ke+SCippoJWc+CFbxfuJaW7A6
% 1x2u449AmUxNK8hspvXs6fiF20cUi5LLMUxunIY83MYUDQNivd/3G4npprCJQby+Fuss5tYYWSGh
% zdVq93GtY1VX+m1qoMCyhxZ8WAEAgtbTiJ2jU84f5giozEqyZnNqZvaBJA3gTu9p21+Cb58FjGwx
% SMPLGKkx2rKzaKpgMOqZLl27P24h8fWVoR1KM30y+YPOYgYOM6trB7WOXAmDlgDrKuRB1SefLe5t
% 1dGtwsfUEh9vzo7E+7F4rDUSt11oV3WvkNXB7pA1PVOLse6ruOYgZcUFMYsC2ri98n7C7S/crhNy
% byIjyjltKNIsbB5W6RkLobYQoWOwZevItT6mNX3vtk43p2OG25mkHc1ZceeO+CzpE0mbCKDaDdlf
% ldLxFvbsQS1oSg1shzq/XR1AdCUyv2duzppOIw+FIfaeusenv6kve8Giu0f/ya7nd/cD+nVP7i+6
% 3b0933vCR/7wxQvw9X/ZLb9QSpVAhB8c9hbHk8Fh7/T4YDDpmfM3la4IE3Y+9vj7SwPPOl9yuoGV
% HgD/aNc8hjZv5DbO8sFKZVCRtGvDXkZ1cn7Eyg55Cd1X/mzgD3t0ZTIYd18dDAL+4A+mY2yKsDsd
% 8cpPSmKF0RGv5tXxaOoPxoeT0fAVD1Q3qkv56ng28w8ORofTwWQ4PJz6wQSz0DTDV0FrqPGRudsf
% jg4PppMJbh/PxvZ23G0OUtNSrfK6K3ClxL8pw5QI0nMaUaw4NJAU9uBNEPwP+GoQDIfdnkvL0sz2
% Z3d2jXwktCZ0BG7P9AnLbGGOAaBNc3eGO+PMOqTczWsDGJOr4BUystwIjLLOIqMwJXckE0Vj6iIx
% /fCTiET7ObhcY8RFromC9loycG2/YhXYXIaQUhQZd0TeaZfReE5saY/rnCji9zNimzgc+6IUIDKd
% U0McxamS9/YG8xRvMrUvVnB38OmZ+hJUpZMiYTbba84Cf1d5VlgrMgA5+X0/dZMxmTopkzSvA3Ai
% wuXCdIBsQkYJ2AyCJIz4nBy5Liti/gBAaMxp3iMoLRMQONvCtK14MC6UZuJhAWJFtxtJZ04MkItX
% VlS8CehWrOCEoNc/YH3DFvzt8fbujz/O6cfbL8E+gh4Xunzlt58hefiyCXxc4gZJU8uV7b52pHNb
% 8MD4d9WKM63Ngm/ffTp705ijloEaXsqBQpPNOA32xWV9wFY73nZUakhcQpLxEcLt3b+GzjjIABmd
% 6NA1e/o6PB62ncOHuvzGQaLLkloorDXNOagOSwsbyrSUw6npwZOypeh8wjZraEuu2xpl5MNe++yu
% SOghp6eAO4IhaG5pBuKqh5O2G5L2zufxDPCeOdoGi9JzRJvUho6g9nCHORApbLOxTqhVqikpwGRI
% GIRPb5krdv1L0MI/sORQ/JvD0KPznjBT9WsgrF+yZusmskzkmtMoaHlzZoR5yYP0vlGrC0dd5GaX
% bsmaQ5SaCXnJ/YSdzPfsEC603daFbUO0oc72MZzj9t9UdnVZOrBy34mvOZZVsd42x5sdUxtQurHV
% OeEtRdGiw+Loe7WxyUL9zfYv9ICaOpraQKemP2FEoTnDZOOIXyaeB1Vq+g9zu6yUigD9oLkeqJU4
% i/qO25x7Lcm+MuLUu7fT1p83NVCrd8v1WWcSzAbj4cxIZX+KXDnf7fa2OxkvNn7nHebrtO1BYRIe
% v/SCgA2rpi3uDrlq4dOcQnZOKAhcEyd1BShbi/zf8LBVUax+DBVSqO4svLMbbAOHYM8eTzi2Qg1E
% /bR6XfWRo31fiBqnrgX2Nstbqr4uUKgMPOkej2azyejAD2bBAeTDwB8fjkfj4XR2MJ4dTHsdW4wV
% xbE/nvnT6Ww6mwz94cF0NggOx7PZaOyPZj6GmAW9jk5BaFz9twLr2D88DEbDYDod+7MxNFdwQLLj
% 0B+NgvEhIH3FJ6J8HoAcD+9+0/LZiVl+s6P5uNPT6tnTgljYd/d6TXejY96Zo9zoWZVMfsTeoH17
% pIZfOlWrK1sjAgtKDjLtNAfHfOplOuV8egw9Tka/Ql7Dgx8yRe93omKDLMlSQvPW4YlfnukYr/B7
% Uk/uTTWUPyrPnUA2S6hbDeKMAdbS7VS0Ajr8FhVbv36fkmxEm4m48CHGph4E1SKuPYs8vtCLWGfI
% xJu1fad03uHenXvZtNWeS5WK7IGGzSek401nyr7/YHut9Zuj/c7/A1BLAwQUAAAACADUqlVdUm8u
% sEQSAAAVJgAAKAAAAGludGVncmF0aW9uX2NoZWNrcy90b3VybmFtZW50c19hdWRpdC50eHSVWl1z
% 2ziyfeevwNuVyqJkSXac2PFUybYy8a4j+8pK7s7d2rggEpIQkwCHIG1ran78Pd0AKSkzL3dra2JR
% INAfp7tPN7S4/zqfTb5MZwuxmE+u/ykmX29uF2IyuxGPv80Wn6ePt4/RpE51JVJZqXMxGor7pLJL
% VYrR8ehd9GjrMlHCqepcVPjbyFyZ6ul4+HR8NihtbdLjYb/v/zgbHI+ecpXqOh8UslBlv1JvYXeX
% 2ALbyyzDXi/KiFxWG4X/6ERm+GBql5S6qFxP4LkuhTYbVepKpSJVhTKpMsk2SjZSm56wpV5rg/dW
% 2mCJeMHKlVYlXpYmpQ2ETHGqtrTIbQ2eOO3oC1uqvB/dNxs4r95KZ8qJV1UqYWwloAIO7kfRzfTu
% 9tt0Prm6m0Z72r/a8nmw++ye2iNY48UGZ+H/Ujj9FhdyrcSdXKh/iVUp1/SCgLVgYGnEcFhUQpaw
% QqYG1qhYm2QT57KEeFFRqhetXvviFm84CJjJpcqcwPOVfoNlcOw5iXHOSzbSQXqR2qTmQ7BM5stM
% 9SGQEoUs6eGLzfAtSafeCpWQfSsrXF0U2Va4CuaTZdoYSijzoktrWMleJHNHXhvgX7fNl4PNFj6G
% LN7qunKNOfEMpjSJcueRf3Q+H09HPTE/m47OM5Wfk4X3PmbWYf+lXmbarktZbLbiWW2dX3B198AH
% 8Idf+9GDN4tIbF7AbzBDneAot6oz6EDaNN54aixYpCvRgTMEOcN1+9HMCgvU0CtAk/cHf1jatwYJ
% pSpsSQbShiFVHJwrCV396JFNlsF17ffkB20q2IzxB5lqUypnsxfs5d0QBzfsDAWwXd0vFvdf4rvb
% 2VR8mSA28Z/b68mdj1kS2dVLclGlX+g0a1cC9of4rzhxRSHYiOpDbGqqciuOR41bOCZdfw+fB0GY
% WAObVRRoQtbsIAoGbZoABKQUYBvUtGUkQ+QlbAwBebBzheik7Zog1dW2L65IOqhPtqsrXu486hSB
% K+KQzbLYlhQYrJtjpyfScERCb4DDbShotAlx38ZzNGfth+eCoM77apxGYY1tEQfpFm8lWZ3CuQxU
% yFA7iAP8JuxJejNaZ3YJ0Rv4s40K6zRbHNZPMukcDODqrILFEcLwLb4rnbffD2xWlwrqyIq90hy+
% kplTF3x0m7ySumQsIDEaSg25Nmt+6VVXm7SUr8ZHbqbNc1yqdZ1JJET4plxzgPdIJU2JUYllna5V
% 5QMR/mGB4g2OhlQOhwJ1Uevfqi8+QXaf5UgePMIfHHiOMo1QlXy6ESWho4Q1yOY5TJroKttGdGSl
% V1sS14MNO2OfuoqN0uvN0pYba1NWgxfQcsSTNgfft14bea/llNJiEqq1f3APIggmrwCryntP1pVl
% 1MKPMOG62pBycK8jGCIDAD9Y9rqxjnI7GTtV61Jx3oNnctiEJWt9Ecwb+WUMBrySKYJGbA0kWOqC
% 8nRFQes04EVmb1yBZxUqE+GKpIw+U53TRv1eE/614lPZVsCdIhcxyrDTxmapPw65lZMRBYC5HEPU
% yCmcSS/4FYA1vluy0bgeSAG0SgiSCj4xqStGAPyF7N74C9FEmYez9VrZXJGMAvJmFFYahlRcP3gd
% F0mRATnkXjN4z2mBkk4vokBkf77JpBKZRZYMAduYaqeT4OBwG1kW+55l4VvPj73n8TLSVQzUonoC
% Tipty3iGOq7K/0LUZGlMXqWUlWwgG6eHyKhX/zhFSUzIKduYcrE3Bzb02UlxBIj9CJh6JRBbjZlJ
% fuRPTrK6qvnzWjcmCZ7bnSN8HFB2o9TxIDO1JRumdULiUKbjzIFdftSOBAk7efMZW+bB8SAR3o6w
% sClKnfMzpEZkbF36NGhr+BMeITVeAKoAicY5sdFUgWrkyCTCZwVMEC4Z7/D3sobwF8HS9JR8iWyG
% s+CqglCQtW458W5BRSnjRJEhoQv8ytGFd90z0j9Ht8Q+5g9VWpTr0qgM9QDkEYWB0pA0z6J6tUFJ
% hEKJs0ht4akYih+iHOUyW3FUfypBO41GWvaLexGs3yYZb37UocKiuAYg8dZLX1vizMq0iXPU7mA6
% Li1GyRLU6gUJlIxVIGiUF6WJADZsCnAgp1I4/15rEhdgKZxPw2zqsuZqsatHHIUe5k3QND4JBYKp
% WtRWkR2AfHlFrDI94MTBJ4JwrZtg4CD7KWpOQ5XjhfjWEY8rkVBeUartKx2bBKjtoRVZkdMRPY/o
% FQIkvMTwakMGsQs7YO1GUsXTfygxKo+G3tAbm9u1MsrWXKMMEMGeZuuWOLmQyTNZQL1xqWWKlcmi
% 4MLmX1Dg6FutMrbUKrO27JjBuutN7HSOwir9CQf29mw6pEOCNLItJ95dMZRVxLETI5D0Cxc2/IXM
% gezt5Q+sofUPNr1+Wl+erL+P4jWUbAz87lx83lM11WUgyoSeNXUKOQJaLFWjKCILh6GQlNvob1bL
% pq0Iq700Rr1VbPTYGz2RTjF3JKtHFHwXnnYCyGzdEEIxoVRy0Sdp/fajEftxeNIXD8EJbGXH8F/B
% Hi2i5Ks/3yean60cigVBt/jlcnj8fXRyEXFpAtKCAZHAKDvBMjYwxBf7TGwMRc7VKw+5iqqRXa38
% YTnSYqYiDjlfwEhcWABpg6ozB966poCCy35QiKPmqkyDFLHHgU2/h2vddObjgNxE6Hc1qhv4Lhhl
% TsyLOxogWZptU96FTdNQsVYaLW6U/3I5OvVm9jQhMA8UHA4RPCFVubkUz1i7aXKnbCkJDBeB+VZM
% Nvzh0BnZz2Wcsn6vbcX22AGCgws6s57gYehJLiLN7JGpPaWuxBJXzeVzUF4npaUmGqWbKqnj1+Jm
% k764aQgOtY1OlS+Es6lMNsI7vhcIMBcRAwr1BhIDwDJX6oWmhemk738DFRIfL0/KoxFgVVpiH8EJ
% DLqVLl3VpF+EA5KVGOfcwJsW7JEkwZ03uBPLrfCVkhG9K7wEByopKMC7lWSD6BAbHk4cNtQ0hjZo
% nyjmzGJ2+vnmx2cqlFF0sgi9V1YbGjeM1F0I3o/VHQqiMI6jgMi0N9LflIHUqgAHZLayCqUdSHzR
% L9QCrUKW9qSGWSngO5v+j7ie3z8+xvP7r7MbcX0/e5z+99fp7HoqbmcCnZ/4NJ/8SlObiIg6IZga
% Skr2EkHJxvLhSzE6OuuJAoQRrjnp0aQG9LH43hnJozF1umotGVW5ymHX80iIT5efnzrFd9n99zX9
% 25FHo273Pz2Rh9dG3Z4oL/mLYbeHFxaIwtpdXj8VnU/YUhBb5FyhuYcNsbVqyCut4Ej/eCn8jqeD
% Uffofdjx6IQWTAINHw758bjbmCpYiVegpGJJZxjzqlh2uyZwex9CbXnD8l/4MJI+PgknxUPxy3E/
% emhqbzBBQOBugMOFSnxiY7DZGiUpKtinlEoCVfNa9ryG7vey6piub7x2UAm6/GyqpUoktVeSYEvR
% BPaErPCidx1ICDA+GdnKFwMhxh8vq4+XHRMPuwOf6nFsXHU/XpIY30eXpid8lwYQ66Xmrju6suBo
% tGBgBMcZWIEyPPFh0oYzCn7i6FFoq7cMez9gkmFYQImaV8EG8/EeY/VElXKKXe7VEYPtqqbGRI1q
% vntGk4blc4D2deP7c2cpz3FiyhgT3pB75CWlHnut+m04oN365RiZJry2V3bgKZy1Vo0GnsLDhH8+
% VmXHI7n7J6A5HArzvTMcjI7Gg86JPHrX7V7ed8Ij7N/tBuULbnQ4dXD0g0Zx89dWgoaA7pY0lkCw
% P06+PNxNxex+/mVyd/u/k8Xt/Uxcf55e/9Mr44d6PMZqSvQeMO96nmgnl0ttbN5xvRFH5J+zp7vO
% ohuPEBVJ13x3pFHSIV8D91DCMVLC6C80JV4+mv20paZBZICuKhTVLGji+QzNWNB4qIjgKANo0h8y
% Ya+sJJF9lB4gGGwPCRuh8OQ8r4VMPUYSj3aSCtQzqo1eUdQ0PLtt+RoQETlhs+whTJuipukj0tA3
% KAmNk3jU9bo+Fej5j/gvdiyCv9GdRfH6ezlFeHM3CFraMKCgeCESK4YhuBbfwACGZL9bpPDvI5yL
% JIRWQGfWdDO7HnWKbs+/y1ykDaIQV9F9p4NlouhS3govIgl2L8gwXrjL4ZHtDAP1RaMEglCzs2js
% yAOyQ57gR9pgcsPBCSya1KVMfGEoUWyg6wQ5cDQY781Dm/aGZoS2RldOY5wtPImOuPLIWnxDCRmM
% AwGnLouBt5uW7Y3/druB3qQ2J0Lv6zTSmAyIv+UQQO8NBXZYKcDZ+dGqtLlogIDNa5pe+HY3IMeD
% IoyZaHiC1JbKgvGK0Cv5CWQIbWvTK1OTmqk3TnzR9F/T668cavPp9f38Jpq0Fw/tzIdnl9t+AYrg
% rx782Jc4FkSqTbIBXePxb2X/Mv1PNip5dn+5AfEDuNp4NT0n1M4yEwytAc3baFIe7jn8Rv0fjgao
% dUVYh8tTn1qySvkAChNcWSYbmKG/pw7lbDCjMG4j/Zk2CdS9+fA8MBL1tpE0g4AFQ2sO7zHNa+cp
% QRA4ClQGzeZ8dC5GZ8G4qV9LLYq74CFPM8bwU6tlZpNniKoMmjHP6MBoBf7XzKFCUm4HRaGHuPAz
% C1Zg/K7JB5yo4Au/BrKMz8Mk6GCm0+PGpk1vTVvS48lFGIb4wUjzFbY6ORfD98fHZ6dBgf3ZBrHR
% XJHNvTEu/E77GZ22QAM+evfh+MPZX7eIS8JM+zof/vPraC/Pzt4PTw+OPjiQbRX7/gnlJ9si9RNH
% CIvYsLzQW//wAHGLIs6MoqZwk/6KxGaKhpgh3T9u84ctb0Pz5lw7qsl9MVk1eAu3RZo6ZS5O5EGb
% 0G2GT3ttePAuPoACdnyFP+cvHn5bfL6fPUwWny8HFDeuQD4YYDmhbpC8HyYfljDFKhkOfo6wYovw
% MU/IPo638p/Hu7CN5ugATz68f7/X1u3PXC/E8N3JsAHOXyaC3ozsIc8HHczLzU2DcB+vh+YN9Xgv
% K/I2gbWVCg2X9PSLc2y7V+NroiShowxDW75qaWdM7aoQjijJ1M4RR/XjjPH7k+HZaU+cjt+dAsQU
% zmfDk/GZb2O3ETVmHpbcJO1pHE7qhVlEmxFy+abzOt/LBNwl7d+SVilSE2RhSEm+G6MJDCnjxH4q
% nM0GvKqP4oc0PPmbW9gnXk9pN8AlJC4SmS44NohAnXCf8jQ6u/z0NP53/Z9Bp/4+jmsU9v0+VqPr
% Q+gHvsu2WNHUQ9AYSVPi4VREniZLn40+0F0E51/u2FvP00CzGefSaiQ+0PMVe3+bUBqhQfFu9kxr
% Tk8OauPfIWAXr6PTXXKJm5MaHzcsidrSlcx1tg3M089bYq8Tl2yegzUdd6rl2oAv66Tdgok6JyJ+
% KUyld0MyT1j3buEaQygiUTTtChOQmBqPdC9xRg/z2y+T+W/x4/3XORrVx8Vk8fVRfJvObz/dXjOt
% je5hDBq8mZ9/RnAeDftiU1WFOx8MZPmmX/q2XA82VZ4NRu+Oz/pIySdnLyOKpd+oOMG5V1L3xJU2
% GX240+zdKwsbwmnlc08s7KssU54LhksgZbgRYYVop3AxdXAptbuk64mXEf/oIRWjE/GPGumRRKWe
% jWcMKPu+sDTjibnKcbAY94fhaljcSBBfcfMHosqqF1DGtoTRJodVrOcbHZzPK7gJk3tUnwYbdeCk
% +wNhVsR3xjvZw/yFLhX74j6MHsWj2uZUFRgkqV7x3XJFnMl5aoatrkr76vgSN9w8UyhDBpWeC/r3
% 7Iwud0EeRjuHrRGU9bIP1AwmlDPM1a96uUQyHDhnksZevt6yadBu23Lr2aS/sTu8pziYoZBp/Q8b
% aCfC4gxr4i/UK163Ggu6Je2HblqGw3psKRIY9ty/y2ajMWVqh+w/X2PH7Yt8nSL/nxYa9hFbbF63
% YWbNt7iU4RqyxoqrbIWiy/ccIEEEz+XP27tIG74moCxZ+69cQX1mP/xsJJC/r/M77/fdjTabq/2l
% zMEPKBCxd7dfbhePYnEvHubTx+n82268dDub3Ilv93dfv0yjG7u7DfJDXGac7a9mnKoqGvk1bXxA
% I6gm/1iABiIcceTaOFxf8BBk//JCHJxC90ZxM+8J2zZxEjXprUU9d60mzFD2fnaDkCqXGp0HYY2G
% j3IJnlPziL+mC0tPj5i6tkPzkO/2Gkwat/aaixo/LgljHS97L+Khw2ETsj3UbmYDJBlHMUS1fFfv
% 1QXJinY3RunP137tKD5HjdAuxAT/TKFpjPfmDlzhm583kVcyP98LIxffFbQ/gQjY2fuBVT/6P1BL
% AwQUAAAACABLq1VdOKihhjkHAADFEAAAJwAAAGludGVncmF0aW9uX2NoZWNrcy9zdHJ1Y3R1cmFs
% X3ZlcmlmeS5weZVXbW/iSBL+7l9Ryul22hnDAAHyspNI0czkNtK+aYl0d5tlo8ZuoDPGdtx2gPz6
% farbNoZJtHdoNHS66+Wpp6uriqOjoy8bGRY014kuFEVaLpLUFDo0NE9zKpaKVnqjIgpVXui5DmWh
% vqckLUgXhmQcd9I8UjlleZrOu0dHR948T1c0z2FUp4khvcrSvKCbaoOkoRsns5LFsj6e4zTNtYy9
% asNsV9mWhY3neZGaU6bylUxUUgjpX3iET0KXJLt5ujb2z7QssHFvupPu7ypPp3RMok8fP1Li1+f3
% vSlEWOKXRNlNjnElzVfSCeUyWSjRD6jScl74Ax/QY7nuTBcPYVoChk8d6jcibJ0FnAO1yWQSCVOu
% hLyHdkCPjKcWoj8raI/+tLHw1ochPu7wJT7pucN8dYWD76jvuwBzVZR5Au8GHMZ6vhXsr9Of+hWF
% i7hUQgY0q0JzjJnuC+gywlFJ72nmFojOb3FUBMAQ6VAZxnLPdmJtCuFQOWXf94O34xGzgPgG3n+r
% GDCW2sR0Rzy71nCs2adKypXKkX+iAtK6oYYoMP34vwjXd6Y1awAU3y/+5JtqswkZsKcyY5n6TXIO
% y5izZDAa+97aXfdWiZHvWbS7mxpVLg8uUIMAqI9agNbOL2yxI+St6d7yLQve9unykkQPaemTio2q
% r8UpImDtci5Mk8dywRFX5nwvbNANv0E3rPw/QsZh8umfNLR7YYNnP+QTvwVOM67eHqTwVThhDUca
% gyJCQxixhB6aHw99+njI8kmj1zoYWCwf8f4knNm8xhsLcR9wFQJZO8VGyEokHiiHi4BOp14Z0DNk
% 1vq+fzFlBf72MpjIwoByfOehffCZaMpOQNZFAM77F7gMq1gv/QZku0xxhNka0cLae5jFKgs97x90
% jQKaFLl+RsGjXyefKUzzXMU2Ovu0L2gpY1AcwZAutpTFpXFbtuImynS9JV+cu90zH/ZNl/ex9ukD
% DTw2g7LNJXJVxoXO4u2DitUKFtfaKLH0vYizFwA7LdTrN6RR0gou4YIZXms8Uy9ijrJwTz38O/XQ
% XhCrx6nhR7VHWNuUC8D3Zqi2HEfUMAnHWK1ryq2hq0tyglfU87KlRFJyev0kQfNG3COTkLYBdbDo
% 2NXhDvIT6fIVHDIoa+DYfbdS2e/evXLRr8dcWfNdGrSi5AxwTdcomI4I/VM/4/KfFaVzhLNgcX6p
% EomxmukE97hW8iv6UFp80EkB+TSHlUjnynbV723j5qeBnpzDTKIQXor2rc1OijR37FzJaIsEBKBY
% v6io67kuMNmuZmks3hXvkP5Kxpd3eanwovh0DQ6yraiKSODK2b0rSngY9jlWZXtdFVGuqEW948qC
% 3cmKvUt3b7fgx+t7mTxICCTJ4mGpjNEysSCb1obMj/R8LjKo4t+AKxNMd6i97x8fD/y3G1LXlDMj
% INjjJ5NJK81EP6yVXiyLfY/1S2ZJ36tAPdTZ2SpOp2ClZyslV7q2vfd/2+z5M3CTS6etuiswbUau
% aA8GJ9YkXJZ5B2mTuVTkgpGGrrKEGOwMt/EwLiOdLOx0x6Iy4eN0Zoq8tJnS9Z7cxLMRYnACPFXN
% RlTHNGj1eMGF/APXoYOiDfgnlY5TagJgy3t8jQbnAQ1752NX0V/g2BXW3tR7Cg+u/WX/NcJFU4S7
% OnkWzPuLjzFj5y9szO7qwT7eoc/TGSBahniKnNpMN3qRcKazTKdfNczZniFXWRpjrMHZiOZw/8o2
% cE6nVa9kV12ZZQqDIsLiUGZvhFDBuqQD3H22+E0oTb/E9Gk1D7rqgIcPbHGe8Q07GU6e2yRSjIfz
% Jq/ESSJXCi5MnC24hzx9Rj1KM56uwFEuI43mhFsxhUwQkh01MQA/lIgtF5uAVhVz9YAKXGJzfPyI
% xGnmfoE552BQWtmhBBHeiJ7PHe6PJu+gXR3v2Wi27PPZ8IK3Bng9NSk7YDfizFJ45lsiQciN6J/b
% t/ua8Ii5xenYSZ4GrZS+ERgGeSJk3zdiDNnBkItPLTpqiZ6M+ZrORie9yqcD4alC4qld1gpcBOyy
% Pxhz8JXiyFUWHA1OUH4HfTDEjjr2ZwDbGP4/NoZsY7Bno47dwWFb5+f909OTszFr9s7G/cGoZsvS
% 1dvTGTqdfu/k7HQ8YlpGw/NhW8lRnOVoZOLdr9eTyQWI5nbofmj+e9TpfBryUMclqukGzTDUffeq
% MrJOh0hUjBmJmaN/oiDif3k4atlfcS0T4hoNGi2v3Xaa05/4Z5YbVCDCi+boNlmiDRcAHKdrOKoa
% AcTs6gCjK8gbzuyvPL+5H3ALtHxT4VdPJfrxm+E5blD8OztGfqjagNqESkWGXKvYR9Qyd33QCcB0
% AKIZ8hNG26fwVce7/oAijDkHrtPSdFptpd03DsHrXUWJt66UAN/TZf8Dkgs/kd1y2CofLbw//khf
% /nP96Y4m//357ocvk9sJfb69/tfPv0zubj9NiD18+Qz5vwBQSwMEFAAAAAgAhqpVXWanGREoAwAA
% LwYAACgAAABpbnRlZ3JhdGlvbl9jaGVja3Mvb3BlcmF0b3JfY29uc3RhbnRzLnB5hVRNj9s4DL3n
% VxC+VFKUjSV/T8eHxWCCAosCBXpcLAzFUWaMieVEtrEpBvPfS9mOk2wLrA+xQ/I9kY8UPc97Pquy
% A6u6qjHqAKrfVR00e+heNbQ/DL7aqoWyMW2nTNd+BtOA6WttqxLDj7bB2LbTx/YPz/MWe9vUsLdI
% iXQtVPWxsR1sJgOoFjZjTK2614t7j97GVuqwWOz0HvT5WPTHo7bkzKGmDwvAR7WtxlAfHnM4wyPU
% sAQ5eKzuemug7WtCzoy9wfpKSN4o7BsLb1AZrNG8aOKAglIOG+JTin8QQ+qloOsravhLxOq83uC3
% pHRMDKVoK2WKbdObHTlNmfX5ycVKdmInOliec+Es4dXyNe9Zz4hc9QNvTxm7S/55tq7E8pnNwWhY
% k8kVMvz4esmlVpUhF234lpf8Kd8QwYVMKd+QiA+vdDYkPKK3QqrlFvIcysGmi0Pzr7aQjyo6nlsJ
% f9Ewoo4S5btlnElylDaOuAzv3Bsio5jLiK6nQKxtjBVhEqJL+FGMbuzt06irG4Fim1/HQbLtpYrR
% Wd44Sy7vzpvg0xEiDZIoCXgQZCJMU3eK0+R3kHKEyDjK4tgPuAiFDHx/wogME71DSabG/is6Hcbj
% YIwoVavb/G8iA5Qry0SSBGnMIySLBZbKHZ3jQyX5AJgeIkPn84M0iSOUJgqz8IpwCdB/hnjXl502
% TV0Zhe3i7u7ostM7DqWuDpV5cU0b0niYD7D5xMJuoHR211hXbpl1Rd2YdafyJ0acakLGdLkhQSyw
% FsrIEMjYCoeXXgGTPI4OlVP/tSOfk+uSMIZcMn6cLrZ7jrYyHdl7J3fvxfr9JuOPh5E7h3f3/kCc
% WOO0fx6oMfwd3876PhF/eHd9u7/NQ9vo1MFAZEma8SAUWRDPfZf+/+CFvBCIJE2kEBlyJiLww2Qc
% bGSep2eszPv25/fvKILbwShFY7Tp8NJBp6pDy93a7SrTN30LB21eutfVtt+96A6GQzn89W3QQGNo
% g7PrZAFtLf7OC5uDMjucglfc2E7nL2PaoM3u2GAKuLdxpVR7KAqjal0UroJPReEWTFF8Gsdm3DaL
% n1BLAwQUAAAACABVq1VdvcD4e0cEAABcCgAAKwAAAGludGVncmF0aW9uX2NoZWNrcy9vcGVyYXRv
% cl9icmlkZ2VfY2hlY2sucHmtVklv2zgUvutXEAYGJhVGlmO7mcmMCswcjN6mQOdmCAIjURFlLQRF
% L0nR/z6PohbLcdtLhUQm3/q9T4/LbDb7pz5UCU8QP7NYozjj8R7VKdIZR6U4g6KoX+6bQ3nPzxJ9
% 4k0jWIVEwist9Ks3m82cVNUlShW4i7pqkChlrTTadgLEGrR1OqH9KcSzd9CisJ6S6Qwkvd9nmPbm
% zWvjOE19UDEPjBxHUSoKHkXEU7ypiyPHxJNMAZhmtwwXcw22FSthHvnLyH9cKFMd/PqbqGTnuQMR
% PZPQE1XDlcY+bbTCNgUhTiN5HExBekYWGaQ2d1HHzNSF50zFmTjyJFKPfBO1zM2pDbWYH7kS6asn
% X+fEOV6HLOvkAKHaoCY8Ni+b3StqlnDl8TNktXb4SCxT7SeKmBI6K7kWcU+ZVKKKhWSF4yQ8RTx5
% 4ZhRQXP6Rp4cBA/biXCXh8FbN8tDEARHL66rHL8R66c4K3KuwTXr3BTXwCfSBwkoSibx0TNG9Ojl
% kqsoka0t6fyhURKmGc47b0kTyoP8MlKykPTB5Qt5n7jJAktXdr4lExXu/E6c7QHbc8GqPd6QVpbW
% CgkkKqRYBdVtOtPe3JTXFvTvIDYe+eghyJMlxpi33OAtXtKHzQdCfWKTjHqfLin0Rm/R6aFV6upl
% gLa+CW19Ac063ADXZrJawILF3ZL8traIVld4OqsLRKsBT8wK3gS7LX4gRmNe42hFwtEqCXa7o9cO
% cU9YHlIbASZuN8pDcs3choTv6behVRPYjoDGgIYoTUcMzdgj3wGHdBX2oOUNH2tJRj5VXevoCJ3e
% LjWTeLdlRcPpf+rAw5Hgd30y0LanJi0w+mAYNbRtegTmEelVjqfezafth1jRR0Lvjfsfl37ZBY+/
% iETz5DLoV5/9XLCsRqUalBfkWjNLLl2HhI667EJMJgtC+5ZL7NteuaCy3RgGOn8nf6JEKN5u4RfS
% ifkPFuV3ViGYDHuRrcBOLpJ12sz+TgIKOKjOwc6nG/qBPoY/wrL+OZa1xWJiQq6wG+WhWaksSfBN
% HR2XNHzlKXk7P4S/wbudUaz9YT33j601GPfbofiriI05ozpzFATdRpzLXRG6qrnLlRnI5k77rh2r
% pq2yGKtcXeU+BeDuG8uFjQuTKThJMxkMe7k0jaBopkaRmkaMaHooikFto05N2lVWwSF+cvHy/kRc
% nMAJoIjrPtwq2MQz5Z7cTN5Zh0zd9UFuefQ69DFAPmJVAu+/0An+lxPzZ3MdCLb4kS59QHFyj15a
% 1arEe9Il6gXjatpfribyXbyQuY2OPiJ/2rUpqmr9brvpfK0PFNuDGhFNohg8Gqfzz39/+QLxqkMl
% NDJ3kQL6g73UFWyPtE2ChiTB12nSb2DAG5EcWBF81f63p+7SZ7ZhVsGlCUE7UhTXA52GyaI+cTXc
% /Fq8HlxqHKgrisxlK4qCYB5F5giPovlTd5Y7/wNQSwMEFAAAAAgAtapVXUKysUv2BAAA7QsAACwA
% AABpbnRlZ3JhdGlvbl9jaGVja3MvY2hlY2tzL3N5bnRoZXNpc19jaGVjay5weXVWW2+jOBR+96+w
% ulrVUEKb0Gk1F1bahx1p1ZeRZt+iLDLgJE7AsLaZJh3Nf99zDCRAE6QQY5/zne/cbN/c3Px14Jml
% adWoXOQ0l3yjKmNlZui60tRuBeV5Lq2sFC+oOSqYMdLgQqVFGRLyD4jUuqrWAVWVxQUjhjgBFcby
% tJAGVhxgzTUvhRV61igJVkqaFVyWJnRYqlKzWstSzNZSFDn9xgtxpBZAaGMA4WuyeKYx/EXLZnXP
% mn+jWTObeyG5ubkha12VgG+3hUypLOtKWwCw23ahnYClsLGy6NdNLbIE15O1LERSVBlHdwNaVnkD
% E24JhVqQtYaAwbrp9b92E6T73hkYE5KLNS0qnjOkA7EBl71PhMKDWODBFbsMJQPnhOfEgQZIT8kw
% fHknvBBNCR2KA4C2sgz+WgEtbKMVQgCvlEN+YhcVlrSmEy/UwlTFD8G8ELIjlCUag+z4O4V7equx
% Rh6e738ILdfHsD7eBt3kc2IrMAC8lU0Gqb/1iI4uw0QXYKLrMG04c7mR1jDexbFza8np7zQKKKf3
% 9zQajD+uOj2oYMYDmo71TFMyxg70jh491PKoTyPfl05m8GAfyICyQ4CCUlGhmlJobgV7kzU7kQp6
% fqkHT2daic2Ur7M7OwxN9jYOY/gT9AmubIqhJ1hGy4cVoHwgA6pXYT6RoVM7cOiiZDqUdHaWEsK0
% W9G7GMB9ejxZyxGAPQbgy1nlDWhly3yFLpIzBkzE9GE4QWc0cqBvk9nFYHYYNiCyeh845KC52ggW
% uUj9Rv/WWuRNJlNZSHuk1dptPBBBqaTa0KxJZfYZo2llXUhsvB+CtvtNthXZPiTcGAHNzIsCasT3
% IyB1gN/clYozexibHSj8PKXJCaZnwXlAF8/eLxpD+wvLhpMemdQcn6p5xPzXQHcaCKMzARV58KZc
% OhM9n0Io1ql5aHYeUa5yV5fgC+7YoNkJkKqxWFE/jw70eAYFRCrbRjoGTvngeQPNX2RA/DDSW5FU
% Wl29IuslJnFOv3yhu7YCnWT12joBA/wGEivYfkLewLnDOmWPvKD+8gFpSHRkR0UBO8ocJ3adHlaH
% m53NyaDOx3ymNeM4DrKHHB0nMPIwJPbiTcVewODSYv+9LHdu5A7NEfYorwDJFk9TH+D4miZ/zO+C
% GyMumGQdhSW32Rbqm0EosDT6hL9Pike688RVO6YGAu5O3XbmHHZyOq1dAlcEwWrX9fPFcwDwc3x9
% hNcHHD31G0FHr4Z+eUQeEblc2E+DjQPOyT2cu1BhAdVgrvZ9aCN4M3638NrRwh+PI6+XGESxM771
% 92i6dCVf+yV+KPehfe1WThobXoPBPbT4IyxBo0+xUOCPwfbVTc/n/t6vXaDnfssjesejvyKwVhr8
% c6kZTDsf+HVFQIW+du6ONYOFd3f6ioJHEHk6w5yTF/K6FiqHLT6zrI6BA48htiqGUOO+mEmhsiPc
% Ql6FTtx1MAaPg+lhOH0MlInaJBko2KQBE73yIBpQbVBayjK8GYV5U9ampeE23KTScHGJsZbMXrwm
% MofjHzbtBP60FAZWfH9xhQjYW4vMJoVU+6Qvf1S5Ii/hlothwBtG1wFa8GzL24Oia4d41BzXTPeh
% 7ZWwCc/xxqo8fcTn4RU0uCLbxsS33/78/v3z6Tou3OV8dJnGHZvDYZjCt22scB3l7t+33gVs6XyN
% F5CD/wFQSwECFAMUAAAACADSqVVdulkEtPhJAABc6AAAHgAAAAAAAAAAAAAApIEAAAAAb3JpZ2lu
% YWxfY2hlY2tzL2ludmVudG9yeS5qc29uUEsBAhQDFAAAAAgA0qlVXfc/jv7gAgAAYQYAACgAAAAA
% AAAAAAAAAKSBNEoAAG9yaWdpbmFsX2NoZWNrcy9yb3VuZDAxLzAxX2xvdy92ZXJpZnkucHlQSwEC
% FAMUAAAACADSqVVdnu9DTKUFAAC+DgAAKwAAAAAAAAAAAAAApIFaTQAAb3JpZ2luYWxfY2hlY2tz
% L3JvdW5kMDEvMDJfbWVkaXVtL3ZlcmlmeS5weVBLAQIUAxQAAAAIANKpVV0N+fMBdAIAAK8GAAAw
% AAAAAAAAAAAAAACkgUhTAABvcmlnaW5hbF9jaGVja3Mvcm91bmQwMS8wM19oaWdoL3ZlcmlmaWNh
% dGlvbi50eHRQSwECFAMUAAAACADSqVVdP2cZXNwJAABbGAAAKQAAAAAAAAAAAAAApIEKVgAAb3Jp
% Z2luYWxfY2hlY2tzL3JvdW5kMDEvMDNfaGlnaC92ZXJpZnkucHlQSwECFAMUAAAACADSqVVdqNQV
% uQcBAABdAQAAMQAAAAAAAAAAAAAApIEtYAAAb3JpZ2luYWxfY2hlY2tzL3JvdW5kMDEvMDRfeGhp
% Z2gvdmVyaWZpY2F0aW9uLnR4dFBLAQIUAxQAAAAIANKpVV2mMxEPnwYAACoNAAAqAAAAAAAAAAAA
% AACkgYNhAABvcmlnaW5hbF9jaGVja3Mvcm91bmQwMS8wNF94aGlnaC92ZXJpZnkucHlQSwECFAMU
% AAAACADSqVVd+4WanKoBAADJAgAALwAAAAAAAAAAAAAApIFqaAAAb3JpZ2luYWxfY2hlY2tzL3Jv
% dW5kMDEvMDVfbWF4L3ZlcmlmaWNhdGlvbi50eHRQSwECFAMUAAAACADSqVVdX6AJT98LAAC1JAAA
% KAAAAAAAAAAAAAAApIFhagAAb3JpZ2luYWxfY2hlY2tzL3JvdW5kMDEvMDVfbWF4L3ZlcmlmeS5w
% eVBLAQIUAxQAAAAIANKpVV1yClWdMgIAAEcEAAAxAAAAAAAAAAAAAACkgYZ2AABvcmlnaW5hbF9j
% aGVja3Mvcm91bmQwMS8wNl91bHRyYS92ZXJpZmljYXRpb24udHh0UEsBAhQDFAAAAAgA0qlVXeIo
% k9d7BwAAqBEAACoAAAAAAAAAAAAAAKSBB3kAAG9yaWdpbmFsX2NoZWNrcy9yb3VuZDAxLzA2X3Vs
% dHJhL3ZlcmlmeS5weVBLAQIUAxQAAAAIANKpVV0OqXraMAEAANoBAAAiAAAAAAAAAAAAAACkgcqA
% AABvcmlnaW5hbF9jaGVja3Mvcm91bmQwMS9SRUFETUUudHh0UEsBAhQDFAAAAAgA0qlVXZ37xMQs
% BAAACAgAACYAAAAAAAAAAAAAAKSBOoIAAG9yaWdpbmFsX2NoZWNrcy9yb3VuZDAxL1NIQTI1NlNV
% TVMudHh0UEsBAhQDFAAAAAgA0qlVXVIEDsgHAQAAbwEAAC0AAAAAAAAAAAAAAKSBqoYAAG9yaWdp
% bmFsX2NoZWNrcy9yb3VuZDAxL1NUQVRVU19DT1JSRUNUSU9OLnR4dFBLAQIUAxQAAAAIANKpVV3c
% jYnJQgEAAPABAAAvAAAAAAAAAAAAAACkgfyHAABvcmlnaW5hbF9jaGVja3Mvcm91bmQwMi8wMV9s
% b3cvdmVyaWZpY2F0aW9uLnR4dFBLAQIUAxQAAAAIANKpVV0u5pt9cAUAAAoMAAAoAAAAAAAAAAAA
% AACkgYuJAABvcmlnaW5hbF9jaGVja3Mvcm91bmQwMi8wMV9sb3cvdmVyaWZ5LnB5UEsBAhQDFAAA
% AAgA0qlVXco+dx6nAwAAhAYAADIAAAAAAAAAAAAAAKSBQY8AAG9yaWdpbmFsX2NoZWNrcy9yb3Vu
% ZDAyLzAyX21lZGl1bS92ZXJpZmljYXRpb24udHh0UEsBAhQDFAAAAAgA0qlVXdgLkjuUCAAAmBUA
% ACsAAAAAAAAAAAAAAKSBOJMAAG9yaWdpbmFsX2NoZWNrcy9yb3VuZDAyLzAyX21lZGl1bS92ZXJp
% ZnkucHlQSwECFAMUAAAACADSqVVd5rWjaW4BAADBAgAAMAAAAAAAAAAAAAAApIEVnAAAb3JpZ2lu
% YWxfY2hlY2tzL3JvdW5kMDIvMDNfaGlnaC92ZXJpZmljYXRpb24udHh0UEsBAhQDFAAAAAgA0qlV
% XTLYfMRBCQAAShcAACkAAAAAAAAAAAAAAKSB0Z0AAG9yaWdpbmFsX2NoZWNrcy9yb3VuZDAyLzAz
% X2hpZ2gvdmVyaWZ5LnB5UEsBAhQDFAAAAAgA0qlVXd5OBhfrAQAAVwQAADEAAAAAAAAAAAAAAKSB
% WacAAG9yaWdpbmFsX2NoZWNrcy9yb3VuZDAyLzA0X3hoaWdoL3ZlcmlmaWNhdGlvbi50eHRQSwEC
% FAMUAAAACADSqVVdmnGR2sIJAAAfFAAAKgAAAAAAAAAAAAAApIGTqQAAb3JpZ2luYWxfY2hlY2tz
% L3JvdW5kMDIvMDRfeGhpZ2gvdmVyaWZ5LnB5UEsBAhQDFAAAAAgA0qlVXSsRN33MAQAADAMAAC8A
% AAAAAAAAAAAAAKSBnbMAAG9yaWdpbmFsX2NoZWNrcy9yb3VuZDAyLzA1X21heC92ZXJpZmljYXRp
% b24udHh0UEsBAhQDFAAAAAgA0qlVXbzQrvbkDgAAyy0AACgAAAAAAAAAAAAAAKSBtrUAAG9yaWdp
% bmFsX2NoZWNrcy9yb3VuZDAyLzA1X21heC92ZXJpZnkucHlQSwECFAMUAAAACADSqVVdu1hLK88C
% AACZBQAAMQAAAAAAAAAAAAAApIHgxAAAb3JpZ2luYWxfY2hlY2tzL3JvdW5kMDIvMDZfdWx0cmEv
% dmVyaWZpY2F0aW9uLnR4dFBLAQIUAxQAAAAIANKpVV0ewASCfgwAAGkjAAAqAAAAAAAAAAAAAACk
% gf7HAABvcmlnaW5hbF9jaGVja3Mvcm91bmQwMi8wNl91bHRyYS92ZXJpZnkucHlQSwECFAMUAAAA
% CADSqVVdbi8WgTABAADaAQAAIgAAAAAAAAAAAAAApIHE1AAAb3JpZ2luYWxfY2hlY2tzL3JvdW5k
% MDIvUkVBRE1FLnR4dFBLAQIUAxQAAAAIANKpVV3wMoU0qgQAAB0JAAAmAAAAAAAAAAAAAACkgTTW
% AABvcmlnaW5hbF9jaGVja3Mvcm91bmQwMi9TSEEyNTZTVU1TLnR4dFBLAQIUAxQAAAAIANKpVV1X
% gC+OJAEAAM4BAAAzAAAAAAAAAAAAAACkgSLbAABvcmlnaW5hbF9jaGVja3Mvcm91bmQwMi9hc3Nl
% c3NtZW50L3ZlcmlmaWNhdGlvbi50eHRQSwECFAMUAAAACADSqVVdcd1NZ/MAAABkAQAALwAAAAAA
% AAAAAAAApIGX3AAAb3JpZ2luYWxfY2hlY2tzL3JvdW5kMDMvMDFfbG93L3ZlcmlmaWNhdGlvbi50
% eHRQSwECFAMUAAAACADSqVVdd8ffKbYEAACdCgAAKAAAAAAAAAAAAAAApIHX3QAAb3JpZ2luYWxf
% Y2hlY2tzL3JvdW5kMDMvMDFfbG93L3ZlcmlmeS5weVBLAQIUAxQAAAAIANKpVV0UMdhAoAQAAAsJ
% AAAyAAAAAAAAAAAAAACkgdPiAABvcmlnaW5hbF9jaGVja3Mvcm91bmQwMy8wMl9tZWRpdW0vdmVy
% aWZpY2F0aW9uLnR4dFBLAQIUAxQAAAAIANKpVV3bbKVRVwsAAIsdAAArAAAAAAAAAAAAAACkgcPn
% AABvcmlnaW5hbF9jaGVja3Mvcm91bmQwMy8wMl9tZWRpdW0vdmVyaWZ5LnB5UEsBAhQDFAAAAAgA
% 0qlVXYAw0dyQAQAAOwIAADAAAAAAAAAAAAAAAKSBY/MAAG9yaWdpbmFsX2NoZWNrcy9yb3VuZDAz
% LzAzX2hpZ2gvdmVyaWZpY2F0aW9uLnR4dFBLAQIUAxQAAAAIANKpVV2jUTZIFgkAAHAVAAApAAAA
% AAAAAAAAAACkgUH1AABvcmlnaW5hbF9jaGVja3Mvcm91bmQwMy8wM19oaWdoL3ZlcmlmeS5weVBL
% AQIUAxQAAAAIANKpVV3+WgUZ7wEAADkDAAAxAAAAAAAAAAAAAACkgZ7+AABvcmlnaW5hbF9jaGVj
% a3Mvcm91bmQwMy8wNF94aGlnaC92ZXJpZmljYXRpb24udHh0UEsBAhQDFAAAAAgA0qlVXTSD9+Y/
% DAAAWh0AACoAAAAAAAAAAAAAAKSB3AABAG9yaWdpbmFsX2NoZWNrcy9yb3VuZDAzLzA0X3hoaWdo
% L3ZlcmlmeS5weVBLAQIUAxQAAAAIANKpVV2HRsw8RQUAALMOAAAyAAAAAAAAAAAAAACkgWMNAQBv
% cmlnaW5hbF9jaGVja3Mvcm91bmQwMy8wNV9tYXgvZXhhY3RfYXJpdGhtZXRpYy5weVBLAQIUAxQA
% AAAIANKpVV3Ouh2YbAIAANwDAAAvAAAAAAAAAAAAAACkgfgSAQBvcmlnaW5hbF9jaGVja3Mvcm91
% bmQwMy8wNV9tYXgvdmVyaWZpY2F0aW9uLnR4dFBLAQIUAxQAAAAIANKpVV07DQdNrBIAAHY6AAAo
% AAAAAAAAAAAAAACkgbEVAQBvcmlnaW5hbF9jaGVja3Mvcm91bmQwMy8wNV9tYXgvdmVyaWZ5LnB5
% UEsBAhQDFAAAAAgA0qlVXTBJgPiqAgAAVgUAADEAAAAAAAAAAAAAAKSBoygBAG9yaWdpbmFsX2No
% ZWNrcy9yb3VuZDAzLzA2X3VsdHJhL3ZlcmlmaWNhdGlvbi50eHRQSwECFAMUAAAACADSqVVdC70L
% T+cLAAApIAAAKgAAAAAAAAAAAAAApIGcKwEAb3JpZ2luYWxfY2hlY2tzL3JvdW5kMDMvMDZfdWx0
% cmEvdmVyaWZ5LnB5UEsBAhQDFAAAAAgA0qlVXWmzJ2HcAQAAGAMAACIAAAAAAAAAAAAAAKSByzcB
% AG9yaWdpbmFsX2NoZWNrcy9yb3VuZDAzL1JFQURNRS50eHRQSwECFAMUAAAACADSqVVd3ikUmhcF
% AADdCQAAJgAAAAAAAAAAAAAApIHnOQEAb3JpZ2luYWxfY2hlY2tzL3JvdW5kMDMvU0hBMjU2U1VN
% Uy50eHRQSwECFAMUAAAACADSqVVdzYRqa4gBAACVAgAAMwAAAAAAAAAAAAAApIFCPwEAb3JpZ2lu
% YWxfY2hlY2tzL3JvdW5kMDMvYXNzZXNzbWVudC92ZXJpZmljYXRpb24udHh0UEsBAhQDFAAAAAgA
% 0qlVXRFKwRguAQAA8QEAAC8AAAAAAAAAAAAAAKSBG0EBAG9yaWdpbmFsX2NoZWNrcy9yb3VuZDA0
% LzAxX2xvdy92ZXJpZmljYXRpb24udHh0UEsBAhQDFAAAAAgA0qlVXT1zkfXlAwAAWwoAACgAAAAA
% AAAAAAAAAKSBlkIBAG9yaWdpbmFsX2NoZWNrcy9yb3VuZDA0LzAxX2xvdy92ZXJpZnkucHlQSwEC
% FAMUAAAACADSqVVd9oNCxnQBAAAuBQAAMgAAAAAAAAAAAAAApIHBRgEAb3JpZ2luYWxfY2hlY2tz
% L3JvdW5kMDQvMDJfbWVkaXVtL3ZlcmlmaWNhdGlvbi50eHRQSwECFAMUAAAACADSqVVdhR/ADLoI
% AABKFAAAKwAAAAAAAAAAAAAApIGFSAEAb3JpZ2luYWxfY2hlY2tzL3JvdW5kMDQvMDJfbWVkaXVt
% L3ZlcmlmeS5weVBLAQIUAxQAAAAIANKpVV2asAm7WwEAAEwCAAAwAAAAAAAAAAAAAACkgYhRAQBv
% cmlnaW5hbF9jaGVja3Mvcm91bmQwNC8wM19oaWdoL3ZlcmlmaWNhdGlvbi50eHRQSwECFAMUAAAA
% CADSqVVdraVJAEkKAACOFwAAKQAAAAAAAAAAAAAApIExUwEAb3JpZ2luYWxfY2hlY2tzL3JvdW5k
% MDQvMDNfaGlnaC92ZXJpZnkucHlQSwECFAMUAAAACADSqVVdznP+g3wDAAAQBgAAMQAAAAAAAAAA
% AAAApIHBXQEAb3JpZ2luYWxfY2hlY2tzL3JvdW5kMDQvMDRfeGhpZ2gvdmVyaWZpY2F0aW9uLnR4
% dFBLAQIUAxQAAAAIANKpVV2BbJVs7REAAMcuAAAqAAAAAAAAAAAAAACkgYxhAQBvcmlnaW5hbF9j
% aGVja3Mvcm91bmQwNC8wNF94aGlnaC92ZXJpZnkucHlQSwECFAMUAAAACADSqVVdKPvWr4QFAABH
% DwAAMgAAAAAAAAAAAAAApIHBcwEAb3JpZ2luYWxfY2hlY2tzL3JvdW5kMDQvMDVfbWF4L2V4YWN0
% X2FyaXRobWV0aWMucHlQSwECFAMUAAAACADSqVVdwuFxq+wCAAB7BQAALwAAAAAAAAAAAAAApIGV
% eQEAb3JpZ2luYWxfY2hlY2tzL3JvdW5kMDQvMDVfbWF4L3ZlcmlmaWNhdGlvbi50eHRQSwECFAMU
% AAAACADSqVVdnLfvImITAACCNQAAKAAAAAAAAAAAAAAApIHOfAEAb3JpZ2luYWxfY2hlY2tzL3Jv
% dW5kMDQvMDVfbWF4L3ZlcmlmeS5weVBLAQIUAxQAAAAIANKpVV2fPecz9gIAAKwGAAAxAAAAAAAA
% AAAAAACkgXaQAQBvcmlnaW5hbF9jaGVja3Mvcm91bmQwNC8wNl91bHRyYS92ZXJpZmljYXRpb24u
% dHh0UEsBAhQDFAAAAAgA0qlVXcG2JqExDAAA7R4AACoAAAAAAAAAAAAAAKSBu5MBAG9yaWdpbmFs
% X2NoZWNrcy9yb3VuZDA0LzA2X3VsdHJhL3ZlcmlmeS5weVBLAQIUAxQAAAAIANKpVV2rLfDX3AEA
% ABgDAAAiAAAAAAAAAAAAAACkgTSgAQBvcmlnaW5hbF9jaGVja3Mvcm91bmQwNC9SRUFETUUudHh0
% UEsBAhQDFAAAAAgA0qlVXTSZo6zeBAAAegkAACYAAAAAAAAAAAAAAKSBUKIBAG9yaWdpbmFsX2No
% ZWNrcy9yb3VuZDA0L1NIQTI1NlNVTVMudHh0UEsBAhQDFAAAAAgA0qlVXb8sjU7aAAAATAEAADMA
% AAAAAAAAAAAAAKSBcqcBAG9yaWdpbmFsX2NoZWNrcy9yb3VuZDA0L2Fzc2Vzc21lbnQvdmVyaWZp
% Y2F0aW9uLnR4dFBLAQIUAxQAAAAIANKpVV2eOsa9lAIAAFsIAAAvAAAAAAAAAAAAAACkgZ2oAQBv
% cmlnaW5hbF9jaGVja3Mvcm91bmQwNS8wMV9sb3cvdmVyaWZpY2F0aW9uLnR4dFBLAQIUAxQAAAAI
% ANKpVV0NHm97TwcAADYSAAAoAAAAAAAAAAAAAACkgX6rAQBvcmlnaW5hbF9jaGVja3Mvcm91bmQw
% NS8wMV9sb3cvdmVyaWZ5LnB5UEsBAhQDFAAAAAgA0qlVXV5hFfsxAgAA2wwAADIAAAAAAAAAAAAA
% AKSBE7MBAG9yaWdpbmFsX2NoZWNrcy9yb3VuZDA1LzAyX21lZGl1bS92ZXJpZmljYXRpb24udHh0
% UEsBAhQDFAAAAAgA0qlVXTTJqRpHCAAAKhIAACsAAAAAAAAAAAAAAKSBlLUBAG9yaWdpbmFsX2No
% ZWNrcy9yb3VuZDA1LzAyX21lZGl1bS92ZXJpZnkucHlQSwECFAMUAAAACADSqVVdXIwfOZsBAABn
% AgAAMAAAAAAAAAAAAAAApIEkvgEAb3JpZ2luYWxfY2hlY2tzL3JvdW5kMDUvMDNfaGlnaC92ZXJp
% ZmljYXRpb24udHh0UEsBAhQDFAAAAAgA0qlVXc9nI/FaCgAA+RYAACkAAAAAAAAAAAAAAKSBDcAB
% AG9yaWdpbmFsX2NoZWNrcy9yb3VuZDA1LzAzX2hpZ2gvdmVyaWZ5LnB5UEsBAhQDFAAAAAgA0qlV
% XQbLPKBIAgAA4AMAADEAAAAAAAAAAAAAAKSBrsoBAG9yaWdpbmFsX2NoZWNrcy9yb3VuZDA1LzA0
% X3hoaWdoL3ZlcmlmaWNhdGlvbi50eHRQSwECFAMUAAAACADSqVVdUMdJPBARAABqKwAAKgAAAAAA
% AAAAAAAApIFFzQEAb3JpZ2luYWxfY2hlY2tzL3JvdW5kMDUvMDRfeGhpZ2gvdmVyaWZ5LnB5UEsB
% AhQDFAAAAAgA0qlVXfbl5C6EBQAARw8AADIAAAAAAAAAAAAAAKSBnd4BAG9yaWdpbmFsX2NoZWNr
% cy9yb3VuZDA1LzA1X21heC9leGFjdF9hcml0aG1ldGljLnB5UEsBAhQDFAAAAAgA0qlVXQwouZxM
% AgAAGQQAAC8AAAAAAAAAAAAAAKSBceQBAG9yaWdpbmFsX2NoZWNrcy9yb3VuZDA1LzA1X21heC92
% ZXJpZmljYXRpb24udHh0UEsBAhQDFAAAAAgA0qlVXcTY5cQFEQAAsC4AACgAAAAAAAAAAAAAAKSB
% CucBAG9yaWdpbmFsX2NoZWNrcy9yb3VuZDA1LzA1X21heC92ZXJpZnkucHlQSwECFAMUAAAACADS
% qVVdrQ9o3p0DAABvCwAAMQAAAAAAAAAAAAAApIFV+AEAb3JpZ2luYWxfY2hlY2tzL3JvdW5kMDUv
% MDZfdWx0cmEvdmVyaWZpY2F0aW9uLnR4dFBLAQIUAxQAAAAIANKpVV29jZ8tlw4AAKwnAAAqAAAA
% AAAAAAAAAACkgUH8AQBvcmlnaW5hbF9jaGVja3Mvcm91bmQwNS8wNl91bHRyYS92ZXJpZnkucHlQ
% SwECFAMUAAAACADSqVVdRnMRrtwBAAAYAwAAIgAAAAAAAAAAAAAApIEgCwIAb3JpZ2luYWxfY2hl
% Y2tzL3JvdW5kMDUvUkVBRE1FLnR4dFBLAQIUAxQAAAAIANKpVV2m174t3gQAAHoJAAAmAAAAAAAA
% AAAAAACkgTwNAgBvcmlnaW5hbF9jaGVja3Mvcm91bmQwNS9TSEEyNTZTVU1TLnR4dFBLAQIUAxQA
% AAAIANKpVV2PjmE/+gAAAHUBAAAzAAAAAAAAAAAAAACkgV4SAgBvcmlnaW5hbF9jaGVja3Mvcm91
% bmQwNS9hc3Nlc3NtZW50L3ZlcmlmaWNhdGlvbi50eHRQSwECFAMUAAAACADSqVVdyAQ7UMsBAADP
% AwAALwAAAAAAAAAAAAAApIGpEwIAb3JpZ2luYWxfY2hlY2tzL3JvdW5kMDYvMDFfbG93L3Zlcmlm
% aWNhdGlvbi50eHRQSwECFAMUAAAACADSqVVdFhb25RsIAACzFAAAKAAAAAAAAAAAAAAApIHBFQIA
% b3JpZ2luYWxfY2hlY2tzL3JvdW5kMDYvMDFfbG93L3ZlcmlmeS5weVBLAQIUAxQAAAAIANKpVV3K
% NA8h6QEAAD8IAAAyAAAAAAAAAAAAAACkgSIeAgBvcmlnaW5hbF9jaGVja3Mvcm91bmQwNi8wMl9t
% ZWRpdW0vdmVyaWZpY2F0aW9uLnR4dFBLAQIUAxQAAAAIANKpVV0+8sTdqAsAAG0aAAArAAAAAAAA
% AAAAAACkgVsgAgBvcmlnaW5hbF9jaGVja3Mvcm91bmQwNi8wMl9tZWRpdW0vdmVyaWZ5LnB5UEsB
% AhQDFAAAAAgA0qlVXdXcdW7QAQAAGAMAADAAAAAAAAAAAAAAAKSBTCwCAG9yaWdpbmFsX2NoZWNr
% cy9yb3VuZDA2LzAzX2hpZ2gvdmVyaWZpY2F0aW9uLnR4dFBLAQIUAxQAAAAIANKpVV249JbfSAoA
% AOsWAAApAAAAAAAAAAAAAACkgWouAgBvcmlnaW5hbF9jaGVja3Mvcm91bmQwNi8wM19oaWdoL3Zl
% cmlmeS5weVBLAQIUAxQAAAAIANKpVV0mOXjZfgIAAFAEAAAxAAAAAAAAAAAAAACkgfk4AgBvcmln
% aW5hbF9jaGVja3Mvcm91bmQwNi8wNF94aGlnaC92ZXJpZmljYXRpb24udHh0UEsBAhQDFAAAAAgA
% 0qlVXfn9BobEFQAAvD8AACoAAAAAAAAAAAAAAKSBxjsCAG9yaWdpbmFsX2NoZWNrcy9yb3VuZDA2
% LzA0X3hoaWdoL3ZlcmlmeS5weVBLAQIUAxQAAAAIANKpVV3VwMN2hAUAAEcPAAAyAAAAAAAAAAAA
% AACkgdJRAgBvcmlnaW5hbF9jaGVja3Mvcm91bmQwNi8wNV9tYXgvZXhhY3RfYXJpdGhtZXRpYy5w
% eVBLAQIUAxQAAAAIANKpVV3hUi4QLAMAAF8GAAAvAAAAAAAAAAAAAACkgaZXAgBvcmlnaW5hbF9j
% aGVja3Mvcm91bmQwNi8wNV9tYXgvdmVyaWZpY2F0aW9uLnR4dFBLAQIUAxQAAAAIANKpVV1DtKAf
% ZxUAAOw6AAAoAAAAAAAAAAAAAACkgR9bAgBvcmlnaW5hbF9jaGVja3Mvcm91bmQwNi8wNV9tYXgv
% dmVyaWZ5LnB5UEsBAhQDFAAAAAgA0qlVXb2Nny2XDgAArCcAAC8AAAAAAAAAAAAAAKSBzHACAG9y
% aWdpbmFsX2NoZWNrcy9yb3VuZDA2LzA2X3VsdHJhL2NvcmVfY2hlY2tzLnB5UEsBAhQDFAAAAAgA
% 0qlVXVf5CJ/nBQAA5RUAADEAAAAAAAAAAAAAAKSBsH8CAG9yaWdpbmFsX2NoZWNrcy9yb3VuZDA2
% LzA2X3VsdHJhL3ZlcmlmaWNhdGlvbi50eHRQSwECFAMUAAAACADSqVVdS4Wez9sRAABoMwAAKgAA
% AAAAAAAAAAAApIHmhQIAb3JpZ2luYWxfY2hlY2tzL3JvdW5kMDYvMDZfdWx0cmEvdmVyaWZ5LnB5
% UEsBAhQDFAAAAAgA0qlVXfkPhGDrAQAAOwMAACIAAAAAAAAAAAAAAKSBCZgCAG9yaWdpbmFsX2No
% ZWNrcy9yb3VuZDA2L1JFQURNRS50eHRQSwECFAMUAAAACADSqVVdcHtKRBUFAADUCQAAJgAAAAAA
% AAAAAAAApIE0mgIAb3JpZ2luYWxfY2hlY2tzL3JvdW5kMDYvU0hBMjU2U1VNUy50eHRQSwECFAMU
% AAAACADSqVVdPl16fOwAAABhAQAAMwAAAAAAAAAAAAAApIGNnwIAb3JpZ2luYWxfY2hlY2tzL3Jv
% dW5kMDYvYXNzZXNzbWVudC92ZXJpZmljYXRpb24udHh0UEsBAhQDFAAAAAgA0qlVXTkbmRR/AgAA
% UgYAAC8AAAAAAAAAAAAAAKSByqACAG9yaWdpbmFsX2NoZWNrcy9yb3VuZDA3LzAxX2xvdy92ZXJp
% ZmljYXRpb24udHh0UEsBAhQDFAAAAAgA0qlVXeMCrlw1CgAAdhsAACgAAAAAAAAAAAAAAKSBlqMC
% AG9yaWdpbmFsX2NoZWNrcy9yb3VuZDA3LzAxX2xvdy92ZXJpZnkucHlQSwECFAMUAAAACADSqVVd
% 9iGM5VICAADABwAAMgAAAAAAAAAAAAAApIERrgIAb3JpZ2luYWxfY2hlY2tzL3JvdW5kMDcvMDJf
% bWVkaXVtL3ZlcmlmaWNhdGlvbi50eHRQSwECFAMUAAAACADSqVVds2R7b9INAAB+IAAAKwAAAAAA
% AAAAAAAApIGzsAIAb3JpZ2luYWxfY2hlY2tzL3JvdW5kMDcvMDJfbWVkaXVtL3ZlcmlmeS5weVBL
% AQIUAxQAAAAIANKpVV1SPLd6SwEAAB4CAAAwAAAAAAAAAAAAAACkgc6+AgBvcmlnaW5hbF9jaGVj
% a3Mvcm91bmQwNy8wM19oaWdoL3ZlcmlmaWNhdGlvbi50eHRQSwECFAMUAAAACADSqVVdsji2woUH
% AABcEQAAKQAAAAAAAAAAAAAApIFnwAIAb3JpZ2luYWxfY2hlY2tzL3JvdW5kMDcvMDNfaGlnaC92
% ZXJpZnkucHlQSwECFAMUAAAACADSqVVd9jOscagCAACHBAAAMQAAAAAAAAAAAAAApIEzyAIAb3Jp
% Z2luYWxfY2hlY2tzL3JvdW5kMDcvMDRfeGhpZ2gvdmVyaWZpY2F0aW9uLnR4dFBLAQIUAxQAAAAI
% ANKpVV164f58KRUAAP05AAAqAAAAAAAAAAAAAACkgSrLAgBvcmlnaW5hbF9jaGVja3Mvcm91bmQw
% Ny8wNF94aGlnaC92ZXJpZnkucHlQSwECFAMUAAAACADSqVVdZP6ko58FAACMDwAAMgAAAAAAAAAA
% AAAApIGb4AIAb3JpZ2luYWxfY2hlY2tzL3JvdW5kMDcvMDVfbWF4L2V4YWN0X2FyaXRobWV0aWMu
% cHlQSwECFAMUAAAACADSqVVdZNmmIIADAADzBgAALwAAAAAAAAAAAAAApIGK5gIAb3JpZ2luYWxf
% Y2hlY2tzL3JvdW5kMDcvMDVfbWF4L3ZlcmlmaWNhdGlvbi50eHRQSwECFAMUAAAACADSqVVdG10G
% FckbAABaVAAAKAAAAAAAAAAAAAAApIFX6gIAb3JpZ2luYWxfY2hlY2tzL3JvdW5kMDcvMDVfbWF4
% L3ZlcmlmeS5weVBLAQIUAxQAAAAIANKpVV0N9VmXcgIAAEEEAAAxAAAAAAAAAAAAAACkgWYGAwBv
% cmlnaW5hbF9jaGVja3Mvcm91bmQwNy8wNl91bHRyYS92ZXJpZmljYXRpb24udHh0UEsBAhQDFAAA
% AAgA0qlVXUrrrY1aEAAA9CwAACoAAAAAAAAAAAAAAKSBJwkDAG9yaWdpbmFsX2NoZWNrcy9yb3Vu
% ZDA3LzA2X3VsdHJhL3ZlcmlmeS5weVBLAQIUAxQAAAAIANKpVV20456k6wEAADsDAAAiAAAAAAAA
% AAAAAACkgckZAwBvcmlnaW5hbF9jaGVja3Mvcm91bmQwNy9SRUFETUUudHh0UEsBAhQDFAAAAAgA
% 0qlVXW5FWBvdBAAAegkAACYAAAAAAAAAAAAAAKSB9BsDAG9yaWdpbmFsX2NoZWNrcy9yb3VuZDA3
% L1NIQTI1NlNVTVMudHh0UEsBAhQDFAAAAAgA0qlVXUKf5+TJAAAAOwEAADMAAAAAAAAAAAAAAKSB
% FSEDAG9yaWdpbmFsX2NoZWNrcy9yb3VuZDA3L2Fzc2Vzc21lbnQvdmVyaWZpY2F0aW9uLnR4dFBL
% AQIUAxQAAAAIAMyqVV0V0NzNZRUAAG8vAAAnAAAAAAAAAAAAAACkgS8iAwBpbnRlZ3JhdGlvbl9j
% aGVja3MvYXJpdGhtZXRpY19hdWRpdC50eHRQSwECFAMUAAAACAAkq1Vd+k88cHwiAADNUAAAJQAA
% AAAAAAAAAAAApIHZNwMAaW50ZWdyYXRpb25fY2hlY2tzL2dlb21ldHJ5X2F1ZGl0LnR4dFBLAQIU
% AxQAAAAIAGWrVV11VcLZ5BMAAFMrAAAlAAAAAAAAAAAAAACkgZhaAwBpbnRlZ3JhdGlvbl9jaGVj
% a3Mvb3BlcmF0b3JfYXVkaXQudHh0UEsBAhQDFAAAAAgAyKpVXVFXjIpWEQAAOyMAACQAAAAAAAAA
% AAAAAKSBv24DAGludGVncmF0aW9uX2NoZWNrcy9zcGVjaWFsX2F1ZGl0LnR4dFBLAQIUAxQAAAAI
% AFGrVV2AJnNDHBMAAB0rAAAnAAAAAAAAAAAAAACkgVeAAwBpbnRlZ3JhdGlvbl9jaGVja3Mvc3Ry
% dWN0dXJhbF9hdWRpdC50eHRQSwECFAMUAAAACADUqlVdUm8usEQSAAAVJgAAKAAAAAAAAAAAAAAA
% pIG4kwMAaW50ZWdyYXRpb25fY2hlY2tzL3RvdXJuYW1lbnRzX2F1ZGl0LnR4dFBLAQIUAxQAAAAI
% AEurVV04qKGGOQcAAMUQAAAnAAAAAAAAAAAAAACkgUKmAwBpbnRlZ3JhdGlvbl9jaGVja3Mvc3Ry
% dWN0dXJhbF92ZXJpZnkucHlQSwECFAMUAAAACACGqlVdZqcZESgDAAAvBgAAKAAAAAAAAAAAAAAA
% pIHArQMAaW50ZWdyYXRpb25fY2hlY2tzL29wZXJhdG9yX2NvbnN0YW50cy5weVBLAQIUAxQAAAAI
% AFWrVV29wPh7RwQAAFwKAAArAAAAAAAAAAAAAACkgS6xAwBpbnRlZ3JhdGlvbl9jaGVja3Mvb3Bl
% cmF0b3JfYnJpZGdlX2NoZWNrLnB5UEsBAhQDFAAAAAgAtapVXUKysUv2BAAA7QsAACwAAAAAAAAA
% AAAAAKSBvrUDAGludGVncmF0aW9uX2NoZWNrcy9jaGVja3Mvc3ludGhlc2lzX2NoZWNrLnB5UEsF
% BgAAAAB4AHgADSoAAP66AwAAAA==
% END INTEGRATED VERIFICATION ZIP BASE64
